% Shared preamble for the GLVbook papers.  This is the single canonical copy;
% a driver one level below GLVbook reaches it as
%
%     \input{../basicFiles/preamble}
%
% What stays in the driver, because it differs from paper to paper:
% \documentclass, font choice, \title, \author, \date, and any macro genuinely
% local to one paper.  Everything else belongs here.  The file is
% self-contained: it does not read any other file.

\usepackage[utf8]{inputenc}
\usepackage[novbox]{pdfsync}   % required: SyncTeX alone does not work here

% Page geometry, set once here for every paper that loads this file.  A driver
% that needs different margins says \geometry{...}; a second
% \usepackage[...]{geometry} raises an option clash and must not be used.
\makeatletter
\@ifundefined{AAdefNnogeometry}{\usepackage[a4paper,margin=2cm]{geometry}}{}
\makeatother

% mathematics
\usepackage{amsmath,amssymb,amsthm,mathtools,bm,mathrsfs,bbm}
\usepackage{centernot}

% floats, tables, lists
\usepackage{graphicx,float}
\makeatletter
\@ifpackageloaded{subfig}{}{\usepackage{subcaption}}
\makeatother
\usepackage{booktabs,longtable,threeparttable,array,multirow}
\usepackage{enumitem}

% colour, code, algorithms
\usepackage{xcolor,listings,tcolorbox}
\makeatletter
\@ifundefined{c@algorithm}{\usepackage{algorithm}}{}
\makeatother
\usepackage{algorithmicx,algpseudocode}

% diagrams
\usepackage{tikz,tikz-cd,pgfplots}
\usetikzlibrary{graphs,decorations.pathmorphing,decorations.markings,
arrows.meta,automata,positioning}
\tikzcdset{every label/.append style = {font = \small}}
\pgfplotsset{compat=1.18}

% chemistry
\usepackage[version=4]{mhchem}

\usepackage{microtype}
\usepackage{lineno}
\usepackage{nomencl}
\makenomenclature

% hyperref first, then bookmark, then glossaries: glossaries must follow
% hyperref or the entry links are built against the wrong anchors.
\makeatletter
\@ifundefined{hyperrefpackage}{%
  \@ifundefined{AAdefNnohyperref}{\usepackage{hyperref}}{}%
}{}
\@ifundefined{hyperrefpackage}{\usepackage{bookmark}}{}
\makeatother
\usepackage[toc,sort=use]{glossaries}
\makeglossaries

% Book/standalone switch.  Material that the book already covers in its
% prelude, but that a standalone paper still has to state, is wrapped as
%     \ifbookbuild\else ... \fi
% Drivers default to standalone; bookM.tex says \bookbuildtrue.
\newif\ifbookbuild
\bookbuildfalse

% Part switches for the gradostat paper.  Both default to true; a driver that
% wants the short submission says \madecfalse and/or \slowfastfalse, and every
% cross-reference into the dropped part disappears with it.
%   \ifmadec     Part III, the Lagasquie--Madec model and its break-even geometry
%   \ifslowfast  Part IV, everything carrying the small parameter \eps
\newif\ifmadec     \madectrue
\newif\ifslowfast  \slowfasttrue

% Part switch for the Barrat/simplicial-contagion project.  Defaults to
% true (the hypernetwork-synchrony part is included); the journal-paper
% driver says \hnetpartfalse to drop it and every cross-reference into it.
\newif\ifhnetpart  \hnetparttrue

%% ========================================
%% PACKAGES
%% ========================================
\RequirePackage{etoolbox}
\usepackage{makecell}
\usepackage[novbox]{pdfsync}% obsolete; SyncTeX is built into modern engines
\usepackage{amssymb, amsfonts,mathtools,bbm,eurosym,graphicx,latexsym, amsmath,amsthm,times,nccmath,url}
\usepackage{listings,verbatim,cases,tabularx}
% hyperref is deliberately NOT in the line above.  Journal classes often
% load it themselves, with a long option list, from a begindocument hook
% (mdpi.cls does: \DeclareHookRule{begindocument}{hyperref}{before}{mdpi}).
% Loading it here first, without options, makes that later load raise
% "Option clash for package hyperref".  So load it only if nobody else
% has, or will:
%   - \@ifpackageloaded  catches a class that loaded it already;
%   - \hyperrefpackage    is mdpi.cls's own deferred loader;
%   - \AAdefNnohyperref   lets any other driver opt out by defining it
%                         before % Shared preamble for the GLVbook papers.  This is the single canonical copy;
% a driver one level below GLVbook reaches it as
%
%     \input{../basicFiles/preamble}
%
% What stays in the driver, because it differs from paper to paper:
% \documentclass, font choice, \title, \author, \date, and any macro genuinely
% local to one paper.  Everything else belongs here.  The file is
% self-contained: it does not read any other file.

\usepackage[utf8]{inputenc}
\usepackage[novbox]{pdfsync}   % required: SyncTeX alone does not work here

% Page geometry, set once here for every paper that loads this file.  A driver
% that needs different margins says \geometry{...}; a second
% \usepackage[...]{geometry} raises an option clash and must not be used.
\makeatletter
\@ifundefined{AAdefNnogeometry}{\usepackage[a4paper,margin=2cm]{geometry}}{}
\makeatother

% mathematics
\usepackage{amsmath,amssymb,amsthm,mathtools,bm,mathrsfs,bbm}
\usepackage{centernot}

% floats, tables, lists
\usepackage{graphicx,float}
\makeatletter
\@ifpackageloaded{subfig}{}{\usepackage{subcaption}}
\makeatother
\usepackage{booktabs,longtable,threeparttable,array,multirow}
\usepackage{enumitem}

% colour, code, algorithms
\usepackage{xcolor,listings,tcolorbox}
\makeatletter
\@ifundefined{c@algorithm}{\usepackage{algorithm}}{}
\makeatother
\usepackage{algorithmicx,algpseudocode}

% diagrams
\usepackage{tikz,tikz-cd,pgfplots}
\usetikzlibrary{graphs,decorations.pathmorphing,decorations.markings,
arrows.meta,automata,positioning}
\tikzcdset{every label/.append style = {font = \small}}
\pgfplotsset{compat=1.18}

% chemistry
\usepackage[version=4]{mhchem}

\usepackage{microtype}
\usepackage{lineno}
\usepackage{nomencl}
\makenomenclature

% hyperref first, then bookmark, then glossaries: glossaries must follow
% hyperref or the entry links are built against the wrong anchors.
\makeatletter
\@ifundefined{hyperrefpackage}{%
  \@ifundefined{AAdefNnohyperref}{\usepackage{hyperref}}{}%
}{}
\@ifundefined{hyperrefpackage}{\usepackage{bookmark}}{}
\makeatother
\usepackage[toc,sort=use]{glossaries}
\makeglossaries

% Book/standalone switch.  Material that the book already covers in its
% prelude, but that a standalone paper still has to state, is wrapped as
%     \ifbookbuild\else ... \fi
% Drivers default to standalone; bookM.tex says \bookbuildtrue.
\newif\ifbookbuild
\bookbuildfalse

% Part switches for the gradostat paper.  Both default to true; a driver that
% wants the short submission says \madecfalse and/or \slowfastfalse, and every
% cross-reference into the dropped part disappears with it.
%   \ifmadec     Part III, the Lagasquie--Madec model and its break-even geometry
%   \ifslowfast  Part IV, everything carrying the small parameter \eps
\newif\ifmadec     \madectrue
\newif\ifslowfast  \slowfasttrue

% Part switch for the Barrat/simplicial-contagion project.  Defaults to
% true (the hypernetwork-synchrony part is included); the journal-paper
% driver says \hnetpartfalse to drop it and every cross-reference into it.
\newif\ifhnetpart  \hnetparttrue

%% ========================================
%% PACKAGES
%% ========================================
\RequirePackage{etoolbox}
\usepackage{makecell}
\usepackage[novbox]{pdfsync}% obsolete; SyncTeX is built into modern engines
\usepackage{amssymb, amsfonts,mathtools,bbm,eurosym,graphicx,latexsym, amsmath,amsthm,times,nccmath,url}
\usepackage{listings,verbatim,cases,tabularx}
% hyperref is deliberately NOT in the line above.  Journal classes often
% load it themselves, with a long option list, from a begindocument hook
% (mdpi.cls does: \DeclareHookRule{begindocument}{hyperref}{before}{mdpi}).
% Loading it here first, without options, makes that later load raise
% "Option clash for package hyperref".  So load it only if nobody else
% has, or will:
%   - \@ifpackageloaded  catches a class that loaded it already;
%   - \hyperrefpackage    is mdpi.cls's own deferred loader;
%   - \AAdefNnohyperref   lets any other driver opt out by defining it
%                         before % Shared preamble for the GLVbook papers.  This is the single canonical copy;
% a driver one level below GLVbook reaches it as
%
%     \input{../basicFiles/preamble}
%
% What stays in the driver, because it differs from paper to paper:
% \documentclass, font choice, \title, \author, \date, and any macro genuinely
% local to one paper.  Everything else belongs here.  The file is
% self-contained: it does not read any other file.

\usepackage[utf8]{inputenc}
\usepackage[novbox]{pdfsync}   % required: SyncTeX alone does not work here

% Page geometry, set once here for every paper that loads this file.  A driver
% that needs different margins says \geometry{...}; a second
% \usepackage[...]{geometry} raises an option clash and must not be used.
\makeatletter
\@ifundefined{AAdefNnogeometry}{\usepackage[a4paper,margin=2cm]{geometry}}{}
\makeatother

% mathematics
\usepackage{amsmath,amssymb,amsthm,mathtools,bm,mathrsfs,bbm}
\usepackage{centernot}

% floats, tables, lists
\usepackage{graphicx,float}
\makeatletter
\@ifpackageloaded{subfig}{}{\usepackage{subcaption}}
\makeatother
\usepackage{booktabs,longtable,threeparttable,array,multirow}
\usepackage{enumitem}

% colour, code, algorithms
\usepackage{xcolor,listings,tcolorbox}
\makeatletter
\@ifundefined{c@algorithm}{\usepackage{algorithm}}{}
\makeatother
\usepackage{algorithmicx,algpseudocode}

% diagrams
\usepackage{tikz,tikz-cd,pgfplots}
\usetikzlibrary{graphs,decorations.pathmorphing,decorations.markings,
arrows.meta,automata,positioning}
\tikzcdset{every label/.append style = {font = \small}}
\pgfplotsset{compat=1.18}

% chemistry
\usepackage[version=4]{mhchem}

\usepackage{microtype}
\usepackage{lineno}
\usepackage{nomencl}
\makenomenclature

% hyperref first, then bookmark, then glossaries: glossaries must follow
% hyperref or the entry links are built against the wrong anchors.
\makeatletter
\@ifundefined{hyperrefpackage}{%
  \@ifundefined{AAdefNnohyperref}{\usepackage{hyperref}}{}%
}{}
\@ifundefined{hyperrefpackage}{\usepackage{bookmark}}{}
\makeatother
\usepackage[toc,sort=use]{glossaries}
\makeglossaries

% Book/standalone switch.  Material that the book already covers in its
% prelude, but that a standalone paper still has to state, is wrapped as
%     \ifbookbuild\else ... \fi
% Drivers default to standalone; bookM.tex says \bookbuildtrue.
\newif\ifbookbuild
\bookbuildfalse

% Part switches for the gradostat paper.  Both default to true; a driver that
% wants the short submission says \madecfalse and/or \slowfastfalse, and every
% cross-reference into the dropped part disappears with it.
%   \ifmadec     Part III, the Lagasquie--Madec model and its break-even geometry
%   \ifslowfast  Part IV, everything carrying the small parameter \eps
\newif\ifmadec     \madectrue
\newif\ifslowfast  \slowfasttrue

% Part switch for the Barrat/simplicial-contagion project.  Defaults to
% true (the hypernetwork-synchrony part is included); the journal-paper
% driver says \hnetpartfalse to drop it and every cross-reference into it.
\newif\ifhnetpart  \hnetparttrue

%% ========================================
%% PACKAGES
%% ========================================
\RequirePackage{etoolbox}
\usepackage{makecell}
\usepackage[novbox]{pdfsync}% obsolete; SyncTeX is built into modern engines
\usepackage{amssymb, amsfonts,mathtools,bbm,eurosym,graphicx,latexsym, amsmath,amsthm,times,nccmath,url}
\usepackage{listings,verbatim,cases,tabularx}
% hyperref is deliberately NOT in the line above.  Journal classes often
% load it themselves, with a long option list, from a begindocument hook
% (mdpi.cls does: \DeclareHookRule{begindocument}{hyperref}{before}{mdpi}).
% Loading it here first, without options, makes that later load raise
% "Option clash for package hyperref".  So load it only if nobody else
% has, or will:
%   - \@ifpackageloaded  catches a class that loaded it already;
%   - \hyperrefpackage    is mdpi.cls's own deferred loader;
%   - \AAdefNnohyperref   lets any other driver opt out by defining it
%                         before % Shared preamble for the GLVbook papers.  This is the single canonical copy;
% a driver one level below GLVbook reaches it as
%
%     \input{../basicFiles/preamble}
%
% What stays in the driver, because it differs from paper to paper:
% \documentclass, font choice, \title, \author, \date, and any macro genuinely
% local to one paper.  Everything else belongs here.  The file is
% self-contained: it does not read any other file.

\usepackage[utf8]{inputenc}
\usepackage[novbox]{pdfsync}   % required: SyncTeX alone does not work here

% Page geometry, set once here for every paper that loads this file.  A driver
% that needs different margins says \geometry{...}; a second
% \usepackage[...]{geometry} raises an option clash and must not be used.
\makeatletter
\@ifundefined{AAdefNnogeometry}{\usepackage[a4paper,margin=2cm]{geometry}}{}
\makeatother

% mathematics
\usepackage{amsmath,amssymb,amsthm,mathtools,bm,mathrsfs,bbm}
\usepackage{centernot}

% floats, tables, lists
\usepackage{graphicx,float}
\makeatletter
\@ifpackageloaded{subfig}{}{\usepackage{subcaption}}
\makeatother
\usepackage{booktabs,longtable,threeparttable,array,multirow}
\usepackage{enumitem}

% colour, code, algorithms
\usepackage{xcolor,listings,tcolorbox}
\makeatletter
\@ifundefined{c@algorithm}{\usepackage{algorithm}}{}
\makeatother
\usepackage{algorithmicx,algpseudocode}

% diagrams
\usepackage{tikz,tikz-cd,pgfplots}
\usetikzlibrary{graphs,decorations.pathmorphing,decorations.markings,
arrows.meta,automata,positioning}
\tikzcdset{every label/.append style = {font = \small}}
\pgfplotsset{compat=1.18}

% chemistry
\usepackage[version=4]{mhchem}

\usepackage{microtype}
\usepackage{lineno}
\usepackage{nomencl}
\makenomenclature

% hyperref first, then bookmark, then glossaries: glossaries must follow
% hyperref or the entry links are built against the wrong anchors.
\makeatletter
\@ifundefined{hyperrefpackage}{%
  \@ifundefined{AAdefNnohyperref}{\usepackage{hyperref}}{}%
}{}
\@ifundefined{hyperrefpackage}{\usepackage{bookmark}}{}
\makeatother
\usepackage[toc,sort=use]{glossaries}
\makeglossaries

% Book/standalone switch.  Material that the book already covers in its
% prelude, but that a standalone paper still has to state, is wrapped as
%     \ifbookbuild\else ... \fi
% Drivers default to standalone; bookM.tex says \bookbuildtrue.
\newif\ifbookbuild
\bookbuildfalse

% Part switches for the gradostat paper.  Both default to true; a driver that
% wants the short submission says \madecfalse and/or \slowfastfalse, and every
% cross-reference into the dropped part disappears with it.
%   \ifmadec     Part III, the Lagasquie--Madec model and its break-even geometry
%   \ifslowfast  Part IV, everything carrying the small parameter \eps
\newif\ifmadec     \madectrue
\newif\ifslowfast  \slowfasttrue

% Part switch for the Barrat/simplicial-contagion project.  Defaults to
% true (the hypernetwork-synchrony part is included); the journal-paper
% driver says \hnetpartfalse to drop it and every cross-reference into it.
\newif\ifhnetpart  \hnetparttrue

%% ========================================
%% PACKAGES
%% ========================================
\RequirePackage{etoolbox}
\usepackage{makecell}
\usepackage[novbox]{pdfsync}% obsolete; SyncTeX is built into modern engines
\usepackage{amssymb, amsfonts,mathtools,bbm,eurosym,graphicx,latexsym, amsmath,amsthm,times,nccmath,url}
\usepackage{listings,verbatim,cases,tabularx}
% hyperref is deliberately NOT in the line above.  Journal classes often
% load it themselves, with a long option list, from a begindocument hook
% (mdpi.cls does: \DeclareHookRule{begindocument}{hyperref}{before}{mdpi}).
% Loading it here first, without options, makes that later load raise
% "Option clash for package hyperref".  So load it only if nobody else
% has, or will:
%   - \@ifpackageloaded  catches a class that loaded it already;
%   - \hyperrefpackage    is mdpi.cls's own deferred loader;
%   - \AAdefNnohyperref   lets any other driver opt out by defining it
%                         before \input{preamble}.
\makeatletter
\@ifpackageloaded{hyperref}{}{%
  \@ifundefined{hyperrefpackage}{%
    \@ifundefined{AAdefNnohyperref}{\usepackage{hyperref}}{}%
  }{}%
}
\makeatother
% linktoc=all makes a table-of-contents page number a link as well as the
% entry text; hyperref's default, linktoc=section, links only the text.
% Deferred to begin-document and guarded: by the block above hyperref may
% be loaded later (mdpi's own hook) or not at all (\AAdefNnohyperref), and
% \hypersetup must exist before it is called.  A driver that wants the
% default back can say \hypersetup{linktoc=section} after \input{preamble}.
\makeatletter
\AtBeginDocument{\@ifundefined{hypersetup}{}{\hypersetup{linktoc=all}}}
\makeatother
\usepackage{graphicx,float,fancybox}
\usepackage{epsfig,epstopdf,tikz}
\usepackage[sans]{dsfont}
\usetikzlibrary{positioning,arrows,arrows.meta,calc}
% Configure listings for Mathematica code
\lstdefinestyle{mathematica}{
 language=Mathematica,
 basicstyle=\small\ttfamily,
 backgroundcolor=\color{gray!10},
 frame=single,
 breaklines=true,
 commentstyle=\color{green!60!black},
 keywordstyle=\color{blue},
 stringstyle=\color{red},
 showstringspaces=false,
 tabsize=2
}
\usepackage{tcolorbox}
% Colored boxes for key results
\newtcolorbox{keyresult}[1][]{
 colback=blue!5!white,
 colframe=blue!75!black,
 title=#1,
 fonttitle=\bfseries
}
\newtcolorbox{examplebox}[1][]{
 colback=green!5!white,
 colframe=green!75!black,
 title=#1,
 fonttitle=\bfseries
}


\makeatletter
\@ifundefined{theorem}   {\newtheorem{theorem}   {Theorem}}   {}
\@ifundefined{definition}{\newtheorem{definition}{Definition}} {}
\@ifundefined{open}      {\newtheorem{open}      {Problem}}   {}
\@ifundefined{example}   {\newtheorem{example}   {Example}}   {}
\makeatother

% COUNTER DEFINITIONS
\newcounter{tabcounter}
\newcounter{defcounter}
\newcounter{opcounter}
\newcounter{thmcounter}
\newcounter{excounter}
\newcounter{ascounter}
\newcounter{cjcounter}
\newcounter{lemcounter}
\newcounter{remcounter}
\newcounter{corcounter}
\newcounter{propcounter}
\newcounter{figcounter}

% UNIFIED GLOSSARY STYLE
\newglossarystyle{unifiedstyle}{%
  \setglossarystyle{list}%
  \renewcommand*{\glossarypreamble}{\gdef\glslastsec{}}%
  \renewcommand{\glossentry}[2]{%
    \edef\glscursec{\glsentryuseriii{##1}}%
    \ifx\glscursec\glslastsec\else
      \item[]\textbf{\S\,\glscursec\enspace\glsentryuseriv{##1}}%
      \global\let\glslastsec\glscursec
    \fi
    \item[\glsentryuserii{##1} \glsentryuseri{##1}:] ``\glossentryname{##1}'' \glsentrydesc{##1}%
  }%
}

\makeatletter
% First-word extractor (used by \beT, \beD and \beXa for the default key)
% \@firstwordof hello world \relax  ->  hello
\def\@firstwordof#1 #2\relax{#1}
% First three letters, used for the heading mnemonic of a theorem given no
% explicit key.  Padding with \@empty keeps it safe for first words shorter
% than three letters: \@empty vanishes under \edef, whereas \relax would
% survive into the label and is illegal inside \csname (it broke every
% two-letter title, e.g. "ME ...", with "Missing \endcsname inserted").
\def\@takethree#1#2#3#4\@nil{#1#2#3}

% EXAMPLE COMMAND -- same signature and behaviour as \beT and \beD.
% Usage: \beXa*[long title][short title][key]
%   star  = suppresses the printed "[key]" tag after the title (as for \beT*
%   and \beD*); the label and glossary entry are unaffected.
%   key = optional label key override; the label becomes e:KEY, and without a
%   key it is e: plus the first three letters of the title's first word.  A
%   mnemonic label lets an example be cited as "Example~\ref{e:KEY}" instead
%   of by section number, which is what the citing prose needs.
\NewDocumentCommand{\beXa}{s O{} O{} o}{%
  \def\@temparg{#2}%
  \ifx\@temparg\@empty
    \begin{example}%
  \else
    \let\@savedSecTitle\@currentsectionname%
    \refstepcounter{excounter}%
    \IfValueTF{#4}{%
      \edef\@exlblkey{#4}%
    }{%
      \edef\@exfirstword{\@firstwordof#2 \relax}%
      \edef\@exlblkey{%
        \expandafter\@takethree\@exfirstword\@empty\@empty\@empty\@nil}%
    }%
    \def\@tempshort{#3}%
    \ifx\@tempshort\@empty
      \expandafter\xdef\csname @xshort@\@exlblkey\endcsname{\@exlblkey}%
    \else
      \expandafter\def\csname @xshort@\@exlblkey\endcsname{#3}%
    \fi
    \IfBooleanTF{#1}{\let\@extag\@empty}{%
      \edef\@extag{\noexpand\ \noexpand\textup{[\noexpand\texttt{e:\@exlblkey}]}}}%
    \begin{example}[#2\ifshowkeys\@extag\fi]%
    \label{e:\@exlblkey}%
    % record the long title so \nref{e:KEY} prints it, as for theorems
    \protected@write\@auxout{}{%
      \string\@thmtitle{e:\@exlblkey}{\unexpanded{#2}}}%
    \expanded{%
      \noexpand\newglossaryentry{e\theexcounter}{%
        name={\unexpanded{#2}%
          \IfBooleanF{#1}{ [\noexpand\texttt{e:\@exlblkey}]}},%
        sort={\theexample},%
        description={p.~\noexpand\pageref{e:\@exlblkey}},%
        user1={\theexample},%
        user2={Example},%
        user3={\thesection},%
        user4={\expandafter\unexpanded\expandafter{\@savedSecTitle}}%
      }%
    }%
    \protected@write\@auxout{}{\string\glsliveentry{e\theexcounter}}%
  \fi
}
\newcommand{\eeXa}{\end{example}}

% DEFINITION COMMAND -- same signature and behaviour as \beT below.
% Usage: \beD[long title][short title][key]
%   long title  = definition heading (plain text only)
%   short title = shown in \dref{word}; defaults to the key
%   key         = optional label key override; the label becomes d:KEY
% Without a key the label is d: plus the first three letters of the title's
% first word, so it is a mnemonic rather than the definition's number: a
% numeric label would silently change whenever a definition is inserted.
\NewDocumentCommand{\beD}{s O{} O{} o}{%
  \def\@temparg{#2}%
  \ifx\@temparg\@empty
    \begin{definition}%
  \else
    \let\@savedSecTitle\@currentsectionname%
    \refstepcounter{defcounter}%
    \IfValueTF{#4}{%
      \edef\@deflblkey{#4}%
    }{%
      \edef\@deffirstword{\@firstwordof#2 \relax}%
      \edef\@deflblkey{%
        \expandafter\@takethree\@deffirstword\@empty\@empty\@empty\@nil}%
    }%
    \def\@tempshort{#3}%
    \ifx\@tempshort\@empty
      \expandafter\xdef\csname @dshort@\@deflblkey\endcsname{\@deflblkey}%
    \else
      \expandafter\def\csname @dshort@\@deflblkey\endcsname{#3}%
    \fi
    \IfBooleanTF{#1}{\let\@deftag\@empty}{%
      \edef\@deftag{\noexpand\ \noexpand\textup{[\noexpand\texttt{d:\@deflblkey}]}}}%
    \begin{definition}[#2\ifshowkeys\@deftag\fi]%
    \label{d:\@deflblkey}%
    % record the long title so \nref{d:KEY} prints it, as for theorems
    \protected@write\@auxout{}{%
      \string\@thmtitle{d:\@deflblkey}{\unexpanded{#2}}}%
    \expanded{%
      \noexpand\newglossaryentry{d\thedefcounter}{%
        name={\unexpanded{#2}%
          \IfBooleanF{#1}{ [\noexpand\texttt{d:\@deflblkey}]}},%
        sort={\thedefinition},%
        description={p.~\noexpand\pageref{d:\@deflblkey}},%
        user1={\thedefinition},%
        user2={Definition},%
        user3={\thesection},%
        user4={\expandafter\unexpanded\expandafter{\@savedSecTitle}}%
      }%
    }%
    \protected@write\@auxout{}{\string\glsliveentry{d\thedefcounter}}%
  \fi
}
\newcommand{\eeD}{\end{definition}}

% \showkeysfalse in a driver hides the [t:KEY] mnemonics in result headings
% (labels and glossary entries are unchanged); used for submission builds.
\newif\ifshowkeys \showkeystrue
% \commonredtrue in a driver that inputs basicFiles/commonRed: tool.tex then
% cites its reduction results instead of restating them.
\newif\ifcommonred \commonredfalse
% TheorEM COMMAND
% Usage: \beT[long title][short title][key]
%   long title  = theorem heading (plain text only)
%   short title = shown in \tref{word}; may contain math; defaults to first word of long title
%   key         = optional label key override; label becomes t:KEY (avoids first-word collisions)
% Label is t:WORD where WORD = first word of long title (or KEY if given); use \tref{WORD} to reference.
\NewDocumentCommand{\beT}{s O{} O{} o}{%
  \def\@temparg{#2}%
  \ifx\@temparg\@empty
    \begin{theorem}%
  \else
    \let\@savedSecTitle\@currentsectionname%
    \refstepcounter{thmcounter}%
    % \@thmtag is the mnemonic shown in the heading; every theorem gets one.
    % With an explicit key the key itself is shown.  Without one the key
    % defaults to the title's first word, and showing that whole word would
    % merely repeat the title ("Dissipativity [Dissipativity]"), so only its
    % first three letters are shown: "Dissipativity [Dis]".  The label is
    % unaffected -- it still uses the full \@thmlblkey.
    \IfValueTF{#4}{%
      \edef\@thmlblkey{#4}%
    }{%
      \edef\@thmfirstword{\@firstwordof#2 \relax}%
      \edef\@thmlblkey{%
        \expandafter\@takethree\@thmfirstword\@empty\@empty\@empty\@nil}%
    }%
    \IfBooleanTF{#1}{\let\@thmtag\@empty}{%
      \edef\@thmtag{\noexpand\ \noexpand\textup{[\noexpand\texttt{t:\@thmlblkey}]}}}%
    \def\@tempshort{#3}%
    \ifx\@tempshort\@empty
      \expandafter\xdef\csname @tshort@\@thmlblkey\endcsname{\@thmlblkey}%
    \else
      \expandafter\def\csname @tshort@\@thmlblkey\endcsname{#3}%
    \fi
    \begin{theorem}[#2\ifshowkeys\@thmtag\fi]%
    \label{t:\@thmlblkey}%
    % record the long title for \nref; via the .aux so forward refs resolve
    \protected@write\@auxout{}{%
      \string\@thmtitle{t:\@thmlblkey}{\unexpanded{#2}}}%
    \expanded{%
      \noexpand\newglossaryentry{t\thethmcounter}{%
        name={\unexpanded{#2}%
          \IfBooleanF{#1}{ [\noexpand\texttt{t:\@thmlblkey}]}},%
        % sort/user1 must be the number the document actually prints, i.e. the
        % theorem environment's counter, not thmcounter.  thmcounter advances
        % only on \beT with a title, so a bare \beT (or \beT with an
        % empty title) makes the two drift and the glossary cites a wrong
        % number.  thmcounter is still fine as the entry key: it is unique.
        sort={\thetheorem},%
        description={p.~\noexpand\pageref{t:\@thmlblkey}},%
        user1={\thetheorem},%
        user2={Theorem},%
        user3={\thesection},%
        user4={\expandafter\unexpanded\expandafter{\@savedSecTitle}}%
      }%
    }%
    \protected@write\@auxout{}{\string\glsliveentry{t\thethmcounter}}%
  \fi
}
\newcommand{\eeT}{\end{theorem}}

% CROSS-PART RELATION PARAGRAPH
% Usage: \xpart{text}  -- place immediately BEFORE \beT, never inside it
%   (anything inside \beT...\eeT would be swallowed into the glossary entry).
% The paragraph is typeset only when Part I is actually present, i.e. in the
% combined book build; in the standalone builds (GranadaDSABNS/A.tex,
% AAHRah-Submission/AAHRah.tex) the label part:glv is undefined and the whole
% paragraph disappears rather than leaving ?? cross-references. Same
% graceful-degradation principle as \GASref in GASt.tex. Resolves on the
% second LaTeX pass, since r@part:glv comes from the .aux file.
% (Already inside the \makeatletter region opened above, so no \makeatletter
% pair here -- adding one would close that region early.)
\newcommand{\xpart}[1]{\@ifundefined{r@part:glv}{}{\paragraph{Relation with Parts~I--II.}#1}}

% \bookref{label}{precise text}{fallback text}: prints the precise cross-reference
% when the label resolves (combined book build), and the fallback text otherwise
% (standalone paper builds, where book-only appendix labels are absent). Same
% graceful-degradation role as \GASref in GASt.tex, but usable in any file since it
% lives in this preamble (loaded by every driver) -- unlike \GASref, which is defined in a
% content file and is therefore undefined in standalone builds that omit that file.
\newcommand{\bookref}[3]{\@ifundefined{r@#1}{#3}{#2}}

% \PVpart and \RBpart: the two parts that a partial build may omit.
% They print "Part~N" in any build that contains the part, and a naming phrase
% in a build that does not, so an excerpt shows no ?? .  Each already supplies
% the word "Part", so write "\PVpart develops ...", not "Part~\PVpart".
\newcommand{\PVpart}{\bookref{part:aahrah}%
  {Part~\ref{part:aahrah}}{the Perron--Volterra part of the full book}}
\newcommand{\RBpart}{\bookref{part:recbrowne}%
  {Part~\ref{part:recbrowne}}{the immuno-virology part of the full book}}

% OPEN PROBLEM COMMANDS
% ASSUMPTION COMMANDS (glossaried)
\newcommand{\beAs}[1][]{%
  \def\@temparg{#1}%
  \ifx\@temparg\@empty
    \begin{assumption}%
  \else
    \let\@savedSecTitle\@currentsectionname%
    \refstepcounter{ascounter}%
    \begin{assumption}[#1]%
    \label{as:\theascounter}%
    \expanded{%
      \noexpand\newglossaryentry{as\theascounter}{%
        name={\unexpanded{#1}},%
        sort={\theascounter},%
        description={p.~\noexpand\pageref{as:\theascounter}},%
        user1={\theascounter},%
        user2={Assumption},%
        user3={\thesection},%
        user4={\expandafter\unexpanded\expandafter{\@savedSecTitle}}%
      }%
    }%
    \protected@write\@auxout{}{\string\glsliveentry{as\theascounter}}%
  \fi
}
\newcommand{\eeAs}{\end{assumption}}
% CONJECTURE COMMANDS (glossaried)
\newcommand{\beCj}[1][]{%
  \def\@temparg{#1}%
  \ifx\@temparg\@empty
    \begin{conjecture}%
  \else
    \let\@savedSecTitle\@currentsectionname%
    \refstepcounter{cjcounter}%
    \begin{conjecture}[#1]%
    \label{cj:\thecjcounter}%
    \expanded{%
      \noexpand\newglossaryentry{cj\thecjcounter}{%
        name={\unexpanded{#1}},%
        sort={\thecjcounter},%
        description={p.~\noexpand\pageref{cj:\thecjcounter}},%
        user1={\thecjcounter},%
        user2={Conjecture},%
        user3={\thesection},%
        user4={\expandafter\unexpanded\expandafter{\@savedSecTitle}}%
      }%
    }%
    \protected@write\@auxout{}{\string\glsliveentry{cj\thecjcounter}}%
  \fi
}
\newcommand{\eeCj}{\end{conjecture}}
% LEMMA COMMANDS
% ---------------------------------------------------------------------
% UNIFORM RESULT WRAPPERS -- \beL, \beC, \beP, \beR share \beT's signature
%
%   \beX[long title][short title][key]  ...  \eeX
%
%   long title  = heading (math allowed)
%   short title = shown in the cross-reference; defaults to the key
%   key         = label key.  With a key the label is X:KEY and the heading
%                 carries the mnemonic [key], exactly as for \beT.  Without
%                 one the label falls back to the running counter, which is
%                 what every pre-existing call site in the book relies on.
%
% A bare \beX -- arguments commented out -- opens the plain environment and
% creates NO label and NO glossary entry.  That is the supported way to keep
% an individual result out of the catalogue.
% ---------------------------------------------------------------------
% #1 wrapper name  #2 environment  #3 label prefix  #4 counter  #5 type name
\newcommand{\@mkresult}[5]{%
  % A leading * suppresses the [key] mnemonic in both the heading and the
  % glossary name, while still making the label X:KEY and the entry:
  %     \beX*[long][short][key]
  \expandafter\NewDocumentCommand\csname be#1\endcsname{s O{} O{} o}{%
    \def\@temparg{##2}%
    \ifx\@temparg\@empty
      \begin{#2}%
    \else
      \let\@savedSecTitle\@currentsectionname%
      \refstepcounter{#4}%
      % With a key: mnemonic label X:KEY and the [key] tag in the heading.
      % Without one: fall back to the running counter, exactly as these
      % wrappers behaved before, so every legacy one-argument call site in
      % the book keeps its label.  The counter is also csname-safe, which a
      % title's first word is not -- titles beginning with math would
      % otherwise build \csname \Omega...\endcsname and fail.
      \IfValueTF{##4}{%
        \edef\@relbl{##4}%
        \IfBooleanTF{##1}{\let\@retag\@empty}{%
          \edef\@retag{\noexpand\ \noexpand\textup{[\noexpand\texttt{#3:##4}]}}}%
      }{%
        \edef\@relbl{\arabic{#4}}%
        \let\@retag\@empty
      }%
      \def\@tempshort{##3}%
      \ifx\@tempshort\@empty
        \expandafter\xdef\csname @sshort@#3\@relbl\endcsname{\@relbl}%
      \else
        \expandafter\def\csname @sshort@#3\@relbl\endcsname{##3}%
      \fi
      \begin{#2}[##2\ifshowkeys\@retag\fi]%
      \edef\@tmplbl{#3:\@relbl}\expandafter\label\expandafter{\@tmplbl}%
      \expanded{%
        \noexpand\newglossaryentry{#3\csname the#4\endcsname}{%
          name={\unexpanded{##2}%
            \IfBooleanF{##1}{\IfValueT{##4}{ [\noexpand\texttt{#3:##4}]}}},%
          sort={\csname the#2\endcsname},%
          description={p.~\noexpand\pageref{#3:\@relbl}},%
          user1={\csname the#2\endcsname},%
          user2={#5},%
          user3={\thesection},%
          user4={\expandafter\unexpanded\expandafter{\@savedSecTitle}}%
        }%
      }%
      \protected@write\@auxout{}{\string\glsliveentry{#3\csname the#4\endcsname}}%
      \protected@write\@auxout{}{\string\@thmtitle{\@tmplbl}{\unexpanded{##2}}}%
    \fi
  }%
  \expandafter\newcommand\csname ee#1\endcsname{\end{#2}}%
}
\@mkresult{L}{lemma}{l}{lemcounter}{Lemma}
% REMARK / COROLLARY / PROPOSITION -- same signature as \beL above
\@mkresult{R}{remark}{r}{remcounter}{Remark}
\@mkresult{C}{corollary}{c}{corcounter}{Corollary}
\@mkresult{P}{proposition}{p}{propcounter}{Proposition}
% OPEN-PROBLEM COMMAND -- generated with the same signature as \beT:
%   \beO[long title][short title][key]   ...   \eeO
% label o:KEY, or o:<counter> when no key is given; a bare \beO makes
% neither a label nor a glossary entry.
\@mkresult{O}{open}{o}{opcounter}{Problem}


% Section-level title tracker: updated only by \section, not \subsection.
% \@currentlabelname is overwritten by every \refstepcounter (including
% \subsection), so we save it separately here.
\newcommand{\@currentsectionname}{}
\AtBeginDocument{%
  \let\@epidcrn@orig@refstepcounter\refstepcounter
  \renewcommand{\refstepcounter}[1]{%
    \@epidcrn@orig@refstepcounter{#1}%
    \ifstrequal{#1}{section}{\global\let\@currentsectionname\@currentlabelname}{}%
  }%
}

% TABLE COMMANDS (glossaried) -- consistent naming with \beT, \beD, etc.
\newcommand{\beTa}[1][]{%
  \def\@temparg{#1}%
  \ifx\@temparg\@empty
    \begin{table}[htbp]%
  \else
    \let\@savedSecTitle\@currentsectionname%
    \refstepcounter{tabcounter}%
    \begin{table}[htbp]%
    \label{A:\thetabcounter}%
    \expanded{%
      \noexpand\newglossaryentry{A\thetabcounter}{%
        name={\unexpanded{#1}},%
        sort={\thetabcounter},%
        description={p.~\noexpand\pageref{A:\thetabcounter}},%
        user1={\thetabcounter},%
        user2={Table},%
        user3={\thesection},%
        user4={\expandafter\unexpanded\expandafter{\@savedSecTitle}}%
      }%
    }%
    \protected@write\@auxout{}{\string\glsliveentry{A\thetabcounter}}%
  \fi
}
\newcommand{\eeTa}{\end{table}}

% FIGURE COMMANDS (glossaried) -- consistent with \beT, \beR, \beTa, etc.
% Usage:
%   \beF[Glossary title]...\caption{Visible caption}\eeF  -> glossaried figure
%   \beF...\caption{Visible caption}\eeF                  -> plain [H] figure, no glossary
\newcommand{\beF}[1][]{%
  \def\@temparg{#1}%
  \ifx\@temparg\@empty
    \begin{figure}[H]%
  \else
    \begin{figure}[H]%
    \let\@savedSecTitle\@currentsectionname%
    \refstepcounter{figcounter}%
    \label{fig:\thefigcounter}%
    \expanded{%
      \noexpand\newglossaryentry{Fig\thefigcounter}{%
        name={\unexpanded{#1}},%
        sort={\thefigcounter},%
        description={p.~\noexpand\pageref{fig:\thefigcounter}},%
        user1={\thefigcounter},%
        user2={Figure},%
        user3={\thesection},%
        user4={\expandafter\unexpanded\expandafter{\@savedSecTitle}}%
      }%
    }%
    \protected@write\@auxout{}{\string\glsliveentry{Fig\thefigcounter}}%
  \fi
}
\newcommand{\eeF}{\end{figure}}
\makeatother
% LIVE-ENTRY REGISTER
% \glsaddall prints every entry the glossaries package currently knows, and
% that set includes entries restored from \jobname.glsdefs.  The package
% rewrites .glsdefs from all *currently defined* entries, so anything it
% restored is written straight back out: a deleted section kept reappearing in
% the catalogue for ever, and the .glsdefs had to be erased by hand.  The
% wrappers therefore record every entry they actually create in the .aux, and
% \printboth adds only those.  The .aux is rebuilt from scratch on each run, so
% this is self-healing: a removed result is gone after one further run.
\makeatletter
\newcommand{\@glslive}{}
\newcommand{\@glsaddlive}[1]{\ifglsentryexists{#1}{\glsadd{#1}}{}}
\newcommand{\glsliveentry}[1]{\g@addto@macro\@glslive{\@glsaddlive{#1}}}
\makeatother
% PRINT COMMAND
% No \clearpage here -- the catalogue continues on the current page.  This has
% been re-added by mistake several times; leave it out.
\makeatletter
\newcommand{\printboth}{%
  \@glslive   % add exactly the entries this document created last run
  \printglossary[style=unifiedstyle, title={0 \hspace{.22em} Definitions, Theorems, Examples and Open Problems}]%
}
\makeatother
% Named theorem reference: \tref{word} -> Thm~N(word-or-short)
% Shows stored short title if available, otherwise shows the key itself.
\newcommand{\tref}[1]{%
  Thm~\ref{t:#1}(\ifcsname @tshort@#1\endcsname\csname @tshort@#1\endcsname\else#1\fi)%
}

% Full-title theorem reference: \nref{t:KEY} prints the number and then the
% theorem's long title, so one writes "Theorem~\nref{t:blocks}" and gets
% "Theorem 3 ``Block factorization and siphons''".  Takes the complete label, like
% \ref, not the bare key like \tref.
%
% \beT writes each title to the .aux (see \@thmtitle below), so \nref resolves
% forward references from the second LaTeX pass, exactly as \ref does.  If the
% title is not yet known -- first pass, or a label that is not a \beT theorem
% -- it degrades to the number rather than producing ??.
% These two use @ in their names, and this point in the file is outside the
% \makeatletter region opened earlier (it closes at \makeatother above), so
% the region has to be reopened here.
\makeatletter
\newcommand{\@thmtitle}[2]{\global\@namedef{@tlong@#1}{#2}}
\newcommand{\nref}[1]{%
  \ref{#1}%
  \ifcsname @tlong@#1\endcsname
    \ ``\csname @tlong@#1\endcsname''%
  \fi
}
% \nm{t:KEY}: the stored long title of any \beT/\beL/\beP/\beR/\beC/\beD/\beE result,
% verbatim and without quotes, so item names follow the macro's first argument.
\newcommand{\nm}[1]{%
  \ifcsname @tlong@#1\endcsname
    \csname @tlong@#1\endcsname%
  \else
    \ref{#1}%
  \fi
}
\makeatother




%% ========================================
%% Other THEOREM ENVIRONMENTS
%% ========================================
\makeatletter
\@ifundefined{lemma}      {\newtheorem{lemma}      {Lemma}}      {}
\@ifundefined{proposition}{\newtheorem{proposition}{Proposition}} {}
\@ifundefined{corollary}  {\newtheorem{corollary}  {Corollary}}  {}
\@ifundefined{conjecture} {\newtheorem{conjecture} {Conjecture}}  {}
\@ifundefined{remark}     {\newtheorem{remark}     {Remark}}     {}
\@ifundefined{question}   {\newtheorem{question}   {Question}}   {}
\@ifundefined{assumption} {\newtheorem{assumption} {Assumption}} {}
\@ifundefined{hypothesis} {\newtheorem{hypothesis} {Hypothesis}} {}
\@ifundefined{fact}       {\newtheorem{fact}       {Fact}}       {}
\makeatother

% \citer: cite the relay-graph threshold theorem of AAHrel. If a document has
% \input the chapter/file containing t:relayRF (so the label resolves), point
% straight at it; otherwise fall back to a bare citation of AAHrel. Shared here
% (not per-paper) so it resolves consistently across every book chapter.
% \citeR: same idea, for the first theorem of AAHRah (t:Pstar).
\makeatletter
\providecommand{\citer}{%
  \@ifundefined{r@t:relayRF}%
    {\cite{AAHrel}}%
    {Theorem~\ref{t:relayRF} of \cite{AAHrel}}%
}
\providecommand{\citeR}{%
  \@ifundefined{r@t:Pstar}%
    {\cite{AAHRah}}%
    {Theorem~\ref{t:Pstar} of \cite{AAHRah}}%
}
\makeatother


% Abbreviations
\def\beQ{\begin{question}}\def\eeQ{\end{question}}
\def\beA{\begin{assumption}}\def\eeA{\end{assumption}}


%% ========================================

%% ========================================
%% COLORS & CUSTOM COMMANDS
%% ========================================
\definecolor{funccolor}{RGB}{25,25,112}
\definecolor{desccolor}{RGB}{64,64,64}
\newcommand{\func}[1]{\texttt{\textcolor{funccolor}{#1}}}
\def\eq{=}
\def\mm{-}

%% ========================================
%% MATHCAL SHORTHANDS
%% ========================================
% Project-wide convention: \mX for a capital X is always \mathcal X.
% Defined for the whole alphabet so the convention never fails to expand.
\newcommand{\mA}{\mathcal A}\newcommand{\mB}{\mathcal B}
\newcommand{\mC}{\mathcal C}\newcommand{\mD}{\mathcal D}
\newcommand{\mE}{\mathcal E}\newcommand{\mF}{\mathcal F}
\newcommand{\mG}{\mathcal G}\newcommand{\mH}{\mathcal H}
\newcommand{\mI}{\mathcal I}\newcommand{\mJ}{\mathcal J}
\newcommand{\mK}{\mathcal K}\newcommand{\mL}{\mathcal L}
\newcommand{\mM}{\mathcal M}\newcommand{\mN}{\mathcal N}
\newcommand{\mO}{\mathcal O}\newcommand{\mP}{\mathcal P}
\newcommand{\mQ}{\mathcal Q}\newcommand{\mR}{\mathcal R}
\newcommand{\mS}{\mathcal S}\newcommand{\mT}{\mathcal T}
\newcommand{\mU}{\mathcal U}\newcommand{\mV}{\mathcal V}
\newcommand{\mW}{\mathcal W}\newcommand{\mX}{\mathcal X}
\newcommand{\mY}{\mathcal Y}\newcommand{\mZ}{\mathcal Z}

%% ========================================
%% BASIC MATH ENVIRONMENTS
%% ========================================
\def\bep{\begin{pmatrix}}\def\eep{\end{pmatrix}}
\def\bev{\begin{vmatrix}}\def\eev{\end{vmatrix}}
\def\bea{\begin{eqnarray*}}\def\eea{\end{eqnarray*}}
\def\bc{\begin{cases}}\def\ec{\end{cases}}
\def\BEN{\begin{enumerate}}\def\EEN{\end{enumerate}}
\def\BI{\begin{itemize}}\def\EI{\end{itemize}}

\newcommand{\be}[1]{\begin{equation}\label{#1}}
\newcommand{\ee}{\end{equation}}
\newcommand{\beq}{\begin{eqnarray}}
\def\eeq{\end{eqnarray}}
\def\eqr{\eqref}\def\fr{\frac}\def\lbl{\label}
\def\Lra{\Longrightarrow}\def\non{\nonumber}\def\cNGM{\cite{Diek,Van,Van08}}
\def\Eq{\Leftrightarrow}\def\b{\beta}\def\uph{\red{UP TO HERE}}
%% ========================================
%% PROBABILITY & STATISTICS
%% ========================================
\renewcommand{\P}{\ensuremath{\mathbb{P}}}  % Note: overrides paragraph symbol
\newcommand{\E}{\ensuremath{\mathbb{E}}}
\newcommand{\ind}[1]{\ensuremath{\mathbbm{1}_{\left\{#1\right\}}}}
\DeclareMathOperator{\Cov}{Cov}
\DeclareMathOperator{\TV}{TV}
\newcommand{\rank}{\operatorname{rank}}

%% ========================================
%% CRN OPERATORS (organizations/ERCs/buffering/siphons framework)
%% ========================================
\def\clos{\mathrm{clos}}
\def\Supply{\mathrm{Supply}}
\def\req{\mathrm{req}}
\def\React{\mathrm{React}}
\def\Prod{\mathrm{Prod}}
\def\Gen{\mathrm{Gen}}
\def\In{\mathrm{In}}
\def\Out{\mathrm{Out}}
\def\Act{\mathrm{Act}}
\def\Min{\mathrm{Min}}

%% ========================================
%% LINE BREAKING
%% ========================================
% Last-resort stretch, used only on paragraphs TeX cannot otherwise break:
% a long \texttt{} token, a comma-separated inline math list (commas are not
% breakpoints in math), or a run of \cite/\ref boxes.  Without it those set as
% overfull \hbox and print into the margin; with it they set slightly loose,
% which is the lesser evil.  Does not affect paragraphs that already break.
\emergencystretch=3em

%% ========================================
%% NUMBER SETS
%% ========================================
\newcommand{\R}{\mathbb{R}}  % Real numbers
\newcommand{\C}{\mathbb{C}}  % Complex numbers
\newcommand{\N}{\mathbb{N}}
\newcommand{\Z}{\mathbb{Z}}
% \providecommand, not \newcommand: some papers define these in their own
% preamble after \input{preamble}, and must not clash.
\providecommand{\Q}{\mathbb{Q}}
\providecommand{\res}{\operatorname{res}}   % resultant; used by epiRes.tex
\newcommand{\MN}{\mathbb{N}}
\newcommand{\MZ}{\mathbb{Z}}
\newcommand{\realpos}{\mathbb{R}_{> 0}}
\newcommand{\realnneg}{\mathbb{R}_{\geq 0}}
\newcommand{\naturalpos}{\mathbb{N}_{>0}}
\newcommand{\Rplus}{\mathbb{R}_+}
\newcommand{\Rp}{\mathbb{R}_{\ge 0}}
\def\Rnne{\mathbb{R}_{\ge 0}}
\def\Se{\bar{S}} \def\Ie{\bar{I}}
\newcommand{\Rpos}{\mathbb{R}_{>0}}
\newcommand{\dd}{\,\mathrm d}
\newcommand{\e}{\mathrm e}
%% ========================================
%% GREEK LETTERS (2+ character abbreviations) greek
%% ========================================
\def\al{\alpha}\def\ba{\bff \alpha}
\def\w{\omega}\def\Lgn{\ga+\mu+\de}\def\om{\omega}
\def\de{\delta}\def\De{\Delta}\def\bde{\bff \delta}
\def\ep{\epsilon}\def\eps{\varepsilon}
\def\f{\varphi}  % Kept as requested
\def\ga{\gamma}\def\Ga{\Gamma}\def\G{\Gamma}\def\bg{\bff \gamma}
\def\up{\upsilon}  % NOT the breve accent: no accent commands anywhere.
                   % Named \up, not \u, so the LaTeX breve stays reachable.
\def\la{\lambda}\def\La{\Lambda}
\def\si{\sigma}\def\ka{\kappa}
\def\th{\theta}\def\vt{\vartheta}
\def\ze{\zeta}\def\Gap{\Gamma^{+}} \def\Gam{\Gamma^{-}}


%% ========================================
%% SPECIAL SHORT COMMANDS (kept as requested)
%% ========================================
\def\I{\infty}
\def\T{\widetilde}
\def\H{\hat}
\newcommand{\relint}{\operatorname{relint}}

%% ========================================
%% BOLD/VECTOR NOTATION
%% ========================================
\newcommand{\bff}[1]{{\mbox{\boldmath$#1$}}}
\newcommand\X{\boldsymbol{X}}
\def\x{\boldsymbol{x}}\def\bx{\bff x}\def\xe{\bar{\x}}
\newcommand\y{\boldsymbol{y}}\def\by{\boldsymbol{y}}\def\ye{\bar{\y}}
\def\yb{\y}\def\xb{\x}
\def\xy{\x-\y}
\newcommand\br{\boldsymbol{r}}\def\re{\bar{\r}}
\newcommand\bd{\boldsymbol{d}}
\newcommand\bQ{\boldsymbol{Q}}
\newcommand{\Yb}{{\mathbf{Y}}}
\newcommand{\Ab}{\mathbf{A}}
\newcommand{\zb}{\mathbf{z}}\def\bz{\bff 0}
\newcommand{\Zb}{\mathbf{Z}}
\newcommand{\bb}{\mathbf{b}}
\def\vbe{\vec {\bff \beta}}
\def\vb{\vec {\bff b}}
\def\v1{\vec {\bff 1}}
\def\vn{\vec \de}
\def\vz{\vec 0}
\def\vi{\overset{\rightarrow}{\mathsf i}}
\def\vv{\overset{\rightarrow}{v}}
\def\vr{\; \vec {\mathsf r}}
\newcommand{\va}{\bff a}
\providecommand{\vect}[1]{\boldsymbol{#1}}   % bold vector (moved from Gradostat.tex)
\providecommand{\one}{\mathbf 1}             % all-ones vector (moved from AliSF)
\providecommand{\spr}{s}                     % spectral bound  (moved from AliSF)
\providecommand{\Ke}{K_{\rm eff}}            % effective carrying capacity
% Critical right/left eigenvectors in the Hopf projection method (AAHOSN/New.tex)
\providecommand{\bq}{\mathbf{q}}
\providecommand{\bp}{\mathbf{p}}


%% ========================================
%% MATHCAL & MATHSCR
%% ========================================
% Mathcal
\newcommand\mc[1]{\mathcal{#1}}
\def\mC{{\mathcal C}}\def\mD{\mathcal D}\def\mE{\mathcal{E}}
\def\mF{\mathcal F}\def\mG{{\mathcal G}}\def\mH{{\mathcal H}}
\def\mI{\mathcal I}\def\mL{\mathcal{L}}\def\mM{{\mathcal M}}
\def\mN{{\mathcal N}}\def\mP{\mathcal{P}}\def\mR{{\mathcal R}}
\def\mS{{\mathcal S}}\def\mT{\mathcal{T}}\def\mV{{\mathcal V}}
\def\mK{\mathcal K}\def\mA{{\mathcal A}}\def\mU{{\mathcal U}}
% Symmetric bilinear and trilinear forms of the Taylor expansion, used in the
% first Lyapunov coefficient (AAHOSN/New.tex)
\providecommand{\Bform}{\mathcal{B}}
\providecommand{\Cform}{\mathcal{C}}

% Blackboard bold.  Kept here, not in the drivers: a \providecommand in a
% driver preamble runs before this file and silently wins, so the same name
% ends up meaning different things in different builds.
\def\IC{\mathbb{C}}\def\IN{\mathbb{N}}\def\IQ{\mathbb{Q}}
\def\IR{\mathbb{R}}\def\IZ{\mathbb{Z}}\def\IG{\mathbb{G}}
\def\II{\mathbb{I}}

% Mathscr
\newcommand{\Rs}{{\mathscr R}}
\newcommand{\Ss}{{\mathscr S}}
\newcommand{\Ns}{{\mathscr N}}
\newcommand{\Cs}{\mathscr C}


%% ========================================
%% OPERATORS & FUNCTIONS
%% ========================================
\DeclareMathOperator{\tr}{tr}
\DeclareMathOperator{\diag}{Diag}
\DeclareMathOperator{\Diag}{Diag}
\DeclareMathOperator{\intr}{\mathrm{Int}}
\DeclareMathOperator{\Inv}{Inv}
\DeclareMathOperator{\spec}{spec}
\DeclareMathOperator{\AUC}{AUC}               % area under the curve
\newcommand{\ii}{\mathrm{i}}                  % imaginary unit
\newcommand\op{\operatorname}
\newcommand{\supp}{\op{supp}}
\newcommand{\diff}{\mathop{}\mathopen{}\mathrm{d}}
\newcommand{\pdd}[1]{\frac{\partial}{\partial{#1}}}
\DeclareMathOperator{\cone}{cone}
\DeclareMathOperator{\proj}{proj}

%% ========================================
%% BRACKETS & NORMS
%% ========================================
\providecommand{\pr}[1]{\left(#1\right)}
\providecommand{\pp}[1]{\left[#1\right]}
\newcommand\croc[1]{\left\langle #1\right\rangle}
\newcommand{\norm}[1]{\|#1\|}
\newcommand{\inner}[2]{\langle #1, #2 \rangle}
% Short synonym of \inner, used by AAHOSN/New.tex as $\ip{u}{v}=\bar u^\top v$
\providecommand{\ip}[2]{\langle #1 , #2 \rangle}


%% ========================================
%% FORMATTING & UTILITIES
%% ========================================
\newcommand*{\Scale}[2][4]{\scalebox{#1}{$#2$}}
\newcommand{\te}[1]{\text{#1}}
\newcommand\steq[1]{\stackrel{\text{\rm #1.}}{=}}
\newcommand\blue[1]{\textcolor{blue}{#1}}
\newcommand{\red}{\textcolor[rgb]{1.00,0.00,0.00}}
\newcommand{\green}{\textcolor[rgb]{0.00,1.00,0.00}}

% Figure command
\newcommand{\figu}[3]{%
\begin{figure}[H]
\centering
\includegraphics[scale=#3]{#1}
\caption{#2\label{f:#1}}
\end{figure}
}

% Footnote
\long\def\symbolfootnote[#1]#2{%
\begingroup
\def\thefootnote{\fnsymbol{footnote}}\footnote[#1]{#2}%
\endgroup}
\def\fn{\symbolfootnote}

\newcommand{\dsum}{\displaystyle\sum}

%% ========================================
%% CRN & EPIDEMIOLOGY SPECIFIC
%% ========================================
% CRN notation
\def\SS{\mathcal S}
\def\CC{\mathcal C}
\def\RR{\mathcal R}

% Reproduction numbers
\def\Rto{R_{1}^{E_2}}
\def\Rot{R_{2}^{E_1}}
\def\Roc{R_{1,c}}
\def\Rtc{R_{2,c}}
\def\rd{r_{0}}
\def\sd{s_{0}}
\def\dd{D_{0}}
\def\vd{v_{0}}
\def\zd{z_{0}}
\def\se{\sE}
\def\xd{x_{0}}
%% ========================================
%% SPECIAL VALUES & EQUILIBRIA
%% ========================================
\def\eE{\bar{e}}
\def\wE{\bar{w}}
\def\sE{s_{*}}
\def\Me{\T M_{ee}}
\def\Rnne{\mathbb{R}_{\ge 0}}

\newcommand{\ones}{\mathbf 1}


%% ========================================
%% ABBREVIATIONS (alphabetical by category)
%% ========================================

% -- General mathematical terms --
\def\adm{admissible}\def\admy{admissibility}
%\def\and{antisymmetric}% disabled 2026-09-30: it overrode LaTeX's author separator \and (printed "antisymmetric" between authors)
\def\ass{assumption}
\def\cadlag{c\`adl\`ag}
\def\ch{characteristic polynomial}
\def\chp{characteristic polynomial}
\def\coe{coefficient}
\def\comp{computation}
\def\con{condition}
\def\coo{coordinate}
\def\corr{corresponding}
\def\dec{decomposition}
\def\den{denominator}
\def\dete{deterministic}
\def\eig{eigenvalue}\def\eval{eigenvalue}
\def\eigv{eigenvector}\def\eiv{eigenvector}\def\evec{eigenvector}
\def\elt{element}
\def\enn{essentially nonnegative/positivity preserving}
\def\equ{equation}
\def\etal{et al.}
\def\fe{for example}
\def\fno{from now on}
\def\for{\forall}
\def\form{formula}
\def\Fr{Furthermore, }
\def\frt{furthermore}
\def\fun{function}
\def\gene{generalization}
\def\heu{heuristic}
\def\how{however}\def\How{However, }
\def\ie{i.e. }
\def\im{\item}
\def\isd{is defined by}
\def\itc{it is easy to check that}
\def\itf{it follows that}\def\Itf{It follows that}
\def\ith{it holds that}
\def\itm{it may be checked that}\def\Itm{It may be checked that}
\def\ito{it turns out that}\def\Ito{It turns out that}\def\Itou{It turns out useful to}
\def\ma{mass action}\def\MAR{mass action representation}
\def\mat{matrix}
\def\mbc{may be checked to be}
\def\mbw{may be written as}\def\mubw{must be written as}
\def\mpr{more precisely, }\def\Mp{More precisely, }
\def\nc{non-catalytic}
\def\nne{non-negative}\def\nny{nonnegativity}
\def\no{\nonumber}
\def\noni{non-infectious}
\def\nota{notation}
\def\np{\newpage}
\def\Nt{Note that}
\def\nt{non trivial}\def\nz{non zero}
\def\oth{otherwise}
\def\Oth{On the other hand}
\def\pa{parameter}\def\para{parameter}
\def\pol{polynomial}
\def\prob{probability}
\def\Prf{{\bf  Proof:}}\def\prf{{\bf  Proof:}}
\def\prop{proportional to}
\def\qu{\quad}
\def\QED{\hfill {$\square$}\goodbreak \medskip}
\def\rep{representation}
\def\resp{respectively}
\def\satd{satisfied}\def\satg{satisfying}\def\sats{satisfies}\def\saty{satisfy}
\def\ssec{\subsection}\def\sssec{\subsubsection}
\def\sef{see for example}
\def\sksy{skew-symmetric}
\def\sub{substrate}
\def\sys{system}
\def\var{variable}
\def\wk{well-known}
\def\wl{without loss of generality}\def\wlo{w.l.o.g.}
\def\wmc{We may check that}\def\Wmc{We may check that}
\def\wmr{we may rewrite}
\def\wms{we must show that}
\def\wrt{with respect to}

% -- CRN & Chemistry --
\newcommand\CRN{chemical reaction network}
\newcommand\CRNs{chemical reaction networks}
\newcommand\RN{reaction network}
\def\ACR{absolute concentration robustness}
\def\AMSD{acyclic siphon decomposition}
\def\DAE{differential algebraic equations}
\def\FHJ{Feinberg-Horn-Jackson}
\def\GMAK{generalized mass action kinetics}
\def\MM{Michaelis-Menten}
\def\sm{stoichiometric matrix}\def\Sm{stoichiometric matrix}
\def\WRZD{weakly reversible, zero deficiency}\def\aut{autocatalytic}

% -- Epidemiology --
\def\brn{basic reproduction number}\def\bzn{basic replacement number}
\def\com{compartment}
\def\dfe{disease free equilibrium}\def\DFE{disease free equilibrium}
\def\dwt{dwell times}\def\edt{expected dwell times}
\def\EM{epidemic model}
\def\endp{endemic fixed point}\def\enp{endemic point}\def\EE{endemic point}
\def\ME{mathematical epidemiology}
\def\MEE{eco-epidemiology}\def\EEP{eco-epidemiological}
\def\ece{eco-epidemiological model}
\def\mgm{multi group model}
\def\mtS{multi-strain model}\def\mSt{multi-strain model}
\def\mSi{minimal siphon}
\def\NGM{next generation matrix}
\def\nim{new infections matrix}
\def\qem{quadratic epidemiological model}
\def\repF{invasion functions}
\def\SPF{SIR-PH-FA model}
\def\sus{susceptible}

% -- Graph theory & Linear algebra --
\def\jac{Jacobi matrix}
\def\jin{Jacobian matrix of the invasion vector field with respect to the invasion variables}
\def\jrn{Jacobian matrix of the resident vector field with respect to the resident variables}
\def\jxn{Jacobian matrix of the $\x$ vector field with respect to the $\x$ variables}
\def\Lm{Laplacian matrix}
\def\MTT{matrix tree theorem}
\def\SGT{spectral graph theory}
\def\SMW{Sherman–Morrison–Woodbury}

% -- Stability & Dynamics --
\def\absC{absorbtion CTMC}
\def\BTb{Bogdanov-Takens bifurcation}
\def\CEP{competitive exclusion partition}
\def\CTMC{continuous time Markov chain}
\def\ds{differential system}
\def\fp{fixed point}\def\bfp{boundary fixed point}\def\fbp{fixed boundary point}
\def\fbl{fixed boundary line}\def\fif{forward invariant face}
\def\gas{global asymptotic stability}\def\GAS{globally asymptotically stable}
\def\GASS{\text{GASS}}
\def\Hb{Hopf bifurcation}
\def\Hur{Hurwitz}\def\MH{Metzler-Hurwitz}\def\RH{Routh-Hurwitz}
\def\IS{\text{IS}}
\def\LAS{locally asymptotically stable}
\def\LdS{Lyapunov diagonally stable}
\def\Lf{Lyapunov function}\def\LyV{Lyapunov Volterra type function}\def\RLF{robust Lyapunov function}
\def\LyM{Lyapunov Malkin stable}
\def\PF{Perron-Frobenius}
\def\spfp{strictly positive fixed point}
\def\sta{stability}
\def\STP{sharp threshold property}\def\STT{sharp threshold theorem}

% -- Models (with \varphi explicitly instead of \f) --
\def\AB{$(A,B)$ Arino-Brauer epidemic models}
\def\ABa{$(A,\bbe,\ba,\varphi)$ balanced bilinear model}
\def\ABP{%$(A,\bbe,P,\varphi)$
\\ balanced bilinear model}
\def\AMG{Anderson May Gupta}
\newcommand\BIN{biological interaction network}
\def\Dt{Descartes type}
\def\gLV{generalized Lotka-Volterra model}
\def\IG{\text{invasion graph}}\def\IGMS{interaction graph on minimal siphons}
\def\LCP{\text{linear complementarity problem}}
\def\LV{Lotka-Volterra}\def\LVK{Lotka-Volterra-Kolmogorov}

% -- Splitting & Structure --
\def\reg{regular splitting}\def\regS{{\bf regular splitting}}
\def\RUR{rational univariate representation}
\def\bds{block-diagonal structure}\def\NGMs{next generation matrices}

% -- Miscellaneous --
\def\Ma{{Mathematica}}
\def\Mm{Max–Min}
\def\num{numerator}
\def\rap{rational approximation} \def\DSR{directed species-reaction bipartite graph}\def\PCP{polynomial complementarity problem}\def\rat{rational substitution}
\def\QSS{quasi steady state} \def\RHS{right hand side} \def\LHS{left hand side}
\def\OSN{online social network}\def\VdD{Van den Driessche}
\def\wli{$\omega$-limit} \def\rk{rich kinetics}\def\sJ{symbolic Jacobian}
\def\cV{\cite{Vas,VasSta,BSV,GSV}} \def\spf{semi-parametric semi-functional}
\def\CRW{Capasso-Ruan-Wang} \def\ABP{$(A,\bbe,P,\f)$ balanced bilinear model}
\def\ABa{%$(A,\bbe,\ba,\f)$ \\
rank-one balanced bilinear model}
\def\bbe{B%{\bff \beta}
} \def\vbe{\vec {\bff \beta}
}
\def\vb{\vec {\bff b}}
\newcommand{\Kh}{\hat{K}}          % K with hat, first mention only
\newcommand{\wS}{w^{\sigma}}       % right Perron J^\perp, full block
\newcommand{\piS}{\vec \pi^{\sigma}}    % left Perron J^\perp, full block
\newcommand{\ws}{w^{\sigma}}       % right Perron J^\perp, sub-block
\newcommand{\pis}{\pi^{\sigma}}    % left Perron J^\perp, sub-block
\newcommand{\hwS}{\hat{w}^{\sigma}}   % right Perron K, full block
\newcommand{\hpiS}{\hat{\pi}^{\sigma}} % left Perron K, full block
\newcommand{\hws}{\hat{w}^{\sigma}}   % right Perron K, sub-block
\newcommand{\hpis}{\hat{\pi}^{\sigma}} % left Perron K, sub-block
\newcommand{\twS}{\tilde{w}^{\sigma}}  % right Perron \tilde K, full block
\newcommand{\tpiS}{\tilde{\pi}^{\sigma}} % left Perron \tilde K
\newcommand{\tws}{\tilde{w}^{\sigma}}
\newcommand{\tpis}{\tilde{\pi}^{\sigma}}
\newcommand{\Mat}[2]{\mathbb{R}_{\ge 0}^{\,#1\times #2}}
\def\tJ{J_\sigma^\perp} \def\lPe{left Perron eigenvector}
% \G is \Gamma (defined at :795).  The Volterra function keeps its own name,
% so that \G no longer has two conflicting meanings in this file.
\def\GV{G}  % Volterra G(u)=u-1-\ln u
\def\PVL{Perron-Volterra Lyapunov minimal siphon ansatz}
\newcommand{\istar}[1]{i_{#1}^{*}}
%\newcommand{\sE}{s_E}
\newcommand{\vE}{v_E}
\def\PV{Perron-Volterra candidate Lyapunov function}
\def\GLV{generalized Lotka-Volterra}
\def\MERN{mathematical epidemiology reaction networks}
\def\CMM{contraction metric method}
\newcommand{\sign}{\operatorname{sign}}
\newcommand{\Var}{\operatorname{Var}} 

% ======================================================================
% IDEAS SHARED BETWEEN THE ABSTRACT AND THE MAIN RESULTS SECTION
%
% Write each idea ONCE, in the abstract; the Main Results section recalls it,
% so the two can never drift apart.
%
%   in the abstract        \ideaitem{1}{The first-order reaction terms ...}
%   in Main Results        \ideahead{1}
%
% \ideaitem stores the text globally and typesets
%   \item \hyperlink{idea:1}{<text>}
% \ideahead typesets
%   \item\hypertarget{idea:1}{}\textbf{<the same text>}
%
% The abstract must come before the Main Results input, which is the normal
% order. If a heading has to read differently from the abstract sentence (for
% instance a long abstract item summarised in one line), override it with the
% optional argument: \ideahead[short heading]{6}.
% ======================================================================
\makeatletter
\newcommand{\ideaitem}[2]{%
  \expandafter\gdef\csname @idea@#1\endcsname{#2}%
  \item\hyperlink{idea:#1}{#2}}
\newcommand{\ideahead}[2][]{%
  \item\hypertarget{idea:#2}{}%
  \def\@ideatmp{#1}%
  \ifx\@ideatmp\@empty
    \@ifundefined{@idea@#2}{%
      \PackageWarning{preamble}{Idea #2 is used in Main Results but was never
        defined by \string\ideaitem\space in the abstract}%
      \textbf{[idea #2 undefined]}%
    }{\textbf{\begingroup\def\ideaterm##1##2{##2}\def\idearest##1{}\@nameuse{@idea@#2}\endgroup}}%
  \else
    \textbf{#1}%
  \fi}
% \ideaterm{<hypertarget>}{<words>}: inside an abstract idea the words become an
% invisible link (the document sets \hypersetup{hidelinks}) jumping to the place
% in Main Results that defines them; inside the Main Results heading itself,
% where the target sits a few lines below, \ideahead strips the link and prints
% the words plainly. Mark the destination with \hypertarget{<name>}{} on the line
% before the display, which adds no equation number and shows nothing.
\newcommand{\ideaterm}[2]{\hyperlink{#1}{#2}}
% \idearest{<text>}: the part of an abstract idea after its first sentence; printed in
% the abstract, omitted from the Main Results heading.
\newcommand{\idearest}[1]{#1}
\makeatother

% ======================================================================
% SHORT-PAPER SWITCH
% \shorttrue in a master that assembles only part of the companion paper
% (Barrat/SipArxiv.tex). \nsh{...} drops text that cites results living in the
% omitted sections, so the short build has no dangling references while the
% companion paper keeps the text verbatim.
% ======================================================================
\newif\ifshort \shortfalse
\newcommand{\nsh}[1]{\ifshort\else#1\fi}
% \shortalt{<short build>}{<companion paper>}: use when a sentence needs
% different wording in the two builds, not merely a deletion.
\newcommand{\shortalt}[2]{\ifshort#1\else#2\fi}
% \Jtr, \JtrS{\Sigma}: transverse Jacobian. J^\perp here; SipArxiv.tex redefines them as J_{\mathrm{tra}}, J_{\mathrm{tra},\Sigma}
\newcommand{\Jtr}{J^\perp}
\newcommand{\JtrS}[1]{J_{#1}^\perp}
.
\makeatletter
\@ifpackageloaded{hyperref}{}{%
  \@ifundefined{hyperrefpackage}{%
    \@ifundefined{AAdefNnohyperref}{\usepackage{hyperref}}{}%
  }{}%
}
\makeatother
% linktoc=all makes a table-of-contents page number a link as well as the
% entry text; hyperref's default, linktoc=section, links only the text.
% Deferred to begin-document and guarded: by the block above hyperref may
% be loaded later (mdpi's own hook) or not at all (\AAdefNnohyperref), and
% \hypersetup must exist before it is called.  A driver that wants the
% default back can say \hypersetup{linktoc=section} after % Shared preamble for the GLVbook papers.  This is the single canonical copy;
% a driver one level below GLVbook reaches it as
%
%     \input{../basicFiles/preamble}
%
% What stays in the driver, because it differs from paper to paper:
% \documentclass, font choice, \title, \author, \date, and any macro genuinely
% local to one paper.  Everything else belongs here.  The file is
% self-contained: it does not read any other file.

\usepackage[utf8]{inputenc}
\usepackage[novbox]{pdfsync}   % required: SyncTeX alone does not work here

% Page geometry, set once here for every paper that loads this file.  A driver
% that needs different margins says \geometry{...}; a second
% \usepackage[...]{geometry} raises an option clash and must not be used.
\makeatletter
\@ifundefined{AAdefNnogeometry}{\usepackage[a4paper,margin=2cm]{geometry}}{}
\makeatother

% mathematics
\usepackage{amsmath,amssymb,amsthm,mathtools,bm,mathrsfs,bbm}
\usepackage{centernot}

% floats, tables, lists
\usepackage{graphicx,float}
\makeatletter
\@ifpackageloaded{subfig}{}{\usepackage{subcaption}}
\makeatother
\usepackage{booktabs,longtable,threeparttable,array,multirow}
\usepackage{enumitem}

% colour, code, algorithms
\usepackage{xcolor,listings,tcolorbox}
\makeatletter
\@ifundefined{c@algorithm}{\usepackage{algorithm}}{}
\makeatother
\usepackage{algorithmicx,algpseudocode}

% diagrams
\usepackage{tikz,tikz-cd,pgfplots}
\usetikzlibrary{graphs,decorations.pathmorphing,decorations.markings,
arrows.meta,automata,positioning}
\tikzcdset{every label/.append style = {font = \small}}
\pgfplotsset{compat=1.18}

% chemistry
\usepackage[version=4]{mhchem}

\usepackage{microtype}
\usepackage{lineno}
\usepackage{nomencl}
\makenomenclature

% hyperref first, then bookmark, then glossaries: glossaries must follow
% hyperref or the entry links are built against the wrong anchors.
\makeatletter
\@ifundefined{hyperrefpackage}{%
  \@ifundefined{AAdefNnohyperref}{\usepackage{hyperref}}{}%
}{}
\@ifundefined{hyperrefpackage}{\usepackage{bookmark}}{}
\makeatother
\usepackage[toc,sort=use]{glossaries}
\makeglossaries

% Book/standalone switch.  Material that the book already covers in its
% prelude, but that a standalone paper still has to state, is wrapped as
%     \ifbookbuild\else ... \fi
% Drivers default to standalone; bookM.tex says \bookbuildtrue.
\newif\ifbookbuild
\bookbuildfalse

% Part switches for the gradostat paper.  Both default to true; a driver that
% wants the short submission says \madecfalse and/or \slowfastfalse, and every
% cross-reference into the dropped part disappears with it.
%   \ifmadec     Part III, the Lagasquie--Madec model and its break-even geometry
%   \ifslowfast  Part IV, everything carrying the small parameter \eps
\newif\ifmadec     \madectrue
\newif\ifslowfast  \slowfasttrue

% Part switch for the Barrat/simplicial-contagion project.  Defaults to
% true (the hypernetwork-synchrony part is included); the journal-paper
% driver says \hnetpartfalse to drop it and every cross-reference into it.
\newif\ifhnetpart  \hnetparttrue

%% ========================================
%% PACKAGES
%% ========================================
\RequirePackage{etoolbox}
\usepackage{makecell}
\usepackage[novbox]{pdfsync}% obsolete; SyncTeX is built into modern engines
\usepackage{amssymb, amsfonts,mathtools,bbm,eurosym,graphicx,latexsym, amsmath,amsthm,times,nccmath,url}
\usepackage{listings,verbatim,cases,tabularx}
% hyperref is deliberately NOT in the line above.  Journal classes often
% load it themselves, with a long option list, from a begindocument hook
% (mdpi.cls does: \DeclareHookRule{begindocument}{hyperref}{before}{mdpi}).
% Loading it here first, without options, makes that later load raise
% "Option clash for package hyperref".  So load it only if nobody else
% has, or will:
%   - \@ifpackageloaded  catches a class that loaded it already;
%   - \hyperrefpackage    is mdpi.cls's own deferred loader;
%   - \AAdefNnohyperref   lets any other driver opt out by defining it
%                         before \input{preamble}.
\makeatletter
\@ifpackageloaded{hyperref}{}{%
  \@ifundefined{hyperrefpackage}{%
    \@ifundefined{AAdefNnohyperref}{\usepackage{hyperref}}{}%
  }{}%
}
\makeatother
% linktoc=all makes a table-of-contents page number a link as well as the
% entry text; hyperref's default, linktoc=section, links only the text.
% Deferred to begin-document and guarded: by the block above hyperref may
% be loaded later (mdpi's own hook) or not at all (\AAdefNnohyperref), and
% \hypersetup must exist before it is called.  A driver that wants the
% default back can say \hypersetup{linktoc=section} after \input{preamble}.
\makeatletter
\AtBeginDocument{\@ifundefined{hypersetup}{}{\hypersetup{linktoc=all}}}
\makeatother
\usepackage{graphicx,float,fancybox}
\usepackage{epsfig,epstopdf,tikz}
\usepackage[sans]{dsfont}
\usetikzlibrary{positioning,arrows,arrows.meta,calc}
% Configure listings for Mathematica code
\lstdefinestyle{mathematica}{
 language=Mathematica,
 basicstyle=\small\ttfamily,
 backgroundcolor=\color{gray!10},
 frame=single,
 breaklines=true,
 commentstyle=\color{green!60!black},
 keywordstyle=\color{blue},
 stringstyle=\color{red},
 showstringspaces=false,
 tabsize=2
}
\usepackage{tcolorbox}
% Colored boxes for key results
\newtcolorbox{keyresult}[1][]{
 colback=blue!5!white,
 colframe=blue!75!black,
 title=#1,
 fonttitle=\bfseries
}
\newtcolorbox{examplebox}[1][]{
 colback=green!5!white,
 colframe=green!75!black,
 title=#1,
 fonttitle=\bfseries
}


\makeatletter
\@ifundefined{theorem}   {\newtheorem{theorem}   {Theorem}}   {}
\@ifundefined{definition}{\newtheorem{definition}{Definition}} {}
\@ifundefined{open}      {\newtheorem{open}      {Problem}}   {}
\@ifundefined{example}   {\newtheorem{example}   {Example}}   {}
\makeatother

% COUNTER DEFINITIONS
\newcounter{tabcounter}
\newcounter{defcounter}
\newcounter{opcounter}
\newcounter{thmcounter}
\newcounter{excounter}
\newcounter{ascounter}
\newcounter{cjcounter}
\newcounter{lemcounter}
\newcounter{remcounter}
\newcounter{corcounter}
\newcounter{propcounter}
\newcounter{figcounter}

% UNIFIED GLOSSARY STYLE
\newglossarystyle{unifiedstyle}{%
  \setglossarystyle{list}%
  \renewcommand*{\glossarypreamble}{\gdef\glslastsec{}}%
  \renewcommand{\glossentry}[2]{%
    \edef\glscursec{\glsentryuseriii{##1}}%
    \ifx\glscursec\glslastsec\else
      \item[]\textbf{\S\,\glscursec\enspace\glsentryuseriv{##1}}%
      \global\let\glslastsec\glscursec
    \fi
    \item[\glsentryuserii{##1} \glsentryuseri{##1}:] ``\glossentryname{##1}'' \glsentrydesc{##1}%
  }%
}

\makeatletter
% First-word extractor (used by \beT, \beD and \beXa for the default key)
% \@firstwordof hello world \relax  ->  hello
\def\@firstwordof#1 #2\relax{#1}
% First three letters, used for the heading mnemonic of a theorem given no
% explicit key.  Padding with \@empty keeps it safe for first words shorter
% than three letters: \@empty vanishes under \edef, whereas \relax would
% survive into the label and is illegal inside \csname (it broke every
% two-letter title, e.g. "ME ...", with "Missing \endcsname inserted").
\def\@takethree#1#2#3#4\@nil{#1#2#3}

% EXAMPLE COMMAND -- same signature and behaviour as \beT and \beD.
% Usage: \beXa*[long title][short title][key]
%   star  = suppresses the printed "[key]" tag after the title (as for \beT*
%   and \beD*); the label and glossary entry are unaffected.
%   key = optional label key override; the label becomes e:KEY, and without a
%   key it is e: plus the first three letters of the title's first word.  A
%   mnemonic label lets an example be cited as "Example~\ref{e:KEY}" instead
%   of by section number, which is what the citing prose needs.
\NewDocumentCommand{\beXa}{s O{} O{} o}{%
  \def\@temparg{#2}%
  \ifx\@temparg\@empty
    \begin{example}%
  \else
    \let\@savedSecTitle\@currentsectionname%
    \refstepcounter{excounter}%
    \IfValueTF{#4}{%
      \edef\@exlblkey{#4}%
    }{%
      \edef\@exfirstword{\@firstwordof#2 \relax}%
      \edef\@exlblkey{%
        \expandafter\@takethree\@exfirstword\@empty\@empty\@empty\@nil}%
    }%
    \def\@tempshort{#3}%
    \ifx\@tempshort\@empty
      \expandafter\xdef\csname @xshort@\@exlblkey\endcsname{\@exlblkey}%
    \else
      \expandafter\def\csname @xshort@\@exlblkey\endcsname{#3}%
    \fi
    \IfBooleanTF{#1}{\let\@extag\@empty}{%
      \edef\@extag{\noexpand\ \noexpand\textup{[\noexpand\texttt{e:\@exlblkey}]}}}%
    \begin{example}[#2\ifshowkeys\@extag\fi]%
    \label{e:\@exlblkey}%
    % record the long title so \nref{e:KEY} prints it, as for theorems
    \protected@write\@auxout{}{%
      \string\@thmtitle{e:\@exlblkey}{\unexpanded{#2}}}%
    \expanded{%
      \noexpand\newglossaryentry{e\theexcounter}{%
        name={\unexpanded{#2}%
          \IfBooleanF{#1}{ [\noexpand\texttt{e:\@exlblkey}]}},%
        sort={\theexample},%
        description={p.~\noexpand\pageref{e:\@exlblkey}},%
        user1={\theexample},%
        user2={Example},%
        user3={\thesection},%
        user4={\expandafter\unexpanded\expandafter{\@savedSecTitle}}%
      }%
    }%
    \protected@write\@auxout{}{\string\glsliveentry{e\theexcounter}}%
  \fi
}
\newcommand{\eeXa}{\end{example}}

% DEFINITION COMMAND -- same signature and behaviour as \beT below.
% Usage: \beD[long title][short title][key]
%   long title  = definition heading (plain text only)
%   short title = shown in \dref{word}; defaults to the key
%   key         = optional label key override; the label becomes d:KEY
% Without a key the label is d: plus the first three letters of the title's
% first word, so it is a mnemonic rather than the definition's number: a
% numeric label would silently change whenever a definition is inserted.
\NewDocumentCommand{\beD}{s O{} O{} o}{%
  \def\@temparg{#2}%
  \ifx\@temparg\@empty
    \begin{definition}%
  \else
    \let\@savedSecTitle\@currentsectionname%
    \refstepcounter{defcounter}%
    \IfValueTF{#4}{%
      \edef\@deflblkey{#4}%
    }{%
      \edef\@deffirstword{\@firstwordof#2 \relax}%
      \edef\@deflblkey{%
        \expandafter\@takethree\@deffirstword\@empty\@empty\@empty\@nil}%
    }%
    \def\@tempshort{#3}%
    \ifx\@tempshort\@empty
      \expandafter\xdef\csname @dshort@\@deflblkey\endcsname{\@deflblkey}%
    \else
      \expandafter\def\csname @dshort@\@deflblkey\endcsname{#3}%
    \fi
    \IfBooleanTF{#1}{\let\@deftag\@empty}{%
      \edef\@deftag{\noexpand\ \noexpand\textup{[\noexpand\texttt{d:\@deflblkey}]}}}%
    \begin{definition}[#2\ifshowkeys\@deftag\fi]%
    \label{d:\@deflblkey}%
    % record the long title so \nref{d:KEY} prints it, as for theorems
    \protected@write\@auxout{}{%
      \string\@thmtitle{d:\@deflblkey}{\unexpanded{#2}}}%
    \expanded{%
      \noexpand\newglossaryentry{d\thedefcounter}{%
        name={\unexpanded{#2}%
          \IfBooleanF{#1}{ [\noexpand\texttt{d:\@deflblkey}]}},%
        sort={\thedefinition},%
        description={p.~\noexpand\pageref{d:\@deflblkey}},%
        user1={\thedefinition},%
        user2={Definition},%
        user3={\thesection},%
        user4={\expandafter\unexpanded\expandafter{\@savedSecTitle}}%
      }%
    }%
    \protected@write\@auxout{}{\string\glsliveentry{d\thedefcounter}}%
  \fi
}
\newcommand{\eeD}{\end{definition}}

% \showkeysfalse in a driver hides the [t:KEY] mnemonics in result headings
% (labels and glossary entries are unchanged); used for submission builds.
\newif\ifshowkeys \showkeystrue
% \commonredtrue in a driver that inputs basicFiles/commonRed: tool.tex then
% cites its reduction results instead of restating them.
\newif\ifcommonred \commonredfalse
% TheorEM COMMAND
% Usage: \beT[long title][short title][key]
%   long title  = theorem heading (plain text only)
%   short title = shown in \tref{word}; may contain math; defaults to first word of long title
%   key         = optional label key override; label becomes t:KEY (avoids first-word collisions)
% Label is t:WORD where WORD = first word of long title (or KEY if given); use \tref{WORD} to reference.
\NewDocumentCommand{\beT}{s O{} O{} o}{%
  \def\@temparg{#2}%
  \ifx\@temparg\@empty
    \begin{theorem}%
  \else
    \let\@savedSecTitle\@currentsectionname%
    \refstepcounter{thmcounter}%
    % \@thmtag is the mnemonic shown in the heading; every theorem gets one.
    % With an explicit key the key itself is shown.  Without one the key
    % defaults to the title's first word, and showing that whole word would
    % merely repeat the title ("Dissipativity [Dissipativity]"), so only its
    % first three letters are shown: "Dissipativity [Dis]".  The label is
    % unaffected -- it still uses the full \@thmlblkey.
    \IfValueTF{#4}{%
      \edef\@thmlblkey{#4}%
    }{%
      \edef\@thmfirstword{\@firstwordof#2 \relax}%
      \edef\@thmlblkey{%
        \expandafter\@takethree\@thmfirstword\@empty\@empty\@empty\@nil}%
    }%
    \IfBooleanTF{#1}{\let\@thmtag\@empty}{%
      \edef\@thmtag{\noexpand\ \noexpand\textup{[\noexpand\texttt{t:\@thmlblkey}]}}}%
    \def\@tempshort{#3}%
    \ifx\@tempshort\@empty
      \expandafter\xdef\csname @tshort@\@thmlblkey\endcsname{\@thmlblkey}%
    \else
      \expandafter\def\csname @tshort@\@thmlblkey\endcsname{#3}%
    \fi
    \begin{theorem}[#2\ifshowkeys\@thmtag\fi]%
    \label{t:\@thmlblkey}%
    % record the long title for \nref; via the .aux so forward refs resolve
    \protected@write\@auxout{}{%
      \string\@thmtitle{t:\@thmlblkey}{\unexpanded{#2}}}%
    \expanded{%
      \noexpand\newglossaryentry{t\thethmcounter}{%
        name={\unexpanded{#2}%
          \IfBooleanF{#1}{ [\noexpand\texttt{t:\@thmlblkey}]}},%
        % sort/user1 must be the number the document actually prints, i.e. the
        % theorem environment's counter, not thmcounter.  thmcounter advances
        % only on \beT with a title, so a bare \beT (or \beT with an
        % empty title) makes the two drift and the glossary cites a wrong
        % number.  thmcounter is still fine as the entry key: it is unique.
        sort={\thetheorem},%
        description={p.~\noexpand\pageref{t:\@thmlblkey}},%
        user1={\thetheorem},%
        user2={Theorem},%
        user3={\thesection},%
        user4={\expandafter\unexpanded\expandafter{\@savedSecTitle}}%
      }%
    }%
    \protected@write\@auxout{}{\string\glsliveentry{t\thethmcounter}}%
  \fi
}
\newcommand{\eeT}{\end{theorem}}

% CROSS-PART RELATION PARAGRAPH
% Usage: \xpart{text}  -- place immediately BEFORE \beT, never inside it
%   (anything inside \beT...\eeT would be swallowed into the glossary entry).
% The paragraph is typeset only when Part I is actually present, i.e. in the
% combined book build; in the standalone builds (GranadaDSABNS/A.tex,
% AAHRah-Submission/AAHRah.tex) the label part:glv is undefined and the whole
% paragraph disappears rather than leaving ?? cross-references. Same
% graceful-degradation principle as \GASref in GASt.tex. Resolves on the
% second LaTeX pass, since r@part:glv comes from the .aux file.
% (Already inside the \makeatletter region opened above, so no \makeatletter
% pair here -- adding one would close that region early.)
\newcommand{\xpart}[1]{\@ifundefined{r@part:glv}{}{\paragraph{Relation with Parts~I--II.}#1}}

% \bookref{label}{precise text}{fallback text}: prints the precise cross-reference
% when the label resolves (combined book build), and the fallback text otherwise
% (standalone paper builds, where book-only appendix labels are absent). Same
% graceful-degradation role as \GASref in GASt.tex, but usable in any file since it
% lives in this preamble (loaded by every driver) -- unlike \GASref, which is defined in a
% content file and is therefore undefined in standalone builds that omit that file.
\newcommand{\bookref}[3]{\@ifundefined{r@#1}{#3}{#2}}

% \PVpart and \RBpart: the two parts that a partial build may omit.
% They print "Part~N" in any build that contains the part, and a naming phrase
% in a build that does not, so an excerpt shows no ?? .  Each already supplies
% the word "Part", so write "\PVpart develops ...", not "Part~\PVpart".
\newcommand{\PVpart}{\bookref{part:aahrah}%
  {Part~\ref{part:aahrah}}{the Perron--Volterra part of the full book}}
\newcommand{\RBpart}{\bookref{part:recbrowne}%
  {Part~\ref{part:recbrowne}}{the immuno-virology part of the full book}}

% OPEN PROBLEM COMMANDS
% ASSUMPTION COMMANDS (glossaried)
\newcommand{\beAs}[1][]{%
  \def\@temparg{#1}%
  \ifx\@temparg\@empty
    \begin{assumption}%
  \else
    \let\@savedSecTitle\@currentsectionname%
    \refstepcounter{ascounter}%
    \begin{assumption}[#1]%
    \label{as:\theascounter}%
    \expanded{%
      \noexpand\newglossaryentry{as\theascounter}{%
        name={\unexpanded{#1}},%
        sort={\theascounter},%
        description={p.~\noexpand\pageref{as:\theascounter}},%
        user1={\theascounter},%
        user2={Assumption},%
        user3={\thesection},%
        user4={\expandafter\unexpanded\expandafter{\@savedSecTitle}}%
      }%
    }%
    \protected@write\@auxout{}{\string\glsliveentry{as\theascounter}}%
  \fi
}
\newcommand{\eeAs}{\end{assumption}}
% CONJECTURE COMMANDS (glossaried)
\newcommand{\beCj}[1][]{%
  \def\@temparg{#1}%
  \ifx\@temparg\@empty
    \begin{conjecture}%
  \else
    \let\@savedSecTitle\@currentsectionname%
    \refstepcounter{cjcounter}%
    \begin{conjecture}[#1]%
    \label{cj:\thecjcounter}%
    \expanded{%
      \noexpand\newglossaryentry{cj\thecjcounter}{%
        name={\unexpanded{#1}},%
        sort={\thecjcounter},%
        description={p.~\noexpand\pageref{cj:\thecjcounter}},%
        user1={\thecjcounter},%
        user2={Conjecture},%
        user3={\thesection},%
        user4={\expandafter\unexpanded\expandafter{\@savedSecTitle}}%
      }%
    }%
    \protected@write\@auxout{}{\string\glsliveentry{cj\thecjcounter}}%
  \fi
}
\newcommand{\eeCj}{\end{conjecture}}
% LEMMA COMMANDS
% ---------------------------------------------------------------------
% UNIFORM RESULT WRAPPERS -- \beL, \beC, \beP, \beR share \beT's signature
%
%   \beX[long title][short title][key]  ...  \eeX
%
%   long title  = heading (math allowed)
%   short title = shown in the cross-reference; defaults to the key
%   key         = label key.  With a key the label is X:KEY and the heading
%                 carries the mnemonic [key], exactly as for \beT.  Without
%                 one the label falls back to the running counter, which is
%                 what every pre-existing call site in the book relies on.
%
% A bare \beX -- arguments commented out -- opens the plain environment and
% creates NO label and NO glossary entry.  That is the supported way to keep
% an individual result out of the catalogue.
% ---------------------------------------------------------------------
% #1 wrapper name  #2 environment  #3 label prefix  #4 counter  #5 type name
\newcommand{\@mkresult}[5]{%
  % A leading * suppresses the [key] mnemonic in both the heading and the
  % glossary name, while still making the label X:KEY and the entry:
  %     \beX*[long][short][key]
  \expandafter\NewDocumentCommand\csname be#1\endcsname{s O{} O{} o}{%
    \def\@temparg{##2}%
    \ifx\@temparg\@empty
      \begin{#2}%
    \else
      \let\@savedSecTitle\@currentsectionname%
      \refstepcounter{#4}%
      % With a key: mnemonic label X:KEY and the [key] tag in the heading.
      % Without one: fall back to the running counter, exactly as these
      % wrappers behaved before, so every legacy one-argument call site in
      % the book keeps its label.  The counter is also csname-safe, which a
      % title's first word is not -- titles beginning with math would
      % otherwise build \csname \Omega...\endcsname and fail.
      \IfValueTF{##4}{%
        \edef\@relbl{##4}%
        \IfBooleanTF{##1}{\let\@retag\@empty}{%
          \edef\@retag{\noexpand\ \noexpand\textup{[\noexpand\texttt{#3:##4}]}}}%
      }{%
        \edef\@relbl{\arabic{#4}}%
        \let\@retag\@empty
      }%
      \def\@tempshort{##3}%
      \ifx\@tempshort\@empty
        \expandafter\xdef\csname @sshort@#3\@relbl\endcsname{\@relbl}%
      \else
        \expandafter\def\csname @sshort@#3\@relbl\endcsname{##3}%
      \fi
      \begin{#2}[##2\ifshowkeys\@retag\fi]%
      \edef\@tmplbl{#3:\@relbl}\expandafter\label\expandafter{\@tmplbl}%
      \expanded{%
        \noexpand\newglossaryentry{#3\csname the#4\endcsname}{%
          name={\unexpanded{##2}%
            \IfBooleanF{##1}{\IfValueT{##4}{ [\noexpand\texttt{#3:##4}]}}},%
          sort={\csname the#2\endcsname},%
          description={p.~\noexpand\pageref{#3:\@relbl}},%
          user1={\csname the#2\endcsname},%
          user2={#5},%
          user3={\thesection},%
          user4={\expandafter\unexpanded\expandafter{\@savedSecTitle}}%
        }%
      }%
      \protected@write\@auxout{}{\string\glsliveentry{#3\csname the#4\endcsname}}%
      \protected@write\@auxout{}{\string\@thmtitle{\@tmplbl}{\unexpanded{##2}}}%
    \fi
  }%
  \expandafter\newcommand\csname ee#1\endcsname{\end{#2}}%
}
\@mkresult{L}{lemma}{l}{lemcounter}{Lemma}
% REMARK / COROLLARY / PROPOSITION -- same signature as \beL above
\@mkresult{R}{remark}{r}{remcounter}{Remark}
\@mkresult{C}{corollary}{c}{corcounter}{Corollary}
\@mkresult{P}{proposition}{p}{propcounter}{Proposition}
% OPEN-PROBLEM COMMAND -- generated with the same signature as \beT:
%   \beO[long title][short title][key]   ...   \eeO
% label o:KEY, or o:<counter> when no key is given; a bare \beO makes
% neither a label nor a glossary entry.
\@mkresult{O}{open}{o}{opcounter}{Problem}


% Section-level title tracker: updated only by \section, not \subsection.
% \@currentlabelname is overwritten by every \refstepcounter (including
% \subsection), so we save it separately here.
\newcommand{\@currentsectionname}{}
\AtBeginDocument{%
  \let\@epidcrn@orig@refstepcounter\refstepcounter
  \renewcommand{\refstepcounter}[1]{%
    \@epidcrn@orig@refstepcounter{#1}%
    \ifstrequal{#1}{section}{\global\let\@currentsectionname\@currentlabelname}{}%
  }%
}

% TABLE COMMANDS (glossaried) -- consistent naming with \beT, \beD, etc.
\newcommand{\beTa}[1][]{%
  \def\@temparg{#1}%
  \ifx\@temparg\@empty
    \begin{table}[htbp]%
  \else
    \let\@savedSecTitle\@currentsectionname%
    \refstepcounter{tabcounter}%
    \begin{table}[htbp]%
    \label{A:\thetabcounter}%
    \expanded{%
      \noexpand\newglossaryentry{A\thetabcounter}{%
        name={\unexpanded{#1}},%
        sort={\thetabcounter},%
        description={p.~\noexpand\pageref{A:\thetabcounter}},%
        user1={\thetabcounter},%
        user2={Table},%
        user3={\thesection},%
        user4={\expandafter\unexpanded\expandafter{\@savedSecTitle}}%
      }%
    }%
    \protected@write\@auxout{}{\string\glsliveentry{A\thetabcounter}}%
  \fi
}
\newcommand{\eeTa}{\end{table}}

% FIGURE COMMANDS (glossaried) -- consistent with \beT, \beR, \beTa, etc.
% Usage:
%   \beF[Glossary title]...\caption{Visible caption}\eeF  -> glossaried figure
%   \beF...\caption{Visible caption}\eeF                  -> plain [H] figure, no glossary
\newcommand{\beF}[1][]{%
  \def\@temparg{#1}%
  \ifx\@temparg\@empty
    \begin{figure}[H]%
  \else
    \begin{figure}[H]%
    \let\@savedSecTitle\@currentsectionname%
    \refstepcounter{figcounter}%
    \label{fig:\thefigcounter}%
    \expanded{%
      \noexpand\newglossaryentry{Fig\thefigcounter}{%
        name={\unexpanded{#1}},%
        sort={\thefigcounter},%
        description={p.~\noexpand\pageref{fig:\thefigcounter}},%
        user1={\thefigcounter},%
        user2={Figure},%
        user3={\thesection},%
        user4={\expandafter\unexpanded\expandafter{\@savedSecTitle}}%
      }%
    }%
    \protected@write\@auxout{}{\string\glsliveentry{Fig\thefigcounter}}%
  \fi
}
\newcommand{\eeF}{\end{figure}}
\makeatother
% LIVE-ENTRY REGISTER
% \glsaddall prints every entry the glossaries package currently knows, and
% that set includes entries restored from \jobname.glsdefs.  The package
% rewrites .glsdefs from all *currently defined* entries, so anything it
% restored is written straight back out: a deleted section kept reappearing in
% the catalogue for ever, and the .glsdefs had to be erased by hand.  The
% wrappers therefore record every entry they actually create in the .aux, and
% \printboth adds only those.  The .aux is rebuilt from scratch on each run, so
% this is self-healing: a removed result is gone after one further run.
\makeatletter
\newcommand{\@glslive}{}
\newcommand{\@glsaddlive}[1]{\ifglsentryexists{#1}{\glsadd{#1}}{}}
\newcommand{\glsliveentry}[1]{\g@addto@macro\@glslive{\@glsaddlive{#1}}}
\makeatother
% PRINT COMMAND
% No \clearpage here -- the catalogue continues on the current page.  This has
% been re-added by mistake several times; leave it out.
\makeatletter
\newcommand{\printboth}{%
  \@glslive   % add exactly the entries this document created last run
  \printglossary[style=unifiedstyle, title={0 \hspace{.22em} Definitions, Theorems, Examples and Open Problems}]%
}
\makeatother
% Named theorem reference: \tref{word} -> Thm~N(word-or-short)
% Shows stored short title if available, otherwise shows the key itself.
\newcommand{\tref}[1]{%
  Thm~\ref{t:#1}(\ifcsname @tshort@#1\endcsname\csname @tshort@#1\endcsname\else#1\fi)%
}

% Full-title theorem reference: \nref{t:KEY} prints the number and then the
% theorem's long title, so one writes "Theorem~\nref{t:blocks}" and gets
% "Theorem 3 ``Block factorization and siphons''".  Takes the complete label, like
% \ref, not the bare key like \tref.
%
% \beT writes each title to the .aux (see \@thmtitle below), so \nref resolves
% forward references from the second LaTeX pass, exactly as \ref does.  If the
% title is not yet known -- first pass, or a label that is not a \beT theorem
% -- it degrades to the number rather than producing ??.
% These two use @ in their names, and this point in the file is outside the
% \makeatletter region opened earlier (it closes at \makeatother above), so
% the region has to be reopened here.
\makeatletter
\newcommand{\@thmtitle}[2]{\global\@namedef{@tlong@#1}{#2}}
\newcommand{\nref}[1]{%
  \ref{#1}%
  \ifcsname @tlong@#1\endcsname
    \ ``\csname @tlong@#1\endcsname''%
  \fi
}
% \nm{t:KEY}: the stored long title of any \beT/\beL/\beP/\beR/\beC/\beD/\beE result,
% verbatim and without quotes, so item names follow the macro's first argument.
\newcommand{\nm}[1]{%
  \ifcsname @tlong@#1\endcsname
    \csname @tlong@#1\endcsname%
  \else
    \ref{#1}%
  \fi
}
\makeatother




%% ========================================
%% Other THEOREM ENVIRONMENTS
%% ========================================
\makeatletter
\@ifundefined{lemma}      {\newtheorem{lemma}      {Lemma}}      {}
\@ifundefined{proposition}{\newtheorem{proposition}{Proposition}} {}
\@ifundefined{corollary}  {\newtheorem{corollary}  {Corollary}}  {}
\@ifundefined{conjecture} {\newtheorem{conjecture} {Conjecture}}  {}
\@ifundefined{remark}     {\newtheorem{remark}     {Remark}}     {}
\@ifundefined{question}   {\newtheorem{question}   {Question}}   {}
\@ifundefined{assumption} {\newtheorem{assumption} {Assumption}} {}
\@ifundefined{hypothesis} {\newtheorem{hypothesis} {Hypothesis}} {}
\@ifundefined{fact}       {\newtheorem{fact}       {Fact}}       {}
\makeatother

% \citer: cite the relay-graph threshold theorem of AAHrel. If a document has
% \input the chapter/file containing t:relayRF (so the label resolves), point
% straight at it; otherwise fall back to a bare citation of AAHrel. Shared here
% (not per-paper) so it resolves consistently across every book chapter.
% \citeR: same idea, for the first theorem of AAHRah (t:Pstar).
\makeatletter
\providecommand{\citer}{%
  \@ifundefined{r@t:relayRF}%
    {\cite{AAHrel}}%
    {Theorem~\ref{t:relayRF} of \cite{AAHrel}}%
}
\providecommand{\citeR}{%
  \@ifundefined{r@t:Pstar}%
    {\cite{AAHRah}}%
    {Theorem~\ref{t:Pstar} of \cite{AAHRah}}%
}
\makeatother


% Abbreviations
\def\beQ{\begin{question}}\def\eeQ{\end{question}}
\def\beA{\begin{assumption}}\def\eeA{\end{assumption}}


%% ========================================

%% ========================================
%% COLORS & CUSTOM COMMANDS
%% ========================================
\definecolor{funccolor}{RGB}{25,25,112}
\definecolor{desccolor}{RGB}{64,64,64}
\newcommand{\func}[1]{\texttt{\textcolor{funccolor}{#1}}}
\def\eq{=}
\def\mm{-}

%% ========================================
%% MATHCAL SHORTHANDS
%% ========================================
% Project-wide convention: \mX for a capital X is always \mathcal X.
% Defined for the whole alphabet so the convention never fails to expand.
\newcommand{\mA}{\mathcal A}\newcommand{\mB}{\mathcal B}
\newcommand{\mC}{\mathcal C}\newcommand{\mD}{\mathcal D}
\newcommand{\mE}{\mathcal E}\newcommand{\mF}{\mathcal F}
\newcommand{\mG}{\mathcal G}\newcommand{\mH}{\mathcal H}
\newcommand{\mI}{\mathcal I}\newcommand{\mJ}{\mathcal J}
\newcommand{\mK}{\mathcal K}\newcommand{\mL}{\mathcal L}
\newcommand{\mM}{\mathcal M}\newcommand{\mN}{\mathcal N}
\newcommand{\mO}{\mathcal O}\newcommand{\mP}{\mathcal P}
\newcommand{\mQ}{\mathcal Q}\newcommand{\mR}{\mathcal R}
\newcommand{\mS}{\mathcal S}\newcommand{\mT}{\mathcal T}
\newcommand{\mU}{\mathcal U}\newcommand{\mV}{\mathcal V}
\newcommand{\mW}{\mathcal W}\newcommand{\mX}{\mathcal X}
\newcommand{\mY}{\mathcal Y}\newcommand{\mZ}{\mathcal Z}

%% ========================================
%% BASIC MATH ENVIRONMENTS
%% ========================================
\def\bep{\begin{pmatrix}}\def\eep{\end{pmatrix}}
\def\bev{\begin{vmatrix}}\def\eev{\end{vmatrix}}
\def\bea{\begin{eqnarray*}}\def\eea{\end{eqnarray*}}
\def\bc{\begin{cases}}\def\ec{\end{cases}}
\def\BEN{\begin{enumerate}}\def\EEN{\end{enumerate}}
\def\BI{\begin{itemize}}\def\EI{\end{itemize}}

\newcommand{\be}[1]{\begin{equation}\label{#1}}
\newcommand{\ee}{\end{equation}}
\newcommand{\beq}{\begin{eqnarray}}
\def\eeq{\end{eqnarray}}
\def\eqr{\eqref}\def\fr{\frac}\def\lbl{\label}
\def\Lra{\Longrightarrow}\def\non{\nonumber}\def\cNGM{\cite{Diek,Van,Van08}}
\def\Eq{\Leftrightarrow}\def\b{\beta}\def\uph{\red{UP TO HERE}}
%% ========================================
%% PROBABILITY & STATISTICS
%% ========================================
\renewcommand{\P}{\ensuremath{\mathbb{P}}}  % Note: overrides paragraph symbol
\newcommand{\E}{\ensuremath{\mathbb{E}}}
\newcommand{\ind}[1]{\ensuremath{\mathbbm{1}_{\left\{#1\right\}}}}
\DeclareMathOperator{\Cov}{Cov}
\DeclareMathOperator{\TV}{TV}
\newcommand{\rank}{\operatorname{rank}}

%% ========================================
%% CRN OPERATORS (organizations/ERCs/buffering/siphons framework)
%% ========================================
\def\clos{\mathrm{clos}}
\def\Supply{\mathrm{Supply}}
\def\req{\mathrm{req}}
\def\React{\mathrm{React}}
\def\Prod{\mathrm{Prod}}
\def\Gen{\mathrm{Gen}}
\def\In{\mathrm{In}}
\def\Out{\mathrm{Out}}
\def\Act{\mathrm{Act}}
\def\Min{\mathrm{Min}}

%% ========================================
%% LINE BREAKING
%% ========================================
% Last-resort stretch, used only on paragraphs TeX cannot otherwise break:
% a long \texttt{} token, a comma-separated inline math list (commas are not
% breakpoints in math), or a run of \cite/\ref boxes.  Without it those set as
% overfull \hbox and print into the margin; with it they set slightly loose,
% which is the lesser evil.  Does not affect paragraphs that already break.
\emergencystretch=3em

%% ========================================
%% NUMBER SETS
%% ========================================
\newcommand{\R}{\mathbb{R}}  % Real numbers
\newcommand{\C}{\mathbb{C}}  % Complex numbers
\newcommand{\N}{\mathbb{N}}
\newcommand{\Z}{\mathbb{Z}}
% \providecommand, not \newcommand: some papers define these in their own
% preamble after \input{preamble}, and must not clash.
\providecommand{\Q}{\mathbb{Q}}
\providecommand{\res}{\operatorname{res}}   % resultant; used by epiRes.tex
\newcommand{\MN}{\mathbb{N}}
\newcommand{\MZ}{\mathbb{Z}}
\newcommand{\realpos}{\mathbb{R}_{> 0}}
\newcommand{\realnneg}{\mathbb{R}_{\geq 0}}
\newcommand{\naturalpos}{\mathbb{N}_{>0}}
\newcommand{\Rplus}{\mathbb{R}_+}
\newcommand{\Rp}{\mathbb{R}_{\ge 0}}
\def\Rnne{\mathbb{R}_{\ge 0}}
\def\Se{\bar{S}} \def\Ie{\bar{I}}
\newcommand{\Rpos}{\mathbb{R}_{>0}}
\newcommand{\dd}{\,\mathrm d}
\newcommand{\e}{\mathrm e}
%% ========================================
%% GREEK LETTERS (2+ character abbreviations) greek
%% ========================================
\def\al{\alpha}\def\ba{\bff \alpha}
\def\w{\omega}\def\Lgn{\ga+\mu+\de}\def\om{\omega}
\def\de{\delta}\def\De{\Delta}\def\bde{\bff \delta}
\def\ep{\epsilon}\def\eps{\varepsilon}
\def\f{\varphi}  % Kept as requested
\def\ga{\gamma}\def\Ga{\Gamma}\def\G{\Gamma}\def\bg{\bff \gamma}
\def\up{\upsilon}  % NOT the breve accent: no accent commands anywhere.
                   % Named \up, not \u, so the LaTeX breve stays reachable.
\def\la{\lambda}\def\La{\Lambda}
\def\si{\sigma}\def\ka{\kappa}
\def\th{\theta}\def\vt{\vartheta}
\def\ze{\zeta}\def\Gap{\Gamma^{+}} \def\Gam{\Gamma^{-}}


%% ========================================
%% SPECIAL SHORT COMMANDS (kept as requested)
%% ========================================
\def\I{\infty}
\def\T{\widetilde}
\def\H{\hat}
\newcommand{\relint}{\operatorname{relint}}

%% ========================================
%% BOLD/VECTOR NOTATION
%% ========================================
\newcommand{\bff}[1]{{\mbox{\boldmath$#1$}}}
\newcommand\X{\boldsymbol{X}}
\def\x{\boldsymbol{x}}\def\bx{\bff x}\def\xe{\bar{\x}}
\newcommand\y{\boldsymbol{y}}\def\by{\boldsymbol{y}}\def\ye{\bar{\y}}
\def\yb{\y}\def\xb{\x}
\def\xy{\x-\y}
\newcommand\br{\boldsymbol{r}}\def\re{\bar{\r}}
\newcommand\bd{\boldsymbol{d}}
\newcommand\bQ{\boldsymbol{Q}}
\newcommand{\Yb}{{\mathbf{Y}}}
\newcommand{\Ab}{\mathbf{A}}
\newcommand{\zb}{\mathbf{z}}\def\bz{\bff 0}
\newcommand{\Zb}{\mathbf{Z}}
\newcommand{\bb}{\mathbf{b}}
\def\vbe{\vec {\bff \beta}}
\def\vb{\vec {\bff b}}
\def\v1{\vec {\bff 1}}
\def\vn{\vec \de}
\def\vz{\vec 0}
\def\vi{\overset{\rightarrow}{\mathsf i}}
\def\vv{\overset{\rightarrow}{v}}
\def\vr{\; \vec {\mathsf r}}
\newcommand{\va}{\bff a}
\providecommand{\vect}[1]{\boldsymbol{#1}}   % bold vector (moved from Gradostat.tex)
\providecommand{\one}{\mathbf 1}             % all-ones vector (moved from AliSF)
\providecommand{\spr}{s}                     % spectral bound  (moved from AliSF)
\providecommand{\Ke}{K_{\rm eff}}            % effective carrying capacity
% Critical right/left eigenvectors in the Hopf projection method (AAHOSN/New.tex)
\providecommand{\bq}{\mathbf{q}}
\providecommand{\bp}{\mathbf{p}}


%% ========================================
%% MATHCAL & MATHSCR
%% ========================================
% Mathcal
\newcommand\mc[1]{\mathcal{#1}}
\def\mC{{\mathcal C}}\def\mD{\mathcal D}\def\mE{\mathcal{E}}
\def\mF{\mathcal F}\def\mG{{\mathcal G}}\def\mH{{\mathcal H}}
\def\mI{\mathcal I}\def\mL{\mathcal{L}}\def\mM{{\mathcal M}}
\def\mN{{\mathcal N}}\def\mP{\mathcal{P}}\def\mR{{\mathcal R}}
\def\mS{{\mathcal S}}\def\mT{\mathcal{T}}\def\mV{{\mathcal V}}
\def\mK{\mathcal K}\def\mA{{\mathcal A}}\def\mU{{\mathcal U}}
% Symmetric bilinear and trilinear forms of the Taylor expansion, used in the
% first Lyapunov coefficient (AAHOSN/New.tex)
\providecommand{\Bform}{\mathcal{B}}
\providecommand{\Cform}{\mathcal{C}}

% Blackboard bold.  Kept here, not in the drivers: a \providecommand in a
% driver preamble runs before this file and silently wins, so the same name
% ends up meaning different things in different builds.
\def\IC{\mathbb{C}}\def\IN{\mathbb{N}}\def\IQ{\mathbb{Q}}
\def\IR{\mathbb{R}}\def\IZ{\mathbb{Z}}\def\IG{\mathbb{G}}
\def\II{\mathbb{I}}

% Mathscr
\newcommand{\Rs}{{\mathscr R}}
\newcommand{\Ss}{{\mathscr S}}
\newcommand{\Ns}{{\mathscr N}}
\newcommand{\Cs}{\mathscr C}


%% ========================================
%% OPERATORS & FUNCTIONS
%% ========================================
\DeclareMathOperator{\tr}{tr}
\DeclareMathOperator{\diag}{Diag}
\DeclareMathOperator{\Diag}{Diag}
\DeclareMathOperator{\intr}{\mathrm{Int}}
\DeclareMathOperator{\Inv}{Inv}
\DeclareMathOperator{\spec}{spec}
\DeclareMathOperator{\AUC}{AUC}               % area under the curve
\newcommand{\ii}{\mathrm{i}}                  % imaginary unit
\newcommand\op{\operatorname}
\newcommand{\supp}{\op{supp}}
\newcommand{\diff}{\mathop{}\mathopen{}\mathrm{d}}
\newcommand{\pdd}[1]{\frac{\partial}{\partial{#1}}}
\DeclareMathOperator{\cone}{cone}
\DeclareMathOperator{\proj}{proj}

%% ========================================
%% BRACKETS & NORMS
%% ========================================
\providecommand{\pr}[1]{\left(#1\right)}
\providecommand{\pp}[1]{\left[#1\right]}
\newcommand\croc[1]{\left\langle #1\right\rangle}
\newcommand{\norm}[1]{\|#1\|}
\newcommand{\inner}[2]{\langle #1, #2 \rangle}
% Short synonym of \inner, used by AAHOSN/New.tex as $\ip{u}{v}=\bar u^\top v$
\providecommand{\ip}[2]{\langle #1 , #2 \rangle}


%% ========================================
%% FORMATTING & UTILITIES
%% ========================================
\newcommand*{\Scale}[2][4]{\scalebox{#1}{$#2$}}
\newcommand{\te}[1]{\text{#1}}
\newcommand\steq[1]{\stackrel{\text{\rm #1.}}{=}}
\newcommand\blue[1]{\textcolor{blue}{#1}}
\newcommand{\red}{\textcolor[rgb]{1.00,0.00,0.00}}
\newcommand{\green}{\textcolor[rgb]{0.00,1.00,0.00}}

% Figure command
\newcommand{\figu}[3]{%
\begin{figure}[H]
\centering
\includegraphics[scale=#3]{#1}
\caption{#2\label{f:#1}}
\end{figure}
}

% Footnote
\long\def\symbolfootnote[#1]#2{%
\begingroup
\def\thefootnote{\fnsymbol{footnote}}\footnote[#1]{#2}%
\endgroup}
\def\fn{\symbolfootnote}

\newcommand{\dsum}{\displaystyle\sum}

%% ========================================
%% CRN & EPIDEMIOLOGY SPECIFIC
%% ========================================
% CRN notation
\def\SS{\mathcal S}
\def\CC{\mathcal C}
\def\RR{\mathcal R}

% Reproduction numbers
\def\Rto{R_{1}^{E_2}}
\def\Rot{R_{2}^{E_1}}
\def\Roc{R_{1,c}}
\def\Rtc{R_{2,c}}
\def\rd{r_{0}}
\def\sd{s_{0}}
\def\dd{D_{0}}
\def\vd{v_{0}}
\def\zd{z_{0}}
\def\se{\sE}
\def\xd{x_{0}}
%% ========================================
%% SPECIAL VALUES & EQUILIBRIA
%% ========================================
\def\eE{\bar{e}}
\def\wE{\bar{w}}
\def\sE{s_{*}}
\def\Me{\T M_{ee}}
\def\Rnne{\mathbb{R}_{\ge 0}}

\newcommand{\ones}{\mathbf 1}


%% ========================================
%% ABBREVIATIONS (alphabetical by category)
%% ========================================

% -- General mathematical terms --
\def\adm{admissible}\def\admy{admissibility}
%\def\and{antisymmetric}% disabled 2026-09-30: it overrode LaTeX's author separator \and (printed "antisymmetric" between authors)
\def\ass{assumption}
\def\cadlag{c\`adl\`ag}
\def\ch{characteristic polynomial}
\def\chp{characteristic polynomial}
\def\coe{coefficient}
\def\comp{computation}
\def\con{condition}
\def\coo{coordinate}
\def\corr{corresponding}
\def\dec{decomposition}
\def\den{denominator}
\def\dete{deterministic}
\def\eig{eigenvalue}\def\eval{eigenvalue}
\def\eigv{eigenvector}\def\eiv{eigenvector}\def\evec{eigenvector}
\def\elt{element}
\def\enn{essentially nonnegative/positivity preserving}
\def\equ{equation}
\def\etal{et al.}
\def\fe{for example}
\def\fno{from now on}
\def\for{\forall}
\def\form{formula}
\def\Fr{Furthermore, }
\def\frt{furthermore}
\def\fun{function}
\def\gene{generalization}
\def\heu{heuristic}
\def\how{however}\def\How{However, }
\def\ie{i.e. }
\def\im{\item}
\def\isd{is defined by}
\def\itc{it is easy to check that}
\def\itf{it follows that}\def\Itf{It follows that}
\def\ith{it holds that}
\def\itm{it may be checked that}\def\Itm{It may be checked that}
\def\ito{it turns out that}\def\Ito{It turns out that}\def\Itou{It turns out useful to}
\def\ma{mass action}\def\MAR{mass action representation}
\def\mat{matrix}
\def\mbc{may be checked to be}
\def\mbw{may be written as}\def\mubw{must be written as}
\def\mpr{more precisely, }\def\Mp{More precisely, }
\def\nc{non-catalytic}
\def\nne{non-negative}\def\nny{nonnegativity}
\def\no{\nonumber}
\def\noni{non-infectious}
\def\nota{notation}
\def\np{\newpage}
\def\Nt{Note that}
\def\nt{non trivial}\def\nz{non zero}
\def\oth{otherwise}
\def\Oth{On the other hand}
\def\pa{parameter}\def\para{parameter}
\def\pol{polynomial}
\def\prob{probability}
\def\Prf{{\bf  Proof:}}\def\prf{{\bf  Proof:}}
\def\prop{proportional to}
\def\qu{\quad}
\def\QED{\hfill {$\square$}\goodbreak \medskip}
\def\rep{representation}
\def\resp{respectively}
\def\satd{satisfied}\def\satg{satisfying}\def\sats{satisfies}\def\saty{satisfy}
\def\ssec{\subsection}\def\sssec{\subsubsection}
\def\sef{see for example}
\def\sksy{skew-symmetric}
\def\sub{substrate}
\def\sys{system}
\def\var{variable}
\def\wk{well-known}
\def\wl{without loss of generality}\def\wlo{w.l.o.g.}
\def\wmc{We may check that}\def\Wmc{We may check that}
\def\wmr{we may rewrite}
\def\wms{we must show that}
\def\wrt{with respect to}

% -- CRN & Chemistry --
\newcommand\CRN{chemical reaction network}
\newcommand\CRNs{chemical reaction networks}
\newcommand\RN{reaction network}
\def\ACR{absolute concentration robustness}
\def\AMSD{acyclic siphon decomposition}
\def\DAE{differential algebraic equations}
\def\FHJ{Feinberg-Horn-Jackson}
\def\GMAK{generalized mass action kinetics}
\def\MM{Michaelis-Menten}
\def\sm{stoichiometric matrix}\def\Sm{stoichiometric matrix}
\def\WRZD{weakly reversible, zero deficiency}\def\aut{autocatalytic}

% -- Epidemiology --
\def\brn{basic reproduction number}\def\bzn{basic replacement number}
\def\com{compartment}
\def\dfe{disease free equilibrium}\def\DFE{disease free equilibrium}
\def\dwt{dwell times}\def\edt{expected dwell times}
\def\EM{epidemic model}
\def\endp{endemic fixed point}\def\enp{endemic point}\def\EE{endemic point}
\def\ME{mathematical epidemiology}
\def\MEE{eco-epidemiology}\def\EEP{eco-epidemiological}
\def\ece{eco-epidemiological model}
\def\mgm{multi group model}
\def\mtS{multi-strain model}\def\mSt{multi-strain model}
\def\mSi{minimal siphon}
\def\NGM{next generation matrix}
\def\nim{new infections matrix}
\def\qem{quadratic epidemiological model}
\def\repF{invasion functions}
\def\SPF{SIR-PH-FA model}
\def\sus{susceptible}

% -- Graph theory & Linear algebra --
\def\jac{Jacobi matrix}
\def\jin{Jacobian matrix of the invasion vector field with respect to the invasion variables}
\def\jrn{Jacobian matrix of the resident vector field with respect to the resident variables}
\def\jxn{Jacobian matrix of the $\x$ vector field with respect to the $\x$ variables}
\def\Lm{Laplacian matrix}
\def\MTT{matrix tree theorem}
\def\SGT{spectral graph theory}
\def\SMW{Sherman–Morrison–Woodbury}

% -- Stability & Dynamics --
\def\absC{absorbtion CTMC}
\def\BTb{Bogdanov-Takens bifurcation}
\def\CEP{competitive exclusion partition}
\def\CTMC{continuous time Markov chain}
\def\ds{differential system}
\def\fp{fixed point}\def\bfp{boundary fixed point}\def\fbp{fixed boundary point}
\def\fbl{fixed boundary line}\def\fif{forward invariant face}
\def\gas{global asymptotic stability}\def\GAS{globally asymptotically stable}
\def\GASS{\text{GASS}}
\def\Hb{Hopf bifurcation}
\def\Hur{Hurwitz}\def\MH{Metzler-Hurwitz}\def\RH{Routh-Hurwitz}
\def\IS{\text{IS}}
\def\LAS{locally asymptotically stable}
\def\LdS{Lyapunov diagonally stable}
\def\Lf{Lyapunov function}\def\LyV{Lyapunov Volterra type function}\def\RLF{robust Lyapunov function}
\def\LyM{Lyapunov Malkin stable}
\def\PF{Perron-Frobenius}
\def\spfp{strictly positive fixed point}
\def\sta{stability}
\def\STP{sharp threshold property}\def\STT{sharp threshold theorem}

% -- Models (with \varphi explicitly instead of \f) --
\def\AB{$(A,B)$ Arino-Brauer epidemic models}
\def\ABa{$(A,\bbe,\ba,\varphi)$ balanced bilinear model}
\def\ABP{%$(A,\bbe,P,\varphi)$
\\ balanced bilinear model}
\def\AMG{Anderson May Gupta}
\newcommand\BIN{biological interaction network}
\def\Dt{Descartes type}
\def\gLV{generalized Lotka-Volterra model}
\def\IG{\text{invasion graph}}\def\IGMS{interaction graph on minimal siphons}
\def\LCP{\text{linear complementarity problem}}
\def\LV{Lotka-Volterra}\def\LVK{Lotka-Volterra-Kolmogorov}

% -- Splitting & Structure --
\def\reg{regular splitting}\def\regS{{\bf regular splitting}}
\def\RUR{rational univariate representation}
\def\bds{block-diagonal structure}\def\NGMs{next generation matrices}

% -- Miscellaneous --
\def\Ma{{Mathematica}}
\def\Mm{Max–Min}
\def\num{numerator}
\def\rap{rational approximation} \def\DSR{directed species-reaction bipartite graph}\def\PCP{polynomial complementarity problem}\def\rat{rational substitution}
\def\QSS{quasi steady state} \def\RHS{right hand side} \def\LHS{left hand side}
\def\OSN{online social network}\def\VdD{Van den Driessche}
\def\wli{$\omega$-limit} \def\rk{rich kinetics}\def\sJ{symbolic Jacobian}
\def\cV{\cite{Vas,VasSta,BSV,GSV}} \def\spf{semi-parametric semi-functional}
\def\CRW{Capasso-Ruan-Wang} \def\ABP{$(A,\bbe,P,\f)$ balanced bilinear model}
\def\ABa{%$(A,\bbe,\ba,\f)$ \\
rank-one balanced bilinear model}
\def\bbe{B%{\bff \beta}
} \def\vbe{\vec {\bff \beta}
}
\def\vb{\vec {\bff b}}
\newcommand{\Kh}{\hat{K}}          % K with hat, first mention only
\newcommand{\wS}{w^{\sigma}}       % right Perron J^\perp, full block
\newcommand{\piS}{\vec \pi^{\sigma}}    % left Perron J^\perp, full block
\newcommand{\ws}{w^{\sigma}}       % right Perron J^\perp, sub-block
\newcommand{\pis}{\pi^{\sigma}}    % left Perron J^\perp, sub-block
\newcommand{\hwS}{\hat{w}^{\sigma}}   % right Perron K, full block
\newcommand{\hpiS}{\hat{\pi}^{\sigma}} % left Perron K, full block
\newcommand{\hws}{\hat{w}^{\sigma}}   % right Perron K, sub-block
\newcommand{\hpis}{\hat{\pi}^{\sigma}} % left Perron K, sub-block
\newcommand{\twS}{\tilde{w}^{\sigma}}  % right Perron \tilde K, full block
\newcommand{\tpiS}{\tilde{\pi}^{\sigma}} % left Perron \tilde K
\newcommand{\tws}{\tilde{w}^{\sigma}}
\newcommand{\tpis}{\tilde{\pi}^{\sigma}}
\newcommand{\Mat}[2]{\mathbb{R}_{\ge 0}^{\,#1\times #2}}
\def\tJ{J_\sigma^\perp} \def\lPe{left Perron eigenvector}
% \G is \Gamma (defined at :795).  The Volterra function keeps its own name,
% so that \G no longer has two conflicting meanings in this file.
\def\GV{G}  % Volterra G(u)=u-1-\ln u
\def\PVL{Perron-Volterra Lyapunov minimal siphon ansatz}
\newcommand{\istar}[1]{i_{#1}^{*}}
%\newcommand{\sE}{s_E}
\newcommand{\vE}{v_E}
\def\PV{Perron-Volterra candidate Lyapunov function}
\def\GLV{generalized Lotka-Volterra}
\def\MERN{mathematical epidemiology reaction networks}
\def\CMM{contraction metric method}
\newcommand{\sign}{\operatorname{sign}}
\newcommand{\Var}{\operatorname{Var}} 

% ======================================================================
% IDEAS SHARED BETWEEN THE ABSTRACT AND THE MAIN RESULTS SECTION
%
% Write each idea ONCE, in the abstract; the Main Results section recalls it,
% so the two can never drift apart.
%
%   in the abstract        \ideaitem{1}{The first-order reaction terms ...}
%   in Main Results        \ideahead{1}
%
% \ideaitem stores the text globally and typesets
%   \item \hyperlink{idea:1}{<text>}
% \ideahead typesets
%   \item\hypertarget{idea:1}{}\textbf{<the same text>}
%
% The abstract must come before the Main Results input, which is the normal
% order. If a heading has to read differently from the abstract sentence (for
% instance a long abstract item summarised in one line), override it with the
% optional argument: \ideahead[short heading]{6}.
% ======================================================================
\makeatletter
\newcommand{\ideaitem}[2]{%
  \expandafter\gdef\csname @idea@#1\endcsname{#2}%
  \item\hyperlink{idea:#1}{#2}}
\newcommand{\ideahead}[2][]{%
  \item\hypertarget{idea:#2}{}%
  \def\@ideatmp{#1}%
  \ifx\@ideatmp\@empty
    \@ifundefined{@idea@#2}{%
      \PackageWarning{preamble}{Idea #2 is used in Main Results but was never
        defined by \string\ideaitem\space in the abstract}%
      \textbf{[idea #2 undefined]}%
    }{\textbf{\begingroup\def\ideaterm##1##2{##2}\def\idearest##1{}\@nameuse{@idea@#2}\endgroup}}%
  \else
    \textbf{#1}%
  \fi}
% \ideaterm{<hypertarget>}{<words>}: inside an abstract idea the words become an
% invisible link (the document sets \hypersetup{hidelinks}) jumping to the place
% in Main Results that defines them; inside the Main Results heading itself,
% where the target sits a few lines below, \ideahead strips the link and prints
% the words plainly. Mark the destination with \hypertarget{<name>}{} on the line
% before the display, which adds no equation number and shows nothing.
\newcommand{\ideaterm}[2]{\hyperlink{#1}{#2}}
% \idearest{<text>}: the part of an abstract idea after its first sentence; printed in
% the abstract, omitted from the Main Results heading.
\newcommand{\idearest}[1]{#1}
\makeatother

% ======================================================================
% SHORT-PAPER SWITCH
% \shorttrue in a master that assembles only part of the companion paper
% (Barrat/SipArxiv.tex). \nsh{...} drops text that cites results living in the
% omitted sections, so the short build has no dangling references while the
% companion paper keeps the text verbatim.
% ======================================================================
\newif\ifshort \shortfalse
\newcommand{\nsh}[1]{\ifshort\else#1\fi}
% \shortalt{<short build>}{<companion paper>}: use when a sentence needs
% different wording in the two builds, not merely a deletion.
\newcommand{\shortalt}[2]{\ifshort#1\else#2\fi}
% \Jtr, \JtrS{\Sigma}: transverse Jacobian. J^\perp here; SipArxiv.tex redefines them as J_{\mathrm{tra}}, J_{\mathrm{tra},\Sigma}
\newcommand{\Jtr}{J^\perp}
\newcommand{\JtrS}[1]{J_{#1}^\perp}
.
\makeatletter
\AtBeginDocument{\@ifundefined{hypersetup}{}{\hypersetup{linktoc=all}}}
\makeatother
\usepackage{graphicx,float,fancybox}
\usepackage{epsfig,epstopdf,tikz}
\usepackage[sans]{dsfont}
\usetikzlibrary{positioning,arrows,arrows.meta,calc}
% Configure listings for Mathematica code
\lstdefinestyle{mathematica}{
 language=Mathematica,
 basicstyle=\small\ttfamily,
 backgroundcolor=\color{gray!10},
 frame=single,
 breaklines=true,
 commentstyle=\color{green!60!black},
 keywordstyle=\color{blue},
 stringstyle=\color{red},
 showstringspaces=false,
 tabsize=2
}
\usepackage{tcolorbox}
% Colored boxes for key results
\newtcolorbox{keyresult}[1][]{
 colback=blue!5!white,
 colframe=blue!75!black,
 title=#1,
 fonttitle=\bfseries
}
\newtcolorbox{examplebox}[1][]{
 colback=green!5!white,
 colframe=green!75!black,
 title=#1,
 fonttitle=\bfseries
}


\makeatletter
\@ifundefined{theorem}   {\newtheorem{theorem}   {Theorem}}   {}
\@ifundefined{definition}{\newtheorem{definition}{Definition}} {}
\@ifundefined{open}      {\newtheorem{open}      {Problem}}   {}
\@ifundefined{example}   {\newtheorem{example}   {Example}}   {}
\makeatother

% COUNTER DEFINITIONS
\newcounter{tabcounter}
\newcounter{defcounter}
\newcounter{opcounter}
\newcounter{thmcounter}
\newcounter{excounter}
\newcounter{ascounter}
\newcounter{cjcounter}
\newcounter{lemcounter}
\newcounter{remcounter}
\newcounter{corcounter}
\newcounter{propcounter}
\newcounter{figcounter}

% UNIFIED GLOSSARY STYLE
\newglossarystyle{unifiedstyle}{%
  \setglossarystyle{list}%
  \renewcommand*{\glossarypreamble}{\gdef\glslastsec{}}%
  \renewcommand{\glossentry}[2]{%
    \edef\glscursec{\glsentryuseriii{##1}}%
    \ifx\glscursec\glslastsec\else
      \item[]\textbf{\S\,\glscursec\enspace\glsentryuseriv{##1}}%
      \global\let\glslastsec\glscursec
    \fi
    \item[\glsentryuserii{##1} \glsentryuseri{##1}:] ``\glossentryname{##1}'' \glsentrydesc{##1}%
  }%
}

\makeatletter
% First-word extractor (used by \beT, \beD and \beXa for the default key)
% \@firstwordof hello world \relax  ->  hello
\def\@firstwordof#1 #2\relax{#1}
% First three letters, used for the heading mnemonic of a theorem given no
% explicit key.  Padding with \@empty keeps it safe for first words shorter
% than three letters: \@empty vanishes under \edef, whereas \relax would
% survive into the label and is illegal inside \csname (it broke every
% two-letter title, e.g. "ME ...", with "Missing \endcsname inserted").
\def\@takethree#1#2#3#4\@nil{#1#2#3}

% EXAMPLE COMMAND -- same signature and behaviour as \beT and \beD.
% Usage: \beXa*[long title][short title][key]
%   star  = suppresses the printed "[key]" tag after the title (as for \beT*
%   and \beD*); the label and glossary entry are unaffected.
%   key = optional label key override; the label becomes e:KEY, and without a
%   key it is e: plus the first three letters of the title's first word.  A
%   mnemonic label lets an example be cited as "Example~\ref{e:KEY}" instead
%   of by section number, which is what the citing prose needs.
\NewDocumentCommand{\beXa}{s O{} O{} o}{%
  \def\@temparg{#2}%
  \ifx\@temparg\@empty
    \begin{example}%
  \else
    \let\@savedSecTitle\@currentsectionname%
    \refstepcounter{excounter}%
    \IfValueTF{#4}{%
      \edef\@exlblkey{#4}%
    }{%
      \edef\@exfirstword{\@firstwordof#2 \relax}%
      \edef\@exlblkey{%
        \expandafter\@takethree\@exfirstword\@empty\@empty\@empty\@nil}%
    }%
    \def\@tempshort{#3}%
    \ifx\@tempshort\@empty
      \expandafter\xdef\csname @xshort@\@exlblkey\endcsname{\@exlblkey}%
    \else
      \expandafter\def\csname @xshort@\@exlblkey\endcsname{#3}%
    \fi
    \IfBooleanTF{#1}{\let\@extag\@empty}{%
      \edef\@extag{\noexpand\ \noexpand\textup{[\noexpand\texttt{e:\@exlblkey}]}}}%
    \begin{example}[#2\ifshowkeys\@extag\fi]%
    \label{e:\@exlblkey}%
    % record the long title so \nref{e:KEY} prints it, as for theorems
    \protected@write\@auxout{}{%
      \string\@thmtitle{e:\@exlblkey}{\unexpanded{#2}}}%
    \expanded{%
      \noexpand\newglossaryentry{e\theexcounter}{%
        name={\unexpanded{#2}%
          \IfBooleanF{#1}{ [\noexpand\texttt{e:\@exlblkey}]}},%
        sort={\theexample},%
        description={p.~\noexpand\pageref{e:\@exlblkey}},%
        user1={\theexample},%
        user2={Example},%
        user3={\thesection},%
        user4={\expandafter\unexpanded\expandafter{\@savedSecTitle}}%
      }%
    }%
    \protected@write\@auxout{}{\string\glsliveentry{e\theexcounter}}%
  \fi
}
\newcommand{\eeXa}{\end{example}}

% DEFINITION COMMAND -- same signature and behaviour as \beT below.
% Usage: \beD[long title][short title][key]
%   long title  = definition heading (plain text only)
%   short title = shown in \dref{word}; defaults to the key
%   key         = optional label key override; the label becomes d:KEY
% Without a key the label is d: plus the first three letters of the title's
% first word, so it is a mnemonic rather than the definition's number: a
% numeric label would silently change whenever a definition is inserted.
\NewDocumentCommand{\beD}{s O{} O{} o}{%
  \def\@temparg{#2}%
  \ifx\@temparg\@empty
    \begin{definition}%
  \else
    \let\@savedSecTitle\@currentsectionname%
    \refstepcounter{defcounter}%
    \IfValueTF{#4}{%
      \edef\@deflblkey{#4}%
    }{%
      \edef\@deffirstword{\@firstwordof#2 \relax}%
      \edef\@deflblkey{%
        \expandafter\@takethree\@deffirstword\@empty\@empty\@empty\@nil}%
    }%
    \def\@tempshort{#3}%
    \ifx\@tempshort\@empty
      \expandafter\xdef\csname @dshort@\@deflblkey\endcsname{\@deflblkey}%
    \else
      \expandafter\def\csname @dshort@\@deflblkey\endcsname{#3}%
    \fi
    \IfBooleanTF{#1}{\let\@deftag\@empty}{%
      \edef\@deftag{\noexpand\ \noexpand\textup{[\noexpand\texttt{d:\@deflblkey}]}}}%
    \begin{definition}[#2\ifshowkeys\@deftag\fi]%
    \label{d:\@deflblkey}%
    % record the long title so \nref{d:KEY} prints it, as for theorems
    \protected@write\@auxout{}{%
      \string\@thmtitle{d:\@deflblkey}{\unexpanded{#2}}}%
    \expanded{%
      \noexpand\newglossaryentry{d\thedefcounter}{%
        name={\unexpanded{#2}%
          \IfBooleanF{#1}{ [\noexpand\texttt{d:\@deflblkey}]}},%
        sort={\thedefinition},%
        description={p.~\noexpand\pageref{d:\@deflblkey}},%
        user1={\thedefinition},%
        user2={Definition},%
        user3={\thesection},%
        user4={\expandafter\unexpanded\expandafter{\@savedSecTitle}}%
      }%
    }%
    \protected@write\@auxout{}{\string\glsliveentry{d\thedefcounter}}%
  \fi
}
\newcommand{\eeD}{\end{definition}}

% \showkeysfalse in a driver hides the [t:KEY] mnemonics in result headings
% (labels and glossary entries are unchanged); used for submission builds.
\newif\ifshowkeys \showkeystrue
% \commonredtrue in a driver that inputs basicFiles/commonRed: tool.tex then
% cites its reduction results instead of restating them.
\newif\ifcommonred \commonredfalse
% TheorEM COMMAND
% Usage: \beT[long title][short title][key]
%   long title  = theorem heading (plain text only)
%   short title = shown in \tref{word}; may contain math; defaults to first word of long title
%   key         = optional label key override; label becomes t:KEY (avoids first-word collisions)
% Label is t:WORD where WORD = first word of long title (or KEY if given); use \tref{WORD} to reference.
\NewDocumentCommand{\beT}{s O{} O{} o}{%
  \def\@temparg{#2}%
  \ifx\@temparg\@empty
    \begin{theorem}%
  \else
    \let\@savedSecTitle\@currentsectionname%
    \refstepcounter{thmcounter}%
    % \@thmtag is the mnemonic shown in the heading; every theorem gets one.
    % With an explicit key the key itself is shown.  Without one the key
    % defaults to the title's first word, and showing that whole word would
    % merely repeat the title ("Dissipativity [Dissipativity]"), so only its
    % first three letters are shown: "Dissipativity [Dis]".  The label is
    % unaffected -- it still uses the full \@thmlblkey.
    \IfValueTF{#4}{%
      \edef\@thmlblkey{#4}%
    }{%
      \edef\@thmfirstword{\@firstwordof#2 \relax}%
      \edef\@thmlblkey{%
        \expandafter\@takethree\@thmfirstword\@empty\@empty\@empty\@nil}%
    }%
    \IfBooleanTF{#1}{\let\@thmtag\@empty}{%
      \edef\@thmtag{\noexpand\ \noexpand\textup{[\noexpand\texttt{t:\@thmlblkey}]}}}%
    \def\@tempshort{#3}%
    \ifx\@tempshort\@empty
      \expandafter\xdef\csname @tshort@\@thmlblkey\endcsname{\@thmlblkey}%
    \else
      \expandafter\def\csname @tshort@\@thmlblkey\endcsname{#3}%
    \fi
    \begin{theorem}[#2\ifshowkeys\@thmtag\fi]%
    \label{t:\@thmlblkey}%
    % record the long title for \nref; via the .aux so forward refs resolve
    \protected@write\@auxout{}{%
      \string\@thmtitle{t:\@thmlblkey}{\unexpanded{#2}}}%
    \expanded{%
      \noexpand\newglossaryentry{t\thethmcounter}{%
        name={\unexpanded{#2}%
          \IfBooleanF{#1}{ [\noexpand\texttt{t:\@thmlblkey}]}},%
        % sort/user1 must be the number the document actually prints, i.e. the
        % theorem environment's counter, not thmcounter.  thmcounter advances
        % only on \beT with a title, so a bare \beT (or \beT with an
        % empty title) makes the two drift and the glossary cites a wrong
        % number.  thmcounter is still fine as the entry key: it is unique.
        sort={\thetheorem},%
        description={p.~\noexpand\pageref{t:\@thmlblkey}},%
        user1={\thetheorem},%
        user2={Theorem},%
        user3={\thesection},%
        user4={\expandafter\unexpanded\expandafter{\@savedSecTitle}}%
      }%
    }%
    \protected@write\@auxout{}{\string\glsliveentry{t\thethmcounter}}%
  \fi
}
\newcommand{\eeT}{\end{theorem}}

% CROSS-PART RELATION PARAGRAPH
% Usage: \xpart{text}  -- place immediately BEFORE \beT, never inside it
%   (anything inside \beT...\eeT would be swallowed into the glossary entry).
% The paragraph is typeset only when Part I is actually present, i.e. in the
% combined book build; in the standalone builds (GranadaDSABNS/A.tex,
% AAHRah-Submission/AAHRah.tex) the label part:glv is undefined and the whole
% paragraph disappears rather than leaving ?? cross-references. Same
% graceful-degradation principle as \GASref in GASt.tex. Resolves on the
% second LaTeX pass, since r@part:glv comes from the .aux file.
% (Already inside the \makeatletter region opened above, so no \makeatletter
% pair here -- adding one would close that region early.)
\newcommand{\xpart}[1]{\@ifundefined{r@part:glv}{}{\paragraph{Relation with Parts~I--II.}#1}}

% \bookref{label}{precise text}{fallback text}: prints the precise cross-reference
% when the label resolves (combined book build), and the fallback text otherwise
% (standalone paper builds, where book-only appendix labels are absent). Same
% graceful-degradation role as \GASref in GASt.tex, but usable in any file since it
% lives in this preamble (loaded by every driver) -- unlike \GASref, which is defined in a
% content file and is therefore undefined in standalone builds that omit that file.
\newcommand{\bookref}[3]{\@ifundefined{r@#1}{#3}{#2}}

% \PVpart and \RBpart: the two parts that a partial build may omit.
% They print "Part~N" in any build that contains the part, and a naming phrase
% in a build that does not, so an excerpt shows no ?? .  Each already supplies
% the word "Part", so write "\PVpart develops ...", not "Part~\PVpart".
\newcommand{\PVpart}{\bookref{part:aahrah}%
  {Part~\ref{part:aahrah}}{the Perron--Volterra part of the full book}}
\newcommand{\RBpart}{\bookref{part:recbrowne}%
  {Part~\ref{part:recbrowne}}{the immuno-virology part of the full book}}

% OPEN PROBLEM COMMANDS
% ASSUMPTION COMMANDS (glossaried)
\newcommand{\beAs}[1][]{%
  \def\@temparg{#1}%
  \ifx\@temparg\@empty
    \begin{assumption}%
  \else
    \let\@savedSecTitle\@currentsectionname%
    \refstepcounter{ascounter}%
    \begin{assumption}[#1]%
    \label{as:\theascounter}%
    \expanded{%
      \noexpand\newglossaryentry{as\theascounter}{%
        name={\unexpanded{#1}},%
        sort={\theascounter},%
        description={p.~\noexpand\pageref{as:\theascounter}},%
        user1={\theascounter},%
        user2={Assumption},%
        user3={\thesection},%
        user4={\expandafter\unexpanded\expandafter{\@savedSecTitle}}%
      }%
    }%
    \protected@write\@auxout{}{\string\glsliveentry{as\theascounter}}%
  \fi
}
\newcommand{\eeAs}{\end{assumption}}
% CONJECTURE COMMANDS (glossaried)
\newcommand{\beCj}[1][]{%
  \def\@temparg{#1}%
  \ifx\@temparg\@empty
    \begin{conjecture}%
  \else
    \let\@savedSecTitle\@currentsectionname%
    \refstepcounter{cjcounter}%
    \begin{conjecture}[#1]%
    \label{cj:\thecjcounter}%
    \expanded{%
      \noexpand\newglossaryentry{cj\thecjcounter}{%
        name={\unexpanded{#1}},%
        sort={\thecjcounter},%
        description={p.~\noexpand\pageref{cj:\thecjcounter}},%
        user1={\thecjcounter},%
        user2={Conjecture},%
        user3={\thesection},%
        user4={\expandafter\unexpanded\expandafter{\@savedSecTitle}}%
      }%
    }%
    \protected@write\@auxout{}{\string\glsliveentry{cj\thecjcounter}}%
  \fi
}
\newcommand{\eeCj}{\end{conjecture}}
% LEMMA COMMANDS
% ---------------------------------------------------------------------
% UNIFORM RESULT WRAPPERS -- \beL, \beC, \beP, \beR share \beT's signature
%
%   \beX[long title][short title][key]  ...  \eeX
%
%   long title  = heading (math allowed)
%   short title = shown in the cross-reference; defaults to the key
%   key         = label key.  With a key the label is X:KEY and the heading
%                 carries the mnemonic [key], exactly as for \beT.  Without
%                 one the label falls back to the running counter, which is
%                 what every pre-existing call site in the book relies on.
%
% A bare \beX -- arguments commented out -- opens the plain environment and
% creates NO label and NO glossary entry.  That is the supported way to keep
% an individual result out of the catalogue.
% ---------------------------------------------------------------------
% #1 wrapper name  #2 environment  #3 label prefix  #4 counter  #5 type name
\newcommand{\@mkresult}[5]{%
  % A leading * suppresses the [key] mnemonic in both the heading and the
  % glossary name, while still making the label X:KEY and the entry:
  %     \beX*[long][short][key]
  \expandafter\NewDocumentCommand\csname be#1\endcsname{s O{} O{} o}{%
    \def\@temparg{##2}%
    \ifx\@temparg\@empty
      \begin{#2}%
    \else
      \let\@savedSecTitle\@currentsectionname%
      \refstepcounter{#4}%
      % With a key: mnemonic label X:KEY and the [key] tag in the heading.
      % Without one: fall back to the running counter, exactly as these
      % wrappers behaved before, so every legacy one-argument call site in
      % the book keeps its label.  The counter is also csname-safe, which a
      % title's first word is not -- titles beginning with math would
      % otherwise build \csname \Omega...\endcsname and fail.
      \IfValueTF{##4}{%
        \edef\@relbl{##4}%
        \IfBooleanTF{##1}{\let\@retag\@empty}{%
          \edef\@retag{\noexpand\ \noexpand\textup{[\noexpand\texttt{#3:##4}]}}}%
      }{%
        \edef\@relbl{\arabic{#4}}%
        \let\@retag\@empty
      }%
      \def\@tempshort{##3}%
      \ifx\@tempshort\@empty
        \expandafter\xdef\csname @sshort@#3\@relbl\endcsname{\@relbl}%
      \else
        \expandafter\def\csname @sshort@#3\@relbl\endcsname{##3}%
      \fi
      \begin{#2}[##2\ifshowkeys\@retag\fi]%
      \edef\@tmplbl{#3:\@relbl}\expandafter\label\expandafter{\@tmplbl}%
      \expanded{%
        \noexpand\newglossaryentry{#3\csname the#4\endcsname}{%
          name={\unexpanded{##2}%
            \IfBooleanF{##1}{\IfValueT{##4}{ [\noexpand\texttt{#3:##4}]}}},%
          sort={\csname the#2\endcsname},%
          description={p.~\noexpand\pageref{#3:\@relbl}},%
          user1={\csname the#2\endcsname},%
          user2={#5},%
          user3={\thesection},%
          user4={\expandafter\unexpanded\expandafter{\@savedSecTitle}}%
        }%
      }%
      \protected@write\@auxout{}{\string\glsliveentry{#3\csname the#4\endcsname}}%
      \protected@write\@auxout{}{\string\@thmtitle{\@tmplbl}{\unexpanded{##2}}}%
    \fi
  }%
  \expandafter\newcommand\csname ee#1\endcsname{\end{#2}}%
}
\@mkresult{L}{lemma}{l}{lemcounter}{Lemma}
% REMARK / COROLLARY / PROPOSITION -- same signature as \beL above
\@mkresult{R}{remark}{r}{remcounter}{Remark}
\@mkresult{C}{corollary}{c}{corcounter}{Corollary}
\@mkresult{P}{proposition}{p}{propcounter}{Proposition}
% OPEN-PROBLEM COMMAND -- generated with the same signature as \beT:
%   \beO[long title][short title][key]   ...   \eeO
% label o:KEY, or o:<counter> when no key is given; a bare \beO makes
% neither a label nor a glossary entry.
\@mkresult{O}{open}{o}{opcounter}{Problem}


% Section-level title tracker: updated only by \section, not \subsection.
% \@currentlabelname is overwritten by every \refstepcounter (including
% \subsection), so we save it separately here.
\newcommand{\@currentsectionname}{}
\AtBeginDocument{%
  \let\@epidcrn@orig@refstepcounter\refstepcounter
  \renewcommand{\refstepcounter}[1]{%
    \@epidcrn@orig@refstepcounter{#1}%
    \ifstrequal{#1}{section}{\global\let\@currentsectionname\@currentlabelname}{}%
  }%
}

% TABLE COMMANDS (glossaried) -- consistent naming with \beT, \beD, etc.
\newcommand{\beTa}[1][]{%
  \def\@temparg{#1}%
  \ifx\@temparg\@empty
    \begin{table}[htbp]%
  \else
    \let\@savedSecTitle\@currentsectionname%
    \refstepcounter{tabcounter}%
    \begin{table}[htbp]%
    \label{A:\thetabcounter}%
    \expanded{%
      \noexpand\newglossaryentry{A\thetabcounter}{%
        name={\unexpanded{#1}},%
        sort={\thetabcounter},%
        description={p.~\noexpand\pageref{A:\thetabcounter}},%
        user1={\thetabcounter},%
        user2={Table},%
        user3={\thesection},%
        user4={\expandafter\unexpanded\expandafter{\@savedSecTitle}}%
      }%
    }%
    \protected@write\@auxout{}{\string\glsliveentry{A\thetabcounter}}%
  \fi
}
\newcommand{\eeTa}{\end{table}}

% FIGURE COMMANDS (glossaried) -- consistent with \beT, \beR, \beTa, etc.
% Usage:
%   \beF[Glossary title]...\caption{Visible caption}\eeF  -> glossaried figure
%   \beF...\caption{Visible caption}\eeF                  -> plain [H] figure, no glossary
\newcommand{\beF}[1][]{%
  \def\@temparg{#1}%
  \ifx\@temparg\@empty
    \begin{figure}[H]%
  \else
    \begin{figure}[H]%
    \let\@savedSecTitle\@currentsectionname%
    \refstepcounter{figcounter}%
    \label{fig:\thefigcounter}%
    \expanded{%
      \noexpand\newglossaryentry{Fig\thefigcounter}{%
        name={\unexpanded{#1}},%
        sort={\thefigcounter},%
        description={p.~\noexpand\pageref{fig:\thefigcounter}},%
        user1={\thefigcounter},%
        user2={Figure},%
        user3={\thesection},%
        user4={\expandafter\unexpanded\expandafter{\@savedSecTitle}}%
      }%
    }%
    \protected@write\@auxout{}{\string\glsliveentry{Fig\thefigcounter}}%
  \fi
}
\newcommand{\eeF}{\end{figure}}
\makeatother
% LIVE-ENTRY REGISTER
% \glsaddall prints every entry the glossaries package currently knows, and
% that set includes entries restored from \jobname.glsdefs.  The package
% rewrites .glsdefs from all *currently defined* entries, so anything it
% restored is written straight back out: a deleted section kept reappearing in
% the catalogue for ever, and the .glsdefs had to be erased by hand.  The
% wrappers therefore record every entry they actually create in the .aux, and
% \printboth adds only those.  The .aux is rebuilt from scratch on each run, so
% this is self-healing: a removed result is gone after one further run.
\makeatletter
\newcommand{\@glslive}{}
\newcommand{\@glsaddlive}[1]{\ifglsentryexists{#1}{\glsadd{#1}}{}}
\newcommand{\glsliveentry}[1]{\g@addto@macro\@glslive{\@glsaddlive{#1}}}
\makeatother
% PRINT COMMAND
% No \clearpage here -- the catalogue continues on the current page.  This has
% been re-added by mistake several times; leave it out.
\makeatletter
\newcommand{\printboth}{%
  \@glslive   % add exactly the entries this document created last run
  \printglossary[style=unifiedstyle, title={0 \hspace{.22em} Definitions, Theorems, Examples and Open Problems}]%
}
\makeatother
% Named theorem reference: \tref{word} -> Thm~N(word-or-short)
% Shows stored short title if available, otherwise shows the key itself.
\newcommand{\tref}[1]{%
  Thm~\ref{t:#1}(\ifcsname @tshort@#1\endcsname\csname @tshort@#1\endcsname\else#1\fi)%
}

% Full-title theorem reference: \nref{t:KEY} prints the number and then the
% theorem's long title, so one writes "Theorem~\nref{t:blocks}" and gets
% "Theorem 3 ``Block factorization and siphons''".  Takes the complete label, like
% \ref, not the bare key like \tref.
%
% \beT writes each title to the .aux (see \@thmtitle below), so \nref resolves
% forward references from the second LaTeX pass, exactly as \ref does.  If the
% title is not yet known -- first pass, or a label that is not a \beT theorem
% -- it degrades to the number rather than producing ??.
% These two use @ in their names, and this point in the file is outside the
% \makeatletter region opened earlier (it closes at \makeatother above), so
% the region has to be reopened here.
\makeatletter
\newcommand{\@thmtitle}[2]{\global\@namedef{@tlong@#1}{#2}}
\newcommand{\nref}[1]{%
  \ref{#1}%
  \ifcsname @tlong@#1\endcsname
    \ ``\csname @tlong@#1\endcsname''%
  \fi
}
% \nm{t:KEY}: the stored long title of any \beT/\beL/\beP/\beR/\beC/\beD/\beE result,
% verbatim and without quotes, so item names follow the macro's first argument.
\newcommand{\nm}[1]{%
  \ifcsname @tlong@#1\endcsname
    \csname @tlong@#1\endcsname%
  \else
    \ref{#1}%
  \fi
}
\makeatother




%% ========================================
%% Other THEOREM ENVIRONMENTS
%% ========================================
\makeatletter
\@ifundefined{lemma}      {\newtheorem{lemma}      {Lemma}}      {}
\@ifundefined{proposition}{\newtheorem{proposition}{Proposition}} {}
\@ifundefined{corollary}  {\newtheorem{corollary}  {Corollary}}  {}
\@ifundefined{conjecture} {\newtheorem{conjecture} {Conjecture}}  {}
\@ifundefined{remark}     {\newtheorem{remark}     {Remark}}     {}
\@ifundefined{question}   {\newtheorem{question}   {Question}}   {}
\@ifundefined{assumption} {\newtheorem{assumption} {Assumption}} {}
\@ifundefined{hypothesis} {\newtheorem{hypothesis} {Hypothesis}} {}
\@ifundefined{fact}       {\newtheorem{fact}       {Fact}}       {}
\makeatother

% \citer: cite the relay-graph threshold theorem of AAHrel. If a document has
% \input the chapter/file containing t:relayRF (so the label resolves), point
% straight at it; otherwise fall back to a bare citation of AAHrel. Shared here
% (not per-paper) so it resolves consistently across every book chapter.
% \citeR: same idea, for the first theorem of AAHRah (t:Pstar).
\makeatletter
\providecommand{\citer}{%
  \@ifundefined{r@t:relayRF}%
    {\cite{AAHrel}}%
    {Theorem~\ref{t:relayRF} of \cite{AAHrel}}%
}
\providecommand{\citeR}{%
  \@ifundefined{r@t:Pstar}%
    {\cite{AAHRah}}%
    {Theorem~\ref{t:Pstar} of \cite{AAHRah}}%
}
\makeatother


% Abbreviations
\def\beQ{\begin{question}}\def\eeQ{\end{question}}
\def\beA{\begin{assumption}}\def\eeA{\end{assumption}}


%% ========================================

%% ========================================
%% COLORS & CUSTOM COMMANDS
%% ========================================
\definecolor{funccolor}{RGB}{25,25,112}
\definecolor{desccolor}{RGB}{64,64,64}
\newcommand{\func}[1]{\texttt{\textcolor{funccolor}{#1}}}
\def\eq{=}
\def\mm{-}

%% ========================================
%% MATHCAL SHORTHANDS
%% ========================================
% Project-wide convention: \mX for a capital X is always \mathcal X.
% Defined for the whole alphabet so the convention never fails to expand.
\newcommand{\mA}{\mathcal A}\newcommand{\mB}{\mathcal B}
\newcommand{\mC}{\mathcal C}\newcommand{\mD}{\mathcal D}
\newcommand{\mE}{\mathcal E}\newcommand{\mF}{\mathcal F}
\newcommand{\mG}{\mathcal G}\newcommand{\mH}{\mathcal H}
\newcommand{\mI}{\mathcal I}\newcommand{\mJ}{\mathcal J}
\newcommand{\mK}{\mathcal K}\newcommand{\mL}{\mathcal L}
\newcommand{\mM}{\mathcal M}\newcommand{\mN}{\mathcal N}
\newcommand{\mO}{\mathcal O}\newcommand{\mP}{\mathcal P}
\newcommand{\mQ}{\mathcal Q}\newcommand{\mR}{\mathcal R}
\newcommand{\mS}{\mathcal S}\newcommand{\mT}{\mathcal T}
\newcommand{\mU}{\mathcal U}\newcommand{\mV}{\mathcal V}
\newcommand{\mW}{\mathcal W}\newcommand{\mX}{\mathcal X}
\newcommand{\mY}{\mathcal Y}\newcommand{\mZ}{\mathcal Z}

%% ========================================
%% BASIC MATH ENVIRONMENTS
%% ========================================
\def\bep{\begin{pmatrix}}\def\eep{\end{pmatrix}}
\def\bev{\begin{vmatrix}}\def\eev{\end{vmatrix}}
\def\bea{\begin{eqnarray*}}\def\eea{\end{eqnarray*}}
\def\bc{\begin{cases}}\def\ec{\end{cases}}
\def\BEN{\begin{enumerate}}\def\EEN{\end{enumerate}}
\def\BI{\begin{itemize}}\def\EI{\end{itemize}}

\newcommand{\be}[1]{\begin{equation}\label{#1}}
\newcommand{\ee}{\end{equation}}
\newcommand{\beq}{\begin{eqnarray}}
\def\eeq{\end{eqnarray}}
\def\eqr{\eqref}\def\fr{\frac}\def\lbl{\label}
\def\Lra{\Longrightarrow}\def\non{\nonumber}\def\cNGM{\cite{Diek,Van,Van08}}
\def\Eq{\Leftrightarrow}\def\b{\beta}\def\uph{\red{UP TO HERE}}
%% ========================================
%% PROBABILITY & STATISTICS
%% ========================================
\renewcommand{\P}{\ensuremath{\mathbb{P}}}  % Note: overrides paragraph symbol
\newcommand{\E}{\ensuremath{\mathbb{E}}}
\newcommand{\ind}[1]{\ensuremath{\mathbbm{1}_{\left\{#1\right\}}}}
\DeclareMathOperator{\Cov}{Cov}
\DeclareMathOperator{\TV}{TV}
\newcommand{\rank}{\operatorname{rank}}

%% ========================================
%% CRN OPERATORS (organizations/ERCs/buffering/siphons framework)
%% ========================================
\def\clos{\mathrm{clos}}
\def\Supply{\mathrm{Supply}}
\def\req{\mathrm{req}}
\def\React{\mathrm{React}}
\def\Prod{\mathrm{Prod}}
\def\Gen{\mathrm{Gen}}
\def\In{\mathrm{In}}
\def\Out{\mathrm{Out}}
\def\Act{\mathrm{Act}}
\def\Min{\mathrm{Min}}

%% ========================================
%% LINE BREAKING
%% ========================================
% Last-resort stretch, used only on paragraphs TeX cannot otherwise break:
% a long \texttt{} token, a comma-separated inline math list (commas are not
% breakpoints in math), or a run of \cite/\ref boxes.  Without it those set as
% overfull \hbox and print into the margin; with it they set slightly loose,
% which is the lesser evil.  Does not affect paragraphs that already break.
\emergencystretch=3em

%% ========================================
%% NUMBER SETS
%% ========================================
\newcommand{\R}{\mathbb{R}}  % Real numbers
\newcommand{\C}{\mathbb{C}}  % Complex numbers
\newcommand{\N}{\mathbb{N}}
\newcommand{\Z}{\mathbb{Z}}
% \providecommand, not \newcommand: some papers define these in their own
% preamble after % Shared preamble for the GLVbook papers.  This is the single canonical copy;
% a driver one level below GLVbook reaches it as
%
%     \input{../basicFiles/preamble}
%
% What stays in the driver, because it differs from paper to paper:
% \documentclass, font choice, \title, \author, \date, and any macro genuinely
% local to one paper.  Everything else belongs here.  The file is
% self-contained: it does not read any other file.

\usepackage[utf8]{inputenc}
\usepackage[novbox]{pdfsync}   % required: SyncTeX alone does not work here

% Page geometry, set once here for every paper that loads this file.  A driver
% that needs different margins says \geometry{...}; a second
% \usepackage[...]{geometry} raises an option clash and must not be used.
\makeatletter
\@ifundefined{AAdefNnogeometry}{\usepackage[a4paper,margin=2cm]{geometry}}{}
\makeatother

% mathematics
\usepackage{amsmath,amssymb,amsthm,mathtools,bm,mathrsfs,bbm}
\usepackage{centernot}

% floats, tables, lists
\usepackage{graphicx,float}
\makeatletter
\@ifpackageloaded{subfig}{}{\usepackage{subcaption}}
\makeatother
\usepackage{booktabs,longtable,threeparttable,array,multirow}
\usepackage{enumitem}

% colour, code, algorithms
\usepackage{xcolor,listings,tcolorbox}
\makeatletter
\@ifundefined{c@algorithm}{\usepackage{algorithm}}{}
\makeatother
\usepackage{algorithmicx,algpseudocode}

% diagrams
\usepackage{tikz,tikz-cd,pgfplots}
\usetikzlibrary{graphs,decorations.pathmorphing,decorations.markings,
arrows.meta,automata,positioning}
\tikzcdset{every label/.append style = {font = \small}}
\pgfplotsset{compat=1.18}

% chemistry
\usepackage[version=4]{mhchem}

\usepackage{microtype}
\usepackage{lineno}
\usepackage{nomencl}
\makenomenclature

% hyperref first, then bookmark, then glossaries: glossaries must follow
% hyperref or the entry links are built against the wrong anchors.
\makeatletter
\@ifundefined{hyperrefpackage}{%
  \@ifundefined{AAdefNnohyperref}{\usepackage{hyperref}}{}%
}{}
\@ifundefined{hyperrefpackage}{\usepackage{bookmark}}{}
\makeatother
\usepackage[toc,sort=use]{glossaries}
\makeglossaries

% Book/standalone switch.  Material that the book already covers in its
% prelude, but that a standalone paper still has to state, is wrapped as
%     \ifbookbuild\else ... \fi
% Drivers default to standalone; bookM.tex says \bookbuildtrue.
\newif\ifbookbuild
\bookbuildfalse

% Part switches for the gradostat paper.  Both default to true; a driver that
% wants the short submission says \madecfalse and/or \slowfastfalse, and every
% cross-reference into the dropped part disappears with it.
%   \ifmadec     Part III, the Lagasquie--Madec model and its break-even geometry
%   \ifslowfast  Part IV, everything carrying the small parameter \eps
\newif\ifmadec     \madectrue
\newif\ifslowfast  \slowfasttrue

% Part switch for the Barrat/simplicial-contagion project.  Defaults to
% true (the hypernetwork-synchrony part is included); the journal-paper
% driver says \hnetpartfalse to drop it and every cross-reference into it.
\newif\ifhnetpart  \hnetparttrue

%% ========================================
%% PACKAGES
%% ========================================
\RequirePackage{etoolbox}
\usepackage{makecell}
\usepackage[novbox]{pdfsync}% obsolete; SyncTeX is built into modern engines
\usepackage{amssymb, amsfonts,mathtools,bbm,eurosym,graphicx,latexsym, amsmath,amsthm,times,nccmath,url}
\usepackage{listings,verbatim,cases,tabularx}
% hyperref is deliberately NOT in the line above.  Journal classes often
% load it themselves, with a long option list, from a begindocument hook
% (mdpi.cls does: \DeclareHookRule{begindocument}{hyperref}{before}{mdpi}).
% Loading it here first, without options, makes that later load raise
% "Option clash for package hyperref".  So load it only if nobody else
% has, or will:
%   - \@ifpackageloaded  catches a class that loaded it already;
%   - \hyperrefpackage    is mdpi.cls's own deferred loader;
%   - \AAdefNnohyperref   lets any other driver opt out by defining it
%                         before \input{preamble}.
\makeatletter
\@ifpackageloaded{hyperref}{}{%
  \@ifundefined{hyperrefpackage}{%
    \@ifundefined{AAdefNnohyperref}{\usepackage{hyperref}}{}%
  }{}%
}
\makeatother
% linktoc=all makes a table-of-contents page number a link as well as the
% entry text; hyperref's default, linktoc=section, links only the text.
% Deferred to begin-document and guarded: by the block above hyperref may
% be loaded later (mdpi's own hook) or not at all (\AAdefNnohyperref), and
% \hypersetup must exist before it is called.  A driver that wants the
% default back can say \hypersetup{linktoc=section} after \input{preamble}.
\makeatletter
\AtBeginDocument{\@ifundefined{hypersetup}{}{\hypersetup{linktoc=all}}}
\makeatother
\usepackage{graphicx,float,fancybox}
\usepackage{epsfig,epstopdf,tikz}
\usepackage[sans]{dsfont}
\usetikzlibrary{positioning,arrows,arrows.meta,calc}
% Configure listings for Mathematica code
\lstdefinestyle{mathematica}{
 language=Mathematica,
 basicstyle=\small\ttfamily,
 backgroundcolor=\color{gray!10},
 frame=single,
 breaklines=true,
 commentstyle=\color{green!60!black},
 keywordstyle=\color{blue},
 stringstyle=\color{red},
 showstringspaces=false,
 tabsize=2
}
\usepackage{tcolorbox}
% Colored boxes for key results
\newtcolorbox{keyresult}[1][]{
 colback=blue!5!white,
 colframe=blue!75!black,
 title=#1,
 fonttitle=\bfseries
}
\newtcolorbox{examplebox}[1][]{
 colback=green!5!white,
 colframe=green!75!black,
 title=#1,
 fonttitle=\bfseries
}


\makeatletter
\@ifundefined{theorem}   {\newtheorem{theorem}   {Theorem}}   {}
\@ifundefined{definition}{\newtheorem{definition}{Definition}} {}
\@ifundefined{open}      {\newtheorem{open}      {Problem}}   {}
\@ifundefined{example}   {\newtheorem{example}   {Example}}   {}
\makeatother

% COUNTER DEFINITIONS
\newcounter{tabcounter}
\newcounter{defcounter}
\newcounter{opcounter}
\newcounter{thmcounter}
\newcounter{excounter}
\newcounter{ascounter}
\newcounter{cjcounter}
\newcounter{lemcounter}
\newcounter{remcounter}
\newcounter{corcounter}
\newcounter{propcounter}
\newcounter{figcounter}

% UNIFIED GLOSSARY STYLE
\newglossarystyle{unifiedstyle}{%
  \setglossarystyle{list}%
  \renewcommand*{\glossarypreamble}{\gdef\glslastsec{}}%
  \renewcommand{\glossentry}[2]{%
    \edef\glscursec{\glsentryuseriii{##1}}%
    \ifx\glscursec\glslastsec\else
      \item[]\textbf{\S\,\glscursec\enspace\glsentryuseriv{##1}}%
      \global\let\glslastsec\glscursec
    \fi
    \item[\glsentryuserii{##1} \glsentryuseri{##1}:] ``\glossentryname{##1}'' \glsentrydesc{##1}%
  }%
}

\makeatletter
% First-word extractor (used by \beT, \beD and \beXa for the default key)
% \@firstwordof hello world \relax  ->  hello
\def\@firstwordof#1 #2\relax{#1}
% First three letters, used for the heading mnemonic of a theorem given no
% explicit key.  Padding with \@empty keeps it safe for first words shorter
% than three letters: \@empty vanishes under \edef, whereas \relax would
% survive into the label and is illegal inside \csname (it broke every
% two-letter title, e.g. "ME ...", with "Missing \endcsname inserted").
\def\@takethree#1#2#3#4\@nil{#1#2#3}

% EXAMPLE COMMAND -- same signature and behaviour as \beT and \beD.
% Usage: \beXa*[long title][short title][key]
%   star  = suppresses the printed "[key]" tag after the title (as for \beT*
%   and \beD*); the label and glossary entry are unaffected.
%   key = optional label key override; the label becomes e:KEY, and without a
%   key it is e: plus the first three letters of the title's first word.  A
%   mnemonic label lets an example be cited as "Example~\ref{e:KEY}" instead
%   of by section number, which is what the citing prose needs.
\NewDocumentCommand{\beXa}{s O{} O{} o}{%
  \def\@temparg{#2}%
  \ifx\@temparg\@empty
    \begin{example}%
  \else
    \let\@savedSecTitle\@currentsectionname%
    \refstepcounter{excounter}%
    \IfValueTF{#4}{%
      \edef\@exlblkey{#4}%
    }{%
      \edef\@exfirstword{\@firstwordof#2 \relax}%
      \edef\@exlblkey{%
        \expandafter\@takethree\@exfirstword\@empty\@empty\@empty\@nil}%
    }%
    \def\@tempshort{#3}%
    \ifx\@tempshort\@empty
      \expandafter\xdef\csname @xshort@\@exlblkey\endcsname{\@exlblkey}%
    \else
      \expandafter\def\csname @xshort@\@exlblkey\endcsname{#3}%
    \fi
    \IfBooleanTF{#1}{\let\@extag\@empty}{%
      \edef\@extag{\noexpand\ \noexpand\textup{[\noexpand\texttt{e:\@exlblkey}]}}}%
    \begin{example}[#2\ifshowkeys\@extag\fi]%
    \label{e:\@exlblkey}%
    % record the long title so \nref{e:KEY} prints it, as for theorems
    \protected@write\@auxout{}{%
      \string\@thmtitle{e:\@exlblkey}{\unexpanded{#2}}}%
    \expanded{%
      \noexpand\newglossaryentry{e\theexcounter}{%
        name={\unexpanded{#2}%
          \IfBooleanF{#1}{ [\noexpand\texttt{e:\@exlblkey}]}},%
        sort={\theexample},%
        description={p.~\noexpand\pageref{e:\@exlblkey}},%
        user1={\theexample},%
        user2={Example},%
        user3={\thesection},%
        user4={\expandafter\unexpanded\expandafter{\@savedSecTitle}}%
      }%
    }%
    \protected@write\@auxout{}{\string\glsliveentry{e\theexcounter}}%
  \fi
}
\newcommand{\eeXa}{\end{example}}

% DEFINITION COMMAND -- same signature and behaviour as \beT below.
% Usage: \beD[long title][short title][key]
%   long title  = definition heading (plain text only)
%   short title = shown in \dref{word}; defaults to the key
%   key         = optional label key override; the label becomes d:KEY
% Without a key the label is d: plus the first three letters of the title's
% first word, so it is a mnemonic rather than the definition's number: a
% numeric label would silently change whenever a definition is inserted.
\NewDocumentCommand{\beD}{s O{} O{} o}{%
  \def\@temparg{#2}%
  \ifx\@temparg\@empty
    \begin{definition}%
  \else
    \let\@savedSecTitle\@currentsectionname%
    \refstepcounter{defcounter}%
    \IfValueTF{#4}{%
      \edef\@deflblkey{#4}%
    }{%
      \edef\@deffirstword{\@firstwordof#2 \relax}%
      \edef\@deflblkey{%
        \expandafter\@takethree\@deffirstword\@empty\@empty\@empty\@nil}%
    }%
    \def\@tempshort{#3}%
    \ifx\@tempshort\@empty
      \expandafter\xdef\csname @dshort@\@deflblkey\endcsname{\@deflblkey}%
    \else
      \expandafter\def\csname @dshort@\@deflblkey\endcsname{#3}%
    \fi
    \IfBooleanTF{#1}{\let\@deftag\@empty}{%
      \edef\@deftag{\noexpand\ \noexpand\textup{[\noexpand\texttt{d:\@deflblkey}]}}}%
    \begin{definition}[#2\ifshowkeys\@deftag\fi]%
    \label{d:\@deflblkey}%
    % record the long title so \nref{d:KEY} prints it, as for theorems
    \protected@write\@auxout{}{%
      \string\@thmtitle{d:\@deflblkey}{\unexpanded{#2}}}%
    \expanded{%
      \noexpand\newglossaryentry{d\thedefcounter}{%
        name={\unexpanded{#2}%
          \IfBooleanF{#1}{ [\noexpand\texttt{d:\@deflblkey}]}},%
        sort={\thedefinition},%
        description={p.~\noexpand\pageref{d:\@deflblkey}},%
        user1={\thedefinition},%
        user2={Definition},%
        user3={\thesection},%
        user4={\expandafter\unexpanded\expandafter{\@savedSecTitle}}%
      }%
    }%
    \protected@write\@auxout{}{\string\glsliveentry{d\thedefcounter}}%
  \fi
}
\newcommand{\eeD}{\end{definition}}

% \showkeysfalse in a driver hides the [t:KEY] mnemonics in result headings
% (labels and glossary entries are unchanged); used for submission builds.
\newif\ifshowkeys \showkeystrue
% \commonredtrue in a driver that inputs basicFiles/commonRed: tool.tex then
% cites its reduction results instead of restating them.
\newif\ifcommonred \commonredfalse
% TheorEM COMMAND
% Usage: \beT[long title][short title][key]
%   long title  = theorem heading (plain text only)
%   short title = shown in \tref{word}; may contain math; defaults to first word of long title
%   key         = optional label key override; label becomes t:KEY (avoids first-word collisions)
% Label is t:WORD where WORD = first word of long title (or KEY if given); use \tref{WORD} to reference.
\NewDocumentCommand{\beT}{s O{} O{} o}{%
  \def\@temparg{#2}%
  \ifx\@temparg\@empty
    \begin{theorem}%
  \else
    \let\@savedSecTitle\@currentsectionname%
    \refstepcounter{thmcounter}%
    % \@thmtag is the mnemonic shown in the heading; every theorem gets one.
    % With an explicit key the key itself is shown.  Without one the key
    % defaults to the title's first word, and showing that whole word would
    % merely repeat the title ("Dissipativity [Dissipativity]"), so only its
    % first three letters are shown: "Dissipativity [Dis]".  The label is
    % unaffected -- it still uses the full \@thmlblkey.
    \IfValueTF{#4}{%
      \edef\@thmlblkey{#4}%
    }{%
      \edef\@thmfirstword{\@firstwordof#2 \relax}%
      \edef\@thmlblkey{%
        \expandafter\@takethree\@thmfirstword\@empty\@empty\@empty\@nil}%
    }%
    \IfBooleanTF{#1}{\let\@thmtag\@empty}{%
      \edef\@thmtag{\noexpand\ \noexpand\textup{[\noexpand\texttt{t:\@thmlblkey}]}}}%
    \def\@tempshort{#3}%
    \ifx\@tempshort\@empty
      \expandafter\xdef\csname @tshort@\@thmlblkey\endcsname{\@thmlblkey}%
    \else
      \expandafter\def\csname @tshort@\@thmlblkey\endcsname{#3}%
    \fi
    \begin{theorem}[#2\ifshowkeys\@thmtag\fi]%
    \label{t:\@thmlblkey}%
    % record the long title for \nref; via the .aux so forward refs resolve
    \protected@write\@auxout{}{%
      \string\@thmtitle{t:\@thmlblkey}{\unexpanded{#2}}}%
    \expanded{%
      \noexpand\newglossaryentry{t\thethmcounter}{%
        name={\unexpanded{#2}%
          \IfBooleanF{#1}{ [\noexpand\texttt{t:\@thmlblkey}]}},%
        % sort/user1 must be the number the document actually prints, i.e. the
        % theorem environment's counter, not thmcounter.  thmcounter advances
        % only on \beT with a title, so a bare \beT (or \beT with an
        % empty title) makes the two drift and the glossary cites a wrong
        % number.  thmcounter is still fine as the entry key: it is unique.
        sort={\thetheorem},%
        description={p.~\noexpand\pageref{t:\@thmlblkey}},%
        user1={\thetheorem},%
        user2={Theorem},%
        user3={\thesection},%
        user4={\expandafter\unexpanded\expandafter{\@savedSecTitle}}%
      }%
    }%
    \protected@write\@auxout{}{\string\glsliveentry{t\thethmcounter}}%
  \fi
}
\newcommand{\eeT}{\end{theorem}}

% CROSS-PART RELATION PARAGRAPH
% Usage: \xpart{text}  -- place immediately BEFORE \beT, never inside it
%   (anything inside \beT...\eeT would be swallowed into the glossary entry).
% The paragraph is typeset only when Part I is actually present, i.e. in the
% combined book build; in the standalone builds (GranadaDSABNS/A.tex,
% AAHRah-Submission/AAHRah.tex) the label part:glv is undefined and the whole
% paragraph disappears rather than leaving ?? cross-references. Same
% graceful-degradation principle as \GASref in GASt.tex. Resolves on the
% second LaTeX pass, since r@part:glv comes from the .aux file.
% (Already inside the \makeatletter region opened above, so no \makeatletter
% pair here -- adding one would close that region early.)
\newcommand{\xpart}[1]{\@ifundefined{r@part:glv}{}{\paragraph{Relation with Parts~I--II.}#1}}

% \bookref{label}{precise text}{fallback text}: prints the precise cross-reference
% when the label resolves (combined book build), and the fallback text otherwise
% (standalone paper builds, where book-only appendix labels are absent). Same
% graceful-degradation role as \GASref in GASt.tex, but usable in any file since it
% lives in this preamble (loaded by every driver) -- unlike \GASref, which is defined in a
% content file and is therefore undefined in standalone builds that omit that file.
\newcommand{\bookref}[3]{\@ifundefined{r@#1}{#3}{#2}}

% \PVpart and \RBpart: the two parts that a partial build may omit.
% They print "Part~N" in any build that contains the part, and a naming phrase
% in a build that does not, so an excerpt shows no ?? .  Each already supplies
% the word "Part", so write "\PVpart develops ...", not "Part~\PVpart".
\newcommand{\PVpart}{\bookref{part:aahrah}%
  {Part~\ref{part:aahrah}}{the Perron--Volterra part of the full book}}
\newcommand{\RBpart}{\bookref{part:recbrowne}%
  {Part~\ref{part:recbrowne}}{the immuno-virology part of the full book}}

% OPEN PROBLEM COMMANDS
% ASSUMPTION COMMANDS (glossaried)
\newcommand{\beAs}[1][]{%
  \def\@temparg{#1}%
  \ifx\@temparg\@empty
    \begin{assumption}%
  \else
    \let\@savedSecTitle\@currentsectionname%
    \refstepcounter{ascounter}%
    \begin{assumption}[#1]%
    \label{as:\theascounter}%
    \expanded{%
      \noexpand\newglossaryentry{as\theascounter}{%
        name={\unexpanded{#1}},%
        sort={\theascounter},%
        description={p.~\noexpand\pageref{as:\theascounter}},%
        user1={\theascounter},%
        user2={Assumption},%
        user3={\thesection},%
        user4={\expandafter\unexpanded\expandafter{\@savedSecTitle}}%
      }%
    }%
    \protected@write\@auxout{}{\string\glsliveentry{as\theascounter}}%
  \fi
}
\newcommand{\eeAs}{\end{assumption}}
% CONJECTURE COMMANDS (glossaried)
\newcommand{\beCj}[1][]{%
  \def\@temparg{#1}%
  \ifx\@temparg\@empty
    \begin{conjecture}%
  \else
    \let\@savedSecTitle\@currentsectionname%
    \refstepcounter{cjcounter}%
    \begin{conjecture}[#1]%
    \label{cj:\thecjcounter}%
    \expanded{%
      \noexpand\newglossaryentry{cj\thecjcounter}{%
        name={\unexpanded{#1}},%
        sort={\thecjcounter},%
        description={p.~\noexpand\pageref{cj:\thecjcounter}},%
        user1={\thecjcounter},%
        user2={Conjecture},%
        user3={\thesection},%
        user4={\expandafter\unexpanded\expandafter{\@savedSecTitle}}%
      }%
    }%
    \protected@write\@auxout{}{\string\glsliveentry{cj\thecjcounter}}%
  \fi
}
\newcommand{\eeCj}{\end{conjecture}}
% LEMMA COMMANDS
% ---------------------------------------------------------------------
% UNIFORM RESULT WRAPPERS -- \beL, \beC, \beP, \beR share \beT's signature
%
%   \beX[long title][short title][key]  ...  \eeX
%
%   long title  = heading (math allowed)
%   short title = shown in the cross-reference; defaults to the key
%   key         = label key.  With a key the label is X:KEY and the heading
%                 carries the mnemonic [key], exactly as for \beT.  Without
%                 one the label falls back to the running counter, which is
%                 what every pre-existing call site in the book relies on.
%
% A bare \beX -- arguments commented out -- opens the plain environment and
% creates NO label and NO glossary entry.  That is the supported way to keep
% an individual result out of the catalogue.
% ---------------------------------------------------------------------
% #1 wrapper name  #2 environment  #3 label prefix  #4 counter  #5 type name
\newcommand{\@mkresult}[5]{%
  % A leading * suppresses the [key] mnemonic in both the heading and the
  % glossary name, while still making the label X:KEY and the entry:
  %     \beX*[long][short][key]
  \expandafter\NewDocumentCommand\csname be#1\endcsname{s O{} O{} o}{%
    \def\@temparg{##2}%
    \ifx\@temparg\@empty
      \begin{#2}%
    \else
      \let\@savedSecTitle\@currentsectionname%
      \refstepcounter{#4}%
      % With a key: mnemonic label X:KEY and the [key] tag in the heading.
      % Without one: fall back to the running counter, exactly as these
      % wrappers behaved before, so every legacy one-argument call site in
      % the book keeps its label.  The counter is also csname-safe, which a
      % title's first word is not -- titles beginning with math would
      % otherwise build \csname \Omega...\endcsname and fail.
      \IfValueTF{##4}{%
        \edef\@relbl{##4}%
        \IfBooleanTF{##1}{\let\@retag\@empty}{%
          \edef\@retag{\noexpand\ \noexpand\textup{[\noexpand\texttt{#3:##4}]}}}%
      }{%
        \edef\@relbl{\arabic{#4}}%
        \let\@retag\@empty
      }%
      \def\@tempshort{##3}%
      \ifx\@tempshort\@empty
        \expandafter\xdef\csname @sshort@#3\@relbl\endcsname{\@relbl}%
      \else
        \expandafter\def\csname @sshort@#3\@relbl\endcsname{##3}%
      \fi
      \begin{#2}[##2\ifshowkeys\@retag\fi]%
      \edef\@tmplbl{#3:\@relbl}\expandafter\label\expandafter{\@tmplbl}%
      \expanded{%
        \noexpand\newglossaryentry{#3\csname the#4\endcsname}{%
          name={\unexpanded{##2}%
            \IfBooleanF{##1}{\IfValueT{##4}{ [\noexpand\texttt{#3:##4}]}}},%
          sort={\csname the#2\endcsname},%
          description={p.~\noexpand\pageref{#3:\@relbl}},%
          user1={\csname the#2\endcsname},%
          user2={#5},%
          user3={\thesection},%
          user4={\expandafter\unexpanded\expandafter{\@savedSecTitle}}%
        }%
      }%
      \protected@write\@auxout{}{\string\glsliveentry{#3\csname the#4\endcsname}}%
      \protected@write\@auxout{}{\string\@thmtitle{\@tmplbl}{\unexpanded{##2}}}%
    \fi
  }%
  \expandafter\newcommand\csname ee#1\endcsname{\end{#2}}%
}
\@mkresult{L}{lemma}{l}{lemcounter}{Lemma}
% REMARK / COROLLARY / PROPOSITION -- same signature as \beL above
\@mkresult{R}{remark}{r}{remcounter}{Remark}
\@mkresult{C}{corollary}{c}{corcounter}{Corollary}
\@mkresult{P}{proposition}{p}{propcounter}{Proposition}
% OPEN-PROBLEM COMMAND -- generated with the same signature as \beT:
%   \beO[long title][short title][key]   ...   \eeO
% label o:KEY, or o:<counter> when no key is given; a bare \beO makes
% neither a label nor a glossary entry.
\@mkresult{O}{open}{o}{opcounter}{Problem}


% Section-level title tracker: updated only by \section, not \subsection.
% \@currentlabelname is overwritten by every \refstepcounter (including
% \subsection), so we save it separately here.
\newcommand{\@currentsectionname}{}
\AtBeginDocument{%
  \let\@epidcrn@orig@refstepcounter\refstepcounter
  \renewcommand{\refstepcounter}[1]{%
    \@epidcrn@orig@refstepcounter{#1}%
    \ifstrequal{#1}{section}{\global\let\@currentsectionname\@currentlabelname}{}%
  }%
}

% TABLE COMMANDS (glossaried) -- consistent naming with \beT, \beD, etc.
\newcommand{\beTa}[1][]{%
  \def\@temparg{#1}%
  \ifx\@temparg\@empty
    \begin{table}[htbp]%
  \else
    \let\@savedSecTitle\@currentsectionname%
    \refstepcounter{tabcounter}%
    \begin{table}[htbp]%
    \label{A:\thetabcounter}%
    \expanded{%
      \noexpand\newglossaryentry{A\thetabcounter}{%
        name={\unexpanded{#1}},%
        sort={\thetabcounter},%
        description={p.~\noexpand\pageref{A:\thetabcounter}},%
        user1={\thetabcounter},%
        user2={Table},%
        user3={\thesection},%
        user4={\expandafter\unexpanded\expandafter{\@savedSecTitle}}%
      }%
    }%
    \protected@write\@auxout{}{\string\glsliveentry{A\thetabcounter}}%
  \fi
}
\newcommand{\eeTa}{\end{table}}

% FIGURE COMMANDS (glossaried) -- consistent with \beT, \beR, \beTa, etc.
% Usage:
%   \beF[Glossary title]...\caption{Visible caption}\eeF  -> glossaried figure
%   \beF...\caption{Visible caption}\eeF                  -> plain [H] figure, no glossary
\newcommand{\beF}[1][]{%
  \def\@temparg{#1}%
  \ifx\@temparg\@empty
    \begin{figure}[H]%
  \else
    \begin{figure}[H]%
    \let\@savedSecTitle\@currentsectionname%
    \refstepcounter{figcounter}%
    \label{fig:\thefigcounter}%
    \expanded{%
      \noexpand\newglossaryentry{Fig\thefigcounter}{%
        name={\unexpanded{#1}},%
        sort={\thefigcounter},%
        description={p.~\noexpand\pageref{fig:\thefigcounter}},%
        user1={\thefigcounter},%
        user2={Figure},%
        user3={\thesection},%
        user4={\expandafter\unexpanded\expandafter{\@savedSecTitle}}%
      }%
    }%
    \protected@write\@auxout{}{\string\glsliveentry{Fig\thefigcounter}}%
  \fi
}
\newcommand{\eeF}{\end{figure}}
\makeatother
% LIVE-ENTRY REGISTER
% \glsaddall prints every entry the glossaries package currently knows, and
% that set includes entries restored from \jobname.glsdefs.  The package
% rewrites .glsdefs from all *currently defined* entries, so anything it
% restored is written straight back out: a deleted section kept reappearing in
% the catalogue for ever, and the .glsdefs had to be erased by hand.  The
% wrappers therefore record every entry they actually create in the .aux, and
% \printboth adds only those.  The .aux is rebuilt from scratch on each run, so
% this is self-healing: a removed result is gone after one further run.
\makeatletter
\newcommand{\@glslive}{}
\newcommand{\@glsaddlive}[1]{\ifglsentryexists{#1}{\glsadd{#1}}{}}
\newcommand{\glsliveentry}[1]{\g@addto@macro\@glslive{\@glsaddlive{#1}}}
\makeatother
% PRINT COMMAND
% No \clearpage here -- the catalogue continues on the current page.  This has
% been re-added by mistake several times; leave it out.
\makeatletter
\newcommand{\printboth}{%
  \@glslive   % add exactly the entries this document created last run
  \printglossary[style=unifiedstyle, title={0 \hspace{.22em} Definitions, Theorems, Examples and Open Problems}]%
}
\makeatother
% Named theorem reference: \tref{word} -> Thm~N(word-or-short)
% Shows stored short title if available, otherwise shows the key itself.
\newcommand{\tref}[1]{%
  Thm~\ref{t:#1}(\ifcsname @tshort@#1\endcsname\csname @tshort@#1\endcsname\else#1\fi)%
}

% Full-title theorem reference: \nref{t:KEY} prints the number and then the
% theorem's long title, so one writes "Theorem~\nref{t:blocks}" and gets
% "Theorem 3 ``Block factorization and siphons''".  Takes the complete label, like
% \ref, not the bare key like \tref.
%
% \beT writes each title to the .aux (see \@thmtitle below), so \nref resolves
% forward references from the second LaTeX pass, exactly as \ref does.  If the
% title is not yet known -- first pass, or a label that is not a \beT theorem
% -- it degrades to the number rather than producing ??.
% These two use @ in their names, and this point in the file is outside the
% \makeatletter region opened earlier (it closes at \makeatother above), so
% the region has to be reopened here.
\makeatletter
\newcommand{\@thmtitle}[2]{\global\@namedef{@tlong@#1}{#2}}
\newcommand{\nref}[1]{%
  \ref{#1}%
  \ifcsname @tlong@#1\endcsname
    \ ``\csname @tlong@#1\endcsname''%
  \fi
}
% \nm{t:KEY}: the stored long title of any \beT/\beL/\beP/\beR/\beC/\beD/\beE result,
% verbatim and without quotes, so item names follow the macro's first argument.
\newcommand{\nm}[1]{%
  \ifcsname @tlong@#1\endcsname
    \csname @tlong@#1\endcsname%
  \else
    \ref{#1}%
  \fi
}
\makeatother




%% ========================================
%% Other THEOREM ENVIRONMENTS
%% ========================================
\makeatletter
\@ifundefined{lemma}      {\newtheorem{lemma}      {Lemma}}      {}
\@ifundefined{proposition}{\newtheorem{proposition}{Proposition}} {}
\@ifundefined{corollary}  {\newtheorem{corollary}  {Corollary}}  {}
\@ifundefined{conjecture} {\newtheorem{conjecture} {Conjecture}}  {}
\@ifundefined{remark}     {\newtheorem{remark}     {Remark}}     {}
\@ifundefined{question}   {\newtheorem{question}   {Question}}   {}
\@ifundefined{assumption} {\newtheorem{assumption} {Assumption}} {}
\@ifundefined{hypothesis} {\newtheorem{hypothesis} {Hypothesis}} {}
\@ifundefined{fact}       {\newtheorem{fact}       {Fact}}       {}
\makeatother

% \citer: cite the relay-graph threshold theorem of AAHrel. If a document has
% \input the chapter/file containing t:relayRF (so the label resolves), point
% straight at it; otherwise fall back to a bare citation of AAHrel. Shared here
% (not per-paper) so it resolves consistently across every book chapter.
% \citeR: same idea, for the first theorem of AAHRah (t:Pstar).
\makeatletter
\providecommand{\citer}{%
  \@ifundefined{r@t:relayRF}%
    {\cite{AAHrel}}%
    {Theorem~\ref{t:relayRF} of \cite{AAHrel}}%
}
\providecommand{\citeR}{%
  \@ifundefined{r@t:Pstar}%
    {\cite{AAHRah}}%
    {Theorem~\ref{t:Pstar} of \cite{AAHRah}}%
}
\makeatother


% Abbreviations
\def\beQ{\begin{question}}\def\eeQ{\end{question}}
\def\beA{\begin{assumption}}\def\eeA{\end{assumption}}


%% ========================================

%% ========================================
%% COLORS & CUSTOM COMMANDS
%% ========================================
\definecolor{funccolor}{RGB}{25,25,112}
\definecolor{desccolor}{RGB}{64,64,64}
\newcommand{\func}[1]{\texttt{\textcolor{funccolor}{#1}}}
\def\eq{=}
\def\mm{-}

%% ========================================
%% MATHCAL SHORTHANDS
%% ========================================
% Project-wide convention: \mX for a capital X is always \mathcal X.
% Defined for the whole alphabet so the convention never fails to expand.
\newcommand{\mA}{\mathcal A}\newcommand{\mB}{\mathcal B}
\newcommand{\mC}{\mathcal C}\newcommand{\mD}{\mathcal D}
\newcommand{\mE}{\mathcal E}\newcommand{\mF}{\mathcal F}
\newcommand{\mG}{\mathcal G}\newcommand{\mH}{\mathcal H}
\newcommand{\mI}{\mathcal I}\newcommand{\mJ}{\mathcal J}
\newcommand{\mK}{\mathcal K}\newcommand{\mL}{\mathcal L}
\newcommand{\mM}{\mathcal M}\newcommand{\mN}{\mathcal N}
\newcommand{\mO}{\mathcal O}\newcommand{\mP}{\mathcal P}
\newcommand{\mQ}{\mathcal Q}\newcommand{\mR}{\mathcal R}
\newcommand{\mS}{\mathcal S}\newcommand{\mT}{\mathcal T}
\newcommand{\mU}{\mathcal U}\newcommand{\mV}{\mathcal V}
\newcommand{\mW}{\mathcal W}\newcommand{\mX}{\mathcal X}
\newcommand{\mY}{\mathcal Y}\newcommand{\mZ}{\mathcal Z}

%% ========================================
%% BASIC MATH ENVIRONMENTS
%% ========================================
\def\bep{\begin{pmatrix}}\def\eep{\end{pmatrix}}
\def\bev{\begin{vmatrix}}\def\eev{\end{vmatrix}}
\def\bea{\begin{eqnarray*}}\def\eea{\end{eqnarray*}}
\def\bc{\begin{cases}}\def\ec{\end{cases}}
\def\BEN{\begin{enumerate}}\def\EEN{\end{enumerate}}
\def\BI{\begin{itemize}}\def\EI{\end{itemize}}

\newcommand{\be}[1]{\begin{equation}\label{#1}}
\newcommand{\ee}{\end{equation}}
\newcommand{\beq}{\begin{eqnarray}}
\def\eeq{\end{eqnarray}}
\def\eqr{\eqref}\def\fr{\frac}\def\lbl{\label}
\def\Lra{\Longrightarrow}\def\non{\nonumber}\def\cNGM{\cite{Diek,Van,Van08}}
\def\Eq{\Leftrightarrow}\def\b{\beta}\def\uph{\red{UP TO HERE}}
%% ========================================
%% PROBABILITY & STATISTICS
%% ========================================
\renewcommand{\P}{\ensuremath{\mathbb{P}}}  % Note: overrides paragraph symbol
\newcommand{\E}{\ensuremath{\mathbb{E}}}
\newcommand{\ind}[1]{\ensuremath{\mathbbm{1}_{\left\{#1\right\}}}}
\DeclareMathOperator{\Cov}{Cov}
\DeclareMathOperator{\TV}{TV}
\newcommand{\rank}{\operatorname{rank}}

%% ========================================
%% CRN OPERATORS (organizations/ERCs/buffering/siphons framework)
%% ========================================
\def\clos{\mathrm{clos}}
\def\Supply{\mathrm{Supply}}
\def\req{\mathrm{req}}
\def\React{\mathrm{React}}
\def\Prod{\mathrm{Prod}}
\def\Gen{\mathrm{Gen}}
\def\In{\mathrm{In}}
\def\Out{\mathrm{Out}}
\def\Act{\mathrm{Act}}
\def\Min{\mathrm{Min}}

%% ========================================
%% LINE BREAKING
%% ========================================
% Last-resort stretch, used only on paragraphs TeX cannot otherwise break:
% a long \texttt{} token, a comma-separated inline math list (commas are not
% breakpoints in math), or a run of \cite/\ref boxes.  Without it those set as
% overfull \hbox and print into the margin; with it they set slightly loose,
% which is the lesser evil.  Does not affect paragraphs that already break.
\emergencystretch=3em

%% ========================================
%% NUMBER SETS
%% ========================================
\newcommand{\R}{\mathbb{R}}  % Real numbers
\newcommand{\C}{\mathbb{C}}  % Complex numbers
\newcommand{\N}{\mathbb{N}}
\newcommand{\Z}{\mathbb{Z}}
% \providecommand, not \newcommand: some papers define these in their own
% preamble after \input{preamble}, and must not clash.
\providecommand{\Q}{\mathbb{Q}}
\providecommand{\res}{\operatorname{res}}   % resultant; used by epiRes.tex
\newcommand{\MN}{\mathbb{N}}
\newcommand{\MZ}{\mathbb{Z}}
\newcommand{\realpos}{\mathbb{R}_{> 0}}
\newcommand{\realnneg}{\mathbb{R}_{\geq 0}}
\newcommand{\naturalpos}{\mathbb{N}_{>0}}
\newcommand{\Rplus}{\mathbb{R}_+}
\newcommand{\Rp}{\mathbb{R}_{\ge 0}}
\def\Rnne{\mathbb{R}_{\ge 0}}
\def\Se{\bar{S}} \def\Ie{\bar{I}}
\newcommand{\Rpos}{\mathbb{R}_{>0}}
\newcommand{\dd}{\,\mathrm d}
\newcommand{\e}{\mathrm e}
%% ========================================
%% GREEK LETTERS (2+ character abbreviations) greek
%% ========================================
\def\al{\alpha}\def\ba{\bff \alpha}
\def\w{\omega}\def\Lgn{\ga+\mu+\de}\def\om{\omega}
\def\de{\delta}\def\De{\Delta}\def\bde{\bff \delta}
\def\ep{\epsilon}\def\eps{\varepsilon}
\def\f{\varphi}  % Kept as requested
\def\ga{\gamma}\def\Ga{\Gamma}\def\G{\Gamma}\def\bg{\bff \gamma}
\def\up{\upsilon}  % NOT the breve accent: no accent commands anywhere.
                   % Named \up, not \u, so the LaTeX breve stays reachable.
\def\la{\lambda}\def\La{\Lambda}
\def\si{\sigma}\def\ka{\kappa}
\def\th{\theta}\def\vt{\vartheta}
\def\ze{\zeta}\def\Gap{\Gamma^{+}} \def\Gam{\Gamma^{-}}


%% ========================================
%% SPECIAL SHORT COMMANDS (kept as requested)
%% ========================================
\def\I{\infty}
\def\T{\widetilde}
\def\H{\hat}
\newcommand{\relint}{\operatorname{relint}}

%% ========================================
%% BOLD/VECTOR NOTATION
%% ========================================
\newcommand{\bff}[1]{{\mbox{\boldmath$#1$}}}
\newcommand\X{\boldsymbol{X}}
\def\x{\boldsymbol{x}}\def\bx{\bff x}\def\xe{\bar{\x}}
\newcommand\y{\boldsymbol{y}}\def\by{\boldsymbol{y}}\def\ye{\bar{\y}}
\def\yb{\y}\def\xb{\x}
\def\xy{\x-\y}
\newcommand\br{\boldsymbol{r}}\def\re{\bar{\r}}
\newcommand\bd{\boldsymbol{d}}
\newcommand\bQ{\boldsymbol{Q}}
\newcommand{\Yb}{{\mathbf{Y}}}
\newcommand{\Ab}{\mathbf{A}}
\newcommand{\zb}{\mathbf{z}}\def\bz{\bff 0}
\newcommand{\Zb}{\mathbf{Z}}
\newcommand{\bb}{\mathbf{b}}
\def\vbe{\vec {\bff \beta}}
\def\vb{\vec {\bff b}}
\def\v1{\vec {\bff 1}}
\def\vn{\vec \de}
\def\vz{\vec 0}
\def\vi{\overset{\rightarrow}{\mathsf i}}
\def\vv{\overset{\rightarrow}{v}}
\def\vr{\; \vec {\mathsf r}}
\newcommand{\va}{\bff a}
\providecommand{\vect}[1]{\boldsymbol{#1}}   % bold vector (moved from Gradostat.tex)
\providecommand{\one}{\mathbf 1}             % all-ones vector (moved from AliSF)
\providecommand{\spr}{s}                     % spectral bound  (moved from AliSF)
\providecommand{\Ke}{K_{\rm eff}}            % effective carrying capacity
% Critical right/left eigenvectors in the Hopf projection method (AAHOSN/New.tex)
\providecommand{\bq}{\mathbf{q}}
\providecommand{\bp}{\mathbf{p}}


%% ========================================
%% MATHCAL & MATHSCR
%% ========================================
% Mathcal
\newcommand\mc[1]{\mathcal{#1}}
\def\mC{{\mathcal C}}\def\mD{\mathcal D}\def\mE{\mathcal{E}}
\def\mF{\mathcal F}\def\mG{{\mathcal G}}\def\mH{{\mathcal H}}
\def\mI{\mathcal I}\def\mL{\mathcal{L}}\def\mM{{\mathcal M}}
\def\mN{{\mathcal N}}\def\mP{\mathcal{P}}\def\mR{{\mathcal R}}
\def\mS{{\mathcal S}}\def\mT{\mathcal{T}}\def\mV{{\mathcal V}}
\def\mK{\mathcal K}\def\mA{{\mathcal A}}\def\mU{{\mathcal U}}
% Symmetric bilinear and trilinear forms of the Taylor expansion, used in the
% first Lyapunov coefficient (AAHOSN/New.tex)
\providecommand{\Bform}{\mathcal{B}}
\providecommand{\Cform}{\mathcal{C}}

% Blackboard bold.  Kept here, not in the drivers: a \providecommand in a
% driver preamble runs before this file and silently wins, so the same name
% ends up meaning different things in different builds.
\def\IC{\mathbb{C}}\def\IN{\mathbb{N}}\def\IQ{\mathbb{Q}}
\def\IR{\mathbb{R}}\def\IZ{\mathbb{Z}}\def\IG{\mathbb{G}}
\def\II{\mathbb{I}}

% Mathscr
\newcommand{\Rs}{{\mathscr R}}
\newcommand{\Ss}{{\mathscr S}}
\newcommand{\Ns}{{\mathscr N}}
\newcommand{\Cs}{\mathscr C}


%% ========================================
%% OPERATORS & FUNCTIONS
%% ========================================
\DeclareMathOperator{\tr}{tr}
\DeclareMathOperator{\diag}{Diag}
\DeclareMathOperator{\Diag}{Diag}
\DeclareMathOperator{\intr}{\mathrm{Int}}
\DeclareMathOperator{\Inv}{Inv}
\DeclareMathOperator{\spec}{spec}
\DeclareMathOperator{\AUC}{AUC}               % area under the curve
\newcommand{\ii}{\mathrm{i}}                  % imaginary unit
\newcommand\op{\operatorname}
\newcommand{\supp}{\op{supp}}
\newcommand{\diff}{\mathop{}\mathopen{}\mathrm{d}}
\newcommand{\pdd}[1]{\frac{\partial}{\partial{#1}}}
\DeclareMathOperator{\cone}{cone}
\DeclareMathOperator{\proj}{proj}

%% ========================================
%% BRACKETS & NORMS
%% ========================================
\providecommand{\pr}[1]{\left(#1\right)}
\providecommand{\pp}[1]{\left[#1\right]}
\newcommand\croc[1]{\left\langle #1\right\rangle}
\newcommand{\norm}[1]{\|#1\|}
\newcommand{\inner}[2]{\langle #1, #2 \rangle}
% Short synonym of \inner, used by AAHOSN/New.tex as $\ip{u}{v}=\bar u^\top v$
\providecommand{\ip}[2]{\langle #1 , #2 \rangle}


%% ========================================
%% FORMATTING & UTILITIES
%% ========================================
\newcommand*{\Scale}[2][4]{\scalebox{#1}{$#2$}}
\newcommand{\te}[1]{\text{#1}}
\newcommand\steq[1]{\stackrel{\text{\rm #1.}}{=}}
\newcommand\blue[1]{\textcolor{blue}{#1}}
\newcommand{\red}{\textcolor[rgb]{1.00,0.00,0.00}}
\newcommand{\green}{\textcolor[rgb]{0.00,1.00,0.00}}

% Figure command
\newcommand{\figu}[3]{%
\begin{figure}[H]
\centering
\includegraphics[scale=#3]{#1}
\caption{#2\label{f:#1}}
\end{figure}
}

% Footnote
\long\def\symbolfootnote[#1]#2{%
\begingroup
\def\thefootnote{\fnsymbol{footnote}}\footnote[#1]{#2}%
\endgroup}
\def\fn{\symbolfootnote}

\newcommand{\dsum}{\displaystyle\sum}

%% ========================================
%% CRN & EPIDEMIOLOGY SPECIFIC
%% ========================================
% CRN notation
\def\SS{\mathcal S}
\def\CC{\mathcal C}
\def\RR{\mathcal R}

% Reproduction numbers
\def\Rto{R_{1}^{E_2}}
\def\Rot{R_{2}^{E_1}}
\def\Roc{R_{1,c}}
\def\Rtc{R_{2,c}}
\def\rd{r_{0}}
\def\sd{s_{0}}
\def\dd{D_{0}}
\def\vd{v_{0}}
\def\zd{z_{0}}
\def\se{\sE}
\def\xd{x_{0}}
%% ========================================
%% SPECIAL VALUES & EQUILIBRIA
%% ========================================
\def\eE{\bar{e}}
\def\wE{\bar{w}}
\def\sE{s_{*}}
\def\Me{\T M_{ee}}
\def\Rnne{\mathbb{R}_{\ge 0}}

\newcommand{\ones}{\mathbf 1}


%% ========================================
%% ABBREVIATIONS (alphabetical by category)
%% ========================================

% -- General mathematical terms --
\def\adm{admissible}\def\admy{admissibility}
%\def\and{antisymmetric}% disabled 2026-09-30: it overrode LaTeX's author separator \and (printed "antisymmetric" between authors)
\def\ass{assumption}
\def\cadlag{c\`adl\`ag}
\def\ch{characteristic polynomial}
\def\chp{characteristic polynomial}
\def\coe{coefficient}
\def\comp{computation}
\def\con{condition}
\def\coo{coordinate}
\def\corr{corresponding}
\def\dec{decomposition}
\def\den{denominator}
\def\dete{deterministic}
\def\eig{eigenvalue}\def\eval{eigenvalue}
\def\eigv{eigenvector}\def\eiv{eigenvector}\def\evec{eigenvector}
\def\elt{element}
\def\enn{essentially nonnegative/positivity preserving}
\def\equ{equation}
\def\etal{et al.}
\def\fe{for example}
\def\fno{from now on}
\def\for{\forall}
\def\form{formula}
\def\Fr{Furthermore, }
\def\frt{furthermore}
\def\fun{function}
\def\gene{generalization}
\def\heu{heuristic}
\def\how{however}\def\How{However, }
\def\ie{i.e. }
\def\im{\item}
\def\isd{is defined by}
\def\itc{it is easy to check that}
\def\itf{it follows that}\def\Itf{It follows that}
\def\ith{it holds that}
\def\itm{it may be checked that}\def\Itm{It may be checked that}
\def\ito{it turns out that}\def\Ito{It turns out that}\def\Itou{It turns out useful to}
\def\ma{mass action}\def\MAR{mass action representation}
\def\mat{matrix}
\def\mbc{may be checked to be}
\def\mbw{may be written as}\def\mubw{must be written as}
\def\mpr{more precisely, }\def\Mp{More precisely, }
\def\nc{non-catalytic}
\def\nne{non-negative}\def\nny{nonnegativity}
\def\no{\nonumber}
\def\noni{non-infectious}
\def\nota{notation}
\def\np{\newpage}
\def\Nt{Note that}
\def\nt{non trivial}\def\nz{non zero}
\def\oth{otherwise}
\def\Oth{On the other hand}
\def\pa{parameter}\def\para{parameter}
\def\pol{polynomial}
\def\prob{probability}
\def\Prf{{\bf  Proof:}}\def\prf{{\bf  Proof:}}
\def\prop{proportional to}
\def\qu{\quad}
\def\QED{\hfill {$\square$}\goodbreak \medskip}
\def\rep{representation}
\def\resp{respectively}
\def\satd{satisfied}\def\satg{satisfying}\def\sats{satisfies}\def\saty{satisfy}
\def\ssec{\subsection}\def\sssec{\subsubsection}
\def\sef{see for example}
\def\sksy{skew-symmetric}
\def\sub{substrate}
\def\sys{system}
\def\var{variable}
\def\wk{well-known}
\def\wl{without loss of generality}\def\wlo{w.l.o.g.}
\def\wmc{We may check that}\def\Wmc{We may check that}
\def\wmr{we may rewrite}
\def\wms{we must show that}
\def\wrt{with respect to}

% -- CRN & Chemistry --
\newcommand\CRN{chemical reaction network}
\newcommand\CRNs{chemical reaction networks}
\newcommand\RN{reaction network}
\def\ACR{absolute concentration robustness}
\def\AMSD{acyclic siphon decomposition}
\def\DAE{differential algebraic equations}
\def\FHJ{Feinberg-Horn-Jackson}
\def\GMAK{generalized mass action kinetics}
\def\MM{Michaelis-Menten}
\def\sm{stoichiometric matrix}\def\Sm{stoichiometric matrix}
\def\WRZD{weakly reversible, zero deficiency}\def\aut{autocatalytic}

% -- Epidemiology --
\def\brn{basic reproduction number}\def\bzn{basic replacement number}
\def\com{compartment}
\def\dfe{disease free equilibrium}\def\DFE{disease free equilibrium}
\def\dwt{dwell times}\def\edt{expected dwell times}
\def\EM{epidemic model}
\def\endp{endemic fixed point}\def\enp{endemic point}\def\EE{endemic point}
\def\ME{mathematical epidemiology}
\def\MEE{eco-epidemiology}\def\EEP{eco-epidemiological}
\def\ece{eco-epidemiological model}
\def\mgm{multi group model}
\def\mtS{multi-strain model}\def\mSt{multi-strain model}
\def\mSi{minimal siphon}
\def\NGM{next generation matrix}
\def\nim{new infections matrix}
\def\qem{quadratic epidemiological model}
\def\repF{invasion functions}
\def\SPF{SIR-PH-FA model}
\def\sus{susceptible}

% -- Graph theory & Linear algebra --
\def\jac{Jacobi matrix}
\def\jin{Jacobian matrix of the invasion vector field with respect to the invasion variables}
\def\jrn{Jacobian matrix of the resident vector field with respect to the resident variables}
\def\jxn{Jacobian matrix of the $\x$ vector field with respect to the $\x$ variables}
\def\Lm{Laplacian matrix}
\def\MTT{matrix tree theorem}
\def\SGT{spectral graph theory}
\def\SMW{Sherman–Morrison–Woodbury}

% -- Stability & Dynamics --
\def\absC{absorbtion CTMC}
\def\BTb{Bogdanov-Takens bifurcation}
\def\CEP{competitive exclusion partition}
\def\CTMC{continuous time Markov chain}
\def\ds{differential system}
\def\fp{fixed point}\def\bfp{boundary fixed point}\def\fbp{fixed boundary point}
\def\fbl{fixed boundary line}\def\fif{forward invariant face}
\def\gas{global asymptotic stability}\def\GAS{globally asymptotically stable}
\def\GASS{\text{GASS}}
\def\Hb{Hopf bifurcation}
\def\Hur{Hurwitz}\def\MH{Metzler-Hurwitz}\def\RH{Routh-Hurwitz}
\def\IS{\text{IS}}
\def\LAS{locally asymptotically stable}
\def\LdS{Lyapunov diagonally stable}
\def\Lf{Lyapunov function}\def\LyV{Lyapunov Volterra type function}\def\RLF{robust Lyapunov function}
\def\LyM{Lyapunov Malkin stable}
\def\PF{Perron-Frobenius}
\def\spfp{strictly positive fixed point}
\def\sta{stability}
\def\STP{sharp threshold property}\def\STT{sharp threshold theorem}

% -- Models (with \varphi explicitly instead of \f) --
\def\AB{$(A,B)$ Arino-Brauer epidemic models}
\def\ABa{$(A,\bbe,\ba,\varphi)$ balanced bilinear model}
\def\ABP{%$(A,\bbe,P,\varphi)$
\\ balanced bilinear model}
\def\AMG{Anderson May Gupta}
\newcommand\BIN{biological interaction network}
\def\Dt{Descartes type}
\def\gLV{generalized Lotka-Volterra model}
\def\IG{\text{invasion graph}}\def\IGMS{interaction graph on minimal siphons}
\def\LCP{\text{linear complementarity problem}}
\def\LV{Lotka-Volterra}\def\LVK{Lotka-Volterra-Kolmogorov}

% -- Splitting & Structure --
\def\reg{regular splitting}\def\regS{{\bf regular splitting}}
\def\RUR{rational univariate representation}
\def\bds{block-diagonal structure}\def\NGMs{next generation matrices}

% -- Miscellaneous --
\def\Ma{{Mathematica}}
\def\Mm{Max–Min}
\def\num{numerator}
\def\rap{rational approximation} \def\DSR{directed species-reaction bipartite graph}\def\PCP{polynomial complementarity problem}\def\rat{rational substitution}
\def\QSS{quasi steady state} \def\RHS{right hand side} \def\LHS{left hand side}
\def\OSN{online social network}\def\VdD{Van den Driessche}
\def\wli{$\omega$-limit} \def\rk{rich kinetics}\def\sJ{symbolic Jacobian}
\def\cV{\cite{Vas,VasSta,BSV,GSV}} \def\spf{semi-parametric semi-functional}
\def\CRW{Capasso-Ruan-Wang} \def\ABP{$(A,\bbe,P,\f)$ balanced bilinear model}
\def\ABa{%$(A,\bbe,\ba,\f)$ \\
rank-one balanced bilinear model}
\def\bbe{B%{\bff \beta}
} \def\vbe{\vec {\bff \beta}
}
\def\vb{\vec {\bff b}}
\newcommand{\Kh}{\hat{K}}          % K with hat, first mention only
\newcommand{\wS}{w^{\sigma}}       % right Perron J^\perp, full block
\newcommand{\piS}{\vec \pi^{\sigma}}    % left Perron J^\perp, full block
\newcommand{\ws}{w^{\sigma}}       % right Perron J^\perp, sub-block
\newcommand{\pis}{\pi^{\sigma}}    % left Perron J^\perp, sub-block
\newcommand{\hwS}{\hat{w}^{\sigma}}   % right Perron K, full block
\newcommand{\hpiS}{\hat{\pi}^{\sigma}} % left Perron K, full block
\newcommand{\hws}{\hat{w}^{\sigma}}   % right Perron K, sub-block
\newcommand{\hpis}{\hat{\pi}^{\sigma}} % left Perron K, sub-block
\newcommand{\twS}{\tilde{w}^{\sigma}}  % right Perron \tilde K, full block
\newcommand{\tpiS}{\tilde{\pi}^{\sigma}} % left Perron \tilde K
\newcommand{\tws}{\tilde{w}^{\sigma}}
\newcommand{\tpis}{\tilde{\pi}^{\sigma}}
\newcommand{\Mat}[2]{\mathbb{R}_{\ge 0}^{\,#1\times #2}}
\def\tJ{J_\sigma^\perp} \def\lPe{left Perron eigenvector}
% \G is \Gamma (defined at :795).  The Volterra function keeps its own name,
% so that \G no longer has two conflicting meanings in this file.
\def\GV{G}  % Volterra G(u)=u-1-\ln u
\def\PVL{Perron-Volterra Lyapunov minimal siphon ansatz}
\newcommand{\istar}[1]{i_{#1}^{*}}
%\newcommand{\sE}{s_E}
\newcommand{\vE}{v_E}
\def\PV{Perron-Volterra candidate Lyapunov function}
\def\GLV{generalized Lotka-Volterra}
\def\MERN{mathematical epidemiology reaction networks}
\def\CMM{contraction metric method}
\newcommand{\sign}{\operatorname{sign}}
\newcommand{\Var}{\operatorname{Var}} 

% ======================================================================
% IDEAS SHARED BETWEEN THE ABSTRACT AND THE MAIN RESULTS SECTION
%
% Write each idea ONCE, in the abstract; the Main Results section recalls it,
% so the two can never drift apart.
%
%   in the abstract        \ideaitem{1}{The first-order reaction terms ...}
%   in Main Results        \ideahead{1}
%
% \ideaitem stores the text globally and typesets
%   \item \hyperlink{idea:1}{<text>}
% \ideahead typesets
%   \item\hypertarget{idea:1}{}\textbf{<the same text>}
%
% The abstract must come before the Main Results input, which is the normal
% order. If a heading has to read differently from the abstract sentence (for
% instance a long abstract item summarised in one line), override it with the
% optional argument: \ideahead[short heading]{6}.
% ======================================================================
\makeatletter
\newcommand{\ideaitem}[2]{%
  \expandafter\gdef\csname @idea@#1\endcsname{#2}%
  \item\hyperlink{idea:#1}{#2}}
\newcommand{\ideahead}[2][]{%
  \item\hypertarget{idea:#2}{}%
  \def\@ideatmp{#1}%
  \ifx\@ideatmp\@empty
    \@ifundefined{@idea@#2}{%
      \PackageWarning{preamble}{Idea #2 is used in Main Results but was never
        defined by \string\ideaitem\space in the abstract}%
      \textbf{[idea #2 undefined]}%
    }{\textbf{\begingroup\def\ideaterm##1##2{##2}\def\idearest##1{}\@nameuse{@idea@#2}\endgroup}}%
  \else
    \textbf{#1}%
  \fi}
% \ideaterm{<hypertarget>}{<words>}: inside an abstract idea the words become an
% invisible link (the document sets \hypersetup{hidelinks}) jumping to the place
% in Main Results that defines them; inside the Main Results heading itself,
% where the target sits a few lines below, \ideahead strips the link and prints
% the words plainly. Mark the destination with \hypertarget{<name>}{} on the line
% before the display, which adds no equation number and shows nothing.
\newcommand{\ideaterm}[2]{\hyperlink{#1}{#2}}
% \idearest{<text>}: the part of an abstract idea after its first sentence; printed in
% the abstract, omitted from the Main Results heading.
\newcommand{\idearest}[1]{#1}
\makeatother

% ======================================================================
% SHORT-PAPER SWITCH
% \shorttrue in a master that assembles only part of the companion paper
% (Barrat/SipArxiv.tex). \nsh{...} drops text that cites results living in the
% omitted sections, so the short build has no dangling references while the
% companion paper keeps the text verbatim.
% ======================================================================
\newif\ifshort \shortfalse
\newcommand{\nsh}[1]{\ifshort\else#1\fi}
% \shortalt{<short build>}{<companion paper>}: use when a sentence needs
% different wording in the two builds, not merely a deletion.
\newcommand{\shortalt}[2]{\ifshort#1\else#2\fi}
% \Jtr, \JtrS{\Sigma}: transverse Jacobian. J^\perp here; SipArxiv.tex redefines them as J_{\mathrm{tra}}, J_{\mathrm{tra},\Sigma}
\newcommand{\Jtr}{J^\perp}
\newcommand{\JtrS}[1]{J_{#1}^\perp}
, and must not clash.
\providecommand{\Q}{\mathbb{Q}}
\providecommand{\res}{\operatorname{res}}   % resultant; used by epiRes.tex
\newcommand{\MN}{\mathbb{N}}
\newcommand{\MZ}{\mathbb{Z}}
\newcommand{\realpos}{\mathbb{R}_{> 0}}
\newcommand{\realnneg}{\mathbb{R}_{\geq 0}}
\newcommand{\naturalpos}{\mathbb{N}_{>0}}
\newcommand{\Rplus}{\mathbb{R}_+}
\newcommand{\Rp}{\mathbb{R}_{\ge 0}}
\def\Rnne{\mathbb{R}_{\ge 0}}
\def\Se{\bar{S}} \def\Ie{\bar{I}}
\newcommand{\Rpos}{\mathbb{R}_{>0}}
\newcommand{\dd}{\,\mathrm d}
\newcommand{\e}{\mathrm e}
%% ========================================
%% GREEK LETTERS (2+ character abbreviations) greek
%% ========================================
\def\al{\alpha}\def\ba{\bff \alpha}
\def\w{\omega}\def\Lgn{\ga+\mu+\de}\def\om{\omega}
\def\de{\delta}\def\De{\Delta}\def\bde{\bff \delta}
\def\ep{\epsilon}\def\eps{\varepsilon}
\def\f{\varphi}  % Kept as requested
\def\ga{\gamma}\def\Ga{\Gamma}\def\G{\Gamma}\def\bg{\bff \gamma}
\def\up{\upsilon}  % NOT the breve accent: no accent commands anywhere.
                   % Named \up, not \u, so the LaTeX breve stays reachable.
\def\la{\lambda}\def\La{\Lambda}
\def\si{\sigma}\def\ka{\kappa}
\def\th{\theta}\def\vt{\vartheta}
\def\ze{\zeta}\def\Gap{\Gamma^{+}} \def\Gam{\Gamma^{-}}


%% ========================================
%% SPECIAL SHORT COMMANDS (kept as requested)
%% ========================================
\def\I{\infty}
\def\T{\widetilde}
\def\H{\hat}
\newcommand{\relint}{\operatorname{relint}}

%% ========================================
%% BOLD/VECTOR NOTATION
%% ========================================
\newcommand{\bff}[1]{{\mbox{\boldmath$#1$}}}
\newcommand\X{\boldsymbol{X}}
\def\x{\boldsymbol{x}}\def\bx{\bff x}\def\xe{\bar{\x}}
\newcommand\y{\boldsymbol{y}}\def\by{\boldsymbol{y}}\def\ye{\bar{\y}}
\def\yb{\y}\def\xb{\x}
\def\xy{\x-\y}
\newcommand\br{\boldsymbol{r}}\def\re{\bar{\r}}
\newcommand\bd{\boldsymbol{d}}
\newcommand\bQ{\boldsymbol{Q}}
\newcommand{\Yb}{{\mathbf{Y}}}
\newcommand{\Ab}{\mathbf{A}}
\newcommand{\zb}{\mathbf{z}}\def\bz{\bff 0}
\newcommand{\Zb}{\mathbf{Z}}
\newcommand{\bb}{\mathbf{b}}
\def\vbe{\vec {\bff \beta}}
\def\vb{\vec {\bff b}}
\def\v1{\vec {\bff 1}}
\def\vn{\vec \de}
\def\vz{\vec 0}
\def\vi{\overset{\rightarrow}{\mathsf i}}
\def\vv{\overset{\rightarrow}{v}}
\def\vr{\; \vec {\mathsf r}}
\newcommand{\va}{\bff a}
\providecommand{\vect}[1]{\boldsymbol{#1}}   % bold vector (moved from Gradostat.tex)
\providecommand{\one}{\mathbf 1}             % all-ones vector (moved from AliSF)
\providecommand{\spr}{s}                     % spectral bound  (moved from AliSF)
\providecommand{\Ke}{K_{\rm eff}}            % effective carrying capacity
% Critical right/left eigenvectors in the Hopf projection method (AAHOSN/New.tex)
\providecommand{\bq}{\mathbf{q}}
\providecommand{\bp}{\mathbf{p}}


%% ========================================
%% MATHCAL & MATHSCR
%% ========================================
% Mathcal
\newcommand\mc[1]{\mathcal{#1}}
\def\mC{{\mathcal C}}\def\mD{\mathcal D}\def\mE{\mathcal{E}}
\def\mF{\mathcal F}\def\mG{{\mathcal G}}\def\mH{{\mathcal H}}
\def\mI{\mathcal I}\def\mL{\mathcal{L}}\def\mM{{\mathcal M}}
\def\mN{{\mathcal N}}\def\mP{\mathcal{P}}\def\mR{{\mathcal R}}
\def\mS{{\mathcal S}}\def\mT{\mathcal{T}}\def\mV{{\mathcal V}}
\def\mK{\mathcal K}\def\mA{{\mathcal A}}\def\mU{{\mathcal U}}
% Symmetric bilinear and trilinear forms of the Taylor expansion, used in the
% first Lyapunov coefficient (AAHOSN/New.tex)
\providecommand{\Bform}{\mathcal{B}}
\providecommand{\Cform}{\mathcal{C}}

% Blackboard bold.  Kept here, not in the drivers: a \providecommand in a
% driver preamble runs before this file and silently wins, so the same name
% ends up meaning different things in different builds.
\def\IC{\mathbb{C}}\def\IN{\mathbb{N}}\def\IQ{\mathbb{Q}}
\def\IR{\mathbb{R}}\def\IZ{\mathbb{Z}}\def\IG{\mathbb{G}}
\def\II{\mathbb{I}}

% Mathscr
\newcommand{\Rs}{{\mathscr R}}
\newcommand{\Ss}{{\mathscr S}}
\newcommand{\Ns}{{\mathscr N}}
\newcommand{\Cs}{\mathscr C}


%% ========================================
%% OPERATORS & FUNCTIONS
%% ========================================
\DeclareMathOperator{\tr}{tr}
\DeclareMathOperator{\diag}{Diag}
\DeclareMathOperator{\Diag}{Diag}
\DeclareMathOperator{\intr}{\mathrm{Int}}
\DeclareMathOperator{\Inv}{Inv}
\DeclareMathOperator{\spec}{spec}
\DeclareMathOperator{\AUC}{AUC}               % area under the curve
\newcommand{\ii}{\mathrm{i}}                  % imaginary unit
\newcommand\op{\operatorname}
\newcommand{\supp}{\op{supp}}
\newcommand{\diff}{\mathop{}\mathopen{}\mathrm{d}}
\newcommand{\pdd}[1]{\frac{\partial}{\partial{#1}}}
\DeclareMathOperator{\cone}{cone}
\DeclareMathOperator{\proj}{proj}

%% ========================================
%% BRACKETS & NORMS
%% ========================================
\providecommand{\pr}[1]{\left(#1\right)}
\providecommand{\pp}[1]{\left[#1\right]}
\newcommand\croc[1]{\left\langle #1\right\rangle}
\newcommand{\norm}[1]{\|#1\|}
\newcommand{\inner}[2]{\langle #1, #2 \rangle}
% Short synonym of \inner, used by AAHOSN/New.tex as $\ip{u}{v}=\bar u^\top v$
\providecommand{\ip}[2]{\langle #1 , #2 \rangle}


%% ========================================
%% FORMATTING & UTILITIES
%% ========================================
\newcommand*{\Scale}[2][4]{\scalebox{#1}{$#2$}}
\newcommand{\te}[1]{\text{#1}}
\newcommand\steq[1]{\stackrel{\text{\rm #1.}}{=}}
\newcommand\blue[1]{\textcolor{blue}{#1}}
\newcommand{\red}{\textcolor[rgb]{1.00,0.00,0.00}}
\newcommand{\green}{\textcolor[rgb]{0.00,1.00,0.00}}

% Figure command
\newcommand{\figu}[3]{%
\begin{figure}[H]
\centering
\includegraphics[scale=#3]{#1}
\caption{#2\label{f:#1}}
\end{figure}
}

% Footnote
\long\def\symbolfootnote[#1]#2{%
\begingroup
\def\thefootnote{\fnsymbol{footnote}}\footnote[#1]{#2}%
\endgroup}
\def\fn{\symbolfootnote}

\newcommand{\dsum}{\displaystyle\sum}

%% ========================================
%% CRN & EPIDEMIOLOGY SPECIFIC
%% ========================================
% CRN notation
\def\SS{\mathcal S}
\def\CC{\mathcal C}
\def\RR{\mathcal R}

% Reproduction numbers
\def\Rto{R_{1}^{E_2}}
\def\Rot{R_{2}^{E_1}}
\def\Roc{R_{1,c}}
\def\Rtc{R_{2,c}}
\def\rd{r_{0}}
\def\sd{s_{0}}
\def\dd{D_{0}}
\def\vd{v_{0}}
\def\zd{z_{0}}
\def\se{\sE}
\def\xd{x_{0}}
%% ========================================
%% SPECIAL VALUES & EQUILIBRIA
%% ========================================
\def\eE{\bar{e}}
\def\wE{\bar{w}}
\def\sE{s_{*}}
\def\Me{\T M_{ee}}
\def\Rnne{\mathbb{R}_{\ge 0}}

\newcommand{\ones}{\mathbf 1}


%% ========================================
%% ABBREVIATIONS (alphabetical by category)
%% ========================================

% -- General mathematical terms --
\def\adm{admissible}\def\admy{admissibility}
%\def\and{antisymmetric}% disabled 2026-09-30: it overrode LaTeX's author separator \and (printed "antisymmetric" between authors)
\def\ass{assumption}
\def\cadlag{c\`adl\`ag}
\def\ch{characteristic polynomial}
\def\chp{characteristic polynomial}
\def\coe{coefficient}
\def\comp{computation}
\def\con{condition}
\def\coo{coordinate}
\def\corr{corresponding}
\def\dec{decomposition}
\def\den{denominator}
\def\dete{deterministic}
\def\eig{eigenvalue}\def\eval{eigenvalue}
\def\eigv{eigenvector}\def\eiv{eigenvector}\def\evec{eigenvector}
\def\elt{element}
\def\enn{essentially nonnegative/positivity preserving}
\def\equ{equation}
\def\etal{et al.}
\def\fe{for example}
\def\fno{from now on}
\def\for{\forall}
\def\form{formula}
\def\Fr{Furthermore, }
\def\frt{furthermore}
\def\fun{function}
\def\gene{generalization}
\def\heu{heuristic}
\def\how{however}\def\How{However, }
\def\ie{i.e. }
\def\im{\item}
\def\isd{is defined by}
\def\itc{it is easy to check that}
\def\itf{it follows that}\def\Itf{It follows that}
\def\ith{it holds that}
\def\itm{it may be checked that}\def\Itm{It may be checked that}
\def\ito{it turns out that}\def\Ito{It turns out that}\def\Itou{It turns out useful to}
\def\ma{mass action}\def\MAR{mass action representation}
\def\mat{matrix}
\def\mbc{may be checked to be}
\def\mbw{may be written as}\def\mubw{must be written as}
\def\mpr{more precisely, }\def\Mp{More precisely, }
\def\nc{non-catalytic}
\def\nne{non-negative}\def\nny{nonnegativity}
\def\no{\nonumber}
\def\noni{non-infectious}
\def\nota{notation}
\def\np{\newpage}
\def\Nt{Note that}
\def\nt{non trivial}\def\nz{non zero}
\def\oth{otherwise}
\def\Oth{On the other hand}
\def\pa{parameter}\def\para{parameter}
\def\pol{polynomial}
\def\prob{probability}
\def\Prf{{\bf  Proof:}}\def\prf{{\bf  Proof:}}
\def\prop{proportional to}
\def\qu{\quad}
\def\QED{\hfill {$\square$}\goodbreak \medskip}
\def\rep{representation}
\def\resp{respectively}
\def\satd{satisfied}\def\satg{satisfying}\def\sats{satisfies}\def\saty{satisfy}
\def\ssec{\subsection}\def\sssec{\subsubsection}
\def\sef{see for example}
\def\sksy{skew-symmetric}
\def\sub{substrate}
\def\sys{system}
\def\var{variable}
\def\wk{well-known}
\def\wl{without loss of generality}\def\wlo{w.l.o.g.}
\def\wmc{We may check that}\def\Wmc{We may check that}
\def\wmr{we may rewrite}
\def\wms{we must show that}
\def\wrt{with respect to}

% -- CRN & Chemistry --
\newcommand\CRN{chemical reaction network}
\newcommand\CRNs{chemical reaction networks}
\newcommand\RN{reaction network}
\def\ACR{absolute concentration robustness}
\def\AMSD{acyclic siphon decomposition}
\def\DAE{differential algebraic equations}
\def\FHJ{Feinberg-Horn-Jackson}
\def\GMAK{generalized mass action kinetics}
\def\MM{Michaelis-Menten}
\def\sm{stoichiometric matrix}\def\Sm{stoichiometric matrix}
\def\WRZD{weakly reversible, zero deficiency}\def\aut{autocatalytic}

% -- Epidemiology --
\def\brn{basic reproduction number}\def\bzn{basic replacement number}
\def\com{compartment}
\def\dfe{disease free equilibrium}\def\DFE{disease free equilibrium}
\def\dwt{dwell times}\def\edt{expected dwell times}
\def\EM{epidemic model}
\def\endp{endemic fixed point}\def\enp{endemic point}\def\EE{endemic point}
\def\ME{mathematical epidemiology}
\def\MEE{eco-epidemiology}\def\EEP{eco-epidemiological}
\def\ece{eco-epidemiological model}
\def\mgm{multi group model}
\def\mtS{multi-strain model}\def\mSt{multi-strain model}
\def\mSi{minimal siphon}
\def\NGM{next generation matrix}
\def\nim{new infections matrix}
\def\qem{quadratic epidemiological model}
\def\repF{invasion functions}
\def\SPF{SIR-PH-FA model}
\def\sus{susceptible}

% -- Graph theory & Linear algebra --
\def\jac{Jacobi matrix}
\def\jin{Jacobian matrix of the invasion vector field with respect to the invasion variables}
\def\jrn{Jacobian matrix of the resident vector field with respect to the resident variables}
\def\jxn{Jacobian matrix of the $\x$ vector field with respect to the $\x$ variables}
\def\Lm{Laplacian matrix}
\def\MTT{matrix tree theorem}
\def\SGT{spectral graph theory}
\def\SMW{Sherman–Morrison–Woodbury}

% -- Stability & Dynamics --
\def\absC{absorbtion CTMC}
\def\BTb{Bogdanov-Takens bifurcation}
\def\CEP{competitive exclusion partition}
\def\CTMC{continuous time Markov chain}
\def\ds{differential system}
\def\fp{fixed point}\def\bfp{boundary fixed point}\def\fbp{fixed boundary point}
\def\fbl{fixed boundary line}\def\fif{forward invariant face}
\def\gas{global asymptotic stability}\def\GAS{globally asymptotically stable}
\def\GASS{\text{GASS}}
\def\Hb{Hopf bifurcation}
\def\Hur{Hurwitz}\def\MH{Metzler-Hurwitz}\def\RH{Routh-Hurwitz}
\def\IS{\text{IS}}
\def\LAS{locally asymptotically stable}
\def\LdS{Lyapunov diagonally stable}
\def\Lf{Lyapunov function}\def\LyV{Lyapunov Volterra type function}\def\RLF{robust Lyapunov function}
\def\LyM{Lyapunov Malkin stable}
\def\PF{Perron-Frobenius}
\def\spfp{strictly positive fixed point}
\def\sta{stability}
\def\STP{sharp threshold property}\def\STT{sharp threshold theorem}

% -- Models (with \varphi explicitly instead of \f) --
\def\AB{$(A,B)$ Arino-Brauer epidemic models}
\def\ABa{$(A,\bbe,\ba,\varphi)$ balanced bilinear model}
\def\ABP{%$(A,\bbe,P,\varphi)$
\\ balanced bilinear model}
\def\AMG{Anderson May Gupta}
\newcommand\BIN{biological interaction network}
\def\Dt{Descartes type}
\def\gLV{generalized Lotka-Volterra model}
\def\IG{\text{invasion graph}}\def\IGMS{interaction graph on minimal siphons}
\def\LCP{\text{linear complementarity problem}}
\def\LV{Lotka-Volterra}\def\LVK{Lotka-Volterra-Kolmogorov}

% -- Splitting & Structure --
\def\reg{regular splitting}\def\regS{{\bf regular splitting}}
\def\RUR{rational univariate representation}
\def\bds{block-diagonal structure}\def\NGMs{next generation matrices}

% -- Miscellaneous --
\def\Ma{{Mathematica}}
\def\Mm{Max–Min}
\def\num{numerator}
\def\rap{rational approximation} \def\DSR{directed species-reaction bipartite graph}\def\PCP{polynomial complementarity problem}\def\rat{rational substitution}
\def\QSS{quasi steady state} \def\RHS{right hand side} \def\LHS{left hand side}
\def\OSN{online social network}\def\VdD{Van den Driessche}
\def\wli{$\omega$-limit} \def\rk{rich kinetics}\def\sJ{symbolic Jacobian}
\def\cV{\cite{Vas,VasSta,BSV,GSV}} \def\spf{semi-parametric semi-functional}
\def\CRW{Capasso-Ruan-Wang} \def\ABP{$(A,\bbe,P,\f)$ balanced bilinear model}
\def\ABa{%$(A,\bbe,\ba,\f)$ \\
rank-one balanced bilinear model}
\def\bbe{B%{\bff \beta}
} \def\vbe{\vec {\bff \beta}
}
\def\vb{\vec {\bff b}}
\newcommand{\Kh}{\hat{K}}          % K with hat, first mention only
\newcommand{\wS}{w^{\sigma}}       % right Perron J^\perp, full block
\newcommand{\piS}{\vec \pi^{\sigma}}    % left Perron J^\perp, full block
\newcommand{\ws}{w^{\sigma}}       % right Perron J^\perp, sub-block
\newcommand{\pis}{\pi^{\sigma}}    % left Perron J^\perp, sub-block
\newcommand{\hwS}{\hat{w}^{\sigma}}   % right Perron K, full block
\newcommand{\hpiS}{\hat{\pi}^{\sigma}} % left Perron K, full block
\newcommand{\hws}{\hat{w}^{\sigma}}   % right Perron K, sub-block
\newcommand{\hpis}{\hat{\pi}^{\sigma}} % left Perron K, sub-block
\newcommand{\twS}{\tilde{w}^{\sigma}}  % right Perron \tilde K, full block
\newcommand{\tpiS}{\tilde{\pi}^{\sigma}} % left Perron \tilde K
\newcommand{\tws}{\tilde{w}^{\sigma}}
\newcommand{\tpis}{\tilde{\pi}^{\sigma}}
\newcommand{\Mat}[2]{\mathbb{R}_{\ge 0}^{\,#1\times #2}}
\def\tJ{J_\sigma^\perp} \def\lPe{left Perron eigenvector}
% \G is \Gamma (defined at :795).  The Volterra function keeps its own name,
% so that \G no longer has two conflicting meanings in this file.
\def\GV{G}  % Volterra G(u)=u-1-\ln u
\def\PVL{Perron-Volterra Lyapunov minimal siphon ansatz}
\newcommand{\istar}[1]{i_{#1}^{*}}
%\newcommand{\sE}{s_E}
\newcommand{\vE}{v_E}
\def\PV{Perron-Volterra candidate Lyapunov function}
\def\GLV{generalized Lotka-Volterra}
\def\MERN{mathematical epidemiology reaction networks}
\def\CMM{contraction metric method}
\newcommand{\sign}{\operatorname{sign}}
\newcommand{\Var}{\operatorname{Var}} 

% ======================================================================
% IDEAS SHARED BETWEEN THE ABSTRACT AND THE MAIN RESULTS SECTION
%
% Write each idea ONCE, in the abstract; the Main Results section recalls it,
% so the two can never drift apart.
%
%   in the abstract        \ideaitem{1}{The first-order reaction terms ...}
%   in Main Results        \ideahead{1}
%
% \ideaitem stores the text globally and typesets
%   \item \hyperlink{idea:1}{<text>}
% \ideahead typesets
%   \item\hypertarget{idea:1}{}\textbf{<the same text>}
%
% The abstract must come before the Main Results input, which is the normal
% order. If a heading has to read differently from the abstract sentence (for
% instance a long abstract item summarised in one line), override it with the
% optional argument: \ideahead[short heading]{6}.
% ======================================================================
\makeatletter
\newcommand{\ideaitem}[2]{%
  \expandafter\gdef\csname @idea@#1\endcsname{#2}%
  \item\hyperlink{idea:#1}{#2}}
\newcommand{\ideahead}[2][]{%
  \item\hypertarget{idea:#2}{}%
  \def\@ideatmp{#1}%
  \ifx\@ideatmp\@empty
    \@ifundefined{@idea@#2}{%
      \PackageWarning{preamble}{Idea #2 is used in Main Results but was never
        defined by \string\ideaitem\space in the abstract}%
      \textbf{[idea #2 undefined]}%
    }{\textbf{\begingroup\def\ideaterm##1##2{##2}\def\idearest##1{}\@nameuse{@idea@#2}\endgroup}}%
  \else
    \textbf{#1}%
  \fi}
% \ideaterm{<hypertarget>}{<words>}: inside an abstract idea the words become an
% invisible link (the document sets \hypersetup{hidelinks}) jumping to the place
% in Main Results that defines them; inside the Main Results heading itself,
% where the target sits a few lines below, \ideahead strips the link and prints
% the words plainly. Mark the destination with \hypertarget{<name>}{} on the line
% before the display, which adds no equation number and shows nothing.
\newcommand{\ideaterm}[2]{\hyperlink{#1}{#2}}
% \idearest{<text>}: the part of an abstract idea after its first sentence; printed in
% the abstract, omitted from the Main Results heading.
\newcommand{\idearest}[1]{#1}
\makeatother

% ======================================================================
% SHORT-PAPER SWITCH
% \shorttrue in a master that assembles only part of the companion paper
% (Barrat/SipArxiv.tex). \nsh{...} drops text that cites results living in the
% omitted sections, so the short build has no dangling references while the
% companion paper keeps the text verbatim.
% ======================================================================
\newif\ifshort \shortfalse
\newcommand{\nsh}[1]{\ifshort\else#1\fi}
% \shortalt{<short build>}{<companion paper>}: use when a sentence needs
% different wording in the two builds, not merely a deletion.
\newcommand{\shortalt}[2]{\ifshort#1\else#2\fi}
% \Jtr, \JtrS{\Sigma}: transverse Jacobian. J^\perp here; SipArxiv.tex redefines them as J_{\mathrm{tra}}, J_{\mathrm{tra},\Sigma}
\newcommand{\Jtr}{J^\perp}
\newcommand{\JtrS}[1]{J_{#1}^\perp}
.
\makeatletter
\@ifpackageloaded{hyperref}{}{%
  \@ifundefined{hyperrefpackage}{%
    \@ifundefined{AAdefNnohyperref}{\usepackage{hyperref}}{}%
  }{}%
}
\makeatother
% linktoc=all makes a table-of-contents page number a link as well as the
% entry text; hyperref's default, linktoc=section, links only the text.
% Deferred to begin-document and guarded: by the block above hyperref may
% be loaded later (mdpi's own hook) or not at all (\AAdefNnohyperref), and
% \hypersetup must exist before it is called.  A driver that wants the
% default back can say \hypersetup{linktoc=section} after % Shared preamble for the GLVbook papers.  This is the single canonical copy;
% a driver one level below GLVbook reaches it as
%
%     \input{../basicFiles/preamble}
%
% What stays in the driver, because it differs from paper to paper:
% \documentclass, font choice, \title, \author, \date, and any macro genuinely
% local to one paper.  Everything else belongs here.  The file is
% self-contained: it does not read any other file.

\usepackage[utf8]{inputenc}
\usepackage[novbox]{pdfsync}   % required: SyncTeX alone does not work here

% Page geometry, set once here for every paper that loads this file.  A driver
% that needs different margins says \geometry{...}; a second
% \usepackage[...]{geometry} raises an option clash and must not be used.
\makeatletter
\@ifundefined{AAdefNnogeometry}{\usepackage[a4paper,margin=2cm]{geometry}}{}
\makeatother

% mathematics
\usepackage{amsmath,amssymb,amsthm,mathtools,bm,mathrsfs,bbm}
\usepackage{centernot}

% floats, tables, lists
\usepackage{graphicx,float}
\makeatletter
\@ifpackageloaded{subfig}{}{\usepackage{subcaption}}
\makeatother
\usepackage{booktabs,longtable,threeparttable,array,multirow}
\usepackage{enumitem}

% colour, code, algorithms
\usepackage{xcolor,listings,tcolorbox}
\makeatletter
\@ifundefined{c@algorithm}{\usepackage{algorithm}}{}
\makeatother
\usepackage{algorithmicx,algpseudocode}

% diagrams
\usepackage{tikz,tikz-cd,pgfplots}
\usetikzlibrary{graphs,decorations.pathmorphing,decorations.markings,
arrows.meta,automata,positioning}
\tikzcdset{every label/.append style = {font = \small}}
\pgfplotsset{compat=1.18}

% chemistry
\usepackage[version=4]{mhchem}

\usepackage{microtype}
\usepackage{lineno}
\usepackage{nomencl}
\makenomenclature

% hyperref first, then bookmark, then glossaries: glossaries must follow
% hyperref or the entry links are built against the wrong anchors.
\makeatletter
\@ifundefined{hyperrefpackage}{%
  \@ifundefined{AAdefNnohyperref}{\usepackage{hyperref}}{}%
}{}
\@ifundefined{hyperrefpackage}{\usepackage{bookmark}}{}
\makeatother
\usepackage[toc,sort=use]{glossaries}
\makeglossaries

% Book/standalone switch.  Material that the book already covers in its
% prelude, but that a standalone paper still has to state, is wrapped as
%     \ifbookbuild\else ... \fi
% Drivers default to standalone; bookM.tex says \bookbuildtrue.
\newif\ifbookbuild
\bookbuildfalse

% Part switches for the gradostat paper.  Both default to true; a driver that
% wants the short submission says \madecfalse and/or \slowfastfalse, and every
% cross-reference into the dropped part disappears with it.
%   \ifmadec     Part III, the Lagasquie--Madec model and its break-even geometry
%   \ifslowfast  Part IV, everything carrying the small parameter \eps
\newif\ifmadec     \madectrue
\newif\ifslowfast  \slowfasttrue

% Part switch for the Barrat/simplicial-contagion project.  Defaults to
% true (the hypernetwork-synchrony part is included); the journal-paper
% driver says \hnetpartfalse to drop it and every cross-reference into it.
\newif\ifhnetpart  \hnetparttrue

%% ========================================
%% PACKAGES
%% ========================================
\RequirePackage{etoolbox}
\usepackage{makecell}
\usepackage[novbox]{pdfsync}% obsolete; SyncTeX is built into modern engines
\usepackage{amssymb, amsfonts,mathtools,bbm,eurosym,graphicx,latexsym, amsmath,amsthm,times,nccmath,url}
\usepackage{listings,verbatim,cases,tabularx}
% hyperref is deliberately NOT in the line above.  Journal classes often
% load it themselves, with a long option list, from a begindocument hook
% (mdpi.cls does: \DeclareHookRule{begindocument}{hyperref}{before}{mdpi}).
% Loading it here first, without options, makes that later load raise
% "Option clash for package hyperref".  So load it only if nobody else
% has, or will:
%   - \@ifpackageloaded  catches a class that loaded it already;
%   - \hyperrefpackage    is mdpi.cls's own deferred loader;
%   - \AAdefNnohyperref   lets any other driver opt out by defining it
%                         before % Shared preamble for the GLVbook papers.  This is the single canonical copy;
% a driver one level below GLVbook reaches it as
%
%     \input{../basicFiles/preamble}
%
% What stays in the driver, because it differs from paper to paper:
% \documentclass, font choice, \title, \author, \date, and any macro genuinely
% local to one paper.  Everything else belongs here.  The file is
% self-contained: it does not read any other file.

\usepackage[utf8]{inputenc}
\usepackage[novbox]{pdfsync}   % required: SyncTeX alone does not work here

% Page geometry, set once here for every paper that loads this file.  A driver
% that needs different margins says \geometry{...}; a second
% \usepackage[...]{geometry} raises an option clash and must not be used.
\makeatletter
\@ifundefined{AAdefNnogeometry}{\usepackage[a4paper,margin=2cm]{geometry}}{}
\makeatother

% mathematics
\usepackage{amsmath,amssymb,amsthm,mathtools,bm,mathrsfs,bbm}
\usepackage{centernot}

% floats, tables, lists
\usepackage{graphicx,float}
\makeatletter
\@ifpackageloaded{subfig}{}{\usepackage{subcaption}}
\makeatother
\usepackage{booktabs,longtable,threeparttable,array,multirow}
\usepackage{enumitem}

% colour, code, algorithms
\usepackage{xcolor,listings,tcolorbox}
\makeatletter
\@ifundefined{c@algorithm}{\usepackage{algorithm}}{}
\makeatother
\usepackage{algorithmicx,algpseudocode}

% diagrams
\usepackage{tikz,tikz-cd,pgfplots}
\usetikzlibrary{graphs,decorations.pathmorphing,decorations.markings,
arrows.meta,automata,positioning}
\tikzcdset{every label/.append style = {font = \small}}
\pgfplotsset{compat=1.18}

% chemistry
\usepackage[version=4]{mhchem}

\usepackage{microtype}
\usepackage{lineno}
\usepackage{nomencl}
\makenomenclature

% hyperref first, then bookmark, then glossaries: glossaries must follow
% hyperref or the entry links are built against the wrong anchors.
\makeatletter
\@ifundefined{hyperrefpackage}{%
  \@ifundefined{AAdefNnohyperref}{\usepackage{hyperref}}{}%
}{}
\@ifundefined{hyperrefpackage}{\usepackage{bookmark}}{}
\makeatother
\usepackage[toc,sort=use]{glossaries}
\makeglossaries

% Book/standalone switch.  Material that the book already covers in its
% prelude, but that a standalone paper still has to state, is wrapped as
%     \ifbookbuild\else ... \fi
% Drivers default to standalone; bookM.tex says \bookbuildtrue.
\newif\ifbookbuild
\bookbuildfalse

% Part switches for the gradostat paper.  Both default to true; a driver that
% wants the short submission says \madecfalse and/or \slowfastfalse, and every
% cross-reference into the dropped part disappears with it.
%   \ifmadec     Part III, the Lagasquie--Madec model and its break-even geometry
%   \ifslowfast  Part IV, everything carrying the small parameter \eps
\newif\ifmadec     \madectrue
\newif\ifslowfast  \slowfasttrue

% Part switch for the Barrat/simplicial-contagion project.  Defaults to
% true (the hypernetwork-synchrony part is included); the journal-paper
% driver says \hnetpartfalse to drop it and every cross-reference into it.
\newif\ifhnetpart  \hnetparttrue

%% ========================================
%% PACKAGES
%% ========================================
\RequirePackage{etoolbox}
\usepackage{makecell}
\usepackage[novbox]{pdfsync}% obsolete; SyncTeX is built into modern engines
\usepackage{amssymb, amsfonts,mathtools,bbm,eurosym,graphicx,latexsym, amsmath,amsthm,times,nccmath,url}
\usepackage{listings,verbatim,cases,tabularx}
% hyperref is deliberately NOT in the line above.  Journal classes often
% load it themselves, with a long option list, from a begindocument hook
% (mdpi.cls does: \DeclareHookRule{begindocument}{hyperref}{before}{mdpi}).
% Loading it here first, without options, makes that later load raise
% "Option clash for package hyperref".  So load it only if nobody else
% has, or will:
%   - \@ifpackageloaded  catches a class that loaded it already;
%   - \hyperrefpackage    is mdpi.cls's own deferred loader;
%   - \AAdefNnohyperref   lets any other driver opt out by defining it
%                         before \input{preamble}.
\makeatletter
\@ifpackageloaded{hyperref}{}{%
  \@ifundefined{hyperrefpackage}{%
    \@ifundefined{AAdefNnohyperref}{\usepackage{hyperref}}{}%
  }{}%
}
\makeatother
% linktoc=all makes a table-of-contents page number a link as well as the
% entry text; hyperref's default, linktoc=section, links only the text.
% Deferred to begin-document and guarded: by the block above hyperref may
% be loaded later (mdpi's own hook) or not at all (\AAdefNnohyperref), and
% \hypersetup must exist before it is called.  A driver that wants the
% default back can say \hypersetup{linktoc=section} after \input{preamble}.
\makeatletter
\AtBeginDocument{\@ifundefined{hypersetup}{}{\hypersetup{linktoc=all}}}
\makeatother
\usepackage{graphicx,float,fancybox}
\usepackage{epsfig,epstopdf,tikz}
\usepackage[sans]{dsfont}
\usetikzlibrary{positioning,arrows,arrows.meta,calc}
% Configure listings for Mathematica code
\lstdefinestyle{mathematica}{
 language=Mathematica,
 basicstyle=\small\ttfamily,
 backgroundcolor=\color{gray!10},
 frame=single,
 breaklines=true,
 commentstyle=\color{green!60!black},
 keywordstyle=\color{blue},
 stringstyle=\color{red},
 showstringspaces=false,
 tabsize=2
}
\usepackage{tcolorbox}
% Colored boxes for key results
\newtcolorbox{keyresult}[1][]{
 colback=blue!5!white,
 colframe=blue!75!black,
 title=#1,
 fonttitle=\bfseries
}
\newtcolorbox{examplebox}[1][]{
 colback=green!5!white,
 colframe=green!75!black,
 title=#1,
 fonttitle=\bfseries
}


\makeatletter
\@ifundefined{theorem}   {\newtheorem{theorem}   {Theorem}}   {}
\@ifundefined{definition}{\newtheorem{definition}{Definition}} {}
\@ifundefined{open}      {\newtheorem{open}      {Problem}}   {}
\@ifundefined{example}   {\newtheorem{example}   {Example}}   {}
\makeatother

% COUNTER DEFINITIONS
\newcounter{tabcounter}
\newcounter{defcounter}
\newcounter{opcounter}
\newcounter{thmcounter}
\newcounter{excounter}
\newcounter{ascounter}
\newcounter{cjcounter}
\newcounter{lemcounter}
\newcounter{remcounter}
\newcounter{corcounter}
\newcounter{propcounter}
\newcounter{figcounter}

% UNIFIED GLOSSARY STYLE
\newglossarystyle{unifiedstyle}{%
  \setglossarystyle{list}%
  \renewcommand*{\glossarypreamble}{\gdef\glslastsec{}}%
  \renewcommand{\glossentry}[2]{%
    \edef\glscursec{\glsentryuseriii{##1}}%
    \ifx\glscursec\glslastsec\else
      \item[]\textbf{\S\,\glscursec\enspace\glsentryuseriv{##1}}%
      \global\let\glslastsec\glscursec
    \fi
    \item[\glsentryuserii{##1} \glsentryuseri{##1}:] ``\glossentryname{##1}'' \glsentrydesc{##1}%
  }%
}

\makeatletter
% First-word extractor (used by \beT, \beD and \beXa for the default key)
% \@firstwordof hello world \relax  ->  hello
\def\@firstwordof#1 #2\relax{#1}
% First three letters, used for the heading mnemonic of a theorem given no
% explicit key.  Padding with \@empty keeps it safe for first words shorter
% than three letters: \@empty vanishes under \edef, whereas \relax would
% survive into the label and is illegal inside \csname (it broke every
% two-letter title, e.g. "ME ...", with "Missing \endcsname inserted").
\def\@takethree#1#2#3#4\@nil{#1#2#3}

% EXAMPLE COMMAND -- same signature and behaviour as \beT and \beD.
% Usage: \beXa*[long title][short title][key]
%   star  = suppresses the printed "[key]" tag after the title (as for \beT*
%   and \beD*); the label and glossary entry are unaffected.
%   key = optional label key override; the label becomes e:KEY, and without a
%   key it is e: plus the first three letters of the title's first word.  A
%   mnemonic label lets an example be cited as "Example~\ref{e:KEY}" instead
%   of by section number, which is what the citing prose needs.
\NewDocumentCommand{\beXa}{s O{} O{} o}{%
  \def\@temparg{#2}%
  \ifx\@temparg\@empty
    \begin{example}%
  \else
    \let\@savedSecTitle\@currentsectionname%
    \refstepcounter{excounter}%
    \IfValueTF{#4}{%
      \edef\@exlblkey{#4}%
    }{%
      \edef\@exfirstword{\@firstwordof#2 \relax}%
      \edef\@exlblkey{%
        \expandafter\@takethree\@exfirstword\@empty\@empty\@empty\@nil}%
    }%
    \def\@tempshort{#3}%
    \ifx\@tempshort\@empty
      \expandafter\xdef\csname @xshort@\@exlblkey\endcsname{\@exlblkey}%
    \else
      \expandafter\def\csname @xshort@\@exlblkey\endcsname{#3}%
    \fi
    \IfBooleanTF{#1}{\let\@extag\@empty}{%
      \edef\@extag{\noexpand\ \noexpand\textup{[\noexpand\texttt{e:\@exlblkey}]}}}%
    \begin{example}[#2\ifshowkeys\@extag\fi]%
    \label{e:\@exlblkey}%
    % record the long title so \nref{e:KEY} prints it, as for theorems
    \protected@write\@auxout{}{%
      \string\@thmtitle{e:\@exlblkey}{\unexpanded{#2}}}%
    \expanded{%
      \noexpand\newglossaryentry{e\theexcounter}{%
        name={\unexpanded{#2}%
          \IfBooleanF{#1}{ [\noexpand\texttt{e:\@exlblkey}]}},%
        sort={\theexample},%
        description={p.~\noexpand\pageref{e:\@exlblkey}},%
        user1={\theexample},%
        user2={Example},%
        user3={\thesection},%
        user4={\expandafter\unexpanded\expandafter{\@savedSecTitle}}%
      }%
    }%
    \protected@write\@auxout{}{\string\glsliveentry{e\theexcounter}}%
  \fi
}
\newcommand{\eeXa}{\end{example}}

% DEFINITION COMMAND -- same signature and behaviour as \beT below.
% Usage: \beD[long title][short title][key]
%   long title  = definition heading (plain text only)
%   short title = shown in \dref{word}; defaults to the key
%   key         = optional label key override; the label becomes d:KEY
% Without a key the label is d: plus the first three letters of the title's
% first word, so it is a mnemonic rather than the definition's number: a
% numeric label would silently change whenever a definition is inserted.
\NewDocumentCommand{\beD}{s O{} O{} o}{%
  \def\@temparg{#2}%
  \ifx\@temparg\@empty
    \begin{definition}%
  \else
    \let\@savedSecTitle\@currentsectionname%
    \refstepcounter{defcounter}%
    \IfValueTF{#4}{%
      \edef\@deflblkey{#4}%
    }{%
      \edef\@deffirstword{\@firstwordof#2 \relax}%
      \edef\@deflblkey{%
        \expandafter\@takethree\@deffirstword\@empty\@empty\@empty\@nil}%
    }%
    \def\@tempshort{#3}%
    \ifx\@tempshort\@empty
      \expandafter\xdef\csname @dshort@\@deflblkey\endcsname{\@deflblkey}%
    \else
      \expandafter\def\csname @dshort@\@deflblkey\endcsname{#3}%
    \fi
    \IfBooleanTF{#1}{\let\@deftag\@empty}{%
      \edef\@deftag{\noexpand\ \noexpand\textup{[\noexpand\texttt{d:\@deflblkey}]}}}%
    \begin{definition}[#2\ifshowkeys\@deftag\fi]%
    \label{d:\@deflblkey}%
    % record the long title so \nref{d:KEY} prints it, as for theorems
    \protected@write\@auxout{}{%
      \string\@thmtitle{d:\@deflblkey}{\unexpanded{#2}}}%
    \expanded{%
      \noexpand\newglossaryentry{d\thedefcounter}{%
        name={\unexpanded{#2}%
          \IfBooleanF{#1}{ [\noexpand\texttt{d:\@deflblkey}]}},%
        sort={\thedefinition},%
        description={p.~\noexpand\pageref{d:\@deflblkey}},%
        user1={\thedefinition},%
        user2={Definition},%
        user3={\thesection},%
        user4={\expandafter\unexpanded\expandafter{\@savedSecTitle}}%
      }%
    }%
    \protected@write\@auxout{}{\string\glsliveentry{d\thedefcounter}}%
  \fi
}
\newcommand{\eeD}{\end{definition}}

% \showkeysfalse in a driver hides the [t:KEY] mnemonics in result headings
% (labels and glossary entries are unchanged); used for submission builds.
\newif\ifshowkeys \showkeystrue
% \commonredtrue in a driver that inputs basicFiles/commonRed: tool.tex then
% cites its reduction results instead of restating them.
\newif\ifcommonred \commonredfalse
% TheorEM COMMAND
% Usage: \beT[long title][short title][key]
%   long title  = theorem heading (plain text only)
%   short title = shown in \tref{word}; may contain math; defaults to first word of long title
%   key         = optional label key override; label becomes t:KEY (avoids first-word collisions)
% Label is t:WORD where WORD = first word of long title (or KEY if given); use \tref{WORD} to reference.
\NewDocumentCommand{\beT}{s O{} O{} o}{%
  \def\@temparg{#2}%
  \ifx\@temparg\@empty
    \begin{theorem}%
  \else
    \let\@savedSecTitle\@currentsectionname%
    \refstepcounter{thmcounter}%
    % \@thmtag is the mnemonic shown in the heading; every theorem gets one.
    % With an explicit key the key itself is shown.  Without one the key
    % defaults to the title's first word, and showing that whole word would
    % merely repeat the title ("Dissipativity [Dissipativity]"), so only its
    % first three letters are shown: "Dissipativity [Dis]".  The label is
    % unaffected -- it still uses the full \@thmlblkey.
    \IfValueTF{#4}{%
      \edef\@thmlblkey{#4}%
    }{%
      \edef\@thmfirstword{\@firstwordof#2 \relax}%
      \edef\@thmlblkey{%
        \expandafter\@takethree\@thmfirstword\@empty\@empty\@empty\@nil}%
    }%
    \IfBooleanTF{#1}{\let\@thmtag\@empty}{%
      \edef\@thmtag{\noexpand\ \noexpand\textup{[\noexpand\texttt{t:\@thmlblkey}]}}}%
    \def\@tempshort{#3}%
    \ifx\@tempshort\@empty
      \expandafter\xdef\csname @tshort@\@thmlblkey\endcsname{\@thmlblkey}%
    \else
      \expandafter\def\csname @tshort@\@thmlblkey\endcsname{#3}%
    \fi
    \begin{theorem}[#2\ifshowkeys\@thmtag\fi]%
    \label{t:\@thmlblkey}%
    % record the long title for \nref; via the .aux so forward refs resolve
    \protected@write\@auxout{}{%
      \string\@thmtitle{t:\@thmlblkey}{\unexpanded{#2}}}%
    \expanded{%
      \noexpand\newglossaryentry{t\thethmcounter}{%
        name={\unexpanded{#2}%
          \IfBooleanF{#1}{ [\noexpand\texttt{t:\@thmlblkey}]}},%
        % sort/user1 must be the number the document actually prints, i.e. the
        % theorem environment's counter, not thmcounter.  thmcounter advances
        % only on \beT with a title, so a bare \beT (or \beT with an
        % empty title) makes the two drift and the glossary cites a wrong
        % number.  thmcounter is still fine as the entry key: it is unique.
        sort={\thetheorem},%
        description={p.~\noexpand\pageref{t:\@thmlblkey}},%
        user1={\thetheorem},%
        user2={Theorem},%
        user3={\thesection},%
        user4={\expandafter\unexpanded\expandafter{\@savedSecTitle}}%
      }%
    }%
    \protected@write\@auxout{}{\string\glsliveentry{t\thethmcounter}}%
  \fi
}
\newcommand{\eeT}{\end{theorem}}

% CROSS-PART RELATION PARAGRAPH
% Usage: \xpart{text}  -- place immediately BEFORE \beT, never inside it
%   (anything inside \beT...\eeT would be swallowed into the glossary entry).
% The paragraph is typeset only when Part I is actually present, i.e. in the
% combined book build; in the standalone builds (GranadaDSABNS/A.tex,
% AAHRah-Submission/AAHRah.tex) the label part:glv is undefined and the whole
% paragraph disappears rather than leaving ?? cross-references. Same
% graceful-degradation principle as \GASref in GASt.tex. Resolves on the
% second LaTeX pass, since r@part:glv comes from the .aux file.
% (Already inside the \makeatletter region opened above, so no \makeatletter
% pair here -- adding one would close that region early.)
\newcommand{\xpart}[1]{\@ifundefined{r@part:glv}{}{\paragraph{Relation with Parts~I--II.}#1}}

% \bookref{label}{precise text}{fallback text}: prints the precise cross-reference
% when the label resolves (combined book build), and the fallback text otherwise
% (standalone paper builds, where book-only appendix labels are absent). Same
% graceful-degradation role as \GASref in GASt.tex, but usable in any file since it
% lives in this preamble (loaded by every driver) -- unlike \GASref, which is defined in a
% content file and is therefore undefined in standalone builds that omit that file.
\newcommand{\bookref}[3]{\@ifundefined{r@#1}{#3}{#2}}

% \PVpart and \RBpart: the two parts that a partial build may omit.
% They print "Part~N" in any build that contains the part, and a naming phrase
% in a build that does not, so an excerpt shows no ?? .  Each already supplies
% the word "Part", so write "\PVpart develops ...", not "Part~\PVpart".
\newcommand{\PVpart}{\bookref{part:aahrah}%
  {Part~\ref{part:aahrah}}{the Perron--Volterra part of the full book}}
\newcommand{\RBpart}{\bookref{part:recbrowne}%
  {Part~\ref{part:recbrowne}}{the immuno-virology part of the full book}}

% OPEN PROBLEM COMMANDS
% ASSUMPTION COMMANDS (glossaried)
\newcommand{\beAs}[1][]{%
  \def\@temparg{#1}%
  \ifx\@temparg\@empty
    \begin{assumption}%
  \else
    \let\@savedSecTitle\@currentsectionname%
    \refstepcounter{ascounter}%
    \begin{assumption}[#1]%
    \label{as:\theascounter}%
    \expanded{%
      \noexpand\newglossaryentry{as\theascounter}{%
        name={\unexpanded{#1}},%
        sort={\theascounter},%
        description={p.~\noexpand\pageref{as:\theascounter}},%
        user1={\theascounter},%
        user2={Assumption},%
        user3={\thesection},%
        user4={\expandafter\unexpanded\expandafter{\@savedSecTitle}}%
      }%
    }%
    \protected@write\@auxout{}{\string\glsliveentry{as\theascounter}}%
  \fi
}
\newcommand{\eeAs}{\end{assumption}}
% CONJECTURE COMMANDS (glossaried)
\newcommand{\beCj}[1][]{%
  \def\@temparg{#1}%
  \ifx\@temparg\@empty
    \begin{conjecture}%
  \else
    \let\@savedSecTitle\@currentsectionname%
    \refstepcounter{cjcounter}%
    \begin{conjecture}[#1]%
    \label{cj:\thecjcounter}%
    \expanded{%
      \noexpand\newglossaryentry{cj\thecjcounter}{%
        name={\unexpanded{#1}},%
        sort={\thecjcounter},%
        description={p.~\noexpand\pageref{cj:\thecjcounter}},%
        user1={\thecjcounter},%
        user2={Conjecture},%
        user3={\thesection},%
        user4={\expandafter\unexpanded\expandafter{\@savedSecTitle}}%
      }%
    }%
    \protected@write\@auxout{}{\string\glsliveentry{cj\thecjcounter}}%
  \fi
}
\newcommand{\eeCj}{\end{conjecture}}
% LEMMA COMMANDS
% ---------------------------------------------------------------------
% UNIFORM RESULT WRAPPERS -- \beL, \beC, \beP, \beR share \beT's signature
%
%   \beX[long title][short title][key]  ...  \eeX
%
%   long title  = heading (math allowed)
%   short title = shown in the cross-reference; defaults to the key
%   key         = label key.  With a key the label is X:KEY and the heading
%                 carries the mnemonic [key], exactly as for \beT.  Without
%                 one the label falls back to the running counter, which is
%                 what every pre-existing call site in the book relies on.
%
% A bare \beX -- arguments commented out -- opens the plain environment and
% creates NO label and NO glossary entry.  That is the supported way to keep
% an individual result out of the catalogue.
% ---------------------------------------------------------------------
% #1 wrapper name  #2 environment  #3 label prefix  #4 counter  #5 type name
\newcommand{\@mkresult}[5]{%
  % A leading * suppresses the [key] mnemonic in both the heading and the
  % glossary name, while still making the label X:KEY and the entry:
  %     \beX*[long][short][key]
  \expandafter\NewDocumentCommand\csname be#1\endcsname{s O{} O{} o}{%
    \def\@temparg{##2}%
    \ifx\@temparg\@empty
      \begin{#2}%
    \else
      \let\@savedSecTitle\@currentsectionname%
      \refstepcounter{#4}%
      % With a key: mnemonic label X:KEY and the [key] tag in the heading.
      % Without one: fall back to the running counter, exactly as these
      % wrappers behaved before, so every legacy one-argument call site in
      % the book keeps its label.  The counter is also csname-safe, which a
      % title's first word is not -- titles beginning with math would
      % otherwise build \csname \Omega...\endcsname and fail.
      \IfValueTF{##4}{%
        \edef\@relbl{##4}%
        \IfBooleanTF{##1}{\let\@retag\@empty}{%
          \edef\@retag{\noexpand\ \noexpand\textup{[\noexpand\texttt{#3:##4}]}}}%
      }{%
        \edef\@relbl{\arabic{#4}}%
        \let\@retag\@empty
      }%
      \def\@tempshort{##3}%
      \ifx\@tempshort\@empty
        \expandafter\xdef\csname @sshort@#3\@relbl\endcsname{\@relbl}%
      \else
        \expandafter\def\csname @sshort@#3\@relbl\endcsname{##3}%
      \fi
      \begin{#2}[##2\ifshowkeys\@retag\fi]%
      \edef\@tmplbl{#3:\@relbl}\expandafter\label\expandafter{\@tmplbl}%
      \expanded{%
        \noexpand\newglossaryentry{#3\csname the#4\endcsname}{%
          name={\unexpanded{##2}%
            \IfBooleanF{##1}{\IfValueT{##4}{ [\noexpand\texttt{#3:##4}]}}},%
          sort={\csname the#2\endcsname},%
          description={p.~\noexpand\pageref{#3:\@relbl}},%
          user1={\csname the#2\endcsname},%
          user2={#5},%
          user3={\thesection},%
          user4={\expandafter\unexpanded\expandafter{\@savedSecTitle}}%
        }%
      }%
      \protected@write\@auxout{}{\string\glsliveentry{#3\csname the#4\endcsname}}%
      \protected@write\@auxout{}{\string\@thmtitle{\@tmplbl}{\unexpanded{##2}}}%
    \fi
  }%
  \expandafter\newcommand\csname ee#1\endcsname{\end{#2}}%
}
\@mkresult{L}{lemma}{l}{lemcounter}{Lemma}
% REMARK / COROLLARY / PROPOSITION -- same signature as \beL above
\@mkresult{R}{remark}{r}{remcounter}{Remark}
\@mkresult{C}{corollary}{c}{corcounter}{Corollary}
\@mkresult{P}{proposition}{p}{propcounter}{Proposition}
% OPEN-PROBLEM COMMAND -- generated with the same signature as \beT:
%   \beO[long title][short title][key]   ...   \eeO
% label o:KEY, or o:<counter> when no key is given; a bare \beO makes
% neither a label nor a glossary entry.
\@mkresult{O}{open}{o}{opcounter}{Problem}


% Section-level title tracker: updated only by \section, not \subsection.
% \@currentlabelname is overwritten by every \refstepcounter (including
% \subsection), so we save it separately here.
\newcommand{\@currentsectionname}{}
\AtBeginDocument{%
  \let\@epidcrn@orig@refstepcounter\refstepcounter
  \renewcommand{\refstepcounter}[1]{%
    \@epidcrn@orig@refstepcounter{#1}%
    \ifstrequal{#1}{section}{\global\let\@currentsectionname\@currentlabelname}{}%
  }%
}

% TABLE COMMANDS (glossaried) -- consistent naming with \beT, \beD, etc.
\newcommand{\beTa}[1][]{%
  \def\@temparg{#1}%
  \ifx\@temparg\@empty
    \begin{table}[htbp]%
  \else
    \let\@savedSecTitle\@currentsectionname%
    \refstepcounter{tabcounter}%
    \begin{table}[htbp]%
    \label{A:\thetabcounter}%
    \expanded{%
      \noexpand\newglossaryentry{A\thetabcounter}{%
        name={\unexpanded{#1}},%
        sort={\thetabcounter},%
        description={p.~\noexpand\pageref{A:\thetabcounter}},%
        user1={\thetabcounter},%
        user2={Table},%
        user3={\thesection},%
        user4={\expandafter\unexpanded\expandafter{\@savedSecTitle}}%
      }%
    }%
    \protected@write\@auxout{}{\string\glsliveentry{A\thetabcounter}}%
  \fi
}
\newcommand{\eeTa}{\end{table}}

% FIGURE COMMANDS (glossaried) -- consistent with \beT, \beR, \beTa, etc.
% Usage:
%   \beF[Glossary title]...\caption{Visible caption}\eeF  -> glossaried figure
%   \beF...\caption{Visible caption}\eeF                  -> plain [H] figure, no glossary
\newcommand{\beF}[1][]{%
  \def\@temparg{#1}%
  \ifx\@temparg\@empty
    \begin{figure}[H]%
  \else
    \begin{figure}[H]%
    \let\@savedSecTitle\@currentsectionname%
    \refstepcounter{figcounter}%
    \label{fig:\thefigcounter}%
    \expanded{%
      \noexpand\newglossaryentry{Fig\thefigcounter}{%
        name={\unexpanded{#1}},%
        sort={\thefigcounter},%
        description={p.~\noexpand\pageref{fig:\thefigcounter}},%
        user1={\thefigcounter},%
        user2={Figure},%
        user3={\thesection},%
        user4={\expandafter\unexpanded\expandafter{\@savedSecTitle}}%
      }%
    }%
    \protected@write\@auxout{}{\string\glsliveentry{Fig\thefigcounter}}%
  \fi
}
\newcommand{\eeF}{\end{figure}}
\makeatother
% LIVE-ENTRY REGISTER
% \glsaddall prints every entry the glossaries package currently knows, and
% that set includes entries restored from \jobname.glsdefs.  The package
% rewrites .glsdefs from all *currently defined* entries, so anything it
% restored is written straight back out: a deleted section kept reappearing in
% the catalogue for ever, and the .glsdefs had to be erased by hand.  The
% wrappers therefore record every entry they actually create in the .aux, and
% \printboth adds only those.  The .aux is rebuilt from scratch on each run, so
% this is self-healing: a removed result is gone after one further run.
\makeatletter
\newcommand{\@glslive}{}
\newcommand{\@glsaddlive}[1]{\ifglsentryexists{#1}{\glsadd{#1}}{}}
\newcommand{\glsliveentry}[1]{\g@addto@macro\@glslive{\@glsaddlive{#1}}}
\makeatother
% PRINT COMMAND
% No \clearpage here -- the catalogue continues on the current page.  This has
% been re-added by mistake several times; leave it out.
\makeatletter
\newcommand{\printboth}{%
  \@glslive   % add exactly the entries this document created last run
  \printglossary[style=unifiedstyle, title={0 \hspace{.22em} Definitions, Theorems, Examples and Open Problems}]%
}
\makeatother
% Named theorem reference: \tref{word} -> Thm~N(word-or-short)
% Shows stored short title if available, otherwise shows the key itself.
\newcommand{\tref}[1]{%
  Thm~\ref{t:#1}(\ifcsname @tshort@#1\endcsname\csname @tshort@#1\endcsname\else#1\fi)%
}

% Full-title theorem reference: \nref{t:KEY} prints the number and then the
% theorem's long title, so one writes "Theorem~\nref{t:blocks}" and gets
% "Theorem 3 ``Block factorization and siphons''".  Takes the complete label, like
% \ref, not the bare key like \tref.
%
% \beT writes each title to the .aux (see \@thmtitle below), so \nref resolves
% forward references from the second LaTeX pass, exactly as \ref does.  If the
% title is not yet known -- first pass, or a label that is not a \beT theorem
% -- it degrades to the number rather than producing ??.
% These two use @ in their names, and this point in the file is outside the
% \makeatletter region opened earlier (it closes at \makeatother above), so
% the region has to be reopened here.
\makeatletter
\newcommand{\@thmtitle}[2]{\global\@namedef{@tlong@#1}{#2}}
\newcommand{\nref}[1]{%
  \ref{#1}%
  \ifcsname @tlong@#1\endcsname
    \ ``\csname @tlong@#1\endcsname''%
  \fi
}
% \nm{t:KEY}: the stored long title of any \beT/\beL/\beP/\beR/\beC/\beD/\beE result,
% verbatim and without quotes, so item names follow the macro's first argument.
\newcommand{\nm}[1]{%
  \ifcsname @tlong@#1\endcsname
    \csname @tlong@#1\endcsname%
  \else
    \ref{#1}%
  \fi
}
\makeatother




%% ========================================
%% Other THEOREM ENVIRONMENTS
%% ========================================
\makeatletter
\@ifundefined{lemma}      {\newtheorem{lemma}      {Lemma}}      {}
\@ifundefined{proposition}{\newtheorem{proposition}{Proposition}} {}
\@ifundefined{corollary}  {\newtheorem{corollary}  {Corollary}}  {}
\@ifundefined{conjecture} {\newtheorem{conjecture} {Conjecture}}  {}
\@ifundefined{remark}     {\newtheorem{remark}     {Remark}}     {}
\@ifundefined{question}   {\newtheorem{question}   {Question}}   {}
\@ifundefined{assumption} {\newtheorem{assumption} {Assumption}} {}
\@ifundefined{hypothesis} {\newtheorem{hypothesis} {Hypothesis}} {}
\@ifundefined{fact}       {\newtheorem{fact}       {Fact}}       {}
\makeatother

% \citer: cite the relay-graph threshold theorem of AAHrel. If a document has
% \input the chapter/file containing t:relayRF (so the label resolves), point
% straight at it; otherwise fall back to a bare citation of AAHrel. Shared here
% (not per-paper) so it resolves consistently across every book chapter.
% \citeR: same idea, for the first theorem of AAHRah (t:Pstar).
\makeatletter
\providecommand{\citer}{%
  \@ifundefined{r@t:relayRF}%
    {\cite{AAHrel}}%
    {Theorem~\ref{t:relayRF} of \cite{AAHrel}}%
}
\providecommand{\citeR}{%
  \@ifundefined{r@t:Pstar}%
    {\cite{AAHRah}}%
    {Theorem~\ref{t:Pstar} of \cite{AAHRah}}%
}
\makeatother


% Abbreviations
\def\beQ{\begin{question}}\def\eeQ{\end{question}}
\def\beA{\begin{assumption}}\def\eeA{\end{assumption}}


%% ========================================

%% ========================================
%% COLORS & CUSTOM COMMANDS
%% ========================================
\definecolor{funccolor}{RGB}{25,25,112}
\definecolor{desccolor}{RGB}{64,64,64}
\newcommand{\func}[1]{\texttt{\textcolor{funccolor}{#1}}}
\def\eq{=}
\def\mm{-}

%% ========================================
%% MATHCAL SHORTHANDS
%% ========================================
% Project-wide convention: \mX for a capital X is always \mathcal X.
% Defined for the whole alphabet so the convention never fails to expand.
\newcommand{\mA}{\mathcal A}\newcommand{\mB}{\mathcal B}
\newcommand{\mC}{\mathcal C}\newcommand{\mD}{\mathcal D}
\newcommand{\mE}{\mathcal E}\newcommand{\mF}{\mathcal F}
\newcommand{\mG}{\mathcal G}\newcommand{\mH}{\mathcal H}
\newcommand{\mI}{\mathcal I}\newcommand{\mJ}{\mathcal J}
\newcommand{\mK}{\mathcal K}\newcommand{\mL}{\mathcal L}
\newcommand{\mM}{\mathcal M}\newcommand{\mN}{\mathcal N}
\newcommand{\mO}{\mathcal O}\newcommand{\mP}{\mathcal P}
\newcommand{\mQ}{\mathcal Q}\newcommand{\mR}{\mathcal R}
\newcommand{\mS}{\mathcal S}\newcommand{\mT}{\mathcal T}
\newcommand{\mU}{\mathcal U}\newcommand{\mV}{\mathcal V}
\newcommand{\mW}{\mathcal W}\newcommand{\mX}{\mathcal X}
\newcommand{\mY}{\mathcal Y}\newcommand{\mZ}{\mathcal Z}

%% ========================================
%% BASIC MATH ENVIRONMENTS
%% ========================================
\def\bep{\begin{pmatrix}}\def\eep{\end{pmatrix}}
\def\bev{\begin{vmatrix}}\def\eev{\end{vmatrix}}
\def\bea{\begin{eqnarray*}}\def\eea{\end{eqnarray*}}
\def\bc{\begin{cases}}\def\ec{\end{cases}}
\def\BEN{\begin{enumerate}}\def\EEN{\end{enumerate}}
\def\BI{\begin{itemize}}\def\EI{\end{itemize}}

\newcommand{\be}[1]{\begin{equation}\label{#1}}
\newcommand{\ee}{\end{equation}}
\newcommand{\beq}{\begin{eqnarray}}
\def\eeq{\end{eqnarray}}
\def\eqr{\eqref}\def\fr{\frac}\def\lbl{\label}
\def\Lra{\Longrightarrow}\def\non{\nonumber}\def\cNGM{\cite{Diek,Van,Van08}}
\def\Eq{\Leftrightarrow}\def\b{\beta}\def\uph{\red{UP TO HERE}}
%% ========================================
%% PROBABILITY & STATISTICS
%% ========================================
\renewcommand{\P}{\ensuremath{\mathbb{P}}}  % Note: overrides paragraph symbol
\newcommand{\E}{\ensuremath{\mathbb{E}}}
\newcommand{\ind}[1]{\ensuremath{\mathbbm{1}_{\left\{#1\right\}}}}
\DeclareMathOperator{\Cov}{Cov}
\DeclareMathOperator{\TV}{TV}
\newcommand{\rank}{\operatorname{rank}}

%% ========================================
%% CRN OPERATORS (organizations/ERCs/buffering/siphons framework)
%% ========================================
\def\clos{\mathrm{clos}}
\def\Supply{\mathrm{Supply}}
\def\req{\mathrm{req}}
\def\React{\mathrm{React}}
\def\Prod{\mathrm{Prod}}
\def\Gen{\mathrm{Gen}}
\def\In{\mathrm{In}}
\def\Out{\mathrm{Out}}
\def\Act{\mathrm{Act}}
\def\Min{\mathrm{Min}}

%% ========================================
%% LINE BREAKING
%% ========================================
% Last-resort stretch, used only on paragraphs TeX cannot otherwise break:
% a long \texttt{} token, a comma-separated inline math list (commas are not
% breakpoints in math), or a run of \cite/\ref boxes.  Without it those set as
% overfull \hbox and print into the margin; with it they set slightly loose,
% which is the lesser evil.  Does not affect paragraphs that already break.
\emergencystretch=3em

%% ========================================
%% NUMBER SETS
%% ========================================
\newcommand{\R}{\mathbb{R}}  % Real numbers
\newcommand{\C}{\mathbb{C}}  % Complex numbers
\newcommand{\N}{\mathbb{N}}
\newcommand{\Z}{\mathbb{Z}}
% \providecommand, not \newcommand: some papers define these in their own
% preamble after \input{preamble}, and must not clash.
\providecommand{\Q}{\mathbb{Q}}
\providecommand{\res}{\operatorname{res}}   % resultant; used by epiRes.tex
\newcommand{\MN}{\mathbb{N}}
\newcommand{\MZ}{\mathbb{Z}}
\newcommand{\realpos}{\mathbb{R}_{> 0}}
\newcommand{\realnneg}{\mathbb{R}_{\geq 0}}
\newcommand{\naturalpos}{\mathbb{N}_{>0}}
\newcommand{\Rplus}{\mathbb{R}_+}
\newcommand{\Rp}{\mathbb{R}_{\ge 0}}
\def\Rnne{\mathbb{R}_{\ge 0}}
\def\Se{\bar{S}} \def\Ie{\bar{I}}
\newcommand{\Rpos}{\mathbb{R}_{>0}}
\newcommand{\dd}{\,\mathrm d}
\newcommand{\e}{\mathrm e}
%% ========================================
%% GREEK LETTERS (2+ character abbreviations) greek
%% ========================================
\def\al{\alpha}\def\ba{\bff \alpha}
\def\w{\omega}\def\Lgn{\ga+\mu+\de}\def\om{\omega}
\def\de{\delta}\def\De{\Delta}\def\bde{\bff \delta}
\def\ep{\epsilon}\def\eps{\varepsilon}
\def\f{\varphi}  % Kept as requested
\def\ga{\gamma}\def\Ga{\Gamma}\def\G{\Gamma}\def\bg{\bff \gamma}
\def\up{\upsilon}  % NOT the breve accent: no accent commands anywhere.
                   % Named \up, not \u, so the LaTeX breve stays reachable.
\def\la{\lambda}\def\La{\Lambda}
\def\si{\sigma}\def\ka{\kappa}
\def\th{\theta}\def\vt{\vartheta}
\def\ze{\zeta}\def\Gap{\Gamma^{+}} \def\Gam{\Gamma^{-}}


%% ========================================
%% SPECIAL SHORT COMMANDS (kept as requested)
%% ========================================
\def\I{\infty}
\def\T{\widetilde}
\def\H{\hat}
\newcommand{\relint}{\operatorname{relint}}

%% ========================================
%% BOLD/VECTOR NOTATION
%% ========================================
\newcommand{\bff}[1]{{\mbox{\boldmath$#1$}}}
\newcommand\X{\boldsymbol{X}}
\def\x{\boldsymbol{x}}\def\bx{\bff x}\def\xe{\bar{\x}}
\newcommand\y{\boldsymbol{y}}\def\by{\boldsymbol{y}}\def\ye{\bar{\y}}
\def\yb{\y}\def\xb{\x}
\def\xy{\x-\y}
\newcommand\br{\boldsymbol{r}}\def\re{\bar{\r}}
\newcommand\bd{\boldsymbol{d}}
\newcommand\bQ{\boldsymbol{Q}}
\newcommand{\Yb}{{\mathbf{Y}}}
\newcommand{\Ab}{\mathbf{A}}
\newcommand{\zb}{\mathbf{z}}\def\bz{\bff 0}
\newcommand{\Zb}{\mathbf{Z}}
\newcommand{\bb}{\mathbf{b}}
\def\vbe{\vec {\bff \beta}}
\def\vb{\vec {\bff b}}
\def\v1{\vec {\bff 1}}
\def\vn{\vec \de}
\def\vz{\vec 0}
\def\vi{\overset{\rightarrow}{\mathsf i}}
\def\vv{\overset{\rightarrow}{v}}
\def\vr{\; \vec {\mathsf r}}
\newcommand{\va}{\bff a}
\providecommand{\vect}[1]{\boldsymbol{#1}}   % bold vector (moved from Gradostat.tex)
\providecommand{\one}{\mathbf 1}             % all-ones vector (moved from AliSF)
\providecommand{\spr}{s}                     % spectral bound  (moved from AliSF)
\providecommand{\Ke}{K_{\rm eff}}            % effective carrying capacity
% Critical right/left eigenvectors in the Hopf projection method (AAHOSN/New.tex)
\providecommand{\bq}{\mathbf{q}}
\providecommand{\bp}{\mathbf{p}}


%% ========================================
%% MATHCAL & MATHSCR
%% ========================================
% Mathcal
\newcommand\mc[1]{\mathcal{#1}}
\def\mC{{\mathcal C}}\def\mD{\mathcal D}\def\mE{\mathcal{E}}
\def\mF{\mathcal F}\def\mG{{\mathcal G}}\def\mH{{\mathcal H}}
\def\mI{\mathcal I}\def\mL{\mathcal{L}}\def\mM{{\mathcal M}}
\def\mN{{\mathcal N}}\def\mP{\mathcal{P}}\def\mR{{\mathcal R}}
\def\mS{{\mathcal S}}\def\mT{\mathcal{T}}\def\mV{{\mathcal V}}
\def\mK{\mathcal K}\def\mA{{\mathcal A}}\def\mU{{\mathcal U}}
% Symmetric bilinear and trilinear forms of the Taylor expansion, used in the
% first Lyapunov coefficient (AAHOSN/New.tex)
\providecommand{\Bform}{\mathcal{B}}
\providecommand{\Cform}{\mathcal{C}}

% Blackboard bold.  Kept here, not in the drivers: a \providecommand in a
% driver preamble runs before this file and silently wins, so the same name
% ends up meaning different things in different builds.
\def\IC{\mathbb{C}}\def\IN{\mathbb{N}}\def\IQ{\mathbb{Q}}
\def\IR{\mathbb{R}}\def\IZ{\mathbb{Z}}\def\IG{\mathbb{G}}
\def\II{\mathbb{I}}

% Mathscr
\newcommand{\Rs}{{\mathscr R}}
\newcommand{\Ss}{{\mathscr S}}
\newcommand{\Ns}{{\mathscr N}}
\newcommand{\Cs}{\mathscr C}


%% ========================================
%% OPERATORS & FUNCTIONS
%% ========================================
\DeclareMathOperator{\tr}{tr}
\DeclareMathOperator{\diag}{Diag}
\DeclareMathOperator{\Diag}{Diag}
\DeclareMathOperator{\intr}{\mathrm{Int}}
\DeclareMathOperator{\Inv}{Inv}
\DeclareMathOperator{\spec}{spec}
\DeclareMathOperator{\AUC}{AUC}               % area under the curve
\newcommand{\ii}{\mathrm{i}}                  % imaginary unit
\newcommand\op{\operatorname}
\newcommand{\supp}{\op{supp}}
\newcommand{\diff}{\mathop{}\mathopen{}\mathrm{d}}
\newcommand{\pdd}[1]{\frac{\partial}{\partial{#1}}}
\DeclareMathOperator{\cone}{cone}
\DeclareMathOperator{\proj}{proj}

%% ========================================
%% BRACKETS & NORMS
%% ========================================
\providecommand{\pr}[1]{\left(#1\right)}
\providecommand{\pp}[1]{\left[#1\right]}
\newcommand\croc[1]{\left\langle #1\right\rangle}
\newcommand{\norm}[1]{\|#1\|}
\newcommand{\inner}[2]{\langle #1, #2 \rangle}
% Short synonym of \inner, used by AAHOSN/New.tex as $\ip{u}{v}=\bar u^\top v$
\providecommand{\ip}[2]{\langle #1 , #2 \rangle}


%% ========================================
%% FORMATTING & UTILITIES
%% ========================================
\newcommand*{\Scale}[2][4]{\scalebox{#1}{$#2$}}
\newcommand{\te}[1]{\text{#1}}
\newcommand\steq[1]{\stackrel{\text{\rm #1.}}{=}}
\newcommand\blue[1]{\textcolor{blue}{#1}}
\newcommand{\red}{\textcolor[rgb]{1.00,0.00,0.00}}
\newcommand{\green}{\textcolor[rgb]{0.00,1.00,0.00}}

% Figure command
\newcommand{\figu}[3]{%
\begin{figure}[H]
\centering
\includegraphics[scale=#3]{#1}
\caption{#2\label{f:#1}}
\end{figure}
}

% Footnote
\long\def\symbolfootnote[#1]#2{%
\begingroup
\def\thefootnote{\fnsymbol{footnote}}\footnote[#1]{#2}%
\endgroup}
\def\fn{\symbolfootnote}

\newcommand{\dsum}{\displaystyle\sum}

%% ========================================
%% CRN & EPIDEMIOLOGY SPECIFIC
%% ========================================
% CRN notation
\def\SS{\mathcal S}
\def\CC{\mathcal C}
\def\RR{\mathcal R}

% Reproduction numbers
\def\Rto{R_{1}^{E_2}}
\def\Rot{R_{2}^{E_1}}
\def\Roc{R_{1,c}}
\def\Rtc{R_{2,c}}
\def\rd{r_{0}}
\def\sd{s_{0}}
\def\dd{D_{0}}
\def\vd{v_{0}}
\def\zd{z_{0}}
\def\se{\sE}
\def\xd{x_{0}}
%% ========================================
%% SPECIAL VALUES & EQUILIBRIA
%% ========================================
\def\eE{\bar{e}}
\def\wE{\bar{w}}
\def\sE{s_{*}}
\def\Me{\T M_{ee}}
\def\Rnne{\mathbb{R}_{\ge 0}}

\newcommand{\ones}{\mathbf 1}


%% ========================================
%% ABBREVIATIONS (alphabetical by category)
%% ========================================

% -- General mathematical terms --
\def\adm{admissible}\def\admy{admissibility}
%\def\and{antisymmetric}% disabled 2026-09-30: it overrode LaTeX's author separator \and (printed "antisymmetric" between authors)
\def\ass{assumption}
\def\cadlag{c\`adl\`ag}
\def\ch{characteristic polynomial}
\def\chp{characteristic polynomial}
\def\coe{coefficient}
\def\comp{computation}
\def\con{condition}
\def\coo{coordinate}
\def\corr{corresponding}
\def\dec{decomposition}
\def\den{denominator}
\def\dete{deterministic}
\def\eig{eigenvalue}\def\eval{eigenvalue}
\def\eigv{eigenvector}\def\eiv{eigenvector}\def\evec{eigenvector}
\def\elt{element}
\def\enn{essentially nonnegative/positivity preserving}
\def\equ{equation}
\def\etal{et al.}
\def\fe{for example}
\def\fno{from now on}
\def\for{\forall}
\def\form{formula}
\def\Fr{Furthermore, }
\def\frt{furthermore}
\def\fun{function}
\def\gene{generalization}
\def\heu{heuristic}
\def\how{however}\def\How{However, }
\def\ie{i.e. }
\def\im{\item}
\def\isd{is defined by}
\def\itc{it is easy to check that}
\def\itf{it follows that}\def\Itf{It follows that}
\def\ith{it holds that}
\def\itm{it may be checked that}\def\Itm{It may be checked that}
\def\ito{it turns out that}\def\Ito{It turns out that}\def\Itou{It turns out useful to}
\def\ma{mass action}\def\MAR{mass action representation}
\def\mat{matrix}
\def\mbc{may be checked to be}
\def\mbw{may be written as}\def\mubw{must be written as}
\def\mpr{more precisely, }\def\Mp{More precisely, }
\def\nc{non-catalytic}
\def\nne{non-negative}\def\nny{nonnegativity}
\def\no{\nonumber}
\def\noni{non-infectious}
\def\nota{notation}
\def\np{\newpage}
\def\Nt{Note that}
\def\nt{non trivial}\def\nz{non zero}
\def\oth{otherwise}
\def\Oth{On the other hand}
\def\pa{parameter}\def\para{parameter}
\def\pol{polynomial}
\def\prob{probability}
\def\Prf{{\bf  Proof:}}\def\prf{{\bf  Proof:}}
\def\prop{proportional to}
\def\qu{\quad}
\def\QED{\hfill {$\square$}\goodbreak \medskip}
\def\rep{representation}
\def\resp{respectively}
\def\satd{satisfied}\def\satg{satisfying}\def\sats{satisfies}\def\saty{satisfy}
\def\ssec{\subsection}\def\sssec{\subsubsection}
\def\sef{see for example}
\def\sksy{skew-symmetric}
\def\sub{substrate}
\def\sys{system}
\def\var{variable}
\def\wk{well-known}
\def\wl{without loss of generality}\def\wlo{w.l.o.g.}
\def\wmc{We may check that}\def\Wmc{We may check that}
\def\wmr{we may rewrite}
\def\wms{we must show that}
\def\wrt{with respect to}

% -- CRN & Chemistry --
\newcommand\CRN{chemical reaction network}
\newcommand\CRNs{chemical reaction networks}
\newcommand\RN{reaction network}
\def\ACR{absolute concentration robustness}
\def\AMSD{acyclic siphon decomposition}
\def\DAE{differential algebraic equations}
\def\FHJ{Feinberg-Horn-Jackson}
\def\GMAK{generalized mass action kinetics}
\def\MM{Michaelis-Menten}
\def\sm{stoichiometric matrix}\def\Sm{stoichiometric matrix}
\def\WRZD{weakly reversible, zero deficiency}\def\aut{autocatalytic}

% -- Epidemiology --
\def\brn{basic reproduction number}\def\bzn{basic replacement number}
\def\com{compartment}
\def\dfe{disease free equilibrium}\def\DFE{disease free equilibrium}
\def\dwt{dwell times}\def\edt{expected dwell times}
\def\EM{epidemic model}
\def\endp{endemic fixed point}\def\enp{endemic point}\def\EE{endemic point}
\def\ME{mathematical epidemiology}
\def\MEE{eco-epidemiology}\def\EEP{eco-epidemiological}
\def\ece{eco-epidemiological model}
\def\mgm{multi group model}
\def\mtS{multi-strain model}\def\mSt{multi-strain model}
\def\mSi{minimal siphon}
\def\NGM{next generation matrix}
\def\nim{new infections matrix}
\def\qem{quadratic epidemiological model}
\def\repF{invasion functions}
\def\SPF{SIR-PH-FA model}
\def\sus{susceptible}

% -- Graph theory & Linear algebra --
\def\jac{Jacobi matrix}
\def\jin{Jacobian matrix of the invasion vector field with respect to the invasion variables}
\def\jrn{Jacobian matrix of the resident vector field with respect to the resident variables}
\def\jxn{Jacobian matrix of the $\x$ vector field with respect to the $\x$ variables}
\def\Lm{Laplacian matrix}
\def\MTT{matrix tree theorem}
\def\SGT{spectral graph theory}
\def\SMW{Sherman–Morrison–Woodbury}

% -- Stability & Dynamics --
\def\absC{absorbtion CTMC}
\def\BTb{Bogdanov-Takens bifurcation}
\def\CEP{competitive exclusion partition}
\def\CTMC{continuous time Markov chain}
\def\ds{differential system}
\def\fp{fixed point}\def\bfp{boundary fixed point}\def\fbp{fixed boundary point}
\def\fbl{fixed boundary line}\def\fif{forward invariant face}
\def\gas{global asymptotic stability}\def\GAS{globally asymptotically stable}
\def\GASS{\text{GASS}}
\def\Hb{Hopf bifurcation}
\def\Hur{Hurwitz}\def\MH{Metzler-Hurwitz}\def\RH{Routh-Hurwitz}
\def\IS{\text{IS}}
\def\LAS{locally asymptotically stable}
\def\LdS{Lyapunov diagonally stable}
\def\Lf{Lyapunov function}\def\LyV{Lyapunov Volterra type function}\def\RLF{robust Lyapunov function}
\def\LyM{Lyapunov Malkin stable}
\def\PF{Perron-Frobenius}
\def\spfp{strictly positive fixed point}
\def\sta{stability}
\def\STP{sharp threshold property}\def\STT{sharp threshold theorem}

% -- Models (with \varphi explicitly instead of \f) --
\def\AB{$(A,B)$ Arino-Brauer epidemic models}
\def\ABa{$(A,\bbe,\ba,\varphi)$ balanced bilinear model}
\def\ABP{%$(A,\bbe,P,\varphi)$
\\ balanced bilinear model}
\def\AMG{Anderson May Gupta}
\newcommand\BIN{biological interaction network}
\def\Dt{Descartes type}
\def\gLV{generalized Lotka-Volterra model}
\def\IG{\text{invasion graph}}\def\IGMS{interaction graph on minimal siphons}
\def\LCP{\text{linear complementarity problem}}
\def\LV{Lotka-Volterra}\def\LVK{Lotka-Volterra-Kolmogorov}

% -- Splitting & Structure --
\def\reg{regular splitting}\def\regS{{\bf regular splitting}}
\def\RUR{rational univariate representation}
\def\bds{block-diagonal structure}\def\NGMs{next generation matrices}

% -- Miscellaneous --
\def\Ma{{Mathematica}}
\def\Mm{Max–Min}
\def\num{numerator}
\def\rap{rational approximation} \def\DSR{directed species-reaction bipartite graph}\def\PCP{polynomial complementarity problem}\def\rat{rational substitution}
\def\QSS{quasi steady state} \def\RHS{right hand side} \def\LHS{left hand side}
\def\OSN{online social network}\def\VdD{Van den Driessche}
\def\wli{$\omega$-limit} \def\rk{rich kinetics}\def\sJ{symbolic Jacobian}
\def\cV{\cite{Vas,VasSta,BSV,GSV}} \def\spf{semi-parametric semi-functional}
\def\CRW{Capasso-Ruan-Wang} \def\ABP{$(A,\bbe,P,\f)$ balanced bilinear model}
\def\ABa{%$(A,\bbe,\ba,\f)$ \\
rank-one balanced bilinear model}
\def\bbe{B%{\bff \beta}
} \def\vbe{\vec {\bff \beta}
}
\def\vb{\vec {\bff b}}
\newcommand{\Kh}{\hat{K}}          % K with hat, first mention only
\newcommand{\wS}{w^{\sigma}}       % right Perron J^\perp, full block
\newcommand{\piS}{\vec \pi^{\sigma}}    % left Perron J^\perp, full block
\newcommand{\ws}{w^{\sigma}}       % right Perron J^\perp, sub-block
\newcommand{\pis}{\pi^{\sigma}}    % left Perron J^\perp, sub-block
\newcommand{\hwS}{\hat{w}^{\sigma}}   % right Perron K, full block
\newcommand{\hpiS}{\hat{\pi}^{\sigma}} % left Perron K, full block
\newcommand{\hws}{\hat{w}^{\sigma}}   % right Perron K, sub-block
\newcommand{\hpis}{\hat{\pi}^{\sigma}} % left Perron K, sub-block
\newcommand{\twS}{\tilde{w}^{\sigma}}  % right Perron \tilde K, full block
\newcommand{\tpiS}{\tilde{\pi}^{\sigma}} % left Perron \tilde K
\newcommand{\tws}{\tilde{w}^{\sigma}}
\newcommand{\tpis}{\tilde{\pi}^{\sigma}}
\newcommand{\Mat}[2]{\mathbb{R}_{\ge 0}^{\,#1\times #2}}
\def\tJ{J_\sigma^\perp} \def\lPe{left Perron eigenvector}
% \G is \Gamma (defined at :795).  The Volterra function keeps its own name,
% so that \G no longer has two conflicting meanings in this file.
\def\GV{G}  % Volterra G(u)=u-1-\ln u
\def\PVL{Perron-Volterra Lyapunov minimal siphon ansatz}
\newcommand{\istar}[1]{i_{#1}^{*}}
%\newcommand{\sE}{s_E}
\newcommand{\vE}{v_E}
\def\PV{Perron-Volterra candidate Lyapunov function}
\def\GLV{generalized Lotka-Volterra}
\def\MERN{mathematical epidemiology reaction networks}
\def\CMM{contraction metric method}
\newcommand{\sign}{\operatorname{sign}}
\newcommand{\Var}{\operatorname{Var}} 

% ======================================================================
% IDEAS SHARED BETWEEN THE ABSTRACT AND THE MAIN RESULTS SECTION
%
% Write each idea ONCE, in the abstract; the Main Results section recalls it,
% so the two can never drift apart.
%
%   in the abstract        \ideaitem{1}{The first-order reaction terms ...}
%   in Main Results        \ideahead{1}
%
% \ideaitem stores the text globally and typesets
%   \item \hyperlink{idea:1}{<text>}
% \ideahead typesets
%   \item\hypertarget{idea:1}{}\textbf{<the same text>}
%
% The abstract must come before the Main Results input, which is the normal
% order. If a heading has to read differently from the abstract sentence (for
% instance a long abstract item summarised in one line), override it with the
% optional argument: \ideahead[short heading]{6}.
% ======================================================================
\makeatletter
\newcommand{\ideaitem}[2]{%
  \expandafter\gdef\csname @idea@#1\endcsname{#2}%
  \item\hyperlink{idea:#1}{#2}}
\newcommand{\ideahead}[2][]{%
  \item\hypertarget{idea:#2}{}%
  \def\@ideatmp{#1}%
  \ifx\@ideatmp\@empty
    \@ifundefined{@idea@#2}{%
      \PackageWarning{preamble}{Idea #2 is used in Main Results but was never
        defined by \string\ideaitem\space in the abstract}%
      \textbf{[idea #2 undefined]}%
    }{\textbf{\begingroup\def\ideaterm##1##2{##2}\def\idearest##1{}\@nameuse{@idea@#2}\endgroup}}%
  \else
    \textbf{#1}%
  \fi}
% \ideaterm{<hypertarget>}{<words>}: inside an abstract idea the words become an
% invisible link (the document sets \hypersetup{hidelinks}) jumping to the place
% in Main Results that defines them; inside the Main Results heading itself,
% where the target sits a few lines below, \ideahead strips the link and prints
% the words plainly. Mark the destination with \hypertarget{<name>}{} on the line
% before the display, which adds no equation number and shows nothing.
\newcommand{\ideaterm}[2]{\hyperlink{#1}{#2}}
% \idearest{<text>}: the part of an abstract idea after its first sentence; printed in
% the abstract, omitted from the Main Results heading.
\newcommand{\idearest}[1]{#1}
\makeatother

% ======================================================================
% SHORT-PAPER SWITCH
% \shorttrue in a master that assembles only part of the companion paper
% (Barrat/SipArxiv.tex). \nsh{...} drops text that cites results living in the
% omitted sections, so the short build has no dangling references while the
% companion paper keeps the text verbatim.
% ======================================================================
\newif\ifshort \shortfalse
\newcommand{\nsh}[1]{\ifshort\else#1\fi}
% \shortalt{<short build>}{<companion paper>}: use when a sentence needs
% different wording in the two builds, not merely a deletion.
\newcommand{\shortalt}[2]{\ifshort#1\else#2\fi}
% \Jtr, \JtrS{\Sigma}: transverse Jacobian. J^\perp here; SipArxiv.tex redefines them as J_{\mathrm{tra}}, J_{\mathrm{tra},\Sigma}
\newcommand{\Jtr}{J^\perp}
\newcommand{\JtrS}[1]{J_{#1}^\perp}
.
\makeatletter
\@ifpackageloaded{hyperref}{}{%
  \@ifundefined{hyperrefpackage}{%
    \@ifundefined{AAdefNnohyperref}{\usepackage{hyperref}}{}%
  }{}%
}
\makeatother
% linktoc=all makes a table-of-contents page number a link as well as the
% entry text; hyperref's default, linktoc=section, links only the text.
% Deferred to begin-document and guarded: by the block above hyperref may
% be loaded later (mdpi's own hook) or not at all (\AAdefNnohyperref), and
% \hypersetup must exist before it is called.  A driver that wants the
% default back can say \hypersetup{linktoc=section} after % Shared preamble for the GLVbook papers.  This is the single canonical copy;
% a driver one level below GLVbook reaches it as
%
%     \input{../basicFiles/preamble}
%
% What stays in the driver, because it differs from paper to paper:
% \documentclass, font choice, \title, \author, \date, and any macro genuinely
% local to one paper.  Everything else belongs here.  The file is
% self-contained: it does not read any other file.

\usepackage[utf8]{inputenc}
\usepackage[novbox]{pdfsync}   % required: SyncTeX alone does not work here

% Page geometry, set once here for every paper that loads this file.  A driver
% that needs different margins says \geometry{...}; a second
% \usepackage[...]{geometry} raises an option clash and must not be used.
\makeatletter
\@ifundefined{AAdefNnogeometry}{\usepackage[a4paper,margin=2cm]{geometry}}{}
\makeatother

% mathematics
\usepackage{amsmath,amssymb,amsthm,mathtools,bm,mathrsfs,bbm}
\usepackage{centernot}

% floats, tables, lists
\usepackage{graphicx,float}
\makeatletter
\@ifpackageloaded{subfig}{}{\usepackage{subcaption}}
\makeatother
\usepackage{booktabs,longtable,threeparttable,array,multirow}
\usepackage{enumitem}

% colour, code, algorithms
\usepackage{xcolor,listings,tcolorbox}
\makeatletter
\@ifundefined{c@algorithm}{\usepackage{algorithm}}{}
\makeatother
\usepackage{algorithmicx,algpseudocode}

% diagrams
\usepackage{tikz,tikz-cd,pgfplots}
\usetikzlibrary{graphs,decorations.pathmorphing,decorations.markings,
arrows.meta,automata,positioning}
\tikzcdset{every label/.append style = {font = \small}}
\pgfplotsset{compat=1.18}

% chemistry
\usepackage[version=4]{mhchem}

\usepackage{microtype}
\usepackage{lineno}
\usepackage{nomencl}
\makenomenclature

% hyperref first, then bookmark, then glossaries: glossaries must follow
% hyperref or the entry links are built against the wrong anchors.
\makeatletter
\@ifundefined{hyperrefpackage}{%
  \@ifundefined{AAdefNnohyperref}{\usepackage{hyperref}}{}%
}{}
\@ifundefined{hyperrefpackage}{\usepackage{bookmark}}{}
\makeatother
\usepackage[toc,sort=use]{glossaries}
\makeglossaries

% Book/standalone switch.  Material that the book already covers in its
% prelude, but that a standalone paper still has to state, is wrapped as
%     \ifbookbuild\else ... \fi
% Drivers default to standalone; bookM.tex says \bookbuildtrue.
\newif\ifbookbuild
\bookbuildfalse

% Part switches for the gradostat paper.  Both default to true; a driver that
% wants the short submission says \madecfalse and/or \slowfastfalse, and every
% cross-reference into the dropped part disappears with it.
%   \ifmadec     Part III, the Lagasquie--Madec model and its break-even geometry
%   \ifslowfast  Part IV, everything carrying the small parameter \eps
\newif\ifmadec     \madectrue
\newif\ifslowfast  \slowfasttrue

% Part switch for the Barrat/simplicial-contagion project.  Defaults to
% true (the hypernetwork-synchrony part is included); the journal-paper
% driver says \hnetpartfalse to drop it and every cross-reference into it.
\newif\ifhnetpart  \hnetparttrue

%% ========================================
%% PACKAGES
%% ========================================
\RequirePackage{etoolbox}
\usepackage{makecell}
\usepackage[novbox]{pdfsync}% obsolete; SyncTeX is built into modern engines
\usepackage{amssymb, amsfonts,mathtools,bbm,eurosym,graphicx,latexsym, amsmath,amsthm,times,nccmath,url}
\usepackage{listings,verbatim,cases,tabularx}
% hyperref is deliberately NOT in the line above.  Journal classes often
% load it themselves, with a long option list, from a begindocument hook
% (mdpi.cls does: \DeclareHookRule{begindocument}{hyperref}{before}{mdpi}).
% Loading it here first, without options, makes that later load raise
% "Option clash for package hyperref".  So load it only if nobody else
% has, or will:
%   - \@ifpackageloaded  catches a class that loaded it already;
%   - \hyperrefpackage    is mdpi.cls's own deferred loader;
%   - \AAdefNnohyperref   lets any other driver opt out by defining it
%                         before \input{preamble}.
\makeatletter
\@ifpackageloaded{hyperref}{}{%
  \@ifundefined{hyperrefpackage}{%
    \@ifundefined{AAdefNnohyperref}{\usepackage{hyperref}}{}%
  }{}%
}
\makeatother
% linktoc=all makes a table-of-contents page number a link as well as the
% entry text; hyperref's default, linktoc=section, links only the text.
% Deferred to begin-document and guarded: by the block above hyperref may
% be loaded later (mdpi's own hook) or not at all (\AAdefNnohyperref), and
% \hypersetup must exist before it is called.  A driver that wants the
% default back can say \hypersetup{linktoc=section} after \input{preamble}.
\makeatletter
\AtBeginDocument{\@ifundefined{hypersetup}{}{\hypersetup{linktoc=all}}}
\makeatother
\usepackage{graphicx,float,fancybox}
\usepackage{epsfig,epstopdf,tikz}
\usepackage[sans]{dsfont}
\usetikzlibrary{positioning,arrows,arrows.meta,calc}
% Configure listings for Mathematica code
\lstdefinestyle{mathematica}{
 language=Mathematica,
 basicstyle=\small\ttfamily,
 backgroundcolor=\color{gray!10},
 frame=single,
 breaklines=true,
 commentstyle=\color{green!60!black},
 keywordstyle=\color{blue},
 stringstyle=\color{red},
 showstringspaces=false,
 tabsize=2
}
\usepackage{tcolorbox}
% Colored boxes for key results
\newtcolorbox{keyresult}[1][]{
 colback=blue!5!white,
 colframe=blue!75!black,
 title=#1,
 fonttitle=\bfseries
}
\newtcolorbox{examplebox}[1][]{
 colback=green!5!white,
 colframe=green!75!black,
 title=#1,
 fonttitle=\bfseries
}


\makeatletter
\@ifundefined{theorem}   {\newtheorem{theorem}   {Theorem}}   {}
\@ifundefined{definition}{\newtheorem{definition}{Definition}} {}
\@ifundefined{open}      {\newtheorem{open}      {Problem}}   {}
\@ifundefined{example}   {\newtheorem{example}   {Example}}   {}
\makeatother

% COUNTER DEFINITIONS
\newcounter{tabcounter}
\newcounter{defcounter}
\newcounter{opcounter}
\newcounter{thmcounter}
\newcounter{excounter}
\newcounter{ascounter}
\newcounter{cjcounter}
\newcounter{lemcounter}
\newcounter{remcounter}
\newcounter{corcounter}
\newcounter{propcounter}
\newcounter{figcounter}

% UNIFIED GLOSSARY STYLE
\newglossarystyle{unifiedstyle}{%
  \setglossarystyle{list}%
  \renewcommand*{\glossarypreamble}{\gdef\glslastsec{}}%
  \renewcommand{\glossentry}[2]{%
    \edef\glscursec{\glsentryuseriii{##1}}%
    \ifx\glscursec\glslastsec\else
      \item[]\textbf{\S\,\glscursec\enspace\glsentryuseriv{##1}}%
      \global\let\glslastsec\glscursec
    \fi
    \item[\glsentryuserii{##1} \glsentryuseri{##1}:] ``\glossentryname{##1}'' \glsentrydesc{##1}%
  }%
}

\makeatletter
% First-word extractor (used by \beT, \beD and \beXa for the default key)
% \@firstwordof hello world \relax  ->  hello
\def\@firstwordof#1 #2\relax{#1}
% First three letters, used for the heading mnemonic of a theorem given no
% explicit key.  Padding with \@empty keeps it safe for first words shorter
% than three letters: \@empty vanishes under \edef, whereas \relax would
% survive into the label and is illegal inside \csname (it broke every
% two-letter title, e.g. "ME ...", with "Missing \endcsname inserted").
\def\@takethree#1#2#3#4\@nil{#1#2#3}

% EXAMPLE COMMAND -- same signature and behaviour as \beT and \beD.
% Usage: \beXa*[long title][short title][key]
%   star  = suppresses the printed "[key]" tag after the title (as for \beT*
%   and \beD*); the label and glossary entry are unaffected.
%   key = optional label key override; the label becomes e:KEY, and without a
%   key it is e: plus the first three letters of the title's first word.  A
%   mnemonic label lets an example be cited as "Example~\ref{e:KEY}" instead
%   of by section number, which is what the citing prose needs.
\NewDocumentCommand{\beXa}{s O{} O{} o}{%
  \def\@temparg{#2}%
  \ifx\@temparg\@empty
    \begin{example}%
  \else
    \let\@savedSecTitle\@currentsectionname%
    \refstepcounter{excounter}%
    \IfValueTF{#4}{%
      \edef\@exlblkey{#4}%
    }{%
      \edef\@exfirstword{\@firstwordof#2 \relax}%
      \edef\@exlblkey{%
        \expandafter\@takethree\@exfirstword\@empty\@empty\@empty\@nil}%
    }%
    \def\@tempshort{#3}%
    \ifx\@tempshort\@empty
      \expandafter\xdef\csname @xshort@\@exlblkey\endcsname{\@exlblkey}%
    \else
      \expandafter\def\csname @xshort@\@exlblkey\endcsname{#3}%
    \fi
    \IfBooleanTF{#1}{\let\@extag\@empty}{%
      \edef\@extag{\noexpand\ \noexpand\textup{[\noexpand\texttt{e:\@exlblkey}]}}}%
    \begin{example}[#2\ifshowkeys\@extag\fi]%
    \label{e:\@exlblkey}%
    % record the long title so \nref{e:KEY} prints it, as for theorems
    \protected@write\@auxout{}{%
      \string\@thmtitle{e:\@exlblkey}{\unexpanded{#2}}}%
    \expanded{%
      \noexpand\newglossaryentry{e\theexcounter}{%
        name={\unexpanded{#2}%
          \IfBooleanF{#1}{ [\noexpand\texttt{e:\@exlblkey}]}},%
        sort={\theexample},%
        description={p.~\noexpand\pageref{e:\@exlblkey}},%
        user1={\theexample},%
        user2={Example},%
        user3={\thesection},%
        user4={\expandafter\unexpanded\expandafter{\@savedSecTitle}}%
      }%
    }%
    \protected@write\@auxout{}{\string\glsliveentry{e\theexcounter}}%
  \fi
}
\newcommand{\eeXa}{\end{example}}

% DEFINITION COMMAND -- same signature and behaviour as \beT below.
% Usage: \beD[long title][short title][key]
%   long title  = definition heading (plain text only)
%   short title = shown in \dref{word}; defaults to the key
%   key         = optional label key override; the label becomes d:KEY
% Without a key the label is d: plus the first three letters of the title's
% first word, so it is a mnemonic rather than the definition's number: a
% numeric label would silently change whenever a definition is inserted.
\NewDocumentCommand{\beD}{s O{} O{} o}{%
  \def\@temparg{#2}%
  \ifx\@temparg\@empty
    \begin{definition}%
  \else
    \let\@savedSecTitle\@currentsectionname%
    \refstepcounter{defcounter}%
    \IfValueTF{#4}{%
      \edef\@deflblkey{#4}%
    }{%
      \edef\@deffirstword{\@firstwordof#2 \relax}%
      \edef\@deflblkey{%
        \expandafter\@takethree\@deffirstword\@empty\@empty\@empty\@nil}%
    }%
    \def\@tempshort{#3}%
    \ifx\@tempshort\@empty
      \expandafter\xdef\csname @dshort@\@deflblkey\endcsname{\@deflblkey}%
    \else
      \expandafter\def\csname @dshort@\@deflblkey\endcsname{#3}%
    \fi
    \IfBooleanTF{#1}{\let\@deftag\@empty}{%
      \edef\@deftag{\noexpand\ \noexpand\textup{[\noexpand\texttt{d:\@deflblkey}]}}}%
    \begin{definition}[#2\ifshowkeys\@deftag\fi]%
    \label{d:\@deflblkey}%
    % record the long title so \nref{d:KEY} prints it, as for theorems
    \protected@write\@auxout{}{%
      \string\@thmtitle{d:\@deflblkey}{\unexpanded{#2}}}%
    \expanded{%
      \noexpand\newglossaryentry{d\thedefcounter}{%
        name={\unexpanded{#2}%
          \IfBooleanF{#1}{ [\noexpand\texttt{d:\@deflblkey}]}},%
        sort={\thedefinition},%
        description={p.~\noexpand\pageref{d:\@deflblkey}},%
        user1={\thedefinition},%
        user2={Definition},%
        user3={\thesection},%
        user4={\expandafter\unexpanded\expandafter{\@savedSecTitle}}%
      }%
    }%
    \protected@write\@auxout{}{\string\glsliveentry{d\thedefcounter}}%
  \fi
}
\newcommand{\eeD}{\end{definition}}

% \showkeysfalse in a driver hides the [t:KEY] mnemonics in result headings
% (labels and glossary entries are unchanged); used for submission builds.
\newif\ifshowkeys \showkeystrue
% \commonredtrue in a driver that inputs basicFiles/commonRed: tool.tex then
% cites its reduction results instead of restating them.
\newif\ifcommonred \commonredfalse
% TheorEM COMMAND
% Usage: \beT[long title][short title][key]
%   long title  = theorem heading (plain text only)
%   short title = shown in \tref{word}; may contain math; defaults to first word of long title
%   key         = optional label key override; label becomes t:KEY (avoids first-word collisions)
% Label is t:WORD where WORD = first word of long title (or KEY if given); use \tref{WORD} to reference.
\NewDocumentCommand{\beT}{s O{} O{} o}{%
  \def\@temparg{#2}%
  \ifx\@temparg\@empty
    \begin{theorem}%
  \else
    \let\@savedSecTitle\@currentsectionname%
    \refstepcounter{thmcounter}%
    % \@thmtag is the mnemonic shown in the heading; every theorem gets one.
    % With an explicit key the key itself is shown.  Without one the key
    % defaults to the title's first word, and showing that whole word would
    % merely repeat the title ("Dissipativity [Dissipativity]"), so only its
    % first three letters are shown: "Dissipativity [Dis]".  The label is
    % unaffected -- it still uses the full \@thmlblkey.
    \IfValueTF{#4}{%
      \edef\@thmlblkey{#4}%
    }{%
      \edef\@thmfirstword{\@firstwordof#2 \relax}%
      \edef\@thmlblkey{%
        \expandafter\@takethree\@thmfirstword\@empty\@empty\@empty\@nil}%
    }%
    \IfBooleanTF{#1}{\let\@thmtag\@empty}{%
      \edef\@thmtag{\noexpand\ \noexpand\textup{[\noexpand\texttt{t:\@thmlblkey}]}}}%
    \def\@tempshort{#3}%
    \ifx\@tempshort\@empty
      \expandafter\xdef\csname @tshort@\@thmlblkey\endcsname{\@thmlblkey}%
    \else
      \expandafter\def\csname @tshort@\@thmlblkey\endcsname{#3}%
    \fi
    \begin{theorem}[#2\ifshowkeys\@thmtag\fi]%
    \label{t:\@thmlblkey}%
    % record the long title for \nref; via the .aux so forward refs resolve
    \protected@write\@auxout{}{%
      \string\@thmtitle{t:\@thmlblkey}{\unexpanded{#2}}}%
    \expanded{%
      \noexpand\newglossaryentry{t\thethmcounter}{%
        name={\unexpanded{#2}%
          \IfBooleanF{#1}{ [\noexpand\texttt{t:\@thmlblkey}]}},%
        % sort/user1 must be the number the document actually prints, i.e. the
        % theorem environment's counter, not thmcounter.  thmcounter advances
        % only on \beT with a title, so a bare \beT (or \beT with an
        % empty title) makes the two drift and the glossary cites a wrong
        % number.  thmcounter is still fine as the entry key: it is unique.
        sort={\thetheorem},%
        description={p.~\noexpand\pageref{t:\@thmlblkey}},%
        user1={\thetheorem},%
        user2={Theorem},%
        user3={\thesection},%
        user4={\expandafter\unexpanded\expandafter{\@savedSecTitle}}%
      }%
    }%
    \protected@write\@auxout{}{\string\glsliveentry{t\thethmcounter}}%
  \fi
}
\newcommand{\eeT}{\end{theorem}}

% CROSS-PART RELATION PARAGRAPH
% Usage: \xpart{text}  -- place immediately BEFORE \beT, never inside it
%   (anything inside \beT...\eeT would be swallowed into the glossary entry).
% The paragraph is typeset only when Part I is actually present, i.e. in the
% combined book build; in the standalone builds (GranadaDSABNS/A.tex,
% AAHRah-Submission/AAHRah.tex) the label part:glv is undefined and the whole
% paragraph disappears rather than leaving ?? cross-references. Same
% graceful-degradation principle as \GASref in GASt.tex. Resolves on the
% second LaTeX pass, since r@part:glv comes from the .aux file.
% (Already inside the \makeatletter region opened above, so no \makeatletter
% pair here -- adding one would close that region early.)
\newcommand{\xpart}[1]{\@ifundefined{r@part:glv}{}{\paragraph{Relation with Parts~I--II.}#1}}

% \bookref{label}{precise text}{fallback text}: prints the precise cross-reference
% when the label resolves (combined book build), and the fallback text otherwise
% (standalone paper builds, where book-only appendix labels are absent). Same
% graceful-degradation role as \GASref in GASt.tex, but usable in any file since it
% lives in this preamble (loaded by every driver) -- unlike \GASref, which is defined in a
% content file and is therefore undefined in standalone builds that omit that file.
\newcommand{\bookref}[3]{\@ifundefined{r@#1}{#3}{#2}}

% \PVpart and \RBpart: the two parts that a partial build may omit.
% They print "Part~N" in any build that contains the part, and a naming phrase
% in a build that does not, so an excerpt shows no ?? .  Each already supplies
% the word "Part", so write "\PVpart develops ...", not "Part~\PVpart".
\newcommand{\PVpart}{\bookref{part:aahrah}%
  {Part~\ref{part:aahrah}}{the Perron--Volterra part of the full book}}
\newcommand{\RBpart}{\bookref{part:recbrowne}%
  {Part~\ref{part:recbrowne}}{the immuno-virology part of the full book}}

% OPEN PROBLEM COMMANDS
% ASSUMPTION COMMANDS (glossaried)
\newcommand{\beAs}[1][]{%
  \def\@temparg{#1}%
  \ifx\@temparg\@empty
    \begin{assumption}%
  \else
    \let\@savedSecTitle\@currentsectionname%
    \refstepcounter{ascounter}%
    \begin{assumption}[#1]%
    \label{as:\theascounter}%
    \expanded{%
      \noexpand\newglossaryentry{as\theascounter}{%
        name={\unexpanded{#1}},%
        sort={\theascounter},%
        description={p.~\noexpand\pageref{as:\theascounter}},%
        user1={\theascounter},%
        user2={Assumption},%
        user3={\thesection},%
        user4={\expandafter\unexpanded\expandafter{\@savedSecTitle}}%
      }%
    }%
    \protected@write\@auxout{}{\string\glsliveentry{as\theascounter}}%
  \fi
}
\newcommand{\eeAs}{\end{assumption}}
% CONJECTURE COMMANDS (glossaried)
\newcommand{\beCj}[1][]{%
  \def\@temparg{#1}%
  \ifx\@temparg\@empty
    \begin{conjecture}%
  \else
    \let\@savedSecTitle\@currentsectionname%
    \refstepcounter{cjcounter}%
    \begin{conjecture}[#1]%
    \label{cj:\thecjcounter}%
    \expanded{%
      \noexpand\newglossaryentry{cj\thecjcounter}{%
        name={\unexpanded{#1}},%
        sort={\thecjcounter},%
        description={p.~\noexpand\pageref{cj:\thecjcounter}},%
        user1={\thecjcounter},%
        user2={Conjecture},%
        user3={\thesection},%
        user4={\expandafter\unexpanded\expandafter{\@savedSecTitle}}%
      }%
    }%
    \protected@write\@auxout{}{\string\glsliveentry{cj\thecjcounter}}%
  \fi
}
\newcommand{\eeCj}{\end{conjecture}}
% LEMMA COMMANDS
% ---------------------------------------------------------------------
% UNIFORM RESULT WRAPPERS -- \beL, \beC, \beP, \beR share \beT's signature
%
%   \beX[long title][short title][key]  ...  \eeX
%
%   long title  = heading (math allowed)
%   short title = shown in the cross-reference; defaults to the key
%   key         = label key.  With a key the label is X:KEY and the heading
%                 carries the mnemonic [key], exactly as for \beT.  Without
%                 one the label falls back to the running counter, which is
%                 what every pre-existing call site in the book relies on.
%
% A bare \beX -- arguments commented out -- opens the plain environment and
% creates NO label and NO glossary entry.  That is the supported way to keep
% an individual result out of the catalogue.
% ---------------------------------------------------------------------
% #1 wrapper name  #2 environment  #3 label prefix  #4 counter  #5 type name
\newcommand{\@mkresult}[5]{%
  % A leading * suppresses the [key] mnemonic in both the heading and the
  % glossary name, while still making the label X:KEY and the entry:
  %     \beX*[long][short][key]
  \expandafter\NewDocumentCommand\csname be#1\endcsname{s O{} O{} o}{%
    \def\@temparg{##2}%
    \ifx\@temparg\@empty
      \begin{#2}%
    \else
      \let\@savedSecTitle\@currentsectionname%
      \refstepcounter{#4}%
      % With a key: mnemonic label X:KEY and the [key] tag in the heading.
      % Without one: fall back to the running counter, exactly as these
      % wrappers behaved before, so every legacy one-argument call site in
      % the book keeps its label.  The counter is also csname-safe, which a
      % title's first word is not -- titles beginning with math would
      % otherwise build \csname \Omega...\endcsname and fail.
      \IfValueTF{##4}{%
        \edef\@relbl{##4}%
        \IfBooleanTF{##1}{\let\@retag\@empty}{%
          \edef\@retag{\noexpand\ \noexpand\textup{[\noexpand\texttt{#3:##4}]}}}%
      }{%
        \edef\@relbl{\arabic{#4}}%
        \let\@retag\@empty
      }%
      \def\@tempshort{##3}%
      \ifx\@tempshort\@empty
        \expandafter\xdef\csname @sshort@#3\@relbl\endcsname{\@relbl}%
      \else
        \expandafter\def\csname @sshort@#3\@relbl\endcsname{##3}%
      \fi
      \begin{#2}[##2\ifshowkeys\@retag\fi]%
      \edef\@tmplbl{#3:\@relbl}\expandafter\label\expandafter{\@tmplbl}%
      \expanded{%
        \noexpand\newglossaryentry{#3\csname the#4\endcsname}{%
          name={\unexpanded{##2}%
            \IfBooleanF{##1}{\IfValueT{##4}{ [\noexpand\texttt{#3:##4}]}}},%
          sort={\csname the#2\endcsname},%
          description={p.~\noexpand\pageref{#3:\@relbl}},%
          user1={\csname the#2\endcsname},%
          user2={#5},%
          user3={\thesection},%
          user4={\expandafter\unexpanded\expandafter{\@savedSecTitle}}%
        }%
      }%
      \protected@write\@auxout{}{\string\glsliveentry{#3\csname the#4\endcsname}}%
      \protected@write\@auxout{}{\string\@thmtitle{\@tmplbl}{\unexpanded{##2}}}%
    \fi
  }%
  \expandafter\newcommand\csname ee#1\endcsname{\end{#2}}%
}
\@mkresult{L}{lemma}{l}{lemcounter}{Lemma}
% REMARK / COROLLARY / PROPOSITION -- same signature as \beL above
\@mkresult{R}{remark}{r}{remcounter}{Remark}
\@mkresult{C}{corollary}{c}{corcounter}{Corollary}
\@mkresult{P}{proposition}{p}{propcounter}{Proposition}
% OPEN-PROBLEM COMMAND -- generated with the same signature as \beT:
%   \beO[long title][short title][key]   ...   \eeO
% label o:KEY, or o:<counter> when no key is given; a bare \beO makes
% neither a label nor a glossary entry.
\@mkresult{O}{open}{o}{opcounter}{Problem}


% Section-level title tracker: updated only by \section, not \subsection.
% \@currentlabelname is overwritten by every \refstepcounter (including
% \subsection), so we save it separately here.
\newcommand{\@currentsectionname}{}
\AtBeginDocument{%
  \let\@epidcrn@orig@refstepcounter\refstepcounter
  \renewcommand{\refstepcounter}[1]{%
    \@epidcrn@orig@refstepcounter{#1}%
    \ifstrequal{#1}{section}{\global\let\@currentsectionname\@currentlabelname}{}%
  }%
}

% TABLE COMMANDS (glossaried) -- consistent naming with \beT, \beD, etc.
\newcommand{\beTa}[1][]{%
  \def\@temparg{#1}%
  \ifx\@temparg\@empty
    \begin{table}[htbp]%
  \else
    \let\@savedSecTitle\@currentsectionname%
    \refstepcounter{tabcounter}%
    \begin{table}[htbp]%
    \label{A:\thetabcounter}%
    \expanded{%
      \noexpand\newglossaryentry{A\thetabcounter}{%
        name={\unexpanded{#1}},%
        sort={\thetabcounter},%
        description={p.~\noexpand\pageref{A:\thetabcounter}},%
        user1={\thetabcounter},%
        user2={Table},%
        user3={\thesection},%
        user4={\expandafter\unexpanded\expandafter{\@savedSecTitle}}%
      }%
    }%
    \protected@write\@auxout{}{\string\glsliveentry{A\thetabcounter}}%
  \fi
}
\newcommand{\eeTa}{\end{table}}

% FIGURE COMMANDS (glossaried) -- consistent with \beT, \beR, \beTa, etc.
% Usage:
%   \beF[Glossary title]...\caption{Visible caption}\eeF  -> glossaried figure
%   \beF...\caption{Visible caption}\eeF                  -> plain [H] figure, no glossary
\newcommand{\beF}[1][]{%
  \def\@temparg{#1}%
  \ifx\@temparg\@empty
    \begin{figure}[H]%
  \else
    \begin{figure}[H]%
    \let\@savedSecTitle\@currentsectionname%
    \refstepcounter{figcounter}%
    \label{fig:\thefigcounter}%
    \expanded{%
      \noexpand\newglossaryentry{Fig\thefigcounter}{%
        name={\unexpanded{#1}},%
        sort={\thefigcounter},%
        description={p.~\noexpand\pageref{fig:\thefigcounter}},%
        user1={\thefigcounter},%
        user2={Figure},%
        user3={\thesection},%
        user4={\expandafter\unexpanded\expandafter{\@savedSecTitle}}%
      }%
    }%
    \protected@write\@auxout{}{\string\glsliveentry{Fig\thefigcounter}}%
  \fi
}
\newcommand{\eeF}{\end{figure}}
\makeatother
% LIVE-ENTRY REGISTER
% \glsaddall prints every entry the glossaries package currently knows, and
% that set includes entries restored from \jobname.glsdefs.  The package
% rewrites .glsdefs from all *currently defined* entries, so anything it
% restored is written straight back out: a deleted section kept reappearing in
% the catalogue for ever, and the .glsdefs had to be erased by hand.  The
% wrappers therefore record every entry they actually create in the .aux, and
% \printboth adds only those.  The .aux is rebuilt from scratch on each run, so
% this is self-healing: a removed result is gone after one further run.
\makeatletter
\newcommand{\@glslive}{}
\newcommand{\@glsaddlive}[1]{\ifglsentryexists{#1}{\glsadd{#1}}{}}
\newcommand{\glsliveentry}[1]{\g@addto@macro\@glslive{\@glsaddlive{#1}}}
\makeatother
% PRINT COMMAND
% No \clearpage here -- the catalogue continues on the current page.  This has
% been re-added by mistake several times; leave it out.
\makeatletter
\newcommand{\printboth}{%
  \@glslive   % add exactly the entries this document created last run
  \printglossary[style=unifiedstyle, title={0 \hspace{.22em} Definitions, Theorems, Examples and Open Problems}]%
}
\makeatother
% Named theorem reference: \tref{word} -> Thm~N(word-or-short)
% Shows stored short title if available, otherwise shows the key itself.
\newcommand{\tref}[1]{%
  Thm~\ref{t:#1}(\ifcsname @tshort@#1\endcsname\csname @tshort@#1\endcsname\else#1\fi)%
}

% Full-title theorem reference: \nref{t:KEY} prints the number and then the
% theorem's long title, so one writes "Theorem~\nref{t:blocks}" and gets
% "Theorem 3 ``Block factorization and siphons''".  Takes the complete label, like
% \ref, not the bare key like \tref.
%
% \beT writes each title to the .aux (see \@thmtitle below), so \nref resolves
% forward references from the second LaTeX pass, exactly as \ref does.  If the
% title is not yet known -- first pass, or a label that is not a \beT theorem
% -- it degrades to the number rather than producing ??.
% These two use @ in their names, and this point in the file is outside the
% \makeatletter region opened earlier (it closes at \makeatother above), so
% the region has to be reopened here.
\makeatletter
\newcommand{\@thmtitle}[2]{\global\@namedef{@tlong@#1}{#2}}
\newcommand{\nref}[1]{%
  \ref{#1}%
  \ifcsname @tlong@#1\endcsname
    \ ``\csname @tlong@#1\endcsname''%
  \fi
}
% \nm{t:KEY}: the stored long title of any \beT/\beL/\beP/\beR/\beC/\beD/\beE result,
% verbatim and without quotes, so item names follow the macro's first argument.
\newcommand{\nm}[1]{%
  \ifcsname @tlong@#1\endcsname
    \csname @tlong@#1\endcsname%
  \else
    \ref{#1}%
  \fi
}
\makeatother




%% ========================================
%% Other THEOREM ENVIRONMENTS
%% ========================================
\makeatletter
\@ifundefined{lemma}      {\newtheorem{lemma}      {Lemma}}      {}
\@ifundefined{proposition}{\newtheorem{proposition}{Proposition}} {}
\@ifundefined{corollary}  {\newtheorem{corollary}  {Corollary}}  {}
\@ifundefined{conjecture} {\newtheorem{conjecture} {Conjecture}}  {}
\@ifundefined{remark}     {\newtheorem{remark}     {Remark}}     {}
\@ifundefined{question}   {\newtheorem{question}   {Question}}   {}
\@ifundefined{assumption} {\newtheorem{assumption} {Assumption}} {}
\@ifundefined{hypothesis} {\newtheorem{hypothesis} {Hypothesis}} {}
\@ifundefined{fact}       {\newtheorem{fact}       {Fact}}       {}
\makeatother

% \citer: cite the relay-graph threshold theorem of AAHrel. If a document has
% \input the chapter/file containing t:relayRF (so the label resolves), point
% straight at it; otherwise fall back to a bare citation of AAHrel. Shared here
% (not per-paper) so it resolves consistently across every book chapter.
% \citeR: same idea, for the first theorem of AAHRah (t:Pstar).
\makeatletter
\providecommand{\citer}{%
  \@ifundefined{r@t:relayRF}%
    {\cite{AAHrel}}%
    {Theorem~\ref{t:relayRF} of \cite{AAHrel}}%
}
\providecommand{\citeR}{%
  \@ifundefined{r@t:Pstar}%
    {\cite{AAHRah}}%
    {Theorem~\ref{t:Pstar} of \cite{AAHRah}}%
}
\makeatother


% Abbreviations
\def\beQ{\begin{question}}\def\eeQ{\end{question}}
\def\beA{\begin{assumption}}\def\eeA{\end{assumption}}


%% ========================================

%% ========================================
%% COLORS & CUSTOM COMMANDS
%% ========================================
\definecolor{funccolor}{RGB}{25,25,112}
\definecolor{desccolor}{RGB}{64,64,64}
\newcommand{\func}[1]{\texttt{\textcolor{funccolor}{#1}}}
\def\eq{=}
\def\mm{-}

%% ========================================
%% MATHCAL SHORTHANDS
%% ========================================
% Project-wide convention: \mX for a capital X is always \mathcal X.
% Defined for the whole alphabet so the convention never fails to expand.
\newcommand{\mA}{\mathcal A}\newcommand{\mB}{\mathcal B}
\newcommand{\mC}{\mathcal C}\newcommand{\mD}{\mathcal D}
\newcommand{\mE}{\mathcal E}\newcommand{\mF}{\mathcal F}
\newcommand{\mG}{\mathcal G}\newcommand{\mH}{\mathcal H}
\newcommand{\mI}{\mathcal I}\newcommand{\mJ}{\mathcal J}
\newcommand{\mK}{\mathcal K}\newcommand{\mL}{\mathcal L}
\newcommand{\mM}{\mathcal M}\newcommand{\mN}{\mathcal N}
\newcommand{\mO}{\mathcal O}\newcommand{\mP}{\mathcal P}
\newcommand{\mQ}{\mathcal Q}\newcommand{\mR}{\mathcal R}
\newcommand{\mS}{\mathcal S}\newcommand{\mT}{\mathcal T}
\newcommand{\mU}{\mathcal U}\newcommand{\mV}{\mathcal V}
\newcommand{\mW}{\mathcal W}\newcommand{\mX}{\mathcal X}
\newcommand{\mY}{\mathcal Y}\newcommand{\mZ}{\mathcal Z}

%% ========================================
%% BASIC MATH ENVIRONMENTS
%% ========================================
\def\bep{\begin{pmatrix}}\def\eep{\end{pmatrix}}
\def\bev{\begin{vmatrix}}\def\eev{\end{vmatrix}}
\def\bea{\begin{eqnarray*}}\def\eea{\end{eqnarray*}}
\def\bc{\begin{cases}}\def\ec{\end{cases}}
\def\BEN{\begin{enumerate}}\def\EEN{\end{enumerate}}
\def\BI{\begin{itemize}}\def\EI{\end{itemize}}

\newcommand{\be}[1]{\begin{equation}\label{#1}}
\newcommand{\ee}{\end{equation}}
\newcommand{\beq}{\begin{eqnarray}}
\def\eeq{\end{eqnarray}}
\def\eqr{\eqref}\def\fr{\frac}\def\lbl{\label}
\def\Lra{\Longrightarrow}\def\non{\nonumber}\def\cNGM{\cite{Diek,Van,Van08}}
\def\Eq{\Leftrightarrow}\def\b{\beta}\def\uph{\red{UP TO HERE}}
%% ========================================
%% PROBABILITY & STATISTICS
%% ========================================
\renewcommand{\P}{\ensuremath{\mathbb{P}}}  % Note: overrides paragraph symbol
\newcommand{\E}{\ensuremath{\mathbb{E}}}
\newcommand{\ind}[1]{\ensuremath{\mathbbm{1}_{\left\{#1\right\}}}}
\DeclareMathOperator{\Cov}{Cov}
\DeclareMathOperator{\TV}{TV}
\newcommand{\rank}{\operatorname{rank}}

%% ========================================
%% CRN OPERATORS (organizations/ERCs/buffering/siphons framework)
%% ========================================
\def\clos{\mathrm{clos}}
\def\Supply{\mathrm{Supply}}
\def\req{\mathrm{req}}
\def\React{\mathrm{React}}
\def\Prod{\mathrm{Prod}}
\def\Gen{\mathrm{Gen}}
\def\In{\mathrm{In}}
\def\Out{\mathrm{Out}}
\def\Act{\mathrm{Act}}
\def\Min{\mathrm{Min}}

%% ========================================
%% LINE BREAKING
%% ========================================
% Last-resort stretch, used only on paragraphs TeX cannot otherwise break:
% a long \texttt{} token, a comma-separated inline math list (commas are not
% breakpoints in math), or a run of \cite/\ref boxes.  Without it those set as
% overfull \hbox and print into the margin; with it they set slightly loose,
% which is the lesser evil.  Does not affect paragraphs that already break.
\emergencystretch=3em

%% ========================================
%% NUMBER SETS
%% ========================================
\newcommand{\R}{\mathbb{R}}  % Real numbers
\newcommand{\C}{\mathbb{C}}  % Complex numbers
\newcommand{\N}{\mathbb{N}}
\newcommand{\Z}{\mathbb{Z}}
% \providecommand, not \newcommand: some papers define these in their own
% preamble after \input{preamble}, and must not clash.
\providecommand{\Q}{\mathbb{Q}}
\providecommand{\res}{\operatorname{res}}   % resultant; used by epiRes.tex
\newcommand{\MN}{\mathbb{N}}
\newcommand{\MZ}{\mathbb{Z}}
\newcommand{\realpos}{\mathbb{R}_{> 0}}
\newcommand{\realnneg}{\mathbb{R}_{\geq 0}}
\newcommand{\naturalpos}{\mathbb{N}_{>0}}
\newcommand{\Rplus}{\mathbb{R}_+}
\newcommand{\Rp}{\mathbb{R}_{\ge 0}}
\def\Rnne{\mathbb{R}_{\ge 0}}
\def\Se{\bar{S}} \def\Ie{\bar{I}}
\newcommand{\Rpos}{\mathbb{R}_{>0}}
\newcommand{\dd}{\,\mathrm d}
\newcommand{\e}{\mathrm e}
%% ========================================
%% GREEK LETTERS (2+ character abbreviations) greek
%% ========================================
\def\al{\alpha}\def\ba{\bff \alpha}
\def\w{\omega}\def\Lgn{\ga+\mu+\de}\def\om{\omega}
\def\de{\delta}\def\De{\Delta}\def\bde{\bff \delta}
\def\ep{\epsilon}\def\eps{\varepsilon}
\def\f{\varphi}  % Kept as requested
\def\ga{\gamma}\def\Ga{\Gamma}\def\G{\Gamma}\def\bg{\bff \gamma}
\def\up{\upsilon}  % NOT the breve accent: no accent commands anywhere.
                   % Named \up, not \u, so the LaTeX breve stays reachable.
\def\la{\lambda}\def\La{\Lambda}
\def\si{\sigma}\def\ka{\kappa}
\def\th{\theta}\def\vt{\vartheta}
\def\ze{\zeta}\def\Gap{\Gamma^{+}} \def\Gam{\Gamma^{-}}


%% ========================================
%% SPECIAL SHORT COMMANDS (kept as requested)
%% ========================================
\def\I{\infty}
\def\T{\widetilde}
\def\H{\hat}
\newcommand{\relint}{\operatorname{relint}}

%% ========================================
%% BOLD/VECTOR NOTATION
%% ========================================
\newcommand{\bff}[1]{{\mbox{\boldmath$#1$}}}
\newcommand\X{\boldsymbol{X}}
\def\x{\boldsymbol{x}}\def\bx{\bff x}\def\xe{\bar{\x}}
\newcommand\y{\boldsymbol{y}}\def\by{\boldsymbol{y}}\def\ye{\bar{\y}}
\def\yb{\y}\def\xb{\x}
\def\xy{\x-\y}
\newcommand\br{\boldsymbol{r}}\def\re{\bar{\r}}
\newcommand\bd{\boldsymbol{d}}
\newcommand\bQ{\boldsymbol{Q}}
\newcommand{\Yb}{{\mathbf{Y}}}
\newcommand{\Ab}{\mathbf{A}}
\newcommand{\zb}{\mathbf{z}}\def\bz{\bff 0}
\newcommand{\Zb}{\mathbf{Z}}
\newcommand{\bb}{\mathbf{b}}
\def\vbe{\vec {\bff \beta}}
\def\vb{\vec {\bff b}}
\def\v1{\vec {\bff 1}}
\def\vn{\vec \de}
\def\vz{\vec 0}
\def\vi{\overset{\rightarrow}{\mathsf i}}
\def\vv{\overset{\rightarrow}{v}}
\def\vr{\; \vec {\mathsf r}}
\newcommand{\va}{\bff a}
\providecommand{\vect}[1]{\boldsymbol{#1}}   % bold vector (moved from Gradostat.tex)
\providecommand{\one}{\mathbf 1}             % all-ones vector (moved from AliSF)
\providecommand{\spr}{s}                     % spectral bound  (moved from AliSF)
\providecommand{\Ke}{K_{\rm eff}}            % effective carrying capacity
% Critical right/left eigenvectors in the Hopf projection method (AAHOSN/New.tex)
\providecommand{\bq}{\mathbf{q}}
\providecommand{\bp}{\mathbf{p}}


%% ========================================
%% MATHCAL & MATHSCR
%% ========================================
% Mathcal
\newcommand\mc[1]{\mathcal{#1}}
\def\mC{{\mathcal C}}\def\mD{\mathcal D}\def\mE{\mathcal{E}}
\def\mF{\mathcal F}\def\mG{{\mathcal G}}\def\mH{{\mathcal H}}
\def\mI{\mathcal I}\def\mL{\mathcal{L}}\def\mM{{\mathcal M}}
\def\mN{{\mathcal N}}\def\mP{\mathcal{P}}\def\mR{{\mathcal R}}
\def\mS{{\mathcal S}}\def\mT{\mathcal{T}}\def\mV{{\mathcal V}}
\def\mK{\mathcal K}\def\mA{{\mathcal A}}\def\mU{{\mathcal U}}
% Symmetric bilinear and trilinear forms of the Taylor expansion, used in the
% first Lyapunov coefficient (AAHOSN/New.tex)
\providecommand{\Bform}{\mathcal{B}}
\providecommand{\Cform}{\mathcal{C}}

% Blackboard bold.  Kept here, not in the drivers: a \providecommand in a
% driver preamble runs before this file and silently wins, so the same name
% ends up meaning different things in different builds.
\def\IC{\mathbb{C}}\def\IN{\mathbb{N}}\def\IQ{\mathbb{Q}}
\def\IR{\mathbb{R}}\def\IZ{\mathbb{Z}}\def\IG{\mathbb{G}}
\def\II{\mathbb{I}}

% Mathscr
\newcommand{\Rs}{{\mathscr R}}
\newcommand{\Ss}{{\mathscr S}}
\newcommand{\Ns}{{\mathscr N}}
\newcommand{\Cs}{\mathscr C}


%% ========================================
%% OPERATORS & FUNCTIONS
%% ========================================
\DeclareMathOperator{\tr}{tr}
\DeclareMathOperator{\diag}{Diag}
\DeclareMathOperator{\Diag}{Diag}
\DeclareMathOperator{\intr}{\mathrm{Int}}
\DeclareMathOperator{\Inv}{Inv}
\DeclareMathOperator{\spec}{spec}
\DeclareMathOperator{\AUC}{AUC}               % area under the curve
\newcommand{\ii}{\mathrm{i}}                  % imaginary unit
\newcommand\op{\operatorname}
\newcommand{\supp}{\op{supp}}
\newcommand{\diff}{\mathop{}\mathopen{}\mathrm{d}}
\newcommand{\pdd}[1]{\frac{\partial}{\partial{#1}}}
\DeclareMathOperator{\cone}{cone}
\DeclareMathOperator{\proj}{proj}

%% ========================================
%% BRACKETS & NORMS
%% ========================================
\providecommand{\pr}[1]{\left(#1\right)}
\providecommand{\pp}[1]{\left[#1\right]}
\newcommand\croc[1]{\left\langle #1\right\rangle}
\newcommand{\norm}[1]{\|#1\|}
\newcommand{\inner}[2]{\langle #1, #2 \rangle}
% Short synonym of \inner, used by AAHOSN/New.tex as $\ip{u}{v}=\bar u^\top v$
\providecommand{\ip}[2]{\langle #1 , #2 \rangle}


%% ========================================
%% FORMATTING & UTILITIES
%% ========================================
\newcommand*{\Scale}[2][4]{\scalebox{#1}{$#2$}}
\newcommand{\te}[1]{\text{#1}}
\newcommand\steq[1]{\stackrel{\text{\rm #1.}}{=}}
\newcommand\blue[1]{\textcolor{blue}{#1}}
\newcommand{\red}{\textcolor[rgb]{1.00,0.00,0.00}}
\newcommand{\green}{\textcolor[rgb]{0.00,1.00,0.00}}

% Figure command
\newcommand{\figu}[3]{%
\begin{figure}[H]
\centering
\includegraphics[scale=#3]{#1}
\caption{#2\label{f:#1}}
\end{figure}
}

% Footnote
\long\def\symbolfootnote[#1]#2{%
\begingroup
\def\thefootnote{\fnsymbol{footnote}}\footnote[#1]{#2}%
\endgroup}
\def\fn{\symbolfootnote}

\newcommand{\dsum}{\displaystyle\sum}

%% ========================================
%% CRN & EPIDEMIOLOGY SPECIFIC
%% ========================================
% CRN notation
\def\SS{\mathcal S}
\def\CC{\mathcal C}
\def\RR{\mathcal R}

% Reproduction numbers
\def\Rto{R_{1}^{E_2}}
\def\Rot{R_{2}^{E_1}}
\def\Roc{R_{1,c}}
\def\Rtc{R_{2,c}}
\def\rd{r_{0}}
\def\sd{s_{0}}
\def\dd{D_{0}}
\def\vd{v_{0}}
\def\zd{z_{0}}
\def\se{\sE}
\def\xd{x_{0}}
%% ========================================
%% SPECIAL VALUES & EQUILIBRIA
%% ========================================
\def\eE{\bar{e}}
\def\wE{\bar{w}}
\def\sE{s_{*}}
\def\Me{\T M_{ee}}
\def\Rnne{\mathbb{R}_{\ge 0}}

\newcommand{\ones}{\mathbf 1}


%% ========================================
%% ABBREVIATIONS (alphabetical by category)
%% ========================================

% -- General mathematical terms --
\def\adm{admissible}\def\admy{admissibility}
%\def\and{antisymmetric}% disabled 2026-09-30: it overrode LaTeX's author separator \and (printed "antisymmetric" between authors)
\def\ass{assumption}
\def\cadlag{c\`adl\`ag}
\def\ch{characteristic polynomial}
\def\chp{characteristic polynomial}
\def\coe{coefficient}
\def\comp{computation}
\def\con{condition}
\def\coo{coordinate}
\def\corr{corresponding}
\def\dec{decomposition}
\def\den{denominator}
\def\dete{deterministic}
\def\eig{eigenvalue}\def\eval{eigenvalue}
\def\eigv{eigenvector}\def\eiv{eigenvector}\def\evec{eigenvector}
\def\elt{element}
\def\enn{essentially nonnegative/positivity preserving}
\def\equ{equation}
\def\etal{et al.}
\def\fe{for example}
\def\fno{from now on}
\def\for{\forall}
\def\form{formula}
\def\Fr{Furthermore, }
\def\frt{furthermore}
\def\fun{function}
\def\gene{generalization}
\def\heu{heuristic}
\def\how{however}\def\How{However, }
\def\ie{i.e. }
\def\im{\item}
\def\isd{is defined by}
\def\itc{it is easy to check that}
\def\itf{it follows that}\def\Itf{It follows that}
\def\ith{it holds that}
\def\itm{it may be checked that}\def\Itm{It may be checked that}
\def\ito{it turns out that}\def\Ito{It turns out that}\def\Itou{It turns out useful to}
\def\ma{mass action}\def\MAR{mass action representation}
\def\mat{matrix}
\def\mbc{may be checked to be}
\def\mbw{may be written as}\def\mubw{must be written as}
\def\mpr{more precisely, }\def\Mp{More precisely, }
\def\nc{non-catalytic}
\def\nne{non-negative}\def\nny{nonnegativity}
\def\no{\nonumber}
\def\noni{non-infectious}
\def\nota{notation}
\def\np{\newpage}
\def\Nt{Note that}
\def\nt{non trivial}\def\nz{non zero}
\def\oth{otherwise}
\def\Oth{On the other hand}
\def\pa{parameter}\def\para{parameter}
\def\pol{polynomial}
\def\prob{probability}
\def\Prf{{\bf  Proof:}}\def\prf{{\bf  Proof:}}
\def\prop{proportional to}
\def\qu{\quad}
\def\QED{\hfill {$\square$}\goodbreak \medskip}
\def\rep{representation}
\def\resp{respectively}
\def\satd{satisfied}\def\satg{satisfying}\def\sats{satisfies}\def\saty{satisfy}
\def\ssec{\subsection}\def\sssec{\subsubsection}
\def\sef{see for example}
\def\sksy{skew-symmetric}
\def\sub{substrate}
\def\sys{system}
\def\var{variable}
\def\wk{well-known}
\def\wl{without loss of generality}\def\wlo{w.l.o.g.}
\def\wmc{We may check that}\def\Wmc{We may check that}
\def\wmr{we may rewrite}
\def\wms{we must show that}
\def\wrt{with respect to}

% -- CRN & Chemistry --
\newcommand\CRN{chemical reaction network}
\newcommand\CRNs{chemical reaction networks}
\newcommand\RN{reaction network}
\def\ACR{absolute concentration robustness}
\def\AMSD{acyclic siphon decomposition}
\def\DAE{differential algebraic equations}
\def\FHJ{Feinberg-Horn-Jackson}
\def\GMAK{generalized mass action kinetics}
\def\MM{Michaelis-Menten}
\def\sm{stoichiometric matrix}\def\Sm{stoichiometric matrix}
\def\WRZD{weakly reversible, zero deficiency}\def\aut{autocatalytic}

% -- Epidemiology --
\def\brn{basic reproduction number}\def\bzn{basic replacement number}
\def\com{compartment}
\def\dfe{disease free equilibrium}\def\DFE{disease free equilibrium}
\def\dwt{dwell times}\def\edt{expected dwell times}
\def\EM{epidemic model}
\def\endp{endemic fixed point}\def\enp{endemic point}\def\EE{endemic point}
\def\ME{mathematical epidemiology}
\def\MEE{eco-epidemiology}\def\EEP{eco-epidemiological}
\def\ece{eco-epidemiological model}
\def\mgm{multi group model}
\def\mtS{multi-strain model}\def\mSt{multi-strain model}
\def\mSi{minimal siphon}
\def\NGM{next generation matrix}
\def\nim{new infections matrix}
\def\qem{quadratic epidemiological model}
\def\repF{invasion functions}
\def\SPF{SIR-PH-FA model}
\def\sus{susceptible}

% -- Graph theory & Linear algebra --
\def\jac{Jacobi matrix}
\def\jin{Jacobian matrix of the invasion vector field with respect to the invasion variables}
\def\jrn{Jacobian matrix of the resident vector field with respect to the resident variables}
\def\jxn{Jacobian matrix of the $\x$ vector field with respect to the $\x$ variables}
\def\Lm{Laplacian matrix}
\def\MTT{matrix tree theorem}
\def\SGT{spectral graph theory}
\def\SMW{Sherman–Morrison–Woodbury}

% -- Stability & Dynamics --
\def\absC{absorbtion CTMC}
\def\BTb{Bogdanov-Takens bifurcation}
\def\CEP{competitive exclusion partition}
\def\CTMC{continuous time Markov chain}
\def\ds{differential system}
\def\fp{fixed point}\def\bfp{boundary fixed point}\def\fbp{fixed boundary point}
\def\fbl{fixed boundary line}\def\fif{forward invariant face}
\def\gas{global asymptotic stability}\def\GAS{globally asymptotically stable}
\def\GASS{\text{GASS}}
\def\Hb{Hopf bifurcation}
\def\Hur{Hurwitz}\def\MH{Metzler-Hurwitz}\def\RH{Routh-Hurwitz}
\def\IS{\text{IS}}
\def\LAS{locally asymptotically stable}
\def\LdS{Lyapunov diagonally stable}
\def\Lf{Lyapunov function}\def\LyV{Lyapunov Volterra type function}\def\RLF{robust Lyapunov function}
\def\LyM{Lyapunov Malkin stable}
\def\PF{Perron-Frobenius}
\def\spfp{strictly positive fixed point}
\def\sta{stability}
\def\STP{sharp threshold property}\def\STT{sharp threshold theorem}

% -- Models (with \varphi explicitly instead of \f) --
\def\AB{$(A,B)$ Arino-Brauer epidemic models}
\def\ABa{$(A,\bbe,\ba,\varphi)$ balanced bilinear model}
\def\ABP{%$(A,\bbe,P,\varphi)$
\\ balanced bilinear model}
\def\AMG{Anderson May Gupta}
\newcommand\BIN{biological interaction network}
\def\Dt{Descartes type}
\def\gLV{generalized Lotka-Volterra model}
\def\IG{\text{invasion graph}}\def\IGMS{interaction graph on minimal siphons}
\def\LCP{\text{linear complementarity problem}}
\def\LV{Lotka-Volterra}\def\LVK{Lotka-Volterra-Kolmogorov}

% -- Splitting & Structure --
\def\reg{regular splitting}\def\regS{{\bf regular splitting}}
\def\RUR{rational univariate representation}
\def\bds{block-diagonal structure}\def\NGMs{next generation matrices}

% -- Miscellaneous --
\def\Ma{{Mathematica}}
\def\Mm{Max–Min}
\def\num{numerator}
\def\rap{rational approximation} \def\DSR{directed species-reaction bipartite graph}\def\PCP{polynomial complementarity problem}\def\rat{rational substitution}
\def\QSS{quasi steady state} \def\RHS{right hand side} \def\LHS{left hand side}
\def\OSN{online social network}\def\VdD{Van den Driessche}
\def\wli{$\omega$-limit} \def\rk{rich kinetics}\def\sJ{symbolic Jacobian}
\def\cV{\cite{Vas,VasSta,BSV,GSV}} \def\spf{semi-parametric semi-functional}
\def\CRW{Capasso-Ruan-Wang} \def\ABP{$(A,\bbe,P,\f)$ balanced bilinear model}
\def\ABa{%$(A,\bbe,\ba,\f)$ \\
rank-one balanced bilinear model}
\def\bbe{B%{\bff \beta}
} \def\vbe{\vec {\bff \beta}
}
\def\vb{\vec {\bff b}}
\newcommand{\Kh}{\hat{K}}          % K with hat, first mention only
\newcommand{\wS}{w^{\sigma}}       % right Perron J^\perp, full block
\newcommand{\piS}{\vec \pi^{\sigma}}    % left Perron J^\perp, full block
\newcommand{\ws}{w^{\sigma}}       % right Perron J^\perp, sub-block
\newcommand{\pis}{\pi^{\sigma}}    % left Perron J^\perp, sub-block
\newcommand{\hwS}{\hat{w}^{\sigma}}   % right Perron K, full block
\newcommand{\hpiS}{\hat{\pi}^{\sigma}} % left Perron K, full block
\newcommand{\hws}{\hat{w}^{\sigma}}   % right Perron K, sub-block
\newcommand{\hpis}{\hat{\pi}^{\sigma}} % left Perron K, sub-block
\newcommand{\twS}{\tilde{w}^{\sigma}}  % right Perron \tilde K, full block
\newcommand{\tpiS}{\tilde{\pi}^{\sigma}} % left Perron \tilde K
\newcommand{\tws}{\tilde{w}^{\sigma}}
\newcommand{\tpis}{\tilde{\pi}^{\sigma}}
\newcommand{\Mat}[2]{\mathbb{R}_{\ge 0}^{\,#1\times #2}}
\def\tJ{J_\sigma^\perp} \def\lPe{left Perron eigenvector}
% \G is \Gamma (defined at :795).  The Volterra function keeps its own name,
% so that \G no longer has two conflicting meanings in this file.
\def\GV{G}  % Volterra G(u)=u-1-\ln u
\def\PVL{Perron-Volterra Lyapunov minimal siphon ansatz}
\newcommand{\istar}[1]{i_{#1}^{*}}
%\newcommand{\sE}{s_E}
\newcommand{\vE}{v_E}
\def\PV{Perron-Volterra candidate Lyapunov function}
\def\GLV{generalized Lotka-Volterra}
\def\MERN{mathematical epidemiology reaction networks}
\def\CMM{contraction metric method}
\newcommand{\sign}{\operatorname{sign}}
\newcommand{\Var}{\operatorname{Var}} 

% ======================================================================
% IDEAS SHARED BETWEEN THE ABSTRACT AND THE MAIN RESULTS SECTION
%
% Write each idea ONCE, in the abstract; the Main Results section recalls it,
% so the two can never drift apart.
%
%   in the abstract        \ideaitem{1}{The first-order reaction terms ...}
%   in Main Results        \ideahead{1}
%
% \ideaitem stores the text globally and typesets
%   \item \hyperlink{idea:1}{<text>}
% \ideahead typesets
%   \item\hypertarget{idea:1}{}\textbf{<the same text>}
%
% The abstract must come before the Main Results input, which is the normal
% order. If a heading has to read differently from the abstract sentence (for
% instance a long abstract item summarised in one line), override it with the
% optional argument: \ideahead[short heading]{6}.
% ======================================================================
\makeatletter
\newcommand{\ideaitem}[2]{%
  \expandafter\gdef\csname @idea@#1\endcsname{#2}%
  \item\hyperlink{idea:#1}{#2}}
\newcommand{\ideahead}[2][]{%
  \item\hypertarget{idea:#2}{}%
  \def\@ideatmp{#1}%
  \ifx\@ideatmp\@empty
    \@ifundefined{@idea@#2}{%
      \PackageWarning{preamble}{Idea #2 is used in Main Results but was never
        defined by \string\ideaitem\space in the abstract}%
      \textbf{[idea #2 undefined]}%
    }{\textbf{\begingroup\def\ideaterm##1##2{##2}\def\idearest##1{}\@nameuse{@idea@#2}\endgroup}}%
  \else
    \textbf{#1}%
  \fi}
% \ideaterm{<hypertarget>}{<words>}: inside an abstract idea the words become an
% invisible link (the document sets \hypersetup{hidelinks}) jumping to the place
% in Main Results that defines them; inside the Main Results heading itself,
% where the target sits a few lines below, \ideahead strips the link and prints
% the words plainly. Mark the destination with \hypertarget{<name>}{} on the line
% before the display, which adds no equation number and shows nothing.
\newcommand{\ideaterm}[2]{\hyperlink{#1}{#2}}
% \idearest{<text>}: the part of an abstract idea after its first sentence; printed in
% the abstract, omitted from the Main Results heading.
\newcommand{\idearest}[1]{#1}
\makeatother

% ======================================================================
% SHORT-PAPER SWITCH
% \shorttrue in a master that assembles only part of the companion paper
% (Barrat/SipArxiv.tex). \nsh{...} drops text that cites results living in the
% omitted sections, so the short build has no dangling references while the
% companion paper keeps the text verbatim.
% ======================================================================
\newif\ifshort \shortfalse
\newcommand{\nsh}[1]{\ifshort\else#1\fi}
% \shortalt{<short build>}{<companion paper>}: use when a sentence needs
% different wording in the two builds, not merely a deletion.
\newcommand{\shortalt}[2]{\ifshort#1\else#2\fi}
% \Jtr, \JtrS{\Sigma}: transverse Jacobian. J^\perp here; SipArxiv.tex redefines them as J_{\mathrm{tra}}, J_{\mathrm{tra},\Sigma}
\newcommand{\Jtr}{J^\perp}
\newcommand{\JtrS}[1]{J_{#1}^\perp}
.
\makeatletter
\AtBeginDocument{\@ifundefined{hypersetup}{}{\hypersetup{linktoc=all}}}
\makeatother
\usepackage{graphicx,float,fancybox}
\usepackage{epsfig,epstopdf,tikz}
\usepackage[sans]{dsfont}
\usetikzlibrary{positioning,arrows,arrows.meta,calc}
% Configure listings for Mathematica code
\lstdefinestyle{mathematica}{
 language=Mathematica,
 basicstyle=\small\ttfamily,
 backgroundcolor=\color{gray!10},
 frame=single,
 breaklines=true,
 commentstyle=\color{green!60!black},
 keywordstyle=\color{blue},
 stringstyle=\color{red},
 showstringspaces=false,
 tabsize=2
}
\usepackage{tcolorbox}
% Colored boxes for key results
\newtcolorbox{keyresult}[1][]{
 colback=blue!5!white,
 colframe=blue!75!black,
 title=#1,
 fonttitle=\bfseries
}
\newtcolorbox{examplebox}[1][]{
 colback=green!5!white,
 colframe=green!75!black,
 title=#1,
 fonttitle=\bfseries
}


\makeatletter
\@ifundefined{theorem}   {\newtheorem{theorem}   {Theorem}}   {}
\@ifundefined{definition}{\newtheorem{definition}{Definition}} {}
\@ifundefined{open}      {\newtheorem{open}      {Problem}}   {}
\@ifundefined{example}   {\newtheorem{example}   {Example}}   {}
\makeatother

% COUNTER DEFINITIONS
\newcounter{tabcounter}
\newcounter{defcounter}
\newcounter{opcounter}
\newcounter{thmcounter}
\newcounter{excounter}
\newcounter{ascounter}
\newcounter{cjcounter}
\newcounter{lemcounter}
\newcounter{remcounter}
\newcounter{corcounter}
\newcounter{propcounter}
\newcounter{figcounter}

% UNIFIED GLOSSARY STYLE
\newglossarystyle{unifiedstyle}{%
  \setglossarystyle{list}%
  \renewcommand*{\glossarypreamble}{\gdef\glslastsec{}}%
  \renewcommand{\glossentry}[2]{%
    \edef\glscursec{\glsentryuseriii{##1}}%
    \ifx\glscursec\glslastsec\else
      \item[]\textbf{\S\,\glscursec\enspace\glsentryuseriv{##1}}%
      \global\let\glslastsec\glscursec
    \fi
    \item[\glsentryuserii{##1} \glsentryuseri{##1}:] ``\glossentryname{##1}'' \glsentrydesc{##1}%
  }%
}

\makeatletter
% First-word extractor (used by \beT, \beD and \beXa for the default key)
% \@firstwordof hello world \relax  ->  hello
\def\@firstwordof#1 #2\relax{#1}
% First three letters, used for the heading mnemonic of a theorem given no
% explicit key.  Padding with \@empty keeps it safe for first words shorter
% than three letters: \@empty vanishes under \edef, whereas \relax would
% survive into the label and is illegal inside \csname (it broke every
% two-letter title, e.g. "ME ...", with "Missing \endcsname inserted").
\def\@takethree#1#2#3#4\@nil{#1#2#3}

% EXAMPLE COMMAND -- same signature and behaviour as \beT and \beD.
% Usage: \beXa*[long title][short title][key]
%   star  = suppresses the printed "[key]" tag after the title (as for \beT*
%   and \beD*); the label and glossary entry are unaffected.
%   key = optional label key override; the label becomes e:KEY, and without a
%   key it is e: plus the first three letters of the title's first word.  A
%   mnemonic label lets an example be cited as "Example~\ref{e:KEY}" instead
%   of by section number, which is what the citing prose needs.
\NewDocumentCommand{\beXa}{s O{} O{} o}{%
  \def\@temparg{#2}%
  \ifx\@temparg\@empty
    \begin{example}%
  \else
    \let\@savedSecTitle\@currentsectionname%
    \refstepcounter{excounter}%
    \IfValueTF{#4}{%
      \edef\@exlblkey{#4}%
    }{%
      \edef\@exfirstword{\@firstwordof#2 \relax}%
      \edef\@exlblkey{%
        \expandafter\@takethree\@exfirstword\@empty\@empty\@empty\@nil}%
    }%
    \def\@tempshort{#3}%
    \ifx\@tempshort\@empty
      \expandafter\xdef\csname @xshort@\@exlblkey\endcsname{\@exlblkey}%
    \else
      \expandafter\def\csname @xshort@\@exlblkey\endcsname{#3}%
    \fi
    \IfBooleanTF{#1}{\let\@extag\@empty}{%
      \edef\@extag{\noexpand\ \noexpand\textup{[\noexpand\texttt{e:\@exlblkey}]}}}%
    \begin{example}[#2\ifshowkeys\@extag\fi]%
    \label{e:\@exlblkey}%
    % record the long title so \nref{e:KEY} prints it, as for theorems
    \protected@write\@auxout{}{%
      \string\@thmtitle{e:\@exlblkey}{\unexpanded{#2}}}%
    \expanded{%
      \noexpand\newglossaryentry{e\theexcounter}{%
        name={\unexpanded{#2}%
          \IfBooleanF{#1}{ [\noexpand\texttt{e:\@exlblkey}]}},%
        sort={\theexample},%
        description={p.~\noexpand\pageref{e:\@exlblkey}},%
        user1={\theexample},%
        user2={Example},%
        user3={\thesection},%
        user4={\expandafter\unexpanded\expandafter{\@savedSecTitle}}%
      }%
    }%
    \protected@write\@auxout{}{\string\glsliveentry{e\theexcounter}}%
  \fi
}
\newcommand{\eeXa}{\end{example}}

% DEFINITION COMMAND -- same signature and behaviour as \beT below.
% Usage: \beD[long title][short title][key]
%   long title  = definition heading (plain text only)
%   short title = shown in \dref{word}; defaults to the key
%   key         = optional label key override; the label becomes d:KEY
% Without a key the label is d: plus the first three letters of the title's
% first word, so it is a mnemonic rather than the definition's number: a
% numeric label would silently change whenever a definition is inserted.
\NewDocumentCommand{\beD}{s O{} O{} o}{%
  \def\@temparg{#2}%
  \ifx\@temparg\@empty
    \begin{definition}%
  \else
    \let\@savedSecTitle\@currentsectionname%
    \refstepcounter{defcounter}%
    \IfValueTF{#4}{%
      \edef\@deflblkey{#4}%
    }{%
      \edef\@deffirstword{\@firstwordof#2 \relax}%
      \edef\@deflblkey{%
        \expandafter\@takethree\@deffirstword\@empty\@empty\@empty\@nil}%
    }%
    \def\@tempshort{#3}%
    \ifx\@tempshort\@empty
      \expandafter\xdef\csname @dshort@\@deflblkey\endcsname{\@deflblkey}%
    \else
      \expandafter\def\csname @dshort@\@deflblkey\endcsname{#3}%
    \fi
    \IfBooleanTF{#1}{\let\@deftag\@empty}{%
      \edef\@deftag{\noexpand\ \noexpand\textup{[\noexpand\texttt{d:\@deflblkey}]}}}%
    \begin{definition}[#2\ifshowkeys\@deftag\fi]%
    \label{d:\@deflblkey}%
    % record the long title so \nref{d:KEY} prints it, as for theorems
    \protected@write\@auxout{}{%
      \string\@thmtitle{d:\@deflblkey}{\unexpanded{#2}}}%
    \expanded{%
      \noexpand\newglossaryentry{d\thedefcounter}{%
        name={\unexpanded{#2}%
          \IfBooleanF{#1}{ [\noexpand\texttt{d:\@deflblkey}]}},%
        sort={\thedefinition},%
        description={p.~\noexpand\pageref{d:\@deflblkey}},%
        user1={\thedefinition},%
        user2={Definition},%
        user3={\thesection},%
        user4={\expandafter\unexpanded\expandafter{\@savedSecTitle}}%
      }%
    }%
    \protected@write\@auxout{}{\string\glsliveentry{d\thedefcounter}}%
  \fi
}
\newcommand{\eeD}{\end{definition}}

% \showkeysfalse in a driver hides the [t:KEY] mnemonics in result headings
% (labels and glossary entries are unchanged); used for submission builds.
\newif\ifshowkeys \showkeystrue
% \commonredtrue in a driver that inputs basicFiles/commonRed: tool.tex then
% cites its reduction results instead of restating them.
\newif\ifcommonred \commonredfalse
% TheorEM COMMAND
% Usage: \beT[long title][short title][key]
%   long title  = theorem heading (plain text only)
%   short title = shown in \tref{word}; may contain math; defaults to first word of long title
%   key         = optional label key override; label becomes t:KEY (avoids first-word collisions)
% Label is t:WORD where WORD = first word of long title (or KEY if given); use \tref{WORD} to reference.
\NewDocumentCommand{\beT}{s O{} O{} o}{%
  \def\@temparg{#2}%
  \ifx\@temparg\@empty
    \begin{theorem}%
  \else
    \let\@savedSecTitle\@currentsectionname%
    \refstepcounter{thmcounter}%
    % \@thmtag is the mnemonic shown in the heading; every theorem gets one.
    % With an explicit key the key itself is shown.  Without one the key
    % defaults to the title's first word, and showing that whole word would
    % merely repeat the title ("Dissipativity [Dissipativity]"), so only its
    % first three letters are shown: "Dissipativity [Dis]".  The label is
    % unaffected -- it still uses the full \@thmlblkey.
    \IfValueTF{#4}{%
      \edef\@thmlblkey{#4}%
    }{%
      \edef\@thmfirstword{\@firstwordof#2 \relax}%
      \edef\@thmlblkey{%
        \expandafter\@takethree\@thmfirstword\@empty\@empty\@empty\@nil}%
    }%
    \IfBooleanTF{#1}{\let\@thmtag\@empty}{%
      \edef\@thmtag{\noexpand\ \noexpand\textup{[\noexpand\texttt{t:\@thmlblkey}]}}}%
    \def\@tempshort{#3}%
    \ifx\@tempshort\@empty
      \expandafter\xdef\csname @tshort@\@thmlblkey\endcsname{\@thmlblkey}%
    \else
      \expandafter\def\csname @tshort@\@thmlblkey\endcsname{#3}%
    \fi
    \begin{theorem}[#2\ifshowkeys\@thmtag\fi]%
    \label{t:\@thmlblkey}%
    % record the long title for \nref; via the .aux so forward refs resolve
    \protected@write\@auxout{}{%
      \string\@thmtitle{t:\@thmlblkey}{\unexpanded{#2}}}%
    \expanded{%
      \noexpand\newglossaryentry{t\thethmcounter}{%
        name={\unexpanded{#2}%
          \IfBooleanF{#1}{ [\noexpand\texttt{t:\@thmlblkey}]}},%
        % sort/user1 must be the number the document actually prints, i.e. the
        % theorem environment's counter, not thmcounter.  thmcounter advances
        % only on \beT with a title, so a bare \beT (or \beT with an
        % empty title) makes the two drift and the glossary cites a wrong
        % number.  thmcounter is still fine as the entry key: it is unique.
        sort={\thetheorem},%
        description={p.~\noexpand\pageref{t:\@thmlblkey}},%
        user1={\thetheorem},%
        user2={Theorem},%
        user3={\thesection},%
        user4={\expandafter\unexpanded\expandafter{\@savedSecTitle}}%
      }%
    }%
    \protected@write\@auxout{}{\string\glsliveentry{t\thethmcounter}}%
  \fi
}
\newcommand{\eeT}{\end{theorem}}

% CROSS-PART RELATION PARAGRAPH
% Usage: \xpart{text}  -- place immediately BEFORE \beT, never inside it
%   (anything inside \beT...\eeT would be swallowed into the glossary entry).
% The paragraph is typeset only when Part I is actually present, i.e. in the
% combined book build; in the standalone builds (GranadaDSABNS/A.tex,
% AAHRah-Submission/AAHRah.tex) the label part:glv is undefined and the whole
% paragraph disappears rather than leaving ?? cross-references. Same
% graceful-degradation principle as \GASref in GASt.tex. Resolves on the
% second LaTeX pass, since r@part:glv comes from the .aux file.
% (Already inside the \makeatletter region opened above, so no \makeatletter
% pair here -- adding one would close that region early.)
\newcommand{\xpart}[1]{\@ifundefined{r@part:glv}{}{\paragraph{Relation with Parts~I--II.}#1}}

% \bookref{label}{precise text}{fallback text}: prints the precise cross-reference
% when the label resolves (combined book build), and the fallback text otherwise
% (standalone paper builds, where book-only appendix labels are absent). Same
% graceful-degradation role as \GASref in GASt.tex, but usable in any file since it
% lives in this preamble (loaded by every driver) -- unlike \GASref, which is defined in a
% content file and is therefore undefined in standalone builds that omit that file.
\newcommand{\bookref}[3]{\@ifundefined{r@#1}{#3}{#2}}

% \PVpart and \RBpart: the two parts that a partial build may omit.
% They print "Part~N" in any build that contains the part, and a naming phrase
% in a build that does not, so an excerpt shows no ?? .  Each already supplies
% the word "Part", so write "\PVpart develops ...", not "Part~\PVpart".
\newcommand{\PVpart}{\bookref{part:aahrah}%
  {Part~\ref{part:aahrah}}{the Perron--Volterra part of the full book}}
\newcommand{\RBpart}{\bookref{part:recbrowne}%
  {Part~\ref{part:recbrowne}}{the immuno-virology part of the full book}}

% OPEN PROBLEM COMMANDS
% ASSUMPTION COMMANDS (glossaried)
\newcommand{\beAs}[1][]{%
  \def\@temparg{#1}%
  \ifx\@temparg\@empty
    \begin{assumption}%
  \else
    \let\@savedSecTitle\@currentsectionname%
    \refstepcounter{ascounter}%
    \begin{assumption}[#1]%
    \label{as:\theascounter}%
    \expanded{%
      \noexpand\newglossaryentry{as\theascounter}{%
        name={\unexpanded{#1}},%
        sort={\theascounter},%
        description={p.~\noexpand\pageref{as:\theascounter}},%
        user1={\theascounter},%
        user2={Assumption},%
        user3={\thesection},%
        user4={\expandafter\unexpanded\expandafter{\@savedSecTitle}}%
      }%
    }%
    \protected@write\@auxout{}{\string\glsliveentry{as\theascounter}}%
  \fi
}
\newcommand{\eeAs}{\end{assumption}}
% CONJECTURE COMMANDS (glossaried)
\newcommand{\beCj}[1][]{%
  \def\@temparg{#1}%
  \ifx\@temparg\@empty
    \begin{conjecture}%
  \else
    \let\@savedSecTitle\@currentsectionname%
    \refstepcounter{cjcounter}%
    \begin{conjecture}[#1]%
    \label{cj:\thecjcounter}%
    \expanded{%
      \noexpand\newglossaryentry{cj\thecjcounter}{%
        name={\unexpanded{#1}},%
        sort={\thecjcounter},%
        description={p.~\noexpand\pageref{cj:\thecjcounter}},%
        user1={\thecjcounter},%
        user2={Conjecture},%
        user3={\thesection},%
        user4={\expandafter\unexpanded\expandafter{\@savedSecTitle}}%
      }%
    }%
    \protected@write\@auxout{}{\string\glsliveentry{cj\thecjcounter}}%
  \fi
}
\newcommand{\eeCj}{\end{conjecture}}
% LEMMA COMMANDS
% ---------------------------------------------------------------------
% UNIFORM RESULT WRAPPERS -- \beL, \beC, \beP, \beR share \beT's signature
%
%   \beX[long title][short title][key]  ...  \eeX
%
%   long title  = heading (math allowed)
%   short title = shown in the cross-reference; defaults to the key
%   key         = label key.  With a key the label is X:KEY and the heading
%                 carries the mnemonic [key], exactly as for \beT.  Without
%                 one the label falls back to the running counter, which is
%                 what every pre-existing call site in the book relies on.
%
% A bare \beX -- arguments commented out -- opens the plain environment and
% creates NO label and NO glossary entry.  That is the supported way to keep
% an individual result out of the catalogue.
% ---------------------------------------------------------------------
% #1 wrapper name  #2 environment  #3 label prefix  #4 counter  #5 type name
\newcommand{\@mkresult}[5]{%
  % A leading * suppresses the [key] mnemonic in both the heading and the
  % glossary name, while still making the label X:KEY and the entry:
  %     \beX*[long][short][key]
  \expandafter\NewDocumentCommand\csname be#1\endcsname{s O{} O{} o}{%
    \def\@temparg{##2}%
    \ifx\@temparg\@empty
      \begin{#2}%
    \else
      \let\@savedSecTitle\@currentsectionname%
      \refstepcounter{#4}%
      % With a key: mnemonic label X:KEY and the [key] tag in the heading.
      % Without one: fall back to the running counter, exactly as these
      % wrappers behaved before, so every legacy one-argument call site in
      % the book keeps its label.  The counter is also csname-safe, which a
      % title's first word is not -- titles beginning with math would
      % otherwise build \csname \Omega...\endcsname and fail.
      \IfValueTF{##4}{%
        \edef\@relbl{##4}%
        \IfBooleanTF{##1}{\let\@retag\@empty}{%
          \edef\@retag{\noexpand\ \noexpand\textup{[\noexpand\texttt{#3:##4}]}}}%
      }{%
        \edef\@relbl{\arabic{#4}}%
        \let\@retag\@empty
      }%
      \def\@tempshort{##3}%
      \ifx\@tempshort\@empty
        \expandafter\xdef\csname @sshort@#3\@relbl\endcsname{\@relbl}%
      \else
        \expandafter\def\csname @sshort@#3\@relbl\endcsname{##3}%
      \fi
      \begin{#2}[##2\ifshowkeys\@retag\fi]%
      \edef\@tmplbl{#3:\@relbl}\expandafter\label\expandafter{\@tmplbl}%
      \expanded{%
        \noexpand\newglossaryentry{#3\csname the#4\endcsname}{%
          name={\unexpanded{##2}%
            \IfBooleanF{##1}{\IfValueT{##4}{ [\noexpand\texttt{#3:##4}]}}},%
          sort={\csname the#2\endcsname},%
          description={p.~\noexpand\pageref{#3:\@relbl}},%
          user1={\csname the#2\endcsname},%
          user2={#5},%
          user3={\thesection},%
          user4={\expandafter\unexpanded\expandafter{\@savedSecTitle}}%
        }%
      }%
      \protected@write\@auxout{}{\string\glsliveentry{#3\csname the#4\endcsname}}%
      \protected@write\@auxout{}{\string\@thmtitle{\@tmplbl}{\unexpanded{##2}}}%
    \fi
  }%
  \expandafter\newcommand\csname ee#1\endcsname{\end{#2}}%
}
\@mkresult{L}{lemma}{l}{lemcounter}{Lemma}
% REMARK / COROLLARY / PROPOSITION -- same signature as \beL above
\@mkresult{R}{remark}{r}{remcounter}{Remark}
\@mkresult{C}{corollary}{c}{corcounter}{Corollary}
\@mkresult{P}{proposition}{p}{propcounter}{Proposition}
% OPEN-PROBLEM COMMAND -- generated with the same signature as \beT:
%   \beO[long title][short title][key]   ...   \eeO
% label o:KEY, or o:<counter> when no key is given; a bare \beO makes
% neither a label nor a glossary entry.
\@mkresult{O}{open}{o}{opcounter}{Problem}


% Section-level title tracker: updated only by \section, not \subsection.
% \@currentlabelname is overwritten by every \refstepcounter (including
% \subsection), so we save it separately here.
\newcommand{\@currentsectionname}{}
\AtBeginDocument{%
  \let\@epidcrn@orig@refstepcounter\refstepcounter
  \renewcommand{\refstepcounter}[1]{%
    \@epidcrn@orig@refstepcounter{#1}%
    \ifstrequal{#1}{section}{\global\let\@currentsectionname\@currentlabelname}{}%
  }%
}

% TABLE COMMANDS (glossaried) -- consistent naming with \beT, \beD, etc.
\newcommand{\beTa}[1][]{%
  \def\@temparg{#1}%
  \ifx\@temparg\@empty
    \begin{table}[htbp]%
  \else
    \let\@savedSecTitle\@currentsectionname%
    \refstepcounter{tabcounter}%
    \begin{table}[htbp]%
    \label{A:\thetabcounter}%
    \expanded{%
      \noexpand\newglossaryentry{A\thetabcounter}{%
        name={\unexpanded{#1}},%
        sort={\thetabcounter},%
        description={p.~\noexpand\pageref{A:\thetabcounter}},%
        user1={\thetabcounter},%
        user2={Table},%
        user3={\thesection},%
        user4={\expandafter\unexpanded\expandafter{\@savedSecTitle}}%
      }%
    }%
    \protected@write\@auxout{}{\string\glsliveentry{A\thetabcounter}}%
  \fi
}
\newcommand{\eeTa}{\end{table}}

% FIGURE COMMANDS (glossaried) -- consistent with \beT, \beR, \beTa, etc.
% Usage:
%   \beF[Glossary title]...\caption{Visible caption}\eeF  -> glossaried figure
%   \beF...\caption{Visible caption}\eeF                  -> plain [H] figure, no glossary
\newcommand{\beF}[1][]{%
  \def\@temparg{#1}%
  \ifx\@temparg\@empty
    \begin{figure}[H]%
  \else
    \begin{figure}[H]%
    \let\@savedSecTitle\@currentsectionname%
    \refstepcounter{figcounter}%
    \label{fig:\thefigcounter}%
    \expanded{%
      \noexpand\newglossaryentry{Fig\thefigcounter}{%
        name={\unexpanded{#1}},%
        sort={\thefigcounter},%
        description={p.~\noexpand\pageref{fig:\thefigcounter}},%
        user1={\thefigcounter},%
        user2={Figure},%
        user3={\thesection},%
        user4={\expandafter\unexpanded\expandafter{\@savedSecTitle}}%
      }%
    }%
    \protected@write\@auxout{}{\string\glsliveentry{Fig\thefigcounter}}%
  \fi
}
\newcommand{\eeF}{\end{figure}}
\makeatother
% LIVE-ENTRY REGISTER
% \glsaddall prints every entry the glossaries package currently knows, and
% that set includes entries restored from \jobname.glsdefs.  The package
% rewrites .glsdefs from all *currently defined* entries, so anything it
% restored is written straight back out: a deleted section kept reappearing in
% the catalogue for ever, and the .glsdefs had to be erased by hand.  The
% wrappers therefore record every entry they actually create in the .aux, and
% \printboth adds only those.  The .aux is rebuilt from scratch on each run, so
% this is self-healing: a removed result is gone after one further run.
\makeatletter
\newcommand{\@glslive}{}
\newcommand{\@glsaddlive}[1]{\ifglsentryexists{#1}{\glsadd{#1}}{}}
\newcommand{\glsliveentry}[1]{\g@addto@macro\@glslive{\@glsaddlive{#1}}}
\makeatother
% PRINT COMMAND
% No \clearpage here -- the catalogue continues on the current page.  This has
% been re-added by mistake several times; leave it out.
\makeatletter
\newcommand{\printboth}{%
  \@glslive   % add exactly the entries this document created last run
  \printglossary[style=unifiedstyle, title={0 \hspace{.22em} Definitions, Theorems, Examples and Open Problems}]%
}
\makeatother
% Named theorem reference: \tref{word} -> Thm~N(word-or-short)
% Shows stored short title if available, otherwise shows the key itself.
\newcommand{\tref}[1]{%
  Thm~\ref{t:#1}(\ifcsname @tshort@#1\endcsname\csname @tshort@#1\endcsname\else#1\fi)%
}

% Full-title theorem reference: \nref{t:KEY} prints the number and then the
% theorem's long title, so one writes "Theorem~\nref{t:blocks}" and gets
% "Theorem 3 ``Block factorization and siphons''".  Takes the complete label, like
% \ref, not the bare key like \tref.
%
% \beT writes each title to the .aux (see \@thmtitle below), so \nref resolves
% forward references from the second LaTeX pass, exactly as \ref does.  If the
% title is not yet known -- first pass, or a label that is not a \beT theorem
% -- it degrades to the number rather than producing ??.
% These two use @ in their names, and this point in the file is outside the
% \makeatletter region opened earlier (it closes at \makeatother above), so
% the region has to be reopened here.
\makeatletter
\newcommand{\@thmtitle}[2]{\global\@namedef{@tlong@#1}{#2}}
\newcommand{\nref}[1]{%
  \ref{#1}%
  \ifcsname @tlong@#1\endcsname
    \ ``\csname @tlong@#1\endcsname''%
  \fi
}
% \nm{t:KEY}: the stored long title of any \beT/\beL/\beP/\beR/\beC/\beD/\beE result,
% verbatim and without quotes, so item names follow the macro's first argument.
\newcommand{\nm}[1]{%
  \ifcsname @tlong@#1\endcsname
    \csname @tlong@#1\endcsname%
  \else
    \ref{#1}%
  \fi
}
\makeatother




%% ========================================
%% Other THEOREM ENVIRONMENTS
%% ========================================
\makeatletter
\@ifundefined{lemma}      {\newtheorem{lemma}      {Lemma}}      {}
\@ifundefined{proposition}{\newtheorem{proposition}{Proposition}} {}
\@ifundefined{corollary}  {\newtheorem{corollary}  {Corollary}}  {}
\@ifundefined{conjecture} {\newtheorem{conjecture} {Conjecture}}  {}
\@ifundefined{remark}     {\newtheorem{remark}     {Remark}}     {}
\@ifundefined{question}   {\newtheorem{question}   {Question}}   {}
\@ifundefined{assumption} {\newtheorem{assumption} {Assumption}} {}
\@ifundefined{hypothesis} {\newtheorem{hypothesis} {Hypothesis}} {}
\@ifundefined{fact}       {\newtheorem{fact}       {Fact}}       {}
\makeatother

% \citer: cite the relay-graph threshold theorem of AAHrel. If a document has
% \input the chapter/file containing t:relayRF (so the label resolves), point
% straight at it; otherwise fall back to a bare citation of AAHrel. Shared here
% (not per-paper) so it resolves consistently across every book chapter.
% \citeR: same idea, for the first theorem of AAHRah (t:Pstar).
\makeatletter
\providecommand{\citer}{%
  \@ifundefined{r@t:relayRF}%
    {\cite{AAHrel}}%
    {Theorem~\ref{t:relayRF} of \cite{AAHrel}}%
}
\providecommand{\citeR}{%
  \@ifundefined{r@t:Pstar}%
    {\cite{AAHRah}}%
    {Theorem~\ref{t:Pstar} of \cite{AAHRah}}%
}
\makeatother


% Abbreviations
\def\beQ{\begin{question}}\def\eeQ{\end{question}}
\def\beA{\begin{assumption}}\def\eeA{\end{assumption}}


%% ========================================

%% ========================================
%% COLORS & CUSTOM COMMANDS
%% ========================================
\definecolor{funccolor}{RGB}{25,25,112}
\definecolor{desccolor}{RGB}{64,64,64}
\newcommand{\func}[1]{\texttt{\textcolor{funccolor}{#1}}}
\def\eq{=}
\def\mm{-}

%% ========================================
%% MATHCAL SHORTHANDS
%% ========================================
% Project-wide convention: \mX for a capital X is always \mathcal X.
% Defined for the whole alphabet so the convention never fails to expand.
\newcommand{\mA}{\mathcal A}\newcommand{\mB}{\mathcal B}
\newcommand{\mC}{\mathcal C}\newcommand{\mD}{\mathcal D}
\newcommand{\mE}{\mathcal E}\newcommand{\mF}{\mathcal F}
\newcommand{\mG}{\mathcal G}\newcommand{\mH}{\mathcal H}
\newcommand{\mI}{\mathcal I}\newcommand{\mJ}{\mathcal J}
\newcommand{\mK}{\mathcal K}\newcommand{\mL}{\mathcal L}
\newcommand{\mM}{\mathcal M}\newcommand{\mN}{\mathcal N}
\newcommand{\mO}{\mathcal O}\newcommand{\mP}{\mathcal P}
\newcommand{\mQ}{\mathcal Q}\newcommand{\mR}{\mathcal R}
\newcommand{\mS}{\mathcal S}\newcommand{\mT}{\mathcal T}
\newcommand{\mU}{\mathcal U}\newcommand{\mV}{\mathcal V}
\newcommand{\mW}{\mathcal W}\newcommand{\mX}{\mathcal X}
\newcommand{\mY}{\mathcal Y}\newcommand{\mZ}{\mathcal Z}

%% ========================================
%% BASIC MATH ENVIRONMENTS
%% ========================================
\def\bep{\begin{pmatrix}}\def\eep{\end{pmatrix}}
\def\bev{\begin{vmatrix}}\def\eev{\end{vmatrix}}
\def\bea{\begin{eqnarray*}}\def\eea{\end{eqnarray*}}
\def\bc{\begin{cases}}\def\ec{\end{cases}}
\def\BEN{\begin{enumerate}}\def\EEN{\end{enumerate}}
\def\BI{\begin{itemize}}\def\EI{\end{itemize}}

\newcommand{\be}[1]{\begin{equation}\label{#1}}
\newcommand{\ee}{\end{equation}}
\newcommand{\beq}{\begin{eqnarray}}
\def\eeq{\end{eqnarray}}
\def\eqr{\eqref}\def\fr{\frac}\def\lbl{\label}
\def\Lra{\Longrightarrow}\def\non{\nonumber}\def\cNGM{\cite{Diek,Van,Van08}}
\def\Eq{\Leftrightarrow}\def\b{\beta}\def\uph{\red{UP TO HERE}}
%% ========================================
%% PROBABILITY & STATISTICS
%% ========================================
\renewcommand{\P}{\ensuremath{\mathbb{P}}}  % Note: overrides paragraph symbol
\newcommand{\E}{\ensuremath{\mathbb{E}}}
\newcommand{\ind}[1]{\ensuremath{\mathbbm{1}_{\left\{#1\right\}}}}
\DeclareMathOperator{\Cov}{Cov}
\DeclareMathOperator{\TV}{TV}
\newcommand{\rank}{\operatorname{rank}}

%% ========================================
%% CRN OPERATORS (organizations/ERCs/buffering/siphons framework)
%% ========================================
\def\clos{\mathrm{clos}}
\def\Supply{\mathrm{Supply}}
\def\req{\mathrm{req}}
\def\React{\mathrm{React}}
\def\Prod{\mathrm{Prod}}
\def\Gen{\mathrm{Gen}}
\def\In{\mathrm{In}}
\def\Out{\mathrm{Out}}
\def\Act{\mathrm{Act}}
\def\Min{\mathrm{Min}}

%% ========================================
%% LINE BREAKING
%% ========================================
% Last-resort stretch, used only on paragraphs TeX cannot otherwise break:
% a long \texttt{} token, a comma-separated inline math list (commas are not
% breakpoints in math), or a run of \cite/\ref boxes.  Without it those set as
% overfull \hbox and print into the margin; with it they set slightly loose,
% which is the lesser evil.  Does not affect paragraphs that already break.
\emergencystretch=3em

%% ========================================
%% NUMBER SETS
%% ========================================
\newcommand{\R}{\mathbb{R}}  % Real numbers
\newcommand{\C}{\mathbb{C}}  % Complex numbers
\newcommand{\N}{\mathbb{N}}
\newcommand{\Z}{\mathbb{Z}}
% \providecommand, not \newcommand: some papers define these in their own
% preamble after % Shared preamble for the GLVbook papers.  This is the single canonical copy;
% a driver one level below GLVbook reaches it as
%
%     \input{../basicFiles/preamble}
%
% What stays in the driver, because it differs from paper to paper:
% \documentclass, font choice, \title, \author, \date, and any macro genuinely
% local to one paper.  Everything else belongs here.  The file is
% self-contained: it does not read any other file.

\usepackage[utf8]{inputenc}
\usepackage[novbox]{pdfsync}   % required: SyncTeX alone does not work here

% Page geometry, set once here for every paper that loads this file.  A driver
% that needs different margins says \geometry{...}; a second
% \usepackage[...]{geometry} raises an option clash and must not be used.
\makeatletter
\@ifundefined{AAdefNnogeometry}{\usepackage[a4paper,margin=2cm]{geometry}}{}
\makeatother

% mathematics
\usepackage{amsmath,amssymb,amsthm,mathtools,bm,mathrsfs,bbm}
\usepackage{centernot}

% floats, tables, lists
\usepackage{graphicx,float}
\makeatletter
\@ifpackageloaded{subfig}{}{\usepackage{subcaption}}
\makeatother
\usepackage{booktabs,longtable,threeparttable,array,multirow}
\usepackage{enumitem}

% colour, code, algorithms
\usepackage{xcolor,listings,tcolorbox}
\makeatletter
\@ifundefined{c@algorithm}{\usepackage{algorithm}}{}
\makeatother
\usepackage{algorithmicx,algpseudocode}

% diagrams
\usepackage{tikz,tikz-cd,pgfplots}
\usetikzlibrary{graphs,decorations.pathmorphing,decorations.markings,
arrows.meta,automata,positioning}
\tikzcdset{every label/.append style = {font = \small}}
\pgfplotsset{compat=1.18}

% chemistry
\usepackage[version=4]{mhchem}

\usepackage{microtype}
\usepackage{lineno}
\usepackage{nomencl}
\makenomenclature

% hyperref first, then bookmark, then glossaries: glossaries must follow
% hyperref or the entry links are built against the wrong anchors.
\makeatletter
\@ifundefined{hyperrefpackage}{%
  \@ifundefined{AAdefNnohyperref}{\usepackage{hyperref}}{}%
}{}
\@ifundefined{hyperrefpackage}{\usepackage{bookmark}}{}
\makeatother
\usepackage[toc,sort=use]{glossaries}
\makeglossaries

% Book/standalone switch.  Material that the book already covers in its
% prelude, but that a standalone paper still has to state, is wrapped as
%     \ifbookbuild\else ... \fi
% Drivers default to standalone; bookM.tex says \bookbuildtrue.
\newif\ifbookbuild
\bookbuildfalse

% Part switches for the gradostat paper.  Both default to true; a driver that
% wants the short submission says \madecfalse and/or \slowfastfalse, and every
% cross-reference into the dropped part disappears with it.
%   \ifmadec     Part III, the Lagasquie--Madec model and its break-even geometry
%   \ifslowfast  Part IV, everything carrying the small parameter \eps
\newif\ifmadec     \madectrue
\newif\ifslowfast  \slowfasttrue

% Part switch for the Barrat/simplicial-contagion project.  Defaults to
% true (the hypernetwork-synchrony part is included); the journal-paper
% driver says \hnetpartfalse to drop it and every cross-reference into it.
\newif\ifhnetpart  \hnetparttrue

%% ========================================
%% PACKAGES
%% ========================================
\RequirePackage{etoolbox}
\usepackage{makecell}
\usepackage[novbox]{pdfsync}% obsolete; SyncTeX is built into modern engines
\usepackage{amssymb, amsfonts,mathtools,bbm,eurosym,graphicx,latexsym, amsmath,amsthm,times,nccmath,url}
\usepackage{listings,verbatim,cases,tabularx}
% hyperref is deliberately NOT in the line above.  Journal classes often
% load it themselves, with a long option list, from a begindocument hook
% (mdpi.cls does: \DeclareHookRule{begindocument}{hyperref}{before}{mdpi}).
% Loading it here first, without options, makes that later load raise
% "Option clash for package hyperref".  So load it only if nobody else
% has, or will:
%   - \@ifpackageloaded  catches a class that loaded it already;
%   - \hyperrefpackage    is mdpi.cls's own deferred loader;
%   - \AAdefNnohyperref   lets any other driver opt out by defining it
%                         before \input{preamble}.
\makeatletter
\@ifpackageloaded{hyperref}{}{%
  \@ifundefined{hyperrefpackage}{%
    \@ifundefined{AAdefNnohyperref}{\usepackage{hyperref}}{}%
  }{}%
}
\makeatother
% linktoc=all makes a table-of-contents page number a link as well as the
% entry text; hyperref's default, linktoc=section, links only the text.
% Deferred to begin-document and guarded: by the block above hyperref may
% be loaded later (mdpi's own hook) or not at all (\AAdefNnohyperref), and
% \hypersetup must exist before it is called.  A driver that wants the
% default back can say \hypersetup{linktoc=section} after \input{preamble}.
\makeatletter
\AtBeginDocument{\@ifundefined{hypersetup}{}{\hypersetup{linktoc=all}}}
\makeatother
\usepackage{graphicx,float,fancybox}
\usepackage{epsfig,epstopdf,tikz}
\usepackage[sans]{dsfont}
\usetikzlibrary{positioning,arrows,arrows.meta,calc}
% Configure listings for Mathematica code
\lstdefinestyle{mathematica}{
 language=Mathematica,
 basicstyle=\small\ttfamily,
 backgroundcolor=\color{gray!10},
 frame=single,
 breaklines=true,
 commentstyle=\color{green!60!black},
 keywordstyle=\color{blue},
 stringstyle=\color{red},
 showstringspaces=false,
 tabsize=2
}
\usepackage{tcolorbox}
% Colored boxes for key results
\newtcolorbox{keyresult}[1][]{
 colback=blue!5!white,
 colframe=blue!75!black,
 title=#1,
 fonttitle=\bfseries
}
\newtcolorbox{examplebox}[1][]{
 colback=green!5!white,
 colframe=green!75!black,
 title=#1,
 fonttitle=\bfseries
}


\makeatletter
\@ifundefined{theorem}   {\newtheorem{theorem}   {Theorem}}   {}
\@ifundefined{definition}{\newtheorem{definition}{Definition}} {}
\@ifundefined{open}      {\newtheorem{open}      {Problem}}   {}
\@ifundefined{example}   {\newtheorem{example}   {Example}}   {}
\makeatother

% COUNTER DEFINITIONS
\newcounter{tabcounter}
\newcounter{defcounter}
\newcounter{opcounter}
\newcounter{thmcounter}
\newcounter{excounter}
\newcounter{ascounter}
\newcounter{cjcounter}
\newcounter{lemcounter}
\newcounter{remcounter}
\newcounter{corcounter}
\newcounter{propcounter}
\newcounter{figcounter}

% UNIFIED GLOSSARY STYLE
\newglossarystyle{unifiedstyle}{%
  \setglossarystyle{list}%
  \renewcommand*{\glossarypreamble}{\gdef\glslastsec{}}%
  \renewcommand{\glossentry}[2]{%
    \edef\glscursec{\glsentryuseriii{##1}}%
    \ifx\glscursec\glslastsec\else
      \item[]\textbf{\S\,\glscursec\enspace\glsentryuseriv{##1}}%
      \global\let\glslastsec\glscursec
    \fi
    \item[\glsentryuserii{##1} \glsentryuseri{##1}:] ``\glossentryname{##1}'' \glsentrydesc{##1}%
  }%
}

\makeatletter
% First-word extractor (used by \beT, \beD and \beXa for the default key)
% \@firstwordof hello world \relax  ->  hello
\def\@firstwordof#1 #2\relax{#1}
% First three letters, used for the heading mnemonic of a theorem given no
% explicit key.  Padding with \@empty keeps it safe for first words shorter
% than three letters: \@empty vanishes under \edef, whereas \relax would
% survive into the label and is illegal inside \csname (it broke every
% two-letter title, e.g. "ME ...", with "Missing \endcsname inserted").
\def\@takethree#1#2#3#4\@nil{#1#2#3}

% EXAMPLE COMMAND -- same signature and behaviour as \beT and \beD.
% Usage: \beXa*[long title][short title][key]
%   star  = suppresses the printed "[key]" tag after the title (as for \beT*
%   and \beD*); the label and glossary entry are unaffected.
%   key = optional label key override; the label becomes e:KEY, and without a
%   key it is e: plus the first three letters of the title's first word.  A
%   mnemonic label lets an example be cited as "Example~\ref{e:KEY}" instead
%   of by section number, which is what the citing prose needs.
\NewDocumentCommand{\beXa}{s O{} O{} o}{%
  \def\@temparg{#2}%
  \ifx\@temparg\@empty
    \begin{example}%
  \else
    \let\@savedSecTitle\@currentsectionname%
    \refstepcounter{excounter}%
    \IfValueTF{#4}{%
      \edef\@exlblkey{#4}%
    }{%
      \edef\@exfirstword{\@firstwordof#2 \relax}%
      \edef\@exlblkey{%
        \expandafter\@takethree\@exfirstword\@empty\@empty\@empty\@nil}%
    }%
    \def\@tempshort{#3}%
    \ifx\@tempshort\@empty
      \expandafter\xdef\csname @xshort@\@exlblkey\endcsname{\@exlblkey}%
    \else
      \expandafter\def\csname @xshort@\@exlblkey\endcsname{#3}%
    \fi
    \IfBooleanTF{#1}{\let\@extag\@empty}{%
      \edef\@extag{\noexpand\ \noexpand\textup{[\noexpand\texttt{e:\@exlblkey}]}}}%
    \begin{example}[#2\ifshowkeys\@extag\fi]%
    \label{e:\@exlblkey}%
    % record the long title so \nref{e:KEY} prints it, as for theorems
    \protected@write\@auxout{}{%
      \string\@thmtitle{e:\@exlblkey}{\unexpanded{#2}}}%
    \expanded{%
      \noexpand\newglossaryentry{e\theexcounter}{%
        name={\unexpanded{#2}%
          \IfBooleanF{#1}{ [\noexpand\texttt{e:\@exlblkey}]}},%
        sort={\theexample},%
        description={p.~\noexpand\pageref{e:\@exlblkey}},%
        user1={\theexample},%
        user2={Example},%
        user3={\thesection},%
        user4={\expandafter\unexpanded\expandafter{\@savedSecTitle}}%
      }%
    }%
    \protected@write\@auxout{}{\string\glsliveentry{e\theexcounter}}%
  \fi
}
\newcommand{\eeXa}{\end{example}}

% DEFINITION COMMAND -- same signature and behaviour as \beT below.
% Usage: \beD[long title][short title][key]
%   long title  = definition heading (plain text only)
%   short title = shown in \dref{word}; defaults to the key
%   key         = optional label key override; the label becomes d:KEY
% Without a key the label is d: plus the first three letters of the title's
% first word, so it is a mnemonic rather than the definition's number: a
% numeric label would silently change whenever a definition is inserted.
\NewDocumentCommand{\beD}{s O{} O{} o}{%
  \def\@temparg{#2}%
  \ifx\@temparg\@empty
    \begin{definition}%
  \else
    \let\@savedSecTitle\@currentsectionname%
    \refstepcounter{defcounter}%
    \IfValueTF{#4}{%
      \edef\@deflblkey{#4}%
    }{%
      \edef\@deffirstword{\@firstwordof#2 \relax}%
      \edef\@deflblkey{%
        \expandafter\@takethree\@deffirstword\@empty\@empty\@empty\@nil}%
    }%
    \def\@tempshort{#3}%
    \ifx\@tempshort\@empty
      \expandafter\xdef\csname @dshort@\@deflblkey\endcsname{\@deflblkey}%
    \else
      \expandafter\def\csname @dshort@\@deflblkey\endcsname{#3}%
    \fi
    \IfBooleanTF{#1}{\let\@deftag\@empty}{%
      \edef\@deftag{\noexpand\ \noexpand\textup{[\noexpand\texttt{d:\@deflblkey}]}}}%
    \begin{definition}[#2\ifshowkeys\@deftag\fi]%
    \label{d:\@deflblkey}%
    % record the long title so \nref{d:KEY} prints it, as for theorems
    \protected@write\@auxout{}{%
      \string\@thmtitle{d:\@deflblkey}{\unexpanded{#2}}}%
    \expanded{%
      \noexpand\newglossaryentry{d\thedefcounter}{%
        name={\unexpanded{#2}%
          \IfBooleanF{#1}{ [\noexpand\texttt{d:\@deflblkey}]}},%
        sort={\thedefinition},%
        description={p.~\noexpand\pageref{d:\@deflblkey}},%
        user1={\thedefinition},%
        user2={Definition},%
        user3={\thesection},%
        user4={\expandafter\unexpanded\expandafter{\@savedSecTitle}}%
      }%
    }%
    \protected@write\@auxout{}{\string\glsliveentry{d\thedefcounter}}%
  \fi
}
\newcommand{\eeD}{\end{definition}}

% \showkeysfalse in a driver hides the [t:KEY] mnemonics in result headings
% (labels and glossary entries are unchanged); used for submission builds.
\newif\ifshowkeys \showkeystrue
% \commonredtrue in a driver that inputs basicFiles/commonRed: tool.tex then
% cites its reduction results instead of restating them.
\newif\ifcommonred \commonredfalse
% TheorEM COMMAND
% Usage: \beT[long title][short title][key]
%   long title  = theorem heading (plain text only)
%   short title = shown in \tref{word}; may contain math; defaults to first word of long title
%   key         = optional label key override; label becomes t:KEY (avoids first-word collisions)
% Label is t:WORD where WORD = first word of long title (or KEY if given); use \tref{WORD} to reference.
\NewDocumentCommand{\beT}{s O{} O{} o}{%
  \def\@temparg{#2}%
  \ifx\@temparg\@empty
    \begin{theorem}%
  \else
    \let\@savedSecTitle\@currentsectionname%
    \refstepcounter{thmcounter}%
    % \@thmtag is the mnemonic shown in the heading; every theorem gets one.
    % With an explicit key the key itself is shown.  Without one the key
    % defaults to the title's first word, and showing that whole word would
    % merely repeat the title ("Dissipativity [Dissipativity]"), so only its
    % first three letters are shown: "Dissipativity [Dis]".  The label is
    % unaffected -- it still uses the full \@thmlblkey.
    \IfValueTF{#4}{%
      \edef\@thmlblkey{#4}%
    }{%
      \edef\@thmfirstword{\@firstwordof#2 \relax}%
      \edef\@thmlblkey{%
        \expandafter\@takethree\@thmfirstword\@empty\@empty\@empty\@nil}%
    }%
    \IfBooleanTF{#1}{\let\@thmtag\@empty}{%
      \edef\@thmtag{\noexpand\ \noexpand\textup{[\noexpand\texttt{t:\@thmlblkey}]}}}%
    \def\@tempshort{#3}%
    \ifx\@tempshort\@empty
      \expandafter\xdef\csname @tshort@\@thmlblkey\endcsname{\@thmlblkey}%
    \else
      \expandafter\def\csname @tshort@\@thmlblkey\endcsname{#3}%
    \fi
    \begin{theorem}[#2\ifshowkeys\@thmtag\fi]%
    \label{t:\@thmlblkey}%
    % record the long title for \nref; via the .aux so forward refs resolve
    \protected@write\@auxout{}{%
      \string\@thmtitle{t:\@thmlblkey}{\unexpanded{#2}}}%
    \expanded{%
      \noexpand\newglossaryentry{t\thethmcounter}{%
        name={\unexpanded{#2}%
          \IfBooleanF{#1}{ [\noexpand\texttt{t:\@thmlblkey}]}},%
        % sort/user1 must be the number the document actually prints, i.e. the
        % theorem environment's counter, not thmcounter.  thmcounter advances
        % only on \beT with a title, so a bare \beT (or \beT with an
        % empty title) makes the two drift and the glossary cites a wrong
        % number.  thmcounter is still fine as the entry key: it is unique.
        sort={\thetheorem},%
        description={p.~\noexpand\pageref{t:\@thmlblkey}},%
        user1={\thetheorem},%
        user2={Theorem},%
        user3={\thesection},%
        user4={\expandafter\unexpanded\expandafter{\@savedSecTitle}}%
      }%
    }%
    \protected@write\@auxout{}{\string\glsliveentry{t\thethmcounter}}%
  \fi
}
\newcommand{\eeT}{\end{theorem}}

% CROSS-PART RELATION PARAGRAPH
% Usage: \xpart{text}  -- place immediately BEFORE \beT, never inside it
%   (anything inside \beT...\eeT would be swallowed into the glossary entry).
% The paragraph is typeset only when Part I is actually present, i.e. in the
% combined book build; in the standalone builds (GranadaDSABNS/A.tex,
% AAHRah-Submission/AAHRah.tex) the label part:glv is undefined and the whole
% paragraph disappears rather than leaving ?? cross-references. Same
% graceful-degradation principle as \GASref in GASt.tex. Resolves on the
% second LaTeX pass, since r@part:glv comes from the .aux file.
% (Already inside the \makeatletter region opened above, so no \makeatletter
% pair here -- adding one would close that region early.)
\newcommand{\xpart}[1]{\@ifundefined{r@part:glv}{}{\paragraph{Relation with Parts~I--II.}#1}}

% \bookref{label}{precise text}{fallback text}: prints the precise cross-reference
% when the label resolves (combined book build), and the fallback text otherwise
% (standalone paper builds, where book-only appendix labels are absent). Same
% graceful-degradation role as \GASref in GASt.tex, but usable in any file since it
% lives in this preamble (loaded by every driver) -- unlike \GASref, which is defined in a
% content file and is therefore undefined in standalone builds that omit that file.
\newcommand{\bookref}[3]{\@ifundefined{r@#1}{#3}{#2}}

% \PVpart and \RBpart: the two parts that a partial build may omit.
% They print "Part~N" in any build that contains the part, and a naming phrase
% in a build that does not, so an excerpt shows no ?? .  Each already supplies
% the word "Part", so write "\PVpart develops ...", not "Part~\PVpart".
\newcommand{\PVpart}{\bookref{part:aahrah}%
  {Part~\ref{part:aahrah}}{the Perron--Volterra part of the full book}}
\newcommand{\RBpart}{\bookref{part:recbrowne}%
  {Part~\ref{part:recbrowne}}{the immuno-virology part of the full book}}

% OPEN PROBLEM COMMANDS
% ASSUMPTION COMMANDS (glossaried)
\newcommand{\beAs}[1][]{%
  \def\@temparg{#1}%
  \ifx\@temparg\@empty
    \begin{assumption}%
  \else
    \let\@savedSecTitle\@currentsectionname%
    \refstepcounter{ascounter}%
    \begin{assumption}[#1]%
    \label{as:\theascounter}%
    \expanded{%
      \noexpand\newglossaryentry{as\theascounter}{%
        name={\unexpanded{#1}},%
        sort={\theascounter},%
        description={p.~\noexpand\pageref{as:\theascounter}},%
        user1={\theascounter},%
        user2={Assumption},%
        user3={\thesection},%
        user4={\expandafter\unexpanded\expandafter{\@savedSecTitle}}%
      }%
    }%
    \protected@write\@auxout{}{\string\glsliveentry{as\theascounter}}%
  \fi
}
\newcommand{\eeAs}{\end{assumption}}
% CONJECTURE COMMANDS (glossaried)
\newcommand{\beCj}[1][]{%
  \def\@temparg{#1}%
  \ifx\@temparg\@empty
    \begin{conjecture}%
  \else
    \let\@savedSecTitle\@currentsectionname%
    \refstepcounter{cjcounter}%
    \begin{conjecture}[#1]%
    \label{cj:\thecjcounter}%
    \expanded{%
      \noexpand\newglossaryentry{cj\thecjcounter}{%
        name={\unexpanded{#1}},%
        sort={\thecjcounter},%
        description={p.~\noexpand\pageref{cj:\thecjcounter}},%
        user1={\thecjcounter},%
        user2={Conjecture},%
        user3={\thesection},%
        user4={\expandafter\unexpanded\expandafter{\@savedSecTitle}}%
      }%
    }%
    \protected@write\@auxout{}{\string\glsliveentry{cj\thecjcounter}}%
  \fi
}
\newcommand{\eeCj}{\end{conjecture}}
% LEMMA COMMANDS
% ---------------------------------------------------------------------
% UNIFORM RESULT WRAPPERS -- \beL, \beC, \beP, \beR share \beT's signature
%
%   \beX[long title][short title][key]  ...  \eeX
%
%   long title  = heading (math allowed)
%   short title = shown in the cross-reference; defaults to the key
%   key         = label key.  With a key the label is X:KEY and the heading
%                 carries the mnemonic [key], exactly as for \beT.  Without
%                 one the label falls back to the running counter, which is
%                 what every pre-existing call site in the book relies on.
%
% A bare \beX -- arguments commented out -- opens the plain environment and
% creates NO label and NO glossary entry.  That is the supported way to keep
% an individual result out of the catalogue.
% ---------------------------------------------------------------------
% #1 wrapper name  #2 environment  #3 label prefix  #4 counter  #5 type name
\newcommand{\@mkresult}[5]{%
  % A leading * suppresses the [key] mnemonic in both the heading and the
  % glossary name, while still making the label X:KEY and the entry:
  %     \beX*[long][short][key]
  \expandafter\NewDocumentCommand\csname be#1\endcsname{s O{} O{} o}{%
    \def\@temparg{##2}%
    \ifx\@temparg\@empty
      \begin{#2}%
    \else
      \let\@savedSecTitle\@currentsectionname%
      \refstepcounter{#4}%
      % With a key: mnemonic label X:KEY and the [key] tag in the heading.
      % Without one: fall back to the running counter, exactly as these
      % wrappers behaved before, so every legacy one-argument call site in
      % the book keeps its label.  The counter is also csname-safe, which a
      % title's first word is not -- titles beginning with math would
      % otherwise build \csname \Omega...\endcsname and fail.
      \IfValueTF{##4}{%
        \edef\@relbl{##4}%
        \IfBooleanTF{##1}{\let\@retag\@empty}{%
          \edef\@retag{\noexpand\ \noexpand\textup{[\noexpand\texttt{#3:##4}]}}}%
      }{%
        \edef\@relbl{\arabic{#4}}%
        \let\@retag\@empty
      }%
      \def\@tempshort{##3}%
      \ifx\@tempshort\@empty
        \expandafter\xdef\csname @sshort@#3\@relbl\endcsname{\@relbl}%
      \else
        \expandafter\def\csname @sshort@#3\@relbl\endcsname{##3}%
      \fi
      \begin{#2}[##2\ifshowkeys\@retag\fi]%
      \edef\@tmplbl{#3:\@relbl}\expandafter\label\expandafter{\@tmplbl}%
      \expanded{%
        \noexpand\newglossaryentry{#3\csname the#4\endcsname}{%
          name={\unexpanded{##2}%
            \IfBooleanF{##1}{\IfValueT{##4}{ [\noexpand\texttt{#3:##4}]}}},%
          sort={\csname the#2\endcsname},%
          description={p.~\noexpand\pageref{#3:\@relbl}},%
          user1={\csname the#2\endcsname},%
          user2={#5},%
          user3={\thesection},%
          user4={\expandafter\unexpanded\expandafter{\@savedSecTitle}}%
        }%
      }%
      \protected@write\@auxout{}{\string\glsliveentry{#3\csname the#4\endcsname}}%
      \protected@write\@auxout{}{\string\@thmtitle{\@tmplbl}{\unexpanded{##2}}}%
    \fi
  }%
  \expandafter\newcommand\csname ee#1\endcsname{\end{#2}}%
}
\@mkresult{L}{lemma}{l}{lemcounter}{Lemma}
% REMARK / COROLLARY / PROPOSITION -- same signature as \beL above
\@mkresult{R}{remark}{r}{remcounter}{Remark}
\@mkresult{C}{corollary}{c}{corcounter}{Corollary}
\@mkresult{P}{proposition}{p}{propcounter}{Proposition}
% OPEN-PROBLEM COMMAND -- generated with the same signature as \beT:
%   \beO[long title][short title][key]   ...   \eeO
% label o:KEY, or o:<counter> when no key is given; a bare \beO makes
% neither a label nor a glossary entry.
\@mkresult{O}{open}{o}{opcounter}{Problem}


% Section-level title tracker: updated only by \section, not \subsection.
% \@currentlabelname is overwritten by every \refstepcounter (including
% \subsection), so we save it separately here.
\newcommand{\@currentsectionname}{}
\AtBeginDocument{%
  \let\@epidcrn@orig@refstepcounter\refstepcounter
  \renewcommand{\refstepcounter}[1]{%
    \@epidcrn@orig@refstepcounter{#1}%
    \ifstrequal{#1}{section}{\global\let\@currentsectionname\@currentlabelname}{}%
  }%
}

% TABLE COMMANDS (glossaried) -- consistent naming with \beT, \beD, etc.
\newcommand{\beTa}[1][]{%
  \def\@temparg{#1}%
  \ifx\@temparg\@empty
    \begin{table}[htbp]%
  \else
    \let\@savedSecTitle\@currentsectionname%
    \refstepcounter{tabcounter}%
    \begin{table}[htbp]%
    \label{A:\thetabcounter}%
    \expanded{%
      \noexpand\newglossaryentry{A\thetabcounter}{%
        name={\unexpanded{#1}},%
        sort={\thetabcounter},%
        description={p.~\noexpand\pageref{A:\thetabcounter}},%
        user1={\thetabcounter},%
        user2={Table},%
        user3={\thesection},%
        user4={\expandafter\unexpanded\expandafter{\@savedSecTitle}}%
      }%
    }%
    \protected@write\@auxout{}{\string\glsliveentry{A\thetabcounter}}%
  \fi
}
\newcommand{\eeTa}{\end{table}}

% FIGURE COMMANDS (glossaried) -- consistent with \beT, \beR, \beTa, etc.
% Usage:
%   \beF[Glossary title]...\caption{Visible caption}\eeF  -> glossaried figure
%   \beF...\caption{Visible caption}\eeF                  -> plain [H] figure, no glossary
\newcommand{\beF}[1][]{%
  \def\@temparg{#1}%
  \ifx\@temparg\@empty
    \begin{figure}[H]%
  \else
    \begin{figure}[H]%
    \let\@savedSecTitle\@currentsectionname%
    \refstepcounter{figcounter}%
    \label{fig:\thefigcounter}%
    \expanded{%
      \noexpand\newglossaryentry{Fig\thefigcounter}{%
        name={\unexpanded{#1}},%
        sort={\thefigcounter},%
        description={p.~\noexpand\pageref{fig:\thefigcounter}},%
        user1={\thefigcounter},%
        user2={Figure},%
        user3={\thesection},%
        user4={\expandafter\unexpanded\expandafter{\@savedSecTitle}}%
      }%
    }%
    \protected@write\@auxout{}{\string\glsliveentry{Fig\thefigcounter}}%
  \fi
}
\newcommand{\eeF}{\end{figure}}
\makeatother
% LIVE-ENTRY REGISTER
% \glsaddall prints every entry the glossaries package currently knows, and
% that set includes entries restored from \jobname.glsdefs.  The package
% rewrites .glsdefs from all *currently defined* entries, so anything it
% restored is written straight back out: a deleted section kept reappearing in
% the catalogue for ever, and the .glsdefs had to be erased by hand.  The
% wrappers therefore record every entry they actually create in the .aux, and
% \printboth adds only those.  The .aux is rebuilt from scratch on each run, so
% this is self-healing: a removed result is gone after one further run.
\makeatletter
\newcommand{\@glslive}{}
\newcommand{\@glsaddlive}[1]{\ifglsentryexists{#1}{\glsadd{#1}}{}}
\newcommand{\glsliveentry}[1]{\g@addto@macro\@glslive{\@glsaddlive{#1}}}
\makeatother
% PRINT COMMAND
% No \clearpage here -- the catalogue continues on the current page.  This has
% been re-added by mistake several times; leave it out.
\makeatletter
\newcommand{\printboth}{%
  \@glslive   % add exactly the entries this document created last run
  \printglossary[style=unifiedstyle, title={0 \hspace{.22em} Definitions, Theorems, Examples and Open Problems}]%
}
\makeatother
% Named theorem reference: \tref{word} -> Thm~N(word-or-short)
% Shows stored short title if available, otherwise shows the key itself.
\newcommand{\tref}[1]{%
  Thm~\ref{t:#1}(\ifcsname @tshort@#1\endcsname\csname @tshort@#1\endcsname\else#1\fi)%
}

% Full-title theorem reference: \nref{t:KEY} prints the number and then the
% theorem's long title, so one writes "Theorem~\nref{t:blocks}" and gets
% "Theorem 3 ``Block factorization and siphons''".  Takes the complete label, like
% \ref, not the bare key like \tref.
%
% \beT writes each title to the .aux (see \@thmtitle below), so \nref resolves
% forward references from the second LaTeX pass, exactly as \ref does.  If the
% title is not yet known -- first pass, or a label that is not a \beT theorem
% -- it degrades to the number rather than producing ??.
% These two use @ in their names, and this point in the file is outside the
% \makeatletter region opened earlier (it closes at \makeatother above), so
% the region has to be reopened here.
\makeatletter
\newcommand{\@thmtitle}[2]{\global\@namedef{@tlong@#1}{#2}}
\newcommand{\nref}[1]{%
  \ref{#1}%
  \ifcsname @tlong@#1\endcsname
    \ ``\csname @tlong@#1\endcsname''%
  \fi
}
% \nm{t:KEY}: the stored long title of any \beT/\beL/\beP/\beR/\beC/\beD/\beE result,
% verbatim and without quotes, so item names follow the macro's first argument.
\newcommand{\nm}[1]{%
  \ifcsname @tlong@#1\endcsname
    \csname @tlong@#1\endcsname%
  \else
    \ref{#1}%
  \fi
}
\makeatother




%% ========================================
%% Other THEOREM ENVIRONMENTS
%% ========================================
\makeatletter
\@ifundefined{lemma}      {\newtheorem{lemma}      {Lemma}}      {}
\@ifundefined{proposition}{\newtheorem{proposition}{Proposition}} {}
\@ifundefined{corollary}  {\newtheorem{corollary}  {Corollary}}  {}
\@ifundefined{conjecture} {\newtheorem{conjecture} {Conjecture}}  {}
\@ifundefined{remark}     {\newtheorem{remark}     {Remark}}     {}
\@ifundefined{question}   {\newtheorem{question}   {Question}}   {}
\@ifundefined{assumption} {\newtheorem{assumption} {Assumption}} {}
\@ifundefined{hypothesis} {\newtheorem{hypothesis} {Hypothesis}} {}
\@ifundefined{fact}       {\newtheorem{fact}       {Fact}}       {}
\makeatother

% \citer: cite the relay-graph threshold theorem of AAHrel. If a document has
% \input the chapter/file containing t:relayRF (so the label resolves), point
% straight at it; otherwise fall back to a bare citation of AAHrel. Shared here
% (not per-paper) so it resolves consistently across every book chapter.
% \citeR: same idea, for the first theorem of AAHRah (t:Pstar).
\makeatletter
\providecommand{\citer}{%
  \@ifundefined{r@t:relayRF}%
    {\cite{AAHrel}}%
    {Theorem~\ref{t:relayRF} of \cite{AAHrel}}%
}
\providecommand{\citeR}{%
  \@ifundefined{r@t:Pstar}%
    {\cite{AAHRah}}%
    {Theorem~\ref{t:Pstar} of \cite{AAHRah}}%
}
\makeatother


% Abbreviations
\def\beQ{\begin{question}}\def\eeQ{\end{question}}
\def\beA{\begin{assumption}}\def\eeA{\end{assumption}}


%% ========================================

%% ========================================
%% COLORS & CUSTOM COMMANDS
%% ========================================
\definecolor{funccolor}{RGB}{25,25,112}
\definecolor{desccolor}{RGB}{64,64,64}
\newcommand{\func}[1]{\texttt{\textcolor{funccolor}{#1}}}
\def\eq{=}
\def\mm{-}

%% ========================================
%% MATHCAL SHORTHANDS
%% ========================================
% Project-wide convention: \mX for a capital X is always \mathcal X.
% Defined for the whole alphabet so the convention never fails to expand.
\newcommand{\mA}{\mathcal A}\newcommand{\mB}{\mathcal B}
\newcommand{\mC}{\mathcal C}\newcommand{\mD}{\mathcal D}
\newcommand{\mE}{\mathcal E}\newcommand{\mF}{\mathcal F}
\newcommand{\mG}{\mathcal G}\newcommand{\mH}{\mathcal H}
\newcommand{\mI}{\mathcal I}\newcommand{\mJ}{\mathcal J}
\newcommand{\mK}{\mathcal K}\newcommand{\mL}{\mathcal L}
\newcommand{\mM}{\mathcal M}\newcommand{\mN}{\mathcal N}
\newcommand{\mO}{\mathcal O}\newcommand{\mP}{\mathcal P}
\newcommand{\mQ}{\mathcal Q}\newcommand{\mR}{\mathcal R}
\newcommand{\mS}{\mathcal S}\newcommand{\mT}{\mathcal T}
\newcommand{\mU}{\mathcal U}\newcommand{\mV}{\mathcal V}
\newcommand{\mW}{\mathcal W}\newcommand{\mX}{\mathcal X}
\newcommand{\mY}{\mathcal Y}\newcommand{\mZ}{\mathcal Z}

%% ========================================
%% BASIC MATH ENVIRONMENTS
%% ========================================
\def\bep{\begin{pmatrix}}\def\eep{\end{pmatrix}}
\def\bev{\begin{vmatrix}}\def\eev{\end{vmatrix}}
\def\bea{\begin{eqnarray*}}\def\eea{\end{eqnarray*}}
\def\bc{\begin{cases}}\def\ec{\end{cases}}
\def\BEN{\begin{enumerate}}\def\EEN{\end{enumerate}}
\def\BI{\begin{itemize}}\def\EI{\end{itemize}}

\newcommand{\be}[1]{\begin{equation}\label{#1}}
\newcommand{\ee}{\end{equation}}
\newcommand{\beq}{\begin{eqnarray}}
\def\eeq{\end{eqnarray}}
\def\eqr{\eqref}\def\fr{\frac}\def\lbl{\label}
\def\Lra{\Longrightarrow}\def\non{\nonumber}\def\cNGM{\cite{Diek,Van,Van08}}
\def\Eq{\Leftrightarrow}\def\b{\beta}\def\uph{\red{UP TO HERE}}
%% ========================================
%% PROBABILITY & STATISTICS
%% ========================================
\renewcommand{\P}{\ensuremath{\mathbb{P}}}  % Note: overrides paragraph symbol
\newcommand{\E}{\ensuremath{\mathbb{E}}}
\newcommand{\ind}[1]{\ensuremath{\mathbbm{1}_{\left\{#1\right\}}}}
\DeclareMathOperator{\Cov}{Cov}
\DeclareMathOperator{\TV}{TV}
\newcommand{\rank}{\operatorname{rank}}

%% ========================================
%% CRN OPERATORS (organizations/ERCs/buffering/siphons framework)
%% ========================================
\def\clos{\mathrm{clos}}
\def\Supply{\mathrm{Supply}}
\def\req{\mathrm{req}}
\def\React{\mathrm{React}}
\def\Prod{\mathrm{Prod}}
\def\Gen{\mathrm{Gen}}
\def\In{\mathrm{In}}
\def\Out{\mathrm{Out}}
\def\Act{\mathrm{Act}}
\def\Min{\mathrm{Min}}

%% ========================================
%% LINE BREAKING
%% ========================================
% Last-resort stretch, used only on paragraphs TeX cannot otherwise break:
% a long \texttt{} token, a comma-separated inline math list (commas are not
% breakpoints in math), or a run of \cite/\ref boxes.  Without it those set as
% overfull \hbox and print into the margin; with it they set slightly loose,
% which is the lesser evil.  Does not affect paragraphs that already break.
\emergencystretch=3em

%% ========================================
%% NUMBER SETS
%% ========================================
\newcommand{\R}{\mathbb{R}}  % Real numbers
\newcommand{\C}{\mathbb{C}}  % Complex numbers
\newcommand{\N}{\mathbb{N}}
\newcommand{\Z}{\mathbb{Z}}
% \providecommand, not \newcommand: some papers define these in their own
% preamble after \input{preamble}, and must not clash.
\providecommand{\Q}{\mathbb{Q}}
\providecommand{\res}{\operatorname{res}}   % resultant; used by epiRes.tex
\newcommand{\MN}{\mathbb{N}}
\newcommand{\MZ}{\mathbb{Z}}
\newcommand{\realpos}{\mathbb{R}_{> 0}}
\newcommand{\realnneg}{\mathbb{R}_{\geq 0}}
\newcommand{\naturalpos}{\mathbb{N}_{>0}}
\newcommand{\Rplus}{\mathbb{R}_+}
\newcommand{\Rp}{\mathbb{R}_{\ge 0}}
\def\Rnne{\mathbb{R}_{\ge 0}}
\def\Se{\bar{S}} \def\Ie{\bar{I}}
\newcommand{\Rpos}{\mathbb{R}_{>0}}
\newcommand{\dd}{\,\mathrm d}
\newcommand{\e}{\mathrm e}
%% ========================================
%% GREEK LETTERS (2+ character abbreviations) greek
%% ========================================
\def\al{\alpha}\def\ba{\bff \alpha}
\def\w{\omega}\def\Lgn{\ga+\mu+\de}\def\om{\omega}
\def\de{\delta}\def\De{\Delta}\def\bde{\bff \delta}
\def\ep{\epsilon}\def\eps{\varepsilon}
\def\f{\varphi}  % Kept as requested
\def\ga{\gamma}\def\Ga{\Gamma}\def\G{\Gamma}\def\bg{\bff \gamma}
\def\up{\upsilon}  % NOT the breve accent: no accent commands anywhere.
                   % Named \up, not \u, so the LaTeX breve stays reachable.
\def\la{\lambda}\def\La{\Lambda}
\def\si{\sigma}\def\ka{\kappa}
\def\th{\theta}\def\vt{\vartheta}
\def\ze{\zeta}\def\Gap{\Gamma^{+}} \def\Gam{\Gamma^{-}}


%% ========================================
%% SPECIAL SHORT COMMANDS (kept as requested)
%% ========================================
\def\I{\infty}
\def\T{\widetilde}
\def\H{\hat}
\newcommand{\relint}{\operatorname{relint}}

%% ========================================
%% BOLD/VECTOR NOTATION
%% ========================================
\newcommand{\bff}[1]{{\mbox{\boldmath$#1$}}}
\newcommand\X{\boldsymbol{X}}
\def\x{\boldsymbol{x}}\def\bx{\bff x}\def\xe{\bar{\x}}
\newcommand\y{\boldsymbol{y}}\def\by{\boldsymbol{y}}\def\ye{\bar{\y}}
\def\yb{\y}\def\xb{\x}
\def\xy{\x-\y}
\newcommand\br{\boldsymbol{r}}\def\re{\bar{\r}}
\newcommand\bd{\boldsymbol{d}}
\newcommand\bQ{\boldsymbol{Q}}
\newcommand{\Yb}{{\mathbf{Y}}}
\newcommand{\Ab}{\mathbf{A}}
\newcommand{\zb}{\mathbf{z}}\def\bz{\bff 0}
\newcommand{\Zb}{\mathbf{Z}}
\newcommand{\bb}{\mathbf{b}}
\def\vbe{\vec {\bff \beta}}
\def\vb{\vec {\bff b}}
\def\v1{\vec {\bff 1}}
\def\vn{\vec \de}
\def\vz{\vec 0}
\def\vi{\overset{\rightarrow}{\mathsf i}}
\def\vv{\overset{\rightarrow}{v}}
\def\vr{\; \vec {\mathsf r}}
\newcommand{\va}{\bff a}
\providecommand{\vect}[1]{\boldsymbol{#1}}   % bold vector (moved from Gradostat.tex)
\providecommand{\one}{\mathbf 1}             % all-ones vector (moved from AliSF)
\providecommand{\spr}{s}                     % spectral bound  (moved from AliSF)
\providecommand{\Ke}{K_{\rm eff}}            % effective carrying capacity
% Critical right/left eigenvectors in the Hopf projection method (AAHOSN/New.tex)
\providecommand{\bq}{\mathbf{q}}
\providecommand{\bp}{\mathbf{p}}


%% ========================================
%% MATHCAL & MATHSCR
%% ========================================
% Mathcal
\newcommand\mc[1]{\mathcal{#1}}
\def\mC{{\mathcal C}}\def\mD{\mathcal D}\def\mE{\mathcal{E}}
\def\mF{\mathcal F}\def\mG{{\mathcal G}}\def\mH{{\mathcal H}}
\def\mI{\mathcal I}\def\mL{\mathcal{L}}\def\mM{{\mathcal M}}
\def\mN{{\mathcal N}}\def\mP{\mathcal{P}}\def\mR{{\mathcal R}}
\def\mS{{\mathcal S}}\def\mT{\mathcal{T}}\def\mV{{\mathcal V}}
\def\mK{\mathcal K}\def\mA{{\mathcal A}}\def\mU{{\mathcal U}}
% Symmetric bilinear and trilinear forms of the Taylor expansion, used in the
% first Lyapunov coefficient (AAHOSN/New.tex)
\providecommand{\Bform}{\mathcal{B}}
\providecommand{\Cform}{\mathcal{C}}

% Blackboard bold.  Kept here, not in the drivers: a \providecommand in a
% driver preamble runs before this file and silently wins, so the same name
% ends up meaning different things in different builds.
\def\IC{\mathbb{C}}\def\IN{\mathbb{N}}\def\IQ{\mathbb{Q}}
\def\IR{\mathbb{R}}\def\IZ{\mathbb{Z}}\def\IG{\mathbb{G}}
\def\II{\mathbb{I}}

% Mathscr
\newcommand{\Rs}{{\mathscr R}}
\newcommand{\Ss}{{\mathscr S}}
\newcommand{\Ns}{{\mathscr N}}
\newcommand{\Cs}{\mathscr C}


%% ========================================
%% OPERATORS & FUNCTIONS
%% ========================================
\DeclareMathOperator{\tr}{tr}
\DeclareMathOperator{\diag}{Diag}
\DeclareMathOperator{\Diag}{Diag}
\DeclareMathOperator{\intr}{\mathrm{Int}}
\DeclareMathOperator{\Inv}{Inv}
\DeclareMathOperator{\spec}{spec}
\DeclareMathOperator{\AUC}{AUC}               % area under the curve
\newcommand{\ii}{\mathrm{i}}                  % imaginary unit
\newcommand\op{\operatorname}
\newcommand{\supp}{\op{supp}}
\newcommand{\diff}{\mathop{}\mathopen{}\mathrm{d}}
\newcommand{\pdd}[1]{\frac{\partial}{\partial{#1}}}
\DeclareMathOperator{\cone}{cone}
\DeclareMathOperator{\proj}{proj}

%% ========================================
%% BRACKETS & NORMS
%% ========================================
\providecommand{\pr}[1]{\left(#1\right)}
\providecommand{\pp}[1]{\left[#1\right]}
\newcommand\croc[1]{\left\langle #1\right\rangle}
\newcommand{\norm}[1]{\|#1\|}
\newcommand{\inner}[2]{\langle #1, #2 \rangle}
% Short synonym of \inner, used by AAHOSN/New.tex as $\ip{u}{v}=\bar u^\top v$
\providecommand{\ip}[2]{\langle #1 , #2 \rangle}


%% ========================================
%% FORMATTING & UTILITIES
%% ========================================
\newcommand*{\Scale}[2][4]{\scalebox{#1}{$#2$}}
\newcommand{\te}[1]{\text{#1}}
\newcommand\steq[1]{\stackrel{\text{\rm #1.}}{=}}
\newcommand\blue[1]{\textcolor{blue}{#1}}
\newcommand{\red}{\textcolor[rgb]{1.00,0.00,0.00}}
\newcommand{\green}{\textcolor[rgb]{0.00,1.00,0.00}}

% Figure command
\newcommand{\figu}[3]{%
\begin{figure}[H]
\centering
\includegraphics[scale=#3]{#1}
\caption{#2\label{f:#1}}
\end{figure}
}

% Footnote
\long\def\symbolfootnote[#1]#2{%
\begingroup
\def\thefootnote{\fnsymbol{footnote}}\footnote[#1]{#2}%
\endgroup}
\def\fn{\symbolfootnote}

\newcommand{\dsum}{\displaystyle\sum}

%% ========================================
%% CRN & EPIDEMIOLOGY SPECIFIC
%% ========================================
% CRN notation
\def\SS{\mathcal S}
\def\CC{\mathcal C}
\def\RR{\mathcal R}

% Reproduction numbers
\def\Rto{R_{1}^{E_2}}
\def\Rot{R_{2}^{E_1}}
\def\Roc{R_{1,c}}
\def\Rtc{R_{2,c}}
\def\rd{r_{0}}
\def\sd{s_{0}}
\def\dd{D_{0}}
\def\vd{v_{0}}
\def\zd{z_{0}}
\def\se{\sE}
\def\xd{x_{0}}
%% ========================================
%% SPECIAL VALUES & EQUILIBRIA
%% ========================================
\def\eE{\bar{e}}
\def\wE{\bar{w}}
\def\sE{s_{*}}
\def\Me{\T M_{ee}}
\def\Rnne{\mathbb{R}_{\ge 0}}

\newcommand{\ones}{\mathbf 1}


%% ========================================
%% ABBREVIATIONS (alphabetical by category)
%% ========================================

% -- General mathematical terms --
\def\adm{admissible}\def\admy{admissibility}
%\def\and{antisymmetric}% disabled 2026-09-30: it overrode LaTeX's author separator \and (printed "antisymmetric" between authors)
\def\ass{assumption}
\def\cadlag{c\`adl\`ag}
\def\ch{characteristic polynomial}
\def\chp{characteristic polynomial}
\def\coe{coefficient}
\def\comp{computation}
\def\con{condition}
\def\coo{coordinate}
\def\corr{corresponding}
\def\dec{decomposition}
\def\den{denominator}
\def\dete{deterministic}
\def\eig{eigenvalue}\def\eval{eigenvalue}
\def\eigv{eigenvector}\def\eiv{eigenvector}\def\evec{eigenvector}
\def\elt{element}
\def\enn{essentially nonnegative/positivity preserving}
\def\equ{equation}
\def\etal{et al.}
\def\fe{for example}
\def\fno{from now on}
\def\for{\forall}
\def\form{formula}
\def\Fr{Furthermore, }
\def\frt{furthermore}
\def\fun{function}
\def\gene{generalization}
\def\heu{heuristic}
\def\how{however}\def\How{However, }
\def\ie{i.e. }
\def\im{\item}
\def\isd{is defined by}
\def\itc{it is easy to check that}
\def\itf{it follows that}\def\Itf{It follows that}
\def\ith{it holds that}
\def\itm{it may be checked that}\def\Itm{It may be checked that}
\def\ito{it turns out that}\def\Ito{It turns out that}\def\Itou{It turns out useful to}
\def\ma{mass action}\def\MAR{mass action representation}
\def\mat{matrix}
\def\mbc{may be checked to be}
\def\mbw{may be written as}\def\mubw{must be written as}
\def\mpr{more precisely, }\def\Mp{More precisely, }
\def\nc{non-catalytic}
\def\nne{non-negative}\def\nny{nonnegativity}
\def\no{\nonumber}
\def\noni{non-infectious}
\def\nota{notation}
\def\np{\newpage}
\def\Nt{Note that}
\def\nt{non trivial}\def\nz{non zero}
\def\oth{otherwise}
\def\Oth{On the other hand}
\def\pa{parameter}\def\para{parameter}
\def\pol{polynomial}
\def\prob{probability}
\def\Prf{{\bf  Proof:}}\def\prf{{\bf  Proof:}}
\def\prop{proportional to}
\def\qu{\quad}
\def\QED{\hfill {$\square$}\goodbreak \medskip}
\def\rep{representation}
\def\resp{respectively}
\def\satd{satisfied}\def\satg{satisfying}\def\sats{satisfies}\def\saty{satisfy}
\def\ssec{\subsection}\def\sssec{\subsubsection}
\def\sef{see for example}
\def\sksy{skew-symmetric}
\def\sub{substrate}
\def\sys{system}
\def\var{variable}
\def\wk{well-known}
\def\wl{without loss of generality}\def\wlo{w.l.o.g.}
\def\wmc{We may check that}\def\Wmc{We may check that}
\def\wmr{we may rewrite}
\def\wms{we must show that}
\def\wrt{with respect to}

% -- CRN & Chemistry --
\newcommand\CRN{chemical reaction network}
\newcommand\CRNs{chemical reaction networks}
\newcommand\RN{reaction network}
\def\ACR{absolute concentration robustness}
\def\AMSD{acyclic siphon decomposition}
\def\DAE{differential algebraic equations}
\def\FHJ{Feinberg-Horn-Jackson}
\def\GMAK{generalized mass action kinetics}
\def\MM{Michaelis-Menten}
\def\sm{stoichiometric matrix}\def\Sm{stoichiometric matrix}
\def\WRZD{weakly reversible, zero deficiency}\def\aut{autocatalytic}

% -- Epidemiology --
\def\brn{basic reproduction number}\def\bzn{basic replacement number}
\def\com{compartment}
\def\dfe{disease free equilibrium}\def\DFE{disease free equilibrium}
\def\dwt{dwell times}\def\edt{expected dwell times}
\def\EM{epidemic model}
\def\endp{endemic fixed point}\def\enp{endemic point}\def\EE{endemic point}
\def\ME{mathematical epidemiology}
\def\MEE{eco-epidemiology}\def\EEP{eco-epidemiological}
\def\ece{eco-epidemiological model}
\def\mgm{multi group model}
\def\mtS{multi-strain model}\def\mSt{multi-strain model}
\def\mSi{minimal siphon}
\def\NGM{next generation matrix}
\def\nim{new infections matrix}
\def\qem{quadratic epidemiological model}
\def\repF{invasion functions}
\def\SPF{SIR-PH-FA model}
\def\sus{susceptible}

% -- Graph theory & Linear algebra --
\def\jac{Jacobi matrix}
\def\jin{Jacobian matrix of the invasion vector field with respect to the invasion variables}
\def\jrn{Jacobian matrix of the resident vector field with respect to the resident variables}
\def\jxn{Jacobian matrix of the $\x$ vector field with respect to the $\x$ variables}
\def\Lm{Laplacian matrix}
\def\MTT{matrix tree theorem}
\def\SGT{spectral graph theory}
\def\SMW{Sherman–Morrison–Woodbury}

% -- Stability & Dynamics --
\def\absC{absorbtion CTMC}
\def\BTb{Bogdanov-Takens bifurcation}
\def\CEP{competitive exclusion partition}
\def\CTMC{continuous time Markov chain}
\def\ds{differential system}
\def\fp{fixed point}\def\bfp{boundary fixed point}\def\fbp{fixed boundary point}
\def\fbl{fixed boundary line}\def\fif{forward invariant face}
\def\gas{global asymptotic stability}\def\GAS{globally asymptotically stable}
\def\GASS{\text{GASS}}
\def\Hb{Hopf bifurcation}
\def\Hur{Hurwitz}\def\MH{Metzler-Hurwitz}\def\RH{Routh-Hurwitz}
\def\IS{\text{IS}}
\def\LAS{locally asymptotically stable}
\def\LdS{Lyapunov diagonally stable}
\def\Lf{Lyapunov function}\def\LyV{Lyapunov Volterra type function}\def\RLF{robust Lyapunov function}
\def\LyM{Lyapunov Malkin stable}
\def\PF{Perron-Frobenius}
\def\spfp{strictly positive fixed point}
\def\sta{stability}
\def\STP{sharp threshold property}\def\STT{sharp threshold theorem}

% -- Models (with \varphi explicitly instead of \f) --
\def\AB{$(A,B)$ Arino-Brauer epidemic models}
\def\ABa{$(A,\bbe,\ba,\varphi)$ balanced bilinear model}
\def\ABP{%$(A,\bbe,P,\varphi)$
\\ balanced bilinear model}
\def\AMG{Anderson May Gupta}
\newcommand\BIN{biological interaction network}
\def\Dt{Descartes type}
\def\gLV{generalized Lotka-Volterra model}
\def\IG{\text{invasion graph}}\def\IGMS{interaction graph on minimal siphons}
\def\LCP{\text{linear complementarity problem}}
\def\LV{Lotka-Volterra}\def\LVK{Lotka-Volterra-Kolmogorov}

% -- Splitting & Structure --
\def\reg{regular splitting}\def\regS{{\bf regular splitting}}
\def\RUR{rational univariate representation}
\def\bds{block-diagonal structure}\def\NGMs{next generation matrices}

% -- Miscellaneous --
\def\Ma{{Mathematica}}
\def\Mm{Max–Min}
\def\num{numerator}
\def\rap{rational approximation} \def\DSR{directed species-reaction bipartite graph}\def\PCP{polynomial complementarity problem}\def\rat{rational substitution}
\def\QSS{quasi steady state} \def\RHS{right hand side} \def\LHS{left hand side}
\def\OSN{online social network}\def\VdD{Van den Driessche}
\def\wli{$\omega$-limit} \def\rk{rich kinetics}\def\sJ{symbolic Jacobian}
\def\cV{\cite{Vas,VasSta,BSV,GSV}} \def\spf{semi-parametric semi-functional}
\def\CRW{Capasso-Ruan-Wang} \def\ABP{$(A,\bbe,P,\f)$ balanced bilinear model}
\def\ABa{%$(A,\bbe,\ba,\f)$ \\
rank-one balanced bilinear model}
\def\bbe{B%{\bff \beta}
} \def\vbe{\vec {\bff \beta}
}
\def\vb{\vec {\bff b}}
\newcommand{\Kh}{\hat{K}}          % K with hat, first mention only
\newcommand{\wS}{w^{\sigma}}       % right Perron J^\perp, full block
\newcommand{\piS}{\vec \pi^{\sigma}}    % left Perron J^\perp, full block
\newcommand{\ws}{w^{\sigma}}       % right Perron J^\perp, sub-block
\newcommand{\pis}{\pi^{\sigma}}    % left Perron J^\perp, sub-block
\newcommand{\hwS}{\hat{w}^{\sigma}}   % right Perron K, full block
\newcommand{\hpiS}{\hat{\pi}^{\sigma}} % left Perron K, full block
\newcommand{\hws}{\hat{w}^{\sigma}}   % right Perron K, sub-block
\newcommand{\hpis}{\hat{\pi}^{\sigma}} % left Perron K, sub-block
\newcommand{\twS}{\tilde{w}^{\sigma}}  % right Perron \tilde K, full block
\newcommand{\tpiS}{\tilde{\pi}^{\sigma}} % left Perron \tilde K
\newcommand{\tws}{\tilde{w}^{\sigma}}
\newcommand{\tpis}{\tilde{\pi}^{\sigma}}
\newcommand{\Mat}[2]{\mathbb{R}_{\ge 0}^{\,#1\times #2}}
\def\tJ{J_\sigma^\perp} \def\lPe{left Perron eigenvector}
% \G is \Gamma (defined at :795).  The Volterra function keeps its own name,
% so that \G no longer has two conflicting meanings in this file.
\def\GV{G}  % Volterra G(u)=u-1-\ln u
\def\PVL{Perron-Volterra Lyapunov minimal siphon ansatz}
\newcommand{\istar}[1]{i_{#1}^{*}}
%\newcommand{\sE}{s_E}
\newcommand{\vE}{v_E}
\def\PV{Perron-Volterra candidate Lyapunov function}
\def\GLV{generalized Lotka-Volterra}
\def\MERN{mathematical epidemiology reaction networks}
\def\CMM{contraction metric method}
\newcommand{\sign}{\operatorname{sign}}
\newcommand{\Var}{\operatorname{Var}} 

% ======================================================================
% IDEAS SHARED BETWEEN THE ABSTRACT AND THE MAIN RESULTS SECTION
%
% Write each idea ONCE, in the abstract; the Main Results section recalls it,
% so the two can never drift apart.
%
%   in the abstract        \ideaitem{1}{The first-order reaction terms ...}
%   in Main Results        \ideahead{1}
%
% \ideaitem stores the text globally and typesets
%   \item \hyperlink{idea:1}{<text>}
% \ideahead typesets
%   \item\hypertarget{idea:1}{}\textbf{<the same text>}
%
% The abstract must come before the Main Results input, which is the normal
% order. If a heading has to read differently from the abstract sentence (for
% instance a long abstract item summarised in one line), override it with the
% optional argument: \ideahead[short heading]{6}.
% ======================================================================
\makeatletter
\newcommand{\ideaitem}[2]{%
  \expandafter\gdef\csname @idea@#1\endcsname{#2}%
  \item\hyperlink{idea:#1}{#2}}
\newcommand{\ideahead}[2][]{%
  \item\hypertarget{idea:#2}{}%
  \def\@ideatmp{#1}%
  \ifx\@ideatmp\@empty
    \@ifundefined{@idea@#2}{%
      \PackageWarning{preamble}{Idea #2 is used in Main Results but was never
        defined by \string\ideaitem\space in the abstract}%
      \textbf{[idea #2 undefined]}%
    }{\textbf{\begingroup\def\ideaterm##1##2{##2}\def\idearest##1{}\@nameuse{@idea@#2}\endgroup}}%
  \else
    \textbf{#1}%
  \fi}
% \ideaterm{<hypertarget>}{<words>}: inside an abstract idea the words become an
% invisible link (the document sets \hypersetup{hidelinks}) jumping to the place
% in Main Results that defines them; inside the Main Results heading itself,
% where the target sits a few lines below, \ideahead strips the link and prints
% the words plainly. Mark the destination with \hypertarget{<name>}{} on the line
% before the display, which adds no equation number and shows nothing.
\newcommand{\ideaterm}[2]{\hyperlink{#1}{#2}}
% \idearest{<text>}: the part of an abstract idea after its first sentence; printed in
% the abstract, omitted from the Main Results heading.
\newcommand{\idearest}[1]{#1}
\makeatother

% ======================================================================
% SHORT-PAPER SWITCH
% \shorttrue in a master that assembles only part of the companion paper
% (Barrat/SipArxiv.tex). \nsh{...} drops text that cites results living in the
% omitted sections, so the short build has no dangling references while the
% companion paper keeps the text verbatim.
% ======================================================================
\newif\ifshort \shortfalse
\newcommand{\nsh}[1]{\ifshort\else#1\fi}
% \shortalt{<short build>}{<companion paper>}: use when a sentence needs
% different wording in the two builds, not merely a deletion.
\newcommand{\shortalt}[2]{\ifshort#1\else#2\fi}
% \Jtr, \JtrS{\Sigma}: transverse Jacobian. J^\perp here; SipArxiv.tex redefines them as J_{\mathrm{tra}}, J_{\mathrm{tra},\Sigma}
\newcommand{\Jtr}{J^\perp}
\newcommand{\JtrS}[1]{J_{#1}^\perp}
, and must not clash.
\providecommand{\Q}{\mathbb{Q}}
\providecommand{\res}{\operatorname{res}}   % resultant; used by epiRes.tex
\newcommand{\MN}{\mathbb{N}}
\newcommand{\MZ}{\mathbb{Z}}
\newcommand{\realpos}{\mathbb{R}_{> 0}}
\newcommand{\realnneg}{\mathbb{R}_{\geq 0}}
\newcommand{\naturalpos}{\mathbb{N}_{>0}}
\newcommand{\Rplus}{\mathbb{R}_+}
\newcommand{\Rp}{\mathbb{R}_{\ge 0}}
\def\Rnne{\mathbb{R}_{\ge 0}}
\def\Se{\bar{S}} \def\Ie{\bar{I}}
\newcommand{\Rpos}{\mathbb{R}_{>0}}
\newcommand{\dd}{\,\mathrm d}
\newcommand{\e}{\mathrm e}
%% ========================================
%% GREEK LETTERS (2+ character abbreviations) greek
%% ========================================
\def\al{\alpha}\def\ba{\bff \alpha}
\def\w{\omega}\def\Lgn{\ga+\mu+\de}\def\om{\omega}
\def\de{\delta}\def\De{\Delta}\def\bde{\bff \delta}
\def\ep{\epsilon}\def\eps{\varepsilon}
\def\f{\varphi}  % Kept as requested
\def\ga{\gamma}\def\Ga{\Gamma}\def\G{\Gamma}\def\bg{\bff \gamma}
\def\up{\upsilon}  % NOT the breve accent: no accent commands anywhere.
                   % Named \up, not \u, so the LaTeX breve stays reachable.
\def\la{\lambda}\def\La{\Lambda}
\def\si{\sigma}\def\ka{\kappa}
\def\th{\theta}\def\vt{\vartheta}
\def\ze{\zeta}\def\Gap{\Gamma^{+}} \def\Gam{\Gamma^{-}}


%% ========================================
%% SPECIAL SHORT COMMANDS (kept as requested)
%% ========================================
\def\I{\infty}
\def\T{\widetilde}
\def\H{\hat}
\newcommand{\relint}{\operatorname{relint}}

%% ========================================
%% BOLD/VECTOR NOTATION
%% ========================================
\newcommand{\bff}[1]{{\mbox{\boldmath$#1$}}}
\newcommand\X{\boldsymbol{X}}
\def\x{\boldsymbol{x}}\def\bx{\bff x}\def\xe{\bar{\x}}
\newcommand\y{\boldsymbol{y}}\def\by{\boldsymbol{y}}\def\ye{\bar{\y}}
\def\yb{\y}\def\xb{\x}
\def\xy{\x-\y}
\newcommand\br{\boldsymbol{r}}\def\re{\bar{\r}}
\newcommand\bd{\boldsymbol{d}}
\newcommand\bQ{\boldsymbol{Q}}
\newcommand{\Yb}{{\mathbf{Y}}}
\newcommand{\Ab}{\mathbf{A}}
\newcommand{\zb}{\mathbf{z}}\def\bz{\bff 0}
\newcommand{\Zb}{\mathbf{Z}}
\newcommand{\bb}{\mathbf{b}}
\def\vbe{\vec {\bff \beta}}
\def\vb{\vec {\bff b}}
\def\v1{\vec {\bff 1}}
\def\vn{\vec \de}
\def\vz{\vec 0}
\def\vi{\overset{\rightarrow}{\mathsf i}}
\def\vv{\overset{\rightarrow}{v}}
\def\vr{\; \vec {\mathsf r}}
\newcommand{\va}{\bff a}
\providecommand{\vect}[1]{\boldsymbol{#1}}   % bold vector (moved from Gradostat.tex)
\providecommand{\one}{\mathbf 1}             % all-ones vector (moved from AliSF)
\providecommand{\spr}{s}                     % spectral bound  (moved from AliSF)
\providecommand{\Ke}{K_{\rm eff}}            % effective carrying capacity
% Critical right/left eigenvectors in the Hopf projection method (AAHOSN/New.tex)
\providecommand{\bq}{\mathbf{q}}
\providecommand{\bp}{\mathbf{p}}


%% ========================================
%% MATHCAL & MATHSCR
%% ========================================
% Mathcal
\newcommand\mc[1]{\mathcal{#1}}
\def\mC{{\mathcal C}}\def\mD{\mathcal D}\def\mE{\mathcal{E}}
\def\mF{\mathcal F}\def\mG{{\mathcal G}}\def\mH{{\mathcal H}}
\def\mI{\mathcal I}\def\mL{\mathcal{L}}\def\mM{{\mathcal M}}
\def\mN{{\mathcal N}}\def\mP{\mathcal{P}}\def\mR{{\mathcal R}}
\def\mS{{\mathcal S}}\def\mT{\mathcal{T}}\def\mV{{\mathcal V}}
\def\mK{\mathcal K}\def\mA{{\mathcal A}}\def\mU{{\mathcal U}}
% Symmetric bilinear and trilinear forms of the Taylor expansion, used in the
% first Lyapunov coefficient (AAHOSN/New.tex)
\providecommand{\Bform}{\mathcal{B}}
\providecommand{\Cform}{\mathcal{C}}

% Blackboard bold.  Kept here, not in the drivers: a \providecommand in a
% driver preamble runs before this file and silently wins, so the same name
% ends up meaning different things in different builds.
\def\IC{\mathbb{C}}\def\IN{\mathbb{N}}\def\IQ{\mathbb{Q}}
\def\IR{\mathbb{R}}\def\IZ{\mathbb{Z}}\def\IG{\mathbb{G}}
\def\II{\mathbb{I}}

% Mathscr
\newcommand{\Rs}{{\mathscr R}}
\newcommand{\Ss}{{\mathscr S}}
\newcommand{\Ns}{{\mathscr N}}
\newcommand{\Cs}{\mathscr C}


%% ========================================
%% OPERATORS & FUNCTIONS
%% ========================================
\DeclareMathOperator{\tr}{tr}
\DeclareMathOperator{\diag}{Diag}
\DeclareMathOperator{\Diag}{Diag}
\DeclareMathOperator{\intr}{\mathrm{Int}}
\DeclareMathOperator{\Inv}{Inv}
\DeclareMathOperator{\spec}{spec}
\DeclareMathOperator{\AUC}{AUC}               % area under the curve
\newcommand{\ii}{\mathrm{i}}                  % imaginary unit
\newcommand\op{\operatorname}
\newcommand{\supp}{\op{supp}}
\newcommand{\diff}{\mathop{}\mathopen{}\mathrm{d}}
\newcommand{\pdd}[1]{\frac{\partial}{\partial{#1}}}
\DeclareMathOperator{\cone}{cone}
\DeclareMathOperator{\proj}{proj}

%% ========================================
%% BRACKETS & NORMS
%% ========================================
\providecommand{\pr}[1]{\left(#1\right)}
\providecommand{\pp}[1]{\left[#1\right]}
\newcommand\croc[1]{\left\langle #1\right\rangle}
\newcommand{\norm}[1]{\|#1\|}
\newcommand{\inner}[2]{\langle #1, #2 \rangle}
% Short synonym of \inner, used by AAHOSN/New.tex as $\ip{u}{v}=\bar u^\top v$
\providecommand{\ip}[2]{\langle #1 , #2 \rangle}


%% ========================================
%% FORMATTING & UTILITIES
%% ========================================
\newcommand*{\Scale}[2][4]{\scalebox{#1}{$#2$}}
\newcommand{\te}[1]{\text{#1}}
\newcommand\steq[1]{\stackrel{\text{\rm #1.}}{=}}
\newcommand\blue[1]{\textcolor{blue}{#1}}
\newcommand{\red}{\textcolor[rgb]{1.00,0.00,0.00}}
\newcommand{\green}{\textcolor[rgb]{0.00,1.00,0.00}}

% Figure command
\newcommand{\figu}[3]{%
\begin{figure}[H]
\centering
\includegraphics[scale=#3]{#1}
\caption{#2\label{f:#1}}
\end{figure}
}

% Footnote
\long\def\symbolfootnote[#1]#2{%
\begingroup
\def\thefootnote{\fnsymbol{footnote}}\footnote[#1]{#2}%
\endgroup}
\def\fn{\symbolfootnote}

\newcommand{\dsum}{\displaystyle\sum}

%% ========================================
%% CRN & EPIDEMIOLOGY SPECIFIC
%% ========================================
% CRN notation
\def\SS{\mathcal S}
\def\CC{\mathcal C}
\def\RR{\mathcal R}

% Reproduction numbers
\def\Rto{R_{1}^{E_2}}
\def\Rot{R_{2}^{E_1}}
\def\Roc{R_{1,c}}
\def\Rtc{R_{2,c}}
\def\rd{r_{0}}
\def\sd{s_{0}}
\def\dd{D_{0}}
\def\vd{v_{0}}
\def\zd{z_{0}}
\def\se{\sE}
\def\xd{x_{0}}
%% ========================================
%% SPECIAL VALUES & EQUILIBRIA
%% ========================================
\def\eE{\bar{e}}
\def\wE{\bar{w}}
\def\sE{s_{*}}
\def\Me{\T M_{ee}}
\def\Rnne{\mathbb{R}_{\ge 0}}

\newcommand{\ones}{\mathbf 1}


%% ========================================
%% ABBREVIATIONS (alphabetical by category)
%% ========================================

% -- General mathematical terms --
\def\adm{admissible}\def\admy{admissibility}
%\def\and{antisymmetric}% disabled 2026-09-30: it overrode LaTeX's author separator \and (printed "antisymmetric" between authors)
\def\ass{assumption}
\def\cadlag{c\`adl\`ag}
\def\ch{characteristic polynomial}
\def\chp{characteristic polynomial}
\def\coe{coefficient}
\def\comp{computation}
\def\con{condition}
\def\coo{coordinate}
\def\corr{corresponding}
\def\dec{decomposition}
\def\den{denominator}
\def\dete{deterministic}
\def\eig{eigenvalue}\def\eval{eigenvalue}
\def\eigv{eigenvector}\def\eiv{eigenvector}\def\evec{eigenvector}
\def\elt{element}
\def\enn{essentially nonnegative/positivity preserving}
\def\equ{equation}
\def\etal{et al.}
\def\fe{for example}
\def\fno{from now on}
\def\for{\forall}
\def\form{formula}
\def\Fr{Furthermore, }
\def\frt{furthermore}
\def\fun{function}
\def\gene{generalization}
\def\heu{heuristic}
\def\how{however}\def\How{However, }
\def\ie{i.e. }
\def\im{\item}
\def\isd{is defined by}
\def\itc{it is easy to check that}
\def\itf{it follows that}\def\Itf{It follows that}
\def\ith{it holds that}
\def\itm{it may be checked that}\def\Itm{It may be checked that}
\def\ito{it turns out that}\def\Ito{It turns out that}\def\Itou{It turns out useful to}
\def\ma{mass action}\def\MAR{mass action representation}
\def\mat{matrix}
\def\mbc{may be checked to be}
\def\mbw{may be written as}\def\mubw{must be written as}
\def\mpr{more precisely, }\def\Mp{More precisely, }
\def\nc{non-catalytic}
\def\nne{non-negative}\def\nny{nonnegativity}
\def\no{\nonumber}
\def\noni{non-infectious}
\def\nota{notation}
\def\np{\newpage}
\def\Nt{Note that}
\def\nt{non trivial}\def\nz{non zero}
\def\oth{otherwise}
\def\Oth{On the other hand}
\def\pa{parameter}\def\para{parameter}
\def\pol{polynomial}
\def\prob{probability}
\def\Prf{{\bf  Proof:}}\def\prf{{\bf  Proof:}}
\def\prop{proportional to}
\def\qu{\quad}
\def\QED{\hfill {$\square$}\goodbreak \medskip}
\def\rep{representation}
\def\resp{respectively}
\def\satd{satisfied}\def\satg{satisfying}\def\sats{satisfies}\def\saty{satisfy}
\def\ssec{\subsection}\def\sssec{\subsubsection}
\def\sef{see for example}
\def\sksy{skew-symmetric}
\def\sub{substrate}
\def\sys{system}
\def\var{variable}
\def\wk{well-known}
\def\wl{without loss of generality}\def\wlo{w.l.o.g.}
\def\wmc{We may check that}\def\Wmc{We may check that}
\def\wmr{we may rewrite}
\def\wms{we must show that}
\def\wrt{with respect to}

% -- CRN & Chemistry --
\newcommand\CRN{chemical reaction network}
\newcommand\CRNs{chemical reaction networks}
\newcommand\RN{reaction network}
\def\ACR{absolute concentration robustness}
\def\AMSD{acyclic siphon decomposition}
\def\DAE{differential algebraic equations}
\def\FHJ{Feinberg-Horn-Jackson}
\def\GMAK{generalized mass action kinetics}
\def\MM{Michaelis-Menten}
\def\sm{stoichiometric matrix}\def\Sm{stoichiometric matrix}
\def\WRZD{weakly reversible, zero deficiency}\def\aut{autocatalytic}

% -- Epidemiology --
\def\brn{basic reproduction number}\def\bzn{basic replacement number}
\def\com{compartment}
\def\dfe{disease free equilibrium}\def\DFE{disease free equilibrium}
\def\dwt{dwell times}\def\edt{expected dwell times}
\def\EM{epidemic model}
\def\endp{endemic fixed point}\def\enp{endemic point}\def\EE{endemic point}
\def\ME{mathematical epidemiology}
\def\MEE{eco-epidemiology}\def\EEP{eco-epidemiological}
\def\ece{eco-epidemiological model}
\def\mgm{multi group model}
\def\mtS{multi-strain model}\def\mSt{multi-strain model}
\def\mSi{minimal siphon}
\def\NGM{next generation matrix}
\def\nim{new infections matrix}
\def\qem{quadratic epidemiological model}
\def\repF{invasion functions}
\def\SPF{SIR-PH-FA model}
\def\sus{susceptible}

% -- Graph theory & Linear algebra --
\def\jac{Jacobi matrix}
\def\jin{Jacobian matrix of the invasion vector field with respect to the invasion variables}
\def\jrn{Jacobian matrix of the resident vector field with respect to the resident variables}
\def\jxn{Jacobian matrix of the $\x$ vector field with respect to the $\x$ variables}
\def\Lm{Laplacian matrix}
\def\MTT{matrix tree theorem}
\def\SGT{spectral graph theory}
\def\SMW{Sherman–Morrison–Woodbury}

% -- Stability & Dynamics --
\def\absC{absorbtion CTMC}
\def\BTb{Bogdanov-Takens bifurcation}
\def\CEP{competitive exclusion partition}
\def\CTMC{continuous time Markov chain}
\def\ds{differential system}
\def\fp{fixed point}\def\bfp{boundary fixed point}\def\fbp{fixed boundary point}
\def\fbl{fixed boundary line}\def\fif{forward invariant face}
\def\gas{global asymptotic stability}\def\GAS{globally asymptotically stable}
\def\GASS{\text{GASS}}
\def\Hb{Hopf bifurcation}
\def\Hur{Hurwitz}\def\MH{Metzler-Hurwitz}\def\RH{Routh-Hurwitz}
\def\IS{\text{IS}}
\def\LAS{locally asymptotically stable}
\def\LdS{Lyapunov diagonally stable}
\def\Lf{Lyapunov function}\def\LyV{Lyapunov Volterra type function}\def\RLF{robust Lyapunov function}
\def\LyM{Lyapunov Malkin stable}
\def\PF{Perron-Frobenius}
\def\spfp{strictly positive fixed point}
\def\sta{stability}
\def\STP{sharp threshold property}\def\STT{sharp threshold theorem}

% -- Models (with \varphi explicitly instead of \f) --
\def\AB{$(A,B)$ Arino-Brauer epidemic models}
\def\ABa{$(A,\bbe,\ba,\varphi)$ balanced bilinear model}
\def\ABP{%$(A,\bbe,P,\varphi)$
\\ balanced bilinear model}
\def\AMG{Anderson May Gupta}
\newcommand\BIN{biological interaction network}
\def\Dt{Descartes type}
\def\gLV{generalized Lotka-Volterra model}
\def\IG{\text{invasion graph}}\def\IGMS{interaction graph on minimal siphons}
\def\LCP{\text{linear complementarity problem}}
\def\LV{Lotka-Volterra}\def\LVK{Lotka-Volterra-Kolmogorov}

% -- Splitting & Structure --
\def\reg{regular splitting}\def\regS{{\bf regular splitting}}
\def\RUR{rational univariate representation}
\def\bds{block-diagonal structure}\def\NGMs{next generation matrices}

% -- Miscellaneous --
\def\Ma{{Mathematica}}
\def\Mm{Max–Min}
\def\num{numerator}
\def\rap{rational approximation} \def\DSR{directed species-reaction bipartite graph}\def\PCP{polynomial complementarity problem}\def\rat{rational substitution}
\def\QSS{quasi steady state} \def\RHS{right hand side} \def\LHS{left hand side}
\def\OSN{online social network}\def\VdD{Van den Driessche}
\def\wli{$\omega$-limit} \def\rk{rich kinetics}\def\sJ{symbolic Jacobian}
\def\cV{\cite{Vas,VasSta,BSV,GSV}} \def\spf{semi-parametric semi-functional}
\def\CRW{Capasso-Ruan-Wang} \def\ABP{$(A,\bbe,P,\f)$ balanced bilinear model}
\def\ABa{%$(A,\bbe,\ba,\f)$ \\
rank-one balanced bilinear model}
\def\bbe{B%{\bff \beta}
} \def\vbe{\vec {\bff \beta}
}
\def\vb{\vec {\bff b}}
\newcommand{\Kh}{\hat{K}}          % K with hat, first mention only
\newcommand{\wS}{w^{\sigma}}       % right Perron J^\perp, full block
\newcommand{\piS}{\vec \pi^{\sigma}}    % left Perron J^\perp, full block
\newcommand{\ws}{w^{\sigma}}       % right Perron J^\perp, sub-block
\newcommand{\pis}{\pi^{\sigma}}    % left Perron J^\perp, sub-block
\newcommand{\hwS}{\hat{w}^{\sigma}}   % right Perron K, full block
\newcommand{\hpiS}{\hat{\pi}^{\sigma}} % left Perron K, full block
\newcommand{\hws}{\hat{w}^{\sigma}}   % right Perron K, sub-block
\newcommand{\hpis}{\hat{\pi}^{\sigma}} % left Perron K, sub-block
\newcommand{\twS}{\tilde{w}^{\sigma}}  % right Perron \tilde K, full block
\newcommand{\tpiS}{\tilde{\pi}^{\sigma}} % left Perron \tilde K
\newcommand{\tws}{\tilde{w}^{\sigma}}
\newcommand{\tpis}{\tilde{\pi}^{\sigma}}
\newcommand{\Mat}[2]{\mathbb{R}_{\ge 0}^{\,#1\times #2}}
\def\tJ{J_\sigma^\perp} \def\lPe{left Perron eigenvector}
% \G is \Gamma (defined at :795).  The Volterra function keeps its own name,
% so that \G no longer has two conflicting meanings in this file.
\def\GV{G}  % Volterra G(u)=u-1-\ln u
\def\PVL{Perron-Volterra Lyapunov minimal siphon ansatz}
\newcommand{\istar}[1]{i_{#1}^{*}}
%\newcommand{\sE}{s_E}
\newcommand{\vE}{v_E}
\def\PV{Perron-Volterra candidate Lyapunov function}
\def\GLV{generalized Lotka-Volterra}
\def\MERN{mathematical epidemiology reaction networks}
\def\CMM{contraction metric method}
\newcommand{\sign}{\operatorname{sign}}
\newcommand{\Var}{\operatorname{Var}} 

% ======================================================================
% IDEAS SHARED BETWEEN THE ABSTRACT AND THE MAIN RESULTS SECTION
%
% Write each idea ONCE, in the abstract; the Main Results section recalls it,
% so the two can never drift apart.
%
%   in the abstract        \ideaitem{1}{The first-order reaction terms ...}
%   in Main Results        \ideahead{1}
%
% \ideaitem stores the text globally and typesets
%   \item \hyperlink{idea:1}{<text>}
% \ideahead typesets
%   \item\hypertarget{idea:1}{}\textbf{<the same text>}
%
% The abstract must come before the Main Results input, which is the normal
% order. If a heading has to read differently from the abstract sentence (for
% instance a long abstract item summarised in one line), override it with the
% optional argument: \ideahead[short heading]{6}.
% ======================================================================
\makeatletter
\newcommand{\ideaitem}[2]{%
  \expandafter\gdef\csname @idea@#1\endcsname{#2}%
  \item\hyperlink{idea:#1}{#2}}
\newcommand{\ideahead}[2][]{%
  \item\hypertarget{idea:#2}{}%
  \def\@ideatmp{#1}%
  \ifx\@ideatmp\@empty
    \@ifundefined{@idea@#2}{%
      \PackageWarning{preamble}{Idea #2 is used in Main Results but was never
        defined by \string\ideaitem\space in the abstract}%
      \textbf{[idea #2 undefined]}%
    }{\textbf{\begingroup\def\ideaterm##1##2{##2}\def\idearest##1{}\@nameuse{@idea@#2}\endgroup}}%
  \else
    \textbf{#1}%
  \fi}
% \ideaterm{<hypertarget>}{<words>}: inside an abstract idea the words become an
% invisible link (the document sets \hypersetup{hidelinks}) jumping to the place
% in Main Results that defines them; inside the Main Results heading itself,
% where the target sits a few lines below, \ideahead strips the link and prints
% the words plainly. Mark the destination with \hypertarget{<name>}{} on the line
% before the display, which adds no equation number and shows nothing.
\newcommand{\ideaterm}[2]{\hyperlink{#1}{#2}}
% \idearest{<text>}: the part of an abstract idea after its first sentence; printed in
% the abstract, omitted from the Main Results heading.
\newcommand{\idearest}[1]{#1}
\makeatother

% ======================================================================
% SHORT-PAPER SWITCH
% \shorttrue in a master that assembles only part of the companion paper
% (Barrat/SipArxiv.tex). \nsh{...} drops text that cites results living in the
% omitted sections, so the short build has no dangling references while the
% companion paper keeps the text verbatim.
% ======================================================================
\newif\ifshort \shortfalse
\newcommand{\nsh}[1]{\ifshort\else#1\fi}
% \shortalt{<short build>}{<companion paper>}: use when a sentence needs
% different wording in the two builds, not merely a deletion.
\newcommand{\shortalt}[2]{\ifshort#1\else#2\fi}
% \Jtr, \JtrS{\Sigma}: transverse Jacobian. J^\perp here; SipArxiv.tex redefines them as J_{\mathrm{tra}}, J_{\mathrm{tra},\Sigma}
\newcommand{\Jtr}{J^\perp}
\newcommand{\JtrS}[1]{J_{#1}^\perp}
.
\makeatletter
\AtBeginDocument{\@ifundefined{hypersetup}{}{\hypersetup{linktoc=all}}}
\makeatother
\usepackage{graphicx,float,fancybox}
\usepackage{epsfig,epstopdf,tikz}
\usepackage[sans]{dsfont}
\usetikzlibrary{positioning,arrows,arrows.meta,calc}
% Configure listings for Mathematica code
\lstdefinestyle{mathematica}{
 language=Mathematica,
 basicstyle=\small\ttfamily,
 backgroundcolor=\color{gray!10},
 frame=single,
 breaklines=true,
 commentstyle=\color{green!60!black},
 keywordstyle=\color{blue},
 stringstyle=\color{red},
 showstringspaces=false,
 tabsize=2
}
\usepackage{tcolorbox}
% Colored boxes for key results
\newtcolorbox{keyresult}[1][]{
 colback=blue!5!white,
 colframe=blue!75!black,
 title=#1,
 fonttitle=\bfseries
}
\newtcolorbox{examplebox}[1][]{
 colback=green!5!white,
 colframe=green!75!black,
 title=#1,
 fonttitle=\bfseries
}


\makeatletter
\@ifundefined{theorem}   {\newtheorem{theorem}   {Theorem}}   {}
\@ifundefined{definition}{\newtheorem{definition}{Definition}} {}
\@ifundefined{open}      {\newtheorem{open}      {Problem}}   {}
\@ifundefined{example}   {\newtheorem{example}   {Example}}   {}
\makeatother

% COUNTER DEFINITIONS
\newcounter{tabcounter}
\newcounter{defcounter}
\newcounter{opcounter}
\newcounter{thmcounter}
\newcounter{excounter}
\newcounter{ascounter}
\newcounter{cjcounter}
\newcounter{lemcounter}
\newcounter{remcounter}
\newcounter{corcounter}
\newcounter{propcounter}
\newcounter{figcounter}

% UNIFIED GLOSSARY STYLE
\newglossarystyle{unifiedstyle}{%
  \setglossarystyle{list}%
  \renewcommand*{\glossarypreamble}{\gdef\glslastsec{}}%
  \renewcommand{\glossentry}[2]{%
    \edef\glscursec{\glsentryuseriii{##1}}%
    \ifx\glscursec\glslastsec\else
      \item[]\textbf{\S\,\glscursec\enspace\glsentryuseriv{##1}}%
      \global\let\glslastsec\glscursec
    \fi
    \item[\glsentryuserii{##1} \glsentryuseri{##1}:] ``\glossentryname{##1}'' \glsentrydesc{##1}%
  }%
}

\makeatletter
% First-word extractor (used by \beT, \beD and \beXa for the default key)
% \@firstwordof hello world \relax  ->  hello
\def\@firstwordof#1 #2\relax{#1}
% First three letters, used for the heading mnemonic of a theorem given no
% explicit key.  Padding with \@empty keeps it safe for first words shorter
% than three letters: \@empty vanishes under \edef, whereas \relax would
% survive into the label and is illegal inside \csname (it broke every
% two-letter title, e.g. "ME ...", with "Missing \endcsname inserted").
\def\@takethree#1#2#3#4\@nil{#1#2#3}

% EXAMPLE COMMAND -- same signature and behaviour as \beT and \beD.
% Usage: \beXa*[long title][short title][key]
%   star  = suppresses the printed "[key]" tag after the title (as for \beT*
%   and \beD*); the label and glossary entry are unaffected.
%   key = optional label key override; the label becomes e:KEY, and without a
%   key it is e: plus the first three letters of the title's first word.  A
%   mnemonic label lets an example be cited as "Example~\ref{e:KEY}" instead
%   of by section number, which is what the citing prose needs.
\NewDocumentCommand{\beXa}{s O{} O{} o}{%
  \def\@temparg{#2}%
  \ifx\@temparg\@empty
    \begin{example}%
  \else
    \let\@savedSecTitle\@currentsectionname%
    \refstepcounter{excounter}%
    \IfValueTF{#4}{%
      \edef\@exlblkey{#4}%
    }{%
      \edef\@exfirstword{\@firstwordof#2 \relax}%
      \edef\@exlblkey{%
        \expandafter\@takethree\@exfirstword\@empty\@empty\@empty\@nil}%
    }%
    \def\@tempshort{#3}%
    \ifx\@tempshort\@empty
      \expandafter\xdef\csname @xshort@\@exlblkey\endcsname{\@exlblkey}%
    \else
      \expandafter\def\csname @xshort@\@exlblkey\endcsname{#3}%
    \fi
    \IfBooleanTF{#1}{\let\@extag\@empty}{%
      \edef\@extag{\noexpand\ \noexpand\textup{[\noexpand\texttt{e:\@exlblkey}]}}}%
    \begin{example}[#2\ifshowkeys\@extag\fi]%
    \label{e:\@exlblkey}%
    % record the long title so \nref{e:KEY} prints it, as for theorems
    \protected@write\@auxout{}{%
      \string\@thmtitle{e:\@exlblkey}{\unexpanded{#2}}}%
    \expanded{%
      \noexpand\newglossaryentry{e\theexcounter}{%
        name={\unexpanded{#2}%
          \IfBooleanF{#1}{ [\noexpand\texttt{e:\@exlblkey}]}},%
        sort={\theexample},%
        description={p.~\noexpand\pageref{e:\@exlblkey}},%
        user1={\theexample},%
        user2={Example},%
        user3={\thesection},%
        user4={\expandafter\unexpanded\expandafter{\@savedSecTitle}}%
      }%
    }%
    \protected@write\@auxout{}{\string\glsliveentry{e\theexcounter}}%
  \fi
}
\newcommand{\eeXa}{\end{example}}

% DEFINITION COMMAND -- same signature and behaviour as \beT below.
% Usage: \beD[long title][short title][key]
%   long title  = definition heading (plain text only)
%   short title = shown in \dref{word}; defaults to the key
%   key         = optional label key override; the label becomes d:KEY
% Without a key the label is d: plus the first three letters of the title's
% first word, so it is a mnemonic rather than the definition's number: a
% numeric label would silently change whenever a definition is inserted.
\NewDocumentCommand{\beD}{s O{} O{} o}{%
  \def\@temparg{#2}%
  \ifx\@temparg\@empty
    \begin{definition}%
  \else
    \let\@savedSecTitle\@currentsectionname%
    \refstepcounter{defcounter}%
    \IfValueTF{#4}{%
      \edef\@deflblkey{#4}%
    }{%
      \edef\@deffirstword{\@firstwordof#2 \relax}%
      \edef\@deflblkey{%
        \expandafter\@takethree\@deffirstword\@empty\@empty\@empty\@nil}%
    }%
    \def\@tempshort{#3}%
    \ifx\@tempshort\@empty
      \expandafter\xdef\csname @dshort@\@deflblkey\endcsname{\@deflblkey}%
    \else
      \expandafter\def\csname @dshort@\@deflblkey\endcsname{#3}%
    \fi
    \IfBooleanTF{#1}{\let\@deftag\@empty}{%
      \edef\@deftag{\noexpand\ \noexpand\textup{[\noexpand\texttt{d:\@deflblkey}]}}}%
    \begin{definition}[#2\ifshowkeys\@deftag\fi]%
    \label{d:\@deflblkey}%
    % record the long title so \nref{d:KEY} prints it, as for theorems
    \protected@write\@auxout{}{%
      \string\@thmtitle{d:\@deflblkey}{\unexpanded{#2}}}%
    \expanded{%
      \noexpand\newglossaryentry{d\thedefcounter}{%
        name={\unexpanded{#2}%
          \IfBooleanF{#1}{ [\noexpand\texttt{d:\@deflblkey}]}},%
        sort={\thedefinition},%
        description={p.~\noexpand\pageref{d:\@deflblkey}},%
        user1={\thedefinition},%
        user2={Definition},%
        user3={\thesection},%
        user4={\expandafter\unexpanded\expandafter{\@savedSecTitle}}%
      }%
    }%
    \protected@write\@auxout{}{\string\glsliveentry{d\thedefcounter}}%
  \fi
}
\newcommand{\eeD}{\end{definition}}

% \showkeysfalse in a driver hides the [t:KEY] mnemonics in result headings
% (labels and glossary entries are unchanged); used for submission builds.
\newif\ifshowkeys \showkeystrue
% \commonredtrue in a driver that inputs basicFiles/commonRed: tool.tex then
% cites its reduction results instead of restating them.
\newif\ifcommonred \commonredfalse
% TheorEM COMMAND
% Usage: \beT[long title][short title][key]
%   long title  = theorem heading (plain text only)
%   short title = shown in \tref{word}; may contain math; defaults to first word of long title
%   key         = optional label key override; label becomes t:KEY (avoids first-word collisions)
% Label is t:WORD where WORD = first word of long title (or KEY if given); use \tref{WORD} to reference.
\NewDocumentCommand{\beT}{s O{} O{} o}{%
  \def\@temparg{#2}%
  \ifx\@temparg\@empty
    \begin{theorem}%
  \else
    \let\@savedSecTitle\@currentsectionname%
    \refstepcounter{thmcounter}%
    % \@thmtag is the mnemonic shown in the heading; every theorem gets one.
    % With an explicit key the key itself is shown.  Without one the key
    % defaults to the title's first word, and showing that whole word would
    % merely repeat the title ("Dissipativity [Dissipativity]"), so only its
    % first three letters are shown: "Dissipativity [Dis]".  The label is
    % unaffected -- it still uses the full \@thmlblkey.
    \IfValueTF{#4}{%
      \edef\@thmlblkey{#4}%
    }{%
      \edef\@thmfirstword{\@firstwordof#2 \relax}%
      \edef\@thmlblkey{%
        \expandafter\@takethree\@thmfirstword\@empty\@empty\@empty\@nil}%
    }%
    \IfBooleanTF{#1}{\let\@thmtag\@empty}{%
      \edef\@thmtag{\noexpand\ \noexpand\textup{[\noexpand\texttt{t:\@thmlblkey}]}}}%
    \def\@tempshort{#3}%
    \ifx\@tempshort\@empty
      \expandafter\xdef\csname @tshort@\@thmlblkey\endcsname{\@thmlblkey}%
    \else
      \expandafter\def\csname @tshort@\@thmlblkey\endcsname{#3}%
    \fi
    \begin{theorem}[#2\ifshowkeys\@thmtag\fi]%
    \label{t:\@thmlblkey}%
    % record the long title for \nref; via the .aux so forward refs resolve
    \protected@write\@auxout{}{%
      \string\@thmtitle{t:\@thmlblkey}{\unexpanded{#2}}}%
    \expanded{%
      \noexpand\newglossaryentry{t\thethmcounter}{%
        name={\unexpanded{#2}%
          \IfBooleanF{#1}{ [\noexpand\texttt{t:\@thmlblkey}]}},%
        % sort/user1 must be the number the document actually prints, i.e. the
        % theorem environment's counter, not thmcounter.  thmcounter advances
        % only on \beT with a title, so a bare \beT (or \beT with an
        % empty title) makes the two drift and the glossary cites a wrong
        % number.  thmcounter is still fine as the entry key: it is unique.
        sort={\thetheorem},%
        description={p.~\noexpand\pageref{t:\@thmlblkey}},%
        user1={\thetheorem},%
        user2={Theorem},%
        user3={\thesection},%
        user4={\expandafter\unexpanded\expandafter{\@savedSecTitle}}%
      }%
    }%
    \protected@write\@auxout{}{\string\glsliveentry{t\thethmcounter}}%
  \fi
}
\newcommand{\eeT}{\end{theorem}}

% CROSS-PART RELATION PARAGRAPH
% Usage: \xpart{text}  -- place immediately BEFORE \beT, never inside it
%   (anything inside \beT...\eeT would be swallowed into the glossary entry).
% The paragraph is typeset only when Part I is actually present, i.e. in the
% combined book build; in the standalone builds (GranadaDSABNS/A.tex,
% AAHRah-Submission/AAHRah.tex) the label part:glv is undefined and the whole
% paragraph disappears rather than leaving ?? cross-references. Same
% graceful-degradation principle as \GASref in GASt.tex. Resolves on the
% second LaTeX pass, since r@part:glv comes from the .aux file.
% (Already inside the \makeatletter region opened above, so no \makeatletter
% pair here -- adding one would close that region early.)
\newcommand{\xpart}[1]{\@ifundefined{r@part:glv}{}{\paragraph{Relation with Parts~I--II.}#1}}

% \bookref{label}{precise text}{fallback text}: prints the precise cross-reference
% when the label resolves (combined book build), and the fallback text otherwise
% (standalone paper builds, where book-only appendix labels are absent). Same
% graceful-degradation role as \GASref in GASt.tex, but usable in any file since it
% lives in this preamble (loaded by every driver) -- unlike \GASref, which is defined in a
% content file and is therefore undefined in standalone builds that omit that file.
\newcommand{\bookref}[3]{\@ifundefined{r@#1}{#3}{#2}}

% \PVpart and \RBpart: the two parts that a partial build may omit.
% They print "Part~N" in any build that contains the part, and a naming phrase
% in a build that does not, so an excerpt shows no ?? .  Each already supplies
% the word "Part", so write "\PVpart develops ...", not "Part~\PVpart".
\newcommand{\PVpart}{\bookref{part:aahrah}%
  {Part~\ref{part:aahrah}}{the Perron--Volterra part of the full book}}
\newcommand{\RBpart}{\bookref{part:recbrowne}%
  {Part~\ref{part:recbrowne}}{the immuno-virology part of the full book}}

% OPEN PROBLEM COMMANDS
% ASSUMPTION COMMANDS (glossaried)
\newcommand{\beAs}[1][]{%
  \def\@temparg{#1}%
  \ifx\@temparg\@empty
    \begin{assumption}%
  \else
    \let\@savedSecTitle\@currentsectionname%
    \refstepcounter{ascounter}%
    \begin{assumption}[#1]%
    \label{as:\theascounter}%
    \expanded{%
      \noexpand\newglossaryentry{as\theascounter}{%
        name={\unexpanded{#1}},%
        sort={\theascounter},%
        description={p.~\noexpand\pageref{as:\theascounter}},%
        user1={\theascounter},%
        user2={Assumption},%
        user3={\thesection},%
        user4={\expandafter\unexpanded\expandafter{\@savedSecTitle}}%
      }%
    }%
    \protected@write\@auxout{}{\string\glsliveentry{as\theascounter}}%
  \fi
}
\newcommand{\eeAs}{\end{assumption}}
% CONJECTURE COMMANDS (glossaried)
\newcommand{\beCj}[1][]{%
  \def\@temparg{#1}%
  \ifx\@temparg\@empty
    \begin{conjecture}%
  \else
    \let\@savedSecTitle\@currentsectionname%
    \refstepcounter{cjcounter}%
    \begin{conjecture}[#1]%
    \label{cj:\thecjcounter}%
    \expanded{%
      \noexpand\newglossaryentry{cj\thecjcounter}{%
        name={\unexpanded{#1}},%
        sort={\thecjcounter},%
        description={p.~\noexpand\pageref{cj:\thecjcounter}},%
        user1={\thecjcounter},%
        user2={Conjecture},%
        user3={\thesection},%
        user4={\expandafter\unexpanded\expandafter{\@savedSecTitle}}%
      }%
    }%
    \protected@write\@auxout{}{\string\glsliveentry{cj\thecjcounter}}%
  \fi
}
\newcommand{\eeCj}{\end{conjecture}}
% LEMMA COMMANDS
% ---------------------------------------------------------------------
% UNIFORM RESULT WRAPPERS -- \beL, \beC, \beP, \beR share \beT's signature
%
%   \beX[long title][short title][key]  ...  \eeX
%
%   long title  = heading (math allowed)
%   short title = shown in the cross-reference; defaults to the key
%   key         = label key.  With a key the label is X:KEY and the heading
%                 carries the mnemonic [key], exactly as for \beT.  Without
%                 one the label falls back to the running counter, which is
%                 what every pre-existing call site in the book relies on.
%
% A bare \beX -- arguments commented out -- opens the plain environment and
% creates NO label and NO glossary entry.  That is the supported way to keep
% an individual result out of the catalogue.
% ---------------------------------------------------------------------
% #1 wrapper name  #2 environment  #3 label prefix  #4 counter  #5 type name
\newcommand{\@mkresult}[5]{%
  % A leading * suppresses the [key] mnemonic in both the heading and the
  % glossary name, while still making the label X:KEY and the entry:
  %     \beX*[long][short][key]
  \expandafter\NewDocumentCommand\csname be#1\endcsname{s O{} O{} o}{%
    \def\@temparg{##2}%
    \ifx\@temparg\@empty
      \begin{#2}%
    \else
      \let\@savedSecTitle\@currentsectionname%
      \refstepcounter{#4}%
      % With a key: mnemonic label X:KEY and the [key] tag in the heading.
      % Without one: fall back to the running counter, exactly as these
      % wrappers behaved before, so every legacy one-argument call site in
      % the book keeps its label.  The counter is also csname-safe, which a
      % title's first word is not -- titles beginning with math would
      % otherwise build \csname \Omega...\endcsname and fail.
      \IfValueTF{##4}{%
        \edef\@relbl{##4}%
        \IfBooleanTF{##1}{\let\@retag\@empty}{%
          \edef\@retag{\noexpand\ \noexpand\textup{[\noexpand\texttt{#3:##4}]}}}%
      }{%
        \edef\@relbl{\arabic{#4}}%
        \let\@retag\@empty
      }%
      \def\@tempshort{##3}%
      \ifx\@tempshort\@empty
        \expandafter\xdef\csname @sshort@#3\@relbl\endcsname{\@relbl}%
      \else
        \expandafter\def\csname @sshort@#3\@relbl\endcsname{##3}%
      \fi
      \begin{#2}[##2\ifshowkeys\@retag\fi]%
      \edef\@tmplbl{#3:\@relbl}\expandafter\label\expandafter{\@tmplbl}%
      \expanded{%
        \noexpand\newglossaryentry{#3\csname the#4\endcsname}{%
          name={\unexpanded{##2}%
            \IfBooleanF{##1}{\IfValueT{##4}{ [\noexpand\texttt{#3:##4}]}}},%
          sort={\csname the#2\endcsname},%
          description={p.~\noexpand\pageref{#3:\@relbl}},%
          user1={\csname the#2\endcsname},%
          user2={#5},%
          user3={\thesection},%
          user4={\expandafter\unexpanded\expandafter{\@savedSecTitle}}%
        }%
      }%
      \protected@write\@auxout{}{\string\glsliveentry{#3\csname the#4\endcsname}}%
      \protected@write\@auxout{}{\string\@thmtitle{\@tmplbl}{\unexpanded{##2}}}%
    \fi
  }%
  \expandafter\newcommand\csname ee#1\endcsname{\end{#2}}%
}
\@mkresult{L}{lemma}{l}{lemcounter}{Lemma}
% REMARK / COROLLARY / PROPOSITION -- same signature as \beL above
\@mkresult{R}{remark}{r}{remcounter}{Remark}
\@mkresult{C}{corollary}{c}{corcounter}{Corollary}
\@mkresult{P}{proposition}{p}{propcounter}{Proposition}
% OPEN-PROBLEM COMMAND -- generated with the same signature as \beT:
%   \beO[long title][short title][key]   ...   \eeO
% label o:KEY, or o:<counter> when no key is given; a bare \beO makes
% neither a label nor a glossary entry.
\@mkresult{O}{open}{o}{opcounter}{Problem}


% Section-level title tracker: updated only by \section, not \subsection.
% \@currentlabelname is overwritten by every \refstepcounter (including
% \subsection), so we save it separately here.
\newcommand{\@currentsectionname}{}
\AtBeginDocument{%
  \let\@epidcrn@orig@refstepcounter\refstepcounter
  \renewcommand{\refstepcounter}[1]{%
    \@epidcrn@orig@refstepcounter{#1}%
    \ifstrequal{#1}{section}{\global\let\@currentsectionname\@currentlabelname}{}%
  }%
}

% TABLE COMMANDS (glossaried) -- consistent naming with \beT, \beD, etc.
\newcommand{\beTa}[1][]{%
  \def\@temparg{#1}%
  \ifx\@temparg\@empty
    \begin{table}[htbp]%
  \else
    \let\@savedSecTitle\@currentsectionname%
    \refstepcounter{tabcounter}%
    \begin{table}[htbp]%
    \label{A:\thetabcounter}%
    \expanded{%
      \noexpand\newglossaryentry{A\thetabcounter}{%
        name={\unexpanded{#1}},%
        sort={\thetabcounter},%
        description={p.~\noexpand\pageref{A:\thetabcounter}},%
        user1={\thetabcounter},%
        user2={Table},%
        user3={\thesection},%
        user4={\expandafter\unexpanded\expandafter{\@savedSecTitle}}%
      }%
    }%
    \protected@write\@auxout{}{\string\glsliveentry{A\thetabcounter}}%
  \fi
}
\newcommand{\eeTa}{\end{table}}

% FIGURE COMMANDS (glossaried) -- consistent with \beT, \beR, \beTa, etc.
% Usage:
%   \beF[Glossary title]...\caption{Visible caption}\eeF  -> glossaried figure
%   \beF...\caption{Visible caption}\eeF                  -> plain [H] figure, no glossary
\newcommand{\beF}[1][]{%
  \def\@temparg{#1}%
  \ifx\@temparg\@empty
    \begin{figure}[H]%
  \else
    \begin{figure}[H]%
    \let\@savedSecTitle\@currentsectionname%
    \refstepcounter{figcounter}%
    \label{fig:\thefigcounter}%
    \expanded{%
      \noexpand\newglossaryentry{Fig\thefigcounter}{%
        name={\unexpanded{#1}},%
        sort={\thefigcounter},%
        description={p.~\noexpand\pageref{fig:\thefigcounter}},%
        user1={\thefigcounter},%
        user2={Figure},%
        user3={\thesection},%
        user4={\expandafter\unexpanded\expandafter{\@savedSecTitle}}%
      }%
    }%
    \protected@write\@auxout{}{\string\glsliveentry{Fig\thefigcounter}}%
  \fi
}
\newcommand{\eeF}{\end{figure}}
\makeatother
% LIVE-ENTRY REGISTER
% \glsaddall prints every entry the glossaries package currently knows, and
% that set includes entries restored from \jobname.glsdefs.  The package
% rewrites .glsdefs from all *currently defined* entries, so anything it
% restored is written straight back out: a deleted section kept reappearing in
% the catalogue for ever, and the .glsdefs had to be erased by hand.  The
% wrappers therefore record every entry they actually create in the .aux, and
% \printboth adds only those.  The .aux is rebuilt from scratch on each run, so
% this is self-healing: a removed result is gone after one further run.
\makeatletter
\newcommand{\@glslive}{}
\newcommand{\@glsaddlive}[1]{\ifglsentryexists{#1}{\glsadd{#1}}{}}
\newcommand{\glsliveentry}[1]{\g@addto@macro\@glslive{\@glsaddlive{#1}}}
\makeatother
% PRINT COMMAND
% No \clearpage here -- the catalogue continues on the current page.  This has
% been re-added by mistake several times; leave it out.
\makeatletter
\newcommand{\printboth}{%
  \@glslive   % add exactly the entries this document created last run
  \printglossary[style=unifiedstyle, title={0 \hspace{.22em} Definitions, Theorems, Examples and Open Problems}]%
}
\makeatother
% Named theorem reference: \tref{word} -> Thm~N(word-or-short)
% Shows stored short title if available, otherwise shows the key itself.
\newcommand{\tref}[1]{%
  Thm~\ref{t:#1}(\ifcsname @tshort@#1\endcsname\csname @tshort@#1\endcsname\else#1\fi)%
}

% Full-title theorem reference: \nref{t:KEY} prints the number and then the
% theorem's long title, so one writes "Theorem~\nref{t:blocks}" and gets
% "Theorem 3 ``Block factorization and siphons''".  Takes the complete label, like
% \ref, not the bare key like \tref.
%
% \beT writes each title to the .aux (see \@thmtitle below), so \nref resolves
% forward references from the second LaTeX pass, exactly as \ref does.  If the
% title is not yet known -- first pass, or a label that is not a \beT theorem
% -- it degrades to the number rather than producing ??.
% These two use @ in their names, and this point in the file is outside the
% \makeatletter region opened earlier (it closes at \makeatother above), so
% the region has to be reopened here.
\makeatletter
\newcommand{\@thmtitle}[2]{\global\@namedef{@tlong@#1}{#2}}
\newcommand{\nref}[1]{%
  \ref{#1}%
  \ifcsname @tlong@#1\endcsname
    \ ``\csname @tlong@#1\endcsname''%
  \fi
}
% \nm{t:KEY}: the stored long title of any \beT/\beL/\beP/\beR/\beC/\beD/\beE result,
% verbatim and without quotes, so item names follow the macro's first argument.
\newcommand{\nm}[1]{%
  \ifcsname @tlong@#1\endcsname
    \csname @tlong@#1\endcsname%
  \else
    \ref{#1}%
  \fi
}
\makeatother




%% ========================================
%% Other THEOREM ENVIRONMENTS
%% ========================================
\makeatletter
\@ifundefined{lemma}      {\newtheorem{lemma}      {Lemma}}      {}
\@ifundefined{proposition}{\newtheorem{proposition}{Proposition}} {}
\@ifundefined{corollary}  {\newtheorem{corollary}  {Corollary}}  {}
\@ifundefined{conjecture} {\newtheorem{conjecture} {Conjecture}}  {}
\@ifundefined{remark}     {\newtheorem{remark}     {Remark}}     {}
\@ifundefined{question}   {\newtheorem{question}   {Question}}   {}
\@ifundefined{assumption} {\newtheorem{assumption} {Assumption}} {}
\@ifundefined{hypothesis} {\newtheorem{hypothesis} {Hypothesis}} {}
\@ifundefined{fact}       {\newtheorem{fact}       {Fact}}       {}
\makeatother

% \citer: cite the relay-graph threshold theorem of AAHrel. If a document has
% \input the chapter/file containing t:relayRF (so the label resolves), point
% straight at it; otherwise fall back to a bare citation of AAHrel. Shared here
% (not per-paper) so it resolves consistently across every book chapter.
% \citeR: same idea, for the first theorem of AAHRah (t:Pstar).
\makeatletter
\providecommand{\citer}{%
  \@ifundefined{r@t:relayRF}%
    {\cite{AAHrel}}%
    {Theorem~\ref{t:relayRF} of \cite{AAHrel}}%
}
\providecommand{\citeR}{%
  \@ifundefined{r@t:Pstar}%
    {\cite{AAHRah}}%
    {Theorem~\ref{t:Pstar} of \cite{AAHRah}}%
}
\makeatother


% Abbreviations
\def\beQ{\begin{question}}\def\eeQ{\end{question}}
\def\beA{\begin{assumption}}\def\eeA{\end{assumption}}


%% ========================================

%% ========================================
%% COLORS & CUSTOM COMMANDS
%% ========================================
\definecolor{funccolor}{RGB}{25,25,112}
\definecolor{desccolor}{RGB}{64,64,64}
\newcommand{\func}[1]{\texttt{\textcolor{funccolor}{#1}}}
\def\eq{=}
\def\mm{-}

%% ========================================
%% MATHCAL SHORTHANDS
%% ========================================
% Project-wide convention: \mX for a capital X is always \mathcal X.
% Defined for the whole alphabet so the convention never fails to expand.
\newcommand{\mA}{\mathcal A}\newcommand{\mB}{\mathcal B}
\newcommand{\mC}{\mathcal C}\newcommand{\mD}{\mathcal D}
\newcommand{\mE}{\mathcal E}\newcommand{\mF}{\mathcal F}
\newcommand{\mG}{\mathcal G}\newcommand{\mH}{\mathcal H}
\newcommand{\mI}{\mathcal I}\newcommand{\mJ}{\mathcal J}
\newcommand{\mK}{\mathcal K}\newcommand{\mL}{\mathcal L}
\newcommand{\mM}{\mathcal M}\newcommand{\mN}{\mathcal N}
\newcommand{\mO}{\mathcal O}\newcommand{\mP}{\mathcal P}
\newcommand{\mQ}{\mathcal Q}\newcommand{\mR}{\mathcal R}
\newcommand{\mS}{\mathcal S}\newcommand{\mT}{\mathcal T}
\newcommand{\mU}{\mathcal U}\newcommand{\mV}{\mathcal V}
\newcommand{\mW}{\mathcal W}\newcommand{\mX}{\mathcal X}
\newcommand{\mY}{\mathcal Y}\newcommand{\mZ}{\mathcal Z}

%% ========================================
%% BASIC MATH ENVIRONMENTS
%% ========================================
\def\bep{\begin{pmatrix}}\def\eep{\end{pmatrix}}
\def\bev{\begin{vmatrix}}\def\eev{\end{vmatrix}}
\def\bea{\begin{eqnarray*}}\def\eea{\end{eqnarray*}}
\def\bc{\begin{cases}}\def\ec{\end{cases}}
\def\BEN{\begin{enumerate}}\def\EEN{\end{enumerate}}
\def\BI{\begin{itemize}}\def\EI{\end{itemize}}

\newcommand{\be}[1]{\begin{equation}\label{#1}}
\newcommand{\ee}{\end{equation}}
\newcommand{\beq}{\begin{eqnarray}}
\def\eeq{\end{eqnarray}}
\def\eqr{\eqref}\def\fr{\frac}\def\lbl{\label}
\def\Lra{\Longrightarrow}\def\non{\nonumber}\def\cNGM{\cite{Diek,Van,Van08}}
\def\Eq{\Leftrightarrow}\def\b{\beta}\def\uph{\red{UP TO HERE}}
%% ========================================
%% PROBABILITY & STATISTICS
%% ========================================
\renewcommand{\P}{\ensuremath{\mathbb{P}}}  % Note: overrides paragraph symbol
\newcommand{\E}{\ensuremath{\mathbb{E}}}
\newcommand{\ind}[1]{\ensuremath{\mathbbm{1}_{\left\{#1\right\}}}}
\DeclareMathOperator{\Cov}{Cov}
\DeclareMathOperator{\TV}{TV}
\newcommand{\rank}{\operatorname{rank}}

%% ========================================
%% CRN OPERATORS (organizations/ERCs/buffering/siphons framework)
%% ========================================
\def\clos{\mathrm{clos}}
\def\Supply{\mathrm{Supply}}
\def\req{\mathrm{req}}
\def\React{\mathrm{React}}
\def\Prod{\mathrm{Prod}}
\def\Gen{\mathrm{Gen}}
\def\In{\mathrm{In}}
\def\Out{\mathrm{Out}}
\def\Act{\mathrm{Act}}
\def\Min{\mathrm{Min}}

%% ========================================
%% LINE BREAKING
%% ========================================
% Last-resort stretch, used only on paragraphs TeX cannot otherwise break:
% a long \texttt{} token, a comma-separated inline math list (commas are not
% breakpoints in math), or a run of \cite/\ref boxes.  Without it those set as
% overfull \hbox and print into the margin; with it they set slightly loose,
% which is the lesser evil.  Does not affect paragraphs that already break.
\emergencystretch=3em

%% ========================================
%% NUMBER SETS
%% ========================================
\newcommand{\R}{\mathbb{R}}  % Real numbers
\newcommand{\C}{\mathbb{C}}  % Complex numbers
\newcommand{\N}{\mathbb{N}}
\newcommand{\Z}{\mathbb{Z}}
% \providecommand, not \newcommand: some papers define these in their own
% preamble after % Shared preamble for the GLVbook papers.  This is the single canonical copy;
% a driver one level below GLVbook reaches it as
%
%     \input{../basicFiles/preamble}
%
% What stays in the driver, because it differs from paper to paper:
% \documentclass, font choice, \title, \author, \date, and any macro genuinely
% local to one paper.  Everything else belongs here.  The file is
% self-contained: it does not read any other file.

\usepackage[utf8]{inputenc}
\usepackage[novbox]{pdfsync}   % required: SyncTeX alone does not work here

% Page geometry, set once here for every paper that loads this file.  A driver
% that needs different margins says \geometry{...}; a second
% \usepackage[...]{geometry} raises an option clash and must not be used.
\makeatletter
\@ifundefined{AAdefNnogeometry}{\usepackage[a4paper,margin=2cm]{geometry}}{}
\makeatother

% mathematics
\usepackage{amsmath,amssymb,amsthm,mathtools,bm,mathrsfs,bbm}
\usepackage{centernot}

% floats, tables, lists
\usepackage{graphicx,float}
\makeatletter
\@ifpackageloaded{subfig}{}{\usepackage{subcaption}}
\makeatother
\usepackage{booktabs,longtable,threeparttable,array,multirow}
\usepackage{enumitem}

% colour, code, algorithms
\usepackage{xcolor,listings,tcolorbox}
\makeatletter
\@ifundefined{c@algorithm}{\usepackage{algorithm}}{}
\makeatother
\usepackage{algorithmicx,algpseudocode}

% diagrams
\usepackage{tikz,tikz-cd,pgfplots}
\usetikzlibrary{graphs,decorations.pathmorphing,decorations.markings,
arrows.meta,automata,positioning}
\tikzcdset{every label/.append style = {font = \small}}
\pgfplotsset{compat=1.18}

% chemistry
\usepackage[version=4]{mhchem}

\usepackage{microtype}
\usepackage{lineno}
\usepackage{nomencl}
\makenomenclature

% hyperref first, then bookmark, then glossaries: glossaries must follow
% hyperref or the entry links are built against the wrong anchors.
\makeatletter
\@ifundefined{hyperrefpackage}{%
  \@ifundefined{AAdefNnohyperref}{\usepackage{hyperref}}{}%
}{}
\@ifundefined{hyperrefpackage}{\usepackage{bookmark}}{}
\makeatother
\usepackage[toc,sort=use]{glossaries}
\makeglossaries

% Book/standalone switch.  Material that the book already covers in its
% prelude, but that a standalone paper still has to state, is wrapped as
%     \ifbookbuild\else ... \fi
% Drivers default to standalone; bookM.tex says \bookbuildtrue.
\newif\ifbookbuild
\bookbuildfalse

% Part switches for the gradostat paper.  Both default to true; a driver that
% wants the short submission says \madecfalse and/or \slowfastfalse, and every
% cross-reference into the dropped part disappears with it.
%   \ifmadec     Part III, the Lagasquie--Madec model and its break-even geometry
%   \ifslowfast  Part IV, everything carrying the small parameter \eps
\newif\ifmadec     \madectrue
\newif\ifslowfast  \slowfasttrue

% Part switch for the Barrat/simplicial-contagion project.  Defaults to
% true (the hypernetwork-synchrony part is included); the journal-paper
% driver says \hnetpartfalse to drop it and every cross-reference into it.
\newif\ifhnetpart  \hnetparttrue

%% ========================================
%% PACKAGES
%% ========================================
\RequirePackage{etoolbox}
\usepackage{makecell}
\usepackage[novbox]{pdfsync}% obsolete; SyncTeX is built into modern engines
\usepackage{amssymb, amsfonts,mathtools,bbm,eurosym,graphicx,latexsym, amsmath,amsthm,times,nccmath,url}
\usepackage{listings,verbatim,cases,tabularx}
% hyperref is deliberately NOT in the line above.  Journal classes often
% load it themselves, with a long option list, from a begindocument hook
% (mdpi.cls does: \DeclareHookRule{begindocument}{hyperref}{before}{mdpi}).
% Loading it here first, without options, makes that later load raise
% "Option clash for package hyperref".  So load it only if nobody else
% has, or will:
%   - \@ifpackageloaded  catches a class that loaded it already;
%   - \hyperrefpackage    is mdpi.cls's own deferred loader;
%   - \AAdefNnohyperref   lets any other driver opt out by defining it
%                         before % Shared preamble for the GLVbook papers.  This is the single canonical copy;
% a driver one level below GLVbook reaches it as
%
%     \input{../basicFiles/preamble}
%
% What stays in the driver, because it differs from paper to paper:
% \documentclass, font choice, \title, \author, \date, and any macro genuinely
% local to one paper.  Everything else belongs here.  The file is
% self-contained: it does not read any other file.

\usepackage[utf8]{inputenc}
\usepackage[novbox]{pdfsync}   % required: SyncTeX alone does not work here

% Page geometry, set once here for every paper that loads this file.  A driver
% that needs different margins says \geometry{...}; a second
% \usepackage[...]{geometry} raises an option clash and must not be used.
\makeatletter
\@ifundefined{AAdefNnogeometry}{\usepackage[a4paper,margin=2cm]{geometry}}{}
\makeatother

% mathematics
\usepackage{amsmath,amssymb,amsthm,mathtools,bm,mathrsfs,bbm}
\usepackage{centernot}

% floats, tables, lists
\usepackage{graphicx,float}
\makeatletter
\@ifpackageloaded{subfig}{}{\usepackage{subcaption}}
\makeatother
\usepackage{booktabs,longtable,threeparttable,array,multirow}
\usepackage{enumitem}

% colour, code, algorithms
\usepackage{xcolor,listings,tcolorbox}
\makeatletter
\@ifundefined{c@algorithm}{\usepackage{algorithm}}{}
\makeatother
\usepackage{algorithmicx,algpseudocode}

% diagrams
\usepackage{tikz,tikz-cd,pgfplots}
\usetikzlibrary{graphs,decorations.pathmorphing,decorations.markings,
arrows.meta,automata,positioning}
\tikzcdset{every label/.append style = {font = \small}}
\pgfplotsset{compat=1.18}

% chemistry
\usepackage[version=4]{mhchem}

\usepackage{microtype}
\usepackage{lineno}
\usepackage{nomencl}
\makenomenclature

% hyperref first, then bookmark, then glossaries: glossaries must follow
% hyperref or the entry links are built against the wrong anchors.
\makeatletter
\@ifundefined{hyperrefpackage}{%
  \@ifundefined{AAdefNnohyperref}{\usepackage{hyperref}}{}%
}{}
\@ifundefined{hyperrefpackage}{\usepackage{bookmark}}{}
\makeatother
\usepackage[toc,sort=use]{glossaries}
\makeglossaries

% Book/standalone switch.  Material that the book already covers in its
% prelude, but that a standalone paper still has to state, is wrapped as
%     \ifbookbuild\else ... \fi
% Drivers default to standalone; bookM.tex says \bookbuildtrue.
\newif\ifbookbuild
\bookbuildfalse

% Part switches for the gradostat paper.  Both default to true; a driver that
% wants the short submission says \madecfalse and/or \slowfastfalse, and every
% cross-reference into the dropped part disappears with it.
%   \ifmadec     Part III, the Lagasquie--Madec model and its break-even geometry
%   \ifslowfast  Part IV, everything carrying the small parameter \eps
\newif\ifmadec     \madectrue
\newif\ifslowfast  \slowfasttrue

% Part switch for the Barrat/simplicial-contagion project.  Defaults to
% true (the hypernetwork-synchrony part is included); the journal-paper
% driver says \hnetpartfalse to drop it and every cross-reference into it.
\newif\ifhnetpart  \hnetparttrue

%% ========================================
%% PACKAGES
%% ========================================
\RequirePackage{etoolbox}
\usepackage{makecell}
\usepackage[novbox]{pdfsync}% obsolete; SyncTeX is built into modern engines
\usepackage{amssymb, amsfonts,mathtools,bbm,eurosym,graphicx,latexsym, amsmath,amsthm,times,nccmath,url}
\usepackage{listings,verbatim,cases,tabularx}
% hyperref is deliberately NOT in the line above.  Journal classes often
% load it themselves, with a long option list, from a begindocument hook
% (mdpi.cls does: \DeclareHookRule{begindocument}{hyperref}{before}{mdpi}).
% Loading it here first, without options, makes that later load raise
% "Option clash for package hyperref".  So load it only if nobody else
% has, or will:
%   - \@ifpackageloaded  catches a class that loaded it already;
%   - \hyperrefpackage    is mdpi.cls's own deferred loader;
%   - \AAdefNnohyperref   lets any other driver opt out by defining it
%                         before \input{preamble}.
\makeatletter
\@ifpackageloaded{hyperref}{}{%
  \@ifundefined{hyperrefpackage}{%
    \@ifundefined{AAdefNnohyperref}{\usepackage{hyperref}}{}%
  }{}%
}
\makeatother
% linktoc=all makes a table-of-contents page number a link as well as the
% entry text; hyperref's default, linktoc=section, links only the text.
% Deferred to begin-document and guarded: by the block above hyperref may
% be loaded later (mdpi's own hook) or not at all (\AAdefNnohyperref), and
% \hypersetup must exist before it is called.  A driver that wants the
% default back can say \hypersetup{linktoc=section} after \input{preamble}.
\makeatletter
\AtBeginDocument{\@ifundefined{hypersetup}{}{\hypersetup{linktoc=all}}}
\makeatother
\usepackage{graphicx,float,fancybox}
\usepackage{epsfig,epstopdf,tikz}
\usepackage[sans]{dsfont}
\usetikzlibrary{positioning,arrows,arrows.meta,calc}
% Configure listings for Mathematica code
\lstdefinestyle{mathematica}{
 language=Mathematica,
 basicstyle=\small\ttfamily,
 backgroundcolor=\color{gray!10},
 frame=single,
 breaklines=true,
 commentstyle=\color{green!60!black},
 keywordstyle=\color{blue},
 stringstyle=\color{red},
 showstringspaces=false,
 tabsize=2
}
\usepackage{tcolorbox}
% Colored boxes for key results
\newtcolorbox{keyresult}[1][]{
 colback=blue!5!white,
 colframe=blue!75!black,
 title=#1,
 fonttitle=\bfseries
}
\newtcolorbox{examplebox}[1][]{
 colback=green!5!white,
 colframe=green!75!black,
 title=#1,
 fonttitle=\bfseries
}


\makeatletter
\@ifundefined{theorem}   {\newtheorem{theorem}   {Theorem}}   {}
\@ifundefined{definition}{\newtheorem{definition}{Definition}} {}
\@ifundefined{open}      {\newtheorem{open}      {Problem}}   {}
\@ifundefined{example}   {\newtheorem{example}   {Example}}   {}
\makeatother

% COUNTER DEFINITIONS
\newcounter{tabcounter}
\newcounter{defcounter}
\newcounter{opcounter}
\newcounter{thmcounter}
\newcounter{excounter}
\newcounter{ascounter}
\newcounter{cjcounter}
\newcounter{lemcounter}
\newcounter{remcounter}
\newcounter{corcounter}
\newcounter{propcounter}
\newcounter{figcounter}

% UNIFIED GLOSSARY STYLE
\newglossarystyle{unifiedstyle}{%
  \setglossarystyle{list}%
  \renewcommand*{\glossarypreamble}{\gdef\glslastsec{}}%
  \renewcommand{\glossentry}[2]{%
    \edef\glscursec{\glsentryuseriii{##1}}%
    \ifx\glscursec\glslastsec\else
      \item[]\textbf{\S\,\glscursec\enspace\glsentryuseriv{##1}}%
      \global\let\glslastsec\glscursec
    \fi
    \item[\glsentryuserii{##1} \glsentryuseri{##1}:] ``\glossentryname{##1}'' \glsentrydesc{##1}%
  }%
}

\makeatletter
% First-word extractor (used by \beT, \beD and \beXa for the default key)
% \@firstwordof hello world \relax  ->  hello
\def\@firstwordof#1 #2\relax{#1}
% First three letters, used for the heading mnemonic of a theorem given no
% explicit key.  Padding with \@empty keeps it safe for first words shorter
% than three letters: \@empty vanishes under \edef, whereas \relax would
% survive into the label and is illegal inside \csname (it broke every
% two-letter title, e.g. "ME ...", with "Missing \endcsname inserted").
\def\@takethree#1#2#3#4\@nil{#1#2#3}

% EXAMPLE COMMAND -- same signature and behaviour as \beT and \beD.
% Usage: \beXa*[long title][short title][key]
%   star  = suppresses the printed "[key]" tag after the title (as for \beT*
%   and \beD*); the label and glossary entry are unaffected.
%   key = optional label key override; the label becomes e:KEY, and without a
%   key it is e: plus the first three letters of the title's first word.  A
%   mnemonic label lets an example be cited as "Example~\ref{e:KEY}" instead
%   of by section number, which is what the citing prose needs.
\NewDocumentCommand{\beXa}{s O{} O{} o}{%
  \def\@temparg{#2}%
  \ifx\@temparg\@empty
    \begin{example}%
  \else
    \let\@savedSecTitle\@currentsectionname%
    \refstepcounter{excounter}%
    \IfValueTF{#4}{%
      \edef\@exlblkey{#4}%
    }{%
      \edef\@exfirstword{\@firstwordof#2 \relax}%
      \edef\@exlblkey{%
        \expandafter\@takethree\@exfirstword\@empty\@empty\@empty\@nil}%
    }%
    \def\@tempshort{#3}%
    \ifx\@tempshort\@empty
      \expandafter\xdef\csname @xshort@\@exlblkey\endcsname{\@exlblkey}%
    \else
      \expandafter\def\csname @xshort@\@exlblkey\endcsname{#3}%
    \fi
    \IfBooleanTF{#1}{\let\@extag\@empty}{%
      \edef\@extag{\noexpand\ \noexpand\textup{[\noexpand\texttt{e:\@exlblkey}]}}}%
    \begin{example}[#2\ifshowkeys\@extag\fi]%
    \label{e:\@exlblkey}%
    % record the long title so \nref{e:KEY} prints it, as for theorems
    \protected@write\@auxout{}{%
      \string\@thmtitle{e:\@exlblkey}{\unexpanded{#2}}}%
    \expanded{%
      \noexpand\newglossaryentry{e\theexcounter}{%
        name={\unexpanded{#2}%
          \IfBooleanF{#1}{ [\noexpand\texttt{e:\@exlblkey}]}},%
        sort={\theexample},%
        description={p.~\noexpand\pageref{e:\@exlblkey}},%
        user1={\theexample},%
        user2={Example},%
        user3={\thesection},%
        user4={\expandafter\unexpanded\expandafter{\@savedSecTitle}}%
      }%
    }%
    \protected@write\@auxout{}{\string\glsliveentry{e\theexcounter}}%
  \fi
}
\newcommand{\eeXa}{\end{example}}

% DEFINITION COMMAND -- same signature and behaviour as \beT below.
% Usage: \beD[long title][short title][key]
%   long title  = definition heading (plain text only)
%   short title = shown in \dref{word}; defaults to the key
%   key         = optional label key override; the label becomes d:KEY
% Without a key the label is d: plus the first three letters of the title's
% first word, so it is a mnemonic rather than the definition's number: a
% numeric label would silently change whenever a definition is inserted.
\NewDocumentCommand{\beD}{s O{} O{} o}{%
  \def\@temparg{#2}%
  \ifx\@temparg\@empty
    \begin{definition}%
  \else
    \let\@savedSecTitle\@currentsectionname%
    \refstepcounter{defcounter}%
    \IfValueTF{#4}{%
      \edef\@deflblkey{#4}%
    }{%
      \edef\@deffirstword{\@firstwordof#2 \relax}%
      \edef\@deflblkey{%
        \expandafter\@takethree\@deffirstword\@empty\@empty\@empty\@nil}%
    }%
    \def\@tempshort{#3}%
    \ifx\@tempshort\@empty
      \expandafter\xdef\csname @dshort@\@deflblkey\endcsname{\@deflblkey}%
    \else
      \expandafter\def\csname @dshort@\@deflblkey\endcsname{#3}%
    \fi
    \IfBooleanTF{#1}{\let\@deftag\@empty}{%
      \edef\@deftag{\noexpand\ \noexpand\textup{[\noexpand\texttt{d:\@deflblkey}]}}}%
    \begin{definition}[#2\ifshowkeys\@deftag\fi]%
    \label{d:\@deflblkey}%
    % record the long title so \nref{d:KEY} prints it, as for theorems
    \protected@write\@auxout{}{%
      \string\@thmtitle{d:\@deflblkey}{\unexpanded{#2}}}%
    \expanded{%
      \noexpand\newglossaryentry{d\thedefcounter}{%
        name={\unexpanded{#2}%
          \IfBooleanF{#1}{ [\noexpand\texttt{d:\@deflblkey}]}},%
        sort={\thedefinition},%
        description={p.~\noexpand\pageref{d:\@deflblkey}},%
        user1={\thedefinition},%
        user2={Definition},%
        user3={\thesection},%
        user4={\expandafter\unexpanded\expandafter{\@savedSecTitle}}%
      }%
    }%
    \protected@write\@auxout{}{\string\glsliveentry{d\thedefcounter}}%
  \fi
}
\newcommand{\eeD}{\end{definition}}

% \showkeysfalse in a driver hides the [t:KEY] mnemonics in result headings
% (labels and glossary entries are unchanged); used for submission builds.
\newif\ifshowkeys \showkeystrue
% \commonredtrue in a driver that inputs basicFiles/commonRed: tool.tex then
% cites its reduction results instead of restating them.
\newif\ifcommonred \commonredfalse
% TheorEM COMMAND
% Usage: \beT[long title][short title][key]
%   long title  = theorem heading (plain text only)
%   short title = shown in \tref{word}; may contain math; defaults to first word of long title
%   key         = optional label key override; label becomes t:KEY (avoids first-word collisions)
% Label is t:WORD where WORD = first word of long title (or KEY if given); use \tref{WORD} to reference.
\NewDocumentCommand{\beT}{s O{} O{} o}{%
  \def\@temparg{#2}%
  \ifx\@temparg\@empty
    \begin{theorem}%
  \else
    \let\@savedSecTitle\@currentsectionname%
    \refstepcounter{thmcounter}%
    % \@thmtag is the mnemonic shown in the heading; every theorem gets one.
    % With an explicit key the key itself is shown.  Without one the key
    % defaults to the title's first word, and showing that whole word would
    % merely repeat the title ("Dissipativity [Dissipativity]"), so only its
    % first three letters are shown: "Dissipativity [Dis]".  The label is
    % unaffected -- it still uses the full \@thmlblkey.
    \IfValueTF{#4}{%
      \edef\@thmlblkey{#4}%
    }{%
      \edef\@thmfirstword{\@firstwordof#2 \relax}%
      \edef\@thmlblkey{%
        \expandafter\@takethree\@thmfirstword\@empty\@empty\@empty\@nil}%
    }%
    \IfBooleanTF{#1}{\let\@thmtag\@empty}{%
      \edef\@thmtag{\noexpand\ \noexpand\textup{[\noexpand\texttt{t:\@thmlblkey}]}}}%
    \def\@tempshort{#3}%
    \ifx\@tempshort\@empty
      \expandafter\xdef\csname @tshort@\@thmlblkey\endcsname{\@thmlblkey}%
    \else
      \expandafter\def\csname @tshort@\@thmlblkey\endcsname{#3}%
    \fi
    \begin{theorem}[#2\ifshowkeys\@thmtag\fi]%
    \label{t:\@thmlblkey}%
    % record the long title for \nref; via the .aux so forward refs resolve
    \protected@write\@auxout{}{%
      \string\@thmtitle{t:\@thmlblkey}{\unexpanded{#2}}}%
    \expanded{%
      \noexpand\newglossaryentry{t\thethmcounter}{%
        name={\unexpanded{#2}%
          \IfBooleanF{#1}{ [\noexpand\texttt{t:\@thmlblkey}]}},%
        % sort/user1 must be the number the document actually prints, i.e. the
        % theorem environment's counter, not thmcounter.  thmcounter advances
        % only on \beT with a title, so a bare \beT (or \beT with an
        % empty title) makes the two drift and the glossary cites a wrong
        % number.  thmcounter is still fine as the entry key: it is unique.
        sort={\thetheorem},%
        description={p.~\noexpand\pageref{t:\@thmlblkey}},%
        user1={\thetheorem},%
        user2={Theorem},%
        user3={\thesection},%
        user4={\expandafter\unexpanded\expandafter{\@savedSecTitle}}%
      }%
    }%
    \protected@write\@auxout{}{\string\glsliveentry{t\thethmcounter}}%
  \fi
}
\newcommand{\eeT}{\end{theorem}}

% CROSS-PART RELATION PARAGRAPH
% Usage: \xpart{text}  -- place immediately BEFORE \beT, never inside it
%   (anything inside \beT...\eeT would be swallowed into the glossary entry).
% The paragraph is typeset only when Part I is actually present, i.e. in the
% combined book build; in the standalone builds (GranadaDSABNS/A.tex,
% AAHRah-Submission/AAHRah.tex) the label part:glv is undefined and the whole
% paragraph disappears rather than leaving ?? cross-references. Same
% graceful-degradation principle as \GASref in GASt.tex. Resolves on the
% second LaTeX pass, since r@part:glv comes from the .aux file.
% (Already inside the \makeatletter region opened above, so no \makeatletter
% pair here -- adding one would close that region early.)
\newcommand{\xpart}[1]{\@ifundefined{r@part:glv}{}{\paragraph{Relation with Parts~I--II.}#1}}

% \bookref{label}{precise text}{fallback text}: prints the precise cross-reference
% when the label resolves (combined book build), and the fallback text otherwise
% (standalone paper builds, where book-only appendix labels are absent). Same
% graceful-degradation role as \GASref in GASt.tex, but usable in any file since it
% lives in this preamble (loaded by every driver) -- unlike \GASref, which is defined in a
% content file and is therefore undefined in standalone builds that omit that file.
\newcommand{\bookref}[3]{\@ifundefined{r@#1}{#3}{#2}}

% \PVpart and \RBpart: the two parts that a partial build may omit.
% They print "Part~N" in any build that contains the part, and a naming phrase
% in a build that does not, so an excerpt shows no ?? .  Each already supplies
% the word "Part", so write "\PVpart develops ...", not "Part~\PVpart".
\newcommand{\PVpart}{\bookref{part:aahrah}%
  {Part~\ref{part:aahrah}}{the Perron--Volterra part of the full book}}
\newcommand{\RBpart}{\bookref{part:recbrowne}%
  {Part~\ref{part:recbrowne}}{the immuno-virology part of the full book}}

% OPEN PROBLEM COMMANDS
% ASSUMPTION COMMANDS (glossaried)
\newcommand{\beAs}[1][]{%
  \def\@temparg{#1}%
  \ifx\@temparg\@empty
    \begin{assumption}%
  \else
    \let\@savedSecTitle\@currentsectionname%
    \refstepcounter{ascounter}%
    \begin{assumption}[#1]%
    \label{as:\theascounter}%
    \expanded{%
      \noexpand\newglossaryentry{as\theascounter}{%
        name={\unexpanded{#1}},%
        sort={\theascounter},%
        description={p.~\noexpand\pageref{as:\theascounter}},%
        user1={\theascounter},%
        user2={Assumption},%
        user3={\thesection},%
        user4={\expandafter\unexpanded\expandafter{\@savedSecTitle}}%
      }%
    }%
    \protected@write\@auxout{}{\string\glsliveentry{as\theascounter}}%
  \fi
}
\newcommand{\eeAs}{\end{assumption}}
% CONJECTURE COMMANDS (glossaried)
\newcommand{\beCj}[1][]{%
  \def\@temparg{#1}%
  \ifx\@temparg\@empty
    \begin{conjecture}%
  \else
    \let\@savedSecTitle\@currentsectionname%
    \refstepcounter{cjcounter}%
    \begin{conjecture}[#1]%
    \label{cj:\thecjcounter}%
    \expanded{%
      \noexpand\newglossaryentry{cj\thecjcounter}{%
        name={\unexpanded{#1}},%
        sort={\thecjcounter},%
        description={p.~\noexpand\pageref{cj:\thecjcounter}},%
        user1={\thecjcounter},%
        user2={Conjecture},%
        user3={\thesection},%
        user4={\expandafter\unexpanded\expandafter{\@savedSecTitle}}%
      }%
    }%
    \protected@write\@auxout{}{\string\glsliveentry{cj\thecjcounter}}%
  \fi
}
\newcommand{\eeCj}{\end{conjecture}}
% LEMMA COMMANDS
% ---------------------------------------------------------------------
% UNIFORM RESULT WRAPPERS -- \beL, \beC, \beP, \beR share \beT's signature
%
%   \beX[long title][short title][key]  ...  \eeX
%
%   long title  = heading (math allowed)
%   short title = shown in the cross-reference; defaults to the key
%   key         = label key.  With a key the label is X:KEY and the heading
%                 carries the mnemonic [key], exactly as for \beT.  Without
%                 one the label falls back to the running counter, which is
%                 what every pre-existing call site in the book relies on.
%
% A bare \beX -- arguments commented out -- opens the plain environment and
% creates NO label and NO glossary entry.  That is the supported way to keep
% an individual result out of the catalogue.
% ---------------------------------------------------------------------
% #1 wrapper name  #2 environment  #3 label prefix  #4 counter  #5 type name
\newcommand{\@mkresult}[5]{%
  % A leading * suppresses the [key] mnemonic in both the heading and the
  % glossary name, while still making the label X:KEY and the entry:
  %     \beX*[long][short][key]
  \expandafter\NewDocumentCommand\csname be#1\endcsname{s O{} O{} o}{%
    \def\@temparg{##2}%
    \ifx\@temparg\@empty
      \begin{#2}%
    \else
      \let\@savedSecTitle\@currentsectionname%
      \refstepcounter{#4}%
      % With a key: mnemonic label X:KEY and the [key] tag in the heading.
      % Without one: fall back to the running counter, exactly as these
      % wrappers behaved before, so every legacy one-argument call site in
      % the book keeps its label.  The counter is also csname-safe, which a
      % title's first word is not -- titles beginning with math would
      % otherwise build \csname \Omega...\endcsname and fail.
      \IfValueTF{##4}{%
        \edef\@relbl{##4}%
        \IfBooleanTF{##1}{\let\@retag\@empty}{%
          \edef\@retag{\noexpand\ \noexpand\textup{[\noexpand\texttt{#3:##4}]}}}%
      }{%
        \edef\@relbl{\arabic{#4}}%
        \let\@retag\@empty
      }%
      \def\@tempshort{##3}%
      \ifx\@tempshort\@empty
        \expandafter\xdef\csname @sshort@#3\@relbl\endcsname{\@relbl}%
      \else
        \expandafter\def\csname @sshort@#3\@relbl\endcsname{##3}%
      \fi
      \begin{#2}[##2\ifshowkeys\@retag\fi]%
      \edef\@tmplbl{#3:\@relbl}\expandafter\label\expandafter{\@tmplbl}%
      \expanded{%
        \noexpand\newglossaryentry{#3\csname the#4\endcsname}{%
          name={\unexpanded{##2}%
            \IfBooleanF{##1}{\IfValueT{##4}{ [\noexpand\texttt{#3:##4}]}}},%
          sort={\csname the#2\endcsname},%
          description={p.~\noexpand\pageref{#3:\@relbl}},%
          user1={\csname the#2\endcsname},%
          user2={#5},%
          user3={\thesection},%
          user4={\expandafter\unexpanded\expandafter{\@savedSecTitle}}%
        }%
      }%
      \protected@write\@auxout{}{\string\glsliveentry{#3\csname the#4\endcsname}}%
      \protected@write\@auxout{}{\string\@thmtitle{\@tmplbl}{\unexpanded{##2}}}%
    \fi
  }%
  \expandafter\newcommand\csname ee#1\endcsname{\end{#2}}%
}
\@mkresult{L}{lemma}{l}{lemcounter}{Lemma}
% REMARK / COROLLARY / PROPOSITION -- same signature as \beL above
\@mkresult{R}{remark}{r}{remcounter}{Remark}
\@mkresult{C}{corollary}{c}{corcounter}{Corollary}
\@mkresult{P}{proposition}{p}{propcounter}{Proposition}
% OPEN-PROBLEM COMMAND -- generated with the same signature as \beT:
%   \beO[long title][short title][key]   ...   \eeO
% label o:KEY, or o:<counter> when no key is given; a bare \beO makes
% neither a label nor a glossary entry.
\@mkresult{O}{open}{o}{opcounter}{Problem}


% Section-level title tracker: updated only by \section, not \subsection.
% \@currentlabelname is overwritten by every \refstepcounter (including
% \subsection), so we save it separately here.
\newcommand{\@currentsectionname}{}
\AtBeginDocument{%
  \let\@epidcrn@orig@refstepcounter\refstepcounter
  \renewcommand{\refstepcounter}[1]{%
    \@epidcrn@orig@refstepcounter{#1}%
    \ifstrequal{#1}{section}{\global\let\@currentsectionname\@currentlabelname}{}%
  }%
}

% TABLE COMMANDS (glossaried) -- consistent naming with \beT, \beD, etc.
\newcommand{\beTa}[1][]{%
  \def\@temparg{#1}%
  \ifx\@temparg\@empty
    \begin{table}[htbp]%
  \else
    \let\@savedSecTitle\@currentsectionname%
    \refstepcounter{tabcounter}%
    \begin{table}[htbp]%
    \label{A:\thetabcounter}%
    \expanded{%
      \noexpand\newglossaryentry{A\thetabcounter}{%
        name={\unexpanded{#1}},%
        sort={\thetabcounter},%
        description={p.~\noexpand\pageref{A:\thetabcounter}},%
        user1={\thetabcounter},%
        user2={Table},%
        user3={\thesection},%
        user4={\expandafter\unexpanded\expandafter{\@savedSecTitle}}%
      }%
    }%
    \protected@write\@auxout{}{\string\glsliveentry{A\thetabcounter}}%
  \fi
}
\newcommand{\eeTa}{\end{table}}

% FIGURE COMMANDS (glossaried) -- consistent with \beT, \beR, \beTa, etc.
% Usage:
%   \beF[Glossary title]...\caption{Visible caption}\eeF  -> glossaried figure
%   \beF...\caption{Visible caption}\eeF                  -> plain [H] figure, no glossary
\newcommand{\beF}[1][]{%
  \def\@temparg{#1}%
  \ifx\@temparg\@empty
    \begin{figure}[H]%
  \else
    \begin{figure}[H]%
    \let\@savedSecTitle\@currentsectionname%
    \refstepcounter{figcounter}%
    \label{fig:\thefigcounter}%
    \expanded{%
      \noexpand\newglossaryentry{Fig\thefigcounter}{%
        name={\unexpanded{#1}},%
        sort={\thefigcounter},%
        description={p.~\noexpand\pageref{fig:\thefigcounter}},%
        user1={\thefigcounter},%
        user2={Figure},%
        user3={\thesection},%
        user4={\expandafter\unexpanded\expandafter{\@savedSecTitle}}%
      }%
    }%
    \protected@write\@auxout{}{\string\glsliveentry{Fig\thefigcounter}}%
  \fi
}
\newcommand{\eeF}{\end{figure}}
\makeatother
% LIVE-ENTRY REGISTER
% \glsaddall prints every entry the glossaries package currently knows, and
% that set includes entries restored from \jobname.glsdefs.  The package
% rewrites .glsdefs from all *currently defined* entries, so anything it
% restored is written straight back out: a deleted section kept reappearing in
% the catalogue for ever, and the .glsdefs had to be erased by hand.  The
% wrappers therefore record every entry they actually create in the .aux, and
% \printboth adds only those.  The .aux is rebuilt from scratch on each run, so
% this is self-healing: a removed result is gone after one further run.
\makeatletter
\newcommand{\@glslive}{}
\newcommand{\@glsaddlive}[1]{\ifglsentryexists{#1}{\glsadd{#1}}{}}
\newcommand{\glsliveentry}[1]{\g@addto@macro\@glslive{\@glsaddlive{#1}}}
\makeatother
% PRINT COMMAND
% No \clearpage here -- the catalogue continues on the current page.  This has
% been re-added by mistake several times; leave it out.
\makeatletter
\newcommand{\printboth}{%
  \@glslive   % add exactly the entries this document created last run
  \printglossary[style=unifiedstyle, title={0 \hspace{.22em} Definitions, Theorems, Examples and Open Problems}]%
}
\makeatother
% Named theorem reference: \tref{word} -> Thm~N(word-or-short)
% Shows stored short title if available, otherwise shows the key itself.
\newcommand{\tref}[1]{%
  Thm~\ref{t:#1}(\ifcsname @tshort@#1\endcsname\csname @tshort@#1\endcsname\else#1\fi)%
}

% Full-title theorem reference: \nref{t:KEY} prints the number and then the
% theorem's long title, so one writes "Theorem~\nref{t:blocks}" and gets
% "Theorem 3 ``Block factorization and siphons''".  Takes the complete label, like
% \ref, not the bare key like \tref.
%
% \beT writes each title to the .aux (see \@thmtitle below), so \nref resolves
% forward references from the second LaTeX pass, exactly as \ref does.  If the
% title is not yet known -- first pass, or a label that is not a \beT theorem
% -- it degrades to the number rather than producing ??.
% These two use @ in their names, and this point in the file is outside the
% \makeatletter region opened earlier (it closes at \makeatother above), so
% the region has to be reopened here.
\makeatletter
\newcommand{\@thmtitle}[2]{\global\@namedef{@tlong@#1}{#2}}
\newcommand{\nref}[1]{%
  \ref{#1}%
  \ifcsname @tlong@#1\endcsname
    \ ``\csname @tlong@#1\endcsname''%
  \fi
}
% \nm{t:KEY}: the stored long title of any \beT/\beL/\beP/\beR/\beC/\beD/\beE result,
% verbatim and without quotes, so item names follow the macro's first argument.
\newcommand{\nm}[1]{%
  \ifcsname @tlong@#1\endcsname
    \csname @tlong@#1\endcsname%
  \else
    \ref{#1}%
  \fi
}
\makeatother




%% ========================================
%% Other THEOREM ENVIRONMENTS
%% ========================================
\makeatletter
\@ifundefined{lemma}      {\newtheorem{lemma}      {Lemma}}      {}
\@ifundefined{proposition}{\newtheorem{proposition}{Proposition}} {}
\@ifundefined{corollary}  {\newtheorem{corollary}  {Corollary}}  {}
\@ifundefined{conjecture} {\newtheorem{conjecture} {Conjecture}}  {}
\@ifundefined{remark}     {\newtheorem{remark}     {Remark}}     {}
\@ifundefined{question}   {\newtheorem{question}   {Question}}   {}
\@ifundefined{assumption} {\newtheorem{assumption} {Assumption}} {}
\@ifundefined{hypothesis} {\newtheorem{hypothesis} {Hypothesis}} {}
\@ifundefined{fact}       {\newtheorem{fact}       {Fact}}       {}
\makeatother

% \citer: cite the relay-graph threshold theorem of AAHrel. If a document has
% \input the chapter/file containing t:relayRF (so the label resolves), point
% straight at it; otherwise fall back to a bare citation of AAHrel. Shared here
% (not per-paper) so it resolves consistently across every book chapter.
% \citeR: same idea, for the first theorem of AAHRah (t:Pstar).
\makeatletter
\providecommand{\citer}{%
  \@ifundefined{r@t:relayRF}%
    {\cite{AAHrel}}%
    {Theorem~\ref{t:relayRF} of \cite{AAHrel}}%
}
\providecommand{\citeR}{%
  \@ifundefined{r@t:Pstar}%
    {\cite{AAHRah}}%
    {Theorem~\ref{t:Pstar} of \cite{AAHRah}}%
}
\makeatother


% Abbreviations
\def\beQ{\begin{question}}\def\eeQ{\end{question}}
\def\beA{\begin{assumption}}\def\eeA{\end{assumption}}


%% ========================================

%% ========================================
%% COLORS & CUSTOM COMMANDS
%% ========================================
\definecolor{funccolor}{RGB}{25,25,112}
\definecolor{desccolor}{RGB}{64,64,64}
\newcommand{\func}[1]{\texttt{\textcolor{funccolor}{#1}}}
\def\eq{=}
\def\mm{-}

%% ========================================
%% MATHCAL SHORTHANDS
%% ========================================
% Project-wide convention: \mX for a capital X is always \mathcal X.
% Defined for the whole alphabet so the convention never fails to expand.
\newcommand{\mA}{\mathcal A}\newcommand{\mB}{\mathcal B}
\newcommand{\mC}{\mathcal C}\newcommand{\mD}{\mathcal D}
\newcommand{\mE}{\mathcal E}\newcommand{\mF}{\mathcal F}
\newcommand{\mG}{\mathcal G}\newcommand{\mH}{\mathcal H}
\newcommand{\mI}{\mathcal I}\newcommand{\mJ}{\mathcal J}
\newcommand{\mK}{\mathcal K}\newcommand{\mL}{\mathcal L}
\newcommand{\mM}{\mathcal M}\newcommand{\mN}{\mathcal N}
\newcommand{\mO}{\mathcal O}\newcommand{\mP}{\mathcal P}
\newcommand{\mQ}{\mathcal Q}\newcommand{\mR}{\mathcal R}
\newcommand{\mS}{\mathcal S}\newcommand{\mT}{\mathcal T}
\newcommand{\mU}{\mathcal U}\newcommand{\mV}{\mathcal V}
\newcommand{\mW}{\mathcal W}\newcommand{\mX}{\mathcal X}
\newcommand{\mY}{\mathcal Y}\newcommand{\mZ}{\mathcal Z}

%% ========================================
%% BASIC MATH ENVIRONMENTS
%% ========================================
\def\bep{\begin{pmatrix}}\def\eep{\end{pmatrix}}
\def\bev{\begin{vmatrix}}\def\eev{\end{vmatrix}}
\def\bea{\begin{eqnarray*}}\def\eea{\end{eqnarray*}}
\def\bc{\begin{cases}}\def\ec{\end{cases}}
\def\BEN{\begin{enumerate}}\def\EEN{\end{enumerate}}
\def\BI{\begin{itemize}}\def\EI{\end{itemize}}

\newcommand{\be}[1]{\begin{equation}\label{#1}}
\newcommand{\ee}{\end{equation}}
\newcommand{\beq}{\begin{eqnarray}}
\def\eeq{\end{eqnarray}}
\def\eqr{\eqref}\def\fr{\frac}\def\lbl{\label}
\def\Lra{\Longrightarrow}\def\non{\nonumber}\def\cNGM{\cite{Diek,Van,Van08}}
\def\Eq{\Leftrightarrow}\def\b{\beta}\def\uph{\red{UP TO HERE}}
%% ========================================
%% PROBABILITY & STATISTICS
%% ========================================
\renewcommand{\P}{\ensuremath{\mathbb{P}}}  % Note: overrides paragraph symbol
\newcommand{\E}{\ensuremath{\mathbb{E}}}
\newcommand{\ind}[1]{\ensuremath{\mathbbm{1}_{\left\{#1\right\}}}}
\DeclareMathOperator{\Cov}{Cov}
\DeclareMathOperator{\TV}{TV}
\newcommand{\rank}{\operatorname{rank}}

%% ========================================
%% CRN OPERATORS (organizations/ERCs/buffering/siphons framework)
%% ========================================
\def\clos{\mathrm{clos}}
\def\Supply{\mathrm{Supply}}
\def\req{\mathrm{req}}
\def\React{\mathrm{React}}
\def\Prod{\mathrm{Prod}}
\def\Gen{\mathrm{Gen}}
\def\In{\mathrm{In}}
\def\Out{\mathrm{Out}}
\def\Act{\mathrm{Act}}
\def\Min{\mathrm{Min}}

%% ========================================
%% LINE BREAKING
%% ========================================
% Last-resort stretch, used only on paragraphs TeX cannot otherwise break:
% a long \texttt{} token, a comma-separated inline math list (commas are not
% breakpoints in math), or a run of \cite/\ref boxes.  Without it those set as
% overfull \hbox and print into the margin; with it they set slightly loose,
% which is the lesser evil.  Does not affect paragraphs that already break.
\emergencystretch=3em

%% ========================================
%% NUMBER SETS
%% ========================================
\newcommand{\R}{\mathbb{R}}  % Real numbers
\newcommand{\C}{\mathbb{C}}  % Complex numbers
\newcommand{\N}{\mathbb{N}}
\newcommand{\Z}{\mathbb{Z}}
% \providecommand, not \newcommand: some papers define these in their own
% preamble after \input{preamble}, and must not clash.
\providecommand{\Q}{\mathbb{Q}}
\providecommand{\res}{\operatorname{res}}   % resultant; used by epiRes.tex
\newcommand{\MN}{\mathbb{N}}
\newcommand{\MZ}{\mathbb{Z}}
\newcommand{\realpos}{\mathbb{R}_{> 0}}
\newcommand{\realnneg}{\mathbb{R}_{\geq 0}}
\newcommand{\naturalpos}{\mathbb{N}_{>0}}
\newcommand{\Rplus}{\mathbb{R}_+}
\newcommand{\Rp}{\mathbb{R}_{\ge 0}}
\def\Rnne{\mathbb{R}_{\ge 0}}
\def\Se{\bar{S}} \def\Ie{\bar{I}}
\newcommand{\Rpos}{\mathbb{R}_{>0}}
\newcommand{\dd}{\,\mathrm d}
\newcommand{\e}{\mathrm e}
%% ========================================
%% GREEK LETTERS (2+ character abbreviations) greek
%% ========================================
\def\al{\alpha}\def\ba{\bff \alpha}
\def\w{\omega}\def\Lgn{\ga+\mu+\de}\def\om{\omega}
\def\de{\delta}\def\De{\Delta}\def\bde{\bff \delta}
\def\ep{\epsilon}\def\eps{\varepsilon}
\def\f{\varphi}  % Kept as requested
\def\ga{\gamma}\def\Ga{\Gamma}\def\G{\Gamma}\def\bg{\bff \gamma}
\def\up{\upsilon}  % NOT the breve accent: no accent commands anywhere.
                   % Named \up, not \u, so the LaTeX breve stays reachable.
\def\la{\lambda}\def\La{\Lambda}
\def\si{\sigma}\def\ka{\kappa}
\def\th{\theta}\def\vt{\vartheta}
\def\ze{\zeta}\def\Gap{\Gamma^{+}} \def\Gam{\Gamma^{-}}


%% ========================================
%% SPECIAL SHORT COMMANDS (kept as requested)
%% ========================================
\def\I{\infty}
\def\T{\widetilde}
\def\H{\hat}
\newcommand{\relint}{\operatorname{relint}}

%% ========================================
%% BOLD/VECTOR NOTATION
%% ========================================
\newcommand{\bff}[1]{{\mbox{\boldmath$#1$}}}
\newcommand\X{\boldsymbol{X}}
\def\x{\boldsymbol{x}}\def\bx{\bff x}\def\xe{\bar{\x}}
\newcommand\y{\boldsymbol{y}}\def\by{\boldsymbol{y}}\def\ye{\bar{\y}}
\def\yb{\y}\def\xb{\x}
\def\xy{\x-\y}
\newcommand\br{\boldsymbol{r}}\def\re{\bar{\r}}
\newcommand\bd{\boldsymbol{d}}
\newcommand\bQ{\boldsymbol{Q}}
\newcommand{\Yb}{{\mathbf{Y}}}
\newcommand{\Ab}{\mathbf{A}}
\newcommand{\zb}{\mathbf{z}}\def\bz{\bff 0}
\newcommand{\Zb}{\mathbf{Z}}
\newcommand{\bb}{\mathbf{b}}
\def\vbe{\vec {\bff \beta}}
\def\vb{\vec {\bff b}}
\def\v1{\vec {\bff 1}}
\def\vn{\vec \de}
\def\vz{\vec 0}
\def\vi{\overset{\rightarrow}{\mathsf i}}
\def\vv{\overset{\rightarrow}{v}}
\def\vr{\; \vec {\mathsf r}}
\newcommand{\va}{\bff a}
\providecommand{\vect}[1]{\boldsymbol{#1}}   % bold vector (moved from Gradostat.tex)
\providecommand{\one}{\mathbf 1}             % all-ones vector (moved from AliSF)
\providecommand{\spr}{s}                     % spectral bound  (moved from AliSF)
\providecommand{\Ke}{K_{\rm eff}}            % effective carrying capacity
% Critical right/left eigenvectors in the Hopf projection method (AAHOSN/New.tex)
\providecommand{\bq}{\mathbf{q}}
\providecommand{\bp}{\mathbf{p}}


%% ========================================
%% MATHCAL & MATHSCR
%% ========================================
% Mathcal
\newcommand\mc[1]{\mathcal{#1}}
\def\mC{{\mathcal C}}\def\mD{\mathcal D}\def\mE{\mathcal{E}}
\def\mF{\mathcal F}\def\mG{{\mathcal G}}\def\mH{{\mathcal H}}
\def\mI{\mathcal I}\def\mL{\mathcal{L}}\def\mM{{\mathcal M}}
\def\mN{{\mathcal N}}\def\mP{\mathcal{P}}\def\mR{{\mathcal R}}
\def\mS{{\mathcal S}}\def\mT{\mathcal{T}}\def\mV{{\mathcal V}}
\def\mK{\mathcal K}\def\mA{{\mathcal A}}\def\mU{{\mathcal U}}
% Symmetric bilinear and trilinear forms of the Taylor expansion, used in the
% first Lyapunov coefficient (AAHOSN/New.tex)
\providecommand{\Bform}{\mathcal{B}}
\providecommand{\Cform}{\mathcal{C}}

% Blackboard bold.  Kept here, not in the drivers: a \providecommand in a
% driver preamble runs before this file and silently wins, so the same name
% ends up meaning different things in different builds.
\def\IC{\mathbb{C}}\def\IN{\mathbb{N}}\def\IQ{\mathbb{Q}}
\def\IR{\mathbb{R}}\def\IZ{\mathbb{Z}}\def\IG{\mathbb{G}}
\def\II{\mathbb{I}}

% Mathscr
\newcommand{\Rs}{{\mathscr R}}
\newcommand{\Ss}{{\mathscr S}}
\newcommand{\Ns}{{\mathscr N}}
\newcommand{\Cs}{\mathscr C}


%% ========================================
%% OPERATORS & FUNCTIONS
%% ========================================
\DeclareMathOperator{\tr}{tr}
\DeclareMathOperator{\diag}{Diag}
\DeclareMathOperator{\Diag}{Diag}
\DeclareMathOperator{\intr}{\mathrm{Int}}
\DeclareMathOperator{\Inv}{Inv}
\DeclareMathOperator{\spec}{spec}
\DeclareMathOperator{\AUC}{AUC}               % area under the curve
\newcommand{\ii}{\mathrm{i}}                  % imaginary unit
\newcommand\op{\operatorname}
\newcommand{\supp}{\op{supp}}
\newcommand{\diff}{\mathop{}\mathopen{}\mathrm{d}}
\newcommand{\pdd}[1]{\frac{\partial}{\partial{#1}}}
\DeclareMathOperator{\cone}{cone}
\DeclareMathOperator{\proj}{proj}

%% ========================================
%% BRACKETS & NORMS
%% ========================================
\providecommand{\pr}[1]{\left(#1\right)}
\providecommand{\pp}[1]{\left[#1\right]}
\newcommand\croc[1]{\left\langle #1\right\rangle}
\newcommand{\norm}[1]{\|#1\|}
\newcommand{\inner}[2]{\langle #1, #2 \rangle}
% Short synonym of \inner, used by AAHOSN/New.tex as $\ip{u}{v}=\bar u^\top v$
\providecommand{\ip}[2]{\langle #1 , #2 \rangle}


%% ========================================
%% FORMATTING & UTILITIES
%% ========================================
\newcommand*{\Scale}[2][4]{\scalebox{#1}{$#2$}}
\newcommand{\te}[1]{\text{#1}}
\newcommand\steq[1]{\stackrel{\text{\rm #1.}}{=}}
\newcommand\blue[1]{\textcolor{blue}{#1}}
\newcommand{\red}{\textcolor[rgb]{1.00,0.00,0.00}}
\newcommand{\green}{\textcolor[rgb]{0.00,1.00,0.00}}

% Figure command
\newcommand{\figu}[3]{%
\begin{figure}[H]
\centering
\includegraphics[scale=#3]{#1}
\caption{#2\label{f:#1}}
\end{figure}
}

% Footnote
\long\def\symbolfootnote[#1]#2{%
\begingroup
\def\thefootnote{\fnsymbol{footnote}}\footnote[#1]{#2}%
\endgroup}
\def\fn{\symbolfootnote}

\newcommand{\dsum}{\displaystyle\sum}

%% ========================================
%% CRN & EPIDEMIOLOGY SPECIFIC
%% ========================================
% CRN notation
\def\SS{\mathcal S}
\def\CC{\mathcal C}
\def\RR{\mathcal R}

% Reproduction numbers
\def\Rto{R_{1}^{E_2}}
\def\Rot{R_{2}^{E_1}}
\def\Roc{R_{1,c}}
\def\Rtc{R_{2,c}}
\def\rd{r_{0}}
\def\sd{s_{0}}
\def\dd{D_{0}}
\def\vd{v_{0}}
\def\zd{z_{0}}
\def\se{\sE}
\def\xd{x_{0}}
%% ========================================
%% SPECIAL VALUES & EQUILIBRIA
%% ========================================
\def\eE{\bar{e}}
\def\wE{\bar{w}}
\def\sE{s_{*}}
\def\Me{\T M_{ee}}
\def\Rnne{\mathbb{R}_{\ge 0}}

\newcommand{\ones}{\mathbf 1}


%% ========================================
%% ABBREVIATIONS (alphabetical by category)
%% ========================================

% -- General mathematical terms --
\def\adm{admissible}\def\admy{admissibility}
%\def\and{antisymmetric}% disabled 2026-09-30: it overrode LaTeX's author separator \and (printed "antisymmetric" between authors)
\def\ass{assumption}
\def\cadlag{c\`adl\`ag}
\def\ch{characteristic polynomial}
\def\chp{characteristic polynomial}
\def\coe{coefficient}
\def\comp{computation}
\def\con{condition}
\def\coo{coordinate}
\def\corr{corresponding}
\def\dec{decomposition}
\def\den{denominator}
\def\dete{deterministic}
\def\eig{eigenvalue}\def\eval{eigenvalue}
\def\eigv{eigenvector}\def\eiv{eigenvector}\def\evec{eigenvector}
\def\elt{element}
\def\enn{essentially nonnegative/positivity preserving}
\def\equ{equation}
\def\etal{et al.}
\def\fe{for example}
\def\fno{from now on}
\def\for{\forall}
\def\form{formula}
\def\Fr{Furthermore, }
\def\frt{furthermore}
\def\fun{function}
\def\gene{generalization}
\def\heu{heuristic}
\def\how{however}\def\How{However, }
\def\ie{i.e. }
\def\im{\item}
\def\isd{is defined by}
\def\itc{it is easy to check that}
\def\itf{it follows that}\def\Itf{It follows that}
\def\ith{it holds that}
\def\itm{it may be checked that}\def\Itm{It may be checked that}
\def\ito{it turns out that}\def\Ito{It turns out that}\def\Itou{It turns out useful to}
\def\ma{mass action}\def\MAR{mass action representation}
\def\mat{matrix}
\def\mbc{may be checked to be}
\def\mbw{may be written as}\def\mubw{must be written as}
\def\mpr{more precisely, }\def\Mp{More precisely, }
\def\nc{non-catalytic}
\def\nne{non-negative}\def\nny{nonnegativity}
\def\no{\nonumber}
\def\noni{non-infectious}
\def\nota{notation}
\def\np{\newpage}
\def\Nt{Note that}
\def\nt{non trivial}\def\nz{non zero}
\def\oth{otherwise}
\def\Oth{On the other hand}
\def\pa{parameter}\def\para{parameter}
\def\pol{polynomial}
\def\prob{probability}
\def\Prf{{\bf  Proof:}}\def\prf{{\bf  Proof:}}
\def\prop{proportional to}
\def\qu{\quad}
\def\QED{\hfill {$\square$}\goodbreak \medskip}
\def\rep{representation}
\def\resp{respectively}
\def\satd{satisfied}\def\satg{satisfying}\def\sats{satisfies}\def\saty{satisfy}
\def\ssec{\subsection}\def\sssec{\subsubsection}
\def\sef{see for example}
\def\sksy{skew-symmetric}
\def\sub{substrate}
\def\sys{system}
\def\var{variable}
\def\wk{well-known}
\def\wl{without loss of generality}\def\wlo{w.l.o.g.}
\def\wmc{We may check that}\def\Wmc{We may check that}
\def\wmr{we may rewrite}
\def\wms{we must show that}
\def\wrt{with respect to}

% -- CRN & Chemistry --
\newcommand\CRN{chemical reaction network}
\newcommand\CRNs{chemical reaction networks}
\newcommand\RN{reaction network}
\def\ACR{absolute concentration robustness}
\def\AMSD{acyclic siphon decomposition}
\def\DAE{differential algebraic equations}
\def\FHJ{Feinberg-Horn-Jackson}
\def\GMAK{generalized mass action kinetics}
\def\MM{Michaelis-Menten}
\def\sm{stoichiometric matrix}\def\Sm{stoichiometric matrix}
\def\WRZD{weakly reversible, zero deficiency}\def\aut{autocatalytic}

% -- Epidemiology --
\def\brn{basic reproduction number}\def\bzn{basic replacement number}
\def\com{compartment}
\def\dfe{disease free equilibrium}\def\DFE{disease free equilibrium}
\def\dwt{dwell times}\def\edt{expected dwell times}
\def\EM{epidemic model}
\def\endp{endemic fixed point}\def\enp{endemic point}\def\EE{endemic point}
\def\ME{mathematical epidemiology}
\def\MEE{eco-epidemiology}\def\EEP{eco-epidemiological}
\def\ece{eco-epidemiological model}
\def\mgm{multi group model}
\def\mtS{multi-strain model}\def\mSt{multi-strain model}
\def\mSi{minimal siphon}
\def\NGM{next generation matrix}
\def\nim{new infections matrix}
\def\qem{quadratic epidemiological model}
\def\repF{invasion functions}
\def\SPF{SIR-PH-FA model}
\def\sus{susceptible}

% -- Graph theory & Linear algebra --
\def\jac{Jacobi matrix}
\def\jin{Jacobian matrix of the invasion vector field with respect to the invasion variables}
\def\jrn{Jacobian matrix of the resident vector field with respect to the resident variables}
\def\jxn{Jacobian matrix of the $\x$ vector field with respect to the $\x$ variables}
\def\Lm{Laplacian matrix}
\def\MTT{matrix tree theorem}
\def\SGT{spectral graph theory}
\def\SMW{Sherman–Morrison–Woodbury}

% -- Stability & Dynamics --
\def\absC{absorbtion CTMC}
\def\BTb{Bogdanov-Takens bifurcation}
\def\CEP{competitive exclusion partition}
\def\CTMC{continuous time Markov chain}
\def\ds{differential system}
\def\fp{fixed point}\def\bfp{boundary fixed point}\def\fbp{fixed boundary point}
\def\fbl{fixed boundary line}\def\fif{forward invariant face}
\def\gas{global asymptotic stability}\def\GAS{globally asymptotically stable}
\def\GASS{\text{GASS}}
\def\Hb{Hopf bifurcation}
\def\Hur{Hurwitz}\def\MH{Metzler-Hurwitz}\def\RH{Routh-Hurwitz}
\def\IS{\text{IS}}
\def\LAS{locally asymptotically stable}
\def\LdS{Lyapunov diagonally stable}
\def\Lf{Lyapunov function}\def\LyV{Lyapunov Volterra type function}\def\RLF{robust Lyapunov function}
\def\LyM{Lyapunov Malkin stable}
\def\PF{Perron-Frobenius}
\def\spfp{strictly positive fixed point}
\def\sta{stability}
\def\STP{sharp threshold property}\def\STT{sharp threshold theorem}

% -- Models (with \varphi explicitly instead of \f) --
\def\AB{$(A,B)$ Arino-Brauer epidemic models}
\def\ABa{$(A,\bbe,\ba,\varphi)$ balanced bilinear model}
\def\ABP{%$(A,\bbe,P,\varphi)$
\\ balanced bilinear model}
\def\AMG{Anderson May Gupta}
\newcommand\BIN{biological interaction network}
\def\Dt{Descartes type}
\def\gLV{generalized Lotka-Volterra model}
\def\IG{\text{invasion graph}}\def\IGMS{interaction graph on minimal siphons}
\def\LCP{\text{linear complementarity problem}}
\def\LV{Lotka-Volterra}\def\LVK{Lotka-Volterra-Kolmogorov}

% -- Splitting & Structure --
\def\reg{regular splitting}\def\regS{{\bf regular splitting}}
\def\RUR{rational univariate representation}
\def\bds{block-diagonal structure}\def\NGMs{next generation matrices}

% -- Miscellaneous --
\def\Ma{{Mathematica}}
\def\Mm{Max–Min}
\def\num{numerator}
\def\rap{rational approximation} \def\DSR{directed species-reaction bipartite graph}\def\PCP{polynomial complementarity problem}\def\rat{rational substitution}
\def\QSS{quasi steady state} \def\RHS{right hand side} \def\LHS{left hand side}
\def\OSN{online social network}\def\VdD{Van den Driessche}
\def\wli{$\omega$-limit} \def\rk{rich kinetics}\def\sJ{symbolic Jacobian}
\def\cV{\cite{Vas,VasSta,BSV,GSV}} \def\spf{semi-parametric semi-functional}
\def\CRW{Capasso-Ruan-Wang} \def\ABP{$(A,\bbe,P,\f)$ balanced bilinear model}
\def\ABa{%$(A,\bbe,\ba,\f)$ \\
rank-one balanced bilinear model}
\def\bbe{B%{\bff \beta}
} \def\vbe{\vec {\bff \beta}
}
\def\vb{\vec {\bff b}}
\newcommand{\Kh}{\hat{K}}          % K with hat, first mention only
\newcommand{\wS}{w^{\sigma}}       % right Perron J^\perp, full block
\newcommand{\piS}{\vec \pi^{\sigma}}    % left Perron J^\perp, full block
\newcommand{\ws}{w^{\sigma}}       % right Perron J^\perp, sub-block
\newcommand{\pis}{\pi^{\sigma}}    % left Perron J^\perp, sub-block
\newcommand{\hwS}{\hat{w}^{\sigma}}   % right Perron K, full block
\newcommand{\hpiS}{\hat{\pi}^{\sigma}} % left Perron K, full block
\newcommand{\hws}{\hat{w}^{\sigma}}   % right Perron K, sub-block
\newcommand{\hpis}{\hat{\pi}^{\sigma}} % left Perron K, sub-block
\newcommand{\twS}{\tilde{w}^{\sigma}}  % right Perron \tilde K, full block
\newcommand{\tpiS}{\tilde{\pi}^{\sigma}} % left Perron \tilde K
\newcommand{\tws}{\tilde{w}^{\sigma}}
\newcommand{\tpis}{\tilde{\pi}^{\sigma}}
\newcommand{\Mat}[2]{\mathbb{R}_{\ge 0}^{\,#1\times #2}}
\def\tJ{J_\sigma^\perp} \def\lPe{left Perron eigenvector}
% \G is \Gamma (defined at :795).  The Volterra function keeps its own name,
% so that \G no longer has two conflicting meanings in this file.
\def\GV{G}  % Volterra G(u)=u-1-\ln u
\def\PVL{Perron-Volterra Lyapunov minimal siphon ansatz}
\newcommand{\istar}[1]{i_{#1}^{*}}
%\newcommand{\sE}{s_E}
\newcommand{\vE}{v_E}
\def\PV{Perron-Volterra candidate Lyapunov function}
\def\GLV{generalized Lotka-Volterra}
\def\MERN{mathematical epidemiology reaction networks}
\def\CMM{contraction metric method}
\newcommand{\sign}{\operatorname{sign}}
\newcommand{\Var}{\operatorname{Var}} 

% ======================================================================
% IDEAS SHARED BETWEEN THE ABSTRACT AND THE MAIN RESULTS SECTION
%
% Write each idea ONCE, in the abstract; the Main Results section recalls it,
% so the two can never drift apart.
%
%   in the abstract        \ideaitem{1}{The first-order reaction terms ...}
%   in Main Results        \ideahead{1}
%
% \ideaitem stores the text globally and typesets
%   \item \hyperlink{idea:1}{<text>}
% \ideahead typesets
%   \item\hypertarget{idea:1}{}\textbf{<the same text>}
%
% The abstract must come before the Main Results input, which is the normal
% order. If a heading has to read differently from the abstract sentence (for
% instance a long abstract item summarised in one line), override it with the
% optional argument: \ideahead[short heading]{6}.
% ======================================================================
\makeatletter
\newcommand{\ideaitem}[2]{%
  \expandafter\gdef\csname @idea@#1\endcsname{#2}%
  \item\hyperlink{idea:#1}{#2}}
\newcommand{\ideahead}[2][]{%
  \item\hypertarget{idea:#2}{}%
  \def\@ideatmp{#1}%
  \ifx\@ideatmp\@empty
    \@ifundefined{@idea@#2}{%
      \PackageWarning{preamble}{Idea #2 is used in Main Results but was never
        defined by \string\ideaitem\space in the abstract}%
      \textbf{[idea #2 undefined]}%
    }{\textbf{\begingroup\def\ideaterm##1##2{##2}\def\idearest##1{}\@nameuse{@idea@#2}\endgroup}}%
  \else
    \textbf{#1}%
  \fi}
% \ideaterm{<hypertarget>}{<words>}: inside an abstract idea the words become an
% invisible link (the document sets \hypersetup{hidelinks}) jumping to the place
% in Main Results that defines them; inside the Main Results heading itself,
% where the target sits a few lines below, \ideahead strips the link and prints
% the words plainly. Mark the destination with \hypertarget{<name>}{} on the line
% before the display, which adds no equation number and shows nothing.
\newcommand{\ideaterm}[2]{\hyperlink{#1}{#2}}
% \idearest{<text>}: the part of an abstract idea after its first sentence; printed in
% the abstract, omitted from the Main Results heading.
\newcommand{\idearest}[1]{#1}
\makeatother

% ======================================================================
% SHORT-PAPER SWITCH
% \shorttrue in a master that assembles only part of the companion paper
% (Barrat/SipArxiv.tex). \nsh{...} drops text that cites results living in the
% omitted sections, so the short build has no dangling references while the
% companion paper keeps the text verbatim.
% ======================================================================
\newif\ifshort \shortfalse
\newcommand{\nsh}[1]{\ifshort\else#1\fi}
% \shortalt{<short build>}{<companion paper>}: use when a sentence needs
% different wording in the two builds, not merely a deletion.
\newcommand{\shortalt}[2]{\ifshort#1\else#2\fi}
% \Jtr, \JtrS{\Sigma}: transverse Jacobian. J^\perp here; SipArxiv.tex redefines them as J_{\mathrm{tra}}, J_{\mathrm{tra},\Sigma}
\newcommand{\Jtr}{J^\perp}
\newcommand{\JtrS}[1]{J_{#1}^\perp}
.
\makeatletter
\@ifpackageloaded{hyperref}{}{%
  \@ifundefined{hyperrefpackage}{%
    \@ifundefined{AAdefNnohyperref}{\usepackage{hyperref}}{}%
  }{}%
}
\makeatother
% linktoc=all makes a table-of-contents page number a link as well as the
% entry text; hyperref's default, linktoc=section, links only the text.
% Deferred to begin-document and guarded: by the block above hyperref may
% be loaded later (mdpi's own hook) or not at all (\AAdefNnohyperref), and
% \hypersetup must exist before it is called.  A driver that wants the
% default back can say \hypersetup{linktoc=section} after % Shared preamble for the GLVbook papers.  This is the single canonical copy;
% a driver one level below GLVbook reaches it as
%
%     \input{../basicFiles/preamble}
%
% What stays in the driver, because it differs from paper to paper:
% \documentclass, font choice, \title, \author, \date, and any macro genuinely
% local to one paper.  Everything else belongs here.  The file is
% self-contained: it does not read any other file.

\usepackage[utf8]{inputenc}
\usepackage[novbox]{pdfsync}   % required: SyncTeX alone does not work here

% Page geometry, set once here for every paper that loads this file.  A driver
% that needs different margins says \geometry{...}; a second
% \usepackage[...]{geometry} raises an option clash and must not be used.
\makeatletter
\@ifundefined{AAdefNnogeometry}{\usepackage[a4paper,margin=2cm]{geometry}}{}
\makeatother

% mathematics
\usepackage{amsmath,amssymb,amsthm,mathtools,bm,mathrsfs,bbm}
\usepackage{centernot}

% floats, tables, lists
\usepackage{graphicx,float}
\makeatletter
\@ifpackageloaded{subfig}{}{\usepackage{subcaption}}
\makeatother
\usepackage{booktabs,longtable,threeparttable,array,multirow}
\usepackage{enumitem}

% colour, code, algorithms
\usepackage{xcolor,listings,tcolorbox}
\makeatletter
\@ifundefined{c@algorithm}{\usepackage{algorithm}}{}
\makeatother
\usepackage{algorithmicx,algpseudocode}

% diagrams
\usepackage{tikz,tikz-cd,pgfplots}
\usetikzlibrary{graphs,decorations.pathmorphing,decorations.markings,
arrows.meta,automata,positioning}
\tikzcdset{every label/.append style = {font = \small}}
\pgfplotsset{compat=1.18}

% chemistry
\usepackage[version=4]{mhchem}

\usepackage{microtype}
\usepackage{lineno}
\usepackage{nomencl}
\makenomenclature

% hyperref first, then bookmark, then glossaries: glossaries must follow
% hyperref or the entry links are built against the wrong anchors.
\makeatletter
\@ifundefined{hyperrefpackage}{%
  \@ifundefined{AAdefNnohyperref}{\usepackage{hyperref}}{}%
}{}
\@ifundefined{hyperrefpackage}{\usepackage{bookmark}}{}
\makeatother
\usepackage[toc,sort=use]{glossaries}
\makeglossaries

% Book/standalone switch.  Material that the book already covers in its
% prelude, but that a standalone paper still has to state, is wrapped as
%     \ifbookbuild\else ... \fi
% Drivers default to standalone; bookM.tex says \bookbuildtrue.
\newif\ifbookbuild
\bookbuildfalse

% Part switches for the gradostat paper.  Both default to true; a driver that
% wants the short submission says \madecfalse and/or \slowfastfalse, and every
% cross-reference into the dropped part disappears with it.
%   \ifmadec     Part III, the Lagasquie--Madec model and its break-even geometry
%   \ifslowfast  Part IV, everything carrying the small parameter \eps
\newif\ifmadec     \madectrue
\newif\ifslowfast  \slowfasttrue

% Part switch for the Barrat/simplicial-contagion project.  Defaults to
% true (the hypernetwork-synchrony part is included); the journal-paper
% driver says \hnetpartfalse to drop it and every cross-reference into it.
\newif\ifhnetpart  \hnetparttrue

%% ========================================
%% PACKAGES
%% ========================================
\RequirePackage{etoolbox}
\usepackage{makecell}
\usepackage[novbox]{pdfsync}% obsolete; SyncTeX is built into modern engines
\usepackage{amssymb, amsfonts,mathtools,bbm,eurosym,graphicx,latexsym, amsmath,amsthm,times,nccmath,url}
\usepackage{listings,verbatim,cases,tabularx}
% hyperref is deliberately NOT in the line above.  Journal classes often
% load it themselves, with a long option list, from a begindocument hook
% (mdpi.cls does: \DeclareHookRule{begindocument}{hyperref}{before}{mdpi}).
% Loading it here first, without options, makes that later load raise
% "Option clash for package hyperref".  So load it only if nobody else
% has, or will:
%   - \@ifpackageloaded  catches a class that loaded it already;
%   - \hyperrefpackage    is mdpi.cls's own deferred loader;
%   - \AAdefNnohyperref   lets any other driver opt out by defining it
%                         before \input{preamble}.
\makeatletter
\@ifpackageloaded{hyperref}{}{%
  \@ifundefined{hyperrefpackage}{%
    \@ifundefined{AAdefNnohyperref}{\usepackage{hyperref}}{}%
  }{}%
}
\makeatother
% linktoc=all makes a table-of-contents page number a link as well as the
% entry text; hyperref's default, linktoc=section, links only the text.
% Deferred to begin-document and guarded: by the block above hyperref may
% be loaded later (mdpi's own hook) or not at all (\AAdefNnohyperref), and
% \hypersetup must exist before it is called.  A driver that wants the
% default back can say \hypersetup{linktoc=section} after \input{preamble}.
\makeatletter
\AtBeginDocument{\@ifundefined{hypersetup}{}{\hypersetup{linktoc=all}}}
\makeatother
\usepackage{graphicx,float,fancybox}
\usepackage{epsfig,epstopdf,tikz}
\usepackage[sans]{dsfont}
\usetikzlibrary{positioning,arrows,arrows.meta,calc}
% Configure listings for Mathematica code
\lstdefinestyle{mathematica}{
 language=Mathematica,
 basicstyle=\small\ttfamily,
 backgroundcolor=\color{gray!10},
 frame=single,
 breaklines=true,
 commentstyle=\color{green!60!black},
 keywordstyle=\color{blue},
 stringstyle=\color{red},
 showstringspaces=false,
 tabsize=2
}
\usepackage{tcolorbox}
% Colored boxes for key results
\newtcolorbox{keyresult}[1][]{
 colback=blue!5!white,
 colframe=blue!75!black,
 title=#1,
 fonttitle=\bfseries
}
\newtcolorbox{examplebox}[1][]{
 colback=green!5!white,
 colframe=green!75!black,
 title=#1,
 fonttitle=\bfseries
}


\makeatletter
\@ifundefined{theorem}   {\newtheorem{theorem}   {Theorem}}   {}
\@ifundefined{definition}{\newtheorem{definition}{Definition}} {}
\@ifundefined{open}      {\newtheorem{open}      {Problem}}   {}
\@ifundefined{example}   {\newtheorem{example}   {Example}}   {}
\makeatother

% COUNTER DEFINITIONS
\newcounter{tabcounter}
\newcounter{defcounter}
\newcounter{opcounter}
\newcounter{thmcounter}
\newcounter{excounter}
\newcounter{ascounter}
\newcounter{cjcounter}
\newcounter{lemcounter}
\newcounter{remcounter}
\newcounter{corcounter}
\newcounter{propcounter}
\newcounter{figcounter}

% UNIFIED GLOSSARY STYLE
\newglossarystyle{unifiedstyle}{%
  \setglossarystyle{list}%
  \renewcommand*{\glossarypreamble}{\gdef\glslastsec{}}%
  \renewcommand{\glossentry}[2]{%
    \edef\glscursec{\glsentryuseriii{##1}}%
    \ifx\glscursec\glslastsec\else
      \item[]\textbf{\S\,\glscursec\enspace\glsentryuseriv{##1}}%
      \global\let\glslastsec\glscursec
    \fi
    \item[\glsentryuserii{##1} \glsentryuseri{##1}:] ``\glossentryname{##1}'' \glsentrydesc{##1}%
  }%
}

\makeatletter
% First-word extractor (used by \beT, \beD and \beXa for the default key)
% \@firstwordof hello world \relax  ->  hello
\def\@firstwordof#1 #2\relax{#1}
% First three letters, used for the heading mnemonic of a theorem given no
% explicit key.  Padding with \@empty keeps it safe for first words shorter
% than three letters: \@empty vanishes under \edef, whereas \relax would
% survive into the label and is illegal inside \csname (it broke every
% two-letter title, e.g. "ME ...", with "Missing \endcsname inserted").
\def\@takethree#1#2#3#4\@nil{#1#2#3}

% EXAMPLE COMMAND -- same signature and behaviour as \beT and \beD.
% Usage: \beXa*[long title][short title][key]
%   star  = suppresses the printed "[key]" tag after the title (as for \beT*
%   and \beD*); the label and glossary entry are unaffected.
%   key = optional label key override; the label becomes e:KEY, and without a
%   key it is e: plus the first three letters of the title's first word.  A
%   mnemonic label lets an example be cited as "Example~\ref{e:KEY}" instead
%   of by section number, which is what the citing prose needs.
\NewDocumentCommand{\beXa}{s O{} O{} o}{%
  \def\@temparg{#2}%
  \ifx\@temparg\@empty
    \begin{example}%
  \else
    \let\@savedSecTitle\@currentsectionname%
    \refstepcounter{excounter}%
    \IfValueTF{#4}{%
      \edef\@exlblkey{#4}%
    }{%
      \edef\@exfirstword{\@firstwordof#2 \relax}%
      \edef\@exlblkey{%
        \expandafter\@takethree\@exfirstword\@empty\@empty\@empty\@nil}%
    }%
    \def\@tempshort{#3}%
    \ifx\@tempshort\@empty
      \expandafter\xdef\csname @xshort@\@exlblkey\endcsname{\@exlblkey}%
    \else
      \expandafter\def\csname @xshort@\@exlblkey\endcsname{#3}%
    \fi
    \IfBooleanTF{#1}{\let\@extag\@empty}{%
      \edef\@extag{\noexpand\ \noexpand\textup{[\noexpand\texttt{e:\@exlblkey}]}}}%
    \begin{example}[#2\ifshowkeys\@extag\fi]%
    \label{e:\@exlblkey}%
    % record the long title so \nref{e:KEY} prints it, as for theorems
    \protected@write\@auxout{}{%
      \string\@thmtitle{e:\@exlblkey}{\unexpanded{#2}}}%
    \expanded{%
      \noexpand\newglossaryentry{e\theexcounter}{%
        name={\unexpanded{#2}%
          \IfBooleanF{#1}{ [\noexpand\texttt{e:\@exlblkey}]}},%
        sort={\theexample},%
        description={p.~\noexpand\pageref{e:\@exlblkey}},%
        user1={\theexample},%
        user2={Example},%
        user3={\thesection},%
        user4={\expandafter\unexpanded\expandafter{\@savedSecTitle}}%
      }%
    }%
    \protected@write\@auxout{}{\string\glsliveentry{e\theexcounter}}%
  \fi
}
\newcommand{\eeXa}{\end{example}}

% DEFINITION COMMAND -- same signature and behaviour as \beT below.
% Usage: \beD[long title][short title][key]
%   long title  = definition heading (plain text only)
%   short title = shown in \dref{word}; defaults to the key
%   key         = optional label key override; the label becomes d:KEY
% Without a key the label is d: plus the first three letters of the title's
% first word, so it is a mnemonic rather than the definition's number: a
% numeric label would silently change whenever a definition is inserted.
\NewDocumentCommand{\beD}{s O{} O{} o}{%
  \def\@temparg{#2}%
  \ifx\@temparg\@empty
    \begin{definition}%
  \else
    \let\@savedSecTitle\@currentsectionname%
    \refstepcounter{defcounter}%
    \IfValueTF{#4}{%
      \edef\@deflblkey{#4}%
    }{%
      \edef\@deffirstword{\@firstwordof#2 \relax}%
      \edef\@deflblkey{%
        \expandafter\@takethree\@deffirstword\@empty\@empty\@empty\@nil}%
    }%
    \def\@tempshort{#3}%
    \ifx\@tempshort\@empty
      \expandafter\xdef\csname @dshort@\@deflblkey\endcsname{\@deflblkey}%
    \else
      \expandafter\def\csname @dshort@\@deflblkey\endcsname{#3}%
    \fi
    \IfBooleanTF{#1}{\let\@deftag\@empty}{%
      \edef\@deftag{\noexpand\ \noexpand\textup{[\noexpand\texttt{d:\@deflblkey}]}}}%
    \begin{definition}[#2\ifshowkeys\@deftag\fi]%
    \label{d:\@deflblkey}%
    % record the long title so \nref{d:KEY} prints it, as for theorems
    \protected@write\@auxout{}{%
      \string\@thmtitle{d:\@deflblkey}{\unexpanded{#2}}}%
    \expanded{%
      \noexpand\newglossaryentry{d\thedefcounter}{%
        name={\unexpanded{#2}%
          \IfBooleanF{#1}{ [\noexpand\texttt{d:\@deflblkey}]}},%
        sort={\thedefinition},%
        description={p.~\noexpand\pageref{d:\@deflblkey}},%
        user1={\thedefinition},%
        user2={Definition},%
        user3={\thesection},%
        user4={\expandafter\unexpanded\expandafter{\@savedSecTitle}}%
      }%
    }%
    \protected@write\@auxout{}{\string\glsliveentry{d\thedefcounter}}%
  \fi
}
\newcommand{\eeD}{\end{definition}}

% \showkeysfalse in a driver hides the [t:KEY] mnemonics in result headings
% (labels and glossary entries are unchanged); used for submission builds.
\newif\ifshowkeys \showkeystrue
% \commonredtrue in a driver that inputs basicFiles/commonRed: tool.tex then
% cites its reduction results instead of restating them.
\newif\ifcommonred \commonredfalse
% TheorEM COMMAND
% Usage: \beT[long title][short title][key]
%   long title  = theorem heading (plain text only)
%   short title = shown in \tref{word}; may contain math; defaults to first word of long title
%   key         = optional label key override; label becomes t:KEY (avoids first-word collisions)
% Label is t:WORD where WORD = first word of long title (or KEY if given); use \tref{WORD} to reference.
\NewDocumentCommand{\beT}{s O{} O{} o}{%
  \def\@temparg{#2}%
  \ifx\@temparg\@empty
    \begin{theorem}%
  \else
    \let\@savedSecTitle\@currentsectionname%
    \refstepcounter{thmcounter}%
    % \@thmtag is the mnemonic shown in the heading; every theorem gets one.
    % With an explicit key the key itself is shown.  Without one the key
    % defaults to the title's first word, and showing that whole word would
    % merely repeat the title ("Dissipativity [Dissipativity]"), so only its
    % first three letters are shown: "Dissipativity [Dis]".  The label is
    % unaffected -- it still uses the full \@thmlblkey.
    \IfValueTF{#4}{%
      \edef\@thmlblkey{#4}%
    }{%
      \edef\@thmfirstword{\@firstwordof#2 \relax}%
      \edef\@thmlblkey{%
        \expandafter\@takethree\@thmfirstword\@empty\@empty\@empty\@nil}%
    }%
    \IfBooleanTF{#1}{\let\@thmtag\@empty}{%
      \edef\@thmtag{\noexpand\ \noexpand\textup{[\noexpand\texttt{t:\@thmlblkey}]}}}%
    \def\@tempshort{#3}%
    \ifx\@tempshort\@empty
      \expandafter\xdef\csname @tshort@\@thmlblkey\endcsname{\@thmlblkey}%
    \else
      \expandafter\def\csname @tshort@\@thmlblkey\endcsname{#3}%
    \fi
    \begin{theorem}[#2\ifshowkeys\@thmtag\fi]%
    \label{t:\@thmlblkey}%
    % record the long title for \nref; via the .aux so forward refs resolve
    \protected@write\@auxout{}{%
      \string\@thmtitle{t:\@thmlblkey}{\unexpanded{#2}}}%
    \expanded{%
      \noexpand\newglossaryentry{t\thethmcounter}{%
        name={\unexpanded{#2}%
          \IfBooleanF{#1}{ [\noexpand\texttt{t:\@thmlblkey}]}},%
        % sort/user1 must be the number the document actually prints, i.e. the
        % theorem environment's counter, not thmcounter.  thmcounter advances
        % only on \beT with a title, so a bare \beT (or \beT with an
        % empty title) makes the two drift and the glossary cites a wrong
        % number.  thmcounter is still fine as the entry key: it is unique.
        sort={\thetheorem},%
        description={p.~\noexpand\pageref{t:\@thmlblkey}},%
        user1={\thetheorem},%
        user2={Theorem},%
        user3={\thesection},%
        user4={\expandafter\unexpanded\expandafter{\@savedSecTitle}}%
      }%
    }%
    \protected@write\@auxout{}{\string\glsliveentry{t\thethmcounter}}%
  \fi
}
\newcommand{\eeT}{\end{theorem}}

% CROSS-PART RELATION PARAGRAPH
% Usage: \xpart{text}  -- place immediately BEFORE \beT, never inside it
%   (anything inside \beT...\eeT would be swallowed into the glossary entry).
% The paragraph is typeset only when Part I is actually present, i.e. in the
% combined book build; in the standalone builds (GranadaDSABNS/A.tex,
% AAHRah-Submission/AAHRah.tex) the label part:glv is undefined and the whole
% paragraph disappears rather than leaving ?? cross-references. Same
% graceful-degradation principle as \GASref in GASt.tex. Resolves on the
% second LaTeX pass, since r@part:glv comes from the .aux file.
% (Already inside the \makeatletter region opened above, so no \makeatletter
% pair here -- adding one would close that region early.)
\newcommand{\xpart}[1]{\@ifundefined{r@part:glv}{}{\paragraph{Relation with Parts~I--II.}#1}}

% \bookref{label}{precise text}{fallback text}: prints the precise cross-reference
% when the label resolves (combined book build), and the fallback text otherwise
% (standalone paper builds, where book-only appendix labels are absent). Same
% graceful-degradation role as \GASref in GASt.tex, but usable in any file since it
% lives in this preamble (loaded by every driver) -- unlike \GASref, which is defined in a
% content file and is therefore undefined in standalone builds that omit that file.
\newcommand{\bookref}[3]{\@ifundefined{r@#1}{#3}{#2}}

% \PVpart and \RBpart: the two parts that a partial build may omit.
% They print "Part~N" in any build that contains the part, and a naming phrase
% in a build that does not, so an excerpt shows no ?? .  Each already supplies
% the word "Part", so write "\PVpart develops ...", not "Part~\PVpart".
\newcommand{\PVpart}{\bookref{part:aahrah}%
  {Part~\ref{part:aahrah}}{the Perron--Volterra part of the full book}}
\newcommand{\RBpart}{\bookref{part:recbrowne}%
  {Part~\ref{part:recbrowne}}{the immuno-virology part of the full book}}

% OPEN PROBLEM COMMANDS
% ASSUMPTION COMMANDS (glossaried)
\newcommand{\beAs}[1][]{%
  \def\@temparg{#1}%
  \ifx\@temparg\@empty
    \begin{assumption}%
  \else
    \let\@savedSecTitle\@currentsectionname%
    \refstepcounter{ascounter}%
    \begin{assumption}[#1]%
    \label{as:\theascounter}%
    \expanded{%
      \noexpand\newglossaryentry{as\theascounter}{%
        name={\unexpanded{#1}},%
        sort={\theascounter},%
        description={p.~\noexpand\pageref{as:\theascounter}},%
        user1={\theascounter},%
        user2={Assumption},%
        user3={\thesection},%
        user4={\expandafter\unexpanded\expandafter{\@savedSecTitle}}%
      }%
    }%
    \protected@write\@auxout{}{\string\glsliveentry{as\theascounter}}%
  \fi
}
\newcommand{\eeAs}{\end{assumption}}
% CONJECTURE COMMANDS (glossaried)
\newcommand{\beCj}[1][]{%
  \def\@temparg{#1}%
  \ifx\@temparg\@empty
    \begin{conjecture}%
  \else
    \let\@savedSecTitle\@currentsectionname%
    \refstepcounter{cjcounter}%
    \begin{conjecture}[#1]%
    \label{cj:\thecjcounter}%
    \expanded{%
      \noexpand\newglossaryentry{cj\thecjcounter}{%
        name={\unexpanded{#1}},%
        sort={\thecjcounter},%
        description={p.~\noexpand\pageref{cj:\thecjcounter}},%
        user1={\thecjcounter},%
        user2={Conjecture},%
        user3={\thesection},%
        user4={\expandafter\unexpanded\expandafter{\@savedSecTitle}}%
      }%
    }%
    \protected@write\@auxout{}{\string\glsliveentry{cj\thecjcounter}}%
  \fi
}
\newcommand{\eeCj}{\end{conjecture}}
% LEMMA COMMANDS
% ---------------------------------------------------------------------
% UNIFORM RESULT WRAPPERS -- \beL, \beC, \beP, \beR share \beT's signature
%
%   \beX[long title][short title][key]  ...  \eeX
%
%   long title  = heading (math allowed)
%   short title = shown in the cross-reference; defaults to the key
%   key         = label key.  With a key the label is X:KEY and the heading
%                 carries the mnemonic [key], exactly as for \beT.  Without
%                 one the label falls back to the running counter, which is
%                 what every pre-existing call site in the book relies on.
%
% A bare \beX -- arguments commented out -- opens the plain environment and
% creates NO label and NO glossary entry.  That is the supported way to keep
% an individual result out of the catalogue.
% ---------------------------------------------------------------------
% #1 wrapper name  #2 environment  #3 label prefix  #4 counter  #5 type name
\newcommand{\@mkresult}[5]{%
  % A leading * suppresses the [key] mnemonic in both the heading and the
  % glossary name, while still making the label X:KEY and the entry:
  %     \beX*[long][short][key]
  \expandafter\NewDocumentCommand\csname be#1\endcsname{s O{} O{} o}{%
    \def\@temparg{##2}%
    \ifx\@temparg\@empty
      \begin{#2}%
    \else
      \let\@savedSecTitle\@currentsectionname%
      \refstepcounter{#4}%
      % With a key: mnemonic label X:KEY and the [key] tag in the heading.
      % Without one: fall back to the running counter, exactly as these
      % wrappers behaved before, so every legacy one-argument call site in
      % the book keeps its label.  The counter is also csname-safe, which a
      % title's first word is not -- titles beginning with math would
      % otherwise build \csname \Omega...\endcsname and fail.
      \IfValueTF{##4}{%
        \edef\@relbl{##4}%
        \IfBooleanTF{##1}{\let\@retag\@empty}{%
          \edef\@retag{\noexpand\ \noexpand\textup{[\noexpand\texttt{#3:##4}]}}}%
      }{%
        \edef\@relbl{\arabic{#4}}%
        \let\@retag\@empty
      }%
      \def\@tempshort{##3}%
      \ifx\@tempshort\@empty
        \expandafter\xdef\csname @sshort@#3\@relbl\endcsname{\@relbl}%
      \else
        \expandafter\def\csname @sshort@#3\@relbl\endcsname{##3}%
      \fi
      \begin{#2}[##2\ifshowkeys\@retag\fi]%
      \edef\@tmplbl{#3:\@relbl}\expandafter\label\expandafter{\@tmplbl}%
      \expanded{%
        \noexpand\newglossaryentry{#3\csname the#4\endcsname}{%
          name={\unexpanded{##2}%
            \IfBooleanF{##1}{\IfValueT{##4}{ [\noexpand\texttt{#3:##4}]}}},%
          sort={\csname the#2\endcsname},%
          description={p.~\noexpand\pageref{#3:\@relbl}},%
          user1={\csname the#2\endcsname},%
          user2={#5},%
          user3={\thesection},%
          user4={\expandafter\unexpanded\expandafter{\@savedSecTitle}}%
        }%
      }%
      \protected@write\@auxout{}{\string\glsliveentry{#3\csname the#4\endcsname}}%
      \protected@write\@auxout{}{\string\@thmtitle{\@tmplbl}{\unexpanded{##2}}}%
    \fi
  }%
  \expandafter\newcommand\csname ee#1\endcsname{\end{#2}}%
}
\@mkresult{L}{lemma}{l}{lemcounter}{Lemma}
% REMARK / COROLLARY / PROPOSITION -- same signature as \beL above
\@mkresult{R}{remark}{r}{remcounter}{Remark}
\@mkresult{C}{corollary}{c}{corcounter}{Corollary}
\@mkresult{P}{proposition}{p}{propcounter}{Proposition}
% OPEN-PROBLEM COMMAND -- generated with the same signature as \beT:
%   \beO[long title][short title][key]   ...   \eeO
% label o:KEY, or o:<counter> when no key is given; a bare \beO makes
% neither a label nor a glossary entry.
\@mkresult{O}{open}{o}{opcounter}{Problem}


% Section-level title tracker: updated only by \section, not \subsection.
% \@currentlabelname is overwritten by every \refstepcounter (including
% \subsection), so we save it separately here.
\newcommand{\@currentsectionname}{}
\AtBeginDocument{%
  \let\@epidcrn@orig@refstepcounter\refstepcounter
  \renewcommand{\refstepcounter}[1]{%
    \@epidcrn@orig@refstepcounter{#1}%
    \ifstrequal{#1}{section}{\global\let\@currentsectionname\@currentlabelname}{}%
  }%
}

% TABLE COMMANDS (glossaried) -- consistent naming with \beT, \beD, etc.
\newcommand{\beTa}[1][]{%
  \def\@temparg{#1}%
  \ifx\@temparg\@empty
    \begin{table}[htbp]%
  \else
    \let\@savedSecTitle\@currentsectionname%
    \refstepcounter{tabcounter}%
    \begin{table}[htbp]%
    \label{A:\thetabcounter}%
    \expanded{%
      \noexpand\newglossaryentry{A\thetabcounter}{%
        name={\unexpanded{#1}},%
        sort={\thetabcounter},%
        description={p.~\noexpand\pageref{A:\thetabcounter}},%
        user1={\thetabcounter},%
        user2={Table},%
        user3={\thesection},%
        user4={\expandafter\unexpanded\expandafter{\@savedSecTitle}}%
      }%
    }%
    \protected@write\@auxout{}{\string\glsliveentry{A\thetabcounter}}%
  \fi
}
\newcommand{\eeTa}{\end{table}}

% FIGURE COMMANDS (glossaried) -- consistent with \beT, \beR, \beTa, etc.
% Usage:
%   \beF[Glossary title]...\caption{Visible caption}\eeF  -> glossaried figure
%   \beF...\caption{Visible caption}\eeF                  -> plain [H] figure, no glossary
\newcommand{\beF}[1][]{%
  \def\@temparg{#1}%
  \ifx\@temparg\@empty
    \begin{figure}[H]%
  \else
    \begin{figure}[H]%
    \let\@savedSecTitle\@currentsectionname%
    \refstepcounter{figcounter}%
    \label{fig:\thefigcounter}%
    \expanded{%
      \noexpand\newglossaryentry{Fig\thefigcounter}{%
        name={\unexpanded{#1}},%
        sort={\thefigcounter},%
        description={p.~\noexpand\pageref{fig:\thefigcounter}},%
        user1={\thefigcounter},%
        user2={Figure},%
        user3={\thesection},%
        user4={\expandafter\unexpanded\expandafter{\@savedSecTitle}}%
      }%
    }%
    \protected@write\@auxout{}{\string\glsliveentry{Fig\thefigcounter}}%
  \fi
}
\newcommand{\eeF}{\end{figure}}
\makeatother
% LIVE-ENTRY REGISTER
% \glsaddall prints every entry the glossaries package currently knows, and
% that set includes entries restored from \jobname.glsdefs.  The package
% rewrites .glsdefs from all *currently defined* entries, so anything it
% restored is written straight back out: a deleted section kept reappearing in
% the catalogue for ever, and the .glsdefs had to be erased by hand.  The
% wrappers therefore record every entry they actually create in the .aux, and
% \printboth adds only those.  The .aux is rebuilt from scratch on each run, so
% this is self-healing: a removed result is gone after one further run.
\makeatletter
\newcommand{\@glslive}{}
\newcommand{\@glsaddlive}[1]{\ifglsentryexists{#1}{\glsadd{#1}}{}}
\newcommand{\glsliveentry}[1]{\g@addto@macro\@glslive{\@glsaddlive{#1}}}
\makeatother
% PRINT COMMAND
% No \clearpage here -- the catalogue continues on the current page.  This has
% been re-added by mistake several times; leave it out.
\makeatletter
\newcommand{\printboth}{%
  \@glslive   % add exactly the entries this document created last run
  \printglossary[style=unifiedstyle, title={0 \hspace{.22em} Definitions, Theorems, Examples and Open Problems}]%
}
\makeatother
% Named theorem reference: \tref{word} -> Thm~N(word-or-short)
% Shows stored short title if available, otherwise shows the key itself.
\newcommand{\tref}[1]{%
  Thm~\ref{t:#1}(\ifcsname @tshort@#1\endcsname\csname @tshort@#1\endcsname\else#1\fi)%
}

% Full-title theorem reference: \nref{t:KEY} prints the number and then the
% theorem's long title, so one writes "Theorem~\nref{t:blocks}" and gets
% "Theorem 3 ``Block factorization and siphons''".  Takes the complete label, like
% \ref, not the bare key like \tref.
%
% \beT writes each title to the .aux (see \@thmtitle below), so \nref resolves
% forward references from the second LaTeX pass, exactly as \ref does.  If the
% title is not yet known -- first pass, or a label that is not a \beT theorem
% -- it degrades to the number rather than producing ??.
% These two use @ in their names, and this point in the file is outside the
% \makeatletter region opened earlier (it closes at \makeatother above), so
% the region has to be reopened here.
\makeatletter
\newcommand{\@thmtitle}[2]{\global\@namedef{@tlong@#1}{#2}}
\newcommand{\nref}[1]{%
  \ref{#1}%
  \ifcsname @tlong@#1\endcsname
    \ ``\csname @tlong@#1\endcsname''%
  \fi
}
% \nm{t:KEY}: the stored long title of any \beT/\beL/\beP/\beR/\beC/\beD/\beE result,
% verbatim and without quotes, so item names follow the macro's first argument.
\newcommand{\nm}[1]{%
  \ifcsname @tlong@#1\endcsname
    \csname @tlong@#1\endcsname%
  \else
    \ref{#1}%
  \fi
}
\makeatother




%% ========================================
%% Other THEOREM ENVIRONMENTS
%% ========================================
\makeatletter
\@ifundefined{lemma}      {\newtheorem{lemma}      {Lemma}}      {}
\@ifundefined{proposition}{\newtheorem{proposition}{Proposition}} {}
\@ifundefined{corollary}  {\newtheorem{corollary}  {Corollary}}  {}
\@ifundefined{conjecture} {\newtheorem{conjecture} {Conjecture}}  {}
\@ifundefined{remark}     {\newtheorem{remark}     {Remark}}     {}
\@ifundefined{question}   {\newtheorem{question}   {Question}}   {}
\@ifundefined{assumption} {\newtheorem{assumption} {Assumption}} {}
\@ifundefined{hypothesis} {\newtheorem{hypothesis} {Hypothesis}} {}
\@ifundefined{fact}       {\newtheorem{fact}       {Fact}}       {}
\makeatother

% \citer: cite the relay-graph threshold theorem of AAHrel. If a document has
% \input the chapter/file containing t:relayRF (so the label resolves), point
% straight at it; otherwise fall back to a bare citation of AAHrel. Shared here
% (not per-paper) so it resolves consistently across every book chapter.
% \citeR: same idea, for the first theorem of AAHRah (t:Pstar).
\makeatletter
\providecommand{\citer}{%
  \@ifundefined{r@t:relayRF}%
    {\cite{AAHrel}}%
    {Theorem~\ref{t:relayRF} of \cite{AAHrel}}%
}
\providecommand{\citeR}{%
  \@ifundefined{r@t:Pstar}%
    {\cite{AAHRah}}%
    {Theorem~\ref{t:Pstar} of \cite{AAHRah}}%
}
\makeatother


% Abbreviations
\def\beQ{\begin{question}}\def\eeQ{\end{question}}
\def\beA{\begin{assumption}}\def\eeA{\end{assumption}}


%% ========================================

%% ========================================
%% COLORS & CUSTOM COMMANDS
%% ========================================
\definecolor{funccolor}{RGB}{25,25,112}
\definecolor{desccolor}{RGB}{64,64,64}
\newcommand{\func}[1]{\texttt{\textcolor{funccolor}{#1}}}
\def\eq{=}
\def\mm{-}

%% ========================================
%% MATHCAL SHORTHANDS
%% ========================================
% Project-wide convention: \mX for a capital X is always \mathcal X.
% Defined for the whole alphabet so the convention never fails to expand.
\newcommand{\mA}{\mathcal A}\newcommand{\mB}{\mathcal B}
\newcommand{\mC}{\mathcal C}\newcommand{\mD}{\mathcal D}
\newcommand{\mE}{\mathcal E}\newcommand{\mF}{\mathcal F}
\newcommand{\mG}{\mathcal G}\newcommand{\mH}{\mathcal H}
\newcommand{\mI}{\mathcal I}\newcommand{\mJ}{\mathcal J}
\newcommand{\mK}{\mathcal K}\newcommand{\mL}{\mathcal L}
\newcommand{\mM}{\mathcal M}\newcommand{\mN}{\mathcal N}
\newcommand{\mO}{\mathcal O}\newcommand{\mP}{\mathcal P}
\newcommand{\mQ}{\mathcal Q}\newcommand{\mR}{\mathcal R}
\newcommand{\mS}{\mathcal S}\newcommand{\mT}{\mathcal T}
\newcommand{\mU}{\mathcal U}\newcommand{\mV}{\mathcal V}
\newcommand{\mW}{\mathcal W}\newcommand{\mX}{\mathcal X}
\newcommand{\mY}{\mathcal Y}\newcommand{\mZ}{\mathcal Z}

%% ========================================
%% BASIC MATH ENVIRONMENTS
%% ========================================
\def\bep{\begin{pmatrix}}\def\eep{\end{pmatrix}}
\def\bev{\begin{vmatrix}}\def\eev{\end{vmatrix}}
\def\bea{\begin{eqnarray*}}\def\eea{\end{eqnarray*}}
\def\bc{\begin{cases}}\def\ec{\end{cases}}
\def\BEN{\begin{enumerate}}\def\EEN{\end{enumerate}}
\def\BI{\begin{itemize}}\def\EI{\end{itemize}}

\newcommand{\be}[1]{\begin{equation}\label{#1}}
\newcommand{\ee}{\end{equation}}
\newcommand{\beq}{\begin{eqnarray}}
\def\eeq{\end{eqnarray}}
\def\eqr{\eqref}\def\fr{\frac}\def\lbl{\label}
\def\Lra{\Longrightarrow}\def\non{\nonumber}\def\cNGM{\cite{Diek,Van,Van08}}
\def\Eq{\Leftrightarrow}\def\b{\beta}\def\uph{\red{UP TO HERE}}
%% ========================================
%% PROBABILITY & STATISTICS
%% ========================================
\renewcommand{\P}{\ensuremath{\mathbb{P}}}  % Note: overrides paragraph symbol
\newcommand{\E}{\ensuremath{\mathbb{E}}}
\newcommand{\ind}[1]{\ensuremath{\mathbbm{1}_{\left\{#1\right\}}}}
\DeclareMathOperator{\Cov}{Cov}
\DeclareMathOperator{\TV}{TV}
\newcommand{\rank}{\operatorname{rank}}

%% ========================================
%% CRN OPERATORS (organizations/ERCs/buffering/siphons framework)
%% ========================================
\def\clos{\mathrm{clos}}
\def\Supply{\mathrm{Supply}}
\def\req{\mathrm{req}}
\def\React{\mathrm{React}}
\def\Prod{\mathrm{Prod}}
\def\Gen{\mathrm{Gen}}
\def\In{\mathrm{In}}
\def\Out{\mathrm{Out}}
\def\Act{\mathrm{Act}}
\def\Min{\mathrm{Min}}

%% ========================================
%% LINE BREAKING
%% ========================================
% Last-resort stretch, used only on paragraphs TeX cannot otherwise break:
% a long \texttt{} token, a comma-separated inline math list (commas are not
% breakpoints in math), or a run of \cite/\ref boxes.  Without it those set as
% overfull \hbox and print into the margin; with it they set slightly loose,
% which is the lesser evil.  Does not affect paragraphs that already break.
\emergencystretch=3em

%% ========================================
%% NUMBER SETS
%% ========================================
\newcommand{\R}{\mathbb{R}}  % Real numbers
\newcommand{\C}{\mathbb{C}}  % Complex numbers
\newcommand{\N}{\mathbb{N}}
\newcommand{\Z}{\mathbb{Z}}
% \providecommand, not \newcommand: some papers define these in their own
% preamble after \input{preamble}, and must not clash.
\providecommand{\Q}{\mathbb{Q}}
\providecommand{\res}{\operatorname{res}}   % resultant; used by epiRes.tex
\newcommand{\MN}{\mathbb{N}}
\newcommand{\MZ}{\mathbb{Z}}
\newcommand{\realpos}{\mathbb{R}_{> 0}}
\newcommand{\realnneg}{\mathbb{R}_{\geq 0}}
\newcommand{\naturalpos}{\mathbb{N}_{>0}}
\newcommand{\Rplus}{\mathbb{R}_+}
\newcommand{\Rp}{\mathbb{R}_{\ge 0}}
\def\Rnne{\mathbb{R}_{\ge 0}}
\def\Se{\bar{S}} \def\Ie{\bar{I}}
\newcommand{\Rpos}{\mathbb{R}_{>0}}
\newcommand{\dd}{\,\mathrm d}
\newcommand{\e}{\mathrm e}
%% ========================================
%% GREEK LETTERS (2+ character abbreviations) greek
%% ========================================
\def\al{\alpha}\def\ba{\bff \alpha}
\def\w{\omega}\def\Lgn{\ga+\mu+\de}\def\om{\omega}
\def\de{\delta}\def\De{\Delta}\def\bde{\bff \delta}
\def\ep{\epsilon}\def\eps{\varepsilon}
\def\f{\varphi}  % Kept as requested
\def\ga{\gamma}\def\Ga{\Gamma}\def\G{\Gamma}\def\bg{\bff \gamma}
\def\up{\upsilon}  % NOT the breve accent: no accent commands anywhere.
                   % Named \up, not \u, so the LaTeX breve stays reachable.
\def\la{\lambda}\def\La{\Lambda}
\def\si{\sigma}\def\ka{\kappa}
\def\th{\theta}\def\vt{\vartheta}
\def\ze{\zeta}\def\Gap{\Gamma^{+}} \def\Gam{\Gamma^{-}}


%% ========================================
%% SPECIAL SHORT COMMANDS (kept as requested)
%% ========================================
\def\I{\infty}
\def\T{\widetilde}
\def\H{\hat}
\newcommand{\relint}{\operatorname{relint}}

%% ========================================
%% BOLD/VECTOR NOTATION
%% ========================================
\newcommand{\bff}[1]{{\mbox{\boldmath$#1$}}}
\newcommand\X{\boldsymbol{X}}
\def\x{\boldsymbol{x}}\def\bx{\bff x}\def\xe{\bar{\x}}
\newcommand\y{\boldsymbol{y}}\def\by{\boldsymbol{y}}\def\ye{\bar{\y}}
\def\yb{\y}\def\xb{\x}
\def\xy{\x-\y}
\newcommand\br{\boldsymbol{r}}\def\re{\bar{\r}}
\newcommand\bd{\boldsymbol{d}}
\newcommand\bQ{\boldsymbol{Q}}
\newcommand{\Yb}{{\mathbf{Y}}}
\newcommand{\Ab}{\mathbf{A}}
\newcommand{\zb}{\mathbf{z}}\def\bz{\bff 0}
\newcommand{\Zb}{\mathbf{Z}}
\newcommand{\bb}{\mathbf{b}}
\def\vbe{\vec {\bff \beta}}
\def\vb{\vec {\bff b}}
\def\v1{\vec {\bff 1}}
\def\vn{\vec \de}
\def\vz{\vec 0}
\def\vi{\overset{\rightarrow}{\mathsf i}}
\def\vv{\overset{\rightarrow}{v}}
\def\vr{\; \vec {\mathsf r}}
\newcommand{\va}{\bff a}
\providecommand{\vect}[1]{\boldsymbol{#1}}   % bold vector (moved from Gradostat.tex)
\providecommand{\one}{\mathbf 1}             % all-ones vector (moved from AliSF)
\providecommand{\spr}{s}                     % spectral bound  (moved from AliSF)
\providecommand{\Ke}{K_{\rm eff}}            % effective carrying capacity
% Critical right/left eigenvectors in the Hopf projection method (AAHOSN/New.tex)
\providecommand{\bq}{\mathbf{q}}
\providecommand{\bp}{\mathbf{p}}


%% ========================================
%% MATHCAL & MATHSCR
%% ========================================
% Mathcal
\newcommand\mc[1]{\mathcal{#1}}
\def\mC{{\mathcal C}}\def\mD{\mathcal D}\def\mE{\mathcal{E}}
\def\mF{\mathcal F}\def\mG{{\mathcal G}}\def\mH{{\mathcal H}}
\def\mI{\mathcal I}\def\mL{\mathcal{L}}\def\mM{{\mathcal M}}
\def\mN{{\mathcal N}}\def\mP{\mathcal{P}}\def\mR{{\mathcal R}}
\def\mS{{\mathcal S}}\def\mT{\mathcal{T}}\def\mV{{\mathcal V}}
\def\mK{\mathcal K}\def\mA{{\mathcal A}}\def\mU{{\mathcal U}}
% Symmetric bilinear and trilinear forms of the Taylor expansion, used in the
% first Lyapunov coefficient (AAHOSN/New.tex)
\providecommand{\Bform}{\mathcal{B}}
\providecommand{\Cform}{\mathcal{C}}

% Blackboard bold.  Kept here, not in the drivers: a \providecommand in a
% driver preamble runs before this file and silently wins, so the same name
% ends up meaning different things in different builds.
\def\IC{\mathbb{C}}\def\IN{\mathbb{N}}\def\IQ{\mathbb{Q}}
\def\IR{\mathbb{R}}\def\IZ{\mathbb{Z}}\def\IG{\mathbb{G}}
\def\II{\mathbb{I}}

% Mathscr
\newcommand{\Rs}{{\mathscr R}}
\newcommand{\Ss}{{\mathscr S}}
\newcommand{\Ns}{{\mathscr N}}
\newcommand{\Cs}{\mathscr C}


%% ========================================
%% OPERATORS & FUNCTIONS
%% ========================================
\DeclareMathOperator{\tr}{tr}
\DeclareMathOperator{\diag}{Diag}
\DeclareMathOperator{\Diag}{Diag}
\DeclareMathOperator{\intr}{\mathrm{Int}}
\DeclareMathOperator{\Inv}{Inv}
\DeclareMathOperator{\spec}{spec}
\DeclareMathOperator{\AUC}{AUC}               % area under the curve
\newcommand{\ii}{\mathrm{i}}                  % imaginary unit
\newcommand\op{\operatorname}
\newcommand{\supp}{\op{supp}}
\newcommand{\diff}{\mathop{}\mathopen{}\mathrm{d}}
\newcommand{\pdd}[1]{\frac{\partial}{\partial{#1}}}
\DeclareMathOperator{\cone}{cone}
\DeclareMathOperator{\proj}{proj}

%% ========================================
%% BRACKETS & NORMS
%% ========================================
\providecommand{\pr}[1]{\left(#1\right)}
\providecommand{\pp}[1]{\left[#1\right]}
\newcommand\croc[1]{\left\langle #1\right\rangle}
\newcommand{\norm}[1]{\|#1\|}
\newcommand{\inner}[2]{\langle #1, #2 \rangle}
% Short synonym of \inner, used by AAHOSN/New.tex as $\ip{u}{v}=\bar u^\top v$
\providecommand{\ip}[2]{\langle #1 , #2 \rangle}


%% ========================================
%% FORMATTING & UTILITIES
%% ========================================
\newcommand*{\Scale}[2][4]{\scalebox{#1}{$#2$}}
\newcommand{\te}[1]{\text{#1}}
\newcommand\steq[1]{\stackrel{\text{\rm #1.}}{=}}
\newcommand\blue[1]{\textcolor{blue}{#1}}
\newcommand{\red}{\textcolor[rgb]{1.00,0.00,0.00}}
\newcommand{\green}{\textcolor[rgb]{0.00,1.00,0.00}}

% Figure command
\newcommand{\figu}[3]{%
\begin{figure}[H]
\centering
\includegraphics[scale=#3]{#1}
\caption{#2\label{f:#1}}
\end{figure}
}

% Footnote
\long\def\symbolfootnote[#1]#2{%
\begingroup
\def\thefootnote{\fnsymbol{footnote}}\footnote[#1]{#2}%
\endgroup}
\def\fn{\symbolfootnote}

\newcommand{\dsum}{\displaystyle\sum}

%% ========================================
%% CRN & EPIDEMIOLOGY SPECIFIC
%% ========================================
% CRN notation
\def\SS{\mathcal S}
\def\CC{\mathcal C}
\def\RR{\mathcal R}

% Reproduction numbers
\def\Rto{R_{1}^{E_2}}
\def\Rot{R_{2}^{E_1}}
\def\Roc{R_{1,c}}
\def\Rtc{R_{2,c}}
\def\rd{r_{0}}
\def\sd{s_{0}}
\def\dd{D_{0}}
\def\vd{v_{0}}
\def\zd{z_{0}}
\def\se{\sE}
\def\xd{x_{0}}
%% ========================================
%% SPECIAL VALUES & EQUILIBRIA
%% ========================================
\def\eE{\bar{e}}
\def\wE{\bar{w}}
\def\sE{s_{*}}
\def\Me{\T M_{ee}}
\def\Rnne{\mathbb{R}_{\ge 0}}

\newcommand{\ones}{\mathbf 1}


%% ========================================
%% ABBREVIATIONS (alphabetical by category)
%% ========================================

% -- General mathematical terms --
\def\adm{admissible}\def\admy{admissibility}
%\def\and{antisymmetric}% disabled 2026-09-30: it overrode LaTeX's author separator \and (printed "antisymmetric" between authors)
\def\ass{assumption}
\def\cadlag{c\`adl\`ag}
\def\ch{characteristic polynomial}
\def\chp{characteristic polynomial}
\def\coe{coefficient}
\def\comp{computation}
\def\con{condition}
\def\coo{coordinate}
\def\corr{corresponding}
\def\dec{decomposition}
\def\den{denominator}
\def\dete{deterministic}
\def\eig{eigenvalue}\def\eval{eigenvalue}
\def\eigv{eigenvector}\def\eiv{eigenvector}\def\evec{eigenvector}
\def\elt{element}
\def\enn{essentially nonnegative/positivity preserving}
\def\equ{equation}
\def\etal{et al.}
\def\fe{for example}
\def\fno{from now on}
\def\for{\forall}
\def\form{formula}
\def\Fr{Furthermore, }
\def\frt{furthermore}
\def\fun{function}
\def\gene{generalization}
\def\heu{heuristic}
\def\how{however}\def\How{However, }
\def\ie{i.e. }
\def\im{\item}
\def\isd{is defined by}
\def\itc{it is easy to check that}
\def\itf{it follows that}\def\Itf{It follows that}
\def\ith{it holds that}
\def\itm{it may be checked that}\def\Itm{It may be checked that}
\def\ito{it turns out that}\def\Ito{It turns out that}\def\Itou{It turns out useful to}
\def\ma{mass action}\def\MAR{mass action representation}
\def\mat{matrix}
\def\mbc{may be checked to be}
\def\mbw{may be written as}\def\mubw{must be written as}
\def\mpr{more precisely, }\def\Mp{More precisely, }
\def\nc{non-catalytic}
\def\nne{non-negative}\def\nny{nonnegativity}
\def\no{\nonumber}
\def\noni{non-infectious}
\def\nota{notation}
\def\np{\newpage}
\def\Nt{Note that}
\def\nt{non trivial}\def\nz{non zero}
\def\oth{otherwise}
\def\Oth{On the other hand}
\def\pa{parameter}\def\para{parameter}
\def\pol{polynomial}
\def\prob{probability}
\def\Prf{{\bf  Proof:}}\def\prf{{\bf  Proof:}}
\def\prop{proportional to}
\def\qu{\quad}
\def\QED{\hfill {$\square$}\goodbreak \medskip}
\def\rep{representation}
\def\resp{respectively}
\def\satd{satisfied}\def\satg{satisfying}\def\sats{satisfies}\def\saty{satisfy}
\def\ssec{\subsection}\def\sssec{\subsubsection}
\def\sef{see for example}
\def\sksy{skew-symmetric}
\def\sub{substrate}
\def\sys{system}
\def\var{variable}
\def\wk{well-known}
\def\wl{without loss of generality}\def\wlo{w.l.o.g.}
\def\wmc{We may check that}\def\Wmc{We may check that}
\def\wmr{we may rewrite}
\def\wms{we must show that}
\def\wrt{with respect to}

% -- CRN & Chemistry --
\newcommand\CRN{chemical reaction network}
\newcommand\CRNs{chemical reaction networks}
\newcommand\RN{reaction network}
\def\ACR{absolute concentration robustness}
\def\AMSD{acyclic siphon decomposition}
\def\DAE{differential algebraic equations}
\def\FHJ{Feinberg-Horn-Jackson}
\def\GMAK{generalized mass action kinetics}
\def\MM{Michaelis-Menten}
\def\sm{stoichiometric matrix}\def\Sm{stoichiometric matrix}
\def\WRZD{weakly reversible, zero deficiency}\def\aut{autocatalytic}

% -- Epidemiology --
\def\brn{basic reproduction number}\def\bzn{basic replacement number}
\def\com{compartment}
\def\dfe{disease free equilibrium}\def\DFE{disease free equilibrium}
\def\dwt{dwell times}\def\edt{expected dwell times}
\def\EM{epidemic model}
\def\endp{endemic fixed point}\def\enp{endemic point}\def\EE{endemic point}
\def\ME{mathematical epidemiology}
\def\MEE{eco-epidemiology}\def\EEP{eco-epidemiological}
\def\ece{eco-epidemiological model}
\def\mgm{multi group model}
\def\mtS{multi-strain model}\def\mSt{multi-strain model}
\def\mSi{minimal siphon}
\def\NGM{next generation matrix}
\def\nim{new infections matrix}
\def\qem{quadratic epidemiological model}
\def\repF{invasion functions}
\def\SPF{SIR-PH-FA model}
\def\sus{susceptible}

% -- Graph theory & Linear algebra --
\def\jac{Jacobi matrix}
\def\jin{Jacobian matrix of the invasion vector field with respect to the invasion variables}
\def\jrn{Jacobian matrix of the resident vector field with respect to the resident variables}
\def\jxn{Jacobian matrix of the $\x$ vector field with respect to the $\x$ variables}
\def\Lm{Laplacian matrix}
\def\MTT{matrix tree theorem}
\def\SGT{spectral graph theory}
\def\SMW{Sherman–Morrison–Woodbury}

% -- Stability & Dynamics --
\def\absC{absorbtion CTMC}
\def\BTb{Bogdanov-Takens bifurcation}
\def\CEP{competitive exclusion partition}
\def\CTMC{continuous time Markov chain}
\def\ds{differential system}
\def\fp{fixed point}\def\bfp{boundary fixed point}\def\fbp{fixed boundary point}
\def\fbl{fixed boundary line}\def\fif{forward invariant face}
\def\gas{global asymptotic stability}\def\GAS{globally asymptotically stable}
\def\GASS{\text{GASS}}
\def\Hb{Hopf bifurcation}
\def\Hur{Hurwitz}\def\MH{Metzler-Hurwitz}\def\RH{Routh-Hurwitz}
\def\IS{\text{IS}}
\def\LAS{locally asymptotically stable}
\def\LdS{Lyapunov diagonally stable}
\def\Lf{Lyapunov function}\def\LyV{Lyapunov Volterra type function}\def\RLF{robust Lyapunov function}
\def\LyM{Lyapunov Malkin stable}
\def\PF{Perron-Frobenius}
\def\spfp{strictly positive fixed point}
\def\sta{stability}
\def\STP{sharp threshold property}\def\STT{sharp threshold theorem}

% -- Models (with \varphi explicitly instead of \f) --
\def\AB{$(A,B)$ Arino-Brauer epidemic models}
\def\ABa{$(A,\bbe,\ba,\varphi)$ balanced bilinear model}
\def\ABP{%$(A,\bbe,P,\varphi)$
\\ balanced bilinear model}
\def\AMG{Anderson May Gupta}
\newcommand\BIN{biological interaction network}
\def\Dt{Descartes type}
\def\gLV{generalized Lotka-Volterra model}
\def\IG{\text{invasion graph}}\def\IGMS{interaction graph on minimal siphons}
\def\LCP{\text{linear complementarity problem}}
\def\LV{Lotka-Volterra}\def\LVK{Lotka-Volterra-Kolmogorov}

% -- Splitting & Structure --
\def\reg{regular splitting}\def\regS{{\bf regular splitting}}
\def\RUR{rational univariate representation}
\def\bds{block-diagonal structure}\def\NGMs{next generation matrices}

% -- Miscellaneous --
\def\Ma{{Mathematica}}
\def\Mm{Max–Min}
\def\num{numerator}
\def\rap{rational approximation} \def\DSR{directed species-reaction bipartite graph}\def\PCP{polynomial complementarity problem}\def\rat{rational substitution}
\def\QSS{quasi steady state} \def\RHS{right hand side} \def\LHS{left hand side}
\def\OSN{online social network}\def\VdD{Van den Driessche}
\def\wli{$\omega$-limit} \def\rk{rich kinetics}\def\sJ{symbolic Jacobian}
\def\cV{\cite{Vas,VasSta,BSV,GSV}} \def\spf{semi-parametric semi-functional}
\def\CRW{Capasso-Ruan-Wang} \def\ABP{$(A,\bbe,P,\f)$ balanced bilinear model}
\def\ABa{%$(A,\bbe,\ba,\f)$ \\
rank-one balanced bilinear model}
\def\bbe{B%{\bff \beta}
} \def\vbe{\vec {\bff \beta}
}
\def\vb{\vec {\bff b}}
\newcommand{\Kh}{\hat{K}}          % K with hat, first mention only
\newcommand{\wS}{w^{\sigma}}       % right Perron J^\perp, full block
\newcommand{\piS}{\vec \pi^{\sigma}}    % left Perron J^\perp, full block
\newcommand{\ws}{w^{\sigma}}       % right Perron J^\perp, sub-block
\newcommand{\pis}{\pi^{\sigma}}    % left Perron J^\perp, sub-block
\newcommand{\hwS}{\hat{w}^{\sigma}}   % right Perron K, full block
\newcommand{\hpiS}{\hat{\pi}^{\sigma}} % left Perron K, full block
\newcommand{\hws}{\hat{w}^{\sigma}}   % right Perron K, sub-block
\newcommand{\hpis}{\hat{\pi}^{\sigma}} % left Perron K, sub-block
\newcommand{\twS}{\tilde{w}^{\sigma}}  % right Perron \tilde K, full block
\newcommand{\tpiS}{\tilde{\pi}^{\sigma}} % left Perron \tilde K
\newcommand{\tws}{\tilde{w}^{\sigma}}
\newcommand{\tpis}{\tilde{\pi}^{\sigma}}
\newcommand{\Mat}[2]{\mathbb{R}_{\ge 0}^{\,#1\times #2}}
\def\tJ{J_\sigma^\perp} \def\lPe{left Perron eigenvector}
% \G is \Gamma (defined at :795).  The Volterra function keeps its own name,
% so that \G no longer has two conflicting meanings in this file.
\def\GV{G}  % Volterra G(u)=u-1-\ln u
\def\PVL{Perron-Volterra Lyapunov minimal siphon ansatz}
\newcommand{\istar}[1]{i_{#1}^{*}}
%\newcommand{\sE}{s_E}
\newcommand{\vE}{v_E}
\def\PV{Perron-Volterra candidate Lyapunov function}
\def\GLV{generalized Lotka-Volterra}
\def\MERN{mathematical epidemiology reaction networks}
\def\CMM{contraction metric method}
\newcommand{\sign}{\operatorname{sign}}
\newcommand{\Var}{\operatorname{Var}} 

% ======================================================================
% IDEAS SHARED BETWEEN THE ABSTRACT AND THE MAIN RESULTS SECTION
%
% Write each idea ONCE, in the abstract; the Main Results section recalls it,
% so the two can never drift apart.
%
%   in the abstract        \ideaitem{1}{The first-order reaction terms ...}
%   in Main Results        \ideahead{1}
%
% \ideaitem stores the text globally and typesets
%   \item \hyperlink{idea:1}{<text>}
% \ideahead typesets
%   \item\hypertarget{idea:1}{}\textbf{<the same text>}
%
% The abstract must come before the Main Results input, which is the normal
% order. If a heading has to read differently from the abstract sentence (for
% instance a long abstract item summarised in one line), override it with the
% optional argument: \ideahead[short heading]{6}.
% ======================================================================
\makeatletter
\newcommand{\ideaitem}[2]{%
  \expandafter\gdef\csname @idea@#1\endcsname{#2}%
  \item\hyperlink{idea:#1}{#2}}
\newcommand{\ideahead}[2][]{%
  \item\hypertarget{idea:#2}{}%
  \def\@ideatmp{#1}%
  \ifx\@ideatmp\@empty
    \@ifundefined{@idea@#2}{%
      \PackageWarning{preamble}{Idea #2 is used in Main Results but was never
        defined by \string\ideaitem\space in the abstract}%
      \textbf{[idea #2 undefined]}%
    }{\textbf{\begingroup\def\ideaterm##1##2{##2}\def\idearest##1{}\@nameuse{@idea@#2}\endgroup}}%
  \else
    \textbf{#1}%
  \fi}
% \ideaterm{<hypertarget>}{<words>}: inside an abstract idea the words become an
% invisible link (the document sets \hypersetup{hidelinks}) jumping to the place
% in Main Results that defines them; inside the Main Results heading itself,
% where the target sits a few lines below, \ideahead strips the link and prints
% the words plainly. Mark the destination with \hypertarget{<name>}{} on the line
% before the display, which adds no equation number and shows nothing.
\newcommand{\ideaterm}[2]{\hyperlink{#1}{#2}}
% \idearest{<text>}: the part of an abstract idea after its first sentence; printed in
% the abstract, omitted from the Main Results heading.
\newcommand{\idearest}[1]{#1}
\makeatother

% ======================================================================
% SHORT-PAPER SWITCH
% \shorttrue in a master that assembles only part of the companion paper
% (Barrat/SipArxiv.tex). \nsh{...} drops text that cites results living in the
% omitted sections, so the short build has no dangling references while the
% companion paper keeps the text verbatim.
% ======================================================================
\newif\ifshort \shortfalse
\newcommand{\nsh}[1]{\ifshort\else#1\fi}
% \shortalt{<short build>}{<companion paper>}: use when a sentence needs
% different wording in the two builds, not merely a deletion.
\newcommand{\shortalt}[2]{\ifshort#1\else#2\fi}
% \Jtr, \JtrS{\Sigma}: transverse Jacobian. J^\perp here; SipArxiv.tex redefines them as J_{\mathrm{tra}}, J_{\mathrm{tra},\Sigma}
\newcommand{\Jtr}{J^\perp}
\newcommand{\JtrS}[1]{J_{#1}^\perp}
.
\makeatletter
\AtBeginDocument{\@ifundefined{hypersetup}{}{\hypersetup{linktoc=all}}}
\makeatother
\usepackage{graphicx,float,fancybox}
\usepackage{epsfig,epstopdf,tikz}
\usepackage[sans]{dsfont}
\usetikzlibrary{positioning,arrows,arrows.meta,calc}
% Configure listings for Mathematica code
\lstdefinestyle{mathematica}{
 language=Mathematica,
 basicstyle=\small\ttfamily,
 backgroundcolor=\color{gray!10},
 frame=single,
 breaklines=true,
 commentstyle=\color{green!60!black},
 keywordstyle=\color{blue},
 stringstyle=\color{red},
 showstringspaces=false,
 tabsize=2
}
\usepackage{tcolorbox}
% Colored boxes for key results
\newtcolorbox{keyresult}[1][]{
 colback=blue!5!white,
 colframe=blue!75!black,
 title=#1,
 fonttitle=\bfseries
}
\newtcolorbox{examplebox}[1][]{
 colback=green!5!white,
 colframe=green!75!black,
 title=#1,
 fonttitle=\bfseries
}


\makeatletter
\@ifundefined{theorem}   {\newtheorem{theorem}   {Theorem}}   {}
\@ifundefined{definition}{\newtheorem{definition}{Definition}} {}
\@ifundefined{open}      {\newtheorem{open}      {Problem}}   {}
\@ifundefined{example}   {\newtheorem{example}   {Example}}   {}
\makeatother

% COUNTER DEFINITIONS
\newcounter{tabcounter}
\newcounter{defcounter}
\newcounter{opcounter}
\newcounter{thmcounter}
\newcounter{excounter}
\newcounter{ascounter}
\newcounter{cjcounter}
\newcounter{lemcounter}
\newcounter{remcounter}
\newcounter{corcounter}
\newcounter{propcounter}
\newcounter{figcounter}

% UNIFIED GLOSSARY STYLE
\newglossarystyle{unifiedstyle}{%
  \setglossarystyle{list}%
  \renewcommand*{\glossarypreamble}{\gdef\glslastsec{}}%
  \renewcommand{\glossentry}[2]{%
    \edef\glscursec{\glsentryuseriii{##1}}%
    \ifx\glscursec\glslastsec\else
      \item[]\textbf{\S\,\glscursec\enspace\glsentryuseriv{##1}}%
      \global\let\glslastsec\glscursec
    \fi
    \item[\glsentryuserii{##1} \glsentryuseri{##1}:] ``\glossentryname{##1}'' \glsentrydesc{##1}%
  }%
}

\makeatletter
% First-word extractor (used by \beT, \beD and \beXa for the default key)
% \@firstwordof hello world \relax  ->  hello
\def\@firstwordof#1 #2\relax{#1}
% First three letters, used for the heading mnemonic of a theorem given no
% explicit key.  Padding with \@empty keeps it safe for first words shorter
% than three letters: \@empty vanishes under \edef, whereas \relax would
% survive into the label and is illegal inside \csname (it broke every
% two-letter title, e.g. "ME ...", with "Missing \endcsname inserted").
\def\@takethree#1#2#3#4\@nil{#1#2#3}

% EXAMPLE COMMAND -- same signature and behaviour as \beT and \beD.
% Usage: \beXa*[long title][short title][key]
%   star  = suppresses the printed "[key]" tag after the title (as for \beT*
%   and \beD*); the label and glossary entry are unaffected.
%   key = optional label key override; the label becomes e:KEY, and without a
%   key it is e: plus the first three letters of the title's first word.  A
%   mnemonic label lets an example be cited as "Example~\ref{e:KEY}" instead
%   of by section number, which is what the citing prose needs.
\NewDocumentCommand{\beXa}{s O{} O{} o}{%
  \def\@temparg{#2}%
  \ifx\@temparg\@empty
    \begin{example}%
  \else
    \let\@savedSecTitle\@currentsectionname%
    \refstepcounter{excounter}%
    \IfValueTF{#4}{%
      \edef\@exlblkey{#4}%
    }{%
      \edef\@exfirstword{\@firstwordof#2 \relax}%
      \edef\@exlblkey{%
        \expandafter\@takethree\@exfirstword\@empty\@empty\@empty\@nil}%
    }%
    \def\@tempshort{#3}%
    \ifx\@tempshort\@empty
      \expandafter\xdef\csname @xshort@\@exlblkey\endcsname{\@exlblkey}%
    \else
      \expandafter\def\csname @xshort@\@exlblkey\endcsname{#3}%
    \fi
    \IfBooleanTF{#1}{\let\@extag\@empty}{%
      \edef\@extag{\noexpand\ \noexpand\textup{[\noexpand\texttt{e:\@exlblkey}]}}}%
    \begin{example}[#2\ifshowkeys\@extag\fi]%
    \label{e:\@exlblkey}%
    % record the long title so \nref{e:KEY} prints it, as for theorems
    \protected@write\@auxout{}{%
      \string\@thmtitle{e:\@exlblkey}{\unexpanded{#2}}}%
    \expanded{%
      \noexpand\newglossaryentry{e\theexcounter}{%
        name={\unexpanded{#2}%
          \IfBooleanF{#1}{ [\noexpand\texttt{e:\@exlblkey}]}},%
        sort={\theexample},%
        description={p.~\noexpand\pageref{e:\@exlblkey}},%
        user1={\theexample},%
        user2={Example},%
        user3={\thesection},%
        user4={\expandafter\unexpanded\expandafter{\@savedSecTitle}}%
      }%
    }%
    \protected@write\@auxout{}{\string\glsliveentry{e\theexcounter}}%
  \fi
}
\newcommand{\eeXa}{\end{example}}

% DEFINITION COMMAND -- same signature and behaviour as \beT below.
% Usage: \beD[long title][short title][key]
%   long title  = definition heading (plain text only)
%   short title = shown in \dref{word}; defaults to the key
%   key         = optional label key override; the label becomes d:KEY
% Without a key the label is d: plus the first three letters of the title's
% first word, so it is a mnemonic rather than the definition's number: a
% numeric label would silently change whenever a definition is inserted.
\NewDocumentCommand{\beD}{s O{} O{} o}{%
  \def\@temparg{#2}%
  \ifx\@temparg\@empty
    \begin{definition}%
  \else
    \let\@savedSecTitle\@currentsectionname%
    \refstepcounter{defcounter}%
    \IfValueTF{#4}{%
      \edef\@deflblkey{#4}%
    }{%
      \edef\@deffirstword{\@firstwordof#2 \relax}%
      \edef\@deflblkey{%
        \expandafter\@takethree\@deffirstword\@empty\@empty\@empty\@nil}%
    }%
    \def\@tempshort{#3}%
    \ifx\@tempshort\@empty
      \expandafter\xdef\csname @dshort@\@deflblkey\endcsname{\@deflblkey}%
    \else
      \expandafter\def\csname @dshort@\@deflblkey\endcsname{#3}%
    \fi
    \IfBooleanTF{#1}{\let\@deftag\@empty}{%
      \edef\@deftag{\noexpand\ \noexpand\textup{[\noexpand\texttt{d:\@deflblkey}]}}}%
    \begin{definition}[#2\ifshowkeys\@deftag\fi]%
    \label{d:\@deflblkey}%
    % record the long title so \nref{d:KEY} prints it, as for theorems
    \protected@write\@auxout{}{%
      \string\@thmtitle{d:\@deflblkey}{\unexpanded{#2}}}%
    \expanded{%
      \noexpand\newglossaryentry{d\thedefcounter}{%
        name={\unexpanded{#2}%
          \IfBooleanF{#1}{ [\noexpand\texttt{d:\@deflblkey}]}},%
        sort={\thedefinition},%
        description={p.~\noexpand\pageref{d:\@deflblkey}},%
        user1={\thedefinition},%
        user2={Definition},%
        user3={\thesection},%
        user4={\expandafter\unexpanded\expandafter{\@savedSecTitle}}%
      }%
    }%
    \protected@write\@auxout{}{\string\glsliveentry{d\thedefcounter}}%
  \fi
}
\newcommand{\eeD}{\end{definition}}

% \showkeysfalse in a driver hides the [t:KEY] mnemonics in result headings
% (labels and glossary entries are unchanged); used for submission builds.
\newif\ifshowkeys \showkeystrue
% \commonredtrue in a driver that inputs basicFiles/commonRed: tool.tex then
% cites its reduction results instead of restating them.
\newif\ifcommonred \commonredfalse
% TheorEM COMMAND
% Usage: \beT[long title][short title][key]
%   long title  = theorem heading (plain text only)
%   short title = shown in \tref{word}; may contain math; defaults to first word of long title
%   key         = optional label key override; label becomes t:KEY (avoids first-word collisions)
% Label is t:WORD where WORD = first word of long title (or KEY if given); use \tref{WORD} to reference.
\NewDocumentCommand{\beT}{s O{} O{} o}{%
  \def\@temparg{#2}%
  \ifx\@temparg\@empty
    \begin{theorem}%
  \else
    \let\@savedSecTitle\@currentsectionname%
    \refstepcounter{thmcounter}%
    % \@thmtag is the mnemonic shown in the heading; every theorem gets one.
    % With an explicit key the key itself is shown.  Without one the key
    % defaults to the title's first word, and showing that whole word would
    % merely repeat the title ("Dissipativity [Dissipativity]"), so only its
    % first three letters are shown: "Dissipativity [Dis]".  The label is
    % unaffected -- it still uses the full \@thmlblkey.
    \IfValueTF{#4}{%
      \edef\@thmlblkey{#4}%
    }{%
      \edef\@thmfirstword{\@firstwordof#2 \relax}%
      \edef\@thmlblkey{%
        \expandafter\@takethree\@thmfirstword\@empty\@empty\@empty\@nil}%
    }%
    \IfBooleanTF{#1}{\let\@thmtag\@empty}{%
      \edef\@thmtag{\noexpand\ \noexpand\textup{[\noexpand\texttt{t:\@thmlblkey}]}}}%
    \def\@tempshort{#3}%
    \ifx\@tempshort\@empty
      \expandafter\xdef\csname @tshort@\@thmlblkey\endcsname{\@thmlblkey}%
    \else
      \expandafter\def\csname @tshort@\@thmlblkey\endcsname{#3}%
    \fi
    \begin{theorem}[#2\ifshowkeys\@thmtag\fi]%
    \label{t:\@thmlblkey}%
    % record the long title for \nref; via the .aux so forward refs resolve
    \protected@write\@auxout{}{%
      \string\@thmtitle{t:\@thmlblkey}{\unexpanded{#2}}}%
    \expanded{%
      \noexpand\newglossaryentry{t\thethmcounter}{%
        name={\unexpanded{#2}%
          \IfBooleanF{#1}{ [\noexpand\texttt{t:\@thmlblkey}]}},%
        % sort/user1 must be the number the document actually prints, i.e. the
        % theorem environment's counter, not thmcounter.  thmcounter advances
        % only on \beT with a title, so a bare \beT (or \beT with an
        % empty title) makes the two drift and the glossary cites a wrong
        % number.  thmcounter is still fine as the entry key: it is unique.
        sort={\thetheorem},%
        description={p.~\noexpand\pageref{t:\@thmlblkey}},%
        user1={\thetheorem},%
        user2={Theorem},%
        user3={\thesection},%
        user4={\expandafter\unexpanded\expandafter{\@savedSecTitle}}%
      }%
    }%
    \protected@write\@auxout{}{\string\glsliveentry{t\thethmcounter}}%
  \fi
}
\newcommand{\eeT}{\end{theorem}}

% CROSS-PART RELATION PARAGRAPH
% Usage: \xpart{text}  -- place immediately BEFORE \beT, never inside it
%   (anything inside \beT...\eeT would be swallowed into the glossary entry).
% The paragraph is typeset only when Part I is actually present, i.e. in the
% combined book build; in the standalone builds (GranadaDSABNS/A.tex,
% AAHRah-Submission/AAHRah.tex) the label part:glv is undefined and the whole
% paragraph disappears rather than leaving ?? cross-references. Same
% graceful-degradation principle as \GASref in GASt.tex. Resolves on the
% second LaTeX pass, since r@part:glv comes from the .aux file.
% (Already inside the \makeatletter region opened above, so no \makeatletter
% pair here -- adding one would close that region early.)
\newcommand{\xpart}[1]{\@ifundefined{r@part:glv}{}{\paragraph{Relation with Parts~I--II.}#1}}

% \bookref{label}{precise text}{fallback text}: prints the precise cross-reference
% when the label resolves (combined book build), and the fallback text otherwise
% (standalone paper builds, where book-only appendix labels are absent). Same
% graceful-degradation role as \GASref in GASt.tex, but usable in any file since it
% lives in this preamble (loaded by every driver) -- unlike \GASref, which is defined in a
% content file and is therefore undefined in standalone builds that omit that file.
\newcommand{\bookref}[3]{\@ifundefined{r@#1}{#3}{#2}}

% \PVpart and \RBpart: the two parts that a partial build may omit.
% They print "Part~N" in any build that contains the part, and a naming phrase
% in a build that does not, so an excerpt shows no ?? .  Each already supplies
% the word "Part", so write "\PVpart develops ...", not "Part~\PVpart".
\newcommand{\PVpart}{\bookref{part:aahrah}%
  {Part~\ref{part:aahrah}}{the Perron--Volterra part of the full book}}
\newcommand{\RBpart}{\bookref{part:recbrowne}%
  {Part~\ref{part:recbrowne}}{the immuno-virology part of the full book}}

% OPEN PROBLEM COMMANDS
% ASSUMPTION COMMANDS (glossaried)
\newcommand{\beAs}[1][]{%
  \def\@temparg{#1}%
  \ifx\@temparg\@empty
    \begin{assumption}%
  \else
    \let\@savedSecTitle\@currentsectionname%
    \refstepcounter{ascounter}%
    \begin{assumption}[#1]%
    \label{as:\theascounter}%
    \expanded{%
      \noexpand\newglossaryentry{as\theascounter}{%
        name={\unexpanded{#1}},%
        sort={\theascounter},%
        description={p.~\noexpand\pageref{as:\theascounter}},%
        user1={\theascounter},%
        user2={Assumption},%
        user3={\thesection},%
        user4={\expandafter\unexpanded\expandafter{\@savedSecTitle}}%
      }%
    }%
    \protected@write\@auxout{}{\string\glsliveentry{as\theascounter}}%
  \fi
}
\newcommand{\eeAs}{\end{assumption}}
% CONJECTURE COMMANDS (glossaried)
\newcommand{\beCj}[1][]{%
  \def\@temparg{#1}%
  \ifx\@temparg\@empty
    \begin{conjecture}%
  \else
    \let\@savedSecTitle\@currentsectionname%
    \refstepcounter{cjcounter}%
    \begin{conjecture}[#1]%
    \label{cj:\thecjcounter}%
    \expanded{%
      \noexpand\newglossaryentry{cj\thecjcounter}{%
        name={\unexpanded{#1}},%
        sort={\thecjcounter},%
        description={p.~\noexpand\pageref{cj:\thecjcounter}},%
        user1={\thecjcounter},%
        user2={Conjecture},%
        user3={\thesection},%
        user4={\expandafter\unexpanded\expandafter{\@savedSecTitle}}%
      }%
    }%
    \protected@write\@auxout{}{\string\glsliveentry{cj\thecjcounter}}%
  \fi
}
\newcommand{\eeCj}{\end{conjecture}}
% LEMMA COMMANDS
% ---------------------------------------------------------------------
% UNIFORM RESULT WRAPPERS -- \beL, \beC, \beP, \beR share \beT's signature
%
%   \beX[long title][short title][key]  ...  \eeX
%
%   long title  = heading (math allowed)
%   short title = shown in the cross-reference; defaults to the key
%   key         = label key.  With a key the label is X:KEY and the heading
%                 carries the mnemonic [key], exactly as for \beT.  Without
%                 one the label falls back to the running counter, which is
%                 what every pre-existing call site in the book relies on.
%
% A bare \beX -- arguments commented out -- opens the plain environment and
% creates NO label and NO glossary entry.  That is the supported way to keep
% an individual result out of the catalogue.
% ---------------------------------------------------------------------
% #1 wrapper name  #2 environment  #3 label prefix  #4 counter  #5 type name
\newcommand{\@mkresult}[5]{%
  % A leading * suppresses the [key] mnemonic in both the heading and the
  % glossary name, while still making the label X:KEY and the entry:
  %     \beX*[long][short][key]
  \expandafter\NewDocumentCommand\csname be#1\endcsname{s O{} O{} o}{%
    \def\@temparg{##2}%
    \ifx\@temparg\@empty
      \begin{#2}%
    \else
      \let\@savedSecTitle\@currentsectionname%
      \refstepcounter{#4}%
      % With a key: mnemonic label X:KEY and the [key] tag in the heading.
      % Without one: fall back to the running counter, exactly as these
      % wrappers behaved before, so every legacy one-argument call site in
      % the book keeps its label.  The counter is also csname-safe, which a
      % title's first word is not -- titles beginning with math would
      % otherwise build \csname \Omega...\endcsname and fail.
      \IfValueTF{##4}{%
        \edef\@relbl{##4}%
        \IfBooleanTF{##1}{\let\@retag\@empty}{%
          \edef\@retag{\noexpand\ \noexpand\textup{[\noexpand\texttt{#3:##4}]}}}%
      }{%
        \edef\@relbl{\arabic{#4}}%
        \let\@retag\@empty
      }%
      \def\@tempshort{##3}%
      \ifx\@tempshort\@empty
        \expandafter\xdef\csname @sshort@#3\@relbl\endcsname{\@relbl}%
      \else
        \expandafter\def\csname @sshort@#3\@relbl\endcsname{##3}%
      \fi
      \begin{#2}[##2\ifshowkeys\@retag\fi]%
      \edef\@tmplbl{#3:\@relbl}\expandafter\label\expandafter{\@tmplbl}%
      \expanded{%
        \noexpand\newglossaryentry{#3\csname the#4\endcsname}{%
          name={\unexpanded{##2}%
            \IfBooleanF{##1}{\IfValueT{##4}{ [\noexpand\texttt{#3:##4}]}}},%
          sort={\csname the#2\endcsname},%
          description={p.~\noexpand\pageref{#3:\@relbl}},%
          user1={\csname the#2\endcsname},%
          user2={#5},%
          user3={\thesection},%
          user4={\expandafter\unexpanded\expandafter{\@savedSecTitle}}%
        }%
      }%
      \protected@write\@auxout{}{\string\glsliveentry{#3\csname the#4\endcsname}}%
      \protected@write\@auxout{}{\string\@thmtitle{\@tmplbl}{\unexpanded{##2}}}%
    \fi
  }%
  \expandafter\newcommand\csname ee#1\endcsname{\end{#2}}%
}
\@mkresult{L}{lemma}{l}{lemcounter}{Lemma}
% REMARK / COROLLARY / PROPOSITION -- same signature as \beL above
\@mkresult{R}{remark}{r}{remcounter}{Remark}
\@mkresult{C}{corollary}{c}{corcounter}{Corollary}
\@mkresult{P}{proposition}{p}{propcounter}{Proposition}
% OPEN-PROBLEM COMMAND -- generated with the same signature as \beT:
%   \beO[long title][short title][key]   ...   \eeO
% label o:KEY, or o:<counter> when no key is given; a bare \beO makes
% neither a label nor a glossary entry.
\@mkresult{O}{open}{o}{opcounter}{Problem}


% Section-level title tracker: updated only by \section, not \subsection.
% \@currentlabelname is overwritten by every \refstepcounter (including
% \subsection), so we save it separately here.
\newcommand{\@currentsectionname}{}
\AtBeginDocument{%
  \let\@epidcrn@orig@refstepcounter\refstepcounter
  \renewcommand{\refstepcounter}[1]{%
    \@epidcrn@orig@refstepcounter{#1}%
    \ifstrequal{#1}{section}{\global\let\@currentsectionname\@currentlabelname}{}%
  }%
}

% TABLE COMMANDS (glossaried) -- consistent naming with \beT, \beD, etc.
\newcommand{\beTa}[1][]{%
  \def\@temparg{#1}%
  \ifx\@temparg\@empty
    \begin{table}[htbp]%
  \else
    \let\@savedSecTitle\@currentsectionname%
    \refstepcounter{tabcounter}%
    \begin{table}[htbp]%
    \label{A:\thetabcounter}%
    \expanded{%
      \noexpand\newglossaryentry{A\thetabcounter}{%
        name={\unexpanded{#1}},%
        sort={\thetabcounter},%
        description={p.~\noexpand\pageref{A:\thetabcounter}},%
        user1={\thetabcounter},%
        user2={Table},%
        user3={\thesection},%
        user4={\expandafter\unexpanded\expandafter{\@savedSecTitle}}%
      }%
    }%
    \protected@write\@auxout{}{\string\glsliveentry{A\thetabcounter}}%
  \fi
}
\newcommand{\eeTa}{\end{table}}

% FIGURE COMMANDS (glossaried) -- consistent with \beT, \beR, \beTa, etc.
% Usage:
%   \beF[Glossary title]...\caption{Visible caption}\eeF  -> glossaried figure
%   \beF...\caption{Visible caption}\eeF                  -> plain [H] figure, no glossary
\newcommand{\beF}[1][]{%
  \def\@temparg{#1}%
  \ifx\@temparg\@empty
    \begin{figure}[H]%
  \else
    \begin{figure}[H]%
    \let\@savedSecTitle\@currentsectionname%
    \refstepcounter{figcounter}%
    \label{fig:\thefigcounter}%
    \expanded{%
      \noexpand\newglossaryentry{Fig\thefigcounter}{%
        name={\unexpanded{#1}},%
        sort={\thefigcounter},%
        description={p.~\noexpand\pageref{fig:\thefigcounter}},%
        user1={\thefigcounter},%
        user2={Figure},%
        user3={\thesection},%
        user4={\expandafter\unexpanded\expandafter{\@savedSecTitle}}%
      }%
    }%
    \protected@write\@auxout{}{\string\glsliveentry{Fig\thefigcounter}}%
  \fi
}
\newcommand{\eeF}{\end{figure}}
\makeatother
% LIVE-ENTRY REGISTER
% \glsaddall prints every entry the glossaries package currently knows, and
% that set includes entries restored from \jobname.glsdefs.  The package
% rewrites .glsdefs from all *currently defined* entries, so anything it
% restored is written straight back out: a deleted section kept reappearing in
% the catalogue for ever, and the .glsdefs had to be erased by hand.  The
% wrappers therefore record every entry they actually create in the .aux, and
% \printboth adds only those.  The .aux is rebuilt from scratch on each run, so
% this is self-healing: a removed result is gone after one further run.
\makeatletter
\newcommand{\@glslive}{}
\newcommand{\@glsaddlive}[1]{\ifglsentryexists{#1}{\glsadd{#1}}{}}
\newcommand{\glsliveentry}[1]{\g@addto@macro\@glslive{\@glsaddlive{#1}}}
\makeatother
% PRINT COMMAND
% No \clearpage here -- the catalogue continues on the current page.  This has
% been re-added by mistake several times; leave it out.
\makeatletter
\newcommand{\printboth}{%
  \@glslive   % add exactly the entries this document created last run
  \printglossary[style=unifiedstyle, title={0 \hspace{.22em} Definitions, Theorems, Examples and Open Problems}]%
}
\makeatother
% Named theorem reference: \tref{word} -> Thm~N(word-or-short)
% Shows stored short title if available, otherwise shows the key itself.
\newcommand{\tref}[1]{%
  Thm~\ref{t:#1}(\ifcsname @tshort@#1\endcsname\csname @tshort@#1\endcsname\else#1\fi)%
}

% Full-title theorem reference: \nref{t:KEY} prints the number and then the
% theorem's long title, so one writes "Theorem~\nref{t:blocks}" and gets
% "Theorem 3 ``Block factorization and siphons''".  Takes the complete label, like
% \ref, not the bare key like \tref.
%
% \beT writes each title to the .aux (see \@thmtitle below), so \nref resolves
% forward references from the second LaTeX pass, exactly as \ref does.  If the
% title is not yet known -- first pass, or a label that is not a \beT theorem
% -- it degrades to the number rather than producing ??.
% These two use @ in their names, and this point in the file is outside the
% \makeatletter region opened earlier (it closes at \makeatother above), so
% the region has to be reopened here.
\makeatletter
\newcommand{\@thmtitle}[2]{\global\@namedef{@tlong@#1}{#2}}
\newcommand{\nref}[1]{%
  \ref{#1}%
  \ifcsname @tlong@#1\endcsname
    \ ``\csname @tlong@#1\endcsname''%
  \fi
}
% \nm{t:KEY}: the stored long title of any \beT/\beL/\beP/\beR/\beC/\beD/\beE result,
% verbatim and without quotes, so item names follow the macro's first argument.
\newcommand{\nm}[1]{%
  \ifcsname @tlong@#1\endcsname
    \csname @tlong@#1\endcsname%
  \else
    \ref{#1}%
  \fi
}
\makeatother




%% ========================================
%% Other THEOREM ENVIRONMENTS
%% ========================================
\makeatletter
\@ifundefined{lemma}      {\newtheorem{lemma}      {Lemma}}      {}
\@ifundefined{proposition}{\newtheorem{proposition}{Proposition}} {}
\@ifundefined{corollary}  {\newtheorem{corollary}  {Corollary}}  {}
\@ifundefined{conjecture} {\newtheorem{conjecture} {Conjecture}}  {}
\@ifundefined{remark}     {\newtheorem{remark}     {Remark}}     {}
\@ifundefined{question}   {\newtheorem{question}   {Question}}   {}
\@ifundefined{assumption} {\newtheorem{assumption} {Assumption}} {}
\@ifundefined{hypothesis} {\newtheorem{hypothesis} {Hypothesis}} {}
\@ifundefined{fact}       {\newtheorem{fact}       {Fact}}       {}
\makeatother

% \citer: cite the relay-graph threshold theorem of AAHrel. If a document has
% \input the chapter/file containing t:relayRF (so the label resolves), point
% straight at it; otherwise fall back to a bare citation of AAHrel. Shared here
% (not per-paper) so it resolves consistently across every book chapter.
% \citeR: same idea, for the first theorem of AAHRah (t:Pstar).
\makeatletter
\providecommand{\citer}{%
  \@ifundefined{r@t:relayRF}%
    {\cite{AAHrel}}%
    {Theorem~\ref{t:relayRF} of \cite{AAHrel}}%
}
\providecommand{\citeR}{%
  \@ifundefined{r@t:Pstar}%
    {\cite{AAHRah}}%
    {Theorem~\ref{t:Pstar} of \cite{AAHRah}}%
}
\makeatother


% Abbreviations
\def\beQ{\begin{question}}\def\eeQ{\end{question}}
\def\beA{\begin{assumption}}\def\eeA{\end{assumption}}


%% ========================================

%% ========================================
%% COLORS & CUSTOM COMMANDS
%% ========================================
\definecolor{funccolor}{RGB}{25,25,112}
\definecolor{desccolor}{RGB}{64,64,64}
\newcommand{\func}[1]{\texttt{\textcolor{funccolor}{#1}}}
\def\eq{=}
\def\mm{-}

%% ========================================
%% MATHCAL SHORTHANDS
%% ========================================
% Project-wide convention: \mX for a capital X is always \mathcal X.
% Defined for the whole alphabet so the convention never fails to expand.
\newcommand{\mA}{\mathcal A}\newcommand{\mB}{\mathcal B}
\newcommand{\mC}{\mathcal C}\newcommand{\mD}{\mathcal D}
\newcommand{\mE}{\mathcal E}\newcommand{\mF}{\mathcal F}
\newcommand{\mG}{\mathcal G}\newcommand{\mH}{\mathcal H}
\newcommand{\mI}{\mathcal I}\newcommand{\mJ}{\mathcal J}
\newcommand{\mK}{\mathcal K}\newcommand{\mL}{\mathcal L}
\newcommand{\mM}{\mathcal M}\newcommand{\mN}{\mathcal N}
\newcommand{\mO}{\mathcal O}\newcommand{\mP}{\mathcal P}
\newcommand{\mQ}{\mathcal Q}\newcommand{\mR}{\mathcal R}
\newcommand{\mS}{\mathcal S}\newcommand{\mT}{\mathcal T}
\newcommand{\mU}{\mathcal U}\newcommand{\mV}{\mathcal V}
\newcommand{\mW}{\mathcal W}\newcommand{\mX}{\mathcal X}
\newcommand{\mY}{\mathcal Y}\newcommand{\mZ}{\mathcal Z}

%% ========================================
%% BASIC MATH ENVIRONMENTS
%% ========================================
\def\bep{\begin{pmatrix}}\def\eep{\end{pmatrix}}
\def\bev{\begin{vmatrix}}\def\eev{\end{vmatrix}}
\def\bea{\begin{eqnarray*}}\def\eea{\end{eqnarray*}}
\def\bc{\begin{cases}}\def\ec{\end{cases}}
\def\BEN{\begin{enumerate}}\def\EEN{\end{enumerate}}
\def\BI{\begin{itemize}}\def\EI{\end{itemize}}

\newcommand{\be}[1]{\begin{equation}\label{#1}}
\newcommand{\ee}{\end{equation}}
\newcommand{\beq}{\begin{eqnarray}}
\def\eeq{\end{eqnarray}}
\def\eqr{\eqref}\def\fr{\frac}\def\lbl{\label}
\def\Lra{\Longrightarrow}\def\non{\nonumber}\def\cNGM{\cite{Diek,Van,Van08}}
\def\Eq{\Leftrightarrow}\def\b{\beta}\def\uph{\red{UP TO HERE}}
%% ========================================
%% PROBABILITY & STATISTICS
%% ========================================
\renewcommand{\P}{\ensuremath{\mathbb{P}}}  % Note: overrides paragraph symbol
\newcommand{\E}{\ensuremath{\mathbb{E}}}
\newcommand{\ind}[1]{\ensuremath{\mathbbm{1}_{\left\{#1\right\}}}}
\DeclareMathOperator{\Cov}{Cov}
\DeclareMathOperator{\TV}{TV}
\newcommand{\rank}{\operatorname{rank}}

%% ========================================
%% CRN OPERATORS (organizations/ERCs/buffering/siphons framework)
%% ========================================
\def\clos{\mathrm{clos}}
\def\Supply{\mathrm{Supply}}
\def\req{\mathrm{req}}
\def\React{\mathrm{React}}
\def\Prod{\mathrm{Prod}}
\def\Gen{\mathrm{Gen}}
\def\In{\mathrm{In}}
\def\Out{\mathrm{Out}}
\def\Act{\mathrm{Act}}
\def\Min{\mathrm{Min}}

%% ========================================
%% LINE BREAKING
%% ========================================
% Last-resort stretch, used only on paragraphs TeX cannot otherwise break:
% a long \texttt{} token, a comma-separated inline math list (commas are not
% breakpoints in math), or a run of \cite/\ref boxes.  Without it those set as
% overfull \hbox and print into the margin; with it they set slightly loose,
% which is the lesser evil.  Does not affect paragraphs that already break.
\emergencystretch=3em

%% ========================================
%% NUMBER SETS
%% ========================================
\newcommand{\R}{\mathbb{R}}  % Real numbers
\newcommand{\C}{\mathbb{C}}  % Complex numbers
\newcommand{\N}{\mathbb{N}}
\newcommand{\Z}{\mathbb{Z}}
% \providecommand, not \newcommand: some papers define these in their own
% preamble after % Shared preamble for the GLVbook papers.  This is the single canonical copy;
% a driver one level below GLVbook reaches it as
%
%     \input{../basicFiles/preamble}
%
% What stays in the driver, because it differs from paper to paper:
% \documentclass, font choice, \title, \author, \date, and any macro genuinely
% local to one paper.  Everything else belongs here.  The file is
% self-contained: it does not read any other file.

\usepackage[utf8]{inputenc}
\usepackage[novbox]{pdfsync}   % required: SyncTeX alone does not work here

% Page geometry, set once here for every paper that loads this file.  A driver
% that needs different margins says \geometry{...}; a second
% \usepackage[...]{geometry} raises an option clash and must not be used.
\makeatletter
\@ifundefined{AAdefNnogeometry}{\usepackage[a4paper,margin=2cm]{geometry}}{}
\makeatother

% mathematics
\usepackage{amsmath,amssymb,amsthm,mathtools,bm,mathrsfs,bbm}
\usepackage{centernot}

% floats, tables, lists
\usepackage{graphicx,float}
\makeatletter
\@ifpackageloaded{subfig}{}{\usepackage{subcaption}}
\makeatother
\usepackage{booktabs,longtable,threeparttable,array,multirow}
\usepackage{enumitem}

% colour, code, algorithms
\usepackage{xcolor,listings,tcolorbox}
\makeatletter
\@ifundefined{c@algorithm}{\usepackage{algorithm}}{}
\makeatother
\usepackage{algorithmicx,algpseudocode}

% diagrams
\usepackage{tikz,tikz-cd,pgfplots}
\usetikzlibrary{graphs,decorations.pathmorphing,decorations.markings,
arrows.meta,automata,positioning}
\tikzcdset{every label/.append style = {font = \small}}
\pgfplotsset{compat=1.18}

% chemistry
\usepackage[version=4]{mhchem}

\usepackage{microtype}
\usepackage{lineno}
\usepackage{nomencl}
\makenomenclature

% hyperref first, then bookmark, then glossaries: glossaries must follow
% hyperref or the entry links are built against the wrong anchors.
\makeatletter
\@ifundefined{hyperrefpackage}{%
  \@ifundefined{AAdefNnohyperref}{\usepackage{hyperref}}{}%
}{}
\@ifundefined{hyperrefpackage}{\usepackage{bookmark}}{}
\makeatother
\usepackage[toc,sort=use]{glossaries}
\makeglossaries

% Book/standalone switch.  Material that the book already covers in its
% prelude, but that a standalone paper still has to state, is wrapped as
%     \ifbookbuild\else ... \fi
% Drivers default to standalone; bookM.tex says \bookbuildtrue.
\newif\ifbookbuild
\bookbuildfalse

% Part switches for the gradostat paper.  Both default to true; a driver that
% wants the short submission says \madecfalse and/or \slowfastfalse, and every
% cross-reference into the dropped part disappears with it.
%   \ifmadec     Part III, the Lagasquie--Madec model and its break-even geometry
%   \ifslowfast  Part IV, everything carrying the small parameter \eps
\newif\ifmadec     \madectrue
\newif\ifslowfast  \slowfasttrue

% Part switch for the Barrat/simplicial-contagion project.  Defaults to
% true (the hypernetwork-synchrony part is included); the journal-paper
% driver says \hnetpartfalse to drop it and every cross-reference into it.
\newif\ifhnetpart  \hnetparttrue

%% ========================================
%% PACKAGES
%% ========================================
\RequirePackage{etoolbox}
\usepackage{makecell}
\usepackage[novbox]{pdfsync}% obsolete; SyncTeX is built into modern engines
\usepackage{amssymb, amsfonts,mathtools,bbm,eurosym,graphicx,latexsym, amsmath,amsthm,times,nccmath,url}
\usepackage{listings,verbatim,cases,tabularx}
% hyperref is deliberately NOT in the line above.  Journal classes often
% load it themselves, with a long option list, from a begindocument hook
% (mdpi.cls does: \DeclareHookRule{begindocument}{hyperref}{before}{mdpi}).
% Loading it here first, without options, makes that later load raise
% "Option clash for package hyperref".  So load it only if nobody else
% has, or will:
%   - \@ifpackageloaded  catches a class that loaded it already;
%   - \hyperrefpackage    is mdpi.cls's own deferred loader;
%   - \AAdefNnohyperref   lets any other driver opt out by defining it
%                         before \input{preamble}.
\makeatletter
\@ifpackageloaded{hyperref}{}{%
  \@ifundefined{hyperrefpackage}{%
    \@ifundefined{AAdefNnohyperref}{\usepackage{hyperref}}{}%
  }{}%
}
\makeatother
% linktoc=all makes a table-of-contents page number a link as well as the
% entry text; hyperref's default, linktoc=section, links only the text.
% Deferred to begin-document and guarded: by the block above hyperref may
% be loaded later (mdpi's own hook) or not at all (\AAdefNnohyperref), and
% \hypersetup must exist before it is called.  A driver that wants the
% default back can say \hypersetup{linktoc=section} after \input{preamble}.
\makeatletter
\AtBeginDocument{\@ifundefined{hypersetup}{}{\hypersetup{linktoc=all}}}
\makeatother
\usepackage{graphicx,float,fancybox}
\usepackage{epsfig,epstopdf,tikz}
\usepackage[sans]{dsfont}
\usetikzlibrary{positioning,arrows,arrows.meta,calc}
% Configure listings for Mathematica code
\lstdefinestyle{mathematica}{
 language=Mathematica,
 basicstyle=\small\ttfamily,
 backgroundcolor=\color{gray!10},
 frame=single,
 breaklines=true,
 commentstyle=\color{green!60!black},
 keywordstyle=\color{blue},
 stringstyle=\color{red},
 showstringspaces=false,
 tabsize=2
}
\usepackage{tcolorbox}
% Colored boxes for key results
\newtcolorbox{keyresult}[1][]{
 colback=blue!5!white,
 colframe=blue!75!black,
 title=#1,
 fonttitle=\bfseries
}
\newtcolorbox{examplebox}[1][]{
 colback=green!5!white,
 colframe=green!75!black,
 title=#1,
 fonttitle=\bfseries
}


\makeatletter
\@ifundefined{theorem}   {\newtheorem{theorem}   {Theorem}}   {}
\@ifundefined{definition}{\newtheorem{definition}{Definition}} {}
\@ifundefined{open}      {\newtheorem{open}      {Problem}}   {}
\@ifundefined{example}   {\newtheorem{example}   {Example}}   {}
\makeatother

% COUNTER DEFINITIONS
\newcounter{tabcounter}
\newcounter{defcounter}
\newcounter{opcounter}
\newcounter{thmcounter}
\newcounter{excounter}
\newcounter{ascounter}
\newcounter{cjcounter}
\newcounter{lemcounter}
\newcounter{remcounter}
\newcounter{corcounter}
\newcounter{propcounter}
\newcounter{figcounter}

% UNIFIED GLOSSARY STYLE
\newglossarystyle{unifiedstyle}{%
  \setglossarystyle{list}%
  \renewcommand*{\glossarypreamble}{\gdef\glslastsec{}}%
  \renewcommand{\glossentry}[2]{%
    \edef\glscursec{\glsentryuseriii{##1}}%
    \ifx\glscursec\glslastsec\else
      \item[]\textbf{\S\,\glscursec\enspace\glsentryuseriv{##1}}%
      \global\let\glslastsec\glscursec
    \fi
    \item[\glsentryuserii{##1} \glsentryuseri{##1}:] ``\glossentryname{##1}'' \glsentrydesc{##1}%
  }%
}

\makeatletter
% First-word extractor (used by \beT, \beD and \beXa for the default key)
% \@firstwordof hello world \relax  ->  hello
\def\@firstwordof#1 #2\relax{#1}
% First three letters, used for the heading mnemonic of a theorem given no
% explicit key.  Padding with \@empty keeps it safe for first words shorter
% than three letters: \@empty vanishes under \edef, whereas \relax would
% survive into the label and is illegal inside \csname (it broke every
% two-letter title, e.g. "ME ...", with "Missing \endcsname inserted").
\def\@takethree#1#2#3#4\@nil{#1#2#3}

% EXAMPLE COMMAND -- same signature and behaviour as \beT and \beD.
% Usage: \beXa*[long title][short title][key]
%   star  = suppresses the printed "[key]" tag after the title (as for \beT*
%   and \beD*); the label and glossary entry are unaffected.
%   key = optional label key override; the label becomes e:KEY, and without a
%   key it is e: plus the first three letters of the title's first word.  A
%   mnemonic label lets an example be cited as "Example~\ref{e:KEY}" instead
%   of by section number, which is what the citing prose needs.
\NewDocumentCommand{\beXa}{s O{} O{} o}{%
  \def\@temparg{#2}%
  \ifx\@temparg\@empty
    \begin{example}%
  \else
    \let\@savedSecTitle\@currentsectionname%
    \refstepcounter{excounter}%
    \IfValueTF{#4}{%
      \edef\@exlblkey{#4}%
    }{%
      \edef\@exfirstword{\@firstwordof#2 \relax}%
      \edef\@exlblkey{%
        \expandafter\@takethree\@exfirstword\@empty\@empty\@empty\@nil}%
    }%
    \def\@tempshort{#3}%
    \ifx\@tempshort\@empty
      \expandafter\xdef\csname @xshort@\@exlblkey\endcsname{\@exlblkey}%
    \else
      \expandafter\def\csname @xshort@\@exlblkey\endcsname{#3}%
    \fi
    \IfBooleanTF{#1}{\let\@extag\@empty}{%
      \edef\@extag{\noexpand\ \noexpand\textup{[\noexpand\texttt{e:\@exlblkey}]}}}%
    \begin{example}[#2\ifshowkeys\@extag\fi]%
    \label{e:\@exlblkey}%
    % record the long title so \nref{e:KEY} prints it, as for theorems
    \protected@write\@auxout{}{%
      \string\@thmtitle{e:\@exlblkey}{\unexpanded{#2}}}%
    \expanded{%
      \noexpand\newglossaryentry{e\theexcounter}{%
        name={\unexpanded{#2}%
          \IfBooleanF{#1}{ [\noexpand\texttt{e:\@exlblkey}]}},%
        sort={\theexample},%
        description={p.~\noexpand\pageref{e:\@exlblkey}},%
        user1={\theexample},%
        user2={Example},%
        user3={\thesection},%
        user4={\expandafter\unexpanded\expandafter{\@savedSecTitle}}%
      }%
    }%
    \protected@write\@auxout{}{\string\glsliveentry{e\theexcounter}}%
  \fi
}
\newcommand{\eeXa}{\end{example}}

% DEFINITION COMMAND -- same signature and behaviour as \beT below.
% Usage: \beD[long title][short title][key]
%   long title  = definition heading (plain text only)
%   short title = shown in \dref{word}; defaults to the key
%   key         = optional label key override; the label becomes d:KEY
% Without a key the label is d: plus the first three letters of the title's
% first word, so it is a mnemonic rather than the definition's number: a
% numeric label would silently change whenever a definition is inserted.
\NewDocumentCommand{\beD}{s O{} O{} o}{%
  \def\@temparg{#2}%
  \ifx\@temparg\@empty
    \begin{definition}%
  \else
    \let\@savedSecTitle\@currentsectionname%
    \refstepcounter{defcounter}%
    \IfValueTF{#4}{%
      \edef\@deflblkey{#4}%
    }{%
      \edef\@deffirstword{\@firstwordof#2 \relax}%
      \edef\@deflblkey{%
        \expandafter\@takethree\@deffirstword\@empty\@empty\@empty\@nil}%
    }%
    \def\@tempshort{#3}%
    \ifx\@tempshort\@empty
      \expandafter\xdef\csname @dshort@\@deflblkey\endcsname{\@deflblkey}%
    \else
      \expandafter\def\csname @dshort@\@deflblkey\endcsname{#3}%
    \fi
    \IfBooleanTF{#1}{\let\@deftag\@empty}{%
      \edef\@deftag{\noexpand\ \noexpand\textup{[\noexpand\texttt{d:\@deflblkey}]}}}%
    \begin{definition}[#2\ifshowkeys\@deftag\fi]%
    \label{d:\@deflblkey}%
    % record the long title so \nref{d:KEY} prints it, as for theorems
    \protected@write\@auxout{}{%
      \string\@thmtitle{d:\@deflblkey}{\unexpanded{#2}}}%
    \expanded{%
      \noexpand\newglossaryentry{d\thedefcounter}{%
        name={\unexpanded{#2}%
          \IfBooleanF{#1}{ [\noexpand\texttt{d:\@deflblkey}]}},%
        sort={\thedefinition},%
        description={p.~\noexpand\pageref{d:\@deflblkey}},%
        user1={\thedefinition},%
        user2={Definition},%
        user3={\thesection},%
        user4={\expandafter\unexpanded\expandafter{\@savedSecTitle}}%
      }%
    }%
    \protected@write\@auxout{}{\string\glsliveentry{d\thedefcounter}}%
  \fi
}
\newcommand{\eeD}{\end{definition}}

% \showkeysfalse in a driver hides the [t:KEY] mnemonics in result headings
% (labels and glossary entries are unchanged); used for submission builds.
\newif\ifshowkeys \showkeystrue
% \commonredtrue in a driver that inputs basicFiles/commonRed: tool.tex then
% cites its reduction results instead of restating them.
\newif\ifcommonred \commonredfalse
% TheorEM COMMAND
% Usage: \beT[long title][short title][key]
%   long title  = theorem heading (plain text only)
%   short title = shown in \tref{word}; may contain math; defaults to first word of long title
%   key         = optional label key override; label becomes t:KEY (avoids first-word collisions)
% Label is t:WORD where WORD = first word of long title (or KEY if given); use \tref{WORD} to reference.
\NewDocumentCommand{\beT}{s O{} O{} o}{%
  \def\@temparg{#2}%
  \ifx\@temparg\@empty
    \begin{theorem}%
  \else
    \let\@savedSecTitle\@currentsectionname%
    \refstepcounter{thmcounter}%
    % \@thmtag is the mnemonic shown in the heading; every theorem gets one.
    % With an explicit key the key itself is shown.  Without one the key
    % defaults to the title's first word, and showing that whole word would
    % merely repeat the title ("Dissipativity [Dissipativity]"), so only its
    % first three letters are shown: "Dissipativity [Dis]".  The label is
    % unaffected -- it still uses the full \@thmlblkey.
    \IfValueTF{#4}{%
      \edef\@thmlblkey{#4}%
    }{%
      \edef\@thmfirstword{\@firstwordof#2 \relax}%
      \edef\@thmlblkey{%
        \expandafter\@takethree\@thmfirstword\@empty\@empty\@empty\@nil}%
    }%
    \IfBooleanTF{#1}{\let\@thmtag\@empty}{%
      \edef\@thmtag{\noexpand\ \noexpand\textup{[\noexpand\texttt{t:\@thmlblkey}]}}}%
    \def\@tempshort{#3}%
    \ifx\@tempshort\@empty
      \expandafter\xdef\csname @tshort@\@thmlblkey\endcsname{\@thmlblkey}%
    \else
      \expandafter\def\csname @tshort@\@thmlblkey\endcsname{#3}%
    \fi
    \begin{theorem}[#2\ifshowkeys\@thmtag\fi]%
    \label{t:\@thmlblkey}%
    % record the long title for \nref; via the .aux so forward refs resolve
    \protected@write\@auxout{}{%
      \string\@thmtitle{t:\@thmlblkey}{\unexpanded{#2}}}%
    \expanded{%
      \noexpand\newglossaryentry{t\thethmcounter}{%
        name={\unexpanded{#2}%
          \IfBooleanF{#1}{ [\noexpand\texttt{t:\@thmlblkey}]}},%
        % sort/user1 must be the number the document actually prints, i.e. the
        % theorem environment's counter, not thmcounter.  thmcounter advances
        % only on \beT with a title, so a bare \beT (or \beT with an
        % empty title) makes the two drift and the glossary cites a wrong
        % number.  thmcounter is still fine as the entry key: it is unique.
        sort={\thetheorem},%
        description={p.~\noexpand\pageref{t:\@thmlblkey}},%
        user1={\thetheorem},%
        user2={Theorem},%
        user3={\thesection},%
        user4={\expandafter\unexpanded\expandafter{\@savedSecTitle}}%
      }%
    }%
    \protected@write\@auxout{}{\string\glsliveentry{t\thethmcounter}}%
  \fi
}
\newcommand{\eeT}{\end{theorem}}

% CROSS-PART RELATION PARAGRAPH
% Usage: \xpart{text}  -- place immediately BEFORE \beT, never inside it
%   (anything inside \beT...\eeT would be swallowed into the glossary entry).
% The paragraph is typeset only when Part I is actually present, i.e. in the
% combined book build; in the standalone builds (GranadaDSABNS/A.tex,
% AAHRah-Submission/AAHRah.tex) the label part:glv is undefined and the whole
% paragraph disappears rather than leaving ?? cross-references. Same
% graceful-degradation principle as \GASref in GASt.tex. Resolves on the
% second LaTeX pass, since r@part:glv comes from the .aux file.
% (Already inside the \makeatletter region opened above, so no \makeatletter
% pair here -- adding one would close that region early.)
\newcommand{\xpart}[1]{\@ifundefined{r@part:glv}{}{\paragraph{Relation with Parts~I--II.}#1}}

% \bookref{label}{precise text}{fallback text}: prints the precise cross-reference
% when the label resolves (combined book build), and the fallback text otherwise
% (standalone paper builds, where book-only appendix labels are absent). Same
% graceful-degradation role as \GASref in GASt.tex, but usable in any file since it
% lives in this preamble (loaded by every driver) -- unlike \GASref, which is defined in a
% content file and is therefore undefined in standalone builds that omit that file.
\newcommand{\bookref}[3]{\@ifundefined{r@#1}{#3}{#2}}

% \PVpart and \RBpart: the two parts that a partial build may omit.
% They print "Part~N" in any build that contains the part, and a naming phrase
% in a build that does not, so an excerpt shows no ?? .  Each already supplies
% the word "Part", so write "\PVpart develops ...", not "Part~\PVpart".
\newcommand{\PVpart}{\bookref{part:aahrah}%
  {Part~\ref{part:aahrah}}{the Perron--Volterra part of the full book}}
\newcommand{\RBpart}{\bookref{part:recbrowne}%
  {Part~\ref{part:recbrowne}}{the immuno-virology part of the full book}}

% OPEN PROBLEM COMMANDS
% ASSUMPTION COMMANDS (glossaried)
\newcommand{\beAs}[1][]{%
  \def\@temparg{#1}%
  \ifx\@temparg\@empty
    \begin{assumption}%
  \else
    \let\@savedSecTitle\@currentsectionname%
    \refstepcounter{ascounter}%
    \begin{assumption}[#1]%
    \label{as:\theascounter}%
    \expanded{%
      \noexpand\newglossaryentry{as\theascounter}{%
        name={\unexpanded{#1}},%
        sort={\theascounter},%
        description={p.~\noexpand\pageref{as:\theascounter}},%
        user1={\theascounter},%
        user2={Assumption},%
        user3={\thesection},%
        user4={\expandafter\unexpanded\expandafter{\@savedSecTitle}}%
      }%
    }%
    \protected@write\@auxout{}{\string\glsliveentry{as\theascounter}}%
  \fi
}
\newcommand{\eeAs}{\end{assumption}}
% CONJECTURE COMMANDS (glossaried)
\newcommand{\beCj}[1][]{%
  \def\@temparg{#1}%
  \ifx\@temparg\@empty
    \begin{conjecture}%
  \else
    \let\@savedSecTitle\@currentsectionname%
    \refstepcounter{cjcounter}%
    \begin{conjecture}[#1]%
    \label{cj:\thecjcounter}%
    \expanded{%
      \noexpand\newglossaryentry{cj\thecjcounter}{%
        name={\unexpanded{#1}},%
        sort={\thecjcounter},%
        description={p.~\noexpand\pageref{cj:\thecjcounter}},%
        user1={\thecjcounter},%
        user2={Conjecture},%
        user3={\thesection},%
        user4={\expandafter\unexpanded\expandafter{\@savedSecTitle}}%
      }%
    }%
    \protected@write\@auxout{}{\string\glsliveentry{cj\thecjcounter}}%
  \fi
}
\newcommand{\eeCj}{\end{conjecture}}
% LEMMA COMMANDS
% ---------------------------------------------------------------------
% UNIFORM RESULT WRAPPERS -- \beL, \beC, \beP, \beR share \beT's signature
%
%   \beX[long title][short title][key]  ...  \eeX
%
%   long title  = heading (math allowed)
%   short title = shown in the cross-reference; defaults to the key
%   key         = label key.  With a key the label is X:KEY and the heading
%                 carries the mnemonic [key], exactly as for \beT.  Without
%                 one the label falls back to the running counter, which is
%                 what every pre-existing call site in the book relies on.
%
% A bare \beX -- arguments commented out -- opens the plain environment and
% creates NO label and NO glossary entry.  That is the supported way to keep
% an individual result out of the catalogue.
% ---------------------------------------------------------------------
% #1 wrapper name  #2 environment  #3 label prefix  #4 counter  #5 type name
\newcommand{\@mkresult}[5]{%
  % A leading * suppresses the [key] mnemonic in both the heading and the
  % glossary name, while still making the label X:KEY and the entry:
  %     \beX*[long][short][key]
  \expandafter\NewDocumentCommand\csname be#1\endcsname{s O{} O{} o}{%
    \def\@temparg{##2}%
    \ifx\@temparg\@empty
      \begin{#2}%
    \else
      \let\@savedSecTitle\@currentsectionname%
      \refstepcounter{#4}%
      % With a key: mnemonic label X:KEY and the [key] tag in the heading.
      % Without one: fall back to the running counter, exactly as these
      % wrappers behaved before, so every legacy one-argument call site in
      % the book keeps its label.  The counter is also csname-safe, which a
      % title's first word is not -- titles beginning with math would
      % otherwise build \csname \Omega...\endcsname and fail.
      \IfValueTF{##4}{%
        \edef\@relbl{##4}%
        \IfBooleanTF{##1}{\let\@retag\@empty}{%
          \edef\@retag{\noexpand\ \noexpand\textup{[\noexpand\texttt{#3:##4}]}}}%
      }{%
        \edef\@relbl{\arabic{#4}}%
        \let\@retag\@empty
      }%
      \def\@tempshort{##3}%
      \ifx\@tempshort\@empty
        \expandafter\xdef\csname @sshort@#3\@relbl\endcsname{\@relbl}%
      \else
        \expandafter\def\csname @sshort@#3\@relbl\endcsname{##3}%
      \fi
      \begin{#2}[##2\ifshowkeys\@retag\fi]%
      \edef\@tmplbl{#3:\@relbl}\expandafter\label\expandafter{\@tmplbl}%
      \expanded{%
        \noexpand\newglossaryentry{#3\csname the#4\endcsname}{%
          name={\unexpanded{##2}%
            \IfBooleanF{##1}{\IfValueT{##4}{ [\noexpand\texttt{#3:##4}]}}},%
          sort={\csname the#2\endcsname},%
          description={p.~\noexpand\pageref{#3:\@relbl}},%
          user1={\csname the#2\endcsname},%
          user2={#5},%
          user3={\thesection},%
          user4={\expandafter\unexpanded\expandafter{\@savedSecTitle}}%
        }%
      }%
      \protected@write\@auxout{}{\string\glsliveentry{#3\csname the#4\endcsname}}%
      \protected@write\@auxout{}{\string\@thmtitle{\@tmplbl}{\unexpanded{##2}}}%
    \fi
  }%
  \expandafter\newcommand\csname ee#1\endcsname{\end{#2}}%
}
\@mkresult{L}{lemma}{l}{lemcounter}{Lemma}
% REMARK / COROLLARY / PROPOSITION -- same signature as \beL above
\@mkresult{R}{remark}{r}{remcounter}{Remark}
\@mkresult{C}{corollary}{c}{corcounter}{Corollary}
\@mkresult{P}{proposition}{p}{propcounter}{Proposition}
% OPEN-PROBLEM COMMAND -- generated with the same signature as \beT:
%   \beO[long title][short title][key]   ...   \eeO
% label o:KEY, or o:<counter> when no key is given; a bare \beO makes
% neither a label nor a glossary entry.
\@mkresult{O}{open}{o}{opcounter}{Problem}


% Section-level title tracker: updated only by \section, not \subsection.
% \@currentlabelname is overwritten by every \refstepcounter (including
% \subsection), so we save it separately here.
\newcommand{\@currentsectionname}{}
\AtBeginDocument{%
  \let\@epidcrn@orig@refstepcounter\refstepcounter
  \renewcommand{\refstepcounter}[1]{%
    \@epidcrn@orig@refstepcounter{#1}%
    \ifstrequal{#1}{section}{\global\let\@currentsectionname\@currentlabelname}{}%
  }%
}

% TABLE COMMANDS (glossaried) -- consistent naming with \beT, \beD, etc.
\newcommand{\beTa}[1][]{%
  \def\@temparg{#1}%
  \ifx\@temparg\@empty
    \begin{table}[htbp]%
  \else
    \let\@savedSecTitle\@currentsectionname%
    \refstepcounter{tabcounter}%
    \begin{table}[htbp]%
    \label{A:\thetabcounter}%
    \expanded{%
      \noexpand\newglossaryentry{A\thetabcounter}{%
        name={\unexpanded{#1}},%
        sort={\thetabcounter},%
        description={p.~\noexpand\pageref{A:\thetabcounter}},%
        user1={\thetabcounter},%
        user2={Table},%
        user3={\thesection},%
        user4={\expandafter\unexpanded\expandafter{\@savedSecTitle}}%
      }%
    }%
    \protected@write\@auxout{}{\string\glsliveentry{A\thetabcounter}}%
  \fi
}
\newcommand{\eeTa}{\end{table}}

% FIGURE COMMANDS (glossaried) -- consistent with \beT, \beR, \beTa, etc.
% Usage:
%   \beF[Glossary title]...\caption{Visible caption}\eeF  -> glossaried figure
%   \beF...\caption{Visible caption}\eeF                  -> plain [H] figure, no glossary
\newcommand{\beF}[1][]{%
  \def\@temparg{#1}%
  \ifx\@temparg\@empty
    \begin{figure}[H]%
  \else
    \begin{figure}[H]%
    \let\@savedSecTitle\@currentsectionname%
    \refstepcounter{figcounter}%
    \label{fig:\thefigcounter}%
    \expanded{%
      \noexpand\newglossaryentry{Fig\thefigcounter}{%
        name={\unexpanded{#1}},%
        sort={\thefigcounter},%
        description={p.~\noexpand\pageref{fig:\thefigcounter}},%
        user1={\thefigcounter},%
        user2={Figure},%
        user3={\thesection},%
        user4={\expandafter\unexpanded\expandafter{\@savedSecTitle}}%
      }%
    }%
    \protected@write\@auxout{}{\string\glsliveentry{Fig\thefigcounter}}%
  \fi
}
\newcommand{\eeF}{\end{figure}}
\makeatother
% LIVE-ENTRY REGISTER
% \glsaddall prints every entry the glossaries package currently knows, and
% that set includes entries restored from \jobname.glsdefs.  The package
% rewrites .glsdefs from all *currently defined* entries, so anything it
% restored is written straight back out: a deleted section kept reappearing in
% the catalogue for ever, and the .glsdefs had to be erased by hand.  The
% wrappers therefore record every entry they actually create in the .aux, and
% \printboth adds only those.  The .aux is rebuilt from scratch on each run, so
% this is self-healing: a removed result is gone after one further run.
\makeatletter
\newcommand{\@glslive}{}
\newcommand{\@glsaddlive}[1]{\ifglsentryexists{#1}{\glsadd{#1}}{}}
\newcommand{\glsliveentry}[1]{\g@addto@macro\@glslive{\@glsaddlive{#1}}}
\makeatother
% PRINT COMMAND
% No \clearpage here -- the catalogue continues on the current page.  This has
% been re-added by mistake several times; leave it out.
\makeatletter
\newcommand{\printboth}{%
  \@glslive   % add exactly the entries this document created last run
  \printglossary[style=unifiedstyle, title={0 \hspace{.22em} Definitions, Theorems, Examples and Open Problems}]%
}
\makeatother
% Named theorem reference: \tref{word} -> Thm~N(word-or-short)
% Shows stored short title if available, otherwise shows the key itself.
\newcommand{\tref}[1]{%
  Thm~\ref{t:#1}(\ifcsname @tshort@#1\endcsname\csname @tshort@#1\endcsname\else#1\fi)%
}

% Full-title theorem reference: \nref{t:KEY} prints the number and then the
% theorem's long title, so one writes "Theorem~\nref{t:blocks}" and gets
% "Theorem 3 ``Block factorization and siphons''".  Takes the complete label, like
% \ref, not the bare key like \tref.
%
% \beT writes each title to the .aux (see \@thmtitle below), so \nref resolves
% forward references from the second LaTeX pass, exactly as \ref does.  If the
% title is not yet known -- first pass, or a label that is not a \beT theorem
% -- it degrades to the number rather than producing ??.
% These two use @ in their names, and this point in the file is outside the
% \makeatletter region opened earlier (it closes at \makeatother above), so
% the region has to be reopened here.
\makeatletter
\newcommand{\@thmtitle}[2]{\global\@namedef{@tlong@#1}{#2}}
\newcommand{\nref}[1]{%
  \ref{#1}%
  \ifcsname @tlong@#1\endcsname
    \ ``\csname @tlong@#1\endcsname''%
  \fi
}
% \nm{t:KEY}: the stored long title of any \beT/\beL/\beP/\beR/\beC/\beD/\beE result,
% verbatim and without quotes, so item names follow the macro's first argument.
\newcommand{\nm}[1]{%
  \ifcsname @tlong@#1\endcsname
    \csname @tlong@#1\endcsname%
  \else
    \ref{#1}%
  \fi
}
\makeatother




%% ========================================
%% Other THEOREM ENVIRONMENTS
%% ========================================
\makeatletter
\@ifundefined{lemma}      {\newtheorem{lemma}      {Lemma}}      {}
\@ifundefined{proposition}{\newtheorem{proposition}{Proposition}} {}
\@ifundefined{corollary}  {\newtheorem{corollary}  {Corollary}}  {}
\@ifundefined{conjecture} {\newtheorem{conjecture} {Conjecture}}  {}
\@ifundefined{remark}     {\newtheorem{remark}     {Remark}}     {}
\@ifundefined{question}   {\newtheorem{question}   {Question}}   {}
\@ifundefined{assumption} {\newtheorem{assumption} {Assumption}} {}
\@ifundefined{hypothesis} {\newtheorem{hypothesis} {Hypothesis}} {}
\@ifundefined{fact}       {\newtheorem{fact}       {Fact}}       {}
\makeatother

% \citer: cite the relay-graph threshold theorem of AAHrel. If a document has
% \input the chapter/file containing t:relayRF (so the label resolves), point
% straight at it; otherwise fall back to a bare citation of AAHrel. Shared here
% (not per-paper) so it resolves consistently across every book chapter.
% \citeR: same idea, for the first theorem of AAHRah (t:Pstar).
\makeatletter
\providecommand{\citer}{%
  \@ifundefined{r@t:relayRF}%
    {\cite{AAHrel}}%
    {Theorem~\ref{t:relayRF} of \cite{AAHrel}}%
}
\providecommand{\citeR}{%
  \@ifundefined{r@t:Pstar}%
    {\cite{AAHRah}}%
    {Theorem~\ref{t:Pstar} of \cite{AAHRah}}%
}
\makeatother


% Abbreviations
\def\beQ{\begin{question}}\def\eeQ{\end{question}}
\def\beA{\begin{assumption}}\def\eeA{\end{assumption}}


%% ========================================

%% ========================================
%% COLORS & CUSTOM COMMANDS
%% ========================================
\definecolor{funccolor}{RGB}{25,25,112}
\definecolor{desccolor}{RGB}{64,64,64}
\newcommand{\func}[1]{\texttt{\textcolor{funccolor}{#1}}}
\def\eq{=}
\def\mm{-}

%% ========================================
%% MATHCAL SHORTHANDS
%% ========================================
% Project-wide convention: \mX for a capital X is always \mathcal X.
% Defined for the whole alphabet so the convention never fails to expand.
\newcommand{\mA}{\mathcal A}\newcommand{\mB}{\mathcal B}
\newcommand{\mC}{\mathcal C}\newcommand{\mD}{\mathcal D}
\newcommand{\mE}{\mathcal E}\newcommand{\mF}{\mathcal F}
\newcommand{\mG}{\mathcal G}\newcommand{\mH}{\mathcal H}
\newcommand{\mI}{\mathcal I}\newcommand{\mJ}{\mathcal J}
\newcommand{\mK}{\mathcal K}\newcommand{\mL}{\mathcal L}
\newcommand{\mM}{\mathcal M}\newcommand{\mN}{\mathcal N}
\newcommand{\mO}{\mathcal O}\newcommand{\mP}{\mathcal P}
\newcommand{\mQ}{\mathcal Q}\newcommand{\mR}{\mathcal R}
\newcommand{\mS}{\mathcal S}\newcommand{\mT}{\mathcal T}
\newcommand{\mU}{\mathcal U}\newcommand{\mV}{\mathcal V}
\newcommand{\mW}{\mathcal W}\newcommand{\mX}{\mathcal X}
\newcommand{\mY}{\mathcal Y}\newcommand{\mZ}{\mathcal Z}

%% ========================================
%% BASIC MATH ENVIRONMENTS
%% ========================================
\def\bep{\begin{pmatrix}}\def\eep{\end{pmatrix}}
\def\bev{\begin{vmatrix}}\def\eev{\end{vmatrix}}
\def\bea{\begin{eqnarray*}}\def\eea{\end{eqnarray*}}
\def\bc{\begin{cases}}\def\ec{\end{cases}}
\def\BEN{\begin{enumerate}}\def\EEN{\end{enumerate}}
\def\BI{\begin{itemize}}\def\EI{\end{itemize}}

\newcommand{\be}[1]{\begin{equation}\label{#1}}
\newcommand{\ee}{\end{equation}}
\newcommand{\beq}{\begin{eqnarray}}
\def\eeq{\end{eqnarray}}
\def\eqr{\eqref}\def\fr{\frac}\def\lbl{\label}
\def\Lra{\Longrightarrow}\def\non{\nonumber}\def\cNGM{\cite{Diek,Van,Van08}}
\def\Eq{\Leftrightarrow}\def\b{\beta}\def\uph{\red{UP TO HERE}}
%% ========================================
%% PROBABILITY & STATISTICS
%% ========================================
\renewcommand{\P}{\ensuremath{\mathbb{P}}}  % Note: overrides paragraph symbol
\newcommand{\E}{\ensuremath{\mathbb{E}}}
\newcommand{\ind}[1]{\ensuremath{\mathbbm{1}_{\left\{#1\right\}}}}
\DeclareMathOperator{\Cov}{Cov}
\DeclareMathOperator{\TV}{TV}
\newcommand{\rank}{\operatorname{rank}}

%% ========================================
%% CRN OPERATORS (organizations/ERCs/buffering/siphons framework)
%% ========================================
\def\clos{\mathrm{clos}}
\def\Supply{\mathrm{Supply}}
\def\req{\mathrm{req}}
\def\React{\mathrm{React}}
\def\Prod{\mathrm{Prod}}
\def\Gen{\mathrm{Gen}}
\def\In{\mathrm{In}}
\def\Out{\mathrm{Out}}
\def\Act{\mathrm{Act}}
\def\Min{\mathrm{Min}}

%% ========================================
%% LINE BREAKING
%% ========================================
% Last-resort stretch, used only on paragraphs TeX cannot otherwise break:
% a long \texttt{} token, a comma-separated inline math list (commas are not
% breakpoints in math), or a run of \cite/\ref boxes.  Without it those set as
% overfull \hbox and print into the margin; with it they set slightly loose,
% which is the lesser evil.  Does not affect paragraphs that already break.
\emergencystretch=3em

%% ========================================
%% NUMBER SETS
%% ========================================
\newcommand{\R}{\mathbb{R}}  % Real numbers
\newcommand{\C}{\mathbb{C}}  % Complex numbers
\newcommand{\N}{\mathbb{N}}
\newcommand{\Z}{\mathbb{Z}}
% \providecommand, not \newcommand: some papers define these in their own
% preamble after \input{preamble}, and must not clash.
\providecommand{\Q}{\mathbb{Q}}
\providecommand{\res}{\operatorname{res}}   % resultant; used by epiRes.tex
\newcommand{\MN}{\mathbb{N}}
\newcommand{\MZ}{\mathbb{Z}}
\newcommand{\realpos}{\mathbb{R}_{> 0}}
\newcommand{\realnneg}{\mathbb{R}_{\geq 0}}
\newcommand{\naturalpos}{\mathbb{N}_{>0}}
\newcommand{\Rplus}{\mathbb{R}_+}
\newcommand{\Rp}{\mathbb{R}_{\ge 0}}
\def\Rnne{\mathbb{R}_{\ge 0}}
\def\Se{\bar{S}} \def\Ie{\bar{I}}
\newcommand{\Rpos}{\mathbb{R}_{>0}}
\newcommand{\dd}{\,\mathrm d}
\newcommand{\e}{\mathrm e}
%% ========================================
%% GREEK LETTERS (2+ character abbreviations) greek
%% ========================================
\def\al{\alpha}\def\ba{\bff \alpha}
\def\w{\omega}\def\Lgn{\ga+\mu+\de}\def\om{\omega}
\def\de{\delta}\def\De{\Delta}\def\bde{\bff \delta}
\def\ep{\epsilon}\def\eps{\varepsilon}
\def\f{\varphi}  % Kept as requested
\def\ga{\gamma}\def\Ga{\Gamma}\def\G{\Gamma}\def\bg{\bff \gamma}
\def\up{\upsilon}  % NOT the breve accent: no accent commands anywhere.
                   % Named \up, not \u, so the LaTeX breve stays reachable.
\def\la{\lambda}\def\La{\Lambda}
\def\si{\sigma}\def\ka{\kappa}
\def\th{\theta}\def\vt{\vartheta}
\def\ze{\zeta}\def\Gap{\Gamma^{+}} \def\Gam{\Gamma^{-}}


%% ========================================
%% SPECIAL SHORT COMMANDS (kept as requested)
%% ========================================
\def\I{\infty}
\def\T{\widetilde}
\def\H{\hat}
\newcommand{\relint}{\operatorname{relint}}

%% ========================================
%% BOLD/VECTOR NOTATION
%% ========================================
\newcommand{\bff}[1]{{\mbox{\boldmath$#1$}}}
\newcommand\X{\boldsymbol{X}}
\def\x{\boldsymbol{x}}\def\bx{\bff x}\def\xe{\bar{\x}}
\newcommand\y{\boldsymbol{y}}\def\by{\boldsymbol{y}}\def\ye{\bar{\y}}
\def\yb{\y}\def\xb{\x}
\def\xy{\x-\y}
\newcommand\br{\boldsymbol{r}}\def\re{\bar{\r}}
\newcommand\bd{\boldsymbol{d}}
\newcommand\bQ{\boldsymbol{Q}}
\newcommand{\Yb}{{\mathbf{Y}}}
\newcommand{\Ab}{\mathbf{A}}
\newcommand{\zb}{\mathbf{z}}\def\bz{\bff 0}
\newcommand{\Zb}{\mathbf{Z}}
\newcommand{\bb}{\mathbf{b}}
\def\vbe{\vec {\bff \beta}}
\def\vb{\vec {\bff b}}
\def\v1{\vec {\bff 1}}
\def\vn{\vec \de}
\def\vz{\vec 0}
\def\vi{\overset{\rightarrow}{\mathsf i}}
\def\vv{\overset{\rightarrow}{v}}
\def\vr{\; \vec {\mathsf r}}
\newcommand{\va}{\bff a}
\providecommand{\vect}[1]{\boldsymbol{#1}}   % bold vector (moved from Gradostat.tex)
\providecommand{\one}{\mathbf 1}             % all-ones vector (moved from AliSF)
\providecommand{\spr}{s}                     % spectral bound  (moved from AliSF)
\providecommand{\Ke}{K_{\rm eff}}            % effective carrying capacity
% Critical right/left eigenvectors in the Hopf projection method (AAHOSN/New.tex)
\providecommand{\bq}{\mathbf{q}}
\providecommand{\bp}{\mathbf{p}}


%% ========================================
%% MATHCAL & MATHSCR
%% ========================================
% Mathcal
\newcommand\mc[1]{\mathcal{#1}}
\def\mC{{\mathcal C}}\def\mD{\mathcal D}\def\mE{\mathcal{E}}
\def\mF{\mathcal F}\def\mG{{\mathcal G}}\def\mH{{\mathcal H}}
\def\mI{\mathcal I}\def\mL{\mathcal{L}}\def\mM{{\mathcal M}}
\def\mN{{\mathcal N}}\def\mP{\mathcal{P}}\def\mR{{\mathcal R}}
\def\mS{{\mathcal S}}\def\mT{\mathcal{T}}\def\mV{{\mathcal V}}
\def\mK{\mathcal K}\def\mA{{\mathcal A}}\def\mU{{\mathcal U}}
% Symmetric bilinear and trilinear forms of the Taylor expansion, used in the
% first Lyapunov coefficient (AAHOSN/New.tex)
\providecommand{\Bform}{\mathcal{B}}
\providecommand{\Cform}{\mathcal{C}}

% Blackboard bold.  Kept here, not in the drivers: a \providecommand in a
% driver preamble runs before this file and silently wins, so the same name
% ends up meaning different things in different builds.
\def\IC{\mathbb{C}}\def\IN{\mathbb{N}}\def\IQ{\mathbb{Q}}
\def\IR{\mathbb{R}}\def\IZ{\mathbb{Z}}\def\IG{\mathbb{G}}
\def\II{\mathbb{I}}

% Mathscr
\newcommand{\Rs}{{\mathscr R}}
\newcommand{\Ss}{{\mathscr S}}
\newcommand{\Ns}{{\mathscr N}}
\newcommand{\Cs}{\mathscr C}


%% ========================================
%% OPERATORS & FUNCTIONS
%% ========================================
\DeclareMathOperator{\tr}{tr}
\DeclareMathOperator{\diag}{Diag}
\DeclareMathOperator{\Diag}{Diag}
\DeclareMathOperator{\intr}{\mathrm{Int}}
\DeclareMathOperator{\Inv}{Inv}
\DeclareMathOperator{\spec}{spec}
\DeclareMathOperator{\AUC}{AUC}               % area under the curve
\newcommand{\ii}{\mathrm{i}}                  % imaginary unit
\newcommand\op{\operatorname}
\newcommand{\supp}{\op{supp}}
\newcommand{\diff}{\mathop{}\mathopen{}\mathrm{d}}
\newcommand{\pdd}[1]{\frac{\partial}{\partial{#1}}}
\DeclareMathOperator{\cone}{cone}
\DeclareMathOperator{\proj}{proj}

%% ========================================
%% BRACKETS & NORMS
%% ========================================
\providecommand{\pr}[1]{\left(#1\right)}
\providecommand{\pp}[1]{\left[#1\right]}
\newcommand\croc[1]{\left\langle #1\right\rangle}
\newcommand{\norm}[1]{\|#1\|}
\newcommand{\inner}[2]{\langle #1, #2 \rangle}
% Short synonym of \inner, used by AAHOSN/New.tex as $\ip{u}{v}=\bar u^\top v$
\providecommand{\ip}[2]{\langle #1 , #2 \rangle}


%% ========================================
%% FORMATTING & UTILITIES
%% ========================================
\newcommand*{\Scale}[2][4]{\scalebox{#1}{$#2$}}
\newcommand{\te}[1]{\text{#1}}
\newcommand\steq[1]{\stackrel{\text{\rm #1.}}{=}}
\newcommand\blue[1]{\textcolor{blue}{#1}}
\newcommand{\red}{\textcolor[rgb]{1.00,0.00,0.00}}
\newcommand{\green}{\textcolor[rgb]{0.00,1.00,0.00}}

% Figure command
\newcommand{\figu}[3]{%
\begin{figure}[H]
\centering
\includegraphics[scale=#3]{#1}
\caption{#2\label{f:#1}}
\end{figure}
}

% Footnote
\long\def\symbolfootnote[#1]#2{%
\begingroup
\def\thefootnote{\fnsymbol{footnote}}\footnote[#1]{#2}%
\endgroup}
\def\fn{\symbolfootnote}

\newcommand{\dsum}{\displaystyle\sum}

%% ========================================
%% CRN & EPIDEMIOLOGY SPECIFIC
%% ========================================
% CRN notation
\def\SS{\mathcal S}
\def\CC{\mathcal C}
\def\RR{\mathcal R}

% Reproduction numbers
\def\Rto{R_{1}^{E_2}}
\def\Rot{R_{2}^{E_1}}
\def\Roc{R_{1,c}}
\def\Rtc{R_{2,c}}
\def\rd{r_{0}}
\def\sd{s_{0}}
\def\dd{D_{0}}
\def\vd{v_{0}}
\def\zd{z_{0}}
\def\se{\sE}
\def\xd{x_{0}}
%% ========================================
%% SPECIAL VALUES & EQUILIBRIA
%% ========================================
\def\eE{\bar{e}}
\def\wE{\bar{w}}
\def\sE{s_{*}}
\def\Me{\T M_{ee}}
\def\Rnne{\mathbb{R}_{\ge 0}}

\newcommand{\ones}{\mathbf 1}


%% ========================================
%% ABBREVIATIONS (alphabetical by category)
%% ========================================

% -- General mathematical terms --
\def\adm{admissible}\def\admy{admissibility}
%\def\and{antisymmetric}% disabled 2026-09-30: it overrode LaTeX's author separator \and (printed "antisymmetric" between authors)
\def\ass{assumption}
\def\cadlag{c\`adl\`ag}
\def\ch{characteristic polynomial}
\def\chp{characteristic polynomial}
\def\coe{coefficient}
\def\comp{computation}
\def\con{condition}
\def\coo{coordinate}
\def\corr{corresponding}
\def\dec{decomposition}
\def\den{denominator}
\def\dete{deterministic}
\def\eig{eigenvalue}\def\eval{eigenvalue}
\def\eigv{eigenvector}\def\eiv{eigenvector}\def\evec{eigenvector}
\def\elt{element}
\def\enn{essentially nonnegative/positivity preserving}
\def\equ{equation}
\def\etal{et al.}
\def\fe{for example}
\def\fno{from now on}
\def\for{\forall}
\def\form{formula}
\def\Fr{Furthermore, }
\def\frt{furthermore}
\def\fun{function}
\def\gene{generalization}
\def\heu{heuristic}
\def\how{however}\def\How{However, }
\def\ie{i.e. }
\def\im{\item}
\def\isd{is defined by}
\def\itc{it is easy to check that}
\def\itf{it follows that}\def\Itf{It follows that}
\def\ith{it holds that}
\def\itm{it may be checked that}\def\Itm{It may be checked that}
\def\ito{it turns out that}\def\Ito{It turns out that}\def\Itou{It turns out useful to}
\def\ma{mass action}\def\MAR{mass action representation}
\def\mat{matrix}
\def\mbc{may be checked to be}
\def\mbw{may be written as}\def\mubw{must be written as}
\def\mpr{more precisely, }\def\Mp{More precisely, }
\def\nc{non-catalytic}
\def\nne{non-negative}\def\nny{nonnegativity}
\def\no{\nonumber}
\def\noni{non-infectious}
\def\nota{notation}
\def\np{\newpage}
\def\Nt{Note that}
\def\nt{non trivial}\def\nz{non zero}
\def\oth{otherwise}
\def\Oth{On the other hand}
\def\pa{parameter}\def\para{parameter}
\def\pol{polynomial}
\def\prob{probability}
\def\Prf{{\bf  Proof:}}\def\prf{{\bf  Proof:}}
\def\prop{proportional to}
\def\qu{\quad}
\def\QED{\hfill {$\square$}\goodbreak \medskip}
\def\rep{representation}
\def\resp{respectively}
\def\satd{satisfied}\def\satg{satisfying}\def\sats{satisfies}\def\saty{satisfy}
\def\ssec{\subsection}\def\sssec{\subsubsection}
\def\sef{see for example}
\def\sksy{skew-symmetric}
\def\sub{substrate}
\def\sys{system}
\def\var{variable}
\def\wk{well-known}
\def\wl{without loss of generality}\def\wlo{w.l.o.g.}
\def\wmc{We may check that}\def\Wmc{We may check that}
\def\wmr{we may rewrite}
\def\wms{we must show that}
\def\wrt{with respect to}

% -- CRN & Chemistry --
\newcommand\CRN{chemical reaction network}
\newcommand\CRNs{chemical reaction networks}
\newcommand\RN{reaction network}
\def\ACR{absolute concentration robustness}
\def\AMSD{acyclic siphon decomposition}
\def\DAE{differential algebraic equations}
\def\FHJ{Feinberg-Horn-Jackson}
\def\GMAK{generalized mass action kinetics}
\def\MM{Michaelis-Menten}
\def\sm{stoichiometric matrix}\def\Sm{stoichiometric matrix}
\def\WRZD{weakly reversible, zero deficiency}\def\aut{autocatalytic}

% -- Epidemiology --
\def\brn{basic reproduction number}\def\bzn{basic replacement number}
\def\com{compartment}
\def\dfe{disease free equilibrium}\def\DFE{disease free equilibrium}
\def\dwt{dwell times}\def\edt{expected dwell times}
\def\EM{epidemic model}
\def\endp{endemic fixed point}\def\enp{endemic point}\def\EE{endemic point}
\def\ME{mathematical epidemiology}
\def\MEE{eco-epidemiology}\def\EEP{eco-epidemiological}
\def\ece{eco-epidemiological model}
\def\mgm{multi group model}
\def\mtS{multi-strain model}\def\mSt{multi-strain model}
\def\mSi{minimal siphon}
\def\NGM{next generation matrix}
\def\nim{new infections matrix}
\def\qem{quadratic epidemiological model}
\def\repF{invasion functions}
\def\SPF{SIR-PH-FA model}
\def\sus{susceptible}

% -- Graph theory & Linear algebra --
\def\jac{Jacobi matrix}
\def\jin{Jacobian matrix of the invasion vector field with respect to the invasion variables}
\def\jrn{Jacobian matrix of the resident vector field with respect to the resident variables}
\def\jxn{Jacobian matrix of the $\x$ vector field with respect to the $\x$ variables}
\def\Lm{Laplacian matrix}
\def\MTT{matrix tree theorem}
\def\SGT{spectral graph theory}
\def\SMW{Sherman–Morrison–Woodbury}

% -- Stability & Dynamics --
\def\absC{absorbtion CTMC}
\def\BTb{Bogdanov-Takens bifurcation}
\def\CEP{competitive exclusion partition}
\def\CTMC{continuous time Markov chain}
\def\ds{differential system}
\def\fp{fixed point}\def\bfp{boundary fixed point}\def\fbp{fixed boundary point}
\def\fbl{fixed boundary line}\def\fif{forward invariant face}
\def\gas{global asymptotic stability}\def\GAS{globally asymptotically stable}
\def\GASS{\text{GASS}}
\def\Hb{Hopf bifurcation}
\def\Hur{Hurwitz}\def\MH{Metzler-Hurwitz}\def\RH{Routh-Hurwitz}
\def\IS{\text{IS}}
\def\LAS{locally asymptotically stable}
\def\LdS{Lyapunov diagonally stable}
\def\Lf{Lyapunov function}\def\LyV{Lyapunov Volterra type function}\def\RLF{robust Lyapunov function}
\def\LyM{Lyapunov Malkin stable}
\def\PF{Perron-Frobenius}
\def\spfp{strictly positive fixed point}
\def\sta{stability}
\def\STP{sharp threshold property}\def\STT{sharp threshold theorem}

% -- Models (with \varphi explicitly instead of \f) --
\def\AB{$(A,B)$ Arino-Brauer epidemic models}
\def\ABa{$(A,\bbe,\ba,\varphi)$ balanced bilinear model}
\def\ABP{%$(A,\bbe,P,\varphi)$
\\ balanced bilinear model}
\def\AMG{Anderson May Gupta}
\newcommand\BIN{biological interaction network}
\def\Dt{Descartes type}
\def\gLV{generalized Lotka-Volterra model}
\def\IG{\text{invasion graph}}\def\IGMS{interaction graph on minimal siphons}
\def\LCP{\text{linear complementarity problem}}
\def\LV{Lotka-Volterra}\def\LVK{Lotka-Volterra-Kolmogorov}

% -- Splitting & Structure --
\def\reg{regular splitting}\def\regS{{\bf regular splitting}}
\def\RUR{rational univariate representation}
\def\bds{block-diagonal structure}\def\NGMs{next generation matrices}

% -- Miscellaneous --
\def\Ma{{Mathematica}}
\def\Mm{Max–Min}
\def\num{numerator}
\def\rap{rational approximation} \def\DSR{directed species-reaction bipartite graph}\def\PCP{polynomial complementarity problem}\def\rat{rational substitution}
\def\QSS{quasi steady state} \def\RHS{right hand side} \def\LHS{left hand side}
\def\OSN{online social network}\def\VdD{Van den Driessche}
\def\wli{$\omega$-limit} \def\rk{rich kinetics}\def\sJ{symbolic Jacobian}
\def\cV{\cite{Vas,VasSta,BSV,GSV}} \def\spf{semi-parametric semi-functional}
\def\CRW{Capasso-Ruan-Wang} \def\ABP{$(A,\bbe,P,\f)$ balanced bilinear model}
\def\ABa{%$(A,\bbe,\ba,\f)$ \\
rank-one balanced bilinear model}
\def\bbe{B%{\bff \beta}
} \def\vbe{\vec {\bff \beta}
}
\def\vb{\vec {\bff b}}
\newcommand{\Kh}{\hat{K}}          % K with hat, first mention only
\newcommand{\wS}{w^{\sigma}}       % right Perron J^\perp, full block
\newcommand{\piS}{\vec \pi^{\sigma}}    % left Perron J^\perp, full block
\newcommand{\ws}{w^{\sigma}}       % right Perron J^\perp, sub-block
\newcommand{\pis}{\pi^{\sigma}}    % left Perron J^\perp, sub-block
\newcommand{\hwS}{\hat{w}^{\sigma}}   % right Perron K, full block
\newcommand{\hpiS}{\hat{\pi}^{\sigma}} % left Perron K, full block
\newcommand{\hws}{\hat{w}^{\sigma}}   % right Perron K, sub-block
\newcommand{\hpis}{\hat{\pi}^{\sigma}} % left Perron K, sub-block
\newcommand{\twS}{\tilde{w}^{\sigma}}  % right Perron \tilde K, full block
\newcommand{\tpiS}{\tilde{\pi}^{\sigma}} % left Perron \tilde K
\newcommand{\tws}{\tilde{w}^{\sigma}}
\newcommand{\tpis}{\tilde{\pi}^{\sigma}}
\newcommand{\Mat}[2]{\mathbb{R}_{\ge 0}^{\,#1\times #2}}
\def\tJ{J_\sigma^\perp} \def\lPe{left Perron eigenvector}
% \G is \Gamma (defined at :795).  The Volterra function keeps its own name,
% so that \G no longer has two conflicting meanings in this file.
\def\GV{G}  % Volterra G(u)=u-1-\ln u
\def\PVL{Perron-Volterra Lyapunov minimal siphon ansatz}
\newcommand{\istar}[1]{i_{#1}^{*}}
%\newcommand{\sE}{s_E}
\newcommand{\vE}{v_E}
\def\PV{Perron-Volterra candidate Lyapunov function}
\def\GLV{generalized Lotka-Volterra}
\def\MERN{mathematical epidemiology reaction networks}
\def\CMM{contraction metric method}
\newcommand{\sign}{\operatorname{sign}}
\newcommand{\Var}{\operatorname{Var}} 

% ======================================================================
% IDEAS SHARED BETWEEN THE ABSTRACT AND THE MAIN RESULTS SECTION
%
% Write each idea ONCE, in the abstract; the Main Results section recalls it,
% so the two can never drift apart.
%
%   in the abstract        \ideaitem{1}{The first-order reaction terms ...}
%   in Main Results        \ideahead{1}
%
% \ideaitem stores the text globally and typesets
%   \item \hyperlink{idea:1}{<text>}
% \ideahead typesets
%   \item\hypertarget{idea:1}{}\textbf{<the same text>}
%
% The abstract must come before the Main Results input, which is the normal
% order. If a heading has to read differently from the abstract sentence (for
% instance a long abstract item summarised in one line), override it with the
% optional argument: \ideahead[short heading]{6}.
% ======================================================================
\makeatletter
\newcommand{\ideaitem}[2]{%
  \expandafter\gdef\csname @idea@#1\endcsname{#2}%
  \item\hyperlink{idea:#1}{#2}}
\newcommand{\ideahead}[2][]{%
  \item\hypertarget{idea:#2}{}%
  \def\@ideatmp{#1}%
  \ifx\@ideatmp\@empty
    \@ifundefined{@idea@#2}{%
      \PackageWarning{preamble}{Idea #2 is used in Main Results but was never
        defined by \string\ideaitem\space in the abstract}%
      \textbf{[idea #2 undefined]}%
    }{\textbf{\begingroup\def\ideaterm##1##2{##2}\def\idearest##1{}\@nameuse{@idea@#2}\endgroup}}%
  \else
    \textbf{#1}%
  \fi}
% \ideaterm{<hypertarget>}{<words>}: inside an abstract idea the words become an
% invisible link (the document sets \hypersetup{hidelinks}) jumping to the place
% in Main Results that defines them; inside the Main Results heading itself,
% where the target sits a few lines below, \ideahead strips the link and prints
% the words plainly. Mark the destination with \hypertarget{<name>}{} on the line
% before the display, which adds no equation number and shows nothing.
\newcommand{\ideaterm}[2]{\hyperlink{#1}{#2}}
% \idearest{<text>}: the part of an abstract idea after its first sentence; printed in
% the abstract, omitted from the Main Results heading.
\newcommand{\idearest}[1]{#1}
\makeatother

% ======================================================================
% SHORT-PAPER SWITCH
% \shorttrue in a master that assembles only part of the companion paper
% (Barrat/SipArxiv.tex). \nsh{...} drops text that cites results living in the
% omitted sections, so the short build has no dangling references while the
% companion paper keeps the text verbatim.
% ======================================================================
\newif\ifshort \shortfalse
\newcommand{\nsh}[1]{\ifshort\else#1\fi}
% \shortalt{<short build>}{<companion paper>}: use when a sentence needs
% different wording in the two builds, not merely a deletion.
\newcommand{\shortalt}[2]{\ifshort#1\else#2\fi}
% \Jtr, \JtrS{\Sigma}: transverse Jacobian. J^\perp here; SipArxiv.tex redefines them as J_{\mathrm{tra}}, J_{\mathrm{tra},\Sigma}
\newcommand{\Jtr}{J^\perp}
\newcommand{\JtrS}[1]{J_{#1}^\perp}
, and must not clash.
\providecommand{\Q}{\mathbb{Q}}
\providecommand{\res}{\operatorname{res}}   % resultant; used by epiRes.tex
\newcommand{\MN}{\mathbb{N}}
\newcommand{\MZ}{\mathbb{Z}}
\newcommand{\realpos}{\mathbb{R}_{> 0}}
\newcommand{\realnneg}{\mathbb{R}_{\geq 0}}
\newcommand{\naturalpos}{\mathbb{N}_{>0}}
\newcommand{\Rplus}{\mathbb{R}_+}
\newcommand{\Rp}{\mathbb{R}_{\ge 0}}
\def\Rnne{\mathbb{R}_{\ge 0}}
\def\Se{\bar{S}} \def\Ie{\bar{I}}
\newcommand{\Rpos}{\mathbb{R}_{>0}}
\newcommand{\dd}{\,\mathrm d}
\newcommand{\e}{\mathrm e}
%% ========================================
%% GREEK LETTERS (2+ character abbreviations) greek
%% ========================================
\def\al{\alpha}\def\ba{\bff \alpha}
\def\w{\omega}\def\Lgn{\ga+\mu+\de}\def\om{\omega}
\def\de{\delta}\def\De{\Delta}\def\bde{\bff \delta}
\def\ep{\epsilon}\def\eps{\varepsilon}
\def\f{\varphi}  % Kept as requested
\def\ga{\gamma}\def\Ga{\Gamma}\def\G{\Gamma}\def\bg{\bff \gamma}
\def\up{\upsilon}  % NOT the breve accent: no accent commands anywhere.
                   % Named \up, not \u, so the LaTeX breve stays reachable.
\def\la{\lambda}\def\La{\Lambda}
\def\si{\sigma}\def\ka{\kappa}
\def\th{\theta}\def\vt{\vartheta}
\def\ze{\zeta}\def\Gap{\Gamma^{+}} \def\Gam{\Gamma^{-}}


%% ========================================
%% SPECIAL SHORT COMMANDS (kept as requested)
%% ========================================
\def\I{\infty}
\def\T{\widetilde}
\def\H{\hat}
\newcommand{\relint}{\operatorname{relint}}

%% ========================================
%% BOLD/VECTOR NOTATION
%% ========================================
\newcommand{\bff}[1]{{\mbox{\boldmath$#1$}}}
\newcommand\X{\boldsymbol{X}}
\def\x{\boldsymbol{x}}\def\bx{\bff x}\def\xe{\bar{\x}}
\newcommand\y{\boldsymbol{y}}\def\by{\boldsymbol{y}}\def\ye{\bar{\y}}
\def\yb{\y}\def\xb{\x}
\def\xy{\x-\y}
\newcommand\br{\boldsymbol{r}}\def\re{\bar{\r}}
\newcommand\bd{\boldsymbol{d}}
\newcommand\bQ{\boldsymbol{Q}}
\newcommand{\Yb}{{\mathbf{Y}}}
\newcommand{\Ab}{\mathbf{A}}
\newcommand{\zb}{\mathbf{z}}\def\bz{\bff 0}
\newcommand{\Zb}{\mathbf{Z}}
\newcommand{\bb}{\mathbf{b}}
\def\vbe{\vec {\bff \beta}}
\def\vb{\vec {\bff b}}
\def\v1{\vec {\bff 1}}
\def\vn{\vec \de}
\def\vz{\vec 0}
\def\vi{\overset{\rightarrow}{\mathsf i}}
\def\vv{\overset{\rightarrow}{v}}
\def\vr{\; \vec {\mathsf r}}
\newcommand{\va}{\bff a}
\providecommand{\vect}[1]{\boldsymbol{#1}}   % bold vector (moved from Gradostat.tex)
\providecommand{\one}{\mathbf 1}             % all-ones vector (moved from AliSF)
\providecommand{\spr}{s}                     % spectral bound  (moved from AliSF)
\providecommand{\Ke}{K_{\rm eff}}            % effective carrying capacity
% Critical right/left eigenvectors in the Hopf projection method (AAHOSN/New.tex)
\providecommand{\bq}{\mathbf{q}}
\providecommand{\bp}{\mathbf{p}}


%% ========================================
%% MATHCAL & MATHSCR
%% ========================================
% Mathcal
\newcommand\mc[1]{\mathcal{#1}}
\def\mC{{\mathcal C}}\def\mD{\mathcal D}\def\mE{\mathcal{E}}
\def\mF{\mathcal F}\def\mG{{\mathcal G}}\def\mH{{\mathcal H}}
\def\mI{\mathcal I}\def\mL{\mathcal{L}}\def\mM{{\mathcal M}}
\def\mN{{\mathcal N}}\def\mP{\mathcal{P}}\def\mR{{\mathcal R}}
\def\mS{{\mathcal S}}\def\mT{\mathcal{T}}\def\mV{{\mathcal V}}
\def\mK{\mathcal K}\def\mA{{\mathcal A}}\def\mU{{\mathcal U}}
% Symmetric bilinear and trilinear forms of the Taylor expansion, used in the
% first Lyapunov coefficient (AAHOSN/New.tex)
\providecommand{\Bform}{\mathcal{B}}
\providecommand{\Cform}{\mathcal{C}}

% Blackboard bold.  Kept here, not in the drivers: a \providecommand in a
% driver preamble runs before this file and silently wins, so the same name
% ends up meaning different things in different builds.
\def\IC{\mathbb{C}}\def\IN{\mathbb{N}}\def\IQ{\mathbb{Q}}
\def\IR{\mathbb{R}}\def\IZ{\mathbb{Z}}\def\IG{\mathbb{G}}
\def\II{\mathbb{I}}

% Mathscr
\newcommand{\Rs}{{\mathscr R}}
\newcommand{\Ss}{{\mathscr S}}
\newcommand{\Ns}{{\mathscr N}}
\newcommand{\Cs}{\mathscr C}


%% ========================================
%% OPERATORS & FUNCTIONS
%% ========================================
\DeclareMathOperator{\tr}{tr}
\DeclareMathOperator{\diag}{Diag}
\DeclareMathOperator{\Diag}{Diag}
\DeclareMathOperator{\intr}{\mathrm{Int}}
\DeclareMathOperator{\Inv}{Inv}
\DeclareMathOperator{\spec}{spec}
\DeclareMathOperator{\AUC}{AUC}               % area under the curve
\newcommand{\ii}{\mathrm{i}}                  % imaginary unit
\newcommand\op{\operatorname}
\newcommand{\supp}{\op{supp}}
\newcommand{\diff}{\mathop{}\mathopen{}\mathrm{d}}
\newcommand{\pdd}[1]{\frac{\partial}{\partial{#1}}}
\DeclareMathOperator{\cone}{cone}
\DeclareMathOperator{\proj}{proj}

%% ========================================
%% BRACKETS & NORMS
%% ========================================
\providecommand{\pr}[1]{\left(#1\right)}
\providecommand{\pp}[1]{\left[#1\right]}
\newcommand\croc[1]{\left\langle #1\right\rangle}
\newcommand{\norm}[1]{\|#1\|}
\newcommand{\inner}[2]{\langle #1, #2 \rangle}
% Short synonym of \inner, used by AAHOSN/New.tex as $\ip{u}{v}=\bar u^\top v$
\providecommand{\ip}[2]{\langle #1 , #2 \rangle}


%% ========================================
%% FORMATTING & UTILITIES
%% ========================================
\newcommand*{\Scale}[2][4]{\scalebox{#1}{$#2$}}
\newcommand{\te}[1]{\text{#1}}
\newcommand\steq[1]{\stackrel{\text{\rm #1.}}{=}}
\newcommand\blue[1]{\textcolor{blue}{#1}}
\newcommand{\red}{\textcolor[rgb]{1.00,0.00,0.00}}
\newcommand{\green}{\textcolor[rgb]{0.00,1.00,0.00}}

% Figure command
\newcommand{\figu}[3]{%
\begin{figure}[H]
\centering
\includegraphics[scale=#3]{#1}
\caption{#2\label{f:#1}}
\end{figure}
}

% Footnote
\long\def\symbolfootnote[#1]#2{%
\begingroup
\def\thefootnote{\fnsymbol{footnote}}\footnote[#1]{#2}%
\endgroup}
\def\fn{\symbolfootnote}

\newcommand{\dsum}{\displaystyle\sum}

%% ========================================
%% CRN & EPIDEMIOLOGY SPECIFIC
%% ========================================
% CRN notation
\def\SS{\mathcal S}
\def\CC{\mathcal C}
\def\RR{\mathcal R}

% Reproduction numbers
\def\Rto{R_{1}^{E_2}}
\def\Rot{R_{2}^{E_1}}
\def\Roc{R_{1,c}}
\def\Rtc{R_{2,c}}
\def\rd{r_{0}}
\def\sd{s_{0}}
\def\dd{D_{0}}
\def\vd{v_{0}}
\def\zd{z_{0}}
\def\se{\sE}
\def\xd{x_{0}}
%% ========================================
%% SPECIAL VALUES & EQUILIBRIA
%% ========================================
\def\eE{\bar{e}}
\def\wE{\bar{w}}
\def\sE{s_{*}}
\def\Me{\T M_{ee}}
\def\Rnne{\mathbb{R}_{\ge 0}}

\newcommand{\ones}{\mathbf 1}


%% ========================================
%% ABBREVIATIONS (alphabetical by category)
%% ========================================

% -- General mathematical terms --
\def\adm{admissible}\def\admy{admissibility}
%\def\and{antisymmetric}% disabled 2026-09-30: it overrode LaTeX's author separator \and (printed "antisymmetric" between authors)
\def\ass{assumption}
\def\cadlag{c\`adl\`ag}
\def\ch{characteristic polynomial}
\def\chp{characteristic polynomial}
\def\coe{coefficient}
\def\comp{computation}
\def\con{condition}
\def\coo{coordinate}
\def\corr{corresponding}
\def\dec{decomposition}
\def\den{denominator}
\def\dete{deterministic}
\def\eig{eigenvalue}\def\eval{eigenvalue}
\def\eigv{eigenvector}\def\eiv{eigenvector}\def\evec{eigenvector}
\def\elt{element}
\def\enn{essentially nonnegative/positivity preserving}
\def\equ{equation}
\def\etal{et al.}
\def\fe{for example}
\def\fno{from now on}
\def\for{\forall}
\def\form{formula}
\def\Fr{Furthermore, }
\def\frt{furthermore}
\def\fun{function}
\def\gene{generalization}
\def\heu{heuristic}
\def\how{however}\def\How{However, }
\def\ie{i.e. }
\def\im{\item}
\def\isd{is defined by}
\def\itc{it is easy to check that}
\def\itf{it follows that}\def\Itf{It follows that}
\def\ith{it holds that}
\def\itm{it may be checked that}\def\Itm{It may be checked that}
\def\ito{it turns out that}\def\Ito{It turns out that}\def\Itou{It turns out useful to}
\def\ma{mass action}\def\MAR{mass action representation}
\def\mat{matrix}
\def\mbc{may be checked to be}
\def\mbw{may be written as}\def\mubw{must be written as}
\def\mpr{more precisely, }\def\Mp{More precisely, }
\def\nc{non-catalytic}
\def\nne{non-negative}\def\nny{nonnegativity}
\def\no{\nonumber}
\def\noni{non-infectious}
\def\nota{notation}
\def\np{\newpage}
\def\Nt{Note that}
\def\nt{non trivial}\def\nz{non zero}
\def\oth{otherwise}
\def\Oth{On the other hand}
\def\pa{parameter}\def\para{parameter}
\def\pol{polynomial}
\def\prob{probability}
\def\Prf{{\bf  Proof:}}\def\prf{{\bf  Proof:}}
\def\prop{proportional to}
\def\qu{\quad}
\def\QED{\hfill {$\square$}\goodbreak \medskip}
\def\rep{representation}
\def\resp{respectively}
\def\satd{satisfied}\def\satg{satisfying}\def\sats{satisfies}\def\saty{satisfy}
\def\ssec{\subsection}\def\sssec{\subsubsection}
\def\sef{see for example}
\def\sksy{skew-symmetric}
\def\sub{substrate}
\def\sys{system}
\def\var{variable}
\def\wk{well-known}
\def\wl{without loss of generality}\def\wlo{w.l.o.g.}
\def\wmc{We may check that}\def\Wmc{We may check that}
\def\wmr{we may rewrite}
\def\wms{we must show that}
\def\wrt{with respect to}

% -- CRN & Chemistry --
\newcommand\CRN{chemical reaction network}
\newcommand\CRNs{chemical reaction networks}
\newcommand\RN{reaction network}
\def\ACR{absolute concentration robustness}
\def\AMSD{acyclic siphon decomposition}
\def\DAE{differential algebraic equations}
\def\FHJ{Feinberg-Horn-Jackson}
\def\GMAK{generalized mass action kinetics}
\def\MM{Michaelis-Menten}
\def\sm{stoichiometric matrix}\def\Sm{stoichiometric matrix}
\def\WRZD{weakly reversible, zero deficiency}\def\aut{autocatalytic}

% -- Epidemiology --
\def\brn{basic reproduction number}\def\bzn{basic replacement number}
\def\com{compartment}
\def\dfe{disease free equilibrium}\def\DFE{disease free equilibrium}
\def\dwt{dwell times}\def\edt{expected dwell times}
\def\EM{epidemic model}
\def\endp{endemic fixed point}\def\enp{endemic point}\def\EE{endemic point}
\def\ME{mathematical epidemiology}
\def\MEE{eco-epidemiology}\def\EEP{eco-epidemiological}
\def\ece{eco-epidemiological model}
\def\mgm{multi group model}
\def\mtS{multi-strain model}\def\mSt{multi-strain model}
\def\mSi{minimal siphon}
\def\NGM{next generation matrix}
\def\nim{new infections matrix}
\def\qem{quadratic epidemiological model}
\def\repF{invasion functions}
\def\SPF{SIR-PH-FA model}
\def\sus{susceptible}

% -- Graph theory & Linear algebra --
\def\jac{Jacobi matrix}
\def\jin{Jacobian matrix of the invasion vector field with respect to the invasion variables}
\def\jrn{Jacobian matrix of the resident vector field with respect to the resident variables}
\def\jxn{Jacobian matrix of the $\x$ vector field with respect to the $\x$ variables}
\def\Lm{Laplacian matrix}
\def\MTT{matrix tree theorem}
\def\SGT{spectral graph theory}
\def\SMW{Sherman–Morrison–Woodbury}

% -- Stability & Dynamics --
\def\absC{absorbtion CTMC}
\def\BTb{Bogdanov-Takens bifurcation}
\def\CEP{competitive exclusion partition}
\def\CTMC{continuous time Markov chain}
\def\ds{differential system}
\def\fp{fixed point}\def\bfp{boundary fixed point}\def\fbp{fixed boundary point}
\def\fbl{fixed boundary line}\def\fif{forward invariant face}
\def\gas{global asymptotic stability}\def\GAS{globally asymptotically stable}
\def\GASS{\text{GASS}}
\def\Hb{Hopf bifurcation}
\def\Hur{Hurwitz}\def\MH{Metzler-Hurwitz}\def\RH{Routh-Hurwitz}
\def\IS{\text{IS}}
\def\LAS{locally asymptotically stable}
\def\LdS{Lyapunov diagonally stable}
\def\Lf{Lyapunov function}\def\LyV{Lyapunov Volterra type function}\def\RLF{robust Lyapunov function}
\def\LyM{Lyapunov Malkin stable}
\def\PF{Perron-Frobenius}
\def\spfp{strictly positive fixed point}
\def\sta{stability}
\def\STP{sharp threshold property}\def\STT{sharp threshold theorem}

% -- Models (with \varphi explicitly instead of \f) --
\def\AB{$(A,B)$ Arino-Brauer epidemic models}
\def\ABa{$(A,\bbe,\ba,\varphi)$ balanced bilinear model}
\def\ABP{%$(A,\bbe,P,\varphi)$
\\ balanced bilinear model}
\def\AMG{Anderson May Gupta}
\newcommand\BIN{biological interaction network}
\def\Dt{Descartes type}
\def\gLV{generalized Lotka-Volterra model}
\def\IG{\text{invasion graph}}\def\IGMS{interaction graph on minimal siphons}
\def\LCP{\text{linear complementarity problem}}
\def\LV{Lotka-Volterra}\def\LVK{Lotka-Volterra-Kolmogorov}

% -- Splitting & Structure --
\def\reg{regular splitting}\def\regS{{\bf regular splitting}}
\def\RUR{rational univariate representation}
\def\bds{block-diagonal structure}\def\NGMs{next generation matrices}

% -- Miscellaneous --
\def\Ma{{Mathematica}}
\def\Mm{Max–Min}
\def\num{numerator}
\def\rap{rational approximation} \def\DSR{directed species-reaction bipartite graph}\def\PCP{polynomial complementarity problem}\def\rat{rational substitution}
\def\QSS{quasi steady state} \def\RHS{right hand side} \def\LHS{left hand side}
\def\OSN{online social network}\def\VdD{Van den Driessche}
\def\wli{$\omega$-limit} \def\rk{rich kinetics}\def\sJ{symbolic Jacobian}
\def\cV{\cite{Vas,VasSta,BSV,GSV}} \def\spf{semi-parametric semi-functional}
\def\CRW{Capasso-Ruan-Wang} \def\ABP{$(A,\bbe,P,\f)$ balanced bilinear model}
\def\ABa{%$(A,\bbe,\ba,\f)$ \\
rank-one balanced bilinear model}
\def\bbe{B%{\bff \beta}
} \def\vbe{\vec {\bff \beta}
}
\def\vb{\vec {\bff b}}
\newcommand{\Kh}{\hat{K}}          % K with hat, first mention only
\newcommand{\wS}{w^{\sigma}}       % right Perron J^\perp, full block
\newcommand{\piS}{\vec \pi^{\sigma}}    % left Perron J^\perp, full block
\newcommand{\ws}{w^{\sigma}}       % right Perron J^\perp, sub-block
\newcommand{\pis}{\pi^{\sigma}}    % left Perron J^\perp, sub-block
\newcommand{\hwS}{\hat{w}^{\sigma}}   % right Perron K, full block
\newcommand{\hpiS}{\hat{\pi}^{\sigma}} % left Perron K, full block
\newcommand{\hws}{\hat{w}^{\sigma}}   % right Perron K, sub-block
\newcommand{\hpis}{\hat{\pi}^{\sigma}} % left Perron K, sub-block
\newcommand{\twS}{\tilde{w}^{\sigma}}  % right Perron \tilde K, full block
\newcommand{\tpiS}{\tilde{\pi}^{\sigma}} % left Perron \tilde K
\newcommand{\tws}{\tilde{w}^{\sigma}}
\newcommand{\tpis}{\tilde{\pi}^{\sigma}}
\newcommand{\Mat}[2]{\mathbb{R}_{\ge 0}^{\,#1\times #2}}
\def\tJ{J_\sigma^\perp} \def\lPe{left Perron eigenvector}
% \G is \Gamma (defined at :795).  The Volterra function keeps its own name,
% so that \G no longer has two conflicting meanings in this file.
\def\GV{G}  % Volterra G(u)=u-1-\ln u
\def\PVL{Perron-Volterra Lyapunov minimal siphon ansatz}
\newcommand{\istar}[1]{i_{#1}^{*}}
%\newcommand{\sE}{s_E}
\newcommand{\vE}{v_E}
\def\PV{Perron-Volterra candidate Lyapunov function}
\def\GLV{generalized Lotka-Volterra}
\def\MERN{mathematical epidemiology reaction networks}
\def\CMM{contraction metric method}
\newcommand{\sign}{\operatorname{sign}}
\newcommand{\Var}{\operatorname{Var}} 

% ======================================================================
% IDEAS SHARED BETWEEN THE ABSTRACT AND THE MAIN RESULTS SECTION
%
% Write each idea ONCE, in the abstract; the Main Results section recalls it,
% so the two can never drift apart.
%
%   in the abstract        \ideaitem{1}{The first-order reaction terms ...}
%   in Main Results        \ideahead{1}
%
% \ideaitem stores the text globally and typesets
%   \item \hyperlink{idea:1}{<text>}
% \ideahead typesets
%   \item\hypertarget{idea:1}{}\textbf{<the same text>}
%
% The abstract must come before the Main Results input, which is the normal
% order. If a heading has to read differently from the abstract sentence (for
% instance a long abstract item summarised in one line), override it with the
% optional argument: \ideahead[short heading]{6}.
% ======================================================================
\makeatletter
\newcommand{\ideaitem}[2]{%
  \expandafter\gdef\csname @idea@#1\endcsname{#2}%
  \item\hyperlink{idea:#1}{#2}}
\newcommand{\ideahead}[2][]{%
  \item\hypertarget{idea:#2}{}%
  \def\@ideatmp{#1}%
  \ifx\@ideatmp\@empty
    \@ifundefined{@idea@#2}{%
      \PackageWarning{preamble}{Idea #2 is used in Main Results but was never
        defined by \string\ideaitem\space in the abstract}%
      \textbf{[idea #2 undefined]}%
    }{\textbf{\begingroup\def\ideaterm##1##2{##2}\def\idearest##1{}\@nameuse{@idea@#2}\endgroup}}%
  \else
    \textbf{#1}%
  \fi}
% \ideaterm{<hypertarget>}{<words>}: inside an abstract idea the words become an
% invisible link (the document sets \hypersetup{hidelinks}) jumping to the place
% in Main Results that defines them; inside the Main Results heading itself,
% where the target sits a few lines below, \ideahead strips the link and prints
% the words plainly. Mark the destination with \hypertarget{<name>}{} on the line
% before the display, which adds no equation number and shows nothing.
\newcommand{\ideaterm}[2]{\hyperlink{#1}{#2}}
% \idearest{<text>}: the part of an abstract idea after its first sentence; printed in
% the abstract, omitted from the Main Results heading.
\newcommand{\idearest}[1]{#1}
\makeatother

% ======================================================================
% SHORT-PAPER SWITCH
% \shorttrue in a master that assembles only part of the companion paper
% (Barrat/SipArxiv.tex). \nsh{...} drops text that cites results living in the
% omitted sections, so the short build has no dangling references while the
% companion paper keeps the text verbatim.
% ======================================================================
\newif\ifshort \shortfalse
\newcommand{\nsh}[1]{\ifshort\else#1\fi}
% \shortalt{<short build>}{<companion paper>}: use when a sentence needs
% different wording in the two builds, not merely a deletion.
\newcommand{\shortalt}[2]{\ifshort#1\else#2\fi}
% \Jtr, \JtrS{\Sigma}: transverse Jacobian. J^\perp here; SipArxiv.tex redefines them as J_{\mathrm{tra}}, J_{\mathrm{tra},\Sigma}
\newcommand{\Jtr}{J^\perp}
\newcommand{\JtrS}[1]{J_{#1}^\perp}
, and must not clash.
\providecommand{\Q}{\mathbb{Q}}
\providecommand{\res}{\operatorname{res}}   % resultant; used by epiRes.tex
\newcommand{\MN}{\mathbb{N}}
\newcommand{\MZ}{\mathbb{Z}}
\newcommand{\realpos}{\mathbb{R}_{> 0}}
\newcommand{\realnneg}{\mathbb{R}_{\geq 0}}
\newcommand{\naturalpos}{\mathbb{N}_{>0}}
\newcommand{\Rplus}{\mathbb{R}_+}
\newcommand{\Rp}{\mathbb{R}_{\ge 0}}
\def\Rnne{\mathbb{R}_{\ge 0}}
\def\Se{\bar{S}} \def\Ie{\bar{I}}
\newcommand{\Rpos}{\mathbb{R}_{>0}}
\newcommand{\dd}{\,\mathrm d}
\newcommand{\e}{\mathrm e}
%% ========================================
%% GREEK LETTERS (2+ character abbreviations) greek
%% ========================================
\def\al{\alpha}\def\ba{\bff \alpha}
\def\w{\omega}\def\Lgn{\ga+\mu+\de}\def\om{\omega}
\def\de{\delta}\def\De{\Delta}\def\bde{\bff \delta}
\def\ep{\epsilon}\def\eps{\varepsilon}
\def\f{\varphi}  % Kept as requested
\def\ga{\gamma}\def\Ga{\Gamma}\def\G{\Gamma}\def\bg{\bff \gamma}
\def\up{\upsilon}  % NOT the breve accent: no accent commands anywhere.
                   % Named \up, not \u, so the LaTeX breve stays reachable.
\def\la{\lambda}\def\La{\Lambda}
\def\si{\sigma}\def\ka{\kappa}
\def\th{\theta}\def\vt{\vartheta}
\def\ze{\zeta}\def\Gap{\Gamma^{+}} \def\Gam{\Gamma^{-}}


%% ========================================
%% SPECIAL SHORT COMMANDS (kept as requested)
%% ========================================
\def\I{\infty}
\def\T{\widetilde}
\def\H{\hat}
\newcommand{\relint}{\operatorname{relint}}

%% ========================================
%% BOLD/VECTOR NOTATION
%% ========================================
\newcommand{\bff}[1]{{\mbox{\boldmath$#1$}}}
\newcommand\X{\boldsymbol{X}}
\def\x{\boldsymbol{x}}\def\bx{\bff x}\def\xe{\bar{\x}}
\newcommand\y{\boldsymbol{y}}\def\by{\boldsymbol{y}}\def\ye{\bar{\y}}
\def\yb{\y}\def\xb{\x}
\def\xy{\x-\y}
\newcommand\br{\boldsymbol{r}}\def\re{\bar{\r}}
\newcommand\bd{\boldsymbol{d}}
\newcommand\bQ{\boldsymbol{Q}}
\newcommand{\Yb}{{\mathbf{Y}}}
\newcommand{\Ab}{\mathbf{A}}
\newcommand{\zb}{\mathbf{z}}\def\bz{\bff 0}
\newcommand{\Zb}{\mathbf{Z}}
\newcommand{\bb}{\mathbf{b}}
\def\vbe{\vec {\bff \beta}}
\def\vb{\vec {\bff b}}
\def\v1{\vec {\bff 1}}
\def\vn{\vec \de}
\def\vz{\vec 0}
\def\vi{\overset{\rightarrow}{\mathsf i}}
\def\vv{\overset{\rightarrow}{v}}
\def\vr{\; \vec {\mathsf r}}
\newcommand{\va}{\bff a}
\providecommand{\vect}[1]{\boldsymbol{#1}}   % bold vector (moved from Gradostat.tex)
\providecommand{\one}{\mathbf 1}             % all-ones vector (moved from AliSF)
\providecommand{\spr}{s}                     % spectral bound  (moved from AliSF)
\providecommand{\Ke}{K_{\rm eff}}            % effective carrying capacity
% Critical right/left eigenvectors in the Hopf projection method (AAHOSN/New.tex)
\providecommand{\bq}{\mathbf{q}}
\providecommand{\bp}{\mathbf{p}}


%% ========================================
%% MATHCAL & MATHSCR
%% ========================================
% Mathcal
\newcommand\mc[1]{\mathcal{#1}}
\def\mC{{\mathcal C}}\def\mD{\mathcal D}\def\mE{\mathcal{E}}
\def\mF{\mathcal F}\def\mG{{\mathcal G}}\def\mH{{\mathcal H}}
\def\mI{\mathcal I}\def\mL{\mathcal{L}}\def\mM{{\mathcal M}}
\def\mN{{\mathcal N}}\def\mP{\mathcal{P}}\def\mR{{\mathcal R}}
\def\mS{{\mathcal S}}\def\mT{\mathcal{T}}\def\mV{{\mathcal V}}
\def\mK{\mathcal K}\def\mA{{\mathcal A}}\def\mU{{\mathcal U}}
% Symmetric bilinear and trilinear forms of the Taylor expansion, used in the
% first Lyapunov coefficient (AAHOSN/New.tex)
\providecommand{\Bform}{\mathcal{B}}
\providecommand{\Cform}{\mathcal{C}}

% Blackboard bold.  Kept here, not in the drivers: a \providecommand in a
% driver preamble runs before this file and silently wins, so the same name
% ends up meaning different things in different builds.
\def\IC{\mathbb{C}}\def\IN{\mathbb{N}}\def\IQ{\mathbb{Q}}
\def\IR{\mathbb{R}}\def\IZ{\mathbb{Z}}\def\IG{\mathbb{G}}
\def\II{\mathbb{I}}

% Mathscr
\newcommand{\Rs}{{\mathscr R}}
\newcommand{\Ss}{{\mathscr S}}
\newcommand{\Ns}{{\mathscr N}}
\newcommand{\Cs}{\mathscr C}


%% ========================================
%% OPERATORS & FUNCTIONS
%% ========================================
\DeclareMathOperator{\tr}{tr}
\DeclareMathOperator{\diag}{Diag}
\DeclareMathOperator{\Diag}{Diag}
\DeclareMathOperator{\intr}{\mathrm{Int}}
\DeclareMathOperator{\Inv}{Inv}
\DeclareMathOperator{\spec}{spec}
\DeclareMathOperator{\AUC}{AUC}               % area under the curve
\newcommand{\ii}{\mathrm{i}}                  % imaginary unit
\newcommand\op{\operatorname}
\newcommand{\supp}{\op{supp}}
\newcommand{\diff}{\mathop{}\mathopen{}\mathrm{d}}
\newcommand{\pdd}[1]{\frac{\partial}{\partial{#1}}}
\DeclareMathOperator{\cone}{cone}
\DeclareMathOperator{\proj}{proj}

%% ========================================
%% BRACKETS & NORMS
%% ========================================
\providecommand{\pr}[1]{\left(#1\right)}
\providecommand{\pp}[1]{\left[#1\right]}
\newcommand\croc[1]{\left\langle #1\right\rangle}
\newcommand{\norm}[1]{\|#1\|}
\newcommand{\inner}[2]{\langle #1, #2 \rangle}
% Short synonym of \inner, used by AAHOSN/New.tex as $\ip{u}{v}=\bar u^\top v$
\providecommand{\ip}[2]{\langle #1 , #2 \rangle}


%% ========================================
%% FORMATTING & UTILITIES
%% ========================================
\newcommand*{\Scale}[2][4]{\scalebox{#1}{$#2$}}
\newcommand{\te}[1]{\text{#1}}
\newcommand\steq[1]{\stackrel{\text{\rm #1.}}{=}}
\newcommand\blue[1]{\textcolor{blue}{#1}}
\newcommand{\red}{\textcolor[rgb]{1.00,0.00,0.00}}
\newcommand{\green}{\textcolor[rgb]{0.00,1.00,0.00}}

% Figure command
\newcommand{\figu}[3]{%
\begin{figure}[H]
\centering
\includegraphics[scale=#3]{#1}
\caption{#2\label{f:#1}}
\end{figure}
}

% Footnote
\long\def\symbolfootnote[#1]#2{%
\begingroup
\def\thefootnote{\fnsymbol{footnote}}\footnote[#1]{#2}%
\endgroup}
\def\fn{\symbolfootnote}

\newcommand{\dsum}{\displaystyle\sum}

%% ========================================
%% CRN & EPIDEMIOLOGY SPECIFIC
%% ========================================
% CRN notation
\def\SS{\mathcal S}
\def\CC{\mathcal C}
\def\RR{\mathcal R}

% Reproduction numbers
\def\Rto{R_{1}^{E_2}}
\def\Rot{R_{2}^{E_1}}
\def\Roc{R_{1,c}}
\def\Rtc{R_{2,c}}
\def\rd{r_{0}}
\def\sd{s_{0}}
\def\dd{D_{0}}
\def\vd{v_{0}}
\def\zd{z_{0}}
\def\se{\sE}
\def\xd{x_{0}}
%% ========================================
%% SPECIAL VALUES & EQUILIBRIA
%% ========================================
\def\eE{\bar{e}}
\def\wE{\bar{w}}
\def\sE{s_{*}}
\def\Me{\T M_{ee}}
\def\Rnne{\mathbb{R}_{\ge 0}}

\newcommand{\ones}{\mathbf 1}


%% ========================================
%% ABBREVIATIONS (alphabetical by category)
%% ========================================

% -- General mathematical terms --
\def\adm{admissible}\def\admy{admissibility}
%\def\and{antisymmetric}% disabled 2026-09-30: it overrode LaTeX's author separator \and (printed "antisymmetric" between authors)
\def\ass{assumption}
\def\cadlag{c\`adl\`ag}
\def\ch{characteristic polynomial}
\def\chp{characteristic polynomial}
\def\coe{coefficient}
\def\comp{computation}
\def\con{condition}
\def\coo{coordinate}
\def\corr{corresponding}
\def\dec{decomposition}
\def\den{denominator}
\def\dete{deterministic}
\def\eig{eigenvalue}\def\eval{eigenvalue}
\def\eigv{eigenvector}\def\eiv{eigenvector}\def\evec{eigenvector}
\def\elt{element}
\def\enn{essentially nonnegative/positivity preserving}
\def\equ{equation}
\def\etal{et al.}
\def\fe{for example}
\def\fno{from now on}
\def\for{\forall}
\def\form{formula}
\def\Fr{Furthermore, }
\def\frt{furthermore}
\def\fun{function}
\def\gene{generalization}
\def\heu{heuristic}
\def\how{however}\def\How{However, }
\def\ie{i.e. }
\def\im{\item}
\def\isd{is defined by}
\def\itc{it is easy to check that}
\def\itf{it follows that}\def\Itf{It follows that}
\def\ith{it holds that}
\def\itm{it may be checked that}\def\Itm{It may be checked that}
\def\ito{it turns out that}\def\Ito{It turns out that}\def\Itou{It turns out useful to}
\def\ma{mass action}\def\MAR{mass action representation}
\def\mat{matrix}
\def\mbc{may be checked to be}
\def\mbw{may be written as}\def\mubw{must be written as}
\def\mpr{more precisely, }\def\Mp{More precisely, }
\def\nc{non-catalytic}
\def\nne{non-negative}\def\nny{nonnegativity}
\def\no{\nonumber}
\def\noni{non-infectious}
\def\nota{notation}
\def\np{\newpage}
\def\Nt{Note that}
\def\nt{non trivial}\def\nz{non zero}
\def\oth{otherwise}
\def\Oth{On the other hand}
\def\pa{parameter}\def\para{parameter}
\def\pol{polynomial}
\def\prob{probability}
\def\Prf{{\bf  Proof:}}\def\prf{{\bf  Proof:}}
\def\prop{proportional to}
\def\qu{\quad}
\def\QED{\hfill {$\square$}\goodbreak \medskip}
\def\rep{representation}
\def\resp{respectively}
\def\satd{satisfied}\def\satg{satisfying}\def\sats{satisfies}\def\saty{satisfy}
\def\ssec{\subsection}\def\sssec{\subsubsection}
\def\sef{see for example}
\def\sksy{skew-symmetric}
\def\sub{substrate}
\def\sys{system}
\def\var{variable}
\def\wk{well-known}
\def\wl{without loss of generality}\def\wlo{w.l.o.g.}
\def\wmc{We may check that}\def\Wmc{We may check that}
\def\wmr{we may rewrite}
\def\wms{we must show that}
\def\wrt{with respect to}

% -- CRN & Chemistry --
\newcommand\CRN{chemical reaction network}
\newcommand\CRNs{chemical reaction networks}
\newcommand\RN{reaction network}
\def\ACR{absolute concentration robustness}
\def\AMSD{acyclic siphon decomposition}
\def\DAE{differential algebraic equations}
\def\FHJ{Feinberg-Horn-Jackson}
\def\GMAK{generalized mass action kinetics}
\def\MM{Michaelis-Menten}
\def\sm{stoichiometric matrix}\def\Sm{stoichiometric matrix}
\def\WRZD{weakly reversible, zero deficiency}\def\aut{autocatalytic}

% -- Epidemiology --
\def\brn{basic reproduction number}\def\bzn{basic replacement number}
\def\com{compartment}
\def\dfe{disease free equilibrium}\def\DFE{disease free equilibrium}
\def\dwt{dwell times}\def\edt{expected dwell times}
\def\EM{epidemic model}
\def\endp{endemic fixed point}\def\enp{endemic point}\def\EE{endemic point}
\def\ME{mathematical epidemiology}
\def\MEE{eco-epidemiology}\def\EEP{eco-epidemiological}
\def\ece{eco-epidemiological model}
\def\mgm{multi group model}
\def\mtS{multi-strain model}\def\mSt{multi-strain model}
\def\mSi{minimal siphon}
\def\NGM{next generation matrix}
\def\nim{new infections matrix}
\def\qem{quadratic epidemiological model}
\def\repF{invasion functions}
\def\SPF{SIR-PH-FA model}
\def\sus{susceptible}

% -- Graph theory & Linear algebra --
\def\jac{Jacobi matrix}
\def\jin{Jacobian matrix of the invasion vector field with respect to the invasion variables}
\def\jrn{Jacobian matrix of the resident vector field with respect to the resident variables}
\def\jxn{Jacobian matrix of the $\x$ vector field with respect to the $\x$ variables}
\def\Lm{Laplacian matrix}
\def\MTT{matrix tree theorem}
\def\SGT{spectral graph theory}
\def\SMW{Sherman–Morrison–Woodbury}

% -- Stability & Dynamics --
\def\absC{absorbtion CTMC}
\def\BTb{Bogdanov-Takens bifurcation}
\def\CEP{competitive exclusion partition}
\def\CTMC{continuous time Markov chain}
\def\ds{differential system}
\def\fp{fixed point}\def\bfp{boundary fixed point}\def\fbp{fixed boundary point}
\def\fbl{fixed boundary line}\def\fif{forward invariant face}
\def\gas{global asymptotic stability}\def\GAS{globally asymptotically stable}
\def\GASS{\text{GASS}}
\def\Hb{Hopf bifurcation}
\def\Hur{Hurwitz}\def\MH{Metzler-Hurwitz}\def\RH{Routh-Hurwitz}
\def\IS{\text{IS}}
\def\LAS{locally asymptotically stable}
\def\LdS{Lyapunov diagonally stable}
\def\Lf{Lyapunov function}\def\LyV{Lyapunov Volterra type function}\def\RLF{robust Lyapunov function}
\def\LyM{Lyapunov Malkin stable}
\def\PF{Perron-Frobenius}
\def\spfp{strictly positive fixed point}
\def\sta{stability}
\def\STP{sharp threshold property}\def\STT{sharp threshold theorem}

% -- Models (with \varphi explicitly instead of \f) --
\def\AB{$(A,B)$ Arino-Brauer epidemic models}
\def\ABa{$(A,\bbe,\ba,\varphi)$ balanced bilinear model}
\def\ABP{%$(A,\bbe,P,\varphi)$
\\ balanced bilinear model}
\def\AMG{Anderson May Gupta}
\newcommand\BIN{biological interaction network}
\def\Dt{Descartes type}
\def\gLV{generalized Lotka-Volterra model}
\def\IG{\text{invasion graph}}\def\IGMS{interaction graph on minimal siphons}
\def\LCP{\text{linear complementarity problem}}
\def\LV{Lotka-Volterra}\def\LVK{Lotka-Volterra-Kolmogorov}

% -- Splitting & Structure --
\def\reg{regular splitting}\def\regS{{\bf regular splitting}}
\def\RUR{rational univariate representation}
\def\bds{block-diagonal structure}\def\NGMs{next generation matrices}

% -- Miscellaneous --
\def\Ma{{Mathematica}}
\def\Mm{Max–Min}
\def\num{numerator}
\def\rap{rational approximation} \def\DSR{directed species-reaction bipartite graph}\def\PCP{polynomial complementarity problem}\def\rat{rational substitution}
\def\QSS{quasi steady state} \def\RHS{right hand side} \def\LHS{left hand side}
\def\OSN{online social network}\def\VdD{Van den Driessche}
\def\wli{$\omega$-limit} \def\rk{rich kinetics}\def\sJ{symbolic Jacobian}
\def\cV{\cite{Vas,VasSta,BSV,GSV}} \def\spf{semi-parametric semi-functional}
\def\CRW{Capasso-Ruan-Wang} \def\ABP{$(A,\bbe,P,\f)$ balanced bilinear model}
\def\ABa{%$(A,\bbe,\ba,\f)$ \\
rank-one balanced bilinear model}
\def\bbe{B%{\bff \beta}
} \def\vbe{\vec {\bff \beta}
}
\def\vb{\vec {\bff b}}
\newcommand{\Kh}{\hat{K}}          % K with hat, first mention only
\newcommand{\wS}{w^{\sigma}}       % right Perron J^\perp, full block
\newcommand{\piS}{\vec \pi^{\sigma}}    % left Perron J^\perp, full block
\newcommand{\ws}{w^{\sigma}}       % right Perron J^\perp, sub-block
\newcommand{\pis}{\pi^{\sigma}}    % left Perron J^\perp, sub-block
\newcommand{\hwS}{\hat{w}^{\sigma}}   % right Perron K, full block
\newcommand{\hpiS}{\hat{\pi}^{\sigma}} % left Perron K, full block
\newcommand{\hws}{\hat{w}^{\sigma}}   % right Perron K, sub-block
\newcommand{\hpis}{\hat{\pi}^{\sigma}} % left Perron K, sub-block
\newcommand{\twS}{\tilde{w}^{\sigma}}  % right Perron \tilde K, full block
\newcommand{\tpiS}{\tilde{\pi}^{\sigma}} % left Perron \tilde K
\newcommand{\tws}{\tilde{w}^{\sigma}}
\newcommand{\tpis}{\tilde{\pi}^{\sigma}}
\newcommand{\Mat}[2]{\mathbb{R}_{\ge 0}^{\,#1\times #2}}
\def\tJ{J_\sigma^\perp} \def\lPe{left Perron eigenvector}
% \G is \Gamma (defined at :795).  The Volterra function keeps its own name,
% so that \G no longer has two conflicting meanings in this file.
\def\GV{G}  % Volterra G(u)=u-1-\ln u
\def\PVL{Perron-Volterra Lyapunov minimal siphon ansatz}
\newcommand{\istar}[1]{i_{#1}^{*}}
%\newcommand{\sE}{s_E}
\newcommand{\vE}{v_E}
\def\PV{Perron-Volterra candidate Lyapunov function}
\def\GLV{generalized Lotka-Volterra}
\def\MERN{mathematical epidemiology reaction networks}
\def\CMM{contraction metric method}
\newcommand{\sign}{\operatorname{sign}}
\newcommand{\Var}{\operatorname{Var}} 

% ======================================================================
% IDEAS SHARED BETWEEN THE ABSTRACT AND THE MAIN RESULTS SECTION
%
% Write each idea ONCE, in the abstract; the Main Results section recalls it,
% so the two can never drift apart.
%
%   in the abstract        \ideaitem{1}{The first-order reaction terms ...}
%   in Main Results        \ideahead{1}
%
% \ideaitem stores the text globally and typesets
%   \item \hyperlink{idea:1}{<text>}
% \ideahead typesets
%   \item\hypertarget{idea:1}{}\textbf{<the same text>}
%
% The abstract must come before the Main Results input, which is the normal
% order. If a heading has to read differently from the abstract sentence (for
% instance a long abstract item summarised in one line), override it with the
% optional argument: \ideahead[short heading]{6}.
% ======================================================================
\makeatletter
\newcommand{\ideaitem}[2]{%
  \expandafter\gdef\csname @idea@#1\endcsname{#2}%
  \item\hyperlink{idea:#1}{#2}}
\newcommand{\ideahead}[2][]{%
  \item\hypertarget{idea:#2}{}%
  \def\@ideatmp{#1}%
  \ifx\@ideatmp\@empty
    \@ifundefined{@idea@#2}{%
      \PackageWarning{preamble}{Idea #2 is used in Main Results but was never
        defined by \string\ideaitem\space in the abstract}%
      \textbf{[idea #2 undefined]}%
    }{\textbf{\begingroup\def\ideaterm##1##2{##2}\def\idearest##1{}\@nameuse{@idea@#2}\endgroup}}%
  \else
    \textbf{#1}%
  \fi}
% \ideaterm{<hypertarget>}{<words>}: inside an abstract idea the words become an
% invisible link (the document sets \hypersetup{hidelinks}) jumping to the place
% in Main Results that defines them; inside the Main Results heading itself,
% where the target sits a few lines below, \ideahead strips the link and prints
% the words plainly. Mark the destination with \hypertarget{<name>}{} on the line
% before the display, which adds no equation number and shows nothing.
\newcommand{\ideaterm}[2]{\hyperlink{#1}{#2}}
% \idearest{<text>}: the part of an abstract idea after its first sentence; printed in
% the abstract, omitted from the Main Results heading.
\newcommand{\idearest}[1]{#1}
\makeatother

% ======================================================================
% SHORT-PAPER SWITCH
% \shorttrue in a master that assembles only part of the companion paper
% (Barrat/SipArxiv.tex). \nsh{...} drops text that cites results living in the
% omitted sections, so the short build has no dangling references while the
% companion paper keeps the text verbatim.
% ======================================================================
\newif\ifshort \shortfalse
\newcommand{\nsh}[1]{\ifshort\else#1\fi}
% \shortalt{<short build>}{<companion paper>}: use when a sentence needs
% different wording in the two builds, not merely a deletion.
\newcommand{\shortalt}[2]{\ifshort#1\else#2\fi}
% \Jtr, \JtrS{\Sigma}: transverse Jacobian. J^\perp here; SipArxiv.tex redefines them as J_{\mathrm{tra}}, J_{\mathrm{tra},\Sigma}
\newcommand{\Jtr}{J^\perp}
\newcommand{\JtrS}[1]{J_{#1}^\perp}
.
\makeatletter
\@ifpackageloaded{hyperref}{}{%
  \@ifundefined{hyperrefpackage}{%
    \@ifundefined{AAdefNnohyperref}{\usepackage{hyperref}}{}%
  }{}%
}
\makeatother
% linktoc=all makes a table-of-contents page number a link as well as the
% entry text; hyperref's default, linktoc=section, links only the text.
% Deferred to begin-document and guarded: by the block above hyperref may
% be loaded later (mdpi's own hook) or not at all (\AAdefNnohyperref), and
% \hypersetup must exist before it is called.  A driver that wants the
% default back can say \hypersetup{linktoc=section} after % Shared preamble for the GLVbook papers.  This is the single canonical copy;
% a driver one level below GLVbook reaches it as
%
%     \input{../basicFiles/preamble}
%
% What stays in the driver, because it differs from paper to paper:
% \documentclass, font choice, \title, \author, \date, and any macro genuinely
% local to one paper.  Everything else belongs here.  The file is
% self-contained: it does not read any other file.

\usepackage[utf8]{inputenc}
\usepackage[novbox]{pdfsync}   % required: SyncTeX alone does not work here

% Page geometry, set once here for every paper that loads this file.  A driver
% that needs different margins says \geometry{...}; a second
% \usepackage[...]{geometry} raises an option clash and must not be used.
\makeatletter
\@ifundefined{AAdefNnogeometry}{\usepackage[a4paper,margin=2cm]{geometry}}{}
\makeatother

% mathematics
\usepackage{amsmath,amssymb,amsthm,mathtools,bm,mathrsfs,bbm}
\usepackage{centernot}

% floats, tables, lists
\usepackage{graphicx,float}
\makeatletter
\@ifpackageloaded{subfig}{}{\usepackage{subcaption}}
\makeatother
\usepackage{booktabs,longtable,threeparttable,array,multirow}
\usepackage{enumitem}

% colour, code, algorithms
\usepackage{xcolor,listings,tcolorbox}
\makeatletter
\@ifundefined{c@algorithm}{\usepackage{algorithm}}{}
\makeatother
\usepackage{algorithmicx,algpseudocode}

% diagrams
\usepackage{tikz,tikz-cd,pgfplots}
\usetikzlibrary{graphs,decorations.pathmorphing,decorations.markings,
arrows.meta,automata,positioning}
\tikzcdset{every label/.append style = {font = \small}}
\pgfplotsset{compat=1.18}

% chemistry
\usepackage[version=4]{mhchem}

\usepackage{microtype}
\usepackage{lineno}
\usepackage{nomencl}
\makenomenclature

% hyperref first, then bookmark, then glossaries: glossaries must follow
% hyperref or the entry links are built against the wrong anchors.
\makeatletter
\@ifundefined{hyperrefpackage}{%
  \@ifundefined{AAdefNnohyperref}{\usepackage{hyperref}}{}%
}{}
\@ifundefined{hyperrefpackage}{\usepackage{bookmark}}{}
\makeatother
\usepackage[toc,sort=use]{glossaries}
\makeglossaries

% Book/standalone switch.  Material that the book already covers in its
% prelude, but that a standalone paper still has to state, is wrapped as
%     \ifbookbuild\else ... \fi
% Drivers default to standalone; bookM.tex says \bookbuildtrue.
\newif\ifbookbuild
\bookbuildfalse

% Part switches for the gradostat paper.  Both default to true; a driver that
% wants the short submission says \madecfalse and/or \slowfastfalse, and every
% cross-reference into the dropped part disappears with it.
%   \ifmadec     Part III, the Lagasquie--Madec model and its break-even geometry
%   \ifslowfast  Part IV, everything carrying the small parameter \eps
\newif\ifmadec     \madectrue
\newif\ifslowfast  \slowfasttrue

% Part switch for the Barrat/simplicial-contagion project.  Defaults to
% true (the hypernetwork-synchrony part is included); the journal-paper
% driver says \hnetpartfalse to drop it and every cross-reference into it.
\newif\ifhnetpart  \hnetparttrue

%% ========================================
%% PACKAGES
%% ========================================
\RequirePackage{etoolbox}
\usepackage{makecell}
\usepackage[novbox]{pdfsync}% obsolete; SyncTeX is built into modern engines
\usepackage{amssymb, amsfonts,mathtools,bbm,eurosym,graphicx,latexsym, amsmath,amsthm,times,nccmath,url}
\usepackage{listings,verbatim,cases,tabularx}
% hyperref is deliberately NOT in the line above.  Journal classes often
% load it themselves, with a long option list, from a begindocument hook
% (mdpi.cls does: \DeclareHookRule{begindocument}{hyperref}{before}{mdpi}).
% Loading it here first, without options, makes that later load raise
% "Option clash for package hyperref".  So load it only if nobody else
% has, or will:
%   - \@ifpackageloaded  catches a class that loaded it already;
%   - \hyperrefpackage    is mdpi.cls's own deferred loader;
%   - \AAdefNnohyperref   lets any other driver opt out by defining it
%                         before % Shared preamble for the GLVbook papers.  This is the single canonical copy;
% a driver one level below GLVbook reaches it as
%
%     \input{../basicFiles/preamble}
%
% What stays in the driver, because it differs from paper to paper:
% \documentclass, font choice, \title, \author, \date, and any macro genuinely
% local to one paper.  Everything else belongs here.  The file is
% self-contained: it does not read any other file.

\usepackage[utf8]{inputenc}
\usepackage[novbox]{pdfsync}   % required: SyncTeX alone does not work here

% Page geometry, set once here for every paper that loads this file.  A driver
% that needs different margins says \geometry{...}; a second
% \usepackage[...]{geometry} raises an option clash and must not be used.
\makeatletter
\@ifundefined{AAdefNnogeometry}{\usepackage[a4paper,margin=2cm]{geometry}}{}
\makeatother

% mathematics
\usepackage{amsmath,amssymb,amsthm,mathtools,bm,mathrsfs,bbm}
\usepackage{centernot}

% floats, tables, lists
\usepackage{graphicx,float}
\makeatletter
\@ifpackageloaded{subfig}{}{\usepackage{subcaption}}
\makeatother
\usepackage{booktabs,longtable,threeparttable,array,multirow}
\usepackage{enumitem}

% colour, code, algorithms
\usepackage{xcolor,listings,tcolorbox}
\makeatletter
\@ifundefined{c@algorithm}{\usepackage{algorithm}}{}
\makeatother
\usepackage{algorithmicx,algpseudocode}

% diagrams
\usepackage{tikz,tikz-cd,pgfplots}
\usetikzlibrary{graphs,decorations.pathmorphing,decorations.markings,
arrows.meta,automata,positioning}
\tikzcdset{every label/.append style = {font = \small}}
\pgfplotsset{compat=1.18}

% chemistry
\usepackage[version=4]{mhchem}

\usepackage{microtype}
\usepackage{lineno}
\usepackage{nomencl}
\makenomenclature

% hyperref first, then bookmark, then glossaries: glossaries must follow
% hyperref or the entry links are built against the wrong anchors.
\makeatletter
\@ifundefined{hyperrefpackage}{%
  \@ifundefined{AAdefNnohyperref}{\usepackage{hyperref}}{}%
}{}
\@ifundefined{hyperrefpackage}{\usepackage{bookmark}}{}
\makeatother
\usepackage[toc,sort=use]{glossaries}
\makeglossaries

% Book/standalone switch.  Material that the book already covers in its
% prelude, but that a standalone paper still has to state, is wrapped as
%     \ifbookbuild\else ... \fi
% Drivers default to standalone; bookM.tex says \bookbuildtrue.
\newif\ifbookbuild
\bookbuildfalse

% Part switches for the gradostat paper.  Both default to true; a driver that
% wants the short submission says \madecfalse and/or \slowfastfalse, and every
% cross-reference into the dropped part disappears with it.
%   \ifmadec     Part III, the Lagasquie--Madec model and its break-even geometry
%   \ifslowfast  Part IV, everything carrying the small parameter \eps
\newif\ifmadec     \madectrue
\newif\ifslowfast  \slowfasttrue

% Part switch for the Barrat/simplicial-contagion project.  Defaults to
% true (the hypernetwork-synchrony part is included); the journal-paper
% driver says \hnetpartfalse to drop it and every cross-reference into it.
\newif\ifhnetpart  \hnetparttrue

%% ========================================
%% PACKAGES
%% ========================================
\RequirePackage{etoolbox}
\usepackage{makecell}
\usepackage[novbox]{pdfsync}% obsolete; SyncTeX is built into modern engines
\usepackage{amssymb, amsfonts,mathtools,bbm,eurosym,graphicx,latexsym, amsmath,amsthm,times,nccmath,url}
\usepackage{listings,verbatim,cases,tabularx}
% hyperref is deliberately NOT in the line above.  Journal classes often
% load it themselves, with a long option list, from a begindocument hook
% (mdpi.cls does: \DeclareHookRule{begindocument}{hyperref}{before}{mdpi}).
% Loading it here first, without options, makes that later load raise
% "Option clash for package hyperref".  So load it only if nobody else
% has, or will:
%   - \@ifpackageloaded  catches a class that loaded it already;
%   - \hyperrefpackage    is mdpi.cls's own deferred loader;
%   - \AAdefNnohyperref   lets any other driver opt out by defining it
%                         before % Shared preamble for the GLVbook papers.  This is the single canonical copy;
% a driver one level below GLVbook reaches it as
%
%     \input{../basicFiles/preamble}
%
% What stays in the driver, because it differs from paper to paper:
% \documentclass, font choice, \title, \author, \date, and any macro genuinely
% local to one paper.  Everything else belongs here.  The file is
% self-contained: it does not read any other file.

\usepackage[utf8]{inputenc}
\usepackage[novbox]{pdfsync}   % required: SyncTeX alone does not work here

% Page geometry, set once here for every paper that loads this file.  A driver
% that needs different margins says \geometry{...}; a second
% \usepackage[...]{geometry} raises an option clash and must not be used.
\makeatletter
\@ifundefined{AAdefNnogeometry}{\usepackage[a4paper,margin=2cm]{geometry}}{}
\makeatother

% mathematics
\usepackage{amsmath,amssymb,amsthm,mathtools,bm,mathrsfs,bbm}
\usepackage{centernot}

% floats, tables, lists
\usepackage{graphicx,float}
\makeatletter
\@ifpackageloaded{subfig}{}{\usepackage{subcaption}}
\makeatother
\usepackage{booktabs,longtable,threeparttable,array,multirow}
\usepackage{enumitem}

% colour, code, algorithms
\usepackage{xcolor,listings,tcolorbox}
\makeatletter
\@ifundefined{c@algorithm}{\usepackage{algorithm}}{}
\makeatother
\usepackage{algorithmicx,algpseudocode}

% diagrams
\usepackage{tikz,tikz-cd,pgfplots}
\usetikzlibrary{graphs,decorations.pathmorphing,decorations.markings,
arrows.meta,automata,positioning}
\tikzcdset{every label/.append style = {font = \small}}
\pgfplotsset{compat=1.18}

% chemistry
\usepackage[version=4]{mhchem}

\usepackage{microtype}
\usepackage{lineno}
\usepackage{nomencl}
\makenomenclature

% hyperref first, then bookmark, then glossaries: glossaries must follow
% hyperref or the entry links are built against the wrong anchors.
\makeatletter
\@ifundefined{hyperrefpackage}{%
  \@ifundefined{AAdefNnohyperref}{\usepackage{hyperref}}{}%
}{}
\@ifundefined{hyperrefpackage}{\usepackage{bookmark}}{}
\makeatother
\usepackage[toc,sort=use]{glossaries}
\makeglossaries

% Book/standalone switch.  Material that the book already covers in its
% prelude, but that a standalone paper still has to state, is wrapped as
%     \ifbookbuild\else ... \fi
% Drivers default to standalone; bookM.tex says \bookbuildtrue.
\newif\ifbookbuild
\bookbuildfalse

% Part switches for the gradostat paper.  Both default to true; a driver that
% wants the short submission says \madecfalse and/or \slowfastfalse, and every
% cross-reference into the dropped part disappears with it.
%   \ifmadec     Part III, the Lagasquie--Madec model and its break-even geometry
%   \ifslowfast  Part IV, everything carrying the small parameter \eps
\newif\ifmadec     \madectrue
\newif\ifslowfast  \slowfasttrue

% Part switch for the Barrat/simplicial-contagion project.  Defaults to
% true (the hypernetwork-synchrony part is included); the journal-paper
% driver says \hnetpartfalse to drop it and every cross-reference into it.
\newif\ifhnetpart  \hnetparttrue

%% ========================================
%% PACKAGES
%% ========================================
\RequirePackage{etoolbox}
\usepackage{makecell}
\usepackage[novbox]{pdfsync}% obsolete; SyncTeX is built into modern engines
\usepackage{amssymb, amsfonts,mathtools,bbm,eurosym,graphicx,latexsym, amsmath,amsthm,times,nccmath,url}
\usepackage{listings,verbatim,cases,tabularx}
% hyperref is deliberately NOT in the line above.  Journal classes often
% load it themselves, with a long option list, from a begindocument hook
% (mdpi.cls does: \DeclareHookRule{begindocument}{hyperref}{before}{mdpi}).
% Loading it here first, without options, makes that later load raise
% "Option clash for package hyperref".  So load it only if nobody else
% has, or will:
%   - \@ifpackageloaded  catches a class that loaded it already;
%   - \hyperrefpackage    is mdpi.cls's own deferred loader;
%   - \AAdefNnohyperref   lets any other driver opt out by defining it
%                         before \input{preamble}.
\makeatletter
\@ifpackageloaded{hyperref}{}{%
  \@ifundefined{hyperrefpackage}{%
    \@ifundefined{AAdefNnohyperref}{\usepackage{hyperref}}{}%
  }{}%
}
\makeatother
% linktoc=all makes a table-of-contents page number a link as well as the
% entry text; hyperref's default, linktoc=section, links only the text.
% Deferred to begin-document and guarded: by the block above hyperref may
% be loaded later (mdpi's own hook) or not at all (\AAdefNnohyperref), and
% \hypersetup must exist before it is called.  A driver that wants the
% default back can say \hypersetup{linktoc=section} after \input{preamble}.
\makeatletter
\AtBeginDocument{\@ifundefined{hypersetup}{}{\hypersetup{linktoc=all}}}
\makeatother
\usepackage{graphicx,float,fancybox}
\usepackage{epsfig,epstopdf,tikz}
\usepackage[sans]{dsfont}
\usetikzlibrary{positioning,arrows,arrows.meta,calc}
% Configure listings for Mathematica code
\lstdefinestyle{mathematica}{
 language=Mathematica,
 basicstyle=\small\ttfamily,
 backgroundcolor=\color{gray!10},
 frame=single,
 breaklines=true,
 commentstyle=\color{green!60!black},
 keywordstyle=\color{blue},
 stringstyle=\color{red},
 showstringspaces=false,
 tabsize=2
}
\usepackage{tcolorbox}
% Colored boxes for key results
\newtcolorbox{keyresult}[1][]{
 colback=blue!5!white,
 colframe=blue!75!black,
 title=#1,
 fonttitle=\bfseries
}
\newtcolorbox{examplebox}[1][]{
 colback=green!5!white,
 colframe=green!75!black,
 title=#1,
 fonttitle=\bfseries
}


\makeatletter
\@ifundefined{theorem}   {\newtheorem{theorem}   {Theorem}}   {}
\@ifundefined{definition}{\newtheorem{definition}{Definition}} {}
\@ifundefined{open}      {\newtheorem{open}      {Problem}}   {}
\@ifundefined{example}   {\newtheorem{example}   {Example}}   {}
\makeatother

% COUNTER DEFINITIONS
\newcounter{tabcounter}
\newcounter{defcounter}
\newcounter{opcounter}
\newcounter{thmcounter}
\newcounter{excounter}
\newcounter{ascounter}
\newcounter{cjcounter}
\newcounter{lemcounter}
\newcounter{remcounter}
\newcounter{corcounter}
\newcounter{propcounter}
\newcounter{figcounter}

% UNIFIED GLOSSARY STYLE
\newglossarystyle{unifiedstyle}{%
  \setglossarystyle{list}%
  \renewcommand*{\glossarypreamble}{\gdef\glslastsec{}}%
  \renewcommand{\glossentry}[2]{%
    \edef\glscursec{\glsentryuseriii{##1}}%
    \ifx\glscursec\glslastsec\else
      \item[]\textbf{\S\,\glscursec\enspace\glsentryuseriv{##1}}%
      \global\let\glslastsec\glscursec
    \fi
    \item[\glsentryuserii{##1} \glsentryuseri{##1}:] ``\glossentryname{##1}'' \glsentrydesc{##1}%
  }%
}

\makeatletter
% First-word extractor (used by \beT, \beD and \beXa for the default key)
% \@firstwordof hello world \relax  ->  hello
\def\@firstwordof#1 #2\relax{#1}
% First three letters, used for the heading mnemonic of a theorem given no
% explicit key.  Padding with \@empty keeps it safe for first words shorter
% than three letters: \@empty vanishes under \edef, whereas \relax would
% survive into the label and is illegal inside \csname (it broke every
% two-letter title, e.g. "ME ...", with "Missing \endcsname inserted").
\def\@takethree#1#2#3#4\@nil{#1#2#3}

% EXAMPLE COMMAND -- same signature and behaviour as \beT and \beD.
% Usage: \beXa*[long title][short title][key]
%   star  = suppresses the printed "[key]" tag after the title (as for \beT*
%   and \beD*); the label and glossary entry are unaffected.
%   key = optional label key override; the label becomes e:KEY, and without a
%   key it is e: plus the first three letters of the title's first word.  A
%   mnemonic label lets an example be cited as "Example~\ref{e:KEY}" instead
%   of by section number, which is what the citing prose needs.
\NewDocumentCommand{\beXa}{s O{} O{} o}{%
  \def\@temparg{#2}%
  \ifx\@temparg\@empty
    \begin{example}%
  \else
    \let\@savedSecTitle\@currentsectionname%
    \refstepcounter{excounter}%
    \IfValueTF{#4}{%
      \edef\@exlblkey{#4}%
    }{%
      \edef\@exfirstword{\@firstwordof#2 \relax}%
      \edef\@exlblkey{%
        \expandafter\@takethree\@exfirstword\@empty\@empty\@empty\@nil}%
    }%
    \def\@tempshort{#3}%
    \ifx\@tempshort\@empty
      \expandafter\xdef\csname @xshort@\@exlblkey\endcsname{\@exlblkey}%
    \else
      \expandafter\def\csname @xshort@\@exlblkey\endcsname{#3}%
    \fi
    \IfBooleanTF{#1}{\let\@extag\@empty}{%
      \edef\@extag{\noexpand\ \noexpand\textup{[\noexpand\texttt{e:\@exlblkey}]}}}%
    \begin{example}[#2\ifshowkeys\@extag\fi]%
    \label{e:\@exlblkey}%
    % record the long title so \nref{e:KEY} prints it, as for theorems
    \protected@write\@auxout{}{%
      \string\@thmtitle{e:\@exlblkey}{\unexpanded{#2}}}%
    \expanded{%
      \noexpand\newglossaryentry{e\theexcounter}{%
        name={\unexpanded{#2}%
          \IfBooleanF{#1}{ [\noexpand\texttt{e:\@exlblkey}]}},%
        sort={\theexample},%
        description={p.~\noexpand\pageref{e:\@exlblkey}},%
        user1={\theexample},%
        user2={Example},%
        user3={\thesection},%
        user4={\expandafter\unexpanded\expandafter{\@savedSecTitle}}%
      }%
    }%
    \protected@write\@auxout{}{\string\glsliveentry{e\theexcounter}}%
  \fi
}
\newcommand{\eeXa}{\end{example}}

% DEFINITION COMMAND -- same signature and behaviour as \beT below.
% Usage: \beD[long title][short title][key]
%   long title  = definition heading (plain text only)
%   short title = shown in \dref{word}; defaults to the key
%   key         = optional label key override; the label becomes d:KEY
% Without a key the label is d: plus the first three letters of the title's
% first word, so it is a mnemonic rather than the definition's number: a
% numeric label would silently change whenever a definition is inserted.
\NewDocumentCommand{\beD}{s O{} O{} o}{%
  \def\@temparg{#2}%
  \ifx\@temparg\@empty
    \begin{definition}%
  \else
    \let\@savedSecTitle\@currentsectionname%
    \refstepcounter{defcounter}%
    \IfValueTF{#4}{%
      \edef\@deflblkey{#4}%
    }{%
      \edef\@deffirstword{\@firstwordof#2 \relax}%
      \edef\@deflblkey{%
        \expandafter\@takethree\@deffirstword\@empty\@empty\@empty\@nil}%
    }%
    \def\@tempshort{#3}%
    \ifx\@tempshort\@empty
      \expandafter\xdef\csname @dshort@\@deflblkey\endcsname{\@deflblkey}%
    \else
      \expandafter\def\csname @dshort@\@deflblkey\endcsname{#3}%
    \fi
    \IfBooleanTF{#1}{\let\@deftag\@empty}{%
      \edef\@deftag{\noexpand\ \noexpand\textup{[\noexpand\texttt{d:\@deflblkey}]}}}%
    \begin{definition}[#2\ifshowkeys\@deftag\fi]%
    \label{d:\@deflblkey}%
    % record the long title so \nref{d:KEY} prints it, as for theorems
    \protected@write\@auxout{}{%
      \string\@thmtitle{d:\@deflblkey}{\unexpanded{#2}}}%
    \expanded{%
      \noexpand\newglossaryentry{d\thedefcounter}{%
        name={\unexpanded{#2}%
          \IfBooleanF{#1}{ [\noexpand\texttt{d:\@deflblkey}]}},%
        sort={\thedefinition},%
        description={p.~\noexpand\pageref{d:\@deflblkey}},%
        user1={\thedefinition},%
        user2={Definition},%
        user3={\thesection},%
        user4={\expandafter\unexpanded\expandafter{\@savedSecTitle}}%
      }%
    }%
    \protected@write\@auxout{}{\string\glsliveentry{d\thedefcounter}}%
  \fi
}
\newcommand{\eeD}{\end{definition}}

% \showkeysfalse in a driver hides the [t:KEY] mnemonics in result headings
% (labels and glossary entries are unchanged); used for submission builds.
\newif\ifshowkeys \showkeystrue
% \commonredtrue in a driver that inputs basicFiles/commonRed: tool.tex then
% cites its reduction results instead of restating them.
\newif\ifcommonred \commonredfalse
% TheorEM COMMAND
% Usage: \beT[long title][short title][key]
%   long title  = theorem heading (plain text only)
%   short title = shown in \tref{word}; may contain math; defaults to first word of long title
%   key         = optional label key override; label becomes t:KEY (avoids first-word collisions)
% Label is t:WORD where WORD = first word of long title (or KEY if given); use \tref{WORD} to reference.
\NewDocumentCommand{\beT}{s O{} O{} o}{%
  \def\@temparg{#2}%
  \ifx\@temparg\@empty
    \begin{theorem}%
  \else
    \let\@savedSecTitle\@currentsectionname%
    \refstepcounter{thmcounter}%
    % \@thmtag is the mnemonic shown in the heading; every theorem gets one.
    % With an explicit key the key itself is shown.  Without one the key
    % defaults to the title's first word, and showing that whole word would
    % merely repeat the title ("Dissipativity [Dissipativity]"), so only its
    % first three letters are shown: "Dissipativity [Dis]".  The label is
    % unaffected -- it still uses the full \@thmlblkey.
    \IfValueTF{#4}{%
      \edef\@thmlblkey{#4}%
    }{%
      \edef\@thmfirstword{\@firstwordof#2 \relax}%
      \edef\@thmlblkey{%
        \expandafter\@takethree\@thmfirstword\@empty\@empty\@empty\@nil}%
    }%
    \IfBooleanTF{#1}{\let\@thmtag\@empty}{%
      \edef\@thmtag{\noexpand\ \noexpand\textup{[\noexpand\texttt{t:\@thmlblkey}]}}}%
    \def\@tempshort{#3}%
    \ifx\@tempshort\@empty
      \expandafter\xdef\csname @tshort@\@thmlblkey\endcsname{\@thmlblkey}%
    \else
      \expandafter\def\csname @tshort@\@thmlblkey\endcsname{#3}%
    \fi
    \begin{theorem}[#2\ifshowkeys\@thmtag\fi]%
    \label{t:\@thmlblkey}%
    % record the long title for \nref; via the .aux so forward refs resolve
    \protected@write\@auxout{}{%
      \string\@thmtitle{t:\@thmlblkey}{\unexpanded{#2}}}%
    \expanded{%
      \noexpand\newglossaryentry{t\thethmcounter}{%
        name={\unexpanded{#2}%
          \IfBooleanF{#1}{ [\noexpand\texttt{t:\@thmlblkey}]}},%
        % sort/user1 must be the number the document actually prints, i.e. the
        % theorem environment's counter, not thmcounter.  thmcounter advances
        % only on \beT with a title, so a bare \beT (or \beT with an
        % empty title) makes the two drift and the glossary cites a wrong
        % number.  thmcounter is still fine as the entry key: it is unique.
        sort={\thetheorem},%
        description={p.~\noexpand\pageref{t:\@thmlblkey}},%
        user1={\thetheorem},%
        user2={Theorem},%
        user3={\thesection},%
        user4={\expandafter\unexpanded\expandafter{\@savedSecTitle}}%
      }%
    }%
    \protected@write\@auxout{}{\string\glsliveentry{t\thethmcounter}}%
  \fi
}
\newcommand{\eeT}{\end{theorem}}

% CROSS-PART RELATION PARAGRAPH
% Usage: \xpart{text}  -- place immediately BEFORE \beT, never inside it
%   (anything inside \beT...\eeT would be swallowed into the glossary entry).
% The paragraph is typeset only when Part I is actually present, i.e. in the
% combined book build; in the standalone builds (GranadaDSABNS/A.tex,
% AAHRah-Submission/AAHRah.tex) the label part:glv is undefined and the whole
% paragraph disappears rather than leaving ?? cross-references. Same
% graceful-degradation principle as \GASref in GASt.tex. Resolves on the
% second LaTeX pass, since r@part:glv comes from the .aux file.
% (Already inside the \makeatletter region opened above, so no \makeatletter
% pair here -- adding one would close that region early.)
\newcommand{\xpart}[1]{\@ifundefined{r@part:glv}{}{\paragraph{Relation with Parts~I--II.}#1}}

% \bookref{label}{precise text}{fallback text}: prints the precise cross-reference
% when the label resolves (combined book build), and the fallback text otherwise
% (standalone paper builds, where book-only appendix labels are absent). Same
% graceful-degradation role as \GASref in GASt.tex, but usable in any file since it
% lives in this preamble (loaded by every driver) -- unlike \GASref, which is defined in a
% content file and is therefore undefined in standalone builds that omit that file.
\newcommand{\bookref}[3]{\@ifundefined{r@#1}{#3}{#2}}

% \PVpart and \RBpart: the two parts that a partial build may omit.
% They print "Part~N" in any build that contains the part, and a naming phrase
% in a build that does not, so an excerpt shows no ?? .  Each already supplies
% the word "Part", so write "\PVpart develops ...", not "Part~\PVpart".
\newcommand{\PVpart}{\bookref{part:aahrah}%
  {Part~\ref{part:aahrah}}{the Perron--Volterra part of the full book}}
\newcommand{\RBpart}{\bookref{part:recbrowne}%
  {Part~\ref{part:recbrowne}}{the immuno-virology part of the full book}}

% OPEN PROBLEM COMMANDS
% ASSUMPTION COMMANDS (glossaried)
\newcommand{\beAs}[1][]{%
  \def\@temparg{#1}%
  \ifx\@temparg\@empty
    \begin{assumption}%
  \else
    \let\@savedSecTitle\@currentsectionname%
    \refstepcounter{ascounter}%
    \begin{assumption}[#1]%
    \label{as:\theascounter}%
    \expanded{%
      \noexpand\newglossaryentry{as\theascounter}{%
        name={\unexpanded{#1}},%
        sort={\theascounter},%
        description={p.~\noexpand\pageref{as:\theascounter}},%
        user1={\theascounter},%
        user2={Assumption},%
        user3={\thesection},%
        user4={\expandafter\unexpanded\expandafter{\@savedSecTitle}}%
      }%
    }%
    \protected@write\@auxout{}{\string\glsliveentry{as\theascounter}}%
  \fi
}
\newcommand{\eeAs}{\end{assumption}}
% CONJECTURE COMMANDS (glossaried)
\newcommand{\beCj}[1][]{%
  \def\@temparg{#1}%
  \ifx\@temparg\@empty
    \begin{conjecture}%
  \else
    \let\@savedSecTitle\@currentsectionname%
    \refstepcounter{cjcounter}%
    \begin{conjecture}[#1]%
    \label{cj:\thecjcounter}%
    \expanded{%
      \noexpand\newglossaryentry{cj\thecjcounter}{%
        name={\unexpanded{#1}},%
        sort={\thecjcounter},%
        description={p.~\noexpand\pageref{cj:\thecjcounter}},%
        user1={\thecjcounter},%
        user2={Conjecture},%
        user3={\thesection},%
        user4={\expandafter\unexpanded\expandafter{\@savedSecTitle}}%
      }%
    }%
    \protected@write\@auxout{}{\string\glsliveentry{cj\thecjcounter}}%
  \fi
}
\newcommand{\eeCj}{\end{conjecture}}
% LEMMA COMMANDS
% ---------------------------------------------------------------------
% UNIFORM RESULT WRAPPERS -- \beL, \beC, \beP, \beR share \beT's signature
%
%   \beX[long title][short title][key]  ...  \eeX
%
%   long title  = heading (math allowed)
%   short title = shown in the cross-reference; defaults to the key
%   key         = label key.  With a key the label is X:KEY and the heading
%                 carries the mnemonic [key], exactly as for \beT.  Without
%                 one the label falls back to the running counter, which is
%                 what every pre-existing call site in the book relies on.
%
% A bare \beX -- arguments commented out -- opens the plain environment and
% creates NO label and NO glossary entry.  That is the supported way to keep
% an individual result out of the catalogue.
% ---------------------------------------------------------------------
% #1 wrapper name  #2 environment  #3 label prefix  #4 counter  #5 type name
\newcommand{\@mkresult}[5]{%
  % A leading * suppresses the [key] mnemonic in both the heading and the
  % glossary name, while still making the label X:KEY and the entry:
  %     \beX*[long][short][key]
  \expandafter\NewDocumentCommand\csname be#1\endcsname{s O{} O{} o}{%
    \def\@temparg{##2}%
    \ifx\@temparg\@empty
      \begin{#2}%
    \else
      \let\@savedSecTitle\@currentsectionname%
      \refstepcounter{#4}%
      % With a key: mnemonic label X:KEY and the [key] tag in the heading.
      % Without one: fall back to the running counter, exactly as these
      % wrappers behaved before, so every legacy one-argument call site in
      % the book keeps its label.  The counter is also csname-safe, which a
      % title's first word is not -- titles beginning with math would
      % otherwise build \csname \Omega...\endcsname and fail.
      \IfValueTF{##4}{%
        \edef\@relbl{##4}%
        \IfBooleanTF{##1}{\let\@retag\@empty}{%
          \edef\@retag{\noexpand\ \noexpand\textup{[\noexpand\texttt{#3:##4}]}}}%
      }{%
        \edef\@relbl{\arabic{#4}}%
        \let\@retag\@empty
      }%
      \def\@tempshort{##3}%
      \ifx\@tempshort\@empty
        \expandafter\xdef\csname @sshort@#3\@relbl\endcsname{\@relbl}%
      \else
        \expandafter\def\csname @sshort@#3\@relbl\endcsname{##3}%
      \fi
      \begin{#2}[##2\ifshowkeys\@retag\fi]%
      \edef\@tmplbl{#3:\@relbl}\expandafter\label\expandafter{\@tmplbl}%
      \expanded{%
        \noexpand\newglossaryentry{#3\csname the#4\endcsname}{%
          name={\unexpanded{##2}%
            \IfBooleanF{##1}{\IfValueT{##4}{ [\noexpand\texttt{#3:##4}]}}},%
          sort={\csname the#2\endcsname},%
          description={p.~\noexpand\pageref{#3:\@relbl}},%
          user1={\csname the#2\endcsname},%
          user2={#5},%
          user3={\thesection},%
          user4={\expandafter\unexpanded\expandafter{\@savedSecTitle}}%
        }%
      }%
      \protected@write\@auxout{}{\string\glsliveentry{#3\csname the#4\endcsname}}%
      \protected@write\@auxout{}{\string\@thmtitle{\@tmplbl}{\unexpanded{##2}}}%
    \fi
  }%
  \expandafter\newcommand\csname ee#1\endcsname{\end{#2}}%
}
\@mkresult{L}{lemma}{l}{lemcounter}{Lemma}
% REMARK / COROLLARY / PROPOSITION -- same signature as \beL above
\@mkresult{R}{remark}{r}{remcounter}{Remark}
\@mkresult{C}{corollary}{c}{corcounter}{Corollary}
\@mkresult{P}{proposition}{p}{propcounter}{Proposition}
% OPEN-PROBLEM COMMAND -- generated with the same signature as \beT:
%   \beO[long title][short title][key]   ...   \eeO
% label o:KEY, or o:<counter> when no key is given; a bare \beO makes
% neither a label nor a glossary entry.
\@mkresult{O}{open}{o}{opcounter}{Problem}


% Section-level title tracker: updated only by \section, not \subsection.
% \@currentlabelname is overwritten by every \refstepcounter (including
% \subsection), so we save it separately here.
\newcommand{\@currentsectionname}{}
\AtBeginDocument{%
  \let\@epidcrn@orig@refstepcounter\refstepcounter
  \renewcommand{\refstepcounter}[1]{%
    \@epidcrn@orig@refstepcounter{#1}%
    \ifstrequal{#1}{section}{\global\let\@currentsectionname\@currentlabelname}{}%
  }%
}

% TABLE COMMANDS (glossaried) -- consistent naming with \beT, \beD, etc.
\newcommand{\beTa}[1][]{%
  \def\@temparg{#1}%
  \ifx\@temparg\@empty
    \begin{table}[htbp]%
  \else
    \let\@savedSecTitle\@currentsectionname%
    \refstepcounter{tabcounter}%
    \begin{table}[htbp]%
    \label{A:\thetabcounter}%
    \expanded{%
      \noexpand\newglossaryentry{A\thetabcounter}{%
        name={\unexpanded{#1}},%
        sort={\thetabcounter},%
        description={p.~\noexpand\pageref{A:\thetabcounter}},%
        user1={\thetabcounter},%
        user2={Table},%
        user3={\thesection},%
        user4={\expandafter\unexpanded\expandafter{\@savedSecTitle}}%
      }%
    }%
    \protected@write\@auxout{}{\string\glsliveentry{A\thetabcounter}}%
  \fi
}
\newcommand{\eeTa}{\end{table}}

% FIGURE COMMANDS (glossaried) -- consistent with \beT, \beR, \beTa, etc.
% Usage:
%   \beF[Glossary title]...\caption{Visible caption}\eeF  -> glossaried figure
%   \beF...\caption{Visible caption}\eeF                  -> plain [H] figure, no glossary
\newcommand{\beF}[1][]{%
  \def\@temparg{#1}%
  \ifx\@temparg\@empty
    \begin{figure}[H]%
  \else
    \begin{figure}[H]%
    \let\@savedSecTitle\@currentsectionname%
    \refstepcounter{figcounter}%
    \label{fig:\thefigcounter}%
    \expanded{%
      \noexpand\newglossaryentry{Fig\thefigcounter}{%
        name={\unexpanded{#1}},%
        sort={\thefigcounter},%
        description={p.~\noexpand\pageref{fig:\thefigcounter}},%
        user1={\thefigcounter},%
        user2={Figure},%
        user3={\thesection},%
        user4={\expandafter\unexpanded\expandafter{\@savedSecTitle}}%
      }%
    }%
    \protected@write\@auxout{}{\string\glsliveentry{Fig\thefigcounter}}%
  \fi
}
\newcommand{\eeF}{\end{figure}}
\makeatother
% LIVE-ENTRY REGISTER
% \glsaddall prints every entry the glossaries package currently knows, and
% that set includes entries restored from \jobname.glsdefs.  The package
% rewrites .glsdefs from all *currently defined* entries, so anything it
% restored is written straight back out: a deleted section kept reappearing in
% the catalogue for ever, and the .glsdefs had to be erased by hand.  The
% wrappers therefore record every entry they actually create in the .aux, and
% \printboth adds only those.  The .aux is rebuilt from scratch on each run, so
% this is self-healing: a removed result is gone after one further run.
\makeatletter
\newcommand{\@glslive}{}
\newcommand{\@glsaddlive}[1]{\ifglsentryexists{#1}{\glsadd{#1}}{}}
\newcommand{\glsliveentry}[1]{\g@addto@macro\@glslive{\@glsaddlive{#1}}}
\makeatother
% PRINT COMMAND
% No \clearpage here -- the catalogue continues on the current page.  This has
% been re-added by mistake several times; leave it out.
\makeatletter
\newcommand{\printboth}{%
  \@glslive   % add exactly the entries this document created last run
  \printglossary[style=unifiedstyle, title={0 \hspace{.22em} Definitions, Theorems, Examples and Open Problems}]%
}
\makeatother
% Named theorem reference: \tref{word} -> Thm~N(word-or-short)
% Shows stored short title if available, otherwise shows the key itself.
\newcommand{\tref}[1]{%
  Thm~\ref{t:#1}(\ifcsname @tshort@#1\endcsname\csname @tshort@#1\endcsname\else#1\fi)%
}

% Full-title theorem reference: \nref{t:KEY} prints the number and then the
% theorem's long title, so one writes "Theorem~\nref{t:blocks}" and gets
% "Theorem 3 ``Block factorization and siphons''".  Takes the complete label, like
% \ref, not the bare key like \tref.
%
% \beT writes each title to the .aux (see \@thmtitle below), so \nref resolves
% forward references from the second LaTeX pass, exactly as \ref does.  If the
% title is not yet known -- first pass, or a label that is not a \beT theorem
% -- it degrades to the number rather than producing ??.
% These two use @ in their names, and this point in the file is outside the
% \makeatletter region opened earlier (it closes at \makeatother above), so
% the region has to be reopened here.
\makeatletter
\newcommand{\@thmtitle}[2]{\global\@namedef{@tlong@#1}{#2}}
\newcommand{\nref}[1]{%
  \ref{#1}%
  \ifcsname @tlong@#1\endcsname
    \ ``\csname @tlong@#1\endcsname''%
  \fi
}
% \nm{t:KEY}: the stored long title of any \beT/\beL/\beP/\beR/\beC/\beD/\beE result,
% verbatim and without quotes, so item names follow the macro's first argument.
\newcommand{\nm}[1]{%
  \ifcsname @tlong@#1\endcsname
    \csname @tlong@#1\endcsname%
  \else
    \ref{#1}%
  \fi
}
\makeatother




%% ========================================
%% Other THEOREM ENVIRONMENTS
%% ========================================
\makeatletter
\@ifundefined{lemma}      {\newtheorem{lemma}      {Lemma}}      {}
\@ifundefined{proposition}{\newtheorem{proposition}{Proposition}} {}
\@ifundefined{corollary}  {\newtheorem{corollary}  {Corollary}}  {}
\@ifundefined{conjecture} {\newtheorem{conjecture} {Conjecture}}  {}
\@ifundefined{remark}     {\newtheorem{remark}     {Remark}}     {}
\@ifundefined{question}   {\newtheorem{question}   {Question}}   {}
\@ifundefined{assumption} {\newtheorem{assumption} {Assumption}} {}
\@ifundefined{hypothesis} {\newtheorem{hypothesis} {Hypothesis}} {}
\@ifundefined{fact}       {\newtheorem{fact}       {Fact}}       {}
\makeatother

% \citer: cite the relay-graph threshold theorem of AAHrel. If a document has
% \input the chapter/file containing t:relayRF (so the label resolves), point
% straight at it; otherwise fall back to a bare citation of AAHrel. Shared here
% (not per-paper) so it resolves consistently across every book chapter.
% \citeR: same idea, for the first theorem of AAHRah (t:Pstar).
\makeatletter
\providecommand{\citer}{%
  \@ifundefined{r@t:relayRF}%
    {\cite{AAHrel}}%
    {Theorem~\ref{t:relayRF} of \cite{AAHrel}}%
}
\providecommand{\citeR}{%
  \@ifundefined{r@t:Pstar}%
    {\cite{AAHRah}}%
    {Theorem~\ref{t:Pstar} of \cite{AAHRah}}%
}
\makeatother


% Abbreviations
\def\beQ{\begin{question}}\def\eeQ{\end{question}}
\def\beA{\begin{assumption}}\def\eeA{\end{assumption}}


%% ========================================

%% ========================================
%% COLORS & CUSTOM COMMANDS
%% ========================================
\definecolor{funccolor}{RGB}{25,25,112}
\definecolor{desccolor}{RGB}{64,64,64}
\newcommand{\func}[1]{\texttt{\textcolor{funccolor}{#1}}}
\def\eq{=}
\def\mm{-}

%% ========================================
%% MATHCAL SHORTHANDS
%% ========================================
% Project-wide convention: \mX for a capital X is always \mathcal X.
% Defined for the whole alphabet so the convention never fails to expand.
\newcommand{\mA}{\mathcal A}\newcommand{\mB}{\mathcal B}
\newcommand{\mC}{\mathcal C}\newcommand{\mD}{\mathcal D}
\newcommand{\mE}{\mathcal E}\newcommand{\mF}{\mathcal F}
\newcommand{\mG}{\mathcal G}\newcommand{\mH}{\mathcal H}
\newcommand{\mI}{\mathcal I}\newcommand{\mJ}{\mathcal J}
\newcommand{\mK}{\mathcal K}\newcommand{\mL}{\mathcal L}
\newcommand{\mM}{\mathcal M}\newcommand{\mN}{\mathcal N}
\newcommand{\mO}{\mathcal O}\newcommand{\mP}{\mathcal P}
\newcommand{\mQ}{\mathcal Q}\newcommand{\mR}{\mathcal R}
\newcommand{\mS}{\mathcal S}\newcommand{\mT}{\mathcal T}
\newcommand{\mU}{\mathcal U}\newcommand{\mV}{\mathcal V}
\newcommand{\mW}{\mathcal W}\newcommand{\mX}{\mathcal X}
\newcommand{\mY}{\mathcal Y}\newcommand{\mZ}{\mathcal Z}

%% ========================================
%% BASIC MATH ENVIRONMENTS
%% ========================================
\def\bep{\begin{pmatrix}}\def\eep{\end{pmatrix}}
\def\bev{\begin{vmatrix}}\def\eev{\end{vmatrix}}
\def\bea{\begin{eqnarray*}}\def\eea{\end{eqnarray*}}
\def\bc{\begin{cases}}\def\ec{\end{cases}}
\def\BEN{\begin{enumerate}}\def\EEN{\end{enumerate}}
\def\BI{\begin{itemize}}\def\EI{\end{itemize}}

\newcommand{\be}[1]{\begin{equation}\label{#1}}
\newcommand{\ee}{\end{equation}}
\newcommand{\beq}{\begin{eqnarray}}
\def\eeq{\end{eqnarray}}
\def\eqr{\eqref}\def\fr{\frac}\def\lbl{\label}
\def\Lra{\Longrightarrow}\def\non{\nonumber}\def\cNGM{\cite{Diek,Van,Van08}}
\def\Eq{\Leftrightarrow}\def\b{\beta}\def\uph{\red{UP TO HERE}}
%% ========================================
%% PROBABILITY & STATISTICS
%% ========================================
\renewcommand{\P}{\ensuremath{\mathbb{P}}}  % Note: overrides paragraph symbol
\newcommand{\E}{\ensuremath{\mathbb{E}}}
\newcommand{\ind}[1]{\ensuremath{\mathbbm{1}_{\left\{#1\right\}}}}
\DeclareMathOperator{\Cov}{Cov}
\DeclareMathOperator{\TV}{TV}
\newcommand{\rank}{\operatorname{rank}}

%% ========================================
%% CRN OPERATORS (organizations/ERCs/buffering/siphons framework)
%% ========================================
\def\clos{\mathrm{clos}}
\def\Supply{\mathrm{Supply}}
\def\req{\mathrm{req}}
\def\React{\mathrm{React}}
\def\Prod{\mathrm{Prod}}
\def\Gen{\mathrm{Gen}}
\def\In{\mathrm{In}}
\def\Out{\mathrm{Out}}
\def\Act{\mathrm{Act}}
\def\Min{\mathrm{Min}}

%% ========================================
%% LINE BREAKING
%% ========================================
% Last-resort stretch, used only on paragraphs TeX cannot otherwise break:
% a long \texttt{} token, a comma-separated inline math list (commas are not
% breakpoints in math), or a run of \cite/\ref boxes.  Without it those set as
% overfull \hbox and print into the margin; with it they set slightly loose,
% which is the lesser evil.  Does not affect paragraphs that already break.
\emergencystretch=3em

%% ========================================
%% NUMBER SETS
%% ========================================
\newcommand{\R}{\mathbb{R}}  % Real numbers
\newcommand{\C}{\mathbb{C}}  % Complex numbers
\newcommand{\N}{\mathbb{N}}
\newcommand{\Z}{\mathbb{Z}}
% \providecommand, not \newcommand: some papers define these in their own
% preamble after \input{preamble}, and must not clash.
\providecommand{\Q}{\mathbb{Q}}
\providecommand{\res}{\operatorname{res}}   % resultant; used by epiRes.tex
\newcommand{\MN}{\mathbb{N}}
\newcommand{\MZ}{\mathbb{Z}}
\newcommand{\realpos}{\mathbb{R}_{> 0}}
\newcommand{\realnneg}{\mathbb{R}_{\geq 0}}
\newcommand{\naturalpos}{\mathbb{N}_{>0}}
\newcommand{\Rplus}{\mathbb{R}_+}
\newcommand{\Rp}{\mathbb{R}_{\ge 0}}
\def\Rnne{\mathbb{R}_{\ge 0}}
\def\Se{\bar{S}} \def\Ie{\bar{I}}
\newcommand{\Rpos}{\mathbb{R}_{>0}}
\newcommand{\dd}{\,\mathrm d}
\newcommand{\e}{\mathrm e}
%% ========================================
%% GREEK LETTERS (2+ character abbreviations) greek
%% ========================================
\def\al{\alpha}\def\ba{\bff \alpha}
\def\w{\omega}\def\Lgn{\ga+\mu+\de}\def\om{\omega}
\def\de{\delta}\def\De{\Delta}\def\bde{\bff \delta}
\def\ep{\epsilon}\def\eps{\varepsilon}
\def\f{\varphi}  % Kept as requested
\def\ga{\gamma}\def\Ga{\Gamma}\def\G{\Gamma}\def\bg{\bff \gamma}
\def\up{\upsilon}  % NOT the breve accent: no accent commands anywhere.
                   % Named \up, not \u, so the LaTeX breve stays reachable.
\def\la{\lambda}\def\La{\Lambda}
\def\si{\sigma}\def\ka{\kappa}
\def\th{\theta}\def\vt{\vartheta}
\def\ze{\zeta}\def\Gap{\Gamma^{+}} \def\Gam{\Gamma^{-}}


%% ========================================
%% SPECIAL SHORT COMMANDS (kept as requested)
%% ========================================
\def\I{\infty}
\def\T{\widetilde}
\def\H{\hat}
\newcommand{\relint}{\operatorname{relint}}

%% ========================================
%% BOLD/VECTOR NOTATION
%% ========================================
\newcommand{\bff}[1]{{\mbox{\boldmath$#1$}}}
\newcommand\X{\boldsymbol{X}}
\def\x{\boldsymbol{x}}\def\bx{\bff x}\def\xe{\bar{\x}}
\newcommand\y{\boldsymbol{y}}\def\by{\boldsymbol{y}}\def\ye{\bar{\y}}
\def\yb{\y}\def\xb{\x}
\def\xy{\x-\y}
\newcommand\br{\boldsymbol{r}}\def\re{\bar{\r}}
\newcommand\bd{\boldsymbol{d}}
\newcommand\bQ{\boldsymbol{Q}}
\newcommand{\Yb}{{\mathbf{Y}}}
\newcommand{\Ab}{\mathbf{A}}
\newcommand{\zb}{\mathbf{z}}\def\bz{\bff 0}
\newcommand{\Zb}{\mathbf{Z}}
\newcommand{\bb}{\mathbf{b}}
\def\vbe{\vec {\bff \beta}}
\def\vb{\vec {\bff b}}
\def\v1{\vec {\bff 1}}
\def\vn{\vec \de}
\def\vz{\vec 0}
\def\vi{\overset{\rightarrow}{\mathsf i}}
\def\vv{\overset{\rightarrow}{v}}
\def\vr{\; \vec {\mathsf r}}
\newcommand{\va}{\bff a}
\providecommand{\vect}[1]{\boldsymbol{#1}}   % bold vector (moved from Gradostat.tex)
\providecommand{\one}{\mathbf 1}             % all-ones vector (moved from AliSF)
\providecommand{\spr}{s}                     % spectral bound  (moved from AliSF)
\providecommand{\Ke}{K_{\rm eff}}            % effective carrying capacity
% Critical right/left eigenvectors in the Hopf projection method (AAHOSN/New.tex)
\providecommand{\bq}{\mathbf{q}}
\providecommand{\bp}{\mathbf{p}}


%% ========================================
%% MATHCAL & MATHSCR
%% ========================================
% Mathcal
\newcommand\mc[1]{\mathcal{#1}}
\def\mC{{\mathcal C}}\def\mD{\mathcal D}\def\mE{\mathcal{E}}
\def\mF{\mathcal F}\def\mG{{\mathcal G}}\def\mH{{\mathcal H}}
\def\mI{\mathcal I}\def\mL{\mathcal{L}}\def\mM{{\mathcal M}}
\def\mN{{\mathcal N}}\def\mP{\mathcal{P}}\def\mR{{\mathcal R}}
\def\mS{{\mathcal S}}\def\mT{\mathcal{T}}\def\mV{{\mathcal V}}
\def\mK{\mathcal K}\def\mA{{\mathcal A}}\def\mU{{\mathcal U}}
% Symmetric bilinear and trilinear forms of the Taylor expansion, used in the
% first Lyapunov coefficient (AAHOSN/New.tex)
\providecommand{\Bform}{\mathcal{B}}
\providecommand{\Cform}{\mathcal{C}}

% Blackboard bold.  Kept here, not in the drivers: a \providecommand in a
% driver preamble runs before this file and silently wins, so the same name
% ends up meaning different things in different builds.
\def\IC{\mathbb{C}}\def\IN{\mathbb{N}}\def\IQ{\mathbb{Q}}
\def\IR{\mathbb{R}}\def\IZ{\mathbb{Z}}\def\IG{\mathbb{G}}
\def\II{\mathbb{I}}

% Mathscr
\newcommand{\Rs}{{\mathscr R}}
\newcommand{\Ss}{{\mathscr S}}
\newcommand{\Ns}{{\mathscr N}}
\newcommand{\Cs}{\mathscr C}


%% ========================================
%% OPERATORS & FUNCTIONS
%% ========================================
\DeclareMathOperator{\tr}{tr}
\DeclareMathOperator{\diag}{Diag}
\DeclareMathOperator{\Diag}{Diag}
\DeclareMathOperator{\intr}{\mathrm{Int}}
\DeclareMathOperator{\Inv}{Inv}
\DeclareMathOperator{\spec}{spec}
\DeclareMathOperator{\AUC}{AUC}               % area under the curve
\newcommand{\ii}{\mathrm{i}}                  % imaginary unit
\newcommand\op{\operatorname}
\newcommand{\supp}{\op{supp}}
\newcommand{\diff}{\mathop{}\mathopen{}\mathrm{d}}
\newcommand{\pdd}[1]{\frac{\partial}{\partial{#1}}}
\DeclareMathOperator{\cone}{cone}
\DeclareMathOperator{\proj}{proj}

%% ========================================
%% BRACKETS & NORMS
%% ========================================
\providecommand{\pr}[1]{\left(#1\right)}
\providecommand{\pp}[1]{\left[#1\right]}
\newcommand\croc[1]{\left\langle #1\right\rangle}
\newcommand{\norm}[1]{\|#1\|}
\newcommand{\inner}[2]{\langle #1, #2 \rangle}
% Short synonym of \inner, used by AAHOSN/New.tex as $\ip{u}{v}=\bar u^\top v$
\providecommand{\ip}[2]{\langle #1 , #2 \rangle}


%% ========================================
%% FORMATTING & UTILITIES
%% ========================================
\newcommand*{\Scale}[2][4]{\scalebox{#1}{$#2$}}
\newcommand{\te}[1]{\text{#1}}
\newcommand\steq[1]{\stackrel{\text{\rm #1.}}{=}}
\newcommand\blue[1]{\textcolor{blue}{#1}}
\newcommand{\red}{\textcolor[rgb]{1.00,0.00,0.00}}
\newcommand{\green}{\textcolor[rgb]{0.00,1.00,0.00}}

% Figure command
\newcommand{\figu}[3]{%
\begin{figure}[H]
\centering
\includegraphics[scale=#3]{#1}
\caption{#2\label{f:#1}}
\end{figure}
}

% Footnote
\long\def\symbolfootnote[#1]#2{%
\begingroup
\def\thefootnote{\fnsymbol{footnote}}\footnote[#1]{#2}%
\endgroup}
\def\fn{\symbolfootnote}

\newcommand{\dsum}{\displaystyle\sum}

%% ========================================
%% CRN & EPIDEMIOLOGY SPECIFIC
%% ========================================
% CRN notation
\def\SS{\mathcal S}
\def\CC{\mathcal C}
\def\RR{\mathcal R}

% Reproduction numbers
\def\Rto{R_{1}^{E_2}}
\def\Rot{R_{2}^{E_1}}
\def\Roc{R_{1,c}}
\def\Rtc{R_{2,c}}
\def\rd{r_{0}}
\def\sd{s_{0}}
\def\dd{D_{0}}
\def\vd{v_{0}}
\def\zd{z_{0}}
\def\se{\sE}
\def\xd{x_{0}}
%% ========================================
%% SPECIAL VALUES & EQUILIBRIA
%% ========================================
\def\eE{\bar{e}}
\def\wE{\bar{w}}
\def\sE{s_{*}}
\def\Me{\T M_{ee}}
\def\Rnne{\mathbb{R}_{\ge 0}}

\newcommand{\ones}{\mathbf 1}


%% ========================================
%% ABBREVIATIONS (alphabetical by category)
%% ========================================

% -- General mathematical terms --
\def\adm{admissible}\def\admy{admissibility}
%\def\and{antisymmetric}% disabled 2026-09-30: it overrode LaTeX's author separator \and (printed "antisymmetric" between authors)
\def\ass{assumption}
\def\cadlag{c\`adl\`ag}
\def\ch{characteristic polynomial}
\def\chp{characteristic polynomial}
\def\coe{coefficient}
\def\comp{computation}
\def\con{condition}
\def\coo{coordinate}
\def\corr{corresponding}
\def\dec{decomposition}
\def\den{denominator}
\def\dete{deterministic}
\def\eig{eigenvalue}\def\eval{eigenvalue}
\def\eigv{eigenvector}\def\eiv{eigenvector}\def\evec{eigenvector}
\def\elt{element}
\def\enn{essentially nonnegative/positivity preserving}
\def\equ{equation}
\def\etal{et al.}
\def\fe{for example}
\def\fno{from now on}
\def\for{\forall}
\def\form{formula}
\def\Fr{Furthermore, }
\def\frt{furthermore}
\def\fun{function}
\def\gene{generalization}
\def\heu{heuristic}
\def\how{however}\def\How{However, }
\def\ie{i.e. }
\def\im{\item}
\def\isd{is defined by}
\def\itc{it is easy to check that}
\def\itf{it follows that}\def\Itf{It follows that}
\def\ith{it holds that}
\def\itm{it may be checked that}\def\Itm{It may be checked that}
\def\ito{it turns out that}\def\Ito{It turns out that}\def\Itou{It turns out useful to}
\def\ma{mass action}\def\MAR{mass action representation}
\def\mat{matrix}
\def\mbc{may be checked to be}
\def\mbw{may be written as}\def\mubw{must be written as}
\def\mpr{more precisely, }\def\Mp{More precisely, }
\def\nc{non-catalytic}
\def\nne{non-negative}\def\nny{nonnegativity}
\def\no{\nonumber}
\def\noni{non-infectious}
\def\nota{notation}
\def\np{\newpage}
\def\Nt{Note that}
\def\nt{non trivial}\def\nz{non zero}
\def\oth{otherwise}
\def\Oth{On the other hand}
\def\pa{parameter}\def\para{parameter}
\def\pol{polynomial}
\def\prob{probability}
\def\Prf{{\bf  Proof:}}\def\prf{{\bf  Proof:}}
\def\prop{proportional to}
\def\qu{\quad}
\def\QED{\hfill {$\square$}\goodbreak \medskip}
\def\rep{representation}
\def\resp{respectively}
\def\satd{satisfied}\def\satg{satisfying}\def\sats{satisfies}\def\saty{satisfy}
\def\ssec{\subsection}\def\sssec{\subsubsection}
\def\sef{see for example}
\def\sksy{skew-symmetric}
\def\sub{substrate}
\def\sys{system}
\def\var{variable}
\def\wk{well-known}
\def\wl{without loss of generality}\def\wlo{w.l.o.g.}
\def\wmc{We may check that}\def\Wmc{We may check that}
\def\wmr{we may rewrite}
\def\wms{we must show that}
\def\wrt{with respect to}

% -- CRN & Chemistry --
\newcommand\CRN{chemical reaction network}
\newcommand\CRNs{chemical reaction networks}
\newcommand\RN{reaction network}
\def\ACR{absolute concentration robustness}
\def\AMSD{acyclic siphon decomposition}
\def\DAE{differential algebraic equations}
\def\FHJ{Feinberg-Horn-Jackson}
\def\GMAK{generalized mass action kinetics}
\def\MM{Michaelis-Menten}
\def\sm{stoichiometric matrix}\def\Sm{stoichiometric matrix}
\def\WRZD{weakly reversible, zero deficiency}\def\aut{autocatalytic}

% -- Epidemiology --
\def\brn{basic reproduction number}\def\bzn{basic replacement number}
\def\com{compartment}
\def\dfe{disease free equilibrium}\def\DFE{disease free equilibrium}
\def\dwt{dwell times}\def\edt{expected dwell times}
\def\EM{epidemic model}
\def\endp{endemic fixed point}\def\enp{endemic point}\def\EE{endemic point}
\def\ME{mathematical epidemiology}
\def\MEE{eco-epidemiology}\def\EEP{eco-epidemiological}
\def\ece{eco-epidemiological model}
\def\mgm{multi group model}
\def\mtS{multi-strain model}\def\mSt{multi-strain model}
\def\mSi{minimal siphon}
\def\NGM{next generation matrix}
\def\nim{new infections matrix}
\def\qem{quadratic epidemiological model}
\def\repF{invasion functions}
\def\SPF{SIR-PH-FA model}
\def\sus{susceptible}

% -- Graph theory & Linear algebra --
\def\jac{Jacobi matrix}
\def\jin{Jacobian matrix of the invasion vector field with respect to the invasion variables}
\def\jrn{Jacobian matrix of the resident vector field with respect to the resident variables}
\def\jxn{Jacobian matrix of the $\x$ vector field with respect to the $\x$ variables}
\def\Lm{Laplacian matrix}
\def\MTT{matrix tree theorem}
\def\SGT{spectral graph theory}
\def\SMW{Sherman–Morrison–Woodbury}

% -- Stability & Dynamics --
\def\absC{absorbtion CTMC}
\def\BTb{Bogdanov-Takens bifurcation}
\def\CEP{competitive exclusion partition}
\def\CTMC{continuous time Markov chain}
\def\ds{differential system}
\def\fp{fixed point}\def\bfp{boundary fixed point}\def\fbp{fixed boundary point}
\def\fbl{fixed boundary line}\def\fif{forward invariant face}
\def\gas{global asymptotic stability}\def\GAS{globally asymptotically stable}
\def\GASS{\text{GASS}}
\def\Hb{Hopf bifurcation}
\def\Hur{Hurwitz}\def\MH{Metzler-Hurwitz}\def\RH{Routh-Hurwitz}
\def\IS{\text{IS}}
\def\LAS{locally asymptotically stable}
\def\LdS{Lyapunov diagonally stable}
\def\Lf{Lyapunov function}\def\LyV{Lyapunov Volterra type function}\def\RLF{robust Lyapunov function}
\def\LyM{Lyapunov Malkin stable}
\def\PF{Perron-Frobenius}
\def\spfp{strictly positive fixed point}
\def\sta{stability}
\def\STP{sharp threshold property}\def\STT{sharp threshold theorem}

% -- Models (with \varphi explicitly instead of \f) --
\def\AB{$(A,B)$ Arino-Brauer epidemic models}
\def\ABa{$(A,\bbe,\ba,\varphi)$ balanced bilinear model}
\def\ABP{%$(A,\bbe,P,\varphi)$
\\ balanced bilinear model}
\def\AMG{Anderson May Gupta}
\newcommand\BIN{biological interaction network}
\def\Dt{Descartes type}
\def\gLV{generalized Lotka-Volterra model}
\def\IG{\text{invasion graph}}\def\IGMS{interaction graph on minimal siphons}
\def\LCP{\text{linear complementarity problem}}
\def\LV{Lotka-Volterra}\def\LVK{Lotka-Volterra-Kolmogorov}

% -- Splitting & Structure --
\def\reg{regular splitting}\def\regS{{\bf regular splitting}}
\def\RUR{rational univariate representation}
\def\bds{block-diagonal structure}\def\NGMs{next generation matrices}

% -- Miscellaneous --
\def\Ma{{Mathematica}}
\def\Mm{Max–Min}
\def\num{numerator}
\def\rap{rational approximation} \def\DSR{directed species-reaction bipartite graph}\def\PCP{polynomial complementarity problem}\def\rat{rational substitution}
\def\QSS{quasi steady state} \def\RHS{right hand side} \def\LHS{left hand side}
\def\OSN{online social network}\def\VdD{Van den Driessche}
\def\wli{$\omega$-limit} \def\rk{rich kinetics}\def\sJ{symbolic Jacobian}
\def\cV{\cite{Vas,VasSta,BSV,GSV}} \def\spf{semi-parametric semi-functional}
\def\CRW{Capasso-Ruan-Wang} \def\ABP{$(A,\bbe,P,\f)$ balanced bilinear model}
\def\ABa{%$(A,\bbe,\ba,\f)$ \\
rank-one balanced bilinear model}
\def\bbe{B%{\bff \beta}
} \def\vbe{\vec {\bff \beta}
}
\def\vb{\vec {\bff b}}
\newcommand{\Kh}{\hat{K}}          % K with hat, first mention only
\newcommand{\wS}{w^{\sigma}}       % right Perron J^\perp, full block
\newcommand{\piS}{\vec \pi^{\sigma}}    % left Perron J^\perp, full block
\newcommand{\ws}{w^{\sigma}}       % right Perron J^\perp, sub-block
\newcommand{\pis}{\pi^{\sigma}}    % left Perron J^\perp, sub-block
\newcommand{\hwS}{\hat{w}^{\sigma}}   % right Perron K, full block
\newcommand{\hpiS}{\hat{\pi}^{\sigma}} % left Perron K, full block
\newcommand{\hws}{\hat{w}^{\sigma}}   % right Perron K, sub-block
\newcommand{\hpis}{\hat{\pi}^{\sigma}} % left Perron K, sub-block
\newcommand{\twS}{\tilde{w}^{\sigma}}  % right Perron \tilde K, full block
\newcommand{\tpiS}{\tilde{\pi}^{\sigma}} % left Perron \tilde K
\newcommand{\tws}{\tilde{w}^{\sigma}}
\newcommand{\tpis}{\tilde{\pi}^{\sigma}}
\newcommand{\Mat}[2]{\mathbb{R}_{\ge 0}^{\,#1\times #2}}
\def\tJ{J_\sigma^\perp} \def\lPe{left Perron eigenvector}
% \G is \Gamma (defined at :795).  The Volterra function keeps its own name,
% so that \G no longer has two conflicting meanings in this file.
\def\GV{G}  % Volterra G(u)=u-1-\ln u
\def\PVL{Perron-Volterra Lyapunov minimal siphon ansatz}
\newcommand{\istar}[1]{i_{#1}^{*}}
%\newcommand{\sE}{s_E}
\newcommand{\vE}{v_E}
\def\PV{Perron-Volterra candidate Lyapunov function}
\def\GLV{generalized Lotka-Volterra}
\def\MERN{mathematical epidemiology reaction networks}
\def\CMM{contraction metric method}
\newcommand{\sign}{\operatorname{sign}}
\newcommand{\Var}{\operatorname{Var}} 

% ======================================================================
% IDEAS SHARED BETWEEN THE ABSTRACT AND THE MAIN RESULTS SECTION
%
% Write each idea ONCE, in the abstract; the Main Results section recalls it,
% so the two can never drift apart.
%
%   in the abstract        \ideaitem{1}{The first-order reaction terms ...}
%   in Main Results        \ideahead{1}
%
% \ideaitem stores the text globally and typesets
%   \item \hyperlink{idea:1}{<text>}
% \ideahead typesets
%   \item\hypertarget{idea:1}{}\textbf{<the same text>}
%
% The abstract must come before the Main Results input, which is the normal
% order. If a heading has to read differently from the abstract sentence (for
% instance a long abstract item summarised in one line), override it with the
% optional argument: \ideahead[short heading]{6}.
% ======================================================================
\makeatletter
\newcommand{\ideaitem}[2]{%
  \expandafter\gdef\csname @idea@#1\endcsname{#2}%
  \item\hyperlink{idea:#1}{#2}}
\newcommand{\ideahead}[2][]{%
  \item\hypertarget{idea:#2}{}%
  \def\@ideatmp{#1}%
  \ifx\@ideatmp\@empty
    \@ifundefined{@idea@#2}{%
      \PackageWarning{preamble}{Idea #2 is used in Main Results but was never
        defined by \string\ideaitem\space in the abstract}%
      \textbf{[idea #2 undefined]}%
    }{\textbf{\begingroup\def\ideaterm##1##2{##2}\def\idearest##1{}\@nameuse{@idea@#2}\endgroup}}%
  \else
    \textbf{#1}%
  \fi}
% \ideaterm{<hypertarget>}{<words>}: inside an abstract idea the words become an
% invisible link (the document sets \hypersetup{hidelinks}) jumping to the place
% in Main Results that defines them; inside the Main Results heading itself,
% where the target sits a few lines below, \ideahead strips the link and prints
% the words plainly. Mark the destination with \hypertarget{<name>}{} on the line
% before the display, which adds no equation number and shows nothing.
\newcommand{\ideaterm}[2]{\hyperlink{#1}{#2}}
% \idearest{<text>}: the part of an abstract idea after its first sentence; printed in
% the abstract, omitted from the Main Results heading.
\newcommand{\idearest}[1]{#1}
\makeatother

% ======================================================================
% SHORT-PAPER SWITCH
% \shorttrue in a master that assembles only part of the companion paper
% (Barrat/SipArxiv.tex). \nsh{...} drops text that cites results living in the
% omitted sections, so the short build has no dangling references while the
% companion paper keeps the text verbatim.
% ======================================================================
\newif\ifshort \shortfalse
\newcommand{\nsh}[1]{\ifshort\else#1\fi}
% \shortalt{<short build>}{<companion paper>}: use when a sentence needs
% different wording in the two builds, not merely a deletion.
\newcommand{\shortalt}[2]{\ifshort#1\else#2\fi}
% \Jtr, \JtrS{\Sigma}: transverse Jacobian. J^\perp here; SipArxiv.tex redefines them as J_{\mathrm{tra}}, J_{\mathrm{tra},\Sigma}
\newcommand{\Jtr}{J^\perp}
\newcommand{\JtrS}[1]{J_{#1}^\perp}
.
\makeatletter
\@ifpackageloaded{hyperref}{}{%
  \@ifundefined{hyperrefpackage}{%
    \@ifundefined{AAdefNnohyperref}{\usepackage{hyperref}}{}%
  }{}%
}
\makeatother
% linktoc=all makes a table-of-contents page number a link as well as the
% entry text; hyperref's default, linktoc=section, links only the text.
% Deferred to begin-document and guarded: by the block above hyperref may
% be loaded later (mdpi's own hook) or not at all (\AAdefNnohyperref), and
% \hypersetup must exist before it is called.  A driver that wants the
% default back can say \hypersetup{linktoc=section} after % Shared preamble for the GLVbook papers.  This is the single canonical copy;
% a driver one level below GLVbook reaches it as
%
%     \input{../basicFiles/preamble}
%
% What stays in the driver, because it differs from paper to paper:
% \documentclass, font choice, \title, \author, \date, and any macro genuinely
% local to one paper.  Everything else belongs here.  The file is
% self-contained: it does not read any other file.

\usepackage[utf8]{inputenc}
\usepackage[novbox]{pdfsync}   % required: SyncTeX alone does not work here

% Page geometry, set once here for every paper that loads this file.  A driver
% that needs different margins says \geometry{...}; a second
% \usepackage[...]{geometry} raises an option clash and must not be used.
\makeatletter
\@ifundefined{AAdefNnogeometry}{\usepackage[a4paper,margin=2cm]{geometry}}{}
\makeatother

% mathematics
\usepackage{amsmath,amssymb,amsthm,mathtools,bm,mathrsfs,bbm}
\usepackage{centernot}

% floats, tables, lists
\usepackage{graphicx,float}
\makeatletter
\@ifpackageloaded{subfig}{}{\usepackage{subcaption}}
\makeatother
\usepackage{booktabs,longtable,threeparttable,array,multirow}
\usepackage{enumitem}

% colour, code, algorithms
\usepackage{xcolor,listings,tcolorbox}
\makeatletter
\@ifundefined{c@algorithm}{\usepackage{algorithm}}{}
\makeatother
\usepackage{algorithmicx,algpseudocode}

% diagrams
\usepackage{tikz,tikz-cd,pgfplots}
\usetikzlibrary{graphs,decorations.pathmorphing,decorations.markings,
arrows.meta,automata,positioning}
\tikzcdset{every label/.append style = {font = \small}}
\pgfplotsset{compat=1.18}

% chemistry
\usepackage[version=4]{mhchem}

\usepackage{microtype}
\usepackage{lineno}
\usepackage{nomencl}
\makenomenclature

% hyperref first, then bookmark, then glossaries: glossaries must follow
% hyperref or the entry links are built against the wrong anchors.
\makeatletter
\@ifundefined{hyperrefpackage}{%
  \@ifundefined{AAdefNnohyperref}{\usepackage{hyperref}}{}%
}{}
\@ifundefined{hyperrefpackage}{\usepackage{bookmark}}{}
\makeatother
\usepackage[toc,sort=use]{glossaries}
\makeglossaries

% Book/standalone switch.  Material that the book already covers in its
% prelude, but that a standalone paper still has to state, is wrapped as
%     \ifbookbuild\else ... \fi
% Drivers default to standalone; bookM.tex says \bookbuildtrue.
\newif\ifbookbuild
\bookbuildfalse

% Part switches for the gradostat paper.  Both default to true; a driver that
% wants the short submission says \madecfalse and/or \slowfastfalse, and every
% cross-reference into the dropped part disappears with it.
%   \ifmadec     Part III, the Lagasquie--Madec model and its break-even geometry
%   \ifslowfast  Part IV, everything carrying the small parameter \eps
\newif\ifmadec     \madectrue
\newif\ifslowfast  \slowfasttrue

% Part switch for the Barrat/simplicial-contagion project.  Defaults to
% true (the hypernetwork-synchrony part is included); the journal-paper
% driver says \hnetpartfalse to drop it and every cross-reference into it.
\newif\ifhnetpart  \hnetparttrue

%% ========================================
%% PACKAGES
%% ========================================
\RequirePackage{etoolbox}
\usepackage{makecell}
\usepackage[novbox]{pdfsync}% obsolete; SyncTeX is built into modern engines
\usepackage{amssymb, amsfonts,mathtools,bbm,eurosym,graphicx,latexsym, amsmath,amsthm,times,nccmath,url}
\usepackage{listings,verbatim,cases,tabularx}
% hyperref is deliberately NOT in the line above.  Journal classes often
% load it themselves, with a long option list, from a begindocument hook
% (mdpi.cls does: \DeclareHookRule{begindocument}{hyperref}{before}{mdpi}).
% Loading it here first, without options, makes that later load raise
% "Option clash for package hyperref".  So load it only if nobody else
% has, or will:
%   - \@ifpackageloaded  catches a class that loaded it already;
%   - \hyperrefpackage    is mdpi.cls's own deferred loader;
%   - \AAdefNnohyperref   lets any other driver opt out by defining it
%                         before \input{preamble}.
\makeatletter
\@ifpackageloaded{hyperref}{}{%
  \@ifundefined{hyperrefpackage}{%
    \@ifundefined{AAdefNnohyperref}{\usepackage{hyperref}}{}%
  }{}%
}
\makeatother
% linktoc=all makes a table-of-contents page number a link as well as the
% entry text; hyperref's default, linktoc=section, links only the text.
% Deferred to begin-document and guarded: by the block above hyperref may
% be loaded later (mdpi's own hook) or not at all (\AAdefNnohyperref), and
% \hypersetup must exist before it is called.  A driver that wants the
% default back can say \hypersetup{linktoc=section} after \input{preamble}.
\makeatletter
\AtBeginDocument{\@ifundefined{hypersetup}{}{\hypersetup{linktoc=all}}}
\makeatother
\usepackage{graphicx,float,fancybox}
\usepackage{epsfig,epstopdf,tikz}
\usepackage[sans]{dsfont}
\usetikzlibrary{positioning,arrows,arrows.meta,calc}
% Configure listings for Mathematica code
\lstdefinestyle{mathematica}{
 language=Mathematica,
 basicstyle=\small\ttfamily,
 backgroundcolor=\color{gray!10},
 frame=single,
 breaklines=true,
 commentstyle=\color{green!60!black},
 keywordstyle=\color{blue},
 stringstyle=\color{red},
 showstringspaces=false,
 tabsize=2
}
\usepackage{tcolorbox}
% Colored boxes for key results
\newtcolorbox{keyresult}[1][]{
 colback=blue!5!white,
 colframe=blue!75!black,
 title=#1,
 fonttitle=\bfseries
}
\newtcolorbox{examplebox}[1][]{
 colback=green!5!white,
 colframe=green!75!black,
 title=#1,
 fonttitle=\bfseries
}


\makeatletter
\@ifundefined{theorem}   {\newtheorem{theorem}   {Theorem}}   {}
\@ifundefined{definition}{\newtheorem{definition}{Definition}} {}
\@ifundefined{open}      {\newtheorem{open}      {Problem}}   {}
\@ifundefined{example}   {\newtheorem{example}   {Example}}   {}
\makeatother

% COUNTER DEFINITIONS
\newcounter{tabcounter}
\newcounter{defcounter}
\newcounter{opcounter}
\newcounter{thmcounter}
\newcounter{excounter}
\newcounter{ascounter}
\newcounter{cjcounter}
\newcounter{lemcounter}
\newcounter{remcounter}
\newcounter{corcounter}
\newcounter{propcounter}
\newcounter{figcounter}

% UNIFIED GLOSSARY STYLE
\newglossarystyle{unifiedstyle}{%
  \setglossarystyle{list}%
  \renewcommand*{\glossarypreamble}{\gdef\glslastsec{}}%
  \renewcommand{\glossentry}[2]{%
    \edef\glscursec{\glsentryuseriii{##1}}%
    \ifx\glscursec\glslastsec\else
      \item[]\textbf{\S\,\glscursec\enspace\glsentryuseriv{##1}}%
      \global\let\glslastsec\glscursec
    \fi
    \item[\glsentryuserii{##1} \glsentryuseri{##1}:] ``\glossentryname{##1}'' \glsentrydesc{##1}%
  }%
}

\makeatletter
% First-word extractor (used by \beT, \beD and \beXa for the default key)
% \@firstwordof hello world \relax  ->  hello
\def\@firstwordof#1 #2\relax{#1}
% First three letters, used for the heading mnemonic of a theorem given no
% explicit key.  Padding with \@empty keeps it safe for first words shorter
% than three letters: \@empty vanishes under \edef, whereas \relax would
% survive into the label and is illegal inside \csname (it broke every
% two-letter title, e.g. "ME ...", with "Missing \endcsname inserted").
\def\@takethree#1#2#3#4\@nil{#1#2#3}

% EXAMPLE COMMAND -- same signature and behaviour as \beT and \beD.
% Usage: \beXa*[long title][short title][key]
%   star  = suppresses the printed "[key]" tag after the title (as for \beT*
%   and \beD*); the label and glossary entry are unaffected.
%   key = optional label key override; the label becomes e:KEY, and without a
%   key it is e: plus the first three letters of the title's first word.  A
%   mnemonic label lets an example be cited as "Example~\ref{e:KEY}" instead
%   of by section number, which is what the citing prose needs.
\NewDocumentCommand{\beXa}{s O{} O{} o}{%
  \def\@temparg{#2}%
  \ifx\@temparg\@empty
    \begin{example}%
  \else
    \let\@savedSecTitle\@currentsectionname%
    \refstepcounter{excounter}%
    \IfValueTF{#4}{%
      \edef\@exlblkey{#4}%
    }{%
      \edef\@exfirstword{\@firstwordof#2 \relax}%
      \edef\@exlblkey{%
        \expandafter\@takethree\@exfirstword\@empty\@empty\@empty\@nil}%
    }%
    \def\@tempshort{#3}%
    \ifx\@tempshort\@empty
      \expandafter\xdef\csname @xshort@\@exlblkey\endcsname{\@exlblkey}%
    \else
      \expandafter\def\csname @xshort@\@exlblkey\endcsname{#3}%
    \fi
    \IfBooleanTF{#1}{\let\@extag\@empty}{%
      \edef\@extag{\noexpand\ \noexpand\textup{[\noexpand\texttt{e:\@exlblkey}]}}}%
    \begin{example}[#2\ifshowkeys\@extag\fi]%
    \label{e:\@exlblkey}%
    % record the long title so \nref{e:KEY} prints it, as for theorems
    \protected@write\@auxout{}{%
      \string\@thmtitle{e:\@exlblkey}{\unexpanded{#2}}}%
    \expanded{%
      \noexpand\newglossaryentry{e\theexcounter}{%
        name={\unexpanded{#2}%
          \IfBooleanF{#1}{ [\noexpand\texttt{e:\@exlblkey}]}},%
        sort={\theexample},%
        description={p.~\noexpand\pageref{e:\@exlblkey}},%
        user1={\theexample},%
        user2={Example},%
        user3={\thesection},%
        user4={\expandafter\unexpanded\expandafter{\@savedSecTitle}}%
      }%
    }%
    \protected@write\@auxout{}{\string\glsliveentry{e\theexcounter}}%
  \fi
}
\newcommand{\eeXa}{\end{example}}

% DEFINITION COMMAND -- same signature and behaviour as \beT below.
% Usage: \beD[long title][short title][key]
%   long title  = definition heading (plain text only)
%   short title = shown in \dref{word}; defaults to the key
%   key         = optional label key override; the label becomes d:KEY
% Without a key the label is d: plus the first three letters of the title's
% first word, so it is a mnemonic rather than the definition's number: a
% numeric label would silently change whenever a definition is inserted.
\NewDocumentCommand{\beD}{s O{} O{} o}{%
  \def\@temparg{#2}%
  \ifx\@temparg\@empty
    \begin{definition}%
  \else
    \let\@savedSecTitle\@currentsectionname%
    \refstepcounter{defcounter}%
    \IfValueTF{#4}{%
      \edef\@deflblkey{#4}%
    }{%
      \edef\@deffirstword{\@firstwordof#2 \relax}%
      \edef\@deflblkey{%
        \expandafter\@takethree\@deffirstword\@empty\@empty\@empty\@nil}%
    }%
    \def\@tempshort{#3}%
    \ifx\@tempshort\@empty
      \expandafter\xdef\csname @dshort@\@deflblkey\endcsname{\@deflblkey}%
    \else
      \expandafter\def\csname @dshort@\@deflblkey\endcsname{#3}%
    \fi
    \IfBooleanTF{#1}{\let\@deftag\@empty}{%
      \edef\@deftag{\noexpand\ \noexpand\textup{[\noexpand\texttt{d:\@deflblkey}]}}}%
    \begin{definition}[#2\ifshowkeys\@deftag\fi]%
    \label{d:\@deflblkey}%
    % record the long title so \nref{d:KEY} prints it, as for theorems
    \protected@write\@auxout{}{%
      \string\@thmtitle{d:\@deflblkey}{\unexpanded{#2}}}%
    \expanded{%
      \noexpand\newglossaryentry{d\thedefcounter}{%
        name={\unexpanded{#2}%
          \IfBooleanF{#1}{ [\noexpand\texttt{d:\@deflblkey}]}},%
        sort={\thedefinition},%
        description={p.~\noexpand\pageref{d:\@deflblkey}},%
        user1={\thedefinition},%
        user2={Definition},%
        user3={\thesection},%
        user4={\expandafter\unexpanded\expandafter{\@savedSecTitle}}%
      }%
    }%
    \protected@write\@auxout{}{\string\glsliveentry{d\thedefcounter}}%
  \fi
}
\newcommand{\eeD}{\end{definition}}

% \showkeysfalse in a driver hides the [t:KEY] mnemonics in result headings
% (labels and glossary entries are unchanged); used for submission builds.
\newif\ifshowkeys \showkeystrue
% \commonredtrue in a driver that inputs basicFiles/commonRed: tool.tex then
% cites its reduction results instead of restating them.
\newif\ifcommonred \commonredfalse
% TheorEM COMMAND
% Usage: \beT[long title][short title][key]
%   long title  = theorem heading (plain text only)
%   short title = shown in \tref{word}; may contain math; defaults to first word of long title
%   key         = optional label key override; label becomes t:KEY (avoids first-word collisions)
% Label is t:WORD where WORD = first word of long title (or KEY if given); use \tref{WORD} to reference.
\NewDocumentCommand{\beT}{s O{} O{} o}{%
  \def\@temparg{#2}%
  \ifx\@temparg\@empty
    \begin{theorem}%
  \else
    \let\@savedSecTitle\@currentsectionname%
    \refstepcounter{thmcounter}%
    % \@thmtag is the mnemonic shown in the heading; every theorem gets one.
    % With an explicit key the key itself is shown.  Without one the key
    % defaults to the title's first word, and showing that whole word would
    % merely repeat the title ("Dissipativity [Dissipativity]"), so only its
    % first three letters are shown: "Dissipativity [Dis]".  The label is
    % unaffected -- it still uses the full \@thmlblkey.
    \IfValueTF{#4}{%
      \edef\@thmlblkey{#4}%
    }{%
      \edef\@thmfirstword{\@firstwordof#2 \relax}%
      \edef\@thmlblkey{%
        \expandafter\@takethree\@thmfirstword\@empty\@empty\@empty\@nil}%
    }%
    \IfBooleanTF{#1}{\let\@thmtag\@empty}{%
      \edef\@thmtag{\noexpand\ \noexpand\textup{[\noexpand\texttt{t:\@thmlblkey}]}}}%
    \def\@tempshort{#3}%
    \ifx\@tempshort\@empty
      \expandafter\xdef\csname @tshort@\@thmlblkey\endcsname{\@thmlblkey}%
    \else
      \expandafter\def\csname @tshort@\@thmlblkey\endcsname{#3}%
    \fi
    \begin{theorem}[#2\ifshowkeys\@thmtag\fi]%
    \label{t:\@thmlblkey}%
    % record the long title for \nref; via the .aux so forward refs resolve
    \protected@write\@auxout{}{%
      \string\@thmtitle{t:\@thmlblkey}{\unexpanded{#2}}}%
    \expanded{%
      \noexpand\newglossaryentry{t\thethmcounter}{%
        name={\unexpanded{#2}%
          \IfBooleanF{#1}{ [\noexpand\texttt{t:\@thmlblkey}]}},%
        % sort/user1 must be the number the document actually prints, i.e. the
        % theorem environment's counter, not thmcounter.  thmcounter advances
        % only on \beT with a title, so a bare \beT (or \beT with an
        % empty title) makes the two drift and the glossary cites a wrong
        % number.  thmcounter is still fine as the entry key: it is unique.
        sort={\thetheorem},%
        description={p.~\noexpand\pageref{t:\@thmlblkey}},%
        user1={\thetheorem},%
        user2={Theorem},%
        user3={\thesection},%
        user4={\expandafter\unexpanded\expandafter{\@savedSecTitle}}%
      }%
    }%
    \protected@write\@auxout{}{\string\glsliveentry{t\thethmcounter}}%
  \fi
}
\newcommand{\eeT}{\end{theorem}}

% CROSS-PART RELATION PARAGRAPH
% Usage: \xpart{text}  -- place immediately BEFORE \beT, never inside it
%   (anything inside \beT...\eeT would be swallowed into the glossary entry).
% The paragraph is typeset only when Part I is actually present, i.e. in the
% combined book build; in the standalone builds (GranadaDSABNS/A.tex,
% AAHRah-Submission/AAHRah.tex) the label part:glv is undefined and the whole
% paragraph disappears rather than leaving ?? cross-references. Same
% graceful-degradation principle as \GASref in GASt.tex. Resolves on the
% second LaTeX pass, since r@part:glv comes from the .aux file.
% (Already inside the \makeatletter region opened above, so no \makeatletter
% pair here -- adding one would close that region early.)
\newcommand{\xpart}[1]{\@ifundefined{r@part:glv}{}{\paragraph{Relation with Parts~I--II.}#1}}

% \bookref{label}{precise text}{fallback text}: prints the precise cross-reference
% when the label resolves (combined book build), and the fallback text otherwise
% (standalone paper builds, where book-only appendix labels are absent). Same
% graceful-degradation role as \GASref in GASt.tex, but usable in any file since it
% lives in this preamble (loaded by every driver) -- unlike \GASref, which is defined in a
% content file and is therefore undefined in standalone builds that omit that file.
\newcommand{\bookref}[3]{\@ifundefined{r@#1}{#3}{#2}}

% \PVpart and \RBpart: the two parts that a partial build may omit.
% They print "Part~N" in any build that contains the part, and a naming phrase
% in a build that does not, so an excerpt shows no ?? .  Each already supplies
% the word "Part", so write "\PVpart develops ...", not "Part~\PVpart".
\newcommand{\PVpart}{\bookref{part:aahrah}%
  {Part~\ref{part:aahrah}}{the Perron--Volterra part of the full book}}
\newcommand{\RBpart}{\bookref{part:recbrowne}%
  {Part~\ref{part:recbrowne}}{the immuno-virology part of the full book}}

% OPEN PROBLEM COMMANDS
% ASSUMPTION COMMANDS (glossaried)
\newcommand{\beAs}[1][]{%
  \def\@temparg{#1}%
  \ifx\@temparg\@empty
    \begin{assumption}%
  \else
    \let\@savedSecTitle\@currentsectionname%
    \refstepcounter{ascounter}%
    \begin{assumption}[#1]%
    \label{as:\theascounter}%
    \expanded{%
      \noexpand\newglossaryentry{as\theascounter}{%
        name={\unexpanded{#1}},%
        sort={\theascounter},%
        description={p.~\noexpand\pageref{as:\theascounter}},%
        user1={\theascounter},%
        user2={Assumption},%
        user3={\thesection},%
        user4={\expandafter\unexpanded\expandafter{\@savedSecTitle}}%
      }%
    }%
    \protected@write\@auxout{}{\string\glsliveentry{as\theascounter}}%
  \fi
}
\newcommand{\eeAs}{\end{assumption}}
% CONJECTURE COMMANDS (glossaried)
\newcommand{\beCj}[1][]{%
  \def\@temparg{#1}%
  \ifx\@temparg\@empty
    \begin{conjecture}%
  \else
    \let\@savedSecTitle\@currentsectionname%
    \refstepcounter{cjcounter}%
    \begin{conjecture}[#1]%
    \label{cj:\thecjcounter}%
    \expanded{%
      \noexpand\newglossaryentry{cj\thecjcounter}{%
        name={\unexpanded{#1}},%
        sort={\thecjcounter},%
        description={p.~\noexpand\pageref{cj:\thecjcounter}},%
        user1={\thecjcounter},%
        user2={Conjecture},%
        user3={\thesection},%
        user4={\expandafter\unexpanded\expandafter{\@savedSecTitle}}%
      }%
    }%
    \protected@write\@auxout{}{\string\glsliveentry{cj\thecjcounter}}%
  \fi
}
\newcommand{\eeCj}{\end{conjecture}}
% LEMMA COMMANDS
% ---------------------------------------------------------------------
% UNIFORM RESULT WRAPPERS -- \beL, \beC, \beP, \beR share \beT's signature
%
%   \beX[long title][short title][key]  ...  \eeX
%
%   long title  = heading (math allowed)
%   short title = shown in the cross-reference; defaults to the key
%   key         = label key.  With a key the label is X:KEY and the heading
%                 carries the mnemonic [key], exactly as for \beT.  Without
%                 one the label falls back to the running counter, which is
%                 what every pre-existing call site in the book relies on.
%
% A bare \beX -- arguments commented out -- opens the plain environment and
% creates NO label and NO glossary entry.  That is the supported way to keep
% an individual result out of the catalogue.
% ---------------------------------------------------------------------
% #1 wrapper name  #2 environment  #3 label prefix  #4 counter  #5 type name
\newcommand{\@mkresult}[5]{%
  % A leading * suppresses the [key] mnemonic in both the heading and the
  % glossary name, while still making the label X:KEY and the entry:
  %     \beX*[long][short][key]
  \expandafter\NewDocumentCommand\csname be#1\endcsname{s O{} O{} o}{%
    \def\@temparg{##2}%
    \ifx\@temparg\@empty
      \begin{#2}%
    \else
      \let\@savedSecTitle\@currentsectionname%
      \refstepcounter{#4}%
      % With a key: mnemonic label X:KEY and the [key] tag in the heading.
      % Without one: fall back to the running counter, exactly as these
      % wrappers behaved before, so every legacy one-argument call site in
      % the book keeps its label.  The counter is also csname-safe, which a
      % title's first word is not -- titles beginning with math would
      % otherwise build \csname \Omega...\endcsname and fail.
      \IfValueTF{##4}{%
        \edef\@relbl{##4}%
        \IfBooleanTF{##1}{\let\@retag\@empty}{%
          \edef\@retag{\noexpand\ \noexpand\textup{[\noexpand\texttt{#3:##4}]}}}%
      }{%
        \edef\@relbl{\arabic{#4}}%
        \let\@retag\@empty
      }%
      \def\@tempshort{##3}%
      \ifx\@tempshort\@empty
        \expandafter\xdef\csname @sshort@#3\@relbl\endcsname{\@relbl}%
      \else
        \expandafter\def\csname @sshort@#3\@relbl\endcsname{##3}%
      \fi
      \begin{#2}[##2\ifshowkeys\@retag\fi]%
      \edef\@tmplbl{#3:\@relbl}\expandafter\label\expandafter{\@tmplbl}%
      \expanded{%
        \noexpand\newglossaryentry{#3\csname the#4\endcsname}{%
          name={\unexpanded{##2}%
            \IfBooleanF{##1}{\IfValueT{##4}{ [\noexpand\texttt{#3:##4}]}}},%
          sort={\csname the#2\endcsname},%
          description={p.~\noexpand\pageref{#3:\@relbl}},%
          user1={\csname the#2\endcsname},%
          user2={#5},%
          user3={\thesection},%
          user4={\expandafter\unexpanded\expandafter{\@savedSecTitle}}%
        }%
      }%
      \protected@write\@auxout{}{\string\glsliveentry{#3\csname the#4\endcsname}}%
      \protected@write\@auxout{}{\string\@thmtitle{\@tmplbl}{\unexpanded{##2}}}%
    \fi
  }%
  \expandafter\newcommand\csname ee#1\endcsname{\end{#2}}%
}
\@mkresult{L}{lemma}{l}{lemcounter}{Lemma}
% REMARK / COROLLARY / PROPOSITION -- same signature as \beL above
\@mkresult{R}{remark}{r}{remcounter}{Remark}
\@mkresult{C}{corollary}{c}{corcounter}{Corollary}
\@mkresult{P}{proposition}{p}{propcounter}{Proposition}
% OPEN-PROBLEM COMMAND -- generated with the same signature as \beT:
%   \beO[long title][short title][key]   ...   \eeO
% label o:KEY, or o:<counter> when no key is given; a bare \beO makes
% neither a label nor a glossary entry.
\@mkresult{O}{open}{o}{opcounter}{Problem}


% Section-level title tracker: updated only by \section, not \subsection.
% \@currentlabelname is overwritten by every \refstepcounter (including
% \subsection), so we save it separately here.
\newcommand{\@currentsectionname}{}
\AtBeginDocument{%
  \let\@epidcrn@orig@refstepcounter\refstepcounter
  \renewcommand{\refstepcounter}[1]{%
    \@epidcrn@orig@refstepcounter{#1}%
    \ifstrequal{#1}{section}{\global\let\@currentsectionname\@currentlabelname}{}%
  }%
}

% TABLE COMMANDS (glossaried) -- consistent naming with \beT, \beD, etc.
\newcommand{\beTa}[1][]{%
  \def\@temparg{#1}%
  \ifx\@temparg\@empty
    \begin{table}[htbp]%
  \else
    \let\@savedSecTitle\@currentsectionname%
    \refstepcounter{tabcounter}%
    \begin{table}[htbp]%
    \label{A:\thetabcounter}%
    \expanded{%
      \noexpand\newglossaryentry{A\thetabcounter}{%
        name={\unexpanded{#1}},%
        sort={\thetabcounter},%
        description={p.~\noexpand\pageref{A:\thetabcounter}},%
        user1={\thetabcounter},%
        user2={Table},%
        user3={\thesection},%
        user4={\expandafter\unexpanded\expandafter{\@savedSecTitle}}%
      }%
    }%
    \protected@write\@auxout{}{\string\glsliveentry{A\thetabcounter}}%
  \fi
}
\newcommand{\eeTa}{\end{table}}

% FIGURE COMMANDS (glossaried) -- consistent with \beT, \beR, \beTa, etc.
% Usage:
%   \beF[Glossary title]...\caption{Visible caption}\eeF  -> glossaried figure
%   \beF...\caption{Visible caption}\eeF                  -> plain [H] figure, no glossary
\newcommand{\beF}[1][]{%
  \def\@temparg{#1}%
  \ifx\@temparg\@empty
    \begin{figure}[H]%
  \else
    \begin{figure}[H]%
    \let\@savedSecTitle\@currentsectionname%
    \refstepcounter{figcounter}%
    \label{fig:\thefigcounter}%
    \expanded{%
      \noexpand\newglossaryentry{Fig\thefigcounter}{%
        name={\unexpanded{#1}},%
        sort={\thefigcounter},%
        description={p.~\noexpand\pageref{fig:\thefigcounter}},%
        user1={\thefigcounter},%
        user2={Figure},%
        user3={\thesection},%
        user4={\expandafter\unexpanded\expandafter{\@savedSecTitle}}%
      }%
    }%
    \protected@write\@auxout{}{\string\glsliveentry{Fig\thefigcounter}}%
  \fi
}
\newcommand{\eeF}{\end{figure}}
\makeatother
% LIVE-ENTRY REGISTER
% \glsaddall prints every entry the glossaries package currently knows, and
% that set includes entries restored from \jobname.glsdefs.  The package
% rewrites .glsdefs from all *currently defined* entries, so anything it
% restored is written straight back out: a deleted section kept reappearing in
% the catalogue for ever, and the .glsdefs had to be erased by hand.  The
% wrappers therefore record every entry they actually create in the .aux, and
% \printboth adds only those.  The .aux is rebuilt from scratch on each run, so
% this is self-healing: a removed result is gone after one further run.
\makeatletter
\newcommand{\@glslive}{}
\newcommand{\@glsaddlive}[1]{\ifglsentryexists{#1}{\glsadd{#1}}{}}
\newcommand{\glsliveentry}[1]{\g@addto@macro\@glslive{\@glsaddlive{#1}}}
\makeatother
% PRINT COMMAND
% No \clearpage here -- the catalogue continues on the current page.  This has
% been re-added by mistake several times; leave it out.
\makeatletter
\newcommand{\printboth}{%
  \@glslive   % add exactly the entries this document created last run
  \printglossary[style=unifiedstyle, title={0 \hspace{.22em} Definitions, Theorems, Examples and Open Problems}]%
}
\makeatother
% Named theorem reference: \tref{word} -> Thm~N(word-or-short)
% Shows stored short title if available, otherwise shows the key itself.
\newcommand{\tref}[1]{%
  Thm~\ref{t:#1}(\ifcsname @tshort@#1\endcsname\csname @tshort@#1\endcsname\else#1\fi)%
}

% Full-title theorem reference: \nref{t:KEY} prints the number and then the
% theorem's long title, so one writes "Theorem~\nref{t:blocks}" and gets
% "Theorem 3 ``Block factorization and siphons''".  Takes the complete label, like
% \ref, not the bare key like \tref.
%
% \beT writes each title to the .aux (see \@thmtitle below), so \nref resolves
% forward references from the second LaTeX pass, exactly as \ref does.  If the
% title is not yet known -- first pass, or a label that is not a \beT theorem
% -- it degrades to the number rather than producing ??.
% These two use @ in their names, and this point in the file is outside the
% \makeatletter region opened earlier (it closes at \makeatother above), so
% the region has to be reopened here.
\makeatletter
\newcommand{\@thmtitle}[2]{\global\@namedef{@tlong@#1}{#2}}
\newcommand{\nref}[1]{%
  \ref{#1}%
  \ifcsname @tlong@#1\endcsname
    \ ``\csname @tlong@#1\endcsname''%
  \fi
}
% \nm{t:KEY}: the stored long title of any \beT/\beL/\beP/\beR/\beC/\beD/\beE result,
% verbatim and without quotes, so item names follow the macro's first argument.
\newcommand{\nm}[1]{%
  \ifcsname @tlong@#1\endcsname
    \csname @tlong@#1\endcsname%
  \else
    \ref{#1}%
  \fi
}
\makeatother




%% ========================================
%% Other THEOREM ENVIRONMENTS
%% ========================================
\makeatletter
\@ifundefined{lemma}      {\newtheorem{lemma}      {Lemma}}      {}
\@ifundefined{proposition}{\newtheorem{proposition}{Proposition}} {}
\@ifundefined{corollary}  {\newtheorem{corollary}  {Corollary}}  {}
\@ifundefined{conjecture} {\newtheorem{conjecture} {Conjecture}}  {}
\@ifundefined{remark}     {\newtheorem{remark}     {Remark}}     {}
\@ifundefined{question}   {\newtheorem{question}   {Question}}   {}
\@ifundefined{assumption} {\newtheorem{assumption} {Assumption}} {}
\@ifundefined{hypothesis} {\newtheorem{hypothesis} {Hypothesis}} {}
\@ifundefined{fact}       {\newtheorem{fact}       {Fact}}       {}
\makeatother

% \citer: cite the relay-graph threshold theorem of AAHrel. If a document has
% \input the chapter/file containing t:relayRF (so the label resolves), point
% straight at it; otherwise fall back to a bare citation of AAHrel. Shared here
% (not per-paper) so it resolves consistently across every book chapter.
% \citeR: same idea, for the first theorem of AAHRah (t:Pstar).
\makeatletter
\providecommand{\citer}{%
  \@ifundefined{r@t:relayRF}%
    {\cite{AAHrel}}%
    {Theorem~\ref{t:relayRF} of \cite{AAHrel}}%
}
\providecommand{\citeR}{%
  \@ifundefined{r@t:Pstar}%
    {\cite{AAHRah}}%
    {Theorem~\ref{t:Pstar} of \cite{AAHRah}}%
}
\makeatother


% Abbreviations
\def\beQ{\begin{question}}\def\eeQ{\end{question}}
\def\beA{\begin{assumption}}\def\eeA{\end{assumption}}


%% ========================================

%% ========================================
%% COLORS & CUSTOM COMMANDS
%% ========================================
\definecolor{funccolor}{RGB}{25,25,112}
\definecolor{desccolor}{RGB}{64,64,64}
\newcommand{\func}[1]{\texttt{\textcolor{funccolor}{#1}}}
\def\eq{=}
\def\mm{-}

%% ========================================
%% MATHCAL SHORTHANDS
%% ========================================
% Project-wide convention: \mX for a capital X is always \mathcal X.
% Defined for the whole alphabet so the convention never fails to expand.
\newcommand{\mA}{\mathcal A}\newcommand{\mB}{\mathcal B}
\newcommand{\mC}{\mathcal C}\newcommand{\mD}{\mathcal D}
\newcommand{\mE}{\mathcal E}\newcommand{\mF}{\mathcal F}
\newcommand{\mG}{\mathcal G}\newcommand{\mH}{\mathcal H}
\newcommand{\mI}{\mathcal I}\newcommand{\mJ}{\mathcal J}
\newcommand{\mK}{\mathcal K}\newcommand{\mL}{\mathcal L}
\newcommand{\mM}{\mathcal M}\newcommand{\mN}{\mathcal N}
\newcommand{\mO}{\mathcal O}\newcommand{\mP}{\mathcal P}
\newcommand{\mQ}{\mathcal Q}\newcommand{\mR}{\mathcal R}
\newcommand{\mS}{\mathcal S}\newcommand{\mT}{\mathcal T}
\newcommand{\mU}{\mathcal U}\newcommand{\mV}{\mathcal V}
\newcommand{\mW}{\mathcal W}\newcommand{\mX}{\mathcal X}
\newcommand{\mY}{\mathcal Y}\newcommand{\mZ}{\mathcal Z}

%% ========================================
%% BASIC MATH ENVIRONMENTS
%% ========================================
\def\bep{\begin{pmatrix}}\def\eep{\end{pmatrix}}
\def\bev{\begin{vmatrix}}\def\eev{\end{vmatrix}}
\def\bea{\begin{eqnarray*}}\def\eea{\end{eqnarray*}}
\def\bc{\begin{cases}}\def\ec{\end{cases}}
\def\BEN{\begin{enumerate}}\def\EEN{\end{enumerate}}
\def\BI{\begin{itemize}}\def\EI{\end{itemize}}

\newcommand{\be}[1]{\begin{equation}\label{#1}}
\newcommand{\ee}{\end{equation}}
\newcommand{\beq}{\begin{eqnarray}}
\def\eeq{\end{eqnarray}}
\def\eqr{\eqref}\def\fr{\frac}\def\lbl{\label}
\def\Lra{\Longrightarrow}\def\non{\nonumber}\def\cNGM{\cite{Diek,Van,Van08}}
\def\Eq{\Leftrightarrow}\def\b{\beta}\def\uph{\red{UP TO HERE}}
%% ========================================
%% PROBABILITY & STATISTICS
%% ========================================
\renewcommand{\P}{\ensuremath{\mathbb{P}}}  % Note: overrides paragraph symbol
\newcommand{\E}{\ensuremath{\mathbb{E}}}
\newcommand{\ind}[1]{\ensuremath{\mathbbm{1}_{\left\{#1\right\}}}}
\DeclareMathOperator{\Cov}{Cov}
\DeclareMathOperator{\TV}{TV}
\newcommand{\rank}{\operatorname{rank}}

%% ========================================
%% CRN OPERATORS (organizations/ERCs/buffering/siphons framework)
%% ========================================
\def\clos{\mathrm{clos}}
\def\Supply{\mathrm{Supply}}
\def\req{\mathrm{req}}
\def\React{\mathrm{React}}
\def\Prod{\mathrm{Prod}}
\def\Gen{\mathrm{Gen}}
\def\In{\mathrm{In}}
\def\Out{\mathrm{Out}}
\def\Act{\mathrm{Act}}
\def\Min{\mathrm{Min}}

%% ========================================
%% LINE BREAKING
%% ========================================
% Last-resort stretch, used only on paragraphs TeX cannot otherwise break:
% a long \texttt{} token, a comma-separated inline math list (commas are not
% breakpoints in math), or a run of \cite/\ref boxes.  Without it those set as
% overfull \hbox and print into the margin; with it they set slightly loose,
% which is the lesser evil.  Does not affect paragraphs that already break.
\emergencystretch=3em

%% ========================================
%% NUMBER SETS
%% ========================================
\newcommand{\R}{\mathbb{R}}  % Real numbers
\newcommand{\C}{\mathbb{C}}  % Complex numbers
\newcommand{\N}{\mathbb{N}}
\newcommand{\Z}{\mathbb{Z}}
% \providecommand, not \newcommand: some papers define these in their own
% preamble after \input{preamble}, and must not clash.
\providecommand{\Q}{\mathbb{Q}}
\providecommand{\res}{\operatorname{res}}   % resultant; used by epiRes.tex
\newcommand{\MN}{\mathbb{N}}
\newcommand{\MZ}{\mathbb{Z}}
\newcommand{\realpos}{\mathbb{R}_{> 0}}
\newcommand{\realnneg}{\mathbb{R}_{\geq 0}}
\newcommand{\naturalpos}{\mathbb{N}_{>0}}
\newcommand{\Rplus}{\mathbb{R}_+}
\newcommand{\Rp}{\mathbb{R}_{\ge 0}}
\def\Rnne{\mathbb{R}_{\ge 0}}
\def\Se{\bar{S}} \def\Ie{\bar{I}}
\newcommand{\Rpos}{\mathbb{R}_{>0}}
\newcommand{\dd}{\,\mathrm d}
\newcommand{\e}{\mathrm e}
%% ========================================
%% GREEK LETTERS (2+ character abbreviations) greek
%% ========================================
\def\al{\alpha}\def\ba{\bff \alpha}
\def\w{\omega}\def\Lgn{\ga+\mu+\de}\def\om{\omega}
\def\de{\delta}\def\De{\Delta}\def\bde{\bff \delta}
\def\ep{\epsilon}\def\eps{\varepsilon}
\def\f{\varphi}  % Kept as requested
\def\ga{\gamma}\def\Ga{\Gamma}\def\G{\Gamma}\def\bg{\bff \gamma}
\def\up{\upsilon}  % NOT the breve accent: no accent commands anywhere.
                   % Named \up, not \u, so the LaTeX breve stays reachable.
\def\la{\lambda}\def\La{\Lambda}
\def\si{\sigma}\def\ka{\kappa}
\def\th{\theta}\def\vt{\vartheta}
\def\ze{\zeta}\def\Gap{\Gamma^{+}} \def\Gam{\Gamma^{-}}


%% ========================================
%% SPECIAL SHORT COMMANDS (kept as requested)
%% ========================================
\def\I{\infty}
\def\T{\widetilde}
\def\H{\hat}
\newcommand{\relint}{\operatorname{relint}}

%% ========================================
%% BOLD/VECTOR NOTATION
%% ========================================
\newcommand{\bff}[1]{{\mbox{\boldmath$#1$}}}
\newcommand\X{\boldsymbol{X}}
\def\x{\boldsymbol{x}}\def\bx{\bff x}\def\xe{\bar{\x}}
\newcommand\y{\boldsymbol{y}}\def\by{\boldsymbol{y}}\def\ye{\bar{\y}}
\def\yb{\y}\def\xb{\x}
\def\xy{\x-\y}
\newcommand\br{\boldsymbol{r}}\def\re{\bar{\r}}
\newcommand\bd{\boldsymbol{d}}
\newcommand\bQ{\boldsymbol{Q}}
\newcommand{\Yb}{{\mathbf{Y}}}
\newcommand{\Ab}{\mathbf{A}}
\newcommand{\zb}{\mathbf{z}}\def\bz{\bff 0}
\newcommand{\Zb}{\mathbf{Z}}
\newcommand{\bb}{\mathbf{b}}
\def\vbe{\vec {\bff \beta}}
\def\vb{\vec {\bff b}}
\def\v1{\vec {\bff 1}}
\def\vn{\vec \de}
\def\vz{\vec 0}
\def\vi{\overset{\rightarrow}{\mathsf i}}
\def\vv{\overset{\rightarrow}{v}}
\def\vr{\; \vec {\mathsf r}}
\newcommand{\va}{\bff a}
\providecommand{\vect}[1]{\boldsymbol{#1}}   % bold vector (moved from Gradostat.tex)
\providecommand{\one}{\mathbf 1}             % all-ones vector (moved from AliSF)
\providecommand{\spr}{s}                     % spectral bound  (moved from AliSF)
\providecommand{\Ke}{K_{\rm eff}}            % effective carrying capacity
% Critical right/left eigenvectors in the Hopf projection method (AAHOSN/New.tex)
\providecommand{\bq}{\mathbf{q}}
\providecommand{\bp}{\mathbf{p}}


%% ========================================
%% MATHCAL & MATHSCR
%% ========================================
% Mathcal
\newcommand\mc[1]{\mathcal{#1}}
\def\mC{{\mathcal C}}\def\mD{\mathcal D}\def\mE{\mathcal{E}}
\def\mF{\mathcal F}\def\mG{{\mathcal G}}\def\mH{{\mathcal H}}
\def\mI{\mathcal I}\def\mL{\mathcal{L}}\def\mM{{\mathcal M}}
\def\mN{{\mathcal N}}\def\mP{\mathcal{P}}\def\mR{{\mathcal R}}
\def\mS{{\mathcal S}}\def\mT{\mathcal{T}}\def\mV{{\mathcal V}}
\def\mK{\mathcal K}\def\mA{{\mathcal A}}\def\mU{{\mathcal U}}
% Symmetric bilinear and trilinear forms of the Taylor expansion, used in the
% first Lyapunov coefficient (AAHOSN/New.tex)
\providecommand{\Bform}{\mathcal{B}}
\providecommand{\Cform}{\mathcal{C}}

% Blackboard bold.  Kept here, not in the drivers: a \providecommand in a
% driver preamble runs before this file and silently wins, so the same name
% ends up meaning different things in different builds.
\def\IC{\mathbb{C}}\def\IN{\mathbb{N}}\def\IQ{\mathbb{Q}}
\def\IR{\mathbb{R}}\def\IZ{\mathbb{Z}}\def\IG{\mathbb{G}}
\def\II{\mathbb{I}}

% Mathscr
\newcommand{\Rs}{{\mathscr R}}
\newcommand{\Ss}{{\mathscr S}}
\newcommand{\Ns}{{\mathscr N}}
\newcommand{\Cs}{\mathscr C}


%% ========================================
%% OPERATORS & FUNCTIONS
%% ========================================
\DeclareMathOperator{\tr}{tr}
\DeclareMathOperator{\diag}{Diag}
\DeclareMathOperator{\Diag}{Diag}
\DeclareMathOperator{\intr}{\mathrm{Int}}
\DeclareMathOperator{\Inv}{Inv}
\DeclareMathOperator{\spec}{spec}
\DeclareMathOperator{\AUC}{AUC}               % area under the curve
\newcommand{\ii}{\mathrm{i}}                  % imaginary unit
\newcommand\op{\operatorname}
\newcommand{\supp}{\op{supp}}
\newcommand{\diff}{\mathop{}\mathopen{}\mathrm{d}}
\newcommand{\pdd}[1]{\frac{\partial}{\partial{#1}}}
\DeclareMathOperator{\cone}{cone}
\DeclareMathOperator{\proj}{proj}

%% ========================================
%% BRACKETS & NORMS
%% ========================================
\providecommand{\pr}[1]{\left(#1\right)}
\providecommand{\pp}[1]{\left[#1\right]}
\newcommand\croc[1]{\left\langle #1\right\rangle}
\newcommand{\norm}[1]{\|#1\|}
\newcommand{\inner}[2]{\langle #1, #2 \rangle}
% Short synonym of \inner, used by AAHOSN/New.tex as $\ip{u}{v}=\bar u^\top v$
\providecommand{\ip}[2]{\langle #1 , #2 \rangle}


%% ========================================
%% FORMATTING & UTILITIES
%% ========================================
\newcommand*{\Scale}[2][4]{\scalebox{#1}{$#2$}}
\newcommand{\te}[1]{\text{#1}}
\newcommand\steq[1]{\stackrel{\text{\rm #1.}}{=}}
\newcommand\blue[1]{\textcolor{blue}{#1}}
\newcommand{\red}{\textcolor[rgb]{1.00,0.00,0.00}}
\newcommand{\green}{\textcolor[rgb]{0.00,1.00,0.00}}

% Figure command
\newcommand{\figu}[3]{%
\begin{figure}[H]
\centering
\includegraphics[scale=#3]{#1}
\caption{#2\label{f:#1}}
\end{figure}
}

% Footnote
\long\def\symbolfootnote[#1]#2{%
\begingroup
\def\thefootnote{\fnsymbol{footnote}}\footnote[#1]{#2}%
\endgroup}
\def\fn{\symbolfootnote}

\newcommand{\dsum}{\displaystyle\sum}

%% ========================================
%% CRN & EPIDEMIOLOGY SPECIFIC
%% ========================================
% CRN notation
\def\SS{\mathcal S}
\def\CC{\mathcal C}
\def\RR{\mathcal R}

% Reproduction numbers
\def\Rto{R_{1}^{E_2}}
\def\Rot{R_{2}^{E_1}}
\def\Roc{R_{1,c}}
\def\Rtc{R_{2,c}}
\def\rd{r_{0}}
\def\sd{s_{0}}
\def\dd{D_{0}}
\def\vd{v_{0}}
\def\zd{z_{0}}
\def\se{\sE}
\def\xd{x_{0}}
%% ========================================
%% SPECIAL VALUES & EQUILIBRIA
%% ========================================
\def\eE{\bar{e}}
\def\wE{\bar{w}}
\def\sE{s_{*}}
\def\Me{\T M_{ee}}
\def\Rnne{\mathbb{R}_{\ge 0}}

\newcommand{\ones}{\mathbf 1}


%% ========================================
%% ABBREVIATIONS (alphabetical by category)
%% ========================================

% -- General mathematical terms --
\def\adm{admissible}\def\admy{admissibility}
%\def\and{antisymmetric}% disabled 2026-09-30: it overrode LaTeX's author separator \and (printed "antisymmetric" between authors)
\def\ass{assumption}
\def\cadlag{c\`adl\`ag}
\def\ch{characteristic polynomial}
\def\chp{characteristic polynomial}
\def\coe{coefficient}
\def\comp{computation}
\def\con{condition}
\def\coo{coordinate}
\def\corr{corresponding}
\def\dec{decomposition}
\def\den{denominator}
\def\dete{deterministic}
\def\eig{eigenvalue}\def\eval{eigenvalue}
\def\eigv{eigenvector}\def\eiv{eigenvector}\def\evec{eigenvector}
\def\elt{element}
\def\enn{essentially nonnegative/positivity preserving}
\def\equ{equation}
\def\etal{et al.}
\def\fe{for example}
\def\fno{from now on}
\def\for{\forall}
\def\form{formula}
\def\Fr{Furthermore, }
\def\frt{furthermore}
\def\fun{function}
\def\gene{generalization}
\def\heu{heuristic}
\def\how{however}\def\How{However, }
\def\ie{i.e. }
\def\im{\item}
\def\isd{is defined by}
\def\itc{it is easy to check that}
\def\itf{it follows that}\def\Itf{It follows that}
\def\ith{it holds that}
\def\itm{it may be checked that}\def\Itm{It may be checked that}
\def\ito{it turns out that}\def\Ito{It turns out that}\def\Itou{It turns out useful to}
\def\ma{mass action}\def\MAR{mass action representation}
\def\mat{matrix}
\def\mbc{may be checked to be}
\def\mbw{may be written as}\def\mubw{must be written as}
\def\mpr{more precisely, }\def\Mp{More precisely, }
\def\nc{non-catalytic}
\def\nne{non-negative}\def\nny{nonnegativity}
\def\no{\nonumber}
\def\noni{non-infectious}
\def\nota{notation}
\def\np{\newpage}
\def\Nt{Note that}
\def\nt{non trivial}\def\nz{non zero}
\def\oth{otherwise}
\def\Oth{On the other hand}
\def\pa{parameter}\def\para{parameter}
\def\pol{polynomial}
\def\prob{probability}
\def\Prf{{\bf  Proof:}}\def\prf{{\bf  Proof:}}
\def\prop{proportional to}
\def\qu{\quad}
\def\QED{\hfill {$\square$}\goodbreak \medskip}
\def\rep{representation}
\def\resp{respectively}
\def\satd{satisfied}\def\satg{satisfying}\def\sats{satisfies}\def\saty{satisfy}
\def\ssec{\subsection}\def\sssec{\subsubsection}
\def\sef{see for example}
\def\sksy{skew-symmetric}
\def\sub{substrate}
\def\sys{system}
\def\var{variable}
\def\wk{well-known}
\def\wl{without loss of generality}\def\wlo{w.l.o.g.}
\def\wmc{We may check that}\def\Wmc{We may check that}
\def\wmr{we may rewrite}
\def\wms{we must show that}
\def\wrt{with respect to}

% -- CRN & Chemistry --
\newcommand\CRN{chemical reaction network}
\newcommand\CRNs{chemical reaction networks}
\newcommand\RN{reaction network}
\def\ACR{absolute concentration robustness}
\def\AMSD{acyclic siphon decomposition}
\def\DAE{differential algebraic equations}
\def\FHJ{Feinberg-Horn-Jackson}
\def\GMAK{generalized mass action kinetics}
\def\MM{Michaelis-Menten}
\def\sm{stoichiometric matrix}\def\Sm{stoichiometric matrix}
\def\WRZD{weakly reversible, zero deficiency}\def\aut{autocatalytic}

% -- Epidemiology --
\def\brn{basic reproduction number}\def\bzn{basic replacement number}
\def\com{compartment}
\def\dfe{disease free equilibrium}\def\DFE{disease free equilibrium}
\def\dwt{dwell times}\def\edt{expected dwell times}
\def\EM{epidemic model}
\def\endp{endemic fixed point}\def\enp{endemic point}\def\EE{endemic point}
\def\ME{mathematical epidemiology}
\def\MEE{eco-epidemiology}\def\EEP{eco-epidemiological}
\def\ece{eco-epidemiological model}
\def\mgm{multi group model}
\def\mtS{multi-strain model}\def\mSt{multi-strain model}
\def\mSi{minimal siphon}
\def\NGM{next generation matrix}
\def\nim{new infections matrix}
\def\qem{quadratic epidemiological model}
\def\repF{invasion functions}
\def\SPF{SIR-PH-FA model}
\def\sus{susceptible}

% -- Graph theory & Linear algebra --
\def\jac{Jacobi matrix}
\def\jin{Jacobian matrix of the invasion vector field with respect to the invasion variables}
\def\jrn{Jacobian matrix of the resident vector field with respect to the resident variables}
\def\jxn{Jacobian matrix of the $\x$ vector field with respect to the $\x$ variables}
\def\Lm{Laplacian matrix}
\def\MTT{matrix tree theorem}
\def\SGT{spectral graph theory}
\def\SMW{Sherman–Morrison–Woodbury}

% -- Stability & Dynamics --
\def\absC{absorbtion CTMC}
\def\BTb{Bogdanov-Takens bifurcation}
\def\CEP{competitive exclusion partition}
\def\CTMC{continuous time Markov chain}
\def\ds{differential system}
\def\fp{fixed point}\def\bfp{boundary fixed point}\def\fbp{fixed boundary point}
\def\fbl{fixed boundary line}\def\fif{forward invariant face}
\def\gas{global asymptotic stability}\def\GAS{globally asymptotically stable}
\def\GASS{\text{GASS}}
\def\Hb{Hopf bifurcation}
\def\Hur{Hurwitz}\def\MH{Metzler-Hurwitz}\def\RH{Routh-Hurwitz}
\def\IS{\text{IS}}
\def\LAS{locally asymptotically stable}
\def\LdS{Lyapunov diagonally stable}
\def\Lf{Lyapunov function}\def\LyV{Lyapunov Volterra type function}\def\RLF{robust Lyapunov function}
\def\LyM{Lyapunov Malkin stable}
\def\PF{Perron-Frobenius}
\def\spfp{strictly positive fixed point}
\def\sta{stability}
\def\STP{sharp threshold property}\def\STT{sharp threshold theorem}

% -- Models (with \varphi explicitly instead of \f) --
\def\AB{$(A,B)$ Arino-Brauer epidemic models}
\def\ABa{$(A,\bbe,\ba,\varphi)$ balanced bilinear model}
\def\ABP{%$(A,\bbe,P,\varphi)$
\\ balanced bilinear model}
\def\AMG{Anderson May Gupta}
\newcommand\BIN{biological interaction network}
\def\Dt{Descartes type}
\def\gLV{generalized Lotka-Volterra model}
\def\IG{\text{invasion graph}}\def\IGMS{interaction graph on minimal siphons}
\def\LCP{\text{linear complementarity problem}}
\def\LV{Lotka-Volterra}\def\LVK{Lotka-Volterra-Kolmogorov}

% -- Splitting & Structure --
\def\reg{regular splitting}\def\regS{{\bf regular splitting}}
\def\RUR{rational univariate representation}
\def\bds{block-diagonal structure}\def\NGMs{next generation matrices}

% -- Miscellaneous --
\def\Ma{{Mathematica}}
\def\Mm{Max–Min}
\def\num{numerator}
\def\rap{rational approximation} \def\DSR{directed species-reaction bipartite graph}\def\PCP{polynomial complementarity problem}\def\rat{rational substitution}
\def\QSS{quasi steady state} \def\RHS{right hand side} \def\LHS{left hand side}
\def\OSN{online social network}\def\VdD{Van den Driessche}
\def\wli{$\omega$-limit} \def\rk{rich kinetics}\def\sJ{symbolic Jacobian}
\def\cV{\cite{Vas,VasSta,BSV,GSV}} \def\spf{semi-parametric semi-functional}
\def\CRW{Capasso-Ruan-Wang} \def\ABP{$(A,\bbe,P,\f)$ balanced bilinear model}
\def\ABa{%$(A,\bbe,\ba,\f)$ \\
rank-one balanced bilinear model}
\def\bbe{B%{\bff \beta}
} \def\vbe{\vec {\bff \beta}
}
\def\vb{\vec {\bff b}}
\newcommand{\Kh}{\hat{K}}          % K with hat, first mention only
\newcommand{\wS}{w^{\sigma}}       % right Perron J^\perp, full block
\newcommand{\piS}{\vec \pi^{\sigma}}    % left Perron J^\perp, full block
\newcommand{\ws}{w^{\sigma}}       % right Perron J^\perp, sub-block
\newcommand{\pis}{\pi^{\sigma}}    % left Perron J^\perp, sub-block
\newcommand{\hwS}{\hat{w}^{\sigma}}   % right Perron K, full block
\newcommand{\hpiS}{\hat{\pi}^{\sigma}} % left Perron K, full block
\newcommand{\hws}{\hat{w}^{\sigma}}   % right Perron K, sub-block
\newcommand{\hpis}{\hat{\pi}^{\sigma}} % left Perron K, sub-block
\newcommand{\twS}{\tilde{w}^{\sigma}}  % right Perron \tilde K, full block
\newcommand{\tpiS}{\tilde{\pi}^{\sigma}} % left Perron \tilde K
\newcommand{\tws}{\tilde{w}^{\sigma}}
\newcommand{\tpis}{\tilde{\pi}^{\sigma}}
\newcommand{\Mat}[2]{\mathbb{R}_{\ge 0}^{\,#1\times #2}}
\def\tJ{J_\sigma^\perp} \def\lPe{left Perron eigenvector}
% \G is \Gamma (defined at :795).  The Volterra function keeps its own name,
% so that \G no longer has two conflicting meanings in this file.
\def\GV{G}  % Volterra G(u)=u-1-\ln u
\def\PVL{Perron-Volterra Lyapunov minimal siphon ansatz}
\newcommand{\istar}[1]{i_{#1}^{*}}
%\newcommand{\sE}{s_E}
\newcommand{\vE}{v_E}
\def\PV{Perron-Volterra candidate Lyapunov function}
\def\GLV{generalized Lotka-Volterra}
\def\MERN{mathematical epidemiology reaction networks}
\def\CMM{contraction metric method}
\newcommand{\sign}{\operatorname{sign}}
\newcommand{\Var}{\operatorname{Var}} 

% ======================================================================
% IDEAS SHARED BETWEEN THE ABSTRACT AND THE MAIN RESULTS SECTION
%
% Write each idea ONCE, in the abstract; the Main Results section recalls it,
% so the two can never drift apart.
%
%   in the abstract        \ideaitem{1}{The first-order reaction terms ...}
%   in Main Results        \ideahead{1}
%
% \ideaitem stores the text globally and typesets
%   \item \hyperlink{idea:1}{<text>}
% \ideahead typesets
%   \item\hypertarget{idea:1}{}\textbf{<the same text>}
%
% The abstract must come before the Main Results input, which is the normal
% order. If a heading has to read differently from the abstract sentence (for
% instance a long abstract item summarised in one line), override it with the
% optional argument: \ideahead[short heading]{6}.
% ======================================================================
\makeatletter
\newcommand{\ideaitem}[2]{%
  \expandafter\gdef\csname @idea@#1\endcsname{#2}%
  \item\hyperlink{idea:#1}{#2}}
\newcommand{\ideahead}[2][]{%
  \item\hypertarget{idea:#2}{}%
  \def\@ideatmp{#1}%
  \ifx\@ideatmp\@empty
    \@ifundefined{@idea@#2}{%
      \PackageWarning{preamble}{Idea #2 is used in Main Results but was never
        defined by \string\ideaitem\space in the abstract}%
      \textbf{[idea #2 undefined]}%
    }{\textbf{\begingroup\def\ideaterm##1##2{##2}\def\idearest##1{}\@nameuse{@idea@#2}\endgroup}}%
  \else
    \textbf{#1}%
  \fi}
% \ideaterm{<hypertarget>}{<words>}: inside an abstract idea the words become an
% invisible link (the document sets \hypersetup{hidelinks}) jumping to the place
% in Main Results that defines them; inside the Main Results heading itself,
% where the target sits a few lines below, \ideahead strips the link and prints
% the words plainly. Mark the destination with \hypertarget{<name>}{} on the line
% before the display, which adds no equation number and shows nothing.
\newcommand{\ideaterm}[2]{\hyperlink{#1}{#2}}
% \idearest{<text>}: the part of an abstract idea after its first sentence; printed in
% the abstract, omitted from the Main Results heading.
\newcommand{\idearest}[1]{#1}
\makeatother

% ======================================================================
% SHORT-PAPER SWITCH
% \shorttrue in a master that assembles only part of the companion paper
% (Barrat/SipArxiv.tex). \nsh{...} drops text that cites results living in the
% omitted sections, so the short build has no dangling references while the
% companion paper keeps the text verbatim.
% ======================================================================
\newif\ifshort \shortfalse
\newcommand{\nsh}[1]{\ifshort\else#1\fi}
% \shortalt{<short build>}{<companion paper>}: use when a sentence needs
% different wording in the two builds, not merely a deletion.
\newcommand{\shortalt}[2]{\ifshort#1\else#2\fi}
% \Jtr, \JtrS{\Sigma}: transverse Jacobian. J^\perp here; SipArxiv.tex redefines them as J_{\mathrm{tra}}, J_{\mathrm{tra},\Sigma}
\newcommand{\Jtr}{J^\perp}
\newcommand{\JtrS}[1]{J_{#1}^\perp}
.
\makeatletter
\AtBeginDocument{\@ifundefined{hypersetup}{}{\hypersetup{linktoc=all}}}
\makeatother
\usepackage{graphicx,float,fancybox}
\usepackage{epsfig,epstopdf,tikz}
\usepackage[sans]{dsfont}
\usetikzlibrary{positioning,arrows,arrows.meta,calc}
% Configure listings for Mathematica code
\lstdefinestyle{mathematica}{
 language=Mathematica,
 basicstyle=\small\ttfamily,
 backgroundcolor=\color{gray!10},
 frame=single,
 breaklines=true,
 commentstyle=\color{green!60!black},
 keywordstyle=\color{blue},
 stringstyle=\color{red},
 showstringspaces=false,
 tabsize=2
}
\usepackage{tcolorbox}
% Colored boxes for key results
\newtcolorbox{keyresult}[1][]{
 colback=blue!5!white,
 colframe=blue!75!black,
 title=#1,
 fonttitle=\bfseries
}
\newtcolorbox{examplebox}[1][]{
 colback=green!5!white,
 colframe=green!75!black,
 title=#1,
 fonttitle=\bfseries
}


\makeatletter
\@ifundefined{theorem}   {\newtheorem{theorem}   {Theorem}}   {}
\@ifundefined{definition}{\newtheorem{definition}{Definition}} {}
\@ifundefined{open}      {\newtheorem{open}      {Problem}}   {}
\@ifundefined{example}   {\newtheorem{example}   {Example}}   {}
\makeatother

% COUNTER DEFINITIONS
\newcounter{tabcounter}
\newcounter{defcounter}
\newcounter{opcounter}
\newcounter{thmcounter}
\newcounter{excounter}
\newcounter{ascounter}
\newcounter{cjcounter}
\newcounter{lemcounter}
\newcounter{remcounter}
\newcounter{corcounter}
\newcounter{propcounter}
\newcounter{figcounter}

% UNIFIED GLOSSARY STYLE
\newglossarystyle{unifiedstyle}{%
  \setglossarystyle{list}%
  \renewcommand*{\glossarypreamble}{\gdef\glslastsec{}}%
  \renewcommand{\glossentry}[2]{%
    \edef\glscursec{\glsentryuseriii{##1}}%
    \ifx\glscursec\glslastsec\else
      \item[]\textbf{\S\,\glscursec\enspace\glsentryuseriv{##1}}%
      \global\let\glslastsec\glscursec
    \fi
    \item[\glsentryuserii{##1} \glsentryuseri{##1}:] ``\glossentryname{##1}'' \glsentrydesc{##1}%
  }%
}

\makeatletter
% First-word extractor (used by \beT, \beD and \beXa for the default key)
% \@firstwordof hello world \relax  ->  hello
\def\@firstwordof#1 #2\relax{#1}
% First three letters, used for the heading mnemonic of a theorem given no
% explicit key.  Padding with \@empty keeps it safe for first words shorter
% than three letters: \@empty vanishes under \edef, whereas \relax would
% survive into the label and is illegal inside \csname (it broke every
% two-letter title, e.g. "ME ...", with "Missing \endcsname inserted").
\def\@takethree#1#2#3#4\@nil{#1#2#3}

% EXAMPLE COMMAND -- same signature and behaviour as \beT and \beD.
% Usage: \beXa*[long title][short title][key]
%   star  = suppresses the printed "[key]" tag after the title (as for \beT*
%   and \beD*); the label and glossary entry are unaffected.
%   key = optional label key override; the label becomes e:KEY, and without a
%   key it is e: plus the first three letters of the title's first word.  A
%   mnemonic label lets an example be cited as "Example~\ref{e:KEY}" instead
%   of by section number, which is what the citing prose needs.
\NewDocumentCommand{\beXa}{s O{} O{} o}{%
  \def\@temparg{#2}%
  \ifx\@temparg\@empty
    \begin{example}%
  \else
    \let\@savedSecTitle\@currentsectionname%
    \refstepcounter{excounter}%
    \IfValueTF{#4}{%
      \edef\@exlblkey{#4}%
    }{%
      \edef\@exfirstword{\@firstwordof#2 \relax}%
      \edef\@exlblkey{%
        \expandafter\@takethree\@exfirstword\@empty\@empty\@empty\@nil}%
    }%
    \def\@tempshort{#3}%
    \ifx\@tempshort\@empty
      \expandafter\xdef\csname @xshort@\@exlblkey\endcsname{\@exlblkey}%
    \else
      \expandafter\def\csname @xshort@\@exlblkey\endcsname{#3}%
    \fi
    \IfBooleanTF{#1}{\let\@extag\@empty}{%
      \edef\@extag{\noexpand\ \noexpand\textup{[\noexpand\texttt{e:\@exlblkey}]}}}%
    \begin{example}[#2\ifshowkeys\@extag\fi]%
    \label{e:\@exlblkey}%
    % record the long title so \nref{e:KEY} prints it, as for theorems
    \protected@write\@auxout{}{%
      \string\@thmtitle{e:\@exlblkey}{\unexpanded{#2}}}%
    \expanded{%
      \noexpand\newglossaryentry{e\theexcounter}{%
        name={\unexpanded{#2}%
          \IfBooleanF{#1}{ [\noexpand\texttt{e:\@exlblkey}]}},%
        sort={\theexample},%
        description={p.~\noexpand\pageref{e:\@exlblkey}},%
        user1={\theexample},%
        user2={Example},%
        user3={\thesection},%
        user4={\expandafter\unexpanded\expandafter{\@savedSecTitle}}%
      }%
    }%
    \protected@write\@auxout{}{\string\glsliveentry{e\theexcounter}}%
  \fi
}
\newcommand{\eeXa}{\end{example}}

% DEFINITION COMMAND -- same signature and behaviour as \beT below.
% Usage: \beD[long title][short title][key]
%   long title  = definition heading (plain text only)
%   short title = shown in \dref{word}; defaults to the key
%   key         = optional label key override; the label becomes d:KEY
% Without a key the label is d: plus the first three letters of the title's
% first word, so it is a mnemonic rather than the definition's number: a
% numeric label would silently change whenever a definition is inserted.
\NewDocumentCommand{\beD}{s O{} O{} o}{%
  \def\@temparg{#2}%
  \ifx\@temparg\@empty
    \begin{definition}%
  \else
    \let\@savedSecTitle\@currentsectionname%
    \refstepcounter{defcounter}%
    \IfValueTF{#4}{%
      \edef\@deflblkey{#4}%
    }{%
      \edef\@deffirstword{\@firstwordof#2 \relax}%
      \edef\@deflblkey{%
        \expandafter\@takethree\@deffirstword\@empty\@empty\@empty\@nil}%
    }%
    \def\@tempshort{#3}%
    \ifx\@tempshort\@empty
      \expandafter\xdef\csname @dshort@\@deflblkey\endcsname{\@deflblkey}%
    \else
      \expandafter\def\csname @dshort@\@deflblkey\endcsname{#3}%
    \fi
    \IfBooleanTF{#1}{\let\@deftag\@empty}{%
      \edef\@deftag{\noexpand\ \noexpand\textup{[\noexpand\texttt{d:\@deflblkey}]}}}%
    \begin{definition}[#2\ifshowkeys\@deftag\fi]%
    \label{d:\@deflblkey}%
    % record the long title so \nref{d:KEY} prints it, as for theorems
    \protected@write\@auxout{}{%
      \string\@thmtitle{d:\@deflblkey}{\unexpanded{#2}}}%
    \expanded{%
      \noexpand\newglossaryentry{d\thedefcounter}{%
        name={\unexpanded{#2}%
          \IfBooleanF{#1}{ [\noexpand\texttt{d:\@deflblkey}]}},%
        sort={\thedefinition},%
        description={p.~\noexpand\pageref{d:\@deflblkey}},%
        user1={\thedefinition},%
        user2={Definition},%
        user3={\thesection},%
        user4={\expandafter\unexpanded\expandafter{\@savedSecTitle}}%
      }%
    }%
    \protected@write\@auxout{}{\string\glsliveentry{d\thedefcounter}}%
  \fi
}
\newcommand{\eeD}{\end{definition}}

% \showkeysfalse in a driver hides the [t:KEY] mnemonics in result headings
% (labels and glossary entries are unchanged); used for submission builds.
\newif\ifshowkeys \showkeystrue
% \commonredtrue in a driver that inputs basicFiles/commonRed: tool.tex then
% cites its reduction results instead of restating them.
\newif\ifcommonred \commonredfalse
% TheorEM COMMAND
% Usage: \beT[long title][short title][key]
%   long title  = theorem heading (plain text only)
%   short title = shown in \tref{word}; may contain math; defaults to first word of long title
%   key         = optional label key override; label becomes t:KEY (avoids first-word collisions)
% Label is t:WORD where WORD = first word of long title (or KEY if given); use \tref{WORD} to reference.
\NewDocumentCommand{\beT}{s O{} O{} o}{%
  \def\@temparg{#2}%
  \ifx\@temparg\@empty
    \begin{theorem}%
  \else
    \let\@savedSecTitle\@currentsectionname%
    \refstepcounter{thmcounter}%
    % \@thmtag is the mnemonic shown in the heading; every theorem gets one.
    % With an explicit key the key itself is shown.  Without one the key
    % defaults to the title's first word, and showing that whole word would
    % merely repeat the title ("Dissipativity [Dissipativity]"), so only its
    % first three letters are shown: "Dissipativity [Dis]".  The label is
    % unaffected -- it still uses the full \@thmlblkey.
    \IfValueTF{#4}{%
      \edef\@thmlblkey{#4}%
    }{%
      \edef\@thmfirstword{\@firstwordof#2 \relax}%
      \edef\@thmlblkey{%
        \expandafter\@takethree\@thmfirstword\@empty\@empty\@empty\@nil}%
    }%
    \IfBooleanTF{#1}{\let\@thmtag\@empty}{%
      \edef\@thmtag{\noexpand\ \noexpand\textup{[\noexpand\texttt{t:\@thmlblkey}]}}}%
    \def\@tempshort{#3}%
    \ifx\@tempshort\@empty
      \expandafter\xdef\csname @tshort@\@thmlblkey\endcsname{\@thmlblkey}%
    \else
      \expandafter\def\csname @tshort@\@thmlblkey\endcsname{#3}%
    \fi
    \begin{theorem}[#2\ifshowkeys\@thmtag\fi]%
    \label{t:\@thmlblkey}%
    % record the long title for \nref; via the .aux so forward refs resolve
    \protected@write\@auxout{}{%
      \string\@thmtitle{t:\@thmlblkey}{\unexpanded{#2}}}%
    \expanded{%
      \noexpand\newglossaryentry{t\thethmcounter}{%
        name={\unexpanded{#2}%
          \IfBooleanF{#1}{ [\noexpand\texttt{t:\@thmlblkey}]}},%
        % sort/user1 must be the number the document actually prints, i.e. the
        % theorem environment's counter, not thmcounter.  thmcounter advances
        % only on \beT with a title, so a bare \beT (or \beT with an
        % empty title) makes the two drift and the glossary cites a wrong
        % number.  thmcounter is still fine as the entry key: it is unique.
        sort={\thetheorem},%
        description={p.~\noexpand\pageref{t:\@thmlblkey}},%
        user1={\thetheorem},%
        user2={Theorem},%
        user3={\thesection},%
        user4={\expandafter\unexpanded\expandafter{\@savedSecTitle}}%
      }%
    }%
    \protected@write\@auxout{}{\string\glsliveentry{t\thethmcounter}}%
  \fi
}
\newcommand{\eeT}{\end{theorem}}

% CROSS-PART RELATION PARAGRAPH
% Usage: \xpart{text}  -- place immediately BEFORE \beT, never inside it
%   (anything inside \beT...\eeT would be swallowed into the glossary entry).
% The paragraph is typeset only when Part I is actually present, i.e. in the
% combined book build; in the standalone builds (GranadaDSABNS/A.tex,
% AAHRah-Submission/AAHRah.tex) the label part:glv is undefined and the whole
% paragraph disappears rather than leaving ?? cross-references. Same
% graceful-degradation principle as \GASref in GASt.tex. Resolves on the
% second LaTeX pass, since r@part:glv comes from the .aux file.
% (Already inside the \makeatletter region opened above, so no \makeatletter
% pair here -- adding one would close that region early.)
\newcommand{\xpart}[1]{\@ifundefined{r@part:glv}{}{\paragraph{Relation with Parts~I--II.}#1}}

% \bookref{label}{precise text}{fallback text}: prints the precise cross-reference
% when the label resolves (combined book build), and the fallback text otherwise
% (standalone paper builds, where book-only appendix labels are absent). Same
% graceful-degradation role as \GASref in GASt.tex, but usable in any file since it
% lives in this preamble (loaded by every driver) -- unlike \GASref, which is defined in a
% content file and is therefore undefined in standalone builds that omit that file.
\newcommand{\bookref}[3]{\@ifundefined{r@#1}{#3}{#2}}

% \PVpart and \RBpart: the two parts that a partial build may omit.
% They print "Part~N" in any build that contains the part, and a naming phrase
% in a build that does not, so an excerpt shows no ?? .  Each already supplies
% the word "Part", so write "\PVpart develops ...", not "Part~\PVpart".
\newcommand{\PVpart}{\bookref{part:aahrah}%
  {Part~\ref{part:aahrah}}{the Perron--Volterra part of the full book}}
\newcommand{\RBpart}{\bookref{part:recbrowne}%
  {Part~\ref{part:recbrowne}}{the immuno-virology part of the full book}}

% OPEN PROBLEM COMMANDS
% ASSUMPTION COMMANDS (glossaried)
\newcommand{\beAs}[1][]{%
  \def\@temparg{#1}%
  \ifx\@temparg\@empty
    \begin{assumption}%
  \else
    \let\@savedSecTitle\@currentsectionname%
    \refstepcounter{ascounter}%
    \begin{assumption}[#1]%
    \label{as:\theascounter}%
    \expanded{%
      \noexpand\newglossaryentry{as\theascounter}{%
        name={\unexpanded{#1}},%
        sort={\theascounter},%
        description={p.~\noexpand\pageref{as:\theascounter}},%
        user1={\theascounter},%
        user2={Assumption},%
        user3={\thesection},%
        user4={\expandafter\unexpanded\expandafter{\@savedSecTitle}}%
      }%
    }%
    \protected@write\@auxout{}{\string\glsliveentry{as\theascounter}}%
  \fi
}
\newcommand{\eeAs}{\end{assumption}}
% CONJECTURE COMMANDS (glossaried)
\newcommand{\beCj}[1][]{%
  \def\@temparg{#1}%
  \ifx\@temparg\@empty
    \begin{conjecture}%
  \else
    \let\@savedSecTitle\@currentsectionname%
    \refstepcounter{cjcounter}%
    \begin{conjecture}[#1]%
    \label{cj:\thecjcounter}%
    \expanded{%
      \noexpand\newglossaryentry{cj\thecjcounter}{%
        name={\unexpanded{#1}},%
        sort={\thecjcounter},%
        description={p.~\noexpand\pageref{cj:\thecjcounter}},%
        user1={\thecjcounter},%
        user2={Conjecture},%
        user3={\thesection},%
        user4={\expandafter\unexpanded\expandafter{\@savedSecTitle}}%
      }%
    }%
    \protected@write\@auxout{}{\string\glsliveentry{cj\thecjcounter}}%
  \fi
}
\newcommand{\eeCj}{\end{conjecture}}
% LEMMA COMMANDS
% ---------------------------------------------------------------------
% UNIFORM RESULT WRAPPERS -- \beL, \beC, \beP, \beR share \beT's signature
%
%   \beX[long title][short title][key]  ...  \eeX
%
%   long title  = heading (math allowed)
%   short title = shown in the cross-reference; defaults to the key
%   key         = label key.  With a key the label is X:KEY and the heading
%                 carries the mnemonic [key], exactly as for \beT.  Without
%                 one the label falls back to the running counter, which is
%                 what every pre-existing call site in the book relies on.
%
% A bare \beX -- arguments commented out -- opens the plain environment and
% creates NO label and NO glossary entry.  That is the supported way to keep
% an individual result out of the catalogue.
% ---------------------------------------------------------------------
% #1 wrapper name  #2 environment  #3 label prefix  #4 counter  #5 type name
\newcommand{\@mkresult}[5]{%
  % A leading * suppresses the [key] mnemonic in both the heading and the
  % glossary name, while still making the label X:KEY and the entry:
  %     \beX*[long][short][key]
  \expandafter\NewDocumentCommand\csname be#1\endcsname{s O{} O{} o}{%
    \def\@temparg{##2}%
    \ifx\@temparg\@empty
      \begin{#2}%
    \else
      \let\@savedSecTitle\@currentsectionname%
      \refstepcounter{#4}%
      % With a key: mnemonic label X:KEY and the [key] tag in the heading.
      % Without one: fall back to the running counter, exactly as these
      % wrappers behaved before, so every legacy one-argument call site in
      % the book keeps its label.  The counter is also csname-safe, which a
      % title's first word is not -- titles beginning with math would
      % otherwise build \csname \Omega...\endcsname and fail.
      \IfValueTF{##4}{%
        \edef\@relbl{##4}%
        \IfBooleanTF{##1}{\let\@retag\@empty}{%
          \edef\@retag{\noexpand\ \noexpand\textup{[\noexpand\texttt{#3:##4}]}}}%
      }{%
        \edef\@relbl{\arabic{#4}}%
        \let\@retag\@empty
      }%
      \def\@tempshort{##3}%
      \ifx\@tempshort\@empty
        \expandafter\xdef\csname @sshort@#3\@relbl\endcsname{\@relbl}%
      \else
        \expandafter\def\csname @sshort@#3\@relbl\endcsname{##3}%
      \fi
      \begin{#2}[##2\ifshowkeys\@retag\fi]%
      \edef\@tmplbl{#3:\@relbl}\expandafter\label\expandafter{\@tmplbl}%
      \expanded{%
        \noexpand\newglossaryentry{#3\csname the#4\endcsname}{%
          name={\unexpanded{##2}%
            \IfBooleanF{##1}{\IfValueT{##4}{ [\noexpand\texttt{#3:##4}]}}},%
          sort={\csname the#2\endcsname},%
          description={p.~\noexpand\pageref{#3:\@relbl}},%
          user1={\csname the#2\endcsname},%
          user2={#5},%
          user3={\thesection},%
          user4={\expandafter\unexpanded\expandafter{\@savedSecTitle}}%
        }%
      }%
      \protected@write\@auxout{}{\string\glsliveentry{#3\csname the#4\endcsname}}%
      \protected@write\@auxout{}{\string\@thmtitle{\@tmplbl}{\unexpanded{##2}}}%
    \fi
  }%
  \expandafter\newcommand\csname ee#1\endcsname{\end{#2}}%
}
\@mkresult{L}{lemma}{l}{lemcounter}{Lemma}
% REMARK / COROLLARY / PROPOSITION -- same signature as \beL above
\@mkresult{R}{remark}{r}{remcounter}{Remark}
\@mkresult{C}{corollary}{c}{corcounter}{Corollary}
\@mkresult{P}{proposition}{p}{propcounter}{Proposition}
% OPEN-PROBLEM COMMAND -- generated with the same signature as \beT:
%   \beO[long title][short title][key]   ...   \eeO
% label o:KEY, or o:<counter> when no key is given; a bare \beO makes
% neither a label nor a glossary entry.
\@mkresult{O}{open}{o}{opcounter}{Problem}


% Section-level title tracker: updated only by \section, not \subsection.
% \@currentlabelname is overwritten by every \refstepcounter (including
% \subsection), so we save it separately here.
\newcommand{\@currentsectionname}{}
\AtBeginDocument{%
  \let\@epidcrn@orig@refstepcounter\refstepcounter
  \renewcommand{\refstepcounter}[1]{%
    \@epidcrn@orig@refstepcounter{#1}%
    \ifstrequal{#1}{section}{\global\let\@currentsectionname\@currentlabelname}{}%
  }%
}

% TABLE COMMANDS (glossaried) -- consistent naming with \beT, \beD, etc.
\newcommand{\beTa}[1][]{%
  \def\@temparg{#1}%
  \ifx\@temparg\@empty
    \begin{table}[htbp]%
  \else
    \let\@savedSecTitle\@currentsectionname%
    \refstepcounter{tabcounter}%
    \begin{table}[htbp]%
    \label{A:\thetabcounter}%
    \expanded{%
      \noexpand\newglossaryentry{A\thetabcounter}{%
        name={\unexpanded{#1}},%
        sort={\thetabcounter},%
        description={p.~\noexpand\pageref{A:\thetabcounter}},%
        user1={\thetabcounter},%
        user2={Table},%
        user3={\thesection},%
        user4={\expandafter\unexpanded\expandafter{\@savedSecTitle}}%
      }%
    }%
    \protected@write\@auxout{}{\string\glsliveentry{A\thetabcounter}}%
  \fi
}
\newcommand{\eeTa}{\end{table}}

% FIGURE COMMANDS (glossaried) -- consistent with \beT, \beR, \beTa, etc.
% Usage:
%   \beF[Glossary title]...\caption{Visible caption}\eeF  -> glossaried figure
%   \beF...\caption{Visible caption}\eeF                  -> plain [H] figure, no glossary
\newcommand{\beF}[1][]{%
  \def\@temparg{#1}%
  \ifx\@temparg\@empty
    \begin{figure}[H]%
  \else
    \begin{figure}[H]%
    \let\@savedSecTitle\@currentsectionname%
    \refstepcounter{figcounter}%
    \label{fig:\thefigcounter}%
    \expanded{%
      \noexpand\newglossaryentry{Fig\thefigcounter}{%
        name={\unexpanded{#1}},%
        sort={\thefigcounter},%
        description={p.~\noexpand\pageref{fig:\thefigcounter}},%
        user1={\thefigcounter},%
        user2={Figure},%
        user3={\thesection},%
        user4={\expandafter\unexpanded\expandafter{\@savedSecTitle}}%
      }%
    }%
    \protected@write\@auxout{}{\string\glsliveentry{Fig\thefigcounter}}%
  \fi
}
\newcommand{\eeF}{\end{figure}}
\makeatother
% LIVE-ENTRY REGISTER
% \glsaddall prints every entry the glossaries package currently knows, and
% that set includes entries restored from \jobname.glsdefs.  The package
% rewrites .glsdefs from all *currently defined* entries, so anything it
% restored is written straight back out: a deleted section kept reappearing in
% the catalogue for ever, and the .glsdefs had to be erased by hand.  The
% wrappers therefore record every entry they actually create in the .aux, and
% \printboth adds only those.  The .aux is rebuilt from scratch on each run, so
% this is self-healing: a removed result is gone after one further run.
\makeatletter
\newcommand{\@glslive}{}
\newcommand{\@glsaddlive}[1]{\ifglsentryexists{#1}{\glsadd{#1}}{}}
\newcommand{\glsliveentry}[1]{\g@addto@macro\@glslive{\@glsaddlive{#1}}}
\makeatother
% PRINT COMMAND
% No \clearpage here -- the catalogue continues on the current page.  This has
% been re-added by mistake several times; leave it out.
\makeatletter
\newcommand{\printboth}{%
  \@glslive   % add exactly the entries this document created last run
  \printglossary[style=unifiedstyle, title={0 \hspace{.22em} Definitions, Theorems, Examples and Open Problems}]%
}
\makeatother
% Named theorem reference: \tref{word} -> Thm~N(word-or-short)
% Shows stored short title if available, otherwise shows the key itself.
\newcommand{\tref}[1]{%
  Thm~\ref{t:#1}(\ifcsname @tshort@#1\endcsname\csname @tshort@#1\endcsname\else#1\fi)%
}

% Full-title theorem reference: \nref{t:KEY} prints the number and then the
% theorem's long title, so one writes "Theorem~\nref{t:blocks}" and gets
% "Theorem 3 ``Block factorization and siphons''".  Takes the complete label, like
% \ref, not the bare key like \tref.
%
% \beT writes each title to the .aux (see \@thmtitle below), so \nref resolves
% forward references from the second LaTeX pass, exactly as \ref does.  If the
% title is not yet known -- first pass, or a label that is not a \beT theorem
% -- it degrades to the number rather than producing ??.
% These two use @ in their names, and this point in the file is outside the
% \makeatletter region opened earlier (it closes at \makeatother above), so
% the region has to be reopened here.
\makeatletter
\newcommand{\@thmtitle}[2]{\global\@namedef{@tlong@#1}{#2}}
\newcommand{\nref}[1]{%
  \ref{#1}%
  \ifcsname @tlong@#1\endcsname
    \ ``\csname @tlong@#1\endcsname''%
  \fi
}
% \nm{t:KEY}: the stored long title of any \beT/\beL/\beP/\beR/\beC/\beD/\beE result,
% verbatim and without quotes, so item names follow the macro's first argument.
\newcommand{\nm}[1]{%
  \ifcsname @tlong@#1\endcsname
    \csname @tlong@#1\endcsname%
  \else
    \ref{#1}%
  \fi
}
\makeatother




%% ========================================
%% Other THEOREM ENVIRONMENTS
%% ========================================
\makeatletter
\@ifundefined{lemma}      {\newtheorem{lemma}      {Lemma}}      {}
\@ifundefined{proposition}{\newtheorem{proposition}{Proposition}} {}
\@ifundefined{corollary}  {\newtheorem{corollary}  {Corollary}}  {}
\@ifundefined{conjecture} {\newtheorem{conjecture} {Conjecture}}  {}
\@ifundefined{remark}     {\newtheorem{remark}     {Remark}}     {}
\@ifundefined{question}   {\newtheorem{question}   {Question}}   {}
\@ifundefined{assumption} {\newtheorem{assumption} {Assumption}} {}
\@ifundefined{hypothesis} {\newtheorem{hypothesis} {Hypothesis}} {}
\@ifundefined{fact}       {\newtheorem{fact}       {Fact}}       {}
\makeatother

% \citer: cite the relay-graph threshold theorem of AAHrel. If a document has
% \input the chapter/file containing t:relayRF (so the label resolves), point
% straight at it; otherwise fall back to a bare citation of AAHrel. Shared here
% (not per-paper) so it resolves consistently across every book chapter.
% \citeR: same idea, for the first theorem of AAHRah (t:Pstar).
\makeatletter
\providecommand{\citer}{%
  \@ifundefined{r@t:relayRF}%
    {\cite{AAHrel}}%
    {Theorem~\ref{t:relayRF} of \cite{AAHrel}}%
}
\providecommand{\citeR}{%
  \@ifundefined{r@t:Pstar}%
    {\cite{AAHRah}}%
    {Theorem~\ref{t:Pstar} of \cite{AAHRah}}%
}
\makeatother


% Abbreviations
\def\beQ{\begin{question}}\def\eeQ{\end{question}}
\def\beA{\begin{assumption}}\def\eeA{\end{assumption}}


%% ========================================

%% ========================================
%% COLORS & CUSTOM COMMANDS
%% ========================================
\definecolor{funccolor}{RGB}{25,25,112}
\definecolor{desccolor}{RGB}{64,64,64}
\newcommand{\func}[1]{\texttt{\textcolor{funccolor}{#1}}}
\def\eq{=}
\def\mm{-}

%% ========================================
%% MATHCAL SHORTHANDS
%% ========================================
% Project-wide convention: \mX for a capital X is always \mathcal X.
% Defined for the whole alphabet so the convention never fails to expand.
\newcommand{\mA}{\mathcal A}\newcommand{\mB}{\mathcal B}
\newcommand{\mC}{\mathcal C}\newcommand{\mD}{\mathcal D}
\newcommand{\mE}{\mathcal E}\newcommand{\mF}{\mathcal F}
\newcommand{\mG}{\mathcal G}\newcommand{\mH}{\mathcal H}
\newcommand{\mI}{\mathcal I}\newcommand{\mJ}{\mathcal J}
\newcommand{\mK}{\mathcal K}\newcommand{\mL}{\mathcal L}
\newcommand{\mM}{\mathcal M}\newcommand{\mN}{\mathcal N}
\newcommand{\mO}{\mathcal O}\newcommand{\mP}{\mathcal P}
\newcommand{\mQ}{\mathcal Q}\newcommand{\mR}{\mathcal R}
\newcommand{\mS}{\mathcal S}\newcommand{\mT}{\mathcal T}
\newcommand{\mU}{\mathcal U}\newcommand{\mV}{\mathcal V}
\newcommand{\mW}{\mathcal W}\newcommand{\mX}{\mathcal X}
\newcommand{\mY}{\mathcal Y}\newcommand{\mZ}{\mathcal Z}

%% ========================================
%% BASIC MATH ENVIRONMENTS
%% ========================================
\def\bep{\begin{pmatrix}}\def\eep{\end{pmatrix}}
\def\bev{\begin{vmatrix}}\def\eev{\end{vmatrix}}
\def\bea{\begin{eqnarray*}}\def\eea{\end{eqnarray*}}
\def\bc{\begin{cases}}\def\ec{\end{cases}}
\def\BEN{\begin{enumerate}}\def\EEN{\end{enumerate}}
\def\BI{\begin{itemize}}\def\EI{\end{itemize}}

\newcommand{\be}[1]{\begin{equation}\label{#1}}
\newcommand{\ee}{\end{equation}}
\newcommand{\beq}{\begin{eqnarray}}
\def\eeq{\end{eqnarray}}
\def\eqr{\eqref}\def\fr{\frac}\def\lbl{\label}
\def\Lra{\Longrightarrow}\def\non{\nonumber}\def\cNGM{\cite{Diek,Van,Van08}}
\def\Eq{\Leftrightarrow}\def\b{\beta}\def\uph{\red{UP TO HERE}}
%% ========================================
%% PROBABILITY & STATISTICS
%% ========================================
\renewcommand{\P}{\ensuremath{\mathbb{P}}}  % Note: overrides paragraph symbol
\newcommand{\E}{\ensuremath{\mathbb{E}}}
\newcommand{\ind}[1]{\ensuremath{\mathbbm{1}_{\left\{#1\right\}}}}
\DeclareMathOperator{\Cov}{Cov}
\DeclareMathOperator{\TV}{TV}
\newcommand{\rank}{\operatorname{rank}}

%% ========================================
%% CRN OPERATORS (organizations/ERCs/buffering/siphons framework)
%% ========================================
\def\clos{\mathrm{clos}}
\def\Supply{\mathrm{Supply}}
\def\req{\mathrm{req}}
\def\React{\mathrm{React}}
\def\Prod{\mathrm{Prod}}
\def\Gen{\mathrm{Gen}}
\def\In{\mathrm{In}}
\def\Out{\mathrm{Out}}
\def\Act{\mathrm{Act}}
\def\Min{\mathrm{Min}}

%% ========================================
%% LINE BREAKING
%% ========================================
% Last-resort stretch, used only on paragraphs TeX cannot otherwise break:
% a long \texttt{} token, a comma-separated inline math list (commas are not
% breakpoints in math), or a run of \cite/\ref boxes.  Without it those set as
% overfull \hbox and print into the margin; with it they set slightly loose,
% which is the lesser evil.  Does not affect paragraphs that already break.
\emergencystretch=3em

%% ========================================
%% NUMBER SETS
%% ========================================
\newcommand{\R}{\mathbb{R}}  % Real numbers
\newcommand{\C}{\mathbb{C}}  % Complex numbers
\newcommand{\N}{\mathbb{N}}
\newcommand{\Z}{\mathbb{Z}}
% \providecommand, not \newcommand: some papers define these in their own
% preamble after % Shared preamble for the GLVbook papers.  This is the single canonical copy;
% a driver one level below GLVbook reaches it as
%
%     \input{../basicFiles/preamble}
%
% What stays in the driver, because it differs from paper to paper:
% \documentclass, font choice, \title, \author, \date, and any macro genuinely
% local to one paper.  Everything else belongs here.  The file is
% self-contained: it does not read any other file.

\usepackage[utf8]{inputenc}
\usepackage[novbox]{pdfsync}   % required: SyncTeX alone does not work here

% Page geometry, set once here for every paper that loads this file.  A driver
% that needs different margins says \geometry{...}; a second
% \usepackage[...]{geometry} raises an option clash and must not be used.
\makeatletter
\@ifundefined{AAdefNnogeometry}{\usepackage[a4paper,margin=2cm]{geometry}}{}
\makeatother

% mathematics
\usepackage{amsmath,amssymb,amsthm,mathtools,bm,mathrsfs,bbm}
\usepackage{centernot}

% floats, tables, lists
\usepackage{graphicx,float}
\makeatletter
\@ifpackageloaded{subfig}{}{\usepackage{subcaption}}
\makeatother
\usepackage{booktabs,longtable,threeparttable,array,multirow}
\usepackage{enumitem}

% colour, code, algorithms
\usepackage{xcolor,listings,tcolorbox}
\makeatletter
\@ifundefined{c@algorithm}{\usepackage{algorithm}}{}
\makeatother
\usepackage{algorithmicx,algpseudocode}

% diagrams
\usepackage{tikz,tikz-cd,pgfplots}
\usetikzlibrary{graphs,decorations.pathmorphing,decorations.markings,
arrows.meta,automata,positioning}
\tikzcdset{every label/.append style = {font = \small}}
\pgfplotsset{compat=1.18}

% chemistry
\usepackage[version=4]{mhchem}

\usepackage{microtype}
\usepackage{lineno}
\usepackage{nomencl}
\makenomenclature

% hyperref first, then bookmark, then glossaries: glossaries must follow
% hyperref or the entry links are built against the wrong anchors.
\makeatletter
\@ifundefined{hyperrefpackage}{%
  \@ifundefined{AAdefNnohyperref}{\usepackage{hyperref}}{}%
}{}
\@ifundefined{hyperrefpackage}{\usepackage{bookmark}}{}
\makeatother
\usepackage[toc,sort=use]{glossaries}
\makeglossaries

% Book/standalone switch.  Material that the book already covers in its
% prelude, but that a standalone paper still has to state, is wrapped as
%     \ifbookbuild\else ... \fi
% Drivers default to standalone; bookM.tex says \bookbuildtrue.
\newif\ifbookbuild
\bookbuildfalse

% Part switches for the gradostat paper.  Both default to true; a driver that
% wants the short submission says \madecfalse and/or \slowfastfalse, and every
% cross-reference into the dropped part disappears with it.
%   \ifmadec     Part III, the Lagasquie--Madec model and its break-even geometry
%   \ifslowfast  Part IV, everything carrying the small parameter \eps
\newif\ifmadec     \madectrue
\newif\ifslowfast  \slowfasttrue

% Part switch for the Barrat/simplicial-contagion project.  Defaults to
% true (the hypernetwork-synchrony part is included); the journal-paper
% driver says \hnetpartfalse to drop it and every cross-reference into it.
\newif\ifhnetpart  \hnetparttrue

%% ========================================
%% PACKAGES
%% ========================================
\RequirePackage{etoolbox}
\usepackage{makecell}
\usepackage[novbox]{pdfsync}% obsolete; SyncTeX is built into modern engines
\usepackage{amssymb, amsfonts,mathtools,bbm,eurosym,graphicx,latexsym, amsmath,amsthm,times,nccmath,url}
\usepackage{listings,verbatim,cases,tabularx}
% hyperref is deliberately NOT in the line above.  Journal classes often
% load it themselves, with a long option list, from a begindocument hook
% (mdpi.cls does: \DeclareHookRule{begindocument}{hyperref}{before}{mdpi}).
% Loading it here first, without options, makes that later load raise
% "Option clash for package hyperref".  So load it only if nobody else
% has, or will:
%   - \@ifpackageloaded  catches a class that loaded it already;
%   - \hyperrefpackage    is mdpi.cls's own deferred loader;
%   - \AAdefNnohyperref   lets any other driver opt out by defining it
%                         before \input{preamble}.
\makeatletter
\@ifpackageloaded{hyperref}{}{%
  \@ifundefined{hyperrefpackage}{%
    \@ifundefined{AAdefNnohyperref}{\usepackage{hyperref}}{}%
  }{}%
}
\makeatother
% linktoc=all makes a table-of-contents page number a link as well as the
% entry text; hyperref's default, linktoc=section, links only the text.
% Deferred to begin-document and guarded: by the block above hyperref may
% be loaded later (mdpi's own hook) or not at all (\AAdefNnohyperref), and
% \hypersetup must exist before it is called.  A driver that wants the
% default back can say \hypersetup{linktoc=section} after \input{preamble}.
\makeatletter
\AtBeginDocument{\@ifundefined{hypersetup}{}{\hypersetup{linktoc=all}}}
\makeatother
\usepackage{graphicx,float,fancybox}
\usepackage{epsfig,epstopdf,tikz}
\usepackage[sans]{dsfont}
\usetikzlibrary{positioning,arrows,arrows.meta,calc}
% Configure listings for Mathematica code
\lstdefinestyle{mathematica}{
 language=Mathematica,
 basicstyle=\small\ttfamily,
 backgroundcolor=\color{gray!10},
 frame=single,
 breaklines=true,
 commentstyle=\color{green!60!black},
 keywordstyle=\color{blue},
 stringstyle=\color{red},
 showstringspaces=false,
 tabsize=2
}
\usepackage{tcolorbox}
% Colored boxes for key results
\newtcolorbox{keyresult}[1][]{
 colback=blue!5!white,
 colframe=blue!75!black,
 title=#1,
 fonttitle=\bfseries
}
\newtcolorbox{examplebox}[1][]{
 colback=green!5!white,
 colframe=green!75!black,
 title=#1,
 fonttitle=\bfseries
}


\makeatletter
\@ifundefined{theorem}   {\newtheorem{theorem}   {Theorem}}   {}
\@ifundefined{definition}{\newtheorem{definition}{Definition}} {}
\@ifundefined{open}      {\newtheorem{open}      {Problem}}   {}
\@ifundefined{example}   {\newtheorem{example}   {Example}}   {}
\makeatother

% COUNTER DEFINITIONS
\newcounter{tabcounter}
\newcounter{defcounter}
\newcounter{opcounter}
\newcounter{thmcounter}
\newcounter{excounter}
\newcounter{ascounter}
\newcounter{cjcounter}
\newcounter{lemcounter}
\newcounter{remcounter}
\newcounter{corcounter}
\newcounter{propcounter}
\newcounter{figcounter}

% UNIFIED GLOSSARY STYLE
\newglossarystyle{unifiedstyle}{%
  \setglossarystyle{list}%
  \renewcommand*{\glossarypreamble}{\gdef\glslastsec{}}%
  \renewcommand{\glossentry}[2]{%
    \edef\glscursec{\glsentryuseriii{##1}}%
    \ifx\glscursec\glslastsec\else
      \item[]\textbf{\S\,\glscursec\enspace\glsentryuseriv{##1}}%
      \global\let\glslastsec\glscursec
    \fi
    \item[\glsentryuserii{##1} \glsentryuseri{##1}:] ``\glossentryname{##1}'' \glsentrydesc{##1}%
  }%
}

\makeatletter
% First-word extractor (used by \beT, \beD and \beXa for the default key)
% \@firstwordof hello world \relax  ->  hello
\def\@firstwordof#1 #2\relax{#1}
% First three letters, used for the heading mnemonic of a theorem given no
% explicit key.  Padding with \@empty keeps it safe for first words shorter
% than three letters: \@empty vanishes under \edef, whereas \relax would
% survive into the label and is illegal inside \csname (it broke every
% two-letter title, e.g. "ME ...", with "Missing \endcsname inserted").
\def\@takethree#1#2#3#4\@nil{#1#2#3}

% EXAMPLE COMMAND -- same signature and behaviour as \beT and \beD.
% Usage: \beXa*[long title][short title][key]
%   star  = suppresses the printed "[key]" tag after the title (as for \beT*
%   and \beD*); the label and glossary entry are unaffected.
%   key = optional label key override; the label becomes e:KEY, and without a
%   key it is e: plus the first three letters of the title's first word.  A
%   mnemonic label lets an example be cited as "Example~\ref{e:KEY}" instead
%   of by section number, which is what the citing prose needs.
\NewDocumentCommand{\beXa}{s O{} O{} o}{%
  \def\@temparg{#2}%
  \ifx\@temparg\@empty
    \begin{example}%
  \else
    \let\@savedSecTitle\@currentsectionname%
    \refstepcounter{excounter}%
    \IfValueTF{#4}{%
      \edef\@exlblkey{#4}%
    }{%
      \edef\@exfirstword{\@firstwordof#2 \relax}%
      \edef\@exlblkey{%
        \expandafter\@takethree\@exfirstword\@empty\@empty\@empty\@nil}%
    }%
    \def\@tempshort{#3}%
    \ifx\@tempshort\@empty
      \expandafter\xdef\csname @xshort@\@exlblkey\endcsname{\@exlblkey}%
    \else
      \expandafter\def\csname @xshort@\@exlblkey\endcsname{#3}%
    \fi
    \IfBooleanTF{#1}{\let\@extag\@empty}{%
      \edef\@extag{\noexpand\ \noexpand\textup{[\noexpand\texttt{e:\@exlblkey}]}}}%
    \begin{example}[#2\ifshowkeys\@extag\fi]%
    \label{e:\@exlblkey}%
    % record the long title so \nref{e:KEY} prints it, as for theorems
    \protected@write\@auxout{}{%
      \string\@thmtitle{e:\@exlblkey}{\unexpanded{#2}}}%
    \expanded{%
      \noexpand\newglossaryentry{e\theexcounter}{%
        name={\unexpanded{#2}%
          \IfBooleanF{#1}{ [\noexpand\texttt{e:\@exlblkey}]}},%
        sort={\theexample},%
        description={p.~\noexpand\pageref{e:\@exlblkey}},%
        user1={\theexample},%
        user2={Example},%
        user3={\thesection},%
        user4={\expandafter\unexpanded\expandafter{\@savedSecTitle}}%
      }%
    }%
    \protected@write\@auxout{}{\string\glsliveentry{e\theexcounter}}%
  \fi
}
\newcommand{\eeXa}{\end{example}}

% DEFINITION COMMAND -- same signature and behaviour as \beT below.
% Usage: \beD[long title][short title][key]
%   long title  = definition heading (plain text only)
%   short title = shown in \dref{word}; defaults to the key
%   key         = optional label key override; the label becomes d:KEY
% Without a key the label is d: plus the first three letters of the title's
% first word, so it is a mnemonic rather than the definition's number: a
% numeric label would silently change whenever a definition is inserted.
\NewDocumentCommand{\beD}{s O{} O{} o}{%
  \def\@temparg{#2}%
  \ifx\@temparg\@empty
    \begin{definition}%
  \else
    \let\@savedSecTitle\@currentsectionname%
    \refstepcounter{defcounter}%
    \IfValueTF{#4}{%
      \edef\@deflblkey{#4}%
    }{%
      \edef\@deffirstword{\@firstwordof#2 \relax}%
      \edef\@deflblkey{%
        \expandafter\@takethree\@deffirstword\@empty\@empty\@empty\@nil}%
    }%
    \def\@tempshort{#3}%
    \ifx\@tempshort\@empty
      \expandafter\xdef\csname @dshort@\@deflblkey\endcsname{\@deflblkey}%
    \else
      \expandafter\def\csname @dshort@\@deflblkey\endcsname{#3}%
    \fi
    \IfBooleanTF{#1}{\let\@deftag\@empty}{%
      \edef\@deftag{\noexpand\ \noexpand\textup{[\noexpand\texttt{d:\@deflblkey}]}}}%
    \begin{definition}[#2\ifshowkeys\@deftag\fi]%
    \label{d:\@deflblkey}%
    % record the long title so \nref{d:KEY} prints it, as for theorems
    \protected@write\@auxout{}{%
      \string\@thmtitle{d:\@deflblkey}{\unexpanded{#2}}}%
    \expanded{%
      \noexpand\newglossaryentry{d\thedefcounter}{%
        name={\unexpanded{#2}%
          \IfBooleanF{#1}{ [\noexpand\texttt{d:\@deflblkey}]}},%
        sort={\thedefinition},%
        description={p.~\noexpand\pageref{d:\@deflblkey}},%
        user1={\thedefinition},%
        user2={Definition},%
        user3={\thesection},%
        user4={\expandafter\unexpanded\expandafter{\@savedSecTitle}}%
      }%
    }%
    \protected@write\@auxout{}{\string\glsliveentry{d\thedefcounter}}%
  \fi
}
\newcommand{\eeD}{\end{definition}}

% \showkeysfalse in a driver hides the [t:KEY] mnemonics in result headings
% (labels and glossary entries are unchanged); used for submission builds.
\newif\ifshowkeys \showkeystrue
% \commonredtrue in a driver that inputs basicFiles/commonRed: tool.tex then
% cites its reduction results instead of restating them.
\newif\ifcommonred \commonredfalse
% TheorEM COMMAND
% Usage: \beT[long title][short title][key]
%   long title  = theorem heading (plain text only)
%   short title = shown in \tref{word}; may contain math; defaults to first word of long title
%   key         = optional label key override; label becomes t:KEY (avoids first-word collisions)
% Label is t:WORD where WORD = first word of long title (or KEY if given); use \tref{WORD} to reference.
\NewDocumentCommand{\beT}{s O{} O{} o}{%
  \def\@temparg{#2}%
  \ifx\@temparg\@empty
    \begin{theorem}%
  \else
    \let\@savedSecTitle\@currentsectionname%
    \refstepcounter{thmcounter}%
    % \@thmtag is the mnemonic shown in the heading; every theorem gets one.
    % With an explicit key the key itself is shown.  Without one the key
    % defaults to the title's first word, and showing that whole word would
    % merely repeat the title ("Dissipativity [Dissipativity]"), so only its
    % first three letters are shown: "Dissipativity [Dis]".  The label is
    % unaffected -- it still uses the full \@thmlblkey.
    \IfValueTF{#4}{%
      \edef\@thmlblkey{#4}%
    }{%
      \edef\@thmfirstword{\@firstwordof#2 \relax}%
      \edef\@thmlblkey{%
        \expandafter\@takethree\@thmfirstword\@empty\@empty\@empty\@nil}%
    }%
    \IfBooleanTF{#1}{\let\@thmtag\@empty}{%
      \edef\@thmtag{\noexpand\ \noexpand\textup{[\noexpand\texttt{t:\@thmlblkey}]}}}%
    \def\@tempshort{#3}%
    \ifx\@tempshort\@empty
      \expandafter\xdef\csname @tshort@\@thmlblkey\endcsname{\@thmlblkey}%
    \else
      \expandafter\def\csname @tshort@\@thmlblkey\endcsname{#3}%
    \fi
    \begin{theorem}[#2\ifshowkeys\@thmtag\fi]%
    \label{t:\@thmlblkey}%
    % record the long title for \nref; via the .aux so forward refs resolve
    \protected@write\@auxout{}{%
      \string\@thmtitle{t:\@thmlblkey}{\unexpanded{#2}}}%
    \expanded{%
      \noexpand\newglossaryentry{t\thethmcounter}{%
        name={\unexpanded{#2}%
          \IfBooleanF{#1}{ [\noexpand\texttt{t:\@thmlblkey}]}},%
        % sort/user1 must be the number the document actually prints, i.e. the
        % theorem environment's counter, not thmcounter.  thmcounter advances
        % only on \beT with a title, so a bare \beT (or \beT with an
        % empty title) makes the two drift and the glossary cites a wrong
        % number.  thmcounter is still fine as the entry key: it is unique.
        sort={\thetheorem},%
        description={p.~\noexpand\pageref{t:\@thmlblkey}},%
        user1={\thetheorem},%
        user2={Theorem},%
        user3={\thesection},%
        user4={\expandafter\unexpanded\expandafter{\@savedSecTitle}}%
      }%
    }%
    \protected@write\@auxout{}{\string\glsliveentry{t\thethmcounter}}%
  \fi
}
\newcommand{\eeT}{\end{theorem}}

% CROSS-PART RELATION PARAGRAPH
% Usage: \xpart{text}  -- place immediately BEFORE \beT, never inside it
%   (anything inside \beT...\eeT would be swallowed into the glossary entry).
% The paragraph is typeset only when Part I is actually present, i.e. in the
% combined book build; in the standalone builds (GranadaDSABNS/A.tex,
% AAHRah-Submission/AAHRah.tex) the label part:glv is undefined and the whole
% paragraph disappears rather than leaving ?? cross-references. Same
% graceful-degradation principle as \GASref in GASt.tex. Resolves on the
% second LaTeX pass, since r@part:glv comes from the .aux file.
% (Already inside the \makeatletter region opened above, so no \makeatletter
% pair here -- adding one would close that region early.)
\newcommand{\xpart}[1]{\@ifundefined{r@part:glv}{}{\paragraph{Relation with Parts~I--II.}#1}}

% \bookref{label}{precise text}{fallback text}: prints the precise cross-reference
% when the label resolves (combined book build), and the fallback text otherwise
% (standalone paper builds, where book-only appendix labels are absent). Same
% graceful-degradation role as \GASref in GASt.tex, but usable in any file since it
% lives in this preamble (loaded by every driver) -- unlike \GASref, which is defined in a
% content file and is therefore undefined in standalone builds that omit that file.
\newcommand{\bookref}[3]{\@ifundefined{r@#1}{#3}{#2}}

% \PVpart and \RBpart: the two parts that a partial build may omit.
% They print "Part~N" in any build that contains the part, and a naming phrase
% in a build that does not, so an excerpt shows no ?? .  Each already supplies
% the word "Part", so write "\PVpart develops ...", not "Part~\PVpart".
\newcommand{\PVpart}{\bookref{part:aahrah}%
  {Part~\ref{part:aahrah}}{the Perron--Volterra part of the full book}}
\newcommand{\RBpart}{\bookref{part:recbrowne}%
  {Part~\ref{part:recbrowne}}{the immuno-virology part of the full book}}

% OPEN PROBLEM COMMANDS
% ASSUMPTION COMMANDS (glossaried)
\newcommand{\beAs}[1][]{%
  \def\@temparg{#1}%
  \ifx\@temparg\@empty
    \begin{assumption}%
  \else
    \let\@savedSecTitle\@currentsectionname%
    \refstepcounter{ascounter}%
    \begin{assumption}[#1]%
    \label{as:\theascounter}%
    \expanded{%
      \noexpand\newglossaryentry{as\theascounter}{%
        name={\unexpanded{#1}},%
        sort={\theascounter},%
        description={p.~\noexpand\pageref{as:\theascounter}},%
        user1={\theascounter},%
        user2={Assumption},%
        user3={\thesection},%
        user4={\expandafter\unexpanded\expandafter{\@savedSecTitle}}%
      }%
    }%
    \protected@write\@auxout{}{\string\glsliveentry{as\theascounter}}%
  \fi
}
\newcommand{\eeAs}{\end{assumption}}
% CONJECTURE COMMANDS (glossaried)
\newcommand{\beCj}[1][]{%
  \def\@temparg{#1}%
  \ifx\@temparg\@empty
    \begin{conjecture}%
  \else
    \let\@savedSecTitle\@currentsectionname%
    \refstepcounter{cjcounter}%
    \begin{conjecture}[#1]%
    \label{cj:\thecjcounter}%
    \expanded{%
      \noexpand\newglossaryentry{cj\thecjcounter}{%
        name={\unexpanded{#1}},%
        sort={\thecjcounter},%
        description={p.~\noexpand\pageref{cj:\thecjcounter}},%
        user1={\thecjcounter},%
        user2={Conjecture},%
        user3={\thesection},%
        user4={\expandafter\unexpanded\expandafter{\@savedSecTitle}}%
      }%
    }%
    \protected@write\@auxout{}{\string\glsliveentry{cj\thecjcounter}}%
  \fi
}
\newcommand{\eeCj}{\end{conjecture}}
% LEMMA COMMANDS
% ---------------------------------------------------------------------
% UNIFORM RESULT WRAPPERS -- \beL, \beC, \beP, \beR share \beT's signature
%
%   \beX[long title][short title][key]  ...  \eeX
%
%   long title  = heading (math allowed)
%   short title = shown in the cross-reference; defaults to the key
%   key         = label key.  With a key the label is X:KEY and the heading
%                 carries the mnemonic [key], exactly as for \beT.  Without
%                 one the label falls back to the running counter, which is
%                 what every pre-existing call site in the book relies on.
%
% A bare \beX -- arguments commented out -- opens the plain environment and
% creates NO label and NO glossary entry.  That is the supported way to keep
% an individual result out of the catalogue.
% ---------------------------------------------------------------------
% #1 wrapper name  #2 environment  #3 label prefix  #4 counter  #5 type name
\newcommand{\@mkresult}[5]{%
  % A leading * suppresses the [key] mnemonic in both the heading and the
  % glossary name, while still making the label X:KEY and the entry:
  %     \beX*[long][short][key]
  \expandafter\NewDocumentCommand\csname be#1\endcsname{s O{} O{} o}{%
    \def\@temparg{##2}%
    \ifx\@temparg\@empty
      \begin{#2}%
    \else
      \let\@savedSecTitle\@currentsectionname%
      \refstepcounter{#4}%
      % With a key: mnemonic label X:KEY and the [key] tag in the heading.
      % Without one: fall back to the running counter, exactly as these
      % wrappers behaved before, so every legacy one-argument call site in
      % the book keeps its label.  The counter is also csname-safe, which a
      % title's first word is not -- titles beginning with math would
      % otherwise build \csname \Omega...\endcsname and fail.
      \IfValueTF{##4}{%
        \edef\@relbl{##4}%
        \IfBooleanTF{##1}{\let\@retag\@empty}{%
          \edef\@retag{\noexpand\ \noexpand\textup{[\noexpand\texttt{#3:##4}]}}}%
      }{%
        \edef\@relbl{\arabic{#4}}%
        \let\@retag\@empty
      }%
      \def\@tempshort{##3}%
      \ifx\@tempshort\@empty
        \expandafter\xdef\csname @sshort@#3\@relbl\endcsname{\@relbl}%
      \else
        \expandafter\def\csname @sshort@#3\@relbl\endcsname{##3}%
      \fi
      \begin{#2}[##2\ifshowkeys\@retag\fi]%
      \edef\@tmplbl{#3:\@relbl}\expandafter\label\expandafter{\@tmplbl}%
      \expanded{%
        \noexpand\newglossaryentry{#3\csname the#4\endcsname}{%
          name={\unexpanded{##2}%
            \IfBooleanF{##1}{\IfValueT{##4}{ [\noexpand\texttt{#3:##4}]}}},%
          sort={\csname the#2\endcsname},%
          description={p.~\noexpand\pageref{#3:\@relbl}},%
          user1={\csname the#2\endcsname},%
          user2={#5},%
          user3={\thesection},%
          user4={\expandafter\unexpanded\expandafter{\@savedSecTitle}}%
        }%
      }%
      \protected@write\@auxout{}{\string\glsliveentry{#3\csname the#4\endcsname}}%
      \protected@write\@auxout{}{\string\@thmtitle{\@tmplbl}{\unexpanded{##2}}}%
    \fi
  }%
  \expandafter\newcommand\csname ee#1\endcsname{\end{#2}}%
}
\@mkresult{L}{lemma}{l}{lemcounter}{Lemma}
% REMARK / COROLLARY / PROPOSITION -- same signature as \beL above
\@mkresult{R}{remark}{r}{remcounter}{Remark}
\@mkresult{C}{corollary}{c}{corcounter}{Corollary}
\@mkresult{P}{proposition}{p}{propcounter}{Proposition}
% OPEN-PROBLEM COMMAND -- generated with the same signature as \beT:
%   \beO[long title][short title][key]   ...   \eeO
% label o:KEY, or o:<counter> when no key is given; a bare \beO makes
% neither a label nor a glossary entry.
\@mkresult{O}{open}{o}{opcounter}{Problem}


% Section-level title tracker: updated only by \section, not \subsection.
% \@currentlabelname is overwritten by every \refstepcounter (including
% \subsection), so we save it separately here.
\newcommand{\@currentsectionname}{}
\AtBeginDocument{%
  \let\@epidcrn@orig@refstepcounter\refstepcounter
  \renewcommand{\refstepcounter}[1]{%
    \@epidcrn@orig@refstepcounter{#1}%
    \ifstrequal{#1}{section}{\global\let\@currentsectionname\@currentlabelname}{}%
  }%
}

% TABLE COMMANDS (glossaried) -- consistent naming with \beT, \beD, etc.
\newcommand{\beTa}[1][]{%
  \def\@temparg{#1}%
  \ifx\@temparg\@empty
    \begin{table}[htbp]%
  \else
    \let\@savedSecTitle\@currentsectionname%
    \refstepcounter{tabcounter}%
    \begin{table}[htbp]%
    \label{A:\thetabcounter}%
    \expanded{%
      \noexpand\newglossaryentry{A\thetabcounter}{%
        name={\unexpanded{#1}},%
        sort={\thetabcounter},%
        description={p.~\noexpand\pageref{A:\thetabcounter}},%
        user1={\thetabcounter},%
        user2={Table},%
        user3={\thesection},%
        user4={\expandafter\unexpanded\expandafter{\@savedSecTitle}}%
      }%
    }%
    \protected@write\@auxout{}{\string\glsliveentry{A\thetabcounter}}%
  \fi
}
\newcommand{\eeTa}{\end{table}}

% FIGURE COMMANDS (glossaried) -- consistent with \beT, \beR, \beTa, etc.
% Usage:
%   \beF[Glossary title]...\caption{Visible caption}\eeF  -> glossaried figure
%   \beF...\caption{Visible caption}\eeF                  -> plain [H] figure, no glossary
\newcommand{\beF}[1][]{%
  \def\@temparg{#1}%
  \ifx\@temparg\@empty
    \begin{figure}[H]%
  \else
    \begin{figure}[H]%
    \let\@savedSecTitle\@currentsectionname%
    \refstepcounter{figcounter}%
    \label{fig:\thefigcounter}%
    \expanded{%
      \noexpand\newglossaryentry{Fig\thefigcounter}{%
        name={\unexpanded{#1}},%
        sort={\thefigcounter},%
        description={p.~\noexpand\pageref{fig:\thefigcounter}},%
        user1={\thefigcounter},%
        user2={Figure},%
        user3={\thesection},%
        user4={\expandafter\unexpanded\expandafter{\@savedSecTitle}}%
      }%
    }%
    \protected@write\@auxout{}{\string\glsliveentry{Fig\thefigcounter}}%
  \fi
}
\newcommand{\eeF}{\end{figure}}
\makeatother
% LIVE-ENTRY REGISTER
% \glsaddall prints every entry the glossaries package currently knows, and
% that set includes entries restored from \jobname.glsdefs.  The package
% rewrites .glsdefs from all *currently defined* entries, so anything it
% restored is written straight back out: a deleted section kept reappearing in
% the catalogue for ever, and the .glsdefs had to be erased by hand.  The
% wrappers therefore record every entry they actually create in the .aux, and
% \printboth adds only those.  The .aux is rebuilt from scratch on each run, so
% this is self-healing: a removed result is gone after one further run.
\makeatletter
\newcommand{\@glslive}{}
\newcommand{\@glsaddlive}[1]{\ifglsentryexists{#1}{\glsadd{#1}}{}}
\newcommand{\glsliveentry}[1]{\g@addto@macro\@glslive{\@glsaddlive{#1}}}
\makeatother
% PRINT COMMAND
% No \clearpage here -- the catalogue continues on the current page.  This has
% been re-added by mistake several times; leave it out.
\makeatletter
\newcommand{\printboth}{%
  \@glslive   % add exactly the entries this document created last run
  \printglossary[style=unifiedstyle, title={0 \hspace{.22em} Definitions, Theorems, Examples and Open Problems}]%
}
\makeatother
% Named theorem reference: \tref{word} -> Thm~N(word-or-short)
% Shows stored short title if available, otherwise shows the key itself.
\newcommand{\tref}[1]{%
  Thm~\ref{t:#1}(\ifcsname @tshort@#1\endcsname\csname @tshort@#1\endcsname\else#1\fi)%
}

% Full-title theorem reference: \nref{t:KEY} prints the number and then the
% theorem's long title, so one writes "Theorem~\nref{t:blocks}" and gets
% "Theorem 3 ``Block factorization and siphons''".  Takes the complete label, like
% \ref, not the bare key like \tref.
%
% \beT writes each title to the .aux (see \@thmtitle below), so \nref resolves
% forward references from the second LaTeX pass, exactly as \ref does.  If the
% title is not yet known -- first pass, or a label that is not a \beT theorem
% -- it degrades to the number rather than producing ??.
% These two use @ in their names, and this point in the file is outside the
% \makeatletter region opened earlier (it closes at \makeatother above), so
% the region has to be reopened here.
\makeatletter
\newcommand{\@thmtitle}[2]{\global\@namedef{@tlong@#1}{#2}}
\newcommand{\nref}[1]{%
  \ref{#1}%
  \ifcsname @tlong@#1\endcsname
    \ ``\csname @tlong@#1\endcsname''%
  \fi
}
% \nm{t:KEY}: the stored long title of any \beT/\beL/\beP/\beR/\beC/\beD/\beE result,
% verbatim and without quotes, so item names follow the macro's first argument.
\newcommand{\nm}[1]{%
  \ifcsname @tlong@#1\endcsname
    \csname @tlong@#1\endcsname%
  \else
    \ref{#1}%
  \fi
}
\makeatother




%% ========================================
%% Other THEOREM ENVIRONMENTS
%% ========================================
\makeatletter
\@ifundefined{lemma}      {\newtheorem{lemma}      {Lemma}}      {}
\@ifundefined{proposition}{\newtheorem{proposition}{Proposition}} {}
\@ifundefined{corollary}  {\newtheorem{corollary}  {Corollary}}  {}
\@ifundefined{conjecture} {\newtheorem{conjecture} {Conjecture}}  {}
\@ifundefined{remark}     {\newtheorem{remark}     {Remark}}     {}
\@ifundefined{question}   {\newtheorem{question}   {Question}}   {}
\@ifundefined{assumption} {\newtheorem{assumption} {Assumption}} {}
\@ifundefined{hypothesis} {\newtheorem{hypothesis} {Hypothesis}} {}
\@ifundefined{fact}       {\newtheorem{fact}       {Fact}}       {}
\makeatother

% \citer: cite the relay-graph threshold theorem of AAHrel. If a document has
% \input the chapter/file containing t:relayRF (so the label resolves), point
% straight at it; otherwise fall back to a bare citation of AAHrel. Shared here
% (not per-paper) so it resolves consistently across every book chapter.
% \citeR: same idea, for the first theorem of AAHRah (t:Pstar).
\makeatletter
\providecommand{\citer}{%
  \@ifundefined{r@t:relayRF}%
    {\cite{AAHrel}}%
    {Theorem~\ref{t:relayRF} of \cite{AAHrel}}%
}
\providecommand{\citeR}{%
  \@ifundefined{r@t:Pstar}%
    {\cite{AAHRah}}%
    {Theorem~\ref{t:Pstar} of \cite{AAHRah}}%
}
\makeatother


% Abbreviations
\def\beQ{\begin{question}}\def\eeQ{\end{question}}
\def\beA{\begin{assumption}}\def\eeA{\end{assumption}}


%% ========================================

%% ========================================
%% COLORS & CUSTOM COMMANDS
%% ========================================
\definecolor{funccolor}{RGB}{25,25,112}
\definecolor{desccolor}{RGB}{64,64,64}
\newcommand{\func}[1]{\texttt{\textcolor{funccolor}{#1}}}
\def\eq{=}
\def\mm{-}

%% ========================================
%% MATHCAL SHORTHANDS
%% ========================================
% Project-wide convention: \mX for a capital X is always \mathcal X.
% Defined for the whole alphabet so the convention never fails to expand.
\newcommand{\mA}{\mathcal A}\newcommand{\mB}{\mathcal B}
\newcommand{\mC}{\mathcal C}\newcommand{\mD}{\mathcal D}
\newcommand{\mE}{\mathcal E}\newcommand{\mF}{\mathcal F}
\newcommand{\mG}{\mathcal G}\newcommand{\mH}{\mathcal H}
\newcommand{\mI}{\mathcal I}\newcommand{\mJ}{\mathcal J}
\newcommand{\mK}{\mathcal K}\newcommand{\mL}{\mathcal L}
\newcommand{\mM}{\mathcal M}\newcommand{\mN}{\mathcal N}
\newcommand{\mO}{\mathcal O}\newcommand{\mP}{\mathcal P}
\newcommand{\mQ}{\mathcal Q}\newcommand{\mR}{\mathcal R}
\newcommand{\mS}{\mathcal S}\newcommand{\mT}{\mathcal T}
\newcommand{\mU}{\mathcal U}\newcommand{\mV}{\mathcal V}
\newcommand{\mW}{\mathcal W}\newcommand{\mX}{\mathcal X}
\newcommand{\mY}{\mathcal Y}\newcommand{\mZ}{\mathcal Z}

%% ========================================
%% BASIC MATH ENVIRONMENTS
%% ========================================
\def\bep{\begin{pmatrix}}\def\eep{\end{pmatrix}}
\def\bev{\begin{vmatrix}}\def\eev{\end{vmatrix}}
\def\bea{\begin{eqnarray*}}\def\eea{\end{eqnarray*}}
\def\bc{\begin{cases}}\def\ec{\end{cases}}
\def\BEN{\begin{enumerate}}\def\EEN{\end{enumerate}}
\def\BI{\begin{itemize}}\def\EI{\end{itemize}}

\newcommand{\be}[1]{\begin{equation}\label{#1}}
\newcommand{\ee}{\end{equation}}
\newcommand{\beq}{\begin{eqnarray}}
\def\eeq{\end{eqnarray}}
\def\eqr{\eqref}\def\fr{\frac}\def\lbl{\label}
\def\Lra{\Longrightarrow}\def\non{\nonumber}\def\cNGM{\cite{Diek,Van,Van08}}
\def\Eq{\Leftrightarrow}\def\b{\beta}\def\uph{\red{UP TO HERE}}
%% ========================================
%% PROBABILITY & STATISTICS
%% ========================================
\renewcommand{\P}{\ensuremath{\mathbb{P}}}  % Note: overrides paragraph symbol
\newcommand{\E}{\ensuremath{\mathbb{E}}}
\newcommand{\ind}[1]{\ensuremath{\mathbbm{1}_{\left\{#1\right\}}}}
\DeclareMathOperator{\Cov}{Cov}
\DeclareMathOperator{\TV}{TV}
\newcommand{\rank}{\operatorname{rank}}

%% ========================================
%% CRN OPERATORS (organizations/ERCs/buffering/siphons framework)
%% ========================================
\def\clos{\mathrm{clos}}
\def\Supply{\mathrm{Supply}}
\def\req{\mathrm{req}}
\def\React{\mathrm{React}}
\def\Prod{\mathrm{Prod}}
\def\Gen{\mathrm{Gen}}
\def\In{\mathrm{In}}
\def\Out{\mathrm{Out}}
\def\Act{\mathrm{Act}}
\def\Min{\mathrm{Min}}

%% ========================================
%% LINE BREAKING
%% ========================================
% Last-resort stretch, used only on paragraphs TeX cannot otherwise break:
% a long \texttt{} token, a comma-separated inline math list (commas are not
% breakpoints in math), or a run of \cite/\ref boxes.  Without it those set as
% overfull \hbox and print into the margin; with it they set slightly loose,
% which is the lesser evil.  Does not affect paragraphs that already break.
\emergencystretch=3em

%% ========================================
%% NUMBER SETS
%% ========================================
\newcommand{\R}{\mathbb{R}}  % Real numbers
\newcommand{\C}{\mathbb{C}}  % Complex numbers
\newcommand{\N}{\mathbb{N}}
\newcommand{\Z}{\mathbb{Z}}
% \providecommand, not \newcommand: some papers define these in their own
% preamble after \input{preamble}, and must not clash.
\providecommand{\Q}{\mathbb{Q}}
\providecommand{\res}{\operatorname{res}}   % resultant; used by epiRes.tex
\newcommand{\MN}{\mathbb{N}}
\newcommand{\MZ}{\mathbb{Z}}
\newcommand{\realpos}{\mathbb{R}_{> 0}}
\newcommand{\realnneg}{\mathbb{R}_{\geq 0}}
\newcommand{\naturalpos}{\mathbb{N}_{>0}}
\newcommand{\Rplus}{\mathbb{R}_+}
\newcommand{\Rp}{\mathbb{R}_{\ge 0}}
\def\Rnne{\mathbb{R}_{\ge 0}}
\def\Se{\bar{S}} \def\Ie{\bar{I}}
\newcommand{\Rpos}{\mathbb{R}_{>0}}
\newcommand{\dd}{\,\mathrm d}
\newcommand{\e}{\mathrm e}
%% ========================================
%% GREEK LETTERS (2+ character abbreviations) greek
%% ========================================
\def\al{\alpha}\def\ba{\bff \alpha}
\def\w{\omega}\def\Lgn{\ga+\mu+\de}\def\om{\omega}
\def\de{\delta}\def\De{\Delta}\def\bde{\bff \delta}
\def\ep{\epsilon}\def\eps{\varepsilon}
\def\f{\varphi}  % Kept as requested
\def\ga{\gamma}\def\Ga{\Gamma}\def\G{\Gamma}\def\bg{\bff \gamma}
\def\up{\upsilon}  % NOT the breve accent: no accent commands anywhere.
                   % Named \up, not \u, so the LaTeX breve stays reachable.
\def\la{\lambda}\def\La{\Lambda}
\def\si{\sigma}\def\ka{\kappa}
\def\th{\theta}\def\vt{\vartheta}
\def\ze{\zeta}\def\Gap{\Gamma^{+}} \def\Gam{\Gamma^{-}}


%% ========================================
%% SPECIAL SHORT COMMANDS (kept as requested)
%% ========================================
\def\I{\infty}
\def\T{\widetilde}
\def\H{\hat}
\newcommand{\relint}{\operatorname{relint}}

%% ========================================
%% BOLD/VECTOR NOTATION
%% ========================================
\newcommand{\bff}[1]{{\mbox{\boldmath$#1$}}}
\newcommand\X{\boldsymbol{X}}
\def\x{\boldsymbol{x}}\def\bx{\bff x}\def\xe{\bar{\x}}
\newcommand\y{\boldsymbol{y}}\def\by{\boldsymbol{y}}\def\ye{\bar{\y}}
\def\yb{\y}\def\xb{\x}
\def\xy{\x-\y}
\newcommand\br{\boldsymbol{r}}\def\re{\bar{\r}}
\newcommand\bd{\boldsymbol{d}}
\newcommand\bQ{\boldsymbol{Q}}
\newcommand{\Yb}{{\mathbf{Y}}}
\newcommand{\Ab}{\mathbf{A}}
\newcommand{\zb}{\mathbf{z}}\def\bz{\bff 0}
\newcommand{\Zb}{\mathbf{Z}}
\newcommand{\bb}{\mathbf{b}}
\def\vbe{\vec {\bff \beta}}
\def\vb{\vec {\bff b}}
\def\v1{\vec {\bff 1}}
\def\vn{\vec \de}
\def\vz{\vec 0}
\def\vi{\overset{\rightarrow}{\mathsf i}}
\def\vv{\overset{\rightarrow}{v}}
\def\vr{\; \vec {\mathsf r}}
\newcommand{\va}{\bff a}
\providecommand{\vect}[1]{\boldsymbol{#1}}   % bold vector (moved from Gradostat.tex)
\providecommand{\one}{\mathbf 1}             % all-ones vector (moved from AliSF)
\providecommand{\spr}{s}                     % spectral bound  (moved from AliSF)
\providecommand{\Ke}{K_{\rm eff}}            % effective carrying capacity
% Critical right/left eigenvectors in the Hopf projection method (AAHOSN/New.tex)
\providecommand{\bq}{\mathbf{q}}
\providecommand{\bp}{\mathbf{p}}


%% ========================================
%% MATHCAL & MATHSCR
%% ========================================
% Mathcal
\newcommand\mc[1]{\mathcal{#1}}
\def\mC{{\mathcal C}}\def\mD{\mathcal D}\def\mE{\mathcal{E}}
\def\mF{\mathcal F}\def\mG{{\mathcal G}}\def\mH{{\mathcal H}}
\def\mI{\mathcal I}\def\mL{\mathcal{L}}\def\mM{{\mathcal M}}
\def\mN{{\mathcal N}}\def\mP{\mathcal{P}}\def\mR{{\mathcal R}}
\def\mS{{\mathcal S}}\def\mT{\mathcal{T}}\def\mV{{\mathcal V}}
\def\mK{\mathcal K}\def\mA{{\mathcal A}}\def\mU{{\mathcal U}}
% Symmetric bilinear and trilinear forms of the Taylor expansion, used in the
% first Lyapunov coefficient (AAHOSN/New.tex)
\providecommand{\Bform}{\mathcal{B}}
\providecommand{\Cform}{\mathcal{C}}

% Blackboard bold.  Kept here, not in the drivers: a \providecommand in a
% driver preamble runs before this file and silently wins, so the same name
% ends up meaning different things in different builds.
\def\IC{\mathbb{C}}\def\IN{\mathbb{N}}\def\IQ{\mathbb{Q}}
\def\IR{\mathbb{R}}\def\IZ{\mathbb{Z}}\def\IG{\mathbb{G}}
\def\II{\mathbb{I}}

% Mathscr
\newcommand{\Rs}{{\mathscr R}}
\newcommand{\Ss}{{\mathscr S}}
\newcommand{\Ns}{{\mathscr N}}
\newcommand{\Cs}{\mathscr C}


%% ========================================
%% OPERATORS & FUNCTIONS
%% ========================================
\DeclareMathOperator{\tr}{tr}
\DeclareMathOperator{\diag}{Diag}
\DeclareMathOperator{\Diag}{Diag}
\DeclareMathOperator{\intr}{\mathrm{Int}}
\DeclareMathOperator{\Inv}{Inv}
\DeclareMathOperator{\spec}{spec}
\DeclareMathOperator{\AUC}{AUC}               % area under the curve
\newcommand{\ii}{\mathrm{i}}                  % imaginary unit
\newcommand\op{\operatorname}
\newcommand{\supp}{\op{supp}}
\newcommand{\diff}{\mathop{}\mathopen{}\mathrm{d}}
\newcommand{\pdd}[1]{\frac{\partial}{\partial{#1}}}
\DeclareMathOperator{\cone}{cone}
\DeclareMathOperator{\proj}{proj}

%% ========================================
%% BRACKETS & NORMS
%% ========================================
\providecommand{\pr}[1]{\left(#1\right)}
\providecommand{\pp}[1]{\left[#1\right]}
\newcommand\croc[1]{\left\langle #1\right\rangle}
\newcommand{\norm}[1]{\|#1\|}
\newcommand{\inner}[2]{\langle #1, #2 \rangle}
% Short synonym of \inner, used by AAHOSN/New.tex as $\ip{u}{v}=\bar u^\top v$
\providecommand{\ip}[2]{\langle #1 , #2 \rangle}


%% ========================================
%% FORMATTING & UTILITIES
%% ========================================
\newcommand*{\Scale}[2][4]{\scalebox{#1}{$#2$}}
\newcommand{\te}[1]{\text{#1}}
\newcommand\steq[1]{\stackrel{\text{\rm #1.}}{=}}
\newcommand\blue[1]{\textcolor{blue}{#1}}
\newcommand{\red}{\textcolor[rgb]{1.00,0.00,0.00}}
\newcommand{\green}{\textcolor[rgb]{0.00,1.00,0.00}}

% Figure command
\newcommand{\figu}[3]{%
\begin{figure}[H]
\centering
\includegraphics[scale=#3]{#1}
\caption{#2\label{f:#1}}
\end{figure}
}

% Footnote
\long\def\symbolfootnote[#1]#2{%
\begingroup
\def\thefootnote{\fnsymbol{footnote}}\footnote[#1]{#2}%
\endgroup}
\def\fn{\symbolfootnote}

\newcommand{\dsum}{\displaystyle\sum}

%% ========================================
%% CRN & EPIDEMIOLOGY SPECIFIC
%% ========================================
% CRN notation
\def\SS{\mathcal S}
\def\CC{\mathcal C}
\def\RR{\mathcal R}

% Reproduction numbers
\def\Rto{R_{1}^{E_2}}
\def\Rot{R_{2}^{E_1}}
\def\Roc{R_{1,c}}
\def\Rtc{R_{2,c}}
\def\rd{r_{0}}
\def\sd{s_{0}}
\def\dd{D_{0}}
\def\vd{v_{0}}
\def\zd{z_{0}}
\def\se{\sE}
\def\xd{x_{0}}
%% ========================================
%% SPECIAL VALUES & EQUILIBRIA
%% ========================================
\def\eE{\bar{e}}
\def\wE{\bar{w}}
\def\sE{s_{*}}
\def\Me{\T M_{ee}}
\def\Rnne{\mathbb{R}_{\ge 0}}

\newcommand{\ones}{\mathbf 1}


%% ========================================
%% ABBREVIATIONS (alphabetical by category)
%% ========================================

% -- General mathematical terms --
\def\adm{admissible}\def\admy{admissibility}
%\def\and{antisymmetric}% disabled 2026-09-30: it overrode LaTeX's author separator \and (printed "antisymmetric" between authors)
\def\ass{assumption}
\def\cadlag{c\`adl\`ag}
\def\ch{characteristic polynomial}
\def\chp{characteristic polynomial}
\def\coe{coefficient}
\def\comp{computation}
\def\con{condition}
\def\coo{coordinate}
\def\corr{corresponding}
\def\dec{decomposition}
\def\den{denominator}
\def\dete{deterministic}
\def\eig{eigenvalue}\def\eval{eigenvalue}
\def\eigv{eigenvector}\def\eiv{eigenvector}\def\evec{eigenvector}
\def\elt{element}
\def\enn{essentially nonnegative/positivity preserving}
\def\equ{equation}
\def\etal{et al.}
\def\fe{for example}
\def\fno{from now on}
\def\for{\forall}
\def\form{formula}
\def\Fr{Furthermore, }
\def\frt{furthermore}
\def\fun{function}
\def\gene{generalization}
\def\heu{heuristic}
\def\how{however}\def\How{However, }
\def\ie{i.e. }
\def\im{\item}
\def\isd{is defined by}
\def\itc{it is easy to check that}
\def\itf{it follows that}\def\Itf{It follows that}
\def\ith{it holds that}
\def\itm{it may be checked that}\def\Itm{It may be checked that}
\def\ito{it turns out that}\def\Ito{It turns out that}\def\Itou{It turns out useful to}
\def\ma{mass action}\def\MAR{mass action representation}
\def\mat{matrix}
\def\mbc{may be checked to be}
\def\mbw{may be written as}\def\mubw{must be written as}
\def\mpr{more precisely, }\def\Mp{More precisely, }
\def\nc{non-catalytic}
\def\nne{non-negative}\def\nny{nonnegativity}
\def\no{\nonumber}
\def\noni{non-infectious}
\def\nota{notation}
\def\np{\newpage}
\def\Nt{Note that}
\def\nt{non trivial}\def\nz{non zero}
\def\oth{otherwise}
\def\Oth{On the other hand}
\def\pa{parameter}\def\para{parameter}
\def\pol{polynomial}
\def\prob{probability}
\def\Prf{{\bf  Proof:}}\def\prf{{\bf  Proof:}}
\def\prop{proportional to}
\def\qu{\quad}
\def\QED{\hfill {$\square$}\goodbreak \medskip}
\def\rep{representation}
\def\resp{respectively}
\def\satd{satisfied}\def\satg{satisfying}\def\sats{satisfies}\def\saty{satisfy}
\def\ssec{\subsection}\def\sssec{\subsubsection}
\def\sef{see for example}
\def\sksy{skew-symmetric}
\def\sub{substrate}
\def\sys{system}
\def\var{variable}
\def\wk{well-known}
\def\wl{without loss of generality}\def\wlo{w.l.o.g.}
\def\wmc{We may check that}\def\Wmc{We may check that}
\def\wmr{we may rewrite}
\def\wms{we must show that}
\def\wrt{with respect to}

% -- CRN & Chemistry --
\newcommand\CRN{chemical reaction network}
\newcommand\CRNs{chemical reaction networks}
\newcommand\RN{reaction network}
\def\ACR{absolute concentration robustness}
\def\AMSD{acyclic siphon decomposition}
\def\DAE{differential algebraic equations}
\def\FHJ{Feinberg-Horn-Jackson}
\def\GMAK{generalized mass action kinetics}
\def\MM{Michaelis-Menten}
\def\sm{stoichiometric matrix}\def\Sm{stoichiometric matrix}
\def\WRZD{weakly reversible, zero deficiency}\def\aut{autocatalytic}

% -- Epidemiology --
\def\brn{basic reproduction number}\def\bzn{basic replacement number}
\def\com{compartment}
\def\dfe{disease free equilibrium}\def\DFE{disease free equilibrium}
\def\dwt{dwell times}\def\edt{expected dwell times}
\def\EM{epidemic model}
\def\endp{endemic fixed point}\def\enp{endemic point}\def\EE{endemic point}
\def\ME{mathematical epidemiology}
\def\MEE{eco-epidemiology}\def\EEP{eco-epidemiological}
\def\ece{eco-epidemiological model}
\def\mgm{multi group model}
\def\mtS{multi-strain model}\def\mSt{multi-strain model}
\def\mSi{minimal siphon}
\def\NGM{next generation matrix}
\def\nim{new infections matrix}
\def\qem{quadratic epidemiological model}
\def\repF{invasion functions}
\def\SPF{SIR-PH-FA model}
\def\sus{susceptible}

% -- Graph theory & Linear algebra --
\def\jac{Jacobi matrix}
\def\jin{Jacobian matrix of the invasion vector field with respect to the invasion variables}
\def\jrn{Jacobian matrix of the resident vector field with respect to the resident variables}
\def\jxn{Jacobian matrix of the $\x$ vector field with respect to the $\x$ variables}
\def\Lm{Laplacian matrix}
\def\MTT{matrix tree theorem}
\def\SGT{spectral graph theory}
\def\SMW{Sherman–Morrison–Woodbury}

% -- Stability & Dynamics --
\def\absC{absorbtion CTMC}
\def\BTb{Bogdanov-Takens bifurcation}
\def\CEP{competitive exclusion partition}
\def\CTMC{continuous time Markov chain}
\def\ds{differential system}
\def\fp{fixed point}\def\bfp{boundary fixed point}\def\fbp{fixed boundary point}
\def\fbl{fixed boundary line}\def\fif{forward invariant face}
\def\gas{global asymptotic stability}\def\GAS{globally asymptotically stable}
\def\GASS{\text{GASS}}
\def\Hb{Hopf bifurcation}
\def\Hur{Hurwitz}\def\MH{Metzler-Hurwitz}\def\RH{Routh-Hurwitz}
\def\IS{\text{IS}}
\def\LAS{locally asymptotically stable}
\def\LdS{Lyapunov diagonally stable}
\def\Lf{Lyapunov function}\def\LyV{Lyapunov Volterra type function}\def\RLF{robust Lyapunov function}
\def\LyM{Lyapunov Malkin stable}
\def\PF{Perron-Frobenius}
\def\spfp{strictly positive fixed point}
\def\sta{stability}
\def\STP{sharp threshold property}\def\STT{sharp threshold theorem}

% -- Models (with \varphi explicitly instead of \f) --
\def\AB{$(A,B)$ Arino-Brauer epidemic models}
\def\ABa{$(A,\bbe,\ba,\varphi)$ balanced bilinear model}
\def\ABP{%$(A,\bbe,P,\varphi)$
\\ balanced bilinear model}
\def\AMG{Anderson May Gupta}
\newcommand\BIN{biological interaction network}
\def\Dt{Descartes type}
\def\gLV{generalized Lotka-Volterra model}
\def\IG{\text{invasion graph}}\def\IGMS{interaction graph on minimal siphons}
\def\LCP{\text{linear complementarity problem}}
\def\LV{Lotka-Volterra}\def\LVK{Lotka-Volterra-Kolmogorov}

% -- Splitting & Structure --
\def\reg{regular splitting}\def\regS{{\bf regular splitting}}
\def\RUR{rational univariate representation}
\def\bds{block-diagonal structure}\def\NGMs{next generation matrices}

% -- Miscellaneous --
\def\Ma{{Mathematica}}
\def\Mm{Max–Min}
\def\num{numerator}
\def\rap{rational approximation} \def\DSR{directed species-reaction bipartite graph}\def\PCP{polynomial complementarity problem}\def\rat{rational substitution}
\def\QSS{quasi steady state} \def\RHS{right hand side} \def\LHS{left hand side}
\def\OSN{online social network}\def\VdD{Van den Driessche}
\def\wli{$\omega$-limit} \def\rk{rich kinetics}\def\sJ{symbolic Jacobian}
\def\cV{\cite{Vas,VasSta,BSV,GSV}} \def\spf{semi-parametric semi-functional}
\def\CRW{Capasso-Ruan-Wang} \def\ABP{$(A,\bbe,P,\f)$ balanced bilinear model}
\def\ABa{%$(A,\bbe,\ba,\f)$ \\
rank-one balanced bilinear model}
\def\bbe{B%{\bff \beta}
} \def\vbe{\vec {\bff \beta}
}
\def\vb{\vec {\bff b}}
\newcommand{\Kh}{\hat{K}}          % K with hat, first mention only
\newcommand{\wS}{w^{\sigma}}       % right Perron J^\perp, full block
\newcommand{\piS}{\vec \pi^{\sigma}}    % left Perron J^\perp, full block
\newcommand{\ws}{w^{\sigma}}       % right Perron J^\perp, sub-block
\newcommand{\pis}{\pi^{\sigma}}    % left Perron J^\perp, sub-block
\newcommand{\hwS}{\hat{w}^{\sigma}}   % right Perron K, full block
\newcommand{\hpiS}{\hat{\pi}^{\sigma}} % left Perron K, full block
\newcommand{\hws}{\hat{w}^{\sigma}}   % right Perron K, sub-block
\newcommand{\hpis}{\hat{\pi}^{\sigma}} % left Perron K, sub-block
\newcommand{\twS}{\tilde{w}^{\sigma}}  % right Perron \tilde K, full block
\newcommand{\tpiS}{\tilde{\pi}^{\sigma}} % left Perron \tilde K
\newcommand{\tws}{\tilde{w}^{\sigma}}
\newcommand{\tpis}{\tilde{\pi}^{\sigma}}
\newcommand{\Mat}[2]{\mathbb{R}_{\ge 0}^{\,#1\times #2}}
\def\tJ{J_\sigma^\perp} \def\lPe{left Perron eigenvector}
% \G is \Gamma (defined at :795).  The Volterra function keeps its own name,
% so that \G no longer has two conflicting meanings in this file.
\def\GV{G}  % Volterra G(u)=u-1-\ln u
\def\PVL{Perron-Volterra Lyapunov minimal siphon ansatz}
\newcommand{\istar}[1]{i_{#1}^{*}}
%\newcommand{\sE}{s_E}
\newcommand{\vE}{v_E}
\def\PV{Perron-Volterra candidate Lyapunov function}
\def\GLV{generalized Lotka-Volterra}
\def\MERN{mathematical epidemiology reaction networks}
\def\CMM{contraction metric method}
\newcommand{\sign}{\operatorname{sign}}
\newcommand{\Var}{\operatorname{Var}} 

% ======================================================================
% IDEAS SHARED BETWEEN THE ABSTRACT AND THE MAIN RESULTS SECTION
%
% Write each idea ONCE, in the abstract; the Main Results section recalls it,
% so the two can never drift apart.
%
%   in the abstract        \ideaitem{1}{The first-order reaction terms ...}
%   in Main Results        \ideahead{1}
%
% \ideaitem stores the text globally and typesets
%   \item \hyperlink{idea:1}{<text>}
% \ideahead typesets
%   \item\hypertarget{idea:1}{}\textbf{<the same text>}
%
% The abstract must come before the Main Results input, which is the normal
% order. If a heading has to read differently from the abstract sentence (for
% instance a long abstract item summarised in one line), override it with the
% optional argument: \ideahead[short heading]{6}.
% ======================================================================
\makeatletter
\newcommand{\ideaitem}[2]{%
  \expandafter\gdef\csname @idea@#1\endcsname{#2}%
  \item\hyperlink{idea:#1}{#2}}
\newcommand{\ideahead}[2][]{%
  \item\hypertarget{idea:#2}{}%
  \def\@ideatmp{#1}%
  \ifx\@ideatmp\@empty
    \@ifundefined{@idea@#2}{%
      \PackageWarning{preamble}{Idea #2 is used in Main Results but was never
        defined by \string\ideaitem\space in the abstract}%
      \textbf{[idea #2 undefined]}%
    }{\textbf{\begingroup\def\ideaterm##1##2{##2}\def\idearest##1{}\@nameuse{@idea@#2}\endgroup}}%
  \else
    \textbf{#1}%
  \fi}
% \ideaterm{<hypertarget>}{<words>}: inside an abstract idea the words become an
% invisible link (the document sets \hypersetup{hidelinks}) jumping to the place
% in Main Results that defines them; inside the Main Results heading itself,
% where the target sits a few lines below, \ideahead strips the link and prints
% the words plainly. Mark the destination with \hypertarget{<name>}{} on the line
% before the display, which adds no equation number and shows nothing.
\newcommand{\ideaterm}[2]{\hyperlink{#1}{#2}}
% \idearest{<text>}: the part of an abstract idea after its first sentence; printed in
% the abstract, omitted from the Main Results heading.
\newcommand{\idearest}[1]{#1}
\makeatother

% ======================================================================
% SHORT-PAPER SWITCH
% \shorttrue in a master that assembles only part of the companion paper
% (Barrat/SipArxiv.tex). \nsh{...} drops text that cites results living in the
% omitted sections, so the short build has no dangling references while the
% companion paper keeps the text verbatim.
% ======================================================================
\newif\ifshort \shortfalse
\newcommand{\nsh}[1]{\ifshort\else#1\fi}
% \shortalt{<short build>}{<companion paper>}: use when a sentence needs
% different wording in the two builds, not merely a deletion.
\newcommand{\shortalt}[2]{\ifshort#1\else#2\fi}
% \Jtr, \JtrS{\Sigma}: transverse Jacobian. J^\perp here; SipArxiv.tex redefines them as J_{\mathrm{tra}}, J_{\mathrm{tra},\Sigma}
\newcommand{\Jtr}{J^\perp}
\newcommand{\JtrS}[1]{J_{#1}^\perp}
, and must not clash.
\providecommand{\Q}{\mathbb{Q}}
\providecommand{\res}{\operatorname{res}}   % resultant; used by epiRes.tex
\newcommand{\MN}{\mathbb{N}}
\newcommand{\MZ}{\mathbb{Z}}
\newcommand{\realpos}{\mathbb{R}_{> 0}}
\newcommand{\realnneg}{\mathbb{R}_{\geq 0}}
\newcommand{\naturalpos}{\mathbb{N}_{>0}}
\newcommand{\Rplus}{\mathbb{R}_+}
\newcommand{\Rp}{\mathbb{R}_{\ge 0}}
\def\Rnne{\mathbb{R}_{\ge 0}}
\def\Se{\bar{S}} \def\Ie{\bar{I}}
\newcommand{\Rpos}{\mathbb{R}_{>0}}
\newcommand{\dd}{\,\mathrm d}
\newcommand{\e}{\mathrm e}
%% ========================================
%% GREEK LETTERS (2+ character abbreviations) greek
%% ========================================
\def\al{\alpha}\def\ba{\bff \alpha}
\def\w{\omega}\def\Lgn{\ga+\mu+\de}\def\om{\omega}
\def\de{\delta}\def\De{\Delta}\def\bde{\bff \delta}
\def\ep{\epsilon}\def\eps{\varepsilon}
\def\f{\varphi}  % Kept as requested
\def\ga{\gamma}\def\Ga{\Gamma}\def\G{\Gamma}\def\bg{\bff \gamma}
\def\up{\upsilon}  % NOT the breve accent: no accent commands anywhere.
                   % Named \up, not \u, so the LaTeX breve stays reachable.
\def\la{\lambda}\def\La{\Lambda}
\def\si{\sigma}\def\ka{\kappa}
\def\th{\theta}\def\vt{\vartheta}
\def\ze{\zeta}\def\Gap{\Gamma^{+}} \def\Gam{\Gamma^{-}}


%% ========================================
%% SPECIAL SHORT COMMANDS (kept as requested)
%% ========================================
\def\I{\infty}
\def\T{\widetilde}
\def\H{\hat}
\newcommand{\relint}{\operatorname{relint}}

%% ========================================
%% BOLD/VECTOR NOTATION
%% ========================================
\newcommand{\bff}[1]{{\mbox{\boldmath$#1$}}}
\newcommand\X{\boldsymbol{X}}
\def\x{\boldsymbol{x}}\def\bx{\bff x}\def\xe{\bar{\x}}
\newcommand\y{\boldsymbol{y}}\def\by{\boldsymbol{y}}\def\ye{\bar{\y}}
\def\yb{\y}\def\xb{\x}
\def\xy{\x-\y}
\newcommand\br{\boldsymbol{r}}\def\re{\bar{\r}}
\newcommand\bd{\boldsymbol{d}}
\newcommand\bQ{\boldsymbol{Q}}
\newcommand{\Yb}{{\mathbf{Y}}}
\newcommand{\Ab}{\mathbf{A}}
\newcommand{\zb}{\mathbf{z}}\def\bz{\bff 0}
\newcommand{\Zb}{\mathbf{Z}}
\newcommand{\bb}{\mathbf{b}}
\def\vbe{\vec {\bff \beta}}
\def\vb{\vec {\bff b}}
\def\v1{\vec {\bff 1}}
\def\vn{\vec \de}
\def\vz{\vec 0}
\def\vi{\overset{\rightarrow}{\mathsf i}}
\def\vv{\overset{\rightarrow}{v}}
\def\vr{\; \vec {\mathsf r}}
\newcommand{\va}{\bff a}
\providecommand{\vect}[1]{\boldsymbol{#1}}   % bold vector (moved from Gradostat.tex)
\providecommand{\one}{\mathbf 1}             % all-ones vector (moved from AliSF)
\providecommand{\spr}{s}                     % spectral bound  (moved from AliSF)
\providecommand{\Ke}{K_{\rm eff}}            % effective carrying capacity
% Critical right/left eigenvectors in the Hopf projection method (AAHOSN/New.tex)
\providecommand{\bq}{\mathbf{q}}
\providecommand{\bp}{\mathbf{p}}


%% ========================================
%% MATHCAL & MATHSCR
%% ========================================
% Mathcal
\newcommand\mc[1]{\mathcal{#1}}
\def\mC{{\mathcal C}}\def\mD{\mathcal D}\def\mE{\mathcal{E}}
\def\mF{\mathcal F}\def\mG{{\mathcal G}}\def\mH{{\mathcal H}}
\def\mI{\mathcal I}\def\mL{\mathcal{L}}\def\mM{{\mathcal M}}
\def\mN{{\mathcal N}}\def\mP{\mathcal{P}}\def\mR{{\mathcal R}}
\def\mS{{\mathcal S}}\def\mT{\mathcal{T}}\def\mV{{\mathcal V}}
\def\mK{\mathcal K}\def\mA{{\mathcal A}}\def\mU{{\mathcal U}}
% Symmetric bilinear and trilinear forms of the Taylor expansion, used in the
% first Lyapunov coefficient (AAHOSN/New.tex)
\providecommand{\Bform}{\mathcal{B}}
\providecommand{\Cform}{\mathcal{C}}

% Blackboard bold.  Kept here, not in the drivers: a \providecommand in a
% driver preamble runs before this file and silently wins, so the same name
% ends up meaning different things in different builds.
\def\IC{\mathbb{C}}\def\IN{\mathbb{N}}\def\IQ{\mathbb{Q}}
\def\IR{\mathbb{R}}\def\IZ{\mathbb{Z}}\def\IG{\mathbb{G}}
\def\II{\mathbb{I}}

% Mathscr
\newcommand{\Rs}{{\mathscr R}}
\newcommand{\Ss}{{\mathscr S}}
\newcommand{\Ns}{{\mathscr N}}
\newcommand{\Cs}{\mathscr C}


%% ========================================
%% OPERATORS & FUNCTIONS
%% ========================================
\DeclareMathOperator{\tr}{tr}
\DeclareMathOperator{\diag}{Diag}
\DeclareMathOperator{\Diag}{Diag}
\DeclareMathOperator{\intr}{\mathrm{Int}}
\DeclareMathOperator{\Inv}{Inv}
\DeclareMathOperator{\spec}{spec}
\DeclareMathOperator{\AUC}{AUC}               % area under the curve
\newcommand{\ii}{\mathrm{i}}                  % imaginary unit
\newcommand\op{\operatorname}
\newcommand{\supp}{\op{supp}}
\newcommand{\diff}{\mathop{}\mathopen{}\mathrm{d}}
\newcommand{\pdd}[1]{\frac{\partial}{\partial{#1}}}
\DeclareMathOperator{\cone}{cone}
\DeclareMathOperator{\proj}{proj}

%% ========================================
%% BRACKETS & NORMS
%% ========================================
\providecommand{\pr}[1]{\left(#1\right)}
\providecommand{\pp}[1]{\left[#1\right]}
\newcommand\croc[1]{\left\langle #1\right\rangle}
\newcommand{\norm}[1]{\|#1\|}
\newcommand{\inner}[2]{\langle #1, #2 \rangle}
% Short synonym of \inner, used by AAHOSN/New.tex as $\ip{u}{v}=\bar u^\top v$
\providecommand{\ip}[2]{\langle #1 , #2 \rangle}


%% ========================================
%% FORMATTING & UTILITIES
%% ========================================
\newcommand*{\Scale}[2][4]{\scalebox{#1}{$#2$}}
\newcommand{\te}[1]{\text{#1}}
\newcommand\steq[1]{\stackrel{\text{\rm #1.}}{=}}
\newcommand\blue[1]{\textcolor{blue}{#1}}
\newcommand{\red}{\textcolor[rgb]{1.00,0.00,0.00}}
\newcommand{\green}{\textcolor[rgb]{0.00,1.00,0.00}}

% Figure command
\newcommand{\figu}[3]{%
\begin{figure}[H]
\centering
\includegraphics[scale=#3]{#1}
\caption{#2\label{f:#1}}
\end{figure}
}

% Footnote
\long\def\symbolfootnote[#1]#2{%
\begingroup
\def\thefootnote{\fnsymbol{footnote}}\footnote[#1]{#2}%
\endgroup}
\def\fn{\symbolfootnote}

\newcommand{\dsum}{\displaystyle\sum}

%% ========================================
%% CRN & EPIDEMIOLOGY SPECIFIC
%% ========================================
% CRN notation
\def\SS{\mathcal S}
\def\CC{\mathcal C}
\def\RR{\mathcal R}

% Reproduction numbers
\def\Rto{R_{1}^{E_2}}
\def\Rot{R_{2}^{E_1}}
\def\Roc{R_{1,c}}
\def\Rtc{R_{2,c}}
\def\rd{r_{0}}
\def\sd{s_{0}}
\def\dd{D_{0}}
\def\vd{v_{0}}
\def\zd{z_{0}}
\def\se{\sE}
\def\xd{x_{0}}
%% ========================================
%% SPECIAL VALUES & EQUILIBRIA
%% ========================================
\def\eE{\bar{e}}
\def\wE{\bar{w}}
\def\sE{s_{*}}
\def\Me{\T M_{ee}}
\def\Rnne{\mathbb{R}_{\ge 0}}

\newcommand{\ones}{\mathbf 1}


%% ========================================
%% ABBREVIATIONS (alphabetical by category)
%% ========================================

% -- General mathematical terms --
\def\adm{admissible}\def\admy{admissibility}
%\def\and{antisymmetric}% disabled 2026-09-30: it overrode LaTeX's author separator \and (printed "antisymmetric" between authors)
\def\ass{assumption}
\def\cadlag{c\`adl\`ag}
\def\ch{characteristic polynomial}
\def\chp{characteristic polynomial}
\def\coe{coefficient}
\def\comp{computation}
\def\con{condition}
\def\coo{coordinate}
\def\corr{corresponding}
\def\dec{decomposition}
\def\den{denominator}
\def\dete{deterministic}
\def\eig{eigenvalue}\def\eval{eigenvalue}
\def\eigv{eigenvector}\def\eiv{eigenvector}\def\evec{eigenvector}
\def\elt{element}
\def\enn{essentially nonnegative/positivity preserving}
\def\equ{equation}
\def\etal{et al.}
\def\fe{for example}
\def\fno{from now on}
\def\for{\forall}
\def\form{formula}
\def\Fr{Furthermore, }
\def\frt{furthermore}
\def\fun{function}
\def\gene{generalization}
\def\heu{heuristic}
\def\how{however}\def\How{However, }
\def\ie{i.e. }
\def\im{\item}
\def\isd{is defined by}
\def\itc{it is easy to check that}
\def\itf{it follows that}\def\Itf{It follows that}
\def\ith{it holds that}
\def\itm{it may be checked that}\def\Itm{It may be checked that}
\def\ito{it turns out that}\def\Ito{It turns out that}\def\Itou{It turns out useful to}
\def\ma{mass action}\def\MAR{mass action representation}
\def\mat{matrix}
\def\mbc{may be checked to be}
\def\mbw{may be written as}\def\mubw{must be written as}
\def\mpr{more precisely, }\def\Mp{More precisely, }
\def\nc{non-catalytic}
\def\nne{non-negative}\def\nny{nonnegativity}
\def\no{\nonumber}
\def\noni{non-infectious}
\def\nota{notation}
\def\np{\newpage}
\def\Nt{Note that}
\def\nt{non trivial}\def\nz{non zero}
\def\oth{otherwise}
\def\Oth{On the other hand}
\def\pa{parameter}\def\para{parameter}
\def\pol{polynomial}
\def\prob{probability}
\def\Prf{{\bf  Proof:}}\def\prf{{\bf  Proof:}}
\def\prop{proportional to}
\def\qu{\quad}
\def\QED{\hfill {$\square$}\goodbreak \medskip}
\def\rep{representation}
\def\resp{respectively}
\def\satd{satisfied}\def\satg{satisfying}\def\sats{satisfies}\def\saty{satisfy}
\def\ssec{\subsection}\def\sssec{\subsubsection}
\def\sef{see for example}
\def\sksy{skew-symmetric}
\def\sub{substrate}
\def\sys{system}
\def\var{variable}
\def\wk{well-known}
\def\wl{without loss of generality}\def\wlo{w.l.o.g.}
\def\wmc{We may check that}\def\Wmc{We may check that}
\def\wmr{we may rewrite}
\def\wms{we must show that}
\def\wrt{with respect to}

% -- CRN & Chemistry --
\newcommand\CRN{chemical reaction network}
\newcommand\CRNs{chemical reaction networks}
\newcommand\RN{reaction network}
\def\ACR{absolute concentration robustness}
\def\AMSD{acyclic siphon decomposition}
\def\DAE{differential algebraic equations}
\def\FHJ{Feinberg-Horn-Jackson}
\def\GMAK{generalized mass action kinetics}
\def\MM{Michaelis-Menten}
\def\sm{stoichiometric matrix}\def\Sm{stoichiometric matrix}
\def\WRZD{weakly reversible, zero deficiency}\def\aut{autocatalytic}

% -- Epidemiology --
\def\brn{basic reproduction number}\def\bzn{basic replacement number}
\def\com{compartment}
\def\dfe{disease free equilibrium}\def\DFE{disease free equilibrium}
\def\dwt{dwell times}\def\edt{expected dwell times}
\def\EM{epidemic model}
\def\endp{endemic fixed point}\def\enp{endemic point}\def\EE{endemic point}
\def\ME{mathematical epidemiology}
\def\MEE{eco-epidemiology}\def\EEP{eco-epidemiological}
\def\ece{eco-epidemiological model}
\def\mgm{multi group model}
\def\mtS{multi-strain model}\def\mSt{multi-strain model}
\def\mSi{minimal siphon}
\def\NGM{next generation matrix}
\def\nim{new infections matrix}
\def\qem{quadratic epidemiological model}
\def\repF{invasion functions}
\def\SPF{SIR-PH-FA model}
\def\sus{susceptible}

% -- Graph theory & Linear algebra --
\def\jac{Jacobi matrix}
\def\jin{Jacobian matrix of the invasion vector field with respect to the invasion variables}
\def\jrn{Jacobian matrix of the resident vector field with respect to the resident variables}
\def\jxn{Jacobian matrix of the $\x$ vector field with respect to the $\x$ variables}
\def\Lm{Laplacian matrix}
\def\MTT{matrix tree theorem}
\def\SGT{spectral graph theory}
\def\SMW{Sherman–Morrison–Woodbury}

% -- Stability & Dynamics --
\def\absC{absorbtion CTMC}
\def\BTb{Bogdanov-Takens bifurcation}
\def\CEP{competitive exclusion partition}
\def\CTMC{continuous time Markov chain}
\def\ds{differential system}
\def\fp{fixed point}\def\bfp{boundary fixed point}\def\fbp{fixed boundary point}
\def\fbl{fixed boundary line}\def\fif{forward invariant face}
\def\gas{global asymptotic stability}\def\GAS{globally asymptotically stable}
\def\GASS{\text{GASS}}
\def\Hb{Hopf bifurcation}
\def\Hur{Hurwitz}\def\MH{Metzler-Hurwitz}\def\RH{Routh-Hurwitz}
\def\IS{\text{IS}}
\def\LAS{locally asymptotically stable}
\def\LdS{Lyapunov diagonally stable}
\def\Lf{Lyapunov function}\def\LyV{Lyapunov Volterra type function}\def\RLF{robust Lyapunov function}
\def\LyM{Lyapunov Malkin stable}
\def\PF{Perron-Frobenius}
\def\spfp{strictly positive fixed point}
\def\sta{stability}
\def\STP{sharp threshold property}\def\STT{sharp threshold theorem}

% -- Models (with \varphi explicitly instead of \f) --
\def\AB{$(A,B)$ Arino-Brauer epidemic models}
\def\ABa{$(A,\bbe,\ba,\varphi)$ balanced bilinear model}
\def\ABP{%$(A,\bbe,P,\varphi)$
\\ balanced bilinear model}
\def\AMG{Anderson May Gupta}
\newcommand\BIN{biological interaction network}
\def\Dt{Descartes type}
\def\gLV{generalized Lotka-Volterra model}
\def\IG{\text{invasion graph}}\def\IGMS{interaction graph on minimal siphons}
\def\LCP{\text{linear complementarity problem}}
\def\LV{Lotka-Volterra}\def\LVK{Lotka-Volterra-Kolmogorov}

% -- Splitting & Structure --
\def\reg{regular splitting}\def\regS{{\bf regular splitting}}
\def\RUR{rational univariate representation}
\def\bds{block-diagonal structure}\def\NGMs{next generation matrices}

% -- Miscellaneous --
\def\Ma{{Mathematica}}
\def\Mm{Max–Min}
\def\num{numerator}
\def\rap{rational approximation} \def\DSR{directed species-reaction bipartite graph}\def\PCP{polynomial complementarity problem}\def\rat{rational substitution}
\def\QSS{quasi steady state} \def\RHS{right hand side} \def\LHS{left hand side}
\def\OSN{online social network}\def\VdD{Van den Driessche}
\def\wli{$\omega$-limit} \def\rk{rich kinetics}\def\sJ{symbolic Jacobian}
\def\cV{\cite{Vas,VasSta,BSV,GSV}} \def\spf{semi-parametric semi-functional}
\def\CRW{Capasso-Ruan-Wang} \def\ABP{$(A,\bbe,P,\f)$ balanced bilinear model}
\def\ABa{%$(A,\bbe,\ba,\f)$ \\
rank-one balanced bilinear model}
\def\bbe{B%{\bff \beta}
} \def\vbe{\vec {\bff \beta}
}
\def\vb{\vec {\bff b}}
\newcommand{\Kh}{\hat{K}}          % K with hat, first mention only
\newcommand{\wS}{w^{\sigma}}       % right Perron J^\perp, full block
\newcommand{\piS}{\vec \pi^{\sigma}}    % left Perron J^\perp, full block
\newcommand{\ws}{w^{\sigma}}       % right Perron J^\perp, sub-block
\newcommand{\pis}{\pi^{\sigma}}    % left Perron J^\perp, sub-block
\newcommand{\hwS}{\hat{w}^{\sigma}}   % right Perron K, full block
\newcommand{\hpiS}{\hat{\pi}^{\sigma}} % left Perron K, full block
\newcommand{\hws}{\hat{w}^{\sigma}}   % right Perron K, sub-block
\newcommand{\hpis}{\hat{\pi}^{\sigma}} % left Perron K, sub-block
\newcommand{\twS}{\tilde{w}^{\sigma}}  % right Perron \tilde K, full block
\newcommand{\tpiS}{\tilde{\pi}^{\sigma}} % left Perron \tilde K
\newcommand{\tws}{\tilde{w}^{\sigma}}
\newcommand{\tpis}{\tilde{\pi}^{\sigma}}
\newcommand{\Mat}[2]{\mathbb{R}_{\ge 0}^{\,#1\times #2}}
\def\tJ{J_\sigma^\perp} \def\lPe{left Perron eigenvector}
% \G is \Gamma (defined at :795).  The Volterra function keeps its own name,
% so that \G no longer has two conflicting meanings in this file.
\def\GV{G}  % Volterra G(u)=u-1-\ln u
\def\PVL{Perron-Volterra Lyapunov minimal siphon ansatz}
\newcommand{\istar}[1]{i_{#1}^{*}}
%\newcommand{\sE}{s_E}
\newcommand{\vE}{v_E}
\def\PV{Perron-Volterra candidate Lyapunov function}
\def\GLV{generalized Lotka-Volterra}
\def\MERN{mathematical epidemiology reaction networks}
\def\CMM{contraction metric method}
\newcommand{\sign}{\operatorname{sign}}
\newcommand{\Var}{\operatorname{Var}} 

% ======================================================================
% IDEAS SHARED BETWEEN THE ABSTRACT AND THE MAIN RESULTS SECTION
%
% Write each idea ONCE, in the abstract; the Main Results section recalls it,
% so the two can never drift apart.
%
%   in the abstract        \ideaitem{1}{The first-order reaction terms ...}
%   in Main Results        \ideahead{1}
%
% \ideaitem stores the text globally and typesets
%   \item \hyperlink{idea:1}{<text>}
% \ideahead typesets
%   \item\hypertarget{idea:1}{}\textbf{<the same text>}
%
% The abstract must come before the Main Results input, which is the normal
% order. If a heading has to read differently from the abstract sentence (for
% instance a long abstract item summarised in one line), override it with the
% optional argument: \ideahead[short heading]{6}.
% ======================================================================
\makeatletter
\newcommand{\ideaitem}[2]{%
  \expandafter\gdef\csname @idea@#1\endcsname{#2}%
  \item\hyperlink{idea:#1}{#2}}
\newcommand{\ideahead}[2][]{%
  \item\hypertarget{idea:#2}{}%
  \def\@ideatmp{#1}%
  \ifx\@ideatmp\@empty
    \@ifundefined{@idea@#2}{%
      \PackageWarning{preamble}{Idea #2 is used in Main Results but was never
        defined by \string\ideaitem\space in the abstract}%
      \textbf{[idea #2 undefined]}%
    }{\textbf{\begingroup\def\ideaterm##1##2{##2}\def\idearest##1{}\@nameuse{@idea@#2}\endgroup}}%
  \else
    \textbf{#1}%
  \fi}
% \ideaterm{<hypertarget>}{<words>}: inside an abstract idea the words become an
% invisible link (the document sets \hypersetup{hidelinks}) jumping to the place
% in Main Results that defines them; inside the Main Results heading itself,
% where the target sits a few lines below, \ideahead strips the link and prints
% the words plainly. Mark the destination with \hypertarget{<name>}{} on the line
% before the display, which adds no equation number and shows nothing.
\newcommand{\ideaterm}[2]{\hyperlink{#1}{#2}}
% \idearest{<text>}: the part of an abstract idea after its first sentence; printed in
% the abstract, omitted from the Main Results heading.
\newcommand{\idearest}[1]{#1}
\makeatother

% ======================================================================
% SHORT-PAPER SWITCH
% \shorttrue in a master that assembles only part of the companion paper
% (Barrat/SipArxiv.tex). \nsh{...} drops text that cites results living in the
% omitted sections, so the short build has no dangling references while the
% companion paper keeps the text verbatim.
% ======================================================================
\newif\ifshort \shortfalse
\newcommand{\nsh}[1]{\ifshort\else#1\fi}
% \shortalt{<short build>}{<companion paper>}: use when a sentence needs
% different wording in the two builds, not merely a deletion.
\newcommand{\shortalt}[2]{\ifshort#1\else#2\fi}
% \Jtr, \JtrS{\Sigma}: transverse Jacobian. J^\perp here; SipArxiv.tex redefines them as J_{\mathrm{tra}}, J_{\mathrm{tra},\Sigma}
\newcommand{\Jtr}{J^\perp}
\newcommand{\JtrS}[1]{J_{#1}^\perp}
.
\makeatletter
\@ifpackageloaded{hyperref}{}{%
  \@ifundefined{hyperrefpackage}{%
    \@ifundefined{AAdefNnohyperref}{\usepackage{hyperref}}{}%
  }{}%
}
\makeatother
% linktoc=all makes a table-of-contents page number a link as well as the
% entry text; hyperref's default, linktoc=section, links only the text.
% Deferred to begin-document and guarded: by the block above hyperref may
% be loaded later (mdpi's own hook) or not at all (\AAdefNnohyperref), and
% \hypersetup must exist before it is called.  A driver that wants the
% default back can say \hypersetup{linktoc=section} after % Shared preamble for the GLVbook papers.  This is the single canonical copy;
% a driver one level below GLVbook reaches it as
%
%     \input{../basicFiles/preamble}
%
% What stays in the driver, because it differs from paper to paper:
% \documentclass, font choice, \title, \author, \date, and any macro genuinely
% local to one paper.  Everything else belongs here.  The file is
% self-contained: it does not read any other file.

\usepackage[utf8]{inputenc}
\usepackage[novbox]{pdfsync}   % required: SyncTeX alone does not work here

% Page geometry, set once here for every paper that loads this file.  A driver
% that needs different margins says \geometry{...}; a second
% \usepackage[...]{geometry} raises an option clash and must not be used.
\makeatletter
\@ifundefined{AAdefNnogeometry}{\usepackage[a4paper,margin=2cm]{geometry}}{}
\makeatother

% mathematics
\usepackage{amsmath,amssymb,amsthm,mathtools,bm,mathrsfs,bbm}
\usepackage{centernot}

% floats, tables, lists
\usepackage{graphicx,float}
\makeatletter
\@ifpackageloaded{subfig}{}{\usepackage{subcaption}}
\makeatother
\usepackage{booktabs,longtable,threeparttable,array,multirow}
\usepackage{enumitem}

% colour, code, algorithms
\usepackage{xcolor,listings,tcolorbox}
\makeatletter
\@ifundefined{c@algorithm}{\usepackage{algorithm}}{}
\makeatother
\usepackage{algorithmicx,algpseudocode}

% diagrams
\usepackage{tikz,tikz-cd,pgfplots}
\usetikzlibrary{graphs,decorations.pathmorphing,decorations.markings,
arrows.meta,automata,positioning}
\tikzcdset{every label/.append style = {font = \small}}
\pgfplotsset{compat=1.18}

% chemistry
\usepackage[version=4]{mhchem}

\usepackage{microtype}
\usepackage{lineno}
\usepackage{nomencl}
\makenomenclature

% hyperref first, then bookmark, then glossaries: glossaries must follow
% hyperref or the entry links are built against the wrong anchors.
\makeatletter
\@ifundefined{hyperrefpackage}{%
  \@ifundefined{AAdefNnohyperref}{\usepackage{hyperref}}{}%
}{}
\@ifundefined{hyperrefpackage}{\usepackage{bookmark}}{}
\makeatother
\usepackage[toc,sort=use]{glossaries}
\makeglossaries

% Book/standalone switch.  Material that the book already covers in its
% prelude, but that a standalone paper still has to state, is wrapped as
%     \ifbookbuild\else ... \fi
% Drivers default to standalone; bookM.tex says \bookbuildtrue.
\newif\ifbookbuild
\bookbuildfalse

% Part switches for the gradostat paper.  Both default to true; a driver that
% wants the short submission says \madecfalse and/or \slowfastfalse, and every
% cross-reference into the dropped part disappears with it.
%   \ifmadec     Part III, the Lagasquie--Madec model and its break-even geometry
%   \ifslowfast  Part IV, everything carrying the small parameter \eps
\newif\ifmadec     \madectrue
\newif\ifslowfast  \slowfasttrue

% Part switch for the Barrat/simplicial-contagion project.  Defaults to
% true (the hypernetwork-synchrony part is included); the journal-paper
% driver says \hnetpartfalse to drop it and every cross-reference into it.
\newif\ifhnetpart  \hnetparttrue

%% ========================================
%% PACKAGES
%% ========================================
\RequirePackage{etoolbox}
\usepackage{makecell}
\usepackage[novbox]{pdfsync}% obsolete; SyncTeX is built into modern engines
\usepackage{amssymb, amsfonts,mathtools,bbm,eurosym,graphicx,latexsym, amsmath,amsthm,times,nccmath,url}
\usepackage{listings,verbatim,cases,tabularx}
% hyperref is deliberately NOT in the line above.  Journal classes often
% load it themselves, with a long option list, from a begindocument hook
% (mdpi.cls does: \DeclareHookRule{begindocument}{hyperref}{before}{mdpi}).
% Loading it here first, without options, makes that later load raise
% "Option clash for package hyperref".  So load it only if nobody else
% has, or will:
%   - \@ifpackageloaded  catches a class that loaded it already;
%   - \hyperrefpackage    is mdpi.cls's own deferred loader;
%   - \AAdefNnohyperref   lets any other driver opt out by defining it
%                         before % Shared preamble for the GLVbook papers.  This is the single canonical copy;
% a driver one level below GLVbook reaches it as
%
%     \input{../basicFiles/preamble}
%
% What stays in the driver, because it differs from paper to paper:
% \documentclass, font choice, \title, \author, \date, and any macro genuinely
% local to one paper.  Everything else belongs here.  The file is
% self-contained: it does not read any other file.

\usepackage[utf8]{inputenc}
\usepackage[novbox]{pdfsync}   % required: SyncTeX alone does not work here

% Page geometry, set once here for every paper that loads this file.  A driver
% that needs different margins says \geometry{...}; a second
% \usepackage[...]{geometry} raises an option clash and must not be used.
\makeatletter
\@ifundefined{AAdefNnogeometry}{\usepackage[a4paper,margin=2cm]{geometry}}{}
\makeatother

% mathematics
\usepackage{amsmath,amssymb,amsthm,mathtools,bm,mathrsfs,bbm}
\usepackage{centernot}

% floats, tables, lists
\usepackage{graphicx,float}
\makeatletter
\@ifpackageloaded{subfig}{}{\usepackage{subcaption}}
\makeatother
\usepackage{booktabs,longtable,threeparttable,array,multirow}
\usepackage{enumitem}

% colour, code, algorithms
\usepackage{xcolor,listings,tcolorbox}
\makeatletter
\@ifundefined{c@algorithm}{\usepackage{algorithm}}{}
\makeatother
\usepackage{algorithmicx,algpseudocode}

% diagrams
\usepackage{tikz,tikz-cd,pgfplots}
\usetikzlibrary{graphs,decorations.pathmorphing,decorations.markings,
arrows.meta,automata,positioning}
\tikzcdset{every label/.append style = {font = \small}}
\pgfplotsset{compat=1.18}

% chemistry
\usepackage[version=4]{mhchem}

\usepackage{microtype}
\usepackage{lineno}
\usepackage{nomencl}
\makenomenclature

% hyperref first, then bookmark, then glossaries: glossaries must follow
% hyperref or the entry links are built against the wrong anchors.
\makeatletter
\@ifundefined{hyperrefpackage}{%
  \@ifundefined{AAdefNnohyperref}{\usepackage{hyperref}}{}%
}{}
\@ifundefined{hyperrefpackage}{\usepackage{bookmark}}{}
\makeatother
\usepackage[toc,sort=use]{glossaries}
\makeglossaries

% Book/standalone switch.  Material that the book already covers in its
% prelude, but that a standalone paper still has to state, is wrapped as
%     \ifbookbuild\else ... \fi
% Drivers default to standalone; bookM.tex says \bookbuildtrue.
\newif\ifbookbuild
\bookbuildfalse

% Part switches for the gradostat paper.  Both default to true; a driver that
% wants the short submission says \madecfalse and/or \slowfastfalse, and every
% cross-reference into the dropped part disappears with it.
%   \ifmadec     Part III, the Lagasquie--Madec model and its break-even geometry
%   \ifslowfast  Part IV, everything carrying the small parameter \eps
\newif\ifmadec     \madectrue
\newif\ifslowfast  \slowfasttrue

% Part switch for the Barrat/simplicial-contagion project.  Defaults to
% true (the hypernetwork-synchrony part is included); the journal-paper
% driver says \hnetpartfalse to drop it and every cross-reference into it.
\newif\ifhnetpart  \hnetparttrue

%% ========================================
%% PACKAGES
%% ========================================
\RequirePackage{etoolbox}
\usepackage{makecell}
\usepackage[novbox]{pdfsync}% obsolete; SyncTeX is built into modern engines
\usepackage{amssymb, amsfonts,mathtools,bbm,eurosym,graphicx,latexsym, amsmath,amsthm,times,nccmath,url}
\usepackage{listings,verbatim,cases,tabularx}
% hyperref is deliberately NOT in the line above.  Journal classes often
% load it themselves, with a long option list, from a begindocument hook
% (mdpi.cls does: \DeclareHookRule{begindocument}{hyperref}{before}{mdpi}).
% Loading it here first, without options, makes that later load raise
% "Option clash for package hyperref".  So load it only if nobody else
% has, or will:
%   - \@ifpackageloaded  catches a class that loaded it already;
%   - \hyperrefpackage    is mdpi.cls's own deferred loader;
%   - \AAdefNnohyperref   lets any other driver opt out by defining it
%                         before \input{preamble}.
\makeatletter
\@ifpackageloaded{hyperref}{}{%
  \@ifundefined{hyperrefpackage}{%
    \@ifundefined{AAdefNnohyperref}{\usepackage{hyperref}}{}%
  }{}%
}
\makeatother
% linktoc=all makes a table-of-contents page number a link as well as the
% entry text; hyperref's default, linktoc=section, links only the text.
% Deferred to begin-document and guarded: by the block above hyperref may
% be loaded later (mdpi's own hook) or not at all (\AAdefNnohyperref), and
% \hypersetup must exist before it is called.  A driver that wants the
% default back can say \hypersetup{linktoc=section} after \input{preamble}.
\makeatletter
\AtBeginDocument{\@ifundefined{hypersetup}{}{\hypersetup{linktoc=all}}}
\makeatother
\usepackage{graphicx,float,fancybox}
\usepackage{epsfig,epstopdf,tikz}
\usepackage[sans]{dsfont}
\usetikzlibrary{positioning,arrows,arrows.meta,calc}
% Configure listings for Mathematica code
\lstdefinestyle{mathematica}{
 language=Mathematica,
 basicstyle=\small\ttfamily,
 backgroundcolor=\color{gray!10},
 frame=single,
 breaklines=true,
 commentstyle=\color{green!60!black},
 keywordstyle=\color{blue},
 stringstyle=\color{red},
 showstringspaces=false,
 tabsize=2
}
\usepackage{tcolorbox}
% Colored boxes for key results
\newtcolorbox{keyresult}[1][]{
 colback=blue!5!white,
 colframe=blue!75!black,
 title=#1,
 fonttitle=\bfseries
}
\newtcolorbox{examplebox}[1][]{
 colback=green!5!white,
 colframe=green!75!black,
 title=#1,
 fonttitle=\bfseries
}


\makeatletter
\@ifundefined{theorem}   {\newtheorem{theorem}   {Theorem}}   {}
\@ifundefined{definition}{\newtheorem{definition}{Definition}} {}
\@ifundefined{open}      {\newtheorem{open}      {Problem}}   {}
\@ifundefined{example}   {\newtheorem{example}   {Example}}   {}
\makeatother

% COUNTER DEFINITIONS
\newcounter{tabcounter}
\newcounter{defcounter}
\newcounter{opcounter}
\newcounter{thmcounter}
\newcounter{excounter}
\newcounter{ascounter}
\newcounter{cjcounter}
\newcounter{lemcounter}
\newcounter{remcounter}
\newcounter{corcounter}
\newcounter{propcounter}
\newcounter{figcounter}

% UNIFIED GLOSSARY STYLE
\newglossarystyle{unifiedstyle}{%
  \setglossarystyle{list}%
  \renewcommand*{\glossarypreamble}{\gdef\glslastsec{}}%
  \renewcommand{\glossentry}[2]{%
    \edef\glscursec{\glsentryuseriii{##1}}%
    \ifx\glscursec\glslastsec\else
      \item[]\textbf{\S\,\glscursec\enspace\glsentryuseriv{##1}}%
      \global\let\glslastsec\glscursec
    \fi
    \item[\glsentryuserii{##1} \glsentryuseri{##1}:] ``\glossentryname{##1}'' \glsentrydesc{##1}%
  }%
}

\makeatletter
% First-word extractor (used by \beT, \beD and \beXa for the default key)
% \@firstwordof hello world \relax  ->  hello
\def\@firstwordof#1 #2\relax{#1}
% First three letters, used for the heading mnemonic of a theorem given no
% explicit key.  Padding with \@empty keeps it safe for first words shorter
% than three letters: \@empty vanishes under \edef, whereas \relax would
% survive into the label and is illegal inside \csname (it broke every
% two-letter title, e.g. "ME ...", with "Missing \endcsname inserted").
\def\@takethree#1#2#3#4\@nil{#1#2#3}

% EXAMPLE COMMAND -- same signature and behaviour as \beT and \beD.
% Usage: \beXa*[long title][short title][key]
%   star  = suppresses the printed "[key]" tag after the title (as for \beT*
%   and \beD*); the label and glossary entry are unaffected.
%   key = optional label key override; the label becomes e:KEY, and without a
%   key it is e: plus the first three letters of the title's first word.  A
%   mnemonic label lets an example be cited as "Example~\ref{e:KEY}" instead
%   of by section number, which is what the citing prose needs.
\NewDocumentCommand{\beXa}{s O{} O{} o}{%
  \def\@temparg{#2}%
  \ifx\@temparg\@empty
    \begin{example}%
  \else
    \let\@savedSecTitle\@currentsectionname%
    \refstepcounter{excounter}%
    \IfValueTF{#4}{%
      \edef\@exlblkey{#4}%
    }{%
      \edef\@exfirstword{\@firstwordof#2 \relax}%
      \edef\@exlblkey{%
        \expandafter\@takethree\@exfirstword\@empty\@empty\@empty\@nil}%
    }%
    \def\@tempshort{#3}%
    \ifx\@tempshort\@empty
      \expandafter\xdef\csname @xshort@\@exlblkey\endcsname{\@exlblkey}%
    \else
      \expandafter\def\csname @xshort@\@exlblkey\endcsname{#3}%
    \fi
    \IfBooleanTF{#1}{\let\@extag\@empty}{%
      \edef\@extag{\noexpand\ \noexpand\textup{[\noexpand\texttt{e:\@exlblkey}]}}}%
    \begin{example}[#2\ifshowkeys\@extag\fi]%
    \label{e:\@exlblkey}%
    % record the long title so \nref{e:KEY} prints it, as for theorems
    \protected@write\@auxout{}{%
      \string\@thmtitle{e:\@exlblkey}{\unexpanded{#2}}}%
    \expanded{%
      \noexpand\newglossaryentry{e\theexcounter}{%
        name={\unexpanded{#2}%
          \IfBooleanF{#1}{ [\noexpand\texttt{e:\@exlblkey}]}},%
        sort={\theexample},%
        description={p.~\noexpand\pageref{e:\@exlblkey}},%
        user1={\theexample},%
        user2={Example},%
        user3={\thesection},%
        user4={\expandafter\unexpanded\expandafter{\@savedSecTitle}}%
      }%
    }%
    \protected@write\@auxout{}{\string\glsliveentry{e\theexcounter}}%
  \fi
}
\newcommand{\eeXa}{\end{example}}

% DEFINITION COMMAND -- same signature and behaviour as \beT below.
% Usage: \beD[long title][short title][key]
%   long title  = definition heading (plain text only)
%   short title = shown in \dref{word}; defaults to the key
%   key         = optional label key override; the label becomes d:KEY
% Without a key the label is d: plus the first three letters of the title's
% first word, so it is a mnemonic rather than the definition's number: a
% numeric label would silently change whenever a definition is inserted.
\NewDocumentCommand{\beD}{s O{} O{} o}{%
  \def\@temparg{#2}%
  \ifx\@temparg\@empty
    \begin{definition}%
  \else
    \let\@savedSecTitle\@currentsectionname%
    \refstepcounter{defcounter}%
    \IfValueTF{#4}{%
      \edef\@deflblkey{#4}%
    }{%
      \edef\@deffirstword{\@firstwordof#2 \relax}%
      \edef\@deflblkey{%
        \expandafter\@takethree\@deffirstword\@empty\@empty\@empty\@nil}%
    }%
    \def\@tempshort{#3}%
    \ifx\@tempshort\@empty
      \expandafter\xdef\csname @dshort@\@deflblkey\endcsname{\@deflblkey}%
    \else
      \expandafter\def\csname @dshort@\@deflblkey\endcsname{#3}%
    \fi
    \IfBooleanTF{#1}{\let\@deftag\@empty}{%
      \edef\@deftag{\noexpand\ \noexpand\textup{[\noexpand\texttt{d:\@deflblkey}]}}}%
    \begin{definition}[#2\ifshowkeys\@deftag\fi]%
    \label{d:\@deflblkey}%
    % record the long title so \nref{d:KEY} prints it, as for theorems
    \protected@write\@auxout{}{%
      \string\@thmtitle{d:\@deflblkey}{\unexpanded{#2}}}%
    \expanded{%
      \noexpand\newglossaryentry{d\thedefcounter}{%
        name={\unexpanded{#2}%
          \IfBooleanF{#1}{ [\noexpand\texttt{d:\@deflblkey}]}},%
        sort={\thedefinition},%
        description={p.~\noexpand\pageref{d:\@deflblkey}},%
        user1={\thedefinition},%
        user2={Definition},%
        user3={\thesection},%
        user4={\expandafter\unexpanded\expandafter{\@savedSecTitle}}%
      }%
    }%
    \protected@write\@auxout{}{\string\glsliveentry{d\thedefcounter}}%
  \fi
}
\newcommand{\eeD}{\end{definition}}

% \showkeysfalse in a driver hides the [t:KEY] mnemonics in result headings
% (labels and glossary entries are unchanged); used for submission builds.
\newif\ifshowkeys \showkeystrue
% \commonredtrue in a driver that inputs basicFiles/commonRed: tool.tex then
% cites its reduction results instead of restating them.
\newif\ifcommonred \commonredfalse
% TheorEM COMMAND
% Usage: \beT[long title][short title][key]
%   long title  = theorem heading (plain text only)
%   short title = shown in \tref{word}; may contain math; defaults to first word of long title
%   key         = optional label key override; label becomes t:KEY (avoids first-word collisions)
% Label is t:WORD where WORD = first word of long title (or KEY if given); use \tref{WORD} to reference.
\NewDocumentCommand{\beT}{s O{} O{} o}{%
  \def\@temparg{#2}%
  \ifx\@temparg\@empty
    \begin{theorem}%
  \else
    \let\@savedSecTitle\@currentsectionname%
    \refstepcounter{thmcounter}%
    % \@thmtag is the mnemonic shown in the heading; every theorem gets one.
    % With an explicit key the key itself is shown.  Without one the key
    % defaults to the title's first word, and showing that whole word would
    % merely repeat the title ("Dissipativity [Dissipativity]"), so only its
    % first three letters are shown: "Dissipativity [Dis]".  The label is
    % unaffected -- it still uses the full \@thmlblkey.
    \IfValueTF{#4}{%
      \edef\@thmlblkey{#4}%
    }{%
      \edef\@thmfirstword{\@firstwordof#2 \relax}%
      \edef\@thmlblkey{%
        \expandafter\@takethree\@thmfirstword\@empty\@empty\@empty\@nil}%
    }%
    \IfBooleanTF{#1}{\let\@thmtag\@empty}{%
      \edef\@thmtag{\noexpand\ \noexpand\textup{[\noexpand\texttt{t:\@thmlblkey}]}}}%
    \def\@tempshort{#3}%
    \ifx\@tempshort\@empty
      \expandafter\xdef\csname @tshort@\@thmlblkey\endcsname{\@thmlblkey}%
    \else
      \expandafter\def\csname @tshort@\@thmlblkey\endcsname{#3}%
    \fi
    \begin{theorem}[#2\ifshowkeys\@thmtag\fi]%
    \label{t:\@thmlblkey}%
    % record the long title for \nref; via the .aux so forward refs resolve
    \protected@write\@auxout{}{%
      \string\@thmtitle{t:\@thmlblkey}{\unexpanded{#2}}}%
    \expanded{%
      \noexpand\newglossaryentry{t\thethmcounter}{%
        name={\unexpanded{#2}%
          \IfBooleanF{#1}{ [\noexpand\texttt{t:\@thmlblkey}]}},%
        % sort/user1 must be the number the document actually prints, i.e. the
        % theorem environment's counter, not thmcounter.  thmcounter advances
        % only on \beT with a title, so a bare \beT (or \beT with an
        % empty title) makes the two drift and the glossary cites a wrong
        % number.  thmcounter is still fine as the entry key: it is unique.
        sort={\thetheorem},%
        description={p.~\noexpand\pageref{t:\@thmlblkey}},%
        user1={\thetheorem},%
        user2={Theorem},%
        user3={\thesection},%
        user4={\expandafter\unexpanded\expandafter{\@savedSecTitle}}%
      }%
    }%
    \protected@write\@auxout{}{\string\glsliveentry{t\thethmcounter}}%
  \fi
}
\newcommand{\eeT}{\end{theorem}}

% CROSS-PART RELATION PARAGRAPH
% Usage: \xpart{text}  -- place immediately BEFORE \beT, never inside it
%   (anything inside \beT...\eeT would be swallowed into the glossary entry).
% The paragraph is typeset only when Part I is actually present, i.e. in the
% combined book build; in the standalone builds (GranadaDSABNS/A.tex,
% AAHRah-Submission/AAHRah.tex) the label part:glv is undefined and the whole
% paragraph disappears rather than leaving ?? cross-references. Same
% graceful-degradation principle as \GASref in GASt.tex. Resolves on the
% second LaTeX pass, since r@part:glv comes from the .aux file.
% (Already inside the \makeatletter region opened above, so no \makeatletter
% pair here -- adding one would close that region early.)
\newcommand{\xpart}[1]{\@ifundefined{r@part:glv}{}{\paragraph{Relation with Parts~I--II.}#1}}

% \bookref{label}{precise text}{fallback text}: prints the precise cross-reference
% when the label resolves (combined book build), and the fallback text otherwise
% (standalone paper builds, where book-only appendix labels are absent). Same
% graceful-degradation role as \GASref in GASt.tex, but usable in any file since it
% lives in this preamble (loaded by every driver) -- unlike \GASref, which is defined in a
% content file and is therefore undefined in standalone builds that omit that file.
\newcommand{\bookref}[3]{\@ifundefined{r@#1}{#3}{#2}}

% \PVpart and \RBpart: the two parts that a partial build may omit.
% They print "Part~N" in any build that contains the part, and a naming phrase
% in a build that does not, so an excerpt shows no ?? .  Each already supplies
% the word "Part", so write "\PVpart develops ...", not "Part~\PVpart".
\newcommand{\PVpart}{\bookref{part:aahrah}%
  {Part~\ref{part:aahrah}}{the Perron--Volterra part of the full book}}
\newcommand{\RBpart}{\bookref{part:recbrowne}%
  {Part~\ref{part:recbrowne}}{the immuno-virology part of the full book}}

% OPEN PROBLEM COMMANDS
% ASSUMPTION COMMANDS (glossaried)
\newcommand{\beAs}[1][]{%
  \def\@temparg{#1}%
  \ifx\@temparg\@empty
    \begin{assumption}%
  \else
    \let\@savedSecTitle\@currentsectionname%
    \refstepcounter{ascounter}%
    \begin{assumption}[#1]%
    \label{as:\theascounter}%
    \expanded{%
      \noexpand\newglossaryentry{as\theascounter}{%
        name={\unexpanded{#1}},%
        sort={\theascounter},%
        description={p.~\noexpand\pageref{as:\theascounter}},%
        user1={\theascounter},%
        user2={Assumption},%
        user3={\thesection},%
        user4={\expandafter\unexpanded\expandafter{\@savedSecTitle}}%
      }%
    }%
    \protected@write\@auxout{}{\string\glsliveentry{as\theascounter}}%
  \fi
}
\newcommand{\eeAs}{\end{assumption}}
% CONJECTURE COMMANDS (glossaried)
\newcommand{\beCj}[1][]{%
  \def\@temparg{#1}%
  \ifx\@temparg\@empty
    \begin{conjecture}%
  \else
    \let\@savedSecTitle\@currentsectionname%
    \refstepcounter{cjcounter}%
    \begin{conjecture}[#1]%
    \label{cj:\thecjcounter}%
    \expanded{%
      \noexpand\newglossaryentry{cj\thecjcounter}{%
        name={\unexpanded{#1}},%
        sort={\thecjcounter},%
        description={p.~\noexpand\pageref{cj:\thecjcounter}},%
        user1={\thecjcounter},%
        user2={Conjecture},%
        user3={\thesection},%
        user4={\expandafter\unexpanded\expandafter{\@savedSecTitle}}%
      }%
    }%
    \protected@write\@auxout{}{\string\glsliveentry{cj\thecjcounter}}%
  \fi
}
\newcommand{\eeCj}{\end{conjecture}}
% LEMMA COMMANDS
% ---------------------------------------------------------------------
% UNIFORM RESULT WRAPPERS -- \beL, \beC, \beP, \beR share \beT's signature
%
%   \beX[long title][short title][key]  ...  \eeX
%
%   long title  = heading (math allowed)
%   short title = shown in the cross-reference; defaults to the key
%   key         = label key.  With a key the label is X:KEY and the heading
%                 carries the mnemonic [key], exactly as for \beT.  Without
%                 one the label falls back to the running counter, which is
%                 what every pre-existing call site in the book relies on.
%
% A bare \beX -- arguments commented out -- opens the plain environment and
% creates NO label and NO glossary entry.  That is the supported way to keep
% an individual result out of the catalogue.
% ---------------------------------------------------------------------
% #1 wrapper name  #2 environment  #3 label prefix  #4 counter  #5 type name
\newcommand{\@mkresult}[5]{%
  % A leading * suppresses the [key] mnemonic in both the heading and the
  % glossary name, while still making the label X:KEY and the entry:
  %     \beX*[long][short][key]
  \expandafter\NewDocumentCommand\csname be#1\endcsname{s O{} O{} o}{%
    \def\@temparg{##2}%
    \ifx\@temparg\@empty
      \begin{#2}%
    \else
      \let\@savedSecTitle\@currentsectionname%
      \refstepcounter{#4}%
      % With a key: mnemonic label X:KEY and the [key] tag in the heading.
      % Without one: fall back to the running counter, exactly as these
      % wrappers behaved before, so every legacy one-argument call site in
      % the book keeps its label.  The counter is also csname-safe, which a
      % title's first word is not -- titles beginning with math would
      % otherwise build \csname \Omega...\endcsname and fail.
      \IfValueTF{##4}{%
        \edef\@relbl{##4}%
        \IfBooleanTF{##1}{\let\@retag\@empty}{%
          \edef\@retag{\noexpand\ \noexpand\textup{[\noexpand\texttt{#3:##4}]}}}%
      }{%
        \edef\@relbl{\arabic{#4}}%
        \let\@retag\@empty
      }%
      \def\@tempshort{##3}%
      \ifx\@tempshort\@empty
        \expandafter\xdef\csname @sshort@#3\@relbl\endcsname{\@relbl}%
      \else
        \expandafter\def\csname @sshort@#3\@relbl\endcsname{##3}%
      \fi
      \begin{#2}[##2\ifshowkeys\@retag\fi]%
      \edef\@tmplbl{#3:\@relbl}\expandafter\label\expandafter{\@tmplbl}%
      \expanded{%
        \noexpand\newglossaryentry{#3\csname the#4\endcsname}{%
          name={\unexpanded{##2}%
            \IfBooleanF{##1}{\IfValueT{##4}{ [\noexpand\texttt{#3:##4}]}}},%
          sort={\csname the#2\endcsname},%
          description={p.~\noexpand\pageref{#3:\@relbl}},%
          user1={\csname the#2\endcsname},%
          user2={#5},%
          user3={\thesection},%
          user4={\expandafter\unexpanded\expandafter{\@savedSecTitle}}%
        }%
      }%
      \protected@write\@auxout{}{\string\glsliveentry{#3\csname the#4\endcsname}}%
      \protected@write\@auxout{}{\string\@thmtitle{\@tmplbl}{\unexpanded{##2}}}%
    \fi
  }%
  \expandafter\newcommand\csname ee#1\endcsname{\end{#2}}%
}
\@mkresult{L}{lemma}{l}{lemcounter}{Lemma}
% REMARK / COROLLARY / PROPOSITION -- same signature as \beL above
\@mkresult{R}{remark}{r}{remcounter}{Remark}
\@mkresult{C}{corollary}{c}{corcounter}{Corollary}
\@mkresult{P}{proposition}{p}{propcounter}{Proposition}
% OPEN-PROBLEM COMMAND -- generated with the same signature as \beT:
%   \beO[long title][short title][key]   ...   \eeO
% label o:KEY, or o:<counter> when no key is given; a bare \beO makes
% neither a label nor a glossary entry.
\@mkresult{O}{open}{o}{opcounter}{Problem}


% Section-level title tracker: updated only by \section, not \subsection.
% \@currentlabelname is overwritten by every \refstepcounter (including
% \subsection), so we save it separately here.
\newcommand{\@currentsectionname}{}
\AtBeginDocument{%
  \let\@epidcrn@orig@refstepcounter\refstepcounter
  \renewcommand{\refstepcounter}[1]{%
    \@epidcrn@orig@refstepcounter{#1}%
    \ifstrequal{#1}{section}{\global\let\@currentsectionname\@currentlabelname}{}%
  }%
}

% TABLE COMMANDS (glossaried) -- consistent naming with \beT, \beD, etc.
\newcommand{\beTa}[1][]{%
  \def\@temparg{#1}%
  \ifx\@temparg\@empty
    \begin{table}[htbp]%
  \else
    \let\@savedSecTitle\@currentsectionname%
    \refstepcounter{tabcounter}%
    \begin{table}[htbp]%
    \label{A:\thetabcounter}%
    \expanded{%
      \noexpand\newglossaryentry{A\thetabcounter}{%
        name={\unexpanded{#1}},%
        sort={\thetabcounter},%
        description={p.~\noexpand\pageref{A:\thetabcounter}},%
        user1={\thetabcounter},%
        user2={Table},%
        user3={\thesection},%
        user4={\expandafter\unexpanded\expandafter{\@savedSecTitle}}%
      }%
    }%
    \protected@write\@auxout{}{\string\glsliveentry{A\thetabcounter}}%
  \fi
}
\newcommand{\eeTa}{\end{table}}

% FIGURE COMMANDS (glossaried) -- consistent with \beT, \beR, \beTa, etc.
% Usage:
%   \beF[Glossary title]...\caption{Visible caption}\eeF  -> glossaried figure
%   \beF...\caption{Visible caption}\eeF                  -> plain [H] figure, no glossary
\newcommand{\beF}[1][]{%
  \def\@temparg{#1}%
  \ifx\@temparg\@empty
    \begin{figure}[H]%
  \else
    \begin{figure}[H]%
    \let\@savedSecTitle\@currentsectionname%
    \refstepcounter{figcounter}%
    \label{fig:\thefigcounter}%
    \expanded{%
      \noexpand\newglossaryentry{Fig\thefigcounter}{%
        name={\unexpanded{#1}},%
        sort={\thefigcounter},%
        description={p.~\noexpand\pageref{fig:\thefigcounter}},%
        user1={\thefigcounter},%
        user2={Figure},%
        user3={\thesection},%
        user4={\expandafter\unexpanded\expandafter{\@savedSecTitle}}%
      }%
    }%
    \protected@write\@auxout{}{\string\glsliveentry{Fig\thefigcounter}}%
  \fi
}
\newcommand{\eeF}{\end{figure}}
\makeatother
% LIVE-ENTRY REGISTER
% \glsaddall prints every entry the glossaries package currently knows, and
% that set includes entries restored from \jobname.glsdefs.  The package
% rewrites .glsdefs from all *currently defined* entries, so anything it
% restored is written straight back out: a deleted section kept reappearing in
% the catalogue for ever, and the .glsdefs had to be erased by hand.  The
% wrappers therefore record every entry they actually create in the .aux, and
% \printboth adds only those.  The .aux is rebuilt from scratch on each run, so
% this is self-healing: a removed result is gone after one further run.
\makeatletter
\newcommand{\@glslive}{}
\newcommand{\@glsaddlive}[1]{\ifglsentryexists{#1}{\glsadd{#1}}{}}
\newcommand{\glsliveentry}[1]{\g@addto@macro\@glslive{\@glsaddlive{#1}}}
\makeatother
% PRINT COMMAND
% No \clearpage here -- the catalogue continues on the current page.  This has
% been re-added by mistake several times; leave it out.
\makeatletter
\newcommand{\printboth}{%
  \@glslive   % add exactly the entries this document created last run
  \printglossary[style=unifiedstyle, title={0 \hspace{.22em} Definitions, Theorems, Examples and Open Problems}]%
}
\makeatother
% Named theorem reference: \tref{word} -> Thm~N(word-or-short)
% Shows stored short title if available, otherwise shows the key itself.
\newcommand{\tref}[1]{%
  Thm~\ref{t:#1}(\ifcsname @tshort@#1\endcsname\csname @tshort@#1\endcsname\else#1\fi)%
}

% Full-title theorem reference: \nref{t:KEY} prints the number and then the
% theorem's long title, so one writes "Theorem~\nref{t:blocks}" and gets
% "Theorem 3 ``Block factorization and siphons''".  Takes the complete label, like
% \ref, not the bare key like \tref.
%
% \beT writes each title to the .aux (see \@thmtitle below), so \nref resolves
% forward references from the second LaTeX pass, exactly as \ref does.  If the
% title is not yet known -- first pass, or a label that is not a \beT theorem
% -- it degrades to the number rather than producing ??.
% These two use @ in their names, and this point in the file is outside the
% \makeatletter region opened earlier (it closes at \makeatother above), so
% the region has to be reopened here.
\makeatletter
\newcommand{\@thmtitle}[2]{\global\@namedef{@tlong@#1}{#2}}
\newcommand{\nref}[1]{%
  \ref{#1}%
  \ifcsname @tlong@#1\endcsname
    \ ``\csname @tlong@#1\endcsname''%
  \fi
}
% \nm{t:KEY}: the stored long title of any \beT/\beL/\beP/\beR/\beC/\beD/\beE result,
% verbatim and without quotes, so item names follow the macro's first argument.
\newcommand{\nm}[1]{%
  \ifcsname @tlong@#1\endcsname
    \csname @tlong@#1\endcsname%
  \else
    \ref{#1}%
  \fi
}
\makeatother




%% ========================================
%% Other THEOREM ENVIRONMENTS
%% ========================================
\makeatletter
\@ifundefined{lemma}      {\newtheorem{lemma}      {Lemma}}      {}
\@ifundefined{proposition}{\newtheorem{proposition}{Proposition}} {}
\@ifundefined{corollary}  {\newtheorem{corollary}  {Corollary}}  {}
\@ifundefined{conjecture} {\newtheorem{conjecture} {Conjecture}}  {}
\@ifundefined{remark}     {\newtheorem{remark}     {Remark}}     {}
\@ifundefined{question}   {\newtheorem{question}   {Question}}   {}
\@ifundefined{assumption} {\newtheorem{assumption} {Assumption}} {}
\@ifundefined{hypothesis} {\newtheorem{hypothesis} {Hypothesis}} {}
\@ifundefined{fact}       {\newtheorem{fact}       {Fact}}       {}
\makeatother

% \citer: cite the relay-graph threshold theorem of AAHrel. If a document has
% \input the chapter/file containing t:relayRF (so the label resolves), point
% straight at it; otherwise fall back to a bare citation of AAHrel. Shared here
% (not per-paper) so it resolves consistently across every book chapter.
% \citeR: same idea, for the first theorem of AAHRah (t:Pstar).
\makeatletter
\providecommand{\citer}{%
  \@ifundefined{r@t:relayRF}%
    {\cite{AAHrel}}%
    {Theorem~\ref{t:relayRF} of \cite{AAHrel}}%
}
\providecommand{\citeR}{%
  \@ifundefined{r@t:Pstar}%
    {\cite{AAHRah}}%
    {Theorem~\ref{t:Pstar} of \cite{AAHRah}}%
}
\makeatother


% Abbreviations
\def\beQ{\begin{question}}\def\eeQ{\end{question}}
\def\beA{\begin{assumption}}\def\eeA{\end{assumption}}


%% ========================================

%% ========================================
%% COLORS & CUSTOM COMMANDS
%% ========================================
\definecolor{funccolor}{RGB}{25,25,112}
\definecolor{desccolor}{RGB}{64,64,64}
\newcommand{\func}[1]{\texttt{\textcolor{funccolor}{#1}}}
\def\eq{=}
\def\mm{-}

%% ========================================
%% MATHCAL SHORTHANDS
%% ========================================
% Project-wide convention: \mX for a capital X is always \mathcal X.
% Defined for the whole alphabet so the convention never fails to expand.
\newcommand{\mA}{\mathcal A}\newcommand{\mB}{\mathcal B}
\newcommand{\mC}{\mathcal C}\newcommand{\mD}{\mathcal D}
\newcommand{\mE}{\mathcal E}\newcommand{\mF}{\mathcal F}
\newcommand{\mG}{\mathcal G}\newcommand{\mH}{\mathcal H}
\newcommand{\mI}{\mathcal I}\newcommand{\mJ}{\mathcal J}
\newcommand{\mK}{\mathcal K}\newcommand{\mL}{\mathcal L}
\newcommand{\mM}{\mathcal M}\newcommand{\mN}{\mathcal N}
\newcommand{\mO}{\mathcal O}\newcommand{\mP}{\mathcal P}
\newcommand{\mQ}{\mathcal Q}\newcommand{\mR}{\mathcal R}
\newcommand{\mS}{\mathcal S}\newcommand{\mT}{\mathcal T}
\newcommand{\mU}{\mathcal U}\newcommand{\mV}{\mathcal V}
\newcommand{\mW}{\mathcal W}\newcommand{\mX}{\mathcal X}
\newcommand{\mY}{\mathcal Y}\newcommand{\mZ}{\mathcal Z}

%% ========================================
%% BASIC MATH ENVIRONMENTS
%% ========================================
\def\bep{\begin{pmatrix}}\def\eep{\end{pmatrix}}
\def\bev{\begin{vmatrix}}\def\eev{\end{vmatrix}}
\def\bea{\begin{eqnarray*}}\def\eea{\end{eqnarray*}}
\def\bc{\begin{cases}}\def\ec{\end{cases}}
\def\BEN{\begin{enumerate}}\def\EEN{\end{enumerate}}
\def\BI{\begin{itemize}}\def\EI{\end{itemize}}

\newcommand{\be}[1]{\begin{equation}\label{#1}}
\newcommand{\ee}{\end{equation}}
\newcommand{\beq}{\begin{eqnarray}}
\def\eeq{\end{eqnarray}}
\def\eqr{\eqref}\def\fr{\frac}\def\lbl{\label}
\def\Lra{\Longrightarrow}\def\non{\nonumber}\def\cNGM{\cite{Diek,Van,Van08}}
\def\Eq{\Leftrightarrow}\def\b{\beta}\def\uph{\red{UP TO HERE}}
%% ========================================
%% PROBABILITY & STATISTICS
%% ========================================
\renewcommand{\P}{\ensuremath{\mathbb{P}}}  % Note: overrides paragraph symbol
\newcommand{\E}{\ensuremath{\mathbb{E}}}
\newcommand{\ind}[1]{\ensuremath{\mathbbm{1}_{\left\{#1\right\}}}}
\DeclareMathOperator{\Cov}{Cov}
\DeclareMathOperator{\TV}{TV}
\newcommand{\rank}{\operatorname{rank}}

%% ========================================
%% CRN OPERATORS (organizations/ERCs/buffering/siphons framework)
%% ========================================
\def\clos{\mathrm{clos}}
\def\Supply{\mathrm{Supply}}
\def\req{\mathrm{req}}
\def\React{\mathrm{React}}
\def\Prod{\mathrm{Prod}}
\def\Gen{\mathrm{Gen}}
\def\In{\mathrm{In}}
\def\Out{\mathrm{Out}}
\def\Act{\mathrm{Act}}
\def\Min{\mathrm{Min}}

%% ========================================
%% LINE BREAKING
%% ========================================
% Last-resort stretch, used only on paragraphs TeX cannot otherwise break:
% a long \texttt{} token, a comma-separated inline math list (commas are not
% breakpoints in math), or a run of \cite/\ref boxes.  Without it those set as
% overfull \hbox and print into the margin; with it they set slightly loose,
% which is the lesser evil.  Does not affect paragraphs that already break.
\emergencystretch=3em

%% ========================================
%% NUMBER SETS
%% ========================================
\newcommand{\R}{\mathbb{R}}  % Real numbers
\newcommand{\C}{\mathbb{C}}  % Complex numbers
\newcommand{\N}{\mathbb{N}}
\newcommand{\Z}{\mathbb{Z}}
% \providecommand, not \newcommand: some papers define these in their own
% preamble after \input{preamble}, and must not clash.
\providecommand{\Q}{\mathbb{Q}}
\providecommand{\res}{\operatorname{res}}   % resultant; used by epiRes.tex
\newcommand{\MN}{\mathbb{N}}
\newcommand{\MZ}{\mathbb{Z}}
\newcommand{\realpos}{\mathbb{R}_{> 0}}
\newcommand{\realnneg}{\mathbb{R}_{\geq 0}}
\newcommand{\naturalpos}{\mathbb{N}_{>0}}
\newcommand{\Rplus}{\mathbb{R}_+}
\newcommand{\Rp}{\mathbb{R}_{\ge 0}}
\def\Rnne{\mathbb{R}_{\ge 0}}
\def\Se{\bar{S}} \def\Ie{\bar{I}}
\newcommand{\Rpos}{\mathbb{R}_{>0}}
\newcommand{\dd}{\,\mathrm d}
\newcommand{\e}{\mathrm e}
%% ========================================
%% GREEK LETTERS (2+ character abbreviations) greek
%% ========================================
\def\al{\alpha}\def\ba{\bff \alpha}
\def\w{\omega}\def\Lgn{\ga+\mu+\de}\def\om{\omega}
\def\de{\delta}\def\De{\Delta}\def\bde{\bff \delta}
\def\ep{\epsilon}\def\eps{\varepsilon}
\def\f{\varphi}  % Kept as requested
\def\ga{\gamma}\def\Ga{\Gamma}\def\G{\Gamma}\def\bg{\bff \gamma}
\def\up{\upsilon}  % NOT the breve accent: no accent commands anywhere.
                   % Named \up, not \u, so the LaTeX breve stays reachable.
\def\la{\lambda}\def\La{\Lambda}
\def\si{\sigma}\def\ka{\kappa}
\def\th{\theta}\def\vt{\vartheta}
\def\ze{\zeta}\def\Gap{\Gamma^{+}} \def\Gam{\Gamma^{-}}


%% ========================================
%% SPECIAL SHORT COMMANDS (kept as requested)
%% ========================================
\def\I{\infty}
\def\T{\widetilde}
\def\H{\hat}
\newcommand{\relint}{\operatorname{relint}}

%% ========================================
%% BOLD/VECTOR NOTATION
%% ========================================
\newcommand{\bff}[1]{{\mbox{\boldmath$#1$}}}
\newcommand\X{\boldsymbol{X}}
\def\x{\boldsymbol{x}}\def\bx{\bff x}\def\xe{\bar{\x}}
\newcommand\y{\boldsymbol{y}}\def\by{\boldsymbol{y}}\def\ye{\bar{\y}}
\def\yb{\y}\def\xb{\x}
\def\xy{\x-\y}
\newcommand\br{\boldsymbol{r}}\def\re{\bar{\r}}
\newcommand\bd{\boldsymbol{d}}
\newcommand\bQ{\boldsymbol{Q}}
\newcommand{\Yb}{{\mathbf{Y}}}
\newcommand{\Ab}{\mathbf{A}}
\newcommand{\zb}{\mathbf{z}}\def\bz{\bff 0}
\newcommand{\Zb}{\mathbf{Z}}
\newcommand{\bb}{\mathbf{b}}
\def\vbe{\vec {\bff \beta}}
\def\vb{\vec {\bff b}}
\def\v1{\vec {\bff 1}}
\def\vn{\vec \de}
\def\vz{\vec 0}
\def\vi{\overset{\rightarrow}{\mathsf i}}
\def\vv{\overset{\rightarrow}{v}}
\def\vr{\; \vec {\mathsf r}}
\newcommand{\va}{\bff a}
\providecommand{\vect}[1]{\boldsymbol{#1}}   % bold vector (moved from Gradostat.tex)
\providecommand{\one}{\mathbf 1}             % all-ones vector (moved from AliSF)
\providecommand{\spr}{s}                     % spectral bound  (moved from AliSF)
\providecommand{\Ke}{K_{\rm eff}}            % effective carrying capacity
% Critical right/left eigenvectors in the Hopf projection method (AAHOSN/New.tex)
\providecommand{\bq}{\mathbf{q}}
\providecommand{\bp}{\mathbf{p}}


%% ========================================
%% MATHCAL & MATHSCR
%% ========================================
% Mathcal
\newcommand\mc[1]{\mathcal{#1}}
\def\mC{{\mathcal C}}\def\mD{\mathcal D}\def\mE{\mathcal{E}}
\def\mF{\mathcal F}\def\mG{{\mathcal G}}\def\mH{{\mathcal H}}
\def\mI{\mathcal I}\def\mL{\mathcal{L}}\def\mM{{\mathcal M}}
\def\mN{{\mathcal N}}\def\mP{\mathcal{P}}\def\mR{{\mathcal R}}
\def\mS{{\mathcal S}}\def\mT{\mathcal{T}}\def\mV{{\mathcal V}}
\def\mK{\mathcal K}\def\mA{{\mathcal A}}\def\mU{{\mathcal U}}
% Symmetric bilinear and trilinear forms of the Taylor expansion, used in the
% first Lyapunov coefficient (AAHOSN/New.tex)
\providecommand{\Bform}{\mathcal{B}}
\providecommand{\Cform}{\mathcal{C}}

% Blackboard bold.  Kept here, not in the drivers: a \providecommand in a
% driver preamble runs before this file and silently wins, so the same name
% ends up meaning different things in different builds.
\def\IC{\mathbb{C}}\def\IN{\mathbb{N}}\def\IQ{\mathbb{Q}}
\def\IR{\mathbb{R}}\def\IZ{\mathbb{Z}}\def\IG{\mathbb{G}}
\def\II{\mathbb{I}}

% Mathscr
\newcommand{\Rs}{{\mathscr R}}
\newcommand{\Ss}{{\mathscr S}}
\newcommand{\Ns}{{\mathscr N}}
\newcommand{\Cs}{\mathscr C}


%% ========================================
%% OPERATORS & FUNCTIONS
%% ========================================
\DeclareMathOperator{\tr}{tr}
\DeclareMathOperator{\diag}{Diag}
\DeclareMathOperator{\Diag}{Diag}
\DeclareMathOperator{\intr}{\mathrm{Int}}
\DeclareMathOperator{\Inv}{Inv}
\DeclareMathOperator{\spec}{spec}
\DeclareMathOperator{\AUC}{AUC}               % area under the curve
\newcommand{\ii}{\mathrm{i}}                  % imaginary unit
\newcommand\op{\operatorname}
\newcommand{\supp}{\op{supp}}
\newcommand{\diff}{\mathop{}\mathopen{}\mathrm{d}}
\newcommand{\pdd}[1]{\frac{\partial}{\partial{#1}}}
\DeclareMathOperator{\cone}{cone}
\DeclareMathOperator{\proj}{proj}

%% ========================================
%% BRACKETS & NORMS
%% ========================================
\providecommand{\pr}[1]{\left(#1\right)}
\providecommand{\pp}[1]{\left[#1\right]}
\newcommand\croc[1]{\left\langle #1\right\rangle}
\newcommand{\norm}[1]{\|#1\|}
\newcommand{\inner}[2]{\langle #1, #2 \rangle}
% Short synonym of \inner, used by AAHOSN/New.tex as $\ip{u}{v}=\bar u^\top v$
\providecommand{\ip}[2]{\langle #1 , #2 \rangle}


%% ========================================
%% FORMATTING & UTILITIES
%% ========================================
\newcommand*{\Scale}[2][4]{\scalebox{#1}{$#2$}}
\newcommand{\te}[1]{\text{#1}}
\newcommand\steq[1]{\stackrel{\text{\rm #1.}}{=}}
\newcommand\blue[1]{\textcolor{blue}{#1}}
\newcommand{\red}{\textcolor[rgb]{1.00,0.00,0.00}}
\newcommand{\green}{\textcolor[rgb]{0.00,1.00,0.00}}

% Figure command
\newcommand{\figu}[3]{%
\begin{figure}[H]
\centering
\includegraphics[scale=#3]{#1}
\caption{#2\label{f:#1}}
\end{figure}
}

% Footnote
\long\def\symbolfootnote[#1]#2{%
\begingroup
\def\thefootnote{\fnsymbol{footnote}}\footnote[#1]{#2}%
\endgroup}
\def\fn{\symbolfootnote}

\newcommand{\dsum}{\displaystyle\sum}

%% ========================================
%% CRN & EPIDEMIOLOGY SPECIFIC
%% ========================================
% CRN notation
\def\SS{\mathcal S}
\def\CC{\mathcal C}
\def\RR{\mathcal R}

% Reproduction numbers
\def\Rto{R_{1}^{E_2}}
\def\Rot{R_{2}^{E_1}}
\def\Roc{R_{1,c}}
\def\Rtc{R_{2,c}}
\def\rd{r_{0}}
\def\sd{s_{0}}
\def\dd{D_{0}}
\def\vd{v_{0}}
\def\zd{z_{0}}
\def\se{\sE}
\def\xd{x_{0}}
%% ========================================
%% SPECIAL VALUES & EQUILIBRIA
%% ========================================
\def\eE{\bar{e}}
\def\wE{\bar{w}}
\def\sE{s_{*}}
\def\Me{\T M_{ee}}
\def\Rnne{\mathbb{R}_{\ge 0}}

\newcommand{\ones}{\mathbf 1}


%% ========================================
%% ABBREVIATIONS (alphabetical by category)
%% ========================================

% -- General mathematical terms --
\def\adm{admissible}\def\admy{admissibility}
%\def\and{antisymmetric}% disabled 2026-09-30: it overrode LaTeX's author separator \and (printed "antisymmetric" between authors)
\def\ass{assumption}
\def\cadlag{c\`adl\`ag}
\def\ch{characteristic polynomial}
\def\chp{characteristic polynomial}
\def\coe{coefficient}
\def\comp{computation}
\def\con{condition}
\def\coo{coordinate}
\def\corr{corresponding}
\def\dec{decomposition}
\def\den{denominator}
\def\dete{deterministic}
\def\eig{eigenvalue}\def\eval{eigenvalue}
\def\eigv{eigenvector}\def\eiv{eigenvector}\def\evec{eigenvector}
\def\elt{element}
\def\enn{essentially nonnegative/positivity preserving}
\def\equ{equation}
\def\etal{et al.}
\def\fe{for example}
\def\fno{from now on}
\def\for{\forall}
\def\form{formula}
\def\Fr{Furthermore, }
\def\frt{furthermore}
\def\fun{function}
\def\gene{generalization}
\def\heu{heuristic}
\def\how{however}\def\How{However, }
\def\ie{i.e. }
\def\im{\item}
\def\isd{is defined by}
\def\itc{it is easy to check that}
\def\itf{it follows that}\def\Itf{It follows that}
\def\ith{it holds that}
\def\itm{it may be checked that}\def\Itm{It may be checked that}
\def\ito{it turns out that}\def\Ito{It turns out that}\def\Itou{It turns out useful to}
\def\ma{mass action}\def\MAR{mass action representation}
\def\mat{matrix}
\def\mbc{may be checked to be}
\def\mbw{may be written as}\def\mubw{must be written as}
\def\mpr{more precisely, }\def\Mp{More precisely, }
\def\nc{non-catalytic}
\def\nne{non-negative}\def\nny{nonnegativity}
\def\no{\nonumber}
\def\noni{non-infectious}
\def\nota{notation}
\def\np{\newpage}
\def\Nt{Note that}
\def\nt{non trivial}\def\nz{non zero}
\def\oth{otherwise}
\def\Oth{On the other hand}
\def\pa{parameter}\def\para{parameter}
\def\pol{polynomial}
\def\prob{probability}
\def\Prf{{\bf  Proof:}}\def\prf{{\bf  Proof:}}
\def\prop{proportional to}
\def\qu{\quad}
\def\QED{\hfill {$\square$}\goodbreak \medskip}
\def\rep{representation}
\def\resp{respectively}
\def\satd{satisfied}\def\satg{satisfying}\def\sats{satisfies}\def\saty{satisfy}
\def\ssec{\subsection}\def\sssec{\subsubsection}
\def\sef{see for example}
\def\sksy{skew-symmetric}
\def\sub{substrate}
\def\sys{system}
\def\var{variable}
\def\wk{well-known}
\def\wl{without loss of generality}\def\wlo{w.l.o.g.}
\def\wmc{We may check that}\def\Wmc{We may check that}
\def\wmr{we may rewrite}
\def\wms{we must show that}
\def\wrt{with respect to}

% -- CRN & Chemistry --
\newcommand\CRN{chemical reaction network}
\newcommand\CRNs{chemical reaction networks}
\newcommand\RN{reaction network}
\def\ACR{absolute concentration robustness}
\def\AMSD{acyclic siphon decomposition}
\def\DAE{differential algebraic equations}
\def\FHJ{Feinberg-Horn-Jackson}
\def\GMAK{generalized mass action kinetics}
\def\MM{Michaelis-Menten}
\def\sm{stoichiometric matrix}\def\Sm{stoichiometric matrix}
\def\WRZD{weakly reversible, zero deficiency}\def\aut{autocatalytic}

% -- Epidemiology --
\def\brn{basic reproduction number}\def\bzn{basic replacement number}
\def\com{compartment}
\def\dfe{disease free equilibrium}\def\DFE{disease free equilibrium}
\def\dwt{dwell times}\def\edt{expected dwell times}
\def\EM{epidemic model}
\def\endp{endemic fixed point}\def\enp{endemic point}\def\EE{endemic point}
\def\ME{mathematical epidemiology}
\def\MEE{eco-epidemiology}\def\EEP{eco-epidemiological}
\def\ece{eco-epidemiological model}
\def\mgm{multi group model}
\def\mtS{multi-strain model}\def\mSt{multi-strain model}
\def\mSi{minimal siphon}
\def\NGM{next generation matrix}
\def\nim{new infections matrix}
\def\qem{quadratic epidemiological model}
\def\repF{invasion functions}
\def\SPF{SIR-PH-FA model}
\def\sus{susceptible}

% -- Graph theory & Linear algebra --
\def\jac{Jacobi matrix}
\def\jin{Jacobian matrix of the invasion vector field with respect to the invasion variables}
\def\jrn{Jacobian matrix of the resident vector field with respect to the resident variables}
\def\jxn{Jacobian matrix of the $\x$ vector field with respect to the $\x$ variables}
\def\Lm{Laplacian matrix}
\def\MTT{matrix tree theorem}
\def\SGT{spectral graph theory}
\def\SMW{Sherman–Morrison–Woodbury}

% -- Stability & Dynamics --
\def\absC{absorbtion CTMC}
\def\BTb{Bogdanov-Takens bifurcation}
\def\CEP{competitive exclusion partition}
\def\CTMC{continuous time Markov chain}
\def\ds{differential system}
\def\fp{fixed point}\def\bfp{boundary fixed point}\def\fbp{fixed boundary point}
\def\fbl{fixed boundary line}\def\fif{forward invariant face}
\def\gas{global asymptotic stability}\def\GAS{globally asymptotically stable}
\def\GASS{\text{GASS}}
\def\Hb{Hopf bifurcation}
\def\Hur{Hurwitz}\def\MH{Metzler-Hurwitz}\def\RH{Routh-Hurwitz}
\def\IS{\text{IS}}
\def\LAS{locally asymptotically stable}
\def\LdS{Lyapunov diagonally stable}
\def\Lf{Lyapunov function}\def\LyV{Lyapunov Volterra type function}\def\RLF{robust Lyapunov function}
\def\LyM{Lyapunov Malkin stable}
\def\PF{Perron-Frobenius}
\def\spfp{strictly positive fixed point}
\def\sta{stability}
\def\STP{sharp threshold property}\def\STT{sharp threshold theorem}

% -- Models (with \varphi explicitly instead of \f) --
\def\AB{$(A,B)$ Arino-Brauer epidemic models}
\def\ABa{$(A,\bbe,\ba,\varphi)$ balanced bilinear model}
\def\ABP{%$(A,\bbe,P,\varphi)$
\\ balanced bilinear model}
\def\AMG{Anderson May Gupta}
\newcommand\BIN{biological interaction network}
\def\Dt{Descartes type}
\def\gLV{generalized Lotka-Volterra model}
\def\IG{\text{invasion graph}}\def\IGMS{interaction graph on minimal siphons}
\def\LCP{\text{linear complementarity problem}}
\def\LV{Lotka-Volterra}\def\LVK{Lotka-Volterra-Kolmogorov}

% -- Splitting & Structure --
\def\reg{regular splitting}\def\regS{{\bf regular splitting}}
\def\RUR{rational univariate representation}
\def\bds{block-diagonal structure}\def\NGMs{next generation matrices}

% -- Miscellaneous --
\def\Ma{{Mathematica}}
\def\Mm{Max–Min}
\def\num{numerator}
\def\rap{rational approximation} \def\DSR{directed species-reaction bipartite graph}\def\PCP{polynomial complementarity problem}\def\rat{rational substitution}
\def\QSS{quasi steady state} \def\RHS{right hand side} \def\LHS{left hand side}
\def\OSN{online social network}\def\VdD{Van den Driessche}
\def\wli{$\omega$-limit} \def\rk{rich kinetics}\def\sJ{symbolic Jacobian}
\def\cV{\cite{Vas,VasSta,BSV,GSV}} \def\spf{semi-parametric semi-functional}
\def\CRW{Capasso-Ruan-Wang} \def\ABP{$(A,\bbe,P,\f)$ balanced bilinear model}
\def\ABa{%$(A,\bbe,\ba,\f)$ \\
rank-one balanced bilinear model}
\def\bbe{B%{\bff \beta}
} \def\vbe{\vec {\bff \beta}
}
\def\vb{\vec {\bff b}}
\newcommand{\Kh}{\hat{K}}          % K with hat, first mention only
\newcommand{\wS}{w^{\sigma}}       % right Perron J^\perp, full block
\newcommand{\piS}{\vec \pi^{\sigma}}    % left Perron J^\perp, full block
\newcommand{\ws}{w^{\sigma}}       % right Perron J^\perp, sub-block
\newcommand{\pis}{\pi^{\sigma}}    % left Perron J^\perp, sub-block
\newcommand{\hwS}{\hat{w}^{\sigma}}   % right Perron K, full block
\newcommand{\hpiS}{\hat{\pi}^{\sigma}} % left Perron K, full block
\newcommand{\hws}{\hat{w}^{\sigma}}   % right Perron K, sub-block
\newcommand{\hpis}{\hat{\pi}^{\sigma}} % left Perron K, sub-block
\newcommand{\twS}{\tilde{w}^{\sigma}}  % right Perron \tilde K, full block
\newcommand{\tpiS}{\tilde{\pi}^{\sigma}} % left Perron \tilde K
\newcommand{\tws}{\tilde{w}^{\sigma}}
\newcommand{\tpis}{\tilde{\pi}^{\sigma}}
\newcommand{\Mat}[2]{\mathbb{R}_{\ge 0}^{\,#1\times #2}}
\def\tJ{J_\sigma^\perp} \def\lPe{left Perron eigenvector}
% \G is \Gamma (defined at :795).  The Volterra function keeps its own name,
% so that \G no longer has two conflicting meanings in this file.
\def\GV{G}  % Volterra G(u)=u-1-\ln u
\def\PVL{Perron-Volterra Lyapunov minimal siphon ansatz}
\newcommand{\istar}[1]{i_{#1}^{*}}
%\newcommand{\sE}{s_E}
\newcommand{\vE}{v_E}
\def\PV{Perron-Volterra candidate Lyapunov function}
\def\GLV{generalized Lotka-Volterra}
\def\MERN{mathematical epidemiology reaction networks}
\def\CMM{contraction metric method}
\newcommand{\sign}{\operatorname{sign}}
\newcommand{\Var}{\operatorname{Var}} 

% ======================================================================
% IDEAS SHARED BETWEEN THE ABSTRACT AND THE MAIN RESULTS SECTION
%
% Write each idea ONCE, in the abstract; the Main Results section recalls it,
% so the two can never drift apart.
%
%   in the abstract        \ideaitem{1}{The first-order reaction terms ...}
%   in Main Results        \ideahead{1}
%
% \ideaitem stores the text globally and typesets
%   \item \hyperlink{idea:1}{<text>}
% \ideahead typesets
%   \item\hypertarget{idea:1}{}\textbf{<the same text>}
%
% The abstract must come before the Main Results input, which is the normal
% order. If a heading has to read differently from the abstract sentence (for
% instance a long abstract item summarised in one line), override it with the
% optional argument: \ideahead[short heading]{6}.
% ======================================================================
\makeatletter
\newcommand{\ideaitem}[2]{%
  \expandafter\gdef\csname @idea@#1\endcsname{#2}%
  \item\hyperlink{idea:#1}{#2}}
\newcommand{\ideahead}[2][]{%
  \item\hypertarget{idea:#2}{}%
  \def\@ideatmp{#1}%
  \ifx\@ideatmp\@empty
    \@ifundefined{@idea@#2}{%
      \PackageWarning{preamble}{Idea #2 is used in Main Results but was never
        defined by \string\ideaitem\space in the abstract}%
      \textbf{[idea #2 undefined]}%
    }{\textbf{\begingroup\def\ideaterm##1##2{##2}\def\idearest##1{}\@nameuse{@idea@#2}\endgroup}}%
  \else
    \textbf{#1}%
  \fi}
% \ideaterm{<hypertarget>}{<words>}: inside an abstract idea the words become an
% invisible link (the document sets \hypersetup{hidelinks}) jumping to the place
% in Main Results that defines them; inside the Main Results heading itself,
% where the target sits a few lines below, \ideahead strips the link and prints
% the words plainly. Mark the destination with \hypertarget{<name>}{} on the line
% before the display, which adds no equation number and shows nothing.
\newcommand{\ideaterm}[2]{\hyperlink{#1}{#2}}
% \idearest{<text>}: the part of an abstract idea after its first sentence; printed in
% the abstract, omitted from the Main Results heading.
\newcommand{\idearest}[1]{#1}
\makeatother

% ======================================================================
% SHORT-PAPER SWITCH
% \shorttrue in a master that assembles only part of the companion paper
% (Barrat/SipArxiv.tex). \nsh{...} drops text that cites results living in the
% omitted sections, so the short build has no dangling references while the
% companion paper keeps the text verbatim.
% ======================================================================
\newif\ifshort \shortfalse
\newcommand{\nsh}[1]{\ifshort\else#1\fi}
% \shortalt{<short build>}{<companion paper>}: use when a sentence needs
% different wording in the two builds, not merely a deletion.
\newcommand{\shortalt}[2]{\ifshort#1\else#2\fi}
% \Jtr, \JtrS{\Sigma}: transverse Jacobian. J^\perp here; SipArxiv.tex redefines them as J_{\mathrm{tra}}, J_{\mathrm{tra},\Sigma}
\newcommand{\Jtr}{J^\perp}
\newcommand{\JtrS}[1]{J_{#1}^\perp}
.
\makeatletter
\@ifpackageloaded{hyperref}{}{%
  \@ifundefined{hyperrefpackage}{%
    \@ifundefined{AAdefNnohyperref}{\usepackage{hyperref}}{}%
  }{}%
}
\makeatother
% linktoc=all makes a table-of-contents page number a link as well as the
% entry text; hyperref's default, linktoc=section, links only the text.
% Deferred to begin-document and guarded: by the block above hyperref may
% be loaded later (mdpi's own hook) or not at all (\AAdefNnohyperref), and
% \hypersetup must exist before it is called.  A driver that wants the
% default back can say \hypersetup{linktoc=section} after % Shared preamble for the GLVbook papers.  This is the single canonical copy;
% a driver one level below GLVbook reaches it as
%
%     \input{../basicFiles/preamble}
%
% What stays in the driver, because it differs from paper to paper:
% \documentclass, font choice, \title, \author, \date, and any macro genuinely
% local to one paper.  Everything else belongs here.  The file is
% self-contained: it does not read any other file.

\usepackage[utf8]{inputenc}
\usepackage[novbox]{pdfsync}   % required: SyncTeX alone does not work here

% Page geometry, set once here for every paper that loads this file.  A driver
% that needs different margins says \geometry{...}; a second
% \usepackage[...]{geometry} raises an option clash and must not be used.
\makeatletter
\@ifundefined{AAdefNnogeometry}{\usepackage[a4paper,margin=2cm]{geometry}}{}
\makeatother

% mathematics
\usepackage{amsmath,amssymb,amsthm,mathtools,bm,mathrsfs,bbm}
\usepackage{centernot}

% floats, tables, lists
\usepackage{graphicx,float}
\makeatletter
\@ifpackageloaded{subfig}{}{\usepackage{subcaption}}
\makeatother
\usepackage{booktabs,longtable,threeparttable,array,multirow}
\usepackage{enumitem}

% colour, code, algorithms
\usepackage{xcolor,listings,tcolorbox}
\makeatletter
\@ifundefined{c@algorithm}{\usepackage{algorithm}}{}
\makeatother
\usepackage{algorithmicx,algpseudocode}

% diagrams
\usepackage{tikz,tikz-cd,pgfplots}
\usetikzlibrary{graphs,decorations.pathmorphing,decorations.markings,
arrows.meta,automata,positioning}
\tikzcdset{every label/.append style = {font = \small}}
\pgfplotsset{compat=1.18}

% chemistry
\usepackage[version=4]{mhchem}

\usepackage{microtype}
\usepackage{lineno}
\usepackage{nomencl}
\makenomenclature

% hyperref first, then bookmark, then glossaries: glossaries must follow
% hyperref or the entry links are built against the wrong anchors.
\makeatletter
\@ifundefined{hyperrefpackage}{%
  \@ifundefined{AAdefNnohyperref}{\usepackage{hyperref}}{}%
}{}
\@ifundefined{hyperrefpackage}{\usepackage{bookmark}}{}
\makeatother
\usepackage[toc,sort=use]{glossaries}
\makeglossaries

% Book/standalone switch.  Material that the book already covers in its
% prelude, but that a standalone paper still has to state, is wrapped as
%     \ifbookbuild\else ... \fi
% Drivers default to standalone; bookM.tex says \bookbuildtrue.
\newif\ifbookbuild
\bookbuildfalse

% Part switches for the gradostat paper.  Both default to true; a driver that
% wants the short submission says \madecfalse and/or \slowfastfalse, and every
% cross-reference into the dropped part disappears with it.
%   \ifmadec     Part III, the Lagasquie--Madec model and its break-even geometry
%   \ifslowfast  Part IV, everything carrying the small parameter \eps
\newif\ifmadec     \madectrue
\newif\ifslowfast  \slowfasttrue

% Part switch for the Barrat/simplicial-contagion project.  Defaults to
% true (the hypernetwork-synchrony part is included); the journal-paper
% driver says \hnetpartfalse to drop it and every cross-reference into it.
\newif\ifhnetpart  \hnetparttrue

%% ========================================
%% PACKAGES
%% ========================================
\RequirePackage{etoolbox}
\usepackage{makecell}
\usepackage[novbox]{pdfsync}% obsolete; SyncTeX is built into modern engines
\usepackage{amssymb, amsfonts,mathtools,bbm,eurosym,graphicx,latexsym, amsmath,amsthm,times,nccmath,url}
\usepackage{listings,verbatim,cases,tabularx}
% hyperref is deliberately NOT in the line above.  Journal classes often
% load it themselves, with a long option list, from a begindocument hook
% (mdpi.cls does: \DeclareHookRule{begindocument}{hyperref}{before}{mdpi}).
% Loading it here first, without options, makes that later load raise
% "Option clash for package hyperref".  So load it only if nobody else
% has, or will:
%   - \@ifpackageloaded  catches a class that loaded it already;
%   - \hyperrefpackage    is mdpi.cls's own deferred loader;
%   - \AAdefNnohyperref   lets any other driver opt out by defining it
%                         before \input{preamble}.
\makeatletter
\@ifpackageloaded{hyperref}{}{%
  \@ifundefined{hyperrefpackage}{%
    \@ifundefined{AAdefNnohyperref}{\usepackage{hyperref}}{}%
  }{}%
}
\makeatother
% linktoc=all makes a table-of-contents page number a link as well as the
% entry text; hyperref's default, linktoc=section, links only the text.
% Deferred to begin-document and guarded: by the block above hyperref may
% be loaded later (mdpi's own hook) or not at all (\AAdefNnohyperref), and
% \hypersetup must exist before it is called.  A driver that wants the
% default back can say \hypersetup{linktoc=section} after \input{preamble}.
\makeatletter
\AtBeginDocument{\@ifundefined{hypersetup}{}{\hypersetup{linktoc=all}}}
\makeatother
\usepackage{graphicx,float,fancybox}
\usepackage{epsfig,epstopdf,tikz}
\usepackage[sans]{dsfont}
\usetikzlibrary{positioning,arrows,arrows.meta,calc}
% Configure listings for Mathematica code
\lstdefinestyle{mathematica}{
 language=Mathematica,
 basicstyle=\small\ttfamily,
 backgroundcolor=\color{gray!10},
 frame=single,
 breaklines=true,
 commentstyle=\color{green!60!black},
 keywordstyle=\color{blue},
 stringstyle=\color{red},
 showstringspaces=false,
 tabsize=2
}
\usepackage{tcolorbox}
% Colored boxes for key results
\newtcolorbox{keyresult}[1][]{
 colback=blue!5!white,
 colframe=blue!75!black,
 title=#1,
 fonttitle=\bfseries
}
\newtcolorbox{examplebox}[1][]{
 colback=green!5!white,
 colframe=green!75!black,
 title=#1,
 fonttitle=\bfseries
}


\makeatletter
\@ifundefined{theorem}   {\newtheorem{theorem}   {Theorem}}   {}
\@ifundefined{definition}{\newtheorem{definition}{Definition}} {}
\@ifundefined{open}      {\newtheorem{open}      {Problem}}   {}
\@ifundefined{example}   {\newtheorem{example}   {Example}}   {}
\makeatother

% COUNTER DEFINITIONS
\newcounter{tabcounter}
\newcounter{defcounter}
\newcounter{opcounter}
\newcounter{thmcounter}
\newcounter{excounter}
\newcounter{ascounter}
\newcounter{cjcounter}
\newcounter{lemcounter}
\newcounter{remcounter}
\newcounter{corcounter}
\newcounter{propcounter}
\newcounter{figcounter}

% UNIFIED GLOSSARY STYLE
\newglossarystyle{unifiedstyle}{%
  \setglossarystyle{list}%
  \renewcommand*{\glossarypreamble}{\gdef\glslastsec{}}%
  \renewcommand{\glossentry}[2]{%
    \edef\glscursec{\glsentryuseriii{##1}}%
    \ifx\glscursec\glslastsec\else
      \item[]\textbf{\S\,\glscursec\enspace\glsentryuseriv{##1}}%
      \global\let\glslastsec\glscursec
    \fi
    \item[\glsentryuserii{##1} \glsentryuseri{##1}:] ``\glossentryname{##1}'' \glsentrydesc{##1}%
  }%
}

\makeatletter
% First-word extractor (used by \beT, \beD and \beXa for the default key)
% \@firstwordof hello world \relax  ->  hello
\def\@firstwordof#1 #2\relax{#1}
% First three letters, used for the heading mnemonic of a theorem given no
% explicit key.  Padding with \@empty keeps it safe for first words shorter
% than three letters: \@empty vanishes under \edef, whereas \relax would
% survive into the label and is illegal inside \csname (it broke every
% two-letter title, e.g. "ME ...", with "Missing \endcsname inserted").
\def\@takethree#1#2#3#4\@nil{#1#2#3}

% EXAMPLE COMMAND -- same signature and behaviour as \beT and \beD.
% Usage: \beXa*[long title][short title][key]
%   star  = suppresses the printed "[key]" tag after the title (as for \beT*
%   and \beD*); the label and glossary entry are unaffected.
%   key = optional label key override; the label becomes e:KEY, and without a
%   key it is e: plus the first three letters of the title's first word.  A
%   mnemonic label lets an example be cited as "Example~\ref{e:KEY}" instead
%   of by section number, which is what the citing prose needs.
\NewDocumentCommand{\beXa}{s O{} O{} o}{%
  \def\@temparg{#2}%
  \ifx\@temparg\@empty
    \begin{example}%
  \else
    \let\@savedSecTitle\@currentsectionname%
    \refstepcounter{excounter}%
    \IfValueTF{#4}{%
      \edef\@exlblkey{#4}%
    }{%
      \edef\@exfirstword{\@firstwordof#2 \relax}%
      \edef\@exlblkey{%
        \expandafter\@takethree\@exfirstword\@empty\@empty\@empty\@nil}%
    }%
    \def\@tempshort{#3}%
    \ifx\@tempshort\@empty
      \expandafter\xdef\csname @xshort@\@exlblkey\endcsname{\@exlblkey}%
    \else
      \expandafter\def\csname @xshort@\@exlblkey\endcsname{#3}%
    \fi
    \IfBooleanTF{#1}{\let\@extag\@empty}{%
      \edef\@extag{\noexpand\ \noexpand\textup{[\noexpand\texttt{e:\@exlblkey}]}}}%
    \begin{example}[#2\ifshowkeys\@extag\fi]%
    \label{e:\@exlblkey}%
    % record the long title so \nref{e:KEY} prints it, as for theorems
    \protected@write\@auxout{}{%
      \string\@thmtitle{e:\@exlblkey}{\unexpanded{#2}}}%
    \expanded{%
      \noexpand\newglossaryentry{e\theexcounter}{%
        name={\unexpanded{#2}%
          \IfBooleanF{#1}{ [\noexpand\texttt{e:\@exlblkey}]}},%
        sort={\theexample},%
        description={p.~\noexpand\pageref{e:\@exlblkey}},%
        user1={\theexample},%
        user2={Example},%
        user3={\thesection},%
        user4={\expandafter\unexpanded\expandafter{\@savedSecTitle}}%
      }%
    }%
    \protected@write\@auxout{}{\string\glsliveentry{e\theexcounter}}%
  \fi
}
\newcommand{\eeXa}{\end{example}}

% DEFINITION COMMAND -- same signature and behaviour as \beT below.
% Usage: \beD[long title][short title][key]
%   long title  = definition heading (plain text only)
%   short title = shown in \dref{word}; defaults to the key
%   key         = optional label key override; the label becomes d:KEY
% Without a key the label is d: plus the first three letters of the title's
% first word, so it is a mnemonic rather than the definition's number: a
% numeric label would silently change whenever a definition is inserted.
\NewDocumentCommand{\beD}{s O{} O{} o}{%
  \def\@temparg{#2}%
  \ifx\@temparg\@empty
    \begin{definition}%
  \else
    \let\@savedSecTitle\@currentsectionname%
    \refstepcounter{defcounter}%
    \IfValueTF{#4}{%
      \edef\@deflblkey{#4}%
    }{%
      \edef\@deffirstword{\@firstwordof#2 \relax}%
      \edef\@deflblkey{%
        \expandafter\@takethree\@deffirstword\@empty\@empty\@empty\@nil}%
    }%
    \def\@tempshort{#3}%
    \ifx\@tempshort\@empty
      \expandafter\xdef\csname @dshort@\@deflblkey\endcsname{\@deflblkey}%
    \else
      \expandafter\def\csname @dshort@\@deflblkey\endcsname{#3}%
    \fi
    \IfBooleanTF{#1}{\let\@deftag\@empty}{%
      \edef\@deftag{\noexpand\ \noexpand\textup{[\noexpand\texttt{d:\@deflblkey}]}}}%
    \begin{definition}[#2\ifshowkeys\@deftag\fi]%
    \label{d:\@deflblkey}%
    % record the long title so \nref{d:KEY} prints it, as for theorems
    \protected@write\@auxout{}{%
      \string\@thmtitle{d:\@deflblkey}{\unexpanded{#2}}}%
    \expanded{%
      \noexpand\newglossaryentry{d\thedefcounter}{%
        name={\unexpanded{#2}%
          \IfBooleanF{#1}{ [\noexpand\texttt{d:\@deflblkey}]}},%
        sort={\thedefinition},%
        description={p.~\noexpand\pageref{d:\@deflblkey}},%
        user1={\thedefinition},%
        user2={Definition},%
        user3={\thesection},%
        user4={\expandafter\unexpanded\expandafter{\@savedSecTitle}}%
      }%
    }%
    \protected@write\@auxout{}{\string\glsliveentry{d\thedefcounter}}%
  \fi
}
\newcommand{\eeD}{\end{definition}}

% \showkeysfalse in a driver hides the [t:KEY] mnemonics in result headings
% (labels and glossary entries are unchanged); used for submission builds.
\newif\ifshowkeys \showkeystrue
% \commonredtrue in a driver that inputs basicFiles/commonRed: tool.tex then
% cites its reduction results instead of restating them.
\newif\ifcommonred \commonredfalse
% TheorEM COMMAND
% Usage: \beT[long title][short title][key]
%   long title  = theorem heading (plain text only)
%   short title = shown in \tref{word}; may contain math; defaults to first word of long title
%   key         = optional label key override; label becomes t:KEY (avoids first-word collisions)
% Label is t:WORD where WORD = first word of long title (or KEY if given); use \tref{WORD} to reference.
\NewDocumentCommand{\beT}{s O{} O{} o}{%
  \def\@temparg{#2}%
  \ifx\@temparg\@empty
    \begin{theorem}%
  \else
    \let\@savedSecTitle\@currentsectionname%
    \refstepcounter{thmcounter}%
    % \@thmtag is the mnemonic shown in the heading; every theorem gets one.
    % With an explicit key the key itself is shown.  Without one the key
    % defaults to the title's first word, and showing that whole word would
    % merely repeat the title ("Dissipativity [Dissipativity]"), so only its
    % first three letters are shown: "Dissipativity [Dis]".  The label is
    % unaffected -- it still uses the full \@thmlblkey.
    \IfValueTF{#4}{%
      \edef\@thmlblkey{#4}%
    }{%
      \edef\@thmfirstword{\@firstwordof#2 \relax}%
      \edef\@thmlblkey{%
        \expandafter\@takethree\@thmfirstword\@empty\@empty\@empty\@nil}%
    }%
    \IfBooleanTF{#1}{\let\@thmtag\@empty}{%
      \edef\@thmtag{\noexpand\ \noexpand\textup{[\noexpand\texttt{t:\@thmlblkey}]}}}%
    \def\@tempshort{#3}%
    \ifx\@tempshort\@empty
      \expandafter\xdef\csname @tshort@\@thmlblkey\endcsname{\@thmlblkey}%
    \else
      \expandafter\def\csname @tshort@\@thmlblkey\endcsname{#3}%
    \fi
    \begin{theorem}[#2\ifshowkeys\@thmtag\fi]%
    \label{t:\@thmlblkey}%
    % record the long title for \nref; via the .aux so forward refs resolve
    \protected@write\@auxout{}{%
      \string\@thmtitle{t:\@thmlblkey}{\unexpanded{#2}}}%
    \expanded{%
      \noexpand\newglossaryentry{t\thethmcounter}{%
        name={\unexpanded{#2}%
          \IfBooleanF{#1}{ [\noexpand\texttt{t:\@thmlblkey}]}},%
        % sort/user1 must be the number the document actually prints, i.e. the
        % theorem environment's counter, not thmcounter.  thmcounter advances
        % only on \beT with a title, so a bare \beT (or \beT with an
        % empty title) makes the two drift and the glossary cites a wrong
        % number.  thmcounter is still fine as the entry key: it is unique.
        sort={\thetheorem},%
        description={p.~\noexpand\pageref{t:\@thmlblkey}},%
        user1={\thetheorem},%
        user2={Theorem},%
        user3={\thesection},%
        user4={\expandafter\unexpanded\expandafter{\@savedSecTitle}}%
      }%
    }%
    \protected@write\@auxout{}{\string\glsliveentry{t\thethmcounter}}%
  \fi
}
\newcommand{\eeT}{\end{theorem}}

% CROSS-PART RELATION PARAGRAPH
% Usage: \xpart{text}  -- place immediately BEFORE \beT, never inside it
%   (anything inside \beT...\eeT would be swallowed into the glossary entry).
% The paragraph is typeset only when Part I is actually present, i.e. in the
% combined book build; in the standalone builds (GranadaDSABNS/A.tex,
% AAHRah-Submission/AAHRah.tex) the label part:glv is undefined and the whole
% paragraph disappears rather than leaving ?? cross-references. Same
% graceful-degradation principle as \GASref in GASt.tex. Resolves on the
% second LaTeX pass, since r@part:glv comes from the .aux file.
% (Already inside the \makeatletter region opened above, so no \makeatletter
% pair here -- adding one would close that region early.)
\newcommand{\xpart}[1]{\@ifundefined{r@part:glv}{}{\paragraph{Relation with Parts~I--II.}#1}}

% \bookref{label}{precise text}{fallback text}: prints the precise cross-reference
% when the label resolves (combined book build), and the fallback text otherwise
% (standalone paper builds, where book-only appendix labels are absent). Same
% graceful-degradation role as \GASref in GASt.tex, but usable in any file since it
% lives in this preamble (loaded by every driver) -- unlike \GASref, which is defined in a
% content file and is therefore undefined in standalone builds that omit that file.
\newcommand{\bookref}[3]{\@ifundefined{r@#1}{#3}{#2}}

% \PVpart and \RBpart: the two parts that a partial build may omit.
% They print "Part~N" in any build that contains the part, and a naming phrase
% in a build that does not, so an excerpt shows no ?? .  Each already supplies
% the word "Part", so write "\PVpart develops ...", not "Part~\PVpart".
\newcommand{\PVpart}{\bookref{part:aahrah}%
  {Part~\ref{part:aahrah}}{the Perron--Volterra part of the full book}}
\newcommand{\RBpart}{\bookref{part:recbrowne}%
  {Part~\ref{part:recbrowne}}{the immuno-virology part of the full book}}

% OPEN PROBLEM COMMANDS
% ASSUMPTION COMMANDS (glossaried)
\newcommand{\beAs}[1][]{%
  \def\@temparg{#1}%
  \ifx\@temparg\@empty
    \begin{assumption}%
  \else
    \let\@savedSecTitle\@currentsectionname%
    \refstepcounter{ascounter}%
    \begin{assumption}[#1]%
    \label{as:\theascounter}%
    \expanded{%
      \noexpand\newglossaryentry{as\theascounter}{%
        name={\unexpanded{#1}},%
        sort={\theascounter},%
        description={p.~\noexpand\pageref{as:\theascounter}},%
        user1={\theascounter},%
        user2={Assumption},%
        user3={\thesection},%
        user4={\expandafter\unexpanded\expandafter{\@savedSecTitle}}%
      }%
    }%
    \protected@write\@auxout{}{\string\glsliveentry{as\theascounter}}%
  \fi
}
\newcommand{\eeAs}{\end{assumption}}
% CONJECTURE COMMANDS (glossaried)
\newcommand{\beCj}[1][]{%
  \def\@temparg{#1}%
  \ifx\@temparg\@empty
    \begin{conjecture}%
  \else
    \let\@savedSecTitle\@currentsectionname%
    \refstepcounter{cjcounter}%
    \begin{conjecture}[#1]%
    \label{cj:\thecjcounter}%
    \expanded{%
      \noexpand\newglossaryentry{cj\thecjcounter}{%
        name={\unexpanded{#1}},%
        sort={\thecjcounter},%
        description={p.~\noexpand\pageref{cj:\thecjcounter}},%
        user1={\thecjcounter},%
        user2={Conjecture},%
        user3={\thesection},%
        user4={\expandafter\unexpanded\expandafter{\@savedSecTitle}}%
      }%
    }%
    \protected@write\@auxout{}{\string\glsliveentry{cj\thecjcounter}}%
  \fi
}
\newcommand{\eeCj}{\end{conjecture}}
% LEMMA COMMANDS
% ---------------------------------------------------------------------
% UNIFORM RESULT WRAPPERS -- \beL, \beC, \beP, \beR share \beT's signature
%
%   \beX[long title][short title][key]  ...  \eeX
%
%   long title  = heading (math allowed)
%   short title = shown in the cross-reference; defaults to the key
%   key         = label key.  With a key the label is X:KEY and the heading
%                 carries the mnemonic [key], exactly as for \beT.  Without
%                 one the label falls back to the running counter, which is
%                 what every pre-existing call site in the book relies on.
%
% A bare \beX -- arguments commented out -- opens the plain environment and
% creates NO label and NO glossary entry.  That is the supported way to keep
% an individual result out of the catalogue.
% ---------------------------------------------------------------------
% #1 wrapper name  #2 environment  #3 label prefix  #4 counter  #5 type name
\newcommand{\@mkresult}[5]{%
  % A leading * suppresses the [key] mnemonic in both the heading and the
  % glossary name, while still making the label X:KEY and the entry:
  %     \beX*[long][short][key]
  \expandafter\NewDocumentCommand\csname be#1\endcsname{s O{} O{} o}{%
    \def\@temparg{##2}%
    \ifx\@temparg\@empty
      \begin{#2}%
    \else
      \let\@savedSecTitle\@currentsectionname%
      \refstepcounter{#4}%
      % With a key: mnemonic label X:KEY and the [key] tag in the heading.
      % Without one: fall back to the running counter, exactly as these
      % wrappers behaved before, so every legacy one-argument call site in
      % the book keeps its label.  The counter is also csname-safe, which a
      % title's first word is not -- titles beginning with math would
      % otherwise build \csname \Omega...\endcsname and fail.
      \IfValueTF{##4}{%
        \edef\@relbl{##4}%
        \IfBooleanTF{##1}{\let\@retag\@empty}{%
          \edef\@retag{\noexpand\ \noexpand\textup{[\noexpand\texttt{#3:##4}]}}}%
      }{%
        \edef\@relbl{\arabic{#4}}%
        \let\@retag\@empty
      }%
      \def\@tempshort{##3}%
      \ifx\@tempshort\@empty
        \expandafter\xdef\csname @sshort@#3\@relbl\endcsname{\@relbl}%
      \else
        \expandafter\def\csname @sshort@#3\@relbl\endcsname{##3}%
      \fi
      \begin{#2}[##2\ifshowkeys\@retag\fi]%
      \edef\@tmplbl{#3:\@relbl}\expandafter\label\expandafter{\@tmplbl}%
      \expanded{%
        \noexpand\newglossaryentry{#3\csname the#4\endcsname}{%
          name={\unexpanded{##2}%
            \IfBooleanF{##1}{\IfValueT{##4}{ [\noexpand\texttt{#3:##4}]}}},%
          sort={\csname the#2\endcsname},%
          description={p.~\noexpand\pageref{#3:\@relbl}},%
          user1={\csname the#2\endcsname},%
          user2={#5},%
          user3={\thesection},%
          user4={\expandafter\unexpanded\expandafter{\@savedSecTitle}}%
        }%
      }%
      \protected@write\@auxout{}{\string\glsliveentry{#3\csname the#4\endcsname}}%
      \protected@write\@auxout{}{\string\@thmtitle{\@tmplbl}{\unexpanded{##2}}}%
    \fi
  }%
  \expandafter\newcommand\csname ee#1\endcsname{\end{#2}}%
}
\@mkresult{L}{lemma}{l}{lemcounter}{Lemma}
% REMARK / COROLLARY / PROPOSITION -- same signature as \beL above
\@mkresult{R}{remark}{r}{remcounter}{Remark}
\@mkresult{C}{corollary}{c}{corcounter}{Corollary}
\@mkresult{P}{proposition}{p}{propcounter}{Proposition}
% OPEN-PROBLEM COMMAND -- generated with the same signature as \beT:
%   \beO[long title][short title][key]   ...   \eeO
% label o:KEY, or o:<counter> when no key is given; a bare \beO makes
% neither a label nor a glossary entry.
\@mkresult{O}{open}{o}{opcounter}{Problem}


% Section-level title tracker: updated only by \section, not \subsection.
% \@currentlabelname is overwritten by every \refstepcounter (including
% \subsection), so we save it separately here.
\newcommand{\@currentsectionname}{}
\AtBeginDocument{%
  \let\@epidcrn@orig@refstepcounter\refstepcounter
  \renewcommand{\refstepcounter}[1]{%
    \@epidcrn@orig@refstepcounter{#1}%
    \ifstrequal{#1}{section}{\global\let\@currentsectionname\@currentlabelname}{}%
  }%
}

% TABLE COMMANDS (glossaried) -- consistent naming with \beT, \beD, etc.
\newcommand{\beTa}[1][]{%
  \def\@temparg{#1}%
  \ifx\@temparg\@empty
    \begin{table}[htbp]%
  \else
    \let\@savedSecTitle\@currentsectionname%
    \refstepcounter{tabcounter}%
    \begin{table}[htbp]%
    \label{A:\thetabcounter}%
    \expanded{%
      \noexpand\newglossaryentry{A\thetabcounter}{%
        name={\unexpanded{#1}},%
        sort={\thetabcounter},%
        description={p.~\noexpand\pageref{A:\thetabcounter}},%
        user1={\thetabcounter},%
        user2={Table},%
        user3={\thesection},%
        user4={\expandafter\unexpanded\expandafter{\@savedSecTitle}}%
      }%
    }%
    \protected@write\@auxout{}{\string\glsliveentry{A\thetabcounter}}%
  \fi
}
\newcommand{\eeTa}{\end{table}}

% FIGURE COMMANDS (glossaried) -- consistent with \beT, \beR, \beTa, etc.
% Usage:
%   \beF[Glossary title]...\caption{Visible caption}\eeF  -> glossaried figure
%   \beF...\caption{Visible caption}\eeF                  -> plain [H] figure, no glossary
\newcommand{\beF}[1][]{%
  \def\@temparg{#1}%
  \ifx\@temparg\@empty
    \begin{figure}[H]%
  \else
    \begin{figure}[H]%
    \let\@savedSecTitle\@currentsectionname%
    \refstepcounter{figcounter}%
    \label{fig:\thefigcounter}%
    \expanded{%
      \noexpand\newglossaryentry{Fig\thefigcounter}{%
        name={\unexpanded{#1}},%
        sort={\thefigcounter},%
        description={p.~\noexpand\pageref{fig:\thefigcounter}},%
        user1={\thefigcounter},%
        user2={Figure},%
        user3={\thesection},%
        user4={\expandafter\unexpanded\expandafter{\@savedSecTitle}}%
      }%
    }%
    \protected@write\@auxout{}{\string\glsliveentry{Fig\thefigcounter}}%
  \fi
}
\newcommand{\eeF}{\end{figure}}
\makeatother
% LIVE-ENTRY REGISTER
% \glsaddall prints every entry the glossaries package currently knows, and
% that set includes entries restored from \jobname.glsdefs.  The package
% rewrites .glsdefs from all *currently defined* entries, so anything it
% restored is written straight back out: a deleted section kept reappearing in
% the catalogue for ever, and the .glsdefs had to be erased by hand.  The
% wrappers therefore record every entry they actually create in the .aux, and
% \printboth adds only those.  The .aux is rebuilt from scratch on each run, so
% this is self-healing: a removed result is gone after one further run.
\makeatletter
\newcommand{\@glslive}{}
\newcommand{\@glsaddlive}[1]{\ifglsentryexists{#1}{\glsadd{#1}}{}}
\newcommand{\glsliveentry}[1]{\g@addto@macro\@glslive{\@glsaddlive{#1}}}
\makeatother
% PRINT COMMAND
% No \clearpage here -- the catalogue continues on the current page.  This has
% been re-added by mistake several times; leave it out.
\makeatletter
\newcommand{\printboth}{%
  \@glslive   % add exactly the entries this document created last run
  \printglossary[style=unifiedstyle, title={0 \hspace{.22em} Definitions, Theorems, Examples and Open Problems}]%
}
\makeatother
% Named theorem reference: \tref{word} -> Thm~N(word-or-short)
% Shows stored short title if available, otherwise shows the key itself.
\newcommand{\tref}[1]{%
  Thm~\ref{t:#1}(\ifcsname @tshort@#1\endcsname\csname @tshort@#1\endcsname\else#1\fi)%
}

% Full-title theorem reference: \nref{t:KEY} prints the number and then the
% theorem's long title, so one writes "Theorem~\nref{t:blocks}" and gets
% "Theorem 3 ``Block factorization and siphons''".  Takes the complete label, like
% \ref, not the bare key like \tref.
%
% \beT writes each title to the .aux (see \@thmtitle below), so \nref resolves
% forward references from the second LaTeX pass, exactly as \ref does.  If the
% title is not yet known -- first pass, or a label that is not a \beT theorem
% -- it degrades to the number rather than producing ??.
% These two use @ in their names, and this point in the file is outside the
% \makeatletter region opened earlier (it closes at \makeatother above), so
% the region has to be reopened here.
\makeatletter
\newcommand{\@thmtitle}[2]{\global\@namedef{@tlong@#1}{#2}}
\newcommand{\nref}[1]{%
  \ref{#1}%
  \ifcsname @tlong@#1\endcsname
    \ ``\csname @tlong@#1\endcsname''%
  \fi
}
% \nm{t:KEY}: the stored long title of any \beT/\beL/\beP/\beR/\beC/\beD/\beE result,
% verbatim and without quotes, so item names follow the macro's first argument.
\newcommand{\nm}[1]{%
  \ifcsname @tlong@#1\endcsname
    \csname @tlong@#1\endcsname%
  \else
    \ref{#1}%
  \fi
}
\makeatother




%% ========================================
%% Other THEOREM ENVIRONMENTS
%% ========================================
\makeatletter
\@ifundefined{lemma}      {\newtheorem{lemma}      {Lemma}}      {}
\@ifundefined{proposition}{\newtheorem{proposition}{Proposition}} {}
\@ifundefined{corollary}  {\newtheorem{corollary}  {Corollary}}  {}
\@ifundefined{conjecture} {\newtheorem{conjecture} {Conjecture}}  {}
\@ifundefined{remark}     {\newtheorem{remark}     {Remark}}     {}
\@ifundefined{question}   {\newtheorem{question}   {Question}}   {}
\@ifundefined{assumption} {\newtheorem{assumption} {Assumption}} {}
\@ifundefined{hypothesis} {\newtheorem{hypothesis} {Hypothesis}} {}
\@ifundefined{fact}       {\newtheorem{fact}       {Fact}}       {}
\makeatother

% \citer: cite the relay-graph threshold theorem of AAHrel. If a document has
% \input the chapter/file containing t:relayRF (so the label resolves), point
% straight at it; otherwise fall back to a bare citation of AAHrel. Shared here
% (not per-paper) so it resolves consistently across every book chapter.
% \citeR: same idea, for the first theorem of AAHRah (t:Pstar).
\makeatletter
\providecommand{\citer}{%
  \@ifundefined{r@t:relayRF}%
    {\cite{AAHrel}}%
    {Theorem~\ref{t:relayRF} of \cite{AAHrel}}%
}
\providecommand{\citeR}{%
  \@ifundefined{r@t:Pstar}%
    {\cite{AAHRah}}%
    {Theorem~\ref{t:Pstar} of \cite{AAHRah}}%
}
\makeatother


% Abbreviations
\def\beQ{\begin{question}}\def\eeQ{\end{question}}
\def\beA{\begin{assumption}}\def\eeA{\end{assumption}}


%% ========================================

%% ========================================
%% COLORS & CUSTOM COMMANDS
%% ========================================
\definecolor{funccolor}{RGB}{25,25,112}
\definecolor{desccolor}{RGB}{64,64,64}
\newcommand{\func}[1]{\texttt{\textcolor{funccolor}{#1}}}
\def\eq{=}
\def\mm{-}

%% ========================================
%% MATHCAL SHORTHANDS
%% ========================================
% Project-wide convention: \mX for a capital X is always \mathcal X.
% Defined for the whole alphabet so the convention never fails to expand.
\newcommand{\mA}{\mathcal A}\newcommand{\mB}{\mathcal B}
\newcommand{\mC}{\mathcal C}\newcommand{\mD}{\mathcal D}
\newcommand{\mE}{\mathcal E}\newcommand{\mF}{\mathcal F}
\newcommand{\mG}{\mathcal G}\newcommand{\mH}{\mathcal H}
\newcommand{\mI}{\mathcal I}\newcommand{\mJ}{\mathcal J}
\newcommand{\mK}{\mathcal K}\newcommand{\mL}{\mathcal L}
\newcommand{\mM}{\mathcal M}\newcommand{\mN}{\mathcal N}
\newcommand{\mO}{\mathcal O}\newcommand{\mP}{\mathcal P}
\newcommand{\mQ}{\mathcal Q}\newcommand{\mR}{\mathcal R}
\newcommand{\mS}{\mathcal S}\newcommand{\mT}{\mathcal T}
\newcommand{\mU}{\mathcal U}\newcommand{\mV}{\mathcal V}
\newcommand{\mW}{\mathcal W}\newcommand{\mX}{\mathcal X}
\newcommand{\mY}{\mathcal Y}\newcommand{\mZ}{\mathcal Z}

%% ========================================
%% BASIC MATH ENVIRONMENTS
%% ========================================
\def\bep{\begin{pmatrix}}\def\eep{\end{pmatrix}}
\def\bev{\begin{vmatrix}}\def\eev{\end{vmatrix}}
\def\bea{\begin{eqnarray*}}\def\eea{\end{eqnarray*}}
\def\bc{\begin{cases}}\def\ec{\end{cases}}
\def\BEN{\begin{enumerate}}\def\EEN{\end{enumerate}}
\def\BI{\begin{itemize}}\def\EI{\end{itemize}}

\newcommand{\be}[1]{\begin{equation}\label{#1}}
\newcommand{\ee}{\end{equation}}
\newcommand{\beq}{\begin{eqnarray}}
\def\eeq{\end{eqnarray}}
\def\eqr{\eqref}\def\fr{\frac}\def\lbl{\label}
\def\Lra{\Longrightarrow}\def\non{\nonumber}\def\cNGM{\cite{Diek,Van,Van08}}
\def\Eq{\Leftrightarrow}\def\b{\beta}\def\uph{\red{UP TO HERE}}
%% ========================================
%% PROBABILITY & STATISTICS
%% ========================================
\renewcommand{\P}{\ensuremath{\mathbb{P}}}  % Note: overrides paragraph symbol
\newcommand{\E}{\ensuremath{\mathbb{E}}}
\newcommand{\ind}[1]{\ensuremath{\mathbbm{1}_{\left\{#1\right\}}}}
\DeclareMathOperator{\Cov}{Cov}
\DeclareMathOperator{\TV}{TV}
\newcommand{\rank}{\operatorname{rank}}

%% ========================================
%% CRN OPERATORS (organizations/ERCs/buffering/siphons framework)
%% ========================================
\def\clos{\mathrm{clos}}
\def\Supply{\mathrm{Supply}}
\def\req{\mathrm{req}}
\def\React{\mathrm{React}}
\def\Prod{\mathrm{Prod}}
\def\Gen{\mathrm{Gen}}
\def\In{\mathrm{In}}
\def\Out{\mathrm{Out}}
\def\Act{\mathrm{Act}}
\def\Min{\mathrm{Min}}

%% ========================================
%% LINE BREAKING
%% ========================================
% Last-resort stretch, used only on paragraphs TeX cannot otherwise break:
% a long \texttt{} token, a comma-separated inline math list (commas are not
% breakpoints in math), or a run of \cite/\ref boxes.  Without it those set as
% overfull \hbox and print into the margin; with it they set slightly loose,
% which is the lesser evil.  Does not affect paragraphs that already break.
\emergencystretch=3em

%% ========================================
%% NUMBER SETS
%% ========================================
\newcommand{\R}{\mathbb{R}}  % Real numbers
\newcommand{\C}{\mathbb{C}}  % Complex numbers
\newcommand{\N}{\mathbb{N}}
\newcommand{\Z}{\mathbb{Z}}
% \providecommand, not \newcommand: some papers define these in their own
% preamble after \input{preamble}, and must not clash.
\providecommand{\Q}{\mathbb{Q}}
\providecommand{\res}{\operatorname{res}}   % resultant; used by epiRes.tex
\newcommand{\MN}{\mathbb{N}}
\newcommand{\MZ}{\mathbb{Z}}
\newcommand{\realpos}{\mathbb{R}_{> 0}}
\newcommand{\realnneg}{\mathbb{R}_{\geq 0}}
\newcommand{\naturalpos}{\mathbb{N}_{>0}}
\newcommand{\Rplus}{\mathbb{R}_+}
\newcommand{\Rp}{\mathbb{R}_{\ge 0}}
\def\Rnne{\mathbb{R}_{\ge 0}}
\def\Se{\bar{S}} \def\Ie{\bar{I}}
\newcommand{\Rpos}{\mathbb{R}_{>0}}
\newcommand{\dd}{\,\mathrm d}
\newcommand{\e}{\mathrm e}
%% ========================================
%% GREEK LETTERS (2+ character abbreviations) greek
%% ========================================
\def\al{\alpha}\def\ba{\bff \alpha}
\def\w{\omega}\def\Lgn{\ga+\mu+\de}\def\om{\omega}
\def\de{\delta}\def\De{\Delta}\def\bde{\bff \delta}
\def\ep{\epsilon}\def\eps{\varepsilon}
\def\f{\varphi}  % Kept as requested
\def\ga{\gamma}\def\Ga{\Gamma}\def\G{\Gamma}\def\bg{\bff \gamma}
\def\up{\upsilon}  % NOT the breve accent: no accent commands anywhere.
                   % Named \up, not \u, so the LaTeX breve stays reachable.
\def\la{\lambda}\def\La{\Lambda}
\def\si{\sigma}\def\ka{\kappa}
\def\th{\theta}\def\vt{\vartheta}
\def\ze{\zeta}\def\Gap{\Gamma^{+}} \def\Gam{\Gamma^{-}}


%% ========================================
%% SPECIAL SHORT COMMANDS (kept as requested)
%% ========================================
\def\I{\infty}
\def\T{\widetilde}
\def\H{\hat}
\newcommand{\relint}{\operatorname{relint}}

%% ========================================
%% BOLD/VECTOR NOTATION
%% ========================================
\newcommand{\bff}[1]{{\mbox{\boldmath$#1$}}}
\newcommand\X{\boldsymbol{X}}
\def\x{\boldsymbol{x}}\def\bx{\bff x}\def\xe{\bar{\x}}
\newcommand\y{\boldsymbol{y}}\def\by{\boldsymbol{y}}\def\ye{\bar{\y}}
\def\yb{\y}\def\xb{\x}
\def\xy{\x-\y}
\newcommand\br{\boldsymbol{r}}\def\re{\bar{\r}}
\newcommand\bd{\boldsymbol{d}}
\newcommand\bQ{\boldsymbol{Q}}
\newcommand{\Yb}{{\mathbf{Y}}}
\newcommand{\Ab}{\mathbf{A}}
\newcommand{\zb}{\mathbf{z}}\def\bz{\bff 0}
\newcommand{\Zb}{\mathbf{Z}}
\newcommand{\bb}{\mathbf{b}}
\def\vbe{\vec {\bff \beta}}
\def\vb{\vec {\bff b}}
\def\v1{\vec {\bff 1}}
\def\vn{\vec \de}
\def\vz{\vec 0}
\def\vi{\overset{\rightarrow}{\mathsf i}}
\def\vv{\overset{\rightarrow}{v}}
\def\vr{\; \vec {\mathsf r}}
\newcommand{\va}{\bff a}
\providecommand{\vect}[1]{\boldsymbol{#1}}   % bold vector (moved from Gradostat.tex)
\providecommand{\one}{\mathbf 1}             % all-ones vector (moved from AliSF)
\providecommand{\spr}{s}                     % spectral bound  (moved from AliSF)
\providecommand{\Ke}{K_{\rm eff}}            % effective carrying capacity
% Critical right/left eigenvectors in the Hopf projection method (AAHOSN/New.tex)
\providecommand{\bq}{\mathbf{q}}
\providecommand{\bp}{\mathbf{p}}


%% ========================================
%% MATHCAL & MATHSCR
%% ========================================
% Mathcal
\newcommand\mc[1]{\mathcal{#1}}
\def\mC{{\mathcal C}}\def\mD{\mathcal D}\def\mE{\mathcal{E}}
\def\mF{\mathcal F}\def\mG{{\mathcal G}}\def\mH{{\mathcal H}}
\def\mI{\mathcal I}\def\mL{\mathcal{L}}\def\mM{{\mathcal M}}
\def\mN{{\mathcal N}}\def\mP{\mathcal{P}}\def\mR{{\mathcal R}}
\def\mS{{\mathcal S}}\def\mT{\mathcal{T}}\def\mV{{\mathcal V}}
\def\mK{\mathcal K}\def\mA{{\mathcal A}}\def\mU{{\mathcal U}}
% Symmetric bilinear and trilinear forms of the Taylor expansion, used in the
% first Lyapunov coefficient (AAHOSN/New.tex)
\providecommand{\Bform}{\mathcal{B}}
\providecommand{\Cform}{\mathcal{C}}

% Blackboard bold.  Kept here, not in the drivers: a \providecommand in a
% driver preamble runs before this file and silently wins, so the same name
% ends up meaning different things in different builds.
\def\IC{\mathbb{C}}\def\IN{\mathbb{N}}\def\IQ{\mathbb{Q}}
\def\IR{\mathbb{R}}\def\IZ{\mathbb{Z}}\def\IG{\mathbb{G}}
\def\II{\mathbb{I}}

% Mathscr
\newcommand{\Rs}{{\mathscr R}}
\newcommand{\Ss}{{\mathscr S}}
\newcommand{\Ns}{{\mathscr N}}
\newcommand{\Cs}{\mathscr C}


%% ========================================
%% OPERATORS & FUNCTIONS
%% ========================================
\DeclareMathOperator{\tr}{tr}
\DeclareMathOperator{\diag}{Diag}
\DeclareMathOperator{\Diag}{Diag}
\DeclareMathOperator{\intr}{\mathrm{Int}}
\DeclareMathOperator{\Inv}{Inv}
\DeclareMathOperator{\spec}{spec}
\DeclareMathOperator{\AUC}{AUC}               % area under the curve
\newcommand{\ii}{\mathrm{i}}                  % imaginary unit
\newcommand\op{\operatorname}
\newcommand{\supp}{\op{supp}}
\newcommand{\diff}{\mathop{}\mathopen{}\mathrm{d}}
\newcommand{\pdd}[1]{\frac{\partial}{\partial{#1}}}
\DeclareMathOperator{\cone}{cone}
\DeclareMathOperator{\proj}{proj}

%% ========================================
%% BRACKETS & NORMS
%% ========================================
\providecommand{\pr}[1]{\left(#1\right)}
\providecommand{\pp}[1]{\left[#1\right]}
\newcommand\croc[1]{\left\langle #1\right\rangle}
\newcommand{\norm}[1]{\|#1\|}
\newcommand{\inner}[2]{\langle #1, #2 \rangle}
% Short synonym of \inner, used by AAHOSN/New.tex as $\ip{u}{v}=\bar u^\top v$
\providecommand{\ip}[2]{\langle #1 , #2 \rangle}


%% ========================================
%% FORMATTING & UTILITIES
%% ========================================
\newcommand*{\Scale}[2][4]{\scalebox{#1}{$#2$}}
\newcommand{\te}[1]{\text{#1}}
\newcommand\steq[1]{\stackrel{\text{\rm #1.}}{=}}
\newcommand\blue[1]{\textcolor{blue}{#1}}
\newcommand{\red}{\textcolor[rgb]{1.00,0.00,0.00}}
\newcommand{\green}{\textcolor[rgb]{0.00,1.00,0.00}}

% Figure command
\newcommand{\figu}[3]{%
\begin{figure}[H]
\centering
\includegraphics[scale=#3]{#1}
\caption{#2\label{f:#1}}
\end{figure}
}

% Footnote
\long\def\symbolfootnote[#1]#2{%
\begingroup
\def\thefootnote{\fnsymbol{footnote}}\footnote[#1]{#2}%
\endgroup}
\def\fn{\symbolfootnote}

\newcommand{\dsum}{\displaystyle\sum}

%% ========================================
%% CRN & EPIDEMIOLOGY SPECIFIC
%% ========================================
% CRN notation
\def\SS{\mathcal S}
\def\CC{\mathcal C}
\def\RR{\mathcal R}

% Reproduction numbers
\def\Rto{R_{1}^{E_2}}
\def\Rot{R_{2}^{E_1}}
\def\Roc{R_{1,c}}
\def\Rtc{R_{2,c}}
\def\rd{r_{0}}
\def\sd{s_{0}}
\def\dd{D_{0}}
\def\vd{v_{0}}
\def\zd{z_{0}}
\def\se{\sE}
\def\xd{x_{0}}
%% ========================================
%% SPECIAL VALUES & EQUILIBRIA
%% ========================================
\def\eE{\bar{e}}
\def\wE{\bar{w}}
\def\sE{s_{*}}
\def\Me{\T M_{ee}}
\def\Rnne{\mathbb{R}_{\ge 0}}

\newcommand{\ones}{\mathbf 1}


%% ========================================
%% ABBREVIATIONS (alphabetical by category)
%% ========================================

% -- General mathematical terms --
\def\adm{admissible}\def\admy{admissibility}
%\def\and{antisymmetric}% disabled 2026-09-30: it overrode LaTeX's author separator \and (printed "antisymmetric" between authors)
\def\ass{assumption}
\def\cadlag{c\`adl\`ag}
\def\ch{characteristic polynomial}
\def\chp{characteristic polynomial}
\def\coe{coefficient}
\def\comp{computation}
\def\con{condition}
\def\coo{coordinate}
\def\corr{corresponding}
\def\dec{decomposition}
\def\den{denominator}
\def\dete{deterministic}
\def\eig{eigenvalue}\def\eval{eigenvalue}
\def\eigv{eigenvector}\def\eiv{eigenvector}\def\evec{eigenvector}
\def\elt{element}
\def\enn{essentially nonnegative/positivity preserving}
\def\equ{equation}
\def\etal{et al.}
\def\fe{for example}
\def\fno{from now on}
\def\for{\forall}
\def\form{formula}
\def\Fr{Furthermore, }
\def\frt{furthermore}
\def\fun{function}
\def\gene{generalization}
\def\heu{heuristic}
\def\how{however}\def\How{However, }
\def\ie{i.e. }
\def\im{\item}
\def\isd{is defined by}
\def\itc{it is easy to check that}
\def\itf{it follows that}\def\Itf{It follows that}
\def\ith{it holds that}
\def\itm{it may be checked that}\def\Itm{It may be checked that}
\def\ito{it turns out that}\def\Ito{It turns out that}\def\Itou{It turns out useful to}
\def\ma{mass action}\def\MAR{mass action representation}
\def\mat{matrix}
\def\mbc{may be checked to be}
\def\mbw{may be written as}\def\mubw{must be written as}
\def\mpr{more precisely, }\def\Mp{More precisely, }
\def\nc{non-catalytic}
\def\nne{non-negative}\def\nny{nonnegativity}
\def\no{\nonumber}
\def\noni{non-infectious}
\def\nota{notation}
\def\np{\newpage}
\def\Nt{Note that}
\def\nt{non trivial}\def\nz{non zero}
\def\oth{otherwise}
\def\Oth{On the other hand}
\def\pa{parameter}\def\para{parameter}
\def\pol{polynomial}
\def\prob{probability}
\def\Prf{{\bf  Proof:}}\def\prf{{\bf  Proof:}}
\def\prop{proportional to}
\def\qu{\quad}
\def\QED{\hfill {$\square$}\goodbreak \medskip}
\def\rep{representation}
\def\resp{respectively}
\def\satd{satisfied}\def\satg{satisfying}\def\sats{satisfies}\def\saty{satisfy}
\def\ssec{\subsection}\def\sssec{\subsubsection}
\def\sef{see for example}
\def\sksy{skew-symmetric}
\def\sub{substrate}
\def\sys{system}
\def\var{variable}
\def\wk{well-known}
\def\wl{without loss of generality}\def\wlo{w.l.o.g.}
\def\wmc{We may check that}\def\Wmc{We may check that}
\def\wmr{we may rewrite}
\def\wms{we must show that}
\def\wrt{with respect to}

% -- CRN & Chemistry --
\newcommand\CRN{chemical reaction network}
\newcommand\CRNs{chemical reaction networks}
\newcommand\RN{reaction network}
\def\ACR{absolute concentration robustness}
\def\AMSD{acyclic siphon decomposition}
\def\DAE{differential algebraic equations}
\def\FHJ{Feinberg-Horn-Jackson}
\def\GMAK{generalized mass action kinetics}
\def\MM{Michaelis-Menten}
\def\sm{stoichiometric matrix}\def\Sm{stoichiometric matrix}
\def\WRZD{weakly reversible, zero deficiency}\def\aut{autocatalytic}

% -- Epidemiology --
\def\brn{basic reproduction number}\def\bzn{basic replacement number}
\def\com{compartment}
\def\dfe{disease free equilibrium}\def\DFE{disease free equilibrium}
\def\dwt{dwell times}\def\edt{expected dwell times}
\def\EM{epidemic model}
\def\endp{endemic fixed point}\def\enp{endemic point}\def\EE{endemic point}
\def\ME{mathematical epidemiology}
\def\MEE{eco-epidemiology}\def\EEP{eco-epidemiological}
\def\ece{eco-epidemiological model}
\def\mgm{multi group model}
\def\mtS{multi-strain model}\def\mSt{multi-strain model}
\def\mSi{minimal siphon}
\def\NGM{next generation matrix}
\def\nim{new infections matrix}
\def\qem{quadratic epidemiological model}
\def\repF{invasion functions}
\def\SPF{SIR-PH-FA model}
\def\sus{susceptible}

% -- Graph theory & Linear algebra --
\def\jac{Jacobi matrix}
\def\jin{Jacobian matrix of the invasion vector field with respect to the invasion variables}
\def\jrn{Jacobian matrix of the resident vector field with respect to the resident variables}
\def\jxn{Jacobian matrix of the $\x$ vector field with respect to the $\x$ variables}
\def\Lm{Laplacian matrix}
\def\MTT{matrix tree theorem}
\def\SGT{spectral graph theory}
\def\SMW{Sherman–Morrison–Woodbury}

% -- Stability & Dynamics --
\def\absC{absorbtion CTMC}
\def\BTb{Bogdanov-Takens bifurcation}
\def\CEP{competitive exclusion partition}
\def\CTMC{continuous time Markov chain}
\def\ds{differential system}
\def\fp{fixed point}\def\bfp{boundary fixed point}\def\fbp{fixed boundary point}
\def\fbl{fixed boundary line}\def\fif{forward invariant face}
\def\gas{global asymptotic stability}\def\GAS{globally asymptotically stable}
\def\GASS{\text{GASS}}
\def\Hb{Hopf bifurcation}
\def\Hur{Hurwitz}\def\MH{Metzler-Hurwitz}\def\RH{Routh-Hurwitz}
\def\IS{\text{IS}}
\def\LAS{locally asymptotically stable}
\def\LdS{Lyapunov diagonally stable}
\def\Lf{Lyapunov function}\def\LyV{Lyapunov Volterra type function}\def\RLF{robust Lyapunov function}
\def\LyM{Lyapunov Malkin stable}
\def\PF{Perron-Frobenius}
\def\spfp{strictly positive fixed point}
\def\sta{stability}
\def\STP{sharp threshold property}\def\STT{sharp threshold theorem}

% -- Models (with \varphi explicitly instead of \f) --
\def\AB{$(A,B)$ Arino-Brauer epidemic models}
\def\ABa{$(A,\bbe,\ba,\varphi)$ balanced bilinear model}
\def\ABP{%$(A,\bbe,P,\varphi)$
\\ balanced bilinear model}
\def\AMG{Anderson May Gupta}
\newcommand\BIN{biological interaction network}
\def\Dt{Descartes type}
\def\gLV{generalized Lotka-Volterra model}
\def\IG{\text{invasion graph}}\def\IGMS{interaction graph on minimal siphons}
\def\LCP{\text{linear complementarity problem}}
\def\LV{Lotka-Volterra}\def\LVK{Lotka-Volterra-Kolmogorov}

% -- Splitting & Structure --
\def\reg{regular splitting}\def\regS{{\bf regular splitting}}
\def\RUR{rational univariate representation}
\def\bds{block-diagonal structure}\def\NGMs{next generation matrices}

% -- Miscellaneous --
\def\Ma{{Mathematica}}
\def\Mm{Max–Min}
\def\num{numerator}
\def\rap{rational approximation} \def\DSR{directed species-reaction bipartite graph}\def\PCP{polynomial complementarity problem}\def\rat{rational substitution}
\def\QSS{quasi steady state} \def\RHS{right hand side} \def\LHS{left hand side}
\def\OSN{online social network}\def\VdD{Van den Driessche}
\def\wli{$\omega$-limit} \def\rk{rich kinetics}\def\sJ{symbolic Jacobian}
\def\cV{\cite{Vas,VasSta,BSV,GSV}} \def\spf{semi-parametric semi-functional}
\def\CRW{Capasso-Ruan-Wang} \def\ABP{$(A,\bbe,P,\f)$ balanced bilinear model}
\def\ABa{%$(A,\bbe,\ba,\f)$ \\
rank-one balanced bilinear model}
\def\bbe{B%{\bff \beta}
} \def\vbe{\vec {\bff \beta}
}
\def\vb{\vec {\bff b}}
\newcommand{\Kh}{\hat{K}}          % K with hat, first mention only
\newcommand{\wS}{w^{\sigma}}       % right Perron J^\perp, full block
\newcommand{\piS}{\vec \pi^{\sigma}}    % left Perron J^\perp, full block
\newcommand{\ws}{w^{\sigma}}       % right Perron J^\perp, sub-block
\newcommand{\pis}{\pi^{\sigma}}    % left Perron J^\perp, sub-block
\newcommand{\hwS}{\hat{w}^{\sigma}}   % right Perron K, full block
\newcommand{\hpiS}{\hat{\pi}^{\sigma}} % left Perron K, full block
\newcommand{\hws}{\hat{w}^{\sigma}}   % right Perron K, sub-block
\newcommand{\hpis}{\hat{\pi}^{\sigma}} % left Perron K, sub-block
\newcommand{\twS}{\tilde{w}^{\sigma}}  % right Perron \tilde K, full block
\newcommand{\tpiS}{\tilde{\pi}^{\sigma}} % left Perron \tilde K
\newcommand{\tws}{\tilde{w}^{\sigma}}
\newcommand{\tpis}{\tilde{\pi}^{\sigma}}
\newcommand{\Mat}[2]{\mathbb{R}_{\ge 0}^{\,#1\times #2}}
\def\tJ{J_\sigma^\perp} \def\lPe{left Perron eigenvector}
% \G is \Gamma (defined at :795).  The Volterra function keeps its own name,
% so that \G no longer has two conflicting meanings in this file.
\def\GV{G}  % Volterra G(u)=u-1-\ln u
\def\PVL{Perron-Volterra Lyapunov minimal siphon ansatz}
\newcommand{\istar}[1]{i_{#1}^{*}}
%\newcommand{\sE}{s_E}
\newcommand{\vE}{v_E}
\def\PV{Perron-Volterra candidate Lyapunov function}
\def\GLV{generalized Lotka-Volterra}
\def\MERN{mathematical epidemiology reaction networks}
\def\CMM{contraction metric method}
\newcommand{\sign}{\operatorname{sign}}
\newcommand{\Var}{\operatorname{Var}} 

% ======================================================================
% IDEAS SHARED BETWEEN THE ABSTRACT AND THE MAIN RESULTS SECTION
%
% Write each idea ONCE, in the abstract; the Main Results section recalls it,
% so the two can never drift apart.
%
%   in the abstract        \ideaitem{1}{The first-order reaction terms ...}
%   in Main Results        \ideahead{1}
%
% \ideaitem stores the text globally and typesets
%   \item \hyperlink{idea:1}{<text>}
% \ideahead typesets
%   \item\hypertarget{idea:1}{}\textbf{<the same text>}
%
% The abstract must come before the Main Results input, which is the normal
% order. If a heading has to read differently from the abstract sentence (for
% instance a long abstract item summarised in one line), override it with the
% optional argument: \ideahead[short heading]{6}.
% ======================================================================
\makeatletter
\newcommand{\ideaitem}[2]{%
  \expandafter\gdef\csname @idea@#1\endcsname{#2}%
  \item\hyperlink{idea:#1}{#2}}
\newcommand{\ideahead}[2][]{%
  \item\hypertarget{idea:#2}{}%
  \def\@ideatmp{#1}%
  \ifx\@ideatmp\@empty
    \@ifundefined{@idea@#2}{%
      \PackageWarning{preamble}{Idea #2 is used in Main Results but was never
        defined by \string\ideaitem\space in the abstract}%
      \textbf{[idea #2 undefined]}%
    }{\textbf{\begingroup\def\ideaterm##1##2{##2}\def\idearest##1{}\@nameuse{@idea@#2}\endgroup}}%
  \else
    \textbf{#1}%
  \fi}
% \ideaterm{<hypertarget>}{<words>}: inside an abstract idea the words become an
% invisible link (the document sets \hypersetup{hidelinks}) jumping to the place
% in Main Results that defines them; inside the Main Results heading itself,
% where the target sits a few lines below, \ideahead strips the link and prints
% the words plainly. Mark the destination with \hypertarget{<name>}{} on the line
% before the display, which adds no equation number and shows nothing.
\newcommand{\ideaterm}[2]{\hyperlink{#1}{#2}}
% \idearest{<text>}: the part of an abstract idea after its first sentence; printed in
% the abstract, omitted from the Main Results heading.
\newcommand{\idearest}[1]{#1}
\makeatother

% ======================================================================
% SHORT-PAPER SWITCH
% \shorttrue in a master that assembles only part of the companion paper
% (Barrat/SipArxiv.tex). \nsh{...} drops text that cites results living in the
% omitted sections, so the short build has no dangling references while the
% companion paper keeps the text verbatim.
% ======================================================================
\newif\ifshort \shortfalse
\newcommand{\nsh}[1]{\ifshort\else#1\fi}
% \shortalt{<short build>}{<companion paper>}: use when a sentence needs
% different wording in the two builds, not merely a deletion.
\newcommand{\shortalt}[2]{\ifshort#1\else#2\fi}
% \Jtr, \JtrS{\Sigma}: transverse Jacobian. J^\perp here; SipArxiv.tex redefines them as J_{\mathrm{tra}}, J_{\mathrm{tra},\Sigma}
\newcommand{\Jtr}{J^\perp}
\newcommand{\JtrS}[1]{J_{#1}^\perp}
.
\makeatletter
\AtBeginDocument{\@ifundefined{hypersetup}{}{\hypersetup{linktoc=all}}}
\makeatother
\usepackage{graphicx,float,fancybox}
\usepackage{epsfig,epstopdf,tikz}
\usepackage[sans]{dsfont}
\usetikzlibrary{positioning,arrows,arrows.meta,calc}
% Configure listings for Mathematica code
\lstdefinestyle{mathematica}{
 language=Mathematica,
 basicstyle=\small\ttfamily,
 backgroundcolor=\color{gray!10},
 frame=single,
 breaklines=true,
 commentstyle=\color{green!60!black},
 keywordstyle=\color{blue},
 stringstyle=\color{red},
 showstringspaces=false,
 tabsize=2
}
\usepackage{tcolorbox}
% Colored boxes for key results
\newtcolorbox{keyresult}[1][]{
 colback=blue!5!white,
 colframe=blue!75!black,
 title=#1,
 fonttitle=\bfseries
}
\newtcolorbox{examplebox}[1][]{
 colback=green!5!white,
 colframe=green!75!black,
 title=#1,
 fonttitle=\bfseries
}


\makeatletter
\@ifundefined{theorem}   {\newtheorem{theorem}   {Theorem}}   {}
\@ifundefined{definition}{\newtheorem{definition}{Definition}} {}
\@ifundefined{open}      {\newtheorem{open}      {Problem}}   {}
\@ifundefined{example}   {\newtheorem{example}   {Example}}   {}
\makeatother

% COUNTER DEFINITIONS
\newcounter{tabcounter}
\newcounter{defcounter}
\newcounter{opcounter}
\newcounter{thmcounter}
\newcounter{excounter}
\newcounter{ascounter}
\newcounter{cjcounter}
\newcounter{lemcounter}
\newcounter{remcounter}
\newcounter{corcounter}
\newcounter{propcounter}
\newcounter{figcounter}

% UNIFIED GLOSSARY STYLE
\newglossarystyle{unifiedstyle}{%
  \setglossarystyle{list}%
  \renewcommand*{\glossarypreamble}{\gdef\glslastsec{}}%
  \renewcommand{\glossentry}[2]{%
    \edef\glscursec{\glsentryuseriii{##1}}%
    \ifx\glscursec\glslastsec\else
      \item[]\textbf{\S\,\glscursec\enspace\glsentryuseriv{##1}}%
      \global\let\glslastsec\glscursec
    \fi
    \item[\glsentryuserii{##1} \glsentryuseri{##1}:] ``\glossentryname{##1}'' \glsentrydesc{##1}%
  }%
}

\makeatletter
% First-word extractor (used by \beT, \beD and \beXa for the default key)
% \@firstwordof hello world \relax  ->  hello
\def\@firstwordof#1 #2\relax{#1}
% First three letters, used for the heading mnemonic of a theorem given no
% explicit key.  Padding with \@empty keeps it safe for first words shorter
% than three letters: \@empty vanishes under \edef, whereas \relax would
% survive into the label and is illegal inside \csname (it broke every
% two-letter title, e.g. "ME ...", with "Missing \endcsname inserted").
\def\@takethree#1#2#3#4\@nil{#1#2#3}

% EXAMPLE COMMAND -- same signature and behaviour as \beT and \beD.
% Usage: \beXa*[long title][short title][key]
%   star  = suppresses the printed "[key]" tag after the title (as for \beT*
%   and \beD*); the label and glossary entry are unaffected.
%   key = optional label key override; the label becomes e:KEY, and without a
%   key it is e: plus the first three letters of the title's first word.  A
%   mnemonic label lets an example be cited as "Example~\ref{e:KEY}" instead
%   of by section number, which is what the citing prose needs.
\NewDocumentCommand{\beXa}{s O{} O{} o}{%
  \def\@temparg{#2}%
  \ifx\@temparg\@empty
    \begin{example}%
  \else
    \let\@savedSecTitle\@currentsectionname%
    \refstepcounter{excounter}%
    \IfValueTF{#4}{%
      \edef\@exlblkey{#4}%
    }{%
      \edef\@exfirstword{\@firstwordof#2 \relax}%
      \edef\@exlblkey{%
        \expandafter\@takethree\@exfirstword\@empty\@empty\@empty\@nil}%
    }%
    \def\@tempshort{#3}%
    \ifx\@tempshort\@empty
      \expandafter\xdef\csname @xshort@\@exlblkey\endcsname{\@exlblkey}%
    \else
      \expandafter\def\csname @xshort@\@exlblkey\endcsname{#3}%
    \fi
    \IfBooleanTF{#1}{\let\@extag\@empty}{%
      \edef\@extag{\noexpand\ \noexpand\textup{[\noexpand\texttt{e:\@exlblkey}]}}}%
    \begin{example}[#2\ifshowkeys\@extag\fi]%
    \label{e:\@exlblkey}%
    % record the long title so \nref{e:KEY} prints it, as for theorems
    \protected@write\@auxout{}{%
      \string\@thmtitle{e:\@exlblkey}{\unexpanded{#2}}}%
    \expanded{%
      \noexpand\newglossaryentry{e\theexcounter}{%
        name={\unexpanded{#2}%
          \IfBooleanF{#1}{ [\noexpand\texttt{e:\@exlblkey}]}},%
        sort={\theexample},%
        description={p.~\noexpand\pageref{e:\@exlblkey}},%
        user1={\theexample},%
        user2={Example},%
        user3={\thesection},%
        user4={\expandafter\unexpanded\expandafter{\@savedSecTitle}}%
      }%
    }%
    \protected@write\@auxout{}{\string\glsliveentry{e\theexcounter}}%
  \fi
}
\newcommand{\eeXa}{\end{example}}

% DEFINITION COMMAND -- same signature and behaviour as \beT below.
% Usage: \beD[long title][short title][key]
%   long title  = definition heading (plain text only)
%   short title = shown in \dref{word}; defaults to the key
%   key         = optional label key override; the label becomes d:KEY
% Without a key the label is d: plus the first three letters of the title's
% first word, so it is a mnemonic rather than the definition's number: a
% numeric label would silently change whenever a definition is inserted.
\NewDocumentCommand{\beD}{s O{} O{} o}{%
  \def\@temparg{#2}%
  \ifx\@temparg\@empty
    \begin{definition}%
  \else
    \let\@savedSecTitle\@currentsectionname%
    \refstepcounter{defcounter}%
    \IfValueTF{#4}{%
      \edef\@deflblkey{#4}%
    }{%
      \edef\@deffirstword{\@firstwordof#2 \relax}%
      \edef\@deflblkey{%
        \expandafter\@takethree\@deffirstword\@empty\@empty\@empty\@nil}%
    }%
    \def\@tempshort{#3}%
    \ifx\@tempshort\@empty
      \expandafter\xdef\csname @dshort@\@deflblkey\endcsname{\@deflblkey}%
    \else
      \expandafter\def\csname @dshort@\@deflblkey\endcsname{#3}%
    \fi
    \IfBooleanTF{#1}{\let\@deftag\@empty}{%
      \edef\@deftag{\noexpand\ \noexpand\textup{[\noexpand\texttt{d:\@deflblkey}]}}}%
    \begin{definition}[#2\ifshowkeys\@deftag\fi]%
    \label{d:\@deflblkey}%
    % record the long title so \nref{d:KEY} prints it, as for theorems
    \protected@write\@auxout{}{%
      \string\@thmtitle{d:\@deflblkey}{\unexpanded{#2}}}%
    \expanded{%
      \noexpand\newglossaryentry{d\thedefcounter}{%
        name={\unexpanded{#2}%
          \IfBooleanF{#1}{ [\noexpand\texttt{d:\@deflblkey}]}},%
        sort={\thedefinition},%
        description={p.~\noexpand\pageref{d:\@deflblkey}},%
        user1={\thedefinition},%
        user2={Definition},%
        user3={\thesection},%
        user4={\expandafter\unexpanded\expandafter{\@savedSecTitle}}%
      }%
    }%
    \protected@write\@auxout{}{\string\glsliveentry{d\thedefcounter}}%
  \fi
}
\newcommand{\eeD}{\end{definition}}

% \showkeysfalse in a driver hides the [t:KEY] mnemonics in result headings
% (labels and glossary entries are unchanged); used for submission builds.
\newif\ifshowkeys \showkeystrue
% \commonredtrue in a driver that inputs basicFiles/commonRed: tool.tex then
% cites its reduction results instead of restating them.
\newif\ifcommonred \commonredfalse
% TheorEM COMMAND
% Usage: \beT[long title][short title][key]
%   long title  = theorem heading (plain text only)
%   short title = shown in \tref{word}; may contain math; defaults to first word of long title
%   key         = optional label key override; label becomes t:KEY (avoids first-word collisions)
% Label is t:WORD where WORD = first word of long title (or KEY if given); use \tref{WORD} to reference.
\NewDocumentCommand{\beT}{s O{} O{} o}{%
  \def\@temparg{#2}%
  \ifx\@temparg\@empty
    \begin{theorem}%
  \else
    \let\@savedSecTitle\@currentsectionname%
    \refstepcounter{thmcounter}%
    % \@thmtag is the mnemonic shown in the heading; every theorem gets one.
    % With an explicit key the key itself is shown.  Without one the key
    % defaults to the title's first word, and showing that whole word would
    % merely repeat the title ("Dissipativity [Dissipativity]"), so only its
    % first three letters are shown: "Dissipativity [Dis]".  The label is
    % unaffected -- it still uses the full \@thmlblkey.
    \IfValueTF{#4}{%
      \edef\@thmlblkey{#4}%
    }{%
      \edef\@thmfirstword{\@firstwordof#2 \relax}%
      \edef\@thmlblkey{%
        \expandafter\@takethree\@thmfirstword\@empty\@empty\@empty\@nil}%
    }%
    \IfBooleanTF{#1}{\let\@thmtag\@empty}{%
      \edef\@thmtag{\noexpand\ \noexpand\textup{[\noexpand\texttt{t:\@thmlblkey}]}}}%
    \def\@tempshort{#3}%
    \ifx\@tempshort\@empty
      \expandafter\xdef\csname @tshort@\@thmlblkey\endcsname{\@thmlblkey}%
    \else
      \expandafter\def\csname @tshort@\@thmlblkey\endcsname{#3}%
    \fi
    \begin{theorem}[#2\ifshowkeys\@thmtag\fi]%
    \label{t:\@thmlblkey}%
    % record the long title for \nref; via the .aux so forward refs resolve
    \protected@write\@auxout{}{%
      \string\@thmtitle{t:\@thmlblkey}{\unexpanded{#2}}}%
    \expanded{%
      \noexpand\newglossaryentry{t\thethmcounter}{%
        name={\unexpanded{#2}%
          \IfBooleanF{#1}{ [\noexpand\texttt{t:\@thmlblkey}]}},%
        % sort/user1 must be the number the document actually prints, i.e. the
        % theorem environment's counter, not thmcounter.  thmcounter advances
        % only on \beT with a title, so a bare \beT (or \beT with an
        % empty title) makes the two drift and the glossary cites a wrong
        % number.  thmcounter is still fine as the entry key: it is unique.
        sort={\thetheorem},%
        description={p.~\noexpand\pageref{t:\@thmlblkey}},%
        user1={\thetheorem},%
        user2={Theorem},%
        user3={\thesection},%
        user4={\expandafter\unexpanded\expandafter{\@savedSecTitle}}%
      }%
    }%
    \protected@write\@auxout{}{\string\glsliveentry{t\thethmcounter}}%
  \fi
}
\newcommand{\eeT}{\end{theorem}}

% CROSS-PART RELATION PARAGRAPH
% Usage: \xpart{text}  -- place immediately BEFORE \beT, never inside it
%   (anything inside \beT...\eeT would be swallowed into the glossary entry).
% The paragraph is typeset only when Part I is actually present, i.e. in the
% combined book build; in the standalone builds (GranadaDSABNS/A.tex,
% AAHRah-Submission/AAHRah.tex) the label part:glv is undefined and the whole
% paragraph disappears rather than leaving ?? cross-references. Same
% graceful-degradation principle as \GASref in GASt.tex. Resolves on the
% second LaTeX pass, since r@part:glv comes from the .aux file.
% (Already inside the \makeatletter region opened above, so no \makeatletter
% pair here -- adding one would close that region early.)
\newcommand{\xpart}[1]{\@ifundefined{r@part:glv}{}{\paragraph{Relation with Parts~I--II.}#1}}

% \bookref{label}{precise text}{fallback text}: prints the precise cross-reference
% when the label resolves (combined book build), and the fallback text otherwise
% (standalone paper builds, where book-only appendix labels are absent). Same
% graceful-degradation role as \GASref in GASt.tex, but usable in any file since it
% lives in this preamble (loaded by every driver) -- unlike \GASref, which is defined in a
% content file and is therefore undefined in standalone builds that omit that file.
\newcommand{\bookref}[3]{\@ifundefined{r@#1}{#3}{#2}}

% \PVpart and \RBpart: the two parts that a partial build may omit.
% They print "Part~N" in any build that contains the part, and a naming phrase
% in a build that does not, so an excerpt shows no ?? .  Each already supplies
% the word "Part", so write "\PVpart develops ...", not "Part~\PVpart".
\newcommand{\PVpart}{\bookref{part:aahrah}%
  {Part~\ref{part:aahrah}}{the Perron--Volterra part of the full book}}
\newcommand{\RBpart}{\bookref{part:recbrowne}%
  {Part~\ref{part:recbrowne}}{the immuno-virology part of the full book}}

% OPEN PROBLEM COMMANDS
% ASSUMPTION COMMANDS (glossaried)
\newcommand{\beAs}[1][]{%
  \def\@temparg{#1}%
  \ifx\@temparg\@empty
    \begin{assumption}%
  \else
    \let\@savedSecTitle\@currentsectionname%
    \refstepcounter{ascounter}%
    \begin{assumption}[#1]%
    \label{as:\theascounter}%
    \expanded{%
      \noexpand\newglossaryentry{as\theascounter}{%
        name={\unexpanded{#1}},%
        sort={\theascounter},%
        description={p.~\noexpand\pageref{as:\theascounter}},%
        user1={\theascounter},%
        user2={Assumption},%
        user3={\thesection},%
        user4={\expandafter\unexpanded\expandafter{\@savedSecTitle}}%
      }%
    }%
    \protected@write\@auxout{}{\string\glsliveentry{as\theascounter}}%
  \fi
}
\newcommand{\eeAs}{\end{assumption}}
% CONJECTURE COMMANDS (glossaried)
\newcommand{\beCj}[1][]{%
  \def\@temparg{#1}%
  \ifx\@temparg\@empty
    \begin{conjecture}%
  \else
    \let\@savedSecTitle\@currentsectionname%
    \refstepcounter{cjcounter}%
    \begin{conjecture}[#1]%
    \label{cj:\thecjcounter}%
    \expanded{%
      \noexpand\newglossaryentry{cj\thecjcounter}{%
        name={\unexpanded{#1}},%
        sort={\thecjcounter},%
        description={p.~\noexpand\pageref{cj:\thecjcounter}},%
        user1={\thecjcounter},%
        user2={Conjecture},%
        user3={\thesection},%
        user4={\expandafter\unexpanded\expandafter{\@savedSecTitle}}%
      }%
    }%
    \protected@write\@auxout{}{\string\glsliveentry{cj\thecjcounter}}%
  \fi
}
\newcommand{\eeCj}{\end{conjecture}}
% LEMMA COMMANDS
% ---------------------------------------------------------------------
% UNIFORM RESULT WRAPPERS -- \beL, \beC, \beP, \beR share \beT's signature
%
%   \beX[long title][short title][key]  ...  \eeX
%
%   long title  = heading (math allowed)
%   short title = shown in the cross-reference; defaults to the key
%   key         = label key.  With a key the label is X:KEY and the heading
%                 carries the mnemonic [key], exactly as for \beT.  Without
%                 one the label falls back to the running counter, which is
%                 what every pre-existing call site in the book relies on.
%
% A bare \beX -- arguments commented out -- opens the plain environment and
% creates NO label and NO glossary entry.  That is the supported way to keep
% an individual result out of the catalogue.
% ---------------------------------------------------------------------
% #1 wrapper name  #2 environment  #3 label prefix  #4 counter  #5 type name
\newcommand{\@mkresult}[5]{%
  % A leading * suppresses the [key] mnemonic in both the heading and the
  % glossary name, while still making the label X:KEY and the entry:
  %     \beX*[long][short][key]
  \expandafter\NewDocumentCommand\csname be#1\endcsname{s O{} O{} o}{%
    \def\@temparg{##2}%
    \ifx\@temparg\@empty
      \begin{#2}%
    \else
      \let\@savedSecTitle\@currentsectionname%
      \refstepcounter{#4}%
      % With a key: mnemonic label X:KEY and the [key] tag in the heading.
      % Without one: fall back to the running counter, exactly as these
      % wrappers behaved before, so every legacy one-argument call site in
      % the book keeps its label.  The counter is also csname-safe, which a
      % title's first word is not -- titles beginning with math would
      % otherwise build \csname \Omega...\endcsname and fail.
      \IfValueTF{##4}{%
        \edef\@relbl{##4}%
        \IfBooleanTF{##1}{\let\@retag\@empty}{%
          \edef\@retag{\noexpand\ \noexpand\textup{[\noexpand\texttt{#3:##4}]}}}%
      }{%
        \edef\@relbl{\arabic{#4}}%
        \let\@retag\@empty
      }%
      \def\@tempshort{##3}%
      \ifx\@tempshort\@empty
        \expandafter\xdef\csname @sshort@#3\@relbl\endcsname{\@relbl}%
      \else
        \expandafter\def\csname @sshort@#3\@relbl\endcsname{##3}%
      \fi
      \begin{#2}[##2\ifshowkeys\@retag\fi]%
      \edef\@tmplbl{#3:\@relbl}\expandafter\label\expandafter{\@tmplbl}%
      \expanded{%
        \noexpand\newglossaryentry{#3\csname the#4\endcsname}{%
          name={\unexpanded{##2}%
            \IfBooleanF{##1}{\IfValueT{##4}{ [\noexpand\texttt{#3:##4}]}}},%
          sort={\csname the#2\endcsname},%
          description={p.~\noexpand\pageref{#3:\@relbl}},%
          user1={\csname the#2\endcsname},%
          user2={#5},%
          user3={\thesection},%
          user4={\expandafter\unexpanded\expandafter{\@savedSecTitle}}%
        }%
      }%
      \protected@write\@auxout{}{\string\glsliveentry{#3\csname the#4\endcsname}}%
      \protected@write\@auxout{}{\string\@thmtitle{\@tmplbl}{\unexpanded{##2}}}%
    \fi
  }%
  \expandafter\newcommand\csname ee#1\endcsname{\end{#2}}%
}
\@mkresult{L}{lemma}{l}{lemcounter}{Lemma}
% REMARK / COROLLARY / PROPOSITION -- same signature as \beL above
\@mkresult{R}{remark}{r}{remcounter}{Remark}
\@mkresult{C}{corollary}{c}{corcounter}{Corollary}
\@mkresult{P}{proposition}{p}{propcounter}{Proposition}
% OPEN-PROBLEM COMMAND -- generated with the same signature as \beT:
%   \beO[long title][short title][key]   ...   \eeO
% label o:KEY, or o:<counter> when no key is given; a bare \beO makes
% neither a label nor a glossary entry.
\@mkresult{O}{open}{o}{opcounter}{Problem}


% Section-level title tracker: updated only by \section, not \subsection.
% \@currentlabelname is overwritten by every \refstepcounter (including
% \subsection), so we save it separately here.
\newcommand{\@currentsectionname}{}
\AtBeginDocument{%
  \let\@epidcrn@orig@refstepcounter\refstepcounter
  \renewcommand{\refstepcounter}[1]{%
    \@epidcrn@orig@refstepcounter{#1}%
    \ifstrequal{#1}{section}{\global\let\@currentsectionname\@currentlabelname}{}%
  }%
}

% TABLE COMMANDS (glossaried) -- consistent naming with \beT, \beD, etc.
\newcommand{\beTa}[1][]{%
  \def\@temparg{#1}%
  \ifx\@temparg\@empty
    \begin{table}[htbp]%
  \else
    \let\@savedSecTitle\@currentsectionname%
    \refstepcounter{tabcounter}%
    \begin{table}[htbp]%
    \label{A:\thetabcounter}%
    \expanded{%
      \noexpand\newglossaryentry{A\thetabcounter}{%
        name={\unexpanded{#1}},%
        sort={\thetabcounter},%
        description={p.~\noexpand\pageref{A:\thetabcounter}},%
        user1={\thetabcounter},%
        user2={Table},%
        user3={\thesection},%
        user4={\expandafter\unexpanded\expandafter{\@savedSecTitle}}%
      }%
    }%
    \protected@write\@auxout{}{\string\glsliveentry{A\thetabcounter}}%
  \fi
}
\newcommand{\eeTa}{\end{table}}

% FIGURE COMMANDS (glossaried) -- consistent with \beT, \beR, \beTa, etc.
% Usage:
%   \beF[Glossary title]...\caption{Visible caption}\eeF  -> glossaried figure
%   \beF...\caption{Visible caption}\eeF                  -> plain [H] figure, no glossary
\newcommand{\beF}[1][]{%
  \def\@temparg{#1}%
  \ifx\@temparg\@empty
    \begin{figure}[H]%
  \else
    \begin{figure}[H]%
    \let\@savedSecTitle\@currentsectionname%
    \refstepcounter{figcounter}%
    \label{fig:\thefigcounter}%
    \expanded{%
      \noexpand\newglossaryentry{Fig\thefigcounter}{%
        name={\unexpanded{#1}},%
        sort={\thefigcounter},%
        description={p.~\noexpand\pageref{fig:\thefigcounter}},%
        user1={\thefigcounter},%
        user2={Figure},%
        user3={\thesection},%
        user4={\expandafter\unexpanded\expandafter{\@savedSecTitle}}%
      }%
    }%
    \protected@write\@auxout{}{\string\glsliveentry{Fig\thefigcounter}}%
  \fi
}
\newcommand{\eeF}{\end{figure}}
\makeatother
% LIVE-ENTRY REGISTER
% \glsaddall prints every entry the glossaries package currently knows, and
% that set includes entries restored from \jobname.glsdefs.  The package
% rewrites .glsdefs from all *currently defined* entries, so anything it
% restored is written straight back out: a deleted section kept reappearing in
% the catalogue for ever, and the .glsdefs had to be erased by hand.  The
% wrappers therefore record every entry they actually create in the .aux, and
% \printboth adds only those.  The .aux is rebuilt from scratch on each run, so
% this is self-healing: a removed result is gone after one further run.
\makeatletter
\newcommand{\@glslive}{}
\newcommand{\@glsaddlive}[1]{\ifglsentryexists{#1}{\glsadd{#1}}{}}
\newcommand{\glsliveentry}[1]{\g@addto@macro\@glslive{\@glsaddlive{#1}}}
\makeatother
% PRINT COMMAND
% No \clearpage here -- the catalogue continues on the current page.  This has
% been re-added by mistake several times; leave it out.
\makeatletter
\newcommand{\printboth}{%
  \@glslive   % add exactly the entries this document created last run
  \printglossary[style=unifiedstyle, title={0 \hspace{.22em} Definitions, Theorems, Examples and Open Problems}]%
}
\makeatother
% Named theorem reference: \tref{word} -> Thm~N(word-or-short)
% Shows stored short title if available, otherwise shows the key itself.
\newcommand{\tref}[1]{%
  Thm~\ref{t:#1}(\ifcsname @tshort@#1\endcsname\csname @tshort@#1\endcsname\else#1\fi)%
}

% Full-title theorem reference: \nref{t:KEY} prints the number and then the
% theorem's long title, so one writes "Theorem~\nref{t:blocks}" and gets
% "Theorem 3 ``Block factorization and siphons''".  Takes the complete label, like
% \ref, not the bare key like \tref.
%
% \beT writes each title to the .aux (see \@thmtitle below), so \nref resolves
% forward references from the second LaTeX pass, exactly as \ref does.  If the
% title is not yet known -- first pass, or a label that is not a \beT theorem
% -- it degrades to the number rather than producing ??.
% These two use @ in their names, and this point in the file is outside the
% \makeatletter region opened earlier (it closes at \makeatother above), so
% the region has to be reopened here.
\makeatletter
\newcommand{\@thmtitle}[2]{\global\@namedef{@tlong@#1}{#2}}
\newcommand{\nref}[1]{%
  \ref{#1}%
  \ifcsname @tlong@#1\endcsname
    \ ``\csname @tlong@#1\endcsname''%
  \fi
}
% \nm{t:KEY}: the stored long title of any \beT/\beL/\beP/\beR/\beC/\beD/\beE result,
% verbatim and without quotes, so item names follow the macro's first argument.
\newcommand{\nm}[1]{%
  \ifcsname @tlong@#1\endcsname
    \csname @tlong@#1\endcsname%
  \else
    \ref{#1}%
  \fi
}
\makeatother




%% ========================================
%% Other THEOREM ENVIRONMENTS
%% ========================================
\makeatletter
\@ifundefined{lemma}      {\newtheorem{lemma}      {Lemma}}      {}
\@ifundefined{proposition}{\newtheorem{proposition}{Proposition}} {}
\@ifundefined{corollary}  {\newtheorem{corollary}  {Corollary}}  {}
\@ifundefined{conjecture} {\newtheorem{conjecture} {Conjecture}}  {}
\@ifundefined{remark}     {\newtheorem{remark}     {Remark}}     {}
\@ifundefined{question}   {\newtheorem{question}   {Question}}   {}
\@ifundefined{assumption} {\newtheorem{assumption} {Assumption}} {}
\@ifundefined{hypothesis} {\newtheorem{hypothesis} {Hypothesis}} {}
\@ifundefined{fact}       {\newtheorem{fact}       {Fact}}       {}
\makeatother

% \citer: cite the relay-graph threshold theorem of AAHrel. If a document has
% \input the chapter/file containing t:relayRF (so the label resolves), point
% straight at it; otherwise fall back to a bare citation of AAHrel. Shared here
% (not per-paper) so it resolves consistently across every book chapter.
% \citeR: same idea, for the first theorem of AAHRah (t:Pstar).
\makeatletter
\providecommand{\citer}{%
  \@ifundefined{r@t:relayRF}%
    {\cite{AAHrel}}%
    {Theorem~\ref{t:relayRF} of \cite{AAHrel}}%
}
\providecommand{\citeR}{%
  \@ifundefined{r@t:Pstar}%
    {\cite{AAHRah}}%
    {Theorem~\ref{t:Pstar} of \cite{AAHRah}}%
}
\makeatother


% Abbreviations
\def\beQ{\begin{question}}\def\eeQ{\end{question}}
\def\beA{\begin{assumption}}\def\eeA{\end{assumption}}


%% ========================================

%% ========================================
%% COLORS & CUSTOM COMMANDS
%% ========================================
\definecolor{funccolor}{RGB}{25,25,112}
\definecolor{desccolor}{RGB}{64,64,64}
\newcommand{\func}[1]{\texttt{\textcolor{funccolor}{#1}}}
\def\eq{=}
\def\mm{-}

%% ========================================
%% MATHCAL SHORTHANDS
%% ========================================
% Project-wide convention: \mX for a capital X is always \mathcal X.
% Defined for the whole alphabet so the convention never fails to expand.
\newcommand{\mA}{\mathcal A}\newcommand{\mB}{\mathcal B}
\newcommand{\mC}{\mathcal C}\newcommand{\mD}{\mathcal D}
\newcommand{\mE}{\mathcal E}\newcommand{\mF}{\mathcal F}
\newcommand{\mG}{\mathcal G}\newcommand{\mH}{\mathcal H}
\newcommand{\mI}{\mathcal I}\newcommand{\mJ}{\mathcal J}
\newcommand{\mK}{\mathcal K}\newcommand{\mL}{\mathcal L}
\newcommand{\mM}{\mathcal M}\newcommand{\mN}{\mathcal N}
\newcommand{\mO}{\mathcal O}\newcommand{\mP}{\mathcal P}
\newcommand{\mQ}{\mathcal Q}\newcommand{\mR}{\mathcal R}
\newcommand{\mS}{\mathcal S}\newcommand{\mT}{\mathcal T}
\newcommand{\mU}{\mathcal U}\newcommand{\mV}{\mathcal V}
\newcommand{\mW}{\mathcal W}\newcommand{\mX}{\mathcal X}
\newcommand{\mY}{\mathcal Y}\newcommand{\mZ}{\mathcal Z}

%% ========================================
%% BASIC MATH ENVIRONMENTS
%% ========================================
\def\bep{\begin{pmatrix}}\def\eep{\end{pmatrix}}
\def\bev{\begin{vmatrix}}\def\eev{\end{vmatrix}}
\def\bea{\begin{eqnarray*}}\def\eea{\end{eqnarray*}}
\def\bc{\begin{cases}}\def\ec{\end{cases}}
\def\BEN{\begin{enumerate}}\def\EEN{\end{enumerate}}
\def\BI{\begin{itemize}}\def\EI{\end{itemize}}

\newcommand{\be}[1]{\begin{equation}\label{#1}}
\newcommand{\ee}{\end{equation}}
\newcommand{\beq}{\begin{eqnarray}}
\def\eeq{\end{eqnarray}}
\def\eqr{\eqref}\def\fr{\frac}\def\lbl{\label}
\def\Lra{\Longrightarrow}\def\non{\nonumber}\def\cNGM{\cite{Diek,Van,Van08}}
\def\Eq{\Leftrightarrow}\def\b{\beta}\def\uph{\red{UP TO HERE}}
%% ========================================
%% PROBABILITY & STATISTICS
%% ========================================
\renewcommand{\P}{\ensuremath{\mathbb{P}}}  % Note: overrides paragraph symbol
\newcommand{\E}{\ensuremath{\mathbb{E}}}
\newcommand{\ind}[1]{\ensuremath{\mathbbm{1}_{\left\{#1\right\}}}}
\DeclareMathOperator{\Cov}{Cov}
\DeclareMathOperator{\TV}{TV}
\newcommand{\rank}{\operatorname{rank}}

%% ========================================
%% CRN OPERATORS (organizations/ERCs/buffering/siphons framework)
%% ========================================
\def\clos{\mathrm{clos}}
\def\Supply{\mathrm{Supply}}
\def\req{\mathrm{req}}
\def\React{\mathrm{React}}
\def\Prod{\mathrm{Prod}}
\def\Gen{\mathrm{Gen}}
\def\In{\mathrm{In}}
\def\Out{\mathrm{Out}}
\def\Act{\mathrm{Act}}
\def\Min{\mathrm{Min}}

%% ========================================
%% LINE BREAKING
%% ========================================
% Last-resort stretch, used only on paragraphs TeX cannot otherwise break:
% a long \texttt{} token, a comma-separated inline math list (commas are not
% breakpoints in math), or a run of \cite/\ref boxes.  Without it those set as
% overfull \hbox and print into the margin; with it they set slightly loose,
% which is the lesser evil.  Does not affect paragraphs that already break.
\emergencystretch=3em

%% ========================================
%% NUMBER SETS
%% ========================================
\newcommand{\R}{\mathbb{R}}  % Real numbers
\newcommand{\C}{\mathbb{C}}  % Complex numbers
\newcommand{\N}{\mathbb{N}}
\newcommand{\Z}{\mathbb{Z}}
% \providecommand, not \newcommand: some papers define these in their own
% preamble after % Shared preamble for the GLVbook papers.  This is the single canonical copy;
% a driver one level below GLVbook reaches it as
%
%     \input{../basicFiles/preamble}
%
% What stays in the driver, because it differs from paper to paper:
% \documentclass, font choice, \title, \author, \date, and any macro genuinely
% local to one paper.  Everything else belongs here.  The file is
% self-contained: it does not read any other file.

\usepackage[utf8]{inputenc}
\usepackage[novbox]{pdfsync}   % required: SyncTeX alone does not work here

% Page geometry, set once here for every paper that loads this file.  A driver
% that needs different margins says \geometry{...}; a second
% \usepackage[...]{geometry} raises an option clash and must not be used.
\makeatletter
\@ifundefined{AAdefNnogeometry}{\usepackage[a4paper,margin=2cm]{geometry}}{}
\makeatother

% mathematics
\usepackage{amsmath,amssymb,amsthm,mathtools,bm,mathrsfs,bbm}
\usepackage{centernot}

% floats, tables, lists
\usepackage{graphicx,float}
\makeatletter
\@ifpackageloaded{subfig}{}{\usepackage{subcaption}}
\makeatother
\usepackage{booktabs,longtable,threeparttable,array,multirow}
\usepackage{enumitem}

% colour, code, algorithms
\usepackage{xcolor,listings,tcolorbox}
\makeatletter
\@ifundefined{c@algorithm}{\usepackage{algorithm}}{}
\makeatother
\usepackage{algorithmicx,algpseudocode}

% diagrams
\usepackage{tikz,tikz-cd,pgfplots}
\usetikzlibrary{graphs,decorations.pathmorphing,decorations.markings,
arrows.meta,automata,positioning}
\tikzcdset{every label/.append style = {font = \small}}
\pgfplotsset{compat=1.18}

% chemistry
\usepackage[version=4]{mhchem}

\usepackage{microtype}
\usepackage{lineno}
\usepackage{nomencl}
\makenomenclature

% hyperref first, then bookmark, then glossaries: glossaries must follow
% hyperref or the entry links are built against the wrong anchors.
\makeatletter
\@ifundefined{hyperrefpackage}{%
  \@ifundefined{AAdefNnohyperref}{\usepackage{hyperref}}{}%
}{}
\@ifundefined{hyperrefpackage}{\usepackage{bookmark}}{}
\makeatother
\usepackage[toc,sort=use]{glossaries}
\makeglossaries

% Book/standalone switch.  Material that the book already covers in its
% prelude, but that a standalone paper still has to state, is wrapped as
%     \ifbookbuild\else ... \fi
% Drivers default to standalone; bookM.tex says \bookbuildtrue.
\newif\ifbookbuild
\bookbuildfalse

% Part switches for the gradostat paper.  Both default to true; a driver that
% wants the short submission says \madecfalse and/or \slowfastfalse, and every
% cross-reference into the dropped part disappears with it.
%   \ifmadec     Part III, the Lagasquie--Madec model and its break-even geometry
%   \ifslowfast  Part IV, everything carrying the small parameter \eps
\newif\ifmadec     \madectrue
\newif\ifslowfast  \slowfasttrue

% Part switch for the Barrat/simplicial-contagion project.  Defaults to
% true (the hypernetwork-synchrony part is included); the journal-paper
% driver says \hnetpartfalse to drop it and every cross-reference into it.
\newif\ifhnetpart  \hnetparttrue

%% ========================================
%% PACKAGES
%% ========================================
\RequirePackage{etoolbox}
\usepackage{makecell}
\usepackage[novbox]{pdfsync}% obsolete; SyncTeX is built into modern engines
\usepackage{amssymb, amsfonts,mathtools,bbm,eurosym,graphicx,latexsym, amsmath,amsthm,times,nccmath,url}
\usepackage{listings,verbatim,cases,tabularx}
% hyperref is deliberately NOT in the line above.  Journal classes often
% load it themselves, with a long option list, from a begindocument hook
% (mdpi.cls does: \DeclareHookRule{begindocument}{hyperref}{before}{mdpi}).
% Loading it here first, without options, makes that later load raise
% "Option clash for package hyperref".  So load it only if nobody else
% has, or will:
%   - \@ifpackageloaded  catches a class that loaded it already;
%   - \hyperrefpackage    is mdpi.cls's own deferred loader;
%   - \AAdefNnohyperref   lets any other driver opt out by defining it
%                         before \input{preamble}.
\makeatletter
\@ifpackageloaded{hyperref}{}{%
  \@ifundefined{hyperrefpackage}{%
    \@ifundefined{AAdefNnohyperref}{\usepackage{hyperref}}{}%
  }{}%
}
\makeatother
% linktoc=all makes a table-of-contents page number a link as well as the
% entry text; hyperref's default, linktoc=section, links only the text.
% Deferred to begin-document and guarded: by the block above hyperref may
% be loaded later (mdpi's own hook) or not at all (\AAdefNnohyperref), and
% \hypersetup must exist before it is called.  A driver that wants the
% default back can say \hypersetup{linktoc=section} after \input{preamble}.
\makeatletter
\AtBeginDocument{\@ifundefined{hypersetup}{}{\hypersetup{linktoc=all}}}
\makeatother
\usepackage{graphicx,float,fancybox}
\usepackage{epsfig,epstopdf,tikz}
\usepackage[sans]{dsfont}
\usetikzlibrary{positioning,arrows,arrows.meta,calc}
% Configure listings for Mathematica code
\lstdefinestyle{mathematica}{
 language=Mathematica,
 basicstyle=\small\ttfamily,
 backgroundcolor=\color{gray!10},
 frame=single,
 breaklines=true,
 commentstyle=\color{green!60!black},
 keywordstyle=\color{blue},
 stringstyle=\color{red},
 showstringspaces=false,
 tabsize=2
}
\usepackage{tcolorbox}
% Colored boxes for key results
\newtcolorbox{keyresult}[1][]{
 colback=blue!5!white,
 colframe=blue!75!black,
 title=#1,
 fonttitle=\bfseries
}
\newtcolorbox{examplebox}[1][]{
 colback=green!5!white,
 colframe=green!75!black,
 title=#1,
 fonttitle=\bfseries
}


\makeatletter
\@ifundefined{theorem}   {\newtheorem{theorem}   {Theorem}}   {}
\@ifundefined{definition}{\newtheorem{definition}{Definition}} {}
\@ifundefined{open}      {\newtheorem{open}      {Problem}}   {}
\@ifundefined{example}   {\newtheorem{example}   {Example}}   {}
\makeatother

% COUNTER DEFINITIONS
\newcounter{tabcounter}
\newcounter{defcounter}
\newcounter{opcounter}
\newcounter{thmcounter}
\newcounter{excounter}
\newcounter{ascounter}
\newcounter{cjcounter}
\newcounter{lemcounter}
\newcounter{remcounter}
\newcounter{corcounter}
\newcounter{propcounter}
\newcounter{figcounter}

% UNIFIED GLOSSARY STYLE
\newglossarystyle{unifiedstyle}{%
  \setglossarystyle{list}%
  \renewcommand*{\glossarypreamble}{\gdef\glslastsec{}}%
  \renewcommand{\glossentry}[2]{%
    \edef\glscursec{\glsentryuseriii{##1}}%
    \ifx\glscursec\glslastsec\else
      \item[]\textbf{\S\,\glscursec\enspace\glsentryuseriv{##1}}%
      \global\let\glslastsec\glscursec
    \fi
    \item[\glsentryuserii{##1} \glsentryuseri{##1}:] ``\glossentryname{##1}'' \glsentrydesc{##1}%
  }%
}

\makeatletter
% First-word extractor (used by \beT, \beD and \beXa for the default key)
% \@firstwordof hello world \relax  ->  hello
\def\@firstwordof#1 #2\relax{#1}
% First three letters, used for the heading mnemonic of a theorem given no
% explicit key.  Padding with \@empty keeps it safe for first words shorter
% than three letters: \@empty vanishes under \edef, whereas \relax would
% survive into the label and is illegal inside \csname (it broke every
% two-letter title, e.g. "ME ...", with "Missing \endcsname inserted").
\def\@takethree#1#2#3#4\@nil{#1#2#3}

% EXAMPLE COMMAND -- same signature and behaviour as \beT and \beD.
% Usage: \beXa*[long title][short title][key]
%   star  = suppresses the printed "[key]" tag after the title (as for \beT*
%   and \beD*); the label and glossary entry are unaffected.
%   key = optional label key override; the label becomes e:KEY, and without a
%   key it is e: plus the first three letters of the title's first word.  A
%   mnemonic label lets an example be cited as "Example~\ref{e:KEY}" instead
%   of by section number, which is what the citing prose needs.
\NewDocumentCommand{\beXa}{s O{} O{} o}{%
  \def\@temparg{#2}%
  \ifx\@temparg\@empty
    \begin{example}%
  \else
    \let\@savedSecTitle\@currentsectionname%
    \refstepcounter{excounter}%
    \IfValueTF{#4}{%
      \edef\@exlblkey{#4}%
    }{%
      \edef\@exfirstword{\@firstwordof#2 \relax}%
      \edef\@exlblkey{%
        \expandafter\@takethree\@exfirstword\@empty\@empty\@empty\@nil}%
    }%
    \def\@tempshort{#3}%
    \ifx\@tempshort\@empty
      \expandafter\xdef\csname @xshort@\@exlblkey\endcsname{\@exlblkey}%
    \else
      \expandafter\def\csname @xshort@\@exlblkey\endcsname{#3}%
    \fi
    \IfBooleanTF{#1}{\let\@extag\@empty}{%
      \edef\@extag{\noexpand\ \noexpand\textup{[\noexpand\texttt{e:\@exlblkey}]}}}%
    \begin{example}[#2\ifshowkeys\@extag\fi]%
    \label{e:\@exlblkey}%
    % record the long title so \nref{e:KEY} prints it, as for theorems
    \protected@write\@auxout{}{%
      \string\@thmtitle{e:\@exlblkey}{\unexpanded{#2}}}%
    \expanded{%
      \noexpand\newglossaryentry{e\theexcounter}{%
        name={\unexpanded{#2}%
          \IfBooleanF{#1}{ [\noexpand\texttt{e:\@exlblkey}]}},%
        sort={\theexample},%
        description={p.~\noexpand\pageref{e:\@exlblkey}},%
        user1={\theexample},%
        user2={Example},%
        user3={\thesection},%
        user4={\expandafter\unexpanded\expandafter{\@savedSecTitle}}%
      }%
    }%
    \protected@write\@auxout{}{\string\glsliveentry{e\theexcounter}}%
  \fi
}
\newcommand{\eeXa}{\end{example}}

% DEFINITION COMMAND -- same signature and behaviour as \beT below.
% Usage: \beD[long title][short title][key]
%   long title  = definition heading (plain text only)
%   short title = shown in \dref{word}; defaults to the key
%   key         = optional label key override; the label becomes d:KEY
% Without a key the label is d: plus the first three letters of the title's
% first word, so it is a mnemonic rather than the definition's number: a
% numeric label would silently change whenever a definition is inserted.
\NewDocumentCommand{\beD}{s O{} O{} o}{%
  \def\@temparg{#2}%
  \ifx\@temparg\@empty
    \begin{definition}%
  \else
    \let\@savedSecTitle\@currentsectionname%
    \refstepcounter{defcounter}%
    \IfValueTF{#4}{%
      \edef\@deflblkey{#4}%
    }{%
      \edef\@deffirstword{\@firstwordof#2 \relax}%
      \edef\@deflblkey{%
        \expandafter\@takethree\@deffirstword\@empty\@empty\@empty\@nil}%
    }%
    \def\@tempshort{#3}%
    \ifx\@tempshort\@empty
      \expandafter\xdef\csname @dshort@\@deflblkey\endcsname{\@deflblkey}%
    \else
      \expandafter\def\csname @dshort@\@deflblkey\endcsname{#3}%
    \fi
    \IfBooleanTF{#1}{\let\@deftag\@empty}{%
      \edef\@deftag{\noexpand\ \noexpand\textup{[\noexpand\texttt{d:\@deflblkey}]}}}%
    \begin{definition}[#2\ifshowkeys\@deftag\fi]%
    \label{d:\@deflblkey}%
    % record the long title so \nref{d:KEY} prints it, as for theorems
    \protected@write\@auxout{}{%
      \string\@thmtitle{d:\@deflblkey}{\unexpanded{#2}}}%
    \expanded{%
      \noexpand\newglossaryentry{d\thedefcounter}{%
        name={\unexpanded{#2}%
          \IfBooleanF{#1}{ [\noexpand\texttt{d:\@deflblkey}]}},%
        sort={\thedefinition},%
        description={p.~\noexpand\pageref{d:\@deflblkey}},%
        user1={\thedefinition},%
        user2={Definition},%
        user3={\thesection},%
        user4={\expandafter\unexpanded\expandafter{\@savedSecTitle}}%
      }%
    }%
    \protected@write\@auxout{}{\string\glsliveentry{d\thedefcounter}}%
  \fi
}
\newcommand{\eeD}{\end{definition}}

% \showkeysfalse in a driver hides the [t:KEY] mnemonics in result headings
% (labels and glossary entries are unchanged); used for submission builds.
\newif\ifshowkeys \showkeystrue
% \commonredtrue in a driver that inputs basicFiles/commonRed: tool.tex then
% cites its reduction results instead of restating them.
\newif\ifcommonred \commonredfalse
% TheorEM COMMAND
% Usage: \beT[long title][short title][key]
%   long title  = theorem heading (plain text only)
%   short title = shown in \tref{word}; may contain math; defaults to first word of long title
%   key         = optional label key override; label becomes t:KEY (avoids first-word collisions)
% Label is t:WORD where WORD = first word of long title (or KEY if given); use \tref{WORD} to reference.
\NewDocumentCommand{\beT}{s O{} O{} o}{%
  \def\@temparg{#2}%
  \ifx\@temparg\@empty
    \begin{theorem}%
  \else
    \let\@savedSecTitle\@currentsectionname%
    \refstepcounter{thmcounter}%
    % \@thmtag is the mnemonic shown in the heading; every theorem gets one.
    % With an explicit key the key itself is shown.  Without one the key
    % defaults to the title's first word, and showing that whole word would
    % merely repeat the title ("Dissipativity [Dissipativity]"), so only its
    % first three letters are shown: "Dissipativity [Dis]".  The label is
    % unaffected -- it still uses the full \@thmlblkey.
    \IfValueTF{#4}{%
      \edef\@thmlblkey{#4}%
    }{%
      \edef\@thmfirstword{\@firstwordof#2 \relax}%
      \edef\@thmlblkey{%
        \expandafter\@takethree\@thmfirstword\@empty\@empty\@empty\@nil}%
    }%
    \IfBooleanTF{#1}{\let\@thmtag\@empty}{%
      \edef\@thmtag{\noexpand\ \noexpand\textup{[\noexpand\texttt{t:\@thmlblkey}]}}}%
    \def\@tempshort{#3}%
    \ifx\@tempshort\@empty
      \expandafter\xdef\csname @tshort@\@thmlblkey\endcsname{\@thmlblkey}%
    \else
      \expandafter\def\csname @tshort@\@thmlblkey\endcsname{#3}%
    \fi
    \begin{theorem}[#2\ifshowkeys\@thmtag\fi]%
    \label{t:\@thmlblkey}%
    % record the long title for \nref; via the .aux so forward refs resolve
    \protected@write\@auxout{}{%
      \string\@thmtitle{t:\@thmlblkey}{\unexpanded{#2}}}%
    \expanded{%
      \noexpand\newglossaryentry{t\thethmcounter}{%
        name={\unexpanded{#2}%
          \IfBooleanF{#1}{ [\noexpand\texttt{t:\@thmlblkey}]}},%
        % sort/user1 must be the number the document actually prints, i.e. the
        % theorem environment's counter, not thmcounter.  thmcounter advances
        % only on \beT with a title, so a bare \beT (or \beT with an
        % empty title) makes the two drift and the glossary cites a wrong
        % number.  thmcounter is still fine as the entry key: it is unique.
        sort={\thetheorem},%
        description={p.~\noexpand\pageref{t:\@thmlblkey}},%
        user1={\thetheorem},%
        user2={Theorem},%
        user3={\thesection},%
        user4={\expandafter\unexpanded\expandafter{\@savedSecTitle}}%
      }%
    }%
    \protected@write\@auxout{}{\string\glsliveentry{t\thethmcounter}}%
  \fi
}
\newcommand{\eeT}{\end{theorem}}

% CROSS-PART RELATION PARAGRAPH
% Usage: \xpart{text}  -- place immediately BEFORE \beT, never inside it
%   (anything inside \beT...\eeT would be swallowed into the glossary entry).
% The paragraph is typeset only when Part I is actually present, i.e. in the
% combined book build; in the standalone builds (GranadaDSABNS/A.tex,
% AAHRah-Submission/AAHRah.tex) the label part:glv is undefined and the whole
% paragraph disappears rather than leaving ?? cross-references. Same
% graceful-degradation principle as \GASref in GASt.tex. Resolves on the
% second LaTeX pass, since r@part:glv comes from the .aux file.
% (Already inside the \makeatletter region opened above, so no \makeatletter
% pair here -- adding one would close that region early.)
\newcommand{\xpart}[1]{\@ifundefined{r@part:glv}{}{\paragraph{Relation with Parts~I--II.}#1}}

% \bookref{label}{precise text}{fallback text}: prints the precise cross-reference
% when the label resolves (combined book build), and the fallback text otherwise
% (standalone paper builds, where book-only appendix labels are absent). Same
% graceful-degradation role as \GASref in GASt.tex, but usable in any file since it
% lives in this preamble (loaded by every driver) -- unlike \GASref, which is defined in a
% content file and is therefore undefined in standalone builds that omit that file.
\newcommand{\bookref}[3]{\@ifundefined{r@#1}{#3}{#2}}

% \PVpart and \RBpart: the two parts that a partial build may omit.
% They print "Part~N" in any build that contains the part, and a naming phrase
% in a build that does not, so an excerpt shows no ?? .  Each already supplies
% the word "Part", so write "\PVpart develops ...", not "Part~\PVpart".
\newcommand{\PVpart}{\bookref{part:aahrah}%
  {Part~\ref{part:aahrah}}{the Perron--Volterra part of the full book}}
\newcommand{\RBpart}{\bookref{part:recbrowne}%
  {Part~\ref{part:recbrowne}}{the immuno-virology part of the full book}}

% OPEN PROBLEM COMMANDS
% ASSUMPTION COMMANDS (glossaried)
\newcommand{\beAs}[1][]{%
  \def\@temparg{#1}%
  \ifx\@temparg\@empty
    \begin{assumption}%
  \else
    \let\@savedSecTitle\@currentsectionname%
    \refstepcounter{ascounter}%
    \begin{assumption}[#1]%
    \label{as:\theascounter}%
    \expanded{%
      \noexpand\newglossaryentry{as\theascounter}{%
        name={\unexpanded{#1}},%
        sort={\theascounter},%
        description={p.~\noexpand\pageref{as:\theascounter}},%
        user1={\theascounter},%
        user2={Assumption},%
        user3={\thesection},%
        user4={\expandafter\unexpanded\expandafter{\@savedSecTitle}}%
      }%
    }%
    \protected@write\@auxout{}{\string\glsliveentry{as\theascounter}}%
  \fi
}
\newcommand{\eeAs}{\end{assumption}}
% CONJECTURE COMMANDS (glossaried)
\newcommand{\beCj}[1][]{%
  \def\@temparg{#1}%
  \ifx\@temparg\@empty
    \begin{conjecture}%
  \else
    \let\@savedSecTitle\@currentsectionname%
    \refstepcounter{cjcounter}%
    \begin{conjecture}[#1]%
    \label{cj:\thecjcounter}%
    \expanded{%
      \noexpand\newglossaryentry{cj\thecjcounter}{%
        name={\unexpanded{#1}},%
        sort={\thecjcounter},%
        description={p.~\noexpand\pageref{cj:\thecjcounter}},%
        user1={\thecjcounter},%
        user2={Conjecture},%
        user3={\thesection},%
        user4={\expandafter\unexpanded\expandafter{\@savedSecTitle}}%
      }%
    }%
    \protected@write\@auxout{}{\string\glsliveentry{cj\thecjcounter}}%
  \fi
}
\newcommand{\eeCj}{\end{conjecture}}
% LEMMA COMMANDS
% ---------------------------------------------------------------------
% UNIFORM RESULT WRAPPERS -- \beL, \beC, \beP, \beR share \beT's signature
%
%   \beX[long title][short title][key]  ...  \eeX
%
%   long title  = heading (math allowed)
%   short title = shown in the cross-reference; defaults to the key
%   key         = label key.  With a key the label is X:KEY and the heading
%                 carries the mnemonic [key], exactly as for \beT.  Without
%                 one the label falls back to the running counter, which is
%                 what every pre-existing call site in the book relies on.
%
% A bare \beX -- arguments commented out -- opens the plain environment and
% creates NO label and NO glossary entry.  That is the supported way to keep
% an individual result out of the catalogue.
% ---------------------------------------------------------------------
% #1 wrapper name  #2 environment  #3 label prefix  #4 counter  #5 type name
\newcommand{\@mkresult}[5]{%
  % A leading * suppresses the [key] mnemonic in both the heading and the
  % glossary name, while still making the label X:KEY and the entry:
  %     \beX*[long][short][key]
  \expandafter\NewDocumentCommand\csname be#1\endcsname{s O{} O{} o}{%
    \def\@temparg{##2}%
    \ifx\@temparg\@empty
      \begin{#2}%
    \else
      \let\@savedSecTitle\@currentsectionname%
      \refstepcounter{#4}%
      % With a key: mnemonic label X:KEY and the [key] tag in the heading.
      % Without one: fall back to the running counter, exactly as these
      % wrappers behaved before, so every legacy one-argument call site in
      % the book keeps its label.  The counter is also csname-safe, which a
      % title's first word is not -- titles beginning with math would
      % otherwise build \csname \Omega...\endcsname and fail.
      \IfValueTF{##4}{%
        \edef\@relbl{##4}%
        \IfBooleanTF{##1}{\let\@retag\@empty}{%
          \edef\@retag{\noexpand\ \noexpand\textup{[\noexpand\texttt{#3:##4}]}}}%
      }{%
        \edef\@relbl{\arabic{#4}}%
        \let\@retag\@empty
      }%
      \def\@tempshort{##3}%
      \ifx\@tempshort\@empty
        \expandafter\xdef\csname @sshort@#3\@relbl\endcsname{\@relbl}%
      \else
        \expandafter\def\csname @sshort@#3\@relbl\endcsname{##3}%
      \fi
      \begin{#2}[##2\ifshowkeys\@retag\fi]%
      \edef\@tmplbl{#3:\@relbl}\expandafter\label\expandafter{\@tmplbl}%
      \expanded{%
        \noexpand\newglossaryentry{#3\csname the#4\endcsname}{%
          name={\unexpanded{##2}%
            \IfBooleanF{##1}{\IfValueT{##4}{ [\noexpand\texttt{#3:##4}]}}},%
          sort={\csname the#2\endcsname},%
          description={p.~\noexpand\pageref{#3:\@relbl}},%
          user1={\csname the#2\endcsname},%
          user2={#5},%
          user3={\thesection},%
          user4={\expandafter\unexpanded\expandafter{\@savedSecTitle}}%
        }%
      }%
      \protected@write\@auxout{}{\string\glsliveentry{#3\csname the#4\endcsname}}%
      \protected@write\@auxout{}{\string\@thmtitle{\@tmplbl}{\unexpanded{##2}}}%
    \fi
  }%
  \expandafter\newcommand\csname ee#1\endcsname{\end{#2}}%
}
\@mkresult{L}{lemma}{l}{lemcounter}{Lemma}
% REMARK / COROLLARY / PROPOSITION -- same signature as \beL above
\@mkresult{R}{remark}{r}{remcounter}{Remark}
\@mkresult{C}{corollary}{c}{corcounter}{Corollary}
\@mkresult{P}{proposition}{p}{propcounter}{Proposition}
% OPEN-PROBLEM COMMAND -- generated with the same signature as \beT:
%   \beO[long title][short title][key]   ...   \eeO
% label o:KEY, or o:<counter> when no key is given; a bare \beO makes
% neither a label nor a glossary entry.
\@mkresult{O}{open}{o}{opcounter}{Problem}


% Section-level title tracker: updated only by \section, not \subsection.
% \@currentlabelname is overwritten by every \refstepcounter (including
% \subsection), so we save it separately here.
\newcommand{\@currentsectionname}{}
\AtBeginDocument{%
  \let\@epidcrn@orig@refstepcounter\refstepcounter
  \renewcommand{\refstepcounter}[1]{%
    \@epidcrn@orig@refstepcounter{#1}%
    \ifstrequal{#1}{section}{\global\let\@currentsectionname\@currentlabelname}{}%
  }%
}

% TABLE COMMANDS (glossaried) -- consistent naming with \beT, \beD, etc.
\newcommand{\beTa}[1][]{%
  \def\@temparg{#1}%
  \ifx\@temparg\@empty
    \begin{table}[htbp]%
  \else
    \let\@savedSecTitle\@currentsectionname%
    \refstepcounter{tabcounter}%
    \begin{table}[htbp]%
    \label{A:\thetabcounter}%
    \expanded{%
      \noexpand\newglossaryentry{A\thetabcounter}{%
        name={\unexpanded{#1}},%
        sort={\thetabcounter},%
        description={p.~\noexpand\pageref{A:\thetabcounter}},%
        user1={\thetabcounter},%
        user2={Table},%
        user3={\thesection},%
        user4={\expandafter\unexpanded\expandafter{\@savedSecTitle}}%
      }%
    }%
    \protected@write\@auxout{}{\string\glsliveentry{A\thetabcounter}}%
  \fi
}
\newcommand{\eeTa}{\end{table}}

% FIGURE COMMANDS (glossaried) -- consistent with \beT, \beR, \beTa, etc.
% Usage:
%   \beF[Glossary title]...\caption{Visible caption}\eeF  -> glossaried figure
%   \beF...\caption{Visible caption}\eeF                  -> plain [H] figure, no glossary
\newcommand{\beF}[1][]{%
  \def\@temparg{#1}%
  \ifx\@temparg\@empty
    \begin{figure}[H]%
  \else
    \begin{figure}[H]%
    \let\@savedSecTitle\@currentsectionname%
    \refstepcounter{figcounter}%
    \label{fig:\thefigcounter}%
    \expanded{%
      \noexpand\newglossaryentry{Fig\thefigcounter}{%
        name={\unexpanded{#1}},%
        sort={\thefigcounter},%
        description={p.~\noexpand\pageref{fig:\thefigcounter}},%
        user1={\thefigcounter},%
        user2={Figure},%
        user3={\thesection},%
        user4={\expandafter\unexpanded\expandafter{\@savedSecTitle}}%
      }%
    }%
    \protected@write\@auxout{}{\string\glsliveentry{Fig\thefigcounter}}%
  \fi
}
\newcommand{\eeF}{\end{figure}}
\makeatother
% LIVE-ENTRY REGISTER
% \glsaddall prints every entry the glossaries package currently knows, and
% that set includes entries restored from \jobname.glsdefs.  The package
% rewrites .glsdefs from all *currently defined* entries, so anything it
% restored is written straight back out: a deleted section kept reappearing in
% the catalogue for ever, and the .glsdefs had to be erased by hand.  The
% wrappers therefore record every entry they actually create in the .aux, and
% \printboth adds only those.  The .aux is rebuilt from scratch on each run, so
% this is self-healing: a removed result is gone after one further run.
\makeatletter
\newcommand{\@glslive}{}
\newcommand{\@glsaddlive}[1]{\ifglsentryexists{#1}{\glsadd{#1}}{}}
\newcommand{\glsliveentry}[1]{\g@addto@macro\@glslive{\@glsaddlive{#1}}}
\makeatother
% PRINT COMMAND
% No \clearpage here -- the catalogue continues on the current page.  This has
% been re-added by mistake several times; leave it out.
\makeatletter
\newcommand{\printboth}{%
  \@glslive   % add exactly the entries this document created last run
  \printglossary[style=unifiedstyle, title={0 \hspace{.22em} Definitions, Theorems, Examples and Open Problems}]%
}
\makeatother
% Named theorem reference: \tref{word} -> Thm~N(word-or-short)
% Shows stored short title if available, otherwise shows the key itself.
\newcommand{\tref}[1]{%
  Thm~\ref{t:#1}(\ifcsname @tshort@#1\endcsname\csname @tshort@#1\endcsname\else#1\fi)%
}

% Full-title theorem reference: \nref{t:KEY} prints the number and then the
% theorem's long title, so one writes "Theorem~\nref{t:blocks}" and gets
% "Theorem 3 ``Block factorization and siphons''".  Takes the complete label, like
% \ref, not the bare key like \tref.
%
% \beT writes each title to the .aux (see \@thmtitle below), so \nref resolves
% forward references from the second LaTeX pass, exactly as \ref does.  If the
% title is not yet known -- first pass, or a label that is not a \beT theorem
% -- it degrades to the number rather than producing ??.
% These two use @ in their names, and this point in the file is outside the
% \makeatletter region opened earlier (it closes at \makeatother above), so
% the region has to be reopened here.
\makeatletter
\newcommand{\@thmtitle}[2]{\global\@namedef{@tlong@#1}{#2}}
\newcommand{\nref}[1]{%
  \ref{#1}%
  \ifcsname @tlong@#1\endcsname
    \ ``\csname @tlong@#1\endcsname''%
  \fi
}
% \nm{t:KEY}: the stored long title of any \beT/\beL/\beP/\beR/\beC/\beD/\beE result,
% verbatim and without quotes, so item names follow the macro's first argument.
\newcommand{\nm}[1]{%
  \ifcsname @tlong@#1\endcsname
    \csname @tlong@#1\endcsname%
  \else
    \ref{#1}%
  \fi
}
\makeatother




%% ========================================
%% Other THEOREM ENVIRONMENTS
%% ========================================
\makeatletter
\@ifundefined{lemma}      {\newtheorem{lemma}      {Lemma}}      {}
\@ifundefined{proposition}{\newtheorem{proposition}{Proposition}} {}
\@ifundefined{corollary}  {\newtheorem{corollary}  {Corollary}}  {}
\@ifundefined{conjecture} {\newtheorem{conjecture} {Conjecture}}  {}
\@ifundefined{remark}     {\newtheorem{remark}     {Remark}}     {}
\@ifundefined{question}   {\newtheorem{question}   {Question}}   {}
\@ifundefined{assumption} {\newtheorem{assumption} {Assumption}} {}
\@ifundefined{hypothesis} {\newtheorem{hypothesis} {Hypothesis}} {}
\@ifundefined{fact}       {\newtheorem{fact}       {Fact}}       {}
\makeatother

% \citer: cite the relay-graph threshold theorem of AAHrel. If a document has
% \input the chapter/file containing t:relayRF (so the label resolves), point
% straight at it; otherwise fall back to a bare citation of AAHrel. Shared here
% (not per-paper) so it resolves consistently across every book chapter.
% \citeR: same idea, for the first theorem of AAHRah (t:Pstar).
\makeatletter
\providecommand{\citer}{%
  \@ifundefined{r@t:relayRF}%
    {\cite{AAHrel}}%
    {Theorem~\ref{t:relayRF} of \cite{AAHrel}}%
}
\providecommand{\citeR}{%
  \@ifundefined{r@t:Pstar}%
    {\cite{AAHRah}}%
    {Theorem~\ref{t:Pstar} of \cite{AAHRah}}%
}
\makeatother


% Abbreviations
\def\beQ{\begin{question}}\def\eeQ{\end{question}}
\def\beA{\begin{assumption}}\def\eeA{\end{assumption}}


%% ========================================

%% ========================================
%% COLORS & CUSTOM COMMANDS
%% ========================================
\definecolor{funccolor}{RGB}{25,25,112}
\definecolor{desccolor}{RGB}{64,64,64}
\newcommand{\func}[1]{\texttt{\textcolor{funccolor}{#1}}}
\def\eq{=}
\def\mm{-}

%% ========================================
%% MATHCAL SHORTHANDS
%% ========================================
% Project-wide convention: \mX for a capital X is always \mathcal X.
% Defined for the whole alphabet so the convention never fails to expand.
\newcommand{\mA}{\mathcal A}\newcommand{\mB}{\mathcal B}
\newcommand{\mC}{\mathcal C}\newcommand{\mD}{\mathcal D}
\newcommand{\mE}{\mathcal E}\newcommand{\mF}{\mathcal F}
\newcommand{\mG}{\mathcal G}\newcommand{\mH}{\mathcal H}
\newcommand{\mI}{\mathcal I}\newcommand{\mJ}{\mathcal J}
\newcommand{\mK}{\mathcal K}\newcommand{\mL}{\mathcal L}
\newcommand{\mM}{\mathcal M}\newcommand{\mN}{\mathcal N}
\newcommand{\mO}{\mathcal O}\newcommand{\mP}{\mathcal P}
\newcommand{\mQ}{\mathcal Q}\newcommand{\mR}{\mathcal R}
\newcommand{\mS}{\mathcal S}\newcommand{\mT}{\mathcal T}
\newcommand{\mU}{\mathcal U}\newcommand{\mV}{\mathcal V}
\newcommand{\mW}{\mathcal W}\newcommand{\mX}{\mathcal X}
\newcommand{\mY}{\mathcal Y}\newcommand{\mZ}{\mathcal Z}

%% ========================================
%% BASIC MATH ENVIRONMENTS
%% ========================================
\def\bep{\begin{pmatrix}}\def\eep{\end{pmatrix}}
\def\bev{\begin{vmatrix}}\def\eev{\end{vmatrix}}
\def\bea{\begin{eqnarray*}}\def\eea{\end{eqnarray*}}
\def\bc{\begin{cases}}\def\ec{\end{cases}}
\def\BEN{\begin{enumerate}}\def\EEN{\end{enumerate}}
\def\BI{\begin{itemize}}\def\EI{\end{itemize}}

\newcommand{\be}[1]{\begin{equation}\label{#1}}
\newcommand{\ee}{\end{equation}}
\newcommand{\beq}{\begin{eqnarray}}
\def\eeq{\end{eqnarray}}
\def\eqr{\eqref}\def\fr{\frac}\def\lbl{\label}
\def\Lra{\Longrightarrow}\def\non{\nonumber}\def\cNGM{\cite{Diek,Van,Van08}}
\def\Eq{\Leftrightarrow}\def\b{\beta}\def\uph{\red{UP TO HERE}}
%% ========================================
%% PROBABILITY & STATISTICS
%% ========================================
\renewcommand{\P}{\ensuremath{\mathbb{P}}}  % Note: overrides paragraph symbol
\newcommand{\E}{\ensuremath{\mathbb{E}}}
\newcommand{\ind}[1]{\ensuremath{\mathbbm{1}_{\left\{#1\right\}}}}
\DeclareMathOperator{\Cov}{Cov}
\DeclareMathOperator{\TV}{TV}
\newcommand{\rank}{\operatorname{rank}}

%% ========================================
%% CRN OPERATORS (organizations/ERCs/buffering/siphons framework)
%% ========================================
\def\clos{\mathrm{clos}}
\def\Supply{\mathrm{Supply}}
\def\req{\mathrm{req}}
\def\React{\mathrm{React}}
\def\Prod{\mathrm{Prod}}
\def\Gen{\mathrm{Gen}}
\def\In{\mathrm{In}}
\def\Out{\mathrm{Out}}
\def\Act{\mathrm{Act}}
\def\Min{\mathrm{Min}}

%% ========================================
%% LINE BREAKING
%% ========================================
% Last-resort stretch, used only on paragraphs TeX cannot otherwise break:
% a long \texttt{} token, a comma-separated inline math list (commas are not
% breakpoints in math), or a run of \cite/\ref boxes.  Without it those set as
% overfull \hbox and print into the margin; with it they set slightly loose,
% which is the lesser evil.  Does not affect paragraphs that already break.
\emergencystretch=3em

%% ========================================
%% NUMBER SETS
%% ========================================
\newcommand{\R}{\mathbb{R}}  % Real numbers
\newcommand{\C}{\mathbb{C}}  % Complex numbers
\newcommand{\N}{\mathbb{N}}
\newcommand{\Z}{\mathbb{Z}}
% \providecommand, not \newcommand: some papers define these in their own
% preamble after \input{preamble}, and must not clash.
\providecommand{\Q}{\mathbb{Q}}
\providecommand{\res}{\operatorname{res}}   % resultant; used by epiRes.tex
\newcommand{\MN}{\mathbb{N}}
\newcommand{\MZ}{\mathbb{Z}}
\newcommand{\realpos}{\mathbb{R}_{> 0}}
\newcommand{\realnneg}{\mathbb{R}_{\geq 0}}
\newcommand{\naturalpos}{\mathbb{N}_{>0}}
\newcommand{\Rplus}{\mathbb{R}_+}
\newcommand{\Rp}{\mathbb{R}_{\ge 0}}
\def\Rnne{\mathbb{R}_{\ge 0}}
\def\Se{\bar{S}} \def\Ie{\bar{I}}
\newcommand{\Rpos}{\mathbb{R}_{>0}}
\newcommand{\dd}{\,\mathrm d}
\newcommand{\e}{\mathrm e}
%% ========================================
%% GREEK LETTERS (2+ character abbreviations) greek
%% ========================================
\def\al{\alpha}\def\ba{\bff \alpha}
\def\w{\omega}\def\Lgn{\ga+\mu+\de}\def\om{\omega}
\def\de{\delta}\def\De{\Delta}\def\bde{\bff \delta}
\def\ep{\epsilon}\def\eps{\varepsilon}
\def\f{\varphi}  % Kept as requested
\def\ga{\gamma}\def\Ga{\Gamma}\def\G{\Gamma}\def\bg{\bff \gamma}
\def\up{\upsilon}  % NOT the breve accent: no accent commands anywhere.
                   % Named \up, not \u, so the LaTeX breve stays reachable.
\def\la{\lambda}\def\La{\Lambda}
\def\si{\sigma}\def\ka{\kappa}
\def\th{\theta}\def\vt{\vartheta}
\def\ze{\zeta}\def\Gap{\Gamma^{+}} \def\Gam{\Gamma^{-}}


%% ========================================
%% SPECIAL SHORT COMMANDS (kept as requested)
%% ========================================
\def\I{\infty}
\def\T{\widetilde}
\def\H{\hat}
\newcommand{\relint}{\operatorname{relint}}

%% ========================================
%% BOLD/VECTOR NOTATION
%% ========================================
\newcommand{\bff}[1]{{\mbox{\boldmath$#1$}}}
\newcommand\X{\boldsymbol{X}}
\def\x{\boldsymbol{x}}\def\bx{\bff x}\def\xe{\bar{\x}}
\newcommand\y{\boldsymbol{y}}\def\by{\boldsymbol{y}}\def\ye{\bar{\y}}
\def\yb{\y}\def\xb{\x}
\def\xy{\x-\y}
\newcommand\br{\boldsymbol{r}}\def\re{\bar{\r}}
\newcommand\bd{\boldsymbol{d}}
\newcommand\bQ{\boldsymbol{Q}}
\newcommand{\Yb}{{\mathbf{Y}}}
\newcommand{\Ab}{\mathbf{A}}
\newcommand{\zb}{\mathbf{z}}\def\bz{\bff 0}
\newcommand{\Zb}{\mathbf{Z}}
\newcommand{\bb}{\mathbf{b}}
\def\vbe{\vec {\bff \beta}}
\def\vb{\vec {\bff b}}
\def\v1{\vec {\bff 1}}
\def\vn{\vec \de}
\def\vz{\vec 0}
\def\vi{\overset{\rightarrow}{\mathsf i}}
\def\vv{\overset{\rightarrow}{v}}
\def\vr{\; \vec {\mathsf r}}
\newcommand{\va}{\bff a}
\providecommand{\vect}[1]{\boldsymbol{#1}}   % bold vector (moved from Gradostat.tex)
\providecommand{\one}{\mathbf 1}             % all-ones vector (moved from AliSF)
\providecommand{\spr}{s}                     % spectral bound  (moved from AliSF)
\providecommand{\Ke}{K_{\rm eff}}            % effective carrying capacity
% Critical right/left eigenvectors in the Hopf projection method (AAHOSN/New.tex)
\providecommand{\bq}{\mathbf{q}}
\providecommand{\bp}{\mathbf{p}}


%% ========================================
%% MATHCAL & MATHSCR
%% ========================================
% Mathcal
\newcommand\mc[1]{\mathcal{#1}}
\def\mC{{\mathcal C}}\def\mD{\mathcal D}\def\mE{\mathcal{E}}
\def\mF{\mathcal F}\def\mG{{\mathcal G}}\def\mH{{\mathcal H}}
\def\mI{\mathcal I}\def\mL{\mathcal{L}}\def\mM{{\mathcal M}}
\def\mN{{\mathcal N}}\def\mP{\mathcal{P}}\def\mR{{\mathcal R}}
\def\mS{{\mathcal S}}\def\mT{\mathcal{T}}\def\mV{{\mathcal V}}
\def\mK{\mathcal K}\def\mA{{\mathcal A}}\def\mU{{\mathcal U}}
% Symmetric bilinear and trilinear forms of the Taylor expansion, used in the
% first Lyapunov coefficient (AAHOSN/New.tex)
\providecommand{\Bform}{\mathcal{B}}
\providecommand{\Cform}{\mathcal{C}}

% Blackboard bold.  Kept here, not in the drivers: a \providecommand in a
% driver preamble runs before this file and silently wins, so the same name
% ends up meaning different things in different builds.
\def\IC{\mathbb{C}}\def\IN{\mathbb{N}}\def\IQ{\mathbb{Q}}
\def\IR{\mathbb{R}}\def\IZ{\mathbb{Z}}\def\IG{\mathbb{G}}
\def\II{\mathbb{I}}

% Mathscr
\newcommand{\Rs}{{\mathscr R}}
\newcommand{\Ss}{{\mathscr S}}
\newcommand{\Ns}{{\mathscr N}}
\newcommand{\Cs}{\mathscr C}


%% ========================================
%% OPERATORS & FUNCTIONS
%% ========================================
\DeclareMathOperator{\tr}{tr}
\DeclareMathOperator{\diag}{Diag}
\DeclareMathOperator{\Diag}{Diag}
\DeclareMathOperator{\intr}{\mathrm{Int}}
\DeclareMathOperator{\Inv}{Inv}
\DeclareMathOperator{\spec}{spec}
\DeclareMathOperator{\AUC}{AUC}               % area under the curve
\newcommand{\ii}{\mathrm{i}}                  % imaginary unit
\newcommand\op{\operatorname}
\newcommand{\supp}{\op{supp}}
\newcommand{\diff}{\mathop{}\mathopen{}\mathrm{d}}
\newcommand{\pdd}[1]{\frac{\partial}{\partial{#1}}}
\DeclareMathOperator{\cone}{cone}
\DeclareMathOperator{\proj}{proj}

%% ========================================
%% BRACKETS & NORMS
%% ========================================
\providecommand{\pr}[1]{\left(#1\right)}
\providecommand{\pp}[1]{\left[#1\right]}
\newcommand\croc[1]{\left\langle #1\right\rangle}
\newcommand{\norm}[1]{\|#1\|}
\newcommand{\inner}[2]{\langle #1, #2 \rangle}
% Short synonym of \inner, used by AAHOSN/New.tex as $\ip{u}{v}=\bar u^\top v$
\providecommand{\ip}[2]{\langle #1 , #2 \rangle}


%% ========================================
%% FORMATTING & UTILITIES
%% ========================================
\newcommand*{\Scale}[2][4]{\scalebox{#1}{$#2$}}
\newcommand{\te}[1]{\text{#1}}
\newcommand\steq[1]{\stackrel{\text{\rm #1.}}{=}}
\newcommand\blue[1]{\textcolor{blue}{#1}}
\newcommand{\red}{\textcolor[rgb]{1.00,0.00,0.00}}
\newcommand{\green}{\textcolor[rgb]{0.00,1.00,0.00}}

% Figure command
\newcommand{\figu}[3]{%
\begin{figure}[H]
\centering
\includegraphics[scale=#3]{#1}
\caption{#2\label{f:#1}}
\end{figure}
}

% Footnote
\long\def\symbolfootnote[#1]#2{%
\begingroup
\def\thefootnote{\fnsymbol{footnote}}\footnote[#1]{#2}%
\endgroup}
\def\fn{\symbolfootnote}

\newcommand{\dsum}{\displaystyle\sum}

%% ========================================
%% CRN & EPIDEMIOLOGY SPECIFIC
%% ========================================
% CRN notation
\def\SS{\mathcal S}
\def\CC{\mathcal C}
\def\RR{\mathcal R}

% Reproduction numbers
\def\Rto{R_{1}^{E_2}}
\def\Rot{R_{2}^{E_1}}
\def\Roc{R_{1,c}}
\def\Rtc{R_{2,c}}
\def\rd{r_{0}}
\def\sd{s_{0}}
\def\dd{D_{0}}
\def\vd{v_{0}}
\def\zd{z_{0}}
\def\se{\sE}
\def\xd{x_{0}}
%% ========================================
%% SPECIAL VALUES & EQUILIBRIA
%% ========================================
\def\eE{\bar{e}}
\def\wE{\bar{w}}
\def\sE{s_{*}}
\def\Me{\T M_{ee}}
\def\Rnne{\mathbb{R}_{\ge 0}}

\newcommand{\ones}{\mathbf 1}


%% ========================================
%% ABBREVIATIONS (alphabetical by category)
%% ========================================

% -- General mathematical terms --
\def\adm{admissible}\def\admy{admissibility}
%\def\and{antisymmetric}% disabled 2026-09-30: it overrode LaTeX's author separator \and (printed "antisymmetric" between authors)
\def\ass{assumption}
\def\cadlag{c\`adl\`ag}
\def\ch{characteristic polynomial}
\def\chp{characteristic polynomial}
\def\coe{coefficient}
\def\comp{computation}
\def\con{condition}
\def\coo{coordinate}
\def\corr{corresponding}
\def\dec{decomposition}
\def\den{denominator}
\def\dete{deterministic}
\def\eig{eigenvalue}\def\eval{eigenvalue}
\def\eigv{eigenvector}\def\eiv{eigenvector}\def\evec{eigenvector}
\def\elt{element}
\def\enn{essentially nonnegative/positivity preserving}
\def\equ{equation}
\def\etal{et al.}
\def\fe{for example}
\def\fno{from now on}
\def\for{\forall}
\def\form{formula}
\def\Fr{Furthermore, }
\def\frt{furthermore}
\def\fun{function}
\def\gene{generalization}
\def\heu{heuristic}
\def\how{however}\def\How{However, }
\def\ie{i.e. }
\def\im{\item}
\def\isd{is defined by}
\def\itc{it is easy to check that}
\def\itf{it follows that}\def\Itf{It follows that}
\def\ith{it holds that}
\def\itm{it may be checked that}\def\Itm{It may be checked that}
\def\ito{it turns out that}\def\Ito{It turns out that}\def\Itou{It turns out useful to}
\def\ma{mass action}\def\MAR{mass action representation}
\def\mat{matrix}
\def\mbc{may be checked to be}
\def\mbw{may be written as}\def\mubw{must be written as}
\def\mpr{more precisely, }\def\Mp{More precisely, }
\def\nc{non-catalytic}
\def\nne{non-negative}\def\nny{nonnegativity}
\def\no{\nonumber}
\def\noni{non-infectious}
\def\nota{notation}
\def\np{\newpage}
\def\Nt{Note that}
\def\nt{non trivial}\def\nz{non zero}
\def\oth{otherwise}
\def\Oth{On the other hand}
\def\pa{parameter}\def\para{parameter}
\def\pol{polynomial}
\def\prob{probability}
\def\Prf{{\bf  Proof:}}\def\prf{{\bf  Proof:}}
\def\prop{proportional to}
\def\qu{\quad}
\def\QED{\hfill {$\square$}\goodbreak \medskip}
\def\rep{representation}
\def\resp{respectively}
\def\satd{satisfied}\def\satg{satisfying}\def\sats{satisfies}\def\saty{satisfy}
\def\ssec{\subsection}\def\sssec{\subsubsection}
\def\sef{see for example}
\def\sksy{skew-symmetric}
\def\sub{substrate}
\def\sys{system}
\def\var{variable}
\def\wk{well-known}
\def\wl{without loss of generality}\def\wlo{w.l.o.g.}
\def\wmc{We may check that}\def\Wmc{We may check that}
\def\wmr{we may rewrite}
\def\wms{we must show that}
\def\wrt{with respect to}

% -- CRN & Chemistry --
\newcommand\CRN{chemical reaction network}
\newcommand\CRNs{chemical reaction networks}
\newcommand\RN{reaction network}
\def\ACR{absolute concentration robustness}
\def\AMSD{acyclic siphon decomposition}
\def\DAE{differential algebraic equations}
\def\FHJ{Feinberg-Horn-Jackson}
\def\GMAK{generalized mass action kinetics}
\def\MM{Michaelis-Menten}
\def\sm{stoichiometric matrix}\def\Sm{stoichiometric matrix}
\def\WRZD{weakly reversible, zero deficiency}\def\aut{autocatalytic}

% -- Epidemiology --
\def\brn{basic reproduction number}\def\bzn{basic replacement number}
\def\com{compartment}
\def\dfe{disease free equilibrium}\def\DFE{disease free equilibrium}
\def\dwt{dwell times}\def\edt{expected dwell times}
\def\EM{epidemic model}
\def\endp{endemic fixed point}\def\enp{endemic point}\def\EE{endemic point}
\def\ME{mathematical epidemiology}
\def\MEE{eco-epidemiology}\def\EEP{eco-epidemiological}
\def\ece{eco-epidemiological model}
\def\mgm{multi group model}
\def\mtS{multi-strain model}\def\mSt{multi-strain model}
\def\mSi{minimal siphon}
\def\NGM{next generation matrix}
\def\nim{new infections matrix}
\def\qem{quadratic epidemiological model}
\def\repF{invasion functions}
\def\SPF{SIR-PH-FA model}
\def\sus{susceptible}

% -- Graph theory & Linear algebra --
\def\jac{Jacobi matrix}
\def\jin{Jacobian matrix of the invasion vector field with respect to the invasion variables}
\def\jrn{Jacobian matrix of the resident vector field with respect to the resident variables}
\def\jxn{Jacobian matrix of the $\x$ vector field with respect to the $\x$ variables}
\def\Lm{Laplacian matrix}
\def\MTT{matrix tree theorem}
\def\SGT{spectral graph theory}
\def\SMW{Sherman–Morrison–Woodbury}

% -- Stability & Dynamics --
\def\absC{absorbtion CTMC}
\def\BTb{Bogdanov-Takens bifurcation}
\def\CEP{competitive exclusion partition}
\def\CTMC{continuous time Markov chain}
\def\ds{differential system}
\def\fp{fixed point}\def\bfp{boundary fixed point}\def\fbp{fixed boundary point}
\def\fbl{fixed boundary line}\def\fif{forward invariant face}
\def\gas{global asymptotic stability}\def\GAS{globally asymptotically stable}
\def\GASS{\text{GASS}}
\def\Hb{Hopf bifurcation}
\def\Hur{Hurwitz}\def\MH{Metzler-Hurwitz}\def\RH{Routh-Hurwitz}
\def\IS{\text{IS}}
\def\LAS{locally asymptotically stable}
\def\LdS{Lyapunov diagonally stable}
\def\Lf{Lyapunov function}\def\LyV{Lyapunov Volterra type function}\def\RLF{robust Lyapunov function}
\def\LyM{Lyapunov Malkin stable}
\def\PF{Perron-Frobenius}
\def\spfp{strictly positive fixed point}
\def\sta{stability}
\def\STP{sharp threshold property}\def\STT{sharp threshold theorem}

% -- Models (with \varphi explicitly instead of \f) --
\def\AB{$(A,B)$ Arino-Brauer epidemic models}
\def\ABa{$(A,\bbe,\ba,\varphi)$ balanced bilinear model}
\def\ABP{%$(A,\bbe,P,\varphi)$
\\ balanced bilinear model}
\def\AMG{Anderson May Gupta}
\newcommand\BIN{biological interaction network}
\def\Dt{Descartes type}
\def\gLV{generalized Lotka-Volterra model}
\def\IG{\text{invasion graph}}\def\IGMS{interaction graph on minimal siphons}
\def\LCP{\text{linear complementarity problem}}
\def\LV{Lotka-Volterra}\def\LVK{Lotka-Volterra-Kolmogorov}

% -- Splitting & Structure --
\def\reg{regular splitting}\def\regS{{\bf regular splitting}}
\def\RUR{rational univariate representation}
\def\bds{block-diagonal structure}\def\NGMs{next generation matrices}

% -- Miscellaneous --
\def\Ma{{Mathematica}}
\def\Mm{Max–Min}
\def\num{numerator}
\def\rap{rational approximation} \def\DSR{directed species-reaction bipartite graph}\def\PCP{polynomial complementarity problem}\def\rat{rational substitution}
\def\QSS{quasi steady state} \def\RHS{right hand side} \def\LHS{left hand side}
\def\OSN{online social network}\def\VdD{Van den Driessche}
\def\wli{$\omega$-limit} \def\rk{rich kinetics}\def\sJ{symbolic Jacobian}
\def\cV{\cite{Vas,VasSta,BSV,GSV}} \def\spf{semi-parametric semi-functional}
\def\CRW{Capasso-Ruan-Wang} \def\ABP{$(A,\bbe,P,\f)$ balanced bilinear model}
\def\ABa{%$(A,\bbe,\ba,\f)$ \\
rank-one balanced bilinear model}
\def\bbe{B%{\bff \beta}
} \def\vbe{\vec {\bff \beta}
}
\def\vb{\vec {\bff b}}
\newcommand{\Kh}{\hat{K}}          % K with hat, first mention only
\newcommand{\wS}{w^{\sigma}}       % right Perron J^\perp, full block
\newcommand{\piS}{\vec \pi^{\sigma}}    % left Perron J^\perp, full block
\newcommand{\ws}{w^{\sigma}}       % right Perron J^\perp, sub-block
\newcommand{\pis}{\pi^{\sigma}}    % left Perron J^\perp, sub-block
\newcommand{\hwS}{\hat{w}^{\sigma}}   % right Perron K, full block
\newcommand{\hpiS}{\hat{\pi}^{\sigma}} % left Perron K, full block
\newcommand{\hws}{\hat{w}^{\sigma}}   % right Perron K, sub-block
\newcommand{\hpis}{\hat{\pi}^{\sigma}} % left Perron K, sub-block
\newcommand{\twS}{\tilde{w}^{\sigma}}  % right Perron \tilde K, full block
\newcommand{\tpiS}{\tilde{\pi}^{\sigma}} % left Perron \tilde K
\newcommand{\tws}{\tilde{w}^{\sigma}}
\newcommand{\tpis}{\tilde{\pi}^{\sigma}}
\newcommand{\Mat}[2]{\mathbb{R}_{\ge 0}^{\,#1\times #2}}
\def\tJ{J_\sigma^\perp} \def\lPe{left Perron eigenvector}
% \G is \Gamma (defined at :795).  The Volterra function keeps its own name,
% so that \G no longer has two conflicting meanings in this file.
\def\GV{G}  % Volterra G(u)=u-1-\ln u
\def\PVL{Perron-Volterra Lyapunov minimal siphon ansatz}
\newcommand{\istar}[1]{i_{#1}^{*}}
%\newcommand{\sE}{s_E}
\newcommand{\vE}{v_E}
\def\PV{Perron-Volterra candidate Lyapunov function}
\def\GLV{generalized Lotka-Volterra}
\def\MERN{mathematical epidemiology reaction networks}
\def\CMM{contraction metric method}
\newcommand{\sign}{\operatorname{sign}}
\newcommand{\Var}{\operatorname{Var}} 

% ======================================================================
% IDEAS SHARED BETWEEN THE ABSTRACT AND THE MAIN RESULTS SECTION
%
% Write each idea ONCE, in the abstract; the Main Results section recalls it,
% so the two can never drift apart.
%
%   in the abstract        \ideaitem{1}{The first-order reaction terms ...}
%   in Main Results        \ideahead{1}
%
% \ideaitem stores the text globally and typesets
%   \item \hyperlink{idea:1}{<text>}
% \ideahead typesets
%   \item\hypertarget{idea:1}{}\textbf{<the same text>}
%
% The abstract must come before the Main Results input, which is the normal
% order. If a heading has to read differently from the abstract sentence (for
% instance a long abstract item summarised in one line), override it with the
% optional argument: \ideahead[short heading]{6}.
% ======================================================================
\makeatletter
\newcommand{\ideaitem}[2]{%
  \expandafter\gdef\csname @idea@#1\endcsname{#2}%
  \item\hyperlink{idea:#1}{#2}}
\newcommand{\ideahead}[2][]{%
  \item\hypertarget{idea:#2}{}%
  \def\@ideatmp{#1}%
  \ifx\@ideatmp\@empty
    \@ifundefined{@idea@#2}{%
      \PackageWarning{preamble}{Idea #2 is used in Main Results but was never
        defined by \string\ideaitem\space in the abstract}%
      \textbf{[idea #2 undefined]}%
    }{\textbf{\begingroup\def\ideaterm##1##2{##2}\def\idearest##1{}\@nameuse{@idea@#2}\endgroup}}%
  \else
    \textbf{#1}%
  \fi}
% \ideaterm{<hypertarget>}{<words>}: inside an abstract idea the words become an
% invisible link (the document sets \hypersetup{hidelinks}) jumping to the place
% in Main Results that defines them; inside the Main Results heading itself,
% where the target sits a few lines below, \ideahead strips the link and prints
% the words plainly. Mark the destination with \hypertarget{<name>}{} on the line
% before the display, which adds no equation number and shows nothing.
\newcommand{\ideaterm}[2]{\hyperlink{#1}{#2}}
% \idearest{<text>}: the part of an abstract idea after its first sentence; printed in
% the abstract, omitted from the Main Results heading.
\newcommand{\idearest}[1]{#1}
\makeatother

% ======================================================================
% SHORT-PAPER SWITCH
% \shorttrue in a master that assembles only part of the companion paper
% (Barrat/SipArxiv.tex). \nsh{...} drops text that cites results living in the
% omitted sections, so the short build has no dangling references while the
% companion paper keeps the text verbatim.
% ======================================================================
\newif\ifshort \shortfalse
\newcommand{\nsh}[1]{\ifshort\else#1\fi}
% \shortalt{<short build>}{<companion paper>}: use when a sentence needs
% different wording in the two builds, not merely a deletion.
\newcommand{\shortalt}[2]{\ifshort#1\else#2\fi}
% \Jtr, \JtrS{\Sigma}: transverse Jacobian. J^\perp here; SipArxiv.tex redefines them as J_{\mathrm{tra}}, J_{\mathrm{tra},\Sigma}
\newcommand{\Jtr}{J^\perp}
\newcommand{\JtrS}[1]{J_{#1}^\perp}
, and must not clash.
\providecommand{\Q}{\mathbb{Q}}
\providecommand{\res}{\operatorname{res}}   % resultant; used by epiRes.tex
\newcommand{\MN}{\mathbb{N}}
\newcommand{\MZ}{\mathbb{Z}}
\newcommand{\realpos}{\mathbb{R}_{> 0}}
\newcommand{\realnneg}{\mathbb{R}_{\geq 0}}
\newcommand{\naturalpos}{\mathbb{N}_{>0}}
\newcommand{\Rplus}{\mathbb{R}_+}
\newcommand{\Rp}{\mathbb{R}_{\ge 0}}
\def\Rnne{\mathbb{R}_{\ge 0}}
\def\Se{\bar{S}} \def\Ie{\bar{I}}
\newcommand{\Rpos}{\mathbb{R}_{>0}}
\newcommand{\dd}{\,\mathrm d}
\newcommand{\e}{\mathrm e}
%% ========================================
%% GREEK LETTERS (2+ character abbreviations) greek
%% ========================================
\def\al{\alpha}\def\ba{\bff \alpha}
\def\w{\omega}\def\Lgn{\ga+\mu+\de}\def\om{\omega}
\def\de{\delta}\def\De{\Delta}\def\bde{\bff \delta}
\def\ep{\epsilon}\def\eps{\varepsilon}
\def\f{\varphi}  % Kept as requested
\def\ga{\gamma}\def\Ga{\Gamma}\def\G{\Gamma}\def\bg{\bff \gamma}
\def\up{\upsilon}  % NOT the breve accent: no accent commands anywhere.
                   % Named \up, not \u, so the LaTeX breve stays reachable.
\def\la{\lambda}\def\La{\Lambda}
\def\si{\sigma}\def\ka{\kappa}
\def\th{\theta}\def\vt{\vartheta}
\def\ze{\zeta}\def\Gap{\Gamma^{+}} \def\Gam{\Gamma^{-}}


%% ========================================
%% SPECIAL SHORT COMMANDS (kept as requested)
%% ========================================
\def\I{\infty}
\def\T{\widetilde}
\def\H{\hat}
\newcommand{\relint}{\operatorname{relint}}

%% ========================================
%% BOLD/VECTOR NOTATION
%% ========================================
\newcommand{\bff}[1]{{\mbox{\boldmath$#1$}}}
\newcommand\X{\boldsymbol{X}}
\def\x{\boldsymbol{x}}\def\bx{\bff x}\def\xe{\bar{\x}}
\newcommand\y{\boldsymbol{y}}\def\by{\boldsymbol{y}}\def\ye{\bar{\y}}
\def\yb{\y}\def\xb{\x}
\def\xy{\x-\y}
\newcommand\br{\boldsymbol{r}}\def\re{\bar{\r}}
\newcommand\bd{\boldsymbol{d}}
\newcommand\bQ{\boldsymbol{Q}}
\newcommand{\Yb}{{\mathbf{Y}}}
\newcommand{\Ab}{\mathbf{A}}
\newcommand{\zb}{\mathbf{z}}\def\bz{\bff 0}
\newcommand{\Zb}{\mathbf{Z}}
\newcommand{\bb}{\mathbf{b}}
\def\vbe{\vec {\bff \beta}}
\def\vb{\vec {\bff b}}
\def\v1{\vec {\bff 1}}
\def\vn{\vec \de}
\def\vz{\vec 0}
\def\vi{\overset{\rightarrow}{\mathsf i}}
\def\vv{\overset{\rightarrow}{v}}
\def\vr{\; \vec {\mathsf r}}
\newcommand{\va}{\bff a}
\providecommand{\vect}[1]{\boldsymbol{#1}}   % bold vector (moved from Gradostat.tex)
\providecommand{\one}{\mathbf 1}             % all-ones vector (moved from AliSF)
\providecommand{\spr}{s}                     % spectral bound  (moved from AliSF)
\providecommand{\Ke}{K_{\rm eff}}            % effective carrying capacity
% Critical right/left eigenvectors in the Hopf projection method (AAHOSN/New.tex)
\providecommand{\bq}{\mathbf{q}}
\providecommand{\bp}{\mathbf{p}}


%% ========================================
%% MATHCAL & MATHSCR
%% ========================================
% Mathcal
\newcommand\mc[1]{\mathcal{#1}}
\def\mC{{\mathcal C}}\def\mD{\mathcal D}\def\mE{\mathcal{E}}
\def\mF{\mathcal F}\def\mG{{\mathcal G}}\def\mH{{\mathcal H}}
\def\mI{\mathcal I}\def\mL{\mathcal{L}}\def\mM{{\mathcal M}}
\def\mN{{\mathcal N}}\def\mP{\mathcal{P}}\def\mR{{\mathcal R}}
\def\mS{{\mathcal S}}\def\mT{\mathcal{T}}\def\mV{{\mathcal V}}
\def\mK{\mathcal K}\def\mA{{\mathcal A}}\def\mU{{\mathcal U}}
% Symmetric bilinear and trilinear forms of the Taylor expansion, used in the
% first Lyapunov coefficient (AAHOSN/New.tex)
\providecommand{\Bform}{\mathcal{B}}
\providecommand{\Cform}{\mathcal{C}}

% Blackboard bold.  Kept here, not in the drivers: a \providecommand in a
% driver preamble runs before this file and silently wins, so the same name
% ends up meaning different things in different builds.
\def\IC{\mathbb{C}}\def\IN{\mathbb{N}}\def\IQ{\mathbb{Q}}
\def\IR{\mathbb{R}}\def\IZ{\mathbb{Z}}\def\IG{\mathbb{G}}
\def\II{\mathbb{I}}

% Mathscr
\newcommand{\Rs}{{\mathscr R}}
\newcommand{\Ss}{{\mathscr S}}
\newcommand{\Ns}{{\mathscr N}}
\newcommand{\Cs}{\mathscr C}


%% ========================================
%% OPERATORS & FUNCTIONS
%% ========================================
\DeclareMathOperator{\tr}{tr}
\DeclareMathOperator{\diag}{Diag}
\DeclareMathOperator{\Diag}{Diag}
\DeclareMathOperator{\intr}{\mathrm{Int}}
\DeclareMathOperator{\Inv}{Inv}
\DeclareMathOperator{\spec}{spec}
\DeclareMathOperator{\AUC}{AUC}               % area under the curve
\newcommand{\ii}{\mathrm{i}}                  % imaginary unit
\newcommand\op{\operatorname}
\newcommand{\supp}{\op{supp}}
\newcommand{\diff}{\mathop{}\mathopen{}\mathrm{d}}
\newcommand{\pdd}[1]{\frac{\partial}{\partial{#1}}}
\DeclareMathOperator{\cone}{cone}
\DeclareMathOperator{\proj}{proj}

%% ========================================
%% BRACKETS & NORMS
%% ========================================
\providecommand{\pr}[1]{\left(#1\right)}
\providecommand{\pp}[1]{\left[#1\right]}
\newcommand\croc[1]{\left\langle #1\right\rangle}
\newcommand{\norm}[1]{\|#1\|}
\newcommand{\inner}[2]{\langle #1, #2 \rangle}
% Short synonym of \inner, used by AAHOSN/New.tex as $\ip{u}{v}=\bar u^\top v$
\providecommand{\ip}[2]{\langle #1 , #2 \rangle}


%% ========================================
%% FORMATTING & UTILITIES
%% ========================================
\newcommand*{\Scale}[2][4]{\scalebox{#1}{$#2$}}
\newcommand{\te}[1]{\text{#1}}
\newcommand\steq[1]{\stackrel{\text{\rm #1.}}{=}}
\newcommand\blue[1]{\textcolor{blue}{#1}}
\newcommand{\red}{\textcolor[rgb]{1.00,0.00,0.00}}
\newcommand{\green}{\textcolor[rgb]{0.00,1.00,0.00}}

% Figure command
\newcommand{\figu}[3]{%
\begin{figure}[H]
\centering
\includegraphics[scale=#3]{#1}
\caption{#2\label{f:#1}}
\end{figure}
}

% Footnote
\long\def\symbolfootnote[#1]#2{%
\begingroup
\def\thefootnote{\fnsymbol{footnote}}\footnote[#1]{#2}%
\endgroup}
\def\fn{\symbolfootnote}

\newcommand{\dsum}{\displaystyle\sum}

%% ========================================
%% CRN & EPIDEMIOLOGY SPECIFIC
%% ========================================
% CRN notation
\def\SS{\mathcal S}
\def\CC{\mathcal C}
\def\RR{\mathcal R}

% Reproduction numbers
\def\Rto{R_{1}^{E_2}}
\def\Rot{R_{2}^{E_1}}
\def\Roc{R_{1,c}}
\def\Rtc{R_{2,c}}
\def\rd{r_{0}}
\def\sd{s_{0}}
\def\dd{D_{0}}
\def\vd{v_{0}}
\def\zd{z_{0}}
\def\se{\sE}
\def\xd{x_{0}}
%% ========================================
%% SPECIAL VALUES & EQUILIBRIA
%% ========================================
\def\eE{\bar{e}}
\def\wE{\bar{w}}
\def\sE{s_{*}}
\def\Me{\T M_{ee}}
\def\Rnne{\mathbb{R}_{\ge 0}}

\newcommand{\ones}{\mathbf 1}


%% ========================================
%% ABBREVIATIONS (alphabetical by category)
%% ========================================

% -- General mathematical terms --
\def\adm{admissible}\def\admy{admissibility}
%\def\and{antisymmetric}% disabled 2026-09-30: it overrode LaTeX's author separator \and (printed "antisymmetric" between authors)
\def\ass{assumption}
\def\cadlag{c\`adl\`ag}
\def\ch{characteristic polynomial}
\def\chp{characteristic polynomial}
\def\coe{coefficient}
\def\comp{computation}
\def\con{condition}
\def\coo{coordinate}
\def\corr{corresponding}
\def\dec{decomposition}
\def\den{denominator}
\def\dete{deterministic}
\def\eig{eigenvalue}\def\eval{eigenvalue}
\def\eigv{eigenvector}\def\eiv{eigenvector}\def\evec{eigenvector}
\def\elt{element}
\def\enn{essentially nonnegative/positivity preserving}
\def\equ{equation}
\def\etal{et al.}
\def\fe{for example}
\def\fno{from now on}
\def\for{\forall}
\def\form{formula}
\def\Fr{Furthermore, }
\def\frt{furthermore}
\def\fun{function}
\def\gene{generalization}
\def\heu{heuristic}
\def\how{however}\def\How{However, }
\def\ie{i.e. }
\def\im{\item}
\def\isd{is defined by}
\def\itc{it is easy to check that}
\def\itf{it follows that}\def\Itf{It follows that}
\def\ith{it holds that}
\def\itm{it may be checked that}\def\Itm{It may be checked that}
\def\ito{it turns out that}\def\Ito{It turns out that}\def\Itou{It turns out useful to}
\def\ma{mass action}\def\MAR{mass action representation}
\def\mat{matrix}
\def\mbc{may be checked to be}
\def\mbw{may be written as}\def\mubw{must be written as}
\def\mpr{more precisely, }\def\Mp{More precisely, }
\def\nc{non-catalytic}
\def\nne{non-negative}\def\nny{nonnegativity}
\def\no{\nonumber}
\def\noni{non-infectious}
\def\nota{notation}
\def\np{\newpage}
\def\Nt{Note that}
\def\nt{non trivial}\def\nz{non zero}
\def\oth{otherwise}
\def\Oth{On the other hand}
\def\pa{parameter}\def\para{parameter}
\def\pol{polynomial}
\def\prob{probability}
\def\Prf{{\bf  Proof:}}\def\prf{{\bf  Proof:}}
\def\prop{proportional to}
\def\qu{\quad}
\def\QED{\hfill {$\square$}\goodbreak \medskip}
\def\rep{representation}
\def\resp{respectively}
\def\satd{satisfied}\def\satg{satisfying}\def\sats{satisfies}\def\saty{satisfy}
\def\ssec{\subsection}\def\sssec{\subsubsection}
\def\sef{see for example}
\def\sksy{skew-symmetric}
\def\sub{substrate}
\def\sys{system}
\def\var{variable}
\def\wk{well-known}
\def\wl{without loss of generality}\def\wlo{w.l.o.g.}
\def\wmc{We may check that}\def\Wmc{We may check that}
\def\wmr{we may rewrite}
\def\wms{we must show that}
\def\wrt{with respect to}

% -- CRN & Chemistry --
\newcommand\CRN{chemical reaction network}
\newcommand\CRNs{chemical reaction networks}
\newcommand\RN{reaction network}
\def\ACR{absolute concentration robustness}
\def\AMSD{acyclic siphon decomposition}
\def\DAE{differential algebraic equations}
\def\FHJ{Feinberg-Horn-Jackson}
\def\GMAK{generalized mass action kinetics}
\def\MM{Michaelis-Menten}
\def\sm{stoichiometric matrix}\def\Sm{stoichiometric matrix}
\def\WRZD{weakly reversible, zero deficiency}\def\aut{autocatalytic}

% -- Epidemiology --
\def\brn{basic reproduction number}\def\bzn{basic replacement number}
\def\com{compartment}
\def\dfe{disease free equilibrium}\def\DFE{disease free equilibrium}
\def\dwt{dwell times}\def\edt{expected dwell times}
\def\EM{epidemic model}
\def\endp{endemic fixed point}\def\enp{endemic point}\def\EE{endemic point}
\def\ME{mathematical epidemiology}
\def\MEE{eco-epidemiology}\def\EEP{eco-epidemiological}
\def\ece{eco-epidemiological model}
\def\mgm{multi group model}
\def\mtS{multi-strain model}\def\mSt{multi-strain model}
\def\mSi{minimal siphon}
\def\NGM{next generation matrix}
\def\nim{new infections matrix}
\def\qem{quadratic epidemiological model}
\def\repF{invasion functions}
\def\SPF{SIR-PH-FA model}
\def\sus{susceptible}

% -- Graph theory & Linear algebra --
\def\jac{Jacobi matrix}
\def\jin{Jacobian matrix of the invasion vector field with respect to the invasion variables}
\def\jrn{Jacobian matrix of the resident vector field with respect to the resident variables}
\def\jxn{Jacobian matrix of the $\x$ vector field with respect to the $\x$ variables}
\def\Lm{Laplacian matrix}
\def\MTT{matrix tree theorem}
\def\SGT{spectral graph theory}
\def\SMW{Sherman–Morrison–Woodbury}

% -- Stability & Dynamics --
\def\absC{absorbtion CTMC}
\def\BTb{Bogdanov-Takens bifurcation}
\def\CEP{competitive exclusion partition}
\def\CTMC{continuous time Markov chain}
\def\ds{differential system}
\def\fp{fixed point}\def\bfp{boundary fixed point}\def\fbp{fixed boundary point}
\def\fbl{fixed boundary line}\def\fif{forward invariant face}
\def\gas{global asymptotic stability}\def\GAS{globally asymptotically stable}
\def\GASS{\text{GASS}}
\def\Hb{Hopf bifurcation}
\def\Hur{Hurwitz}\def\MH{Metzler-Hurwitz}\def\RH{Routh-Hurwitz}
\def\IS{\text{IS}}
\def\LAS{locally asymptotically stable}
\def\LdS{Lyapunov diagonally stable}
\def\Lf{Lyapunov function}\def\LyV{Lyapunov Volterra type function}\def\RLF{robust Lyapunov function}
\def\LyM{Lyapunov Malkin stable}
\def\PF{Perron-Frobenius}
\def\spfp{strictly positive fixed point}
\def\sta{stability}
\def\STP{sharp threshold property}\def\STT{sharp threshold theorem}

% -- Models (with \varphi explicitly instead of \f) --
\def\AB{$(A,B)$ Arino-Brauer epidemic models}
\def\ABa{$(A,\bbe,\ba,\varphi)$ balanced bilinear model}
\def\ABP{%$(A,\bbe,P,\varphi)$
\\ balanced bilinear model}
\def\AMG{Anderson May Gupta}
\newcommand\BIN{biological interaction network}
\def\Dt{Descartes type}
\def\gLV{generalized Lotka-Volterra model}
\def\IG{\text{invasion graph}}\def\IGMS{interaction graph on minimal siphons}
\def\LCP{\text{linear complementarity problem}}
\def\LV{Lotka-Volterra}\def\LVK{Lotka-Volterra-Kolmogorov}

% -- Splitting & Structure --
\def\reg{regular splitting}\def\regS{{\bf regular splitting}}
\def\RUR{rational univariate representation}
\def\bds{block-diagonal structure}\def\NGMs{next generation matrices}

% -- Miscellaneous --
\def\Ma{{Mathematica}}
\def\Mm{Max–Min}
\def\num{numerator}
\def\rap{rational approximation} \def\DSR{directed species-reaction bipartite graph}\def\PCP{polynomial complementarity problem}\def\rat{rational substitution}
\def\QSS{quasi steady state} \def\RHS{right hand side} \def\LHS{left hand side}
\def\OSN{online social network}\def\VdD{Van den Driessche}
\def\wli{$\omega$-limit} \def\rk{rich kinetics}\def\sJ{symbolic Jacobian}
\def\cV{\cite{Vas,VasSta,BSV,GSV}} \def\spf{semi-parametric semi-functional}
\def\CRW{Capasso-Ruan-Wang} \def\ABP{$(A,\bbe,P,\f)$ balanced bilinear model}
\def\ABa{%$(A,\bbe,\ba,\f)$ \\
rank-one balanced bilinear model}
\def\bbe{B%{\bff \beta}
} \def\vbe{\vec {\bff \beta}
}
\def\vb{\vec {\bff b}}
\newcommand{\Kh}{\hat{K}}          % K with hat, first mention only
\newcommand{\wS}{w^{\sigma}}       % right Perron J^\perp, full block
\newcommand{\piS}{\vec \pi^{\sigma}}    % left Perron J^\perp, full block
\newcommand{\ws}{w^{\sigma}}       % right Perron J^\perp, sub-block
\newcommand{\pis}{\pi^{\sigma}}    % left Perron J^\perp, sub-block
\newcommand{\hwS}{\hat{w}^{\sigma}}   % right Perron K, full block
\newcommand{\hpiS}{\hat{\pi}^{\sigma}} % left Perron K, full block
\newcommand{\hws}{\hat{w}^{\sigma}}   % right Perron K, sub-block
\newcommand{\hpis}{\hat{\pi}^{\sigma}} % left Perron K, sub-block
\newcommand{\twS}{\tilde{w}^{\sigma}}  % right Perron \tilde K, full block
\newcommand{\tpiS}{\tilde{\pi}^{\sigma}} % left Perron \tilde K
\newcommand{\tws}{\tilde{w}^{\sigma}}
\newcommand{\tpis}{\tilde{\pi}^{\sigma}}
\newcommand{\Mat}[2]{\mathbb{R}_{\ge 0}^{\,#1\times #2}}
\def\tJ{J_\sigma^\perp} \def\lPe{left Perron eigenvector}
% \G is \Gamma (defined at :795).  The Volterra function keeps its own name,
% so that \G no longer has two conflicting meanings in this file.
\def\GV{G}  % Volterra G(u)=u-1-\ln u
\def\PVL{Perron-Volterra Lyapunov minimal siphon ansatz}
\newcommand{\istar}[1]{i_{#1}^{*}}
%\newcommand{\sE}{s_E}
\newcommand{\vE}{v_E}
\def\PV{Perron-Volterra candidate Lyapunov function}
\def\GLV{generalized Lotka-Volterra}
\def\MERN{mathematical epidemiology reaction networks}
\def\CMM{contraction metric method}
\newcommand{\sign}{\operatorname{sign}}
\newcommand{\Var}{\operatorname{Var}} 

% ======================================================================
% IDEAS SHARED BETWEEN THE ABSTRACT AND THE MAIN RESULTS SECTION
%
% Write each idea ONCE, in the abstract; the Main Results section recalls it,
% so the two can never drift apart.
%
%   in the abstract        \ideaitem{1}{The first-order reaction terms ...}
%   in Main Results        \ideahead{1}
%
% \ideaitem stores the text globally and typesets
%   \item \hyperlink{idea:1}{<text>}
% \ideahead typesets
%   \item\hypertarget{idea:1}{}\textbf{<the same text>}
%
% The abstract must come before the Main Results input, which is the normal
% order. If a heading has to read differently from the abstract sentence (for
% instance a long abstract item summarised in one line), override it with the
% optional argument: \ideahead[short heading]{6}.
% ======================================================================
\makeatletter
\newcommand{\ideaitem}[2]{%
  \expandafter\gdef\csname @idea@#1\endcsname{#2}%
  \item\hyperlink{idea:#1}{#2}}
\newcommand{\ideahead}[2][]{%
  \item\hypertarget{idea:#2}{}%
  \def\@ideatmp{#1}%
  \ifx\@ideatmp\@empty
    \@ifundefined{@idea@#2}{%
      \PackageWarning{preamble}{Idea #2 is used in Main Results but was never
        defined by \string\ideaitem\space in the abstract}%
      \textbf{[idea #2 undefined]}%
    }{\textbf{\begingroup\def\ideaterm##1##2{##2}\def\idearest##1{}\@nameuse{@idea@#2}\endgroup}}%
  \else
    \textbf{#1}%
  \fi}
% \ideaterm{<hypertarget>}{<words>}: inside an abstract idea the words become an
% invisible link (the document sets \hypersetup{hidelinks}) jumping to the place
% in Main Results that defines them; inside the Main Results heading itself,
% where the target sits a few lines below, \ideahead strips the link and prints
% the words plainly. Mark the destination with \hypertarget{<name>}{} on the line
% before the display, which adds no equation number and shows nothing.
\newcommand{\ideaterm}[2]{\hyperlink{#1}{#2}}
% \idearest{<text>}: the part of an abstract idea after its first sentence; printed in
% the abstract, omitted from the Main Results heading.
\newcommand{\idearest}[1]{#1}
\makeatother

% ======================================================================
% SHORT-PAPER SWITCH
% \shorttrue in a master that assembles only part of the companion paper
% (Barrat/SipArxiv.tex). \nsh{...} drops text that cites results living in the
% omitted sections, so the short build has no dangling references while the
% companion paper keeps the text verbatim.
% ======================================================================
\newif\ifshort \shortfalse
\newcommand{\nsh}[1]{\ifshort\else#1\fi}
% \shortalt{<short build>}{<companion paper>}: use when a sentence needs
% different wording in the two builds, not merely a deletion.
\newcommand{\shortalt}[2]{\ifshort#1\else#2\fi}
% \Jtr, \JtrS{\Sigma}: transverse Jacobian. J^\perp here; SipArxiv.tex redefines them as J_{\mathrm{tra}}, J_{\mathrm{tra},\Sigma}
\newcommand{\Jtr}{J^\perp}
\newcommand{\JtrS}[1]{J_{#1}^\perp}
.
\makeatletter
\AtBeginDocument{\@ifundefined{hypersetup}{}{\hypersetup{linktoc=all}}}
\makeatother
\usepackage{graphicx,float,fancybox}
\usepackage{epsfig,epstopdf,tikz}
\usepackage[sans]{dsfont}
\usetikzlibrary{positioning,arrows,arrows.meta,calc}
% Configure listings for Mathematica code
\lstdefinestyle{mathematica}{
 language=Mathematica,
 basicstyle=\small\ttfamily,
 backgroundcolor=\color{gray!10},
 frame=single,
 breaklines=true,
 commentstyle=\color{green!60!black},
 keywordstyle=\color{blue},
 stringstyle=\color{red},
 showstringspaces=false,
 tabsize=2
}
\usepackage{tcolorbox}
% Colored boxes for key results
\newtcolorbox{keyresult}[1][]{
 colback=blue!5!white,
 colframe=blue!75!black,
 title=#1,
 fonttitle=\bfseries
}
\newtcolorbox{examplebox}[1][]{
 colback=green!5!white,
 colframe=green!75!black,
 title=#1,
 fonttitle=\bfseries
}


\makeatletter
\@ifundefined{theorem}   {\newtheorem{theorem}   {Theorem}}   {}
\@ifundefined{definition}{\newtheorem{definition}{Definition}} {}
\@ifundefined{open}      {\newtheorem{open}      {Problem}}   {}
\@ifundefined{example}   {\newtheorem{example}   {Example}}   {}
\makeatother

% COUNTER DEFINITIONS
\newcounter{tabcounter}
\newcounter{defcounter}
\newcounter{opcounter}
\newcounter{thmcounter}
\newcounter{excounter}
\newcounter{ascounter}
\newcounter{cjcounter}
\newcounter{lemcounter}
\newcounter{remcounter}
\newcounter{corcounter}
\newcounter{propcounter}
\newcounter{figcounter}

% UNIFIED GLOSSARY STYLE
\newglossarystyle{unifiedstyle}{%
  \setglossarystyle{list}%
  \renewcommand*{\glossarypreamble}{\gdef\glslastsec{}}%
  \renewcommand{\glossentry}[2]{%
    \edef\glscursec{\glsentryuseriii{##1}}%
    \ifx\glscursec\glslastsec\else
      \item[]\textbf{\S\,\glscursec\enspace\glsentryuseriv{##1}}%
      \global\let\glslastsec\glscursec
    \fi
    \item[\glsentryuserii{##1} \glsentryuseri{##1}:] ``\glossentryname{##1}'' \glsentrydesc{##1}%
  }%
}

\makeatletter
% First-word extractor (used by \beT, \beD and \beXa for the default key)
% \@firstwordof hello world \relax  ->  hello
\def\@firstwordof#1 #2\relax{#1}
% First three letters, used for the heading mnemonic of a theorem given no
% explicit key.  Padding with \@empty keeps it safe for first words shorter
% than three letters: \@empty vanishes under \edef, whereas \relax would
% survive into the label and is illegal inside \csname (it broke every
% two-letter title, e.g. "ME ...", with "Missing \endcsname inserted").
\def\@takethree#1#2#3#4\@nil{#1#2#3}

% EXAMPLE COMMAND -- same signature and behaviour as \beT and \beD.
% Usage: \beXa*[long title][short title][key]
%   star  = suppresses the printed "[key]" tag after the title (as for \beT*
%   and \beD*); the label and glossary entry are unaffected.
%   key = optional label key override; the label becomes e:KEY, and without a
%   key it is e: plus the first three letters of the title's first word.  A
%   mnemonic label lets an example be cited as "Example~\ref{e:KEY}" instead
%   of by section number, which is what the citing prose needs.
\NewDocumentCommand{\beXa}{s O{} O{} o}{%
  \def\@temparg{#2}%
  \ifx\@temparg\@empty
    \begin{example}%
  \else
    \let\@savedSecTitle\@currentsectionname%
    \refstepcounter{excounter}%
    \IfValueTF{#4}{%
      \edef\@exlblkey{#4}%
    }{%
      \edef\@exfirstword{\@firstwordof#2 \relax}%
      \edef\@exlblkey{%
        \expandafter\@takethree\@exfirstword\@empty\@empty\@empty\@nil}%
    }%
    \def\@tempshort{#3}%
    \ifx\@tempshort\@empty
      \expandafter\xdef\csname @xshort@\@exlblkey\endcsname{\@exlblkey}%
    \else
      \expandafter\def\csname @xshort@\@exlblkey\endcsname{#3}%
    \fi
    \IfBooleanTF{#1}{\let\@extag\@empty}{%
      \edef\@extag{\noexpand\ \noexpand\textup{[\noexpand\texttt{e:\@exlblkey}]}}}%
    \begin{example}[#2\ifshowkeys\@extag\fi]%
    \label{e:\@exlblkey}%
    % record the long title so \nref{e:KEY} prints it, as for theorems
    \protected@write\@auxout{}{%
      \string\@thmtitle{e:\@exlblkey}{\unexpanded{#2}}}%
    \expanded{%
      \noexpand\newglossaryentry{e\theexcounter}{%
        name={\unexpanded{#2}%
          \IfBooleanF{#1}{ [\noexpand\texttt{e:\@exlblkey}]}},%
        sort={\theexample},%
        description={p.~\noexpand\pageref{e:\@exlblkey}},%
        user1={\theexample},%
        user2={Example},%
        user3={\thesection},%
        user4={\expandafter\unexpanded\expandafter{\@savedSecTitle}}%
      }%
    }%
    \protected@write\@auxout{}{\string\glsliveentry{e\theexcounter}}%
  \fi
}
\newcommand{\eeXa}{\end{example}}

% DEFINITION COMMAND -- same signature and behaviour as \beT below.
% Usage: \beD[long title][short title][key]
%   long title  = definition heading (plain text only)
%   short title = shown in \dref{word}; defaults to the key
%   key         = optional label key override; the label becomes d:KEY
% Without a key the label is d: plus the first three letters of the title's
% first word, so it is a mnemonic rather than the definition's number: a
% numeric label would silently change whenever a definition is inserted.
\NewDocumentCommand{\beD}{s O{} O{} o}{%
  \def\@temparg{#2}%
  \ifx\@temparg\@empty
    \begin{definition}%
  \else
    \let\@savedSecTitle\@currentsectionname%
    \refstepcounter{defcounter}%
    \IfValueTF{#4}{%
      \edef\@deflblkey{#4}%
    }{%
      \edef\@deffirstword{\@firstwordof#2 \relax}%
      \edef\@deflblkey{%
        \expandafter\@takethree\@deffirstword\@empty\@empty\@empty\@nil}%
    }%
    \def\@tempshort{#3}%
    \ifx\@tempshort\@empty
      \expandafter\xdef\csname @dshort@\@deflblkey\endcsname{\@deflblkey}%
    \else
      \expandafter\def\csname @dshort@\@deflblkey\endcsname{#3}%
    \fi
    \IfBooleanTF{#1}{\let\@deftag\@empty}{%
      \edef\@deftag{\noexpand\ \noexpand\textup{[\noexpand\texttt{d:\@deflblkey}]}}}%
    \begin{definition}[#2\ifshowkeys\@deftag\fi]%
    \label{d:\@deflblkey}%
    % record the long title so \nref{d:KEY} prints it, as for theorems
    \protected@write\@auxout{}{%
      \string\@thmtitle{d:\@deflblkey}{\unexpanded{#2}}}%
    \expanded{%
      \noexpand\newglossaryentry{d\thedefcounter}{%
        name={\unexpanded{#2}%
          \IfBooleanF{#1}{ [\noexpand\texttt{d:\@deflblkey}]}},%
        sort={\thedefinition},%
        description={p.~\noexpand\pageref{d:\@deflblkey}},%
        user1={\thedefinition},%
        user2={Definition},%
        user3={\thesection},%
        user4={\expandafter\unexpanded\expandafter{\@savedSecTitle}}%
      }%
    }%
    \protected@write\@auxout{}{\string\glsliveentry{d\thedefcounter}}%
  \fi
}
\newcommand{\eeD}{\end{definition}}

% \showkeysfalse in a driver hides the [t:KEY] mnemonics in result headings
% (labels and glossary entries are unchanged); used for submission builds.
\newif\ifshowkeys \showkeystrue
% \commonredtrue in a driver that inputs basicFiles/commonRed: tool.tex then
% cites its reduction results instead of restating them.
\newif\ifcommonred \commonredfalse
% TheorEM COMMAND
% Usage: \beT[long title][short title][key]
%   long title  = theorem heading (plain text only)
%   short title = shown in \tref{word}; may contain math; defaults to first word of long title
%   key         = optional label key override; label becomes t:KEY (avoids first-word collisions)
% Label is t:WORD where WORD = first word of long title (or KEY if given); use \tref{WORD} to reference.
\NewDocumentCommand{\beT}{s O{} O{} o}{%
  \def\@temparg{#2}%
  \ifx\@temparg\@empty
    \begin{theorem}%
  \else
    \let\@savedSecTitle\@currentsectionname%
    \refstepcounter{thmcounter}%
    % \@thmtag is the mnemonic shown in the heading; every theorem gets one.
    % With an explicit key the key itself is shown.  Without one the key
    % defaults to the title's first word, and showing that whole word would
    % merely repeat the title ("Dissipativity [Dissipativity]"), so only its
    % first three letters are shown: "Dissipativity [Dis]".  The label is
    % unaffected -- it still uses the full \@thmlblkey.
    \IfValueTF{#4}{%
      \edef\@thmlblkey{#4}%
    }{%
      \edef\@thmfirstword{\@firstwordof#2 \relax}%
      \edef\@thmlblkey{%
        \expandafter\@takethree\@thmfirstword\@empty\@empty\@empty\@nil}%
    }%
    \IfBooleanTF{#1}{\let\@thmtag\@empty}{%
      \edef\@thmtag{\noexpand\ \noexpand\textup{[\noexpand\texttt{t:\@thmlblkey}]}}}%
    \def\@tempshort{#3}%
    \ifx\@tempshort\@empty
      \expandafter\xdef\csname @tshort@\@thmlblkey\endcsname{\@thmlblkey}%
    \else
      \expandafter\def\csname @tshort@\@thmlblkey\endcsname{#3}%
    \fi
    \begin{theorem}[#2\ifshowkeys\@thmtag\fi]%
    \label{t:\@thmlblkey}%
    % record the long title for \nref; via the .aux so forward refs resolve
    \protected@write\@auxout{}{%
      \string\@thmtitle{t:\@thmlblkey}{\unexpanded{#2}}}%
    \expanded{%
      \noexpand\newglossaryentry{t\thethmcounter}{%
        name={\unexpanded{#2}%
          \IfBooleanF{#1}{ [\noexpand\texttt{t:\@thmlblkey}]}},%
        % sort/user1 must be the number the document actually prints, i.e. the
        % theorem environment's counter, not thmcounter.  thmcounter advances
        % only on \beT with a title, so a bare \beT (or \beT with an
        % empty title) makes the two drift and the glossary cites a wrong
        % number.  thmcounter is still fine as the entry key: it is unique.
        sort={\thetheorem},%
        description={p.~\noexpand\pageref{t:\@thmlblkey}},%
        user1={\thetheorem},%
        user2={Theorem},%
        user3={\thesection},%
        user4={\expandafter\unexpanded\expandafter{\@savedSecTitle}}%
      }%
    }%
    \protected@write\@auxout{}{\string\glsliveentry{t\thethmcounter}}%
  \fi
}
\newcommand{\eeT}{\end{theorem}}

% CROSS-PART RELATION PARAGRAPH
% Usage: \xpart{text}  -- place immediately BEFORE \beT, never inside it
%   (anything inside \beT...\eeT would be swallowed into the glossary entry).
% The paragraph is typeset only when Part I is actually present, i.e. in the
% combined book build; in the standalone builds (GranadaDSABNS/A.tex,
% AAHRah-Submission/AAHRah.tex) the label part:glv is undefined and the whole
% paragraph disappears rather than leaving ?? cross-references. Same
% graceful-degradation principle as \GASref in GASt.tex. Resolves on the
% second LaTeX pass, since r@part:glv comes from the .aux file.
% (Already inside the \makeatletter region opened above, so no \makeatletter
% pair here -- adding one would close that region early.)
\newcommand{\xpart}[1]{\@ifundefined{r@part:glv}{}{\paragraph{Relation with Parts~I--II.}#1}}

% \bookref{label}{precise text}{fallback text}: prints the precise cross-reference
% when the label resolves (combined book build), and the fallback text otherwise
% (standalone paper builds, where book-only appendix labels are absent). Same
% graceful-degradation role as \GASref in GASt.tex, but usable in any file since it
% lives in this preamble (loaded by every driver) -- unlike \GASref, which is defined in a
% content file and is therefore undefined in standalone builds that omit that file.
\newcommand{\bookref}[3]{\@ifundefined{r@#1}{#3}{#2}}

% \PVpart and \RBpart: the two parts that a partial build may omit.
% They print "Part~N" in any build that contains the part, and a naming phrase
% in a build that does not, so an excerpt shows no ?? .  Each already supplies
% the word "Part", so write "\PVpart develops ...", not "Part~\PVpart".
\newcommand{\PVpart}{\bookref{part:aahrah}%
  {Part~\ref{part:aahrah}}{the Perron--Volterra part of the full book}}
\newcommand{\RBpart}{\bookref{part:recbrowne}%
  {Part~\ref{part:recbrowne}}{the immuno-virology part of the full book}}

% OPEN PROBLEM COMMANDS
% ASSUMPTION COMMANDS (glossaried)
\newcommand{\beAs}[1][]{%
  \def\@temparg{#1}%
  \ifx\@temparg\@empty
    \begin{assumption}%
  \else
    \let\@savedSecTitle\@currentsectionname%
    \refstepcounter{ascounter}%
    \begin{assumption}[#1]%
    \label{as:\theascounter}%
    \expanded{%
      \noexpand\newglossaryentry{as\theascounter}{%
        name={\unexpanded{#1}},%
        sort={\theascounter},%
        description={p.~\noexpand\pageref{as:\theascounter}},%
        user1={\theascounter},%
        user2={Assumption},%
        user3={\thesection},%
        user4={\expandafter\unexpanded\expandafter{\@savedSecTitle}}%
      }%
    }%
    \protected@write\@auxout{}{\string\glsliveentry{as\theascounter}}%
  \fi
}
\newcommand{\eeAs}{\end{assumption}}
% CONJECTURE COMMANDS (glossaried)
\newcommand{\beCj}[1][]{%
  \def\@temparg{#1}%
  \ifx\@temparg\@empty
    \begin{conjecture}%
  \else
    \let\@savedSecTitle\@currentsectionname%
    \refstepcounter{cjcounter}%
    \begin{conjecture}[#1]%
    \label{cj:\thecjcounter}%
    \expanded{%
      \noexpand\newglossaryentry{cj\thecjcounter}{%
        name={\unexpanded{#1}},%
        sort={\thecjcounter},%
        description={p.~\noexpand\pageref{cj:\thecjcounter}},%
        user1={\thecjcounter},%
        user2={Conjecture},%
        user3={\thesection},%
        user4={\expandafter\unexpanded\expandafter{\@savedSecTitle}}%
      }%
    }%
    \protected@write\@auxout{}{\string\glsliveentry{cj\thecjcounter}}%
  \fi
}
\newcommand{\eeCj}{\end{conjecture}}
% LEMMA COMMANDS
% ---------------------------------------------------------------------
% UNIFORM RESULT WRAPPERS -- \beL, \beC, \beP, \beR share \beT's signature
%
%   \beX[long title][short title][key]  ...  \eeX
%
%   long title  = heading (math allowed)
%   short title = shown in the cross-reference; defaults to the key
%   key         = label key.  With a key the label is X:KEY and the heading
%                 carries the mnemonic [key], exactly as for \beT.  Without
%                 one the label falls back to the running counter, which is
%                 what every pre-existing call site in the book relies on.
%
% A bare \beX -- arguments commented out -- opens the plain environment and
% creates NO label and NO glossary entry.  That is the supported way to keep
% an individual result out of the catalogue.
% ---------------------------------------------------------------------
% #1 wrapper name  #2 environment  #3 label prefix  #4 counter  #5 type name
\newcommand{\@mkresult}[5]{%
  % A leading * suppresses the [key] mnemonic in both the heading and the
  % glossary name, while still making the label X:KEY and the entry:
  %     \beX*[long][short][key]
  \expandafter\NewDocumentCommand\csname be#1\endcsname{s O{} O{} o}{%
    \def\@temparg{##2}%
    \ifx\@temparg\@empty
      \begin{#2}%
    \else
      \let\@savedSecTitle\@currentsectionname%
      \refstepcounter{#4}%
      % With a key: mnemonic label X:KEY and the [key] tag in the heading.
      % Without one: fall back to the running counter, exactly as these
      % wrappers behaved before, so every legacy one-argument call site in
      % the book keeps its label.  The counter is also csname-safe, which a
      % title's first word is not -- titles beginning with math would
      % otherwise build \csname \Omega...\endcsname and fail.
      \IfValueTF{##4}{%
        \edef\@relbl{##4}%
        \IfBooleanTF{##1}{\let\@retag\@empty}{%
          \edef\@retag{\noexpand\ \noexpand\textup{[\noexpand\texttt{#3:##4}]}}}%
      }{%
        \edef\@relbl{\arabic{#4}}%
        \let\@retag\@empty
      }%
      \def\@tempshort{##3}%
      \ifx\@tempshort\@empty
        \expandafter\xdef\csname @sshort@#3\@relbl\endcsname{\@relbl}%
      \else
        \expandafter\def\csname @sshort@#3\@relbl\endcsname{##3}%
      \fi
      \begin{#2}[##2\ifshowkeys\@retag\fi]%
      \edef\@tmplbl{#3:\@relbl}\expandafter\label\expandafter{\@tmplbl}%
      \expanded{%
        \noexpand\newglossaryentry{#3\csname the#4\endcsname}{%
          name={\unexpanded{##2}%
            \IfBooleanF{##1}{\IfValueT{##4}{ [\noexpand\texttt{#3:##4}]}}},%
          sort={\csname the#2\endcsname},%
          description={p.~\noexpand\pageref{#3:\@relbl}},%
          user1={\csname the#2\endcsname},%
          user2={#5},%
          user3={\thesection},%
          user4={\expandafter\unexpanded\expandafter{\@savedSecTitle}}%
        }%
      }%
      \protected@write\@auxout{}{\string\glsliveentry{#3\csname the#4\endcsname}}%
      \protected@write\@auxout{}{\string\@thmtitle{\@tmplbl}{\unexpanded{##2}}}%
    \fi
  }%
  \expandafter\newcommand\csname ee#1\endcsname{\end{#2}}%
}
\@mkresult{L}{lemma}{l}{lemcounter}{Lemma}
% REMARK / COROLLARY / PROPOSITION -- same signature as \beL above
\@mkresult{R}{remark}{r}{remcounter}{Remark}
\@mkresult{C}{corollary}{c}{corcounter}{Corollary}
\@mkresult{P}{proposition}{p}{propcounter}{Proposition}
% OPEN-PROBLEM COMMAND -- generated with the same signature as \beT:
%   \beO[long title][short title][key]   ...   \eeO
% label o:KEY, or o:<counter> when no key is given; a bare \beO makes
% neither a label nor a glossary entry.
\@mkresult{O}{open}{o}{opcounter}{Problem}


% Section-level title tracker: updated only by \section, not \subsection.
% \@currentlabelname is overwritten by every \refstepcounter (including
% \subsection), so we save it separately here.
\newcommand{\@currentsectionname}{}
\AtBeginDocument{%
  \let\@epidcrn@orig@refstepcounter\refstepcounter
  \renewcommand{\refstepcounter}[1]{%
    \@epidcrn@orig@refstepcounter{#1}%
    \ifstrequal{#1}{section}{\global\let\@currentsectionname\@currentlabelname}{}%
  }%
}

% TABLE COMMANDS (glossaried) -- consistent naming with \beT, \beD, etc.
\newcommand{\beTa}[1][]{%
  \def\@temparg{#1}%
  \ifx\@temparg\@empty
    \begin{table}[htbp]%
  \else
    \let\@savedSecTitle\@currentsectionname%
    \refstepcounter{tabcounter}%
    \begin{table}[htbp]%
    \label{A:\thetabcounter}%
    \expanded{%
      \noexpand\newglossaryentry{A\thetabcounter}{%
        name={\unexpanded{#1}},%
        sort={\thetabcounter},%
        description={p.~\noexpand\pageref{A:\thetabcounter}},%
        user1={\thetabcounter},%
        user2={Table},%
        user3={\thesection},%
        user4={\expandafter\unexpanded\expandafter{\@savedSecTitle}}%
      }%
    }%
    \protected@write\@auxout{}{\string\glsliveentry{A\thetabcounter}}%
  \fi
}
\newcommand{\eeTa}{\end{table}}

% FIGURE COMMANDS (glossaried) -- consistent with \beT, \beR, \beTa, etc.
% Usage:
%   \beF[Glossary title]...\caption{Visible caption}\eeF  -> glossaried figure
%   \beF...\caption{Visible caption}\eeF                  -> plain [H] figure, no glossary
\newcommand{\beF}[1][]{%
  \def\@temparg{#1}%
  \ifx\@temparg\@empty
    \begin{figure}[H]%
  \else
    \begin{figure}[H]%
    \let\@savedSecTitle\@currentsectionname%
    \refstepcounter{figcounter}%
    \label{fig:\thefigcounter}%
    \expanded{%
      \noexpand\newglossaryentry{Fig\thefigcounter}{%
        name={\unexpanded{#1}},%
        sort={\thefigcounter},%
        description={p.~\noexpand\pageref{fig:\thefigcounter}},%
        user1={\thefigcounter},%
        user2={Figure},%
        user3={\thesection},%
        user4={\expandafter\unexpanded\expandafter{\@savedSecTitle}}%
      }%
    }%
    \protected@write\@auxout{}{\string\glsliveentry{Fig\thefigcounter}}%
  \fi
}
\newcommand{\eeF}{\end{figure}}
\makeatother
% LIVE-ENTRY REGISTER
% \glsaddall prints every entry the glossaries package currently knows, and
% that set includes entries restored from \jobname.glsdefs.  The package
% rewrites .glsdefs from all *currently defined* entries, so anything it
% restored is written straight back out: a deleted section kept reappearing in
% the catalogue for ever, and the .glsdefs had to be erased by hand.  The
% wrappers therefore record every entry they actually create in the .aux, and
% \printboth adds only those.  The .aux is rebuilt from scratch on each run, so
% this is self-healing: a removed result is gone after one further run.
\makeatletter
\newcommand{\@glslive}{}
\newcommand{\@glsaddlive}[1]{\ifglsentryexists{#1}{\glsadd{#1}}{}}
\newcommand{\glsliveentry}[1]{\g@addto@macro\@glslive{\@glsaddlive{#1}}}
\makeatother
% PRINT COMMAND
% No \clearpage here -- the catalogue continues on the current page.  This has
% been re-added by mistake several times; leave it out.
\makeatletter
\newcommand{\printboth}{%
  \@glslive   % add exactly the entries this document created last run
  \printglossary[style=unifiedstyle, title={0 \hspace{.22em} Definitions, Theorems, Examples and Open Problems}]%
}
\makeatother
% Named theorem reference: \tref{word} -> Thm~N(word-or-short)
% Shows stored short title if available, otherwise shows the key itself.
\newcommand{\tref}[1]{%
  Thm~\ref{t:#1}(\ifcsname @tshort@#1\endcsname\csname @tshort@#1\endcsname\else#1\fi)%
}

% Full-title theorem reference: \nref{t:KEY} prints the number and then the
% theorem's long title, so one writes "Theorem~\nref{t:blocks}" and gets
% "Theorem 3 ``Block factorization and siphons''".  Takes the complete label, like
% \ref, not the bare key like \tref.
%
% \beT writes each title to the .aux (see \@thmtitle below), so \nref resolves
% forward references from the second LaTeX pass, exactly as \ref does.  If the
% title is not yet known -- first pass, or a label that is not a \beT theorem
% -- it degrades to the number rather than producing ??.
% These two use @ in their names, and this point in the file is outside the
% \makeatletter region opened earlier (it closes at \makeatother above), so
% the region has to be reopened here.
\makeatletter
\newcommand{\@thmtitle}[2]{\global\@namedef{@tlong@#1}{#2}}
\newcommand{\nref}[1]{%
  \ref{#1}%
  \ifcsname @tlong@#1\endcsname
    \ ``\csname @tlong@#1\endcsname''%
  \fi
}
% \nm{t:KEY}: the stored long title of any \beT/\beL/\beP/\beR/\beC/\beD/\beE result,
% verbatim and without quotes, so item names follow the macro's first argument.
\newcommand{\nm}[1]{%
  \ifcsname @tlong@#1\endcsname
    \csname @tlong@#1\endcsname%
  \else
    \ref{#1}%
  \fi
}
\makeatother




%% ========================================
%% Other THEOREM ENVIRONMENTS
%% ========================================
\makeatletter
\@ifundefined{lemma}      {\newtheorem{lemma}      {Lemma}}      {}
\@ifundefined{proposition}{\newtheorem{proposition}{Proposition}} {}
\@ifundefined{corollary}  {\newtheorem{corollary}  {Corollary}}  {}
\@ifundefined{conjecture} {\newtheorem{conjecture} {Conjecture}}  {}
\@ifundefined{remark}     {\newtheorem{remark}     {Remark}}     {}
\@ifundefined{question}   {\newtheorem{question}   {Question}}   {}
\@ifundefined{assumption} {\newtheorem{assumption} {Assumption}} {}
\@ifundefined{hypothesis} {\newtheorem{hypothesis} {Hypothesis}} {}
\@ifundefined{fact}       {\newtheorem{fact}       {Fact}}       {}
\makeatother

% \citer: cite the relay-graph threshold theorem of AAHrel. If a document has
% \input the chapter/file containing t:relayRF (so the label resolves), point
% straight at it; otherwise fall back to a bare citation of AAHrel. Shared here
% (not per-paper) so it resolves consistently across every book chapter.
% \citeR: same idea, for the first theorem of AAHRah (t:Pstar).
\makeatletter
\providecommand{\citer}{%
  \@ifundefined{r@t:relayRF}%
    {\cite{AAHrel}}%
    {Theorem~\ref{t:relayRF} of \cite{AAHrel}}%
}
\providecommand{\citeR}{%
  \@ifundefined{r@t:Pstar}%
    {\cite{AAHRah}}%
    {Theorem~\ref{t:Pstar} of \cite{AAHRah}}%
}
\makeatother


% Abbreviations
\def\beQ{\begin{question}}\def\eeQ{\end{question}}
\def\beA{\begin{assumption}}\def\eeA{\end{assumption}}


%% ========================================

%% ========================================
%% COLORS & CUSTOM COMMANDS
%% ========================================
\definecolor{funccolor}{RGB}{25,25,112}
\definecolor{desccolor}{RGB}{64,64,64}
\newcommand{\func}[1]{\texttt{\textcolor{funccolor}{#1}}}
\def\eq{=}
\def\mm{-}

%% ========================================
%% MATHCAL SHORTHANDS
%% ========================================
% Project-wide convention: \mX for a capital X is always \mathcal X.
% Defined for the whole alphabet so the convention never fails to expand.
\newcommand{\mA}{\mathcal A}\newcommand{\mB}{\mathcal B}
\newcommand{\mC}{\mathcal C}\newcommand{\mD}{\mathcal D}
\newcommand{\mE}{\mathcal E}\newcommand{\mF}{\mathcal F}
\newcommand{\mG}{\mathcal G}\newcommand{\mH}{\mathcal H}
\newcommand{\mI}{\mathcal I}\newcommand{\mJ}{\mathcal J}
\newcommand{\mK}{\mathcal K}\newcommand{\mL}{\mathcal L}
\newcommand{\mM}{\mathcal M}\newcommand{\mN}{\mathcal N}
\newcommand{\mO}{\mathcal O}\newcommand{\mP}{\mathcal P}
\newcommand{\mQ}{\mathcal Q}\newcommand{\mR}{\mathcal R}
\newcommand{\mS}{\mathcal S}\newcommand{\mT}{\mathcal T}
\newcommand{\mU}{\mathcal U}\newcommand{\mV}{\mathcal V}
\newcommand{\mW}{\mathcal W}\newcommand{\mX}{\mathcal X}
\newcommand{\mY}{\mathcal Y}\newcommand{\mZ}{\mathcal Z}

%% ========================================
%% BASIC MATH ENVIRONMENTS
%% ========================================
\def\bep{\begin{pmatrix}}\def\eep{\end{pmatrix}}
\def\bev{\begin{vmatrix}}\def\eev{\end{vmatrix}}
\def\bea{\begin{eqnarray*}}\def\eea{\end{eqnarray*}}
\def\bc{\begin{cases}}\def\ec{\end{cases}}
\def\BEN{\begin{enumerate}}\def\EEN{\end{enumerate}}
\def\BI{\begin{itemize}}\def\EI{\end{itemize}}

\newcommand{\be}[1]{\begin{equation}\label{#1}}
\newcommand{\ee}{\end{equation}}
\newcommand{\beq}{\begin{eqnarray}}
\def\eeq{\end{eqnarray}}
\def\eqr{\eqref}\def\fr{\frac}\def\lbl{\label}
\def\Lra{\Longrightarrow}\def\non{\nonumber}\def\cNGM{\cite{Diek,Van,Van08}}
\def\Eq{\Leftrightarrow}\def\b{\beta}\def\uph{\red{UP TO HERE}}
%% ========================================
%% PROBABILITY & STATISTICS
%% ========================================
\renewcommand{\P}{\ensuremath{\mathbb{P}}}  % Note: overrides paragraph symbol
\newcommand{\E}{\ensuremath{\mathbb{E}}}
\newcommand{\ind}[1]{\ensuremath{\mathbbm{1}_{\left\{#1\right\}}}}
\DeclareMathOperator{\Cov}{Cov}
\DeclareMathOperator{\TV}{TV}
\newcommand{\rank}{\operatorname{rank}}

%% ========================================
%% CRN OPERATORS (organizations/ERCs/buffering/siphons framework)
%% ========================================
\def\clos{\mathrm{clos}}
\def\Supply{\mathrm{Supply}}
\def\req{\mathrm{req}}
\def\React{\mathrm{React}}
\def\Prod{\mathrm{Prod}}
\def\Gen{\mathrm{Gen}}
\def\In{\mathrm{In}}
\def\Out{\mathrm{Out}}
\def\Act{\mathrm{Act}}
\def\Min{\mathrm{Min}}

%% ========================================
%% LINE BREAKING
%% ========================================
% Last-resort stretch, used only on paragraphs TeX cannot otherwise break:
% a long \texttt{} token, a comma-separated inline math list (commas are not
% breakpoints in math), or a run of \cite/\ref boxes.  Without it those set as
% overfull \hbox and print into the margin; with it they set slightly loose,
% which is the lesser evil.  Does not affect paragraphs that already break.
\emergencystretch=3em

%% ========================================
%% NUMBER SETS
%% ========================================
\newcommand{\R}{\mathbb{R}}  % Real numbers
\newcommand{\C}{\mathbb{C}}  % Complex numbers
\newcommand{\N}{\mathbb{N}}
\newcommand{\Z}{\mathbb{Z}}
% \providecommand, not \newcommand: some papers define these in their own
% preamble after % Shared preamble for the GLVbook papers.  This is the single canonical copy;
% a driver one level below GLVbook reaches it as
%
%     \input{../basicFiles/preamble}
%
% What stays in the driver, because it differs from paper to paper:
% \documentclass, font choice, \title, \author, \date, and any macro genuinely
% local to one paper.  Everything else belongs here.  The file is
% self-contained: it does not read any other file.

\usepackage[utf8]{inputenc}
\usepackage[novbox]{pdfsync}   % required: SyncTeX alone does not work here

% Page geometry, set once here for every paper that loads this file.  A driver
% that needs different margins says \geometry{...}; a second
% \usepackage[...]{geometry} raises an option clash and must not be used.
\makeatletter
\@ifundefined{AAdefNnogeometry}{\usepackage[a4paper,margin=2cm]{geometry}}{}
\makeatother

% mathematics
\usepackage{amsmath,amssymb,amsthm,mathtools,bm,mathrsfs,bbm}
\usepackage{centernot}

% floats, tables, lists
\usepackage{graphicx,float}
\makeatletter
\@ifpackageloaded{subfig}{}{\usepackage{subcaption}}
\makeatother
\usepackage{booktabs,longtable,threeparttable,array,multirow}
\usepackage{enumitem}

% colour, code, algorithms
\usepackage{xcolor,listings,tcolorbox}
\makeatletter
\@ifundefined{c@algorithm}{\usepackage{algorithm}}{}
\makeatother
\usepackage{algorithmicx,algpseudocode}

% diagrams
\usepackage{tikz,tikz-cd,pgfplots}
\usetikzlibrary{graphs,decorations.pathmorphing,decorations.markings,
arrows.meta,automata,positioning}
\tikzcdset{every label/.append style = {font = \small}}
\pgfplotsset{compat=1.18}

% chemistry
\usepackage[version=4]{mhchem}

\usepackage{microtype}
\usepackage{lineno}
\usepackage{nomencl}
\makenomenclature

% hyperref first, then bookmark, then glossaries: glossaries must follow
% hyperref or the entry links are built against the wrong anchors.
\makeatletter
\@ifundefined{hyperrefpackage}{%
  \@ifundefined{AAdefNnohyperref}{\usepackage{hyperref}}{}%
}{}
\@ifundefined{hyperrefpackage}{\usepackage{bookmark}}{}
\makeatother
\usepackage[toc,sort=use]{glossaries}
\makeglossaries

% Book/standalone switch.  Material that the book already covers in its
% prelude, but that a standalone paper still has to state, is wrapped as
%     \ifbookbuild\else ... \fi
% Drivers default to standalone; bookM.tex says \bookbuildtrue.
\newif\ifbookbuild
\bookbuildfalse

% Part switches for the gradostat paper.  Both default to true; a driver that
% wants the short submission says \madecfalse and/or \slowfastfalse, and every
% cross-reference into the dropped part disappears with it.
%   \ifmadec     Part III, the Lagasquie--Madec model and its break-even geometry
%   \ifslowfast  Part IV, everything carrying the small parameter \eps
\newif\ifmadec     \madectrue
\newif\ifslowfast  \slowfasttrue

% Part switch for the Barrat/simplicial-contagion project.  Defaults to
% true (the hypernetwork-synchrony part is included); the journal-paper
% driver says \hnetpartfalse to drop it and every cross-reference into it.
\newif\ifhnetpart  \hnetparttrue

%% ========================================
%% PACKAGES
%% ========================================
\RequirePackage{etoolbox}
\usepackage{makecell}
\usepackage[novbox]{pdfsync}% obsolete; SyncTeX is built into modern engines
\usepackage{amssymb, amsfonts,mathtools,bbm,eurosym,graphicx,latexsym, amsmath,amsthm,times,nccmath,url}
\usepackage{listings,verbatim,cases,tabularx}
% hyperref is deliberately NOT in the line above.  Journal classes often
% load it themselves, with a long option list, from a begindocument hook
% (mdpi.cls does: \DeclareHookRule{begindocument}{hyperref}{before}{mdpi}).
% Loading it here first, without options, makes that later load raise
% "Option clash for package hyperref".  So load it only if nobody else
% has, or will:
%   - \@ifpackageloaded  catches a class that loaded it already;
%   - \hyperrefpackage    is mdpi.cls's own deferred loader;
%   - \AAdefNnohyperref   lets any other driver opt out by defining it
%                         before % Shared preamble for the GLVbook papers.  This is the single canonical copy;
% a driver one level below GLVbook reaches it as
%
%     \input{../basicFiles/preamble}
%
% What stays in the driver, because it differs from paper to paper:
% \documentclass, font choice, \title, \author, \date, and any macro genuinely
% local to one paper.  Everything else belongs here.  The file is
% self-contained: it does not read any other file.

\usepackage[utf8]{inputenc}
\usepackage[novbox]{pdfsync}   % required: SyncTeX alone does not work here

% Page geometry, set once here for every paper that loads this file.  A driver
% that needs different margins says \geometry{...}; a second
% \usepackage[...]{geometry} raises an option clash and must not be used.
\makeatletter
\@ifundefined{AAdefNnogeometry}{\usepackage[a4paper,margin=2cm]{geometry}}{}
\makeatother

% mathematics
\usepackage{amsmath,amssymb,amsthm,mathtools,bm,mathrsfs,bbm}
\usepackage{centernot}

% floats, tables, lists
\usepackage{graphicx,float}
\makeatletter
\@ifpackageloaded{subfig}{}{\usepackage{subcaption}}
\makeatother
\usepackage{booktabs,longtable,threeparttable,array,multirow}
\usepackage{enumitem}

% colour, code, algorithms
\usepackage{xcolor,listings,tcolorbox}
\makeatletter
\@ifundefined{c@algorithm}{\usepackage{algorithm}}{}
\makeatother
\usepackage{algorithmicx,algpseudocode}

% diagrams
\usepackage{tikz,tikz-cd,pgfplots}
\usetikzlibrary{graphs,decorations.pathmorphing,decorations.markings,
arrows.meta,automata,positioning}
\tikzcdset{every label/.append style = {font = \small}}
\pgfplotsset{compat=1.18}

% chemistry
\usepackage[version=4]{mhchem}

\usepackage{microtype}
\usepackage{lineno}
\usepackage{nomencl}
\makenomenclature

% hyperref first, then bookmark, then glossaries: glossaries must follow
% hyperref or the entry links are built against the wrong anchors.
\makeatletter
\@ifundefined{hyperrefpackage}{%
  \@ifundefined{AAdefNnohyperref}{\usepackage{hyperref}}{}%
}{}
\@ifundefined{hyperrefpackage}{\usepackage{bookmark}}{}
\makeatother
\usepackage[toc,sort=use]{glossaries}
\makeglossaries

% Book/standalone switch.  Material that the book already covers in its
% prelude, but that a standalone paper still has to state, is wrapped as
%     \ifbookbuild\else ... \fi
% Drivers default to standalone; bookM.tex says \bookbuildtrue.
\newif\ifbookbuild
\bookbuildfalse

% Part switches for the gradostat paper.  Both default to true; a driver that
% wants the short submission says \madecfalse and/or \slowfastfalse, and every
% cross-reference into the dropped part disappears with it.
%   \ifmadec     Part III, the Lagasquie--Madec model and its break-even geometry
%   \ifslowfast  Part IV, everything carrying the small parameter \eps
\newif\ifmadec     \madectrue
\newif\ifslowfast  \slowfasttrue

% Part switch for the Barrat/simplicial-contagion project.  Defaults to
% true (the hypernetwork-synchrony part is included); the journal-paper
% driver says \hnetpartfalse to drop it and every cross-reference into it.
\newif\ifhnetpart  \hnetparttrue

%% ========================================
%% PACKAGES
%% ========================================
\RequirePackage{etoolbox}
\usepackage{makecell}
\usepackage[novbox]{pdfsync}% obsolete; SyncTeX is built into modern engines
\usepackage{amssymb, amsfonts,mathtools,bbm,eurosym,graphicx,latexsym, amsmath,amsthm,times,nccmath,url}
\usepackage{listings,verbatim,cases,tabularx}
% hyperref is deliberately NOT in the line above.  Journal classes often
% load it themselves, with a long option list, from a begindocument hook
% (mdpi.cls does: \DeclareHookRule{begindocument}{hyperref}{before}{mdpi}).
% Loading it here first, without options, makes that later load raise
% "Option clash for package hyperref".  So load it only if nobody else
% has, or will:
%   - \@ifpackageloaded  catches a class that loaded it already;
%   - \hyperrefpackage    is mdpi.cls's own deferred loader;
%   - \AAdefNnohyperref   lets any other driver opt out by defining it
%                         before \input{preamble}.
\makeatletter
\@ifpackageloaded{hyperref}{}{%
  \@ifundefined{hyperrefpackage}{%
    \@ifundefined{AAdefNnohyperref}{\usepackage{hyperref}}{}%
  }{}%
}
\makeatother
% linktoc=all makes a table-of-contents page number a link as well as the
% entry text; hyperref's default, linktoc=section, links only the text.
% Deferred to begin-document and guarded: by the block above hyperref may
% be loaded later (mdpi's own hook) or not at all (\AAdefNnohyperref), and
% \hypersetup must exist before it is called.  A driver that wants the
% default back can say \hypersetup{linktoc=section} after \input{preamble}.
\makeatletter
\AtBeginDocument{\@ifundefined{hypersetup}{}{\hypersetup{linktoc=all}}}
\makeatother
\usepackage{graphicx,float,fancybox}
\usepackage{epsfig,epstopdf,tikz}
\usepackage[sans]{dsfont}
\usetikzlibrary{positioning,arrows,arrows.meta,calc}
% Configure listings for Mathematica code
\lstdefinestyle{mathematica}{
 language=Mathematica,
 basicstyle=\small\ttfamily,
 backgroundcolor=\color{gray!10},
 frame=single,
 breaklines=true,
 commentstyle=\color{green!60!black},
 keywordstyle=\color{blue},
 stringstyle=\color{red},
 showstringspaces=false,
 tabsize=2
}
\usepackage{tcolorbox}
% Colored boxes for key results
\newtcolorbox{keyresult}[1][]{
 colback=blue!5!white,
 colframe=blue!75!black,
 title=#1,
 fonttitle=\bfseries
}
\newtcolorbox{examplebox}[1][]{
 colback=green!5!white,
 colframe=green!75!black,
 title=#1,
 fonttitle=\bfseries
}


\makeatletter
\@ifundefined{theorem}   {\newtheorem{theorem}   {Theorem}}   {}
\@ifundefined{definition}{\newtheorem{definition}{Definition}} {}
\@ifundefined{open}      {\newtheorem{open}      {Problem}}   {}
\@ifundefined{example}   {\newtheorem{example}   {Example}}   {}
\makeatother

% COUNTER DEFINITIONS
\newcounter{tabcounter}
\newcounter{defcounter}
\newcounter{opcounter}
\newcounter{thmcounter}
\newcounter{excounter}
\newcounter{ascounter}
\newcounter{cjcounter}
\newcounter{lemcounter}
\newcounter{remcounter}
\newcounter{corcounter}
\newcounter{propcounter}
\newcounter{figcounter}

% UNIFIED GLOSSARY STYLE
\newglossarystyle{unifiedstyle}{%
  \setglossarystyle{list}%
  \renewcommand*{\glossarypreamble}{\gdef\glslastsec{}}%
  \renewcommand{\glossentry}[2]{%
    \edef\glscursec{\glsentryuseriii{##1}}%
    \ifx\glscursec\glslastsec\else
      \item[]\textbf{\S\,\glscursec\enspace\glsentryuseriv{##1}}%
      \global\let\glslastsec\glscursec
    \fi
    \item[\glsentryuserii{##1} \glsentryuseri{##1}:] ``\glossentryname{##1}'' \glsentrydesc{##1}%
  }%
}

\makeatletter
% First-word extractor (used by \beT, \beD and \beXa for the default key)
% \@firstwordof hello world \relax  ->  hello
\def\@firstwordof#1 #2\relax{#1}
% First three letters, used for the heading mnemonic of a theorem given no
% explicit key.  Padding with \@empty keeps it safe for first words shorter
% than three letters: \@empty vanishes under \edef, whereas \relax would
% survive into the label and is illegal inside \csname (it broke every
% two-letter title, e.g. "ME ...", with "Missing \endcsname inserted").
\def\@takethree#1#2#3#4\@nil{#1#2#3}

% EXAMPLE COMMAND -- same signature and behaviour as \beT and \beD.
% Usage: \beXa*[long title][short title][key]
%   star  = suppresses the printed "[key]" tag after the title (as for \beT*
%   and \beD*); the label and glossary entry are unaffected.
%   key = optional label key override; the label becomes e:KEY, and without a
%   key it is e: plus the first three letters of the title's first word.  A
%   mnemonic label lets an example be cited as "Example~\ref{e:KEY}" instead
%   of by section number, which is what the citing prose needs.
\NewDocumentCommand{\beXa}{s O{} O{} o}{%
  \def\@temparg{#2}%
  \ifx\@temparg\@empty
    \begin{example}%
  \else
    \let\@savedSecTitle\@currentsectionname%
    \refstepcounter{excounter}%
    \IfValueTF{#4}{%
      \edef\@exlblkey{#4}%
    }{%
      \edef\@exfirstword{\@firstwordof#2 \relax}%
      \edef\@exlblkey{%
        \expandafter\@takethree\@exfirstword\@empty\@empty\@empty\@nil}%
    }%
    \def\@tempshort{#3}%
    \ifx\@tempshort\@empty
      \expandafter\xdef\csname @xshort@\@exlblkey\endcsname{\@exlblkey}%
    \else
      \expandafter\def\csname @xshort@\@exlblkey\endcsname{#3}%
    \fi
    \IfBooleanTF{#1}{\let\@extag\@empty}{%
      \edef\@extag{\noexpand\ \noexpand\textup{[\noexpand\texttt{e:\@exlblkey}]}}}%
    \begin{example}[#2\ifshowkeys\@extag\fi]%
    \label{e:\@exlblkey}%
    % record the long title so \nref{e:KEY} prints it, as for theorems
    \protected@write\@auxout{}{%
      \string\@thmtitle{e:\@exlblkey}{\unexpanded{#2}}}%
    \expanded{%
      \noexpand\newglossaryentry{e\theexcounter}{%
        name={\unexpanded{#2}%
          \IfBooleanF{#1}{ [\noexpand\texttt{e:\@exlblkey}]}},%
        sort={\theexample},%
        description={p.~\noexpand\pageref{e:\@exlblkey}},%
        user1={\theexample},%
        user2={Example},%
        user3={\thesection},%
        user4={\expandafter\unexpanded\expandafter{\@savedSecTitle}}%
      }%
    }%
    \protected@write\@auxout{}{\string\glsliveentry{e\theexcounter}}%
  \fi
}
\newcommand{\eeXa}{\end{example}}

% DEFINITION COMMAND -- same signature and behaviour as \beT below.
% Usage: \beD[long title][short title][key]
%   long title  = definition heading (plain text only)
%   short title = shown in \dref{word}; defaults to the key
%   key         = optional label key override; the label becomes d:KEY
% Without a key the label is d: plus the first three letters of the title's
% first word, so it is a mnemonic rather than the definition's number: a
% numeric label would silently change whenever a definition is inserted.
\NewDocumentCommand{\beD}{s O{} O{} o}{%
  \def\@temparg{#2}%
  \ifx\@temparg\@empty
    \begin{definition}%
  \else
    \let\@savedSecTitle\@currentsectionname%
    \refstepcounter{defcounter}%
    \IfValueTF{#4}{%
      \edef\@deflblkey{#4}%
    }{%
      \edef\@deffirstword{\@firstwordof#2 \relax}%
      \edef\@deflblkey{%
        \expandafter\@takethree\@deffirstword\@empty\@empty\@empty\@nil}%
    }%
    \def\@tempshort{#3}%
    \ifx\@tempshort\@empty
      \expandafter\xdef\csname @dshort@\@deflblkey\endcsname{\@deflblkey}%
    \else
      \expandafter\def\csname @dshort@\@deflblkey\endcsname{#3}%
    \fi
    \IfBooleanTF{#1}{\let\@deftag\@empty}{%
      \edef\@deftag{\noexpand\ \noexpand\textup{[\noexpand\texttt{d:\@deflblkey}]}}}%
    \begin{definition}[#2\ifshowkeys\@deftag\fi]%
    \label{d:\@deflblkey}%
    % record the long title so \nref{d:KEY} prints it, as for theorems
    \protected@write\@auxout{}{%
      \string\@thmtitle{d:\@deflblkey}{\unexpanded{#2}}}%
    \expanded{%
      \noexpand\newglossaryentry{d\thedefcounter}{%
        name={\unexpanded{#2}%
          \IfBooleanF{#1}{ [\noexpand\texttt{d:\@deflblkey}]}},%
        sort={\thedefinition},%
        description={p.~\noexpand\pageref{d:\@deflblkey}},%
        user1={\thedefinition},%
        user2={Definition},%
        user3={\thesection},%
        user4={\expandafter\unexpanded\expandafter{\@savedSecTitle}}%
      }%
    }%
    \protected@write\@auxout{}{\string\glsliveentry{d\thedefcounter}}%
  \fi
}
\newcommand{\eeD}{\end{definition}}

% \showkeysfalse in a driver hides the [t:KEY] mnemonics in result headings
% (labels and glossary entries are unchanged); used for submission builds.
\newif\ifshowkeys \showkeystrue
% \commonredtrue in a driver that inputs basicFiles/commonRed: tool.tex then
% cites its reduction results instead of restating them.
\newif\ifcommonred \commonredfalse
% TheorEM COMMAND
% Usage: \beT[long title][short title][key]
%   long title  = theorem heading (plain text only)
%   short title = shown in \tref{word}; may contain math; defaults to first word of long title
%   key         = optional label key override; label becomes t:KEY (avoids first-word collisions)
% Label is t:WORD where WORD = first word of long title (or KEY if given); use \tref{WORD} to reference.
\NewDocumentCommand{\beT}{s O{} O{} o}{%
  \def\@temparg{#2}%
  \ifx\@temparg\@empty
    \begin{theorem}%
  \else
    \let\@savedSecTitle\@currentsectionname%
    \refstepcounter{thmcounter}%
    % \@thmtag is the mnemonic shown in the heading; every theorem gets one.
    % With an explicit key the key itself is shown.  Without one the key
    % defaults to the title's first word, and showing that whole word would
    % merely repeat the title ("Dissipativity [Dissipativity]"), so only its
    % first three letters are shown: "Dissipativity [Dis]".  The label is
    % unaffected -- it still uses the full \@thmlblkey.
    \IfValueTF{#4}{%
      \edef\@thmlblkey{#4}%
    }{%
      \edef\@thmfirstword{\@firstwordof#2 \relax}%
      \edef\@thmlblkey{%
        \expandafter\@takethree\@thmfirstword\@empty\@empty\@empty\@nil}%
    }%
    \IfBooleanTF{#1}{\let\@thmtag\@empty}{%
      \edef\@thmtag{\noexpand\ \noexpand\textup{[\noexpand\texttt{t:\@thmlblkey}]}}}%
    \def\@tempshort{#3}%
    \ifx\@tempshort\@empty
      \expandafter\xdef\csname @tshort@\@thmlblkey\endcsname{\@thmlblkey}%
    \else
      \expandafter\def\csname @tshort@\@thmlblkey\endcsname{#3}%
    \fi
    \begin{theorem}[#2\ifshowkeys\@thmtag\fi]%
    \label{t:\@thmlblkey}%
    % record the long title for \nref; via the .aux so forward refs resolve
    \protected@write\@auxout{}{%
      \string\@thmtitle{t:\@thmlblkey}{\unexpanded{#2}}}%
    \expanded{%
      \noexpand\newglossaryentry{t\thethmcounter}{%
        name={\unexpanded{#2}%
          \IfBooleanF{#1}{ [\noexpand\texttt{t:\@thmlblkey}]}},%
        % sort/user1 must be the number the document actually prints, i.e. the
        % theorem environment's counter, not thmcounter.  thmcounter advances
        % only on \beT with a title, so a bare \beT (or \beT with an
        % empty title) makes the two drift and the glossary cites a wrong
        % number.  thmcounter is still fine as the entry key: it is unique.
        sort={\thetheorem},%
        description={p.~\noexpand\pageref{t:\@thmlblkey}},%
        user1={\thetheorem},%
        user2={Theorem},%
        user3={\thesection},%
        user4={\expandafter\unexpanded\expandafter{\@savedSecTitle}}%
      }%
    }%
    \protected@write\@auxout{}{\string\glsliveentry{t\thethmcounter}}%
  \fi
}
\newcommand{\eeT}{\end{theorem}}

% CROSS-PART RELATION PARAGRAPH
% Usage: \xpart{text}  -- place immediately BEFORE \beT, never inside it
%   (anything inside \beT...\eeT would be swallowed into the glossary entry).
% The paragraph is typeset only when Part I is actually present, i.e. in the
% combined book build; in the standalone builds (GranadaDSABNS/A.tex,
% AAHRah-Submission/AAHRah.tex) the label part:glv is undefined and the whole
% paragraph disappears rather than leaving ?? cross-references. Same
% graceful-degradation principle as \GASref in GASt.tex. Resolves on the
% second LaTeX pass, since r@part:glv comes from the .aux file.
% (Already inside the \makeatletter region opened above, so no \makeatletter
% pair here -- adding one would close that region early.)
\newcommand{\xpart}[1]{\@ifundefined{r@part:glv}{}{\paragraph{Relation with Parts~I--II.}#1}}

% \bookref{label}{precise text}{fallback text}: prints the precise cross-reference
% when the label resolves (combined book build), and the fallback text otherwise
% (standalone paper builds, where book-only appendix labels are absent). Same
% graceful-degradation role as \GASref in GASt.tex, but usable in any file since it
% lives in this preamble (loaded by every driver) -- unlike \GASref, which is defined in a
% content file and is therefore undefined in standalone builds that omit that file.
\newcommand{\bookref}[3]{\@ifundefined{r@#1}{#3}{#2}}

% \PVpart and \RBpart: the two parts that a partial build may omit.
% They print "Part~N" in any build that contains the part, and a naming phrase
% in a build that does not, so an excerpt shows no ?? .  Each already supplies
% the word "Part", so write "\PVpart develops ...", not "Part~\PVpart".
\newcommand{\PVpart}{\bookref{part:aahrah}%
  {Part~\ref{part:aahrah}}{the Perron--Volterra part of the full book}}
\newcommand{\RBpart}{\bookref{part:recbrowne}%
  {Part~\ref{part:recbrowne}}{the immuno-virology part of the full book}}

% OPEN PROBLEM COMMANDS
% ASSUMPTION COMMANDS (glossaried)
\newcommand{\beAs}[1][]{%
  \def\@temparg{#1}%
  \ifx\@temparg\@empty
    \begin{assumption}%
  \else
    \let\@savedSecTitle\@currentsectionname%
    \refstepcounter{ascounter}%
    \begin{assumption}[#1]%
    \label{as:\theascounter}%
    \expanded{%
      \noexpand\newglossaryentry{as\theascounter}{%
        name={\unexpanded{#1}},%
        sort={\theascounter},%
        description={p.~\noexpand\pageref{as:\theascounter}},%
        user1={\theascounter},%
        user2={Assumption},%
        user3={\thesection},%
        user4={\expandafter\unexpanded\expandafter{\@savedSecTitle}}%
      }%
    }%
    \protected@write\@auxout{}{\string\glsliveentry{as\theascounter}}%
  \fi
}
\newcommand{\eeAs}{\end{assumption}}
% CONJECTURE COMMANDS (glossaried)
\newcommand{\beCj}[1][]{%
  \def\@temparg{#1}%
  \ifx\@temparg\@empty
    \begin{conjecture}%
  \else
    \let\@savedSecTitle\@currentsectionname%
    \refstepcounter{cjcounter}%
    \begin{conjecture}[#1]%
    \label{cj:\thecjcounter}%
    \expanded{%
      \noexpand\newglossaryentry{cj\thecjcounter}{%
        name={\unexpanded{#1}},%
        sort={\thecjcounter},%
        description={p.~\noexpand\pageref{cj:\thecjcounter}},%
        user1={\thecjcounter},%
        user2={Conjecture},%
        user3={\thesection},%
        user4={\expandafter\unexpanded\expandafter{\@savedSecTitle}}%
      }%
    }%
    \protected@write\@auxout{}{\string\glsliveentry{cj\thecjcounter}}%
  \fi
}
\newcommand{\eeCj}{\end{conjecture}}
% LEMMA COMMANDS
% ---------------------------------------------------------------------
% UNIFORM RESULT WRAPPERS -- \beL, \beC, \beP, \beR share \beT's signature
%
%   \beX[long title][short title][key]  ...  \eeX
%
%   long title  = heading (math allowed)
%   short title = shown in the cross-reference; defaults to the key
%   key         = label key.  With a key the label is X:KEY and the heading
%                 carries the mnemonic [key], exactly as for \beT.  Without
%                 one the label falls back to the running counter, which is
%                 what every pre-existing call site in the book relies on.
%
% A bare \beX -- arguments commented out -- opens the plain environment and
% creates NO label and NO glossary entry.  That is the supported way to keep
% an individual result out of the catalogue.
% ---------------------------------------------------------------------
% #1 wrapper name  #2 environment  #3 label prefix  #4 counter  #5 type name
\newcommand{\@mkresult}[5]{%
  % A leading * suppresses the [key] mnemonic in both the heading and the
  % glossary name, while still making the label X:KEY and the entry:
  %     \beX*[long][short][key]
  \expandafter\NewDocumentCommand\csname be#1\endcsname{s O{} O{} o}{%
    \def\@temparg{##2}%
    \ifx\@temparg\@empty
      \begin{#2}%
    \else
      \let\@savedSecTitle\@currentsectionname%
      \refstepcounter{#4}%
      % With a key: mnemonic label X:KEY and the [key] tag in the heading.
      % Without one: fall back to the running counter, exactly as these
      % wrappers behaved before, so every legacy one-argument call site in
      % the book keeps its label.  The counter is also csname-safe, which a
      % title's first word is not -- titles beginning with math would
      % otherwise build \csname \Omega...\endcsname and fail.
      \IfValueTF{##4}{%
        \edef\@relbl{##4}%
        \IfBooleanTF{##1}{\let\@retag\@empty}{%
          \edef\@retag{\noexpand\ \noexpand\textup{[\noexpand\texttt{#3:##4}]}}}%
      }{%
        \edef\@relbl{\arabic{#4}}%
        \let\@retag\@empty
      }%
      \def\@tempshort{##3}%
      \ifx\@tempshort\@empty
        \expandafter\xdef\csname @sshort@#3\@relbl\endcsname{\@relbl}%
      \else
        \expandafter\def\csname @sshort@#3\@relbl\endcsname{##3}%
      \fi
      \begin{#2}[##2\ifshowkeys\@retag\fi]%
      \edef\@tmplbl{#3:\@relbl}\expandafter\label\expandafter{\@tmplbl}%
      \expanded{%
        \noexpand\newglossaryentry{#3\csname the#4\endcsname}{%
          name={\unexpanded{##2}%
            \IfBooleanF{##1}{\IfValueT{##4}{ [\noexpand\texttt{#3:##4}]}}},%
          sort={\csname the#2\endcsname},%
          description={p.~\noexpand\pageref{#3:\@relbl}},%
          user1={\csname the#2\endcsname},%
          user2={#5},%
          user3={\thesection},%
          user4={\expandafter\unexpanded\expandafter{\@savedSecTitle}}%
        }%
      }%
      \protected@write\@auxout{}{\string\glsliveentry{#3\csname the#4\endcsname}}%
      \protected@write\@auxout{}{\string\@thmtitle{\@tmplbl}{\unexpanded{##2}}}%
    \fi
  }%
  \expandafter\newcommand\csname ee#1\endcsname{\end{#2}}%
}
\@mkresult{L}{lemma}{l}{lemcounter}{Lemma}
% REMARK / COROLLARY / PROPOSITION -- same signature as \beL above
\@mkresult{R}{remark}{r}{remcounter}{Remark}
\@mkresult{C}{corollary}{c}{corcounter}{Corollary}
\@mkresult{P}{proposition}{p}{propcounter}{Proposition}
% OPEN-PROBLEM COMMAND -- generated with the same signature as \beT:
%   \beO[long title][short title][key]   ...   \eeO
% label o:KEY, or o:<counter> when no key is given; a bare \beO makes
% neither a label nor a glossary entry.
\@mkresult{O}{open}{o}{opcounter}{Problem}


% Section-level title tracker: updated only by \section, not \subsection.
% \@currentlabelname is overwritten by every \refstepcounter (including
% \subsection), so we save it separately here.
\newcommand{\@currentsectionname}{}
\AtBeginDocument{%
  \let\@epidcrn@orig@refstepcounter\refstepcounter
  \renewcommand{\refstepcounter}[1]{%
    \@epidcrn@orig@refstepcounter{#1}%
    \ifstrequal{#1}{section}{\global\let\@currentsectionname\@currentlabelname}{}%
  }%
}

% TABLE COMMANDS (glossaried) -- consistent naming with \beT, \beD, etc.
\newcommand{\beTa}[1][]{%
  \def\@temparg{#1}%
  \ifx\@temparg\@empty
    \begin{table}[htbp]%
  \else
    \let\@savedSecTitle\@currentsectionname%
    \refstepcounter{tabcounter}%
    \begin{table}[htbp]%
    \label{A:\thetabcounter}%
    \expanded{%
      \noexpand\newglossaryentry{A\thetabcounter}{%
        name={\unexpanded{#1}},%
        sort={\thetabcounter},%
        description={p.~\noexpand\pageref{A:\thetabcounter}},%
        user1={\thetabcounter},%
        user2={Table},%
        user3={\thesection},%
        user4={\expandafter\unexpanded\expandafter{\@savedSecTitle}}%
      }%
    }%
    \protected@write\@auxout{}{\string\glsliveentry{A\thetabcounter}}%
  \fi
}
\newcommand{\eeTa}{\end{table}}

% FIGURE COMMANDS (glossaried) -- consistent with \beT, \beR, \beTa, etc.
% Usage:
%   \beF[Glossary title]...\caption{Visible caption}\eeF  -> glossaried figure
%   \beF...\caption{Visible caption}\eeF                  -> plain [H] figure, no glossary
\newcommand{\beF}[1][]{%
  \def\@temparg{#1}%
  \ifx\@temparg\@empty
    \begin{figure}[H]%
  \else
    \begin{figure}[H]%
    \let\@savedSecTitle\@currentsectionname%
    \refstepcounter{figcounter}%
    \label{fig:\thefigcounter}%
    \expanded{%
      \noexpand\newglossaryentry{Fig\thefigcounter}{%
        name={\unexpanded{#1}},%
        sort={\thefigcounter},%
        description={p.~\noexpand\pageref{fig:\thefigcounter}},%
        user1={\thefigcounter},%
        user2={Figure},%
        user3={\thesection},%
        user4={\expandafter\unexpanded\expandafter{\@savedSecTitle}}%
      }%
    }%
    \protected@write\@auxout{}{\string\glsliveentry{Fig\thefigcounter}}%
  \fi
}
\newcommand{\eeF}{\end{figure}}
\makeatother
% LIVE-ENTRY REGISTER
% \glsaddall prints every entry the glossaries package currently knows, and
% that set includes entries restored from \jobname.glsdefs.  The package
% rewrites .glsdefs from all *currently defined* entries, so anything it
% restored is written straight back out: a deleted section kept reappearing in
% the catalogue for ever, and the .glsdefs had to be erased by hand.  The
% wrappers therefore record every entry they actually create in the .aux, and
% \printboth adds only those.  The .aux is rebuilt from scratch on each run, so
% this is self-healing: a removed result is gone after one further run.
\makeatletter
\newcommand{\@glslive}{}
\newcommand{\@glsaddlive}[1]{\ifglsentryexists{#1}{\glsadd{#1}}{}}
\newcommand{\glsliveentry}[1]{\g@addto@macro\@glslive{\@glsaddlive{#1}}}
\makeatother
% PRINT COMMAND
% No \clearpage here -- the catalogue continues on the current page.  This has
% been re-added by mistake several times; leave it out.
\makeatletter
\newcommand{\printboth}{%
  \@glslive   % add exactly the entries this document created last run
  \printglossary[style=unifiedstyle, title={0 \hspace{.22em} Definitions, Theorems, Examples and Open Problems}]%
}
\makeatother
% Named theorem reference: \tref{word} -> Thm~N(word-or-short)
% Shows stored short title if available, otherwise shows the key itself.
\newcommand{\tref}[1]{%
  Thm~\ref{t:#1}(\ifcsname @tshort@#1\endcsname\csname @tshort@#1\endcsname\else#1\fi)%
}

% Full-title theorem reference: \nref{t:KEY} prints the number and then the
% theorem's long title, so one writes "Theorem~\nref{t:blocks}" and gets
% "Theorem 3 ``Block factorization and siphons''".  Takes the complete label, like
% \ref, not the bare key like \tref.
%
% \beT writes each title to the .aux (see \@thmtitle below), so \nref resolves
% forward references from the second LaTeX pass, exactly as \ref does.  If the
% title is not yet known -- first pass, or a label that is not a \beT theorem
% -- it degrades to the number rather than producing ??.
% These two use @ in their names, and this point in the file is outside the
% \makeatletter region opened earlier (it closes at \makeatother above), so
% the region has to be reopened here.
\makeatletter
\newcommand{\@thmtitle}[2]{\global\@namedef{@tlong@#1}{#2}}
\newcommand{\nref}[1]{%
  \ref{#1}%
  \ifcsname @tlong@#1\endcsname
    \ ``\csname @tlong@#1\endcsname''%
  \fi
}
% \nm{t:KEY}: the stored long title of any \beT/\beL/\beP/\beR/\beC/\beD/\beE result,
% verbatim and without quotes, so item names follow the macro's first argument.
\newcommand{\nm}[1]{%
  \ifcsname @tlong@#1\endcsname
    \csname @tlong@#1\endcsname%
  \else
    \ref{#1}%
  \fi
}
\makeatother




%% ========================================
%% Other THEOREM ENVIRONMENTS
%% ========================================
\makeatletter
\@ifundefined{lemma}      {\newtheorem{lemma}      {Lemma}}      {}
\@ifundefined{proposition}{\newtheorem{proposition}{Proposition}} {}
\@ifundefined{corollary}  {\newtheorem{corollary}  {Corollary}}  {}
\@ifundefined{conjecture} {\newtheorem{conjecture} {Conjecture}}  {}
\@ifundefined{remark}     {\newtheorem{remark}     {Remark}}     {}
\@ifundefined{question}   {\newtheorem{question}   {Question}}   {}
\@ifundefined{assumption} {\newtheorem{assumption} {Assumption}} {}
\@ifundefined{hypothesis} {\newtheorem{hypothesis} {Hypothesis}} {}
\@ifundefined{fact}       {\newtheorem{fact}       {Fact}}       {}
\makeatother

% \citer: cite the relay-graph threshold theorem of AAHrel. If a document has
% \input the chapter/file containing t:relayRF (so the label resolves), point
% straight at it; otherwise fall back to a bare citation of AAHrel. Shared here
% (not per-paper) so it resolves consistently across every book chapter.
% \citeR: same idea, for the first theorem of AAHRah (t:Pstar).
\makeatletter
\providecommand{\citer}{%
  \@ifundefined{r@t:relayRF}%
    {\cite{AAHrel}}%
    {Theorem~\ref{t:relayRF} of \cite{AAHrel}}%
}
\providecommand{\citeR}{%
  \@ifundefined{r@t:Pstar}%
    {\cite{AAHRah}}%
    {Theorem~\ref{t:Pstar} of \cite{AAHRah}}%
}
\makeatother


% Abbreviations
\def\beQ{\begin{question}}\def\eeQ{\end{question}}
\def\beA{\begin{assumption}}\def\eeA{\end{assumption}}


%% ========================================

%% ========================================
%% COLORS & CUSTOM COMMANDS
%% ========================================
\definecolor{funccolor}{RGB}{25,25,112}
\definecolor{desccolor}{RGB}{64,64,64}
\newcommand{\func}[1]{\texttt{\textcolor{funccolor}{#1}}}
\def\eq{=}
\def\mm{-}

%% ========================================
%% MATHCAL SHORTHANDS
%% ========================================
% Project-wide convention: \mX for a capital X is always \mathcal X.
% Defined for the whole alphabet so the convention never fails to expand.
\newcommand{\mA}{\mathcal A}\newcommand{\mB}{\mathcal B}
\newcommand{\mC}{\mathcal C}\newcommand{\mD}{\mathcal D}
\newcommand{\mE}{\mathcal E}\newcommand{\mF}{\mathcal F}
\newcommand{\mG}{\mathcal G}\newcommand{\mH}{\mathcal H}
\newcommand{\mI}{\mathcal I}\newcommand{\mJ}{\mathcal J}
\newcommand{\mK}{\mathcal K}\newcommand{\mL}{\mathcal L}
\newcommand{\mM}{\mathcal M}\newcommand{\mN}{\mathcal N}
\newcommand{\mO}{\mathcal O}\newcommand{\mP}{\mathcal P}
\newcommand{\mQ}{\mathcal Q}\newcommand{\mR}{\mathcal R}
\newcommand{\mS}{\mathcal S}\newcommand{\mT}{\mathcal T}
\newcommand{\mU}{\mathcal U}\newcommand{\mV}{\mathcal V}
\newcommand{\mW}{\mathcal W}\newcommand{\mX}{\mathcal X}
\newcommand{\mY}{\mathcal Y}\newcommand{\mZ}{\mathcal Z}

%% ========================================
%% BASIC MATH ENVIRONMENTS
%% ========================================
\def\bep{\begin{pmatrix}}\def\eep{\end{pmatrix}}
\def\bev{\begin{vmatrix}}\def\eev{\end{vmatrix}}
\def\bea{\begin{eqnarray*}}\def\eea{\end{eqnarray*}}
\def\bc{\begin{cases}}\def\ec{\end{cases}}
\def\BEN{\begin{enumerate}}\def\EEN{\end{enumerate}}
\def\BI{\begin{itemize}}\def\EI{\end{itemize}}

\newcommand{\be}[1]{\begin{equation}\label{#1}}
\newcommand{\ee}{\end{equation}}
\newcommand{\beq}{\begin{eqnarray}}
\def\eeq{\end{eqnarray}}
\def\eqr{\eqref}\def\fr{\frac}\def\lbl{\label}
\def\Lra{\Longrightarrow}\def\non{\nonumber}\def\cNGM{\cite{Diek,Van,Van08}}
\def\Eq{\Leftrightarrow}\def\b{\beta}\def\uph{\red{UP TO HERE}}
%% ========================================
%% PROBABILITY & STATISTICS
%% ========================================
\renewcommand{\P}{\ensuremath{\mathbb{P}}}  % Note: overrides paragraph symbol
\newcommand{\E}{\ensuremath{\mathbb{E}}}
\newcommand{\ind}[1]{\ensuremath{\mathbbm{1}_{\left\{#1\right\}}}}
\DeclareMathOperator{\Cov}{Cov}
\DeclareMathOperator{\TV}{TV}
\newcommand{\rank}{\operatorname{rank}}

%% ========================================
%% CRN OPERATORS (organizations/ERCs/buffering/siphons framework)
%% ========================================
\def\clos{\mathrm{clos}}
\def\Supply{\mathrm{Supply}}
\def\req{\mathrm{req}}
\def\React{\mathrm{React}}
\def\Prod{\mathrm{Prod}}
\def\Gen{\mathrm{Gen}}
\def\In{\mathrm{In}}
\def\Out{\mathrm{Out}}
\def\Act{\mathrm{Act}}
\def\Min{\mathrm{Min}}

%% ========================================
%% LINE BREAKING
%% ========================================
% Last-resort stretch, used only on paragraphs TeX cannot otherwise break:
% a long \texttt{} token, a comma-separated inline math list (commas are not
% breakpoints in math), or a run of \cite/\ref boxes.  Without it those set as
% overfull \hbox and print into the margin; with it they set slightly loose,
% which is the lesser evil.  Does not affect paragraphs that already break.
\emergencystretch=3em

%% ========================================
%% NUMBER SETS
%% ========================================
\newcommand{\R}{\mathbb{R}}  % Real numbers
\newcommand{\C}{\mathbb{C}}  % Complex numbers
\newcommand{\N}{\mathbb{N}}
\newcommand{\Z}{\mathbb{Z}}
% \providecommand, not \newcommand: some papers define these in their own
% preamble after \input{preamble}, and must not clash.
\providecommand{\Q}{\mathbb{Q}}
\providecommand{\res}{\operatorname{res}}   % resultant; used by epiRes.tex
\newcommand{\MN}{\mathbb{N}}
\newcommand{\MZ}{\mathbb{Z}}
\newcommand{\realpos}{\mathbb{R}_{> 0}}
\newcommand{\realnneg}{\mathbb{R}_{\geq 0}}
\newcommand{\naturalpos}{\mathbb{N}_{>0}}
\newcommand{\Rplus}{\mathbb{R}_+}
\newcommand{\Rp}{\mathbb{R}_{\ge 0}}
\def\Rnne{\mathbb{R}_{\ge 0}}
\def\Se{\bar{S}} \def\Ie{\bar{I}}
\newcommand{\Rpos}{\mathbb{R}_{>0}}
\newcommand{\dd}{\,\mathrm d}
\newcommand{\e}{\mathrm e}
%% ========================================
%% GREEK LETTERS (2+ character abbreviations) greek
%% ========================================
\def\al{\alpha}\def\ba{\bff \alpha}
\def\w{\omega}\def\Lgn{\ga+\mu+\de}\def\om{\omega}
\def\de{\delta}\def\De{\Delta}\def\bde{\bff \delta}
\def\ep{\epsilon}\def\eps{\varepsilon}
\def\f{\varphi}  % Kept as requested
\def\ga{\gamma}\def\Ga{\Gamma}\def\G{\Gamma}\def\bg{\bff \gamma}
\def\up{\upsilon}  % NOT the breve accent: no accent commands anywhere.
                   % Named \up, not \u, so the LaTeX breve stays reachable.
\def\la{\lambda}\def\La{\Lambda}
\def\si{\sigma}\def\ka{\kappa}
\def\th{\theta}\def\vt{\vartheta}
\def\ze{\zeta}\def\Gap{\Gamma^{+}} \def\Gam{\Gamma^{-}}


%% ========================================
%% SPECIAL SHORT COMMANDS (kept as requested)
%% ========================================
\def\I{\infty}
\def\T{\widetilde}
\def\H{\hat}
\newcommand{\relint}{\operatorname{relint}}

%% ========================================
%% BOLD/VECTOR NOTATION
%% ========================================
\newcommand{\bff}[1]{{\mbox{\boldmath$#1$}}}
\newcommand\X{\boldsymbol{X}}
\def\x{\boldsymbol{x}}\def\bx{\bff x}\def\xe{\bar{\x}}
\newcommand\y{\boldsymbol{y}}\def\by{\boldsymbol{y}}\def\ye{\bar{\y}}
\def\yb{\y}\def\xb{\x}
\def\xy{\x-\y}
\newcommand\br{\boldsymbol{r}}\def\re{\bar{\r}}
\newcommand\bd{\boldsymbol{d}}
\newcommand\bQ{\boldsymbol{Q}}
\newcommand{\Yb}{{\mathbf{Y}}}
\newcommand{\Ab}{\mathbf{A}}
\newcommand{\zb}{\mathbf{z}}\def\bz{\bff 0}
\newcommand{\Zb}{\mathbf{Z}}
\newcommand{\bb}{\mathbf{b}}
\def\vbe{\vec {\bff \beta}}
\def\vb{\vec {\bff b}}
\def\v1{\vec {\bff 1}}
\def\vn{\vec \de}
\def\vz{\vec 0}
\def\vi{\overset{\rightarrow}{\mathsf i}}
\def\vv{\overset{\rightarrow}{v}}
\def\vr{\; \vec {\mathsf r}}
\newcommand{\va}{\bff a}
\providecommand{\vect}[1]{\boldsymbol{#1}}   % bold vector (moved from Gradostat.tex)
\providecommand{\one}{\mathbf 1}             % all-ones vector (moved from AliSF)
\providecommand{\spr}{s}                     % spectral bound  (moved from AliSF)
\providecommand{\Ke}{K_{\rm eff}}            % effective carrying capacity
% Critical right/left eigenvectors in the Hopf projection method (AAHOSN/New.tex)
\providecommand{\bq}{\mathbf{q}}
\providecommand{\bp}{\mathbf{p}}


%% ========================================
%% MATHCAL & MATHSCR
%% ========================================
% Mathcal
\newcommand\mc[1]{\mathcal{#1}}
\def\mC{{\mathcal C}}\def\mD{\mathcal D}\def\mE{\mathcal{E}}
\def\mF{\mathcal F}\def\mG{{\mathcal G}}\def\mH{{\mathcal H}}
\def\mI{\mathcal I}\def\mL{\mathcal{L}}\def\mM{{\mathcal M}}
\def\mN{{\mathcal N}}\def\mP{\mathcal{P}}\def\mR{{\mathcal R}}
\def\mS{{\mathcal S}}\def\mT{\mathcal{T}}\def\mV{{\mathcal V}}
\def\mK{\mathcal K}\def\mA{{\mathcal A}}\def\mU{{\mathcal U}}
% Symmetric bilinear and trilinear forms of the Taylor expansion, used in the
% first Lyapunov coefficient (AAHOSN/New.tex)
\providecommand{\Bform}{\mathcal{B}}
\providecommand{\Cform}{\mathcal{C}}

% Blackboard bold.  Kept here, not in the drivers: a \providecommand in a
% driver preamble runs before this file and silently wins, so the same name
% ends up meaning different things in different builds.
\def\IC{\mathbb{C}}\def\IN{\mathbb{N}}\def\IQ{\mathbb{Q}}
\def\IR{\mathbb{R}}\def\IZ{\mathbb{Z}}\def\IG{\mathbb{G}}
\def\II{\mathbb{I}}

% Mathscr
\newcommand{\Rs}{{\mathscr R}}
\newcommand{\Ss}{{\mathscr S}}
\newcommand{\Ns}{{\mathscr N}}
\newcommand{\Cs}{\mathscr C}


%% ========================================
%% OPERATORS & FUNCTIONS
%% ========================================
\DeclareMathOperator{\tr}{tr}
\DeclareMathOperator{\diag}{Diag}
\DeclareMathOperator{\Diag}{Diag}
\DeclareMathOperator{\intr}{\mathrm{Int}}
\DeclareMathOperator{\Inv}{Inv}
\DeclareMathOperator{\spec}{spec}
\DeclareMathOperator{\AUC}{AUC}               % area under the curve
\newcommand{\ii}{\mathrm{i}}                  % imaginary unit
\newcommand\op{\operatorname}
\newcommand{\supp}{\op{supp}}
\newcommand{\diff}{\mathop{}\mathopen{}\mathrm{d}}
\newcommand{\pdd}[1]{\frac{\partial}{\partial{#1}}}
\DeclareMathOperator{\cone}{cone}
\DeclareMathOperator{\proj}{proj}

%% ========================================
%% BRACKETS & NORMS
%% ========================================
\providecommand{\pr}[1]{\left(#1\right)}
\providecommand{\pp}[1]{\left[#1\right]}
\newcommand\croc[1]{\left\langle #1\right\rangle}
\newcommand{\norm}[1]{\|#1\|}
\newcommand{\inner}[2]{\langle #1, #2 \rangle}
% Short synonym of \inner, used by AAHOSN/New.tex as $\ip{u}{v}=\bar u^\top v$
\providecommand{\ip}[2]{\langle #1 , #2 \rangle}


%% ========================================
%% FORMATTING & UTILITIES
%% ========================================
\newcommand*{\Scale}[2][4]{\scalebox{#1}{$#2$}}
\newcommand{\te}[1]{\text{#1}}
\newcommand\steq[1]{\stackrel{\text{\rm #1.}}{=}}
\newcommand\blue[1]{\textcolor{blue}{#1}}
\newcommand{\red}{\textcolor[rgb]{1.00,0.00,0.00}}
\newcommand{\green}{\textcolor[rgb]{0.00,1.00,0.00}}

% Figure command
\newcommand{\figu}[3]{%
\begin{figure}[H]
\centering
\includegraphics[scale=#3]{#1}
\caption{#2\label{f:#1}}
\end{figure}
}

% Footnote
\long\def\symbolfootnote[#1]#2{%
\begingroup
\def\thefootnote{\fnsymbol{footnote}}\footnote[#1]{#2}%
\endgroup}
\def\fn{\symbolfootnote}

\newcommand{\dsum}{\displaystyle\sum}

%% ========================================
%% CRN & EPIDEMIOLOGY SPECIFIC
%% ========================================
% CRN notation
\def\SS{\mathcal S}
\def\CC{\mathcal C}
\def\RR{\mathcal R}

% Reproduction numbers
\def\Rto{R_{1}^{E_2}}
\def\Rot{R_{2}^{E_1}}
\def\Roc{R_{1,c}}
\def\Rtc{R_{2,c}}
\def\rd{r_{0}}
\def\sd{s_{0}}
\def\dd{D_{0}}
\def\vd{v_{0}}
\def\zd{z_{0}}
\def\se{\sE}
\def\xd{x_{0}}
%% ========================================
%% SPECIAL VALUES & EQUILIBRIA
%% ========================================
\def\eE{\bar{e}}
\def\wE{\bar{w}}
\def\sE{s_{*}}
\def\Me{\T M_{ee}}
\def\Rnne{\mathbb{R}_{\ge 0}}

\newcommand{\ones}{\mathbf 1}


%% ========================================
%% ABBREVIATIONS (alphabetical by category)
%% ========================================

% -- General mathematical terms --
\def\adm{admissible}\def\admy{admissibility}
%\def\and{antisymmetric}% disabled 2026-09-30: it overrode LaTeX's author separator \and (printed "antisymmetric" between authors)
\def\ass{assumption}
\def\cadlag{c\`adl\`ag}
\def\ch{characteristic polynomial}
\def\chp{characteristic polynomial}
\def\coe{coefficient}
\def\comp{computation}
\def\con{condition}
\def\coo{coordinate}
\def\corr{corresponding}
\def\dec{decomposition}
\def\den{denominator}
\def\dete{deterministic}
\def\eig{eigenvalue}\def\eval{eigenvalue}
\def\eigv{eigenvector}\def\eiv{eigenvector}\def\evec{eigenvector}
\def\elt{element}
\def\enn{essentially nonnegative/positivity preserving}
\def\equ{equation}
\def\etal{et al.}
\def\fe{for example}
\def\fno{from now on}
\def\for{\forall}
\def\form{formula}
\def\Fr{Furthermore, }
\def\frt{furthermore}
\def\fun{function}
\def\gene{generalization}
\def\heu{heuristic}
\def\how{however}\def\How{However, }
\def\ie{i.e. }
\def\im{\item}
\def\isd{is defined by}
\def\itc{it is easy to check that}
\def\itf{it follows that}\def\Itf{It follows that}
\def\ith{it holds that}
\def\itm{it may be checked that}\def\Itm{It may be checked that}
\def\ito{it turns out that}\def\Ito{It turns out that}\def\Itou{It turns out useful to}
\def\ma{mass action}\def\MAR{mass action representation}
\def\mat{matrix}
\def\mbc{may be checked to be}
\def\mbw{may be written as}\def\mubw{must be written as}
\def\mpr{more precisely, }\def\Mp{More precisely, }
\def\nc{non-catalytic}
\def\nne{non-negative}\def\nny{nonnegativity}
\def\no{\nonumber}
\def\noni{non-infectious}
\def\nota{notation}
\def\np{\newpage}
\def\Nt{Note that}
\def\nt{non trivial}\def\nz{non zero}
\def\oth{otherwise}
\def\Oth{On the other hand}
\def\pa{parameter}\def\para{parameter}
\def\pol{polynomial}
\def\prob{probability}
\def\Prf{{\bf  Proof:}}\def\prf{{\bf  Proof:}}
\def\prop{proportional to}
\def\qu{\quad}
\def\QED{\hfill {$\square$}\goodbreak \medskip}
\def\rep{representation}
\def\resp{respectively}
\def\satd{satisfied}\def\satg{satisfying}\def\sats{satisfies}\def\saty{satisfy}
\def\ssec{\subsection}\def\sssec{\subsubsection}
\def\sef{see for example}
\def\sksy{skew-symmetric}
\def\sub{substrate}
\def\sys{system}
\def\var{variable}
\def\wk{well-known}
\def\wl{without loss of generality}\def\wlo{w.l.o.g.}
\def\wmc{We may check that}\def\Wmc{We may check that}
\def\wmr{we may rewrite}
\def\wms{we must show that}
\def\wrt{with respect to}

% -- CRN & Chemistry --
\newcommand\CRN{chemical reaction network}
\newcommand\CRNs{chemical reaction networks}
\newcommand\RN{reaction network}
\def\ACR{absolute concentration robustness}
\def\AMSD{acyclic siphon decomposition}
\def\DAE{differential algebraic equations}
\def\FHJ{Feinberg-Horn-Jackson}
\def\GMAK{generalized mass action kinetics}
\def\MM{Michaelis-Menten}
\def\sm{stoichiometric matrix}\def\Sm{stoichiometric matrix}
\def\WRZD{weakly reversible, zero deficiency}\def\aut{autocatalytic}

% -- Epidemiology --
\def\brn{basic reproduction number}\def\bzn{basic replacement number}
\def\com{compartment}
\def\dfe{disease free equilibrium}\def\DFE{disease free equilibrium}
\def\dwt{dwell times}\def\edt{expected dwell times}
\def\EM{epidemic model}
\def\endp{endemic fixed point}\def\enp{endemic point}\def\EE{endemic point}
\def\ME{mathematical epidemiology}
\def\MEE{eco-epidemiology}\def\EEP{eco-epidemiological}
\def\ece{eco-epidemiological model}
\def\mgm{multi group model}
\def\mtS{multi-strain model}\def\mSt{multi-strain model}
\def\mSi{minimal siphon}
\def\NGM{next generation matrix}
\def\nim{new infections matrix}
\def\qem{quadratic epidemiological model}
\def\repF{invasion functions}
\def\SPF{SIR-PH-FA model}
\def\sus{susceptible}

% -- Graph theory & Linear algebra --
\def\jac{Jacobi matrix}
\def\jin{Jacobian matrix of the invasion vector field with respect to the invasion variables}
\def\jrn{Jacobian matrix of the resident vector field with respect to the resident variables}
\def\jxn{Jacobian matrix of the $\x$ vector field with respect to the $\x$ variables}
\def\Lm{Laplacian matrix}
\def\MTT{matrix tree theorem}
\def\SGT{spectral graph theory}
\def\SMW{Sherman–Morrison–Woodbury}

% -- Stability & Dynamics --
\def\absC{absorbtion CTMC}
\def\BTb{Bogdanov-Takens bifurcation}
\def\CEP{competitive exclusion partition}
\def\CTMC{continuous time Markov chain}
\def\ds{differential system}
\def\fp{fixed point}\def\bfp{boundary fixed point}\def\fbp{fixed boundary point}
\def\fbl{fixed boundary line}\def\fif{forward invariant face}
\def\gas{global asymptotic stability}\def\GAS{globally asymptotically stable}
\def\GASS{\text{GASS}}
\def\Hb{Hopf bifurcation}
\def\Hur{Hurwitz}\def\MH{Metzler-Hurwitz}\def\RH{Routh-Hurwitz}
\def\IS{\text{IS}}
\def\LAS{locally asymptotically stable}
\def\LdS{Lyapunov diagonally stable}
\def\Lf{Lyapunov function}\def\LyV{Lyapunov Volterra type function}\def\RLF{robust Lyapunov function}
\def\LyM{Lyapunov Malkin stable}
\def\PF{Perron-Frobenius}
\def\spfp{strictly positive fixed point}
\def\sta{stability}
\def\STP{sharp threshold property}\def\STT{sharp threshold theorem}

% -- Models (with \varphi explicitly instead of \f) --
\def\AB{$(A,B)$ Arino-Brauer epidemic models}
\def\ABa{$(A,\bbe,\ba,\varphi)$ balanced bilinear model}
\def\ABP{%$(A,\bbe,P,\varphi)$
\\ balanced bilinear model}
\def\AMG{Anderson May Gupta}
\newcommand\BIN{biological interaction network}
\def\Dt{Descartes type}
\def\gLV{generalized Lotka-Volterra model}
\def\IG{\text{invasion graph}}\def\IGMS{interaction graph on minimal siphons}
\def\LCP{\text{linear complementarity problem}}
\def\LV{Lotka-Volterra}\def\LVK{Lotka-Volterra-Kolmogorov}

% -- Splitting & Structure --
\def\reg{regular splitting}\def\regS{{\bf regular splitting}}
\def\RUR{rational univariate representation}
\def\bds{block-diagonal structure}\def\NGMs{next generation matrices}

% -- Miscellaneous --
\def\Ma{{Mathematica}}
\def\Mm{Max–Min}
\def\num{numerator}
\def\rap{rational approximation} \def\DSR{directed species-reaction bipartite graph}\def\PCP{polynomial complementarity problem}\def\rat{rational substitution}
\def\QSS{quasi steady state} \def\RHS{right hand side} \def\LHS{left hand side}
\def\OSN{online social network}\def\VdD{Van den Driessche}
\def\wli{$\omega$-limit} \def\rk{rich kinetics}\def\sJ{symbolic Jacobian}
\def\cV{\cite{Vas,VasSta,BSV,GSV}} \def\spf{semi-parametric semi-functional}
\def\CRW{Capasso-Ruan-Wang} \def\ABP{$(A,\bbe,P,\f)$ balanced bilinear model}
\def\ABa{%$(A,\bbe,\ba,\f)$ \\
rank-one balanced bilinear model}
\def\bbe{B%{\bff \beta}
} \def\vbe{\vec {\bff \beta}
}
\def\vb{\vec {\bff b}}
\newcommand{\Kh}{\hat{K}}          % K with hat, first mention only
\newcommand{\wS}{w^{\sigma}}       % right Perron J^\perp, full block
\newcommand{\piS}{\vec \pi^{\sigma}}    % left Perron J^\perp, full block
\newcommand{\ws}{w^{\sigma}}       % right Perron J^\perp, sub-block
\newcommand{\pis}{\pi^{\sigma}}    % left Perron J^\perp, sub-block
\newcommand{\hwS}{\hat{w}^{\sigma}}   % right Perron K, full block
\newcommand{\hpiS}{\hat{\pi}^{\sigma}} % left Perron K, full block
\newcommand{\hws}{\hat{w}^{\sigma}}   % right Perron K, sub-block
\newcommand{\hpis}{\hat{\pi}^{\sigma}} % left Perron K, sub-block
\newcommand{\twS}{\tilde{w}^{\sigma}}  % right Perron \tilde K, full block
\newcommand{\tpiS}{\tilde{\pi}^{\sigma}} % left Perron \tilde K
\newcommand{\tws}{\tilde{w}^{\sigma}}
\newcommand{\tpis}{\tilde{\pi}^{\sigma}}
\newcommand{\Mat}[2]{\mathbb{R}_{\ge 0}^{\,#1\times #2}}
\def\tJ{J_\sigma^\perp} \def\lPe{left Perron eigenvector}
% \G is \Gamma (defined at :795).  The Volterra function keeps its own name,
% so that \G no longer has two conflicting meanings in this file.
\def\GV{G}  % Volterra G(u)=u-1-\ln u
\def\PVL{Perron-Volterra Lyapunov minimal siphon ansatz}
\newcommand{\istar}[1]{i_{#1}^{*}}
%\newcommand{\sE}{s_E}
\newcommand{\vE}{v_E}
\def\PV{Perron-Volterra candidate Lyapunov function}
\def\GLV{generalized Lotka-Volterra}
\def\MERN{mathematical epidemiology reaction networks}
\def\CMM{contraction metric method}
\newcommand{\sign}{\operatorname{sign}}
\newcommand{\Var}{\operatorname{Var}} 

% ======================================================================
% IDEAS SHARED BETWEEN THE ABSTRACT AND THE MAIN RESULTS SECTION
%
% Write each idea ONCE, in the abstract; the Main Results section recalls it,
% so the two can never drift apart.
%
%   in the abstract        \ideaitem{1}{The first-order reaction terms ...}
%   in Main Results        \ideahead{1}
%
% \ideaitem stores the text globally and typesets
%   \item \hyperlink{idea:1}{<text>}
% \ideahead typesets
%   \item\hypertarget{idea:1}{}\textbf{<the same text>}
%
% The abstract must come before the Main Results input, which is the normal
% order. If a heading has to read differently from the abstract sentence (for
% instance a long abstract item summarised in one line), override it with the
% optional argument: \ideahead[short heading]{6}.
% ======================================================================
\makeatletter
\newcommand{\ideaitem}[2]{%
  \expandafter\gdef\csname @idea@#1\endcsname{#2}%
  \item\hyperlink{idea:#1}{#2}}
\newcommand{\ideahead}[2][]{%
  \item\hypertarget{idea:#2}{}%
  \def\@ideatmp{#1}%
  \ifx\@ideatmp\@empty
    \@ifundefined{@idea@#2}{%
      \PackageWarning{preamble}{Idea #2 is used in Main Results but was never
        defined by \string\ideaitem\space in the abstract}%
      \textbf{[idea #2 undefined]}%
    }{\textbf{\begingroup\def\ideaterm##1##2{##2}\def\idearest##1{}\@nameuse{@idea@#2}\endgroup}}%
  \else
    \textbf{#1}%
  \fi}
% \ideaterm{<hypertarget>}{<words>}: inside an abstract idea the words become an
% invisible link (the document sets \hypersetup{hidelinks}) jumping to the place
% in Main Results that defines them; inside the Main Results heading itself,
% where the target sits a few lines below, \ideahead strips the link and prints
% the words plainly. Mark the destination with \hypertarget{<name>}{} on the line
% before the display, which adds no equation number and shows nothing.
\newcommand{\ideaterm}[2]{\hyperlink{#1}{#2}}
% \idearest{<text>}: the part of an abstract idea after its first sentence; printed in
% the abstract, omitted from the Main Results heading.
\newcommand{\idearest}[1]{#1}
\makeatother

% ======================================================================
% SHORT-PAPER SWITCH
% \shorttrue in a master that assembles only part of the companion paper
% (Barrat/SipArxiv.tex). \nsh{...} drops text that cites results living in the
% omitted sections, so the short build has no dangling references while the
% companion paper keeps the text verbatim.
% ======================================================================
\newif\ifshort \shortfalse
\newcommand{\nsh}[1]{\ifshort\else#1\fi}
% \shortalt{<short build>}{<companion paper>}: use when a sentence needs
% different wording in the two builds, not merely a deletion.
\newcommand{\shortalt}[2]{\ifshort#1\else#2\fi}
% \Jtr, \JtrS{\Sigma}: transverse Jacobian. J^\perp here; SipArxiv.tex redefines them as J_{\mathrm{tra}}, J_{\mathrm{tra},\Sigma}
\newcommand{\Jtr}{J^\perp}
\newcommand{\JtrS}[1]{J_{#1}^\perp}
.
\makeatletter
\@ifpackageloaded{hyperref}{}{%
  \@ifundefined{hyperrefpackage}{%
    \@ifundefined{AAdefNnohyperref}{\usepackage{hyperref}}{}%
  }{}%
}
\makeatother
% linktoc=all makes a table-of-contents page number a link as well as the
% entry text; hyperref's default, linktoc=section, links only the text.
% Deferred to begin-document and guarded: by the block above hyperref may
% be loaded later (mdpi's own hook) or not at all (\AAdefNnohyperref), and
% \hypersetup must exist before it is called.  A driver that wants the
% default back can say \hypersetup{linktoc=section} after % Shared preamble for the GLVbook papers.  This is the single canonical copy;
% a driver one level below GLVbook reaches it as
%
%     \input{../basicFiles/preamble}
%
% What stays in the driver, because it differs from paper to paper:
% \documentclass, font choice, \title, \author, \date, and any macro genuinely
% local to one paper.  Everything else belongs here.  The file is
% self-contained: it does not read any other file.

\usepackage[utf8]{inputenc}
\usepackage[novbox]{pdfsync}   % required: SyncTeX alone does not work here

% Page geometry, set once here for every paper that loads this file.  A driver
% that needs different margins says \geometry{...}; a second
% \usepackage[...]{geometry} raises an option clash and must not be used.
\makeatletter
\@ifundefined{AAdefNnogeometry}{\usepackage[a4paper,margin=2cm]{geometry}}{}
\makeatother

% mathematics
\usepackage{amsmath,amssymb,amsthm,mathtools,bm,mathrsfs,bbm}
\usepackage{centernot}

% floats, tables, lists
\usepackage{graphicx,float}
\makeatletter
\@ifpackageloaded{subfig}{}{\usepackage{subcaption}}
\makeatother
\usepackage{booktabs,longtable,threeparttable,array,multirow}
\usepackage{enumitem}

% colour, code, algorithms
\usepackage{xcolor,listings,tcolorbox}
\makeatletter
\@ifundefined{c@algorithm}{\usepackage{algorithm}}{}
\makeatother
\usepackage{algorithmicx,algpseudocode}

% diagrams
\usepackage{tikz,tikz-cd,pgfplots}
\usetikzlibrary{graphs,decorations.pathmorphing,decorations.markings,
arrows.meta,automata,positioning}
\tikzcdset{every label/.append style = {font = \small}}
\pgfplotsset{compat=1.18}

% chemistry
\usepackage[version=4]{mhchem}

\usepackage{microtype}
\usepackage{lineno}
\usepackage{nomencl}
\makenomenclature

% hyperref first, then bookmark, then glossaries: glossaries must follow
% hyperref or the entry links are built against the wrong anchors.
\makeatletter
\@ifundefined{hyperrefpackage}{%
  \@ifundefined{AAdefNnohyperref}{\usepackage{hyperref}}{}%
}{}
\@ifundefined{hyperrefpackage}{\usepackage{bookmark}}{}
\makeatother
\usepackage[toc,sort=use]{glossaries}
\makeglossaries

% Book/standalone switch.  Material that the book already covers in its
% prelude, but that a standalone paper still has to state, is wrapped as
%     \ifbookbuild\else ... \fi
% Drivers default to standalone; bookM.tex says \bookbuildtrue.
\newif\ifbookbuild
\bookbuildfalse

% Part switches for the gradostat paper.  Both default to true; a driver that
% wants the short submission says \madecfalse and/or \slowfastfalse, and every
% cross-reference into the dropped part disappears with it.
%   \ifmadec     Part III, the Lagasquie--Madec model and its break-even geometry
%   \ifslowfast  Part IV, everything carrying the small parameter \eps
\newif\ifmadec     \madectrue
\newif\ifslowfast  \slowfasttrue

% Part switch for the Barrat/simplicial-contagion project.  Defaults to
% true (the hypernetwork-synchrony part is included); the journal-paper
% driver says \hnetpartfalse to drop it and every cross-reference into it.
\newif\ifhnetpart  \hnetparttrue

%% ========================================
%% PACKAGES
%% ========================================
\RequirePackage{etoolbox}
\usepackage{makecell}
\usepackage[novbox]{pdfsync}% obsolete; SyncTeX is built into modern engines
\usepackage{amssymb, amsfonts,mathtools,bbm,eurosym,graphicx,latexsym, amsmath,amsthm,times,nccmath,url}
\usepackage{listings,verbatim,cases,tabularx}
% hyperref is deliberately NOT in the line above.  Journal classes often
% load it themselves, with a long option list, from a begindocument hook
% (mdpi.cls does: \DeclareHookRule{begindocument}{hyperref}{before}{mdpi}).
% Loading it here first, without options, makes that later load raise
% "Option clash for package hyperref".  So load it only if nobody else
% has, or will:
%   - \@ifpackageloaded  catches a class that loaded it already;
%   - \hyperrefpackage    is mdpi.cls's own deferred loader;
%   - \AAdefNnohyperref   lets any other driver opt out by defining it
%                         before \input{preamble}.
\makeatletter
\@ifpackageloaded{hyperref}{}{%
  \@ifundefined{hyperrefpackage}{%
    \@ifundefined{AAdefNnohyperref}{\usepackage{hyperref}}{}%
  }{}%
}
\makeatother
% linktoc=all makes a table-of-contents page number a link as well as the
% entry text; hyperref's default, linktoc=section, links only the text.
% Deferred to begin-document and guarded: by the block above hyperref may
% be loaded later (mdpi's own hook) or not at all (\AAdefNnohyperref), and
% \hypersetup must exist before it is called.  A driver that wants the
% default back can say \hypersetup{linktoc=section} after \input{preamble}.
\makeatletter
\AtBeginDocument{\@ifundefined{hypersetup}{}{\hypersetup{linktoc=all}}}
\makeatother
\usepackage{graphicx,float,fancybox}
\usepackage{epsfig,epstopdf,tikz}
\usepackage[sans]{dsfont}
\usetikzlibrary{positioning,arrows,arrows.meta,calc}
% Configure listings for Mathematica code
\lstdefinestyle{mathematica}{
 language=Mathematica,
 basicstyle=\small\ttfamily,
 backgroundcolor=\color{gray!10},
 frame=single,
 breaklines=true,
 commentstyle=\color{green!60!black},
 keywordstyle=\color{blue},
 stringstyle=\color{red},
 showstringspaces=false,
 tabsize=2
}
\usepackage{tcolorbox}
% Colored boxes for key results
\newtcolorbox{keyresult}[1][]{
 colback=blue!5!white,
 colframe=blue!75!black,
 title=#1,
 fonttitle=\bfseries
}
\newtcolorbox{examplebox}[1][]{
 colback=green!5!white,
 colframe=green!75!black,
 title=#1,
 fonttitle=\bfseries
}


\makeatletter
\@ifundefined{theorem}   {\newtheorem{theorem}   {Theorem}}   {}
\@ifundefined{definition}{\newtheorem{definition}{Definition}} {}
\@ifundefined{open}      {\newtheorem{open}      {Problem}}   {}
\@ifundefined{example}   {\newtheorem{example}   {Example}}   {}
\makeatother

% COUNTER DEFINITIONS
\newcounter{tabcounter}
\newcounter{defcounter}
\newcounter{opcounter}
\newcounter{thmcounter}
\newcounter{excounter}
\newcounter{ascounter}
\newcounter{cjcounter}
\newcounter{lemcounter}
\newcounter{remcounter}
\newcounter{corcounter}
\newcounter{propcounter}
\newcounter{figcounter}

% UNIFIED GLOSSARY STYLE
\newglossarystyle{unifiedstyle}{%
  \setglossarystyle{list}%
  \renewcommand*{\glossarypreamble}{\gdef\glslastsec{}}%
  \renewcommand{\glossentry}[2]{%
    \edef\glscursec{\glsentryuseriii{##1}}%
    \ifx\glscursec\glslastsec\else
      \item[]\textbf{\S\,\glscursec\enspace\glsentryuseriv{##1}}%
      \global\let\glslastsec\glscursec
    \fi
    \item[\glsentryuserii{##1} \glsentryuseri{##1}:] ``\glossentryname{##1}'' \glsentrydesc{##1}%
  }%
}

\makeatletter
% First-word extractor (used by \beT, \beD and \beXa for the default key)
% \@firstwordof hello world \relax  ->  hello
\def\@firstwordof#1 #2\relax{#1}
% First three letters, used for the heading mnemonic of a theorem given no
% explicit key.  Padding with \@empty keeps it safe for first words shorter
% than three letters: \@empty vanishes under \edef, whereas \relax would
% survive into the label and is illegal inside \csname (it broke every
% two-letter title, e.g. "ME ...", with "Missing \endcsname inserted").
\def\@takethree#1#2#3#4\@nil{#1#2#3}

% EXAMPLE COMMAND -- same signature and behaviour as \beT and \beD.
% Usage: \beXa*[long title][short title][key]
%   star  = suppresses the printed "[key]" tag after the title (as for \beT*
%   and \beD*); the label and glossary entry are unaffected.
%   key = optional label key override; the label becomes e:KEY, and without a
%   key it is e: plus the first three letters of the title's first word.  A
%   mnemonic label lets an example be cited as "Example~\ref{e:KEY}" instead
%   of by section number, which is what the citing prose needs.
\NewDocumentCommand{\beXa}{s O{} O{} o}{%
  \def\@temparg{#2}%
  \ifx\@temparg\@empty
    \begin{example}%
  \else
    \let\@savedSecTitle\@currentsectionname%
    \refstepcounter{excounter}%
    \IfValueTF{#4}{%
      \edef\@exlblkey{#4}%
    }{%
      \edef\@exfirstword{\@firstwordof#2 \relax}%
      \edef\@exlblkey{%
        \expandafter\@takethree\@exfirstword\@empty\@empty\@empty\@nil}%
    }%
    \def\@tempshort{#3}%
    \ifx\@tempshort\@empty
      \expandafter\xdef\csname @xshort@\@exlblkey\endcsname{\@exlblkey}%
    \else
      \expandafter\def\csname @xshort@\@exlblkey\endcsname{#3}%
    \fi
    \IfBooleanTF{#1}{\let\@extag\@empty}{%
      \edef\@extag{\noexpand\ \noexpand\textup{[\noexpand\texttt{e:\@exlblkey}]}}}%
    \begin{example}[#2\ifshowkeys\@extag\fi]%
    \label{e:\@exlblkey}%
    % record the long title so \nref{e:KEY} prints it, as for theorems
    \protected@write\@auxout{}{%
      \string\@thmtitle{e:\@exlblkey}{\unexpanded{#2}}}%
    \expanded{%
      \noexpand\newglossaryentry{e\theexcounter}{%
        name={\unexpanded{#2}%
          \IfBooleanF{#1}{ [\noexpand\texttt{e:\@exlblkey}]}},%
        sort={\theexample},%
        description={p.~\noexpand\pageref{e:\@exlblkey}},%
        user1={\theexample},%
        user2={Example},%
        user3={\thesection},%
        user4={\expandafter\unexpanded\expandafter{\@savedSecTitle}}%
      }%
    }%
    \protected@write\@auxout{}{\string\glsliveentry{e\theexcounter}}%
  \fi
}
\newcommand{\eeXa}{\end{example}}

% DEFINITION COMMAND -- same signature and behaviour as \beT below.
% Usage: \beD[long title][short title][key]
%   long title  = definition heading (plain text only)
%   short title = shown in \dref{word}; defaults to the key
%   key         = optional label key override; the label becomes d:KEY
% Without a key the label is d: plus the first three letters of the title's
% first word, so it is a mnemonic rather than the definition's number: a
% numeric label would silently change whenever a definition is inserted.
\NewDocumentCommand{\beD}{s O{} O{} o}{%
  \def\@temparg{#2}%
  \ifx\@temparg\@empty
    \begin{definition}%
  \else
    \let\@savedSecTitle\@currentsectionname%
    \refstepcounter{defcounter}%
    \IfValueTF{#4}{%
      \edef\@deflblkey{#4}%
    }{%
      \edef\@deffirstword{\@firstwordof#2 \relax}%
      \edef\@deflblkey{%
        \expandafter\@takethree\@deffirstword\@empty\@empty\@empty\@nil}%
    }%
    \def\@tempshort{#3}%
    \ifx\@tempshort\@empty
      \expandafter\xdef\csname @dshort@\@deflblkey\endcsname{\@deflblkey}%
    \else
      \expandafter\def\csname @dshort@\@deflblkey\endcsname{#3}%
    \fi
    \IfBooleanTF{#1}{\let\@deftag\@empty}{%
      \edef\@deftag{\noexpand\ \noexpand\textup{[\noexpand\texttt{d:\@deflblkey}]}}}%
    \begin{definition}[#2\ifshowkeys\@deftag\fi]%
    \label{d:\@deflblkey}%
    % record the long title so \nref{d:KEY} prints it, as for theorems
    \protected@write\@auxout{}{%
      \string\@thmtitle{d:\@deflblkey}{\unexpanded{#2}}}%
    \expanded{%
      \noexpand\newglossaryentry{d\thedefcounter}{%
        name={\unexpanded{#2}%
          \IfBooleanF{#1}{ [\noexpand\texttt{d:\@deflblkey}]}},%
        sort={\thedefinition},%
        description={p.~\noexpand\pageref{d:\@deflblkey}},%
        user1={\thedefinition},%
        user2={Definition},%
        user3={\thesection},%
        user4={\expandafter\unexpanded\expandafter{\@savedSecTitle}}%
      }%
    }%
    \protected@write\@auxout{}{\string\glsliveentry{d\thedefcounter}}%
  \fi
}
\newcommand{\eeD}{\end{definition}}

% \showkeysfalse in a driver hides the [t:KEY] mnemonics in result headings
% (labels and glossary entries are unchanged); used for submission builds.
\newif\ifshowkeys \showkeystrue
% \commonredtrue in a driver that inputs basicFiles/commonRed: tool.tex then
% cites its reduction results instead of restating them.
\newif\ifcommonred \commonredfalse
% TheorEM COMMAND
% Usage: \beT[long title][short title][key]
%   long title  = theorem heading (plain text only)
%   short title = shown in \tref{word}; may contain math; defaults to first word of long title
%   key         = optional label key override; label becomes t:KEY (avoids first-word collisions)
% Label is t:WORD where WORD = first word of long title (or KEY if given); use \tref{WORD} to reference.
\NewDocumentCommand{\beT}{s O{} O{} o}{%
  \def\@temparg{#2}%
  \ifx\@temparg\@empty
    \begin{theorem}%
  \else
    \let\@savedSecTitle\@currentsectionname%
    \refstepcounter{thmcounter}%
    % \@thmtag is the mnemonic shown in the heading; every theorem gets one.
    % With an explicit key the key itself is shown.  Without one the key
    % defaults to the title's first word, and showing that whole word would
    % merely repeat the title ("Dissipativity [Dissipativity]"), so only its
    % first three letters are shown: "Dissipativity [Dis]".  The label is
    % unaffected -- it still uses the full \@thmlblkey.
    \IfValueTF{#4}{%
      \edef\@thmlblkey{#4}%
    }{%
      \edef\@thmfirstword{\@firstwordof#2 \relax}%
      \edef\@thmlblkey{%
        \expandafter\@takethree\@thmfirstword\@empty\@empty\@empty\@nil}%
    }%
    \IfBooleanTF{#1}{\let\@thmtag\@empty}{%
      \edef\@thmtag{\noexpand\ \noexpand\textup{[\noexpand\texttt{t:\@thmlblkey}]}}}%
    \def\@tempshort{#3}%
    \ifx\@tempshort\@empty
      \expandafter\xdef\csname @tshort@\@thmlblkey\endcsname{\@thmlblkey}%
    \else
      \expandafter\def\csname @tshort@\@thmlblkey\endcsname{#3}%
    \fi
    \begin{theorem}[#2\ifshowkeys\@thmtag\fi]%
    \label{t:\@thmlblkey}%
    % record the long title for \nref; via the .aux so forward refs resolve
    \protected@write\@auxout{}{%
      \string\@thmtitle{t:\@thmlblkey}{\unexpanded{#2}}}%
    \expanded{%
      \noexpand\newglossaryentry{t\thethmcounter}{%
        name={\unexpanded{#2}%
          \IfBooleanF{#1}{ [\noexpand\texttt{t:\@thmlblkey}]}},%
        % sort/user1 must be the number the document actually prints, i.e. the
        % theorem environment's counter, not thmcounter.  thmcounter advances
        % only on \beT with a title, so a bare \beT (or \beT with an
        % empty title) makes the two drift and the glossary cites a wrong
        % number.  thmcounter is still fine as the entry key: it is unique.
        sort={\thetheorem},%
        description={p.~\noexpand\pageref{t:\@thmlblkey}},%
        user1={\thetheorem},%
        user2={Theorem},%
        user3={\thesection},%
        user4={\expandafter\unexpanded\expandafter{\@savedSecTitle}}%
      }%
    }%
    \protected@write\@auxout{}{\string\glsliveentry{t\thethmcounter}}%
  \fi
}
\newcommand{\eeT}{\end{theorem}}

% CROSS-PART RELATION PARAGRAPH
% Usage: \xpart{text}  -- place immediately BEFORE \beT, never inside it
%   (anything inside \beT...\eeT would be swallowed into the glossary entry).
% The paragraph is typeset only when Part I is actually present, i.e. in the
% combined book build; in the standalone builds (GranadaDSABNS/A.tex,
% AAHRah-Submission/AAHRah.tex) the label part:glv is undefined and the whole
% paragraph disappears rather than leaving ?? cross-references. Same
% graceful-degradation principle as \GASref in GASt.tex. Resolves on the
% second LaTeX pass, since r@part:glv comes from the .aux file.
% (Already inside the \makeatletter region opened above, so no \makeatletter
% pair here -- adding one would close that region early.)
\newcommand{\xpart}[1]{\@ifundefined{r@part:glv}{}{\paragraph{Relation with Parts~I--II.}#1}}

% \bookref{label}{precise text}{fallback text}: prints the precise cross-reference
% when the label resolves (combined book build), and the fallback text otherwise
% (standalone paper builds, where book-only appendix labels are absent). Same
% graceful-degradation role as \GASref in GASt.tex, but usable in any file since it
% lives in this preamble (loaded by every driver) -- unlike \GASref, which is defined in a
% content file and is therefore undefined in standalone builds that omit that file.
\newcommand{\bookref}[3]{\@ifundefined{r@#1}{#3}{#2}}

% \PVpart and \RBpart: the two parts that a partial build may omit.
% They print "Part~N" in any build that contains the part, and a naming phrase
% in a build that does not, so an excerpt shows no ?? .  Each already supplies
% the word "Part", so write "\PVpart develops ...", not "Part~\PVpart".
\newcommand{\PVpart}{\bookref{part:aahrah}%
  {Part~\ref{part:aahrah}}{the Perron--Volterra part of the full book}}
\newcommand{\RBpart}{\bookref{part:recbrowne}%
  {Part~\ref{part:recbrowne}}{the immuno-virology part of the full book}}

% OPEN PROBLEM COMMANDS
% ASSUMPTION COMMANDS (glossaried)
\newcommand{\beAs}[1][]{%
  \def\@temparg{#1}%
  \ifx\@temparg\@empty
    \begin{assumption}%
  \else
    \let\@savedSecTitle\@currentsectionname%
    \refstepcounter{ascounter}%
    \begin{assumption}[#1]%
    \label{as:\theascounter}%
    \expanded{%
      \noexpand\newglossaryentry{as\theascounter}{%
        name={\unexpanded{#1}},%
        sort={\theascounter},%
        description={p.~\noexpand\pageref{as:\theascounter}},%
        user1={\theascounter},%
        user2={Assumption},%
        user3={\thesection},%
        user4={\expandafter\unexpanded\expandafter{\@savedSecTitle}}%
      }%
    }%
    \protected@write\@auxout{}{\string\glsliveentry{as\theascounter}}%
  \fi
}
\newcommand{\eeAs}{\end{assumption}}
% CONJECTURE COMMANDS (glossaried)
\newcommand{\beCj}[1][]{%
  \def\@temparg{#1}%
  \ifx\@temparg\@empty
    \begin{conjecture}%
  \else
    \let\@savedSecTitle\@currentsectionname%
    \refstepcounter{cjcounter}%
    \begin{conjecture}[#1]%
    \label{cj:\thecjcounter}%
    \expanded{%
      \noexpand\newglossaryentry{cj\thecjcounter}{%
        name={\unexpanded{#1}},%
        sort={\thecjcounter},%
        description={p.~\noexpand\pageref{cj:\thecjcounter}},%
        user1={\thecjcounter},%
        user2={Conjecture},%
        user3={\thesection},%
        user4={\expandafter\unexpanded\expandafter{\@savedSecTitle}}%
      }%
    }%
    \protected@write\@auxout{}{\string\glsliveentry{cj\thecjcounter}}%
  \fi
}
\newcommand{\eeCj}{\end{conjecture}}
% LEMMA COMMANDS
% ---------------------------------------------------------------------
% UNIFORM RESULT WRAPPERS -- \beL, \beC, \beP, \beR share \beT's signature
%
%   \beX[long title][short title][key]  ...  \eeX
%
%   long title  = heading (math allowed)
%   short title = shown in the cross-reference; defaults to the key
%   key         = label key.  With a key the label is X:KEY and the heading
%                 carries the mnemonic [key], exactly as for \beT.  Without
%                 one the label falls back to the running counter, which is
%                 what every pre-existing call site in the book relies on.
%
% A bare \beX -- arguments commented out -- opens the plain environment and
% creates NO label and NO glossary entry.  That is the supported way to keep
% an individual result out of the catalogue.
% ---------------------------------------------------------------------
% #1 wrapper name  #2 environment  #3 label prefix  #4 counter  #5 type name
\newcommand{\@mkresult}[5]{%
  % A leading * suppresses the [key] mnemonic in both the heading and the
  % glossary name, while still making the label X:KEY and the entry:
  %     \beX*[long][short][key]
  \expandafter\NewDocumentCommand\csname be#1\endcsname{s O{} O{} o}{%
    \def\@temparg{##2}%
    \ifx\@temparg\@empty
      \begin{#2}%
    \else
      \let\@savedSecTitle\@currentsectionname%
      \refstepcounter{#4}%
      % With a key: mnemonic label X:KEY and the [key] tag in the heading.
      % Without one: fall back to the running counter, exactly as these
      % wrappers behaved before, so every legacy one-argument call site in
      % the book keeps its label.  The counter is also csname-safe, which a
      % title's first word is not -- titles beginning with math would
      % otherwise build \csname \Omega...\endcsname and fail.
      \IfValueTF{##4}{%
        \edef\@relbl{##4}%
        \IfBooleanTF{##1}{\let\@retag\@empty}{%
          \edef\@retag{\noexpand\ \noexpand\textup{[\noexpand\texttt{#3:##4}]}}}%
      }{%
        \edef\@relbl{\arabic{#4}}%
        \let\@retag\@empty
      }%
      \def\@tempshort{##3}%
      \ifx\@tempshort\@empty
        \expandafter\xdef\csname @sshort@#3\@relbl\endcsname{\@relbl}%
      \else
        \expandafter\def\csname @sshort@#3\@relbl\endcsname{##3}%
      \fi
      \begin{#2}[##2\ifshowkeys\@retag\fi]%
      \edef\@tmplbl{#3:\@relbl}\expandafter\label\expandafter{\@tmplbl}%
      \expanded{%
        \noexpand\newglossaryentry{#3\csname the#4\endcsname}{%
          name={\unexpanded{##2}%
            \IfBooleanF{##1}{\IfValueT{##4}{ [\noexpand\texttt{#3:##4}]}}},%
          sort={\csname the#2\endcsname},%
          description={p.~\noexpand\pageref{#3:\@relbl}},%
          user1={\csname the#2\endcsname},%
          user2={#5},%
          user3={\thesection},%
          user4={\expandafter\unexpanded\expandafter{\@savedSecTitle}}%
        }%
      }%
      \protected@write\@auxout{}{\string\glsliveentry{#3\csname the#4\endcsname}}%
      \protected@write\@auxout{}{\string\@thmtitle{\@tmplbl}{\unexpanded{##2}}}%
    \fi
  }%
  \expandafter\newcommand\csname ee#1\endcsname{\end{#2}}%
}
\@mkresult{L}{lemma}{l}{lemcounter}{Lemma}
% REMARK / COROLLARY / PROPOSITION -- same signature as \beL above
\@mkresult{R}{remark}{r}{remcounter}{Remark}
\@mkresult{C}{corollary}{c}{corcounter}{Corollary}
\@mkresult{P}{proposition}{p}{propcounter}{Proposition}
% OPEN-PROBLEM COMMAND -- generated with the same signature as \beT:
%   \beO[long title][short title][key]   ...   \eeO
% label o:KEY, or o:<counter> when no key is given; a bare \beO makes
% neither a label nor a glossary entry.
\@mkresult{O}{open}{o}{opcounter}{Problem}


% Section-level title tracker: updated only by \section, not \subsection.
% \@currentlabelname is overwritten by every \refstepcounter (including
% \subsection), so we save it separately here.
\newcommand{\@currentsectionname}{}
\AtBeginDocument{%
  \let\@epidcrn@orig@refstepcounter\refstepcounter
  \renewcommand{\refstepcounter}[1]{%
    \@epidcrn@orig@refstepcounter{#1}%
    \ifstrequal{#1}{section}{\global\let\@currentsectionname\@currentlabelname}{}%
  }%
}

% TABLE COMMANDS (glossaried) -- consistent naming with \beT, \beD, etc.
\newcommand{\beTa}[1][]{%
  \def\@temparg{#1}%
  \ifx\@temparg\@empty
    \begin{table}[htbp]%
  \else
    \let\@savedSecTitle\@currentsectionname%
    \refstepcounter{tabcounter}%
    \begin{table}[htbp]%
    \label{A:\thetabcounter}%
    \expanded{%
      \noexpand\newglossaryentry{A\thetabcounter}{%
        name={\unexpanded{#1}},%
        sort={\thetabcounter},%
        description={p.~\noexpand\pageref{A:\thetabcounter}},%
        user1={\thetabcounter},%
        user2={Table},%
        user3={\thesection},%
        user4={\expandafter\unexpanded\expandafter{\@savedSecTitle}}%
      }%
    }%
    \protected@write\@auxout{}{\string\glsliveentry{A\thetabcounter}}%
  \fi
}
\newcommand{\eeTa}{\end{table}}

% FIGURE COMMANDS (glossaried) -- consistent with \beT, \beR, \beTa, etc.
% Usage:
%   \beF[Glossary title]...\caption{Visible caption}\eeF  -> glossaried figure
%   \beF...\caption{Visible caption}\eeF                  -> plain [H] figure, no glossary
\newcommand{\beF}[1][]{%
  \def\@temparg{#1}%
  \ifx\@temparg\@empty
    \begin{figure}[H]%
  \else
    \begin{figure}[H]%
    \let\@savedSecTitle\@currentsectionname%
    \refstepcounter{figcounter}%
    \label{fig:\thefigcounter}%
    \expanded{%
      \noexpand\newglossaryentry{Fig\thefigcounter}{%
        name={\unexpanded{#1}},%
        sort={\thefigcounter},%
        description={p.~\noexpand\pageref{fig:\thefigcounter}},%
        user1={\thefigcounter},%
        user2={Figure},%
        user3={\thesection},%
        user4={\expandafter\unexpanded\expandafter{\@savedSecTitle}}%
      }%
    }%
    \protected@write\@auxout{}{\string\glsliveentry{Fig\thefigcounter}}%
  \fi
}
\newcommand{\eeF}{\end{figure}}
\makeatother
% LIVE-ENTRY REGISTER
% \glsaddall prints every entry the glossaries package currently knows, and
% that set includes entries restored from \jobname.glsdefs.  The package
% rewrites .glsdefs from all *currently defined* entries, so anything it
% restored is written straight back out: a deleted section kept reappearing in
% the catalogue for ever, and the .glsdefs had to be erased by hand.  The
% wrappers therefore record every entry they actually create in the .aux, and
% \printboth adds only those.  The .aux is rebuilt from scratch on each run, so
% this is self-healing: a removed result is gone after one further run.
\makeatletter
\newcommand{\@glslive}{}
\newcommand{\@glsaddlive}[1]{\ifglsentryexists{#1}{\glsadd{#1}}{}}
\newcommand{\glsliveentry}[1]{\g@addto@macro\@glslive{\@glsaddlive{#1}}}
\makeatother
% PRINT COMMAND
% No \clearpage here -- the catalogue continues on the current page.  This has
% been re-added by mistake several times; leave it out.
\makeatletter
\newcommand{\printboth}{%
  \@glslive   % add exactly the entries this document created last run
  \printglossary[style=unifiedstyle, title={0 \hspace{.22em} Definitions, Theorems, Examples and Open Problems}]%
}
\makeatother
% Named theorem reference: \tref{word} -> Thm~N(word-or-short)
% Shows stored short title if available, otherwise shows the key itself.
\newcommand{\tref}[1]{%
  Thm~\ref{t:#1}(\ifcsname @tshort@#1\endcsname\csname @tshort@#1\endcsname\else#1\fi)%
}

% Full-title theorem reference: \nref{t:KEY} prints the number and then the
% theorem's long title, so one writes "Theorem~\nref{t:blocks}" and gets
% "Theorem 3 ``Block factorization and siphons''".  Takes the complete label, like
% \ref, not the bare key like \tref.
%
% \beT writes each title to the .aux (see \@thmtitle below), so \nref resolves
% forward references from the second LaTeX pass, exactly as \ref does.  If the
% title is not yet known -- first pass, or a label that is not a \beT theorem
% -- it degrades to the number rather than producing ??.
% These two use @ in their names, and this point in the file is outside the
% \makeatletter region opened earlier (it closes at \makeatother above), so
% the region has to be reopened here.
\makeatletter
\newcommand{\@thmtitle}[2]{\global\@namedef{@tlong@#1}{#2}}
\newcommand{\nref}[1]{%
  \ref{#1}%
  \ifcsname @tlong@#1\endcsname
    \ ``\csname @tlong@#1\endcsname''%
  \fi
}
% \nm{t:KEY}: the stored long title of any \beT/\beL/\beP/\beR/\beC/\beD/\beE result,
% verbatim and without quotes, so item names follow the macro's first argument.
\newcommand{\nm}[1]{%
  \ifcsname @tlong@#1\endcsname
    \csname @tlong@#1\endcsname%
  \else
    \ref{#1}%
  \fi
}
\makeatother




%% ========================================
%% Other THEOREM ENVIRONMENTS
%% ========================================
\makeatletter
\@ifundefined{lemma}      {\newtheorem{lemma}      {Lemma}}      {}
\@ifundefined{proposition}{\newtheorem{proposition}{Proposition}} {}
\@ifundefined{corollary}  {\newtheorem{corollary}  {Corollary}}  {}
\@ifundefined{conjecture} {\newtheorem{conjecture} {Conjecture}}  {}
\@ifundefined{remark}     {\newtheorem{remark}     {Remark}}     {}
\@ifundefined{question}   {\newtheorem{question}   {Question}}   {}
\@ifundefined{assumption} {\newtheorem{assumption} {Assumption}} {}
\@ifundefined{hypothesis} {\newtheorem{hypothesis} {Hypothesis}} {}
\@ifundefined{fact}       {\newtheorem{fact}       {Fact}}       {}
\makeatother

% \citer: cite the relay-graph threshold theorem of AAHrel. If a document has
% \input the chapter/file containing t:relayRF (so the label resolves), point
% straight at it; otherwise fall back to a bare citation of AAHrel. Shared here
% (not per-paper) so it resolves consistently across every book chapter.
% \citeR: same idea, for the first theorem of AAHRah (t:Pstar).
\makeatletter
\providecommand{\citer}{%
  \@ifundefined{r@t:relayRF}%
    {\cite{AAHrel}}%
    {Theorem~\ref{t:relayRF} of \cite{AAHrel}}%
}
\providecommand{\citeR}{%
  \@ifundefined{r@t:Pstar}%
    {\cite{AAHRah}}%
    {Theorem~\ref{t:Pstar} of \cite{AAHRah}}%
}
\makeatother


% Abbreviations
\def\beQ{\begin{question}}\def\eeQ{\end{question}}
\def\beA{\begin{assumption}}\def\eeA{\end{assumption}}


%% ========================================

%% ========================================
%% COLORS & CUSTOM COMMANDS
%% ========================================
\definecolor{funccolor}{RGB}{25,25,112}
\definecolor{desccolor}{RGB}{64,64,64}
\newcommand{\func}[1]{\texttt{\textcolor{funccolor}{#1}}}
\def\eq{=}
\def\mm{-}

%% ========================================
%% MATHCAL SHORTHANDS
%% ========================================
% Project-wide convention: \mX for a capital X is always \mathcal X.
% Defined for the whole alphabet so the convention never fails to expand.
\newcommand{\mA}{\mathcal A}\newcommand{\mB}{\mathcal B}
\newcommand{\mC}{\mathcal C}\newcommand{\mD}{\mathcal D}
\newcommand{\mE}{\mathcal E}\newcommand{\mF}{\mathcal F}
\newcommand{\mG}{\mathcal G}\newcommand{\mH}{\mathcal H}
\newcommand{\mI}{\mathcal I}\newcommand{\mJ}{\mathcal J}
\newcommand{\mK}{\mathcal K}\newcommand{\mL}{\mathcal L}
\newcommand{\mM}{\mathcal M}\newcommand{\mN}{\mathcal N}
\newcommand{\mO}{\mathcal O}\newcommand{\mP}{\mathcal P}
\newcommand{\mQ}{\mathcal Q}\newcommand{\mR}{\mathcal R}
\newcommand{\mS}{\mathcal S}\newcommand{\mT}{\mathcal T}
\newcommand{\mU}{\mathcal U}\newcommand{\mV}{\mathcal V}
\newcommand{\mW}{\mathcal W}\newcommand{\mX}{\mathcal X}
\newcommand{\mY}{\mathcal Y}\newcommand{\mZ}{\mathcal Z}

%% ========================================
%% BASIC MATH ENVIRONMENTS
%% ========================================
\def\bep{\begin{pmatrix}}\def\eep{\end{pmatrix}}
\def\bev{\begin{vmatrix}}\def\eev{\end{vmatrix}}
\def\bea{\begin{eqnarray*}}\def\eea{\end{eqnarray*}}
\def\bc{\begin{cases}}\def\ec{\end{cases}}
\def\BEN{\begin{enumerate}}\def\EEN{\end{enumerate}}
\def\BI{\begin{itemize}}\def\EI{\end{itemize}}

\newcommand{\be}[1]{\begin{equation}\label{#1}}
\newcommand{\ee}{\end{equation}}
\newcommand{\beq}{\begin{eqnarray}}
\def\eeq{\end{eqnarray}}
\def\eqr{\eqref}\def\fr{\frac}\def\lbl{\label}
\def\Lra{\Longrightarrow}\def\non{\nonumber}\def\cNGM{\cite{Diek,Van,Van08}}
\def\Eq{\Leftrightarrow}\def\b{\beta}\def\uph{\red{UP TO HERE}}
%% ========================================
%% PROBABILITY & STATISTICS
%% ========================================
\renewcommand{\P}{\ensuremath{\mathbb{P}}}  % Note: overrides paragraph symbol
\newcommand{\E}{\ensuremath{\mathbb{E}}}
\newcommand{\ind}[1]{\ensuremath{\mathbbm{1}_{\left\{#1\right\}}}}
\DeclareMathOperator{\Cov}{Cov}
\DeclareMathOperator{\TV}{TV}
\newcommand{\rank}{\operatorname{rank}}

%% ========================================
%% CRN OPERATORS (organizations/ERCs/buffering/siphons framework)
%% ========================================
\def\clos{\mathrm{clos}}
\def\Supply{\mathrm{Supply}}
\def\req{\mathrm{req}}
\def\React{\mathrm{React}}
\def\Prod{\mathrm{Prod}}
\def\Gen{\mathrm{Gen}}
\def\In{\mathrm{In}}
\def\Out{\mathrm{Out}}
\def\Act{\mathrm{Act}}
\def\Min{\mathrm{Min}}

%% ========================================
%% LINE BREAKING
%% ========================================
% Last-resort stretch, used only on paragraphs TeX cannot otherwise break:
% a long \texttt{} token, a comma-separated inline math list (commas are not
% breakpoints in math), or a run of \cite/\ref boxes.  Without it those set as
% overfull \hbox and print into the margin; with it they set slightly loose,
% which is the lesser evil.  Does not affect paragraphs that already break.
\emergencystretch=3em

%% ========================================
%% NUMBER SETS
%% ========================================
\newcommand{\R}{\mathbb{R}}  % Real numbers
\newcommand{\C}{\mathbb{C}}  % Complex numbers
\newcommand{\N}{\mathbb{N}}
\newcommand{\Z}{\mathbb{Z}}
% \providecommand, not \newcommand: some papers define these in their own
% preamble after \input{preamble}, and must not clash.
\providecommand{\Q}{\mathbb{Q}}
\providecommand{\res}{\operatorname{res}}   % resultant; used by epiRes.tex
\newcommand{\MN}{\mathbb{N}}
\newcommand{\MZ}{\mathbb{Z}}
\newcommand{\realpos}{\mathbb{R}_{> 0}}
\newcommand{\realnneg}{\mathbb{R}_{\geq 0}}
\newcommand{\naturalpos}{\mathbb{N}_{>0}}
\newcommand{\Rplus}{\mathbb{R}_+}
\newcommand{\Rp}{\mathbb{R}_{\ge 0}}
\def\Rnne{\mathbb{R}_{\ge 0}}
\def\Se{\bar{S}} \def\Ie{\bar{I}}
\newcommand{\Rpos}{\mathbb{R}_{>0}}
\newcommand{\dd}{\,\mathrm d}
\newcommand{\e}{\mathrm e}
%% ========================================
%% GREEK LETTERS (2+ character abbreviations) greek
%% ========================================
\def\al{\alpha}\def\ba{\bff \alpha}
\def\w{\omega}\def\Lgn{\ga+\mu+\de}\def\om{\omega}
\def\de{\delta}\def\De{\Delta}\def\bde{\bff \delta}
\def\ep{\epsilon}\def\eps{\varepsilon}
\def\f{\varphi}  % Kept as requested
\def\ga{\gamma}\def\Ga{\Gamma}\def\G{\Gamma}\def\bg{\bff \gamma}
\def\up{\upsilon}  % NOT the breve accent: no accent commands anywhere.
                   % Named \up, not \u, so the LaTeX breve stays reachable.
\def\la{\lambda}\def\La{\Lambda}
\def\si{\sigma}\def\ka{\kappa}
\def\th{\theta}\def\vt{\vartheta}
\def\ze{\zeta}\def\Gap{\Gamma^{+}} \def\Gam{\Gamma^{-}}


%% ========================================
%% SPECIAL SHORT COMMANDS (kept as requested)
%% ========================================
\def\I{\infty}
\def\T{\widetilde}
\def\H{\hat}
\newcommand{\relint}{\operatorname{relint}}

%% ========================================
%% BOLD/VECTOR NOTATION
%% ========================================
\newcommand{\bff}[1]{{\mbox{\boldmath$#1$}}}
\newcommand\X{\boldsymbol{X}}
\def\x{\boldsymbol{x}}\def\bx{\bff x}\def\xe{\bar{\x}}
\newcommand\y{\boldsymbol{y}}\def\by{\boldsymbol{y}}\def\ye{\bar{\y}}
\def\yb{\y}\def\xb{\x}
\def\xy{\x-\y}
\newcommand\br{\boldsymbol{r}}\def\re{\bar{\r}}
\newcommand\bd{\boldsymbol{d}}
\newcommand\bQ{\boldsymbol{Q}}
\newcommand{\Yb}{{\mathbf{Y}}}
\newcommand{\Ab}{\mathbf{A}}
\newcommand{\zb}{\mathbf{z}}\def\bz{\bff 0}
\newcommand{\Zb}{\mathbf{Z}}
\newcommand{\bb}{\mathbf{b}}
\def\vbe{\vec {\bff \beta}}
\def\vb{\vec {\bff b}}
\def\v1{\vec {\bff 1}}
\def\vn{\vec \de}
\def\vz{\vec 0}
\def\vi{\overset{\rightarrow}{\mathsf i}}
\def\vv{\overset{\rightarrow}{v}}
\def\vr{\; \vec {\mathsf r}}
\newcommand{\va}{\bff a}
\providecommand{\vect}[1]{\boldsymbol{#1}}   % bold vector (moved from Gradostat.tex)
\providecommand{\one}{\mathbf 1}             % all-ones vector (moved from AliSF)
\providecommand{\spr}{s}                     % spectral bound  (moved from AliSF)
\providecommand{\Ke}{K_{\rm eff}}            % effective carrying capacity
% Critical right/left eigenvectors in the Hopf projection method (AAHOSN/New.tex)
\providecommand{\bq}{\mathbf{q}}
\providecommand{\bp}{\mathbf{p}}


%% ========================================
%% MATHCAL & MATHSCR
%% ========================================
% Mathcal
\newcommand\mc[1]{\mathcal{#1}}
\def\mC{{\mathcal C}}\def\mD{\mathcal D}\def\mE{\mathcal{E}}
\def\mF{\mathcal F}\def\mG{{\mathcal G}}\def\mH{{\mathcal H}}
\def\mI{\mathcal I}\def\mL{\mathcal{L}}\def\mM{{\mathcal M}}
\def\mN{{\mathcal N}}\def\mP{\mathcal{P}}\def\mR{{\mathcal R}}
\def\mS{{\mathcal S}}\def\mT{\mathcal{T}}\def\mV{{\mathcal V}}
\def\mK{\mathcal K}\def\mA{{\mathcal A}}\def\mU{{\mathcal U}}
% Symmetric bilinear and trilinear forms of the Taylor expansion, used in the
% first Lyapunov coefficient (AAHOSN/New.tex)
\providecommand{\Bform}{\mathcal{B}}
\providecommand{\Cform}{\mathcal{C}}

% Blackboard bold.  Kept here, not in the drivers: a \providecommand in a
% driver preamble runs before this file and silently wins, so the same name
% ends up meaning different things in different builds.
\def\IC{\mathbb{C}}\def\IN{\mathbb{N}}\def\IQ{\mathbb{Q}}
\def\IR{\mathbb{R}}\def\IZ{\mathbb{Z}}\def\IG{\mathbb{G}}
\def\II{\mathbb{I}}

% Mathscr
\newcommand{\Rs}{{\mathscr R}}
\newcommand{\Ss}{{\mathscr S}}
\newcommand{\Ns}{{\mathscr N}}
\newcommand{\Cs}{\mathscr C}


%% ========================================
%% OPERATORS & FUNCTIONS
%% ========================================
\DeclareMathOperator{\tr}{tr}
\DeclareMathOperator{\diag}{Diag}
\DeclareMathOperator{\Diag}{Diag}
\DeclareMathOperator{\intr}{\mathrm{Int}}
\DeclareMathOperator{\Inv}{Inv}
\DeclareMathOperator{\spec}{spec}
\DeclareMathOperator{\AUC}{AUC}               % area under the curve
\newcommand{\ii}{\mathrm{i}}                  % imaginary unit
\newcommand\op{\operatorname}
\newcommand{\supp}{\op{supp}}
\newcommand{\diff}{\mathop{}\mathopen{}\mathrm{d}}
\newcommand{\pdd}[1]{\frac{\partial}{\partial{#1}}}
\DeclareMathOperator{\cone}{cone}
\DeclareMathOperator{\proj}{proj}

%% ========================================
%% BRACKETS & NORMS
%% ========================================
\providecommand{\pr}[1]{\left(#1\right)}
\providecommand{\pp}[1]{\left[#1\right]}
\newcommand\croc[1]{\left\langle #1\right\rangle}
\newcommand{\norm}[1]{\|#1\|}
\newcommand{\inner}[2]{\langle #1, #2 \rangle}
% Short synonym of \inner, used by AAHOSN/New.tex as $\ip{u}{v}=\bar u^\top v$
\providecommand{\ip}[2]{\langle #1 , #2 \rangle}


%% ========================================
%% FORMATTING & UTILITIES
%% ========================================
\newcommand*{\Scale}[2][4]{\scalebox{#1}{$#2$}}
\newcommand{\te}[1]{\text{#1}}
\newcommand\steq[1]{\stackrel{\text{\rm #1.}}{=}}
\newcommand\blue[1]{\textcolor{blue}{#1}}
\newcommand{\red}{\textcolor[rgb]{1.00,0.00,0.00}}
\newcommand{\green}{\textcolor[rgb]{0.00,1.00,0.00}}

% Figure command
\newcommand{\figu}[3]{%
\begin{figure}[H]
\centering
\includegraphics[scale=#3]{#1}
\caption{#2\label{f:#1}}
\end{figure}
}

% Footnote
\long\def\symbolfootnote[#1]#2{%
\begingroup
\def\thefootnote{\fnsymbol{footnote}}\footnote[#1]{#2}%
\endgroup}
\def\fn{\symbolfootnote}

\newcommand{\dsum}{\displaystyle\sum}

%% ========================================
%% CRN & EPIDEMIOLOGY SPECIFIC
%% ========================================
% CRN notation
\def\SS{\mathcal S}
\def\CC{\mathcal C}
\def\RR{\mathcal R}

% Reproduction numbers
\def\Rto{R_{1}^{E_2}}
\def\Rot{R_{2}^{E_1}}
\def\Roc{R_{1,c}}
\def\Rtc{R_{2,c}}
\def\rd{r_{0}}
\def\sd{s_{0}}
\def\dd{D_{0}}
\def\vd{v_{0}}
\def\zd{z_{0}}
\def\se{\sE}
\def\xd{x_{0}}
%% ========================================
%% SPECIAL VALUES & EQUILIBRIA
%% ========================================
\def\eE{\bar{e}}
\def\wE{\bar{w}}
\def\sE{s_{*}}
\def\Me{\T M_{ee}}
\def\Rnne{\mathbb{R}_{\ge 0}}

\newcommand{\ones}{\mathbf 1}


%% ========================================
%% ABBREVIATIONS (alphabetical by category)
%% ========================================

% -- General mathematical terms --
\def\adm{admissible}\def\admy{admissibility}
%\def\and{antisymmetric}% disabled 2026-09-30: it overrode LaTeX's author separator \and (printed "antisymmetric" between authors)
\def\ass{assumption}
\def\cadlag{c\`adl\`ag}
\def\ch{characteristic polynomial}
\def\chp{characteristic polynomial}
\def\coe{coefficient}
\def\comp{computation}
\def\con{condition}
\def\coo{coordinate}
\def\corr{corresponding}
\def\dec{decomposition}
\def\den{denominator}
\def\dete{deterministic}
\def\eig{eigenvalue}\def\eval{eigenvalue}
\def\eigv{eigenvector}\def\eiv{eigenvector}\def\evec{eigenvector}
\def\elt{element}
\def\enn{essentially nonnegative/positivity preserving}
\def\equ{equation}
\def\etal{et al.}
\def\fe{for example}
\def\fno{from now on}
\def\for{\forall}
\def\form{formula}
\def\Fr{Furthermore, }
\def\frt{furthermore}
\def\fun{function}
\def\gene{generalization}
\def\heu{heuristic}
\def\how{however}\def\How{However, }
\def\ie{i.e. }
\def\im{\item}
\def\isd{is defined by}
\def\itc{it is easy to check that}
\def\itf{it follows that}\def\Itf{It follows that}
\def\ith{it holds that}
\def\itm{it may be checked that}\def\Itm{It may be checked that}
\def\ito{it turns out that}\def\Ito{It turns out that}\def\Itou{It turns out useful to}
\def\ma{mass action}\def\MAR{mass action representation}
\def\mat{matrix}
\def\mbc{may be checked to be}
\def\mbw{may be written as}\def\mubw{must be written as}
\def\mpr{more precisely, }\def\Mp{More precisely, }
\def\nc{non-catalytic}
\def\nne{non-negative}\def\nny{nonnegativity}
\def\no{\nonumber}
\def\noni{non-infectious}
\def\nota{notation}
\def\np{\newpage}
\def\Nt{Note that}
\def\nt{non trivial}\def\nz{non zero}
\def\oth{otherwise}
\def\Oth{On the other hand}
\def\pa{parameter}\def\para{parameter}
\def\pol{polynomial}
\def\prob{probability}
\def\Prf{{\bf  Proof:}}\def\prf{{\bf  Proof:}}
\def\prop{proportional to}
\def\qu{\quad}
\def\QED{\hfill {$\square$}\goodbreak \medskip}
\def\rep{representation}
\def\resp{respectively}
\def\satd{satisfied}\def\satg{satisfying}\def\sats{satisfies}\def\saty{satisfy}
\def\ssec{\subsection}\def\sssec{\subsubsection}
\def\sef{see for example}
\def\sksy{skew-symmetric}
\def\sub{substrate}
\def\sys{system}
\def\var{variable}
\def\wk{well-known}
\def\wl{without loss of generality}\def\wlo{w.l.o.g.}
\def\wmc{We may check that}\def\Wmc{We may check that}
\def\wmr{we may rewrite}
\def\wms{we must show that}
\def\wrt{with respect to}

% -- CRN & Chemistry --
\newcommand\CRN{chemical reaction network}
\newcommand\CRNs{chemical reaction networks}
\newcommand\RN{reaction network}
\def\ACR{absolute concentration robustness}
\def\AMSD{acyclic siphon decomposition}
\def\DAE{differential algebraic equations}
\def\FHJ{Feinberg-Horn-Jackson}
\def\GMAK{generalized mass action kinetics}
\def\MM{Michaelis-Menten}
\def\sm{stoichiometric matrix}\def\Sm{stoichiometric matrix}
\def\WRZD{weakly reversible, zero deficiency}\def\aut{autocatalytic}

% -- Epidemiology --
\def\brn{basic reproduction number}\def\bzn{basic replacement number}
\def\com{compartment}
\def\dfe{disease free equilibrium}\def\DFE{disease free equilibrium}
\def\dwt{dwell times}\def\edt{expected dwell times}
\def\EM{epidemic model}
\def\endp{endemic fixed point}\def\enp{endemic point}\def\EE{endemic point}
\def\ME{mathematical epidemiology}
\def\MEE{eco-epidemiology}\def\EEP{eco-epidemiological}
\def\ece{eco-epidemiological model}
\def\mgm{multi group model}
\def\mtS{multi-strain model}\def\mSt{multi-strain model}
\def\mSi{minimal siphon}
\def\NGM{next generation matrix}
\def\nim{new infections matrix}
\def\qem{quadratic epidemiological model}
\def\repF{invasion functions}
\def\SPF{SIR-PH-FA model}
\def\sus{susceptible}

% -- Graph theory & Linear algebra --
\def\jac{Jacobi matrix}
\def\jin{Jacobian matrix of the invasion vector field with respect to the invasion variables}
\def\jrn{Jacobian matrix of the resident vector field with respect to the resident variables}
\def\jxn{Jacobian matrix of the $\x$ vector field with respect to the $\x$ variables}
\def\Lm{Laplacian matrix}
\def\MTT{matrix tree theorem}
\def\SGT{spectral graph theory}
\def\SMW{Sherman–Morrison–Woodbury}

% -- Stability & Dynamics --
\def\absC{absorbtion CTMC}
\def\BTb{Bogdanov-Takens bifurcation}
\def\CEP{competitive exclusion partition}
\def\CTMC{continuous time Markov chain}
\def\ds{differential system}
\def\fp{fixed point}\def\bfp{boundary fixed point}\def\fbp{fixed boundary point}
\def\fbl{fixed boundary line}\def\fif{forward invariant face}
\def\gas{global asymptotic stability}\def\GAS{globally asymptotically stable}
\def\GASS{\text{GASS}}
\def\Hb{Hopf bifurcation}
\def\Hur{Hurwitz}\def\MH{Metzler-Hurwitz}\def\RH{Routh-Hurwitz}
\def\IS{\text{IS}}
\def\LAS{locally asymptotically stable}
\def\LdS{Lyapunov diagonally stable}
\def\Lf{Lyapunov function}\def\LyV{Lyapunov Volterra type function}\def\RLF{robust Lyapunov function}
\def\LyM{Lyapunov Malkin stable}
\def\PF{Perron-Frobenius}
\def\spfp{strictly positive fixed point}
\def\sta{stability}
\def\STP{sharp threshold property}\def\STT{sharp threshold theorem}

% -- Models (with \varphi explicitly instead of \f) --
\def\AB{$(A,B)$ Arino-Brauer epidemic models}
\def\ABa{$(A,\bbe,\ba,\varphi)$ balanced bilinear model}
\def\ABP{%$(A,\bbe,P,\varphi)$
\\ balanced bilinear model}
\def\AMG{Anderson May Gupta}
\newcommand\BIN{biological interaction network}
\def\Dt{Descartes type}
\def\gLV{generalized Lotka-Volterra model}
\def\IG{\text{invasion graph}}\def\IGMS{interaction graph on minimal siphons}
\def\LCP{\text{linear complementarity problem}}
\def\LV{Lotka-Volterra}\def\LVK{Lotka-Volterra-Kolmogorov}

% -- Splitting & Structure --
\def\reg{regular splitting}\def\regS{{\bf regular splitting}}
\def\RUR{rational univariate representation}
\def\bds{block-diagonal structure}\def\NGMs{next generation matrices}

% -- Miscellaneous --
\def\Ma{{Mathematica}}
\def\Mm{Max–Min}
\def\num{numerator}
\def\rap{rational approximation} \def\DSR{directed species-reaction bipartite graph}\def\PCP{polynomial complementarity problem}\def\rat{rational substitution}
\def\QSS{quasi steady state} \def\RHS{right hand side} \def\LHS{left hand side}
\def\OSN{online social network}\def\VdD{Van den Driessche}
\def\wli{$\omega$-limit} \def\rk{rich kinetics}\def\sJ{symbolic Jacobian}
\def\cV{\cite{Vas,VasSta,BSV,GSV}} \def\spf{semi-parametric semi-functional}
\def\CRW{Capasso-Ruan-Wang} \def\ABP{$(A,\bbe,P,\f)$ balanced bilinear model}
\def\ABa{%$(A,\bbe,\ba,\f)$ \\
rank-one balanced bilinear model}
\def\bbe{B%{\bff \beta}
} \def\vbe{\vec {\bff \beta}
}
\def\vb{\vec {\bff b}}
\newcommand{\Kh}{\hat{K}}          % K with hat, first mention only
\newcommand{\wS}{w^{\sigma}}       % right Perron J^\perp, full block
\newcommand{\piS}{\vec \pi^{\sigma}}    % left Perron J^\perp, full block
\newcommand{\ws}{w^{\sigma}}       % right Perron J^\perp, sub-block
\newcommand{\pis}{\pi^{\sigma}}    % left Perron J^\perp, sub-block
\newcommand{\hwS}{\hat{w}^{\sigma}}   % right Perron K, full block
\newcommand{\hpiS}{\hat{\pi}^{\sigma}} % left Perron K, full block
\newcommand{\hws}{\hat{w}^{\sigma}}   % right Perron K, sub-block
\newcommand{\hpis}{\hat{\pi}^{\sigma}} % left Perron K, sub-block
\newcommand{\twS}{\tilde{w}^{\sigma}}  % right Perron \tilde K, full block
\newcommand{\tpiS}{\tilde{\pi}^{\sigma}} % left Perron \tilde K
\newcommand{\tws}{\tilde{w}^{\sigma}}
\newcommand{\tpis}{\tilde{\pi}^{\sigma}}
\newcommand{\Mat}[2]{\mathbb{R}_{\ge 0}^{\,#1\times #2}}
\def\tJ{J_\sigma^\perp} \def\lPe{left Perron eigenvector}
% \G is \Gamma (defined at :795).  The Volterra function keeps its own name,
% so that \G no longer has two conflicting meanings in this file.
\def\GV{G}  % Volterra G(u)=u-1-\ln u
\def\PVL{Perron-Volterra Lyapunov minimal siphon ansatz}
\newcommand{\istar}[1]{i_{#1}^{*}}
%\newcommand{\sE}{s_E}
\newcommand{\vE}{v_E}
\def\PV{Perron-Volterra candidate Lyapunov function}
\def\GLV{generalized Lotka-Volterra}
\def\MERN{mathematical epidemiology reaction networks}
\def\CMM{contraction metric method}
\newcommand{\sign}{\operatorname{sign}}
\newcommand{\Var}{\operatorname{Var}} 

% ======================================================================
% IDEAS SHARED BETWEEN THE ABSTRACT AND THE MAIN RESULTS SECTION
%
% Write each idea ONCE, in the abstract; the Main Results section recalls it,
% so the two can never drift apart.
%
%   in the abstract        \ideaitem{1}{The first-order reaction terms ...}
%   in Main Results        \ideahead{1}
%
% \ideaitem stores the text globally and typesets
%   \item \hyperlink{idea:1}{<text>}
% \ideahead typesets
%   \item\hypertarget{idea:1}{}\textbf{<the same text>}
%
% The abstract must come before the Main Results input, which is the normal
% order. If a heading has to read differently from the abstract sentence (for
% instance a long abstract item summarised in one line), override it with the
% optional argument: \ideahead[short heading]{6}.
% ======================================================================
\makeatletter
\newcommand{\ideaitem}[2]{%
  \expandafter\gdef\csname @idea@#1\endcsname{#2}%
  \item\hyperlink{idea:#1}{#2}}
\newcommand{\ideahead}[2][]{%
  \item\hypertarget{idea:#2}{}%
  \def\@ideatmp{#1}%
  \ifx\@ideatmp\@empty
    \@ifundefined{@idea@#2}{%
      \PackageWarning{preamble}{Idea #2 is used in Main Results but was never
        defined by \string\ideaitem\space in the abstract}%
      \textbf{[idea #2 undefined]}%
    }{\textbf{\begingroup\def\ideaterm##1##2{##2}\def\idearest##1{}\@nameuse{@idea@#2}\endgroup}}%
  \else
    \textbf{#1}%
  \fi}
% \ideaterm{<hypertarget>}{<words>}: inside an abstract idea the words become an
% invisible link (the document sets \hypersetup{hidelinks}) jumping to the place
% in Main Results that defines them; inside the Main Results heading itself,
% where the target sits a few lines below, \ideahead strips the link and prints
% the words plainly. Mark the destination with \hypertarget{<name>}{} on the line
% before the display, which adds no equation number and shows nothing.
\newcommand{\ideaterm}[2]{\hyperlink{#1}{#2}}
% \idearest{<text>}: the part of an abstract idea after its first sentence; printed in
% the abstract, omitted from the Main Results heading.
\newcommand{\idearest}[1]{#1}
\makeatother

% ======================================================================
% SHORT-PAPER SWITCH
% \shorttrue in a master that assembles only part of the companion paper
% (Barrat/SipArxiv.tex). \nsh{...} drops text that cites results living in the
% omitted sections, so the short build has no dangling references while the
% companion paper keeps the text verbatim.
% ======================================================================
\newif\ifshort \shortfalse
\newcommand{\nsh}[1]{\ifshort\else#1\fi}
% \shortalt{<short build>}{<companion paper>}: use when a sentence needs
% different wording in the two builds, not merely a deletion.
\newcommand{\shortalt}[2]{\ifshort#1\else#2\fi}
% \Jtr, \JtrS{\Sigma}: transverse Jacobian. J^\perp here; SipArxiv.tex redefines them as J_{\mathrm{tra}}, J_{\mathrm{tra},\Sigma}
\newcommand{\Jtr}{J^\perp}
\newcommand{\JtrS}[1]{J_{#1}^\perp}
.
\makeatletter
\AtBeginDocument{\@ifundefined{hypersetup}{}{\hypersetup{linktoc=all}}}
\makeatother
\usepackage{graphicx,float,fancybox}
\usepackage{epsfig,epstopdf,tikz}
\usepackage[sans]{dsfont}
\usetikzlibrary{positioning,arrows,arrows.meta,calc}
% Configure listings for Mathematica code
\lstdefinestyle{mathematica}{
 language=Mathematica,
 basicstyle=\small\ttfamily,
 backgroundcolor=\color{gray!10},
 frame=single,
 breaklines=true,
 commentstyle=\color{green!60!black},
 keywordstyle=\color{blue},
 stringstyle=\color{red},
 showstringspaces=false,
 tabsize=2
}
\usepackage{tcolorbox}
% Colored boxes for key results
\newtcolorbox{keyresult}[1][]{
 colback=blue!5!white,
 colframe=blue!75!black,
 title=#1,
 fonttitle=\bfseries
}
\newtcolorbox{examplebox}[1][]{
 colback=green!5!white,
 colframe=green!75!black,
 title=#1,
 fonttitle=\bfseries
}


\makeatletter
\@ifundefined{theorem}   {\newtheorem{theorem}   {Theorem}}   {}
\@ifundefined{definition}{\newtheorem{definition}{Definition}} {}
\@ifundefined{open}      {\newtheorem{open}      {Problem}}   {}
\@ifundefined{example}   {\newtheorem{example}   {Example}}   {}
\makeatother

% COUNTER DEFINITIONS
\newcounter{tabcounter}
\newcounter{defcounter}
\newcounter{opcounter}
\newcounter{thmcounter}
\newcounter{excounter}
\newcounter{ascounter}
\newcounter{cjcounter}
\newcounter{lemcounter}
\newcounter{remcounter}
\newcounter{corcounter}
\newcounter{propcounter}
\newcounter{figcounter}

% UNIFIED GLOSSARY STYLE
\newglossarystyle{unifiedstyle}{%
  \setglossarystyle{list}%
  \renewcommand*{\glossarypreamble}{\gdef\glslastsec{}}%
  \renewcommand{\glossentry}[2]{%
    \edef\glscursec{\glsentryuseriii{##1}}%
    \ifx\glscursec\glslastsec\else
      \item[]\textbf{\S\,\glscursec\enspace\glsentryuseriv{##1}}%
      \global\let\glslastsec\glscursec
    \fi
    \item[\glsentryuserii{##1} \glsentryuseri{##1}:] ``\glossentryname{##1}'' \glsentrydesc{##1}%
  }%
}

\makeatletter
% First-word extractor (used by \beT, \beD and \beXa for the default key)
% \@firstwordof hello world \relax  ->  hello
\def\@firstwordof#1 #2\relax{#1}
% First three letters, used for the heading mnemonic of a theorem given no
% explicit key.  Padding with \@empty keeps it safe for first words shorter
% than three letters: \@empty vanishes under \edef, whereas \relax would
% survive into the label and is illegal inside \csname (it broke every
% two-letter title, e.g. "ME ...", with "Missing \endcsname inserted").
\def\@takethree#1#2#3#4\@nil{#1#2#3}

% EXAMPLE COMMAND -- same signature and behaviour as \beT and \beD.
% Usage: \beXa*[long title][short title][key]
%   star  = suppresses the printed "[key]" tag after the title (as for \beT*
%   and \beD*); the label and glossary entry are unaffected.
%   key = optional label key override; the label becomes e:KEY, and without a
%   key it is e: plus the first three letters of the title's first word.  A
%   mnemonic label lets an example be cited as "Example~\ref{e:KEY}" instead
%   of by section number, which is what the citing prose needs.
\NewDocumentCommand{\beXa}{s O{} O{} o}{%
  \def\@temparg{#2}%
  \ifx\@temparg\@empty
    \begin{example}%
  \else
    \let\@savedSecTitle\@currentsectionname%
    \refstepcounter{excounter}%
    \IfValueTF{#4}{%
      \edef\@exlblkey{#4}%
    }{%
      \edef\@exfirstword{\@firstwordof#2 \relax}%
      \edef\@exlblkey{%
        \expandafter\@takethree\@exfirstword\@empty\@empty\@empty\@nil}%
    }%
    \def\@tempshort{#3}%
    \ifx\@tempshort\@empty
      \expandafter\xdef\csname @xshort@\@exlblkey\endcsname{\@exlblkey}%
    \else
      \expandafter\def\csname @xshort@\@exlblkey\endcsname{#3}%
    \fi
    \IfBooleanTF{#1}{\let\@extag\@empty}{%
      \edef\@extag{\noexpand\ \noexpand\textup{[\noexpand\texttt{e:\@exlblkey}]}}}%
    \begin{example}[#2\ifshowkeys\@extag\fi]%
    \label{e:\@exlblkey}%
    % record the long title so \nref{e:KEY} prints it, as for theorems
    \protected@write\@auxout{}{%
      \string\@thmtitle{e:\@exlblkey}{\unexpanded{#2}}}%
    \expanded{%
      \noexpand\newglossaryentry{e\theexcounter}{%
        name={\unexpanded{#2}%
          \IfBooleanF{#1}{ [\noexpand\texttt{e:\@exlblkey}]}},%
        sort={\theexample},%
        description={p.~\noexpand\pageref{e:\@exlblkey}},%
        user1={\theexample},%
        user2={Example},%
        user3={\thesection},%
        user4={\expandafter\unexpanded\expandafter{\@savedSecTitle}}%
      }%
    }%
    \protected@write\@auxout{}{\string\glsliveentry{e\theexcounter}}%
  \fi
}
\newcommand{\eeXa}{\end{example}}

% DEFINITION COMMAND -- same signature and behaviour as \beT below.
% Usage: \beD[long title][short title][key]
%   long title  = definition heading (plain text only)
%   short title = shown in \dref{word}; defaults to the key
%   key         = optional label key override; the label becomes d:KEY
% Without a key the label is d: plus the first three letters of the title's
% first word, so it is a mnemonic rather than the definition's number: a
% numeric label would silently change whenever a definition is inserted.
\NewDocumentCommand{\beD}{s O{} O{} o}{%
  \def\@temparg{#2}%
  \ifx\@temparg\@empty
    \begin{definition}%
  \else
    \let\@savedSecTitle\@currentsectionname%
    \refstepcounter{defcounter}%
    \IfValueTF{#4}{%
      \edef\@deflblkey{#4}%
    }{%
      \edef\@deffirstword{\@firstwordof#2 \relax}%
      \edef\@deflblkey{%
        \expandafter\@takethree\@deffirstword\@empty\@empty\@empty\@nil}%
    }%
    \def\@tempshort{#3}%
    \ifx\@tempshort\@empty
      \expandafter\xdef\csname @dshort@\@deflblkey\endcsname{\@deflblkey}%
    \else
      \expandafter\def\csname @dshort@\@deflblkey\endcsname{#3}%
    \fi
    \IfBooleanTF{#1}{\let\@deftag\@empty}{%
      \edef\@deftag{\noexpand\ \noexpand\textup{[\noexpand\texttt{d:\@deflblkey}]}}}%
    \begin{definition}[#2\ifshowkeys\@deftag\fi]%
    \label{d:\@deflblkey}%
    % record the long title so \nref{d:KEY} prints it, as for theorems
    \protected@write\@auxout{}{%
      \string\@thmtitle{d:\@deflblkey}{\unexpanded{#2}}}%
    \expanded{%
      \noexpand\newglossaryentry{d\thedefcounter}{%
        name={\unexpanded{#2}%
          \IfBooleanF{#1}{ [\noexpand\texttt{d:\@deflblkey}]}},%
        sort={\thedefinition},%
        description={p.~\noexpand\pageref{d:\@deflblkey}},%
        user1={\thedefinition},%
        user2={Definition},%
        user3={\thesection},%
        user4={\expandafter\unexpanded\expandafter{\@savedSecTitle}}%
      }%
    }%
    \protected@write\@auxout{}{\string\glsliveentry{d\thedefcounter}}%
  \fi
}
\newcommand{\eeD}{\end{definition}}

% \showkeysfalse in a driver hides the [t:KEY] mnemonics in result headings
% (labels and glossary entries are unchanged); used for submission builds.
\newif\ifshowkeys \showkeystrue
% \commonredtrue in a driver that inputs basicFiles/commonRed: tool.tex then
% cites its reduction results instead of restating them.
\newif\ifcommonred \commonredfalse
% TheorEM COMMAND
% Usage: \beT[long title][short title][key]
%   long title  = theorem heading (plain text only)
%   short title = shown in \tref{word}; may contain math; defaults to first word of long title
%   key         = optional label key override; label becomes t:KEY (avoids first-word collisions)
% Label is t:WORD where WORD = first word of long title (or KEY if given); use \tref{WORD} to reference.
\NewDocumentCommand{\beT}{s O{} O{} o}{%
  \def\@temparg{#2}%
  \ifx\@temparg\@empty
    \begin{theorem}%
  \else
    \let\@savedSecTitle\@currentsectionname%
    \refstepcounter{thmcounter}%
    % \@thmtag is the mnemonic shown in the heading; every theorem gets one.
    % With an explicit key the key itself is shown.  Without one the key
    % defaults to the title's first word, and showing that whole word would
    % merely repeat the title ("Dissipativity [Dissipativity]"), so only its
    % first three letters are shown: "Dissipativity [Dis]".  The label is
    % unaffected -- it still uses the full \@thmlblkey.
    \IfValueTF{#4}{%
      \edef\@thmlblkey{#4}%
    }{%
      \edef\@thmfirstword{\@firstwordof#2 \relax}%
      \edef\@thmlblkey{%
        \expandafter\@takethree\@thmfirstword\@empty\@empty\@empty\@nil}%
    }%
    \IfBooleanTF{#1}{\let\@thmtag\@empty}{%
      \edef\@thmtag{\noexpand\ \noexpand\textup{[\noexpand\texttt{t:\@thmlblkey}]}}}%
    \def\@tempshort{#3}%
    \ifx\@tempshort\@empty
      \expandafter\xdef\csname @tshort@\@thmlblkey\endcsname{\@thmlblkey}%
    \else
      \expandafter\def\csname @tshort@\@thmlblkey\endcsname{#3}%
    \fi
    \begin{theorem}[#2\ifshowkeys\@thmtag\fi]%
    \label{t:\@thmlblkey}%
    % record the long title for \nref; via the .aux so forward refs resolve
    \protected@write\@auxout{}{%
      \string\@thmtitle{t:\@thmlblkey}{\unexpanded{#2}}}%
    \expanded{%
      \noexpand\newglossaryentry{t\thethmcounter}{%
        name={\unexpanded{#2}%
          \IfBooleanF{#1}{ [\noexpand\texttt{t:\@thmlblkey}]}},%
        % sort/user1 must be the number the document actually prints, i.e. the
        % theorem environment's counter, not thmcounter.  thmcounter advances
        % only on \beT with a title, so a bare \beT (or \beT with an
        % empty title) makes the two drift and the glossary cites a wrong
        % number.  thmcounter is still fine as the entry key: it is unique.
        sort={\thetheorem},%
        description={p.~\noexpand\pageref{t:\@thmlblkey}},%
        user1={\thetheorem},%
        user2={Theorem},%
        user3={\thesection},%
        user4={\expandafter\unexpanded\expandafter{\@savedSecTitle}}%
      }%
    }%
    \protected@write\@auxout{}{\string\glsliveentry{t\thethmcounter}}%
  \fi
}
\newcommand{\eeT}{\end{theorem}}

% CROSS-PART RELATION PARAGRAPH
% Usage: \xpart{text}  -- place immediately BEFORE \beT, never inside it
%   (anything inside \beT...\eeT would be swallowed into the glossary entry).
% The paragraph is typeset only when Part I is actually present, i.e. in the
% combined book build; in the standalone builds (GranadaDSABNS/A.tex,
% AAHRah-Submission/AAHRah.tex) the label part:glv is undefined and the whole
% paragraph disappears rather than leaving ?? cross-references. Same
% graceful-degradation principle as \GASref in GASt.tex. Resolves on the
% second LaTeX pass, since r@part:glv comes from the .aux file.
% (Already inside the \makeatletter region opened above, so no \makeatletter
% pair here -- adding one would close that region early.)
\newcommand{\xpart}[1]{\@ifundefined{r@part:glv}{}{\paragraph{Relation with Parts~I--II.}#1}}

% \bookref{label}{precise text}{fallback text}: prints the precise cross-reference
% when the label resolves (combined book build), and the fallback text otherwise
% (standalone paper builds, where book-only appendix labels are absent). Same
% graceful-degradation role as \GASref in GASt.tex, but usable in any file since it
% lives in this preamble (loaded by every driver) -- unlike \GASref, which is defined in a
% content file and is therefore undefined in standalone builds that omit that file.
\newcommand{\bookref}[3]{\@ifundefined{r@#1}{#3}{#2}}

% \PVpart and \RBpart: the two parts that a partial build may omit.
% They print "Part~N" in any build that contains the part, and a naming phrase
% in a build that does not, so an excerpt shows no ?? .  Each already supplies
% the word "Part", so write "\PVpart develops ...", not "Part~\PVpart".
\newcommand{\PVpart}{\bookref{part:aahrah}%
  {Part~\ref{part:aahrah}}{the Perron--Volterra part of the full book}}
\newcommand{\RBpart}{\bookref{part:recbrowne}%
  {Part~\ref{part:recbrowne}}{the immuno-virology part of the full book}}

% OPEN PROBLEM COMMANDS
% ASSUMPTION COMMANDS (glossaried)
\newcommand{\beAs}[1][]{%
  \def\@temparg{#1}%
  \ifx\@temparg\@empty
    \begin{assumption}%
  \else
    \let\@savedSecTitle\@currentsectionname%
    \refstepcounter{ascounter}%
    \begin{assumption}[#1]%
    \label{as:\theascounter}%
    \expanded{%
      \noexpand\newglossaryentry{as\theascounter}{%
        name={\unexpanded{#1}},%
        sort={\theascounter},%
        description={p.~\noexpand\pageref{as:\theascounter}},%
        user1={\theascounter},%
        user2={Assumption},%
        user3={\thesection},%
        user4={\expandafter\unexpanded\expandafter{\@savedSecTitle}}%
      }%
    }%
    \protected@write\@auxout{}{\string\glsliveentry{as\theascounter}}%
  \fi
}
\newcommand{\eeAs}{\end{assumption}}
% CONJECTURE COMMANDS (glossaried)
\newcommand{\beCj}[1][]{%
  \def\@temparg{#1}%
  \ifx\@temparg\@empty
    \begin{conjecture}%
  \else
    \let\@savedSecTitle\@currentsectionname%
    \refstepcounter{cjcounter}%
    \begin{conjecture}[#1]%
    \label{cj:\thecjcounter}%
    \expanded{%
      \noexpand\newglossaryentry{cj\thecjcounter}{%
        name={\unexpanded{#1}},%
        sort={\thecjcounter},%
        description={p.~\noexpand\pageref{cj:\thecjcounter}},%
        user1={\thecjcounter},%
        user2={Conjecture},%
        user3={\thesection},%
        user4={\expandafter\unexpanded\expandafter{\@savedSecTitle}}%
      }%
    }%
    \protected@write\@auxout{}{\string\glsliveentry{cj\thecjcounter}}%
  \fi
}
\newcommand{\eeCj}{\end{conjecture}}
% LEMMA COMMANDS
% ---------------------------------------------------------------------
% UNIFORM RESULT WRAPPERS -- \beL, \beC, \beP, \beR share \beT's signature
%
%   \beX[long title][short title][key]  ...  \eeX
%
%   long title  = heading (math allowed)
%   short title = shown in the cross-reference; defaults to the key
%   key         = label key.  With a key the label is X:KEY and the heading
%                 carries the mnemonic [key], exactly as for \beT.  Without
%                 one the label falls back to the running counter, which is
%                 what every pre-existing call site in the book relies on.
%
% A bare \beX -- arguments commented out -- opens the plain environment and
% creates NO label and NO glossary entry.  That is the supported way to keep
% an individual result out of the catalogue.
% ---------------------------------------------------------------------
% #1 wrapper name  #2 environment  #3 label prefix  #4 counter  #5 type name
\newcommand{\@mkresult}[5]{%
  % A leading * suppresses the [key] mnemonic in both the heading and the
  % glossary name, while still making the label X:KEY and the entry:
  %     \beX*[long][short][key]
  \expandafter\NewDocumentCommand\csname be#1\endcsname{s O{} O{} o}{%
    \def\@temparg{##2}%
    \ifx\@temparg\@empty
      \begin{#2}%
    \else
      \let\@savedSecTitle\@currentsectionname%
      \refstepcounter{#4}%
      % With a key: mnemonic label X:KEY and the [key] tag in the heading.
      % Without one: fall back to the running counter, exactly as these
      % wrappers behaved before, so every legacy one-argument call site in
      % the book keeps its label.  The counter is also csname-safe, which a
      % title's first word is not -- titles beginning with math would
      % otherwise build \csname \Omega...\endcsname and fail.
      \IfValueTF{##4}{%
        \edef\@relbl{##4}%
        \IfBooleanTF{##1}{\let\@retag\@empty}{%
          \edef\@retag{\noexpand\ \noexpand\textup{[\noexpand\texttt{#3:##4}]}}}%
      }{%
        \edef\@relbl{\arabic{#4}}%
        \let\@retag\@empty
      }%
      \def\@tempshort{##3}%
      \ifx\@tempshort\@empty
        \expandafter\xdef\csname @sshort@#3\@relbl\endcsname{\@relbl}%
      \else
        \expandafter\def\csname @sshort@#3\@relbl\endcsname{##3}%
      \fi
      \begin{#2}[##2\ifshowkeys\@retag\fi]%
      \edef\@tmplbl{#3:\@relbl}\expandafter\label\expandafter{\@tmplbl}%
      \expanded{%
        \noexpand\newglossaryentry{#3\csname the#4\endcsname}{%
          name={\unexpanded{##2}%
            \IfBooleanF{##1}{\IfValueT{##4}{ [\noexpand\texttt{#3:##4}]}}},%
          sort={\csname the#2\endcsname},%
          description={p.~\noexpand\pageref{#3:\@relbl}},%
          user1={\csname the#2\endcsname},%
          user2={#5},%
          user3={\thesection},%
          user4={\expandafter\unexpanded\expandafter{\@savedSecTitle}}%
        }%
      }%
      \protected@write\@auxout{}{\string\glsliveentry{#3\csname the#4\endcsname}}%
      \protected@write\@auxout{}{\string\@thmtitle{\@tmplbl}{\unexpanded{##2}}}%
    \fi
  }%
  \expandafter\newcommand\csname ee#1\endcsname{\end{#2}}%
}
\@mkresult{L}{lemma}{l}{lemcounter}{Lemma}
% REMARK / COROLLARY / PROPOSITION -- same signature as \beL above
\@mkresult{R}{remark}{r}{remcounter}{Remark}
\@mkresult{C}{corollary}{c}{corcounter}{Corollary}
\@mkresult{P}{proposition}{p}{propcounter}{Proposition}
% OPEN-PROBLEM COMMAND -- generated with the same signature as \beT:
%   \beO[long title][short title][key]   ...   \eeO
% label o:KEY, or o:<counter> when no key is given; a bare \beO makes
% neither a label nor a glossary entry.
\@mkresult{O}{open}{o}{opcounter}{Problem}


% Section-level title tracker: updated only by \section, not \subsection.
% \@currentlabelname is overwritten by every \refstepcounter (including
% \subsection), so we save it separately here.
\newcommand{\@currentsectionname}{}
\AtBeginDocument{%
  \let\@epidcrn@orig@refstepcounter\refstepcounter
  \renewcommand{\refstepcounter}[1]{%
    \@epidcrn@orig@refstepcounter{#1}%
    \ifstrequal{#1}{section}{\global\let\@currentsectionname\@currentlabelname}{}%
  }%
}

% TABLE COMMANDS (glossaried) -- consistent naming with \beT, \beD, etc.
\newcommand{\beTa}[1][]{%
  \def\@temparg{#1}%
  \ifx\@temparg\@empty
    \begin{table}[htbp]%
  \else
    \let\@savedSecTitle\@currentsectionname%
    \refstepcounter{tabcounter}%
    \begin{table}[htbp]%
    \label{A:\thetabcounter}%
    \expanded{%
      \noexpand\newglossaryentry{A\thetabcounter}{%
        name={\unexpanded{#1}},%
        sort={\thetabcounter},%
        description={p.~\noexpand\pageref{A:\thetabcounter}},%
        user1={\thetabcounter},%
        user2={Table},%
        user3={\thesection},%
        user4={\expandafter\unexpanded\expandafter{\@savedSecTitle}}%
      }%
    }%
    \protected@write\@auxout{}{\string\glsliveentry{A\thetabcounter}}%
  \fi
}
\newcommand{\eeTa}{\end{table}}

% FIGURE COMMANDS (glossaried) -- consistent with \beT, \beR, \beTa, etc.
% Usage:
%   \beF[Glossary title]...\caption{Visible caption}\eeF  -> glossaried figure
%   \beF...\caption{Visible caption}\eeF                  -> plain [H] figure, no glossary
\newcommand{\beF}[1][]{%
  \def\@temparg{#1}%
  \ifx\@temparg\@empty
    \begin{figure}[H]%
  \else
    \begin{figure}[H]%
    \let\@savedSecTitle\@currentsectionname%
    \refstepcounter{figcounter}%
    \label{fig:\thefigcounter}%
    \expanded{%
      \noexpand\newglossaryentry{Fig\thefigcounter}{%
        name={\unexpanded{#1}},%
        sort={\thefigcounter},%
        description={p.~\noexpand\pageref{fig:\thefigcounter}},%
        user1={\thefigcounter},%
        user2={Figure},%
        user3={\thesection},%
        user4={\expandafter\unexpanded\expandafter{\@savedSecTitle}}%
      }%
    }%
    \protected@write\@auxout{}{\string\glsliveentry{Fig\thefigcounter}}%
  \fi
}
\newcommand{\eeF}{\end{figure}}
\makeatother
% LIVE-ENTRY REGISTER
% \glsaddall prints every entry the glossaries package currently knows, and
% that set includes entries restored from \jobname.glsdefs.  The package
% rewrites .glsdefs from all *currently defined* entries, so anything it
% restored is written straight back out: a deleted section kept reappearing in
% the catalogue for ever, and the .glsdefs had to be erased by hand.  The
% wrappers therefore record every entry they actually create in the .aux, and
% \printboth adds only those.  The .aux is rebuilt from scratch on each run, so
% this is self-healing: a removed result is gone after one further run.
\makeatletter
\newcommand{\@glslive}{}
\newcommand{\@glsaddlive}[1]{\ifglsentryexists{#1}{\glsadd{#1}}{}}
\newcommand{\glsliveentry}[1]{\g@addto@macro\@glslive{\@glsaddlive{#1}}}
\makeatother
% PRINT COMMAND
% No \clearpage here -- the catalogue continues on the current page.  This has
% been re-added by mistake several times; leave it out.
\makeatletter
\newcommand{\printboth}{%
  \@glslive   % add exactly the entries this document created last run
  \printglossary[style=unifiedstyle, title={0 \hspace{.22em} Definitions, Theorems, Examples and Open Problems}]%
}
\makeatother
% Named theorem reference: \tref{word} -> Thm~N(word-or-short)
% Shows stored short title if available, otherwise shows the key itself.
\newcommand{\tref}[1]{%
  Thm~\ref{t:#1}(\ifcsname @tshort@#1\endcsname\csname @tshort@#1\endcsname\else#1\fi)%
}

% Full-title theorem reference: \nref{t:KEY} prints the number and then the
% theorem's long title, so one writes "Theorem~\nref{t:blocks}" and gets
% "Theorem 3 ``Block factorization and siphons''".  Takes the complete label, like
% \ref, not the bare key like \tref.
%
% \beT writes each title to the .aux (see \@thmtitle below), so \nref resolves
% forward references from the second LaTeX pass, exactly as \ref does.  If the
% title is not yet known -- first pass, or a label that is not a \beT theorem
% -- it degrades to the number rather than producing ??.
% These two use @ in their names, and this point in the file is outside the
% \makeatletter region opened earlier (it closes at \makeatother above), so
% the region has to be reopened here.
\makeatletter
\newcommand{\@thmtitle}[2]{\global\@namedef{@tlong@#1}{#2}}
\newcommand{\nref}[1]{%
  \ref{#1}%
  \ifcsname @tlong@#1\endcsname
    \ ``\csname @tlong@#1\endcsname''%
  \fi
}
% \nm{t:KEY}: the stored long title of any \beT/\beL/\beP/\beR/\beC/\beD/\beE result,
% verbatim and without quotes, so item names follow the macro's first argument.
\newcommand{\nm}[1]{%
  \ifcsname @tlong@#1\endcsname
    \csname @tlong@#1\endcsname%
  \else
    \ref{#1}%
  \fi
}
\makeatother




%% ========================================
%% Other THEOREM ENVIRONMENTS
%% ========================================
\makeatletter
\@ifundefined{lemma}      {\newtheorem{lemma}      {Lemma}}      {}
\@ifundefined{proposition}{\newtheorem{proposition}{Proposition}} {}
\@ifundefined{corollary}  {\newtheorem{corollary}  {Corollary}}  {}
\@ifundefined{conjecture} {\newtheorem{conjecture} {Conjecture}}  {}
\@ifundefined{remark}     {\newtheorem{remark}     {Remark}}     {}
\@ifundefined{question}   {\newtheorem{question}   {Question}}   {}
\@ifundefined{assumption} {\newtheorem{assumption} {Assumption}} {}
\@ifundefined{hypothesis} {\newtheorem{hypothesis} {Hypothesis}} {}
\@ifundefined{fact}       {\newtheorem{fact}       {Fact}}       {}
\makeatother

% \citer: cite the relay-graph threshold theorem of AAHrel. If a document has
% \input the chapter/file containing t:relayRF (so the label resolves), point
% straight at it; otherwise fall back to a bare citation of AAHrel. Shared here
% (not per-paper) so it resolves consistently across every book chapter.
% \citeR: same idea, for the first theorem of AAHRah (t:Pstar).
\makeatletter
\providecommand{\citer}{%
  \@ifundefined{r@t:relayRF}%
    {\cite{AAHrel}}%
    {Theorem~\ref{t:relayRF} of \cite{AAHrel}}%
}
\providecommand{\citeR}{%
  \@ifundefined{r@t:Pstar}%
    {\cite{AAHRah}}%
    {Theorem~\ref{t:Pstar} of \cite{AAHRah}}%
}
\makeatother


% Abbreviations
\def\beQ{\begin{question}}\def\eeQ{\end{question}}
\def\beA{\begin{assumption}}\def\eeA{\end{assumption}}


%% ========================================

%% ========================================
%% COLORS & CUSTOM COMMANDS
%% ========================================
\definecolor{funccolor}{RGB}{25,25,112}
\definecolor{desccolor}{RGB}{64,64,64}
\newcommand{\func}[1]{\texttt{\textcolor{funccolor}{#1}}}
\def\eq{=}
\def\mm{-}

%% ========================================
%% MATHCAL SHORTHANDS
%% ========================================
% Project-wide convention: \mX for a capital X is always \mathcal X.
% Defined for the whole alphabet so the convention never fails to expand.
\newcommand{\mA}{\mathcal A}\newcommand{\mB}{\mathcal B}
\newcommand{\mC}{\mathcal C}\newcommand{\mD}{\mathcal D}
\newcommand{\mE}{\mathcal E}\newcommand{\mF}{\mathcal F}
\newcommand{\mG}{\mathcal G}\newcommand{\mH}{\mathcal H}
\newcommand{\mI}{\mathcal I}\newcommand{\mJ}{\mathcal J}
\newcommand{\mK}{\mathcal K}\newcommand{\mL}{\mathcal L}
\newcommand{\mM}{\mathcal M}\newcommand{\mN}{\mathcal N}
\newcommand{\mO}{\mathcal O}\newcommand{\mP}{\mathcal P}
\newcommand{\mQ}{\mathcal Q}\newcommand{\mR}{\mathcal R}
\newcommand{\mS}{\mathcal S}\newcommand{\mT}{\mathcal T}
\newcommand{\mU}{\mathcal U}\newcommand{\mV}{\mathcal V}
\newcommand{\mW}{\mathcal W}\newcommand{\mX}{\mathcal X}
\newcommand{\mY}{\mathcal Y}\newcommand{\mZ}{\mathcal Z}

%% ========================================
%% BASIC MATH ENVIRONMENTS
%% ========================================
\def\bep{\begin{pmatrix}}\def\eep{\end{pmatrix}}
\def\bev{\begin{vmatrix}}\def\eev{\end{vmatrix}}
\def\bea{\begin{eqnarray*}}\def\eea{\end{eqnarray*}}
\def\bc{\begin{cases}}\def\ec{\end{cases}}
\def\BEN{\begin{enumerate}}\def\EEN{\end{enumerate}}
\def\BI{\begin{itemize}}\def\EI{\end{itemize}}

\newcommand{\be}[1]{\begin{equation}\label{#1}}
\newcommand{\ee}{\end{equation}}
\newcommand{\beq}{\begin{eqnarray}}
\def\eeq{\end{eqnarray}}
\def\eqr{\eqref}\def\fr{\frac}\def\lbl{\label}
\def\Lra{\Longrightarrow}\def\non{\nonumber}\def\cNGM{\cite{Diek,Van,Van08}}
\def\Eq{\Leftrightarrow}\def\b{\beta}\def\uph{\red{UP TO HERE}}
%% ========================================
%% PROBABILITY & STATISTICS
%% ========================================
\renewcommand{\P}{\ensuremath{\mathbb{P}}}  % Note: overrides paragraph symbol
\newcommand{\E}{\ensuremath{\mathbb{E}}}
\newcommand{\ind}[1]{\ensuremath{\mathbbm{1}_{\left\{#1\right\}}}}
\DeclareMathOperator{\Cov}{Cov}
\DeclareMathOperator{\TV}{TV}
\newcommand{\rank}{\operatorname{rank}}

%% ========================================
%% CRN OPERATORS (organizations/ERCs/buffering/siphons framework)
%% ========================================
\def\clos{\mathrm{clos}}
\def\Supply{\mathrm{Supply}}
\def\req{\mathrm{req}}
\def\React{\mathrm{React}}
\def\Prod{\mathrm{Prod}}
\def\Gen{\mathrm{Gen}}
\def\In{\mathrm{In}}
\def\Out{\mathrm{Out}}
\def\Act{\mathrm{Act}}
\def\Min{\mathrm{Min}}

%% ========================================
%% LINE BREAKING
%% ========================================
% Last-resort stretch, used only on paragraphs TeX cannot otherwise break:
% a long \texttt{} token, a comma-separated inline math list (commas are not
% breakpoints in math), or a run of \cite/\ref boxes.  Without it those set as
% overfull \hbox and print into the margin; with it they set slightly loose,
% which is the lesser evil.  Does not affect paragraphs that already break.
\emergencystretch=3em

%% ========================================
%% NUMBER SETS
%% ========================================
\newcommand{\R}{\mathbb{R}}  % Real numbers
\newcommand{\C}{\mathbb{C}}  % Complex numbers
\newcommand{\N}{\mathbb{N}}
\newcommand{\Z}{\mathbb{Z}}
% \providecommand, not \newcommand: some papers define these in their own
% preamble after % Shared preamble for the GLVbook papers.  This is the single canonical copy;
% a driver one level below GLVbook reaches it as
%
%     \input{../basicFiles/preamble}
%
% What stays in the driver, because it differs from paper to paper:
% \documentclass, font choice, \title, \author, \date, and any macro genuinely
% local to one paper.  Everything else belongs here.  The file is
% self-contained: it does not read any other file.

\usepackage[utf8]{inputenc}
\usepackage[novbox]{pdfsync}   % required: SyncTeX alone does not work here

% Page geometry, set once here for every paper that loads this file.  A driver
% that needs different margins says \geometry{...}; a second
% \usepackage[...]{geometry} raises an option clash and must not be used.
\makeatletter
\@ifundefined{AAdefNnogeometry}{\usepackage[a4paper,margin=2cm]{geometry}}{}
\makeatother

% mathematics
\usepackage{amsmath,amssymb,amsthm,mathtools,bm,mathrsfs,bbm}
\usepackage{centernot}

% floats, tables, lists
\usepackage{graphicx,float}
\makeatletter
\@ifpackageloaded{subfig}{}{\usepackage{subcaption}}
\makeatother
\usepackage{booktabs,longtable,threeparttable,array,multirow}
\usepackage{enumitem}

% colour, code, algorithms
\usepackage{xcolor,listings,tcolorbox}
\makeatletter
\@ifundefined{c@algorithm}{\usepackage{algorithm}}{}
\makeatother
\usepackage{algorithmicx,algpseudocode}

% diagrams
\usepackage{tikz,tikz-cd,pgfplots}
\usetikzlibrary{graphs,decorations.pathmorphing,decorations.markings,
arrows.meta,automata,positioning}
\tikzcdset{every label/.append style = {font = \small}}
\pgfplotsset{compat=1.18}

% chemistry
\usepackage[version=4]{mhchem}

\usepackage{microtype}
\usepackage{lineno}
\usepackage{nomencl}
\makenomenclature

% hyperref first, then bookmark, then glossaries: glossaries must follow
% hyperref or the entry links are built against the wrong anchors.
\makeatletter
\@ifundefined{hyperrefpackage}{%
  \@ifundefined{AAdefNnohyperref}{\usepackage{hyperref}}{}%
}{}
\@ifundefined{hyperrefpackage}{\usepackage{bookmark}}{}
\makeatother
\usepackage[toc,sort=use]{glossaries}
\makeglossaries

% Book/standalone switch.  Material that the book already covers in its
% prelude, but that a standalone paper still has to state, is wrapped as
%     \ifbookbuild\else ... \fi
% Drivers default to standalone; bookM.tex says \bookbuildtrue.
\newif\ifbookbuild
\bookbuildfalse

% Part switches for the gradostat paper.  Both default to true; a driver that
% wants the short submission says \madecfalse and/or \slowfastfalse, and every
% cross-reference into the dropped part disappears with it.
%   \ifmadec     Part III, the Lagasquie--Madec model and its break-even geometry
%   \ifslowfast  Part IV, everything carrying the small parameter \eps
\newif\ifmadec     \madectrue
\newif\ifslowfast  \slowfasttrue

% Part switch for the Barrat/simplicial-contagion project.  Defaults to
% true (the hypernetwork-synchrony part is included); the journal-paper
% driver says \hnetpartfalse to drop it and every cross-reference into it.
\newif\ifhnetpart  \hnetparttrue

%% ========================================
%% PACKAGES
%% ========================================
\RequirePackage{etoolbox}
\usepackage{makecell}
\usepackage[novbox]{pdfsync}% obsolete; SyncTeX is built into modern engines
\usepackage{amssymb, amsfonts,mathtools,bbm,eurosym,graphicx,latexsym, amsmath,amsthm,times,nccmath,url}
\usepackage{listings,verbatim,cases,tabularx}
% hyperref is deliberately NOT in the line above.  Journal classes often
% load it themselves, with a long option list, from a begindocument hook
% (mdpi.cls does: \DeclareHookRule{begindocument}{hyperref}{before}{mdpi}).
% Loading it here first, without options, makes that later load raise
% "Option clash for package hyperref".  So load it only if nobody else
% has, or will:
%   - \@ifpackageloaded  catches a class that loaded it already;
%   - \hyperrefpackage    is mdpi.cls's own deferred loader;
%   - \AAdefNnohyperref   lets any other driver opt out by defining it
%                         before \input{preamble}.
\makeatletter
\@ifpackageloaded{hyperref}{}{%
  \@ifundefined{hyperrefpackage}{%
    \@ifundefined{AAdefNnohyperref}{\usepackage{hyperref}}{}%
  }{}%
}
\makeatother
% linktoc=all makes a table-of-contents page number a link as well as the
% entry text; hyperref's default, linktoc=section, links only the text.
% Deferred to begin-document and guarded: by the block above hyperref may
% be loaded later (mdpi's own hook) or not at all (\AAdefNnohyperref), and
% \hypersetup must exist before it is called.  A driver that wants the
% default back can say \hypersetup{linktoc=section} after \input{preamble}.
\makeatletter
\AtBeginDocument{\@ifundefined{hypersetup}{}{\hypersetup{linktoc=all}}}
\makeatother
\usepackage{graphicx,float,fancybox}
\usepackage{epsfig,epstopdf,tikz}
\usepackage[sans]{dsfont}
\usetikzlibrary{positioning,arrows,arrows.meta,calc}
% Configure listings for Mathematica code
\lstdefinestyle{mathematica}{
 language=Mathematica,
 basicstyle=\small\ttfamily,
 backgroundcolor=\color{gray!10},
 frame=single,
 breaklines=true,
 commentstyle=\color{green!60!black},
 keywordstyle=\color{blue},
 stringstyle=\color{red},
 showstringspaces=false,
 tabsize=2
}
\usepackage{tcolorbox}
% Colored boxes for key results
\newtcolorbox{keyresult}[1][]{
 colback=blue!5!white,
 colframe=blue!75!black,
 title=#1,
 fonttitle=\bfseries
}
\newtcolorbox{examplebox}[1][]{
 colback=green!5!white,
 colframe=green!75!black,
 title=#1,
 fonttitle=\bfseries
}


\makeatletter
\@ifundefined{theorem}   {\newtheorem{theorem}   {Theorem}}   {}
\@ifundefined{definition}{\newtheorem{definition}{Definition}} {}
\@ifundefined{open}      {\newtheorem{open}      {Problem}}   {}
\@ifundefined{example}   {\newtheorem{example}   {Example}}   {}
\makeatother

% COUNTER DEFINITIONS
\newcounter{tabcounter}
\newcounter{defcounter}
\newcounter{opcounter}
\newcounter{thmcounter}
\newcounter{excounter}
\newcounter{ascounter}
\newcounter{cjcounter}
\newcounter{lemcounter}
\newcounter{remcounter}
\newcounter{corcounter}
\newcounter{propcounter}
\newcounter{figcounter}

% UNIFIED GLOSSARY STYLE
\newglossarystyle{unifiedstyle}{%
  \setglossarystyle{list}%
  \renewcommand*{\glossarypreamble}{\gdef\glslastsec{}}%
  \renewcommand{\glossentry}[2]{%
    \edef\glscursec{\glsentryuseriii{##1}}%
    \ifx\glscursec\glslastsec\else
      \item[]\textbf{\S\,\glscursec\enspace\glsentryuseriv{##1}}%
      \global\let\glslastsec\glscursec
    \fi
    \item[\glsentryuserii{##1} \glsentryuseri{##1}:] ``\glossentryname{##1}'' \glsentrydesc{##1}%
  }%
}

\makeatletter
% First-word extractor (used by \beT, \beD and \beXa for the default key)
% \@firstwordof hello world \relax  ->  hello
\def\@firstwordof#1 #2\relax{#1}
% First three letters, used for the heading mnemonic of a theorem given no
% explicit key.  Padding with \@empty keeps it safe for first words shorter
% than three letters: \@empty vanishes under \edef, whereas \relax would
% survive into the label and is illegal inside \csname (it broke every
% two-letter title, e.g. "ME ...", with "Missing \endcsname inserted").
\def\@takethree#1#2#3#4\@nil{#1#2#3}

% EXAMPLE COMMAND -- same signature and behaviour as \beT and \beD.
% Usage: \beXa*[long title][short title][key]
%   star  = suppresses the printed "[key]" tag after the title (as for \beT*
%   and \beD*); the label and glossary entry are unaffected.
%   key = optional label key override; the label becomes e:KEY, and without a
%   key it is e: plus the first three letters of the title's first word.  A
%   mnemonic label lets an example be cited as "Example~\ref{e:KEY}" instead
%   of by section number, which is what the citing prose needs.
\NewDocumentCommand{\beXa}{s O{} O{} o}{%
  \def\@temparg{#2}%
  \ifx\@temparg\@empty
    \begin{example}%
  \else
    \let\@savedSecTitle\@currentsectionname%
    \refstepcounter{excounter}%
    \IfValueTF{#4}{%
      \edef\@exlblkey{#4}%
    }{%
      \edef\@exfirstword{\@firstwordof#2 \relax}%
      \edef\@exlblkey{%
        \expandafter\@takethree\@exfirstword\@empty\@empty\@empty\@nil}%
    }%
    \def\@tempshort{#3}%
    \ifx\@tempshort\@empty
      \expandafter\xdef\csname @xshort@\@exlblkey\endcsname{\@exlblkey}%
    \else
      \expandafter\def\csname @xshort@\@exlblkey\endcsname{#3}%
    \fi
    \IfBooleanTF{#1}{\let\@extag\@empty}{%
      \edef\@extag{\noexpand\ \noexpand\textup{[\noexpand\texttt{e:\@exlblkey}]}}}%
    \begin{example}[#2\ifshowkeys\@extag\fi]%
    \label{e:\@exlblkey}%
    % record the long title so \nref{e:KEY} prints it, as for theorems
    \protected@write\@auxout{}{%
      \string\@thmtitle{e:\@exlblkey}{\unexpanded{#2}}}%
    \expanded{%
      \noexpand\newglossaryentry{e\theexcounter}{%
        name={\unexpanded{#2}%
          \IfBooleanF{#1}{ [\noexpand\texttt{e:\@exlblkey}]}},%
        sort={\theexample},%
        description={p.~\noexpand\pageref{e:\@exlblkey}},%
        user1={\theexample},%
        user2={Example},%
        user3={\thesection},%
        user4={\expandafter\unexpanded\expandafter{\@savedSecTitle}}%
      }%
    }%
    \protected@write\@auxout{}{\string\glsliveentry{e\theexcounter}}%
  \fi
}
\newcommand{\eeXa}{\end{example}}

% DEFINITION COMMAND -- same signature and behaviour as \beT below.
% Usage: \beD[long title][short title][key]
%   long title  = definition heading (plain text only)
%   short title = shown in \dref{word}; defaults to the key
%   key         = optional label key override; the label becomes d:KEY
% Without a key the label is d: plus the first three letters of the title's
% first word, so it is a mnemonic rather than the definition's number: a
% numeric label would silently change whenever a definition is inserted.
\NewDocumentCommand{\beD}{s O{} O{} o}{%
  \def\@temparg{#2}%
  \ifx\@temparg\@empty
    \begin{definition}%
  \else
    \let\@savedSecTitle\@currentsectionname%
    \refstepcounter{defcounter}%
    \IfValueTF{#4}{%
      \edef\@deflblkey{#4}%
    }{%
      \edef\@deffirstword{\@firstwordof#2 \relax}%
      \edef\@deflblkey{%
        \expandafter\@takethree\@deffirstword\@empty\@empty\@empty\@nil}%
    }%
    \def\@tempshort{#3}%
    \ifx\@tempshort\@empty
      \expandafter\xdef\csname @dshort@\@deflblkey\endcsname{\@deflblkey}%
    \else
      \expandafter\def\csname @dshort@\@deflblkey\endcsname{#3}%
    \fi
    \IfBooleanTF{#1}{\let\@deftag\@empty}{%
      \edef\@deftag{\noexpand\ \noexpand\textup{[\noexpand\texttt{d:\@deflblkey}]}}}%
    \begin{definition}[#2\ifshowkeys\@deftag\fi]%
    \label{d:\@deflblkey}%
    % record the long title so \nref{d:KEY} prints it, as for theorems
    \protected@write\@auxout{}{%
      \string\@thmtitle{d:\@deflblkey}{\unexpanded{#2}}}%
    \expanded{%
      \noexpand\newglossaryentry{d\thedefcounter}{%
        name={\unexpanded{#2}%
          \IfBooleanF{#1}{ [\noexpand\texttt{d:\@deflblkey}]}},%
        sort={\thedefinition},%
        description={p.~\noexpand\pageref{d:\@deflblkey}},%
        user1={\thedefinition},%
        user2={Definition},%
        user3={\thesection},%
        user4={\expandafter\unexpanded\expandafter{\@savedSecTitle}}%
      }%
    }%
    \protected@write\@auxout{}{\string\glsliveentry{d\thedefcounter}}%
  \fi
}
\newcommand{\eeD}{\end{definition}}

% \showkeysfalse in a driver hides the [t:KEY] mnemonics in result headings
% (labels and glossary entries are unchanged); used for submission builds.
\newif\ifshowkeys \showkeystrue
% \commonredtrue in a driver that inputs basicFiles/commonRed: tool.tex then
% cites its reduction results instead of restating them.
\newif\ifcommonred \commonredfalse
% TheorEM COMMAND
% Usage: \beT[long title][short title][key]
%   long title  = theorem heading (plain text only)
%   short title = shown in \tref{word}; may contain math; defaults to first word of long title
%   key         = optional label key override; label becomes t:KEY (avoids first-word collisions)
% Label is t:WORD where WORD = first word of long title (or KEY if given); use \tref{WORD} to reference.
\NewDocumentCommand{\beT}{s O{} O{} o}{%
  \def\@temparg{#2}%
  \ifx\@temparg\@empty
    \begin{theorem}%
  \else
    \let\@savedSecTitle\@currentsectionname%
    \refstepcounter{thmcounter}%
    % \@thmtag is the mnemonic shown in the heading; every theorem gets one.
    % With an explicit key the key itself is shown.  Without one the key
    % defaults to the title's first word, and showing that whole word would
    % merely repeat the title ("Dissipativity [Dissipativity]"), so only its
    % first three letters are shown: "Dissipativity [Dis]".  The label is
    % unaffected -- it still uses the full \@thmlblkey.
    \IfValueTF{#4}{%
      \edef\@thmlblkey{#4}%
    }{%
      \edef\@thmfirstword{\@firstwordof#2 \relax}%
      \edef\@thmlblkey{%
        \expandafter\@takethree\@thmfirstword\@empty\@empty\@empty\@nil}%
    }%
    \IfBooleanTF{#1}{\let\@thmtag\@empty}{%
      \edef\@thmtag{\noexpand\ \noexpand\textup{[\noexpand\texttt{t:\@thmlblkey}]}}}%
    \def\@tempshort{#3}%
    \ifx\@tempshort\@empty
      \expandafter\xdef\csname @tshort@\@thmlblkey\endcsname{\@thmlblkey}%
    \else
      \expandafter\def\csname @tshort@\@thmlblkey\endcsname{#3}%
    \fi
    \begin{theorem}[#2\ifshowkeys\@thmtag\fi]%
    \label{t:\@thmlblkey}%
    % record the long title for \nref; via the .aux so forward refs resolve
    \protected@write\@auxout{}{%
      \string\@thmtitle{t:\@thmlblkey}{\unexpanded{#2}}}%
    \expanded{%
      \noexpand\newglossaryentry{t\thethmcounter}{%
        name={\unexpanded{#2}%
          \IfBooleanF{#1}{ [\noexpand\texttt{t:\@thmlblkey}]}},%
        % sort/user1 must be the number the document actually prints, i.e. the
        % theorem environment's counter, not thmcounter.  thmcounter advances
        % only on \beT with a title, so a bare \beT (or \beT with an
        % empty title) makes the two drift and the glossary cites a wrong
        % number.  thmcounter is still fine as the entry key: it is unique.
        sort={\thetheorem},%
        description={p.~\noexpand\pageref{t:\@thmlblkey}},%
        user1={\thetheorem},%
        user2={Theorem},%
        user3={\thesection},%
        user4={\expandafter\unexpanded\expandafter{\@savedSecTitle}}%
      }%
    }%
    \protected@write\@auxout{}{\string\glsliveentry{t\thethmcounter}}%
  \fi
}
\newcommand{\eeT}{\end{theorem}}

% CROSS-PART RELATION PARAGRAPH
% Usage: \xpart{text}  -- place immediately BEFORE \beT, never inside it
%   (anything inside \beT...\eeT would be swallowed into the glossary entry).
% The paragraph is typeset only when Part I is actually present, i.e. in the
% combined book build; in the standalone builds (GranadaDSABNS/A.tex,
% AAHRah-Submission/AAHRah.tex) the label part:glv is undefined and the whole
% paragraph disappears rather than leaving ?? cross-references. Same
% graceful-degradation principle as \GASref in GASt.tex. Resolves on the
% second LaTeX pass, since r@part:glv comes from the .aux file.
% (Already inside the \makeatletter region opened above, so no \makeatletter
% pair here -- adding one would close that region early.)
\newcommand{\xpart}[1]{\@ifundefined{r@part:glv}{}{\paragraph{Relation with Parts~I--II.}#1}}

% \bookref{label}{precise text}{fallback text}: prints the precise cross-reference
% when the label resolves (combined book build), and the fallback text otherwise
% (standalone paper builds, where book-only appendix labels are absent). Same
% graceful-degradation role as \GASref in GASt.tex, but usable in any file since it
% lives in this preamble (loaded by every driver) -- unlike \GASref, which is defined in a
% content file and is therefore undefined in standalone builds that omit that file.
\newcommand{\bookref}[3]{\@ifundefined{r@#1}{#3}{#2}}

% \PVpart and \RBpart: the two parts that a partial build may omit.
% They print "Part~N" in any build that contains the part, and a naming phrase
% in a build that does not, so an excerpt shows no ?? .  Each already supplies
% the word "Part", so write "\PVpart develops ...", not "Part~\PVpart".
\newcommand{\PVpart}{\bookref{part:aahrah}%
  {Part~\ref{part:aahrah}}{the Perron--Volterra part of the full book}}
\newcommand{\RBpart}{\bookref{part:recbrowne}%
  {Part~\ref{part:recbrowne}}{the immuno-virology part of the full book}}

% OPEN PROBLEM COMMANDS
% ASSUMPTION COMMANDS (glossaried)
\newcommand{\beAs}[1][]{%
  \def\@temparg{#1}%
  \ifx\@temparg\@empty
    \begin{assumption}%
  \else
    \let\@savedSecTitle\@currentsectionname%
    \refstepcounter{ascounter}%
    \begin{assumption}[#1]%
    \label{as:\theascounter}%
    \expanded{%
      \noexpand\newglossaryentry{as\theascounter}{%
        name={\unexpanded{#1}},%
        sort={\theascounter},%
        description={p.~\noexpand\pageref{as:\theascounter}},%
        user1={\theascounter},%
        user2={Assumption},%
        user3={\thesection},%
        user4={\expandafter\unexpanded\expandafter{\@savedSecTitle}}%
      }%
    }%
    \protected@write\@auxout{}{\string\glsliveentry{as\theascounter}}%
  \fi
}
\newcommand{\eeAs}{\end{assumption}}
% CONJECTURE COMMANDS (glossaried)
\newcommand{\beCj}[1][]{%
  \def\@temparg{#1}%
  \ifx\@temparg\@empty
    \begin{conjecture}%
  \else
    \let\@savedSecTitle\@currentsectionname%
    \refstepcounter{cjcounter}%
    \begin{conjecture}[#1]%
    \label{cj:\thecjcounter}%
    \expanded{%
      \noexpand\newglossaryentry{cj\thecjcounter}{%
        name={\unexpanded{#1}},%
        sort={\thecjcounter},%
        description={p.~\noexpand\pageref{cj:\thecjcounter}},%
        user1={\thecjcounter},%
        user2={Conjecture},%
        user3={\thesection},%
        user4={\expandafter\unexpanded\expandafter{\@savedSecTitle}}%
      }%
    }%
    \protected@write\@auxout{}{\string\glsliveentry{cj\thecjcounter}}%
  \fi
}
\newcommand{\eeCj}{\end{conjecture}}
% LEMMA COMMANDS
% ---------------------------------------------------------------------
% UNIFORM RESULT WRAPPERS -- \beL, \beC, \beP, \beR share \beT's signature
%
%   \beX[long title][short title][key]  ...  \eeX
%
%   long title  = heading (math allowed)
%   short title = shown in the cross-reference; defaults to the key
%   key         = label key.  With a key the label is X:KEY and the heading
%                 carries the mnemonic [key], exactly as for \beT.  Without
%                 one the label falls back to the running counter, which is
%                 what every pre-existing call site in the book relies on.
%
% A bare \beX -- arguments commented out -- opens the plain environment and
% creates NO label and NO glossary entry.  That is the supported way to keep
% an individual result out of the catalogue.
% ---------------------------------------------------------------------
% #1 wrapper name  #2 environment  #3 label prefix  #4 counter  #5 type name
\newcommand{\@mkresult}[5]{%
  % A leading * suppresses the [key] mnemonic in both the heading and the
  % glossary name, while still making the label X:KEY and the entry:
  %     \beX*[long][short][key]
  \expandafter\NewDocumentCommand\csname be#1\endcsname{s O{} O{} o}{%
    \def\@temparg{##2}%
    \ifx\@temparg\@empty
      \begin{#2}%
    \else
      \let\@savedSecTitle\@currentsectionname%
      \refstepcounter{#4}%
      % With a key: mnemonic label X:KEY and the [key] tag in the heading.
      % Without one: fall back to the running counter, exactly as these
      % wrappers behaved before, so every legacy one-argument call site in
      % the book keeps its label.  The counter is also csname-safe, which a
      % title's first word is not -- titles beginning with math would
      % otherwise build \csname \Omega...\endcsname and fail.
      \IfValueTF{##4}{%
        \edef\@relbl{##4}%
        \IfBooleanTF{##1}{\let\@retag\@empty}{%
          \edef\@retag{\noexpand\ \noexpand\textup{[\noexpand\texttt{#3:##4}]}}}%
      }{%
        \edef\@relbl{\arabic{#4}}%
        \let\@retag\@empty
      }%
      \def\@tempshort{##3}%
      \ifx\@tempshort\@empty
        \expandafter\xdef\csname @sshort@#3\@relbl\endcsname{\@relbl}%
      \else
        \expandafter\def\csname @sshort@#3\@relbl\endcsname{##3}%
      \fi
      \begin{#2}[##2\ifshowkeys\@retag\fi]%
      \edef\@tmplbl{#3:\@relbl}\expandafter\label\expandafter{\@tmplbl}%
      \expanded{%
        \noexpand\newglossaryentry{#3\csname the#4\endcsname}{%
          name={\unexpanded{##2}%
            \IfBooleanF{##1}{\IfValueT{##4}{ [\noexpand\texttt{#3:##4}]}}},%
          sort={\csname the#2\endcsname},%
          description={p.~\noexpand\pageref{#3:\@relbl}},%
          user1={\csname the#2\endcsname},%
          user2={#5},%
          user3={\thesection},%
          user4={\expandafter\unexpanded\expandafter{\@savedSecTitle}}%
        }%
      }%
      \protected@write\@auxout{}{\string\glsliveentry{#3\csname the#4\endcsname}}%
      \protected@write\@auxout{}{\string\@thmtitle{\@tmplbl}{\unexpanded{##2}}}%
    \fi
  }%
  \expandafter\newcommand\csname ee#1\endcsname{\end{#2}}%
}
\@mkresult{L}{lemma}{l}{lemcounter}{Lemma}
% REMARK / COROLLARY / PROPOSITION -- same signature as \beL above
\@mkresult{R}{remark}{r}{remcounter}{Remark}
\@mkresult{C}{corollary}{c}{corcounter}{Corollary}
\@mkresult{P}{proposition}{p}{propcounter}{Proposition}
% OPEN-PROBLEM COMMAND -- generated with the same signature as \beT:
%   \beO[long title][short title][key]   ...   \eeO
% label o:KEY, or o:<counter> when no key is given; a bare \beO makes
% neither a label nor a glossary entry.
\@mkresult{O}{open}{o}{opcounter}{Problem}


% Section-level title tracker: updated only by \section, not \subsection.
% \@currentlabelname is overwritten by every \refstepcounter (including
% \subsection), so we save it separately here.
\newcommand{\@currentsectionname}{}
\AtBeginDocument{%
  \let\@epidcrn@orig@refstepcounter\refstepcounter
  \renewcommand{\refstepcounter}[1]{%
    \@epidcrn@orig@refstepcounter{#1}%
    \ifstrequal{#1}{section}{\global\let\@currentsectionname\@currentlabelname}{}%
  }%
}

% TABLE COMMANDS (glossaried) -- consistent naming with \beT, \beD, etc.
\newcommand{\beTa}[1][]{%
  \def\@temparg{#1}%
  \ifx\@temparg\@empty
    \begin{table}[htbp]%
  \else
    \let\@savedSecTitle\@currentsectionname%
    \refstepcounter{tabcounter}%
    \begin{table}[htbp]%
    \label{A:\thetabcounter}%
    \expanded{%
      \noexpand\newglossaryentry{A\thetabcounter}{%
        name={\unexpanded{#1}},%
        sort={\thetabcounter},%
        description={p.~\noexpand\pageref{A:\thetabcounter}},%
        user1={\thetabcounter},%
        user2={Table},%
        user3={\thesection},%
        user4={\expandafter\unexpanded\expandafter{\@savedSecTitle}}%
      }%
    }%
    \protected@write\@auxout{}{\string\glsliveentry{A\thetabcounter}}%
  \fi
}
\newcommand{\eeTa}{\end{table}}

% FIGURE COMMANDS (glossaried) -- consistent with \beT, \beR, \beTa, etc.
% Usage:
%   \beF[Glossary title]...\caption{Visible caption}\eeF  -> glossaried figure
%   \beF...\caption{Visible caption}\eeF                  -> plain [H] figure, no glossary
\newcommand{\beF}[1][]{%
  \def\@temparg{#1}%
  \ifx\@temparg\@empty
    \begin{figure}[H]%
  \else
    \begin{figure}[H]%
    \let\@savedSecTitle\@currentsectionname%
    \refstepcounter{figcounter}%
    \label{fig:\thefigcounter}%
    \expanded{%
      \noexpand\newglossaryentry{Fig\thefigcounter}{%
        name={\unexpanded{#1}},%
        sort={\thefigcounter},%
        description={p.~\noexpand\pageref{fig:\thefigcounter}},%
        user1={\thefigcounter},%
        user2={Figure},%
        user3={\thesection},%
        user4={\expandafter\unexpanded\expandafter{\@savedSecTitle}}%
      }%
    }%
    \protected@write\@auxout{}{\string\glsliveentry{Fig\thefigcounter}}%
  \fi
}
\newcommand{\eeF}{\end{figure}}
\makeatother
% LIVE-ENTRY REGISTER
% \glsaddall prints every entry the glossaries package currently knows, and
% that set includes entries restored from \jobname.glsdefs.  The package
% rewrites .glsdefs from all *currently defined* entries, so anything it
% restored is written straight back out: a deleted section kept reappearing in
% the catalogue for ever, and the .glsdefs had to be erased by hand.  The
% wrappers therefore record every entry they actually create in the .aux, and
% \printboth adds only those.  The .aux is rebuilt from scratch on each run, so
% this is self-healing: a removed result is gone after one further run.
\makeatletter
\newcommand{\@glslive}{}
\newcommand{\@glsaddlive}[1]{\ifglsentryexists{#1}{\glsadd{#1}}{}}
\newcommand{\glsliveentry}[1]{\g@addto@macro\@glslive{\@glsaddlive{#1}}}
\makeatother
% PRINT COMMAND
% No \clearpage here -- the catalogue continues on the current page.  This has
% been re-added by mistake several times; leave it out.
\makeatletter
\newcommand{\printboth}{%
  \@glslive   % add exactly the entries this document created last run
  \printglossary[style=unifiedstyle, title={0 \hspace{.22em} Definitions, Theorems, Examples and Open Problems}]%
}
\makeatother
% Named theorem reference: \tref{word} -> Thm~N(word-or-short)
% Shows stored short title if available, otherwise shows the key itself.
\newcommand{\tref}[1]{%
  Thm~\ref{t:#1}(\ifcsname @tshort@#1\endcsname\csname @tshort@#1\endcsname\else#1\fi)%
}

% Full-title theorem reference: \nref{t:KEY} prints the number and then the
% theorem's long title, so one writes "Theorem~\nref{t:blocks}" and gets
% "Theorem 3 ``Block factorization and siphons''".  Takes the complete label, like
% \ref, not the bare key like \tref.
%
% \beT writes each title to the .aux (see \@thmtitle below), so \nref resolves
% forward references from the second LaTeX pass, exactly as \ref does.  If the
% title is not yet known -- first pass, or a label that is not a \beT theorem
% -- it degrades to the number rather than producing ??.
% These two use @ in their names, and this point in the file is outside the
% \makeatletter region opened earlier (it closes at \makeatother above), so
% the region has to be reopened here.
\makeatletter
\newcommand{\@thmtitle}[2]{\global\@namedef{@tlong@#1}{#2}}
\newcommand{\nref}[1]{%
  \ref{#1}%
  \ifcsname @tlong@#1\endcsname
    \ ``\csname @tlong@#1\endcsname''%
  \fi
}
% \nm{t:KEY}: the stored long title of any \beT/\beL/\beP/\beR/\beC/\beD/\beE result,
% verbatim and without quotes, so item names follow the macro's first argument.
\newcommand{\nm}[1]{%
  \ifcsname @tlong@#1\endcsname
    \csname @tlong@#1\endcsname%
  \else
    \ref{#1}%
  \fi
}
\makeatother




%% ========================================
%% Other THEOREM ENVIRONMENTS
%% ========================================
\makeatletter
\@ifundefined{lemma}      {\newtheorem{lemma}      {Lemma}}      {}
\@ifundefined{proposition}{\newtheorem{proposition}{Proposition}} {}
\@ifundefined{corollary}  {\newtheorem{corollary}  {Corollary}}  {}
\@ifundefined{conjecture} {\newtheorem{conjecture} {Conjecture}}  {}
\@ifundefined{remark}     {\newtheorem{remark}     {Remark}}     {}
\@ifundefined{question}   {\newtheorem{question}   {Question}}   {}
\@ifundefined{assumption} {\newtheorem{assumption} {Assumption}} {}
\@ifundefined{hypothesis} {\newtheorem{hypothesis} {Hypothesis}} {}
\@ifundefined{fact}       {\newtheorem{fact}       {Fact}}       {}
\makeatother

% \citer: cite the relay-graph threshold theorem of AAHrel. If a document has
% \input the chapter/file containing t:relayRF (so the label resolves), point
% straight at it; otherwise fall back to a bare citation of AAHrel. Shared here
% (not per-paper) so it resolves consistently across every book chapter.
% \citeR: same idea, for the first theorem of AAHRah (t:Pstar).
\makeatletter
\providecommand{\citer}{%
  \@ifundefined{r@t:relayRF}%
    {\cite{AAHrel}}%
    {Theorem~\ref{t:relayRF} of \cite{AAHrel}}%
}
\providecommand{\citeR}{%
  \@ifundefined{r@t:Pstar}%
    {\cite{AAHRah}}%
    {Theorem~\ref{t:Pstar} of \cite{AAHRah}}%
}
\makeatother


% Abbreviations
\def\beQ{\begin{question}}\def\eeQ{\end{question}}
\def\beA{\begin{assumption}}\def\eeA{\end{assumption}}


%% ========================================

%% ========================================
%% COLORS & CUSTOM COMMANDS
%% ========================================
\definecolor{funccolor}{RGB}{25,25,112}
\definecolor{desccolor}{RGB}{64,64,64}
\newcommand{\func}[1]{\texttt{\textcolor{funccolor}{#1}}}
\def\eq{=}
\def\mm{-}

%% ========================================
%% MATHCAL SHORTHANDS
%% ========================================
% Project-wide convention: \mX for a capital X is always \mathcal X.
% Defined for the whole alphabet so the convention never fails to expand.
\newcommand{\mA}{\mathcal A}\newcommand{\mB}{\mathcal B}
\newcommand{\mC}{\mathcal C}\newcommand{\mD}{\mathcal D}
\newcommand{\mE}{\mathcal E}\newcommand{\mF}{\mathcal F}
\newcommand{\mG}{\mathcal G}\newcommand{\mH}{\mathcal H}
\newcommand{\mI}{\mathcal I}\newcommand{\mJ}{\mathcal J}
\newcommand{\mK}{\mathcal K}\newcommand{\mL}{\mathcal L}
\newcommand{\mM}{\mathcal M}\newcommand{\mN}{\mathcal N}
\newcommand{\mO}{\mathcal O}\newcommand{\mP}{\mathcal P}
\newcommand{\mQ}{\mathcal Q}\newcommand{\mR}{\mathcal R}
\newcommand{\mS}{\mathcal S}\newcommand{\mT}{\mathcal T}
\newcommand{\mU}{\mathcal U}\newcommand{\mV}{\mathcal V}
\newcommand{\mW}{\mathcal W}\newcommand{\mX}{\mathcal X}
\newcommand{\mY}{\mathcal Y}\newcommand{\mZ}{\mathcal Z}

%% ========================================
%% BASIC MATH ENVIRONMENTS
%% ========================================
\def\bep{\begin{pmatrix}}\def\eep{\end{pmatrix}}
\def\bev{\begin{vmatrix}}\def\eev{\end{vmatrix}}
\def\bea{\begin{eqnarray*}}\def\eea{\end{eqnarray*}}
\def\bc{\begin{cases}}\def\ec{\end{cases}}
\def\BEN{\begin{enumerate}}\def\EEN{\end{enumerate}}
\def\BI{\begin{itemize}}\def\EI{\end{itemize}}

\newcommand{\be}[1]{\begin{equation}\label{#1}}
\newcommand{\ee}{\end{equation}}
\newcommand{\beq}{\begin{eqnarray}}
\def\eeq{\end{eqnarray}}
\def\eqr{\eqref}\def\fr{\frac}\def\lbl{\label}
\def\Lra{\Longrightarrow}\def\non{\nonumber}\def\cNGM{\cite{Diek,Van,Van08}}
\def\Eq{\Leftrightarrow}\def\b{\beta}\def\uph{\red{UP TO HERE}}
%% ========================================
%% PROBABILITY & STATISTICS
%% ========================================
\renewcommand{\P}{\ensuremath{\mathbb{P}}}  % Note: overrides paragraph symbol
\newcommand{\E}{\ensuremath{\mathbb{E}}}
\newcommand{\ind}[1]{\ensuremath{\mathbbm{1}_{\left\{#1\right\}}}}
\DeclareMathOperator{\Cov}{Cov}
\DeclareMathOperator{\TV}{TV}
\newcommand{\rank}{\operatorname{rank}}

%% ========================================
%% CRN OPERATORS (organizations/ERCs/buffering/siphons framework)
%% ========================================
\def\clos{\mathrm{clos}}
\def\Supply{\mathrm{Supply}}
\def\req{\mathrm{req}}
\def\React{\mathrm{React}}
\def\Prod{\mathrm{Prod}}
\def\Gen{\mathrm{Gen}}
\def\In{\mathrm{In}}
\def\Out{\mathrm{Out}}
\def\Act{\mathrm{Act}}
\def\Min{\mathrm{Min}}

%% ========================================
%% LINE BREAKING
%% ========================================
% Last-resort stretch, used only on paragraphs TeX cannot otherwise break:
% a long \texttt{} token, a comma-separated inline math list (commas are not
% breakpoints in math), or a run of \cite/\ref boxes.  Without it those set as
% overfull \hbox and print into the margin; with it they set slightly loose,
% which is the lesser evil.  Does not affect paragraphs that already break.
\emergencystretch=3em

%% ========================================
%% NUMBER SETS
%% ========================================
\newcommand{\R}{\mathbb{R}}  % Real numbers
\newcommand{\C}{\mathbb{C}}  % Complex numbers
\newcommand{\N}{\mathbb{N}}
\newcommand{\Z}{\mathbb{Z}}
% \providecommand, not \newcommand: some papers define these in their own
% preamble after \input{preamble}, and must not clash.
\providecommand{\Q}{\mathbb{Q}}
\providecommand{\res}{\operatorname{res}}   % resultant; used by epiRes.tex
\newcommand{\MN}{\mathbb{N}}
\newcommand{\MZ}{\mathbb{Z}}
\newcommand{\realpos}{\mathbb{R}_{> 0}}
\newcommand{\realnneg}{\mathbb{R}_{\geq 0}}
\newcommand{\naturalpos}{\mathbb{N}_{>0}}
\newcommand{\Rplus}{\mathbb{R}_+}
\newcommand{\Rp}{\mathbb{R}_{\ge 0}}
\def\Rnne{\mathbb{R}_{\ge 0}}
\def\Se{\bar{S}} \def\Ie{\bar{I}}
\newcommand{\Rpos}{\mathbb{R}_{>0}}
\newcommand{\dd}{\,\mathrm d}
\newcommand{\e}{\mathrm e}
%% ========================================
%% GREEK LETTERS (2+ character abbreviations) greek
%% ========================================
\def\al{\alpha}\def\ba{\bff \alpha}
\def\w{\omega}\def\Lgn{\ga+\mu+\de}\def\om{\omega}
\def\de{\delta}\def\De{\Delta}\def\bde{\bff \delta}
\def\ep{\epsilon}\def\eps{\varepsilon}
\def\f{\varphi}  % Kept as requested
\def\ga{\gamma}\def\Ga{\Gamma}\def\G{\Gamma}\def\bg{\bff \gamma}
\def\up{\upsilon}  % NOT the breve accent: no accent commands anywhere.
                   % Named \up, not \u, so the LaTeX breve stays reachable.
\def\la{\lambda}\def\La{\Lambda}
\def\si{\sigma}\def\ka{\kappa}
\def\th{\theta}\def\vt{\vartheta}
\def\ze{\zeta}\def\Gap{\Gamma^{+}} \def\Gam{\Gamma^{-}}


%% ========================================
%% SPECIAL SHORT COMMANDS (kept as requested)
%% ========================================
\def\I{\infty}
\def\T{\widetilde}
\def\H{\hat}
\newcommand{\relint}{\operatorname{relint}}

%% ========================================
%% BOLD/VECTOR NOTATION
%% ========================================
\newcommand{\bff}[1]{{\mbox{\boldmath$#1$}}}
\newcommand\X{\boldsymbol{X}}
\def\x{\boldsymbol{x}}\def\bx{\bff x}\def\xe{\bar{\x}}
\newcommand\y{\boldsymbol{y}}\def\by{\boldsymbol{y}}\def\ye{\bar{\y}}
\def\yb{\y}\def\xb{\x}
\def\xy{\x-\y}
\newcommand\br{\boldsymbol{r}}\def\re{\bar{\r}}
\newcommand\bd{\boldsymbol{d}}
\newcommand\bQ{\boldsymbol{Q}}
\newcommand{\Yb}{{\mathbf{Y}}}
\newcommand{\Ab}{\mathbf{A}}
\newcommand{\zb}{\mathbf{z}}\def\bz{\bff 0}
\newcommand{\Zb}{\mathbf{Z}}
\newcommand{\bb}{\mathbf{b}}
\def\vbe{\vec {\bff \beta}}
\def\vb{\vec {\bff b}}
\def\v1{\vec {\bff 1}}
\def\vn{\vec \de}
\def\vz{\vec 0}
\def\vi{\overset{\rightarrow}{\mathsf i}}
\def\vv{\overset{\rightarrow}{v}}
\def\vr{\; \vec {\mathsf r}}
\newcommand{\va}{\bff a}
\providecommand{\vect}[1]{\boldsymbol{#1}}   % bold vector (moved from Gradostat.tex)
\providecommand{\one}{\mathbf 1}             % all-ones vector (moved from AliSF)
\providecommand{\spr}{s}                     % spectral bound  (moved from AliSF)
\providecommand{\Ke}{K_{\rm eff}}            % effective carrying capacity
% Critical right/left eigenvectors in the Hopf projection method (AAHOSN/New.tex)
\providecommand{\bq}{\mathbf{q}}
\providecommand{\bp}{\mathbf{p}}


%% ========================================
%% MATHCAL & MATHSCR
%% ========================================
% Mathcal
\newcommand\mc[1]{\mathcal{#1}}
\def\mC{{\mathcal C}}\def\mD{\mathcal D}\def\mE{\mathcal{E}}
\def\mF{\mathcal F}\def\mG{{\mathcal G}}\def\mH{{\mathcal H}}
\def\mI{\mathcal I}\def\mL{\mathcal{L}}\def\mM{{\mathcal M}}
\def\mN{{\mathcal N}}\def\mP{\mathcal{P}}\def\mR{{\mathcal R}}
\def\mS{{\mathcal S}}\def\mT{\mathcal{T}}\def\mV{{\mathcal V}}
\def\mK{\mathcal K}\def\mA{{\mathcal A}}\def\mU{{\mathcal U}}
% Symmetric bilinear and trilinear forms of the Taylor expansion, used in the
% first Lyapunov coefficient (AAHOSN/New.tex)
\providecommand{\Bform}{\mathcal{B}}
\providecommand{\Cform}{\mathcal{C}}

% Blackboard bold.  Kept here, not in the drivers: a \providecommand in a
% driver preamble runs before this file and silently wins, so the same name
% ends up meaning different things in different builds.
\def\IC{\mathbb{C}}\def\IN{\mathbb{N}}\def\IQ{\mathbb{Q}}
\def\IR{\mathbb{R}}\def\IZ{\mathbb{Z}}\def\IG{\mathbb{G}}
\def\II{\mathbb{I}}

% Mathscr
\newcommand{\Rs}{{\mathscr R}}
\newcommand{\Ss}{{\mathscr S}}
\newcommand{\Ns}{{\mathscr N}}
\newcommand{\Cs}{\mathscr C}


%% ========================================
%% OPERATORS & FUNCTIONS
%% ========================================
\DeclareMathOperator{\tr}{tr}
\DeclareMathOperator{\diag}{Diag}
\DeclareMathOperator{\Diag}{Diag}
\DeclareMathOperator{\intr}{\mathrm{Int}}
\DeclareMathOperator{\Inv}{Inv}
\DeclareMathOperator{\spec}{spec}
\DeclareMathOperator{\AUC}{AUC}               % area under the curve
\newcommand{\ii}{\mathrm{i}}                  % imaginary unit
\newcommand\op{\operatorname}
\newcommand{\supp}{\op{supp}}
\newcommand{\diff}{\mathop{}\mathopen{}\mathrm{d}}
\newcommand{\pdd}[1]{\frac{\partial}{\partial{#1}}}
\DeclareMathOperator{\cone}{cone}
\DeclareMathOperator{\proj}{proj}

%% ========================================
%% BRACKETS & NORMS
%% ========================================
\providecommand{\pr}[1]{\left(#1\right)}
\providecommand{\pp}[1]{\left[#1\right]}
\newcommand\croc[1]{\left\langle #1\right\rangle}
\newcommand{\norm}[1]{\|#1\|}
\newcommand{\inner}[2]{\langle #1, #2 \rangle}
% Short synonym of \inner, used by AAHOSN/New.tex as $\ip{u}{v}=\bar u^\top v$
\providecommand{\ip}[2]{\langle #1 , #2 \rangle}


%% ========================================
%% FORMATTING & UTILITIES
%% ========================================
\newcommand*{\Scale}[2][4]{\scalebox{#1}{$#2$}}
\newcommand{\te}[1]{\text{#1}}
\newcommand\steq[1]{\stackrel{\text{\rm #1.}}{=}}
\newcommand\blue[1]{\textcolor{blue}{#1}}
\newcommand{\red}{\textcolor[rgb]{1.00,0.00,0.00}}
\newcommand{\green}{\textcolor[rgb]{0.00,1.00,0.00}}

% Figure command
\newcommand{\figu}[3]{%
\begin{figure}[H]
\centering
\includegraphics[scale=#3]{#1}
\caption{#2\label{f:#1}}
\end{figure}
}

% Footnote
\long\def\symbolfootnote[#1]#2{%
\begingroup
\def\thefootnote{\fnsymbol{footnote}}\footnote[#1]{#2}%
\endgroup}
\def\fn{\symbolfootnote}

\newcommand{\dsum}{\displaystyle\sum}

%% ========================================
%% CRN & EPIDEMIOLOGY SPECIFIC
%% ========================================
% CRN notation
\def\SS{\mathcal S}
\def\CC{\mathcal C}
\def\RR{\mathcal R}

% Reproduction numbers
\def\Rto{R_{1}^{E_2}}
\def\Rot{R_{2}^{E_1}}
\def\Roc{R_{1,c}}
\def\Rtc{R_{2,c}}
\def\rd{r_{0}}
\def\sd{s_{0}}
\def\dd{D_{0}}
\def\vd{v_{0}}
\def\zd{z_{0}}
\def\se{\sE}
\def\xd{x_{0}}
%% ========================================
%% SPECIAL VALUES & EQUILIBRIA
%% ========================================
\def\eE{\bar{e}}
\def\wE{\bar{w}}
\def\sE{s_{*}}
\def\Me{\T M_{ee}}
\def\Rnne{\mathbb{R}_{\ge 0}}

\newcommand{\ones}{\mathbf 1}


%% ========================================
%% ABBREVIATIONS (alphabetical by category)
%% ========================================

% -- General mathematical terms --
\def\adm{admissible}\def\admy{admissibility}
%\def\and{antisymmetric}% disabled 2026-09-30: it overrode LaTeX's author separator \and (printed "antisymmetric" between authors)
\def\ass{assumption}
\def\cadlag{c\`adl\`ag}
\def\ch{characteristic polynomial}
\def\chp{characteristic polynomial}
\def\coe{coefficient}
\def\comp{computation}
\def\con{condition}
\def\coo{coordinate}
\def\corr{corresponding}
\def\dec{decomposition}
\def\den{denominator}
\def\dete{deterministic}
\def\eig{eigenvalue}\def\eval{eigenvalue}
\def\eigv{eigenvector}\def\eiv{eigenvector}\def\evec{eigenvector}
\def\elt{element}
\def\enn{essentially nonnegative/positivity preserving}
\def\equ{equation}
\def\etal{et al.}
\def\fe{for example}
\def\fno{from now on}
\def\for{\forall}
\def\form{formula}
\def\Fr{Furthermore, }
\def\frt{furthermore}
\def\fun{function}
\def\gene{generalization}
\def\heu{heuristic}
\def\how{however}\def\How{However, }
\def\ie{i.e. }
\def\im{\item}
\def\isd{is defined by}
\def\itc{it is easy to check that}
\def\itf{it follows that}\def\Itf{It follows that}
\def\ith{it holds that}
\def\itm{it may be checked that}\def\Itm{It may be checked that}
\def\ito{it turns out that}\def\Ito{It turns out that}\def\Itou{It turns out useful to}
\def\ma{mass action}\def\MAR{mass action representation}
\def\mat{matrix}
\def\mbc{may be checked to be}
\def\mbw{may be written as}\def\mubw{must be written as}
\def\mpr{more precisely, }\def\Mp{More precisely, }
\def\nc{non-catalytic}
\def\nne{non-negative}\def\nny{nonnegativity}
\def\no{\nonumber}
\def\noni{non-infectious}
\def\nota{notation}
\def\np{\newpage}
\def\Nt{Note that}
\def\nt{non trivial}\def\nz{non zero}
\def\oth{otherwise}
\def\Oth{On the other hand}
\def\pa{parameter}\def\para{parameter}
\def\pol{polynomial}
\def\prob{probability}
\def\Prf{{\bf  Proof:}}\def\prf{{\bf  Proof:}}
\def\prop{proportional to}
\def\qu{\quad}
\def\QED{\hfill {$\square$}\goodbreak \medskip}
\def\rep{representation}
\def\resp{respectively}
\def\satd{satisfied}\def\satg{satisfying}\def\sats{satisfies}\def\saty{satisfy}
\def\ssec{\subsection}\def\sssec{\subsubsection}
\def\sef{see for example}
\def\sksy{skew-symmetric}
\def\sub{substrate}
\def\sys{system}
\def\var{variable}
\def\wk{well-known}
\def\wl{without loss of generality}\def\wlo{w.l.o.g.}
\def\wmc{We may check that}\def\Wmc{We may check that}
\def\wmr{we may rewrite}
\def\wms{we must show that}
\def\wrt{with respect to}

% -- CRN & Chemistry --
\newcommand\CRN{chemical reaction network}
\newcommand\CRNs{chemical reaction networks}
\newcommand\RN{reaction network}
\def\ACR{absolute concentration robustness}
\def\AMSD{acyclic siphon decomposition}
\def\DAE{differential algebraic equations}
\def\FHJ{Feinberg-Horn-Jackson}
\def\GMAK{generalized mass action kinetics}
\def\MM{Michaelis-Menten}
\def\sm{stoichiometric matrix}\def\Sm{stoichiometric matrix}
\def\WRZD{weakly reversible, zero deficiency}\def\aut{autocatalytic}

% -- Epidemiology --
\def\brn{basic reproduction number}\def\bzn{basic replacement number}
\def\com{compartment}
\def\dfe{disease free equilibrium}\def\DFE{disease free equilibrium}
\def\dwt{dwell times}\def\edt{expected dwell times}
\def\EM{epidemic model}
\def\endp{endemic fixed point}\def\enp{endemic point}\def\EE{endemic point}
\def\ME{mathematical epidemiology}
\def\MEE{eco-epidemiology}\def\EEP{eco-epidemiological}
\def\ece{eco-epidemiological model}
\def\mgm{multi group model}
\def\mtS{multi-strain model}\def\mSt{multi-strain model}
\def\mSi{minimal siphon}
\def\NGM{next generation matrix}
\def\nim{new infections matrix}
\def\qem{quadratic epidemiological model}
\def\repF{invasion functions}
\def\SPF{SIR-PH-FA model}
\def\sus{susceptible}

% -- Graph theory & Linear algebra --
\def\jac{Jacobi matrix}
\def\jin{Jacobian matrix of the invasion vector field with respect to the invasion variables}
\def\jrn{Jacobian matrix of the resident vector field with respect to the resident variables}
\def\jxn{Jacobian matrix of the $\x$ vector field with respect to the $\x$ variables}
\def\Lm{Laplacian matrix}
\def\MTT{matrix tree theorem}
\def\SGT{spectral graph theory}
\def\SMW{Sherman–Morrison–Woodbury}

% -- Stability & Dynamics --
\def\absC{absorbtion CTMC}
\def\BTb{Bogdanov-Takens bifurcation}
\def\CEP{competitive exclusion partition}
\def\CTMC{continuous time Markov chain}
\def\ds{differential system}
\def\fp{fixed point}\def\bfp{boundary fixed point}\def\fbp{fixed boundary point}
\def\fbl{fixed boundary line}\def\fif{forward invariant face}
\def\gas{global asymptotic stability}\def\GAS{globally asymptotically stable}
\def\GASS{\text{GASS}}
\def\Hb{Hopf bifurcation}
\def\Hur{Hurwitz}\def\MH{Metzler-Hurwitz}\def\RH{Routh-Hurwitz}
\def\IS{\text{IS}}
\def\LAS{locally asymptotically stable}
\def\LdS{Lyapunov diagonally stable}
\def\Lf{Lyapunov function}\def\LyV{Lyapunov Volterra type function}\def\RLF{robust Lyapunov function}
\def\LyM{Lyapunov Malkin stable}
\def\PF{Perron-Frobenius}
\def\spfp{strictly positive fixed point}
\def\sta{stability}
\def\STP{sharp threshold property}\def\STT{sharp threshold theorem}

% -- Models (with \varphi explicitly instead of \f) --
\def\AB{$(A,B)$ Arino-Brauer epidemic models}
\def\ABa{$(A,\bbe,\ba,\varphi)$ balanced bilinear model}
\def\ABP{%$(A,\bbe,P,\varphi)$
\\ balanced bilinear model}
\def\AMG{Anderson May Gupta}
\newcommand\BIN{biological interaction network}
\def\Dt{Descartes type}
\def\gLV{generalized Lotka-Volterra model}
\def\IG{\text{invasion graph}}\def\IGMS{interaction graph on minimal siphons}
\def\LCP{\text{linear complementarity problem}}
\def\LV{Lotka-Volterra}\def\LVK{Lotka-Volterra-Kolmogorov}

% -- Splitting & Structure --
\def\reg{regular splitting}\def\regS{{\bf regular splitting}}
\def\RUR{rational univariate representation}
\def\bds{block-diagonal structure}\def\NGMs{next generation matrices}

% -- Miscellaneous --
\def\Ma{{Mathematica}}
\def\Mm{Max–Min}
\def\num{numerator}
\def\rap{rational approximation} \def\DSR{directed species-reaction bipartite graph}\def\PCP{polynomial complementarity problem}\def\rat{rational substitution}
\def\QSS{quasi steady state} \def\RHS{right hand side} \def\LHS{left hand side}
\def\OSN{online social network}\def\VdD{Van den Driessche}
\def\wli{$\omega$-limit} \def\rk{rich kinetics}\def\sJ{symbolic Jacobian}
\def\cV{\cite{Vas,VasSta,BSV,GSV}} \def\spf{semi-parametric semi-functional}
\def\CRW{Capasso-Ruan-Wang} \def\ABP{$(A,\bbe,P,\f)$ balanced bilinear model}
\def\ABa{%$(A,\bbe,\ba,\f)$ \\
rank-one balanced bilinear model}
\def\bbe{B%{\bff \beta}
} \def\vbe{\vec {\bff \beta}
}
\def\vb{\vec {\bff b}}
\newcommand{\Kh}{\hat{K}}          % K with hat, first mention only
\newcommand{\wS}{w^{\sigma}}       % right Perron J^\perp, full block
\newcommand{\piS}{\vec \pi^{\sigma}}    % left Perron J^\perp, full block
\newcommand{\ws}{w^{\sigma}}       % right Perron J^\perp, sub-block
\newcommand{\pis}{\pi^{\sigma}}    % left Perron J^\perp, sub-block
\newcommand{\hwS}{\hat{w}^{\sigma}}   % right Perron K, full block
\newcommand{\hpiS}{\hat{\pi}^{\sigma}} % left Perron K, full block
\newcommand{\hws}{\hat{w}^{\sigma}}   % right Perron K, sub-block
\newcommand{\hpis}{\hat{\pi}^{\sigma}} % left Perron K, sub-block
\newcommand{\twS}{\tilde{w}^{\sigma}}  % right Perron \tilde K, full block
\newcommand{\tpiS}{\tilde{\pi}^{\sigma}} % left Perron \tilde K
\newcommand{\tws}{\tilde{w}^{\sigma}}
\newcommand{\tpis}{\tilde{\pi}^{\sigma}}
\newcommand{\Mat}[2]{\mathbb{R}_{\ge 0}^{\,#1\times #2}}
\def\tJ{J_\sigma^\perp} \def\lPe{left Perron eigenvector}
% \G is \Gamma (defined at :795).  The Volterra function keeps its own name,
% so that \G no longer has two conflicting meanings in this file.
\def\GV{G}  % Volterra G(u)=u-1-\ln u
\def\PVL{Perron-Volterra Lyapunov minimal siphon ansatz}
\newcommand{\istar}[1]{i_{#1}^{*}}
%\newcommand{\sE}{s_E}
\newcommand{\vE}{v_E}
\def\PV{Perron-Volterra candidate Lyapunov function}
\def\GLV{generalized Lotka-Volterra}
\def\MERN{mathematical epidemiology reaction networks}
\def\CMM{contraction metric method}
\newcommand{\sign}{\operatorname{sign}}
\newcommand{\Var}{\operatorname{Var}} 

% ======================================================================
% IDEAS SHARED BETWEEN THE ABSTRACT AND THE MAIN RESULTS SECTION
%
% Write each idea ONCE, in the abstract; the Main Results section recalls it,
% so the two can never drift apart.
%
%   in the abstract        \ideaitem{1}{The first-order reaction terms ...}
%   in Main Results        \ideahead{1}
%
% \ideaitem stores the text globally and typesets
%   \item \hyperlink{idea:1}{<text>}
% \ideahead typesets
%   \item\hypertarget{idea:1}{}\textbf{<the same text>}
%
% The abstract must come before the Main Results input, which is the normal
% order. If a heading has to read differently from the abstract sentence (for
% instance a long abstract item summarised in one line), override it with the
% optional argument: \ideahead[short heading]{6}.
% ======================================================================
\makeatletter
\newcommand{\ideaitem}[2]{%
  \expandafter\gdef\csname @idea@#1\endcsname{#2}%
  \item\hyperlink{idea:#1}{#2}}
\newcommand{\ideahead}[2][]{%
  \item\hypertarget{idea:#2}{}%
  \def\@ideatmp{#1}%
  \ifx\@ideatmp\@empty
    \@ifundefined{@idea@#2}{%
      \PackageWarning{preamble}{Idea #2 is used in Main Results but was never
        defined by \string\ideaitem\space in the abstract}%
      \textbf{[idea #2 undefined]}%
    }{\textbf{\begingroup\def\ideaterm##1##2{##2}\def\idearest##1{}\@nameuse{@idea@#2}\endgroup}}%
  \else
    \textbf{#1}%
  \fi}
% \ideaterm{<hypertarget>}{<words>}: inside an abstract idea the words become an
% invisible link (the document sets \hypersetup{hidelinks}) jumping to the place
% in Main Results that defines them; inside the Main Results heading itself,
% where the target sits a few lines below, \ideahead strips the link and prints
% the words plainly. Mark the destination with \hypertarget{<name>}{} on the line
% before the display, which adds no equation number and shows nothing.
\newcommand{\ideaterm}[2]{\hyperlink{#1}{#2}}
% \idearest{<text>}: the part of an abstract idea after its first sentence; printed in
% the abstract, omitted from the Main Results heading.
\newcommand{\idearest}[1]{#1}
\makeatother

% ======================================================================
% SHORT-PAPER SWITCH
% \shorttrue in a master that assembles only part of the companion paper
% (Barrat/SipArxiv.tex). \nsh{...} drops text that cites results living in the
% omitted sections, so the short build has no dangling references while the
% companion paper keeps the text verbatim.
% ======================================================================
\newif\ifshort \shortfalse
\newcommand{\nsh}[1]{\ifshort\else#1\fi}
% \shortalt{<short build>}{<companion paper>}: use when a sentence needs
% different wording in the two builds, not merely a deletion.
\newcommand{\shortalt}[2]{\ifshort#1\else#2\fi}
% \Jtr, \JtrS{\Sigma}: transverse Jacobian. J^\perp here; SipArxiv.tex redefines them as J_{\mathrm{tra}}, J_{\mathrm{tra},\Sigma}
\newcommand{\Jtr}{J^\perp}
\newcommand{\JtrS}[1]{J_{#1}^\perp}
, and must not clash.
\providecommand{\Q}{\mathbb{Q}}
\providecommand{\res}{\operatorname{res}}   % resultant; used by epiRes.tex
\newcommand{\MN}{\mathbb{N}}
\newcommand{\MZ}{\mathbb{Z}}
\newcommand{\realpos}{\mathbb{R}_{> 0}}
\newcommand{\realnneg}{\mathbb{R}_{\geq 0}}
\newcommand{\naturalpos}{\mathbb{N}_{>0}}
\newcommand{\Rplus}{\mathbb{R}_+}
\newcommand{\Rp}{\mathbb{R}_{\ge 0}}
\def\Rnne{\mathbb{R}_{\ge 0}}
\def\Se{\bar{S}} \def\Ie{\bar{I}}
\newcommand{\Rpos}{\mathbb{R}_{>0}}
\newcommand{\dd}{\,\mathrm d}
\newcommand{\e}{\mathrm e}
%% ========================================
%% GREEK LETTERS (2+ character abbreviations) greek
%% ========================================
\def\al{\alpha}\def\ba{\bff \alpha}
\def\w{\omega}\def\Lgn{\ga+\mu+\de}\def\om{\omega}
\def\de{\delta}\def\De{\Delta}\def\bde{\bff \delta}
\def\ep{\epsilon}\def\eps{\varepsilon}
\def\f{\varphi}  % Kept as requested
\def\ga{\gamma}\def\Ga{\Gamma}\def\G{\Gamma}\def\bg{\bff \gamma}
\def\up{\upsilon}  % NOT the breve accent: no accent commands anywhere.
                   % Named \up, not \u, so the LaTeX breve stays reachable.
\def\la{\lambda}\def\La{\Lambda}
\def\si{\sigma}\def\ka{\kappa}
\def\th{\theta}\def\vt{\vartheta}
\def\ze{\zeta}\def\Gap{\Gamma^{+}} \def\Gam{\Gamma^{-}}


%% ========================================
%% SPECIAL SHORT COMMANDS (kept as requested)
%% ========================================
\def\I{\infty}
\def\T{\widetilde}
\def\H{\hat}
\newcommand{\relint}{\operatorname{relint}}

%% ========================================
%% BOLD/VECTOR NOTATION
%% ========================================
\newcommand{\bff}[1]{{\mbox{\boldmath$#1$}}}
\newcommand\X{\boldsymbol{X}}
\def\x{\boldsymbol{x}}\def\bx{\bff x}\def\xe{\bar{\x}}
\newcommand\y{\boldsymbol{y}}\def\by{\boldsymbol{y}}\def\ye{\bar{\y}}
\def\yb{\y}\def\xb{\x}
\def\xy{\x-\y}
\newcommand\br{\boldsymbol{r}}\def\re{\bar{\r}}
\newcommand\bd{\boldsymbol{d}}
\newcommand\bQ{\boldsymbol{Q}}
\newcommand{\Yb}{{\mathbf{Y}}}
\newcommand{\Ab}{\mathbf{A}}
\newcommand{\zb}{\mathbf{z}}\def\bz{\bff 0}
\newcommand{\Zb}{\mathbf{Z}}
\newcommand{\bb}{\mathbf{b}}
\def\vbe{\vec {\bff \beta}}
\def\vb{\vec {\bff b}}
\def\v1{\vec {\bff 1}}
\def\vn{\vec \de}
\def\vz{\vec 0}
\def\vi{\overset{\rightarrow}{\mathsf i}}
\def\vv{\overset{\rightarrow}{v}}
\def\vr{\; \vec {\mathsf r}}
\newcommand{\va}{\bff a}
\providecommand{\vect}[1]{\boldsymbol{#1}}   % bold vector (moved from Gradostat.tex)
\providecommand{\one}{\mathbf 1}             % all-ones vector (moved from AliSF)
\providecommand{\spr}{s}                     % spectral bound  (moved from AliSF)
\providecommand{\Ke}{K_{\rm eff}}            % effective carrying capacity
% Critical right/left eigenvectors in the Hopf projection method (AAHOSN/New.tex)
\providecommand{\bq}{\mathbf{q}}
\providecommand{\bp}{\mathbf{p}}


%% ========================================
%% MATHCAL & MATHSCR
%% ========================================
% Mathcal
\newcommand\mc[1]{\mathcal{#1}}
\def\mC{{\mathcal C}}\def\mD{\mathcal D}\def\mE{\mathcal{E}}
\def\mF{\mathcal F}\def\mG{{\mathcal G}}\def\mH{{\mathcal H}}
\def\mI{\mathcal I}\def\mL{\mathcal{L}}\def\mM{{\mathcal M}}
\def\mN{{\mathcal N}}\def\mP{\mathcal{P}}\def\mR{{\mathcal R}}
\def\mS{{\mathcal S}}\def\mT{\mathcal{T}}\def\mV{{\mathcal V}}
\def\mK{\mathcal K}\def\mA{{\mathcal A}}\def\mU{{\mathcal U}}
% Symmetric bilinear and trilinear forms of the Taylor expansion, used in the
% first Lyapunov coefficient (AAHOSN/New.tex)
\providecommand{\Bform}{\mathcal{B}}
\providecommand{\Cform}{\mathcal{C}}

% Blackboard bold.  Kept here, not in the drivers: a \providecommand in a
% driver preamble runs before this file and silently wins, so the same name
% ends up meaning different things in different builds.
\def\IC{\mathbb{C}}\def\IN{\mathbb{N}}\def\IQ{\mathbb{Q}}
\def\IR{\mathbb{R}}\def\IZ{\mathbb{Z}}\def\IG{\mathbb{G}}
\def\II{\mathbb{I}}

% Mathscr
\newcommand{\Rs}{{\mathscr R}}
\newcommand{\Ss}{{\mathscr S}}
\newcommand{\Ns}{{\mathscr N}}
\newcommand{\Cs}{\mathscr C}


%% ========================================
%% OPERATORS & FUNCTIONS
%% ========================================
\DeclareMathOperator{\tr}{tr}
\DeclareMathOperator{\diag}{Diag}
\DeclareMathOperator{\Diag}{Diag}
\DeclareMathOperator{\intr}{\mathrm{Int}}
\DeclareMathOperator{\Inv}{Inv}
\DeclareMathOperator{\spec}{spec}
\DeclareMathOperator{\AUC}{AUC}               % area under the curve
\newcommand{\ii}{\mathrm{i}}                  % imaginary unit
\newcommand\op{\operatorname}
\newcommand{\supp}{\op{supp}}
\newcommand{\diff}{\mathop{}\mathopen{}\mathrm{d}}
\newcommand{\pdd}[1]{\frac{\partial}{\partial{#1}}}
\DeclareMathOperator{\cone}{cone}
\DeclareMathOperator{\proj}{proj}

%% ========================================
%% BRACKETS & NORMS
%% ========================================
\providecommand{\pr}[1]{\left(#1\right)}
\providecommand{\pp}[1]{\left[#1\right]}
\newcommand\croc[1]{\left\langle #1\right\rangle}
\newcommand{\norm}[1]{\|#1\|}
\newcommand{\inner}[2]{\langle #1, #2 \rangle}
% Short synonym of \inner, used by AAHOSN/New.tex as $\ip{u}{v}=\bar u^\top v$
\providecommand{\ip}[2]{\langle #1 , #2 \rangle}


%% ========================================
%% FORMATTING & UTILITIES
%% ========================================
\newcommand*{\Scale}[2][4]{\scalebox{#1}{$#2$}}
\newcommand{\te}[1]{\text{#1}}
\newcommand\steq[1]{\stackrel{\text{\rm #1.}}{=}}
\newcommand\blue[1]{\textcolor{blue}{#1}}
\newcommand{\red}{\textcolor[rgb]{1.00,0.00,0.00}}
\newcommand{\green}{\textcolor[rgb]{0.00,1.00,0.00}}

% Figure command
\newcommand{\figu}[3]{%
\begin{figure}[H]
\centering
\includegraphics[scale=#3]{#1}
\caption{#2\label{f:#1}}
\end{figure}
}

% Footnote
\long\def\symbolfootnote[#1]#2{%
\begingroup
\def\thefootnote{\fnsymbol{footnote}}\footnote[#1]{#2}%
\endgroup}
\def\fn{\symbolfootnote}

\newcommand{\dsum}{\displaystyle\sum}

%% ========================================
%% CRN & EPIDEMIOLOGY SPECIFIC
%% ========================================
% CRN notation
\def\SS{\mathcal S}
\def\CC{\mathcal C}
\def\RR{\mathcal R}

% Reproduction numbers
\def\Rto{R_{1}^{E_2}}
\def\Rot{R_{2}^{E_1}}
\def\Roc{R_{1,c}}
\def\Rtc{R_{2,c}}
\def\rd{r_{0}}
\def\sd{s_{0}}
\def\dd{D_{0}}
\def\vd{v_{0}}
\def\zd{z_{0}}
\def\se{\sE}
\def\xd{x_{0}}
%% ========================================
%% SPECIAL VALUES & EQUILIBRIA
%% ========================================
\def\eE{\bar{e}}
\def\wE{\bar{w}}
\def\sE{s_{*}}
\def\Me{\T M_{ee}}
\def\Rnne{\mathbb{R}_{\ge 0}}

\newcommand{\ones}{\mathbf 1}


%% ========================================
%% ABBREVIATIONS (alphabetical by category)
%% ========================================

% -- General mathematical terms --
\def\adm{admissible}\def\admy{admissibility}
%\def\and{antisymmetric}% disabled 2026-09-30: it overrode LaTeX's author separator \and (printed "antisymmetric" between authors)
\def\ass{assumption}
\def\cadlag{c\`adl\`ag}
\def\ch{characteristic polynomial}
\def\chp{characteristic polynomial}
\def\coe{coefficient}
\def\comp{computation}
\def\con{condition}
\def\coo{coordinate}
\def\corr{corresponding}
\def\dec{decomposition}
\def\den{denominator}
\def\dete{deterministic}
\def\eig{eigenvalue}\def\eval{eigenvalue}
\def\eigv{eigenvector}\def\eiv{eigenvector}\def\evec{eigenvector}
\def\elt{element}
\def\enn{essentially nonnegative/positivity preserving}
\def\equ{equation}
\def\etal{et al.}
\def\fe{for example}
\def\fno{from now on}
\def\for{\forall}
\def\form{formula}
\def\Fr{Furthermore, }
\def\frt{furthermore}
\def\fun{function}
\def\gene{generalization}
\def\heu{heuristic}
\def\how{however}\def\How{However, }
\def\ie{i.e. }
\def\im{\item}
\def\isd{is defined by}
\def\itc{it is easy to check that}
\def\itf{it follows that}\def\Itf{It follows that}
\def\ith{it holds that}
\def\itm{it may be checked that}\def\Itm{It may be checked that}
\def\ito{it turns out that}\def\Ito{It turns out that}\def\Itou{It turns out useful to}
\def\ma{mass action}\def\MAR{mass action representation}
\def\mat{matrix}
\def\mbc{may be checked to be}
\def\mbw{may be written as}\def\mubw{must be written as}
\def\mpr{more precisely, }\def\Mp{More precisely, }
\def\nc{non-catalytic}
\def\nne{non-negative}\def\nny{nonnegativity}
\def\no{\nonumber}
\def\noni{non-infectious}
\def\nota{notation}
\def\np{\newpage}
\def\Nt{Note that}
\def\nt{non trivial}\def\nz{non zero}
\def\oth{otherwise}
\def\Oth{On the other hand}
\def\pa{parameter}\def\para{parameter}
\def\pol{polynomial}
\def\prob{probability}
\def\Prf{{\bf  Proof:}}\def\prf{{\bf  Proof:}}
\def\prop{proportional to}
\def\qu{\quad}
\def\QED{\hfill {$\square$}\goodbreak \medskip}
\def\rep{representation}
\def\resp{respectively}
\def\satd{satisfied}\def\satg{satisfying}\def\sats{satisfies}\def\saty{satisfy}
\def\ssec{\subsection}\def\sssec{\subsubsection}
\def\sef{see for example}
\def\sksy{skew-symmetric}
\def\sub{substrate}
\def\sys{system}
\def\var{variable}
\def\wk{well-known}
\def\wl{without loss of generality}\def\wlo{w.l.o.g.}
\def\wmc{We may check that}\def\Wmc{We may check that}
\def\wmr{we may rewrite}
\def\wms{we must show that}
\def\wrt{with respect to}

% -- CRN & Chemistry --
\newcommand\CRN{chemical reaction network}
\newcommand\CRNs{chemical reaction networks}
\newcommand\RN{reaction network}
\def\ACR{absolute concentration robustness}
\def\AMSD{acyclic siphon decomposition}
\def\DAE{differential algebraic equations}
\def\FHJ{Feinberg-Horn-Jackson}
\def\GMAK{generalized mass action kinetics}
\def\MM{Michaelis-Menten}
\def\sm{stoichiometric matrix}\def\Sm{stoichiometric matrix}
\def\WRZD{weakly reversible, zero deficiency}\def\aut{autocatalytic}

% -- Epidemiology --
\def\brn{basic reproduction number}\def\bzn{basic replacement number}
\def\com{compartment}
\def\dfe{disease free equilibrium}\def\DFE{disease free equilibrium}
\def\dwt{dwell times}\def\edt{expected dwell times}
\def\EM{epidemic model}
\def\endp{endemic fixed point}\def\enp{endemic point}\def\EE{endemic point}
\def\ME{mathematical epidemiology}
\def\MEE{eco-epidemiology}\def\EEP{eco-epidemiological}
\def\ece{eco-epidemiological model}
\def\mgm{multi group model}
\def\mtS{multi-strain model}\def\mSt{multi-strain model}
\def\mSi{minimal siphon}
\def\NGM{next generation matrix}
\def\nim{new infections matrix}
\def\qem{quadratic epidemiological model}
\def\repF{invasion functions}
\def\SPF{SIR-PH-FA model}
\def\sus{susceptible}

% -- Graph theory & Linear algebra --
\def\jac{Jacobi matrix}
\def\jin{Jacobian matrix of the invasion vector field with respect to the invasion variables}
\def\jrn{Jacobian matrix of the resident vector field with respect to the resident variables}
\def\jxn{Jacobian matrix of the $\x$ vector field with respect to the $\x$ variables}
\def\Lm{Laplacian matrix}
\def\MTT{matrix tree theorem}
\def\SGT{spectral graph theory}
\def\SMW{Sherman–Morrison–Woodbury}

% -- Stability & Dynamics --
\def\absC{absorbtion CTMC}
\def\BTb{Bogdanov-Takens bifurcation}
\def\CEP{competitive exclusion partition}
\def\CTMC{continuous time Markov chain}
\def\ds{differential system}
\def\fp{fixed point}\def\bfp{boundary fixed point}\def\fbp{fixed boundary point}
\def\fbl{fixed boundary line}\def\fif{forward invariant face}
\def\gas{global asymptotic stability}\def\GAS{globally asymptotically stable}
\def\GASS{\text{GASS}}
\def\Hb{Hopf bifurcation}
\def\Hur{Hurwitz}\def\MH{Metzler-Hurwitz}\def\RH{Routh-Hurwitz}
\def\IS{\text{IS}}
\def\LAS{locally asymptotically stable}
\def\LdS{Lyapunov diagonally stable}
\def\Lf{Lyapunov function}\def\LyV{Lyapunov Volterra type function}\def\RLF{robust Lyapunov function}
\def\LyM{Lyapunov Malkin stable}
\def\PF{Perron-Frobenius}
\def\spfp{strictly positive fixed point}
\def\sta{stability}
\def\STP{sharp threshold property}\def\STT{sharp threshold theorem}

% -- Models (with \varphi explicitly instead of \f) --
\def\AB{$(A,B)$ Arino-Brauer epidemic models}
\def\ABa{$(A,\bbe,\ba,\varphi)$ balanced bilinear model}
\def\ABP{%$(A,\bbe,P,\varphi)$
\\ balanced bilinear model}
\def\AMG{Anderson May Gupta}
\newcommand\BIN{biological interaction network}
\def\Dt{Descartes type}
\def\gLV{generalized Lotka-Volterra model}
\def\IG{\text{invasion graph}}\def\IGMS{interaction graph on minimal siphons}
\def\LCP{\text{linear complementarity problem}}
\def\LV{Lotka-Volterra}\def\LVK{Lotka-Volterra-Kolmogorov}

% -- Splitting & Structure --
\def\reg{regular splitting}\def\regS{{\bf regular splitting}}
\def\RUR{rational univariate representation}
\def\bds{block-diagonal structure}\def\NGMs{next generation matrices}

% -- Miscellaneous --
\def\Ma{{Mathematica}}
\def\Mm{Max–Min}
\def\num{numerator}
\def\rap{rational approximation} \def\DSR{directed species-reaction bipartite graph}\def\PCP{polynomial complementarity problem}\def\rat{rational substitution}
\def\QSS{quasi steady state} \def\RHS{right hand side} \def\LHS{left hand side}
\def\OSN{online social network}\def\VdD{Van den Driessche}
\def\wli{$\omega$-limit} \def\rk{rich kinetics}\def\sJ{symbolic Jacobian}
\def\cV{\cite{Vas,VasSta,BSV,GSV}} \def\spf{semi-parametric semi-functional}
\def\CRW{Capasso-Ruan-Wang} \def\ABP{$(A,\bbe,P,\f)$ balanced bilinear model}
\def\ABa{%$(A,\bbe,\ba,\f)$ \\
rank-one balanced bilinear model}
\def\bbe{B%{\bff \beta}
} \def\vbe{\vec {\bff \beta}
}
\def\vb{\vec {\bff b}}
\newcommand{\Kh}{\hat{K}}          % K with hat, first mention only
\newcommand{\wS}{w^{\sigma}}       % right Perron J^\perp, full block
\newcommand{\piS}{\vec \pi^{\sigma}}    % left Perron J^\perp, full block
\newcommand{\ws}{w^{\sigma}}       % right Perron J^\perp, sub-block
\newcommand{\pis}{\pi^{\sigma}}    % left Perron J^\perp, sub-block
\newcommand{\hwS}{\hat{w}^{\sigma}}   % right Perron K, full block
\newcommand{\hpiS}{\hat{\pi}^{\sigma}} % left Perron K, full block
\newcommand{\hws}{\hat{w}^{\sigma}}   % right Perron K, sub-block
\newcommand{\hpis}{\hat{\pi}^{\sigma}} % left Perron K, sub-block
\newcommand{\twS}{\tilde{w}^{\sigma}}  % right Perron \tilde K, full block
\newcommand{\tpiS}{\tilde{\pi}^{\sigma}} % left Perron \tilde K
\newcommand{\tws}{\tilde{w}^{\sigma}}
\newcommand{\tpis}{\tilde{\pi}^{\sigma}}
\newcommand{\Mat}[2]{\mathbb{R}_{\ge 0}^{\,#1\times #2}}
\def\tJ{J_\sigma^\perp} \def\lPe{left Perron eigenvector}
% \G is \Gamma (defined at :795).  The Volterra function keeps its own name,
% so that \G no longer has two conflicting meanings in this file.
\def\GV{G}  % Volterra G(u)=u-1-\ln u
\def\PVL{Perron-Volterra Lyapunov minimal siphon ansatz}
\newcommand{\istar}[1]{i_{#1}^{*}}
%\newcommand{\sE}{s_E}
\newcommand{\vE}{v_E}
\def\PV{Perron-Volterra candidate Lyapunov function}
\def\GLV{generalized Lotka-Volterra}
\def\MERN{mathematical epidemiology reaction networks}
\def\CMM{contraction metric method}
\newcommand{\sign}{\operatorname{sign}}
\newcommand{\Var}{\operatorname{Var}} 

% ======================================================================
% IDEAS SHARED BETWEEN THE ABSTRACT AND THE MAIN RESULTS SECTION
%
% Write each idea ONCE, in the abstract; the Main Results section recalls it,
% so the two can never drift apart.
%
%   in the abstract        \ideaitem{1}{The first-order reaction terms ...}
%   in Main Results        \ideahead{1}
%
% \ideaitem stores the text globally and typesets
%   \item \hyperlink{idea:1}{<text>}
% \ideahead typesets
%   \item\hypertarget{idea:1}{}\textbf{<the same text>}
%
% The abstract must come before the Main Results input, which is the normal
% order. If a heading has to read differently from the abstract sentence (for
% instance a long abstract item summarised in one line), override it with the
% optional argument: \ideahead[short heading]{6}.
% ======================================================================
\makeatletter
\newcommand{\ideaitem}[2]{%
  \expandafter\gdef\csname @idea@#1\endcsname{#2}%
  \item\hyperlink{idea:#1}{#2}}
\newcommand{\ideahead}[2][]{%
  \item\hypertarget{idea:#2}{}%
  \def\@ideatmp{#1}%
  \ifx\@ideatmp\@empty
    \@ifundefined{@idea@#2}{%
      \PackageWarning{preamble}{Idea #2 is used in Main Results but was never
        defined by \string\ideaitem\space in the abstract}%
      \textbf{[idea #2 undefined]}%
    }{\textbf{\begingroup\def\ideaterm##1##2{##2}\def\idearest##1{}\@nameuse{@idea@#2}\endgroup}}%
  \else
    \textbf{#1}%
  \fi}
% \ideaterm{<hypertarget>}{<words>}: inside an abstract idea the words become an
% invisible link (the document sets \hypersetup{hidelinks}) jumping to the place
% in Main Results that defines them; inside the Main Results heading itself,
% where the target sits a few lines below, \ideahead strips the link and prints
% the words plainly. Mark the destination with \hypertarget{<name>}{} on the line
% before the display, which adds no equation number and shows nothing.
\newcommand{\ideaterm}[2]{\hyperlink{#1}{#2}}
% \idearest{<text>}: the part of an abstract idea after its first sentence; printed in
% the abstract, omitted from the Main Results heading.
\newcommand{\idearest}[1]{#1}
\makeatother

% ======================================================================
% SHORT-PAPER SWITCH
% \shorttrue in a master that assembles only part of the companion paper
% (Barrat/SipArxiv.tex). \nsh{...} drops text that cites results living in the
% omitted sections, so the short build has no dangling references while the
% companion paper keeps the text verbatim.
% ======================================================================
\newif\ifshort \shortfalse
\newcommand{\nsh}[1]{\ifshort\else#1\fi}
% \shortalt{<short build>}{<companion paper>}: use when a sentence needs
% different wording in the two builds, not merely a deletion.
\newcommand{\shortalt}[2]{\ifshort#1\else#2\fi}
% \Jtr, \JtrS{\Sigma}: transverse Jacobian. J^\perp here; SipArxiv.tex redefines them as J_{\mathrm{tra}}, J_{\mathrm{tra},\Sigma}
\newcommand{\Jtr}{J^\perp}
\newcommand{\JtrS}[1]{J_{#1}^\perp}
, and must not clash.
\providecommand{\Q}{\mathbb{Q}}
\providecommand{\res}{\operatorname{res}}   % resultant; used by epiRes.tex
\newcommand{\MN}{\mathbb{N}}
\newcommand{\MZ}{\mathbb{Z}}
\newcommand{\realpos}{\mathbb{R}_{> 0}}
\newcommand{\realnneg}{\mathbb{R}_{\geq 0}}
\newcommand{\naturalpos}{\mathbb{N}_{>0}}
\newcommand{\Rplus}{\mathbb{R}_+}
\newcommand{\Rp}{\mathbb{R}_{\ge 0}}
\def\Rnne{\mathbb{R}_{\ge 0}}
\def\Se{\bar{S}} \def\Ie{\bar{I}}
\newcommand{\Rpos}{\mathbb{R}_{>0}}
\newcommand{\dd}{\,\mathrm d}
\newcommand{\e}{\mathrm e}
%% ========================================
%% GREEK LETTERS (2+ character abbreviations) greek
%% ========================================
\def\al{\alpha}\def\ba{\bff \alpha}
\def\w{\omega}\def\Lgn{\ga+\mu+\de}\def\om{\omega}
\def\de{\delta}\def\De{\Delta}\def\bde{\bff \delta}
\def\ep{\epsilon}\def\eps{\varepsilon}
\def\f{\varphi}  % Kept as requested
\def\ga{\gamma}\def\Ga{\Gamma}\def\G{\Gamma}\def\bg{\bff \gamma}
\def\up{\upsilon}  % NOT the breve accent: no accent commands anywhere.
                   % Named \up, not \u, so the LaTeX breve stays reachable.
\def\la{\lambda}\def\La{\Lambda}
\def\si{\sigma}\def\ka{\kappa}
\def\th{\theta}\def\vt{\vartheta}
\def\ze{\zeta}\def\Gap{\Gamma^{+}} \def\Gam{\Gamma^{-}}


%% ========================================
%% SPECIAL SHORT COMMANDS (kept as requested)
%% ========================================
\def\I{\infty}
\def\T{\widetilde}
\def\H{\hat}
\newcommand{\relint}{\operatorname{relint}}

%% ========================================
%% BOLD/VECTOR NOTATION
%% ========================================
\newcommand{\bff}[1]{{\mbox{\boldmath$#1$}}}
\newcommand\X{\boldsymbol{X}}
\def\x{\boldsymbol{x}}\def\bx{\bff x}\def\xe{\bar{\x}}
\newcommand\y{\boldsymbol{y}}\def\by{\boldsymbol{y}}\def\ye{\bar{\y}}
\def\yb{\y}\def\xb{\x}
\def\xy{\x-\y}
\newcommand\br{\boldsymbol{r}}\def\re{\bar{\r}}
\newcommand\bd{\boldsymbol{d}}
\newcommand\bQ{\boldsymbol{Q}}
\newcommand{\Yb}{{\mathbf{Y}}}
\newcommand{\Ab}{\mathbf{A}}
\newcommand{\zb}{\mathbf{z}}\def\bz{\bff 0}
\newcommand{\Zb}{\mathbf{Z}}
\newcommand{\bb}{\mathbf{b}}
\def\vbe{\vec {\bff \beta}}
\def\vb{\vec {\bff b}}
\def\v1{\vec {\bff 1}}
\def\vn{\vec \de}
\def\vz{\vec 0}
\def\vi{\overset{\rightarrow}{\mathsf i}}
\def\vv{\overset{\rightarrow}{v}}
\def\vr{\; \vec {\mathsf r}}
\newcommand{\va}{\bff a}
\providecommand{\vect}[1]{\boldsymbol{#1}}   % bold vector (moved from Gradostat.tex)
\providecommand{\one}{\mathbf 1}             % all-ones vector (moved from AliSF)
\providecommand{\spr}{s}                     % spectral bound  (moved from AliSF)
\providecommand{\Ke}{K_{\rm eff}}            % effective carrying capacity
% Critical right/left eigenvectors in the Hopf projection method (AAHOSN/New.tex)
\providecommand{\bq}{\mathbf{q}}
\providecommand{\bp}{\mathbf{p}}


%% ========================================
%% MATHCAL & MATHSCR
%% ========================================
% Mathcal
\newcommand\mc[1]{\mathcal{#1}}
\def\mC{{\mathcal C}}\def\mD{\mathcal D}\def\mE{\mathcal{E}}
\def\mF{\mathcal F}\def\mG{{\mathcal G}}\def\mH{{\mathcal H}}
\def\mI{\mathcal I}\def\mL{\mathcal{L}}\def\mM{{\mathcal M}}
\def\mN{{\mathcal N}}\def\mP{\mathcal{P}}\def\mR{{\mathcal R}}
\def\mS{{\mathcal S}}\def\mT{\mathcal{T}}\def\mV{{\mathcal V}}
\def\mK{\mathcal K}\def\mA{{\mathcal A}}\def\mU{{\mathcal U}}
% Symmetric bilinear and trilinear forms of the Taylor expansion, used in the
% first Lyapunov coefficient (AAHOSN/New.tex)
\providecommand{\Bform}{\mathcal{B}}
\providecommand{\Cform}{\mathcal{C}}

% Blackboard bold.  Kept here, not in the drivers: a \providecommand in a
% driver preamble runs before this file and silently wins, so the same name
% ends up meaning different things in different builds.
\def\IC{\mathbb{C}}\def\IN{\mathbb{N}}\def\IQ{\mathbb{Q}}
\def\IR{\mathbb{R}}\def\IZ{\mathbb{Z}}\def\IG{\mathbb{G}}
\def\II{\mathbb{I}}

% Mathscr
\newcommand{\Rs}{{\mathscr R}}
\newcommand{\Ss}{{\mathscr S}}
\newcommand{\Ns}{{\mathscr N}}
\newcommand{\Cs}{\mathscr C}


%% ========================================
%% OPERATORS & FUNCTIONS
%% ========================================
\DeclareMathOperator{\tr}{tr}
\DeclareMathOperator{\diag}{Diag}
\DeclareMathOperator{\Diag}{Diag}
\DeclareMathOperator{\intr}{\mathrm{Int}}
\DeclareMathOperator{\Inv}{Inv}
\DeclareMathOperator{\spec}{spec}
\DeclareMathOperator{\AUC}{AUC}               % area under the curve
\newcommand{\ii}{\mathrm{i}}                  % imaginary unit
\newcommand\op{\operatorname}
\newcommand{\supp}{\op{supp}}
\newcommand{\diff}{\mathop{}\mathopen{}\mathrm{d}}
\newcommand{\pdd}[1]{\frac{\partial}{\partial{#1}}}
\DeclareMathOperator{\cone}{cone}
\DeclareMathOperator{\proj}{proj}

%% ========================================
%% BRACKETS & NORMS
%% ========================================
\providecommand{\pr}[1]{\left(#1\right)}
\providecommand{\pp}[1]{\left[#1\right]}
\newcommand\croc[1]{\left\langle #1\right\rangle}
\newcommand{\norm}[1]{\|#1\|}
\newcommand{\inner}[2]{\langle #1, #2 \rangle}
% Short synonym of \inner, used by AAHOSN/New.tex as $\ip{u}{v}=\bar u^\top v$
\providecommand{\ip}[2]{\langle #1 , #2 \rangle}


%% ========================================
%% FORMATTING & UTILITIES
%% ========================================
\newcommand*{\Scale}[2][4]{\scalebox{#1}{$#2$}}
\newcommand{\te}[1]{\text{#1}}
\newcommand\steq[1]{\stackrel{\text{\rm #1.}}{=}}
\newcommand\blue[1]{\textcolor{blue}{#1}}
\newcommand{\red}{\textcolor[rgb]{1.00,0.00,0.00}}
\newcommand{\green}{\textcolor[rgb]{0.00,1.00,0.00}}

% Figure command
\newcommand{\figu}[3]{%
\begin{figure}[H]
\centering
\includegraphics[scale=#3]{#1}
\caption{#2\label{f:#1}}
\end{figure}
}

% Footnote
\long\def\symbolfootnote[#1]#2{%
\begingroup
\def\thefootnote{\fnsymbol{footnote}}\footnote[#1]{#2}%
\endgroup}
\def\fn{\symbolfootnote}

\newcommand{\dsum}{\displaystyle\sum}

%% ========================================
%% CRN & EPIDEMIOLOGY SPECIFIC
%% ========================================
% CRN notation
\def\SS{\mathcal S}
\def\CC{\mathcal C}
\def\RR{\mathcal R}

% Reproduction numbers
\def\Rto{R_{1}^{E_2}}
\def\Rot{R_{2}^{E_1}}
\def\Roc{R_{1,c}}
\def\Rtc{R_{2,c}}
\def\rd{r_{0}}
\def\sd{s_{0}}
\def\dd{D_{0}}
\def\vd{v_{0}}
\def\zd{z_{0}}
\def\se{\sE}
\def\xd{x_{0}}
%% ========================================
%% SPECIAL VALUES & EQUILIBRIA
%% ========================================
\def\eE{\bar{e}}
\def\wE{\bar{w}}
\def\sE{s_{*}}
\def\Me{\T M_{ee}}
\def\Rnne{\mathbb{R}_{\ge 0}}

\newcommand{\ones}{\mathbf 1}


%% ========================================
%% ABBREVIATIONS (alphabetical by category)
%% ========================================

% -- General mathematical terms --
\def\adm{admissible}\def\admy{admissibility}
%\def\and{antisymmetric}% disabled 2026-09-30: it overrode LaTeX's author separator \and (printed "antisymmetric" between authors)
\def\ass{assumption}
\def\cadlag{c\`adl\`ag}
\def\ch{characteristic polynomial}
\def\chp{characteristic polynomial}
\def\coe{coefficient}
\def\comp{computation}
\def\con{condition}
\def\coo{coordinate}
\def\corr{corresponding}
\def\dec{decomposition}
\def\den{denominator}
\def\dete{deterministic}
\def\eig{eigenvalue}\def\eval{eigenvalue}
\def\eigv{eigenvector}\def\eiv{eigenvector}\def\evec{eigenvector}
\def\elt{element}
\def\enn{essentially nonnegative/positivity preserving}
\def\equ{equation}
\def\etal{et al.}
\def\fe{for example}
\def\fno{from now on}
\def\for{\forall}
\def\form{formula}
\def\Fr{Furthermore, }
\def\frt{furthermore}
\def\fun{function}
\def\gene{generalization}
\def\heu{heuristic}
\def\how{however}\def\How{However, }
\def\ie{i.e. }
\def\im{\item}
\def\isd{is defined by}
\def\itc{it is easy to check that}
\def\itf{it follows that}\def\Itf{It follows that}
\def\ith{it holds that}
\def\itm{it may be checked that}\def\Itm{It may be checked that}
\def\ito{it turns out that}\def\Ito{It turns out that}\def\Itou{It turns out useful to}
\def\ma{mass action}\def\MAR{mass action representation}
\def\mat{matrix}
\def\mbc{may be checked to be}
\def\mbw{may be written as}\def\mubw{must be written as}
\def\mpr{more precisely, }\def\Mp{More precisely, }
\def\nc{non-catalytic}
\def\nne{non-negative}\def\nny{nonnegativity}
\def\no{\nonumber}
\def\noni{non-infectious}
\def\nota{notation}
\def\np{\newpage}
\def\Nt{Note that}
\def\nt{non trivial}\def\nz{non zero}
\def\oth{otherwise}
\def\Oth{On the other hand}
\def\pa{parameter}\def\para{parameter}
\def\pol{polynomial}
\def\prob{probability}
\def\Prf{{\bf  Proof:}}\def\prf{{\bf  Proof:}}
\def\prop{proportional to}
\def\qu{\quad}
\def\QED{\hfill {$\square$}\goodbreak \medskip}
\def\rep{representation}
\def\resp{respectively}
\def\satd{satisfied}\def\satg{satisfying}\def\sats{satisfies}\def\saty{satisfy}
\def\ssec{\subsection}\def\sssec{\subsubsection}
\def\sef{see for example}
\def\sksy{skew-symmetric}
\def\sub{substrate}
\def\sys{system}
\def\var{variable}
\def\wk{well-known}
\def\wl{without loss of generality}\def\wlo{w.l.o.g.}
\def\wmc{We may check that}\def\Wmc{We may check that}
\def\wmr{we may rewrite}
\def\wms{we must show that}
\def\wrt{with respect to}

% -- CRN & Chemistry --
\newcommand\CRN{chemical reaction network}
\newcommand\CRNs{chemical reaction networks}
\newcommand\RN{reaction network}
\def\ACR{absolute concentration robustness}
\def\AMSD{acyclic siphon decomposition}
\def\DAE{differential algebraic equations}
\def\FHJ{Feinberg-Horn-Jackson}
\def\GMAK{generalized mass action kinetics}
\def\MM{Michaelis-Menten}
\def\sm{stoichiometric matrix}\def\Sm{stoichiometric matrix}
\def\WRZD{weakly reversible, zero deficiency}\def\aut{autocatalytic}

% -- Epidemiology --
\def\brn{basic reproduction number}\def\bzn{basic replacement number}
\def\com{compartment}
\def\dfe{disease free equilibrium}\def\DFE{disease free equilibrium}
\def\dwt{dwell times}\def\edt{expected dwell times}
\def\EM{epidemic model}
\def\endp{endemic fixed point}\def\enp{endemic point}\def\EE{endemic point}
\def\ME{mathematical epidemiology}
\def\MEE{eco-epidemiology}\def\EEP{eco-epidemiological}
\def\ece{eco-epidemiological model}
\def\mgm{multi group model}
\def\mtS{multi-strain model}\def\mSt{multi-strain model}
\def\mSi{minimal siphon}
\def\NGM{next generation matrix}
\def\nim{new infections matrix}
\def\qem{quadratic epidemiological model}
\def\repF{invasion functions}
\def\SPF{SIR-PH-FA model}
\def\sus{susceptible}

% -- Graph theory & Linear algebra --
\def\jac{Jacobi matrix}
\def\jin{Jacobian matrix of the invasion vector field with respect to the invasion variables}
\def\jrn{Jacobian matrix of the resident vector field with respect to the resident variables}
\def\jxn{Jacobian matrix of the $\x$ vector field with respect to the $\x$ variables}
\def\Lm{Laplacian matrix}
\def\MTT{matrix tree theorem}
\def\SGT{spectral graph theory}
\def\SMW{Sherman–Morrison–Woodbury}

% -- Stability & Dynamics --
\def\absC{absorbtion CTMC}
\def\BTb{Bogdanov-Takens bifurcation}
\def\CEP{competitive exclusion partition}
\def\CTMC{continuous time Markov chain}
\def\ds{differential system}
\def\fp{fixed point}\def\bfp{boundary fixed point}\def\fbp{fixed boundary point}
\def\fbl{fixed boundary line}\def\fif{forward invariant face}
\def\gas{global asymptotic stability}\def\GAS{globally asymptotically stable}
\def\GASS{\text{GASS}}
\def\Hb{Hopf bifurcation}
\def\Hur{Hurwitz}\def\MH{Metzler-Hurwitz}\def\RH{Routh-Hurwitz}
\def\IS{\text{IS}}
\def\LAS{locally asymptotically stable}
\def\LdS{Lyapunov diagonally stable}
\def\Lf{Lyapunov function}\def\LyV{Lyapunov Volterra type function}\def\RLF{robust Lyapunov function}
\def\LyM{Lyapunov Malkin stable}
\def\PF{Perron-Frobenius}
\def\spfp{strictly positive fixed point}
\def\sta{stability}
\def\STP{sharp threshold property}\def\STT{sharp threshold theorem}

% -- Models (with \varphi explicitly instead of \f) --
\def\AB{$(A,B)$ Arino-Brauer epidemic models}
\def\ABa{$(A,\bbe,\ba,\varphi)$ balanced bilinear model}
\def\ABP{%$(A,\bbe,P,\varphi)$
\\ balanced bilinear model}
\def\AMG{Anderson May Gupta}
\newcommand\BIN{biological interaction network}
\def\Dt{Descartes type}
\def\gLV{generalized Lotka-Volterra model}
\def\IG{\text{invasion graph}}\def\IGMS{interaction graph on minimal siphons}
\def\LCP{\text{linear complementarity problem}}
\def\LV{Lotka-Volterra}\def\LVK{Lotka-Volterra-Kolmogorov}

% -- Splitting & Structure --
\def\reg{regular splitting}\def\regS{{\bf regular splitting}}
\def\RUR{rational univariate representation}
\def\bds{block-diagonal structure}\def\NGMs{next generation matrices}

% -- Miscellaneous --
\def\Ma{{Mathematica}}
\def\Mm{Max–Min}
\def\num{numerator}
\def\rap{rational approximation} \def\DSR{directed species-reaction bipartite graph}\def\PCP{polynomial complementarity problem}\def\rat{rational substitution}
\def\QSS{quasi steady state} \def\RHS{right hand side} \def\LHS{left hand side}
\def\OSN{online social network}\def\VdD{Van den Driessche}
\def\wli{$\omega$-limit} \def\rk{rich kinetics}\def\sJ{symbolic Jacobian}
\def\cV{\cite{Vas,VasSta,BSV,GSV}} \def\spf{semi-parametric semi-functional}
\def\CRW{Capasso-Ruan-Wang} \def\ABP{$(A,\bbe,P,\f)$ balanced bilinear model}
\def\ABa{%$(A,\bbe,\ba,\f)$ \\
rank-one balanced bilinear model}
\def\bbe{B%{\bff \beta}
} \def\vbe{\vec {\bff \beta}
}
\def\vb{\vec {\bff b}}
\newcommand{\Kh}{\hat{K}}          % K with hat, first mention only
\newcommand{\wS}{w^{\sigma}}       % right Perron J^\perp, full block
\newcommand{\piS}{\vec \pi^{\sigma}}    % left Perron J^\perp, full block
\newcommand{\ws}{w^{\sigma}}       % right Perron J^\perp, sub-block
\newcommand{\pis}{\pi^{\sigma}}    % left Perron J^\perp, sub-block
\newcommand{\hwS}{\hat{w}^{\sigma}}   % right Perron K, full block
\newcommand{\hpiS}{\hat{\pi}^{\sigma}} % left Perron K, full block
\newcommand{\hws}{\hat{w}^{\sigma}}   % right Perron K, sub-block
\newcommand{\hpis}{\hat{\pi}^{\sigma}} % left Perron K, sub-block
\newcommand{\twS}{\tilde{w}^{\sigma}}  % right Perron \tilde K, full block
\newcommand{\tpiS}{\tilde{\pi}^{\sigma}} % left Perron \tilde K
\newcommand{\tws}{\tilde{w}^{\sigma}}
\newcommand{\tpis}{\tilde{\pi}^{\sigma}}
\newcommand{\Mat}[2]{\mathbb{R}_{\ge 0}^{\,#1\times #2}}
\def\tJ{J_\sigma^\perp} \def\lPe{left Perron eigenvector}
% \G is \Gamma (defined at :795).  The Volterra function keeps its own name,
% so that \G no longer has two conflicting meanings in this file.
\def\GV{G}  % Volterra G(u)=u-1-\ln u
\def\PVL{Perron-Volterra Lyapunov minimal siphon ansatz}
\newcommand{\istar}[1]{i_{#1}^{*}}
%\newcommand{\sE}{s_E}
\newcommand{\vE}{v_E}
\def\PV{Perron-Volterra candidate Lyapunov function}
\def\GLV{generalized Lotka-Volterra}
\def\MERN{mathematical epidemiology reaction networks}
\def\CMM{contraction metric method}
\newcommand{\sign}{\operatorname{sign}}
\newcommand{\Var}{\operatorname{Var}} 

% ======================================================================
% IDEAS SHARED BETWEEN THE ABSTRACT AND THE MAIN RESULTS SECTION
%
% Write each idea ONCE, in the abstract; the Main Results section recalls it,
% so the two can never drift apart.
%
%   in the abstract        \ideaitem{1}{The first-order reaction terms ...}
%   in Main Results        \ideahead{1}
%
% \ideaitem stores the text globally and typesets
%   \item \hyperlink{idea:1}{<text>}
% \ideahead typesets
%   \item\hypertarget{idea:1}{}\textbf{<the same text>}
%
% The abstract must come before the Main Results input, which is the normal
% order. If a heading has to read differently from the abstract sentence (for
% instance a long abstract item summarised in one line), override it with the
% optional argument: \ideahead[short heading]{6}.
% ======================================================================
\makeatletter
\newcommand{\ideaitem}[2]{%
  \expandafter\gdef\csname @idea@#1\endcsname{#2}%
  \item\hyperlink{idea:#1}{#2}}
\newcommand{\ideahead}[2][]{%
  \item\hypertarget{idea:#2}{}%
  \def\@ideatmp{#1}%
  \ifx\@ideatmp\@empty
    \@ifundefined{@idea@#2}{%
      \PackageWarning{preamble}{Idea #2 is used in Main Results but was never
        defined by \string\ideaitem\space in the abstract}%
      \textbf{[idea #2 undefined]}%
    }{\textbf{\begingroup\def\ideaterm##1##2{##2}\def\idearest##1{}\@nameuse{@idea@#2}\endgroup}}%
  \else
    \textbf{#1}%
  \fi}
% \ideaterm{<hypertarget>}{<words>}: inside an abstract idea the words become an
% invisible link (the document sets \hypersetup{hidelinks}) jumping to the place
% in Main Results that defines them; inside the Main Results heading itself,
% where the target sits a few lines below, \ideahead strips the link and prints
% the words plainly. Mark the destination with \hypertarget{<name>}{} on the line
% before the display, which adds no equation number and shows nothing.
\newcommand{\ideaterm}[2]{\hyperlink{#1}{#2}}
% \idearest{<text>}: the part of an abstract idea after its first sentence; printed in
% the abstract, omitted from the Main Results heading.
\newcommand{\idearest}[1]{#1}
\makeatother

% ======================================================================
% SHORT-PAPER SWITCH
% \shorttrue in a master that assembles only part of the companion paper
% (Barrat/SipArxiv.tex). \nsh{...} drops text that cites results living in the
% omitted sections, so the short build has no dangling references while the
% companion paper keeps the text verbatim.
% ======================================================================
\newif\ifshort \shortfalse
\newcommand{\nsh}[1]{\ifshort\else#1\fi}
% \shortalt{<short build>}{<companion paper>}: use when a sentence needs
% different wording in the two builds, not merely a deletion.
\newcommand{\shortalt}[2]{\ifshort#1\else#2\fi}
% \Jtr, \JtrS{\Sigma}: transverse Jacobian. J^\perp here; SipArxiv.tex redefines them as J_{\mathrm{tra}}, J_{\mathrm{tra},\Sigma}
\newcommand{\Jtr}{J^\perp}
\newcommand{\JtrS}[1]{J_{#1}^\perp}
.
\makeatletter
\AtBeginDocument{\@ifundefined{hypersetup}{}{\hypersetup{linktoc=all}}}
\makeatother
\usepackage{graphicx,float,fancybox}
\usepackage{epsfig,epstopdf,tikz}
\usepackage[sans]{dsfont}
\usetikzlibrary{positioning,arrows,arrows.meta,calc}
% Configure listings for Mathematica code
\lstdefinestyle{mathematica}{
 language=Mathematica,
 basicstyle=\small\ttfamily,
 backgroundcolor=\color{gray!10},
 frame=single,
 breaklines=true,
 commentstyle=\color{green!60!black},
 keywordstyle=\color{blue},
 stringstyle=\color{red},
 showstringspaces=false,
 tabsize=2
}
\usepackage{tcolorbox}
% Colored boxes for key results
\newtcolorbox{keyresult}[1][]{
 colback=blue!5!white,
 colframe=blue!75!black,
 title=#1,
 fonttitle=\bfseries
}
\newtcolorbox{examplebox}[1][]{
 colback=green!5!white,
 colframe=green!75!black,
 title=#1,
 fonttitle=\bfseries
}


\makeatletter
\@ifundefined{theorem}   {\newtheorem{theorem}   {Theorem}}   {}
\@ifundefined{definition}{\newtheorem{definition}{Definition}} {}
\@ifundefined{open}      {\newtheorem{open}      {Problem}}   {}
\@ifundefined{example}   {\newtheorem{example}   {Example}}   {}
\makeatother

% COUNTER DEFINITIONS
\newcounter{tabcounter}
\newcounter{defcounter}
\newcounter{opcounter}
\newcounter{thmcounter}
\newcounter{excounter}
\newcounter{ascounter}
\newcounter{cjcounter}
\newcounter{lemcounter}
\newcounter{remcounter}
\newcounter{corcounter}
\newcounter{propcounter}
\newcounter{figcounter}

% UNIFIED GLOSSARY STYLE
\newglossarystyle{unifiedstyle}{%
  \setglossarystyle{list}%
  \renewcommand*{\glossarypreamble}{\gdef\glslastsec{}}%
  \renewcommand{\glossentry}[2]{%
    \edef\glscursec{\glsentryuseriii{##1}}%
    \ifx\glscursec\glslastsec\else
      \item[]\textbf{\S\,\glscursec\enspace\glsentryuseriv{##1}}%
      \global\let\glslastsec\glscursec
    \fi
    \item[\glsentryuserii{##1} \glsentryuseri{##1}:] ``\glossentryname{##1}'' \glsentrydesc{##1}%
  }%
}

\makeatletter
% First-word extractor (used by \beT, \beD and \beXa for the default key)
% \@firstwordof hello world \relax  ->  hello
\def\@firstwordof#1 #2\relax{#1}
% First three letters, used for the heading mnemonic of a theorem given no
% explicit key.  Padding with \@empty keeps it safe for first words shorter
% than three letters: \@empty vanishes under \edef, whereas \relax would
% survive into the label and is illegal inside \csname (it broke every
% two-letter title, e.g. "ME ...", with "Missing \endcsname inserted").
\def\@takethree#1#2#3#4\@nil{#1#2#3}

% EXAMPLE COMMAND -- same signature and behaviour as \beT and \beD.
% Usage: \beXa*[long title][short title][key]
%   star  = suppresses the printed "[key]" tag after the title (as for \beT*
%   and \beD*); the label and glossary entry are unaffected.
%   key = optional label key override; the label becomes e:KEY, and without a
%   key it is e: plus the first three letters of the title's first word.  A
%   mnemonic label lets an example be cited as "Example~\ref{e:KEY}" instead
%   of by section number, which is what the citing prose needs.
\NewDocumentCommand{\beXa}{s O{} O{} o}{%
  \def\@temparg{#2}%
  \ifx\@temparg\@empty
    \begin{example}%
  \else
    \let\@savedSecTitle\@currentsectionname%
    \refstepcounter{excounter}%
    \IfValueTF{#4}{%
      \edef\@exlblkey{#4}%
    }{%
      \edef\@exfirstword{\@firstwordof#2 \relax}%
      \edef\@exlblkey{%
        \expandafter\@takethree\@exfirstword\@empty\@empty\@empty\@nil}%
    }%
    \def\@tempshort{#3}%
    \ifx\@tempshort\@empty
      \expandafter\xdef\csname @xshort@\@exlblkey\endcsname{\@exlblkey}%
    \else
      \expandafter\def\csname @xshort@\@exlblkey\endcsname{#3}%
    \fi
    \IfBooleanTF{#1}{\let\@extag\@empty}{%
      \edef\@extag{\noexpand\ \noexpand\textup{[\noexpand\texttt{e:\@exlblkey}]}}}%
    \begin{example}[#2\ifshowkeys\@extag\fi]%
    \label{e:\@exlblkey}%
    % record the long title so \nref{e:KEY} prints it, as for theorems
    \protected@write\@auxout{}{%
      \string\@thmtitle{e:\@exlblkey}{\unexpanded{#2}}}%
    \expanded{%
      \noexpand\newglossaryentry{e\theexcounter}{%
        name={\unexpanded{#2}%
          \IfBooleanF{#1}{ [\noexpand\texttt{e:\@exlblkey}]}},%
        sort={\theexample},%
        description={p.~\noexpand\pageref{e:\@exlblkey}},%
        user1={\theexample},%
        user2={Example},%
        user3={\thesection},%
        user4={\expandafter\unexpanded\expandafter{\@savedSecTitle}}%
      }%
    }%
    \protected@write\@auxout{}{\string\glsliveentry{e\theexcounter}}%
  \fi
}
\newcommand{\eeXa}{\end{example}}

% DEFINITION COMMAND -- same signature and behaviour as \beT below.
% Usage: \beD[long title][short title][key]
%   long title  = definition heading (plain text only)
%   short title = shown in \dref{word}; defaults to the key
%   key         = optional label key override; the label becomes d:KEY
% Without a key the label is d: plus the first three letters of the title's
% first word, so it is a mnemonic rather than the definition's number: a
% numeric label would silently change whenever a definition is inserted.
\NewDocumentCommand{\beD}{s O{} O{} o}{%
  \def\@temparg{#2}%
  \ifx\@temparg\@empty
    \begin{definition}%
  \else
    \let\@savedSecTitle\@currentsectionname%
    \refstepcounter{defcounter}%
    \IfValueTF{#4}{%
      \edef\@deflblkey{#4}%
    }{%
      \edef\@deffirstword{\@firstwordof#2 \relax}%
      \edef\@deflblkey{%
        \expandafter\@takethree\@deffirstword\@empty\@empty\@empty\@nil}%
    }%
    \def\@tempshort{#3}%
    \ifx\@tempshort\@empty
      \expandafter\xdef\csname @dshort@\@deflblkey\endcsname{\@deflblkey}%
    \else
      \expandafter\def\csname @dshort@\@deflblkey\endcsname{#3}%
    \fi
    \IfBooleanTF{#1}{\let\@deftag\@empty}{%
      \edef\@deftag{\noexpand\ \noexpand\textup{[\noexpand\texttt{d:\@deflblkey}]}}}%
    \begin{definition}[#2\ifshowkeys\@deftag\fi]%
    \label{d:\@deflblkey}%
    % record the long title so \nref{d:KEY} prints it, as for theorems
    \protected@write\@auxout{}{%
      \string\@thmtitle{d:\@deflblkey}{\unexpanded{#2}}}%
    \expanded{%
      \noexpand\newglossaryentry{d\thedefcounter}{%
        name={\unexpanded{#2}%
          \IfBooleanF{#1}{ [\noexpand\texttt{d:\@deflblkey}]}},%
        sort={\thedefinition},%
        description={p.~\noexpand\pageref{d:\@deflblkey}},%
        user1={\thedefinition},%
        user2={Definition},%
        user3={\thesection},%
        user4={\expandafter\unexpanded\expandafter{\@savedSecTitle}}%
      }%
    }%
    \protected@write\@auxout{}{\string\glsliveentry{d\thedefcounter}}%
  \fi
}
\newcommand{\eeD}{\end{definition}}

% \showkeysfalse in a driver hides the [t:KEY] mnemonics in result headings
% (labels and glossary entries are unchanged); used for submission builds.
\newif\ifshowkeys \showkeystrue
% \commonredtrue in a driver that inputs basicFiles/commonRed: tool.tex then
% cites its reduction results instead of restating them.
\newif\ifcommonred \commonredfalse
% TheorEM COMMAND
% Usage: \beT[long title][short title][key]
%   long title  = theorem heading (plain text only)
%   short title = shown in \tref{word}; may contain math; defaults to first word of long title
%   key         = optional label key override; label becomes t:KEY (avoids first-word collisions)
% Label is t:WORD where WORD = first word of long title (or KEY if given); use \tref{WORD} to reference.
\NewDocumentCommand{\beT}{s O{} O{} o}{%
  \def\@temparg{#2}%
  \ifx\@temparg\@empty
    \begin{theorem}%
  \else
    \let\@savedSecTitle\@currentsectionname%
    \refstepcounter{thmcounter}%
    % \@thmtag is the mnemonic shown in the heading; every theorem gets one.
    % With an explicit key the key itself is shown.  Without one the key
    % defaults to the title's first word, and showing that whole word would
    % merely repeat the title ("Dissipativity [Dissipativity]"), so only its
    % first three letters are shown: "Dissipativity [Dis]".  The label is
    % unaffected -- it still uses the full \@thmlblkey.
    \IfValueTF{#4}{%
      \edef\@thmlblkey{#4}%
    }{%
      \edef\@thmfirstword{\@firstwordof#2 \relax}%
      \edef\@thmlblkey{%
        \expandafter\@takethree\@thmfirstword\@empty\@empty\@empty\@nil}%
    }%
    \IfBooleanTF{#1}{\let\@thmtag\@empty}{%
      \edef\@thmtag{\noexpand\ \noexpand\textup{[\noexpand\texttt{t:\@thmlblkey}]}}}%
    \def\@tempshort{#3}%
    \ifx\@tempshort\@empty
      \expandafter\xdef\csname @tshort@\@thmlblkey\endcsname{\@thmlblkey}%
    \else
      \expandafter\def\csname @tshort@\@thmlblkey\endcsname{#3}%
    \fi
    \begin{theorem}[#2\ifshowkeys\@thmtag\fi]%
    \label{t:\@thmlblkey}%
    % record the long title for \nref; via the .aux so forward refs resolve
    \protected@write\@auxout{}{%
      \string\@thmtitle{t:\@thmlblkey}{\unexpanded{#2}}}%
    \expanded{%
      \noexpand\newglossaryentry{t\thethmcounter}{%
        name={\unexpanded{#2}%
          \IfBooleanF{#1}{ [\noexpand\texttt{t:\@thmlblkey}]}},%
        % sort/user1 must be the number the document actually prints, i.e. the
        % theorem environment's counter, not thmcounter.  thmcounter advances
        % only on \beT with a title, so a bare \beT (or \beT with an
        % empty title) makes the two drift and the glossary cites a wrong
        % number.  thmcounter is still fine as the entry key: it is unique.
        sort={\thetheorem},%
        description={p.~\noexpand\pageref{t:\@thmlblkey}},%
        user1={\thetheorem},%
        user2={Theorem},%
        user3={\thesection},%
        user4={\expandafter\unexpanded\expandafter{\@savedSecTitle}}%
      }%
    }%
    \protected@write\@auxout{}{\string\glsliveentry{t\thethmcounter}}%
  \fi
}
\newcommand{\eeT}{\end{theorem}}

% CROSS-PART RELATION PARAGRAPH
% Usage: \xpart{text}  -- place immediately BEFORE \beT, never inside it
%   (anything inside \beT...\eeT would be swallowed into the glossary entry).
% The paragraph is typeset only when Part I is actually present, i.e. in the
% combined book build; in the standalone builds (GranadaDSABNS/A.tex,
% AAHRah-Submission/AAHRah.tex) the label part:glv is undefined and the whole
% paragraph disappears rather than leaving ?? cross-references. Same
% graceful-degradation principle as \GASref in GASt.tex. Resolves on the
% second LaTeX pass, since r@part:glv comes from the .aux file.
% (Already inside the \makeatletter region opened above, so no \makeatletter
% pair here -- adding one would close that region early.)
\newcommand{\xpart}[1]{\@ifundefined{r@part:glv}{}{\paragraph{Relation with Parts~I--II.}#1}}

% \bookref{label}{precise text}{fallback text}: prints the precise cross-reference
% when the label resolves (combined book build), and the fallback text otherwise
% (standalone paper builds, where book-only appendix labels are absent). Same
% graceful-degradation role as \GASref in GASt.tex, but usable in any file since it
% lives in this preamble (loaded by every driver) -- unlike \GASref, which is defined in a
% content file and is therefore undefined in standalone builds that omit that file.
\newcommand{\bookref}[3]{\@ifundefined{r@#1}{#3}{#2}}

% \PVpart and \RBpart: the two parts that a partial build may omit.
% They print "Part~N" in any build that contains the part, and a naming phrase
% in a build that does not, so an excerpt shows no ?? .  Each already supplies
% the word "Part", so write "\PVpart develops ...", not "Part~\PVpart".
\newcommand{\PVpart}{\bookref{part:aahrah}%
  {Part~\ref{part:aahrah}}{the Perron--Volterra part of the full book}}
\newcommand{\RBpart}{\bookref{part:recbrowne}%
  {Part~\ref{part:recbrowne}}{the immuno-virology part of the full book}}

% OPEN PROBLEM COMMANDS
% ASSUMPTION COMMANDS (glossaried)
\newcommand{\beAs}[1][]{%
  \def\@temparg{#1}%
  \ifx\@temparg\@empty
    \begin{assumption}%
  \else
    \let\@savedSecTitle\@currentsectionname%
    \refstepcounter{ascounter}%
    \begin{assumption}[#1]%
    \label{as:\theascounter}%
    \expanded{%
      \noexpand\newglossaryentry{as\theascounter}{%
        name={\unexpanded{#1}},%
        sort={\theascounter},%
        description={p.~\noexpand\pageref{as:\theascounter}},%
        user1={\theascounter},%
        user2={Assumption},%
        user3={\thesection},%
        user4={\expandafter\unexpanded\expandafter{\@savedSecTitle}}%
      }%
    }%
    \protected@write\@auxout{}{\string\glsliveentry{as\theascounter}}%
  \fi
}
\newcommand{\eeAs}{\end{assumption}}
% CONJECTURE COMMANDS (glossaried)
\newcommand{\beCj}[1][]{%
  \def\@temparg{#1}%
  \ifx\@temparg\@empty
    \begin{conjecture}%
  \else
    \let\@savedSecTitle\@currentsectionname%
    \refstepcounter{cjcounter}%
    \begin{conjecture}[#1]%
    \label{cj:\thecjcounter}%
    \expanded{%
      \noexpand\newglossaryentry{cj\thecjcounter}{%
        name={\unexpanded{#1}},%
        sort={\thecjcounter},%
        description={p.~\noexpand\pageref{cj:\thecjcounter}},%
        user1={\thecjcounter},%
        user2={Conjecture},%
        user3={\thesection},%
        user4={\expandafter\unexpanded\expandafter{\@savedSecTitle}}%
      }%
    }%
    \protected@write\@auxout{}{\string\glsliveentry{cj\thecjcounter}}%
  \fi
}
\newcommand{\eeCj}{\end{conjecture}}
% LEMMA COMMANDS
% ---------------------------------------------------------------------
% UNIFORM RESULT WRAPPERS -- \beL, \beC, \beP, \beR share \beT's signature
%
%   \beX[long title][short title][key]  ...  \eeX
%
%   long title  = heading (math allowed)
%   short title = shown in the cross-reference; defaults to the key
%   key         = label key.  With a key the label is X:KEY and the heading
%                 carries the mnemonic [key], exactly as for \beT.  Without
%                 one the label falls back to the running counter, which is
%                 what every pre-existing call site in the book relies on.
%
% A bare \beX -- arguments commented out -- opens the plain environment and
% creates NO label and NO glossary entry.  That is the supported way to keep
% an individual result out of the catalogue.
% ---------------------------------------------------------------------
% #1 wrapper name  #2 environment  #3 label prefix  #4 counter  #5 type name
\newcommand{\@mkresult}[5]{%
  % A leading * suppresses the [key] mnemonic in both the heading and the
  % glossary name, while still making the label X:KEY and the entry:
  %     \beX*[long][short][key]
  \expandafter\NewDocumentCommand\csname be#1\endcsname{s O{} O{} o}{%
    \def\@temparg{##2}%
    \ifx\@temparg\@empty
      \begin{#2}%
    \else
      \let\@savedSecTitle\@currentsectionname%
      \refstepcounter{#4}%
      % With a key: mnemonic label X:KEY and the [key] tag in the heading.
      % Without one: fall back to the running counter, exactly as these
      % wrappers behaved before, so every legacy one-argument call site in
      % the book keeps its label.  The counter is also csname-safe, which a
      % title's first word is not -- titles beginning with math would
      % otherwise build \csname \Omega...\endcsname and fail.
      \IfValueTF{##4}{%
        \edef\@relbl{##4}%
        \IfBooleanTF{##1}{\let\@retag\@empty}{%
          \edef\@retag{\noexpand\ \noexpand\textup{[\noexpand\texttt{#3:##4}]}}}%
      }{%
        \edef\@relbl{\arabic{#4}}%
        \let\@retag\@empty
      }%
      \def\@tempshort{##3}%
      \ifx\@tempshort\@empty
        \expandafter\xdef\csname @sshort@#3\@relbl\endcsname{\@relbl}%
      \else
        \expandafter\def\csname @sshort@#3\@relbl\endcsname{##3}%
      \fi
      \begin{#2}[##2\ifshowkeys\@retag\fi]%
      \edef\@tmplbl{#3:\@relbl}\expandafter\label\expandafter{\@tmplbl}%
      \expanded{%
        \noexpand\newglossaryentry{#3\csname the#4\endcsname}{%
          name={\unexpanded{##2}%
            \IfBooleanF{##1}{\IfValueT{##4}{ [\noexpand\texttt{#3:##4}]}}},%
          sort={\csname the#2\endcsname},%
          description={p.~\noexpand\pageref{#3:\@relbl}},%
          user1={\csname the#2\endcsname},%
          user2={#5},%
          user3={\thesection},%
          user4={\expandafter\unexpanded\expandafter{\@savedSecTitle}}%
        }%
      }%
      \protected@write\@auxout{}{\string\glsliveentry{#3\csname the#4\endcsname}}%
      \protected@write\@auxout{}{\string\@thmtitle{\@tmplbl}{\unexpanded{##2}}}%
    \fi
  }%
  \expandafter\newcommand\csname ee#1\endcsname{\end{#2}}%
}
\@mkresult{L}{lemma}{l}{lemcounter}{Lemma}
% REMARK / COROLLARY / PROPOSITION -- same signature as \beL above
\@mkresult{R}{remark}{r}{remcounter}{Remark}
\@mkresult{C}{corollary}{c}{corcounter}{Corollary}
\@mkresult{P}{proposition}{p}{propcounter}{Proposition}
% OPEN-PROBLEM COMMAND -- generated with the same signature as \beT:
%   \beO[long title][short title][key]   ...   \eeO
% label o:KEY, or o:<counter> when no key is given; a bare \beO makes
% neither a label nor a glossary entry.
\@mkresult{O}{open}{o}{opcounter}{Problem}


% Section-level title tracker: updated only by \section, not \subsection.
% \@currentlabelname is overwritten by every \refstepcounter (including
% \subsection), so we save it separately here.
\newcommand{\@currentsectionname}{}
\AtBeginDocument{%
  \let\@epidcrn@orig@refstepcounter\refstepcounter
  \renewcommand{\refstepcounter}[1]{%
    \@epidcrn@orig@refstepcounter{#1}%
    \ifstrequal{#1}{section}{\global\let\@currentsectionname\@currentlabelname}{}%
  }%
}

% TABLE COMMANDS (glossaried) -- consistent naming with \beT, \beD, etc.
\newcommand{\beTa}[1][]{%
  \def\@temparg{#1}%
  \ifx\@temparg\@empty
    \begin{table}[htbp]%
  \else
    \let\@savedSecTitle\@currentsectionname%
    \refstepcounter{tabcounter}%
    \begin{table}[htbp]%
    \label{A:\thetabcounter}%
    \expanded{%
      \noexpand\newglossaryentry{A\thetabcounter}{%
        name={\unexpanded{#1}},%
        sort={\thetabcounter},%
        description={p.~\noexpand\pageref{A:\thetabcounter}},%
        user1={\thetabcounter},%
        user2={Table},%
        user3={\thesection},%
        user4={\expandafter\unexpanded\expandafter{\@savedSecTitle}}%
      }%
    }%
    \protected@write\@auxout{}{\string\glsliveentry{A\thetabcounter}}%
  \fi
}
\newcommand{\eeTa}{\end{table}}

% FIGURE COMMANDS (glossaried) -- consistent with \beT, \beR, \beTa, etc.
% Usage:
%   \beF[Glossary title]...\caption{Visible caption}\eeF  -> glossaried figure
%   \beF...\caption{Visible caption}\eeF                  -> plain [H] figure, no glossary
\newcommand{\beF}[1][]{%
  \def\@temparg{#1}%
  \ifx\@temparg\@empty
    \begin{figure}[H]%
  \else
    \begin{figure}[H]%
    \let\@savedSecTitle\@currentsectionname%
    \refstepcounter{figcounter}%
    \label{fig:\thefigcounter}%
    \expanded{%
      \noexpand\newglossaryentry{Fig\thefigcounter}{%
        name={\unexpanded{#1}},%
        sort={\thefigcounter},%
        description={p.~\noexpand\pageref{fig:\thefigcounter}},%
        user1={\thefigcounter},%
        user2={Figure},%
        user3={\thesection},%
        user4={\expandafter\unexpanded\expandafter{\@savedSecTitle}}%
      }%
    }%
    \protected@write\@auxout{}{\string\glsliveentry{Fig\thefigcounter}}%
  \fi
}
\newcommand{\eeF}{\end{figure}}
\makeatother
% LIVE-ENTRY REGISTER
% \glsaddall prints every entry the glossaries package currently knows, and
% that set includes entries restored from \jobname.glsdefs.  The package
% rewrites .glsdefs from all *currently defined* entries, so anything it
% restored is written straight back out: a deleted section kept reappearing in
% the catalogue for ever, and the .glsdefs had to be erased by hand.  The
% wrappers therefore record every entry they actually create in the .aux, and
% \printboth adds only those.  The .aux is rebuilt from scratch on each run, so
% this is self-healing: a removed result is gone after one further run.
\makeatletter
\newcommand{\@glslive}{}
\newcommand{\@glsaddlive}[1]{\ifglsentryexists{#1}{\glsadd{#1}}{}}
\newcommand{\glsliveentry}[1]{\g@addto@macro\@glslive{\@glsaddlive{#1}}}
\makeatother
% PRINT COMMAND
% No \clearpage here -- the catalogue continues on the current page.  This has
% been re-added by mistake several times; leave it out.
\makeatletter
\newcommand{\printboth}{%
  \@glslive   % add exactly the entries this document created last run
  \printglossary[style=unifiedstyle, title={0 \hspace{.22em} Definitions, Theorems, Examples and Open Problems}]%
}
\makeatother
% Named theorem reference: \tref{word} -> Thm~N(word-or-short)
% Shows stored short title if available, otherwise shows the key itself.
\newcommand{\tref}[1]{%
  Thm~\ref{t:#1}(\ifcsname @tshort@#1\endcsname\csname @tshort@#1\endcsname\else#1\fi)%
}

% Full-title theorem reference: \nref{t:KEY} prints the number and then the
% theorem's long title, so one writes "Theorem~\nref{t:blocks}" and gets
% "Theorem 3 ``Block factorization and siphons''".  Takes the complete label, like
% \ref, not the bare key like \tref.
%
% \beT writes each title to the .aux (see \@thmtitle below), so \nref resolves
% forward references from the second LaTeX pass, exactly as \ref does.  If the
% title is not yet known -- first pass, or a label that is not a \beT theorem
% -- it degrades to the number rather than producing ??.
% These two use @ in their names, and this point in the file is outside the
% \makeatletter region opened earlier (it closes at \makeatother above), so
% the region has to be reopened here.
\makeatletter
\newcommand{\@thmtitle}[2]{\global\@namedef{@tlong@#1}{#2}}
\newcommand{\nref}[1]{%
  \ref{#1}%
  \ifcsname @tlong@#1\endcsname
    \ ``\csname @tlong@#1\endcsname''%
  \fi
}
% \nm{t:KEY}: the stored long title of any \beT/\beL/\beP/\beR/\beC/\beD/\beE result,
% verbatim and without quotes, so item names follow the macro's first argument.
\newcommand{\nm}[1]{%
  \ifcsname @tlong@#1\endcsname
    \csname @tlong@#1\endcsname%
  \else
    \ref{#1}%
  \fi
}
\makeatother




%% ========================================
%% Other THEOREM ENVIRONMENTS
%% ========================================
\makeatletter
\@ifundefined{lemma}      {\newtheorem{lemma}      {Lemma}}      {}
\@ifundefined{proposition}{\newtheorem{proposition}{Proposition}} {}
\@ifundefined{corollary}  {\newtheorem{corollary}  {Corollary}}  {}
\@ifundefined{conjecture} {\newtheorem{conjecture} {Conjecture}}  {}
\@ifundefined{remark}     {\newtheorem{remark}     {Remark}}     {}
\@ifundefined{question}   {\newtheorem{question}   {Question}}   {}
\@ifundefined{assumption} {\newtheorem{assumption} {Assumption}} {}
\@ifundefined{hypothesis} {\newtheorem{hypothesis} {Hypothesis}} {}
\@ifundefined{fact}       {\newtheorem{fact}       {Fact}}       {}
\makeatother

% \citer: cite the relay-graph threshold theorem of AAHrel. If a document has
% \input the chapter/file containing t:relayRF (so the label resolves), point
% straight at it; otherwise fall back to a bare citation of AAHrel. Shared here
% (not per-paper) so it resolves consistently across every book chapter.
% \citeR: same idea, for the first theorem of AAHRah (t:Pstar).
\makeatletter
\providecommand{\citer}{%
  \@ifundefined{r@t:relayRF}%
    {\cite{AAHrel}}%
    {Theorem~\ref{t:relayRF} of \cite{AAHrel}}%
}
\providecommand{\citeR}{%
  \@ifundefined{r@t:Pstar}%
    {\cite{AAHRah}}%
    {Theorem~\ref{t:Pstar} of \cite{AAHRah}}%
}
\makeatother


% Abbreviations
\def\beQ{\begin{question}}\def\eeQ{\end{question}}
\def\beA{\begin{assumption}}\def\eeA{\end{assumption}}


%% ========================================

%% ========================================
%% COLORS & CUSTOM COMMANDS
%% ========================================
\definecolor{funccolor}{RGB}{25,25,112}
\definecolor{desccolor}{RGB}{64,64,64}
\newcommand{\func}[1]{\texttt{\textcolor{funccolor}{#1}}}
\def\eq{=}
\def\mm{-}

%% ========================================
%% MATHCAL SHORTHANDS
%% ========================================
% Project-wide convention: \mX for a capital X is always \mathcal X.
% Defined for the whole alphabet so the convention never fails to expand.
\newcommand{\mA}{\mathcal A}\newcommand{\mB}{\mathcal B}
\newcommand{\mC}{\mathcal C}\newcommand{\mD}{\mathcal D}
\newcommand{\mE}{\mathcal E}\newcommand{\mF}{\mathcal F}
\newcommand{\mG}{\mathcal G}\newcommand{\mH}{\mathcal H}
\newcommand{\mI}{\mathcal I}\newcommand{\mJ}{\mathcal J}
\newcommand{\mK}{\mathcal K}\newcommand{\mL}{\mathcal L}
\newcommand{\mM}{\mathcal M}\newcommand{\mN}{\mathcal N}
\newcommand{\mO}{\mathcal O}\newcommand{\mP}{\mathcal P}
\newcommand{\mQ}{\mathcal Q}\newcommand{\mR}{\mathcal R}
\newcommand{\mS}{\mathcal S}\newcommand{\mT}{\mathcal T}
\newcommand{\mU}{\mathcal U}\newcommand{\mV}{\mathcal V}
\newcommand{\mW}{\mathcal W}\newcommand{\mX}{\mathcal X}
\newcommand{\mY}{\mathcal Y}\newcommand{\mZ}{\mathcal Z}

%% ========================================
%% BASIC MATH ENVIRONMENTS
%% ========================================
\def\bep{\begin{pmatrix}}\def\eep{\end{pmatrix}}
\def\bev{\begin{vmatrix}}\def\eev{\end{vmatrix}}
\def\bea{\begin{eqnarray*}}\def\eea{\end{eqnarray*}}
\def\bc{\begin{cases}}\def\ec{\end{cases}}
\def\BEN{\begin{enumerate}}\def\EEN{\end{enumerate}}
\def\BI{\begin{itemize}}\def\EI{\end{itemize}}

\newcommand{\be}[1]{\begin{equation}\label{#1}}
\newcommand{\ee}{\end{equation}}
\newcommand{\beq}{\begin{eqnarray}}
\def\eeq{\end{eqnarray}}
\def\eqr{\eqref}\def\fr{\frac}\def\lbl{\label}
\def\Lra{\Longrightarrow}\def\non{\nonumber}\def\cNGM{\cite{Diek,Van,Van08}}
\def\Eq{\Leftrightarrow}\def\b{\beta}\def\uph{\red{UP TO HERE}}
%% ========================================
%% PROBABILITY & STATISTICS
%% ========================================
\renewcommand{\P}{\ensuremath{\mathbb{P}}}  % Note: overrides paragraph symbol
\newcommand{\E}{\ensuremath{\mathbb{E}}}
\newcommand{\ind}[1]{\ensuremath{\mathbbm{1}_{\left\{#1\right\}}}}
\DeclareMathOperator{\Cov}{Cov}
\DeclareMathOperator{\TV}{TV}
\newcommand{\rank}{\operatorname{rank}}

%% ========================================
%% CRN OPERATORS (organizations/ERCs/buffering/siphons framework)
%% ========================================
\def\clos{\mathrm{clos}}
\def\Supply{\mathrm{Supply}}
\def\req{\mathrm{req}}
\def\React{\mathrm{React}}
\def\Prod{\mathrm{Prod}}
\def\Gen{\mathrm{Gen}}
\def\In{\mathrm{In}}
\def\Out{\mathrm{Out}}
\def\Act{\mathrm{Act}}
\def\Min{\mathrm{Min}}

%% ========================================
%% LINE BREAKING
%% ========================================
% Last-resort stretch, used only on paragraphs TeX cannot otherwise break:
% a long \texttt{} token, a comma-separated inline math list (commas are not
% breakpoints in math), or a run of \cite/\ref boxes.  Without it those set as
% overfull \hbox and print into the margin; with it they set slightly loose,
% which is the lesser evil.  Does not affect paragraphs that already break.
\emergencystretch=3em

%% ========================================
%% NUMBER SETS
%% ========================================
\newcommand{\R}{\mathbb{R}}  % Real numbers
\newcommand{\C}{\mathbb{C}}  % Complex numbers
\newcommand{\N}{\mathbb{N}}
\newcommand{\Z}{\mathbb{Z}}
% \providecommand, not \newcommand: some papers define these in their own
% preamble after % Shared preamble for the GLVbook papers.  This is the single canonical copy;
% a driver one level below GLVbook reaches it as
%
%     \input{../basicFiles/preamble}
%
% What stays in the driver, because it differs from paper to paper:
% \documentclass, font choice, \title, \author, \date, and any macro genuinely
% local to one paper.  Everything else belongs here.  The file is
% self-contained: it does not read any other file.

\usepackage[utf8]{inputenc}
\usepackage[novbox]{pdfsync}   % required: SyncTeX alone does not work here

% Page geometry, set once here for every paper that loads this file.  A driver
% that needs different margins says \geometry{...}; a second
% \usepackage[...]{geometry} raises an option clash and must not be used.
\makeatletter
\@ifundefined{AAdefNnogeometry}{\usepackage[a4paper,margin=2cm]{geometry}}{}
\makeatother

% mathematics
\usepackage{amsmath,amssymb,amsthm,mathtools,bm,mathrsfs,bbm}
\usepackage{centernot}

% floats, tables, lists
\usepackage{graphicx,float}
\makeatletter
\@ifpackageloaded{subfig}{}{\usepackage{subcaption}}
\makeatother
\usepackage{booktabs,longtable,threeparttable,array,multirow}
\usepackage{enumitem}

% colour, code, algorithms
\usepackage{xcolor,listings,tcolorbox}
\makeatletter
\@ifundefined{c@algorithm}{\usepackage{algorithm}}{}
\makeatother
\usepackage{algorithmicx,algpseudocode}

% diagrams
\usepackage{tikz,tikz-cd,pgfplots}
\usetikzlibrary{graphs,decorations.pathmorphing,decorations.markings,
arrows.meta,automata,positioning}
\tikzcdset{every label/.append style = {font = \small}}
\pgfplotsset{compat=1.18}

% chemistry
\usepackage[version=4]{mhchem}

\usepackage{microtype}
\usepackage{lineno}
\usepackage{nomencl}
\makenomenclature

% hyperref first, then bookmark, then glossaries: glossaries must follow
% hyperref or the entry links are built against the wrong anchors.
\makeatletter
\@ifundefined{hyperrefpackage}{%
  \@ifundefined{AAdefNnohyperref}{\usepackage{hyperref}}{}%
}{}
\@ifundefined{hyperrefpackage}{\usepackage{bookmark}}{}
\makeatother
\usepackage[toc,sort=use]{glossaries}
\makeglossaries

% Book/standalone switch.  Material that the book already covers in its
% prelude, but that a standalone paper still has to state, is wrapped as
%     \ifbookbuild\else ... \fi
% Drivers default to standalone; bookM.tex says \bookbuildtrue.
\newif\ifbookbuild
\bookbuildfalse

% Part switches for the gradostat paper.  Both default to true; a driver that
% wants the short submission says \madecfalse and/or \slowfastfalse, and every
% cross-reference into the dropped part disappears with it.
%   \ifmadec     Part III, the Lagasquie--Madec model and its break-even geometry
%   \ifslowfast  Part IV, everything carrying the small parameter \eps
\newif\ifmadec     \madectrue
\newif\ifslowfast  \slowfasttrue

% Part switch for the Barrat/simplicial-contagion project.  Defaults to
% true (the hypernetwork-synchrony part is included); the journal-paper
% driver says \hnetpartfalse to drop it and every cross-reference into it.
\newif\ifhnetpart  \hnetparttrue

%% ========================================
%% PACKAGES
%% ========================================
\RequirePackage{etoolbox}
\usepackage{makecell}
\usepackage[novbox]{pdfsync}% obsolete; SyncTeX is built into modern engines
\usepackage{amssymb, amsfonts,mathtools,bbm,eurosym,graphicx,latexsym, amsmath,amsthm,times,nccmath,url}
\usepackage{listings,verbatim,cases,tabularx}
% hyperref is deliberately NOT in the line above.  Journal classes often
% load it themselves, with a long option list, from a begindocument hook
% (mdpi.cls does: \DeclareHookRule{begindocument}{hyperref}{before}{mdpi}).
% Loading it here first, without options, makes that later load raise
% "Option clash for package hyperref".  So load it only if nobody else
% has, or will:
%   - \@ifpackageloaded  catches a class that loaded it already;
%   - \hyperrefpackage    is mdpi.cls's own deferred loader;
%   - \AAdefNnohyperref   lets any other driver opt out by defining it
%                         before % Shared preamble for the GLVbook papers.  This is the single canonical copy;
% a driver one level below GLVbook reaches it as
%
%     \input{../basicFiles/preamble}
%
% What stays in the driver, because it differs from paper to paper:
% \documentclass, font choice, \title, \author, \date, and any macro genuinely
% local to one paper.  Everything else belongs here.  The file is
% self-contained: it does not read any other file.

\usepackage[utf8]{inputenc}
\usepackage[novbox]{pdfsync}   % required: SyncTeX alone does not work here

% Page geometry, set once here for every paper that loads this file.  A driver
% that needs different margins says \geometry{...}; a second
% \usepackage[...]{geometry} raises an option clash and must not be used.
\makeatletter
\@ifundefined{AAdefNnogeometry}{\usepackage[a4paper,margin=2cm]{geometry}}{}
\makeatother

% mathematics
\usepackage{amsmath,amssymb,amsthm,mathtools,bm,mathrsfs,bbm}
\usepackage{centernot}

% floats, tables, lists
\usepackage{graphicx,float}
\makeatletter
\@ifpackageloaded{subfig}{}{\usepackage{subcaption}}
\makeatother
\usepackage{booktabs,longtable,threeparttable,array,multirow}
\usepackage{enumitem}

% colour, code, algorithms
\usepackage{xcolor,listings,tcolorbox}
\makeatletter
\@ifundefined{c@algorithm}{\usepackage{algorithm}}{}
\makeatother
\usepackage{algorithmicx,algpseudocode}

% diagrams
\usepackage{tikz,tikz-cd,pgfplots}
\usetikzlibrary{graphs,decorations.pathmorphing,decorations.markings,
arrows.meta,automata,positioning}
\tikzcdset{every label/.append style = {font = \small}}
\pgfplotsset{compat=1.18}

% chemistry
\usepackage[version=4]{mhchem}

\usepackage{microtype}
\usepackage{lineno}
\usepackage{nomencl}
\makenomenclature

% hyperref first, then bookmark, then glossaries: glossaries must follow
% hyperref or the entry links are built against the wrong anchors.
\makeatletter
\@ifundefined{hyperrefpackage}{%
  \@ifundefined{AAdefNnohyperref}{\usepackage{hyperref}}{}%
}{}
\@ifundefined{hyperrefpackage}{\usepackage{bookmark}}{}
\makeatother
\usepackage[toc,sort=use]{glossaries}
\makeglossaries

% Book/standalone switch.  Material that the book already covers in its
% prelude, but that a standalone paper still has to state, is wrapped as
%     \ifbookbuild\else ... \fi
% Drivers default to standalone; bookM.tex says \bookbuildtrue.
\newif\ifbookbuild
\bookbuildfalse

% Part switches for the gradostat paper.  Both default to true; a driver that
% wants the short submission says \madecfalse and/or \slowfastfalse, and every
% cross-reference into the dropped part disappears with it.
%   \ifmadec     Part III, the Lagasquie--Madec model and its break-even geometry
%   \ifslowfast  Part IV, everything carrying the small parameter \eps
\newif\ifmadec     \madectrue
\newif\ifslowfast  \slowfasttrue

% Part switch for the Barrat/simplicial-contagion project.  Defaults to
% true (the hypernetwork-synchrony part is included); the journal-paper
% driver says \hnetpartfalse to drop it and every cross-reference into it.
\newif\ifhnetpart  \hnetparttrue

%% ========================================
%% PACKAGES
%% ========================================
\RequirePackage{etoolbox}
\usepackage{makecell}
\usepackage[novbox]{pdfsync}% obsolete; SyncTeX is built into modern engines
\usepackage{amssymb, amsfonts,mathtools,bbm,eurosym,graphicx,latexsym, amsmath,amsthm,times,nccmath,url}
\usepackage{listings,verbatim,cases,tabularx}
% hyperref is deliberately NOT in the line above.  Journal classes often
% load it themselves, with a long option list, from a begindocument hook
% (mdpi.cls does: \DeclareHookRule{begindocument}{hyperref}{before}{mdpi}).
% Loading it here first, without options, makes that later load raise
% "Option clash for package hyperref".  So load it only if nobody else
% has, or will:
%   - \@ifpackageloaded  catches a class that loaded it already;
%   - \hyperrefpackage    is mdpi.cls's own deferred loader;
%   - \AAdefNnohyperref   lets any other driver opt out by defining it
%                         before % Shared preamble for the GLVbook papers.  This is the single canonical copy;
% a driver one level below GLVbook reaches it as
%
%     \input{../basicFiles/preamble}
%
% What stays in the driver, because it differs from paper to paper:
% \documentclass, font choice, \title, \author, \date, and any macro genuinely
% local to one paper.  Everything else belongs here.  The file is
% self-contained: it does not read any other file.

\usepackage[utf8]{inputenc}
\usepackage[novbox]{pdfsync}   % required: SyncTeX alone does not work here

% Page geometry, set once here for every paper that loads this file.  A driver
% that needs different margins says \geometry{...}; a second
% \usepackage[...]{geometry} raises an option clash and must not be used.
\makeatletter
\@ifundefined{AAdefNnogeometry}{\usepackage[a4paper,margin=2cm]{geometry}}{}
\makeatother

% mathematics
\usepackage{amsmath,amssymb,amsthm,mathtools,bm,mathrsfs,bbm}
\usepackage{centernot}

% floats, tables, lists
\usepackage{graphicx,float}
\makeatletter
\@ifpackageloaded{subfig}{}{\usepackage{subcaption}}
\makeatother
\usepackage{booktabs,longtable,threeparttable,array,multirow}
\usepackage{enumitem}

% colour, code, algorithms
\usepackage{xcolor,listings,tcolorbox}
\makeatletter
\@ifundefined{c@algorithm}{\usepackage{algorithm}}{}
\makeatother
\usepackage{algorithmicx,algpseudocode}

% diagrams
\usepackage{tikz,tikz-cd,pgfplots}
\usetikzlibrary{graphs,decorations.pathmorphing,decorations.markings,
arrows.meta,automata,positioning}
\tikzcdset{every label/.append style = {font = \small}}
\pgfplotsset{compat=1.18}

% chemistry
\usepackage[version=4]{mhchem}

\usepackage{microtype}
\usepackage{lineno}
\usepackage{nomencl}
\makenomenclature

% hyperref first, then bookmark, then glossaries: glossaries must follow
% hyperref or the entry links are built against the wrong anchors.
\makeatletter
\@ifundefined{hyperrefpackage}{%
  \@ifundefined{AAdefNnohyperref}{\usepackage{hyperref}}{}%
}{}
\@ifundefined{hyperrefpackage}{\usepackage{bookmark}}{}
\makeatother
\usepackage[toc,sort=use]{glossaries}
\makeglossaries

% Book/standalone switch.  Material that the book already covers in its
% prelude, but that a standalone paper still has to state, is wrapped as
%     \ifbookbuild\else ... \fi
% Drivers default to standalone; bookM.tex says \bookbuildtrue.
\newif\ifbookbuild
\bookbuildfalse

% Part switches for the gradostat paper.  Both default to true; a driver that
% wants the short submission says \madecfalse and/or \slowfastfalse, and every
% cross-reference into the dropped part disappears with it.
%   \ifmadec     Part III, the Lagasquie--Madec model and its break-even geometry
%   \ifslowfast  Part IV, everything carrying the small parameter \eps
\newif\ifmadec     \madectrue
\newif\ifslowfast  \slowfasttrue

% Part switch for the Barrat/simplicial-contagion project.  Defaults to
% true (the hypernetwork-synchrony part is included); the journal-paper
% driver says \hnetpartfalse to drop it and every cross-reference into it.
\newif\ifhnetpart  \hnetparttrue

%% ========================================
%% PACKAGES
%% ========================================
\RequirePackage{etoolbox}
\usepackage{makecell}
\usepackage[novbox]{pdfsync}% obsolete; SyncTeX is built into modern engines
\usepackage{amssymb, amsfonts,mathtools,bbm,eurosym,graphicx,latexsym, amsmath,amsthm,times,nccmath,url}
\usepackage{listings,verbatim,cases,tabularx}
% hyperref is deliberately NOT in the line above.  Journal classes often
% load it themselves, with a long option list, from a begindocument hook
% (mdpi.cls does: \DeclareHookRule{begindocument}{hyperref}{before}{mdpi}).
% Loading it here first, without options, makes that later load raise
% "Option clash for package hyperref".  So load it only if nobody else
% has, or will:
%   - \@ifpackageloaded  catches a class that loaded it already;
%   - \hyperrefpackage    is mdpi.cls's own deferred loader;
%   - \AAdefNnohyperref   lets any other driver opt out by defining it
%                         before \input{preamble}.
\makeatletter
\@ifpackageloaded{hyperref}{}{%
  \@ifundefined{hyperrefpackage}{%
    \@ifundefined{AAdefNnohyperref}{\usepackage{hyperref}}{}%
  }{}%
}
\makeatother
% linktoc=all makes a table-of-contents page number a link as well as the
% entry text; hyperref's default, linktoc=section, links only the text.
% Deferred to begin-document and guarded: by the block above hyperref may
% be loaded later (mdpi's own hook) or not at all (\AAdefNnohyperref), and
% \hypersetup must exist before it is called.  A driver that wants the
% default back can say \hypersetup{linktoc=section} after \input{preamble}.
\makeatletter
\AtBeginDocument{\@ifundefined{hypersetup}{}{\hypersetup{linktoc=all}}}
\makeatother
\usepackage{graphicx,float,fancybox}
\usepackage{epsfig,epstopdf,tikz}
\usepackage[sans]{dsfont}
\usetikzlibrary{positioning,arrows,arrows.meta,calc}
% Configure listings for Mathematica code
\lstdefinestyle{mathematica}{
 language=Mathematica,
 basicstyle=\small\ttfamily,
 backgroundcolor=\color{gray!10},
 frame=single,
 breaklines=true,
 commentstyle=\color{green!60!black},
 keywordstyle=\color{blue},
 stringstyle=\color{red},
 showstringspaces=false,
 tabsize=2
}
\usepackage{tcolorbox}
% Colored boxes for key results
\newtcolorbox{keyresult}[1][]{
 colback=blue!5!white,
 colframe=blue!75!black,
 title=#1,
 fonttitle=\bfseries
}
\newtcolorbox{examplebox}[1][]{
 colback=green!5!white,
 colframe=green!75!black,
 title=#1,
 fonttitle=\bfseries
}


\makeatletter
\@ifundefined{theorem}   {\newtheorem{theorem}   {Theorem}}   {}
\@ifundefined{definition}{\newtheorem{definition}{Definition}} {}
\@ifundefined{open}      {\newtheorem{open}      {Problem}}   {}
\@ifundefined{example}   {\newtheorem{example}   {Example}}   {}
\makeatother

% COUNTER DEFINITIONS
\newcounter{tabcounter}
\newcounter{defcounter}
\newcounter{opcounter}
\newcounter{thmcounter}
\newcounter{excounter}
\newcounter{ascounter}
\newcounter{cjcounter}
\newcounter{lemcounter}
\newcounter{remcounter}
\newcounter{corcounter}
\newcounter{propcounter}
\newcounter{figcounter}

% UNIFIED GLOSSARY STYLE
\newglossarystyle{unifiedstyle}{%
  \setglossarystyle{list}%
  \renewcommand*{\glossarypreamble}{\gdef\glslastsec{}}%
  \renewcommand{\glossentry}[2]{%
    \edef\glscursec{\glsentryuseriii{##1}}%
    \ifx\glscursec\glslastsec\else
      \item[]\textbf{\S\,\glscursec\enspace\glsentryuseriv{##1}}%
      \global\let\glslastsec\glscursec
    \fi
    \item[\glsentryuserii{##1} \glsentryuseri{##1}:] ``\glossentryname{##1}'' \glsentrydesc{##1}%
  }%
}

\makeatletter
% First-word extractor (used by \beT, \beD and \beXa for the default key)
% \@firstwordof hello world \relax  ->  hello
\def\@firstwordof#1 #2\relax{#1}
% First three letters, used for the heading mnemonic of a theorem given no
% explicit key.  Padding with \@empty keeps it safe for first words shorter
% than three letters: \@empty vanishes under \edef, whereas \relax would
% survive into the label and is illegal inside \csname (it broke every
% two-letter title, e.g. "ME ...", with "Missing \endcsname inserted").
\def\@takethree#1#2#3#4\@nil{#1#2#3}

% EXAMPLE COMMAND -- same signature and behaviour as \beT and \beD.
% Usage: \beXa*[long title][short title][key]
%   star  = suppresses the printed "[key]" tag after the title (as for \beT*
%   and \beD*); the label and glossary entry are unaffected.
%   key = optional label key override; the label becomes e:KEY, and without a
%   key it is e: plus the first three letters of the title's first word.  A
%   mnemonic label lets an example be cited as "Example~\ref{e:KEY}" instead
%   of by section number, which is what the citing prose needs.
\NewDocumentCommand{\beXa}{s O{} O{} o}{%
  \def\@temparg{#2}%
  \ifx\@temparg\@empty
    \begin{example}%
  \else
    \let\@savedSecTitle\@currentsectionname%
    \refstepcounter{excounter}%
    \IfValueTF{#4}{%
      \edef\@exlblkey{#4}%
    }{%
      \edef\@exfirstword{\@firstwordof#2 \relax}%
      \edef\@exlblkey{%
        \expandafter\@takethree\@exfirstword\@empty\@empty\@empty\@nil}%
    }%
    \def\@tempshort{#3}%
    \ifx\@tempshort\@empty
      \expandafter\xdef\csname @xshort@\@exlblkey\endcsname{\@exlblkey}%
    \else
      \expandafter\def\csname @xshort@\@exlblkey\endcsname{#3}%
    \fi
    \IfBooleanTF{#1}{\let\@extag\@empty}{%
      \edef\@extag{\noexpand\ \noexpand\textup{[\noexpand\texttt{e:\@exlblkey}]}}}%
    \begin{example}[#2\ifshowkeys\@extag\fi]%
    \label{e:\@exlblkey}%
    % record the long title so \nref{e:KEY} prints it, as for theorems
    \protected@write\@auxout{}{%
      \string\@thmtitle{e:\@exlblkey}{\unexpanded{#2}}}%
    \expanded{%
      \noexpand\newglossaryentry{e\theexcounter}{%
        name={\unexpanded{#2}%
          \IfBooleanF{#1}{ [\noexpand\texttt{e:\@exlblkey}]}},%
        sort={\theexample},%
        description={p.~\noexpand\pageref{e:\@exlblkey}},%
        user1={\theexample},%
        user2={Example},%
        user3={\thesection},%
        user4={\expandafter\unexpanded\expandafter{\@savedSecTitle}}%
      }%
    }%
    \protected@write\@auxout{}{\string\glsliveentry{e\theexcounter}}%
  \fi
}
\newcommand{\eeXa}{\end{example}}

% DEFINITION COMMAND -- same signature and behaviour as \beT below.
% Usage: \beD[long title][short title][key]
%   long title  = definition heading (plain text only)
%   short title = shown in \dref{word}; defaults to the key
%   key         = optional label key override; the label becomes d:KEY
% Without a key the label is d: plus the first three letters of the title's
% first word, so it is a mnemonic rather than the definition's number: a
% numeric label would silently change whenever a definition is inserted.
\NewDocumentCommand{\beD}{s O{} O{} o}{%
  \def\@temparg{#2}%
  \ifx\@temparg\@empty
    \begin{definition}%
  \else
    \let\@savedSecTitle\@currentsectionname%
    \refstepcounter{defcounter}%
    \IfValueTF{#4}{%
      \edef\@deflblkey{#4}%
    }{%
      \edef\@deffirstword{\@firstwordof#2 \relax}%
      \edef\@deflblkey{%
        \expandafter\@takethree\@deffirstword\@empty\@empty\@empty\@nil}%
    }%
    \def\@tempshort{#3}%
    \ifx\@tempshort\@empty
      \expandafter\xdef\csname @dshort@\@deflblkey\endcsname{\@deflblkey}%
    \else
      \expandafter\def\csname @dshort@\@deflblkey\endcsname{#3}%
    \fi
    \IfBooleanTF{#1}{\let\@deftag\@empty}{%
      \edef\@deftag{\noexpand\ \noexpand\textup{[\noexpand\texttt{d:\@deflblkey}]}}}%
    \begin{definition}[#2\ifshowkeys\@deftag\fi]%
    \label{d:\@deflblkey}%
    % record the long title so \nref{d:KEY} prints it, as for theorems
    \protected@write\@auxout{}{%
      \string\@thmtitle{d:\@deflblkey}{\unexpanded{#2}}}%
    \expanded{%
      \noexpand\newglossaryentry{d\thedefcounter}{%
        name={\unexpanded{#2}%
          \IfBooleanF{#1}{ [\noexpand\texttt{d:\@deflblkey}]}},%
        sort={\thedefinition},%
        description={p.~\noexpand\pageref{d:\@deflblkey}},%
        user1={\thedefinition},%
        user2={Definition},%
        user3={\thesection},%
        user4={\expandafter\unexpanded\expandafter{\@savedSecTitle}}%
      }%
    }%
    \protected@write\@auxout{}{\string\glsliveentry{d\thedefcounter}}%
  \fi
}
\newcommand{\eeD}{\end{definition}}

% \showkeysfalse in a driver hides the [t:KEY] mnemonics in result headings
% (labels and glossary entries are unchanged); used for submission builds.
\newif\ifshowkeys \showkeystrue
% \commonredtrue in a driver that inputs basicFiles/commonRed: tool.tex then
% cites its reduction results instead of restating them.
\newif\ifcommonred \commonredfalse
% TheorEM COMMAND
% Usage: \beT[long title][short title][key]
%   long title  = theorem heading (plain text only)
%   short title = shown in \tref{word}; may contain math; defaults to first word of long title
%   key         = optional label key override; label becomes t:KEY (avoids first-word collisions)
% Label is t:WORD where WORD = first word of long title (or KEY if given); use \tref{WORD} to reference.
\NewDocumentCommand{\beT}{s O{} O{} o}{%
  \def\@temparg{#2}%
  \ifx\@temparg\@empty
    \begin{theorem}%
  \else
    \let\@savedSecTitle\@currentsectionname%
    \refstepcounter{thmcounter}%
    % \@thmtag is the mnemonic shown in the heading; every theorem gets one.
    % With an explicit key the key itself is shown.  Without one the key
    % defaults to the title's first word, and showing that whole word would
    % merely repeat the title ("Dissipativity [Dissipativity]"), so only its
    % first three letters are shown: "Dissipativity [Dis]".  The label is
    % unaffected -- it still uses the full \@thmlblkey.
    \IfValueTF{#4}{%
      \edef\@thmlblkey{#4}%
    }{%
      \edef\@thmfirstword{\@firstwordof#2 \relax}%
      \edef\@thmlblkey{%
        \expandafter\@takethree\@thmfirstword\@empty\@empty\@empty\@nil}%
    }%
    \IfBooleanTF{#1}{\let\@thmtag\@empty}{%
      \edef\@thmtag{\noexpand\ \noexpand\textup{[\noexpand\texttt{t:\@thmlblkey}]}}}%
    \def\@tempshort{#3}%
    \ifx\@tempshort\@empty
      \expandafter\xdef\csname @tshort@\@thmlblkey\endcsname{\@thmlblkey}%
    \else
      \expandafter\def\csname @tshort@\@thmlblkey\endcsname{#3}%
    \fi
    \begin{theorem}[#2\ifshowkeys\@thmtag\fi]%
    \label{t:\@thmlblkey}%
    % record the long title for \nref; via the .aux so forward refs resolve
    \protected@write\@auxout{}{%
      \string\@thmtitle{t:\@thmlblkey}{\unexpanded{#2}}}%
    \expanded{%
      \noexpand\newglossaryentry{t\thethmcounter}{%
        name={\unexpanded{#2}%
          \IfBooleanF{#1}{ [\noexpand\texttt{t:\@thmlblkey}]}},%
        % sort/user1 must be the number the document actually prints, i.e. the
        % theorem environment's counter, not thmcounter.  thmcounter advances
        % only on \beT with a title, so a bare \beT (or \beT with an
        % empty title) makes the two drift and the glossary cites a wrong
        % number.  thmcounter is still fine as the entry key: it is unique.
        sort={\thetheorem},%
        description={p.~\noexpand\pageref{t:\@thmlblkey}},%
        user1={\thetheorem},%
        user2={Theorem},%
        user3={\thesection},%
        user4={\expandafter\unexpanded\expandafter{\@savedSecTitle}}%
      }%
    }%
    \protected@write\@auxout{}{\string\glsliveentry{t\thethmcounter}}%
  \fi
}
\newcommand{\eeT}{\end{theorem}}

% CROSS-PART RELATION PARAGRAPH
% Usage: \xpart{text}  -- place immediately BEFORE \beT, never inside it
%   (anything inside \beT...\eeT would be swallowed into the glossary entry).
% The paragraph is typeset only when Part I is actually present, i.e. in the
% combined book build; in the standalone builds (GranadaDSABNS/A.tex,
% AAHRah-Submission/AAHRah.tex) the label part:glv is undefined and the whole
% paragraph disappears rather than leaving ?? cross-references. Same
% graceful-degradation principle as \GASref in GASt.tex. Resolves on the
% second LaTeX pass, since r@part:glv comes from the .aux file.
% (Already inside the \makeatletter region opened above, so no \makeatletter
% pair here -- adding one would close that region early.)
\newcommand{\xpart}[1]{\@ifundefined{r@part:glv}{}{\paragraph{Relation with Parts~I--II.}#1}}

% \bookref{label}{precise text}{fallback text}: prints the precise cross-reference
% when the label resolves (combined book build), and the fallback text otherwise
% (standalone paper builds, where book-only appendix labels are absent). Same
% graceful-degradation role as \GASref in GASt.tex, but usable in any file since it
% lives in this preamble (loaded by every driver) -- unlike \GASref, which is defined in a
% content file and is therefore undefined in standalone builds that omit that file.
\newcommand{\bookref}[3]{\@ifundefined{r@#1}{#3}{#2}}

% \PVpart and \RBpart: the two parts that a partial build may omit.
% They print "Part~N" in any build that contains the part, and a naming phrase
% in a build that does not, so an excerpt shows no ?? .  Each already supplies
% the word "Part", so write "\PVpart develops ...", not "Part~\PVpart".
\newcommand{\PVpart}{\bookref{part:aahrah}%
  {Part~\ref{part:aahrah}}{the Perron--Volterra part of the full book}}
\newcommand{\RBpart}{\bookref{part:recbrowne}%
  {Part~\ref{part:recbrowne}}{the immuno-virology part of the full book}}

% OPEN PROBLEM COMMANDS
% ASSUMPTION COMMANDS (glossaried)
\newcommand{\beAs}[1][]{%
  \def\@temparg{#1}%
  \ifx\@temparg\@empty
    \begin{assumption}%
  \else
    \let\@savedSecTitle\@currentsectionname%
    \refstepcounter{ascounter}%
    \begin{assumption}[#1]%
    \label{as:\theascounter}%
    \expanded{%
      \noexpand\newglossaryentry{as\theascounter}{%
        name={\unexpanded{#1}},%
        sort={\theascounter},%
        description={p.~\noexpand\pageref{as:\theascounter}},%
        user1={\theascounter},%
        user2={Assumption},%
        user3={\thesection},%
        user4={\expandafter\unexpanded\expandafter{\@savedSecTitle}}%
      }%
    }%
    \protected@write\@auxout{}{\string\glsliveentry{as\theascounter}}%
  \fi
}
\newcommand{\eeAs}{\end{assumption}}
% CONJECTURE COMMANDS (glossaried)
\newcommand{\beCj}[1][]{%
  \def\@temparg{#1}%
  \ifx\@temparg\@empty
    \begin{conjecture}%
  \else
    \let\@savedSecTitle\@currentsectionname%
    \refstepcounter{cjcounter}%
    \begin{conjecture}[#1]%
    \label{cj:\thecjcounter}%
    \expanded{%
      \noexpand\newglossaryentry{cj\thecjcounter}{%
        name={\unexpanded{#1}},%
        sort={\thecjcounter},%
        description={p.~\noexpand\pageref{cj:\thecjcounter}},%
        user1={\thecjcounter},%
        user2={Conjecture},%
        user3={\thesection},%
        user4={\expandafter\unexpanded\expandafter{\@savedSecTitle}}%
      }%
    }%
    \protected@write\@auxout{}{\string\glsliveentry{cj\thecjcounter}}%
  \fi
}
\newcommand{\eeCj}{\end{conjecture}}
% LEMMA COMMANDS
% ---------------------------------------------------------------------
% UNIFORM RESULT WRAPPERS -- \beL, \beC, \beP, \beR share \beT's signature
%
%   \beX[long title][short title][key]  ...  \eeX
%
%   long title  = heading (math allowed)
%   short title = shown in the cross-reference; defaults to the key
%   key         = label key.  With a key the label is X:KEY and the heading
%                 carries the mnemonic [key], exactly as for \beT.  Without
%                 one the label falls back to the running counter, which is
%                 what every pre-existing call site in the book relies on.
%
% A bare \beX -- arguments commented out -- opens the plain environment and
% creates NO label and NO glossary entry.  That is the supported way to keep
% an individual result out of the catalogue.
% ---------------------------------------------------------------------
% #1 wrapper name  #2 environment  #3 label prefix  #4 counter  #5 type name
\newcommand{\@mkresult}[5]{%
  % A leading * suppresses the [key] mnemonic in both the heading and the
  % glossary name, while still making the label X:KEY and the entry:
  %     \beX*[long][short][key]
  \expandafter\NewDocumentCommand\csname be#1\endcsname{s O{} O{} o}{%
    \def\@temparg{##2}%
    \ifx\@temparg\@empty
      \begin{#2}%
    \else
      \let\@savedSecTitle\@currentsectionname%
      \refstepcounter{#4}%
      % With a key: mnemonic label X:KEY and the [key] tag in the heading.
      % Without one: fall back to the running counter, exactly as these
      % wrappers behaved before, so every legacy one-argument call site in
      % the book keeps its label.  The counter is also csname-safe, which a
      % title's first word is not -- titles beginning with math would
      % otherwise build \csname \Omega...\endcsname and fail.
      \IfValueTF{##4}{%
        \edef\@relbl{##4}%
        \IfBooleanTF{##1}{\let\@retag\@empty}{%
          \edef\@retag{\noexpand\ \noexpand\textup{[\noexpand\texttt{#3:##4}]}}}%
      }{%
        \edef\@relbl{\arabic{#4}}%
        \let\@retag\@empty
      }%
      \def\@tempshort{##3}%
      \ifx\@tempshort\@empty
        \expandafter\xdef\csname @sshort@#3\@relbl\endcsname{\@relbl}%
      \else
        \expandafter\def\csname @sshort@#3\@relbl\endcsname{##3}%
      \fi
      \begin{#2}[##2\ifshowkeys\@retag\fi]%
      \edef\@tmplbl{#3:\@relbl}\expandafter\label\expandafter{\@tmplbl}%
      \expanded{%
        \noexpand\newglossaryentry{#3\csname the#4\endcsname}{%
          name={\unexpanded{##2}%
            \IfBooleanF{##1}{\IfValueT{##4}{ [\noexpand\texttt{#3:##4}]}}},%
          sort={\csname the#2\endcsname},%
          description={p.~\noexpand\pageref{#3:\@relbl}},%
          user1={\csname the#2\endcsname},%
          user2={#5},%
          user3={\thesection},%
          user4={\expandafter\unexpanded\expandafter{\@savedSecTitle}}%
        }%
      }%
      \protected@write\@auxout{}{\string\glsliveentry{#3\csname the#4\endcsname}}%
      \protected@write\@auxout{}{\string\@thmtitle{\@tmplbl}{\unexpanded{##2}}}%
    \fi
  }%
  \expandafter\newcommand\csname ee#1\endcsname{\end{#2}}%
}
\@mkresult{L}{lemma}{l}{lemcounter}{Lemma}
% REMARK / COROLLARY / PROPOSITION -- same signature as \beL above
\@mkresult{R}{remark}{r}{remcounter}{Remark}
\@mkresult{C}{corollary}{c}{corcounter}{Corollary}
\@mkresult{P}{proposition}{p}{propcounter}{Proposition}
% OPEN-PROBLEM COMMAND -- generated with the same signature as \beT:
%   \beO[long title][short title][key]   ...   \eeO
% label o:KEY, or o:<counter> when no key is given; a bare \beO makes
% neither a label nor a glossary entry.
\@mkresult{O}{open}{o}{opcounter}{Problem}


% Section-level title tracker: updated only by \section, not \subsection.
% \@currentlabelname is overwritten by every \refstepcounter (including
% \subsection), so we save it separately here.
\newcommand{\@currentsectionname}{}
\AtBeginDocument{%
  \let\@epidcrn@orig@refstepcounter\refstepcounter
  \renewcommand{\refstepcounter}[1]{%
    \@epidcrn@orig@refstepcounter{#1}%
    \ifstrequal{#1}{section}{\global\let\@currentsectionname\@currentlabelname}{}%
  }%
}

% TABLE COMMANDS (glossaried) -- consistent naming with \beT, \beD, etc.
\newcommand{\beTa}[1][]{%
  \def\@temparg{#1}%
  \ifx\@temparg\@empty
    \begin{table}[htbp]%
  \else
    \let\@savedSecTitle\@currentsectionname%
    \refstepcounter{tabcounter}%
    \begin{table}[htbp]%
    \label{A:\thetabcounter}%
    \expanded{%
      \noexpand\newglossaryentry{A\thetabcounter}{%
        name={\unexpanded{#1}},%
        sort={\thetabcounter},%
        description={p.~\noexpand\pageref{A:\thetabcounter}},%
        user1={\thetabcounter},%
        user2={Table},%
        user3={\thesection},%
        user4={\expandafter\unexpanded\expandafter{\@savedSecTitle}}%
      }%
    }%
    \protected@write\@auxout{}{\string\glsliveentry{A\thetabcounter}}%
  \fi
}
\newcommand{\eeTa}{\end{table}}

% FIGURE COMMANDS (glossaried) -- consistent with \beT, \beR, \beTa, etc.
% Usage:
%   \beF[Glossary title]...\caption{Visible caption}\eeF  -> glossaried figure
%   \beF...\caption{Visible caption}\eeF                  -> plain [H] figure, no glossary
\newcommand{\beF}[1][]{%
  \def\@temparg{#1}%
  \ifx\@temparg\@empty
    \begin{figure}[H]%
  \else
    \begin{figure}[H]%
    \let\@savedSecTitle\@currentsectionname%
    \refstepcounter{figcounter}%
    \label{fig:\thefigcounter}%
    \expanded{%
      \noexpand\newglossaryentry{Fig\thefigcounter}{%
        name={\unexpanded{#1}},%
        sort={\thefigcounter},%
        description={p.~\noexpand\pageref{fig:\thefigcounter}},%
        user1={\thefigcounter},%
        user2={Figure},%
        user3={\thesection},%
        user4={\expandafter\unexpanded\expandafter{\@savedSecTitle}}%
      }%
    }%
    \protected@write\@auxout{}{\string\glsliveentry{Fig\thefigcounter}}%
  \fi
}
\newcommand{\eeF}{\end{figure}}
\makeatother
% LIVE-ENTRY REGISTER
% \glsaddall prints every entry the glossaries package currently knows, and
% that set includes entries restored from \jobname.glsdefs.  The package
% rewrites .glsdefs from all *currently defined* entries, so anything it
% restored is written straight back out: a deleted section kept reappearing in
% the catalogue for ever, and the .glsdefs had to be erased by hand.  The
% wrappers therefore record every entry they actually create in the .aux, and
% \printboth adds only those.  The .aux is rebuilt from scratch on each run, so
% this is self-healing: a removed result is gone after one further run.
\makeatletter
\newcommand{\@glslive}{}
\newcommand{\@glsaddlive}[1]{\ifglsentryexists{#1}{\glsadd{#1}}{}}
\newcommand{\glsliveentry}[1]{\g@addto@macro\@glslive{\@glsaddlive{#1}}}
\makeatother
% PRINT COMMAND
% No \clearpage here -- the catalogue continues on the current page.  This has
% been re-added by mistake several times; leave it out.
\makeatletter
\newcommand{\printboth}{%
  \@glslive   % add exactly the entries this document created last run
  \printglossary[style=unifiedstyle, title={0 \hspace{.22em} Definitions, Theorems, Examples and Open Problems}]%
}
\makeatother
% Named theorem reference: \tref{word} -> Thm~N(word-or-short)
% Shows stored short title if available, otherwise shows the key itself.
\newcommand{\tref}[1]{%
  Thm~\ref{t:#1}(\ifcsname @tshort@#1\endcsname\csname @tshort@#1\endcsname\else#1\fi)%
}

% Full-title theorem reference: \nref{t:KEY} prints the number and then the
% theorem's long title, so one writes "Theorem~\nref{t:blocks}" and gets
% "Theorem 3 ``Block factorization and siphons''".  Takes the complete label, like
% \ref, not the bare key like \tref.
%
% \beT writes each title to the .aux (see \@thmtitle below), so \nref resolves
% forward references from the second LaTeX pass, exactly as \ref does.  If the
% title is not yet known -- first pass, or a label that is not a \beT theorem
% -- it degrades to the number rather than producing ??.
% These two use @ in their names, and this point in the file is outside the
% \makeatletter region opened earlier (it closes at \makeatother above), so
% the region has to be reopened here.
\makeatletter
\newcommand{\@thmtitle}[2]{\global\@namedef{@tlong@#1}{#2}}
\newcommand{\nref}[1]{%
  \ref{#1}%
  \ifcsname @tlong@#1\endcsname
    \ ``\csname @tlong@#1\endcsname''%
  \fi
}
% \nm{t:KEY}: the stored long title of any \beT/\beL/\beP/\beR/\beC/\beD/\beE result,
% verbatim and without quotes, so item names follow the macro's first argument.
\newcommand{\nm}[1]{%
  \ifcsname @tlong@#1\endcsname
    \csname @tlong@#1\endcsname%
  \else
    \ref{#1}%
  \fi
}
\makeatother




%% ========================================
%% Other THEOREM ENVIRONMENTS
%% ========================================
\makeatletter
\@ifundefined{lemma}      {\newtheorem{lemma}      {Lemma}}      {}
\@ifundefined{proposition}{\newtheorem{proposition}{Proposition}} {}
\@ifundefined{corollary}  {\newtheorem{corollary}  {Corollary}}  {}
\@ifundefined{conjecture} {\newtheorem{conjecture} {Conjecture}}  {}
\@ifundefined{remark}     {\newtheorem{remark}     {Remark}}     {}
\@ifundefined{question}   {\newtheorem{question}   {Question}}   {}
\@ifundefined{assumption} {\newtheorem{assumption} {Assumption}} {}
\@ifundefined{hypothesis} {\newtheorem{hypothesis} {Hypothesis}} {}
\@ifundefined{fact}       {\newtheorem{fact}       {Fact}}       {}
\makeatother

% \citer: cite the relay-graph threshold theorem of AAHrel. If a document has
% \input the chapter/file containing t:relayRF (so the label resolves), point
% straight at it; otherwise fall back to a bare citation of AAHrel. Shared here
% (not per-paper) so it resolves consistently across every book chapter.
% \citeR: same idea, for the first theorem of AAHRah (t:Pstar).
\makeatletter
\providecommand{\citer}{%
  \@ifundefined{r@t:relayRF}%
    {\cite{AAHrel}}%
    {Theorem~\ref{t:relayRF} of \cite{AAHrel}}%
}
\providecommand{\citeR}{%
  \@ifundefined{r@t:Pstar}%
    {\cite{AAHRah}}%
    {Theorem~\ref{t:Pstar} of \cite{AAHRah}}%
}
\makeatother


% Abbreviations
\def\beQ{\begin{question}}\def\eeQ{\end{question}}
\def\beA{\begin{assumption}}\def\eeA{\end{assumption}}


%% ========================================

%% ========================================
%% COLORS & CUSTOM COMMANDS
%% ========================================
\definecolor{funccolor}{RGB}{25,25,112}
\definecolor{desccolor}{RGB}{64,64,64}
\newcommand{\func}[1]{\texttt{\textcolor{funccolor}{#1}}}
\def\eq{=}
\def\mm{-}

%% ========================================
%% MATHCAL SHORTHANDS
%% ========================================
% Project-wide convention: \mX for a capital X is always \mathcal X.
% Defined for the whole alphabet so the convention never fails to expand.
\newcommand{\mA}{\mathcal A}\newcommand{\mB}{\mathcal B}
\newcommand{\mC}{\mathcal C}\newcommand{\mD}{\mathcal D}
\newcommand{\mE}{\mathcal E}\newcommand{\mF}{\mathcal F}
\newcommand{\mG}{\mathcal G}\newcommand{\mH}{\mathcal H}
\newcommand{\mI}{\mathcal I}\newcommand{\mJ}{\mathcal J}
\newcommand{\mK}{\mathcal K}\newcommand{\mL}{\mathcal L}
\newcommand{\mM}{\mathcal M}\newcommand{\mN}{\mathcal N}
\newcommand{\mO}{\mathcal O}\newcommand{\mP}{\mathcal P}
\newcommand{\mQ}{\mathcal Q}\newcommand{\mR}{\mathcal R}
\newcommand{\mS}{\mathcal S}\newcommand{\mT}{\mathcal T}
\newcommand{\mU}{\mathcal U}\newcommand{\mV}{\mathcal V}
\newcommand{\mW}{\mathcal W}\newcommand{\mX}{\mathcal X}
\newcommand{\mY}{\mathcal Y}\newcommand{\mZ}{\mathcal Z}

%% ========================================
%% BASIC MATH ENVIRONMENTS
%% ========================================
\def\bep{\begin{pmatrix}}\def\eep{\end{pmatrix}}
\def\bev{\begin{vmatrix}}\def\eev{\end{vmatrix}}
\def\bea{\begin{eqnarray*}}\def\eea{\end{eqnarray*}}
\def\bc{\begin{cases}}\def\ec{\end{cases}}
\def\BEN{\begin{enumerate}}\def\EEN{\end{enumerate}}
\def\BI{\begin{itemize}}\def\EI{\end{itemize}}

\newcommand{\be}[1]{\begin{equation}\label{#1}}
\newcommand{\ee}{\end{equation}}
\newcommand{\beq}{\begin{eqnarray}}
\def\eeq{\end{eqnarray}}
\def\eqr{\eqref}\def\fr{\frac}\def\lbl{\label}
\def\Lra{\Longrightarrow}\def\non{\nonumber}\def\cNGM{\cite{Diek,Van,Van08}}
\def\Eq{\Leftrightarrow}\def\b{\beta}\def\uph{\red{UP TO HERE}}
%% ========================================
%% PROBABILITY & STATISTICS
%% ========================================
\renewcommand{\P}{\ensuremath{\mathbb{P}}}  % Note: overrides paragraph symbol
\newcommand{\E}{\ensuremath{\mathbb{E}}}
\newcommand{\ind}[1]{\ensuremath{\mathbbm{1}_{\left\{#1\right\}}}}
\DeclareMathOperator{\Cov}{Cov}
\DeclareMathOperator{\TV}{TV}
\newcommand{\rank}{\operatorname{rank}}

%% ========================================
%% CRN OPERATORS (organizations/ERCs/buffering/siphons framework)
%% ========================================
\def\clos{\mathrm{clos}}
\def\Supply{\mathrm{Supply}}
\def\req{\mathrm{req}}
\def\React{\mathrm{React}}
\def\Prod{\mathrm{Prod}}
\def\Gen{\mathrm{Gen}}
\def\In{\mathrm{In}}
\def\Out{\mathrm{Out}}
\def\Act{\mathrm{Act}}
\def\Min{\mathrm{Min}}

%% ========================================
%% LINE BREAKING
%% ========================================
% Last-resort stretch, used only on paragraphs TeX cannot otherwise break:
% a long \texttt{} token, a comma-separated inline math list (commas are not
% breakpoints in math), or a run of \cite/\ref boxes.  Without it those set as
% overfull \hbox and print into the margin; with it they set slightly loose,
% which is the lesser evil.  Does not affect paragraphs that already break.
\emergencystretch=3em

%% ========================================
%% NUMBER SETS
%% ========================================
\newcommand{\R}{\mathbb{R}}  % Real numbers
\newcommand{\C}{\mathbb{C}}  % Complex numbers
\newcommand{\N}{\mathbb{N}}
\newcommand{\Z}{\mathbb{Z}}
% \providecommand, not \newcommand: some papers define these in their own
% preamble after \input{preamble}, and must not clash.
\providecommand{\Q}{\mathbb{Q}}
\providecommand{\res}{\operatorname{res}}   % resultant; used by epiRes.tex
\newcommand{\MN}{\mathbb{N}}
\newcommand{\MZ}{\mathbb{Z}}
\newcommand{\realpos}{\mathbb{R}_{> 0}}
\newcommand{\realnneg}{\mathbb{R}_{\geq 0}}
\newcommand{\naturalpos}{\mathbb{N}_{>0}}
\newcommand{\Rplus}{\mathbb{R}_+}
\newcommand{\Rp}{\mathbb{R}_{\ge 0}}
\def\Rnne{\mathbb{R}_{\ge 0}}
\def\Se{\bar{S}} \def\Ie{\bar{I}}
\newcommand{\Rpos}{\mathbb{R}_{>0}}
\newcommand{\dd}{\,\mathrm d}
\newcommand{\e}{\mathrm e}
%% ========================================
%% GREEK LETTERS (2+ character abbreviations) greek
%% ========================================
\def\al{\alpha}\def\ba{\bff \alpha}
\def\w{\omega}\def\Lgn{\ga+\mu+\de}\def\om{\omega}
\def\de{\delta}\def\De{\Delta}\def\bde{\bff \delta}
\def\ep{\epsilon}\def\eps{\varepsilon}
\def\f{\varphi}  % Kept as requested
\def\ga{\gamma}\def\Ga{\Gamma}\def\G{\Gamma}\def\bg{\bff \gamma}
\def\up{\upsilon}  % NOT the breve accent: no accent commands anywhere.
                   % Named \up, not \u, so the LaTeX breve stays reachable.
\def\la{\lambda}\def\La{\Lambda}
\def\si{\sigma}\def\ka{\kappa}
\def\th{\theta}\def\vt{\vartheta}
\def\ze{\zeta}\def\Gap{\Gamma^{+}} \def\Gam{\Gamma^{-}}


%% ========================================
%% SPECIAL SHORT COMMANDS (kept as requested)
%% ========================================
\def\I{\infty}
\def\T{\widetilde}
\def\H{\hat}
\newcommand{\relint}{\operatorname{relint}}

%% ========================================
%% BOLD/VECTOR NOTATION
%% ========================================
\newcommand{\bff}[1]{{\mbox{\boldmath$#1$}}}
\newcommand\X{\boldsymbol{X}}
\def\x{\boldsymbol{x}}\def\bx{\bff x}\def\xe{\bar{\x}}
\newcommand\y{\boldsymbol{y}}\def\by{\boldsymbol{y}}\def\ye{\bar{\y}}
\def\yb{\y}\def\xb{\x}
\def\xy{\x-\y}
\newcommand\br{\boldsymbol{r}}\def\re{\bar{\r}}
\newcommand\bd{\boldsymbol{d}}
\newcommand\bQ{\boldsymbol{Q}}
\newcommand{\Yb}{{\mathbf{Y}}}
\newcommand{\Ab}{\mathbf{A}}
\newcommand{\zb}{\mathbf{z}}\def\bz{\bff 0}
\newcommand{\Zb}{\mathbf{Z}}
\newcommand{\bb}{\mathbf{b}}
\def\vbe{\vec {\bff \beta}}
\def\vb{\vec {\bff b}}
\def\v1{\vec {\bff 1}}
\def\vn{\vec \de}
\def\vz{\vec 0}
\def\vi{\overset{\rightarrow}{\mathsf i}}
\def\vv{\overset{\rightarrow}{v}}
\def\vr{\; \vec {\mathsf r}}
\newcommand{\va}{\bff a}
\providecommand{\vect}[1]{\boldsymbol{#1}}   % bold vector (moved from Gradostat.tex)
\providecommand{\one}{\mathbf 1}             % all-ones vector (moved from AliSF)
\providecommand{\spr}{s}                     % spectral bound  (moved from AliSF)
\providecommand{\Ke}{K_{\rm eff}}            % effective carrying capacity
% Critical right/left eigenvectors in the Hopf projection method (AAHOSN/New.tex)
\providecommand{\bq}{\mathbf{q}}
\providecommand{\bp}{\mathbf{p}}


%% ========================================
%% MATHCAL & MATHSCR
%% ========================================
% Mathcal
\newcommand\mc[1]{\mathcal{#1}}
\def\mC{{\mathcal C}}\def\mD{\mathcal D}\def\mE{\mathcal{E}}
\def\mF{\mathcal F}\def\mG{{\mathcal G}}\def\mH{{\mathcal H}}
\def\mI{\mathcal I}\def\mL{\mathcal{L}}\def\mM{{\mathcal M}}
\def\mN{{\mathcal N}}\def\mP{\mathcal{P}}\def\mR{{\mathcal R}}
\def\mS{{\mathcal S}}\def\mT{\mathcal{T}}\def\mV{{\mathcal V}}
\def\mK{\mathcal K}\def\mA{{\mathcal A}}\def\mU{{\mathcal U}}
% Symmetric bilinear and trilinear forms of the Taylor expansion, used in the
% first Lyapunov coefficient (AAHOSN/New.tex)
\providecommand{\Bform}{\mathcal{B}}
\providecommand{\Cform}{\mathcal{C}}

% Blackboard bold.  Kept here, not in the drivers: a \providecommand in a
% driver preamble runs before this file and silently wins, so the same name
% ends up meaning different things in different builds.
\def\IC{\mathbb{C}}\def\IN{\mathbb{N}}\def\IQ{\mathbb{Q}}
\def\IR{\mathbb{R}}\def\IZ{\mathbb{Z}}\def\IG{\mathbb{G}}
\def\II{\mathbb{I}}

% Mathscr
\newcommand{\Rs}{{\mathscr R}}
\newcommand{\Ss}{{\mathscr S}}
\newcommand{\Ns}{{\mathscr N}}
\newcommand{\Cs}{\mathscr C}


%% ========================================
%% OPERATORS & FUNCTIONS
%% ========================================
\DeclareMathOperator{\tr}{tr}
\DeclareMathOperator{\diag}{Diag}
\DeclareMathOperator{\Diag}{Diag}
\DeclareMathOperator{\intr}{\mathrm{Int}}
\DeclareMathOperator{\Inv}{Inv}
\DeclareMathOperator{\spec}{spec}
\DeclareMathOperator{\AUC}{AUC}               % area under the curve
\newcommand{\ii}{\mathrm{i}}                  % imaginary unit
\newcommand\op{\operatorname}
\newcommand{\supp}{\op{supp}}
\newcommand{\diff}{\mathop{}\mathopen{}\mathrm{d}}
\newcommand{\pdd}[1]{\frac{\partial}{\partial{#1}}}
\DeclareMathOperator{\cone}{cone}
\DeclareMathOperator{\proj}{proj}

%% ========================================
%% BRACKETS & NORMS
%% ========================================
\providecommand{\pr}[1]{\left(#1\right)}
\providecommand{\pp}[1]{\left[#1\right]}
\newcommand\croc[1]{\left\langle #1\right\rangle}
\newcommand{\norm}[1]{\|#1\|}
\newcommand{\inner}[2]{\langle #1, #2 \rangle}
% Short synonym of \inner, used by AAHOSN/New.tex as $\ip{u}{v}=\bar u^\top v$
\providecommand{\ip}[2]{\langle #1 , #2 \rangle}


%% ========================================
%% FORMATTING & UTILITIES
%% ========================================
\newcommand*{\Scale}[2][4]{\scalebox{#1}{$#2$}}
\newcommand{\te}[1]{\text{#1}}
\newcommand\steq[1]{\stackrel{\text{\rm #1.}}{=}}
\newcommand\blue[1]{\textcolor{blue}{#1}}
\newcommand{\red}{\textcolor[rgb]{1.00,0.00,0.00}}
\newcommand{\green}{\textcolor[rgb]{0.00,1.00,0.00}}

% Figure command
\newcommand{\figu}[3]{%
\begin{figure}[H]
\centering
\includegraphics[scale=#3]{#1}
\caption{#2\label{f:#1}}
\end{figure}
}

% Footnote
\long\def\symbolfootnote[#1]#2{%
\begingroup
\def\thefootnote{\fnsymbol{footnote}}\footnote[#1]{#2}%
\endgroup}
\def\fn{\symbolfootnote}

\newcommand{\dsum}{\displaystyle\sum}

%% ========================================
%% CRN & EPIDEMIOLOGY SPECIFIC
%% ========================================
% CRN notation
\def\SS{\mathcal S}
\def\CC{\mathcal C}
\def\RR{\mathcal R}

% Reproduction numbers
\def\Rto{R_{1}^{E_2}}
\def\Rot{R_{2}^{E_1}}
\def\Roc{R_{1,c}}
\def\Rtc{R_{2,c}}
\def\rd{r_{0}}
\def\sd{s_{0}}
\def\dd{D_{0}}
\def\vd{v_{0}}
\def\zd{z_{0}}
\def\se{\sE}
\def\xd{x_{0}}
%% ========================================
%% SPECIAL VALUES & EQUILIBRIA
%% ========================================
\def\eE{\bar{e}}
\def\wE{\bar{w}}
\def\sE{s_{*}}
\def\Me{\T M_{ee}}
\def\Rnne{\mathbb{R}_{\ge 0}}

\newcommand{\ones}{\mathbf 1}


%% ========================================
%% ABBREVIATIONS (alphabetical by category)
%% ========================================

% -- General mathematical terms --
\def\adm{admissible}\def\admy{admissibility}
%\def\and{antisymmetric}% disabled 2026-09-30: it overrode LaTeX's author separator \and (printed "antisymmetric" between authors)
\def\ass{assumption}
\def\cadlag{c\`adl\`ag}
\def\ch{characteristic polynomial}
\def\chp{characteristic polynomial}
\def\coe{coefficient}
\def\comp{computation}
\def\con{condition}
\def\coo{coordinate}
\def\corr{corresponding}
\def\dec{decomposition}
\def\den{denominator}
\def\dete{deterministic}
\def\eig{eigenvalue}\def\eval{eigenvalue}
\def\eigv{eigenvector}\def\eiv{eigenvector}\def\evec{eigenvector}
\def\elt{element}
\def\enn{essentially nonnegative/positivity preserving}
\def\equ{equation}
\def\etal{et al.}
\def\fe{for example}
\def\fno{from now on}
\def\for{\forall}
\def\form{formula}
\def\Fr{Furthermore, }
\def\frt{furthermore}
\def\fun{function}
\def\gene{generalization}
\def\heu{heuristic}
\def\how{however}\def\How{However, }
\def\ie{i.e. }
\def\im{\item}
\def\isd{is defined by}
\def\itc{it is easy to check that}
\def\itf{it follows that}\def\Itf{It follows that}
\def\ith{it holds that}
\def\itm{it may be checked that}\def\Itm{It may be checked that}
\def\ito{it turns out that}\def\Ito{It turns out that}\def\Itou{It turns out useful to}
\def\ma{mass action}\def\MAR{mass action representation}
\def\mat{matrix}
\def\mbc{may be checked to be}
\def\mbw{may be written as}\def\mubw{must be written as}
\def\mpr{more precisely, }\def\Mp{More precisely, }
\def\nc{non-catalytic}
\def\nne{non-negative}\def\nny{nonnegativity}
\def\no{\nonumber}
\def\noni{non-infectious}
\def\nota{notation}
\def\np{\newpage}
\def\Nt{Note that}
\def\nt{non trivial}\def\nz{non zero}
\def\oth{otherwise}
\def\Oth{On the other hand}
\def\pa{parameter}\def\para{parameter}
\def\pol{polynomial}
\def\prob{probability}
\def\Prf{{\bf  Proof:}}\def\prf{{\bf  Proof:}}
\def\prop{proportional to}
\def\qu{\quad}
\def\QED{\hfill {$\square$}\goodbreak \medskip}
\def\rep{representation}
\def\resp{respectively}
\def\satd{satisfied}\def\satg{satisfying}\def\sats{satisfies}\def\saty{satisfy}
\def\ssec{\subsection}\def\sssec{\subsubsection}
\def\sef{see for example}
\def\sksy{skew-symmetric}
\def\sub{substrate}
\def\sys{system}
\def\var{variable}
\def\wk{well-known}
\def\wl{without loss of generality}\def\wlo{w.l.o.g.}
\def\wmc{We may check that}\def\Wmc{We may check that}
\def\wmr{we may rewrite}
\def\wms{we must show that}
\def\wrt{with respect to}

% -- CRN & Chemistry --
\newcommand\CRN{chemical reaction network}
\newcommand\CRNs{chemical reaction networks}
\newcommand\RN{reaction network}
\def\ACR{absolute concentration robustness}
\def\AMSD{acyclic siphon decomposition}
\def\DAE{differential algebraic equations}
\def\FHJ{Feinberg-Horn-Jackson}
\def\GMAK{generalized mass action kinetics}
\def\MM{Michaelis-Menten}
\def\sm{stoichiometric matrix}\def\Sm{stoichiometric matrix}
\def\WRZD{weakly reversible, zero deficiency}\def\aut{autocatalytic}

% -- Epidemiology --
\def\brn{basic reproduction number}\def\bzn{basic replacement number}
\def\com{compartment}
\def\dfe{disease free equilibrium}\def\DFE{disease free equilibrium}
\def\dwt{dwell times}\def\edt{expected dwell times}
\def\EM{epidemic model}
\def\endp{endemic fixed point}\def\enp{endemic point}\def\EE{endemic point}
\def\ME{mathematical epidemiology}
\def\MEE{eco-epidemiology}\def\EEP{eco-epidemiological}
\def\ece{eco-epidemiological model}
\def\mgm{multi group model}
\def\mtS{multi-strain model}\def\mSt{multi-strain model}
\def\mSi{minimal siphon}
\def\NGM{next generation matrix}
\def\nim{new infections matrix}
\def\qem{quadratic epidemiological model}
\def\repF{invasion functions}
\def\SPF{SIR-PH-FA model}
\def\sus{susceptible}

% -- Graph theory & Linear algebra --
\def\jac{Jacobi matrix}
\def\jin{Jacobian matrix of the invasion vector field with respect to the invasion variables}
\def\jrn{Jacobian matrix of the resident vector field with respect to the resident variables}
\def\jxn{Jacobian matrix of the $\x$ vector field with respect to the $\x$ variables}
\def\Lm{Laplacian matrix}
\def\MTT{matrix tree theorem}
\def\SGT{spectral graph theory}
\def\SMW{Sherman–Morrison–Woodbury}

% -- Stability & Dynamics --
\def\absC{absorbtion CTMC}
\def\BTb{Bogdanov-Takens bifurcation}
\def\CEP{competitive exclusion partition}
\def\CTMC{continuous time Markov chain}
\def\ds{differential system}
\def\fp{fixed point}\def\bfp{boundary fixed point}\def\fbp{fixed boundary point}
\def\fbl{fixed boundary line}\def\fif{forward invariant face}
\def\gas{global asymptotic stability}\def\GAS{globally asymptotically stable}
\def\GASS{\text{GASS}}
\def\Hb{Hopf bifurcation}
\def\Hur{Hurwitz}\def\MH{Metzler-Hurwitz}\def\RH{Routh-Hurwitz}
\def\IS{\text{IS}}
\def\LAS{locally asymptotically stable}
\def\LdS{Lyapunov diagonally stable}
\def\Lf{Lyapunov function}\def\LyV{Lyapunov Volterra type function}\def\RLF{robust Lyapunov function}
\def\LyM{Lyapunov Malkin stable}
\def\PF{Perron-Frobenius}
\def\spfp{strictly positive fixed point}
\def\sta{stability}
\def\STP{sharp threshold property}\def\STT{sharp threshold theorem}

% -- Models (with \varphi explicitly instead of \f) --
\def\AB{$(A,B)$ Arino-Brauer epidemic models}
\def\ABa{$(A,\bbe,\ba,\varphi)$ balanced bilinear model}
\def\ABP{%$(A,\bbe,P,\varphi)$
\\ balanced bilinear model}
\def\AMG{Anderson May Gupta}
\newcommand\BIN{biological interaction network}
\def\Dt{Descartes type}
\def\gLV{generalized Lotka-Volterra model}
\def\IG{\text{invasion graph}}\def\IGMS{interaction graph on minimal siphons}
\def\LCP{\text{linear complementarity problem}}
\def\LV{Lotka-Volterra}\def\LVK{Lotka-Volterra-Kolmogorov}

% -- Splitting & Structure --
\def\reg{regular splitting}\def\regS{{\bf regular splitting}}
\def\RUR{rational univariate representation}
\def\bds{block-diagonal structure}\def\NGMs{next generation matrices}

% -- Miscellaneous --
\def\Ma{{Mathematica}}
\def\Mm{Max–Min}
\def\num{numerator}
\def\rap{rational approximation} \def\DSR{directed species-reaction bipartite graph}\def\PCP{polynomial complementarity problem}\def\rat{rational substitution}
\def\QSS{quasi steady state} \def\RHS{right hand side} \def\LHS{left hand side}
\def\OSN{online social network}\def\VdD{Van den Driessche}
\def\wli{$\omega$-limit} \def\rk{rich kinetics}\def\sJ{symbolic Jacobian}
\def\cV{\cite{Vas,VasSta,BSV,GSV}} \def\spf{semi-parametric semi-functional}
\def\CRW{Capasso-Ruan-Wang} \def\ABP{$(A,\bbe,P,\f)$ balanced bilinear model}
\def\ABa{%$(A,\bbe,\ba,\f)$ \\
rank-one balanced bilinear model}
\def\bbe{B%{\bff \beta}
} \def\vbe{\vec {\bff \beta}
}
\def\vb{\vec {\bff b}}
\newcommand{\Kh}{\hat{K}}          % K with hat, first mention only
\newcommand{\wS}{w^{\sigma}}       % right Perron J^\perp, full block
\newcommand{\piS}{\vec \pi^{\sigma}}    % left Perron J^\perp, full block
\newcommand{\ws}{w^{\sigma}}       % right Perron J^\perp, sub-block
\newcommand{\pis}{\pi^{\sigma}}    % left Perron J^\perp, sub-block
\newcommand{\hwS}{\hat{w}^{\sigma}}   % right Perron K, full block
\newcommand{\hpiS}{\hat{\pi}^{\sigma}} % left Perron K, full block
\newcommand{\hws}{\hat{w}^{\sigma}}   % right Perron K, sub-block
\newcommand{\hpis}{\hat{\pi}^{\sigma}} % left Perron K, sub-block
\newcommand{\twS}{\tilde{w}^{\sigma}}  % right Perron \tilde K, full block
\newcommand{\tpiS}{\tilde{\pi}^{\sigma}} % left Perron \tilde K
\newcommand{\tws}{\tilde{w}^{\sigma}}
\newcommand{\tpis}{\tilde{\pi}^{\sigma}}
\newcommand{\Mat}[2]{\mathbb{R}_{\ge 0}^{\,#1\times #2}}
\def\tJ{J_\sigma^\perp} \def\lPe{left Perron eigenvector}
% \G is \Gamma (defined at :795).  The Volterra function keeps its own name,
% so that \G no longer has two conflicting meanings in this file.
\def\GV{G}  % Volterra G(u)=u-1-\ln u
\def\PVL{Perron-Volterra Lyapunov minimal siphon ansatz}
\newcommand{\istar}[1]{i_{#1}^{*}}
%\newcommand{\sE}{s_E}
\newcommand{\vE}{v_E}
\def\PV{Perron-Volterra candidate Lyapunov function}
\def\GLV{generalized Lotka-Volterra}
\def\MERN{mathematical epidemiology reaction networks}
\def\CMM{contraction metric method}
\newcommand{\sign}{\operatorname{sign}}
\newcommand{\Var}{\operatorname{Var}} 

% ======================================================================
% IDEAS SHARED BETWEEN THE ABSTRACT AND THE MAIN RESULTS SECTION
%
% Write each idea ONCE, in the abstract; the Main Results section recalls it,
% so the two can never drift apart.
%
%   in the abstract        \ideaitem{1}{The first-order reaction terms ...}
%   in Main Results        \ideahead{1}
%
% \ideaitem stores the text globally and typesets
%   \item \hyperlink{idea:1}{<text>}
% \ideahead typesets
%   \item\hypertarget{idea:1}{}\textbf{<the same text>}
%
% The abstract must come before the Main Results input, which is the normal
% order. If a heading has to read differently from the abstract sentence (for
% instance a long abstract item summarised in one line), override it with the
% optional argument: \ideahead[short heading]{6}.
% ======================================================================
\makeatletter
\newcommand{\ideaitem}[2]{%
  \expandafter\gdef\csname @idea@#1\endcsname{#2}%
  \item\hyperlink{idea:#1}{#2}}
\newcommand{\ideahead}[2][]{%
  \item\hypertarget{idea:#2}{}%
  \def\@ideatmp{#1}%
  \ifx\@ideatmp\@empty
    \@ifundefined{@idea@#2}{%
      \PackageWarning{preamble}{Idea #2 is used in Main Results but was never
        defined by \string\ideaitem\space in the abstract}%
      \textbf{[idea #2 undefined]}%
    }{\textbf{\begingroup\def\ideaterm##1##2{##2}\def\idearest##1{}\@nameuse{@idea@#2}\endgroup}}%
  \else
    \textbf{#1}%
  \fi}
% \ideaterm{<hypertarget>}{<words>}: inside an abstract idea the words become an
% invisible link (the document sets \hypersetup{hidelinks}) jumping to the place
% in Main Results that defines them; inside the Main Results heading itself,
% where the target sits a few lines below, \ideahead strips the link and prints
% the words plainly. Mark the destination with \hypertarget{<name>}{} on the line
% before the display, which adds no equation number and shows nothing.
\newcommand{\ideaterm}[2]{\hyperlink{#1}{#2}}
% \idearest{<text>}: the part of an abstract idea after its first sentence; printed in
% the abstract, omitted from the Main Results heading.
\newcommand{\idearest}[1]{#1}
\makeatother

% ======================================================================
% SHORT-PAPER SWITCH
% \shorttrue in a master that assembles only part of the companion paper
% (Barrat/SipArxiv.tex). \nsh{...} drops text that cites results living in the
% omitted sections, so the short build has no dangling references while the
% companion paper keeps the text verbatim.
% ======================================================================
\newif\ifshort \shortfalse
\newcommand{\nsh}[1]{\ifshort\else#1\fi}
% \shortalt{<short build>}{<companion paper>}: use when a sentence needs
% different wording in the two builds, not merely a deletion.
\newcommand{\shortalt}[2]{\ifshort#1\else#2\fi}
% \Jtr, \JtrS{\Sigma}: transverse Jacobian. J^\perp here; SipArxiv.tex redefines them as J_{\mathrm{tra}}, J_{\mathrm{tra},\Sigma}
\newcommand{\Jtr}{J^\perp}
\newcommand{\JtrS}[1]{J_{#1}^\perp}
.
\makeatletter
\@ifpackageloaded{hyperref}{}{%
  \@ifundefined{hyperrefpackage}{%
    \@ifundefined{AAdefNnohyperref}{\usepackage{hyperref}}{}%
  }{}%
}
\makeatother
% linktoc=all makes a table-of-contents page number a link as well as the
% entry text; hyperref's default, linktoc=section, links only the text.
% Deferred to begin-document and guarded: by the block above hyperref may
% be loaded later (mdpi's own hook) or not at all (\AAdefNnohyperref), and
% \hypersetup must exist before it is called.  A driver that wants the
% default back can say \hypersetup{linktoc=section} after % Shared preamble for the GLVbook papers.  This is the single canonical copy;
% a driver one level below GLVbook reaches it as
%
%     \input{../basicFiles/preamble}
%
% What stays in the driver, because it differs from paper to paper:
% \documentclass, font choice, \title, \author, \date, and any macro genuinely
% local to one paper.  Everything else belongs here.  The file is
% self-contained: it does not read any other file.

\usepackage[utf8]{inputenc}
\usepackage[novbox]{pdfsync}   % required: SyncTeX alone does not work here

% Page geometry, set once here for every paper that loads this file.  A driver
% that needs different margins says \geometry{...}; a second
% \usepackage[...]{geometry} raises an option clash and must not be used.
\makeatletter
\@ifundefined{AAdefNnogeometry}{\usepackage[a4paper,margin=2cm]{geometry}}{}
\makeatother

% mathematics
\usepackage{amsmath,amssymb,amsthm,mathtools,bm,mathrsfs,bbm}
\usepackage{centernot}

% floats, tables, lists
\usepackage{graphicx,float}
\makeatletter
\@ifpackageloaded{subfig}{}{\usepackage{subcaption}}
\makeatother
\usepackage{booktabs,longtable,threeparttable,array,multirow}
\usepackage{enumitem}

% colour, code, algorithms
\usepackage{xcolor,listings,tcolorbox}
\makeatletter
\@ifundefined{c@algorithm}{\usepackage{algorithm}}{}
\makeatother
\usepackage{algorithmicx,algpseudocode}

% diagrams
\usepackage{tikz,tikz-cd,pgfplots}
\usetikzlibrary{graphs,decorations.pathmorphing,decorations.markings,
arrows.meta,automata,positioning}
\tikzcdset{every label/.append style = {font = \small}}
\pgfplotsset{compat=1.18}

% chemistry
\usepackage[version=4]{mhchem}

\usepackage{microtype}
\usepackage{lineno}
\usepackage{nomencl}
\makenomenclature

% hyperref first, then bookmark, then glossaries: glossaries must follow
% hyperref or the entry links are built against the wrong anchors.
\makeatletter
\@ifundefined{hyperrefpackage}{%
  \@ifundefined{AAdefNnohyperref}{\usepackage{hyperref}}{}%
}{}
\@ifundefined{hyperrefpackage}{\usepackage{bookmark}}{}
\makeatother
\usepackage[toc,sort=use]{glossaries}
\makeglossaries

% Book/standalone switch.  Material that the book already covers in its
% prelude, but that a standalone paper still has to state, is wrapped as
%     \ifbookbuild\else ... \fi
% Drivers default to standalone; bookM.tex says \bookbuildtrue.
\newif\ifbookbuild
\bookbuildfalse

% Part switches for the gradostat paper.  Both default to true; a driver that
% wants the short submission says \madecfalse and/or \slowfastfalse, and every
% cross-reference into the dropped part disappears with it.
%   \ifmadec     Part III, the Lagasquie--Madec model and its break-even geometry
%   \ifslowfast  Part IV, everything carrying the small parameter \eps
\newif\ifmadec     \madectrue
\newif\ifslowfast  \slowfasttrue

% Part switch for the Barrat/simplicial-contagion project.  Defaults to
% true (the hypernetwork-synchrony part is included); the journal-paper
% driver says \hnetpartfalse to drop it and every cross-reference into it.
\newif\ifhnetpart  \hnetparttrue

%% ========================================
%% PACKAGES
%% ========================================
\RequirePackage{etoolbox}
\usepackage{makecell}
\usepackage[novbox]{pdfsync}% obsolete; SyncTeX is built into modern engines
\usepackage{amssymb, amsfonts,mathtools,bbm,eurosym,graphicx,latexsym, amsmath,amsthm,times,nccmath,url}
\usepackage{listings,verbatim,cases,tabularx}
% hyperref is deliberately NOT in the line above.  Journal classes often
% load it themselves, with a long option list, from a begindocument hook
% (mdpi.cls does: \DeclareHookRule{begindocument}{hyperref}{before}{mdpi}).
% Loading it here first, without options, makes that later load raise
% "Option clash for package hyperref".  So load it only if nobody else
% has, or will:
%   - \@ifpackageloaded  catches a class that loaded it already;
%   - \hyperrefpackage    is mdpi.cls's own deferred loader;
%   - \AAdefNnohyperref   lets any other driver opt out by defining it
%                         before \input{preamble}.
\makeatletter
\@ifpackageloaded{hyperref}{}{%
  \@ifundefined{hyperrefpackage}{%
    \@ifundefined{AAdefNnohyperref}{\usepackage{hyperref}}{}%
  }{}%
}
\makeatother
% linktoc=all makes a table-of-contents page number a link as well as the
% entry text; hyperref's default, linktoc=section, links only the text.
% Deferred to begin-document and guarded: by the block above hyperref may
% be loaded later (mdpi's own hook) or not at all (\AAdefNnohyperref), and
% \hypersetup must exist before it is called.  A driver that wants the
% default back can say \hypersetup{linktoc=section} after \input{preamble}.
\makeatletter
\AtBeginDocument{\@ifundefined{hypersetup}{}{\hypersetup{linktoc=all}}}
\makeatother
\usepackage{graphicx,float,fancybox}
\usepackage{epsfig,epstopdf,tikz}
\usepackage[sans]{dsfont}
\usetikzlibrary{positioning,arrows,arrows.meta,calc}
% Configure listings for Mathematica code
\lstdefinestyle{mathematica}{
 language=Mathematica,
 basicstyle=\small\ttfamily,
 backgroundcolor=\color{gray!10},
 frame=single,
 breaklines=true,
 commentstyle=\color{green!60!black},
 keywordstyle=\color{blue},
 stringstyle=\color{red},
 showstringspaces=false,
 tabsize=2
}
\usepackage{tcolorbox}
% Colored boxes for key results
\newtcolorbox{keyresult}[1][]{
 colback=blue!5!white,
 colframe=blue!75!black,
 title=#1,
 fonttitle=\bfseries
}
\newtcolorbox{examplebox}[1][]{
 colback=green!5!white,
 colframe=green!75!black,
 title=#1,
 fonttitle=\bfseries
}


\makeatletter
\@ifundefined{theorem}   {\newtheorem{theorem}   {Theorem}}   {}
\@ifundefined{definition}{\newtheorem{definition}{Definition}} {}
\@ifundefined{open}      {\newtheorem{open}      {Problem}}   {}
\@ifundefined{example}   {\newtheorem{example}   {Example}}   {}
\makeatother

% COUNTER DEFINITIONS
\newcounter{tabcounter}
\newcounter{defcounter}
\newcounter{opcounter}
\newcounter{thmcounter}
\newcounter{excounter}
\newcounter{ascounter}
\newcounter{cjcounter}
\newcounter{lemcounter}
\newcounter{remcounter}
\newcounter{corcounter}
\newcounter{propcounter}
\newcounter{figcounter}

% UNIFIED GLOSSARY STYLE
\newglossarystyle{unifiedstyle}{%
  \setglossarystyle{list}%
  \renewcommand*{\glossarypreamble}{\gdef\glslastsec{}}%
  \renewcommand{\glossentry}[2]{%
    \edef\glscursec{\glsentryuseriii{##1}}%
    \ifx\glscursec\glslastsec\else
      \item[]\textbf{\S\,\glscursec\enspace\glsentryuseriv{##1}}%
      \global\let\glslastsec\glscursec
    \fi
    \item[\glsentryuserii{##1} \glsentryuseri{##1}:] ``\glossentryname{##1}'' \glsentrydesc{##1}%
  }%
}

\makeatletter
% First-word extractor (used by \beT, \beD and \beXa for the default key)
% \@firstwordof hello world \relax  ->  hello
\def\@firstwordof#1 #2\relax{#1}
% First three letters, used for the heading mnemonic of a theorem given no
% explicit key.  Padding with \@empty keeps it safe for first words shorter
% than three letters: \@empty vanishes under \edef, whereas \relax would
% survive into the label and is illegal inside \csname (it broke every
% two-letter title, e.g. "ME ...", with "Missing \endcsname inserted").
\def\@takethree#1#2#3#4\@nil{#1#2#3}

% EXAMPLE COMMAND -- same signature and behaviour as \beT and \beD.
% Usage: \beXa*[long title][short title][key]
%   star  = suppresses the printed "[key]" tag after the title (as for \beT*
%   and \beD*); the label and glossary entry are unaffected.
%   key = optional label key override; the label becomes e:KEY, and without a
%   key it is e: plus the first three letters of the title's first word.  A
%   mnemonic label lets an example be cited as "Example~\ref{e:KEY}" instead
%   of by section number, which is what the citing prose needs.
\NewDocumentCommand{\beXa}{s O{} O{} o}{%
  \def\@temparg{#2}%
  \ifx\@temparg\@empty
    \begin{example}%
  \else
    \let\@savedSecTitle\@currentsectionname%
    \refstepcounter{excounter}%
    \IfValueTF{#4}{%
      \edef\@exlblkey{#4}%
    }{%
      \edef\@exfirstword{\@firstwordof#2 \relax}%
      \edef\@exlblkey{%
        \expandafter\@takethree\@exfirstword\@empty\@empty\@empty\@nil}%
    }%
    \def\@tempshort{#3}%
    \ifx\@tempshort\@empty
      \expandafter\xdef\csname @xshort@\@exlblkey\endcsname{\@exlblkey}%
    \else
      \expandafter\def\csname @xshort@\@exlblkey\endcsname{#3}%
    \fi
    \IfBooleanTF{#1}{\let\@extag\@empty}{%
      \edef\@extag{\noexpand\ \noexpand\textup{[\noexpand\texttt{e:\@exlblkey}]}}}%
    \begin{example}[#2\ifshowkeys\@extag\fi]%
    \label{e:\@exlblkey}%
    % record the long title so \nref{e:KEY} prints it, as for theorems
    \protected@write\@auxout{}{%
      \string\@thmtitle{e:\@exlblkey}{\unexpanded{#2}}}%
    \expanded{%
      \noexpand\newglossaryentry{e\theexcounter}{%
        name={\unexpanded{#2}%
          \IfBooleanF{#1}{ [\noexpand\texttt{e:\@exlblkey}]}},%
        sort={\theexample},%
        description={p.~\noexpand\pageref{e:\@exlblkey}},%
        user1={\theexample},%
        user2={Example},%
        user3={\thesection},%
        user4={\expandafter\unexpanded\expandafter{\@savedSecTitle}}%
      }%
    }%
    \protected@write\@auxout{}{\string\glsliveentry{e\theexcounter}}%
  \fi
}
\newcommand{\eeXa}{\end{example}}

% DEFINITION COMMAND -- same signature and behaviour as \beT below.
% Usage: \beD[long title][short title][key]
%   long title  = definition heading (plain text only)
%   short title = shown in \dref{word}; defaults to the key
%   key         = optional label key override; the label becomes d:KEY
% Without a key the label is d: plus the first three letters of the title's
% first word, so it is a mnemonic rather than the definition's number: a
% numeric label would silently change whenever a definition is inserted.
\NewDocumentCommand{\beD}{s O{} O{} o}{%
  \def\@temparg{#2}%
  \ifx\@temparg\@empty
    \begin{definition}%
  \else
    \let\@savedSecTitle\@currentsectionname%
    \refstepcounter{defcounter}%
    \IfValueTF{#4}{%
      \edef\@deflblkey{#4}%
    }{%
      \edef\@deffirstword{\@firstwordof#2 \relax}%
      \edef\@deflblkey{%
        \expandafter\@takethree\@deffirstword\@empty\@empty\@empty\@nil}%
    }%
    \def\@tempshort{#3}%
    \ifx\@tempshort\@empty
      \expandafter\xdef\csname @dshort@\@deflblkey\endcsname{\@deflblkey}%
    \else
      \expandafter\def\csname @dshort@\@deflblkey\endcsname{#3}%
    \fi
    \IfBooleanTF{#1}{\let\@deftag\@empty}{%
      \edef\@deftag{\noexpand\ \noexpand\textup{[\noexpand\texttt{d:\@deflblkey}]}}}%
    \begin{definition}[#2\ifshowkeys\@deftag\fi]%
    \label{d:\@deflblkey}%
    % record the long title so \nref{d:KEY} prints it, as for theorems
    \protected@write\@auxout{}{%
      \string\@thmtitle{d:\@deflblkey}{\unexpanded{#2}}}%
    \expanded{%
      \noexpand\newglossaryentry{d\thedefcounter}{%
        name={\unexpanded{#2}%
          \IfBooleanF{#1}{ [\noexpand\texttt{d:\@deflblkey}]}},%
        sort={\thedefinition},%
        description={p.~\noexpand\pageref{d:\@deflblkey}},%
        user1={\thedefinition},%
        user2={Definition},%
        user3={\thesection},%
        user4={\expandafter\unexpanded\expandafter{\@savedSecTitle}}%
      }%
    }%
    \protected@write\@auxout{}{\string\glsliveentry{d\thedefcounter}}%
  \fi
}
\newcommand{\eeD}{\end{definition}}

% \showkeysfalse in a driver hides the [t:KEY] mnemonics in result headings
% (labels and glossary entries are unchanged); used for submission builds.
\newif\ifshowkeys \showkeystrue
% \commonredtrue in a driver that inputs basicFiles/commonRed: tool.tex then
% cites its reduction results instead of restating them.
\newif\ifcommonred \commonredfalse
% TheorEM COMMAND
% Usage: \beT[long title][short title][key]
%   long title  = theorem heading (plain text only)
%   short title = shown in \tref{word}; may contain math; defaults to first word of long title
%   key         = optional label key override; label becomes t:KEY (avoids first-word collisions)
% Label is t:WORD where WORD = first word of long title (or KEY if given); use \tref{WORD} to reference.
\NewDocumentCommand{\beT}{s O{} O{} o}{%
  \def\@temparg{#2}%
  \ifx\@temparg\@empty
    \begin{theorem}%
  \else
    \let\@savedSecTitle\@currentsectionname%
    \refstepcounter{thmcounter}%
    % \@thmtag is the mnemonic shown in the heading; every theorem gets one.
    % With an explicit key the key itself is shown.  Without one the key
    % defaults to the title's first word, and showing that whole word would
    % merely repeat the title ("Dissipativity [Dissipativity]"), so only its
    % first three letters are shown: "Dissipativity [Dis]".  The label is
    % unaffected -- it still uses the full \@thmlblkey.
    \IfValueTF{#4}{%
      \edef\@thmlblkey{#4}%
    }{%
      \edef\@thmfirstword{\@firstwordof#2 \relax}%
      \edef\@thmlblkey{%
        \expandafter\@takethree\@thmfirstword\@empty\@empty\@empty\@nil}%
    }%
    \IfBooleanTF{#1}{\let\@thmtag\@empty}{%
      \edef\@thmtag{\noexpand\ \noexpand\textup{[\noexpand\texttt{t:\@thmlblkey}]}}}%
    \def\@tempshort{#3}%
    \ifx\@tempshort\@empty
      \expandafter\xdef\csname @tshort@\@thmlblkey\endcsname{\@thmlblkey}%
    \else
      \expandafter\def\csname @tshort@\@thmlblkey\endcsname{#3}%
    \fi
    \begin{theorem}[#2\ifshowkeys\@thmtag\fi]%
    \label{t:\@thmlblkey}%
    % record the long title for \nref; via the .aux so forward refs resolve
    \protected@write\@auxout{}{%
      \string\@thmtitle{t:\@thmlblkey}{\unexpanded{#2}}}%
    \expanded{%
      \noexpand\newglossaryentry{t\thethmcounter}{%
        name={\unexpanded{#2}%
          \IfBooleanF{#1}{ [\noexpand\texttt{t:\@thmlblkey}]}},%
        % sort/user1 must be the number the document actually prints, i.e. the
        % theorem environment's counter, not thmcounter.  thmcounter advances
        % only on \beT with a title, so a bare \beT (or \beT with an
        % empty title) makes the two drift and the glossary cites a wrong
        % number.  thmcounter is still fine as the entry key: it is unique.
        sort={\thetheorem},%
        description={p.~\noexpand\pageref{t:\@thmlblkey}},%
        user1={\thetheorem},%
        user2={Theorem},%
        user3={\thesection},%
        user4={\expandafter\unexpanded\expandafter{\@savedSecTitle}}%
      }%
    }%
    \protected@write\@auxout{}{\string\glsliveentry{t\thethmcounter}}%
  \fi
}
\newcommand{\eeT}{\end{theorem}}

% CROSS-PART RELATION PARAGRAPH
% Usage: \xpart{text}  -- place immediately BEFORE \beT, never inside it
%   (anything inside \beT...\eeT would be swallowed into the glossary entry).
% The paragraph is typeset only when Part I is actually present, i.e. in the
% combined book build; in the standalone builds (GranadaDSABNS/A.tex,
% AAHRah-Submission/AAHRah.tex) the label part:glv is undefined and the whole
% paragraph disappears rather than leaving ?? cross-references. Same
% graceful-degradation principle as \GASref in GASt.tex. Resolves on the
% second LaTeX pass, since r@part:glv comes from the .aux file.
% (Already inside the \makeatletter region opened above, so no \makeatletter
% pair here -- adding one would close that region early.)
\newcommand{\xpart}[1]{\@ifundefined{r@part:glv}{}{\paragraph{Relation with Parts~I--II.}#1}}

% \bookref{label}{precise text}{fallback text}: prints the precise cross-reference
% when the label resolves (combined book build), and the fallback text otherwise
% (standalone paper builds, where book-only appendix labels are absent). Same
% graceful-degradation role as \GASref in GASt.tex, but usable in any file since it
% lives in this preamble (loaded by every driver) -- unlike \GASref, which is defined in a
% content file and is therefore undefined in standalone builds that omit that file.
\newcommand{\bookref}[3]{\@ifundefined{r@#1}{#3}{#2}}

% \PVpart and \RBpart: the two parts that a partial build may omit.
% They print "Part~N" in any build that contains the part, and a naming phrase
% in a build that does not, so an excerpt shows no ?? .  Each already supplies
% the word "Part", so write "\PVpart develops ...", not "Part~\PVpart".
\newcommand{\PVpart}{\bookref{part:aahrah}%
  {Part~\ref{part:aahrah}}{the Perron--Volterra part of the full book}}
\newcommand{\RBpart}{\bookref{part:recbrowne}%
  {Part~\ref{part:recbrowne}}{the immuno-virology part of the full book}}

% OPEN PROBLEM COMMANDS
% ASSUMPTION COMMANDS (glossaried)
\newcommand{\beAs}[1][]{%
  \def\@temparg{#1}%
  \ifx\@temparg\@empty
    \begin{assumption}%
  \else
    \let\@savedSecTitle\@currentsectionname%
    \refstepcounter{ascounter}%
    \begin{assumption}[#1]%
    \label{as:\theascounter}%
    \expanded{%
      \noexpand\newglossaryentry{as\theascounter}{%
        name={\unexpanded{#1}},%
        sort={\theascounter},%
        description={p.~\noexpand\pageref{as:\theascounter}},%
        user1={\theascounter},%
        user2={Assumption},%
        user3={\thesection},%
        user4={\expandafter\unexpanded\expandafter{\@savedSecTitle}}%
      }%
    }%
    \protected@write\@auxout{}{\string\glsliveentry{as\theascounter}}%
  \fi
}
\newcommand{\eeAs}{\end{assumption}}
% CONJECTURE COMMANDS (glossaried)
\newcommand{\beCj}[1][]{%
  \def\@temparg{#1}%
  \ifx\@temparg\@empty
    \begin{conjecture}%
  \else
    \let\@savedSecTitle\@currentsectionname%
    \refstepcounter{cjcounter}%
    \begin{conjecture}[#1]%
    \label{cj:\thecjcounter}%
    \expanded{%
      \noexpand\newglossaryentry{cj\thecjcounter}{%
        name={\unexpanded{#1}},%
        sort={\thecjcounter},%
        description={p.~\noexpand\pageref{cj:\thecjcounter}},%
        user1={\thecjcounter},%
        user2={Conjecture},%
        user3={\thesection},%
        user4={\expandafter\unexpanded\expandafter{\@savedSecTitle}}%
      }%
    }%
    \protected@write\@auxout{}{\string\glsliveentry{cj\thecjcounter}}%
  \fi
}
\newcommand{\eeCj}{\end{conjecture}}
% LEMMA COMMANDS
% ---------------------------------------------------------------------
% UNIFORM RESULT WRAPPERS -- \beL, \beC, \beP, \beR share \beT's signature
%
%   \beX[long title][short title][key]  ...  \eeX
%
%   long title  = heading (math allowed)
%   short title = shown in the cross-reference; defaults to the key
%   key         = label key.  With a key the label is X:KEY and the heading
%                 carries the mnemonic [key], exactly as for \beT.  Without
%                 one the label falls back to the running counter, which is
%                 what every pre-existing call site in the book relies on.
%
% A bare \beX -- arguments commented out -- opens the plain environment and
% creates NO label and NO glossary entry.  That is the supported way to keep
% an individual result out of the catalogue.
% ---------------------------------------------------------------------
% #1 wrapper name  #2 environment  #3 label prefix  #4 counter  #5 type name
\newcommand{\@mkresult}[5]{%
  % A leading * suppresses the [key] mnemonic in both the heading and the
  % glossary name, while still making the label X:KEY and the entry:
  %     \beX*[long][short][key]
  \expandafter\NewDocumentCommand\csname be#1\endcsname{s O{} O{} o}{%
    \def\@temparg{##2}%
    \ifx\@temparg\@empty
      \begin{#2}%
    \else
      \let\@savedSecTitle\@currentsectionname%
      \refstepcounter{#4}%
      % With a key: mnemonic label X:KEY and the [key] tag in the heading.
      % Without one: fall back to the running counter, exactly as these
      % wrappers behaved before, so every legacy one-argument call site in
      % the book keeps its label.  The counter is also csname-safe, which a
      % title's first word is not -- titles beginning with math would
      % otherwise build \csname \Omega...\endcsname and fail.
      \IfValueTF{##4}{%
        \edef\@relbl{##4}%
        \IfBooleanTF{##1}{\let\@retag\@empty}{%
          \edef\@retag{\noexpand\ \noexpand\textup{[\noexpand\texttt{#3:##4}]}}}%
      }{%
        \edef\@relbl{\arabic{#4}}%
        \let\@retag\@empty
      }%
      \def\@tempshort{##3}%
      \ifx\@tempshort\@empty
        \expandafter\xdef\csname @sshort@#3\@relbl\endcsname{\@relbl}%
      \else
        \expandafter\def\csname @sshort@#3\@relbl\endcsname{##3}%
      \fi
      \begin{#2}[##2\ifshowkeys\@retag\fi]%
      \edef\@tmplbl{#3:\@relbl}\expandafter\label\expandafter{\@tmplbl}%
      \expanded{%
        \noexpand\newglossaryentry{#3\csname the#4\endcsname}{%
          name={\unexpanded{##2}%
            \IfBooleanF{##1}{\IfValueT{##4}{ [\noexpand\texttt{#3:##4}]}}},%
          sort={\csname the#2\endcsname},%
          description={p.~\noexpand\pageref{#3:\@relbl}},%
          user1={\csname the#2\endcsname},%
          user2={#5},%
          user3={\thesection},%
          user4={\expandafter\unexpanded\expandafter{\@savedSecTitle}}%
        }%
      }%
      \protected@write\@auxout{}{\string\glsliveentry{#3\csname the#4\endcsname}}%
      \protected@write\@auxout{}{\string\@thmtitle{\@tmplbl}{\unexpanded{##2}}}%
    \fi
  }%
  \expandafter\newcommand\csname ee#1\endcsname{\end{#2}}%
}
\@mkresult{L}{lemma}{l}{lemcounter}{Lemma}
% REMARK / COROLLARY / PROPOSITION -- same signature as \beL above
\@mkresult{R}{remark}{r}{remcounter}{Remark}
\@mkresult{C}{corollary}{c}{corcounter}{Corollary}
\@mkresult{P}{proposition}{p}{propcounter}{Proposition}
% OPEN-PROBLEM COMMAND -- generated with the same signature as \beT:
%   \beO[long title][short title][key]   ...   \eeO
% label o:KEY, or o:<counter> when no key is given; a bare \beO makes
% neither a label nor a glossary entry.
\@mkresult{O}{open}{o}{opcounter}{Problem}


% Section-level title tracker: updated only by \section, not \subsection.
% \@currentlabelname is overwritten by every \refstepcounter (including
% \subsection), so we save it separately here.
\newcommand{\@currentsectionname}{}
\AtBeginDocument{%
  \let\@epidcrn@orig@refstepcounter\refstepcounter
  \renewcommand{\refstepcounter}[1]{%
    \@epidcrn@orig@refstepcounter{#1}%
    \ifstrequal{#1}{section}{\global\let\@currentsectionname\@currentlabelname}{}%
  }%
}

% TABLE COMMANDS (glossaried) -- consistent naming with \beT, \beD, etc.
\newcommand{\beTa}[1][]{%
  \def\@temparg{#1}%
  \ifx\@temparg\@empty
    \begin{table}[htbp]%
  \else
    \let\@savedSecTitle\@currentsectionname%
    \refstepcounter{tabcounter}%
    \begin{table}[htbp]%
    \label{A:\thetabcounter}%
    \expanded{%
      \noexpand\newglossaryentry{A\thetabcounter}{%
        name={\unexpanded{#1}},%
        sort={\thetabcounter},%
        description={p.~\noexpand\pageref{A:\thetabcounter}},%
        user1={\thetabcounter},%
        user2={Table},%
        user3={\thesection},%
        user4={\expandafter\unexpanded\expandafter{\@savedSecTitle}}%
      }%
    }%
    \protected@write\@auxout{}{\string\glsliveentry{A\thetabcounter}}%
  \fi
}
\newcommand{\eeTa}{\end{table}}

% FIGURE COMMANDS (glossaried) -- consistent with \beT, \beR, \beTa, etc.
% Usage:
%   \beF[Glossary title]...\caption{Visible caption}\eeF  -> glossaried figure
%   \beF...\caption{Visible caption}\eeF                  -> plain [H] figure, no glossary
\newcommand{\beF}[1][]{%
  \def\@temparg{#1}%
  \ifx\@temparg\@empty
    \begin{figure}[H]%
  \else
    \begin{figure}[H]%
    \let\@savedSecTitle\@currentsectionname%
    \refstepcounter{figcounter}%
    \label{fig:\thefigcounter}%
    \expanded{%
      \noexpand\newglossaryentry{Fig\thefigcounter}{%
        name={\unexpanded{#1}},%
        sort={\thefigcounter},%
        description={p.~\noexpand\pageref{fig:\thefigcounter}},%
        user1={\thefigcounter},%
        user2={Figure},%
        user3={\thesection},%
        user4={\expandafter\unexpanded\expandafter{\@savedSecTitle}}%
      }%
    }%
    \protected@write\@auxout{}{\string\glsliveentry{Fig\thefigcounter}}%
  \fi
}
\newcommand{\eeF}{\end{figure}}
\makeatother
% LIVE-ENTRY REGISTER
% \glsaddall prints every entry the glossaries package currently knows, and
% that set includes entries restored from \jobname.glsdefs.  The package
% rewrites .glsdefs from all *currently defined* entries, so anything it
% restored is written straight back out: a deleted section kept reappearing in
% the catalogue for ever, and the .glsdefs had to be erased by hand.  The
% wrappers therefore record every entry they actually create in the .aux, and
% \printboth adds only those.  The .aux is rebuilt from scratch on each run, so
% this is self-healing: a removed result is gone after one further run.
\makeatletter
\newcommand{\@glslive}{}
\newcommand{\@glsaddlive}[1]{\ifglsentryexists{#1}{\glsadd{#1}}{}}
\newcommand{\glsliveentry}[1]{\g@addto@macro\@glslive{\@glsaddlive{#1}}}
\makeatother
% PRINT COMMAND
% No \clearpage here -- the catalogue continues on the current page.  This has
% been re-added by mistake several times; leave it out.
\makeatletter
\newcommand{\printboth}{%
  \@glslive   % add exactly the entries this document created last run
  \printglossary[style=unifiedstyle, title={0 \hspace{.22em} Definitions, Theorems, Examples and Open Problems}]%
}
\makeatother
% Named theorem reference: \tref{word} -> Thm~N(word-or-short)
% Shows stored short title if available, otherwise shows the key itself.
\newcommand{\tref}[1]{%
  Thm~\ref{t:#1}(\ifcsname @tshort@#1\endcsname\csname @tshort@#1\endcsname\else#1\fi)%
}

% Full-title theorem reference: \nref{t:KEY} prints the number and then the
% theorem's long title, so one writes "Theorem~\nref{t:blocks}" and gets
% "Theorem 3 ``Block factorization and siphons''".  Takes the complete label, like
% \ref, not the bare key like \tref.
%
% \beT writes each title to the .aux (see \@thmtitle below), so \nref resolves
% forward references from the second LaTeX pass, exactly as \ref does.  If the
% title is not yet known -- first pass, or a label that is not a \beT theorem
% -- it degrades to the number rather than producing ??.
% These two use @ in their names, and this point in the file is outside the
% \makeatletter region opened earlier (it closes at \makeatother above), so
% the region has to be reopened here.
\makeatletter
\newcommand{\@thmtitle}[2]{\global\@namedef{@tlong@#1}{#2}}
\newcommand{\nref}[1]{%
  \ref{#1}%
  \ifcsname @tlong@#1\endcsname
    \ ``\csname @tlong@#1\endcsname''%
  \fi
}
% \nm{t:KEY}: the stored long title of any \beT/\beL/\beP/\beR/\beC/\beD/\beE result,
% verbatim and without quotes, so item names follow the macro's first argument.
\newcommand{\nm}[1]{%
  \ifcsname @tlong@#1\endcsname
    \csname @tlong@#1\endcsname%
  \else
    \ref{#1}%
  \fi
}
\makeatother




%% ========================================
%% Other THEOREM ENVIRONMENTS
%% ========================================
\makeatletter
\@ifundefined{lemma}      {\newtheorem{lemma}      {Lemma}}      {}
\@ifundefined{proposition}{\newtheorem{proposition}{Proposition}} {}
\@ifundefined{corollary}  {\newtheorem{corollary}  {Corollary}}  {}
\@ifundefined{conjecture} {\newtheorem{conjecture} {Conjecture}}  {}
\@ifundefined{remark}     {\newtheorem{remark}     {Remark}}     {}
\@ifundefined{question}   {\newtheorem{question}   {Question}}   {}
\@ifundefined{assumption} {\newtheorem{assumption} {Assumption}} {}
\@ifundefined{hypothesis} {\newtheorem{hypothesis} {Hypothesis}} {}
\@ifundefined{fact}       {\newtheorem{fact}       {Fact}}       {}
\makeatother

% \citer: cite the relay-graph threshold theorem of AAHrel. If a document has
% \input the chapter/file containing t:relayRF (so the label resolves), point
% straight at it; otherwise fall back to a bare citation of AAHrel. Shared here
% (not per-paper) so it resolves consistently across every book chapter.
% \citeR: same idea, for the first theorem of AAHRah (t:Pstar).
\makeatletter
\providecommand{\citer}{%
  \@ifundefined{r@t:relayRF}%
    {\cite{AAHrel}}%
    {Theorem~\ref{t:relayRF} of \cite{AAHrel}}%
}
\providecommand{\citeR}{%
  \@ifundefined{r@t:Pstar}%
    {\cite{AAHRah}}%
    {Theorem~\ref{t:Pstar} of \cite{AAHRah}}%
}
\makeatother


% Abbreviations
\def\beQ{\begin{question}}\def\eeQ{\end{question}}
\def\beA{\begin{assumption}}\def\eeA{\end{assumption}}


%% ========================================

%% ========================================
%% COLORS & CUSTOM COMMANDS
%% ========================================
\definecolor{funccolor}{RGB}{25,25,112}
\definecolor{desccolor}{RGB}{64,64,64}
\newcommand{\func}[1]{\texttt{\textcolor{funccolor}{#1}}}
\def\eq{=}
\def\mm{-}

%% ========================================
%% MATHCAL SHORTHANDS
%% ========================================
% Project-wide convention: \mX for a capital X is always \mathcal X.
% Defined for the whole alphabet so the convention never fails to expand.
\newcommand{\mA}{\mathcal A}\newcommand{\mB}{\mathcal B}
\newcommand{\mC}{\mathcal C}\newcommand{\mD}{\mathcal D}
\newcommand{\mE}{\mathcal E}\newcommand{\mF}{\mathcal F}
\newcommand{\mG}{\mathcal G}\newcommand{\mH}{\mathcal H}
\newcommand{\mI}{\mathcal I}\newcommand{\mJ}{\mathcal J}
\newcommand{\mK}{\mathcal K}\newcommand{\mL}{\mathcal L}
\newcommand{\mM}{\mathcal M}\newcommand{\mN}{\mathcal N}
\newcommand{\mO}{\mathcal O}\newcommand{\mP}{\mathcal P}
\newcommand{\mQ}{\mathcal Q}\newcommand{\mR}{\mathcal R}
\newcommand{\mS}{\mathcal S}\newcommand{\mT}{\mathcal T}
\newcommand{\mU}{\mathcal U}\newcommand{\mV}{\mathcal V}
\newcommand{\mW}{\mathcal W}\newcommand{\mX}{\mathcal X}
\newcommand{\mY}{\mathcal Y}\newcommand{\mZ}{\mathcal Z}

%% ========================================
%% BASIC MATH ENVIRONMENTS
%% ========================================
\def\bep{\begin{pmatrix}}\def\eep{\end{pmatrix}}
\def\bev{\begin{vmatrix}}\def\eev{\end{vmatrix}}
\def\bea{\begin{eqnarray*}}\def\eea{\end{eqnarray*}}
\def\bc{\begin{cases}}\def\ec{\end{cases}}
\def\BEN{\begin{enumerate}}\def\EEN{\end{enumerate}}
\def\BI{\begin{itemize}}\def\EI{\end{itemize}}

\newcommand{\be}[1]{\begin{equation}\label{#1}}
\newcommand{\ee}{\end{equation}}
\newcommand{\beq}{\begin{eqnarray}}
\def\eeq{\end{eqnarray}}
\def\eqr{\eqref}\def\fr{\frac}\def\lbl{\label}
\def\Lra{\Longrightarrow}\def\non{\nonumber}\def\cNGM{\cite{Diek,Van,Van08}}
\def\Eq{\Leftrightarrow}\def\b{\beta}\def\uph{\red{UP TO HERE}}
%% ========================================
%% PROBABILITY & STATISTICS
%% ========================================
\renewcommand{\P}{\ensuremath{\mathbb{P}}}  % Note: overrides paragraph symbol
\newcommand{\E}{\ensuremath{\mathbb{E}}}
\newcommand{\ind}[1]{\ensuremath{\mathbbm{1}_{\left\{#1\right\}}}}
\DeclareMathOperator{\Cov}{Cov}
\DeclareMathOperator{\TV}{TV}
\newcommand{\rank}{\operatorname{rank}}

%% ========================================
%% CRN OPERATORS (organizations/ERCs/buffering/siphons framework)
%% ========================================
\def\clos{\mathrm{clos}}
\def\Supply{\mathrm{Supply}}
\def\req{\mathrm{req}}
\def\React{\mathrm{React}}
\def\Prod{\mathrm{Prod}}
\def\Gen{\mathrm{Gen}}
\def\In{\mathrm{In}}
\def\Out{\mathrm{Out}}
\def\Act{\mathrm{Act}}
\def\Min{\mathrm{Min}}

%% ========================================
%% LINE BREAKING
%% ========================================
% Last-resort stretch, used only on paragraphs TeX cannot otherwise break:
% a long \texttt{} token, a comma-separated inline math list (commas are not
% breakpoints in math), or a run of \cite/\ref boxes.  Without it those set as
% overfull \hbox and print into the margin; with it they set slightly loose,
% which is the lesser evil.  Does not affect paragraphs that already break.
\emergencystretch=3em

%% ========================================
%% NUMBER SETS
%% ========================================
\newcommand{\R}{\mathbb{R}}  % Real numbers
\newcommand{\C}{\mathbb{C}}  % Complex numbers
\newcommand{\N}{\mathbb{N}}
\newcommand{\Z}{\mathbb{Z}}
% \providecommand, not \newcommand: some papers define these in their own
% preamble after \input{preamble}, and must not clash.
\providecommand{\Q}{\mathbb{Q}}
\providecommand{\res}{\operatorname{res}}   % resultant; used by epiRes.tex
\newcommand{\MN}{\mathbb{N}}
\newcommand{\MZ}{\mathbb{Z}}
\newcommand{\realpos}{\mathbb{R}_{> 0}}
\newcommand{\realnneg}{\mathbb{R}_{\geq 0}}
\newcommand{\naturalpos}{\mathbb{N}_{>0}}
\newcommand{\Rplus}{\mathbb{R}_+}
\newcommand{\Rp}{\mathbb{R}_{\ge 0}}
\def\Rnne{\mathbb{R}_{\ge 0}}
\def\Se{\bar{S}} \def\Ie{\bar{I}}
\newcommand{\Rpos}{\mathbb{R}_{>0}}
\newcommand{\dd}{\,\mathrm d}
\newcommand{\e}{\mathrm e}
%% ========================================
%% GREEK LETTERS (2+ character abbreviations) greek
%% ========================================
\def\al{\alpha}\def\ba{\bff \alpha}
\def\w{\omega}\def\Lgn{\ga+\mu+\de}\def\om{\omega}
\def\de{\delta}\def\De{\Delta}\def\bde{\bff \delta}
\def\ep{\epsilon}\def\eps{\varepsilon}
\def\f{\varphi}  % Kept as requested
\def\ga{\gamma}\def\Ga{\Gamma}\def\G{\Gamma}\def\bg{\bff \gamma}
\def\up{\upsilon}  % NOT the breve accent: no accent commands anywhere.
                   % Named \up, not \u, so the LaTeX breve stays reachable.
\def\la{\lambda}\def\La{\Lambda}
\def\si{\sigma}\def\ka{\kappa}
\def\th{\theta}\def\vt{\vartheta}
\def\ze{\zeta}\def\Gap{\Gamma^{+}} \def\Gam{\Gamma^{-}}


%% ========================================
%% SPECIAL SHORT COMMANDS (kept as requested)
%% ========================================
\def\I{\infty}
\def\T{\widetilde}
\def\H{\hat}
\newcommand{\relint}{\operatorname{relint}}

%% ========================================
%% BOLD/VECTOR NOTATION
%% ========================================
\newcommand{\bff}[1]{{\mbox{\boldmath$#1$}}}
\newcommand\X{\boldsymbol{X}}
\def\x{\boldsymbol{x}}\def\bx{\bff x}\def\xe{\bar{\x}}
\newcommand\y{\boldsymbol{y}}\def\by{\boldsymbol{y}}\def\ye{\bar{\y}}
\def\yb{\y}\def\xb{\x}
\def\xy{\x-\y}
\newcommand\br{\boldsymbol{r}}\def\re{\bar{\r}}
\newcommand\bd{\boldsymbol{d}}
\newcommand\bQ{\boldsymbol{Q}}
\newcommand{\Yb}{{\mathbf{Y}}}
\newcommand{\Ab}{\mathbf{A}}
\newcommand{\zb}{\mathbf{z}}\def\bz{\bff 0}
\newcommand{\Zb}{\mathbf{Z}}
\newcommand{\bb}{\mathbf{b}}
\def\vbe{\vec {\bff \beta}}
\def\vb{\vec {\bff b}}
\def\v1{\vec {\bff 1}}
\def\vn{\vec \de}
\def\vz{\vec 0}
\def\vi{\overset{\rightarrow}{\mathsf i}}
\def\vv{\overset{\rightarrow}{v}}
\def\vr{\; \vec {\mathsf r}}
\newcommand{\va}{\bff a}
\providecommand{\vect}[1]{\boldsymbol{#1}}   % bold vector (moved from Gradostat.tex)
\providecommand{\one}{\mathbf 1}             % all-ones vector (moved from AliSF)
\providecommand{\spr}{s}                     % spectral bound  (moved from AliSF)
\providecommand{\Ke}{K_{\rm eff}}            % effective carrying capacity
% Critical right/left eigenvectors in the Hopf projection method (AAHOSN/New.tex)
\providecommand{\bq}{\mathbf{q}}
\providecommand{\bp}{\mathbf{p}}


%% ========================================
%% MATHCAL & MATHSCR
%% ========================================
% Mathcal
\newcommand\mc[1]{\mathcal{#1}}
\def\mC{{\mathcal C}}\def\mD{\mathcal D}\def\mE{\mathcal{E}}
\def\mF{\mathcal F}\def\mG{{\mathcal G}}\def\mH{{\mathcal H}}
\def\mI{\mathcal I}\def\mL{\mathcal{L}}\def\mM{{\mathcal M}}
\def\mN{{\mathcal N}}\def\mP{\mathcal{P}}\def\mR{{\mathcal R}}
\def\mS{{\mathcal S}}\def\mT{\mathcal{T}}\def\mV{{\mathcal V}}
\def\mK{\mathcal K}\def\mA{{\mathcal A}}\def\mU{{\mathcal U}}
% Symmetric bilinear and trilinear forms of the Taylor expansion, used in the
% first Lyapunov coefficient (AAHOSN/New.tex)
\providecommand{\Bform}{\mathcal{B}}
\providecommand{\Cform}{\mathcal{C}}

% Blackboard bold.  Kept here, not in the drivers: a \providecommand in a
% driver preamble runs before this file and silently wins, so the same name
% ends up meaning different things in different builds.
\def\IC{\mathbb{C}}\def\IN{\mathbb{N}}\def\IQ{\mathbb{Q}}
\def\IR{\mathbb{R}}\def\IZ{\mathbb{Z}}\def\IG{\mathbb{G}}
\def\II{\mathbb{I}}

% Mathscr
\newcommand{\Rs}{{\mathscr R}}
\newcommand{\Ss}{{\mathscr S}}
\newcommand{\Ns}{{\mathscr N}}
\newcommand{\Cs}{\mathscr C}


%% ========================================
%% OPERATORS & FUNCTIONS
%% ========================================
\DeclareMathOperator{\tr}{tr}
\DeclareMathOperator{\diag}{Diag}
\DeclareMathOperator{\Diag}{Diag}
\DeclareMathOperator{\intr}{\mathrm{Int}}
\DeclareMathOperator{\Inv}{Inv}
\DeclareMathOperator{\spec}{spec}
\DeclareMathOperator{\AUC}{AUC}               % area under the curve
\newcommand{\ii}{\mathrm{i}}                  % imaginary unit
\newcommand\op{\operatorname}
\newcommand{\supp}{\op{supp}}
\newcommand{\diff}{\mathop{}\mathopen{}\mathrm{d}}
\newcommand{\pdd}[1]{\frac{\partial}{\partial{#1}}}
\DeclareMathOperator{\cone}{cone}
\DeclareMathOperator{\proj}{proj}

%% ========================================
%% BRACKETS & NORMS
%% ========================================
\providecommand{\pr}[1]{\left(#1\right)}
\providecommand{\pp}[1]{\left[#1\right]}
\newcommand\croc[1]{\left\langle #1\right\rangle}
\newcommand{\norm}[1]{\|#1\|}
\newcommand{\inner}[2]{\langle #1, #2 \rangle}
% Short synonym of \inner, used by AAHOSN/New.tex as $\ip{u}{v}=\bar u^\top v$
\providecommand{\ip}[2]{\langle #1 , #2 \rangle}


%% ========================================
%% FORMATTING & UTILITIES
%% ========================================
\newcommand*{\Scale}[2][4]{\scalebox{#1}{$#2$}}
\newcommand{\te}[1]{\text{#1}}
\newcommand\steq[1]{\stackrel{\text{\rm #1.}}{=}}
\newcommand\blue[1]{\textcolor{blue}{#1}}
\newcommand{\red}{\textcolor[rgb]{1.00,0.00,0.00}}
\newcommand{\green}{\textcolor[rgb]{0.00,1.00,0.00}}

% Figure command
\newcommand{\figu}[3]{%
\begin{figure}[H]
\centering
\includegraphics[scale=#3]{#1}
\caption{#2\label{f:#1}}
\end{figure}
}

% Footnote
\long\def\symbolfootnote[#1]#2{%
\begingroup
\def\thefootnote{\fnsymbol{footnote}}\footnote[#1]{#2}%
\endgroup}
\def\fn{\symbolfootnote}

\newcommand{\dsum}{\displaystyle\sum}

%% ========================================
%% CRN & EPIDEMIOLOGY SPECIFIC
%% ========================================
% CRN notation
\def\SS{\mathcal S}
\def\CC{\mathcal C}
\def\RR{\mathcal R}

% Reproduction numbers
\def\Rto{R_{1}^{E_2}}
\def\Rot{R_{2}^{E_1}}
\def\Roc{R_{1,c}}
\def\Rtc{R_{2,c}}
\def\rd{r_{0}}
\def\sd{s_{0}}
\def\dd{D_{0}}
\def\vd{v_{0}}
\def\zd{z_{0}}
\def\se{\sE}
\def\xd{x_{0}}
%% ========================================
%% SPECIAL VALUES & EQUILIBRIA
%% ========================================
\def\eE{\bar{e}}
\def\wE{\bar{w}}
\def\sE{s_{*}}
\def\Me{\T M_{ee}}
\def\Rnne{\mathbb{R}_{\ge 0}}

\newcommand{\ones}{\mathbf 1}


%% ========================================
%% ABBREVIATIONS (alphabetical by category)
%% ========================================

% -- General mathematical terms --
\def\adm{admissible}\def\admy{admissibility}
%\def\and{antisymmetric}% disabled 2026-09-30: it overrode LaTeX's author separator \and (printed "antisymmetric" between authors)
\def\ass{assumption}
\def\cadlag{c\`adl\`ag}
\def\ch{characteristic polynomial}
\def\chp{characteristic polynomial}
\def\coe{coefficient}
\def\comp{computation}
\def\con{condition}
\def\coo{coordinate}
\def\corr{corresponding}
\def\dec{decomposition}
\def\den{denominator}
\def\dete{deterministic}
\def\eig{eigenvalue}\def\eval{eigenvalue}
\def\eigv{eigenvector}\def\eiv{eigenvector}\def\evec{eigenvector}
\def\elt{element}
\def\enn{essentially nonnegative/positivity preserving}
\def\equ{equation}
\def\etal{et al.}
\def\fe{for example}
\def\fno{from now on}
\def\for{\forall}
\def\form{formula}
\def\Fr{Furthermore, }
\def\frt{furthermore}
\def\fun{function}
\def\gene{generalization}
\def\heu{heuristic}
\def\how{however}\def\How{However, }
\def\ie{i.e. }
\def\im{\item}
\def\isd{is defined by}
\def\itc{it is easy to check that}
\def\itf{it follows that}\def\Itf{It follows that}
\def\ith{it holds that}
\def\itm{it may be checked that}\def\Itm{It may be checked that}
\def\ito{it turns out that}\def\Ito{It turns out that}\def\Itou{It turns out useful to}
\def\ma{mass action}\def\MAR{mass action representation}
\def\mat{matrix}
\def\mbc{may be checked to be}
\def\mbw{may be written as}\def\mubw{must be written as}
\def\mpr{more precisely, }\def\Mp{More precisely, }
\def\nc{non-catalytic}
\def\nne{non-negative}\def\nny{nonnegativity}
\def\no{\nonumber}
\def\noni{non-infectious}
\def\nota{notation}
\def\np{\newpage}
\def\Nt{Note that}
\def\nt{non trivial}\def\nz{non zero}
\def\oth{otherwise}
\def\Oth{On the other hand}
\def\pa{parameter}\def\para{parameter}
\def\pol{polynomial}
\def\prob{probability}
\def\Prf{{\bf  Proof:}}\def\prf{{\bf  Proof:}}
\def\prop{proportional to}
\def\qu{\quad}
\def\QED{\hfill {$\square$}\goodbreak \medskip}
\def\rep{representation}
\def\resp{respectively}
\def\satd{satisfied}\def\satg{satisfying}\def\sats{satisfies}\def\saty{satisfy}
\def\ssec{\subsection}\def\sssec{\subsubsection}
\def\sef{see for example}
\def\sksy{skew-symmetric}
\def\sub{substrate}
\def\sys{system}
\def\var{variable}
\def\wk{well-known}
\def\wl{without loss of generality}\def\wlo{w.l.o.g.}
\def\wmc{We may check that}\def\Wmc{We may check that}
\def\wmr{we may rewrite}
\def\wms{we must show that}
\def\wrt{with respect to}

% -- CRN & Chemistry --
\newcommand\CRN{chemical reaction network}
\newcommand\CRNs{chemical reaction networks}
\newcommand\RN{reaction network}
\def\ACR{absolute concentration robustness}
\def\AMSD{acyclic siphon decomposition}
\def\DAE{differential algebraic equations}
\def\FHJ{Feinberg-Horn-Jackson}
\def\GMAK{generalized mass action kinetics}
\def\MM{Michaelis-Menten}
\def\sm{stoichiometric matrix}\def\Sm{stoichiometric matrix}
\def\WRZD{weakly reversible, zero deficiency}\def\aut{autocatalytic}

% -- Epidemiology --
\def\brn{basic reproduction number}\def\bzn{basic replacement number}
\def\com{compartment}
\def\dfe{disease free equilibrium}\def\DFE{disease free equilibrium}
\def\dwt{dwell times}\def\edt{expected dwell times}
\def\EM{epidemic model}
\def\endp{endemic fixed point}\def\enp{endemic point}\def\EE{endemic point}
\def\ME{mathematical epidemiology}
\def\MEE{eco-epidemiology}\def\EEP{eco-epidemiological}
\def\ece{eco-epidemiological model}
\def\mgm{multi group model}
\def\mtS{multi-strain model}\def\mSt{multi-strain model}
\def\mSi{minimal siphon}
\def\NGM{next generation matrix}
\def\nim{new infections matrix}
\def\qem{quadratic epidemiological model}
\def\repF{invasion functions}
\def\SPF{SIR-PH-FA model}
\def\sus{susceptible}

% -- Graph theory & Linear algebra --
\def\jac{Jacobi matrix}
\def\jin{Jacobian matrix of the invasion vector field with respect to the invasion variables}
\def\jrn{Jacobian matrix of the resident vector field with respect to the resident variables}
\def\jxn{Jacobian matrix of the $\x$ vector field with respect to the $\x$ variables}
\def\Lm{Laplacian matrix}
\def\MTT{matrix tree theorem}
\def\SGT{spectral graph theory}
\def\SMW{Sherman–Morrison–Woodbury}

% -- Stability & Dynamics --
\def\absC{absorbtion CTMC}
\def\BTb{Bogdanov-Takens bifurcation}
\def\CEP{competitive exclusion partition}
\def\CTMC{continuous time Markov chain}
\def\ds{differential system}
\def\fp{fixed point}\def\bfp{boundary fixed point}\def\fbp{fixed boundary point}
\def\fbl{fixed boundary line}\def\fif{forward invariant face}
\def\gas{global asymptotic stability}\def\GAS{globally asymptotically stable}
\def\GASS{\text{GASS}}
\def\Hb{Hopf bifurcation}
\def\Hur{Hurwitz}\def\MH{Metzler-Hurwitz}\def\RH{Routh-Hurwitz}
\def\IS{\text{IS}}
\def\LAS{locally asymptotically stable}
\def\LdS{Lyapunov diagonally stable}
\def\Lf{Lyapunov function}\def\LyV{Lyapunov Volterra type function}\def\RLF{robust Lyapunov function}
\def\LyM{Lyapunov Malkin stable}
\def\PF{Perron-Frobenius}
\def\spfp{strictly positive fixed point}
\def\sta{stability}
\def\STP{sharp threshold property}\def\STT{sharp threshold theorem}

% -- Models (with \varphi explicitly instead of \f) --
\def\AB{$(A,B)$ Arino-Brauer epidemic models}
\def\ABa{$(A,\bbe,\ba,\varphi)$ balanced bilinear model}
\def\ABP{%$(A,\bbe,P,\varphi)$
\\ balanced bilinear model}
\def\AMG{Anderson May Gupta}
\newcommand\BIN{biological interaction network}
\def\Dt{Descartes type}
\def\gLV{generalized Lotka-Volterra model}
\def\IG{\text{invasion graph}}\def\IGMS{interaction graph on minimal siphons}
\def\LCP{\text{linear complementarity problem}}
\def\LV{Lotka-Volterra}\def\LVK{Lotka-Volterra-Kolmogorov}

% -- Splitting & Structure --
\def\reg{regular splitting}\def\regS{{\bf regular splitting}}
\def\RUR{rational univariate representation}
\def\bds{block-diagonal structure}\def\NGMs{next generation matrices}

% -- Miscellaneous --
\def\Ma{{Mathematica}}
\def\Mm{Max–Min}
\def\num{numerator}
\def\rap{rational approximation} \def\DSR{directed species-reaction bipartite graph}\def\PCP{polynomial complementarity problem}\def\rat{rational substitution}
\def\QSS{quasi steady state} \def\RHS{right hand side} \def\LHS{left hand side}
\def\OSN{online social network}\def\VdD{Van den Driessche}
\def\wli{$\omega$-limit} \def\rk{rich kinetics}\def\sJ{symbolic Jacobian}
\def\cV{\cite{Vas,VasSta,BSV,GSV}} \def\spf{semi-parametric semi-functional}
\def\CRW{Capasso-Ruan-Wang} \def\ABP{$(A,\bbe,P,\f)$ balanced bilinear model}
\def\ABa{%$(A,\bbe,\ba,\f)$ \\
rank-one balanced bilinear model}
\def\bbe{B%{\bff \beta}
} \def\vbe{\vec {\bff \beta}
}
\def\vb{\vec {\bff b}}
\newcommand{\Kh}{\hat{K}}          % K with hat, first mention only
\newcommand{\wS}{w^{\sigma}}       % right Perron J^\perp, full block
\newcommand{\piS}{\vec \pi^{\sigma}}    % left Perron J^\perp, full block
\newcommand{\ws}{w^{\sigma}}       % right Perron J^\perp, sub-block
\newcommand{\pis}{\pi^{\sigma}}    % left Perron J^\perp, sub-block
\newcommand{\hwS}{\hat{w}^{\sigma}}   % right Perron K, full block
\newcommand{\hpiS}{\hat{\pi}^{\sigma}} % left Perron K, full block
\newcommand{\hws}{\hat{w}^{\sigma}}   % right Perron K, sub-block
\newcommand{\hpis}{\hat{\pi}^{\sigma}} % left Perron K, sub-block
\newcommand{\twS}{\tilde{w}^{\sigma}}  % right Perron \tilde K, full block
\newcommand{\tpiS}{\tilde{\pi}^{\sigma}} % left Perron \tilde K
\newcommand{\tws}{\tilde{w}^{\sigma}}
\newcommand{\tpis}{\tilde{\pi}^{\sigma}}
\newcommand{\Mat}[2]{\mathbb{R}_{\ge 0}^{\,#1\times #2}}
\def\tJ{J_\sigma^\perp} \def\lPe{left Perron eigenvector}
% \G is \Gamma (defined at :795).  The Volterra function keeps its own name,
% so that \G no longer has two conflicting meanings in this file.
\def\GV{G}  % Volterra G(u)=u-1-\ln u
\def\PVL{Perron-Volterra Lyapunov minimal siphon ansatz}
\newcommand{\istar}[1]{i_{#1}^{*}}
%\newcommand{\sE}{s_E}
\newcommand{\vE}{v_E}
\def\PV{Perron-Volterra candidate Lyapunov function}
\def\GLV{generalized Lotka-Volterra}
\def\MERN{mathematical epidemiology reaction networks}
\def\CMM{contraction metric method}
\newcommand{\sign}{\operatorname{sign}}
\newcommand{\Var}{\operatorname{Var}} 

% ======================================================================
% IDEAS SHARED BETWEEN THE ABSTRACT AND THE MAIN RESULTS SECTION
%
% Write each idea ONCE, in the abstract; the Main Results section recalls it,
% so the two can never drift apart.
%
%   in the abstract        \ideaitem{1}{The first-order reaction terms ...}
%   in Main Results        \ideahead{1}
%
% \ideaitem stores the text globally and typesets
%   \item \hyperlink{idea:1}{<text>}
% \ideahead typesets
%   \item\hypertarget{idea:1}{}\textbf{<the same text>}
%
% The abstract must come before the Main Results input, which is the normal
% order. If a heading has to read differently from the abstract sentence (for
% instance a long abstract item summarised in one line), override it with the
% optional argument: \ideahead[short heading]{6}.
% ======================================================================
\makeatletter
\newcommand{\ideaitem}[2]{%
  \expandafter\gdef\csname @idea@#1\endcsname{#2}%
  \item\hyperlink{idea:#1}{#2}}
\newcommand{\ideahead}[2][]{%
  \item\hypertarget{idea:#2}{}%
  \def\@ideatmp{#1}%
  \ifx\@ideatmp\@empty
    \@ifundefined{@idea@#2}{%
      \PackageWarning{preamble}{Idea #2 is used in Main Results but was never
        defined by \string\ideaitem\space in the abstract}%
      \textbf{[idea #2 undefined]}%
    }{\textbf{\begingroup\def\ideaterm##1##2{##2}\def\idearest##1{}\@nameuse{@idea@#2}\endgroup}}%
  \else
    \textbf{#1}%
  \fi}
% \ideaterm{<hypertarget>}{<words>}: inside an abstract idea the words become an
% invisible link (the document sets \hypersetup{hidelinks}) jumping to the place
% in Main Results that defines them; inside the Main Results heading itself,
% where the target sits a few lines below, \ideahead strips the link and prints
% the words plainly. Mark the destination with \hypertarget{<name>}{} on the line
% before the display, which adds no equation number and shows nothing.
\newcommand{\ideaterm}[2]{\hyperlink{#1}{#2}}
% \idearest{<text>}: the part of an abstract idea after its first sentence; printed in
% the abstract, omitted from the Main Results heading.
\newcommand{\idearest}[1]{#1}
\makeatother

% ======================================================================
% SHORT-PAPER SWITCH
% \shorttrue in a master that assembles only part of the companion paper
% (Barrat/SipArxiv.tex). \nsh{...} drops text that cites results living in the
% omitted sections, so the short build has no dangling references while the
% companion paper keeps the text verbatim.
% ======================================================================
\newif\ifshort \shortfalse
\newcommand{\nsh}[1]{\ifshort\else#1\fi}
% \shortalt{<short build>}{<companion paper>}: use when a sentence needs
% different wording in the two builds, not merely a deletion.
\newcommand{\shortalt}[2]{\ifshort#1\else#2\fi}
% \Jtr, \JtrS{\Sigma}: transverse Jacobian. J^\perp here; SipArxiv.tex redefines them as J_{\mathrm{tra}}, J_{\mathrm{tra},\Sigma}
\newcommand{\Jtr}{J^\perp}
\newcommand{\JtrS}[1]{J_{#1}^\perp}
.
\makeatletter
\AtBeginDocument{\@ifundefined{hypersetup}{}{\hypersetup{linktoc=all}}}
\makeatother
\usepackage{graphicx,float,fancybox}
\usepackage{epsfig,epstopdf,tikz}
\usepackage[sans]{dsfont}
\usetikzlibrary{positioning,arrows,arrows.meta,calc}
% Configure listings for Mathematica code
\lstdefinestyle{mathematica}{
 language=Mathematica,
 basicstyle=\small\ttfamily,
 backgroundcolor=\color{gray!10},
 frame=single,
 breaklines=true,
 commentstyle=\color{green!60!black},
 keywordstyle=\color{blue},
 stringstyle=\color{red},
 showstringspaces=false,
 tabsize=2
}
\usepackage{tcolorbox}
% Colored boxes for key results
\newtcolorbox{keyresult}[1][]{
 colback=blue!5!white,
 colframe=blue!75!black,
 title=#1,
 fonttitle=\bfseries
}
\newtcolorbox{examplebox}[1][]{
 colback=green!5!white,
 colframe=green!75!black,
 title=#1,
 fonttitle=\bfseries
}


\makeatletter
\@ifundefined{theorem}   {\newtheorem{theorem}   {Theorem}}   {}
\@ifundefined{definition}{\newtheorem{definition}{Definition}} {}
\@ifundefined{open}      {\newtheorem{open}      {Problem}}   {}
\@ifundefined{example}   {\newtheorem{example}   {Example}}   {}
\makeatother

% COUNTER DEFINITIONS
\newcounter{tabcounter}
\newcounter{defcounter}
\newcounter{opcounter}
\newcounter{thmcounter}
\newcounter{excounter}
\newcounter{ascounter}
\newcounter{cjcounter}
\newcounter{lemcounter}
\newcounter{remcounter}
\newcounter{corcounter}
\newcounter{propcounter}
\newcounter{figcounter}

% UNIFIED GLOSSARY STYLE
\newglossarystyle{unifiedstyle}{%
  \setglossarystyle{list}%
  \renewcommand*{\glossarypreamble}{\gdef\glslastsec{}}%
  \renewcommand{\glossentry}[2]{%
    \edef\glscursec{\glsentryuseriii{##1}}%
    \ifx\glscursec\glslastsec\else
      \item[]\textbf{\S\,\glscursec\enspace\glsentryuseriv{##1}}%
      \global\let\glslastsec\glscursec
    \fi
    \item[\glsentryuserii{##1} \glsentryuseri{##1}:] ``\glossentryname{##1}'' \glsentrydesc{##1}%
  }%
}

\makeatletter
% First-word extractor (used by \beT, \beD and \beXa for the default key)
% \@firstwordof hello world \relax  ->  hello
\def\@firstwordof#1 #2\relax{#1}
% First three letters, used for the heading mnemonic of a theorem given no
% explicit key.  Padding with \@empty keeps it safe for first words shorter
% than three letters: \@empty vanishes under \edef, whereas \relax would
% survive into the label and is illegal inside \csname (it broke every
% two-letter title, e.g. "ME ...", with "Missing \endcsname inserted").
\def\@takethree#1#2#3#4\@nil{#1#2#3}

% EXAMPLE COMMAND -- same signature and behaviour as \beT and \beD.
% Usage: \beXa*[long title][short title][key]
%   star  = suppresses the printed "[key]" tag after the title (as for \beT*
%   and \beD*); the label and glossary entry are unaffected.
%   key = optional label key override; the label becomes e:KEY, and without a
%   key it is e: plus the first three letters of the title's first word.  A
%   mnemonic label lets an example be cited as "Example~\ref{e:KEY}" instead
%   of by section number, which is what the citing prose needs.
\NewDocumentCommand{\beXa}{s O{} O{} o}{%
  \def\@temparg{#2}%
  \ifx\@temparg\@empty
    \begin{example}%
  \else
    \let\@savedSecTitle\@currentsectionname%
    \refstepcounter{excounter}%
    \IfValueTF{#4}{%
      \edef\@exlblkey{#4}%
    }{%
      \edef\@exfirstword{\@firstwordof#2 \relax}%
      \edef\@exlblkey{%
        \expandafter\@takethree\@exfirstword\@empty\@empty\@empty\@nil}%
    }%
    \def\@tempshort{#3}%
    \ifx\@tempshort\@empty
      \expandafter\xdef\csname @xshort@\@exlblkey\endcsname{\@exlblkey}%
    \else
      \expandafter\def\csname @xshort@\@exlblkey\endcsname{#3}%
    \fi
    \IfBooleanTF{#1}{\let\@extag\@empty}{%
      \edef\@extag{\noexpand\ \noexpand\textup{[\noexpand\texttt{e:\@exlblkey}]}}}%
    \begin{example}[#2\ifshowkeys\@extag\fi]%
    \label{e:\@exlblkey}%
    % record the long title so \nref{e:KEY} prints it, as for theorems
    \protected@write\@auxout{}{%
      \string\@thmtitle{e:\@exlblkey}{\unexpanded{#2}}}%
    \expanded{%
      \noexpand\newglossaryentry{e\theexcounter}{%
        name={\unexpanded{#2}%
          \IfBooleanF{#1}{ [\noexpand\texttt{e:\@exlblkey}]}},%
        sort={\theexample},%
        description={p.~\noexpand\pageref{e:\@exlblkey}},%
        user1={\theexample},%
        user2={Example},%
        user3={\thesection},%
        user4={\expandafter\unexpanded\expandafter{\@savedSecTitle}}%
      }%
    }%
    \protected@write\@auxout{}{\string\glsliveentry{e\theexcounter}}%
  \fi
}
\newcommand{\eeXa}{\end{example}}

% DEFINITION COMMAND -- same signature and behaviour as \beT below.
% Usage: \beD[long title][short title][key]
%   long title  = definition heading (plain text only)
%   short title = shown in \dref{word}; defaults to the key
%   key         = optional label key override; the label becomes d:KEY
% Without a key the label is d: plus the first three letters of the title's
% first word, so it is a mnemonic rather than the definition's number: a
% numeric label would silently change whenever a definition is inserted.
\NewDocumentCommand{\beD}{s O{} O{} o}{%
  \def\@temparg{#2}%
  \ifx\@temparg\@empty
    \begin{definition}%
  \else
    \let\@savedSecTitle\@currentsectionname%
    \refstepcounter{defcounter}%
    \IfValueTF{#4}{%
      \edef\@deflblkey{#4}%
    }{%
      \edef\@deffirstword{\@firstwordof#2 \relax}%
      \edef\@deflblkey{%
        \expandafter\@takethree\@deffirstword\@empty\@empty\@empty\@nil}%
    }%
    \def\@tempshort{#3}%
    \ifx\@tempshort\@empty
      \expandafter\xdef\csname @dshort@\@deflblkey\endcsname{\@deflblkey}%
    \else
      \expandafter\def\csname @dshort@\@deflblkey\endcsname{#3}%
    \fi
    \IfBooleanTF{#1}{\let\@deftag\@empty}{%
      \edef\@deftag{\noexpand\ \noexpand\textup{[\noexpand\texttt{d:\@deflblkey}]}}}%
    \begin{definition}[#2\ifshowkeys\@deftag\fi]%
    \label{d:\@deflblkey}%
    % record the long title so \nref{d:KEY} prints it, as for theorems
    \protected@write\@auxout{}{%
      \string\@thmtitle{d:\@deflblkey}{\unexpanded{#2}}}%
    \expanded{%
      \noexpand\newglossaryentry{d\thedefcounter}{%
        name={\unexpanded{#2}%
          \IfBooleanF{#1}{ [\noexpand\texttt{d:\@deflblkey}]}},%
        sort={\thedefinition},%
        description={p.~\noexpand\pageref{d:\@deflblkey}},%
        user1={\thedefinition},%
        user2={Definition},%
        user3={\thesection},%
        user4={\expandafter\unexpanded\expandafter{\@savedSecTitle}}%
      }%
    }%
    \protected@write\@auxout{}{\string\glsliveentry{d\thedefcounter}}%
  \fi
}
\newcommand{\eeD}{\end{definition}}

% \showkeysfalse in a driver hides the [t:KEY] mnemonics in result headings
% (labels and glossary entries are unchanged); used for submission builds.
\newif\ifshowkeys \showkeystrue
% \commonredtrue in a driver that inputs basicFiles/commonRed: tool.tex then
% cites its reduction results instead of restating them.
\newif\ifcommonred \commonredfalse
% TheorEM COMMAND
% Usage: \beT[long title][short title][key]
%   long title  = theorem heading (plain text only)
%   short title = shown in \tref{word}; may contain math; defaults to first word of long title
%   key         = optional label key override; label becomes t:KEY (avoids first-word collisions)
% Label is t:WORD where WORD = first word of long title (or KEY if given); use \tref{WORD} to reference.
\NewDocumentCommand{\beT}{s O{} O{} o}{%
  \def\@temparg{#2}%
  \ifx\@temparg\@empty
    \begin{theorem}%
  \else
    \let\@savedSecTitle\@currentsectionname%
    \refstepcounter{thmcounter}%
    % \@thmtag is the mnemonic shown in the heading; every theorem gets one.
    % With an explicit key the key itself is shown.  Without one the key
    % defaults to the title's first word, and showing that whole word would
    % merely repeat the title ("Dissipativity [Dissipativity]"), so only its
    % first three letters are shown: "Dissipativity [Dis]".  The label is
    % unaffected -- it still uses the full \@thmlblkey.
    \IfValueTF{#4}{%
      \edef\@thmlblkey{#4}%
    }{%
      \edef\@thmfirstword{\@firstwordof#2 \relax}%
      \edef\@thmlblkey{%
        \expandafter\@takethree\@thmfirstword\@empty\@empty\@empty\@nil}%
    }%
    \IfBooleanTF{#1}{\let\@thmtag\@empty}{%
      \edef\@thmtag{\noexpand\ \noexpand\textup{[\noexpand\texttt{t:\@thmlblkey}]}}}%
    \def\@tempshort{#3}%
    \ifx\@tempshort\@empty
      \expandafter\xdef\csname @tshort@\@thmlblkey\endcsname{\@thmlblkey}%
    \else
      \expandafter\def\csname @tshort@\@thmlblkey\endcsname{#3}%
    \fi
    \begin{theorem}[#2\ifshowkeys\@thmtag\fi]%
    \label{t:\@thmlblkey}%
    % record the long title for \nref; via the .aux so forward refs resolve
    \protected@write\@auxout{}{%
      \string\@thmtitle{t:\@thmlblkey}{\unexpanded{#2}}}%
    \expanded{%
      \noexpand\newglossaryentry{t\thethmcounter}{%
        name={\unexpanded{#2}%
          \IfBooleanF{#1}{ [\noexpand\texttt{t:\@thmlblkey}]}},%
        % sort/user1 must be the number the document actually prints, i.e. the
        % theorem environment's counter, not thmcounter.  thmcounter advances
        % only on \beT with a title, so a bare \beT (or \beT with an
        % empty title) makes the two drift and the glossary cites a wrong
        % number.  thmcounter is still fine as the entry key: it is unique.
        sort={\thetheorem},%
        description={p.~\noexpand\pageref{t:\@thmlblkey}},%
        user1={\thetheorem},%
        user2={Theorem},%
        user3={\thesection},%
        user4={\expandafter\unexpanded\expandafter{\@savedSecTitle}}%
      }%
    }%
    \protected@write\@auxout{}{\string\glsliveentry{t\thethmcounter}}%
  \fi
}
\newcommand{\eeT}{\end{theorem}}

% CROSS-PART RELATION PARAGRAPH
% Usage: \xpart{text}  -- place immediately BEFORE \beT, never inside it
%   (anything inside \beT...\eeT would be swallowed into the glossary entry).
% The paragraph is typeset only when Part I is actually present, i.e. in the
% combined book build; in the standalone builds (GranadaDSABNS/A.tex,
% AAHRah-Submission/AAHRah.tex) the label part:glv is undefined and the whole
% paragraph disappears rather than leaving ?? cross-references. Same
% graceful-degradation principle as \GASref in GASt.tex. Resolves on the
% second LaTeX pass, since r@part:glv comes from the .aux file.
% (Already inside the \makeatletter region opened above, so no \makeatletter
% pair here -- adding one would close that region early.)
\newcommand{\xpart}[1]{\@ifundefined{r@part:glv}{}{\paragraph{Relation with Parts~I--II.}#1}}

% \bookref{label}{precise text}{fallback text}: prints the precise cross-reference
% when the label resolves (combined book build), and the fallback text otherwise
% (standalone paper builds, where book-only appendix labels are absent). Same
% graceful-degradation role as \GASref in GASt.tex, but usable in any file since it
% lives in this preamble (loaded by every driver) -- unlike \GASref, which is defined in a
% content file and is therefore undefined in standalone builds that omit that file.
\newcommand{\bookref}[3]{\@ifundefined{r@#1}{#3}{#2}}

% \PVpart and \RBpart: the two parts that a partial build may omit.
% They print "Part~N" in any build that contains the part, and a naming phrase
% in a build that does not, so an excerpt shows no ?? .  Each already supplies
% the word "Part", so write "\PVpart develops ...", not "Part~\PVpart".
\newcommand{\PVpart}{\bookref{part:aahrah}%
  {Part~\ref{part:aahrah}}{the Perron--Volterra part of the full book}}
\newcommand{\RBpart}{\bookref{part:recbrowne}%
  {Part~\ref{part:recbrowne}}{the immuno-virology part of the full book}}

% OPEN PROBLEM COMMANDS
% ASSUMPTION COMMANDS (glossaried)
\newcommand{\beAs}[1][]{%
  \def\@temparg{#1}%
  \ifx\@temparg\@empty
    \begin{assumption}%
  \else
    \let\@savedSecTitle\@currentsectionname%
    \refstepcounter{ascounter}%
    \begin{assumption}[#1]%
    \label{as:\theascounter}%
    \expanded{%
      \noexpand\newglossaryentry{as\theascounter}{%
        name={\unexpanded{#1}},%
        sort={\theascounter},%
        description={p.~\noexpand\pageref{as:\theascounter}},%
        user1={\theascounter},%
        user2={Assumption},%
        user3={\thesection},%
        user4={\expandafter\unexpanded\expandafter{\@savedSecTitle}}%
      }%
    }%
    \protected@write\@auxout{}{\string\glsliveentry{as\theascounter}}%
  \fi
}
\newcommand{\eeAs}{\end{assumption}}
% CONJECTURE COMMANDS (glossaried)
\newcommand{\beCj}[1][]{%
  \def\@temparg{#1}%
  \ifx\@temparg\@empty
    \begin{conjecture}%
  \else
    \let\@savedSecTitle\@currentsectionname%
    \refstepcounter{cjcounter}%
    \begin{conjecture}[#1]%
    \label{cj:\thecjcounter}%
    \expanded{%
      \noexpand\newglossaryentry{cj\thecjcounter}{%
        name={\unexpanded{#1}},%
        sort={\thecjcounter},%
        description={p.~\noexpand\pageref{cj:\thecjcounter}},%
        user1={\thecjcounter},%
        user2={Conjecture},%
        user3={\thesection},%
        user4={\expandafter\unexpanded\expandafter{\@savedSecTitle}}%
      }%
    }%
    \protected@write\@auxout{}{\string\glsliveentry{cj\thecjcounter}}%
  \fi
}
\newcommand{\eeCj}{\end{conjecture}}
% LEMMA COMMANDS
% ---------------------------------------------------------------------
% UNIFORM RESULT WRAPPERS -- \beL, \beC, \beP, \beR share \beT's signature
%
%   \beX[long title][short title][key]  ...  \eeX
%
%   long title  = heading (math allowed)
%   short title = shown in the cross-reference; defaults to the key
%   key         = label key.  With a key the label is X:KEY and the heading
%                 carries the mnemonic [key], exactly as for \beT.  Without
%                 one the label falls back to the running counter, which is
%                 what every pre-existing call site in the book relies on.
%
% A bare \beX -- arguments commented out -- opens the plain environment and
% creates NO label and NO glossary entry.  That is the supported way to keep
% an individual result out of the catalogue.
% ---------------------------------------------------------------------
% #1 wrapper name  #2 environment  #3 label prefix  #4 counter  #5 type name
\newcommand{\@mkresult}[5]{%
  % A leading * suppresses the [key] mnemonic in both the heading and the
  % glossary name, while still making the label X:KEY and the entry:
  %     \beX*[long][short][key]
  \expandafter\NewDocumentCommand\csname be#1\endcsname{s O{} O{} o}{%
    \def\@temparg{##2}%
    \ifx\@temparg\@empty
      \begin{#2}%
    \else
      \let\@savedSecTitle\@currentsectionname%
      \refstepcounter{#4}%
      % With a key: mnemonic label X:KEY and the [key] tag in the heading.
      % Without one: fall back to the running counter, exactly as these
      % wrappers behaved before, so every legacy one-argument call site in
      % the book keeps its label.  The counter is also csname-safe, which a
      % title's first word is not -- titles beginning with math would
      % otherwise build \csname \Omega...\endcsname and fail.
      \IfValueTF{##4}{%
        \edef\@relbl{##4}%
        \IfBooleanTF{##1}{\let\@retag\@empty}{%
          \edef\@retag{\noexpand\ \noexpand\textup{[\noexpand\texttt{#3:##4}]}}}%
      }{%
        \edef\@relbl{\arabic{#4}}%
        \let\@retag\@empty
      }%
      \def\@tempshort{##3}%
      \ifx\@tempshort\@empty
        \expandafter\xdef\csname @sshort@#3\@relbl\endcsname{\@relbl}%
      \else
        \expandafter\def\csname @sshort@#3\@relbl\endcsname{##3}%
      \fi
      \begin{#2}[##2\ifshowkeys\@retag\fi]%
      \edef\@tmplbl{#3:\@relbl}\expandafter\label\expandafter{\@tmplbl}%
      \expanded{%
        \noexpand\newglossaryentry{#3\csname the#4\endcsname}{%
          name={\unexpanded{##2}%
            \IfBooleanF{##1}{\IfValueT{##4}{ [\noexpand\texttt{#3:##4}]}}},%
          sort={\csname the#2\endcsname},%
          description={p.~\noexpand\pageref{#3:\@relbl}},%
          user1={\csname the#2\endcsname},%
          user2={#5},%
          user3={\thesection},%
          user4={\expandafter\unexpanded\expandafter{\@savedSecTitle}}%
        }%
      }%
      \protected@write\@auxout{}{\string\glsliveentry{#3\csname the#4\endcsname}}%
      \protected@write\@auxout{}{\string\@thmtitle{\@tmplbl}{\unexpanded{##2}}}%
    \fi
  }%
  \expandafter\newcommand\csname ee#1\endcsname{\end{#2}}%
}
\@mkresult{L}{lemma}{l}{lemcounter}{Lemma}
% REMARK / COROLLARY / PROPOSITION -- same signature as \beL above
\@mkresult{R}{remark}{r}{remcounter}{Remark}
\@mkresult{C}{corollary}{c}{corcounter}{Corollary}
\@mkresult{P}{proposition}{p}{propcounter}{Proposition}
% OPEN-PROBLEM COMMAND -- generated with the same signature as \beT:
%   \beO[long title][short title][key]   ...   \eeO
% label o:KEY, or o:<counter> when no key is given; a bare \beO makes
% neither a label nor a glossary entry.
\@mkresult{O}{open}{o}{opcounter}{Problem}


% Section-level title tracker: updated only by \section, not \subsection.
% \@currentlabelname is overwritten by every \refstepcounter (including
% \subsection), so we save it separately here.
\newcommand{\@currentsectionname}{}
\AtBeginDocument{%
  \let\@epidcrn@orig@refstepcounter\refstepcounter
  \renewcommand{\refstepcounter}[1]{%
    \@epidcrn@orig@refstepcounter{#1}%
    \ifstrequal{#1}{section}{\global\let\@currentsectionname\@currentlabelname}{}%
  }%
}

% TABLE COMMANDS (glossaried) -- consistent naming with \beT, \beD, etc.
\newcommand{\beTa}[1][]{%
  \def\@temparg{#1}%
  \ifx\@temparg\@empty
    \begin{table}[htbp]%
  \else
    \let\@savedSecTitle\@currentsectionname%
    \refstepcounter{tabcounter}%
    \begin{table}[htbp]%
    \label{A:\thetabcounter}%
    \expanded{%
      \noexpand\newglossaryentry{A\thetabcounter}{%
        name={\unexpanded{#1}},%
        sort={\thetabcounter},%
        description={p.~\noexpand\pageref{A:\thetabcounter}},%
        user1={\thetabcounter},%
        user2={Table},%
        user3={\thesection},%
        user4={\expandafter\unexpanded\expandafter{\@savedSecTitle}}%
      }%
    }%
    \protected@write\@auxout{}{\string\glsliveentry{A\thetabcounter}}%
  \fi
}
\newcommand{\eeTa}{\end{table}}

% FIGURE COMMANDS (glossaried) -- consistent with \beT, \beR, \beTa, etc.
% Usage:
%   \beF[Glossary title]...\caption{Visible caption}\eeF  -> glossaried figure
%   \beF...\caption{Visible caption}\eeF                  -> plain [H] figure, no glossary
\newcommand{\beF}[1][]{%
  \def\@temparg{#1}%
  \ifx\@temparg\@empty
    \begin{figure}[H]%
  \else
    \begin{figure}[H]%
    \let\@savedSecTitle\@currentsectionname%
    \refstepcounter{figcounter}%
    \label{fig:\thefigcounter}%
    \expanded{%
      \noexpand\newglossaryentry{Fig\thefigcounter}{%
        name={\unexpanded{#1}},%
        sort={\thefigcounter},%
        description={p.~\noexpand\pageref{fig:\thefigcounter}},%
        user1={\thefigcounter},%
        user2={Figure},%
        user3={\thesection},%
        user4={\expandafter\unexpanded\expandafter{\@savedSecTitle}}%
      }%
    }%
    \protected@write\@auxout{}{\string\glsliveentry{Fig\thefigcounter}}%
  \fi
}
\newcommand{\eeF}{\end{figure}}
\makeatother
% LIVE-ENTRY REGISTER
% \glsaddall prints every entry the glossaries package currently knows, and
% that set includes entries restored from \jobname.glsdefs.  The package
% rewrites .glsdefs from all *currently defined* entries, so anything it
% restored is written straight back out: a deleted section kept reappearing in
% the catalogue for ever, and the .glsdefs had to be erased by hand.  The
% wrappers therefore record every entry they actually create in the .aux, and
% \printboth adds only those.  The .aux is rebuilt from scratch on each run, so
% this is self-healing: a removed result is gone after one further run.
\makeatletter
\newcommand{\@glslive}{}
\newcommand{\@glsaddlive}[1]{\ifglsentryexists{#1}{\glsadd{#1}}{}}
\newcommand{\glsliveentry}[1]{\g@addto@macro\@glslive{\@glsaddlive{#1}}}
\makeatother
% PRINT COMMAND
% No \clearpage here -- the catalogue continues on the current page.  This has
% been re-added by mistake several times; leave it out.
\makeatletter
\newcommand{\printboth}{%
  \@glslive   % add exactly the entries this document created last run
  \printglossary[style=unifiedstyle, title={0 \hspace{.22em} Definitions, Theorems, Examples and Open Problems}]%
}
\makeatother
% Named theorem reference: \tref{word} -> Thm~N(word-or-short)
% Shows stored short title if available, otherwise shows the key itself.
\newcommand{\tref}[1]{%
  Thm~\ref{t:#1}(\ifcsname @tshort@#1\endcsname\csname @tshort@#1\endcsname\else#1\fi)%
}

% Full-title theorem reference: \nref{t:KEY} prints the number and then the
% theorem's long title, so one writes "Theorem~\nref{t:blocks}" and gets
% "Theorem 3 ``Block factorization and siphons''".  Takes the complete label, like
% \ref, not the bare key like \tref.
%
% \beT writes each title to the .aux (see \@thmtitle below), so \nref resolves
% forward references from the second LaTeX pass, exactly as \ref does.  If the
% title is not yet known -- first pass, or a label that is not a \beT theorem
% -- it degrades to the number rather than producing ??.
% These two use @ in their names, and this point in the file is outside the
% \makeatletter region opened earlier (it closes at \makeatother above), so
% the region has to be reopened here.
\makeatletter
\newcommand{\@thmtitle}[2]{\global\@namedef{@tlong@#1}{#2}}
\newcommand{\nref}[1]{%
  \ref{#1}%
  \ifcsname @tlong@#1\endcsname
    \ ``\csname @tlong@#1\endcsname''%
  \fi
}
% \nm{t:KEY}: the stored long title of any \beT/\beL/\beP/\beR/\beC/\beD/\beE result,
% verbatim and without quotes, so item names follow the macro's first argument.
\newcommand{\nm}[1]{%
  \ifcsname @tlong@#1\endcsname
    \csname @tlong@#1\endcsname%
  \else
    \ref{#1}%
  \fi
}
\makeatother




%% ========================================
%% Other THEOREM ENVIRONMENTS
%% ========================================
\makeatletter
\@ifundefined{lemma}      {\newtheorem{lemma}      {Lemma}}      {}
\@ifundefined{proposition}{\newtheorem{proposition}{Proposition}} {}
\@ifundefined{corollary}  {\newtheorem{corollary}  {Corollary}}  {}
\@ifundefined{conjecture} {\newtheorem{conjecture} {Conjecture}}  {}
\@ifundefined{remark}     {\newtheorem{remark}     {Remark}}     {}
\@ifundefined{question}   {\newtheorem{question}   {Question}}   {}
\@ifundefined{assumption} {\newtheorem{assumption} {Assumption}} {}
\@ifundefined{hypothesis} {\newtheorem{hypothesis} {Hypothesis}} {}
\@ifundefined{fact}       {\newtheorem{fact}       {Fact}}       {}
\makeatother

% \citer: cite the relay-graph threshold theorem of AAHrel. If a document has
% \input the chapter/file containing t:relayRF (so the label resolves), point
% straight at it; otherwise fall back to a bare citation of AAHrel. Shared here
% (not per-paper) so it resolves consistently across every book chapter.
% \citeR: same idea, for the first theorem of AAHRah (t:Pstar).
\makeatletter
\providecommand{\citer}{%
  \@ifundefined{r@t:relayRF}%
    {\cite{AAHrel}}%
    {Theorem~\ref{t:relayRF} of \cite{AAHrel}}%
}
\providecommand{\citeR}{%
  \@ifundefined{r@t:Pstar}%
    {\cite{AAHRah}}%
    {Theorem~\ref{t:Pstar} of \cite{AAHRah}}%
}
\makeatother


% Abbreviations
\def\beQ{\begin{question}}\def\eeQ{\end{question}}
\def\beA{\begin{assumption}}\def\eeA{\end{assumption}}


%% ========================================

%% ========================================
%% COLORS & CUSTOM COMMANDS
%% ========================================
\definecolor{funccolor}{RGB}{25,25,112}
\definecolor{desccolor}{RGB}{64,64,64}
\newcommand{\func}[1]{\texttt{\textcolor{funccolor}{#1}}}
\def\eq{=}
\def\mm{-}

%% ========================================
%% MATHCAL SHORTHANDS
%% ========================================
% Project-wide convention: \mX for a capital X is always \mathcal X.
% Defined for the whole alphabet so the convention never fails to expand.
\newcommand{\mA}{\mathcal A}\newcommand{\mB}{\mathcal B}
\newcommand{\mC}{\mathcal C}\newcommand{\mD}{\mathcal D}
\newcommand{\mE}{\mathcal E}\newcommand{\mF}{\mathcal F}
\newcommand{\mG}{\mathcal G}\newcommand{\mH}{\mathcal H}
\newcommand{\mI}{\mathcal I}\newcommand{\mJ}{\mathcal J}
\newcommand{\mK}{\mathcal K}\newcommand{\mL}{\mathcal L}
\newcommand{\mM}{\mathcal M}\newcommand{\mN}{\mathcal N}
\newcommand{\mO}{\mathcal O}\newcommand{\mP}{\mathcal P}
\newcommand{\mQ}{\mathcal Q}\newcommand{\mR}{\mathcal R}
\newcommand{\mS}{\mathcal S}\newcommand{\mT}{\mathcal T}
\newcommand{\mU}{\mathcal U}\newcommand{\mV}{\mathcal V}
\newcommand{\mW}{\mathcal W}\newcommand{\mX}{\mathcal X}
\newcommand{\mY}{\mathcal Y}\newcommand{\mZ}{\mathcal Z}

%% ========================================
%% BASIC MATH ENVIRONMENTS
%% ========================================
\def\bep{\begin{pmatrix}}\def\eep{\end{pmatrix}}
\def\bev{\begin{vmatrix}}\def\eev{\end{vmatrix}}
\def\bea{\begin{eqnarray*}}\def\eea{\end{eqnarray*}}
\def\bc{\begin{cases}}\def\ec{\end{cases}}
\def\BEN{\begin{enumerate}}\def\EEN{\end{enumerate}}
\def\BI{\begin{itemize}}\def\EI{\end{itemize}}

\newcommand{\be}[1]{\begin{equation}\label{#1}}
\newcommand{\ee}{\end{equation}}
\newcommand{\beq}{\begin{eqnarray}}
\def\eeq{\end{eqnarray}}
\def\eqr{\eqref}\def\fr{\frac}\def\lbl{\label}
\def\Lra{\Longrightarrow}\def\non{\nonumber}\def\cNGM{\cite{Diek,Van,Van08}}
\def\Eq{\Leftrightarrow}\def\b{\beta}\def\uph{\red{UP TO HERE}}
%% ========================================
%% PROBABILITY & STATISTICS
%% ========================================
\renewcommand{\P}{\ensuremath{\mathbb{P}}}  % Note: overrides paragraph symbol
\newcommand{\E}{\ensuremath{\mathbb{E}}}
\newcommand{\ind}[1]{\ensuremath{\mathbbm{1}_{\left\{#1\right\}}}}
\DeclareMathOperator{\Cov}{Cov}
\DeclareMathOperator{\TV}{TV}
\newcommand{\rank}{\operatorname{rank}}

%% ========================================
%% CRN OPERATORS (organizations/ERCs/buffering/siphons framework)
%% ========================================
\def\clos{\mathrm{clos}}
\def\Supply{\mathrm{Supply}}
\def\req{\mathrm{req}}
\def\React{\mathrm{React}}
\def\Prod{\mathrm{Prod}}
\def\Gen{\mathrm{Gen}}
\def\In{\mathrm{In}}
\def\Out{\mathrm{Out}}
\def\Act{\mathrm{Act}}
\def\Min{\mathrm{Min}}

%% ========================================
%% LINE BREAKING
%% ========================================
% Last-resort stretch, used only on paragraphs TeX cannot otherwise break:
% a long \texttt{} token, a comma-separated inline math list (commas are not
% breakpoints in math), or a run of \cite/\ref boxes.  Without it those set as
% overfull \hbox and print into the margin; with it they set slightly loose,
% which is the lesser evil.  Does not affect paragraphs that already break.
\emergencystretch=3em

%% ========================================
%% NUMBER SETS
%% ========================================
\newcommand{\R}{\mathbb{R}}  % Real numbers
\newcommand{\C}{\mathbb{C}}  % Complex numbers
\newcommand{\N}{\mathbb{N}}
\newcommand{\Z}{\mathbb{Z}}
% \providecommand, not \newcommand: some papers define these in their own
% preamble after % Shared preamble for the GLVbook papers.  This is the single canonical copy;
% a driver one level below GLVbook reaches it as
%
%     \input{../basicFiles/preamble}
%
% What stays in the driver, because it differs from paper to paper:
% \documentclass, font choice, \title, \author, \date, and any macro genuinely
% local to one paper.  Everything else belongs here.  The file is
% self-contained: it does not read any other file.

\usepackage[utf8]{inputenc}
\usepackage[novbox]{pdfsync}   % required: SyncTeX alone does not work here

% Page geometry, set once here for every paper that loads this file.  A driver
% that needs different margins says \geometry{...}; a second
% \usepackage[...]{geometry} raises an option clash and must not be used.
\makeatletter
\@ifundefined{AAdefNnogeometry}{\usepackage[a4paper,margin=2cm]{geometry}}{}
\makeatother

% mathematics
\usepackage{amsmath,amssymb,amsthm,mathtools,bm,mathrsfs,bbm}
\usepackage{centernot}

% floats, tables, lists
\usepackage{graphicx,float}
\makeatletter
\@ifpackageloaded{subfig}{}{\usepackage{subcaption}}
\makeatother
\usepackage{booktabs,longtable,threeparttable,array,multirow}
\usepackage{enumitem}

% colour, code, algorithms
\usepackage{xcolor,listings,tcolorbox}
\makeatletter
\@ifundefined{c@algorithm}{\usepackage{algorithm}}{}
\makeatother
\usepackage{algorithmicx,algpseudocode}

% diagrams
\usepackage{tikz,tikz-cd,pgfplots}
\usetikzlibrary{graphs,decorations.pathmorphing,decorations.markings,
arrows.meta,automata,positioning}
\tikzcdset{every label/.append style = {font = \small}}
\pgfplotsset{compat=1.18}

% chemistry
\usepackage[version=4]{mhchem}

\usepackage{microtype}
\usepackage{lineno}
\usepackage{nomencl}
\makenomenclature

% hyperref first, then bookmark, then glossaries: glossaries must follow
% hyperref or the entry links are built against the wrong anchors.
\makeatletter
\@ifundefined{hyperrefpackage}{%
  \@ifundefined{AAdefNnohyperref}{\usepackage{hyperref}}{}%
}{}
\@ifundefined{hyperrefpackage}{\usepackage{bookmark}}{}
\makeatother
\usepackage[toc,sort=use]{glossaries}
\makeglossaries

% Book/standalone switch.  Material that the book already covers in its
% prelude, but that a standalone paper still has to state, is wrapped as
%     \ifbookbuild\else ... \fi
% Drivers default to standalone; bookM.tex says \bookbuildtrue.
\newif\ifbookbuild
\bookbuildfalse

% Part switches for the gradostat paper.  Both default to true; a driver that
% wants the short submission says \madecfalse and/or \slowfastfalse, and every
% cross-reference into the dropped part disappears with it.
%   \ifmadec     Part III, the Lagasquie--Madec model and its break-even geometry
%   \ifslowfast  Part IV, everything carrying the small parameter \eps
\newif\ifmadec     \madectrue
\newif\ifslowfast  \slowfasttrue

% Part switch for the Barrat/simplicial-contagion project.  Defaults to
% true (the hypernetwork-synchrony part is included); the journal-paper
% driver says \hnetpartfalse to drop it and every cross-reference into it.
\newif\ifhnetpart  \hnetparttrue

%% ========================================
%% PACKAGES
%% ========================================
\RequirePackage{etoolbox}
\usepackage{makecell}
\usepackage[novbox]{pdfsync}% obsolete; SyncTeX is built into modern engines
\usepackage{amssymb, amsfonts,mathtools,bbm,eurosym,graphicx,latexsym, amsmath,amsthm,times,nccmath,url}
\usepackage{listings,verbatim,cases,tabularx}
% hyperref is deliberately NOT in the line above.  Journal classes often
% load it themselves, with a long option list, from a begindocument hook
% (mdpi.cls does: \DeclareHookRule{begindocument}{hyperref}{before}{mdpi}).
% Loading it here first, without options, makes that later load raise
% "Option clash for package hyperref".  So load it only if nobody else
% has, or will:
%   - \@ifpackageloaded  catches a class that loaded it already;
%   - \hyperrefpackage    is mdpi.cls's own deferred loader;
%   - \AAdefNnohyperref   lets any other driver opt out by defining it
%                         before \input{preamble}.
\makeatletter
\@ifpackageloaded{hyperref}{}{%
  \@ifundefined{hyperrefpackage}{%
    \@ifundefined{AAdefNnohyperref}{\usepackage{hyperref}}{}%
  }{}%
}
\makeatother
% linktoc=all makes a table-of-contents page number a link as well as the
% entry text; hyperref's default, linktoc=section, links only the text.
% Deferred to begin-document and guarded: by the block above hyperref may
% be loaded later (mdpi's own hook) or not at all (\AAdefNnohyperref), and
% \hypersetup must exist before it is called.  A driver that wants the
% default back can say \hypersetup{linktoc=section} after \input{preamble}.
\makeatletter
\AtBeginDocument{\@ifundefined{hypersetup}{}{\hypersetup{linktoc=all}}}
\makeatother
\usepackage{graphicx,float,fancybox}
\usepackage{epsfig,epstopdf,tikz}
\usepackage[sans]{dsfont}
\usetikzlibrary{positioning,arrows,arrows.meta,calc}
% Configure listings for Mathematica code
\lstdefinestyle{mathematica}{
 language=Mathematica,
 basicstyle=\small\ttfamily,
 backgroundcolor=\color{gray!10},
 frame=single,
 breaklines=true,
 commentstyle=\color{green!60!black},
 keywordstyle=\color{blue},
 stringstyle=\color{red},
 showstringspaces=false,
 tabsize=2
}
\usepackage{tcolorbox}
% Colored boxes for key results
\newtcolorbox{keyresult}[1][]{
 colback=blue!5!white,
 colframe=blue!75!black,
 title=#1,
 fonttitle=\bfseries
}
\newtcolorbox{examplebox}[1][]{
 colback=green!5!white,
 colframe=green!75!black,
 title=#1,
 fonttitle=\bfseries
}


\makeatletter
\@ifundefined{theorem}   {\newtheorem{theorem}   {Theorem}}   {}
\@ifundefined{definition}{\newtheorem{definition}{Definition}} {}
\@ifundefined{open}      {\newtheorem{open}      {Problem}}   {}
\@ifundefined{example}   {\newtheorem{example}   {Example}}   {}
\makeatother

% COUNTER DEFINITIONS
\newcounter{tabcounter}
\newcounter{defcounter}
\newcounter{opcounter}
\newcounter{thmcounter}
\newcounter{excounter}
\newcounter{ascounter}
\newcounter{cjcounter}
\newcounter{lemcounter}
\newcounter{remcounter}
\newcounter{corcounter}
\newcounter{propcounter}
\newcounter{figcounter}

% UNIFIED GLOSSARY STYLE
\newglossarystyle{unifiedstyle}{%
  \setglossarystyle{list}%
  \renewcommand*{\glossarypreamble}{\gdef\glslastsec{}}%
  \renewcommand{\glossentry}[2]{%
    \edef\glscursec{\glsentryuseriii{##1}}%
    \ifx\glscursec\glslastsec\else
      \item[]\textbf{\S\,\glscursec\enspace\glsentryuseriv{##1}}%
      \global\let\glslastsec\glscursec
    \fi
    \item[\glsentryuserii{##1} \glsentryuseri{##1}:] ``\glossentryname{##1}'' \glsentrydesc{##1}%
  }%
}

\makeatletter
% First-word extractor (used by \beT, \beD and \beXa for the default key)
% \@firstwordof hello world \relax  ->  hello
\def\@firstwordof#1 #2\relax{#1}
% First three letters, used for the heading mnemonic of a theorem given no
% explicit key.  Padding with \@empty keeps it safe for first words shorter
% than three letters: \@empty vanishes under \edef, whereas \relax would
% survive into the label and is illegal inside \csname (it broke every
% two-letter title, e.g. "ME ...", with "Missing \endcsname inserted").
\def\@takethree#1#2#3#4\@nil{#1#2#3}

% EXAMPLE COMMAND -- same signature and behaviour as \beT and \beD.
% Usage: \beXa*[long title][short title][key]
%   star  = suppresses the printed "[key]" tag after the title (as for \beT*
%   and \beD*); the label and glossary entry are unaffected.
%   key = optional label key override; the label becomes e:KEY, and without a
%   key it is e: plus the first three letters of the title's first word.  A
%   mnemonic label lets an example be cited as "Example~\ref{e:KEY}" instead
%   of by section number, which is what the citing prose needs.
\NewDocumentCommand{\beXa}{s O{} O{} o}{%
  \def\@temparg{#2}%
  \ifx\@temparg\@empty
    \begin{example}%
  \else
    \let\@savedSecTitle\@currentsectionname%
    \refstepcounter{excounter}%
    \IfValueTF{#4}{%
      \edef\@exlblkey{#4}%
    }{%
      \edef\@exfirstword{\@firstwordof#2 \relax}%
      \edef\@exlblkey{%
        \expandafter\@takethree\@exfirstword\@empty\@empty\@empty\@nil}%
    }%
    \def\@tempshort{#3}%
    \ifx\@tempshort\@empty
      \expandafter\xdef\csname @xshort@\@exlblkey\endcsname{\@exlblkey}%
    \else
      \expandafter\def\csname @xshort@\@exlblkey\endcsname{#3}%
    \fi
    \IfBooleanTF{#1}{\let\@extag\@empty}{%
      \edef\@extag{\noexpand\ \noexpand\textup{[\noexpand\texttt{e:\@exlblkey}]}}}%
    \begin{example}[#2\ifshowkeys\@extag\fi]%
    \label{e:\@exlblkey}%
    % record the long title so \nref{e:KEY} prints it, as for theorems
    \protected@write\@auxout{}{%
      \string\@thmtitle{e:\@exlblkey}{\unexpanded{#2}}}%
    \expanded{%
      \noexpand\newglossaryentry{e\theexcounter}{%
        name={\unexpanded{#2}%
          \IfBooleanF{#1}{ [\noexpand\texttt{e:\@exlblkey}]}},%
        sort={\theexample},%
        description={p.~\noexpand\pageref{e:\@exlblkey}},%
        user1={\theexample},%
        user2={Example},%
        user3={\thesection},%
        user4={\expandafter\unexpanded\expandafter{\@savedSecTitle}}%
      }%
    }%
    \protected@write\@auxout{}{\string\glsliveentry{e\theexcounter}}%
  \fi
}
\newcommand{\eeXa}{\end{example}}

% DEFINITION COMMAND -- same signature and behaviour as \beT below.
% Usage: \beD[long title][short title][key]
%   long title  = definition heading (plain text only)
%   short title = shown in \dref{word}; defaults to the key
%   key         = optional label key override; the label becomes d:KEY
% Without a key the label is d: plus the first three letters of the title's
% first word, so it is a mnemonic rather than the definition's number: a
% numeric label would silently change whenever a definition is inserted.
\NewDocumentCommand{\beD}{s O{} O{} o}{%
  \def\@temparg{#2}%
  \ifx\@temparg\@empty
    \begin{definition}%
  \else
    \let\@savedSecTitle\@currentsectionname%
    \refstepcounter{defcounter}%
    \IfValueTF{#4}{%
      \edef\@deflblkey{#4}%
    }{%
      \edef\@deffirstword{\@firstwordof#2 \relax}%
      \edef\@deflblkey{%
        \expandafter\@takethree\@deffirstword\@empty\@empty\@empty\@nil}%
    }%
    \def\@tempshort{#3}%
    \ifx\@tempshort\@empty
      \expandafter\xdef\csname @dshort@\@deflblkey\endcsname{\@deflblkey}%
    \else
      \expandafter\def\csname @dshort@\@deflblkey\endcsname{#3}%
    \fi
    \IfBooleanTF{#1}{\let\@deftag\@empty}{%
      \edef\@deftag{\noexpand\ \noexpand\textup{[\noexpand\texttt{d:\@deflblkey}]}}}%
    \begin{definition}[#2\ifshowkeys\@deftag\fi]%
    \label{d:\@deflblkey}%
    % record the long title so \nref{d:KEY} prints it, as for theorems
    \protected@write\@auxout{}{%
      \string\@thmtitle{d:\@deflblkey}{\unexpanded{#2}}}%
    \expanded{%
      \noexpand\newglossaryentry{d\thedefcounter}{%
        name={\unexpanded{#2}%
          \IfBooleanF{#1}{ [\noexpand\texttt{d:\@deflblkey}]}},%
        sort={\thedefinition},%
        description={p.~\noexpand\pageref{d:\@deflblkey}},%
        user1={\thedefinition},%
        user2={Definition},%
        user3={\thesection},%
        user4={\expandafter\unexpanded\expandafter{\@savedSecTitle}}%
      }%
    }%
    \protected@write\@auxout{}{\string\glsliveentry{d\thedefcounter}}%
  \fi
}
\newcommand{\eeD}{\end{definition}}

% \showkeysfalse in a driver hides the [t:KEY] mnemonics in result headings
% (labels and glossary entries are unchanged); used for submission builds.
\newif\ifshowkeys \showkeystrue
% \commonredtrue in a driver that inputs basicFiles/commonRed: tool.tex then
% cites its reduction results instead of restating them.
\newif\ifcommonred \commonredfalse
% TheorEM COMMAND
% Usage: \beT[long title][short title][key]
%   long title  = theorem heading (plain text only)
%   short title = shown in \tref{word}; may contain math; defaults to first word of long title
%   key         = optional label key override; label becomes t:KEY (avoids first-word collisions)
% Label is t:WORD where WORD = first word of long title (or KEY if given); use \tref{WORD} to reference.
\NewDocumentCommand{\beT}{s O{} O{} o}{%
  \def\@temparg{#2}%
  \ifx\@temparg\@empty
    \begin{theorem}%
  \else
    \let\@savedSecTitle\@currentsectionname%
    \refstepcounter{thmcounter}%
    % \@thmtag is the mnemonic shown in the heading; every theorem gets one.
    % With an explicit key the key itself is shown.  Without one the key
    % defaults to the title's first word, and showing that whole word would
    % merely repeat the title ("Dissipativity [Dissipativity]"), so only its
    % first three letters are shown: "Dissipativity [Dis]".  The label is
    % unaffected -- it still uses the full \@thmlblkey.
    \IfValueTF{#4}{%
      \edef\@thmlblkey{#4}%
    }{%
      \edef\@thmfirstword{\@firstwordof#2 \relax}%
      \edef\@thmlblkey{%
        \expandafter\@takethree\@thmfirstword\@empty\@empty\@empty\@nil}%
    }%
    \IfBooleanTF{#1}{\let\@thmtag\@empty}{%
      \edef\@thmtag{\noexpand\ \noexpand\textup{[\noexpand\texttt{t:\@thmlblkey}]}}}%
    \def\@tempshort{#3}%
    \ifx\@tempshort\@empty
      \expandafter\xdef\csname @tshort@\@thmlblkey\endcsname{\@thmlblkey}%
    \else
      \expandafter\def\csname @tshort@\@thmlblkey\endcsname{#3}%
    \fi
    \begin{theorem}[#2\ifshowkeys\@thmtag\fi]%
    \label{t:\@thmlblkey}%
    % record the long title for \nref; via the .aux so forward refs resolve
    \protected@write\@auxout{}{%
      \string\@thmtitle{t:\@thmlblkey}{\unexpanded{#2}}}%
    \expanded{%
      \noexpand\newglossaryentry{t\thethmcounter}{%
        name={\unexpanded{#2}%
          \IfBooleanF{#1}{ [\noexpand\texttt{t:\@thmlblkey}]}},%
        % sort/user1 must be the number the document actually prints, i.e. the
        % theorem environment's counter, not thmcounter.  thmcounter advances
        % only on \beT with a title, so a bare \beT (or \beT with an
        % empty title) makes the two drift and the glossary cites a wrong
        % number.  thmcounter is still fine as the entry key: it is unique.
        sort={\thetheorem},%
        description={p.~\noexpand\pageref{t:\@thmlblkey}},%
        user1={\thetheorem},%
        user2={Theorem},%
        user3={\thesection},%
        user4={\expandafter\unexpanded\expandafter{\@savedSecTitle}}%
      }%
    }%
    \protected@write\@auxout{}{\string\glsliveentry{t\thethmcounter}}%
  \fi
}
\newcommand{\eeT}{\end{theorem}}

% CROSS-PART RELATION PARAGRAPH
% Usage: \xpart{text}  -- place immediately BEFORE \beT, never inside it
%   (anything inside \beT...\eeT would be swallowed into the glossary entry).
% The paragraph is typeset only when Part I is actually present, i.e. in the
% combined book build; in the standalone builds (GranadaDSABNS/A.tex,
% AAHRah-Submission/AAHRah.tex) the label part:glv is undefined and the whole
% paragraph disappears rather than leaving ?? cross-references. Same
% graceful-degradation principle as \GASref in GASt.tex. Resolves on the
% second LaTeX pass, since r@part:glv comes from the .aux file.
% (Already inside the \makeatletter region opened above, so no \makeatletter
% pair here -- adding one would close that region early.)
\newcommand{\xpart}[1]{\@ifundefined{r@part:glv}{}{\paragraph{Relation with Parts~I--II.}#1}}

% \bookref{label}{precise text}{fallback text}: prints the precise cross-reference
% when the label resolves (combined book build), and the fallback text otherwise
% (standalone paper builds, where book-only appendix labels are absent). Same
% graceful-degradation role as \GASref in GASt.tex, but usable in any file since it
% lives in this preamble (loaded by every driver) -- unlike \GASref, which is defined in a
% content file and is therefore undefined in standalone builds that omit that file.
\newcommand{\bookref}[3]{\@ifundefined{r@#1}{#3}{#2}}

% \PVpart and \RBpart: the two parts that a partial build may omit.
% They print "Part~N" in any build that contains the part, and a naming phrase
% in a build that does not, so an excerpt shows no ?? .  Each already supplies
% the word "Part", so write "\PVpart develops ...", not "Part~\PVpart".
\newcommand{\PVpart}{\bookref{part:aahrah}%
  {Part~\ref{part:aahrah}}{the Perron--Volterra part of the full book}}
\newcommand{\RBpart}{\bookref{part:recbrowne}%
  {Part~\ref{part:recbrowne}}{the immuno-virology part of the full book}}

% OPEN PROBLEM COMMANDS
% ASSUMPTION COMMANDS (glossaried)
\newcommand{\beAs}[1][]{%
  \def\@temparg{#1}%
  \ifx\@temparg\@empty
    \begin{assumption}%
  \else
    \let\@savedSecTitle\@currentsectionname%
    \refstepcounter{ascounter}%
    \begin{assumption}[#1]%
    \label{as:\theascounter}%
    \expanded{%
      \noexpand\newglossaryentry{as\theascounter}{%
        name={\unexpanded{#1}},%
        sort={\theascounter},%
        description={p.~\noexpand\pageref{as:\theascounter}},%
        user1={\theascounter},%
        user2={Assumption},%
        user3={\thesection},%
        user4={\expandafter\unexpanded\expandafter{\@savedSecTitle}}%
      }%
    }%
    \protected@write\@auxout{}{\string\glsliveentry{as\theascounter}}%
  \fi
}
\newcommand{\eeAs}{\end{assumption}}
% CONJECTURE COMMANDS (glossaried)
\newcommand{\beCj}[1][]{%
  \def\@temparg{#1}%
  \ifx\@temparg\@empty
    \begin{conjecture}%
  \else
    \let\@savedSecTitle\@currentsectionname%
    \refstepcounter{cjcounter}%
    \begin{conjecture}[#1]%
    \label{cj:\thecjcounter}%
    \expanded{%
      \noexpand\newglossaryentry{cj\thecjcounter}{%
        name={\unexpanded{#1}},%
        sort={\thecjcounter},%
        description={p.~\noexpand\pageref{cj:\thecjcounter}},%
        user1={\thecjcounter},%
        user2={Conjecture},%
        user3={\thesection},%
        user4={\expandafter\unexpanded\expandafter{\@savedSecTitle}}%
      }%
    }%
    \protected@write\@auxout{}{\string\glsliveentry{cj\thecjcounter}}%
  \fi
}
\newcommand{\eeCj}{\end{conjecture}}
% LEMMA COMMANDS
% ---------------------------------------------------------------------
% UNIFORM RESULT WRAPPERS -- \beL, \beC, \beP, \beR share \beT's signature
%
%   \beX[long title][short title][key]  ...  \eeX
%
%   long title  = heading (math allowed)
%   short title = shown in the cross-reference; defaults to the key
%   key         = label key.  With a key the label is X:KEY and the heading
%                 carries the mnemonic [key], exactly as for \beT.  Without
%                 one the label falls back to the running counter, which is
%                 what every pre-existing call site in the book relies on.
%
% A bare \beX -- arguments commented out -- opens the plain environment and
% creates NO label and NO glossary entry.  That is the supported way to keep
% an individual result out of the catalogue.
% ---------------------------------------------------------------------
% #1 wrapper name  #2 environment  #3 label prefix  #4 counter  #5 type name
\newcommand{\@mkresult}[5]{%
  % A leading * suppresses the [key] mnemonic in both the heading and the
  % glossary name, while still making the label X:KEY and the entry:
  %     \beX*[long][short][key]
  \expandafter\NewDocumentCommand\csname be#1\endcsname{s O{} O{} o}{%
    \def\@temparg{##2}%
    \ifx\@temparg\@empty
      \begin{#2}%
    \else
      \let\@savedSecTitle\@currentsectionname%
      \refstepcounter{#4}%
      % With a key: mnemonic label X:KEY and the [key] tag in the heading.
      % Without one: fall back to the running counter, exactly as these
      % wrappers behaved before, so every legacy one-argument call site in
      % the book keeps its label.  The counter is also csname-safe, which a
      % title's first word is not -- titles beginning with math would
      % otherwise build \csname \Omega...\endcsname and fail.
      \IfValueTF{##4}{%
        \edef\@relbl{##4}%
        \IfBooleanTF{##1}{\let\@retag\@empty}{%
          \edef\@retag{\noexpand\ \noexpand\textup{[\noexpand\texttt{#3:##4}]}}}%
      }{%
        \edef\@relbl{\arabic{#4}}%
        \let\@retag\@empty
      }%
      \def\@tempshort{##3}%
      \ifx\@tempshort\@empty
        \expandafter\xdef\csname @sshort@#3\@relbl\endcsname{\@relbl}%
      \else
        \expandafter\def\csname @sshort@#3\@relbl\endcsname{##3}%
      \fi
      \begin{#2}[##2\ifshowkeys\@retag\fi]%
      \edef\@tmplbl{#3:\@relbl}\expandafter\label\expandafter{\@tmplbl}%
      \expanded{%
        \noexpand\newglossaryentry{#3\csname the#4\endcsname}{%
          name={\unexpanded{##2}%
            \IfBooleanF{##1}{\IfValueT{##4}{ [\noexpand\texttt{#3:##4}]}}},%
          sort={\csname the#2\endcsname},%
          description={p.~\noexpand\pageref{#3:\@relbl}},%
          user1={\csname the#2\endcsname},%
          user2={#5},%
          user3={\thesection},%
          user4={\expandafter\unexpanded\expandafter{\@savedSecTitle}}%
        }%
      }%
      \protected@write\@auxout{}{\string\glsliveentry{#3\csname the#4\endcsname}}%
      \protected@write\@auxout{}{\string\@thmtitle{\@tmplbl}{\unexpanded{##2}}}%
    \fi
  }%
  \expandafter\newcommand\csname ee#1\endcsname{\end{#2}}%
}
\@mkresult{L}{lemma}{l}{lemcounter}{Lemma}
% REMARK / COROLLARY / PROPOSITION -- same signature as \beL above
\@mkresult{R}{remark}{r}{remcounter}{Remark}
\@mkresult{C}{corollary}{c}{corcounter}{Corollary}
\@mkresult{P}{proposition}{p}{propcounter}{Proposition}
% OPEN-PROBLEM COMMAND -- generated with the same signature as \beT:
%   \beO[long title][short title][key]   ...   \eeO
% label o:KEY, or o:<counter> when no key is given; a bare \beO makes
% neither a label nor a glossary entry.
\@mkresult{O}{open}{o}{opcounter}{Problem}


% Section-level title tracker: updated only by \section, not \subsection.
% \@currentlabelname is overwritten by every \refstepcounter (including
% \subsection), so we save it separately here.
\newcommand{\@currentsectionname}{}
\AtBeginDocument{%
  \let\@epidcrn@orig@refstepcounter\refstepcounter
  \renewcommand{\refstepcounter}[1]{%
    \@epidcrn@orig@refstepcounter{#1}%
    \ifstrequal{#1}{section}{\global\let\@currentsectionname\@currentlabelname}{}%
  }%
}

% TABLE COMMANDS (glossaried) -- consistent naming with \beT, \beD, etc.
\newcommand{\beTa}[1][]{%
  \def\@temparg{#1}%
  \ifx\@temparg\@empty
    \begin{table}[htbp]%
  \else
    \let\@savedSecTitle\@currentsectionname%
    \refstepcounter{tabcounter}%
    \begin{table}[htbp]%
    \label{A:\thetabcounter}%
    \expanded{%
      \noexpand\newglossaryentry{A\thetabcounter}{%
        name={\unexpanded{#1}},%
        sort={\thetabcounter},%
        description={p.~\noexpand\pageref{A:\thetabcounter}},%
        user1={\thetabcounter},%
        user2={Table},%
        user3={\thesection},%
        user4={\expandafter\unexpanded\expandafter{\@savedSecTitle}}%
      }%
    }%
    \protected@write\@auxout{}{\string\glsliveentry{A\thetabcounter}}%
  \fi
}
\newcommand{\eeTa}{\end{table}}

% FIGURE COMMANDS (glossaried) -- consistent with \beT, \beR, \beTa, etc.
% Usage:
%   \beF[Glossary title]...\caption{Visible caption}\eeF  -> glossaried figure
%   \beF...\caption{Visible caption}\eeF                  -> plain [H] figure, no glossary
\newcommand{\beF}[1][]{%
  \def\@temparg{#1}%
  \ifx\@temparg\@empty
    \begin{figure}[H]%
  \else
    \begin{figure}[H]%
    \let\@savedSecTitle\@currentsectionname%
    \refstepcounter{figcounter}%
    \label{fig:\thefigcounter}%
    \expanded{%
      \noexpand\newglossaryentry{Fig\thefigcounter}{%
        name={\unexpanded{#1}},%
        sort={\thefigcounter},%
        description={p.~\noexpand\pageref{fig:\thefigcounter}},%
        user1={\thefigcounter},%
        user2={Figure},%
        user3={\thesection},%
        user4={\expandafter\unexpanded\expandafter{\@savedSecTitle}}%
      }%
    }%
    \protected@write\@auxout{}{\string\glsliveentry{Fig\thefigcounter}}%
  \fi
}
\newcommand{\eeF}{\end{figure}}
\makeatother
% LIVE-ENTRY REGISTER
% \glsaddall prints every entry the glossaries package currently knows, and
% that set includes entries restored from \jobname.glsdefs.  The package
% rewrites .glsdefs from all *currently defined* entries, so anything it
% restored is written straight back out: a deleted section kept reappearing in
% the catalogue for ever, and the .glsdefs had to be erased by hand.  The
% wrappers therefore record every entry they actually create in the .aux, and
% \printboth adds only those.  The .aux is rebuilt from scratch on each run, so
% this is self-healing: a removed result is gone after one further run.
\makeatletter
\newcommand{\@glslive}{}
\newcommand{\@glsaddlive}[1]{\ifglsentryexists{#1}{\glsadd{#1}}{}}
\newcommand{\glsliveentry}[1]{\g@addto@macro\@glslive{\@glsaddlive{#1}}}
\makeatother
% PRINT COMMAND
% No \clearpage here -- the catalogue continues on the current page.  This has
% been re-added by mistake several times; leave it out.
\makeatletter
\newcommand{\printboth}{%
  \@glslive   % add exactly the entries this document created last run
  \printglossary[style=unifiedstyle, title={0 \hspace{.22em} Definitions, Theorems, Examples and Open Problems}]%
}
\makeatother
% Named theorem reference: \tref{word} -> Thm~N(word-or-short)
% Shows stored short title if available, otherwise shows the key itself.
\newcommand{\tref}[1]{%
  Thm~\ref{t:#1}(\ifcsname @tshort@#1\endcsname\csname @tshort@#1\endcsname\else#1\fi)%
}

% Full-title theorem reference: \nref{t:KEY} prints the number and then the
% theorem's long title, so one writes "Theorem~\nref{t:blocks}" and gets
% "Theorem 3 ``Block factorization and siphons''".  Takes the complete label, like
% \ref, not the bare key like \tref.
%
% \beT writes each title to the .aux (see \@thmtitle below), so \nref resolves
% forward references from the second LaTeX pass, exactly as \ref does.  If the
% title is not yet known -- first pass, or a label that is not a \beT theorem
% -- it degrades to the number rather than producing ??.
% These two use @ in their names, and this point in the file is outside the
% \makeatletter region opened earlier (it closes at \makeatother above), so
% the region has to be reopened here.
\makeatletter
\newcommand{\@thmtitle}[2]{\global\@namedef{@tlong@#1}{#2}}
\newcommand{\nref}[1]{%
  \ref{#1}%
  \ifcsname @tlong@#1\endcsname
    \ ``\csname @tlong@#1\endcsname''%
  \fi
}
% \nm{t:KEY}: the stored long title of any \beT/\beL/\beP/\beR/\beC/\beD/\beE result,
% verbatim and without quotes, so item names follow the macro's first argument.
\newcommand{\nm}[1]{%
  \ifcsname @tlong@#1\endcsname
    \csname @tlong@#1\endcsname%
  \else
    \ref{#1}%
  \fi
}
\makeatother




%% ========================================
%% Other THEOREM ENVIRONMENTS
%% ========================================
\makeatletter
\@ifundefined{lemma}      {\newtheorem{lemma}      {Lemma}}      {}
\@ifundefined{proposition}{\newtheorem{proposition}{Proposition}} {}
\@ifundefined{corollary}  {\newtheorem{corollary}  {Corollary}}  {}
\@ifundefined{conjecture} {\newtheorem{conjecture} {Conjecture}}  {}
\@ifundefined{remark}     {\newtheorem{remark}     {Remark}}     {}
\@ifundefined{question}   {\newtheorem{question}   {Question}}   {}
\@ifundefined{assumption} {\newtheorem{assumption} {Assumption}} {}
\@ifundefined{hypothesis} {\newtheorem{hypothesis} {Hypothesis}} {}
\@ifundefined{fact}       {\newtheorem{fact}       {Fact}}       {}
\makeatother

% \citer: cite the relay-graph threshold theorem of AAHrel. If a document has
% \input the chapter/file containing t:relayRF (so the label resolves), point
% straight at it; otherwise fall back to a bare citation of AAHrel. Shared here
% (not per-paper) so it resolves consistently across every book chapter.
% \citeR: same idea, for the first theorem of AAHRah (t:Pstar).
\makeatletter
\providecommand{\citer}{%
  \@ifundefined{r@t:relayRF}%
    {\cite{AAHrel}}%
    {Theorem~\ref{t:relayRF} of \cite{AAHrel}}%
}
\providecommand{\citeR}{%
  \@ifundefined{r@t:Pstar}%
    {\cite{AAHRah}}%
    {Theorem~\ref{t:Pstar} of \cite{AAHRah}}%
}
\makeatother


% Abbreviations
\def\beQ{\begin{question}}\def\eeQ{\end{question}}
\def\beA{\begin{assumption}}\def\eeA{\end{assumption}}


%% ========================================

%% ========================================
%% COLORS & CUSTOM COMMANDS
%% ========================================
\definecolor{funccolor}{RGB}{25,25,112}
\definecolor{desccolor}{RGB}{64,64,64}
\newcommand{\func}[1]{\texttt{\textcolor{funccolor}{#1}}}
\def\eq{=}
\def\mm{-}

%% ========================================
%% MATHCAL SHORTHANDS
%% ========================================
% Project-wide convention: \mX for a capital X is always \mathcal X.
% Defined for the whole alphabet so the convention never fails to expand.
\newcommand{\mA}{\mathcal A}\newcommand{\mB}{\mathcal B}
\newcommand{\mC}{\mathcal C}\newcommand{\mD}{\mathcal D}
\newcommand{\mE}{\mathcal E}\newcommand{\mF}{\mathcal F}
\newcommand{\mG}{\mathcal G}\newcommand{\mH}{\mathcal H}
\newcommand{\mI}{\mathcal I}\newcommand{\mJ}{\mathcal J}
\newcommand{\mK}{\mathcal K}\newcommand{\mL}{\mathcal L}
\newcommand{\mM}{\mathcal M}\newcommand{\mN}{\mathcal N}
\newcommand{\mO}{\mathcal O}\newcommand{\mP}{\mathcal P}
\newcommand{\mQ}{\mathcal Q}\newcommand{\mR}{\mathcal R}
\newcommand{\mS}{\mathcal S}\newcommand{\mT}{\mathcal T}
\newcommand{\mU}{\mathcal U}\newcommand{\mV}{\mathcal V}
\newcommand{\mW}{\mathcal W}\newcommand{\mX}{\mathcal X}
\newcommand{\mY}{\mathcal Y}\newcommand{\mZ}{\mathcal Z}

%% ========================================
%% BASIC MATH ENVIRONMENTS
%% ========================================
\def\bep{\begin{pmatrix}}\def\eep{\end{pmatrix}}
\def\bev{\begin{vmatrix}}\def\eev{\end{vmatrix}}
\def\bea{\begin{eqnarray*}}\def\eea{\end{eqnarray*}}
\def\bc{\begin{cases}}\def\ec{\end{cases}}
\def\BEN{\begin{enumerate}}\def\EEN{\end{enumerate}}
\def\BI{\begin{itemize}}\def\EI{\end{itemize}}

\newcommand{\be}[1]{\begin{equation}\label{#1}}
\newcommand{\ee}{\end{equation}}
\newcommand{\beq}{\begin{eqnarray}}
\def\eeq{\end{eqnarray}}
\def\eqr{\eqref}\def\fr{\frac}\def\lbl{\label}
\def\Lra{\Longrightarrow}\def\non{\nonumber}\def\cNGM{\cite{Diek,Van,Van08}}
\def\Eq{\Leftrightarrow}\def\b{\beta}\def\uph{\red{UP TO HERE}}
%% ========================================
%% PROBABILITY & STATISTICS
%% ========================================
\renewcommand{\P}{\ensuremath{\mathbb{P}}}  % Note: overrides paragraph symbol
\newcommand{\E}{\ensuremath{\mathbb{E}}}
\newcommand{\ind}[1]{\ensuremath{\mathbbm{1}_{\left\{#1\right\}}}}
\DeclareMathOperator{\Cov}{Cov}
\DeclareMathOperator{\TV}{TV}
\newcommand{\rank}{\operatorname{rank}}

%% ========================================
%% CRN OPERATORS (organizations/ERCs/buffering/siphons framework)
%% ========================================
\def\clos{\mathrm{clos}}
\def\Supply{\mathrm{Supply}}
\def\req{\mathrm{req}}
\def\React{\mathrm{React}}
\def\Prod{\mathrm{Prod}}
\def\Gen{\mathrm{Gen}}
\def\In{\mathrm{In}}
\def\Out{\mathrm{Out}}
\def\Act{\mathrm{Act}}
\def\Min{\mathrm{Min}}

%% ========================================
%% LINE BREAKING
%% ========================================
% Last-resort stretch, used only on paragraphs TeX cannot otherwise break:
% a long \texttt{} token, a comma-separated inline math list (commas are not
% breakpoints in math), or a run of \cite/\ref boxes.  Without it those set as
% overfull \hbox and print into the margin; with it they set slightly loose,
% which is the lesser evil.  Does not affect paragraphs that already break.
\emergencystretch=3em

%% ========================================
%% NUMBER SETS
%% ========================================
\newcommand{\R}{\mathbb{R}}  % Real numbers
\newcommand{\C}{\mathbb{C}}  % Complex numbers
\newcommand{\N}{\mathbb{N}}
\newcommand{\Z}{\mathbb{Z}}
% \providecommand, not \newcommand: some papers define these in their own
% preamble after \input{preamble}, and must not clash.
\providecommand{\Q}{\mathbb{Q}}
\providecommand{\res}{\operatorname{res}}   % resultant; used by epiRes.tex
\newcommand{\MN}{\mathbb{N}}
\newcommand{\MZ}{\mathbb{Z}}
\newcommand{\realpos}{\mathbb{R}_{> 0}}
\newcommand{\realnneg}{\mathbb{R}_{\geq 0}}
\newcommand{\naturalpos}{\mathbb{N}_{>0}}
\newcommand{\Rplus}{\mathbb{R}_+}
\newcommand{\Rp}{\mathbb{R}_{\ge 0}}
\def\Rnne{\mathbb{R}_{\ge 0}}
\def\Se{\bar{S}} \def\Ie{\bar{I}}
\newcommand{\Rpos}{\mathbb{R}_{>0}}
\newcommand{\dd}{\,\mathrm d}
\newcommand{\e}{\mathrm e}
%% ========================================
%% GREEK LETTERS (2+ character abbreviations) greek
%% ========================================
\def\al{\alpha}\def\ba{\bff \alpha}
\def\w{\omega}\def\Lgn{\ga+\mu+\de}\def\om{\omega}
\def\de{\delta}\def\De{\Delta}\def\bde{\bff \delta}
\def\ep{\epsilon}\def\eps{\varepsilon}
\def\f{\varphi}  % Kept as requested
\def\ga{\gamma}\def\Ga{\Gamma}\def\G{\Gamma}\def\bg{\bff \gamma}
\def\up{\upsilon}  % NOT the breve accent: no accent commands anywhere.
                   % Named \up, not \u, so the LaTeX breve stays reachable.
\def\la{\lambda}\def\La{\Lambda}
\def\si{\sigma}\def\ka{\kappa}
\def\th{\theta}\def\vt{\vartheta}
\def\ze{\zeta}\def\Gap{\Gamma^{+}} \def\Gam{\Gamma^{-}}


%% ========================================
%% SPECIAL SHORT COMMANDS (kept as requested)
%% ========================================
\def\I{\infty}
\def\T{\widetilde}
\def\H{\hat}
\newcommand{\relint}{\operatorname{relint}}

%% ========================================
%% BOLD/VECTOR NOTATION
%% ========================================
\newcommand{\bff}[1]{{\mbox{\boldmath$#1$}}}
\newcommand\X{\boldsymbol{X}}
\def\x{\boldsymbol{x}}\def\bx{\bff x}\def\xe{\bar{\x}}
\newcommand\y{\boldsymbol{y}}\def\by{\boldsymbol{y}}\def\ye{\bar{\y}}
\def\yb{\y}\def\xb{\x}
\def\xy{\x-\y}
\newcommand\br{\boldsymbol{r}}\def\re{\bar{\r}}
\newcommand\bd{\boldsymbol{d}}
\newcommand\bQ{\boldsymbol{Q}}
\newcommand{\Yb}{{\mathbf{Y}}}
\newcommand{\Ab}{\mathbf{A}}
\newcommand{\zb}{\mathbf{z}}\def\bz{\bff 0}
\newcommand{\Zb}{\mathbf{Z}}
\newcommand{\bb}{\mathbf{b}}
\def\vbe{\vec {\bff \beta}}
\def\vb{\vec {\bff b}}
\def\v1{\vec {\bff 1}}
\def\vn{\vec \de}
\def\vz{\vec 0}
\def\vi{\overset{\rightarrow}{\mathsf i}}
\def\vv{\overset{\rightarrow}{v}}
\def\vr{\; \vec {\mathsf r}}
\newcommand{\va}{\bff a}
\providecommand{\vect}[1]{\boldsymbol{#1}}   % bold vector (moved from Gradostat.tex)
\providecommand{\one}{\mathbf 1}             % all-ones vector (moved from AliSF)
\providecommand{\spr}{s}                     % spectral bound  (moved from AliSF)
\providecommand{\Ke}{K_{\rm eff}}            % effective carrying capacity
% Critical right/left eigenvectors in the Hopf projection method (AAHOSN/New.tex)
\providecommand{\bq}{\mathbf{q}}
\providecommand{\bp}{\mathbf{p}}


%% ========================================
%% MATHCAL & MATHSCR
%% ========================================
% Mathcal
\newcommand\mc[1]{\mathcal{#1}}
\def\mC{{\mathcal C}}\def\mD{\mathcal D}\def\mE{\mathcal{E}}
\def\mF{\mathcal F}\def\mG{{\mathcal G}}\def\mH{{\mathcal H}}
\def\mI{\mathcal I}\def\mL{\mathcal{L}}\def\mM{{\mathcal M}}
\def\mN{{\mathcal N}}\def\mP{\mathcal{P}}\def\mR{{\mathcal R}}
\def\mS{{\mathcal S}}\def\mT{\mathcal{T}}\def\mV{{\mathcal V}}
\def\mK{\mathcal K}\def\mA{{\mathcal A}}\def\mU{{\mathcal U}}
% Symmetric bilinear and trilinear forms of the Taylor expansion, used in the
% first Lyapunov coefficient (AAHOSN/New.tex)
\providecommand{\Bform}{\mathcal{B}}
\providecommand{\Cform}{\mathcal{C}}

% Blackboard bold.  Kept here, not in the drivers: a \providecommand in a
% driver preamble runs before this file and silently wins, so the same name
% ends up meaning different things in different builds.
\def\IC{\mathbb{C}}\def\IN{\mathbb{N}}\def\IQ{\mathbb{Q}}
\def\IR{\mathbb{R}}\def\IZ{\mathbb{Z}}\def\IG{\mathbb{G}}
\def\II{\mathbb{I}}

% Mathscr
\newcommand{\Rs}{{\mathscr R}}
\newcommand{\Ss}{{\mathscr S}}
\newcommand{\Ns}{{\mathscr N}}
\newcommand{\Cs}{\mathscr C}


%% ========================================
%% OPERATORS & FUNCTIONS
%% ========================================
\DeclareMathOperator{\tr}{tr}
\DeclareMathOperator{\diag}{Diag}
\DeclareMathOperator{\Diag}{Diag}
\DeclareMathOperator{\intr}{\mathrm{Int}}
\DeclareMathOperator{\Inv}{Inv}
\DeclareMathOperator{\spec}{spec}
\DeclareMathOperator{\AUC}{AUC}               % area under the curve
\newcommand{\ii}{\mathrm{i}}                  % imaginary unit
\newcommand\op{\operatorname}
\newcommand{\supp}{\op{supp}}
\newcommand{\diff}{\mathop{}\mathopen{}\mathrm{d}}
\newcommand{\pdd}[1]{\frac{\partial}{\partial{#1}}}
\DeclareMathOperator{\cone}{cone}
\DeclareMathOperator{\proj}{proj}

%% ========================================
%% BRACKETS & NORMS
%% ========================================
\providecommand{\pr}[1]{\left(#1\right)}
\providecommand{\pp}[1]{\left[#1\right]}
\newcommand\croc[1]{\left\langle #1\right\rangle}
\newcommand{\norm}[1]{\|#1\|}
\newcommand{\inner}[2]{\langle #1, #2 \rangle}
% Short synonym of \inner, used by AAHOSN/New.tex as $\ip{u}{v}=\bar u^\top v$
\providecommand{\ip}[2]{\langle #1 , #2 \rangle}


%% ========================================
%% FORMATTING & UTILITIES
%% ========================================
\newcommand*{\Scale}[2][4]{\scalebox{#1}{$#2$}}
\newcommand{\te}[1]{\text{#1}}
\newcommand\steq[1]{\stackrel{\text{\rm #1.}}{=}}
\newcommand\blue[1]{\textcolor{blue}{#1}}
\newcommand{\red}{\textcolor[rgb]{1.00,0.00,0.00}}
\newcommand{\green}{\textcolor[rgb]{0.00,1.00,0.00}}

% Figure command
\newcommand{\figu}[3]{%
\begin{figure}[H]
\centering
\includegraphics[scale=#3]{#1}
\caption{#2\label{f:#1}}
\end{figure}
}

% Footnote
\long\def\symbolfootnote[#1]#2{%
\begingroup
\def\thefootnote{\fnsymbol{footnote}}\footnote[#1]{#2}%
\endgroup}
\def\fn{\symbolfootnote}

\newcommand{\dsum}{\displaystyle\sum}

%% ========================================
%% CRN & EPIDEMIOLOGY SPECIFIC
%% ========================================
% CRN notation
\def\SS{\mathcal S}
\def\CC{\mathcal C}
\def\RR{\mathcal R}

% Reproduction numbers
\def\Rto{R_{1}^{E_2}}
\def\Rot{R_{2}^{E_1}}
\def\Roc{R_{1,c}}
\def\Rtc{R_{2,c}}
\def\rd{r_{0}}
\def\sd{s_{0}}
\def\dd{D_{0}}
\def\vd{v_{0}}
\def\zd{z_{0}}
\def\se{\sE}
\def\xd{x_{0}}
%% ========================================
%% SPECIAL VALUES & EQUILIBRIA
%% ========================================
\def\eE{\bar{e}}
\def\wE{\bar{w}}
\def\sE{s_{*}}
\def\Me{\T M_{ee}}
\def\Rnne{\mathbb{R}_{\ge 0}}

\newcommand{\ones}{\mathbf 1}


%% ========================================
%% ABBREVIATIONS (alphabetical by category)
%% ========================================

% -- General mathematical terms --
\def\adm{admissible}\def\admy{admissibility}
%\def\and{antisymmetric}% disabled 2026-09-30: it overrode LaTeX's author separator \and (printed "antisymmetric" between authors)
\def\ass{assumption}
\def\cadlag{c\`adl\`ag}
\def\ch{characteristic polynomial}
\def\chp{characteristic polynomial}
\def\coe{coefficient}
\def\comp{computation}
\def\con{condition}
\def\coo{coordinate}
\def\corr{corresponding}
\def\dec{decomposition}
\def\den{denominator}
\def\dete{deterministic}
\def\eig{eigenvalue}\def\eval{eigenvalue}
\def\eigv{eigenvector}\def\eiv{eigenvector}\def\evec{eigenvector}
\def\elt{element}
\def\enn{essentially nonnegative/positivity preserving}
\def\equ{equation}
\def\etal{et al.}
\def\fe{for example}
\def\fno{from now on}
\def\for{\forall}
\def\form{formula}
\def\Fr{Furthermore, }
\def\frt{furthermore}
\def\fun{function}
\def\gene{generalization}
\def\heu{heuristic}
\def\how{however}\def\How{However, }
\def\ie{i.e. }
\def\im{\item}
\def\isd{is defined by}
\def\itc{it is easy to check that}
\def\itf{it follows that}\def\Itf{It follows that}
\def\ith{it holds that}
\def\itm{it may be checked that}\def\Itm{It may be checked that}
\def\ito{it turns out that}\def\Ito{It turns out that}\def\Itou{It turns out useful to}
\def\ma{mass action}\def\MAR{mass action representation}
\def\mat{matrix}
\def\mbc{may be checked to be}
\def\mbw{may be written as}\def\mubw{must be written as}
\def\mpr{more precisely, }\def\Mp{More precisely, }
\def\nc{non-catalytic}
\def\nne{non-negative}\def\nny{nonnegativity}
\def\no{\nonumber}
\def\noni{non-infectious}
\def\nota{notation}
\def\np{\newpage}
\def\Nt{Note that}
\def\nt{non trivial}\def\nz{non zero}
\def\oth{otherwise}
\def\Oth{On the other hand}
\def\pa{parameter}\def\para{parameter}
\def\pol{polynomial}
\def\prob{probability}
\def\Prf{{\bf  Proof:}}\def\prf{{\bf  Proof:}}
\def\prop{proportional to}
\def\qu{\quad}
\def\QED{\hfill {$\square$}\goodbreak \medskip}
\def\rep{representation}
\def\resp{respectively}
\def\satd{satisfied}\def\satg{satisfying}\def\sats{satisfies}\def\saty{satisfy}
\def\ssec{\subsection}\def\sssec{\subsubsection}
\def\sef{see for example}
\def\sksy{skew-symmetric}
\def\sub{substrate}
\def\sys{system}
\def\var{variable}
\def\wk{well-known}
\def\wl{without loss of generality}\def\wlo{w.l.o.g.}
\def\wmc{We may check that}\def\Wmc{We may check that}
\def\wmr{we may rewrite}
\def\wms{we must show that}
\def\wrt{with respect to}

% -- CRN & Chemistry --
\newcommand\CRN{chemical reaction network}
\newcommand\CRNs{chemical reaction networks}
\newcommand\RN{reaction network}
\def\ACR{absolute concentration robustness}
\def\AMSD{acyclic siphon decomposition}
\def\DAE{differential algebraic equations}
\def\FHJ{Feinberg-Horn-Jackson}
\def\GMAK{generalized mass action kinetics}
\def\MM{Michaelis-Menten}
\def\sm{stoichiometric matrix}\def\Sm{stoichiometric matrix}
\def\WRZD{weakly reversible, zero deficiency}\def\aut{autocatalytic}

% -- Epidemiology --
\def\brn{basic reproduction number}\def\bzn{basic replacement number}
\def\com{compartment}
\def\dfe{disease free equilibrium}\def\DFE{disease free equilibrium}
\def\dwt{dwell times}\def\edt{expected dwell times}
\def\EM{epidemic model}
\def\endp{endemic fixed point}\def\enp{endemic point}\def\EE{endemic point}
\def\ME{mathematical epidemiology}
\def\MEE{eco-epidemiology}\def\EEP{eco-epidemiological}
\def\ece{eco-epidemiological model}
\def\mgm{multi group model}
\def\mtS{multi-strain model}\def\mSt{multi-strain model}
\def\mSi{minimal siphon}
\def\NGM{next generation matrix}
\def\nim{new infections matrix}
\def\qem{quadratic epidemiological model}
\def\repF{invasion functions}
\def\SPF{SIR-PH-FA model}
\def\sus{susceptible}

% -- Graph theory & Linear algebra --
\def\jac{Jacobi matrix}
\def\jin{Jacobian matrix of the invasion vector field with respect to the invasion variables}
\def\jrn{Jacobian matrix of the resident vector field with respect to the resident variables}
\def\jxn{Jacobian matrix of the $\x$ vector field with respect to the $\x$ variables}
\def\Lm{Laplacian matrix}
\def\MTT{matrix tree theorem}
\def\SGT{spectral graph theory}
\def\SMW{Sherman–Morrison–Woodbury}

% -- Stability & Dynamics --
\def\absC{absorbtion CTMC}
\def\BTb{Bogdanov-Takens bifurcation}
\def\CEP{competitive exclusion partition}
\def\CTMC{continuous time Markov chain}
\def\ds{differential system}
\def\fp{fixed point}\def\bfp{boundary fixed point}\def\fbp{fixed boundary point}
\def\fbl{fixed boundary line}\def\fif{forward invariant face}
\def\gas{global asymptotic stability}\def\GAS{globally asymptotically stable}
\def\GASS{\text{GASS}}
\def\Hb{Hopf bifurcation}
\def\Hur{Hurwitz}\def\MH{Metzler-Hurwitz}\def\RH{Routh-Hurwitz}
\def\IS{\text{IS}}
\def\LAS{locally asymptotically stable}
\def\LdS{Lyapunov diagonally stable}
\def\Lf{Lyapunov function}\def\LyV{Lyapunov Volterra type function}\def\RLF{robust Lyapunov function}
\def\LyM{Lyapunov Malkin stable}
\def\PF{Perron-Frobenius}
\def\spfp{strictly positive fixed point}
\def\sta{stability}
\def\STP{sharp threshold property}\def\STT{sharp threshold theorem}

% -- Models (with \varphi explicitly instead of \f) --
\def\AB{$(A,B)$ Arino-Brauer epidemic models}
\def\ABa{$(A,\bbe,\ba,\varphi)$ balanced bilinear model}
\def\ABP{%$(A,\bbe,P,\varphi)$
\\ balanced bilinear model}
\def\AMG{Anderson May Gupta}
\newcommand\BIN{biological interaction network}
\def\Dt{Descartes type}
\def\gLV{generalized Lotka-Volterra model}
\def\IG{\text{invasion graph}}\def\IGMS{interaction graph on minimal siphons}
\def\LCP{\text{linear complementarity problem}}
\def\LV{Lotka-Volterra}\def\LVK{Lotka-Volterra-Kolmogorov}

% -- Splitting & Structure --
\def\reg{regular splitting}\def\regS{{\bf regular splitting}}
\def\RUR{rational univariate representation}
\def\bds{block-diagonal structure}\def\NGMs{next generation matrices}

% -- Miscellaneous --
\def\Ma{{Mathematica}}
\def\Mm{Max–Min}
\def\num{numerator}
\def\rap{rational approximation} \def\DSR{directed species-reaction bipartite graph}\def\PCP{polynomial complementarity problem}\def\rat{rational substitution}
\def\QSS{quasi steady state} \def\RHS{right hand side} \def\LHS{left hand side}
\def\OSN{online social network}\def\VdD{Van den Driessche}
\def\wli{$\omega$-limit} \def\rk{rich kinetics}\def\sJ{symbolic Jacobian}
\def\cV{\cite{Vas,VasSta,BSV,GSV}} \def\spf{semi-parametric semi-functional}
\def\CRW{Capasso-Ruan-Wang} \def\ABP{$(A,\bbe,P,\f)$ balanced bilinear model}
\def\ABa{%$(A,\bbe,\ba,\f)$ \\
rank-one balanced bilinear model}
\def\bbe{B%{\bff \beta}
} \def\vbe{\vec {\bff \beta}
}
\def\vb{\vec {\bff b}}
\newcommand{\Kh}{\hat{K}}          % K with hat, first mention only
\newcommand{\wS}{w^{\sigma}}       % right Perron J^\perp, full block
\newcommand{\piS}{\vec \pi^{\sigma}}    % left Perron J^\perp, full block
\newcommand{\ws}{w^{\sigma}}       % right Perron J^\perp, sub-block
\newcommand{\pis}{\pi^{\sigma}}    % left Perron J^\perp, sub-block
\newcommand{\hwS}{\hat{w}^{\sigma}}   % right Perron K, full block
\newcommand{\hpiS}{\hat{\pi}^{\sigma}} % left Perron K, full block
\newcommand{\hws}{\hat{w}^{\sigma}}   % right Perron K, sub-block
\newcommand{\hpis}{\hat{\pi}^{\sigma}} % left Perron K, sub-block
\newcommand{\twS}{\tilde{w}^{\sigma}}  % right Perron \tilde K, full block
\newcommand{\tpiS}{\tilde{\pi}^{\sigma}} % left Perron \tilde K
\newcommand{\tws}{\tilde{w}^{\sigma}}
\newcommand{\tpis}{\tilde{\pi}^{\sigma}}
\newcommand{\Mat}[2]{\mathbb{R}_{\ge 0}^{\,#1\times #2}}
\def\tJ{J_\sigma^\perp} \def\lPe{left Perron eigenvector}
% \G is \Gamma (defined at :795).  The Volterra function keeps its own name,
% so that \G no longer has two conflicting meanings in this file.
\def\GV{G}  % Volterra G(u)=u-1-\ln u
\def\PVL{Perron-Volterra Lyapunov minimal siphon ansatz}
\newcommand{\istar}[1]{i_{#1}^{*}}
%\newcommand{\sE}{s_E}
\newcommand{\vE}{v_E}
\def\PV{Perron-Volterra candidate Lyapunov function}
\def\GLV{generalized Lotka-Volterra}
\def\MERN{mathematical epidemiology reaction networks}
\def\CMM{contraction metric method}
\newcommand{\sign}{\operatorname{sign}}
\newcommand{\Var}{\operatorname{Var}} 

% ======================================================================
% IDEAS SHARED BETWEEN THE ABSTRACT AND THE MAIN RESULTS SECTION
%
% Write each idea ONCE, in the abstract; the Main Results section recalls it,
% so the two can never drift apart.
%
%   in the abstract        \ideaitem{1}{The first-order reaction terms ...}
%   in Main Results        \ideahead{1}
%
% \ideaitem stores the text globally and typesets
%   \item \hyperlink{idea:1}{<text>}
% \ideahead typesets
%   \item\hypertarget{idea:1}{}\textbf{<the same text>}
%
% The abstract must come before the Main Results input, which is the normal
% order. If a heading has to read differently from the abstract sentence (for
% instance a long abstract item summarised in one line), override it with the
% optional argument: \ideahead[short heading]{6}.
% ======================================================================
\makeatletter
\newcommand{\ideaitem}[2]{%
  \expandafter\gdef\csname @idea@#1\endcsname{#2}%
  \item\hyperlink{idea:#1}{#2}}
\newcommand{\ideahead}[2][]{%
  \item\hypertarget{idea:#2}{}%
  \def\@ideatmp{#1}%
  \ifx\@ideatmp\@empty
    \@ifundefined{@idea@#2}{%
      \PackageWarning{preamble}{Idea #2 is used in Main Results but was never
        defined by \string\ideaitem\space in the abstract}%
      \textbf{[idea #2 undefined]}%
    }{\textbf{\begingroup\def\ideaterm##1##2{##2}\def\idearest##1{}\@nameuse{@idea@#2}\endgroup}}%
  \else
    \textbf{#1}%
  \fi}
% \ideaterm{<hypertarget>}{<words>}: inside an abstract idea the words become an
% invisible link (the document sets \hypersetup{hidelinks}) jumping to the place
% in Main Results that defines them; inside the Main Results heading itself,
% where the target sits a few lines below, \ideahead strips the link and prints
% the words plainly. Mark the destination with \hypertarget{<name>}{} on the line
% before the display, which adds no equation number and shows nothing.
\newcommand{\ideaterm}[2]{\hyperlink{#1}{#2}}
% \idearest{<text>}: the part of an abstract idea after its first sentence; printed in
% the abstract, omitted from the Main Results heading.
\newcommand{\idearest}[1]{#1}
\makeatother

% ======================================================================
% SHORT-PAPER SWITCH
% \shorttrue in a master that assembles only part of the companion paper
% (Barrat/SipArxiv.tex). \nsh{...} drops text that cites results living in the
% omitted sections, so the short build has no dangling references while the
% companion paper keeps the text verbatim.
% ======================================================================
\newif\ifshort \shortfalse
\newcommand{\nsh}[1]{\ifshort\else#1\fi}
% \shortalt{<short build>}{<companion paper>}: use when a sentence needs
% different wording in the two builds, not merely a deletion.
\newcommand{\shortalt}[2]{\ifshort#1\else#2\fi}
% \Jtr, \JtrS{\Sigma}: transverse Jacobian. J^\perp here; SipArxiv.tex redefines them as J_{\mathrm{tra}}, J_{\mathrm{tra},\Sigma}
\newcommand{\Jtr}{J^\perp}
\newcommand{\JtrS}[1]{J_{#1}^\perp}
, and must not clash.
\providecommand{\Q}{\mathbb{Q}}
\providecommand{\res}{\operatorname{res}}   % resultant; used by epiRes.tex
\newcommand{\MN}{\mathbb{N}}
\newcommand{\MZ}{\mathbb{Z}}
\newcommand{\realpos}{\mathbb{R}_{> 0}}
\newcommand{\realnneg}{\mathbb{R}_{\geq 0}}
\newcommand{\naturalpos}{\mathbb{N}_{>0}}
\newcommand{\Rplus}{\mathbb{R}_+}
\newcommand{\Rp}{\mathbb{R}_{\ge 0}}
\def\Rnne{\mathbb{R}_{\ge 0}}
\def\Se{\bar{S}} \def\Ie{\bar{I}}
\newcommand{\Rpos}{\mathbb{R}_{>0}}
\newcommand{\dd}{\,\mathrm d}
\newcommand{\e}{\mathrm e}
%% ========================================
%% GREEK LETTERS (2+ character abbreviations) greek
%% ========================================
\def\al{\alpha}\def\ba{\bff \alpha}
\def\w{\omega}\def\Lgn{\ga+\mu+\de}\def\om{\omega}
\def\de{\delta}\def\De{\Delta}\def\bde{\bff \delta}
\def\ep{\epsilon}\def\eps{\varepsilon}
\def\f{\varphi}  % Kept as requested
\def\ga{\gamma}\def\Ga{\Gamma}\def\G{\Gamma}\def\bg{\bff \gamma}
\def\up{\upsilon}  % NOT the breve accent: no accent commands anywhere.
                   % Named \up, not \u, so the LaTeX breve stays reachable.
\def\la{\lambda}\def\La{\Lambda}
\def\si{\sigma}\def\ka{\kappa}
\def\th{\theta}\def\vt{\vartheta}
\def\ze{\zeta}\def\Gap{\Gamma^{+}} \def\Gam{\Gamma^{-}}


%% ========================================
%% SPECIAL SHORT COMMANDS (kept as requested)
%% ========================================
\def\I{\infty}
\def\T{\widetilde}
\def\H{\hat}
\newcommand{\relint}{\operatorname{relint}}

%% ========================================
%% BOLD/VECTOR NOTATION
%% ========================================
\newcommand{\bff}[1]{{\mbox{\boldmath$#1$}}}
\newcommand\X{\boldsymbol{X}}
\def\x{\boldsymbol{x}}\def\bx{\bff x}\def\xe{\bar{\x}}
\newcommand\y{\boldsymbol{y}}\def\by{\boldsymbol{y}}\def\ye{\bar{\y}}
\def\yb{\y}\def\xb{\x}
\def\xy{\x-\y}
\newcommand\br{\boldsymbol{r}}\def\re{\bar{\r}}
\newcommand\bd{\boldsymbol{d}}
\newcommand\bQ{\boldsymbol{Q}}
\newcommand{\Yb}{{\mathbf{Y}}}
\newcommand{\Ab}{\mathbf{A}}
\newcommand{\zb}{\mathbf{z}}\def\bz{\bff 0}
\newcommand{\Zb}{\mathbf{Z}}
\newcommand{\bb}{\mathbf{b}}
\def\vbe{\vec {\bff \beta}}
\def\vb{\vec {\bff b}}
\def\v1{\vec {\bff 1}}
\def\vn{\vec \de}
\def\vz{\vec 0}
\def\vi{\overset{\rightarrow}{\mathsf i}}
\def\vv{\overset{\rightarrow}{v}}
\def\vr{\; \vec {\mathsf r}}
\newcommand{\va}{\bff a}
\providecommand{\vect}[1]{\boldsymbol{#1}}   % bold vector (moved from Gradostat.tex)
\providecommand{\one}{\mathbf 1}             % all-ones vector (moved from AliSF)
\providecommand{\spr}{s}                     % spectral bound  (moved from AliSF)
\providecommand{\Ke}{K_{\rm eff}}            % effective carrying capacity
% Critical right/left eigenvectors in the Hopf projection method (AAHOSN/New.tex)
\providecommand{\bq}{\mathbf{q}}
\providecommand{\bp}{\mathbf{p}}


%% ========================================
%% MATHCAL & MATHSCR
%% ========================================
% Mathcal
\newcommand\mc[1]{\mathcal{#1}}
\def\mC{{\mathcal C}}\def\mD{\mathcal D}\def\mE{\mathcal{E}}
\def\mF{\mathcal F}\def\mG{{\mathcal G}}\def\mH{{\mathcal H}}
\def\mI{\mathcal I}\def\mL{\mathcal{L}}\def\mM{{\mathcal M}}
\def\mN{{\mathcal N}}\def\mP{\mathcal{P}}\def\mR{{\mathcal R}}
\def\mS{{\mathcal S}}\def\mT{\mathcal{T}}\def\mV{{\mathcal V}}
\def\mK{\mathcal K}\def\mA{{\mathcal A}}\def\mU{{\mathcal U}}
% Symmetric bilinear and trilinear forms of the Taylor expansion, used in the
% first Lyapunov coefficient (AAHOSN/New.tex)
\providecommand{\Bform}{\mathcal{B}}
\providecommand{\Cform}{\mathcal{C}}

% Blackboard bold.  Kept here, not in the drivers: a \providecommand in a
% driver preamble runs before this file and silently wins, so the same name
% ends up meaning different things in different builds.
\def\IC{\mathbb{C}}\def\IN{\mathbb{N}}\def\IQ{\mathbb{Q}}
\def\IR{\mathbb{R}}\def\IZ{\mathbb{Z}}\def\IG{\mathbb{G}}
\def\II{\mathbb{I}}

% Mathscr
\newcommand{\Rs}{{\mathscr R}}
\newcommand{\Ss}{{\mathscr S}}
\newcommand{\Ns}{{\mathscr N}}
\newcommand{\Cs}{\mathscr C}


%% ========================================
%% OPERATORS & FUNCTIONS
%% ========================================
\DeclareMathOperator{\tr}{tr}
\DeclareMathOperator{\diag}{Diag}
\DeclareMathOperator{\Diag}{Diag}
\DeclareMathOperator{\intr}{\mathrm{Int}}
\DeclareMathOperator{\Inv}{Inv}
\DeclareMathOperator{\spec}{spec}
\DeclareMathOperator{\AUC}{AUC}               % area under the curve
\newcommand{\ii}{\mathrm{i}}                  % imaginary unit
\newcommand\op{\operatorname}
\newcommand{\supp}{\op{supp}}
\newcommand{\diff}{\mathop{}\mathopen{}\mathrm{d}}
\newcommand{\pdd}[1]{\frac{\partial}{\partial{#1}}}
\DeclareMathOperator{\cone}{cone}
\DeclareMathOperator{\proj}{proj}

%% ========================================
%% BRACKETS & NORMS
%% ========================================
\providecommand{\pr}[1]{\left(#1\right)}
\providecommand{\pp}[1]{\left[#1\right]}
\newcommand\croc[1]{\left\langle #1\right\rangle}
\newcommand{\norm}[1]{\|#1\|}
\newcommand{\inner}[2]{\langle #1, #2 \rangle}
% Short synonym of \inner, used by AAHOSN/New.tex as $\ip{u}{v}=\bar u^\top v$
\providecommand{\ip}[2]{\langle #1 , #2 \rangle}


%% ========================================
%% FORMATTING & UTILITIES
%% ========================================
\newcommand*{\Scale}[2][4]{\scalebox{#1}{$#2$}}
\newcommand{\te}[1]{\text{#1}}
\newcommand\steq[1]{\stackrel{\text{\rm #1.}}{=}}
\newcommand\blue[1]{\textcolor{blue}{#1}}
\newcommand{\red}{\textcolor[rgb]{1.00,0.00,0.00}}
\newcommand{\green}{\textcolor[rgb]{0.00,1.00,0.00}}

% Figure command
\newcommand{\figu}[3]{%
\begin{figure}[H]
\centering
\includegraphics[scale=#3]{#1}
\caption{#2\label{f:#1}}
\end{figure}
}

% Footnote
\long\def\symbolfootnote[#1]#2{%
\begingroup
\def\thefootnote{\fnsymbol{footnote}}\footnote[#1]{#2}%
\endgroup}
\def\fn{\symbolfootnote}

\newcommand{\dsum}{\displaystyle\sum}

%% ========================================
%% CRN & EPIDEMIOLOGY SPECIFIC
%% ========================================
% CRN notation
\def\SS{\mathcal S}
\def\CC{\mathcal C}
\def\RR{\mathcal R}

% Reproduction numbers
\def\Rto{R_{1}^{E_2}}
\def\Rot{R_{2}^{E_1}}
\def\Roc{R_{1,c}}
\def\Rtc{R_{2,c}}
\def\rd{r_{0}}
\def\sd{s_{0}}
\def\dd{D_{0}}
\def\vd{v_{0}}
\def\zd{z_{0}}
\def\se{\sE}
\def\xd{x_{0}}
%% ========================================
%% SPECIAL VALUES & EQUILIBRIA
%% ========================================
\def\eE{\bar{e}}
\def\wE{\bar{w}}
\def\sE{s_{*}}
\def\Me{\T M_{ee}}
\def\Rnne{\mathbb{R}_{\ge 0}}

\newcommand{\ones}{\mathbf 1}


%% ========================================
%% ABBREVIATIONS (alphabetical by category)
%% ========================================

% -- General mathematical terms --
\def\adm{admissible}\def\admy{admissibility}
%\def\and{antisymmetric}% disabled 2026-09-30: it overrode LaTeX's author separator \and (printed "antisymmetric" between authors)
\def\ass{assumption}
\def\cadlag{c\`adl\`ag}
\def\ch{characteristic polynomial}
\def\chp{characteristic polynomial}
\def\coe{coefficient}
\def\comp{computation}
\def\con{condition}
\def\coo{coordinate}
\def\corr{corresponding}
\def\dec{decomposition}
\def\den{denominator}
\def\dete{deterministic}
\def\eig{eigenvalue}\def\eval{eigenvalue}
\def\eigv{eigenvector}\def\eiv{eigenvector}\def\evec{eigenvector}
\def\elt{element}
\def\enn{essentially nonnegative/positivity preserving}
\def\equ{equation}
\def\etal{et al.}
\def\fe{for example}
\def\fno{from now on}
\def\for{\forall}
\def\form{formula}
\def\Fr{Furthermore, }
\def\frt{furthermore}
\def\fun{function}
\def\gene{generalization}
\def\heu{heuristic}
\def\how{however}\def\How{However, }
\def\ie{i.e. }
\def\im{\item}
\def\isd{is defined by}
\def\itc{it is easy to check that}
\def\itf{it follows that}\def\Itf{It follows that}
\def\ith{it holds that}
\def\itm{it may be checked that}\def\Itm{It may be checked that}
\def\ito{it turns out that}\def\Ito{It turns out that}\def\Itou{It turns out useful to}
\def\ma{mass action}\def\MAR{mass action representation}
\def\mat{matrix}
\def\mbc{may be checked to be}
\def\mbw{may be written as}\def\mubw{must be written as}
\def\mpr{more precisely, }\def\Mp{More precisely, }
\def\nc{non-catalytic}
\def\nne{non-negative}\def\nny{nonnegativity}
\def\no{\nonumber}
\def\noni{non-infectious}
\def\nota{notation}
\def\np{\newpage}
\def\Nt{Note that}
\def\nt{non trivial}\def\nz{non zero}
\def\oth{otherwise}
\def\Oth{On the other hand}
\def\pa{parameter}\def\para{parameter}
\def\pol{polynomial}
\def\prob{probability}
\def\Prf{{\bf  Proof:}}\def\prf{{\bf  Proof:}}
\def\prop{proportional to}
\def\qu{\quad}
\def\QED{\hfill {$\square$}\goodbreak \medskip}
\def\rep{representation}
\def\resp{respectively}
\def\satd{satisfied}\def\satg{satisfying}\def\sats{satisfies}\def\saty{satisfy}
\def\ssec{\subsection}\def\sssec{\subsubsection}
\def\sef{see for example}
\def\sksy{skew-symmetric}
\def\sub{substrate}
\def\sys{system}
\def\var{variable}
\def\wk{well-known}
\def\wl{without loss of generality}\def\wlo{w.l.o.g.}
\def\wmc{We may check that}\def\Wmc{We may check that}
\def\wmr{we may rewrite}
\def\wms{we must show that}
\def\wrt{with respect to}

% -- CRN & Chemistry --
\newcommand\CRN{chemical reaction network}
\newcommand\CRNs{chemical reaction networks}
\newcommand\RN{reaction network}
\def\ACR{absolute concentration robustness}
\def\AMSD{acyclic siphon decomposition}
\def\DAE{differential algebraic equations}
\def\FHJ{Feinberg-Horn-Jackson}
\def\GMAK{generalized mass action kinetics}
\def\MM{Michaelis-Menten}
\def\sm{stoichiometric matrix}\def\Sm{stoichiometric matrix}
\def\WRZD{weakly reversible, zero deficiency}\def\aut{autocatalytic}

% -- Epidemiology --
\def\brn{basic reproduction number}\def\bzn{basic replacement number}
\def\com{compartment}
\def\dfe{disease free equilibrium}\def\DFE{disease free equilibrium}
\def\dwt{dwell times}\def\edt{expected dwell times}
\def\EM{epidemic model}
\def\endp{endemic fixed point}\def\enp{endemic point}\def\EE{endemic point}
\def\ME{mathematical epidemiology}
\def\MEE{eco-epidemiology}\def\EEP{eco-epidemiological}
\def\ece{eco-epidemiological model}
\def\mgm{multi group model}
\def\mtS{multi-strain model}\def\mSt{multi-strain model}
\def\mSi{minimal siphon}
\def\NGM{next generation matrix}
\def\nim{new infections matrix}
\def\qem{quadratic epidemiological model}
\def\repF{invasion functions}
\def\SPF{SIR-PH-FA model}
\def\sus{susceptible}

% -- Graph theory & Linear algebra --
\def\jac{Jacobi matrix}
\def\jin{Jacobian matrix of the invasion vector field with respect to the invasion variables}
\def\jrn{Jacobian matrix of the resident vector field with respect to the resident variables}
\def\jxn{Jacobian matrix of the $\x$ vector field with respect to the $\x$ variables}
\def\Lm{Laplacian matrix}
\def\MTT{matrix tree theorem}
\def\SGT{spectral graph theory}
\def\SMW{Sherman–Morrison–Woodbury}

% -- Stability & Dynamics --
\def\absC{absorbtion CTMC}
\def\BTb{Bogdanov-Takens bifurcation}
\def\CEP{competitive exclusion partition}
\def\CTMC{continuous time Markov chain}
\def\ds{differential system}
\def\fp{fixed point}\def\bfp{boundary fixed point}\def\fbp{fixed boundary point}
\def\fbl{fixed boundary line}\def\fif{forward invariant face}
\def\gas{global asymptotic stability}\def\GAS{globally asymptotically stable}
\def\GASS{\text{GASS}}
\def\Hb{Hopf bifurcation}
\def\Hur{Hurwitz}\def\MH{Metzler-Hurwitz}\def\RH{Routh-Hurwitz}
\def\IS{\text{IS}}
\def\LAS{locally asymptotically stable}
\def\LdS{Lyapunov diagonally stable}
\def\Lf{Lyapunov function}\def\LyV{Lyapunov Volterra type function}\def\RLF{robust Lyapunov function}
\def\LyM{Lyapunov Malkin stable}
\def\PF{Perron-Frobenius}
\def\spfp{strictly positive fixed point}
\def\sta{stability}
\def\STP{sharp threshold property}\def\STT{sharp threshold theorem}

% -- Models (with \varphi explicitly instead of \f) --
\def\AB{$(A,B)$ Arino-Brauer epidemic models}
\def\ABa{$(A,\bbe,\ba,\varphi)$ balanced bilinear model}
\def\ABP{%$(A,\bbe,P,\varphi)$
\\ balanced bilinear model}
\def\AMG{Anderson May Gupta}
\newcommand\BIN{biological interaction network}
\def\Dt{Descartes type}
\def\gLV{generalized Lotka-Volterra model}
\def\IG{\text{invasion graph}}\def\IGMS{interaction graph on minimal siphons}
\def\LCP{\text{linear complementarity problem}}
\def\LV{Lotka-Volterra}\def\LVK{Lotka-Volterra-Kolmogorov}

% -- Splitting & Structure --
\def\reg{regular splitting}\def\regS{{\bf regular splitting}}
\def\RUR{rational univariate representation}
\def\bds{block-diagonal structure}\def\NGMs{next generation matrices}

% -- Miscellaneous --
\def\Ma{{Mathematica}}
\def\Mm{Max–Min}
\def\num{numerator}
\def\rap{rational approximation} \def\DSR{directed species-reaction bipartite graph}\def\PCP{polynomial complementarity problem}\def\rat{rational substitution}
\def\QSS{quasi steady state} \def\RHS{right hand side} \def\LHS{left hand side}
\def\OSN{online social network}\def\VdD{Van den Driessche}
\def\wli{$\omega$-limit} \def\rk{rich kinetics}\def\sJ{symbolic Jacobian}
\def\cV{\cite{Vas,VasSta,BSV,GSV}} \def\spf{semi-parametric semi-functional}
\def\CRW{Capasso-Ruan-Wang} \def\ABP{$(A,\bbe,P,\f)$ balanced bilinear model}
\def\ABa{%$(A,\bbe,\ba,\f)$ \\
rank-one balanced bilinear model}
\def\bbe{B%{\bff \beta}
} \def\vbe{\vec {\bff \beta}
}
\def\vb{\vec {\bff b}}
\newcommand{\Kh}{\hat{K}}          % K with hat, first mention only
\newcommand{\wS}{w^{\sigma}}       % right Perron J^\perp, full block
\newcommand{\piS}{\vec \pi^{\sigma}}    % left Perron J^\perp, full block
\newcommand{\ws}{w^{\sigma}}       % right Perron J^\perp, sub-block
\newcommand{\pis}{\pi^{\sigma}}    % left Perron J^\perp, sub-block
\newcommand{\hwS}{\hat{w}^{\sigma}}   % right Perron K, full block
\newcommand{\hpiS}{\hat{\pi}^{\sigma}} % left Perron K, full block
\newcommand{\hws}{\hat{w}^{\sigma}}   % right Perron K, sub-block
\newcommand{\hpis}{\hat{\pi}^{\sigma}} % left Perron K, sub-block
\newcommand{\twS}{\tilde{w}^{\sigma}}  % right Perron \tilde K, full block
\newcommand{\tpiS}{\tilde{\pi}^{\sigma}} % left Perron \tilde K
\newcommand{\tws}{\tilde{w}^{\sigma}}
\newcommand{\tpis}{\tilde{\pi}^{\sigma}}
\newcommand{\Mat}[2]{\mathbb{R}_{\ge 0}^{\,#1\times #2}}
\def\tJ{J_\sigma^\perp} \def\lPe{left Perron eigenvector}
% \G is \Gamma (defined at :795).  The Volterra function keeps its own name,
% so that \G no longer has two conflicting meanings in this file.
\def\GV{G}  % Volterra G(u)=u-1-\ln u
\def\PVL{Perron-Volterra Lyapunov minimal siphon ansatz}
\newcommand{\istar}[1]{i_{#1}^{*}}
%\newcommand{\sE}{s_E}
\newcommand{\vE}{v_E}
\def\PV{Perron-Volterra candidate Lyapunov function}
\def\GLV{generalized Lotka-Volterra}
\def\MERN{mathematical epidemiology reaction networks}
\def\CMM{contraction metric method}
\newcommand{\sign}{\operatorname{sign}}
\newcommand{\Var}{\operatorname{Var}} 

% ======================================================================
% IDEAS SHARED BETWEEN THE ABSTRACT AND THE MAIN RESULTS SECTION
%
% Write each idea ONCE, in the abstract; the Main Results section recalls it,
% so the two can never drift apart.
%
%   in the abstract        \ideaitem{1}{The first-order reaction terms ...}
%   in Main Results        \ideahead{1}
%
% \ideaitem stores the text globally and typesets
%   \item \hyperlink{idea:1}{<text>}
% \ideahead typesets
%   \item\hypertarget{idea:1}{}\textbf{<the same text>}
%
% The abstract must come before the Main Results input, which is the normal
% order. If a heading has to read differently from the abstract sentence (for
% instance a long abstract item summarised in one line), override it with the
% optional argument: \ideahead[short heading]{6}.
% ======================================================================
\makeatletter
\newcommand{\ideaitem}[2]{%
  \expandafter\gdef\csname @idea@#1\endcsname{#2}%
  \item\hyperlink{idea:#1}{#2}}
\newcommand{\ideahead}[2][]{%
  \item\hypertarget{idea:#2}{}%
  \def\@ideatmp{#1}%
  \ifx\@ideatmp\@empty
    \@ifundefined{@idea@#2}{%
      \PackageWarning{preamble}{Idea #2 is used in Main Results but was never
        defined by \string\ideaitem\space in the abstract}%
      \textbf{[idea #2 undefined]}%
    }{\textbf{\begingroup\def\ideaterm##1##2{##2}\def\idearest##1{}\@nameuse{@idea@#2}\endgroup}}%
  \else
    \textbf{#1}%
  \fi}
% \ideaterm{<hypertarget>}{<words>}: inside an abstract idea the words become an
% invisible link (the document sets \hypersetup{hidelinks}) jumping to the place
% in Main Results that defines them; inside the Main Results heading itself,
% where the target sits a few lines below, \ideahead strips the link and prints
% the words plainly. Mark the destination with \hypertarget{<name>}{} on the line
% before the display, which adds no equation number and shows nothing.
\newcommand{\ideaterm}[2]{\hyperlink{#1}{#2}}
% \idearest{<text>}: the part of an abstract idea after its first sentence; printed in
% the abstract, omitted from the Main Results heading.
\newcommand{\idearest}[1]{#1}
\makeatother

% ======================================================================
% SHORT-PAPER SWITCH
% \shorttrue in a master that assembles only part of the companion paper
% (Barrat/SipArxiv.tex). \nsh{...} drops text that cites results living in the
% omitted sections, so the short build has no dangling references while the
% companion paper keeps the text verbatim.
% ======================================================================
\newif\ifshort \shortfalse
\newcommand{\nsh}[1]{\ifshort\else#1\fi}
% \shortalt{<short build>}{<companion paper>}: use when a sentence needs
% different wording in the two builds, not merely a deletion.
\newcommand{\shortalt}[2]{\ifshort#1\else#2\fi}
% \Jtr, \JtrS{\Sigma}: transverse Jacobian. J^\perp here; SipArxiv.tex redefines them as J_{\mathrm{tra}}, J_{\mathrm{tra},\Sigma}
\newcommand{\Jtr}{J^\perp}
\newcommand{\JtrS}[1]{J_{#1}^\perp}
.
\makeatletter
\@ifpackageloaded{hyperref}{}{%
  \@ifundefined{hyperrefpackage}{%
    \@ifundefined{AAdefNnohyperref}{\usepackage{hyperref}}{}%
  }{}%
}
\makeatother
% linktoc=all makes a table-of-contents page number a link as well as the
% entry text; hyperref's default, linktoc=section, links only the text.
% Deferred to begin-document and guarded: by the block above hyperref may
% be loaded later (mdpi's own hook) or not at all (\AAdefNnohyperref), and
% \hypersetup must exist before it is called.  A driver that wants the
% default back can say \hypersetup{linktoc=section} after % Shared preamble for the GLVbook papers.  This is the single canonical copy;
% a driver one level below GLVbook reaches it as
%
%     \input{../basicFiles/preamble}
%
% What stays in the driver, because it differs from paper to paper:
% \documentclass, font choice, \title, \author, \date, and any macro genuinely
% local to one paper.  Everything else belongs here.  The file is
% self-contained: it does not read any other file.

\usepackage[utf8]{inputenc}
\usepackage[novbox]{pdfsync}   % required: SyncTeX alone does not work here

% Page geometry, set once here for every paper that loads this file.  A driver
% that needs different margins says \geometry{...}; a second
% \usepackage[...]{geometry} raises an option clash and must not be used.
\makeatletter
\@ifundefined{AAdefNnogeometry}{\usepackage[a4paper,margin=2cm]{geometry}}{}
\makeatother

% mathematics
\usepackage{amsmath,amssymb,amsthm,mathtools,bm,mathrsfs,bbm}
\usepackage{centernot}

% floats, tables, lists
\usepackage{graphicx,float}
\makeatletter
\@ifpackageloaded{subfig}{}{\usepackage{subcaption}}
\makeatother
\usepackage{booktabs,longtable,threeparttable,array,multirow}
\usepackage{enumitem}

% colour, code, algorithms
\usepackage{xcolor,listings,tcolorbox}
\makeatletter
\@ifundefined{c@algorithm}{\usepackage{algorithm}}{}
\makeatother
\usepackage{algorithmicx,algpseudocode}

% diagrams
\usepackage{tikz,tikz-cd,pgfplots}
\usetikzlibrary{graphs,decorations.pathmorphing,decorations.markings,
arrows.meta,automata,positioning}
\tikzcdset{every label/.append style = {font = \small}}
\pgfplotsset{compat=1.18}

% chemistry
\usepackage[version=4]{mhchem}

\usepackage{microtype}
\usepackage{lineno}
\usepackage{nomencl}
\makenomenclature

% hyperref first, then bookmark, then glossaries: glossaries must follow
% hyperref or the entry links are built against the wrong anchors.
\makeatletter
\@ifundefined{hyperrefpackage}{%
  \@ifundefined{AAdefNnohyperref}{\usepackage{hyperref}}{}%
}{}
\@ifundefined{hyperrefpackage}{\usepackage{bookmark}}{}
\makeatother
\usepackage[toc,sort=use]{glossaries}
\makeglossaries

% Book/standalone switch.  Material that the book already covers in its
% prelude, but that a standalone paper still has to state, is wrapped as
%     \ifbookbuild\else ... \fi
% Drivers default to standalone; bookM.tex says \bookbuildtrue.
\newif\ifbookbuild
\bookbuildfalse

% Part switches for the gradostat paper.  Both default to true; a driver that
% wants the short submission says \madecfalse and/or \slowfastfalse, and every
% cross-reference into the dropped part disappears with it.
%   \ifmadec     Part III, the Lagasquie--Madec model and its break-even geometry
%   \ifslowfast  Part IV, everything carrying the small parameter \eps
\newif\ifmadec     \madectrue
\newif\ifslowfast  \slowfasttrue

% Part switch for the Barrat/simplicial-contagion project.  Defaults to
% true (the hypernetwork-synchrony part is included); the journal-paper
% driver says \hnetpartfalse to drop it and every cross-reference into it.
\newif\ifhnetpart  \hnetparttrue

%% ========================================
%% PACKAGES
%% ========================================
\RequirePackage{etoolbox}
\usepackage{makecell}
\usepackage[novbox]{pdfsync}% obsolete; SyncTeX is built into modern engines
\usepackage{amssymb, amsfonts,mathtools,bbm,eurosym,graphicx,latexsym, amsmath,amsthm,times,nccmath,url}
\usepackage{listings,verbatim,cases,tabularx}
% hyperref is deliberately NOT in the line above.  Journal classes often
% load it themselves, with a long option list, from a begindocument hook
% (mdpi.cls does: \DeclareHookRule{begindocument}{hyperref}{before}{mdpi}).
% Loading it here first, without options, makes that later load raise
% "Option clash for package hyperref".  So load it only if nobody else
% has, or will:
%   - \@ifpackageloaded  catches a class that loaded it already;
%   - \hyperrefpackage    is mdpi.cls's own deferred loader;
%   - \AAdefNnohyperref   lets any other driver opt out by defining it
%                         before % Shared preamble for the GLVbook papers.  This is the single canonical copy;
% a driver one level below GLVbook reaches it as
%
%     \input{../basicFiles/preamble}
%
% What stays in the driver, because it differs from paper to paper:
% \documentclass, font choice, \title, \author, \date, and any macro genuinely
% local to one paper.  Everything else belongs here.  The file is
% self-contained: it does not read any other file.

\usepackage[utf8]{inputenc}
\usepackage[novbox]{pdfsync}   % required: SyncTeX alone does not work here

% Page geometry, set once here for every paper that loads this file.  A driver
% that needs different margins says \geometry{...}; a second
% \usepackage[...]{geometry} raises an option clash and must not be used.
\makeatletter
\@ifundefined{AAdefNnogeometry}{\usepackage[a4paper,margin=2cm]{geometry}}{}
\makeatother

% mathematics
\usepackage{amsmath,amssymb,amsthm,mathtools,bm,mathrsfs,bbm}
\usepackage{centernot}

% floats, tables, lists
\usepackage{graphicx,float}
\makeatletter
\@ifpackageloaded{subfig}{}{\usepackage{subcaption}}
\makeatother
\usepackage{booktabs,longtable,threeparttable,array,multirow}
\usepackage{enumitem}

% colour, code, algorithms
\usepackage{xcolor,listings,tcolorbox}
\makeatletter
\@ifundefined{c@algorithm}{\usepackage{algorithm}}{}
\makeatother
\usepackage{algorithmicx,algpseudocode}

% diagrams
\usepackage{tikz,tikz-cd,pgfplots}
\usetikzlibrary{graphs,decorations.pathmorphing,decorations.markings,
arrows.meta,automata,positioning}
\tikzcdset{every label/.append style = {font = \small}}
\pgfplotsset{compat=1.18}

% chemistry
\usepackage[version=4]{mhchem}

\usepackage{microtype}
\usepackage{lineno}
\usepackage{nomencl}
\makenomenclature

% hyperref first, then bookmark, then glossaries: glossaries must follow
% hyperref or the entry links are built against the wrong anchors.
\makeatletter
\@ifundefined{hyperrefpackage}{%
  \@ifundefined{AAdefNnohyperref}{\usepackage{hyperref}}{}%
}{}
\@ifundefined{hyperrefpackage}{\usepackage{bookmark}}{}
\makeatother
\usepackage[toc,sort=use]{glossaries}
\makeglossaries

% Book/standalone switch.  Material that the book already covers in its
% prelude, but that a standalone paper still has to state, is wrapped as
%     \ifbookbuild\else ... \fi
% Drivers default to standalone; bookM.tex says \bookbuildtrue.
\newif\ifbookbuild
\bookbuildfalse

% Part switches for the gradostat paper.  Both default to true; a driver that
% wants the short submission says \madecfalse and/or \slowfastfalse, and every
% cross-reference into the dropped part disappears with it.
%   \ifmadec     Part III, the Lagasquie--Madec model and its break-even geometry
%   \ifslowfast  Part IV, everything carrying the small parameter \eps
\newif\ifmadec     \madectrue
\newif\ifslowfast  \slowfasttrue

% Part switch for the Barrat/simplicial-contagion project.  Defaults to
% true (the hypernetwork-synchrony part is included); the journal-paper
% driver says \hnetpartfalse to drop it and every cross-reference into it.
\newif\ifhnetpart  \hnetparttrue

%% ========================================
%% PACKAGES
%% ========================================
\RequirePackage{etoolbox}
\usepackage{makecell}
\usepackage[novbox]{pdfsync}% obsolete; SyncTeX is built into modern engines
\usepackage{amssymb, amsfonts,mathtools,bbm,eurosym,graphicx,latexsym, amsmath,amsthm,times,nccmath,url}
\usepackage{listings,verbatim,cases,tabularx}
% hyperref is deliberately NOT in the line above.  Journal classes often
% load it themselves, with a long option list, from a begindocument hook
% (mdpi.cls does: \DeclareHookRule{begindocument}{hyperref}{before}{mdpi}).
% Loading it here first, without options, makes that later load raise
% "Option clash for package hyperref".  So load it only if nobody else
% has, or will:
%   - \@ifpackageloaded  catches a class that loaded it already;
%   - \hyperrefpackage    is mdpi.cls's own deferred loader;
%   - \AAdefNnohyperref   lets any other driver opt out by defining it
%                         before \input{preamble}.
\makeatletter
\@ifpackageloaded{hyperref}{}{%
  \@ifundefined{hyperrefpackage}{%
    \@ifundefined{AAdefNnohyperref}{\usepackage{hyperref}}{}%
  }{}%
}
\makeatother
% linktoc=all makes a table-of-contents page number a link as well as the
% entry text; hyperref's default, linktoc=section, links only the text.
% Deferred to begin-document and guarded: by the block above hyperref may
% be loaded later (mdpi's own hook) or not at all (\AAdefNnohyperref), and
% \hypersetup must exist before it is called.  A driver that wants the
% default back can say \hypersetup{linktoc=section} after \input{preamble}.
\makeatletter
\AtBeginDocument{\@ifundefined{hypersetup}{}{\hypersetup{linktoc=all}}}
\makeatother
\usepackage{graphicx,float,fancybox}
\usepackage{epsfig,epstopdf,tikz}
\usepackage[sans]{dsfont}
\usetikzlibrary{positioning,arrows,arrows.meta,calc}
% Configure listings for Mathematica code
\lstdefinestyle{mathematica}{
 language=Mathematica,
 basicstyle=\small\ttfamily,
 backgroundcolor=\color{gray!10},
 frame=single,
 breaklines=true,
 commentstyle=\color{green!60!black},
 keywordstyle=\color{blue},
 stringstyle=\color{red},
 showstringspaces=false,
 tabsize=2
}
\usepackage{tcolorbox}
% Colored boxes for key results
\newtcolorbox{keyresult}[1][]{
 colback=blue!5!white,
 colframe=blue!75!black,
 title=#1,
 fonttitle=\bfseries
}
\newtcolorbox{examplebox}[1][]{
 colback=green!5!white,
 colframe=green!75!black,
 title=#1,
 fonttitle=\bfseries
}


\makeatletter
\@ifundefined{theorem}   {\newtheorem{theorem}   {Theorem}}   {}
\@ifundefined{definition}{\newtheorem{definition}{Definition}} {}
\@ifundefined{open}      {\newtheorem{open}      {Problem}}   {}
\@ifundefined{example}   {\newtheorem{example}   {Example}}   {}
\makeatother

% COUNTER DEFINITIONS
\newcounter{tabcounter}
\newcounter{defcounter}
\newcounter{opcounter}
\newcounter{thmcounter}
\newcounter{excounter}
\newcounter{ascounter}
\newcounter{cjcounter}
\newcounter{lemcounter}
\newcounter{remcounter}
\newcounter{corcounter}
\newcounter{propcounter}
\newcounter{figcounter}

% UNIFIED GLOSSARY STYLE
\newglossarystyle{unifiedstyle}{%
  \setglossarystyle{list}%
  \renewcommand*{\glossarypreamble}{\gdef\glslastsec{}}%
  \renewcommand{\glossentry}[2]{%
    \edef\glscursec{\glsentryuseriii{##1}}%
    \ifx\glscursec\glslastsec\else
      \item[]\textbf{\S\,\glscursec\enspace\glsentryuseriv{##1}}%
      \global\let\glslastsec\glscursec
    \fi
    \item[\glsentryuserii{##1} \glsentryuseri{##1}:] ``\glossentryname{##1}'' \glsentrydesc{##1}%
  }%
}

\makeatletter
% First-word extractor (used by \beT, \beD and \beXa for the default key)
% \@firstwordof hello world \relax  ->  hello
\def\@firstwordof#1 #2\relax{#1}
% First three letters, used for the heading mnemonic of a theorem given no
% explicit key.  Padding with \@empty keeps it safe for first words shorter
% than three letters: \@empty vanishes under \edef, whereas \relax would
% survive into the label and is illegal inside \csname (it broke every
% two-letter title, e.g. "ME ...", with "Missing \endcsname inserted").
\def\@takethree#1#2#3#4\@nil{#1#2#3}

% EXAMPLE COMMAND -- same signature and behaviour as \beT and \beD.
% Usage: \beXa*[long title][short title][key]
%   star  = suppresses the printed "[key]" tag after the title (as for \beT*
%   and \beD*); the label and glossary entry are unaffected.
%   key = optional label key override; the label becomes e:KEY, and without a
%   key it is e: plus the first three letters of the title's first word.  A
%   mnemonic label lets an example be cited as "Example~\ref{e:KEY}" instead
%   of by section number, which is what the citing prose needs.
\NewDocumentCommand{\beXa}{s O{} O{} o}{%
  \def\@temparg{#2}%
  \ifx\@temparg\@empty
    \begin{example}%
  \else
    \let\@savedSecTitle\@currentsectionname%
    \refstepcounter{excounter}%
    \IfValueTF{#4}{%
      \edef\@exlblkey{#4}%
    }{%
      \edef\@exfirstword{\@firstwordof#2 \relax}%
      \edef\@exlblkey{%
        \expandafter\@takethree\@exfirstword\@empty\@empty\@empty\@nil}%
    }%
    \def\@tempshort{#3}%
    \ifx\@tempshort\@empty
      \expandafter\xdef\csname @xshort@\@exlblkey\endcsname{\@exlblkey}%
    \else
      \expandafter\def\csname @xshort@\@exlblkey\endcsname{#3}%
    \fi
    \IfBooleanTF{#1}{\let\@extag\@empty}{%
      \edef\@extag{\noexpand\ \noexpand\textup{[\noexpand\texttt{e:\@exlblkey}]}}}%
    \begin{example}[#2\ifshowkeys\@extag\fi]%
    \label{e:\@exlblkey}%
    % record the long title so \nref{e:KEY} prints it, as for theorems
    \protected@write\@auxout{}{%
      \string\@thmtitle{e:\@exlblkey}{\unexpanded{#2}}}%
    \expanded{%
      \noexpand\newglossaryentry{e\theexcounter}{%
        name={\unexpanded{#2}%
          \IfBooleanF{#1}{ [\noexpand\texttt{e:\@exlblkey}]}},%
        sort={\theexample},%
        description={p.~\noexpand\pageref{e:\@exlblkey}},%
        user1={\theexample},%
        user2={Example},%
        user3={\thesection},%
        user4={\expandafter\unexpanded\expandafter{\@savedSecTitle}}%
      }%
    }%
    \protected@write\@auxout{}{\string\glsliveentry{e\theexcounter}}%
  \fi
}
\newcommand{\eeXa}{\end{example}}

% DEFINITION COMMAND -- same signature and behaviour as \beT below.
% Usage: \beD[long title][short title][key]
%   long title  = definition heading (plain text only)
%   short title = shown in \dref{word}; defaults to the key
%   key         = optional label key override; the label becomes d:KEY
% Without a key the label is d: plus the first three letters of the title's
% first word, so it is a mnemonic rather than the definition's number: a
% numeric label would silently change whenever a definition is inserted.
\NewDocumentCommand{\beD}{s O{} O{} o}{%
  \def\@temparg{#2}%
  \ifx\@temparg\@empty
    \begin{definition}%
  \else
    \let\@savedSecTitle\@currentsectionname%
    \refstepcounter{defcounter}%
    \IfValueTF{#4}{%
      \edef\@deflblkey{#4}%
    }{%
      \edef\@deffirstword{\@firstwordof#2 \relax}%
      \edef\@deflblkey{%
        \expandafter\@takethree\@deffirstword\@empty\@empty\@empty\@nil}%
    }%
    \def\@tempshort{#3}%
    \ifx\@tempshort\@empty
      \expandafter\xdef\csname @dshort@\@deflblkey\endcsname{\@deflblkey}%
    \else
      \expandafter\def\csname @dshort@\@deflblkey\endcsname{#3}%
    \fi
    \IfBooleanTF{#1}{\let\@deftag\@empty}{%
      \edef\@deftag{\noexpand\ \noexpand\textup{[\noexpand\texttt{d:\@deflblkey}]}}}%
    \begin{definition}[#2\ifshowkeys\@deftag\fi]%
    \label{d:\@deflblkey}%
    % record the long title so \nref{d:KEY} prints it, as for theorems
    \protected@write\@auxout{}{%
      \string\@thmtitle{d:\@deflblkey}{\unexpanded{#2}}}%
    \expanded{%
      \noexpand\newglossaryentry{d\thedefcounter}{%
        name={\unexpanded{#2}%
          \IfBooleanF{#1}{ [\noexpand\texttt{d:\@deflblkey}]}},%
        sort={\thedefinition},%
        description={p.~\noexpand\pageref{d:\@deflblkey}},%
        user1={\thedefinition},%
        user2={Definition},%
        user3={\thesection},%
        user4={\expandafter\unexpanded\expandafter{\@savedSecTitle}}%
      }%
    }%
    \protected@write\@auxout{}{\string\glsliveentry{d\thedefcounter}}%
  \fi
}
\newcommand{\eeD}{\end{definition}}

% \showkeysfalse in a driver hides the [t:KEY] mnemonics in result headings
% (labels and glossary entries are unchanged); used for submission builds.
\newif\ifshowkeys \showkeystrue
% \commonredtrue in a driver that inputs basicFiles/commonRed: tool.tex then
% cites its reduction results instead of restating them.
\newif\ifcommonred \commonredfalse
% TheorEM COMMAND
% Usage: \beT[long title][short title][key]
%   long title  = theorem heading (plain text only)
%   short title = shown in \tref{word}; may contain math; defaults to first word of long title
%   key         = optional label key override; label becomes t:KEY (avoids first-word collisions)
% Label is t:WORD where WORD = first word of long title (or KEY if given); use \tref{WORD} to reference.
\NewDocumentCommand{\beT}{s O{} O{} o}{%
  \def\@temparg{#2}%
  \ifx\@temparg\@empty
    \begin{theorem}%
  \else
    \let\@savedSecTitle\@currentsectionname%
    \refstepcounter{thmcounter}%
    % \@thmtag is the mnemonic shown in the heading; every theorem gets one.
    % With an explicit key the key itself is shown.  Without one the key
    % defaults to the title's first word, and showing that whole word would
    % merely repeat the title ("Dissipativity [Dissipativity]"), so only its
    % first three letters are shown: "Dissipativity [Dis]".  The label is
    % unaffected -- it still uses the full \@thmlblkey.
    \IfValueTF{#4}{%
      \edef\@thmlblkey{#4}%
    }{%
      \edef\@thmfirstword{\@firstwordof#2 \relax}%
      \edef\@thmlblkey{%
        \expandafter\@takethree\@thmfirstword\@empty\@empty\@empty\@nil}%
    }%
    \IfBooleanTF{#1}{\let\@thmtag\@empty}{%
      \edef\@thmtag{\noexpand\ \noexpand\textup{[\noexpand\texttt{t:\@thmlblkey}]}}}%
    \def\@tempshort{#3}%
    \ifx\@tempshort\@empty
      \expandafter\xdef\csname @tshort@\@thmlblkey\endcsname{\@thmlblkey}%
    \else
      \expandafter\def\csname @tshort@\@thmlblkey\endcsname{#3}%
    \fi
    \begin{theorem}[#2\ifshowkeys\@thmtag\fi]%
    \label{t:\@thmlblkey}%
    % record the long title for \nref; via the .aux so forward refs resolve
    \protected@write\@auxout{}{%
      \string\@thmtitle{t:\@thmlblkey}{\unexpanded{#2}}}%
    \expanded{%
      \noexpand\newglossaryentry{t\thethmcounter}{%
        name={\unexpanded{#2}%
          \IfBooleanF{#1}{ [\noexpand\texttt{t:\@thmlblkey}]}},%
        % sort/user1 must be the number the document actually prints, i.e. the
        % theorem environment's counter, not thmcounter.  thmcounter advances
        % only on \beT with a title, so a bare \beT (or \beT with an
        % empty title) makes the two drift and the glossary cites a wrong
        % number.  thmcounter is still fine as the entry key: it is unique.
        sort={\thetheorem},%
        description={p.~\noexpand\pageref{t:\@thmlblkey}},%
        user1={\thetheorem},%
        user2={Theorem},%
        user3={\thesection},%
        user4={\expandafter\unexpanded\expandafter{\@savedSecTitle}}%
      }%
    }%
    \protected@write\@auxout{}{\string\glsliveentry{t\thethmcounter}}%
  \fi
}
\newcommand{\eeT}{\end{theorem}}

% CROSS-PART RELATION PARAGRAPH
% Usage: \xpart{text}  -- place immediately BEFORE \beT, never inside it
%   (anything inside \beT...\eeT would be swallowed into the glossary entry).
% The paragraph is typeset only when Part I is actually present, i.e. in the
% combined book build; in the standalone builds (GranadaDSABNS/A.tex,
% AAHRah-Submission/AAHRah.tex) the label part:glv is undefined and the whole
% paragraph disappears rather than leaving ?? cross-references. Same
% graceful-degradation principle as \GASref in GASt.tex. Resolves on the
% second LaTeX pass, since r@part:glv comes from the .aux file.
% (Already inside the \makeatletter region opened above, so no \makeatletter
% pair here -- adding one would close that region early.)
\newcommand{\xpart}[1]{\@ifundefined{r@part:glv}{}{\paragraph{Relation with Parts~I--II.}#1}}

% \bookref{label}{precise text}{fallback text}: prints the precise cross-reference
% when the label resolves (combined book build), and the fallback text otherwise
% (standalone paper builds, where book-only appendix labels are absent). Same
% graceful-degradation role as \GASref in GASt.tex, but usable in any file since it
% lives in this preamble (loaded by every driver) -- unlike \GASref, which is defined in a
% content file and is therefore undefined in standalone builds that omit that file.
\newcommand{\bookref}[3]{\@ifundefined{r@#1}{#3}{#2}}

% \PVpart and \RBpart: the two parts that a partial build may omit.
% They print "Part~N" in any build that contains the part, and a naming phrase
% in a build that does not, so an excerpt shows no ?? .  Each already supplies
% the word "Part", so write "\PVpart develops ...", not "Part~\PVpart".
\newcommand{\PVpart}{\bookref{part:aahrah}%
  {Part~\ref{part:aahrah}}{the Perron--Volterra part of the full book}}
\newcommand{\RBpart}{\bookref{part:recbrowne}%
  {Part~\ref{part:recbrowne}}{the immuno-virology part of the full book}}

% OPEN PROBLEM COMMANDS
% ASSUMPTION COMMANDS (glossaried)
\newcommand{\beAs}[1][]{%
  \def\@temparg{#1}%
  \ifx\@temparg\@empty
    \begin{assumption}%
  \else
    \let\@savedSecTitle\@currentsectionname%
    \refstepcounter{ascounter}%
    \begin{assumption}[#1]%
    \label{as:\theascounter}%
    \expanded{%
      \noexpand\newglossaryentry{as\theascounter}{%
        name={\unexpanded{#1}},%
        sort={\theascounter},%
        description={p.~\noexpand\pageref{as:\theascounter}},%
        user1={\theascounter},%
        user2={Assumption},%
        user3={\thesection},%
        user4={\expandafter\unexpanded\expandafter{\@savedSecTitle}}%
      }%
    }%
    \protected@write\@auxout{}{\string\glsliveentry{as\theascounter}}%
  \fi
}
\newcommand{\eeAs}{\end{assumption}}
% CONJECTURE COMMANDS (glossaried)
\newcommand{\beCj}[1][]{%
  \def\@temparg{#1}%
  \ifx\@temparg\@empty
    \begin{conjecture}%
  \else
    \let\@savedSecTitle\@currentsectionname%
    \refstepcounter{cjcounter}%
    \begin{conjecture}[#1]%
    \label{cj:\thecjcounter}%
    \expanded{%
      \noexpand\newglossaryentry{cj\thecjcounter}{%
        name={\unexpanded{#1}},%
        sort={\thecjcounter},%
        description={p.~\noexpand\pageref{cj:\thecjcounter}},%
        user1={\thecjcounter},%
        user2={Conjecture},%
        user3={\thesection},%
        user4={\expandafter\unexpanded\expandafter{\@savedSecTitle}}%
      }%
    }%
    \protected@write\@auxout{}{\string\glsliveentry{cj\thecjcounter}}%
  \fi
}
\newcommand{\eeCj}{\end{conjecture}}
% LEMMA COMMANDS
% ---------------------------------------------------------------------
% UNIFORM RESULT WRAPPERS -- \beL, \beC, \beP, \beR share \beT's signature
%
%   \beX[long title][short title][key]  ...  \eeX
%
%   long title  = heading (math allowed)
%   short title = shown in the cross-reference; defaults to the key
%   key         = label key.  With a key the label is X:KEY and the heading
%                 carries the mnemonic [key], exactly as for \beT.  Without
%                 one the label falls back to the running counter, which is
%                 what every pre-existing call site in the book relies on.
%
% A bare \beX -- arguments commented out -- opens the plain environment and
% creates NO label and NO glossary entry.  That is the supported way to keep
% an individual result out of the catalogue.
% ---------------------------------------------------------------------
% #1 wrapper name  #2 environment  #3 label prefix  #4 counter  #5 type name
\newcommand{\@mkresult}[5]{%
  % A leading * suppresses the [key] mnemonic in both the heading and the
  % glossary name, while still making the label X:KEY and the entry:
  %     \beX*[long][short][key]
  \expandafter\NewDocumentCommand\csname be#1\endcsname{s O{} O{} o}{%
    \def\@temparg{##2}%
    \ifx\@temparg\@empty
      \begin{#2}%
    \else
      \let\@savedSecTitle\@currentsectionname%
      \refstepcounter{#4}%
      % With a key: mnemonic label X:KEY and the [key] tag in the heading.
      % Without one: fall back to the running counter, exactly as these
      % wrappers behaved before, so every legacy one-argument call site in
      % the book keeps its label.  The counter is also csname-safe, which a
      % title's first word is not -- titles beginning with math would
      % otherwise build \csname \Omega...\endcsname and fail.
      \IfValueTF{##4}{%
        \edef\@relbl{##4}%
        \IfBooleanTF{##1}{\let\@retag\@empty}{%
          \edef\@retag{\noexpand\ \noexpand\textup{[\noexpand\texttt{#3:##4}]}}}%
      }{%
        \edef\@relbl{\arabic{#4}}%
        \let\@retag\@empty
      }%
      \def\@tempshort{##3}%
      \ifx\@tempshort\@empty
        \expandafter\xdef\csname @sshort@#3\@relbl\endcsname{\@relbl}%
      \else
        \expandafter\def\csname @sshort@#3\@relbl\endcsname{##3}%
      \fi
      \begin{#2}[##2\ifshowkeys\@retag\fi]%
      \edef\@tmplbl{#3:\@relbl}\expandafter\label\expandafter{\@tmplbl}%
      \expanded{%
        \noexpand\newglossaryentry{#3\csname the#4\endcsname}{%
          name={\unexpanded{##2}%
            \IfBooleanF{##1}{\IfValueT{##4}{ [\noexpand\texttt{#3:##4}]}}},%
          sort={\csname the#2\endcsname},%
          description={p.~\noexpand\pageref{#3:\@relbl}},%
          user1={\csname the#2\endcsname},%
          user2={#5},%
          user3={\thesection},%
          user4={\expandafter\unexpanded\expandafter{\@savedSecTitle}}%
        }%
      }%
      \protected@write\@auxout{}{\string\glsliveentry{#3\csname the#4\endcsname}}%
      \protected@write\@auxout{}{\string\@thmtitle{\@tmplbl}{\unexpanded{##2}}}%
    \fi
  }%
  \expandafter\newcommand\csname ee#1\endcsname{\end{#2}}%
}
\@mkresult{L}{lemma}{l}{lemcounter}{Lemma}
% REMARK / COROLLARY / PROPOSITION -- same signature as \beL above
\@mkresult{R}{remark}{r}{remcounter}{Remark}
\@mkresult{C}{corollary}{c}{corcounter}{Corollary}
\@mkresult{P}{proposition}{p}{propcounter}{Proposition}
% OPEN-PROBLEM COMMAND -- generated with the same signature as \beT:
%   \beO[long title][short title][key]   ...   \eeO
% label o:KEY, or o:<counter> when no key is given; a bare \beO makes
% neither a label nor a glossary entry.
\@mkresult{O}{open}{o}{opcounter}{Problem}


% Section-level title tracker: updated only by \section, not \subsection.
% \@currentlabelname is overwritten by every \refstepcounter (including
% \subsection), so we save it separately here.
\newcommand{\@currentsectionname}{}
\AtBeginDocument{%
  \let\@epidcrn@orig@refstepcounter\refstepcounter
  \renewcommand{\refstepcounter}[1]{%
    \@epidcrn@orig@refstepcounter{#1}%
    \ifstrequal{#1}{section}{\global\let\@currentsectionname\@currentlabelname}{}%
  }%
}

% TABLE COMMANDS (glossaried) -- consistent naming with \beT, \beD, etc.
\newcommand{\beTa}[1][]{%
  \def\@temparg{#1}%
  \ifx\@temparg\@empty
    \begin{table}[htbp]%
  \else
    \let\@savedSecTitle\@currentsectionname%
    \refstepcounter{tabcounter}%
    \begin{table}[htbp]%
    \label{A:\thetabcounter}%
    \expanded{%
      \noexpand\newglossaryentry{A\thetabcounter}{%
        name={\unexpanded{#1}},%
        sort={\thetabcounter},%
        description={p.~\noexpand\pageref{A:\thetabcounter}},%
        user1={\thetabcounter},%
        user2={Table},%
        user3={\thesection},%
        user4={\expandafter\unexpanded\expandafter{\@savedSecTitle}}%
      }%
    }%
    \protected@write\@auxout{}{\string\glsliveentry{A\thetabcounter}}%
  \fi
}
\newcommand{\eeTa}{\end{table}}

% FIGURE COMMANDS (glossaried) -- consistent with \beT, \beR, \beTa, etc.
% Usage:
%   \beF[Glossary title]...\caption{Visible caption}\eeF  -> glossaried figure
%   \beF...\caption{Visible caption}\eeF                  -> plain [H] figure, no glossary
\newcommand{\beF}[1][]{%
  \def\@temparg{#1}%
  \ifx\@temparg\@empty
    \begin{figure}[H]%
  \else
    \begin{figure}[H]%
    \let\@savedSecTitle\@currentsectionname%
    \refstepcounter{figcounter}%
    \label{fig:\thefigcounter}%
    \expanded{%
      \noexpand\newglossaryentry{Fig\thefigcounter}{%
        name={\unexpanded{#1}},%
        sort={\thefigcounter},%
        description={p.~\noexpand\pageref{fig:\thefigcounter}},%
        user1={\thefigcounter},%
        user2={Figure},%
        user3={\thesection},%
        user4={\expandafter\unexpanded\expandafter{\@savedSecTitle}}%
      }%
    }%
    \protected@write\@auxout{}{\string\glsliveentry{Fig\thefigcounter}}%
  \fi
}
\newcommand{\eeF}{\end{figure}}
\makeatother
% LIVE-ENTRY REGISTER
% \glsaddall prints every entry the glossaries package currently knows, and
% that set includes entries restored from \jobname.glsdefs.  The package
% rewrites .glsdefs from all *currently defined* entries, so anything it
% restored is written straight back out: a deleted section kept reappearing in
% the catalogue for ever, and the .glsdefs had to be erased by hand.  The
% wrappers therefore record every entry they actually create in the .aux, and
% \printboth adds only those.  The .aux is rebuilt from scratch on each run, so
% this is self-healing: a removed result is gone after one further run.
\makeatletter
\newcommand{\@glslive}{}
\newcommand{\@glsaddlive}[1]{\ifglsentryexists{#1}{\glsadd{#1}}{}}
\newcommand{\glsliveentry}[1]{\g@addto@macro\@glslive{\@glsaddlive{#1}}}
\makeatother
% PRINT COMMAND
% No \clearpage here -- the catalogue continues on the current page.  This has
% been re-added by mistake several times; leave it out.
\makeatletter
\newcommand{\printboth}{%
  \@glslive   % add exactly the entries this document created last run
  \printglossary[style=unifiedstyle, title={0 \hspace{.22em} Definitions, Theorems, Examples and Open Problems}]%
}
\makeatother
% Named theorem reference: \tref{word} -> Thm~N(word-or-short)
% Shows stored short title if available, otherwise shows the key itself.
\newcommand{\tref}[1]{%
  Thm~\ref{t:#1}(\ifcsname @tshort@#1\endcsname\csname @tshort@#1\endcsname\else#1\fi)%
}

% Full-title theorem reference: \nref{t:KEY} prints the number and then the
% theorem's long title, so one writes "Theorem~\nref{t:blocks}" and gets
% "Theorem 3 ``Block factorization and siphons''".  Takes the complete label, like
% \ref, not the bare key like \tref.
%
% \beT writes each title to the .aux (see \@thmtitle below), so \nref resolves
% forward references from the second LaTeX pass, exactly as \ref does.  If the
% title is not yet known -- first pass, or a label that is not a \beT theorem
% -- it degrades to the number rather than producing ??.
% These two use @ in their names, and this point in the file is outside the
% \makeatletter region opened earlier (it closes at \makeatother above), so
% the region has to be reopened here.
\makeatletter
\newcommand{\@thmtitle}[2]{\global\@namedef{@tlong@#1}{#2}}
\newcommand{\nref}[1]{%
  \ref{#1}%
  \ifcsname @tlong@#1\endcsname
    \ ``\csname @tlong@#1\endcsname''%
  \fi
}
% \nm{t:KEY}: the stored long title of any \beT/\beL/\beP/\beR/\beC/\beD/\beE result,
% verbatim and without quotes, so item names follow the macro's first argument.
\newcommand{\nm}[1]{%
  \ifcsname @tlong@#1\endcsname
    \csname @tlong@#1\endcsname%
  \else
    \ref{#1}%
  \fi
}
\makeatother




%% ========================================
%% Other THEOREM ENVIRONMENTS
%% ========================================
\makeatletter
\@ifundefined{lemma}      {\newtheorem{lemma}      {Lemma}}      {}
\@ifundefined{proposition}{\newtheorem{proposition}{Proposition}} {}
\@ifundefined{corollary}  {\newtheorem{corollary}  {Corollary}}  {}
\@ifundefined{conjecture} {\newtheorem{conjecture} {Conjecture}}  {}
\@ifundefined{remark}     {\newtheorem{remark}     {Remark}}     {}
\@ifundefined{question}   {\newtheorem{question}   {Question}}   {}
\@ifundefined{assumption} {\newtheorem{assumption} {Assumption}} {}
\@ifundefined{hypothesis} {\newtheorem{hypothesis} {Hypothesis}} {}
\@ifundefined{fact}       {\newtheorem{fact}       {Fact}}       {}
\makeatother

% \citer: cite the relay-graph threshold theorem of AAHrel. If a document has
% \input the chapter/file containing t:relayRF (so the label resolves), point
% straight at it; otherwise fall back to a bare citation of AAHrel. Shared here
% (not per-paper) so it resolves consistently across every book chapter.
% \citeR: same idea, for the first theorem of AAHRah (t:Pstar).
\makeatletter
\providecommand{\citer}{%
  \@ifundefined{r@t:relayRF}%
    {\cite{AAHrel}}%
    {Theorem~\ref{t:relayRF} of \cite{AAHrel}}%
}
\providecommand{\citeR}{%
  \@ifundefined{r@t:Pstar}%
    {\cite{AAHRah}}%
    {Theorem~\ref{t:Pstar} of \cite{AAHRah}}%
}
\makeatother


% Abbreviations
\def\beQ{\begin{question}}\def\eeQ{\end{question}}
\def\beA{\begin{assumption}}\def\eeA{\end{assumption}}


%% ========================================

%% ========================================
%% COLORS & CUSTOM COMMANDS
%% ========================================
\definecolor{funccolor}{RGB}{25,25,112}
\definecolor{desccolor}{RGB}{64,64,64}
\newcommand{\func}[1]{\texttt{\textcolor{funccolor}{#1}}}
\def\eq{=}
\def\mm{-}

%% ========================================
%% MATHCAL SHORTHANDS
%% ========================================
% Project-wide convention: \mX for a capital X is always \mathcal X.
% Defined for the whole alphabet so the convention never fails to expand.
\newcommand{\mA}{\mathcal A}\newcommand{\mB}{\mathcal B}
\newcommand{\mC}{\mathcal C}\newcommand{\mD}{\mathcal D}
\newcommand{\mE}{\mathcal E}\newcommand{\mF}{\mathcal F}
\newcommand{\mG}{\mathcal G}\newcommand{\mH}{\mathcal H}
\newcommand{\mI}{\mathcal I}\newcommand{\mJ}{\mathcal J}
\newcommand{\mK}{\mathcal K}\newcommand{\mL}{\mathcal L}
\newcommand{\mM}{\mathcal M}\newcommand{\mN}{\mathcal N}
\newcommand{\mO}{\mathcal O}\newcommand{\mP}{\mathcal P}
\newcommand{\mQ}{\mathcal Q}\newcommand{\mR}{\mathcal R}
\newcommand{\mS}{\mathcal S}\newcommand{\mT}{\mathcal T}
\newcommand{\mU}{\mathcal U}\newcommand{\mV}{\mathcal V}
\newcommand{\mW}{\mathcal W}\newcommand{\mX}{\mathcal X}
\newcommand{\mY}{\mathcal Y}\newcommand{\mZ}{\mathcal Z}

%% ========================================
%% BASIC MATH ENVIRONMENTS
%% ========================================
\def\bep{\begin{pmatrix}}\def\eep{\end{pmatrix}}
\def\bev{\begin{vmatrix}}\def\eev{\end{vmatrix}}
\def\bea{\begin{eqnarray*}}\def\eea{\end{eqnarray*}}
\def\bc{\begin{cases}}\def\ec{\end{cases}}
\def\BEN{\begin{enumerate}}\def\EEN{\end{enumerate}}
\def\BI{\begin{itemize}}\def\EI{\end{itemize}}

\newcommand{\be}[1]{\begin{equation}\label{#1}}
\newcommand{\ee}{\end{equation}}
\newcommand{\beq}{\begin{eqnarray}}
\def\eeq{\end{eqnarray}}
\def\eqr{\eqref}\def\fr{\frac}\def\lbl{\label}
\def\Lra{\Longrightarrow}\def\non{\nonumber}\def\cNGM{\cite{Diek,Van,Van08}}
\def\Eq{\Leftrightarrow}\def\b{\beta}\def\uph{\red{UP TO HERE}}
%% ========================================
%% PROBABILITY & STATISTICS
%% ========================================
\renewcommand{\P}{\ensuremath{\mathbb{P}}}  % Note: overrides paragraph symbol
\newcommand{\E}{\ensuremath{\mathbb{E}}}
\newcommand{\ind}[1]{\ensuremath{\mathbbm{1}_{\left\{#1\right\}}}}
\DeclareMathOperator{\Cov}{Cov}
\DeclareMathOperator{\TV}{TV}
\newcommand{\rank}{\operatorname{rank}}

%% ========================================
%% CRN OPERATORS (organizations/ERCs/buffering/siphons framework)
%% ========================================
\def\clos{\mathrm{clos}}
\def\Supply{\mathrm{Supply}}
\def\req{\mathrm{req}}
\def\React{\mathrm{React}}
\def\Prod{\mathrm{Prod}}
\def\Gen{\mathrm{Gen}}
\def\In{\mathrm{In}}
\def\Out{\mathrm{Out}}
\def\Act{\mathrm{Act}}
\def\Min{\mathrm{Min}}

%% ========================================
%% LINE BREAKING
%% ========================================
% Last-resort stretch, used only on paragraphs TeX cannot otherwise break:
% a long \texttt{} token, a comma-separated inline math list (commas are not
% breakpoints in math), or a run of \cite/\ref boxes.  Without it those set as
% overfull \hbox and print into the margin; with it they set slightly loose,
% which is the lesser evil.  Does not affect paragraphs that already break.
\emergencystretch=3em

%% ========================================
%% NUMBER SETS
%% ========================================
\newcommand{\R}{\mathbb{R}}  % Real numbers
\newcommand{\C}{\mathbb{C}}  % Complex numbers
\newcommand{\N}{\mathbb{N}}
\newcommand{\Z}{\mathbb{Z}}
% \providecommand, not \newcommand: some papers define these in their own
% preamble after \input{preamble}, and must not clash.
\providecommand{\Q}{\mathbb{Q}}
\providecommand{\res}{\operatorname{res}}   % resultant; used by epiRes.tex
\newcommand{\MN}{\mathbb{N}}
\newcommand{\MZ}{\mathbb{Z}}
\newcommand{\realpos}{\mathbb{R}_{> 0}}
\newcommand{\realnneg}{\mathbb{R}_{\geq 0}}
\newcommand{\naturalpos}{\mathbb{N}_{>0}}
\newcommand{\Rplus}{\mathbb{R}_+}
\newcommand{\Rp}{\mathbb{R}_{\ge 0}}
\def\Rnne{\mathbb{R}_{\ge 0}}
\def\Se{\bar{S}} \def\Ie{\bar{I}}
\newcommand{\Rpos}{\mathbb{R}_{>0}}
\newcommand{\dd}{\,\mathrm d}
\newcommand{\e}{\mathrm e}
%% ========================================
%% GREEK LETTERS (2+ character abbreviations) greek
%% ========================================
\def\al{\alpha}\def\ba{\bff \alpha}
\def\w{\omega}\def\Lgn{\ga+\mu+\de}\def\om{\omega}
\def\de{\delta}\def\De{\Delta}\def\bde{\bff \delta}
\def\ep{\epsilon}\def\eps{\varepsilon}
\def\f{\varphi}  % Kept as requested
\def\ga{\gamma}\def\Ga{\Gamma}\def\G{\Gamma}\def\bg{\bff \gamma}
\def\up{\upsilon}  % NOT the breve accent: no accent commands anywhere.
                   % Named \up, not \u, so the LaTeX breve stays reachable.
\def\la{\lambda}\def\La{\Lambda}
\def\si{\sigma}\def\ka{\kappa}
\def\th{\theta}\def\vt{\vartheta}
\def\ze{\zeta}\def\Gap{\Gamma^{+}} \def\Gam{\Gamma^{-}}


%% ========================================
%% SPECIAL SHORT COMMANDS (kept as requested)
%% ========================================
\def\I{\infty}
\def\T{\widetilde}
\def\H{\hat}
\newcommand{\relint}{\operatorname{relint}}

%% ========================================
%% BOLD/VECTOR NOTATION
%% ========================================
\newcommand{\bff}[1]{{\mbox{\boldmath$#1$}}}
\newcommand\X{\boldsymbol{X}}
\def\x{\boldsymbol{x}}\def\bx{\bff x}\def\xe{\bar{\x}}
\newcommand\y{\boldsymbol{y}}\def\by{\boldsymbol{y}}\def\ye{\bar{\y}}
\def\yb{\y}\def\xb{\x}
\def\xy{\x-\y}
\newcommand\br{\boldsymbol{r}}\def\re{\bar{\r}}
\newcommand\bd{\boldsymbol{d}}
\newcommand\bQ{\boldsymbol{Q}}
\newcommand{\Yb}{{\mathbf{Y}}}
\newcommand{\Ab}{\mathbf{A}}
\newcommand{\zb}{\mathbf{z}}\def\bz{\bff 0}
\newcommand{\Zb}{\mathbf{Z}}
\newcommand{\bb}{\mathbf{b}}
\def\vbe{\vec {\bff \beta}}
\def\vb{\vec {\bff b}}
\def\v1{\vec {\bff 1}}
\def\vn{\vec \de}
\def\vz{\vec 0}
\def\vi{\overset{\rightarrow}{\mathsf i}}
\def\vv{\overset{\rightarrow}{v}}
\def\vr{\; \vec {\mathsf r}}
\newcommand{\va}{\bff a}
\providecommand{\vect}[1]{\boldsymbol{#1}}   % bold vector (moved from Gradostat.tex)
\providecommand{\one}{\mathbf 1}             % all-ones vector (moved from AliSF)
\providecommand{\spr}{s}                     % spectral bound  (moved from AliSF)
\providecommand{\Ke}{K_{\rm eff}}            % effective carrying capacity
% Critical right/left eigenvectors in the Hopf projection method (AAHOSN/New.tex)
\providecommand{\bq}{\mathbf{q}}
\providecommand{\bp}{\mathbf{p}}


%% ========================================
%% MATHCAL & MATHSCR
%% ========================================
% Mathcal
\newcommand\mc[1]{\mathcal{#1}}
\def\mC{{\mathcal C}}\def\mD{\mathcal D}\def\mE{\mathcal{E}}
\def\mF{\mathcal F}\def\mG{{\mathcal G}}\def\mH{{\mathcal H}}
\def\mI{\mathcal I}\def\mL{\mathcal{L}}\def\mM{{\mathcal M}}
\def\mN{{\mathcal N}}\def\mP{\mathcal{P}}\def\mR{{\mathcal R}}
\def\mS{{\mathcal S}}\def\mT{\mathcal{T}}\def\mV{{\mathcal V}}
\def\mK{\mathcal K}\def\mA{{\mathcal A}}\def\mU{{\mathcal U}}
% Symmetric bilinear and trilinear forms of the Taylor expansion, used in the
% first Lyapunov coefficient (AAHOSN/New.tex)
\providecommand{\Bform}{\mathcal{B}}
\providecommand{\Cform}{\mathcal{C}}

% Blackboard bold.  Kept here, not in the drivers: a \providecommand in a
% driver preamble runs before this file and silently wins, so the same name
% ends up meaning different things in different builds.
\def\IC{\mathbb{C}}\def\IN{\mathbb{N}}\def\IQ{\mathbb{Q}}
\def\IR{\mathbb{R}}\def\IZ{\mathbb{Z}}\def\IG{\mathbb{G}}
\def\II{\mathbb{I}}

% Mathscr
\newcommand{\Rs}{{\mathscr R}}
\newcommand{\Ss}{{\mathscr S}}
\newcommand{\Ns}{{\mathscr N}}
\newcommand{\Cs}{\mathscr C}


%% ========================================
%% OPERATORS & FUNCTIONS
%% ========================================
\DeclareMathOperator{\tr}{tr}
\DeclareMathOperator{\diag}{Diag}
\DeclareMathOperator{\Diag}{Diag}
\DeclareMathOperator{\intr}{\mathrm{Int}}
\DeclareMathOperator{\Inv}{Inv}
\DeclareMathOperator{\spec}{spec}
\DeclareMathOperator{\AUC}{AUC}               % area under the curve
\newcommand{\ii}{\mathrm{i}}                  % imaginary unit
\newcommand\op{\operatorname}
\newcommand{\supp}{\op{supp}}
\newcommand{\diff}{\mathop{}\mathopen{}\mathrm{d}}
\newcommand{\pdd}[1]{\frac{\partial}{\partial{#1}}}
\DeclareMathOperator{\cone}{cone}
\DeclareMathOperator{\proj}{proj}

%% ========================================
%% BRACKETS & NORMS
%% ========================================
\providecommand{\pr}[1]{\left(#1\right)}
\providecommand{\pp}[1]{\left[#1\right]}
\newcommand\croc[1]{\left\langle #1\right\rangle}
\newcommand{\norm}[1]{\|#1\|}
\newcommand{\inner}[2]{\langle #1, #2 \rangle}
% Short synonym of \inner, used by AAHOSN/New.tex as $\ip{u}{v}=\bar u^\top v$
\providecommand{\ip}[2]{\langle #1 , #2 \rangle}


%% ========================================
%% FORMATTING & UTILITIES
%% ========================================
\newcommand*{\Scale}[2][4]{\scalebox{#1}{$#2$}}
\newcommand{\te}[1]{\text{#1}}
\newcommand\steq[1]{\stackrel{\text{\rm #1.}}{=}}
\newcommand\blue[1]{\textcolor{blue}{#1}}
\newcommand{\red}{\textcolor[rgb]{1.00,0.00,0.00}}
\newcommand{\green}{\textcolor[rgb]{0.00,1.00,0.00}}

% Figure command
\newcommand{\figu}[3]{%
\begin{figure}[H]
\centering
\includegraphics[scale=#3]{#1}
\caption{#2\label{f:#1}}
\end{figure}
}

% Footnote
\long\def\symbolfootnote[#1]#2{%
\begingroup
\def\thefootnote{\fnsymbol{footnote}}\footnote[#1]{#2}%
\endgroup}
\def\fn{\symbolfootnote}

\newcommand{\dsum}{\displaystyle\sum}

%% ========================================
%% CRN & EPIDEMIOLOGY SPECIFIC
%% ========================================
% CRN notation
\def\SS{\mathcal S}
\def\CC{\mathcal C}
\def\RR{\mathcal R}

% Reproduction numbers
\def\Rto{R_{1}^{E_2}}
\def\Rot{R_{2}^{E_1}}
\def\Roc{R_{1,c}}
\def\Rtc{R_{2,c}}
\def\rd{r_{0}}
\def\sd{s_{0}}
\def\dd{D_{0}}
\def\vd{v_{0}}
\def\zd{z_{0}}
\def\se{\sE}
\def\xd{x_{0}}
%% ========================================
%% SPECIAL VALUES & EQUILIBRIA
%% ========================================
\def\eE{\bar{e}}
\def\wE{\bar{w}}
\def\sE{s_{*}}
\def\Me{\T M_{ee}}
\def\Rnne{\mathbb{R}_{\ge 0}}

\newcommand{\ones}{\mathbf 1}


%% ========================================
%% ABBREVIATIONS (alphabetical by category)
%% ========================================

% -- General mathematical terms --
\def\adm{admissible}\def\admy{admissibility}
%\def\and{antisymmetric}% disabled 2026-09-30: it overrode LaTeX's author separator \and (printed "antisymmetric" between authors)
\def\ass{assumption}
\def\cadlag{c\`adl\`ag}
\def\ch{characteristic polynomial}
\def\chp{characteristic polynomial}
\def\coe{coefficient}
\def\comp{computation}
\def\con{condition}
\def\coo{coordinate}
\def\corr{corresponding}
\def\dec{decomposition}
\def\den{denominator}
\def\dete{deterministic}
\def\eig{eigenvalue}\def\eval{eigenvalue}
\def\eigv{eigenvector}\def\eiv{eigenvector}\def\evec{eigenvector}
\def\elt{element}
\def\enn{essentially nonnegative/positivity preserving}
\def\equ{equation}
\def\etal{et al.}
\def\fe{for example}
\def\fno{from now on}
\def\for{\forall}
\def\form{formula}
\def\Fr{Furthermore, }
\def\frt{furthermore}
\def\fun{function}
\def\gene{generalization}
\def\heu{heuristic}
\def\how{however}\def\How{However, }
\def\ie{i.e. }
\def\im{\item}
\def\isd{is defined by}
\def\itc{it is easy to check that}
\def\itf{it follows that}\def\Itf{It follows that}
\def\ith{it holds that}
\def\itm{it may be checked that}\def\Itm{It may be checked that}
\def\ito{it turns out that}\def\Ito{It turns out that}\def\Itou{It turns out useful to}
\def\ma{mass action}\def\MAR{mass action representation}
\def\mat{matrix}
\def\mbc{may be checked to be}
\def\mbw{may be written as}\def\mubw{must be written as}
\def\mpr{more precisely, }\def\Mp{More precisely, }
\def\nc{non-catalytic}
\def\nne{non-negative}\def\nny{nonnegativity}
\def\no{\nonumber}
\def\noni{non-infectious}
\def\nota{notation}
\def\np{\newpage}
\def\Nt{Note that}
\def\nt{non trivial}\def\nz{non zero}
\def\oth{otherwise}
\def\Oth{On the other hand}
\def\pa{parameter}\def\para{parameter}
\def\pol{polynomial}
\def\prob{probability}
\def\Prf{{\bf  Proof:}}\def\prf{{\bf  Proof:}}
\def\prop{proportional to}
\def\qu{\quad}
\def\QED{\hfill {$\square$}\goodbreak \medskip}
\def\rep{representation}
\def\resp{respectively}
\def\satd{satisfied}\def\satg{satisfying}\def\sats{satisfies}\def\saty{satisfy}
\def\ssec{\subsection}\def\sssec{\subsubsection}
\def\sef{see for example}
\def\sksy{skew-symmetric}
\def\sub{substrate}
\def\sys{system}
\def\var{variable}
\def\wk{well-known}
\def\wl{without loss of generality}\def\wlo{w.l.o.g.}
\def\wmc{We may check that}\def\Wmc{We may check that}
\def\wmr{we may rewrite}
\def\wms{we must show that}
\def\wrt{with respect to}

% -- CRN & Chemistry --
\newcommand\CRN{chemical reaction network}
\newcommand\CRNs{chemical reaction networks}
\newcommand\RN{reaction network}
\def\ACR{absolute concentration robustness}
\def\AMSD{acyclic siphon decomposition}
\def\DAE{differential algebraic equations}
\def\FHJ{Feinberg-Horn-Jackson}
\def\GMAK{generalized mass action kinetics}
\def\MM{Michaelis-Menten}
\def\sm{stoichiometric matrix}\def\Sm{stoichiometric matrix}
\def\WRZD{weakly reversible, zero deficiency}\def\aut{autocatalytic}

% -- Epidemiology --
\def\brn{basic reproduction number}\def\bzn{basic replacement number}
\def\com{compartment}
\def\dfe{disease free equilibrium}\def\DFE{disease free equilibrium}
\def\dwt{dwell times}\def\edt{expected dwell times}
\def\EM{epidemic model}
\def\endp{endemic fixed point}\def\enp{endemic point}\def\EE{endemic point}
\def\ME{mathematical epidemiology}
\def\MEE{eco-epidemiology}\def\EEP{eco-epidemiological}
\def\ece{eco-epidemiological model}
\def\mgm{multi group model}
\def\mtS{multi-strain model}\def\mSt{multi-strain model}
\def\mSi{minimal siphon}
\def\NGM{next generation matrix}
\def\nim{new infections matrix}
\def\qem{quadratic epidemiological model}
\def\repF{invasion functions}
\def\SPF{SIR-PH-FA model}
\def\sus{susceptible}

% -- Graph theory & Linear algebra --
\def\jac{Jacobi matrix}
\def\jin{Jacobian matrix of the invasion vector field with respect to the invasion variables}
\def\jrn{Jacobian matrix of the resident vector field with respect to the resident variables}
\def\jxn{Jacobian matrix of the $\x$ vector field with respect to the $\x$ variables}
\def\Lm{Laplacian matrix}
\def\MTT{matrix tree theorem}
\def\SGT{spectral graph theory}
\def\SMW{Sherman–Morrison–Woodbury}

% -- Stability & Dynamics --
\def\absC{absorbtion CTMC}
\def\BTb{Bogdanov-Takens bifurcation}
\def\CEP{competitive exclusion partition}
\def\CTMC{continuous time Markov chain}
\def\ds{differential system}
\def\fp{fixed point}\def\bfp{boundary fixed point}\def\fbp{fixed boundary point}
\def\fbl{fixed boundary line}\def\fif{forward invariant face}
\def\gas{global asymptotic stability}\def\GAS{globally asymptotically stable}
\def\GASS{\text{GASS}}
\def\Hb{Hopf bifurcation}
\def\Hur{Hurwitz}\def\MH{Metzler-Hurwitz}\def\RH{Routh-Hurwitz}
\def\IS{\text{IS}}
\def\LAS{locally asymptotically stable}
\def\LdS{Lyapunov diagonally stable}
\def\Lf{Lyapunov function}\def\LyV{Lyapunov Volterra type function}\def\RLF{robust Lyapunov function}
\def\LyM{Lyapunov Malkin stable}
\def\PF{Perron-Frobenius}
\def\spfp{strictly positive fixed point}
\def\sta{stability}
\def\STP{sharp threshold property}\def\STT{sharp threshold theorem}

% -- Models (with \varphi explicitly instead of \f) --
\def\AB{$(A,B)$ Arino-Brauer epidemic models}
\def\ABa{$(A,\bbe,\ba,\varphi)$ balanced bilinear model}
\def\ABP{%$(A,\bbe,P,\varphi)$
\\ balanced bilinear model}
\def\AMG{Anderson May Gupta}
\newcommand\BIN{biological interaction network}
\def\Dt{Descartes type}
\def\gLV{generalized Lotka-Volterra model}
\def\IG{\text{invasion graph}}\def\IGMS{interaction graph on minimal siphons}
\def\LCP{\text{linear complementarity problem}}
\def\LV{Lotka-Volterra}\def\LVK{Lotka-Volterra-Kolmogorov}

% -- Splitting & Structure --
\def\reg{regular splitting}\def\regS{{\bf regular splitting}}
\def\RUR{rational univariate representation}
\def\bds{block-diagonal structure}\def\NGMs{next generation matrices}

% -- Miscellaneous --
\def\Ma{{Mathematica}}
\def\Mm{Max–Min}
\def\num{numerator}
\def\rap{rational approximation} \def\DSR{directed species-reaction bipartite graph}\def\PCP{polynomial complementarity problem}\def\rat{rational substitution}
\def\QSS{quasi steady state} \def\RHS{right hand side} \def\LHS{left hand side}
\def\OSN{online social network}\def\VdD{Van den Driessche}
\def\wli{$\omega$-limit} \def\rk{rich kinetics}\def\sJ{symbolic Jacobian}
\def\cV{\cite{Vas,VasSta,BSV,GSV}} \def\spf{semi-parametric semi-functional}
\def\CRW{Capasso-Ruan-Wang} \def\ABP{$(A,\bbe,P,\f)$ balanced bilinear model}
\def\ABa{%$(A,\bbe,\ba,\f)$ \\
rank-one balanced bilinear model}
\def\bbe{B%{\bff \beta}
} \def\vbe{\vec {\bff \beta}
}
\def\vb{\vec {\bff b}}
\newcommand{\Kh}{\hat{K}}          % K with hat, first mention only
\newcommand{\wS}{w^{\sigma}}       % right Perron J^\perp, full block
\newcommand{\piS}{\vec \pi^{\sigma}}    % left Perron J^\perp, full block
\newcommand{\ws}{w^{\sigma}}       % right Perron J^\perp, sub-block
\newcommand{\pis}{\pi^{\sigma}}    % left Perron J^\perp, sub-block
\newcommand{\hwS}{\hat{w}^{\sigma}}   % right Perron K, full block
\newcommand{\hpiS}{\hat{\pi}^{\sigma}} % left Perron K, full block
\newcommand{\hws}{\hat{w}^{\sigma}}   % right Perron K, sub-block
\newcommand{\hpis}{\hat{\pi}^{\sigma}} % left Perron K, sub-block
\newcommand{\twS}{\tilde{w}^{\sigma}}  % right Perron \tilde K, full block
\newcommand{\tpiS}{\tilde{\pi}^{\sigma}} % left Perron \tilde K
\newcommand{\tws}{\tilde{w}^{\sigma}}
\newcommand{\tpis}{\tilde{\pi}^{\sigma}}
\newcommand{\Mat}[2]{\mathbb{R}_{\ge 0}^{\,#1\times #2}}
\def\tJ{J_\sigma^\perp} \def\lPe{left Perron eigenvector}
% \G is \Gamma (defined at :795).  The Volterra function keeps its own name,
% so that \G no longer has two conflicting meanings in this file.
\def\GV{G}  % Volterra G(u)=u-1-\ln u
\def\PVL{Perron-Volterra Lyapunov minimal siphon ansatz}
\newcommand{\istar}[1]{i_{#1}^{*}}
%\newcommand{\sE}{s_E}
\newcommand{\vE}{v_E}
\def\PV{Perron-Volterra candidate Lyapunov function}
\def\GLV{generalized Lotka-Volterra}
\def\MERN{mathematical epidemiology reaction networks}
\def\CMM{contraction metric method}
\newcommand{\sign}{\operatorname{sign}}
\newcommand{\Var}{\operatorname{Var}} 

% ======================================================================
% IDEAS SHARED BETWEEN THE ABSTRACT AND THE MAIN RESULTS SECTION
%
% Write each idea ONCE, in the abstract; the Main Results section recalls it,
% so the two can never drift apart.
%
%   in the abstract        \ideaitem{1}{The first-order reaction terms ...}
%   in Main Results        \ideahead{1}
%
% \ideaitem stores the text globally and typesets
%   \item \hyperlink{idea:1}{<text>}
% \ideahead typesets
%   \item\hypertarget{idea:1}{}\textbf{<the same text>}
%
% The abstract must come before the Main Results input, which is the normal
% order. If a heading has to read differently from the abstract sentence (for
% instance a long abstract item summarised in one line), override it with the
% optional argument: \ideahead[short heading]{6}.
% ======================================================================
\makeatletter
\newcommand{\ideaitem}[2]{%
  \expandafter\gdef\csname @idea@#1\endcsname{#2}%
  \item\hyperlink{idea:#1}{#2}}
\newcommand{\ideahead}[2][]{%
  \item\hypertarget{idea:#2}{}%
  \def\@ideatmp{#1}%
  \ifx\@ideatmp\@empty
    \@ifundefined{@idea@#2}{%
      \PackageWarning{preamble}{Idea #2 is used in Main Results but was never
        defined by \string\ideaitem\space in the abstract}%
      \textbf{[idea #2 undefined]}%
    }{\textbf{\begingroup\def\ideaterm##1##2{##2}\def\idearest##1{}\@nameuse{@idea@#2}\endgroup}}%
  \else
    \textbf{#1}%
  \fi}
% \ideaterm{<hypertarget>}{<words>}: inside an abstract idea the words become an
% invisible link (the document sets \hypersetup{hidelinks}) jumping to the place
% in Main Results that defines them; inside the Main Results heading itself,
% where the target sits a few lines below, \ideahead strips the link and prints
% the words plainly. Mark the destination with \hypertarget{<name>}{} on the line
% before the display, which adds no equation number and shows nothing.
\newcommand{\ideaterm}[2]{\hyperlink{#1}{#2}}
% \idearest{<text>}: the part of an abstract idea after its first sentence; printed in
% the abstract, omitted from the Main Results heading.
\newcommand{\idearest}[1]{#1}
\makeatother

% ======================================================================
% SHORT-PAPER SWITCH
% \shorttrue in a master that assembles only part of the companion paper
% (Barrat/SipArxiv.tex). \nsh{...} drops text that cites results living in the
% omitted sections, so the short build has no dangling references while the
% companion paper keeps the text verbatim.
% ======================================================================
\newif\ifshort \shortfalse
\newcommand{\nsh}[1]{\ifshort\else#1\fi}
% \shortalt{<short build>}{<companion paper>}: use when a sentence needs
% different wording in the two builds, not merely a deletion.
\newcommand{\shortalt}[2]{\ifshort#1\else#2\fi}
% \Jtr, \JtrS{\Sigma}: transverse Jacobian. J^\perp here; SipArxiv.tex redefines them as J_{\mathrm{tra}}, J_{\mathrm{tra},\Sigma}
\newcommand{\Jtr}{J^\perp}
\newcommand{\JtrS}[1]{J_{#1}^\perp}
.
\makeatletter
\@ifpackageloaded{hyperref}{}{%
  \@ifundefined{hyperrefpackage}{%
    \@ifundefined{AAdefNnohyperref}{\usepackage{hyperref}}{}%
  }{}%
}
\makeatother
% linktoc=all makes a table-of-contents page number a link as well as the
% entry text; hyperref's default, linktoc=section, links only the text.
% Deferred to begin-document and guarded: by the block above hyperref may
% be loaded later (mdpi's own hook) or not at all (\AAdefNnohyperref), and
% \hypersetup must exist before it is called.  A driver that wants the
% default back can say \hypersetup{linktoc=section} after % Shared preamble for the GLVbook papers.  This is the single canonical copy;
% a driver one level below GLVbook reaches it as
%
%     \input{../basicFiles/preamble}
%
% What stays in the driver, because it differs from paper to paper:
% \documentclass, font choice, \title, \author, \date, and any macro genuinely
% local to one paper.  Everything else belongs here.  The file is
% self-contained: it does not read any other file.

\usepackage[utf8]{inputenc}
\usepackage[novbox]{pdfsync}   % required: SyncTeX alone does not work here

% Page geometry, set once here for every paper that loads this file.  A driver
% that needs different margins says \geometry{...}; a second
% \usepackage[...]{geometry} raises an option clash and must not be used.
\makeatletter
\@ifundefined{AAdefNnogeometry}{\usepackage[a4paper,margin=2cm]{geometry}}{}
\makeatother

% mathematics
\usepackage{amsmath,amssymb,amsthm,mathtools,bm,mathrsfs,bbm}
\usepackage{centernot}

% floats, tables, lists
\usepackage{graphicx,float}
\makeatletter
\@ifpackageloaded{subfig}{}{\usepackage{subcaption}}
\makeatother
\usepackage{booktabs,longtable,threeparttable,array,multirow}
\usepackage{enumitem}

% colour, code, algorithms
\usepackage{xcolor,listings,tcolorbox}
\makeatletter
\@ifundefined{c@algorithm}{\usepackage{algorithm}}{}
\makeatother
\usepackage{algorithmicx,algpseudocode}

% diagrams
\usepackage{tikz,tikz-cd,pgfplots}
\usetikzlibrary{graphs,decorations.pathmorphing,decorations.markings,
arrows.meta,automata,positioning}
\tikzcdset{every label/.append style = {font = \small}}
\pgfplotsset{compat=1.18}

% chemistry
\usepackage[version=4]{mhchem}

\usepackage{microtype}
\usepackage{lineno}
\usepackage{nomencl}
\makenomenclature

% hyperref first, then bookmark, then glossaries: glossaries must follow
% hyperref or the entry links are built against the wrong anchors.
\makeatletter
\@ifundefined{hyperrefpackage}{%
  \@ifundefined{AAdefNnohyperref}{\usepackage{hyperref}}{}%
}{}
\@ifundefined{hyperrefpackage}{\usepackage{bookmark}}{}
\makeatother
\usepackage[toc,sort=use]{glossaries}
\makeglossaries

% Book/standalone switch.  Material that the book already covers in its
% prelude, but that a standalone paper still has to state, is wrapped as
%     \ifbookbuild\else ... \fi
% Drivers default to standalone; bookM.tex says \bookbuildtrue.
\newif\ifbookbuild
\bookbuildfalse

% Part switches for the gradostat paper.  Both default to true; a driver that
% wants the short submission says \madecfalse and/or \slowfastfalse, and every
% cross-reference into the dropped part disappears with it.
%   \ifmadec     Part III, the Lagasquie--Madec model and its break-even geometry
%   \ifslowfast  Part IV, everything carrying the small parameter \eps
\newif\ifmadec     \madectrue
\newif\ifslowfast  \slowfasttrue

% Part switch for the Barrat/simplicial-contagion project.  Defaults to
% true (the hypernetwork-synchrony part is included); the journal-paper
% driver says \hnetpartfalse to drop it and every cross-reference into it.
\newif\ifhnetpart  \hnetparttrue

%% ========================================
%% PACKAGES
%% ========================================
\RequirePackage{etoolbox}
\usepackage{makecell}
\usepackage[novbox]{pdfsync}% obsolete; SyncTeX is built into modern engines
\usepackage{amssymb, amsfonts,mathtools,bbm,eurosym,graphicx,latexsym, amsmath,amsthm,times,nccmath,url}
\usepackage{listings,verbatim,cases,tabularx}
% hyperref is deliberately NOT in the line above.  Journal classes often
% load it themselves, with a long option list, from a begindocument hook
% (mdpi.cls does: \DeclareHookRule{begindocument}{hyperref}{before}{mdpi}).
% Loading it here first, without options, makes that later load raise
% "Option clash for package hyperref".  So load it only if nobody else
% has, or will:
%   - \@ifpackageloaded  catches a class that loaded it already;
%   - \hyperrefpackage    is mdpi.cls's own deferred loader;
%   - \AAdefNnohyperref   lets any other driver opt out by defining it
%                         before \input{preamble}.
\makeatletter
\@ifpackageloaded{hyperref}{}{%
  \@ifundefined{hyperrefpackage}{%
    \@ifundefined{AAdefNnohyperref}{\usepackage{hyperref}}{}%
  }{}%
}
\makeatother
% linktoc=all makes a table-of-contents page number a link as well as the
% entry text; hyperref's default, linktoc=section, links only the text.
% Deferred to begin-document and guarded: by the block above hyperref may
% be loaded later (mdpi's own hook) or not at all (\AAdefNnohyperref), and
% \hypersetup must exist before it is called.  A driver that wants the
% default back can say \hypersetup{linktoc=section} after \input{preamble}.
\makeatletter
\AtBeginDocument{\@ifundefined{hypersetup}{}{\hypersetup{linktoc=all}}}
\makeatother
\usepackage{graphicx,float,fancybox}
\usepackage{epsfig,epstopdf,tikz}
\usepackage[sans]{dsfont}
\usetikzlibrary{positioning,arrows,arrows.meta,calc}
% Configure listings for Mathematica code
\lstdefinestyle{mathematica}{
 language=Mathematica,
 basicstyle=\small\ttfamily,
 backgroundcolor=\color{gray!10},
 frame=single,
 breaklines=true,
 commentstyle=\color{green!60!black},
 keywordstyle=\color{blue},
 stringstyle=\color{red},
 showstringspaces=false,
 tabsize=2
}
\usepackage{tcolorbox}
% Colored boxes for key results
\newtcolorbox{keyresult}[1][]{
 colback=blue!5!white,
 colframe=blue!75!black,
 title=#1,
 fonttitle=\bfseries
}
\newtcolorbox{examplebox}[1][]{
 colback=green!5!white,
 colframe=green!75!black,
 title=#1,
 fonttitle=\bfseries
}


\makeatletter
\@ifundefined{theorem}   {\newtheorem{theorem}   {Theorem}}   {}
\@ifundefined{definition}{\newtheorem{definition}{Definition}} {}
\@ifundefined{open}      {\newtheorem{open}      {Problem}}   {}
\@ifundefined{example}   {\newtheorem{example}   {Example}}   {}
\makeatother

% COUNTER DEFINITIONS
\newcounter{tabcounter}
\newcounter{defcounter}
\newcounter{opcounter}
\newcounter{thmcounter}
\newcounter{excounter}
\newcounter{ascounter}
\newcounter{cjcounter}
\newcounter{lemcounter}
\newcounter{remcounter}
\newcounter{corcounter}
\newcounter{propcounter}
\newcounter{figcounter}

% UNIFIED GLOSSARY STYLE
\newglossarystyle{unifiedstyle}{%
  \setglossarystyle{list}%
  \renewcommand*{\glossarypreamble}{\gdef\glslastsec{}}%
  \renewcommand{\glossentry}[2]{%
    \edef\glscursec{\glsentryuseriii{##1}}%
    \ifx\glscursec\glslastsec\else
      \item[]\textbf{\S\,\glscursec\enspace\glsentryuseriv{##1}}%
      \global\let\glslastsec\glscursec
    \fi
    \item[\glsentryuserii{##1} \glsentryuseri{##1}:] ``\glossentryname{##1}'' \glsentrydesc{##1}%
  }%
}

\makeatletter
% First-word extractor (used by \beT, \beD and \beXa for the default key)
% \@firstwordof hello world \relax  ->  hello
\def\@firstwordof#1 #2\relax{#1}
% First three letters, used for the heading mnemonic of a theorem given no
% explicit key.  Padding with \@empty keeps it safe for first words shorter
% than three letters: \@empty vanishes under \edef, whereas \relax would
% survive into the label and is illegal inside \csname (it broke every
% two-letter title, e.g. "ME ...", with "Missing \endcsname inserted").
\def\@takethree#1#2#3#4\@nil{#1#2#3}

% EXAMPLE COMMAND -- same signature and behaviour as \beT and \beD.
% Usage: \beXa*[long title][short title][key]
%   star  = suppresses the printed "[key]" tag after the title (as for \beT*
%   and \beD*); the label and glossary entry are unaffected.
%   key = optional label key override; the label becomes e:KEY, and without a
%   key it is e: plus the first three letters of the title's first word.  A
%   mnemonic label lets an example be cited as "Example~\ref{e:KEY}" instead
%   of by section number, which is what the citing prose needs.
\NewDocumentCommand{\beXa}{s O{} O{} o}{%
  \def\@temparg{#2}%
  \ifx\@temparg\@empty
    \begin{example}%
  \else
    \let\@savedSecTitle\@currentsectionname%
    \refstepcounter{excounter}%
    \IfValueTF{#4}{%
      \edef\@exlblkey{#4}%
    }{%
      \edef\@exfirstword{\@firstwordof#2 \relax}%
      \edef\@exlblkey{%
        \expandafter\@takethree\@exfirstword\@empty\@empty\@empty\@nil}%
    }%
    \def\@tempshort{#3}%
    \ifx\@tempshort\@empty
      \expandafter\xdef\csname @xshort@\@exlblkey\endcsname{\@exlblkey}%
    \else
      \expandafter\def\csname @xshort@\@exlblkey\endcsname{#3}%
    \fi
    \IfBooleanTF{#1}{\let\@extag\@empty}{%
      \edef\@extag{\noexpand\ \noexpand\textup{[\noexpand\texttt{e:\@exlblkey}]}}}%
    \begin{example}[#2\ifshowkeys\@extag\fi]%
    \label{e:\@exlblkey}%
    % record the long title so \nref{e:KEY} prints it, as for theorems
    \protected@write\@auxout{}{%
      \string\@thmtitle{e:\@exlblkey}{\unexpanded{#2}}}%
    \expanded{%
      \noexpand\newglossaryentry{e\theexcounter}{%
        name={\unexpanded{#2}%
          \IfBooleanF{#1}{ [\noexpand\texttt{e:\@exlblkey}]}},%
        sort={\theexample},%
        description={p.~\noexpand\pageref{e:\@exlblkey}},%
        user1={\theexample},%
        user2={Example},%
        user3={\thesection},%
        user4={\expandafter\unexpanded\expandafter{\@savedSecTitle}}%
      }%
    }%
    \protected@write\@auxout{}{\string\glsliveentry{e\theexcounter}}%
  \fi
}
\newcommand{\eeXa}{\end{example}}

% DEFINITION COMMAND -- same signature and behaviour as \beT below.
% Usage: \beD[long title][short title][key]
%   long title  = definition heading (plain text only)
%   short title = shown in \dref{word}; defaults to the key
%   key         = optional label key override; the label becomes d:KEY
% Without a key the label is d: plus the first three letters of the title's
% first word, so it is a mnemonic rather than the definition's number: a
% numeric label would silently change whenever a definition is inserted.
\NewDocumentCommand{\beD}{s O{} O{} o}{%
  \def\@temparg{#2}%
  \ifx\@temparg\@empty
    \begin{definition}%
  \else
    \let\@savedSecTitle\@currentsectionname%
    \refstepcounter{defcounter}%
    \IfValueTF{#4}{%
      \edef\@deflblkey{#4}%
    }{%
      \edef\@deffirstword{\@firstwordof#2 \relax}%
      \edef\@deflblkey{%
        \expandafter\@takethree\@deffirstword\@empty\@empty\@empty\@nil}%
    }%
    \def\@tempshort{#3}%
    \ifx\@tempshort\@empty
      \expandafter\xdef\csname @dshort@\@deflblkey\endcsname{\@deflblkey}%
    \else
      \expandafter\def\csname @dshort@\@deflblkey\endcsname{#3}%
    \fi
    \IfBooleanTF{#1}{\let\@deftag\@empty}{%
      \edef\@deftag{\noexpand\ \noexpand\textup{[\noexpand\texttt{d:\@deflblkey}]}}}%
    \begin{definition}[#2\ifshowkeys\@deftag\fi]%
    \label{d:\@deflblkey}%
    % record the long title so \nref{d:KEY} prints it, as for theorems
    \protected@write\@auxout{}{%
      \string\@thmtitle{d:\@deflblkey}{\unexpanded{#2}}}%
    \expanded{%
      \noexpand\newglossaryentry{d\thedefcounter}{%
        name={\unexpanded{#2}%
          \IfBooleanF{#1}{ [\noexpand\texttt{d:\@deflblkey}]}},%
        sort={\thedefinition},%
        description={p.~\noexpand\pageref{d:\@deflblkey}},%
        user1={\thedefinition},%
        user2={Definition},%
        user3={\thesection},%
        user4={\expandafter\unexpanded\expandafter{\@savedSecTitle}}%
      }%
    }%
    \protected@write\@auxout{}{\string\glsliveentry{d\thedefcounter}}%
  \fi
}
\newcommand{\eeD}{\end{definition}}

% \showkeysfalse in a driver hides the [t:KEY] mnemonics in result headings
% (labels and glossary entries are unchanged); used for submission builds.
\newif\ifshowkeys \showkeystrue
% \commonredtrue in a driver that inputs basicFiles/commonRed: tool.tex then
% cites its reduction results instead of restating them.
\newif\ifcommonred \commonredfalse
% TheorEM COMMAND
% Usage: \beT[long title][short title][key]
%   long title  = theorem heading (plain text only)
%   short title = shown in \tref{word}; may contain math; defaults to first word of long title
%   key         = optional label key override; label becomes t:KEY (avoids first-word collisions)
% Label is t:WORD where WORD = first word of long title (or KEY if given); use \tref{WORD} to reference.
\NewDocumentCommand{\beT}{s O{} O{} o}{%
  \def\@temparg{#2}%
  \ifx\@temparg\@empty
    \begin{theorem}%
  \else
    \let\@savedSecTitle\@currentsectionname%
    \refstepcounter{thmcounter}%
    % \@thmtag is the mnemonic shown in the heading; every theorem gets one.
    % With an explicit key the key itself is shown.  Without one the key
    % defaults to the title's first word, and showing that whole word would
    % merely repeat the title ("Dissipativity [Dissipativity]"), so only its
    % first three letters are shown: "Dissipativity [Dis]".  The label is
    % unaffected -- it still uses the full \@thmlblkey.
    \IfValueTF{#4}{%
      \edef\@thmlblkey{#4}%
    }{%
      \edef\@thmfirstword{\@firstwordof#2 \relax}%
      \edef\@thmlblkey{%
        \expandafter\@takethree\@thmfirstword\@empty\@empty\@empty\@nil}%
    }%
    \IfBooleanTF{#1}{\let\@thmtag\@empty}{%
      \edef\@thmtag{\noexpand\ \noexpand\textup{[\noexpand\texttt{t:\@thmlblkey}]}}}%
    \def\@tempshort{#3}%
    \ifx\@tempshort\@empty
      \expandafter\xdef\csname @tshort@\@thmlblkey\endcsname{\@thmlblkey}%
    \else
      \expandafter\def\csname @tshort@\@thmlblkey\endcsname{#3}%
    \fi
    \begin{theorem}[#2\ifshowkeys\@thmtag\fi]%
    \label{t:\@thmlblkey}%
    % record the long title for \nref; via the .aux so forward refs resolve
    \protected@write\@auxout{}{%
      \string\@thmtitle{t:\@thmlblkey}{\unexpanded{#2}}}%
    \expanded{%
      \noexpand\newglossaryentry{t\thethmcounter}{%
        name={\unexpanded{#2}%
          \IfBooleanF{#1}{ [\noexpand\texttt{t:\@thmlblkey}]}},%
        % sort/user1 must be the number the document actually prints, i.e. the
        % theorem environment's counter, not thmcounter.  thmcounter advances
        % only on \beT with a title, so a bare \beT (or \beT with an
        % empty title) makes the two drift and the glossary cites a wrong
        % number.  thmcounter is still fine as the entry key: it is unique.
        sort={\thetheorem},%
        description={p.~\noexpand\pageref{t:\@thmlblkey}},%
        user1={\thetheorem},%
        user2={Theorem},%
        user3={\thesection},%
        user4={\expandafter\unexpanded\expandafter{\@savedSecTitle}}%
      }%
    }%
    \protected@write\@auxout{}{\string\glsliveentry{t\thethmcounter}}%
  \fi
}
\newcommand{\eeT}{\end{theorem}}

% CROSS-PART RELATION PARAGRAPH
% Usage: \xpart{text}  -- place immediately BEFORE \beT, never inside it
%   (anything inside \beT...\eeT would be swallowed into the glossary entry).
% The paragraph is typeset only when Part I is actually present, i.e. in the
% combined book build; in the standalone builds (GranadaDSABNS/A.tex,
% AAHRah-Submission/AAHRah.tex) the label part:glv is undefined and the whole
% paragraph disappears rather than leaving ?? cross-references. Same
% graceful-degradation principle as \GASref in GASt.tex. Resolves on the
% second LaTeX pass, since r@part:glv comes from the .aux file.
% (Already inside the \makeatletter region opened above, so no \makeatletter
% pair here -- adding one would close that region early.)
\newcommand{\xpart}[1]{\@ifundefined{r@part:glv}{}{\paragraph{Relation with Parts~I--II.}#1}}

% \bookref{label}{precise text}{fallback text}: prints the precise cross-reference
% when the label resolves (combined book build), and the fallback text otherwise
% (standalone paper builds, where book-only appendix labels are absent). Same
% graceful-degradation role as \GASref in GASt.tex, but usable in any file since it
% lives in this preamble (loaded by every driver) -- unlike \GASref, which is defined in a
% content file and is therefore undefined in standalone builds that omit that file.
\newcommand{\bookref}[3]{\@ifundefined{r@#1}{#3}{#2}}

% \PVpart and \RBpart: the two parts that a partial build may omit.
% They print "Part~N" in any build that contains the part, and a naming phrase
% in a build that does not, so an excerpt shows no ?? .  Each already supplies
% the word "Part", so write "\PVpart develops ...", not "Part~\PVpart".
\newcommand{\PVpart}{\bookref{part:aahrah}%
  {Part~\ref{part:aahrah}}{the Perron--Volterra part of the full book}}
\newcommand{\RBpart}{\bookref{part:recbrowne}%
  {Part~\ref{part:recbrowne}}{the immuno-virology part of the full book}}

% OPEN PROBLEM COMMANDS
% ASSUMPTION COMMANDS (glossaried)
\newcommand{\beAs}[1][]{%
  \def\@temparg{#1}%
  \ifx\@temparg\@empty
    \begin{assumption}%
  \else
    \let\@savedSecTitle\@currentsectionname%
    \refstepcounter{ascounter}%
    \begin{assumption}[#1]%
    \label{as:\theascounter}%
    \expanded{%
      \noexpand\newglossaryentry{as\theascounter}{%
        name={\unexpanded{#1}},%
        sort={\theascounter},%
        description={p.~\noexpand\pageref{as:\theascounter}},%
        user1={\theascounter},%
        user2={Assumption},%
        user3={\thesection},%
        user4={\expandafter\unexpanded\expandafter{\@savedSecTitle}}%
      }%
    }%
    \protected@write\@auxout{}{\string\glsliveentry{as\theascounter}}%
  \fi
}
\newcommand{\eeAs}{\end{assumption}}
% CONJECTURE COMMANDS (glossaried)
\newcommand{\beCj}[1][]{%
  \def\@temparg{#1}%
  \ifx\@temparg\@empty
    \begin{conjecture}%
  \else
    \let\@savedSecTitle\@currentsectionname%
    \refstepcounter{cjcounter}%
    \begin{conjecture}[#1]%
    \label{cj:\thecjcounter}%
    \expanded{%
      \noexpand\newglossaryentry{cj\thecjcounter}{%
        name={\unexpanded{#1}},%
        sort={\thecjcounter},%
        description={p.~\noexpand\pageref{cj:\thecjcounter}},%
        user1={\thecjcounter},%
        user2={Conjecture},%
        user3={\thesection},%
        user4={\expandafter\unexpanded\expandafter{\@savedSecTitle}}%
      }%
    }%
    \protected@write\@auxout{}{\string\glsliveentry{cj\thecjcounter}}%
  \fi
}
\newcommand{\eeCj}{\end{conjecture}}
% LEMMA COMMANDS
% ---------------------------------------------------------------------
% UNIFORM RESULT WRAPPERS -- \beL, \beC, \beP, \beR share \beT's signature
%
%   \beX[long title][short title][key]  ...  \eeX
%
%   long title  = heading (math allowed)
%   short title = shown in the cross-reference; defaults to the key
%   key         = label key.  With a key the label is X:KEY and the heading
%                 carries the mnemonic [key], exactly as for \beT.  Without
%                 one the label falls back to the running counter, which is
%                 what every pre-existing call site in the book relies on.
%
% A bare \beX -- arguments commented out -- opens the plain environment and
% creates NO label and NO glossary entry.  That is the supported way to keep
% an individual result out of the catalogue.
% ---------------------------------------------------------------------
% #1 wrapper name  #2 environment  #3 label prefix  #4 counter  #5 type name
\newcommand{\@mkresult}[5]{%
  % A leading * suppresses the [key] mnemonic in both the heading and the
  % glossary name, while still making the label X:KEY and the entry:
  %     \beX*[long][short][key]
  \expandafter\NewDocumentCommand\csname be#1\endcsname{s O{} O{} o}{%
    \def\@temparg{##2}%
    \ifx\@temparg\@empty
      \begin{#2}%
    \else
      \let\@savedSecTitle\@currentsectionname%
      \refstepcounter{#4}%
      % With a key: mnemonic label X:KEY and the [key] tag in the heading.
      % Without one: fall back to the running counter, exactly as these
      % wrappers behaved before, so every legacy one-argument call site in
      % the book keeps its label.  The counter is also csname-safe, which a
      % title's first word is not -- titles beginning with math would
      % otherwise build \csname \Omega...\endcsname and fail.
      \IfValueTF{##4}{%
        \edef\@relbl{##4}%
        \IfBooleanTF{##1}{\let\@retag\@empty}{%
          \edef\@retag{\noexpand\ \noexpand\textup{[\noexpand\texttt{#3:##4}]}}}%
      }{%
        \edef\@relbl{\arabic{#4}}%
        \let\@retag\@empty
      }%
      \def\@tempshort{##3}%
      \ifx\@tempshort\@empty
        \expandafter\xdef\csname @sshort@#3\@relbl\endcsname{\@relbl}%
      \else
        \expandafter\def\csname @sshort@#3\@relbl\endcsname{##3}%
      \fi
      \begin{#2}[##2\ifshowkeys\@retag\fi]%
      \edef\@tmplbl{#3:\@relbl}\expandafter\label\expandafter{\@tmplbl}%
      \expanded{%
        \noexpand\newglossaryentry{#3\csname the#4\endcsname}{%
          name={\unexpanded{##2}%
            \IfBooleanF{##1}{\IfValueT{##4}{ [\noexpand\texttt{#3:##4}]}}},%
          sort={\csname the#2\endcsname},%
          description={p.~\noexpand\pageref{#3:\@relbl}},%
          user1={\csname the#2\endcsname},%
          user2={#5},%
          user3={\thesection},%
          user4={\expandafter\unexpanded\expandafter{\@savedSecTitle}}%
        }%
      }%
      \protected@write\@auxout{}{\string\glsliveentry{#3\csname the#4\endcsname}}%
      \protected@write\@auxout{}{\string\@thmtitle{\@tmplbl}{\unexpanded{##2}}}%
    \fi
  }%
  \expandafter\newcommand\csname ee#1\endcsname{\end{#2}}%
}
\@mkresult{L}{lemma}{l}{lemcounter}{Lemma}
% REMARK / COROLLARY / PROPOSITION -- same signature as \beL above
\@mkresult{R}{remark}{r}{remcounter}{Remark}
\@mkresult{C}{corollary}{c}{corcounter}{Corollary}
\@mkresult{P}{proposition}{p}{propcounter}{Proposition}
% OPEN-PROBLEM COMMAND -- generated with the same signature as \beT:
%   \beO[long title][short title][key]   ...   \eeO
% label o:KEY, or o:<counter> when no key is given; a bare \beO makes
% neither a label nor a glossary entry.
\@mkresult{O}{open}{o}{opcounter}{Problem}


% Section-level title tracker: updated only by \section, not \subsection.
% \@currentlabelname is overwritten by every \refstepcounter (including
% \subsection), so we save it separately here.
\newcommand{\@currentsectionname}{}
\AtBeginDocument{%
  \let\@epidcrn@orig@refstepcounter\refstepcounter
  \renewcommand{\refstepcounter}[1]{%
    \@epidcrn@orig@refstepcounter{#1}%
    \ifstrequal{#1}{section}{\global\let\@currentsectionname\@currentlabelname}{}%
  }%
}

% TABLE COMMANDS (glossaried) -- consistent naming with \beT, \beD, etc.
\newcommand{\beTa}[1][]{%
  \def\@temparg{#1}%
  \ifx\@temparg\@empty
    \begin{table}[htbp]%
  \else
    \let\@savedSecTitle\@currentsectionname%
    \refstepcounter{tabcounter}%
    \begin{table}[htbp]%
    \label{A:\thetabcounter}%
    \expanded{%
      \noexpand\newglossaryentry{A\thetabcounter}{%
        name={\unexpanded{#1}},%
        sort={\thetabcounter},%
        description={p.~\noexpand\pageref{A:\thetabcounter}},%
        user1={\thetabcounter},%
        user2={Table},%
        user3={\thesection},%
        user4={\expandafter\unexpanded\expandafter{\@savedSecTitle}}%
      }%
    }%
    \protected@write\@auxout{}{\string\glsliveentry{A\thetabcounter}}%
  \fi
}
\newcommand{\eeTa}{\end{table}}

% FIGURE COMMANDS (glossaried) -- consistent with \beT, \beR, \beTa, etc.
% Usage:
%   \beF[Glossary title]...\caption{Visible caption}\eeF  -> glossaried figure
%   \beF...\caption{Visible caption}\eeF                  -> plain [H] figure, no glossary
\newcommand{\beF}[1][]{%
  \def\@temparg{#1}%
  \ifx\@temparg\@empty
    \begin{figure}[H]%
  \else
    \begin{figure}[H]%
    \let\@savedSecTitle\@currentsectionname%
    \refstepcounter{figcounter}%
    \label{fig:\thefigcounter}%
    \expanded{%
      \noexpand\newglossaryentry{Fig\thefigcounter}{%
        name={\unexpanded{#1}},%
        sort={\thefigcounter},%
        description={p.~\noexpand\pageref{fig:\thefigcounter}},%
        user1={\thefigcounter},%
        user2={Figure},%
        user3={\thesection},%
        user4={\expandafter\unexpanded\expandafter{\@savedSecTitle}}%
      }%
    }%
    \protected@write\@auxout{}{\string\glsliveentry{Fig\thefigcounter}}%
  \fi
}
\newcommand{\eeF}{\end{figure}}
\makeatother
% LIVE-ENTRY REGISTER
% \glsaddall prints every entry the glossaries package currently knows, and
% that set includes entries restored from \jobname.glsdefs.  The package
% rewrites .glsdefs from all *currently defined* entries, so anything it
% restored is written straight back out: a deleted section kept reappearing in
% the catalogue for ever, and the .glsdefs had to be erased by hand.  The
% wrappers therefore record every entry they actually create in the .aux, and
% \printboth adds only those.  The .aux is rebuilt from scratch on each run, so
% this is self-healing: a removed result is gone after one further run.
\makeatletter
\newcommand{\@glslive}{}
\newcommand{\@glsaddlive}[1]{\ifglsentryexists{#1}{\glsadd{#1}}{}}
\newcommand{\glsliveentry}[1]{\g@addto@macro\@glslive{\@glsaddlive{#1}}}
\makeatother
% PRINT COMMAND
% No \clearpage here -- the catalogue continues on the current page.  This has
% been re-added by mistake several times; leave it out.
\makeatletter
\newcommand{\printboth}{%
  \@glslive   % add exactly the entries this document created last run
  \printglossary[style=unifiedstyle, title={0 \hspace{.22em} Definitions, Theorems, Examples and Open Problems}]%
}
\makeatother
% Named theorem reference: \tref{word} -> Thm~N(word-or-short)
% Shows stored short title if available, otherwise shows the key itself.
\newcommand{\tref}[1]{%
  Thm~\ref{t:#1}(\ifcsname @tshort@#1\endcsname\csname @tshort@#1\endcsname\else#1\fi)%
}

% Full-title theorem reference: \nref{t:KEY} prints the number and then the
% theorem's long title, so one writes "Theorem~\nref{t:blocks}" and gets
% "Theorem 3 ``Block factorization and siphons''".  Takes the complete label, like
% \ref, not the bare key like \tref.
%
% \beT writes each title to the .aux (see \@thmtitle below), so \nref resolves
% forward references from the second LaTeX pass, exactly as \ref does.  If the
% title is not yet known -- first pass, or a label that is not a \beT theorem
% -- it degrades to the number rather than producing ??.
% These two use @ in their names, and this point in the file is outside the
% \makeatletter region opened earlier (it closes at \makeatother above), so
% the region has to be reopened here.
\makeatletter
\newcommand{\@thmtitle}[2]{\global\@namedef{@tlong@#1}{#2}}
\newcommand{\nref}[1]{%
  \ref{#1}%
  \ifcsname @tlong@#1\endcsname
    \ ``\csname @tlong@#1\endcsname''%
  \fi
}
% \nm{t:KEY}: the stored long title of any \beT/\beL/\beP/\beR/\beC/\beD/\beE result,
% verbatim and without quotes, so item names follow the macro's first argument.
\newcommand{\nm}[1]{%
  \ifcsname @tlong@#1\endcsname
    \csname @tlong@#1\endcsname%
  \else
    \ref{#1}%
  \fi
}
\makeatother




%% ========================================
%% Other THEOREM ENVIRONMENTS
%% ========================================
\makeatletter
\@ifundefined{lemma}      {\newtheorem{lemma}      {Lemma}}      {}
\@ifundefined{proposition}{\newtheorem{proposition}{Proposition}} {}
\@ifundefined{corollary}  {\newtheorem{corollary}  {Corollary}}  {}
\@ifundefined{conjecture} {\newtheorem{conjecture} {Conjecture}}  {}
\@ifundefined{remark}     {\newtheorem{remark}     {Remark}}     {}
\@ifundefined{question}   {\newtheorem{question}   {Question}}   {}
\@ifundefined{assumption} {\newtheorem{assumption} {Assumption}} {}
\@ifundefined{hypothesis} {\newtheorem{hypothesis} {Hypothesis}} {}
\@ifundefined{fact}       {\newtheorem{fact}       {Fact}}       {}
\makeatother

% \citer: cite the relay-graph threshold theorem of AAHrel. If a document has
% \input the chapter/file containing t:relayRF (so the label resolves), point
% straight at it; otherwise fall back to a bare citation of AAHrel. Shared here
% (not per-paper) so it resolves consistently across every book chapter.
% \citeR: same idea, for the first theorem of AAHRah (t:Pstar).
\makeatletter
\providecommand{\citer}{%
  \@ifundefined{r@t:relayRF}%
    {\cite{AAHrel}}%
    {Theorem~\ref{t:relayRF} of \cite{AAHrel}}%
}
\providecommand{\citeR}{%
  \@ifundefined{r@t:Pstar}%
    {\cite{AAHRah}}%
    {Theorem~\ref{t:Pstar} of \cite{AAHRah}}%
}
\makeatother


% Abbreviations
\def\beQ{\begin{question}}\def\eeQ{\end{question}}
\def\beA{\begin{assumption}}\def\eeA{\end{assumption}}


%% ========================================

%% ========================================
%% COLORS & CUSTOM COMMANDS
%% ========================================
\definecolor{funccolor}{RGB}{25,25,112}
\definecolor{desccolor}{RGB}{64,64,64}
\newcommand{\func}[1]{\texttt{\textcolor{funccolor}{#1}}}
\def\eq{=}
\def\mm{-}

%% ========================================
%% MATHCAL SHORTHANDS
%% ========================================
% Project-wide convention: \mX for a capital X is always \mathcal X.
% Defined for the whole alphabet so the convention never fails to expand.
\newcommand{\mA}{\mathcal A}\newcommand{\mB}{\mathcal B}
\newcommand{\mC}{\mathcal C}\newcommand{\mD}{\mathcal D}
\newcommand{\mE}{\mathcal E}\newcommand{\mF}{\mathcal F}
\newcommand{\mG}{\mathcal G}\newcommand{\mH}{\mathcal H}
\newcommand{\mI}{\mathcal I}\newcommand{\mJ}{\mathcal J}
\newcommand{\mK}{\mathcal K}\newcommand{\mL}{\mathcal L}
\newcommand{\mM}{\mathcal M}\newcommand{\mN}{\mathcal N}
\newcommand{\mO}{\mathcal O}\newcommand{\mP}{\mathcal P}
\newcommand{\mQ}{\mathcal Q}\newcommand{\mR}{\mathcal R}
\newcommand{\mS}{\mathcal S}\newcommand{\mT}{\mathcal T}
\newcommand{\mU}{\mathcal U}\newcommand{\mV}{\mathcal V}
\newcommand{\mW}{\mathcal W}\newcommand{\mX}{\mathcal X}
\newcommand{\mY}{\mathcal Y}\newcommand{\mZ}{\mathcal Z}

%% ========================================
%% BASIC MATH ENVIRONMENTS
%% ========================================
\def\bep{\begin{pmatrix}}\def\eep{\end{pmatrix}}
\def\bev{\begin{vmatrix}}\def\eev{\end{vmatrix}}
\def\bea{\begin{eqnarray*}}\def\eea{\end{eqnarray*}}
\def\bc{\begin{cases}}\def\ec{\end{cases}}
\def\BEN{\begin{enumerate}}\def\EEN{\end{enumerate}}
\def\BI{\begin{itemize}}\def\EI{\end{itemize}}

\newcommand{\be}[1]{\begin{equation}\label{#1}}
\newcommand{\ee}{\end{equation}}
\newcommand{\beq}{\begin{eqnarray}}
\def\eeq{\end{eqnarray}}
\def\eqr{\eqref}\def\fr{\frac}\def\lbl{\label}
\def\Lra{\Longrightarrow}\def\non{\nonumber}\def\cNGM{\cite{Diek,Van,Van08}}
\def\Eq{\Leftrightarrow}\def\b{\beta}\def\uph{\red{UP TO HERE}}
%% ========================================
%% PROBABILITY & STATISTICS
%% ========================================
\renewcommand{\P}{\ensuremath{\mathbb{P}}}  % Note: overrides paragraph symbol
\newcommand{\E}{\ensuremath{\mathbb{E}}}
\newcommand{\ind}[1]{\ensuremath{\mathbbm{1}_{\left\{#1\right\}}}}
\DeclareMathOperator{\Cov}{Cov}
\DeclareMathOperator{\TV}{TV}
\newcommand{\rank}{\operatorname{rank}}

%% ========================================
%% CRN OPERATORS (organizations/ERCs/buffering/siphons framework)
%% ========================================
\def\clos{\mathrm{clos}}
\def\Supply{\mathrm{Supply}}
\def\req{\mathrm{req}}
\def\React{\mathrm{React}}
\def\Prod{\mathrm{Prod}}
\def\Gen{\mathrm{Gen}}
\def\In{\mathrm{In}}
\def\Out{\mathrm{Out}}
\def\Act{\mathrm{Act}}
\def\Min{\mathrm{Min}}

%% ========================================
%% LINE BREAKING
%% ========================================
% Last-resort stretch, used only on paragraphs TeX cannot otherwise break:
% a long \texttt{} token, a comma-separated inline math list (commas are not
% breakpoints in math), or a run of \cite/\ref boxes.  Without it those set as
% overfull \hbox and print into the margin; with it they set slightly loose,
% which is the lesser evil.  Does not affect paragraphs that already break.
\emergencystretch=3em

%% ========================================
%% NUMBER SETS
%% ========================================
\newcommand{\R}{\mathbb{R}}  % Real numbers
\newcommand{\C}{\mathbb{C}}  % Complex numbers
\newcommand{\N}{\mathbb{N}}
\newcommand{\Z}{\mathbb{Z}}
% \providecommand, not \newcommand: some papers define these in their own
% preamble after \input{preamble}, and must not clash.
\providecommand{\Q}{\mathbb{Q}}
\providecommand{\res}{\operatorname{res}}   % resultant; used by epiRes.tex
\newcommand{\MN}{\mathbb{N}}
\newcommand{\MZ}{\mathbb{Z}}
\newcommand{\realpos}{\mathbb{R}_{> 0}}
\newcommand{\realnneg}{\mathbb{R}_{\geq 0}}
\newcommand{\naturalpos}{\mathbb{N}_{>0}}
\newcommand{\Rplus}{\mathbb{R}_+}
\newcommand{\Rp}{\mathbb{R}_{\ge 0}}
\def\Rnne{\mathbb{R}_{\ge 0}}
\def\Se{\bar{S}} \def\Ie{\bar{I}}
\newcommand{\Rpos}{\mathbb{R}_{>0}}
\newcommand{\dd}{\,\mathrm d}
\newcommand{\e}{\mathrm e}
%% ========================================
%% GREEK LETTERS (2+ character abbreviations) greek
%% ========================================
\def\al{\alpha}\def\ba{\bff \alpha}
\def\w{\omega}\def\Lgn{\ga+\mu+\de}\def\om{\omega}
\def\de{\delta}\def\De{\Delta}\def\bde{\bff \delta}
\def\ep{\epsilon}\def\eps{\varepsilon}
\def\f{\varphi}  % Kept as requested
\def\ga{\gamma}\def\Ga{\Gamma}\def\G{\Gamma}\def\bg{\bff \gamma}
\def\up{\upsilon}  % NOT the breve accent: no accent commands anywhere.
                   % Named \up, not \u, so the LaTeX breve stays reachable.
\def\la{\lambda}\def\La{\Lambda}
\def\si{\sigma}\def\ka{\kappa}
\def\th{\theta}\def\vt{\vartheta}
\def\ze{\zeta}\def\Gap{\Gamma^{+}} \def\Gam{\Gamma^{-}}


%% ========================================
%% SPECIAL SHORT COMMANDS (kept as requested)
%% ========================================
\def\I{\infty}
\def\T{\widetilde}
\def\H{\hat}
\newcommand{\relint}{\operatorname{relint}}

%% ========================================
%% BOLD/VECTOR NOTATION
%% ========================================
\newcommand{\bff}[1]{{\mbox{\boldmath$#1$}}}
\newcommand\X{\boldsymbol{X}}
\def\x{\boldsymbol{x}}\def\bx{\bff x}\def\xe{\bar{\x}}
\newcommand\y{\boldsymbol{y}}\def\by{\boldsymbol{y}}\def\ye{\bar{\y}}
\def\yb{\y}\def\xb{\x}
\def\xy{\x-\y}
\newcommand\br{\boldsymbol{r}}\def\re{\bar{\r}}
\newcommand\bd{\boldsymbol{d}}
\newcommand\bQ{\boldsymbol{Q}}
\newcommand{\Yb}{{\mathbf{Y}}}
\newcommand{\Ab}{\mathbf{A}}
\newcommand{\zb}{\mathbf{z}}\def\bz{\bff 0}
\newcommand{\Zb}{\mathbf{Z}}
\newcommand{\bb}{\mathbf{b}}
\def\vbe{\vec {\bff \beta}}
\def\vb{\vec {\bff b}}
\def\v1{\vec {\bff 1}}
\def\vn{\vec \de}
\def\vz{\vec 0}
\def\vi{\overset{\rightarrow}{\mathsf i}}
\def\vv{\overset{\rightarrow}{v}}
\def\vr{\; \vec {\mathsf r}}
\newcommand{\va}{\bff a}
\providecommand{\vect}[1]{\boldsymbol{#1}}   % bold vector (moved from Gradostat.tex)
\providecommand{\one}{\mathbf 1}             % all-ones vector (moved from AliSF)
\providecommand{\spr}{s}                     % spectral bound  (moved from AliSF)
\providecommand{\Ke}{K_{\rm eff}}            % effective carrying capacity
% Critical right/left eigenvectors in the Hopf projection method (AAHOSN/New.tex)
\providecommand{\bq}{\mathbf{q}}
\providecommand{\bp}{\mathbf{p}}


%% ========================================
%% MATHCAL & MATHSCR
%% ========================================
% Mathcal
\newcommand\mc[1]{\mathcal{#1}}
\def\mC{{\mathcal C}}\def\mD{\mathcal D}\def\mE{\mathcal{E}}
\def\mF{\mathcal F}\def\mG{{\mathcal G}}\def\mH{{\mathcal H}}
\def\mI{\mathcal I}\def\mL{\mathcal{L}}\def\mM{{\mathcal M}}
\def\mN{{\mathcal N}}\def\mP{\mathcal{P}}\def\mR{{\mathcal R}}
\def\mS{{\mathcal S}}\def\mT{\mathcal{T}}\def\mV{{\mathcal V}}
\def\mK{\mathcal K}\def\mA{{\mathcal A}}\def\mU{{\mathcal U}}
% Symmetric bilinear and trilinear forms of the Taylor expansion, used in the
% first Lyapunov coefficient (AAHOSN/New.tex)
\providecommand{\Bform}{\mathcal{B}}
\providecommand{\Cform}{\mathcal{C}}

% Blackboard bold.  Kept here, not in the drivers: a \providecommand in a
% driver preamble runs before this file and silently wins, so the same name
% ends up meaning different things in different builds.
\def\IC{\mathbb{C}}\def\IN{\mathbb{N}}\def\IQ{\mathbb{Q}}
\def\IR{\mathbb{R}}\def\IZ{\mathbb{Z}}\def\IG{\mathbb{G}}
\def\II{\mathbb{I}}

% Mathscr
\newcommand{\Rs}{{\mathscr R}}
\newcommand{\Ss}{{\mathscr S}}
\newcommand{\Ns}{{\mathscr N}}
\newcommand{\Cs}{\mathscr C}


%% ========================================
%% OPERATORS & FUNCTIONS
%% ========================================
\DeclareMathOperator{\tr}{tr}
\DeclareMathOperator{\diag}{Diag}
\DeclareMathOperator{\Diag}{Diag}
\DeclareMathOperator{\intr}{\mathrm{Int}}
\DeclareMathOperator{\Inv}{Inv}
\DeclareMathOperator{\spec}{spec}
\DeclareMathOperator{\AUC}{AUC}               % area under the curve
\newcommand{\ii}{\mathrm{i}}                  % imaginary unit
\newcommand\op{\operatorname}
\newcommand{\supp}{\op{supp}}
\newcommand{\diff}{\mathop{}\mathopen{}\mathrm{d}}
\newcommand{\pdd}[1]{\frac{\partial}{\partial{#1}}}
\DeclareMathOperator{\cone}{cone}
\DeclareMathOperator{\proj}{proj}

%% ========================================
%% BRACKETS & NORMS
%% ========================================
\providecommand{\pr}[1]{\left(#1\right)}
\providecommand{\pp}[1]{\left[#1\right]}
\newcommand\croc[1]{\left\langle #1\right\rangle}
\newcommand{\norm}[1]{\|#1\|}
\newcommand{\inner}[2]{\langle #1, #2 \rangle}
% Short synonym of \inner, used by AAHOSN/New.tex as $\ip{u}{v}=\bar u^\top v$
\providecommand{\ip}[2]{\langle #1 , #2 \rangle}


%% ========================================
%% FORMATTING & UTILITIES
%% ========================================
\newcommand*{\Scale}[2][4]{\scalebox{#1}{$#2$}}
\newcommand{\te}[1]{\text{#1}}
\newcommand\steq[1]{\stackrel{\text{\rm #1.}}{=}}
\newcommand\blue[1]{\textcolor{blue}{#1}}
\newcommand{\red}{\textcolor[rgb]{1.00,0.00,0.00}}
\newcommand{\green}{\textcolor[rgb]{0.00,1.00,0.00}}

% Figure command
\newcommand{\figu}[3]{%
\begin{figure}[H]
\centering
\includegraphics[scale=#3]{#1}
\caption{#2\label{f:#1}}
\end{figure}
}

% Footnote
\long\def\symbolfootnote[#1]#2{%
\begingroup
\def\thefootnote{\fnsymbol{footnote}}\footnote[#1]{#2}%
\endgroup}
\def\fn{\symbolfootnote}

\newcommand{\dsum}{\displaystyle\sum}

%% ========================================
%% CRN & EPIDEMIOLOGY SPECIFIC
%% ========================================
% CRN notation
\def\SS{\mathcal S}
\def\CC{\mathcal C}
\def\RR{\mathcal R}

% Reproduction numbers
\def\Rto{R_{1}^{E_2}}
\def\Rot{R_{2}^{E_1}}
\def\Roc{R_{1,c}}
\def\Rtc{R_{2,c}}
\def\rd{r_{0}}
\def\sd{s_{0}}
\def\dd{D_{0}}
\def\vd{v_{0}}
\def\zd{z_{0}}
\def\se{\sE}
\def\xd{x_{0}}
%% ========================================
%% SPECIAL VALUES & EQUILIBRIA
%% ========================================
\def\eE{\bar{e}}
\def\wE{\bar{w}}
\def\sE{s_{*}}
\def\Me{\T M_{ee}}
\def\Rnne{\mathbb{R}_{\ge 0}}

\newcommand{\ones}{\mathbf 1}


%% ========================================
%% ABBREVIATIONS (alphabetical by category)
%% ========================================

% -- General mathematical terms --
\def\adm{admissible}\def\admy{admissibility}
%\def\and{antisymmetric}% disabled 2026-09-30: it overrode LaTeX's author separator \and (printed "antisymmetric" between authors)
\def\ass{assumption}
\def\cadlag{c\`adl\`ag}
\def\ch{characteristic polynomial}
\def\chp{characteristic polynomial}
\def\coe{coefficient}
\def\comp{computation}
\def\con{condition}
\def\coo{coordinate}
\def\corr{corresponding}
\def\dec{decomposition}
\def\den{denominator}
\def\dete{deterministic}
\def\eig{eigenvalue}\def\eval{eigenvalue}
\def\eigv{eigenvector}\def\eiv{eigenvector}\def\evec{eigenvector}
\def\elt{element}
\def\enn{essentially nonnegative/positivity preserving}
\def\equ{equation}
\def\etal{et al.}
\def\fe{for example}
\def\fno{from now on}
\def\for{\forall}
\def\form{formula}
\def\Fr{Furthermore, }
\def\frt{furthermore}
\def\fun{function}
\def\gene{generalization}
\def\heu{heuristic}
\def\how{however}\def\How{However, }
\def\ie{i.e. }
\def\im{\item}
\def\isd{is defined by}
\def\itc{it is easy to check that}
\def\itf{it follows that}\def\Itf{It follows that}
\def\ith{it holds that}
\def\itm{it may be checked that}\def\Itm{It may be checked that}
\def\ito{it turns out that}\def\Ito{It turns out that}\def\Itou{It turns out useful to}
\def\ma{mass action}\def\MAR{mass action representation}
\def\mat{matrix}
\def\mbc{may be checked to be}
\def\mbw{may be written as}\def\mubw{must be written as}
\def\mpr{more precisely, }\def\Mp{More precisely, }
\def\nc{non-catalytic}
\def\nne{non-negative}\def\nny{nonnegativity}
\def\no{\nonumber}
\def\noni{non-infectious}
\def\nota{notation}
\def\np{\newpage}
\def\Nt{Note that}
\def\nt{non trivial}\def\nz{non zero}
\def\oth{otherwise}
\def\Oth{On the other hand}
\def\pa{parameter}\def\para{parameter}
\def\pol{polynomial}
\def\prob{probability}
\def\Prf{{\bf  Proof:}}\def\prf{{\bf  Proof:}}
\def\prop{proportional to}
\def\qu{\quad}
\def\QED{\hfill {$\square$}\goodbreak \medskip}
\def\rep{representation}
\def\resp{respectively}
\def\satd{satisfied}\def\satg{satisfying}\def\sats{satisfies}\def\saty{satisfy}
\def\ssec{\subsection}\def\sssec{\subsubsection}
\def\sef{see for example}
\def\sksy{skew-symmetric}
\def\sub{substrate}
\def\sys{system}
\def\var{variable}
\def\wk{well-known}
\def\wl{without loss of generality}\def\wlo{w.l.o.g.}
\def\wmc{We may check that}\def\Wmc{We may check that}
\def\wmr{we may rewrite}
\def\wms{we must show that}
\def\wrt{with respect to}

% -- CRN & Chemistry --
\newcommand\CRN{chemical reaction network}
\newcommand\CRNs{chemical reaction networks}
\newcommand\RN{reaction network}
\def\ACR{absolute concentration robustness}
\def\AMSD{acyclic siphon decomposition}
\def\DAE{differential algebraic equations}
\def\FHJ{Feinberg-Horn-Jackson}
\def\GMAK{generalized mass action kinetics}
\def\MM{Michaelis-Menten}
\def\sm{stoichiometric matrix}\def\Sm{stoichiometric matrix}
\def\WRZD{weakly reversible, zero deficiency}\def\aut{autocatalytic}

% -- Epidemiology --
\def\brn{basic reproduction number}\def\bzn{basic replacement number}
\def\com{compartment}
\def\dfe{disease free equilibrium}\def\DFE{disease free equilibrium}
\def\dwt{dwell times}\def\edt{expected dwell times}
\def\EM{epidemic model}
\def\endp{endemic fixed point}\def\enp{endemic point}\def\EE{endemic point}
\def\ME{mathematical epidemiology}
\def\MEE{eco-epidemiology}\def\EEP{eco-epidemiological}
\def\ece{eco-epidemiological model}
\def\mgm{multi group model}
\def\mtS{multi-strain model}\def\mSt{multi-strain model}
\def\mSi{minimal siphon}
\def\NGM{next generation matrix}
\def\nim{new infections matrix}
\def\qem{quadratic epidemiological model}
\def\repF{invasion functions}
\def\SPF{SIR-PH-FA model}
\def\sus{susceptible}

% -- Graph theory & Linear algebra --
\def\jac{Jacobi matrix}
\def\jin{Jacobian matrix of the invasion vector field with respect to the invasion variables}
\def\jrn{Jacobian matrix of the resident vector field with respect to the resident variables}
\def\jxn{Jacobian matrix of the $\x$ vector field with respect to the $\x$ variables}
\def\Lm{Laplacian matrix}
\def\MTT{matrix tree theorem}
\def\SGT{spectral graph theory}
\def\SMW{Sherman–Morrison–Woodbury}

% -- Stability & Dynamics --
\def\absC{absorbtion CTMC}
\def\BTb{Bogdanov-Takens bifurcation}
\def\CEP{competitive exclusion partition}
\def\CTMC{continuous time Markov chain}
\def\ds{differential system}
\def\fp{fixed point}\def\bfp{boundary fixed point}\def\fbp{fixed boundary point}
\def\fbl{fixed boundary line}\def\fif{forward invariant face}
\def\gas{global asymptotic stability}\def\GAS{globally asymptotically stable}
\def\GASS{\text{GASS}}
\def\Hb{Hopf bifurcation}
\def\Hur{Hurwitz}\def\MH{Metzler-Hurwitz}\def\RH{Routh-Hurwitz}
\def\IS{\text{IS}}
\def\LAS{locally asymptotically stable}
\def\LdS{Lyapunov diagonally stable}
\def\Lf{Lyapunov function}\def\LyV{Lyapunov Volterra type function}\def\RLF{robust Lyapunov function}
\def\LyM{Lyapunov Malkin stable}
\def\PF{Perron-Frobenius}
\def\spfp{strictly positive fixed point}
\def\sta{stability}
\def\STP{sharp threshold property}\def\STT{sharp threshold theorem}

% -- Models (with \varphi explicitly instead of \f) --
\def\AB{$(A,B)$ Arino-Brauer epidemic models}
\def\ABa{$(A,\bbe,\ba,\varphi)$ balanced bilinear model}
\def\ABP{%$(A,\bbe,P,\varphi)$
\\ balanced bilinear model}
\def\AMG{Anderson May Gupta}
\newcommand\BIN{biological interaction network}
\def\Dt{Descartes type}
\def\gLV{generalized Lotka-Volterra model}
\def\IG{\text{invasion graph}}\def\IGMS{interaction graph on minimal siphons}
\def\LCP{\text{linear complementarity problem}}
\def\LV{Lotka-Volterra}\def\LVK{Lotka-Volterra-Kolmogorov}

% -- Splitting & Structure --
\def\reg{regular splitting}\def\regS{{\bf regular splitting}}
\def\RUR{rational univariate representation}
\def\bds{block-diagonal structure}\def\NGMs{next generation matrices}

% -- Miscellaneous --
\def\Ma{{Mathematica}}
\def\Mm{Max–Min}
\def\num{numerator}
\def\rap{rational approximation} \def\DSR{directed species-reaction bipartite graph}\def\PCP{polynomial complementarity problem}\def\rat{rational substitution}
\def\QSS{quasi steady state} \def\RHS{right hand side} \def\LHS{left hand side}
\def\OSN{online social network}\def\VdD{Van den Driessche}
\def\wli{$\omega$-limit} \def\rk{rich kinetics}\def\sJ{symbolic Jacobian}
\def\cV{\cite{Vas,VasSta,BSV,GSV}} \def\spf{semi-parametric semi-functional}
\def\CRW{Capasso-Ruan-Wang} \def\ABP{$(A,\bbe,P,\f)$ balanced bilinear model}
\def\ABa{%$(A,\bbe,\ba,\f)$ \\
rank-one balanced bilinear model}
\def\bbe{B%{\bff \beta}
} \def\vbe{\vec {\bff \beta}
}
\def\vb{\vec {\bff b}}
\newcommand{\Kh}{\hat{K}}          % K with hat, first mention only
\newcommand{\wS}{w^{\sigma}}       % right Perron J^\perp, full block
\newcommand{\piS}{\vec \pi^{\sigma}}    % left Perron J^\perp, full block
\newcommand{\ws}{w^{\sigma}}       % right Perron J^\perp, sub-block
\newcommand{\pis}{\pi^{\sigma}}    % left Perron J^\perp, sub-block
\newcommand{\hwS}{\hat{w}^{\sigma}}   % right Perron K, full block
\newcommand{\hpiS}{\hat{\pi}^{\sigma}} % left Perron K, full block
\newcommand{\hws}{\hat{w}^{\sigma}}   % right Perron K, sub-block
\newcommand{\hpis}{\hat{\pi}^{\sigma}} % left Perron K, sub-block
\newcommand{\twS}{\tilde{w}^{\sigma}}  % right Perron \tilde K, full block
\newcommand{\tpiS}{\tilde{\pi}^{\sigma}} % left Perron \tilde K
\newcommand{\tws}{\tilde{w}^{\sigma}}
\newcommand{\tpis}{\tilde{\pi}^{\sigma}}
\newcommand{\Mat}[2]{\mathbb{R}_{\ge 0}^{\,#1\times #2}}
\def\tJ{J_\sigma^\perp} \def\lPe{left Perron eigenvector}
% \G is \Gamma (defined at :795).  The Volterra function keeps its own name,
% so that \G no longer has two conflicting meanings in this file.
\def\GV{G}  % Volterra G(u)=u-1-\ln u
\def\PVL{Perron-Volterra Lyapunov minimal siphon ansatz}
\newcommand{\istar}[1]{i_{#1}^{*}}
%\newcommand{\sE}{s_E}
\newcommand{\vE}{v_E}
\def\PV{Perron-Volterra candidate Lyapunov function}
\def\GLV{generalized Lotka-Volterra}
\def\MERN{mathematical epidemiology reaction networks}
\def\CMM{contraction metric method}
\newcommand{\sign}{\operatorname{sign}}
\newcommand{\Var}{\operatorname{Var}} 

% ======================================================================
% IDEAS SHARED BETWEEN THE ABSTRACT AND THE MAIN RESULTS SECTION
%
% Write each idea ONCE, in the abstract; the Main Results section recalls it,
% so the two can never drift apart.
%
%   in the abstract        \ideaitem{1}{The first-order reaction terms ...}
%   in Main Results        \ideahead{1}
%
% \ideaitem stores the text globally and typesets
%   \item \hyperlink{idea:1}{<text>}
% \ideahead typesets
%   \item\hypertarget{idea:1}{}\textbf{<the same text>}
%
% The abstract must come before the Main Results input, which is the normal
% order. If a heading has to read differently from the abstract sentence (for
% instance a long abstract item summarised in one line), override it with the
% optional argument: \ideahead[short heading]{6}.
% ======================================================================
\makeatletter
\newcommand{\ideaitem}[2]{%
  \expandafter\gdef\csname @idea@#1\endcsname{#2}%
  \item\hyperlink{idea:#1}{#2}}
\newcommand{\ideahead}[2][]{%
  \item\hypertarget{idea:#2}{}%
  \def\@ideatmp{#1}%
  \ifx\@ideatmp\@empty
    \@ifundefined{@idea@#2}{%
      \PackageWarning{preamble}{Idea #2 is used in Main Results but was never
        defined by \string\ideaitem\space in the abstract}%
      \textbf{[idea #2 undefined]}%
    }{\textbf{\begingroup\def\ideaterm##1##2{##2}\def\idearest##1{}\@nameuse{@idea@#2}\endgroup}}%
  \else
    \textbf{#1}%
  \fi}
% \ideaterm{<hypertarget>}{<words>}: inside an abstract idea the words become an
% invisible link (the document sets \hypersetup{hidelinks}) jumping to the place
% in Main Results that defines them; inside the Main Results heading itself,
% where the target sits a few lines below, \ideahead strips the link and prints
% the words plainly. Mark the destination with \hypertarget{<name>}{} on the line
% before the display, which adds no equation number and shows nothing.
\newcommand{\ideaterm}[2]{\hyperlink{#1}{#2}}
% \idearest{<text>}: the part of an abstract idea after its first sentence; printed in
% the abstract, omitted from the Main Results heading.
\newcommand{\idearest}[1]{#1}
\makeatother

% ======================================================================
% SHORT-PAPER SWITCH
% \shorttrue in a master that assembles only part of the companion paper
% (Barrat/SipArxiv.tex). \nsh{...} drops text that cites results living in the
% omitted sections, so the short build has no dangling references while the
% companion paper keeps the text verbatim.
% ======================================================================
\newif\ifshort \shortfalse
\newcommand{\nsh}[1]{\ifshort\else#1\fi}
% \shortalt{<short build>}{<companion paper>}: use when a sentence needs
% different wording in the two builds, not merely a deletion.
\newcommand{\shortalt}[2]{\ifshort#1\else#2\fi}
% \Jtr, \JtrS{\Sigma}: transverse Jacobian. J^\perp here; SipArxiv.tex redefines them as J_{\mathrm{tra}}, J_{\mathrm{tra},\Sigma}
\newcommand{\Jtr}{J^\perp}
\newcommand{\JtrS}[1]{J_{#1}^\perp}
.
\makeatletter
\AtBeginDocument{\@ifundefined{hypersetup}{}{\hypersetup{linktoc=all}}}
\makeatother
\usepackage{graphicx,float,fancybox}
\usepackage{epsfig,epstopdf,tikz}
\usepackage[sans]{dsfont}
\usetikzlibrary{positioning,arrows,arrows.meta,calc}
% Configure listings for Mathematica code
\lstdefinestyle{mathematica}{
 language=Mathematica,
 basicstyle=\small\ttfamily,
 backgroundcolor=\color{gray!10},
 frame=single,
 breaklines=true,
 commentstyle=\color{green!60!black},
 keywordstyle=\color{blue},
 stringstyle=\color{red},
 showstringspaces=false,
 tabsize=2
}
\usepackage{tcolorbox}
% Colored boxes for key results
\newtcolorbox{keyresult}[1][]{
 colback=blue!5!white,
 colframe=blue!75!black,
 title=#1,
 fonttitle=\bfseries
}
\newtcolorbox{examplebox}[1][]{
 colback=green!5!white,
 colframe=green!75!black,
 title=#1,
 fonttitle=\bfseries
}


\makeatletter
\@ifundefined{theorem}   {\newtheorem{theorem}   {Theorem}}   {}
\@ifundefined{definition}{\newtheorem{definition}{Definition}} {}
\@ifundefined{open}      {\newtheorem{open}      {Problem}}   {}
\@ifundefined{example}   {\newtheorem{example}   {Example}}   {}
\makeatother

% COUNTER DEFINITIONS
\newcounter{tabcounter}
\newcounter{defcounter}
\newcounter{opcounter}
\newcounter{thmcounter}
\newcounter{excounter}
\newcounter{ascounter}
\newcounter{cjcounter}
\newcounter{lemcounter}
\newcounter{remcounter}
\newcounter{corcounter}
\newcounter{propcounter}
\newcounter{figcounter}

% UNIFIED GLOSSARY STYLE
\newglossarystyle{unifiedstyle}{%
  \setglossarystyle{list}%
  \renewcommand*{\glossarypreamble}{\gdef\glslastsec{}}%
  \renewcommand{\glossentry}[2]{%
    \edef\glscursec{\glsentryuseriii{##1}}%
    \ifx\glscursec\glslastsec\else
      \item[]\textbf{\S\,\glscursec\enspace\glsentryuseriv{##1}}%
      \global\let\glslastsec\glscursec
    \fi
    \item[\glsentryuserii{##1} \glsentryuseri{##1}:] ``\glossentryname{##1}'' \glsentrydesc{##1}%
  }%
}

\makeatletter
% First-word extractor (used by \beT, \beD and \beXa for the default key)
% \@firstwordof hello world \relax  ->  hello
\def\@firstwordof#1 #2\relax{#1}
% First three letters, used for the heading mnemonic of a theorem given no
% explicit key.  Padding with \@empty keeps it safe for first words shorter
% than three letters: \@empty vanishes under \edef, whereas \relax would
% survive into the label and is illegal inside \csname (it broke every
% two-letter title, e.g. "ME ...", with "Missing \endcsname inserted").
\def\@takethree#1#2#3#4\@nil{#1#2#3}

% EXAMPLE COMMAND -- same signature and behaviour as \beT and \beD.
% Usage: \beXa*[long title][short title][key]
%   star  = suppresses the printed "[key]" tag after the title (as for \beT*
%   and \beD*); the label and glossary entry are unaffected.
%   key = optional label key override; the label becomes e:KEY, and without a
%   key it is e: plus the first three letters of the title's first word.  A
%   mnemonic label lets an example be cited as "Example~\ref{e:KEY}" instead
%   of by section number, which is what the citing prose needs.
\NewDocumentCommand{\beXa}{s O{} O{} o}{%
  \def\@temparg{#2}%
  \ifx\@temparg\@empty
    \begin{example}%
  \else
    \let\@savedSecTitle\@currentsectionname%
    \refstepcounter{excounter}%
    \IfValueTF{#4}{%
      \edef\@exlblkey{#4}%
    }{%
      \edef\@exfirstword{\@firstwordof#2 \relax}%
      \edef\@exlblkey{%
        \expandafter\@takethree\@exfirstword\@empty\@empty\@empty\@nil}%
    }%
    \def\@tempshort{#3}%
    \ifx\@tempshort\@empty
      \expandafter\xdef\csname @xshort@\@exlblkey\endcsname{\@exlblkey}%
    \else
      \expandafter\def\csname @xshort@\@exlblkey\endcsname{#3}%
    \fi
    \IfBooleanTF{#1}{\let\@extag\@empty}{%
      \edef\@extag{\noexpand\ \noexpand\textup{[\noexpand\texttt{e:\@exlblkey}]}}}%
    \begin{example}[#2\ifshowkeys\@extag\fi]%
    \label{e:\@exlblkey}%
    % record the long title so \nref{e:KEY} prints it, as for theorems
    \protected@write\@auxout{}{%
      \string\@thmtitle{e:\@exlblkey}{\unexpanded{#2}}}%
    \expanded{%
      \noexpand\newglossaryentry{e\theexcounter}{%
        name={\unexpanded{#2}%
          \IfBooleanF{#1}{ [\noexpand\texttt{e:\@exlblkey}]}},%
        sort={\theexample},%
        description={p.~\noexpand\pageref{e:\@exlblkey}},%
        user1={\theexample},%
        user2={Example},%
        user3={\thesection},%
        user4={\expandafter\unexpanded\expandafter{\@savedSecTitle}}%
      }%
    }%
    \protected@write\@auxout{}{\string\glsliveentry{e\theexcounter}}%
  \fi
}
\newcommand{\eeXa}{\end{example}}

% DEFINITION COMMAND -- same signature and behaviour as \beT below.
% Usage: \beD[long title][short title][key]
%   long title  = definition heading (plain text only)
%   short title = shown in \dref{word}; defaults to the key
%   key         = optional label key override; the label becomes d:KEY
% Without a key the label is d: plus the first three letters of the title's
% first word, so it is a mnemonic rather than the definition's number: a
% numeric label would silently change whenever a definition is inserted.
\NewDocumentCommand{\beD}{s O{} O{} o}{%
  \def\@temparg{#2}%
  \ifx\@temparg\@empty
    \begin{definition}%
  \else
    \let\@savedSecTitle\@currentsectionname%
    \refstepcounter{defcounter}%
    \IfValueTF{#4}{%
      \edef\@deflblkey{#4}%
    }{%
      \edef\@deffirstword{\@firstwordof#2 \relax}%
      \edef\@deflblkey{%
        \expandafter\@takethree\@deffirstword\@empty\@empty\@empty\@nil}%
    }%
    \def\@tempshort{#3}%
    \ifx\@tempshort\@empty
      \expandafter\xdef\csname @dshort@\@deflblkey\endcsname{\@deflblkey}%
    \else
      \expandafter\def\csname @dshort@\@deflblkey\endcsname{#3}%
    \fi
    \IfBooleanTF{#1}{\let\@deftag\@empty}{%
      \edef\@deftag{\noexpand\ \noexpand\textup{[\noexpand\texttt{d:\@deflblkey}]}}}%
    \begin{definition}[#2\ifshowkeys\@deftag\fi]%
    \label{d:\@deflblkey}%
    % record the long title so \nref{d:KEY} prints it, as for theorems
    \protected@write\@auxout{}{%
      \string\@thmtitle{d:\@deflblkey}{\unexpanded{#2}}}%
    \expanded{%
      \noexpand\newglossaryentry{d\thedefcounter}{%
        name={\unexpanded{#2}%
          \IfBooleanF{#1}{ [\noexpand\texttt{d:\@deflblkey}]}},%
        sort={\thedefinition},%
        description={p.~\noexpand\pageref{d:\@deflblkey}},%
        user1={\thedefinition},%
        user2={Definition},%
        user3={\thesection},%
        user4={\expandafter\unexpanded\expandafter{\@savedSecTitle}}%
      }%
    }%
    \protected@write\@auxout{}{\string\glsliveentry{d\thedefcounter}}%
  \fi
}
\newcommand{\eeD}{\end{definition}}

% \showkeysfalse in a driver hides the [t:KEY] mnemonics in result headings
% (labels and glossary entries are unchanged); used for submission builds.
\newif\ifshowkeys \showkeystrue
% \commonredtrue in a driver that inputs basicFiles/commonRed: tool.tex then
% cites its reduction results instead of restating them.
\newif\ifcommonred \commonredfalse
% TheorEM COMMAND
% Usage: \beT[long title][short title][key]
%   long title  = theorem heading (plain text only)
%   short title = shown in \tref{word}; may contain math; defaults to first word of long title
%   key         = optional label key override; label becomes t:KEY (avoids first-word collisions)
% Label is t:WORD where WORD = first word of long title (or KEY if given); use \tref{WORD} to reference.
\NewDocumentCommand{\beT}{s O{} O{} o}{%
  \def\@temparg{#2}%
  \ifx\@temparg\@empty
    \begin{theorem}%
  \else
    \let\@savedSecTitle\@currentsectionname%
    \refstepcounter{thmcounter}%
    % \@thmtag is the mnemonic shown in the heading; every theorem gets one.
    % With an explicit key the key itself is shown.  Without one the key
    % defaults to the title's first word, and showing that whole word would
    % merely repeat the title ("Dissipativity [Dissipativity]"), so only its
    % first three letters are shown: "Dissipativity [Dis]".  The label is
    % unaffected -- it still uses the full \@thmlblkey.
    \IfValueTF{#4}{%
      \edef\@thmlblkey{#4}%
    }{%
      \edef\@thmfirstword{\@firstwordof#2 \relax}%
      \edef\@thmlblkey{%
        \expandafter\@takethree\@thmfirstword\@empty\@empty\@empty\@nil}%
    }%
    \IfBooleanTF{#1}{\let\@thmtag\@empty}{%
      \edef\@thmtag{\noexpand\ \noexpand\textup{[\noexpand\texttt{t:\@thmlblkey}]}}}%
    \def\@tempshort{#3}%
    \ifx\@tempshort\@empty
      \expandafter\xdef\csname @tshort@\@thmlblkey\endcsname{\@thmlblkey}%
    \else
      \expandafter\def\csname @tshort@\@thmlblkey\endcsname{#3}%
    \fi
    \begin{theorem}[#2\ifshowkeys\@thmtag\fi]%
    \label{t:\@thmlblkey}%
    % record the long title for \nref; via the .aux so forward refs resolve
    \protected@write\@auxout{}{%
      \string\@thmtitle{t:\@thmlblkey}{\unexpanded{#2}}}%
    \expanded{%
      \noexpand\newglossaryentry{t\thethmcounter}{%
        name={\unexpanded{#2}%
          \IfBooleanF{#1}{ [\noexpand\texttt{t:\@thmlblkey}]}},%
        % sort/user1 must be the number the document actually prints, i.e. the
        % theorem environment's counter, not thmcounter.  thmcounter advances
        % only on \beT with a title, so a bare \beT (or \beT with an
        % empty title) makes the two drift and the glossary cites a wrong
        % number.  thmcounter is still fine as the entry key: it is unique.
        sort={\thetheorem},%
        description={p.~\noexpand\pageref{t:\@thmlblkey}},%
        user1={\thetheorem},%
        user2={Theorem},%
        user3={\thesection},%
        user4={\expandafter\unexpanded\expandafter{\@savedSecTitle}}%
      }%
    }%
    \protected@write\@auxout{}{\string\glsliveentry{t\thethmcounter}}%
  \fi
}
\newcommand{\eeT}{\end{theorem}}

% CROSS-PART RELATION PARAGRAPH
% Usage: \xpart{text}  -- place immediately BEFORE \beT, never inside it
%   (anything inside \beT...\eeT would be swallowed into the glossary entry).
% The paragraph is typeset only when Part I is actually present, i.e. in the
% combined book build; in the standalone builds (GranadaDSABNS/A.tex,
% AAHRah-Submission/AAHRah.tex) the label part:glv is undefined and the whole
% paragraph disappears rather than leaving ?? cross-references. Same
% graceful-degradation principle as \GASref in GASt.tex. Resolves on the
% second LaTeX pass, since r@part:glv comes from the .aux file.
% (Already inside the \makeatletter region opened above, so no \makeatletter
% pair here -- adding one would close that region early.)
\newcommand{\xpart}[1]{\@ifundefined{r@part:glv}{}{\paragraph{Relation with Parts~I--II.}#1}}

% \bookref{label}{precise text}{fallback text}: prints the precise cross-reference
% when the label resolves (combined book build), and the fallback text otherwise
% (standalone paper builds, where book-only appendix labels are absent). Same
% graceful-degradation role as \GASref in GASt.tex, but usable in any file since it
% lives in this preamble (loaded by every driver) -- unlike \GASref, which is defined in a
% content file and is therefore undefined in standalone builds that omit that file.
\newcommand{\bookref}[3]{\@ifundefined{r@#1}{#3}{#2}}

% \PVpart and \RBpart: the two parts that a partial build may omit.
% They print "Part~N" in any build that contains the part, and a naming phrase
% in a build that does not, so an excerpt shows no ?? .  Each already supplies
% the word "Part", so write "\PVpart develops ...", not "Part~\PVpart".
\newcommand{\PVpart}{\bookref{part:aahrah}%
  {Part~\ref{part:aahrah}}{the Perron--Volterra part of the full book}}
\newcommand{\RBpart}{\bookref{part:recbrowne}%
  {Part~\ref{part:recbrowne}}{the immuno-virology part of the full book}}

% OPEN PROBLEM COMMANDS
% ASSUMPTION COMMANDS (glossaried)
\newcommand{\beAs}[1][]{%
  \def\@temparg{#1}%
  \ifx\@temparg\@empty
    \begin{assumption}%
  \else
    \let\@savedSecTitle\@currentsectionname%
    \refstepcounter{ascounter}%
    \begin{assumption}[#1]%
    \label{as:\theascounter}%
    \expanded{%
      \noexpand\newglossaryentry{as\theascounter}{%
        name={\unexpanded{#1}},%
        sort={\theascounter},%
        description={p.~\noexpand\pageref{as:\theascounter}},%
        user1={\theascounter},%
        user2={Assumption},%
        user3={\thesection},%
        user4={\expandafter\unexpanded\expandafter{\@savedSecTitle}}%
      }%
    }%
    \protected@write\@auxout{}{\string\glsliveentry{as\theascounter}}%
  \fi
}
\newcommand{\eeAs}{\end{assumption}}
% CONJECTURE COMMANDS (glossaried)
\newcommand{\beCj}[1][]{%
  \def\@temparg{#1}%
  \ifx\@temparg\@empty
    \begin{conjecture}%
  \else
    \let\@savedSecTitle\@currentsectionname%
    \refstepcounter{cjcounter}%
    \begin{conjecture}[#1]%
    \label{cj:\thecjcounter}%
    \expanded{%
      \noexpand\newglossaryentry{cj\thecjcounter}{%
        name={\unexpanded{#1}},%
        sort={\thecjcounter},%
        description={p.~\noexpand\pageref{cj:\thecjcounter}},%
        user1={\thecjcounter},%
        user2={Conjecture},%
        user3={\thesection},%
        user4={\expandafter\unexpanded\expandafter{\@savedSecTitle}}%
      }%
    }%
    \protected@write\@auxout{}{\string\glsliveentry{cj\thecjcounter}}%
  \fi
}
\newcommand{\eeCj}{\end{conjecture}}
% LEMMA COMMANDS
% ---------------------------------------------------------------------
% UNIFORM RESULT WRAPPERS -- \beL, \beC, \beP, \beR share \beT's signature
%
%   \beX[long title][short title][key]  ...  \eeX
%
%   long title  = heading (math allowed)
%   short title = shown in the cross-reference; defaults to the key
%   key         = label key.  With a key the label is X:KEY and the heading
%                 carries the mnemonic [key], exactly as for \beT.  Without
%                 one the label falls back to the running counter, which is
%                 what every pre-existing call site in the book relies on.
%
% A bare \beX -- arguments commented out -- opens the plain environment and
% creates NO label and NO glossary entry.  That is the supported way to keep
% an individual result out of the catalogue.
% ---------------------------------------------------------------------
% #1 wrapper name  #2 environment  #3 label prefix  #4 counter  #5 type name
\newcommand{\@mkresult}[5]{%
  % A leading * suppresses the [key] mnemonic in both the heading and the
  % glossary name, while still making the label X:KEY and the entry:
  %     \beX*[long][short][key]
  \expandafter\NewDocumentCommand\csname be#1\endcsname{s O{} O{} o}{%
    \def\@temparg{##2}%
    \ifx\@temparg\@empty
      \begin{#2}%
    \else
      \let\@savedSecTitle\@currentsectionname%
      \refstepcounter{#4}%
      % With a key: mnemonic label X:KEY and the [key] tag in the heading.
      % Without one: fall back to the running counter, exactly as these
      % wrappers behaved before, so every legacy one-argument call site in
      % the book keeps its label.  The counter is also csname-safe, which a
      % title's first word is not -- titles beginning with math would
      % otherwise build \csname \Omega...\endcsname and fail.
      \IfValueTF{##4}{%
        \edef\@relbl{##4}%
        \IfBooleanTF{##1}{\let\@retag\@empty}{%
          \edef\@retag{\noexpand\ \noexpand\textup{[\noexpand\texttt{#3:##4}]}}}%
      }{%
        \edef\@relbl{\arabic{#4}}%
        \let\@retag\@empty
      }%
      \def\@tempshort{##3}%
      \ifx\@tempshort\@empty
        \expandafter\xdef\csname @sshort@#3\@relbl\endcsname{\@relbl}%
      \else
        \expandafter\def\csname @sshort@#3\@relbl\endcsname{##3}%
      \fi
      \begin{#2}[##2\ifshowkeys\@retag\fi]%
      \edef\@tmplbl{#3:\@relbl}\expandafter\label\expandafter{\@tmplbl}%
      \expanded{%
        \noexpand\newglossaryentry{#3\csname the#4\endcsname}{%
          name={\unexpanded{##2}%
            \IfBooleanF{##1}{\IfValueT{##4}{ [\noexpand\texttt{#3:##4}]}}},%
          sort={\csname the#2\endcsname},%
          description={p.~\noexpand\pageref{#3:\@relbl}},%
          user1={\csname the#2\endcsname},%
          user2={#5},%
          user3={\thesection},%
          user4={\expandafter\unexpanded\expandafter{\@savedSecTitle}}%
        }%
      }%
      \protected@write\@auxout{}{\string\glsliveentry{#3\csname the#4\endcsname}}%
      \protected@write\@auxout{}{\string\@thmtitle{\@tmplbl}{\unexpanded{##2}}}%
    \fi
  }%
  \expandafter\newcommand\csname ee#1\endcsname{\end{#2}}%
}
\@mkresult{L}{lemma}{l}{lemcounter}{Lemma}
% REMARK / COROLLARY / PROPOSITION -- same signature as \beL above
\@mkresult{R}{remark}{r}{remcounter}{Remark}
\@mkresult{C}{corollary}{c}{corcounter}{Corollary}
\@mkresult{P}{proposition}{p}{propcounter}{Proposition}
% OPEN-PROBLEM COMMAND -- generated with the same signature as \beT:
%   \beO[long title][short title][key]   ...   \eeO
% label o:KEY, or o:<counter> when no key is given; a bare \beO makes
% neither a label nor a glossary entry.
\@mkresult{O}{open}{o}{opcounter}{Problem}


% Section-level title tracker: updated only by \section, not \subsection.
% \@currentlabelname is overwritten by every \refstepcounter (including
% \subsection), so we save it separately here.
\newcommand{\@currentsectionname}{}
\AtBeginDocument{%
  \let\@epidcrn@orig@refstepcounter\refstepcounter
  \renewcommand{\refstepcounter}[1]{%
    \@epidcrn@orig@refstepcounter{#1}%
    \ifstrequal{#1}{section}{\global\let\@currentsectionname\@currentlabelname}{}%
  }%
}

% TABLE COMMANDS (glossaried) -- consistent naming with \beT, \beD, etc.
\newcommand{\beTa}[1][]{%
  \def\@temparg{#1}%
  \ifx\@temparg\@empty
    \begin{table}[htbp]%
  \else
    \let\@savedSecTitle\@currentsectionname%
    \refstepcounter{tabcounter}%
    \begin{table}[htbp]%
    \label{A:\thetabcounter}%
    \expanded{%
      \noexpand\newglossaryentry{A\thetabcounter}{%
        name={\unexpanded{#1}},%
        sort={\thetabcounter},%
        description={p.~\noexpand\pageref{A:\thetabcounter}},%
        user1={\thetabcounter},%
        user2={Table},%
        user3={\thesection},%
        user4={\expandafter\unexpanded\expandafter{\@savedSecTitle}}%
      }%
    }%
    \protected@write\@auxout{}{\string\glsliveentry{A\thetabcounter}}%
  \fi
}
\newcommand{\eeTa}{\end{table}}

% FIGURE COMMANDS (glossaried) -- consistent with \beT, \beR, \beTa, etc.
% Usage:
%   \beF[Glossary title]...\caption{Visible caption}\eeF  -> glossaried figure
%   \beF...\caption{Visible caption}\eeF                  -> plain [H] figure, no glossary
\newcommand{\beF}[1][]{%
  \def\@temparg{#1}%
  \ifx\@temparg\@empty
    \begin{figure}[H]%
  \else
    \begin{figure}[H]%
    \let\@savedSecTitle\@currentsectionname%
    \refstepcounter{figcounter}%
    \label{fig:\thefigcounter}%
    \expanded{%
      \noexpand\newglossaryentry{Fig\thefigcounter}{%
        name={\unexpanded{#1}},%
        sort={\thefigcounter},%
        description={p.~\noexpand\pageref{fig:\thefigcounter}},%
        user1={\thefigcounter},%
        user2={Figure},%
        user3={\thesection},%
        user4={\expandafter\unexpanded\expandafter{\@savedSecTitle}}%
      }%
    }%
    \protected@write\@auxout{}{\string\glsliveentry{Fig\thefigcounter}}%
  \fi
}
\newcommand{\eeF}{\end{figure}}
\makeatother
% LIVE-ENTRY REGISTER
% \glsaddall prints every entry the glossaries package currently knows, and
% that set includes entries restored from \jobname.glsdefs.  The package
% rewrites .glsdefs from all *currently defined* entries, so anything it
% restored is written straight back out: a deleted section kept reappearing in
% the catalogue for ever, and the .glsdefs had to be erased by hand.  The
% wrappers therefore record every entry they actually create in the .aux, and
% \printboth adds only those.  The .aux is rebuilt from scratch on each run, so
% this is self-healing: a removed result is gone after one further run.
\makeatletter
\newcommand{\@glslive}{}
\newcommand{\@glsaddlive}[1]{\ifglsentryexists{#1}{\glsadd{#1}}{}}
\newcommand{\glsliveentry}[1]{\g@addto@macro\@glslive{\@glsaddlive{#1}}}
\makeatother
% PRINT COMMAND
% No \clearpage here -- the catalogue continues on the current page.  This has
% been re-added by mistake several times; leave it out.
\makeatletter
\newcommand{\printboth}{%
  \@glslive   % add exactly the entries this document created last run
  \printglossary[style=unifiedstyle, title={0 \hspace{.22em} Definitions, Theorems, Examples and Open Problems}]%
}
\makeatother
% Named theorem reference: \tref{word} -> Thm~N(word-or-short)
% Shows stored short title if available, otherwise shows the key itself.
\newcommand{\tref}[1]{%
  Thm~\ref{t:#1}(\ifcsname @tshort@#1\endcsname\csname @tshort@#1\endcsname\else#1\fi)%
}

% Full-title theorem reference: \nref{t:KEY} prints the number and then the
% theorem's long title, so one writes "Theorem~\nref{t:blocks}" and gets
% "Theorem 3 ``Block factorization and siphons''".  Takes the complete label, like
% \ref, not the bare key like \tref.
%
% \beT writes each title to the .aux (see \@thmtitle below), so \nref resolves
% forward references from the second LaTeX pass, exactly as \ref does.  If the
% title is not yet known -- first pass, or a label that is not a \beT theorem
% -- it degrades to the number rather than producing ??.
% These two use @ in their names, and this point in the file is outside the
% \makeatletter region opened earlier (it closes at \makeatother above), so
% the region has to be reopened here.
\makeatletter
\newcommand{\@thmtitle}[2]{\global\@namedef{@tlong@#1}{#2}}
\newcommand{\nref}[1]{%
  \ref{#1}%
  \ifcsname @tlong@#1\endcsname
    \ ``\csname @tlong@#1\endcsname''%
  \fi
}
% \nm{t:KEY}: the stored long title of any \beT/\beL/\beP/\beR/\beC/\beD/\beE result,
% verbatim and without quotes, so item names follow the macro's first argument.
\newcommand{\nm}[1]{%
  \ifcsname @tlong@#1\endcsname
    \csname @tlong@#1\endcsname%
  \else
    \ref{#1}%
  \fi
}
\makeatother




%% ========================================
%% Other THEOREM ENVIRONMENTS
%% ========================================
\makeatletter
\@ifundefined{lemma}      {\newtheorem{lemma}      {Lemma}}      {}
\@ifundefined{proposition}{\newtheorem{proposition}{Proposition}} {}
\@ifundefined{corollary}  {\newtheorem{corollary}  {Corollary}}  {}
\@ifundefined{conjecture} {\newtheorem{conjecture} {Conjecture}}  {}
\@ifundefined{remark}     {\newtheorem{remark}     {Remark}}     {}
\@ifundefined{question}   {\newtheorem{question}   {Question}}   {}
\@ifundefined{assumption} {\newtheorem{assumption} {Assumption}} {}
\@ifundefined{hypothesis} {\newtheorem{hypothesis} {Hypothesis}} {}
\@ifundefined{fact}       {\newtheorem{fact}       {Fact}}       {}
\makeatother

% \citer: cite the relay-graph threshold theorem of AAHrel. If a document has
% \input the chapter/file containing t:relayRF (so the label resolves), point
% straight at it; otherwise fall back to a bare citation of AAHrel. Shared here
% (not per-paper) so it resolves consistently across every book chapter.
% \citeR: same idea, for the first theorem of AAHRah (t:Pstar).
\makeatletter
\providecommand{\citer}{%
  \@ifundefined{r@t:relayRF}%
    {\cite{AAHrel}}%
    {Theorem~\ref{t:relayRF} of \cite{AAHrel}}%
}
\providecommand{\citeR}{%
  \@ifundefined{r@t:Pstar}%
    {\cite{AAHRah}}%
    {Theorem~\ref{t:Pstar} of \cite{AAHRah}}%
}
\makeatother


% Abbreviations
\def\beQ{\begin{question}}\def\eeQ{\end{question}}
\def\beA{\begin{assumption}}\def\eeA{\end{assumption}}


%% ========================================

%% ========================================
%% COLORS & CUSTOM COMMANDS
%% ========================================
\definecolor{funccolor}{RGB}{25,25,112}
\definecolor{desccolor}{RGB}{64,64,64}
\newcommand{\func}[1]{\texttt{\textcolor{funccolor}{#1}}}
\def\eq{=}
\def\mm{-}

%% ========================================
%% MATHCAL SHORTHANDS
%% ========================================
% Project-wide convention: \mX for a capital X is always \mathcal X.
% Defined for the whole alphabet so the convention never fails to expand.
\newcommand{\mA}{\mathcal A}\newcommand{\mB}{\mathcal B}
\newcommand{\mC}{\mathcal C}\newcommand{\mD}{\mathcal D}
\newcommand{\mE}{\mathcal E}\newcommand{\mF}{\mathcal F}
\newcommand{\mG}{\mathcal G}\newcommand{\mH}{\mathcal H}
\newcommand{\mI}{\mathcal I}\newcommand{\mJ}{\mathcal J}
\newcommand{\mK}{\mathcal K}\newcommand{\mL}{\mathcal L}
\newcommand{\mM}{\mathcal M}\newcommand{\mN}{\mathcal N}
\newcommand{\mO}{\mathcal O}\newcommand{\mP}{\mathcal P}
\newcommand{\mQ}{\mathcal Q}\newcommand{\mR}{\mathcal R}
\newcommand{\mS}{\mathcal S}\newcommand{\mT}{\mathcal T}
\newcommand{\mU}{\mathcal U}\newcommand{\mV}{\mathcal V}
\newcommand{\mW}{\mathcal W}\newcommand{\mX}{\mathcal X}
\newcommand{\mY}{\mathcal Y}\newcommand{\mZ}{\mathcal Z}

%% ========================================
%% BASIC MATH ENVIRONMENTS
%% ========================================
\def\bep{\begin{pmatrix}}\def\eep{\end{pmatrix}}
\def\bev{\begin{vmatrix}}\def\eev{\end{vmatrix}}
\def\bea{\begin{eqnarray*}}\def\eea{\end{eqnarray*}}
\def\bc{\begin{cases}}\def\ec{\end{cases}}
\def\BEN{\begin{enumerate}}\def\EEN{\end{enumerate}}
\def\BI{\begin{itemize}}\def\EI{\end{itemize}}

\newcommand{\be}[1]{\begin{equation}\label{#1}}
\newcommand{\ee}{\end{equation}}
\newcommand{\beq}{\begin{eqnarray}}
\def\eeq{\end{eqnarray}}
\def\eqr{\eqref}\def\fr{\frac}\def\lbl{\label}
\def\Lra{\Longrightarrow}\def\non{\nonumber}\def\cNGM{\cite{Diek,Van,Van08}}
\def\Eq{\Leftrightarrow}\def\b{\beta}\def\uph{\red{UP TO HERE}}
%% ========================================
%% PROBABILITY & STATISTICS
%% ========================================
\renewcommand{\P}{\ensuremath{\mathbb{P}}}  % Note: overrides paragraph symbol
\newcommand{\E}{\ensuremath{\mathbb{E}}}
\newcommand{\ind}[1]{\ensuremath{\mathbbm{1}_{\left\{#1\right\}}}}
\DeclareMathOperator{\Cov}{Cov}
\DeclareMathOperator{\TV}{TV}
\newcommand{\rank}{\operatorname{rank}}

%% ========================================
%% CRN OPERATORS (organizations/ERCs/buffering/siphons framework)
%% ========================================
\def\clos{\mathrm{clos}}
\def\Supply{\mathrm{Supply}}
\def\req{\mathrm{req}}
\def\React{\mathrm{React}}
\def\Prod{\mathrm{Prod}}
\def\Gen{\mathrm{Gen}}
\def\In{\mathrm{In}}
\def\Out{\mathrm{Out}}
\def\Act{\mathrm{Act}}
\def\Min{\mathrm{Min}}

%% ========================================
%% LINE BREAKING
%% ========================================
% Last-resort stretch, used only on paragraphs TeX cannot otherwise break:
% a long \texttt{} token, a comma-separated inline math list (commas are not
% breakpoints in math), or a run of \cite/\ref boxes.  Without it those set as
% overfull \hbox and print into the margin; with it they set slightly loose,
% which is the lesser evil.  Does not affect paragraphs that already break.
\emergencystretch=3em

%% ========================================
%% NUMBER SETS
%% ========================================
\newcommand{\R}{\mathbb{R}}  % Real numbers
\newcommand{\C}{\mathbb{C}}  % Complex numbers
\newcommand{\N}{\mathbb{N}}
\newcommand{\Z}{\mathbb{Z}}
% \providecommand, not \newcommand: some papers define these in their own
% preamble after % Shared preamble for the GLVbook papers.  This is the single canonical copy;
% a driver one level below GLVbook reaches it as
%
%     \input{../basicFiles/preamble}
%
% What stays in the driver, because it differs from paper to paper:
% \documentclass, font choice, \title, \author, \date, and any macro genuinely
% local to one paper.  Everything else belongs here.  The file is
% self-contained: it does not read any other file.

\usepackage[utf8]{inputenc}
\usepackage[novbox]{pdfsync}   % required: SyncTeX alone does not work here

% Page geometry, set once here for every paper that loads this file.  A driver
% that needs different margins says \geometry{...}; a second
% \usepackage[...]{geometry} raises an option clash and must not be used.
\makeatletter
\@ifundefined{AAdefNnogeometry}{\usepackage[a4paper,margin=2cm]{geometry}}{}
\makeatother

% mathematics
\usepackage{amsmath,amssymb,amsthm,mathtools,bm,mathrsfs,bbm}
\usepackage{centernot}

% floats, tables, lists
\usepackage{graphicx,float}
\makeatletter
\@ifpackageloaded{subfig}{}{\usepackage{subcaption}}
\makeatother
\usepackage{booktabs,longtable,threeparttable,array,multirow}
\usepackage{enumitem}

% colour, code, algorithms
\usepackage{xcolor,listings,tcolorbox}
\makeatletter
\@ifundefined{c@algorithm}{\usepackage{algorithm}}{}
\makeatother
\usepackage{algorithmicx,algpseudocode}

% diagrams
\usepackage{tikz,tikz-cd,pgfplots}
\usetikzlibrary{graphs,decorations.pathmorphing,decorations.markings,
arrows.meta,automata,positioning}
\tikzcdset{every label/.append style = {font = \small}}
\pgfplotsset{compat=1.18}

% chemistry
\usepackage[version=4]{mhchem}

\usepackage{microtype}
\usepackage{lineno}
\usepackage{nomencl}
\makenomenclature

% hyperref first, then bookmark, then glossaries: glossaries must follow
% hyperref or the entry links are built against the wrong anchors.
\makeatletter
\@ifundefined{hyperrefpackage}{%
  \@ifundefined{AAdefNnohyperref}{\usepackage{hyperref}}{}%
}{}
\@ifundefined{hyperrefpackage}{\usepackage{bookmark}}{}
\makeatother
\usepackage[toc,sort=use]{glossaries}
\makeglossaries

% Book/standalone switch.  Material that the book already covers in its
% prelude, but that a standalone paper still has to state, is wrapped as
%     \ifbookbuild\else ... \fi
% Drivers default to standalone; bookM.tex says \bookbuildtrue.
\newif\ifbookbuild
\bookbuildfalse

% Part switches for the gradostat paper.  Both default to true; a driver that
% wants the short submission says \madecfalse and/or \slowfastfalse, and every
% cross-reference into the dropped part disappears with it.
%   \ifmadec     Part III, the Lagasquie--Madec model and its break-even geometry
%   \ifslowfast  Part IV, everything carrying the small parameter \eps
\newif\ifmadec     \madectrue
\newif\ifslowfast  \slowfasttrue

% Part switch for the Barrat/simplicial-contagion project.  Defaults to
% true (the hypernetwork-synchrony part is included); the journal-paper
% driver says \hnetpartfalse to drop it and every cross-reference into it.
\newif\ifhnetpart  \hnetparttrue

%% ========================================
%% PACKAGES
%% ========================================
\RequirePackage{etoolbox}
\usepackage{makecell}
\usepackage[novbox]{pdfsync}% obsolete; SyncTeX is built into modern engines
\usepackage{amssymb, amsfonts,mathtools,bbm,eurosym,graphicx,latexsym, amsmath,amsthm,times,nccmath,url}
\usepackage{listings,verbatim,cases,tabularx}
% hyperref is deliberately NOT in the line above.  Journal classes often
% load it themselves, with a long option list, from a begindocument hook
% (mdpi.cls does: \DeclareHookRule{begindocument}{hyperref}{before}{mdpi}).
% Loading it here first, without options, makes that later load raise
% "Option clash for package hyperref".  So load it only if nobody else
% has, or will:
%   - \@ifpackageloaded  catches a class that loaded it already;
%   - \hyperrefpackage    is mdpi.cls's own deferred loader;
%   - \AAdefNnohyperref   lets any other driver opt out by defining it
%                         before \input{preamble}.
\makeatletter
\@ifpackageloaded{hyperref}{}{%
  \@ifundefined{hyperrefpackage}{%
    \@ifundefined{AAdefNnohyperref}{\usepackage{hyperref}}{}%
  }{}%
}
\makeatother
% linktoc=all makes a table-of-contents page number a link as well as the
% entry text; hyperref's default, linktoc=section, links only the text.
% Deferred to begin-document and guarded: by the block above hyperref may
% be loaded later (mdpi's own hook) or not at all (\AAdefNnohyperref), and
% \hypersetup must exist before it is called.  A driver that wants the
% default back can say \hypersetup{linktoc=section} after \input{preamble}.
\makeatletter
\AtBeginDocument{\@ifundefined{hypersetup}{}{\hypersetup{linktoc=all}}}
\makeatother
\usepackage{graphicx,float,fancybox}
\usepackage{epsfig,epstopdf,tikz}
\usepackage[sans]{dsfont}
\usetikzlibrary{positioning,arrows,arrows.meta,calc}
% Configure listings for Mathematica code
\lstdefinestyle{mathematica}{
 language=Mathematica,
 basicstyle=\small\ttfamily,
 backgroundcolor=\color{gray!10},
 frame=single,
 breaklines=true,
 commentstyle=\color{green!60!black},
 keywordstyle=\color{blue},
 stringstyle=\color{red},
 showstringspaces=false,
 tabsize=2
}
\usepackage{tcolorbox}
% Colored boxes for key results
\newtcolorbox{keyresult}[1][]{
 colback=blue!5!white,
 colframe=blue!75!black,
 title=#1,
 fonttitle=\bfseries
}
\newtcolorbox{examplebox}[1][]{
 colback=green!5!white,
 colframe=green!75!black,
 title=#1,
 fonttitle=\bfseries
}


\makeatletter
\@ifundefined{theorem}   {\newtheorem{theorem}   {Theorem}}   {}
\@ifundefined{definition}{\newtheorem{definition}{Definition}} {}
\@ifundefined{open}      {\newtheorem{open}      {Problem}}   {}
\@ifundefined{example}   {\newtheorem{example}   {Example}}   {}
\makeatother

% COUNTER DEFINITIONS
\newcounter{tabcounter}
\newcounter{defcounter}
\newcounter{opcounter}
\newcounter{thmcounter}
\newcounter{excounter}
\newcounter{ascounter}
\newcounter{cjcounter}
\newcounter{lemcounter}
\newcounter{remcounter}
\newcounter{corcounter}
\newcounter{propcounter}
\newcounter{figcounter}

% UNIFIED GLOSSARY STYLE
\newglossarystyle{unifiedstyle}{%
  \setglossarystyle{list}%
  \renewcommand*{\glossarypreamble}{\gdef\glslastsec{}}%
  \renewcommand{\glossentry}[2]{%
    \edef\glscursec{\glsentryuseriii{##1}}%
    \ifx\glscursec\glslastsec\else
      \item[]\textbf{\S\,\glscursec\enspace\glsentryuseriv{##1}}%
      \global\let\glslastsec\glscursec
    \fi
    \item[\glsentryuserii{##1} \glsentryuseri{##1}:] ``\glossentryname{##1}'' \glsentrydesc{##1}%
  }%
}

\makeatletter
% First-word extractor (used by \beT, \beD and \beXa for the default key)
% \@firstwordof hello world \relax  ->  hello
\def\@firstwordof#1 #2\relax{#1}
% First three letters, used for the heading mnemonic of a theorem given no
% explicit key.  Padding with \@empty keeps it safe for first words shorter
% than three letters: \@empty vanishes under \edef, whereas \relax would
% survive into the label and is illegal inside \csname (it broke every
% two-letter title, e.g. "ME ...", with "Missing \endcsname inserted").
\def\@takethree#1#2#3#4\@nil{#1#2#3}

% EXAMPLE COMMAND -- same signature and behaviour as \beT and \beD.
% Usage: \beXa*[long title][short title][key]
%   star  = suppresses the printed "[key]" tag after the title (as for \beT*
%   and \beD*); the label and glossary entry are unaffected.
%   key = optional label key override; the label becomes e:KEY, and without a
%   key it is e: plus the first three letters of the title's first word.  A
%   mnemonic label lets an example be cited as "Example~\ref{e:KEY}" instead
%   of by section number, which is what the citing prose needs.
\NewDocumentCommand{\beXa}{s O{} O{} o}{%
  \def\@temparg{#2}%
  \ifx\@temparg\@empty
    \begin{example}%
  \else
    \let\@savedSecTitle\@currentsectionname%
    \refstepcounter{excounter}%
    \IfValueTF{#4}{%
      \edef\@exlblkey{#4}%
    }{%
      \edef\@exfirstword{\@firstwordof#2 \relax}%
      \edef\@exlblkey{%
        \expandafter\@takethree\@exfirstword\@empty\@empty\@empty\@nil}%
    }%
    \def\@tempshort{#3}%
    \ifx\@tempshort\@empty
      \expandafter\xdef\csname @xshort@\@exlblkey\endcsname{\@exlblkey}%
    \else
      \expandafter\def\csname @xshort@\@exlblkey\endcsname{#3}%
    \fi
    \IfBooleanTF{#1}{\let\@extag\@empty}{%
      \edef\@extag{\noexpand\ \noexpand\textup{[\noexpand\texttt{e:\@exlblkey}]}}}%
    \begin{example}[#2\ifshowkeys\@extag\fi]%
    \label{e:\@exlblkey}%
    % record the long title so \nref{e:KEY} prints it, as for theorems
    \protected@write\@auxout{}{%
      \string\@thmtitle{e:\@exlblkey}{\unexpanded{#2}}}%
    \expanded{%
      \noexpand\newglossaryentry{e\theexcounter}{%
        name={\unexpanded{#2}%
          \IfBooleanF{#1}{ [\noexpand\texttt{e:\@exlblkey}]}},%
        sort={\theexample},%
        description={p.~\noexpand\pageref{e:\@exlblkey}},%
        user1={\theexample},%
        user2={Example},%
        user3={\thesection},%
        user4={\expandafter\unexpanded\expandafter{\@savedSecTitle}}%
      }%
    }%
    \protected@write\@auxout{}{\string\glsliveentry{e\theexcounter}}%
  \fi
}
\newcommand{\eeXa}{\end{example}}

% DEFINITION COMMAND -- same signature and behaviour as \beT below.
% Usage: \beD[long title][short title][key]
%   long title  = definition heading (plain text only)
%   short title = shown in \dref{word}; defaults to the key
%   key         = optional label key override; the label becomes d:KEY
% Without a key the label is d: plus the first three letters of the title's
% first word, so it is a mnemonic rather than the definition's number: a
% numeric label would silently change whenever a definition is inserted.
\NewDocumentCommand{\beD}{s O{} O{} o}{%
  \def\@temparg{#2}%
  \ifx\@temparg\@empty
    \begin{definition}%
  \else
    \let\@savedSecTitle\@currentsectionname%
    \refstepcounter{defcounter}%
    \IfValueTF{#4}{%
      \edef\@deflblkey{#4}%
    }{%
      \edef\@deffirstword{\@firstwordof#2 \relax}%
      \edef\@deflblkey{%
        \expandafter\@takethree\@deffirstword\@empty\@empty\@empty\@nil}%
    }%
    \def\@tempshort{#3}%
    \ifx\@tempshort\@empty
      \expandafter\xdef\csname @dshort@\@deflblkey\endcsname{\@deflblkey}%
    \else
      \expandafter\def\csname @dshort@\@deflblkey\endcsname{#3}%
    \fi
    \IfBooleanTF{#1}{\let\@deftag\@empty}{%
      \edef\@deftag{\noexpand\ \noexpand\textup{[\noexpand\texttt{d:\@deflblkey}]}}}%
    \begin{definition}[#2\ifshowkeys\@deftag\fi]%
    \label{d:\@deflblkey}%
    % record the long title so \nref{d:KEY} prints it, as for theorems
    \protected@write\@auxout{}{%
      \string\@thmtitle{d:\@deflblkey}{\unexpanded{#2}}}%
    \expanded{%
      \noexpand\newglossaryentry{d\thedefcounter}{%
        name={\unexpanded{#2}%
          \IfBooleanF{#1}{ [\noexpand\texttt{d:\@deflblkey}]}},%
        sort={\thedefinition},%
        description={p.~\noexpand\pageref{d:\@deflblkey}},%
        user1={\thedefinition},%
        user2={Definition},%
        user3={\thesection},%
        user4={\expandafter\unexpanded\expandafter{\@savedSecTitle}}%
      }%
    }%
    \protected@write\@auxout{}{\string\glsliveentry{d\thedefcounter}}%
  \fi
}
\newcommand{\eeD}{\end{definition}}

% \showkeysfalse in a driver hides the [t:KEY] mnemonics in result headings
% (labels and glossary entries are unchanged); used for submission builds.
\newif\ifshowkeys \showkeystrue
% \commonredtrue in a driver that inputs basicFiles/commonRed: tool.tex then
% cites its reduction results instead of restating them.
\newif\ifcommonred \commonredfalse
% TheorEM COMMAND
% Usage: \beT[long title][short title][key]
%   long title  = theorem heading (plain text only)
%   short title = shown in \tref{word}; may contain math; defaults to first word of long title
%   key         = optional label key override; label becomes t:KEY (avoids first-word collisions)
% Label is t:WORD where WORD = first word of long title (or KEY if given); use \tref{WORD} to reference.
\NewDocumentCommand{\beT}{s O{} O{} o}{%
  \def\@temparg{#2}%
  \ifx\@temparg\@empty
    \begin{theorem}%
  \else
    \let\@savedSecTitle\@currentsectionname%
    \refstepcounter{thmcounter}%
    % \@thmtag is the mnemonic shown in the heading; every theorem gets one.
    % With an explicit key the key itself is shown.  Without one the key
    % defaults to the title's first word, and showing that whole word would
    % merely repeat the title ("Dissipativity [Dissipativity]"), so only its
    % first three letters are shown: "Dissipativity [Dis]".  The label is
    % unaffected -- it still uses the full \@thmlblkey.
    \IfValueTF{#4}{%
      \edef\@thmlblkey{#4}%
    }{%
      \edef\@thmfirstword{\@firstwordof#2 \relax}%
      \edef\@thmlblkey{%
        \expandafter\@takethree\@thmfirstword\@empty\@empty\@empty\@nil}%
    }%
    \IfBooleanTF{#1}{\let\@thmtag\@empty}{%
      \edef\@thmtag{\noexpand\ \noexpand\textup{[\noexpand\texttt{t:\@thmlblkey}]}}}%
    \def\@tempshort{#3}%
    \ifx\@tempshort\@empty
      \expandafter\xdef\csname @tshort@\@thmlblkey\endcsname{\@thmlblkey}%
    \else
      \expandafter\def\csname @tshort@\@thmlblkey\endcsname{#3}%
    \fi
    \begin{theorem}[#2\ifshowkeys\@thmtag\fi]%
    \label{t:\@thmlblkey}%
    % record the long title for \nref; via the .aux so forward refs resolve
    \protected@write\@auxout{}{%
      \string\@thmtitle{t:\@thmlblkey}{\unexpanded{#2}}}%
    \expanded{%
      \noexpand\newglossaryentry{t\thethmcounter}{%
        name={\unexpanded{#2}%
          \IfBooleanF{#1}{ [\noexpand\texttt{t:\@thmlblkey}]}},%
        % sort/user1 must be the number the document actually prints, i.e. the
        % theorem environment's counter, not thmcounter.  thmcounter advances
        % only on \beT with a title, so a bare \beT (or \beT with an
        % empty title) makes the two drift and the glossary cites a wrong
        % number.  thmcounter is still fine as the entry key: it is unique.
        sort={\thetheorem},%
        description={p.~\noexpand\pageref{t:\@thmlblkey}},%
        user1={\thetheorem},%
        user2={Theorem},%
        user3={\thesection},%
        user4={\expandafter\unexpanded\expandafter{\@savedSecTitle}}%
      }%
    }%
    \protected@write\@auxout{}{\string\glsliveentry{t\thethmcounter}}%
  \fi
}
\newcommand{\eeT}{\end{theorem}}

% CROSS-PART RELATION PARAGRAPH
% Usage: \xpart{text}  -- place immediately BEFORE \beT, never inside it
%   (anything inside \beT...\eeT would be swallowed into the glossary entry).
% The paragraph is typeset only when Part I is actually present, i.e. in the
% combined book build; in the standalone builds (GranadaDSABNS/A.tex,
% AAHRah-Submission/AAHRah.tex) the label part:glv is undefined and the whole
% paragraph disappears rather than leaving ?? cross-references. Same
% graceful-degradation principle as \GASref in GASt.tex. Resolves on the
% second LaTeX pass, since r@part:glv comes from the .aux file.
% (Already inside the \makeatletter region opened above, so no \makeatletter
% pair here -- adding one would close that region early.)
\newcommand{\xpart}[1]{\@ifundefined{r@part:glv}{}{\paragraph{Relation with Parts~I--II.}#1}}

% \bookref{label}{precise text}{fallback text}: prints the precise cross-reference
% when the label resolves (combined book build), and the fallback text otherwise
% (standalone paper builds, where book-only appendix labels are absent). Same
% graceful-degradation role as \GASref in GASt.tex, but usable in any file since it
% lives in this preamble (loaded by every driver) -- unlike \GASref, which is defined in a
% content file and is therefore undefined in standalone builds that omit that file.
\newcommand{\bookref}[3]{\@ifundefined{r@#1}{#3}{#2}}

% \PVpart and \RBpart: the two parts that a partial build may omit.
% They print "Part~N" in any build that contains the part, and a naming phrase
% in a build that does not, so an excerpt shows no ?? .  Each already supplies
% the word "Part", so write "\PVpart develops ...", not "Part~\PVpart".
\newcommand{\PVpart}{\bookref{part:aahrah}%
  {Part~\ref{part:aahrah}}{the Perron--Volterra part of the full book}}
\newcommand{\RBpart}{\bookref{part:recbrowne}%
  {Part~\ref{part:recbrowne}}{the immuno-virology part of the full book}}

% OPEN PROBLEM COMMANDS
% ASSUMPTION COMMANDS (glossaried)
\newcommand{\beAs}[1][]{%
  \def\@temparg{#1}%
  \ifx\@temparg\@empty
    \begin{assumption}%
  \else
    \let\@savedSecTitle\@currentsectionname%
    \refstepcounter{ascounter}%
    \begin{assumption}[#1]%
    \label{as:\theascounter}%
    \expanded{%
      \noexpand\newglossaryentry{as\theascounter}{%
        name={\unexpanded{#1}},%
        sort={\theascounter},%
        description={p.~\noexpand\pageref{as:\theascounter}},%
        user1={\theascounter},%
        user2={Assumption},%
        user3={\thesection},%
        user4={\expandafter\unexpanded\expandafter{\@savedSecTitle}}%
      }%
    }%
    \protected@write\@auxout{}{\string\glsliveentry{as\theascounter}}%
  \fi
}
\newcommand{\eeAs}{\end{assumption}}
% CONJECTURE COMMANDS (glossaried)
\newcommand{\beCj}[1][]{%
  \def\@temparg{#1}%
  \ifx\@temparg\@empty
    \begin{conjecture}%
  \else
    \let\@savedSecTitle\@currentsectionname%
    \refstepcounter{cjcounter}%
    \begin{conjecture}[#1]%
    \label{cj:\thecjcounter}%
    \expanded{%
      \noexpand\newglossaryentry{cj\thecjcounter}{%
        name={\unexpanded{#1}},%
        sort={\thecjcounter},%
        description={p.~\noexpand\pageref{cj:\thecjcounter}},%
        user1={\thecjcounter},%
        user2={Conjecture},%
        user3={\thesection},%
        user4={\expandafter\unexpanded\expandafter{\@savedSecTitle}}%
      }%
    }%
    \protected@write\@auxout{}{\string\glsliveentry{cj\thecjcounter}}%
  \fi
}
\newcommand{\eeCj}{\end{conjecture}}
% LEMMA COMMANDS
% ---------------------------------------------------------------------
% UNIFORM RESULT WRAPPERS -- \beL, \beC, \beP, \beR share \beT's signature
%
%   \beX[long title][short title][key]  ...  \eeX
%
%   long title  = heading (math allowed)
%   short title = shown in the cross-reference; defaults to the key
%   key         = label key.  With a key the label is X:KEY and the heading
%                 carries the mnemonic [key], exactly as for \beT.  Without
%                 one the label falls back to the running counter, which is
%                 what every pre-existing call site in the book relies on.
%
% A bare \beX -- arguments commented out -- opens the plain environment and
% creates NO label and NO glossary entry.  That is the supported way to keep
% an individual result out of the catalogue.
% ---------------------------------------------------------------------
% #1 wrapper name  #2 environment  #3 label prefix  #4 counter  #5 type name
\newcommand{\@mkresult}[5]{%
  % A leading * suppresses the [key] mnemonic in both the heading and the
  % glossary name, while still making the label X:KEY and the entry:
  %     \beX*[long][short][key]
  \expandafter\NewDocumentCommand\csname be#1\endcsname{s O{} O{} o}{%
    \def\@temparg{##2}%
    \ifx\@temparg\@empty
      \begin{#2}%
    \else
      \let\@savedSecTitle\@currentsectionname%
      \refstepcounter{#4}%
      % With a key: mnemonic label X:KEY and the [key] tag in the heading.
      % Without one: fall back to the running counter, exactly as these
      % wrappers behaved before, so every legacy one-argument call site in
      % the book keeps its label.  The counter is also csname-safe, which a
      % title's first word is not -- titles beginning with math would
      % otherwise build \csname \Omega...\endcsname and fail.
      \IfValueTF{##4}{%
        \edef\@relbl{##4}%
        \IfBooleanTF{##1}{\let\@retag\@empty}{%
          \edef\@retag{\noexpand\ \noexpand\textup{[\noexpand\texttt{#3:##4}]}}}%
      }{%
        \edef\@relbl{\arabic{#4}}%
        \let\@retag\@empty
      }%
      \def\@tempshort{##3}%
      \ifx\@tempshort\@empty
        \expandafter\xdef\csname @sshort@#3\@relbl\endcsname{\@relbl}%
      \else
        \expandafter\def\csname @sshort@#3\@relbl\endcsname{##3}%
      \fi
      \begin{#2}[##2\ifshowkeys\@retag\fi]%
      \edef\@tmplbl{#3:\@relbl}\expandafter\label\expandafter{\@tmplbl}%
      \expanded{%
        \noexpand\newglossaryentry{#3\csname the#4\endcsname}{%
          name={\unexpanded{##2}%
            \IfBooleanF{##1}{\IfValueT{##4}{ [\noexpand\texttt{#3:##4}]}}},%
          sort={\csname the#2\endcsname},%
          description={p.~\noexpand\pageref{#3:\@relbl}},%
          user1={\csname the#2\endcsname},%
          user2={#5},%
          user3={\thesection},%
          user4={\expandafter\unexpanded\expandafter{\@savedSecTitle}}%
        }%
      }%
      \protected@write\@auxout{}{\string\glsliveentry{#3\csname the#4\endcsname}}%
      \protected@write\@auxout{}{\string\@thmtitle{\@tmplbl}{\unexpanded{##2}}}%
    \fi
  }%
  \expandafter\newcommand\csname ee#1\endcsname{\end{#2}}%
}
\@mkresult{L}{lemma}{l}{lemcounter}{Lemma}
% REMARK / COROLLARY / PROPOSITION -- same signature as \beL above
\@mkresult{R}{remark}{r}{remcounter}{Remark}
\@mkresult{C}{corollary}{c}{corcounter}{Corollary}
\@mkresult{P}{proposition}{p}{propcounter}{Proposition}
% OPEN-PROBLEM COMMAND -- generated with the same signature as \beT:
%   \beO[long title][short title][key]   ...   \eeO
% label o:KEY, or o:<counter> when no key is given; a bare \beO makes
% neither a label nor a glossary entry.
\@mkresult{O}{open}{o}{opcounter}{Problem}


% Section-level title tracker: updated only by \section, not \subsection.
% \@currentlabelname is overwritten by every \refstepcounter (including
% \subsection), so we save it separately here.
\newcommand{\@currentsectionname}{}
\AtBeginDocument{%
  \let\@epidcrn@orig@refstepcounter\refstepcounter
  \renewcommand{\refstepcounter}[1]{%
    \@epidcrn@orig@refstepcounter{#1}%
    \ifstrequal{#1}{section}{\global\let\@currentsectionname\@currentlabelname}{}%
  }%
}

% TABLE COMMANDS (glossaried) -- consistent naming with \beT, \beD, etc.
\newcommand{\beTa}[1][]{%
  \def\@temparg{#1}%
  \ifx\@temparg\@empty
    \begin{table}[htbp]%
  \else
    \let\@savedSecTitle\@currentsectionname%
    \refstepcounter{tabcounter}%
    \begin{table}[htbp]%
    \label{A:\thetabcounter}%
    \expanded{%
      \noexpand\newglossaryentry{A\thetabcounter}{%
        name={\unexpanded{#1}},%
        sort={\thetabcounter},%
        description={p.~\noexpand\pageref{A:\thetabcounter}},%
        user1={\thetabcounter},%
        user2={Table},%
        user3={\thesection},%
        user4={\expandafter\unexpanded\expandafter{\@savedSecTitle}}%
      }%
    }%
    \protected@write\@auxout{}{\string\glsliveentry{A\thetabcounter}}%
  \fi
}
\newcommand{\eeTa}{\end{table}}

% FIGURE COMMANDS (glossaried) -- consistent with \beT, \beR, \beTa, etc.
% Usage:
%   \beF[Glossary title]...\caption{Visible caption}\eeF  -> glossaried figure
%   \beF...\caption{Visible caption}\eeF                  -> plain [H] figure, no glossary
\newcommand{\beF}[1][]{%
  \def\@temparg{#1}%
  \ifx\@temparg\@empty
    \begin{figure}[H]%
  \else
    \begin{figure}[H]%
    \let\@savedSecTitle\@currentsectionname%
    \refstepcounter{figcounter}%
    \label{fig:\thefigcounter}%
    \expanded{%
      \noexpand\newglossaryentry{Fig\thefigcounter}{%
        name={\unexpanded{#1}},%
        sort={\thefigcounter},%
        description={p.~\noexpand\pageref{fig:\thefigcounter}},%
        user1={\thefigcounter},%
        user2={Figure},%
        user3={\thesection},%
        user4={\expandafter\unexpanded\expandafter{\@savedSecTitle}}%
      }%
    }%
    \protected@write\@auxout{}{\string\glsliveentry{Fig\thefigcounter}}%
  \fi
}
\newcommand{\eeF}{\end{figure}}
\makeatother
% LIVE-ENTRY REGISTER
% \glsaddall prints every entry the glossaries package currently knows, and
% that set includes entries restored from \jobname.glsdefs.  The package
% rewrites .glsdefs from all *currently defined* entries, so anything it
% restored is written straight back out: a deleted section kept reappearing in
% the catalogue for ever, and the .glsdefs had to be erased by hand.  The
% wrappers therefore record every entry they actually create in the .aux, and
% \printboth adds only those.  The .aux is rebuilt from scratch on each run, so
% this is self-healing: a removed result is gone after one further run.
\makeatletter
\newcommand{\@glslive}{}
\newcommand{\@glsaddlive}[1]{\ifglsentryexists{#1}{\glsadd{#1}}{}}
\newcommand{\glsliveentry}[1]{\g@addto@macro\@glslive{\@glsaddlive{#1}}}
\makeatother
% PRINT COMMAND
% No \clearpage here -- the catalogue continues on the current page.  This has
% been re-added by mistake several times; leave it out.
\makeatletter
\newcommand{\printboth}{%
  \@glslive   % add exactly the entries this document created last run
  \printglossary[style=unifiedstyle, title={0 \hspace{.22em} Definitions, Theorems, Examples and Open Problems}]%
}
\makeatother
% Named theorem reference: \tref{word} -> Thm~N(word-or-short)
% Shows stored short title if available, otherwise shows the key itself.
\newcommand{\tref}[1]{%
  Thm~\ref{t:#1}(\ifcsname @tshort@#1\endcsname\csname @tshort@#1\endcsname\else#1\fi)%
}

% Full-title theorem reference: \nref{t:KEY} prints the number and then the
% theorem's long title, so one writes "Theorem~\nref{t:blocks}" and gets
% "Theorem 3 ``Block factorization and siphons''".  Takes the complete label, like
% \ref, not the bare key like \tref.
%
% \beT writes each title to the .aux (see \@thmtitle below), so \nref resolves
% forward references from the second LaTeX pass, exactly as \ref does.  If the
% title is not yet known -- first pass, or a label that is not a \beT theorem
% -- it degrades to the number rather than producing ??.
% These two use @ in their names, and this point in the file is outside the
% \makeatletter region opened earlier (it closes at \makeatother above), so
% the region has to be reopened here.
\makeatletter
\newcommand{\@thmtitle}[2]{\global\@namedef{@tlong@#1}{#2}}
\newcommand{\nref}[1]{%
  \ref{#1}%
  \ifcsname @tlong@#1\endcsname
    \ ``\csname @tlong@#1\endcsname''%
  \fi
}
% \nm{t:KEY}: the stored long title of any \beT/\beL/\beP/\beR/\beC/\beD/\beE result,
% verbatim and without quotes, so item names follow the macro's first argument.
\newcommand{\nm}[1]{%
  \ifcsname @tlong@#1\endcsname
    \csname @tlong@#1\endcsname%
  \else
    \ref{#1}%
  \fi
}
\makeatother




%% ========================================
%% Other THEOREM ENVIRONMENTS
%% ========================================
\makeatletter
\@ifundefined{lemma}      {\newtheorem{lemma}      {Lemma}}      {}
\@ifundefined{proposition}{\newtheorem{proposition}{Proposition}} {}
\@ifundefined{corollary}  {\newtheorem{corollary}  {Corollary}}  {}
\@ifundefined{conjecture} {\newtheorem{conjecture} {Conjecture}}  {}
\@ifundefined{remark}     {\newtheorem{remark}     {Remark}}     {}
\@ifundefined{question}   {\newtheorem{question}   {Question}}   {}
\@ifundefined{assumption} {\newtheorem{assumption} {Assumption}} {}
\@ifundefined{hypothesis} {\newtheorem{hypothesis} {Hypothesis}} {}
\@ifundefined{fact}       {\newtheorem{fact}       {Fact}}       {}
\makeatother

% \citer: cite the relay-graph threshold theorem of AAHrel. If a document has
% \input the chapter/file containing t:relayRF (so the label resolves), point
% straight at it; otherwise fall back to a bare citation of AAHrel. Shared here
% (not per-paper) so it resolves consistently across every book chapter.
% \citeR: same idea, for the first theorem of AAHRah (t:Pstar).
\makeatletter
\providecommand{\citer}{%
  \@ifundefined{r@t:relayRF}%
    {\cite{AAHrel}}%
    {Theorem~\ref{t:relayRF} of \cite{AAHrel}}%
}
\providecommand{\citeR}{%
  \@ifundefined{r@t:Pstar}%
    {\cite{AAHRah}}%
    {Theorem~\ref{t:Pstar} of \cite{AAHRah}}%
}
\makeatother


% Abbreviations
\def\beQ{\begin{question}}\def\eeQ{\end{question}}
\def\beA{\begin{assumption}}\def\eeA{\end{assumption}}


%% ========================================

%% ========================================
%% COLORS & CUSTOM COMMANDS
%% ========================================
\definecolor{funccolor}{RGB}{25,25,112}
\definecolor{desccolor}{RGB}{64,64,64}
\newcommand{\func}[1]{\texttt{\textcolor{funccolor}{#1}}}
\def\eq{=}
\def\mm{-}

%% ========================================
%% MATHCAL SHORTHANDS
%% ========================================
% Project-wide convention: \mX for a capital X is always \mathcal X.
% Defined for the whole alphabet so the convention never fails to expand.
\newcommand{\mA}{\mathcal A}\newcommand{\mB}{\mathcal B}
\newcommand{\mC}{\mathcal C}\newcommand{\mD}{\mathcal D}
\newcommand{\mE}{\mathcal E}\newcommand{\mF}{\mathcal F}
\newcommand{\mG}{\mathcal G}\newcommand{\mH}{\mathcal H}
\newcommand{\mI}{\mathcal I}\newcommand{\mJ}{\mathcal J}
\newcommand{\mK}{\mathcal K}\newcommand{\mL}{\mathcal L}
\newcommand{\mM}{\mathcal M}\newcommand{\mN}{\mathcal N}
\newcommand{\mO}{\mathcal O}\newcommand{\mP}{\mathcal P}
\newcommand{\mQ}{\mathcal Q}\newcommand{\mR}{\mathcal R}
\newcommand{\mS}{\mathcal S}\newcommand{\mT}{\mathcal T}
\newcommand{\mU}{\mathcal U}\newcommand{\mV}{\mathcal V}
\newcommand{\mW}{\mathcal W}\newcommand{\mX}{\mathcal X}
\newcommand{\mY}{\mathcal Y}\newcommand{\mZ}{\mathcal Z}

%% ========================================
%% BASIC MATH ENVIRONMENTS
%% ========================================
\def\bep{\begin{pmatrix}}\def\eep{\end{pmatrix}}
\def\bev{\begin{vmatrix}}\def\eev{\end{vmatrix}}
\def\bea{\begin{eqnarray*}}\def\eea{\end{eqnarray*}}
\def\bc{\begin{cases}}\def\ec{\end{cases}}
\def\BEN{\begin{enumerate}}\def\EEN{\end{enumerate}}
\def\BI{\begin{itemize}}\def\EI{\end{itemize}}

\newcommand{\be}[1]{\begin{equation}\label{#1}}
\newcommand{\ee}{\end{equation}}
\newcommand{\beq}{\begin{eqnarray}}
\def\eeq{\end{eqnarray}}
\def\eqr{\eqref}\def\fr{\frac}\def\lbl{\label}
\def\Lra{\Longrightarrow}\def\non{\nonumber}\def\cNGM{\cite{Diek,Van,Van08}}
\def\Eq{\Leftrightarrow}\def\b{\beta}\def\uph{\red{UP TO HERE}}
%% ========================================
%% PROBABILITY & STATISTICS
%% ========================================
\renewcommand{\P}{\ensuremath{\mathbb{P}}}  % Note: overrides paragraph symbol
\newcommand{\E}{\ensuremath{\mathbb{E}}}
\newcommand{\ind}[1]{\ensuremath{\mathbbm{1}_{\left\{#1\right\}}}}
\DeclareMathOperator{\Cov}{Cov}
\DeclareMathOperator{\TV}{TV}
\newcommand{\rank}{\operatorname{rank}}

%% ========================================
%% CRN OPERATORS (organizations/ERCs/buffering/siphons framework)
%% ========================================
\def\clos{\mathrm{clos}}
\def\Supply{\mathrm{Supply}}
\def\req{\mathrm{req}}
\def\React{\mathrm{React}}
\def\Prod{\mathrm{Prod}}
\def\Gen{\mathrm{Gen}}
\def\In{\mathrm{In}}
\def\Out{\mathrm{Out}}
\def\Act{\mathrm{Act}}
\def\Min{\mathrm{Min}}

%% ========================================
%% LINE BREAKING
%% ========================================
% Last-resort stretch, used only on paragraphs TeX cannot otherwise break:
% a long \texttt{} token, a comma-separated inline math list (commas are not
% breakpoints in math), or a run of \cite/\ref boxes.  Without it those set as
% overfull \hbox and print into the margin; with it they set slightly loose,
% which is the lesser evil.  Does not affect paragraphs that already break.
\emergencystretch=3em

%% ========================================
%% NUMBER SETS
%% ========================================
\newcommand{\R}{\mathbb{R}}  % Real numbers
\newcommand{\C}{\mathbb{C}}  % Complex numbers
\newcommand{\N}{\mathbb{N}}
\newcommand{\Z}{\mathbb{Z}}
% \providecommand, not \newcommand: some papers define these in their own
% preamble after \input{preamble}, and must not clash.
\providecommand{\Q}{\mathbb{Q}}
\providecommand{\res}{\operatorname{res}}   % resultant; used by epiRes.tex
\newcommand{\MN}{\mathbb{N}}
\newcommand{\MZ}{\mathbb{Z}}
\newcommand{\realpos}{\mathbb{R}_{> 0}}
\newcommand{\realnneg}{\mathbb{R}_{\geq 0}}
\newcommand{\naturalpos}{\mathbb{N}_{>0}}
\newcommand{\Rplus}{\mathbb{R}_+}
\newcommand{\Rp}{\mathbb{R}_{\ge 0}}
\def\Rnne{\mathbb{R}_{\ge 0}}
\def\Se{\bar{S}} \def\Ie{\bar{I}}
\newcommand{\Rpos}{\mathbb{R}_{>0}}
\newcommand{\dd}{\,\mathrm d}
\newcommand{\e}{\mathrm e}
%% ========================================
%% GREEK LETTERS (2+ character abbreviations) greek
%% ========================================
\def\al{\alpha}\def\ba{\bff \alpha}
\def\w{\omega}\def\Lgn{\ga+\mu+\de}\def\om{\omega}
\def\de{\delta}\def\De{\Delta}\def\bde{\bff \delta}
\def\ep{\epsilon}\def\eps{\varepsilon}
\def\f{\varphi}  % Kept as requested
\def\ga{\gamma}\def\Ga{\Gamma}\def\G{\Gamma}\def\bg{\bff \gamma}
\def\up{\upsilon}  % NOT the breve accent: no accent commands anywhere.
                   % Named \up, not \u, so the LaTeX breve stays reachable.
\def\la{\lambda}\def\La{\Lambda}
\def\si{\sigma}\def\ka{\kappa}
\def\th{\theta}\def\vt{\vartheta}
\def\ze{\zeta}\def\Gap{\Gamma^{+}} \def\Gam{\Gamma^{-}}


%% ========================================
%% SPECIAL SHORT COMMANDS (kept as requested)
%% ========================================
\def\I{\infty}
\def\T{\widetilde}
\def\H{\hat}
\newcommand{\relint}{\operatorname{relint}}

%% ========================================
%% BOLD/VECTOR NOTATION
%% ========================================
\newcommand{\bff}[1]{{\mbox{\boldmath$#1$}}}
\newcommand\X{\boldsymbol{X}}
\def\x{\boldsymbol{x}}\def\bx{\bff x}\def\xe{\bar{\x}}
\newcommand\y{\boldsymbol{y}}\def\by{\boldsymbol{y}}\def\ye{\bar{\y}}
\def\yb{\y}\def\xb{\x}
\def\xy{\x-\y}
\newcommand\br{\boldsymbol{r}}\def\re{\bar{\r}}
\newcommand\bd{\boldsymbol{d}}
\newcommand\bQ{\boldsymbol{Q}}
\newcommand{\Yb}{{\mathbf{Y}}}
\newcommand{\Ab}{\mathbf{A}}
\newcommand{\zb}{\mathbf{z}}\def\bz{\bff 0}
\newcommand{\Zb}{\mathbf{Z}}
\newcommand{\bb}{\mathbf{b}}
\def\vbe{\vec {\bff \beta}}
\def\vb{\vec {\bff b}}
\def\v1{\vec {\bff 1}}
\def\vn{\vec \de}
\def\vz{\vec 0}
\def\vi{\overset{\rightarrow}{\mathsf i}}
\def\vv{\overset{\rightarrow}{v}}
\def\vr{\; \vec {\mathsf r}}
\newcommand{\va}{\bff a}
\providecommand{\vect}[1]{\boldsymbol{#1}}   % bold vector (moved from Gradostat.tex)
\providecommand{\one}{\mathbf 1}             % all-ones vector (moved from AliSF)
\providecommand{\spr}{s}                     % spectral bound  (moved from AliSF)
\providecommand{\Ke}{K_{\rm eff}}            % effective carrying capacity
% Critical right/left eigenvectors in the Hopf projection method (AAHOSN/New.tex)
\providecommand{\bq}{\mathbf{q}}
\providecommand{\bp}{\mathbf{p}}


%% ========================================
%% MATHCAL & MATHSCR
%% ========================================
% Mathcal
\newcommand\mc[1]{\mathcal{#1}}
\def\mC{{\mathcal C}}\def\mD{\mathcal D}\def\mE{\mathcal{E}}
\def\mF{\mathcal F}\def\mG{{\mathcal G}}\def\mH{{\mathcal H}}
\def\mI{\mathcal I}\def\mL{\mathcal{L}}\def\mM{{\mathcal M}}
\def\mN{{\mathcal N}}\def\mP{\mathcal{P}}\def\mR{{\mathcal R}}
\def\mS{{\mathcal S}}\def\mT{\mathcal{T}}\def\mV{{\mathcal V}}
\def\mK{\mathcal K}\def\mA{{\mathcal A}}\def\mU{{\mathcal U}}
% Symmetric bilinear and trilinear forms of the Taylor expansion, used in the
% first Lyapunov coefficient (AAHOSN/New.tex)
\providecommand{\Bform}{\mathcal{B}}
\providecommand{\Cform}{\mathcal{C}}

% Blackboard bold.  Kept here, not in the drivers: a \providecommand in a
% driver preamble runs before this file and silently wins, so the same name
% ends up meaning different things in different builds.
\def\IC{\mathbb{C}}\def\IN{\mathbb{N}}\def\IQ{\mathbb{Q}}
\def\IR{\mathbb{R}}\def\IZ{\mathbb{Z}}\def\IG{\mathbb{G}}
\def\II{\mathbb{I}}

% Mathscr
\newcommand{\Rs}{{\mathscr R}}
\newcommand{\Ss}{{\mathscr S}}
\newcommand{\Ns}{{\mathscr N}}
\newcommand{\Cs}{\mathscr C}


%% ========================================
%% OPERATORS & FUNCTIONS
%% ========================================
\DeclareMathOperator{\tr}{tr}
\DeclareMathOperator{\diag}{Diag}
\DeclareMathOperator{\Diag}{Diag}
\DeclareMathOperator{\intr}{\mathrm{Int}}
\DeclareMathOperator{\Inv}{Inv}
\DeclareMathOperator{\spec}{spec}
\DeclareMathOperator{\AUC}{AUC}               % area under the curve
\newcommand{\ii}{\mathrm{i}}                  % imaginary unit
\newcommand\op{\operatorname}
\newcommand{\supp}{\op{supp}}
\newcommand{\diff}{\mathop{}\mathopen{}\mathrm{d}}
\newcommand{\pdd}[1]{\frac{\partial}{\partial{#1}}}
\DeclareMathOperator{\cone}{cone}
\DeclareMathOperator{\proj}{proj}

%% ========================================
%% BRACKETS & NORMS
%% ========================================
\providecommand{\pr}[1]{\left(#1\right)}
\providecommand{\pp}[1]{\left[#1\right]}
\newcommand\croc[1]{\left\langle #1\right\rangle}
\newcommand{\norm}[1]{\|#1\|}
\newcommand{\inner}[2]{\langle #1, #2 \rangle}
% Short synonym of \inner, used by AAHOSN/New.tex as $\ip{u}{v}=\bar u^\top v$
\providecommand{\ip}[2]{\langle #1 , #2 \rangle}


%% ========================================
%% FORMATTING & UTILITIES
%% ========================================
\newcommand*{\Scale}[2][4]{\scalebox{#1}{$#2$}}
\newcommand{\te}[1]{\text{#1}}
\newcommand\steq[1]{\stackrel{\text{\rm #1.}}{=}}
\newcommand\blue[1]{\textcolor{blue}{#1}}
\newcommand{\red}{\textcolor[rgb]{1.00,0.00,0.00}}
\newcommand{\green}{\textcolor[rgb]{0.00,1.00,0.00}}

% Figure command
\newcommand{\figu}[3]{%
\begin{figure}[H]
\centering
\includegraphics[scale=#3]{#1}
\caption{#2\label{f:#1}}
\end{figure}
}

% Footnote
\long\def\symbolfootnote[#1]#2{%
\begingroup
\def\thefootnote{\fnsymbol{footnote}}\footnote[#1]{#2}%
\endgroup}
\def\fn{\symbolfootnote}

\newcommand{\dsum}{\displaystyle\sum}

%% ========================================
%% CRN & EPIDEMIOLOGY SPECIFIC
%% ========================================
% CRN notation
\def\SS{\mathcal S}
\def\CC{\mathcal C}
\def\RR{\mathcal R}

% Reproduction numbers
\def\Rto{R_{1}^{E_2}}
\def\Rot{R_{2}^{E_1}}
\def\Roc{R_{1,c}}
\def\Rtc{R_{2,c}}
\def\rd{r_{0}}
\def\sd{s_{0}}
\def\dd{D_{0}}
\def\vd{v_{0}}
\def\zd{z_{0}}
\def\se{\sE}
\def\xd{x_{0}}
%% ========================================
%% SPECIAL VALUES & EQUILIBRIA
%% ========================================
\def\eE{\bar{e}}
\def\wE{\bar{w}}
\def\sE{s_{*}}
\def\Me{\T M_{ee}}
\def\Rnne{\mathbb{R}_{\ge 0}}

\newcommand{\ones}{\mathbf 1}


%% ========================================
%% ABBREVIATIONS (alphabetical by category)
%% ========================================

% -- General mathematical terms --
\def\adm{admissible}\def\admy{admissibility}
%\def\and{antisymmetric}% disabled 2026-09-30: it overrode LaTeX's author separator \and (printed "antisymmetric" between authors)
\def\ass{assumption}
\def\cadlag{c\`adl\`ag}
\def\ch{characteristic polynomial}
\def\chp{characteristic polynomial}
\def\coe{coefficient}
\def\comp{computation}
\def\con{condition}
\def\coo{coordinate}
\def\corr{corresponding}
\def\dec{decomposition}
\def\den{denominator}
\def\dete{deterministic}
\def\eig{eigenvalue}\def\eval{eigenvalue}
\def\eigv{eigenvector}\def\eiv{eigenvector}\def\evec{eigenvector}
\def\elt{element}
\def\enn{essentially nonnegative/positivity preserving}
\def\equ{equation}
\def\etal{et al.}
\def\fe{for example}
\def\fno{from now on}
\def\for{\forall}
\def\form{formula}
\def\Fr{Furthermore, }
\def\frt{furthermore}
\def\fun{function}
\def\gene{generalization}
\def\heu{heuristic}
\def\how{however}\def\How{However, }
\def\ie{i.e. }
\def\im{\item}
\def\isd{is defined by}
\def\itc{it is easy to check that}
\def\itf{it follows that}\def\Itf{It follows that}
\def\ith{it holds that}
\def\itm{it may be checked that}\def\Itm{It may be checked that}
\def\ito{it turns out that}\def\Ito{It turns out that}\def\Itou{It turns out useful to}
\def\ma{mass action}\def\MAR{mass action representation}
\def\mat{matrix}
\def\mbc{may be checked to be}
\def\mbw{may be written as}\def\mubw{must be written as}
\def\mpr{more precisely, }\def\Mp{More precisely, }
\def\nc{non-catalytic}
\def\nne{non-negative}\def\nny{nonnegativity}
\def\no{\nonumber}
\def\noni{non-infectious}
\def\nota{notation}
\def\np{\newpage}
\def\Nt{Note that}
\def\nt{non trivial}\def\nz{non zero}
\def\oth{otherwise}
\def\Oth{On the other hand}
\def\pa{parameter}\def\para{parameter}
\def\pol{polynomial}
\def\prob{probability}
\def\Prf{{\bf  Proof:}}\def\prf{{\bf  Proof:}}
\def\prop{proportional to}
\def\qu{\quad}
\def\QED{\hfill {$\square$}\goodbreak \medskip}
\def\rep{representation}
\def\resp{respectively}
\def\satd{satisfied}\def\satg{satisfying}\def\sats{satisfies}\def\saty{satisfy}
\def\ssec{\subsection}\def\sssec{\subsubsection}
\def\sef{see for example}
\def\sksy{skew-symmetric}
\def\sub{substrate}
\def\sys{system}
\def\var{variable}
\def\wk{well-known}
\def\wl{without loss of generality}\def\wlo{w.l.o.g.}
\def\wmc{We may check that}\def\Wmc{We may check that}
\def\wmr{we may rewrite}
\def\wms{we must show that}
\def\wrt{with respect to}

% -- CRN & Chemistry --
\newcommand\CRN{chemical reaction network}
\newcommand\CRNs{chemical reaction networks}
\newcommand\RN{reaction network}
\def\ACR{absolute concentration robustness}
\def\AMSD{acyclic siphon decomposition}
\def\DAE{differential algebraic equations}
\def\FHJ{Feinberg-Horn-Jackson}
\def\GMAK{generalized mass action kinetics}
\def\MM{Michaelis-Menten}
\def\sm{stoichiometric matrix}\def\Sm{stoichiometric matrix}
\def\WRZD{weakly reversible, zero deficiency}\def\aut{autocatalytic}

% -- Epidemiology --
\def\brn{basic reproduction number}\def\bzn{basic replacement number}
\def\com{compartment}
\def\dfe{disease free equilibrium}\def\DFE{disease free equilibrium}
\def\dwt{dwell times}\def\edt{expected dwell times}
\def\EM{epidemic model}
\def\endp{endemic fixed point}\def\enp{endemic point}\def\EE{endemic point}
\def\ME{mathematical epidemiology}
\def\MEE{eco-epidemiology}\def\EEP{eco-epidemiological}
\def\ece{eco-epidemiological model}
\def\mgm{multi group model}
\def\mtS{multi-strain model}\def\mSt{multi-strain model}
\def\mSi{minimal siphon}
\def\NGM{next generation matrix}
\def\nim{new infections matrix}
\def\qem{quadratic epidemiological model}
\def\repF{invasion functions}
\def\SPF{SIR-PH-FA model}
\def\sus{susceptible}

% -- Graph theory & Linear algebra --
\def\jac{Jacobi matrix}
\def\jin{Jacobian matrix of the invasion vector field with respect to the invasion variables}
\def\jrn{Jacobian matrix of the resident vector field with respect to the resident variables}
\def\jxn{Jacobian matrix of the $\x$ vector field with respect to the $\x$ variables}
\def\Lm{Laplacian matrix}
\def\MTT{matrix tree theorem}
\def\SGT{spectral graph theory}
\def\SMW{Sherman–Morrison–Woodbury}

% -- Stability & Dynamics --
\def\absC{absorbtion CTMC}
\def\BTb{Bogdanov-Takens bifurcation}
\def\CEP{competitive exclusion partition}
\def\CTMC{continuous time Markov chain}
\def\ds{differential system}
\def\fp{fixed point}\def\bfp{boundary fixed point}\def\fbp{fixed boundary point}
\def\fbl{fixed boundary line}\def\fif{forward invariant face}
\def\gas{global asymptotic stability}\def\GAS{globally asymptotically stable}
\def\GASS{\text{GASS}}
\def\Hb{Hopf bifurcation}
\def\Hur{Hurwitz}\def\MH{Metzler-Hurwitz}\def\RH{Routh-Hurwitz}
\def\IS{\text{IS}}
\def\LAS{locally asymptotically stable}
\def\LdS{Lyapunov diagonally stable}
\def\Lf{Lyapunov function}\def\LyV{Lyapunov Volterra type function}\def\RLF{robust Lyapunov function}
\def\LyM{Lyapunov Malkin stable}
\def\PF{Perron-Frobenius}
\def\spfp{strictly positive fixed point}
\def\sta{stability}
\def\STP{sharp threshold property}\def\STT{sharp threshold theorem}

% -- Models (with \varphi explicitly instead of \f) --
\def\AB{$(A,B)$ Arino-Brauer epidemic models}
\def\ABa{$(A,\bbe,\ba,\varphi)$ balanced bilinear model}
\def\ABP{%$(A,\bbe,P,\varphi)$
\\ balanced bilinear model}
\def\AMG{Anderson May Gupta}
\newcommand\BIN{biological interaction network}
\def\Dt{Descartes type}
\def\gLV{generalized Lotka-Volterra model}
\def\IG{\text{invasion graph}}\def\IGMS{interaction graph on minimal siphons}
\def\LCP{\text{linear complementarity problem}}
\def\LV{Lotka-Volterra}\def\LVK{Lotka-Volterra-Kolmogorov}

% -- Splitting & Structure --
\def\reg{regular splitting}\def\regS{{\bf regular splitting}}
\def\RUR{rational univariate representation}
\def\bds{block-diagonal structure}\def\NGMs{next generation matrices}

% -- Miscellaneous --
\def\Ma{{Mathematica}}
\def\Mm{Max–Min}
\def\num{numerator}
\def\rap{rational approximation} \def\DSR{directed species-reaction bipartite graph}\def\PCP{polynomial complementarity problem}\def\rat{rational substitution}
\def\QSS{quasi steady state} \def\RHS{right hand side} \def\LHS{left hand side}
\def\OSN{online social network}\def\VdD{Van den Driessche}
\def\wli{$\omega$-limit} \def\rk{rich kinetics}\def\sJ{symbolic Jacobian}
\def\cV{\cite{Vas,VasSta,BSV,GSV}} \def\spf{semi-parametric semi-functional}
\def\CRW{Capasso-Ruan-Wang} \def\ABP{$(A,\bbe,P,\f)$ balanced bilinear model}
\def\ABa{%$(A,\bbe,\ba,\f)$ \\
rank-one balanced bilinear model}
\def\bbe{B%{\bff \beta}
} \def\vbe{\vec {\bff \beta}
}
\def\vb{\vec {\bff b}}
\newcommand{\Kh}{\hat{K}}          % K with hat, first mention only
\newcommand{\wS}{w^{\sigma}}       % right Perron J^\perp, full block
\newcommand{\piS}{\vec \pi^{\sigma}}    % left Perron J^\perp, full block
\newcommand{\ws}{w^{\sigma}}       % right Perron J^\perp, sub-block
\newcommand{\pis}{\pi^{\sigma}}    % left Perron J^\perp, sub-block
\newcommand{\hwS}{\hat{w}^{\sigma}}   % right Perron K, full block
\newcommand{\hpiS}{\hat{\pi}^{\sigma}} % left Perron K, full block
\newcommand{\hws}{\hat{w}^{\sigma}}   % right Perron K, sub-block
\newcommand{\hpis}{\hat{\pi}^{\sigma}} % left Perron K, sub-block
\newcommand{\twS}{\tilde{w}^{\sigma}}  % right Perron \tilde K, full block
\newcommand{\tpiS}{\tilde{\pi}^{\sigma}} % left Perron \tilde K
\newcommand{\tws}{\tilde{w}^{\sigma}}
\newcommand{\tpis}{\tilde{\pi}^{\sigma}}
\newcommand{\Mat}[2]{\mathbb{R}_{\ge 0}^{\,#1\times #2}}
\def\tJ{J_\sigma^\perp} \def\lPe{left Perron eigenvector}
% \G is \Gamma (defined at :795).  The Volterra function keeps its own name,
% so that \G no longer has two conflicting meanings in this file.
\def\GV{G}  % Volterra G(u)=u-1-\ln u
\def\PVL{Perron-Volterra Lyapunov minimal siphon ansatz}
\newcommand{\istar}[1]{i_{#1}^{*}}
%\newcommand{\sE}{s_E}
\newcommand{\vE}{v_E}
\def\PV{Perron-Volterra candidate Lyapunov function}
\def\GLV{generalized Lotka-Volterra}
\def\MERN{mathematical epidemiology reaction networks}
\def\CMM{contraction metric method}
\newcommand{\sign}{\operatorname{sign}}
\newcommand{\Var}{\operatorname{Var}} 

% ======================================================================
% IDEAS SHARED BETWEEN THE ABSTRACT AND THE MAIN RESULTS SECTION
%
% Write each idea ONCE, in the abstract; the Main Results section recalls it,
% so the two can never drift apart.
%
%   in the abstract        \ideaitem{1}{The first-order reaction terms ...}
%   in Main Results        \ideahead{1}
%
% \ideaitem stores the text globally and typesets
%   \item \hyperlink{idea:1}{<text>}
% \ideahead typesets
%   \item\hypertarget{idea:1}{}\textbf{<the same text>}
%
% The abstract must come before the Main Results input, which is the normal
% order. If a heading has to read differently from the abstract sentence (for
% instance a long abstract item summarised in one line), override it with the
% optional argument: \ideahead[short heading]{6}.
% ======================================================================
\makeatletter
\newcommand{\ideaitem}[2]{%
  \expandafter\gdef\csname @idea@#1\endcsname{#2}%
  \item\hyperlink{idea:#1}{#2}}
\newcommand{\ideahead}[2][]{%
  \item\hypertarget{idea:#2}{}%
  \def\@ideatmp{#1}%
  \ifx\@ideatmp\@empty
    \@ifundefined{@idea@#2}{%
      \PackageWarning{preamble}{Idea #2 is used in Main Results but was never
        defined by \string\ideaitem\space in the abstract}%
      \textbf{[idea #2 undefined]}%
    }{\textbf{\begingroup\def\ideaterm##1##2{##2}\def\idearest##1{}\@nameuse{@idea@#2}\endgroup}}%
  \else
    \textbf{#1}%
  \fi}
% \ideaterm{<hypertarget>}{<words>}: inside an abstract idea the words become an
% invisible link (the document sets \hypersetup{hidelinks}) jumping to the place
% in Main Results that defines them; inside the Main Results heading itself,
% where the target sits a few lines below, \ideahead strips the link and prints
% the words plainly. Mark the destination with \hypertarget{<name>}{} on the line
% before the display, which adds no equation number and shows nothing.
\newcommand{\ideaterm}[2]{\hyperlink{#1}{#2}}
% \idearest{<text>}: the part of an abstract idea after its first sentence; printed in
% the abstract, omitted from the Main Results heading.
\newcommand{\idearest}[1]{#1}
\makeatother

% ======================================================================
% SHORT-PAPER SWITCH
% \shorttrue in a master that assembles only part of the companion paper
% (Barrat/SipArxiv.tex). \nsh{...} drops text that cites results living in the
% omitted sections, so the short build has no dangling references while the
% companion paper keeps the text verbatim.
% ======================================================================
\newif\ifshort \shortfalse
\newcommand{\nsh}[1]{\ifshort\else#1\fi}
% \shortalt{<short build>}{<companion paper>}: use when a sentence needs
% different wording in the two builds, not merely a deletion.
\newcommand{\shortalt}[2]{\ifshort#1\else#2\fi}
% \Jtr, \JtrS{\Sigma}: transverse Jacobian. J^\perp here; SipArxiv.tex redefines them as J_{\mathrm{tra}}, J_{\mathrm{tra},\Sigma}
\newcommand{\Jtr}{J^\perp}
\newcommand{\JtrS}[1]{J_{#1}^\perp}
, and must not clash.
\providecommand{\Q}{\mathbb{Q}}
\providecommand{\res}{\operatorname{res}}   % resultant; used by epiRes.tex
\newcommand{\MN}{\mathbb{N}}
\newcommand{\MZ}{\mathbb{Z}}
\newcommand{\realpos}{\mathbb{R}_{> 0}}
\newcommand{\realnneg}{\mathbb{R}_{\geq 0}}
\newcommand{\naturalpos}{\mathbb{N}_{>0}}
\newcommand{\Rplus}{\mathbb{R}_+}
\newcommand{\Rp}{\mathbb{R}_{\ge 0}}
\def\Rnne{\mathbb{R}_{\ge 0}}
\def\Se{\bar{S}} \def\Ie{\bar{I}}
\newcommand{\Rpos}{\mathbb{R}_{>0}}
\newcommand{\dd}{\,\mathrm d}
\newcommand{\e}{\mathrm e}
%% ========================================
%% GREEK LETTERS (2+ character abbreviations) greek
%% ========================================
\def\al{\alpha}\def\ba{\bff \alpha}
\def\w{\omega}\def\Lgn{\ga+\mu+\de}\def\om{\omega}
\def\de{\delta}\def\De{\Delta}\def\bde{\bff \delta}
\def\ep{\epsilon}\def\eps{\varepsilon}
\def\f{\varphi}  % Kept as requested
\def\ga{\gamma}\def\Ga{\Gamma}\def\G{\Gamma}\def\bg{\bff \gamma}
\def\up{\upsilon}  % NOT the breve accent: no accent commands anywhere.
                   % Named \up, not \u, so the LaTeX breve stays reachable.
\def\la{\lambda}\def\La{\Lambda}
\def\si{\sigma}\def\ka{\kappa}
\def\th{\theta}\def\vt{\vartheta}
\def\ze{\zeta}\def\Gap{\Gamma^{+}} \def\Gam{\Gamma^{-}}


%% ========================================
%% SPECIAL SHORT COMMANDS (kept as requested)
%% ========================================
\def\I{\infty}
\def\T{\widetilde}
\def\H{\hat}
\newcommand{\relint}{\operatorname{relint}}

%% ========================================
%% BOLD/VECTOR NOTATION
%% ========================================
\newcommand{\bff}[1]{{\mbox{\boldmath$#1$}}}
\newcommand\X{\boldsymbol{X}}
\def\x{\boldsymbol{x}}\def\bx{\bff x}\def\xe{\bar{\x}}
\newcommand\y{\boldsymbol{y}}\def\by{\boldsymbol{y}}\def\ye{\bar{\y}}
\def\yb{\y}\def\xb{\x}
\def\xy{\x-\y}
\newcommand\br{\boldsymbol{r}}\def\re{\bar{\r}}
\newcommand\bd{\boldsymbol{d}}
\newcommand\bQ{\boldsymbol{Q}}
\newcommand{\Yb}{{\mathbf{Y}}}
\newcommand{\Ab}{\mathbf{A}}
\newcommand{\zb}{\mathbf{z}}\def\bz{\bff 0}
\newcommand{\Zb}{\mathbf{Z}}
\newcommand{\bb}{\mathbf{b}}
\def\vbe{\vec {\bff \beta}}
\def\vb{\vec {\bff b}}
\def\v1{\vec {\bff 1}}
\def\vn{\vec \de}
\def\vz{\vec 0}
\def\vi{\overset{\rightarrow}{\mathsf i}}
\def\vv{\overset{\rightarrow}{v}}
\def\vr{\; \vec {\mathsf r}}
\newcommand{\va}{\bff a}
\providecommand{\vect}[1]{\boldsymbol{#1}}   % bold vector (moved from Gradostat.tex)
\providecommand{\one}{\mathbf 1}             % all-ones vector (moved from AliSF)
\providecommand{\spr}{s}                     % spectral bound  (moved from AliSF)
\providecommand{\Ke}{K_{\rm eff}}            % effective carrying capacity
% Critical right/left eigenvectors in the Hopf projection method (AAHOSN/New.tex)
\providecommand{\bq}{\mathbf{q}}
\providecommand{\bp}{\mathbf{p}}


%% ========================================
%% MATHCAL & MATHSCR
%% ========================================
% Mathcal
\newcommand\mc[1]{\mathcal{#1}}
\def\mC{{\mathcal C}}\def\mD{\mathcal D}\def\mE{\mathcal{E}}
\def\mF{\mathcal F}\def\mG{{\mathcal G}}\def\mH{{\mathcal H}}
\def\mI{\mathcal I}\def\mL{\mathcal{L}}\def\mM{{\mathcal M}}
\def\mN{{\mathcal N}}\def\mP{\mathcal{P}}\def\mR{{\mathcal R}}
\def\mS{{\mathcal S}}\def\mT{\mathcal{T}}\def\mV{{\mathcal V}}
\def\mK{\mathcal K}\def\mA{{\mathcal A}}\def\mU{{\mathcal U}}
% Symmetric bilinear and trilinear forms of the Taylor expansion, used in the
% first Lyapunov coefficient (AAHOSN/New.tex)
\providecommand{\Bform}{\mathcal{B}}
\providecommand{\Cform}{\mathcal{C}}

% Blackboard bold.  Kept here, not in the drivers: a \providecommand in a
% driver preamble runs before this file and silently wins, so the same name
% ends up meaning different things in different builds.
\def\IC{\mathbb{C}}\def\IN{\mathbb{N}}\def\IQ{\mathbb{Q}}
\def\IR{\mathbb{R}}\def\IZ{\mathbb{Z}}\def\IG{\mathbb{G}}
\def\II{\mathbb{I}}

% Mathscr
\newcommand{\Rs}{{\mathscr R}}
\newcommand{\Ss}{{\mathscr S}}
\newcommand{\Ns}{{\mathscr N}}
\newcommand{\Cs}{\mathscr C}


%% ========================================
%% OPERATORS & FUNCTIONS
%% ========================================
\DeclareMathOperator{\tr}{tr}
\DeclareMathOperator{\diag}{Diag}
\DeclareMathOperator{\Diag}{Diag}
\DeclareMathOperator{\intr}{\mathrm{Int}}
\DeclareMathOperator{\Inv}{Inv}
\DeclareMathOperator{\spec}{spec}
\DeclareMathOperator{\AUC}{AUC}               % area under the curve
\newcommand{\ii}{\mathrm{i}}                  % imaginary unit
\newcommand\op{\operatorname}
\newcommand{\supp}{\op{supp}}
\newcommand{\diff}{\mathop{}\mathopen{}\mathrm{d}}
\newcommand{\pdd}[1]{\frac{\partial}{\partial{#1}}}
\DeclareMathOperator{\cone}{cone}
\DeclareMathOperator{\proj}{proj}

%% ========================================
%% BRACKETS & NORMS
%% ========================================
\providecommand{\pr}[1]{\left(#1\right)}
\providecommand{\pp}[1]{\left[#1\right]}
\newcommand\croc[1]{\left\langle #1\right\rangle}
\newcommand{\norm}[1]{\|#1\|}
\newcommand{\inner}[2]{\langle #1, #2 \rangle}
% Short synonym of \inner, used by AAHOSN/New.tex as $\ip{u}{v}=\bar u^\top v$
\providecommand{\ip}[2]{\langle #1 , #2 \rangle}


%% ========================================
%% FORMATTING & UTILITIES
%% ========================================
\newcommand*{\Scale}[2][4]{\scalebox{#1}{$#2$}}
\newcommand{\te}[1]{\text{#1}}
\newcommand\steq[1]{\stackrel{\text{\rm #1.}}{=}}
\newcommand\blue[1]{\textcolor{blue}{#1}}
\newcommand{\red}{\textcolor[rgb]{1.00,0.00,0.00}}
\newcommand{\green}{\textcolor[rgb]{0.00,1.00,0.00}}

% Figure command
\newcommand{\figu}[3]{%
\begin{figure}[H]
\centering
\includegraphics[scale=#3]{#1}
\caption{#2\label{f:#1}}
\end{figure}
}

% Footnote
\long\def\symbolfootnote[#1]#2{%
\begingroup
\def\thefootnote{\fnsymbol{footnote}}\footnote[#1]{#2}%
\endgroup}
\def\fn{\symbolfootnote}

\newcommand{\dsum}{\displaystyle\sum}

%% ========================================
%% CRN & EPIDEMIOLOGY SPECIFIC
%% ========================================
% CRN notation
\def\SS{\mathcal S}
\def\CC{\mathcal C}
\def\RR{\mathcal R}

% Reproduction numbers
\def\Rto{R_{1}^{E_2}}
\def\Rot{R_{2}^{E_1}}
\def\Roc{R_{1,c}}
\def\Rtc{R_{2,c}}
\def\rd{r_{0}}
\def\sd{s_{0}}
\def\dd{D_{0}}
\def\vd{v_{0}}
\def\zd{z_{0}}
\def\se{\sE}
\def\xd{x_{0}}
%% ========================================
%% SPECIAL VALUES & EQUILIBRIA
%% ========================================
\def\eE{\bar{e}}
\def\wE{\bar{w}}
\def\sE{s_{*}}
\def\Me{\T M_{ee}}
\def\Rnne{\mathbb{R}_{\ge 0}}

\newcommand{\ones}{\mathbf 1}


%% ========================================
%% ABBREVIATIONS (alphabetical by category)
%% ========================================

% -- General mathematical terms --
\def\adm{admissible}\def\admy{admissibility}
%\def\and{antisymmetric}% disabled 2026-09-30: it overrode LaTeX's author separator \and (printed "antisymmetric" between authors)
\def\ass{assumption}
\def\cadlag{c\`adl\`ag}
\def\ch{characteristic polynomial}
\def\chp{characteristic polynomial}
\def\coe{coefficient}
\def\comp{computation}
\def\con{condition}
\def\coo{coordinate}
\def\corr{corresponding}
\def\dec{decomposition}
\def\den{denominator}
\def\dete{deterministic}
\def\eig{eigenvalue}\def\eval{eigenvalue}
\def\eigv{eigenvector}\def\eiv{eigenvector}\def\evec{eigenvector}
\def\elt{element}
\def\enn{essentially nonnegative/positivity preserving}
\def\equ{equation}
\def\etal{et al.}
\def\fe{for example}
\def\fno{from now on}
\def\for{\forall}
\def\form{formula}
\def\Fr{Furthermore, }
\def\frt{furthermore}
\def\fun{function}
\def\gene{generalization}
\def\heu{heuristic}
\def\how{however}\def\How{However, }
\def\ie{i.e. }
\def\im{\item}
\def\isd{is defined by}
\def\itc{it is easy to check that}
\def\itf{it follows that}\def\Itf{It follows that}
\def\ith{it holds that}
\def\itm{it may be checked that}\def\Itm{It may be checked that}
\def\ito{it turns out that}\def\Ito{It turns out that}\def\Itou{It turns out useful to}
\def\ma{mass action}\def\MAR{mass action representation}
\def\mat{matrix}
\def\mbc{may be checked to be}
\def\mbw{may be written as}\def\mubw{must be written as}
\def\mpr{more precisely, }\def\Mp{More precisely, }
\def\nc{non-catalytic}
\def\nne{non-negative}\def\nny{nonnegativity}
\def\no{\nonumber}
\def\noni{non-infectious}
\def\nota{notation}
\def\np{\newpage}
\def\Nt{Note that}
\def\nt{non trivial}\def\nz{non zero}
\def\oth{otherwise}
\def\Oth{On the other hand}
\def\pa{parameter}\def\para{parameter}
\def\pol{polynomial}
\def\prob{probability}
\def\Prf{{\bf  Proof:}}\def\prf{{\bf  Proof:}}
\def\prop{proportional to}
\def\qu{\quad}
\def\QED{\hfill {$\square$}\goodbreak \medskip}
\def\rep{representation}
\def\resp{respectively}
\def\satd{satisfied}\def\satg{satisfying}\def\sats{satisfies}\def\saty{satisfy}
\def\ssec{\subsection}\def\sssec{\subsubsection}
\def\sef{see for example}
\def\sksy{skew-symmetric}
\def\sub{substrate}
\def\sys{system}
\def\var{variable}
\def\wk{well-known}
\def\wl{without loss of generality}\def\wlo{w.l.o.g.}
\def\wmc{We may check that}\def\Wmc{We may check that}
\def\wmr{we may rewrite}
\def\wms{we must show that}
\def\wrt{with respect to}

% -- CRN & Chemistry --
\newcommand\CRN{chemical reaction network}
\newcommand\CRNs{chemical reaction networks}
\newcommand\RN{reaction network}
\def\ACR{absolute concentration robustness}
\def\AMSD{acyclic siphon decomposition}
\def\DAE{differential algebraic equations}
\def\FHJ{Feinberg-Horn-Jackson}
\def\GMAK{generalized mass action kinetics}
\def\MM{Michaelis-Menten}
\def\sm{stoichiometric matrix}\def\Sm{stoichiometric matrix}
\def\WRZD{weakly reversible, zero deficiency}\def\aut{autocatalytic}

% -- Epidemiology --
\def\brn{basic reproduction number}\def\bzn{basic replacement number}
\def\com{compartment}
\def\dfe{disease free equilibrium}\def\DFE{disease free equilibrium}
\def\dwt{dwell times}\def\edt{expected dwell times}
\def\EM{epidemic model}
\def\endp{endemic fixed point}\def\enp{endemic point}\def\EE{endemic point}
\def\ME{mathematical epidemiology}
\def\MEE{eco-epidemiology}\def\EEP{eco-epidemiological}
\def\ece{eco-epidemiological model}
\def\mgm{multi group model}
\def\mtS{multi-strain model}\def\mSt{multi-strain model}
\def\mSi{minimal siphon}
\def\NGM{next generation matrix}
\def\nim{new infections matrix}
\def\qem{quadratic epidemiological model}
\def\repF{invasion functions}
\def\SPF{SIR-PH-FA model}
\def\sus{susceptible}

% -- Graph theory & Linear algebra --
\def\jac{Jacobi matrix}
\def\jin{Jacobian matrix of the invasion vector field with respect to the invasion variables}
\def\jrn{Jacobian matrix of the resident vector field with respect to the resident variables}
\def\jxn{Jacobian matrix of the $\x$ vector field with respect to the $\x$ variables}
\def\Lm{Laplacian matrix}
\def\MTT{matrix tree theorem}
\def\SGT{spectral graph theory}
\def\SMW{Sherman–Morrison–Woodbury}

% -- Stability & Dynamics --
\def\absC{absorbtion CTMC}
\def\BTb{Bogdanov-Takens bifurcation}
\def\CEP{competitive exclusion partition}
\def\CTMC{continuous time Markov chain}
\def\ds{differential system}
\def\fp{fixed point}\def\bfp{boundary fixed point}\def\fbp{fixed boundary point}
\def\fbl{fixed boundary line}\def\fif{forward invariant face}
\def\gas{global asymptotic stability}\def\GAS{globally asymptotically stable}
\def\GASS{\text{GASS}}
\def\Hb{Hopf bifurcation}
\def\Hur{Hurwitz}\def\MH{Metzler-Hurwitz}\def\RH{Routh-Hurwitz}
\def\IS{\text{IS}}
\def\LAS{locally asymptotically stable}
\def\LdS{Lyapunov diagonally stable}
\def\Lf{Lyapunov function}\def\LyV{Lyapunov Volterra type function}\def\RLF{robust Lyapunov function}
\def\LyM{Lyapunov Malkin stable}
\def\PF{Perron-Frobenius}
\def\spfp{strictly positive fixed point}
\def\sta{stability}
\def\STP{sharp threshold property}\def\STT{sharp threshold theorem}

% -- Models (with \varphi explicitly instead of \f) --
\def\AB{$(A,B)$ Arino-Brauer epidemic models}
\def\ABa{$(A,\bbe,\ba,\varphi)$ balanced bilinear model}
\def\ABP{%$(A,\bbe,P,\varphi)$
\\ balanced bilinear model}
\def\AMG{Anderson May Gupta}
\newcommand\BIN{biological interaction network}
\def\Dt{Descartes type}
\def\gLV{generalized Lotka-Volterra model}
\def\IG{\text{invasion graph}}\def\IGMS{interaction graph on minimal siphons}
\def\LCP{\text{linear complementarity problem}}
\def\LV{Lotka-Volterra}\def\LVK{Lotka-Volterra-Kolmogorov}

% -- Splitting & Structure --
\def\reg{regular splitting}\def\regS{{\bf regular splitting}}
\def\RUR{rational univariate representation}
\def\bds{block-diagonal structure}\def\NGMs{next generation matrices}

% -- Miscellaneous --
\def\Ma{{Mathematica}}
\def\Mm{Max–Min}
\def\num{numerator}
\def\rap{rational approximation} \def\DSR{directed species-reaction bipartite graph}\def\PCP{polynomial complementarity problem}\def\rat{rational substitution}
\def\QSS{quasi steady state} \def\RHS{right hand side} \def\LHS{left hand side}
\def\OSN{online social network}\def\VdD{Van den Driessche}
\def\wli{$\omega$-limit} \def\rk{rich kinetics}\def\sJ{symbolic Jacobian}
\def\cV{\cite{Vas,VasSta,BSV,GSV}} \def\spf{semi-parametric semi-functional}
\def\CRW{Capasso-Ruan-Wang} \def\ABP{$(A,\bbe,P,\f)$ balanced bilinear model}
\def\ABa{%$(A,\bbe,\ba,\f)$ \\
rank-one balanced bilinear model}
\def\bbe{B%{\bff \beta}
} \def\vbe{\vec {\bff \beta}
}
\def\vb{\vec {\bff b}}
\newcommand{\Kh}{\hat{K}}          % K with hat, first mention only
\newcommand{\wS}{w^{\sigma}}       % right Perron J^\perp, full block
\newcommand{\piS}{\vec \pi^{\sigma}}    % left Perron J^\perp, full block
\newcommand{\ws}{w^{\sigma}}       % right Perron J^\perp, sub-block
\newcommand{\pis}{\pi^{\sigma}}    % left Perron J^\perp, sub-block
\newcommand{\hwS}{\hat{w}^{\sigma}}   % right Perron K, full block
\newcommand{\hpiS}{\hat{\pi}^{\sigma}} % left Perron K, full block
\newcommand{\hws}{\hat{w}^{\sigma}}   % right Perron K, sub-block
\newcommand{\hpis}{\hat{\pi}^{\sigma}} % left Perron K, sub-block
\newcommand{\twS}{\tilde{w}^{\sigma}}  % right Perron \tilde K, full block
\newcommand{\tpiS}{\tilde{\pi}^{\sigma}} % left Perron \tilde K
\newcommand{\tws}{\tilde{w}^{\sigma}}
\newcommand{\tpis}{\tilde{\pi}^{\sigma}}
\newcommand{\Mat}[2]{\mathbb{R}_{\ge 0}^{\,#1\times #2}}
\def\tJ{J_\sigma^\perp} \def\lPe{left Perron eigenvector}
% \G is \Gamma (defined at :795).  The Volterra function keeps its own name,
% so that \G no longer has two conflicting meanings in this file.
\def\GV{G}  % Volterra G(u)=u-1-\ln u
\def\PVL{Perron-Volterra Lyapunov minimal siphon ansatz}
\newcommand{\istar}[1]{i_{#1}^{*}}
%\newcommand{\sE}{s_E}
\newcommand{\vE}{v_E}
\def\PV{Perron-Volterra candidate Lyapunov function}
\def\GLV{generalized Lotka-Volterra}
\def\MERN{mathematical epidemiology reaction networks}
\def\CMM{contraction metric method}
\newcommand{\sign}{\operatorname{sign}}
\newcommand{\Var}{\operatorname{Var}} 

% ======================================================================
% IDEAS SHARED BETWEEN THE ABSTRACT AND THE MAIN RESULTS SECTION
%
% Write each idea ONCE, in the abstract; the Main Results section recalls it,
% so the two can never drift apart.
%
%   in the abstract        \ideaitem{1}{The first-order reaction terms ...}
%   in Main Results        \ideahead{1}
%
% \ideaitem stores the text globally and typesets
%   \item \hyperlink{idea:1}{<text>}
% \ideahead typesets
%   \item\hypertarget{idea:1}{}\textbf{<the same text>}
%
% The abstract must come before the Main Results input, which is the normal
% order. If a heading has to read differently from the abstract sentence (for
% instance a long abstract item summarised in one line), override it with the
% optional argument: \ideahead[short heading]{6}.
% ======================================================================
\makeatletter
\newcommand{\ideaitem}[2]{%
  \expandafter\gdef\csname @idea@#1\endcsname{#2}%
  \item\hyperlink{idea:#1}{#2}}
\newcommand{\ideahead}[2][]{%
  \item\hypertarget{idea:#2}{}%
  \def\@ideatmp{#1}%
  \ifx\@ideatmp\@empty
    \@ifundefined{@idea@#2}{%
      \PackageWarning{preamble}{Idea #2 is used in Main Results but was never
        defined by \string\ideaitem\space in the abstract}%
      \textbf{[idea #2 undefined]}%
    }{\textbf{\begingroup\def\ideaterm##1##2{##2}\def\idearest##1{}\@nameuse{@idea@#2}\endgroup}}%
  \else
    \textbf{#1}%
  \fi}
% \ideaterm{<hypertarget>}{<words>}: inside an abstract idea the words become an
% invisible link (the document sets \hypersetup{hidelinks}) jumping to the place
% in Main Results that defines them; inside the Main Results heading itself,
% where the target sits a few lines below, \ideahead strips the link and prints
% the words plainly. Mark the destination with \hypertarget{<name>}{} on the line
% before the display, which adds no equation number and shows nothing.
\newcommand{\ideaterm}[2]{\hyperlink{#1}{#2}}
% \idearest{<text>}: the part of an abstract idea after its first sentence; printed in
% the abstract, omitted from the Main Results heading.
\newcommand{\idearest}[1]{#1}
\makeatother

% ======================================================================
% SHORT-PAPER SWITCH
% \shorttrue in a master that assembles only part of the companion paper
% (Barrat/SipArxiv.tex). \nsh{...} drops text that cites results living in the
% omitted sections, so the short build has no dangling references while the
% companion paper keeps the text verbatim.
% ======================================================================
\newif\ifshort \shortfalse
\newcommand{\nsh}[1]{\ifshort\else#1\fi}
% \shortalt{<short build>}{<companion paper>}: use when a sentence needs
% different wording in the two builds, not merely a deletion.
\newcommand{\shortalt}[2]{\ifshort#1\else#2\fi}
% \Jtr, \JtrS{\Sigma}: transverse Jacobian. J^\perp here; SipArxiv.tex redefines them as J_{\mathrm{tra}}, J_{\mathrm{tra},\Sigma}
\newcommand{\Jtr}{J^\perp}
\newcommand{\JtrS}[1]{J_{#1}^\perp}
.
\makeatletter
\AtBeginDocument{\@ifundefined{hypersetup}{}{\hypersetup{linktoc=all}}}
\makeatother
\usepackage{graphicx,float,fancybox}
\usepackage{epsfig,epstopdf,tikz}
\usepackage[sans]{dsfont}
\usetikzlibrary{positioning,arrows,arrows.meta,calc}
% Configure listings for Mathematica code
\lstdefinestyle{mathematica}{
 language=Mathematica,
 basicstyle=\small\ttfamily,
 backgroundcolor=\color{gray!10},
 frame=single,
 breaklines=true,
 commentstyle=\color{green!60!black},
 keywordstyle=\color{blue},
 stringstyle=\color{red},
 showstringspaces=false,
 tabsize=2
}
\usepackage{tcolorbox}
% Colored boxes for key results
\newtcolorbox{keyresult}[1][]{
 colback=blue!5!white,
 colframe=blue!75!black,
 title=#1,
 fonttitle=\bfseries
}
\newtcolorbox{examplebox}[1][]{
 colback=green!5!white,
 colframe=green!75!black,
 title=#1,
 fonttitle=\bfseries
}


\makeatletter
\@ifundefined{theorem}   {\newtheorem{theorem}   {Theorem}}   {}
\@ifundefined{definition}{\newtheorem{definition}{Definition}} {}
\@ifundefined{open}      {\newtheorem{open}      {Problem}}   {}
\@ifundefined{example}   {\newtheorem{example}   {Example}}   {}
\makeatother

% COUNTER DEFINITIONS
\newcounter{tabcounter}
\newcounter{defcounter}
\newcounter{opcounter}
\newcounter{thmcounter}
\newcounter{excounter}
\newcounter{ascounter}
\newcounter{cjcounter}
\newcounter{lemcounter}
\newcounter{remcounter}
\newcounter{corcounter}
\newcounter{propcounter}
\newcounter{figcounter}

% UNIFIED GLOSSARY STYLE
\newglossarystyle{unifiedstyle}{%
  \setglossarystyle{list}%
  \renewcommand*{\glossarypreamble}{\gdef\glslastsec{}}%
  \renewcommand{\glossentry}[2]{%
    \edef\glscursec{\glsentryuseriii{##1}}%
    \ifx\glscursec\glslastsec\else
      \item[]\textbf{\S\,\glscursec\enspace\glsentryuseriv{##1}}%
      \global\let\glslastsec\glscursec
    \fi
    \item[\glsentryuserii{##1} \glsentryuseri{##1}:] ``\glossentryname{##1}'' \glsentrydesc{##1}%
  }%
}

\makeatletter
% First-word extractor (used by \beT, \beD and \beXa for the default key)
% \@firstwordof hello world \relax  ->  hello
\def\@firstwordof#1 #2\relax{#1}
% First three letters, used for the heading mnemonic of a theorem given no
% explicit key.  Padding with \@empty keeps it safe for first words shorter
% than three letters: \@empty vanishes under \edef, whereas \relax would
% survive into the label and is illegal inside \csname (it broke every
% two-letter title, e.g. "ME ...", with "Missing \endcsname inserted").
\def\@takethree#1#2#3#4\@nil{#1#2#3}

% EXAMPLE COMMAND -- same signature and behaviour as \beT and \beD.
% Usage: \beXa*[long title][short title][key]
%   star  = suppresses the printed "[key]" tag after the title (as for \beT*
%   and \beD*); the label and glossary entry are unaffected.
%   key = optional label key override; the label becomes e:KEY, and without a
%   key it is e: plus the first three letters of the title's first word.  A
%   mnemonic label lets an example be cited as "Example~\ref{e:KEY}" instead
%   of by section number, which is what the citing prose needs.
\NewDocumentCommand{\beXa}{s O{} O{} o}{%
  \def\@temparg{#2}%
  \ifx\@temparg\@empty
    \begin{example}%
  \else
    \let\@savedSecTitle\@currentsectionname%
    \refstepcounter{excounter}%
    \IfValueTF{#4}{%
      \edef\@exlblkey{#4}%
    }{%
      \edef\@exfirstword{\@firstwordof#2 \relax}%
      \edef\@exlblkey{%
        \expandafter\@takethree\@exfirstword\@empty\@empty\@empty\@nil}%
    }%
    \def\@tempshort{#3}%
    \ifx\@tempshort\@empty
      \expandafter\xdef\csname @xshort@\@exlblkey\endcsname{\@exlblkey}%
    \else
      \expandafter\def\csname @xshort@\@exlblkey\endcsname{#3}%
    \fi
    \IfBooleanTF{#1}{\let\@extag\@empty}{%
      \edef\@extag{\noexpand\ \noexpand\textup{[\noexpand\texttt{e:\@exlblkey}]}}}%
    \begin{example}[#2\ifshowkeys\@extag\fi]%
    \label{e:\@exlblkey}%
    % record the long title so \nref{e:KEY} prints it, as for theorems
    \protected@write\@auxout{}{%
      \string\@thmtitle{e:\@exlblkey}{\unexpanded{#2}}}%
    \expanded{%
      \noexpand\newglossaryentry{e\theexcounter}{%
        name={\unexpanded{#2}%
          \IfBooleanF{#1}{ [\noexpand\texttt{e:\@exlblkey}]}},%
        sort={\theexample},%
        description={p.~\noexpand\pageref{e:\@exlblkey}},%
        user1={\theexample},%
        user2={Example},%
        user3={\thesection},%
        user4={\expandafter\unexpanded\expandafter{\@savedSecTitle}}%
      }%
    }%
    \protected@write\@auxout{}{\string\glsliveentry{e\theexcounter}}%
  \fi
}
\newcommand{\eeXa}{\end{example}}

% DEFINITION COMMAND -- same signature and behaviour as \beT below.
% Usage: \beD[long title][short title][key]
%   long title  = definition heading (plain text only)
%   short title = shown in \dref{word}; defaults to the key
%   key         = optional label key override; the label becomes d:KEY
% Without a key the label is d: plus the first three letters of the title's
% first word, so it is a mnemonic rather than the definition's number: a
% numeric label would silently change whenever a definition is inserted.
\NewDocumentCommand{\beD}{s O{} O{} o}{%
  \def\@temparg{#2}%
  \ifx\@temparg\@empty
    \begin{definition}%
  \else
    \let\@savedSecTitle\@currentsectionname%
    \refstepcounter{defcounter}%
    \IfValueTF{#4}{%
      \edef\@deflblkey{#4}%
    }{%
      \edef\@deffirstword{\@firstwordof#2 \relax}%
      \edef\@deflblkey{%
        \expandafter\@takethree\@deffirstword\@empty\@empty\@empty\@nil}%
    }%
    \def\@tempshort{#3}%
    \ifx\@tempshort\@empty
      \expandafter\xdef\csname @dshort@\@deflblkey\endcsname{\@deflblkey}%
    \else
      \expandafter\def\csname @dshort@\@deflblkey\endcsname{#3}%
    \fi
    \IfBooleanTF{#1}{\let\@deftag\@empty}{%
      \edef\@deftag{\noexpand\ \noexpand\textup{[\noexpand\texttt{d:\@deflblkey}]}}}%
    \begin{definition}[#2\ifshowkeys\@deftag\fi]%
    \label{d:\@deflblkey}%
    % record the long title so \nref{d:KEY} prints it, as for theorems
    \protected@write\@auxout{}{%
      \string\@thmtitle{d:\@deflblkey}{\unexpanded{#2}}}%
    \expanded{%
      \noexpand\newglossaryentry{d\thedefcounter}{%
        name={\unexpanded{#2}%
          \IfBooleanF{#1}{ [\noexpand\texttt{d:\@deflblkey}]}},%
        sort={\thedefinition},%
        description={p.~\noexpand\pageref{d:\@deflblkey}},%
        user1={\thedefinition},%
        user2={Definition},%
        user3={\thesection},%
        user4={\expandafter\unexpanded\expandafter{\@savedSecTitle}}%
      }%
    }%
    \protected@write\@auxout{}{\string\glsliveentry{d\thedefcounter}}%
  \fi
}
\newcommand{\eeD}{\end{definition}}

% \showkeysfalse in a driver hides the [t:KEY] mnemonics in result headings
% (labels and glossary entries are unchanged); used for submission builds.
\newif\ifshowkeys \showkeystrue
% \commonredtrue in a driver that inputs basicFiles/commonRed: tool.tex then
% cites its reduction results instead of restating them.
\newif\ifcommonred \commonredfalse
% TheorEM COMMAND
% Usage: \beT[long title][short title][key]
%   long title  = theorem heading (plain text only)
%   short title = shown in \tref{word}; may contain math; defaults to first word of long title
%   key         = optional label key override; label becomes t:KEY (avoids first-word collisions)
% Label is t:WORD where WORD = first word of long title (or KEY if given); use \tref{WORD} to reference.
\NewDocumentCommand{\beT}{s O{} O{} o}{%
  \def\@temparg{#2}%
  \ifx\@temparg\@empty
    \begin{theorem}%
  \else
    \let\@savedSecTitle\@currentsectionname%
    \refstepcounter{thmcounter}%
    % \@thmtag is the mnemonic shown in the heading; every theorem gets one.
    % With an explicit key the key itself is shown.  Without one the key
    % defaults to the title's first word, and showing that whole word would
    % merely repeat the title ("Dissipativity [Dissipativity]"), so only its
    % first three letters are shown: "Dissipativity [Dis]".  The label is
    % unaffected -- it still uses the full \@thmlblkey.
    \IfValueTF{#4}{%
      \edef\@thmlblkey{#4}%
    }{%
      \edef\@thmfirstword{\@firstwordof#2 \relax}%
      \edef\@thmlblkey{%
        \expandafter\@takethree\@thmfirstword\@empty\@empty\@empty\@nil}%
    }%
    \IfBooleanTF{#1}{\let\@thmtag\@empty}{%
      \edef\@thmtag{\noexpand\ \noexpand\textup{[\noexpand\texttt{t:\@thmlblkey}]}}}%
    \def\@tempshort{#3}%
    \ifx\@tempshort\@empty
      \expandafter\xdef\csname @tshort@\@thmlblkey\endcsname{\@thmlblkey}%
    \else
      \expandafter\def\csname @tshort@\@thmlblkey\endcsname{#3}%
    \fi
    \begin{theorem}[#2\ifshowkeys\@thmtag\fi]%
    \label{t:\@thmlblkey}%
    % record the long title for \nref; via the .aux so forward refs resolve
    \protected@write\@auxout{}{%
      \string\@thmtitle{t:\@thmlblkey}{\unexpanded{#2}}}%
    \expanded{%
      \noexpand\newglossaryentry{t\thethmcounter}{%
        name={\unexpanded{#2}%
          \IfBooleanF{#1}{ [\noexpand\texttt{t:\@thmlblkey}]}},%
        % sort/user1 must be the number the document actually prints, i.e. the
        % theorem environment's counter, not thmcounter.  thmcounter advances
        % only on \beT with a title, so a bare \beT (or \beT with an
        % empty title) makes the two drift and the glossary cites a wrong
        % number.  thmcounter is still fine as the entry key: it is unique.
        sort={\thetheorem},%
        description={p.~\noexpand\pageref{t:\@thmlblkey}},%
        user1={\thetheorem},%
        user2={Theorem},%
        user3={\thesection},%
        user4={\expandafter\unexpanded\expandafter{\@savedSecTitle}}%
      }%
    }%
    \protected@write\@auxout{}{\string\glsliveentry{t\thethmcounter}}%
  \fi
}
\newcommand{\eeT}{\end{theorem}}

% CROSS-PART RELATION PARAGRAPH
% Usage: \xpart{text}  -- place immediately BEFORE \beT, never inside it
%   (anything inside \beT...\eeT would be swallowed into the glossary entry).
% The paragraph is typeset only when Part I is actually present, i.e. in the
% combined book build; in the standalone builds (GranadaDSABNS/A.tex,
% AAHRah-Submission/AAHRah.tex) the label part:glv is undefined and the whole
% paragraph disappears rather than leaving ?? cross-references. Same
% graceful-degradation principle as \GASref in GASt.tex. Resolves on the
% second LaTeX pass, since r@part:glv comes from the .aux file.
% (Already inside the \makeatletter region opened above, so no \makeatletter
% pair here -- adding one would close that region early.)
\newcommand{\xpart}[1]{\@ifundefined{r@part:glv}{}{\paragraph{Relation with Parts~I--II.}#1}}

% \bookref{label}{precise text}{fallback text}: prints the precise cross-reference
% when the label resolves (combined book build), and the fallback text otherwise
% (standalone paper builds, where book-only appendix labels are absent). Same
% graceful-degradation role as \GASref in GASt.tex, but usable in any file since it
% lives in this preamble (loaded by every driver) -- unlike \GASref, which is defined in a
% content file and is therefore undefined in standalone builds that omit that file.
\newcommand{\bookref}[3]{\@ifundefined{r@#1}{#3}{#2}}

% \PVpart and \RBpart: the two parts that a partial build may omit.
% They print "Part~N" in any build that contains the part, and a naming phrase
% in a build that does not, so an excerpt shows no ?? .  Each already supplies
% the word "Part", so write "\PVpart develops ...", not "Part~\PVpart".
\newcommand{\PVpart}{\bookref{part:aahrah}%
  {Part~\ref{part:aahrah}}{the Perron--Volterra part of the full book}}
\newcommand{\RBpart}{\bookref{part:recbrowne}%
  {Part~\ref{part:recbrowne}}{the immuno-virology part of the full book}}

% OPEN PROBLEM COMMANDS
% ASSUMPTION COMMANDS (glossaried)
\newcommand{\beAs}[1][]{%
  \def\@temparg{#1}%
  \ifx\@temparg\@empty
    \begin{assumption}%
  \else
    \let\@savedSecTitle\@currentsectionname%
    \refstepcounter{ascounter}%
    \begin{assumption}[#1]%
    \label{as:\theascounter}%
    \expanded{%
      \noexpand\newglossaryentry{as\theascounter}{%
        name={\unexpanded{#1}},%
        sort={\theascounter},%
        description={p.~\noexpand\pageref{as:\theascounter}},%
        user1={\theascounter},%
        user2={Assumption},%
        user3={\thesection},%
        user4={\expandafter\unexpanded\expandafter{\@savedSecTitle}}%
      }%
    }%
    \protected@write\@auxout{}{\string\glsliveentry{as\theascounter}}%
  \fi
}
\newcommand{\eeAs}{\end{assumption}}
% CONJECTURE COMMANDS (glossaried)
\newcommand{\beCj}[1][]{%
  \def\@temparg{#1}%
  \ifx\@temparg\@empty
    \begin{conjecture}%
  \else
    \let\@savedSecTitle\@currentsectionname%
    \refstepcounter{cjcounter}%
    \begin{conjecture}[#1]%
    \label{cj:\thecjcounter}%
    \expanded{%
      \noexpand\newglossaryentry{cj\thecjcounter}{%
        name={\unexpanded{#1}},%
        sort={\thecjcounter},%
        description={p.~\noexpand\pageref{cj:\thecjcounter}},%
        user1={\thecjcounter},%
        user2={Conjecture},%
        user3={\thesection},%
        user4={\expandafter\unexpanded\expandafter{\@savedSecTitle}}%
      }%
    }%
    \protected@write\@auxout{}{\string\glsliveentry{cj\thecjcounter}}%
  \fi
}
\newcommand{\eeCj}{\end{conjecture}}
% LEMMA COMMANDS
% ---------------------------------------------------------------------
% UNIFORM RESULT WRAPPERS -- \beL, \beC, \beP, \beR share \beT's signature
%
%   \beX[long title][short title][key]  ...  \eeX
%
%   long title  = heading (math allowed)
%   short title = shown in the cross-reference; defaults to the key
%   key         = label key.  With a key the label is X:KEY and the heading
%                 carries the mnemonic [key], exactly as for \beT.  Without
%                 one the label falls back to the running counter, which is
%                 what every pre-existing call site in the book relies on.
%
% A bare \beX -- arguments commented out -- opens the plain environment and
% creates NO label and NO glossary entry.  That is the supported way to keep
% an individual result out of the catalogue.
% ---------------------------------------------------------------------
% #1 wrapper name  #2 environment  #3 label prefix  #4 counter  #5 type name
\newcommand{\@mkresult}[5]{%
  % A leading * suppresses the [key] mnemonic in both the heading and the
  % glossary name, while still making the label X:KEY and the entry:
  %     \beX*[long][short][key]
  \expandafter\NewDocumentCommand\csname be#1\endcsname{s O{} O{} o}{%
    \def\@temparg{##2}%
    \ifx\@temparg\@empty
      \begin{#2}%
    \else
      \let\@savedSecTitle\@currentsectionname%
      \refstepcounter{#4}%
      % With a key: mnemonic label X:KEY and the [key] tag in the heading.
      % Without one: fall back to the running counter, exactly as these
      % wrappers behaved before, so every legacy one-argument call site in
      % the book keeps its label.  The counter is also csname-safe, which a
      % title's first word is not -- titles beginning with math would
      % otherwise build \csname \Omega...\endcsname and fail.
      \IfValueTF{##4}{%
        \edef\@relbl{##4}%
        \IfBooleanTF{##1}{\let\@retag\@empty}{%
          \edef\@retag{\noexpand\ \noexpand\textup{[\noexpand\texttt{#3:##4}]}}}%
      }{%
        \edef\@relbl{\arabic{#4}}%
        \let\@retag\@empty
      }%
      \def\@tempshort{##3}%
      \ifx\@tempshort\@empty
        \expandafter\xdef\csname @sshort@#3\@relbl\endcsname{\@relbl}%
      \else
        \expandafter\def\csname @sshort@#3\@relbl\endcsname{##3}%
      \fi
      \begin{#2}[##2\ifshowkeys\@retag\fi]%
      \edef\@tmplbl{#3:\@relbl}\expandafter\label\expandafter{\@tmplbl}%
      \expanded{%
        \noexpand\newglossaryentry{#3\csname the#4\endcsname}{%
          name={\unexpanded{##2}%
            \IfBooleanF{##1}{\IfValueT{##4}{ [\noexpand\texttt{#3:##4}]}}},%
          sort={\csname the#2\endcsname},%
          description={p.~\noexpand\pageref{#3:\@relbl}},%
          user1={\csname the#2\endcsname},%
          user2={#5},%
          user3={\thesection},%
          user4={\expandafter\unexpanded\expandafter{\@savedSecTitle}}%
        }%
      }%
      \protected@write\@auxout{}{\string\glsliveentry{#3\csname the#4\endcsname}}%
      \protected@write\@auxout{}{\string\@thmtitle{\@tmplbl}{\unexpanded{##2}}}%
    \fi
  }%
  \expandafter\newcommand\csname ee#1\endcsname{\end{#2}}%
}
\@mkresult{L}{lemma}{l}{lemcounter}{Lemma}
% REMARK / COROLLARY / PROPOSITION -- same signature as \beL above
\@mkresult{R}{remark}{r}{remcounter}{Remark}
\@mkresult{C}{corollary}{c}{corcounter}{Corollary}
\@mkresult{P}{proposition}{p}{propcounter}{Proposition}
% OPEN-PROBLEM COMMAND -- generated with the same signature as \beT:
%   \beO[long title][short title][key]   ...   \eeO
% label o:KEY, or o:<counter> when no key is given; a bare \beO makes
% neither a label nor a glossary entry.
\@mkresult{O}{open}{o}{opcounter}{Problem}


% Section-level title tracker: updated only by \section, not \subsection.
% \@currentlabelname is overwritten by every \refstepcounter (including
% \subsection), so we save it separately here.
\newcommand{\@currentsectionname}{}
\AtBeginDocument{%
  \let\@epidcrn@orig@refstepcounter\refstepcounter
  \renewcommand{\refstepcounter}[1]{%
    \@epidcrn@orig@refstepcounter{#1}%
    \ifstrequal{#1}{section}{\global\let\@currentsectionname\@currentlabelname}{}%
  }%
}

% TABLE COMMANDS (glossaried) -- consistent naming with \beT, \beD, etc.
\newcommand{\beTa}[1][]{%
  \def\@temparg{#1}%
  \ifx\@temparg\@empty
    \begin{table}[htbp]%
  \else
    \let\@savedSecTitle\@currentsectionname%
    \refstepcounter{tabcounter}%
    \begin{table}[htbp]%
    \label{A:\thetabcounter}%
    \expanded{%
      \noexpand\newglossaryentry{A\thetabcounter}{%
        name={\unexpanded{#1}},%
        sort={\thetabcounter},%
        description={p.~\noexpand\pageref{A:\thetabcounter}},%
        user1={\thetabcounter},%
        user2={Table},%
        user3={\thesection},%
        user4={\expandafter\unexpanded\expandafter{\@savedSecTitle}}%
      }%
    }%
    \protected@write\@auxout{}{\string\glsliveentry{A\thetabcounter}}%
  \fi
}
\newcommand{\eeTa}{\end{table}}

% FIGURE COMMANDS (glossaried) -- consistent with \beT, \beR, \beTa, etc.
% Usage:
%   \beF[Glossary title]...\caption{Visible caption}\eeF  -> glossaried figure
%   \beF...\caption{Visible caption}\eeF                  -> plain [H] figure, no glossary
\newcommand{\beF}[1][]{%
  \def\@temparg{#1}%
  \ifx\@temparg\@empty
    \begin{figure}[H]%
  \else
    \begin{figure}[H]%
    \let\@savedSecTitle\@currentsectionname%
    \refstepcounter{figcounter}%
    \label{fig:\thefigcounter}%
    \expanded{%
      \noexpand\newglossaryentry{Fig\thefigcounter}{%
        name={\unexpanded{#1}},%
        sort={\thefigcounter},%
        description={p.~\noexpand\pageref{fig:\thefigcounter}},%
        user1={\thefigcounter},%
        user2={Figure},%
        user3={\thesection},%
        user4={\expandafter\unexpanded\expandafter{\@savedSecTitle}}%
      }%
    }%
    \protected@write\@auxout{}{\string\glsliveentry{Fig\thefigcounter}}%
  \fi
}
\newcommand{\eeF}{\end{figure}}
\makeatother
% LIVE-ENTRY REGISTER
% \glsaddall prints every entry the glossaries package currently knows, and
% that set includes entries restored from \jobname.glsdefs.  The package
% rewrites .glsdefs from all *currently defined* entries, so anything it
% restored is written straight back out: a deleted section kept reappearing in
% the catalogue for ever, and the .glsdefs had to be erased by hand.  The
% wrappers therefore record every entry they actually create in the .aux, and
% \printboth adds only those.  The .aux is rebuilt from scratch on each run, so
% this is self-healing: a removed result is gone after one further run.
\makeatletter
\newcommand{\@glslive}{}
\newcommand{\@glsaddlive}[1]{\ifglsentryexists{#1}{\glsadd{#1}}{}}
\newcommand{\glsliveentry}[1]{\g@addto@macro\@glslive{\@glsaddlive{#1}}}
\makeatother
% PRINT COMMAND
% No \clearpage here -- the catalogue continues on the current page.  This has
% been re-added by mistake several times; leave it out.
\makeatletter
\newcommand{\printboth}{%
  \@glslive   % add exactly the entries this document created last run
  \printglossary[style=unifiedstyle, title={0 \hspace{.22em} Definitions, Theorems, Examples and Open Problems}]%
}
\makeatother
% Named theorem reference: \tref{word} -> Thm~N(word-or-short)
% Shows stored short title if available, otherwise shows the key itself.
\newcommand{\tref}[1]{%
  Thm~\ref{t:#1}(\ifcsname @tshort@#1\endcsname\csname @tshort@#1\endcsname\else#1\fi)%
}

% Full-title theorem reference: \nref{t:KEY} prints the number and then the
% theorem's long title, so one writes "Theorem~\nref{t:blocks}" and gets
% "Theorem 3 ``Block factorization and siphons''".  Takes the complete label, like
% \ref, not the bare key like \tref.
%
% \beT writes each title to the .aux (see \@thmtitle below), so \nref resolves
% forward references from the second LaTeX pass, exactly as \ref does.  If the
% title is not yet known -- first pass, or a label that is not a \beT theorem
% -- it degrades to the number rather than producing ??.
% These two use @ in their names, and this point in the file is outside the
% \makeatletter region opened earlier (it closes at \makeatother above), so
% the region has to be reopened here.
\makeatletter
\newcommand{\@thmtitle}[2]{\global\@namedef{@tlong@#1}{#2}}
\newcommand{\nref}[1]{%
  \ref{#1}%
  \ifcsname @tlong@#1\endcsname
    \ ``\csname @tlong@#1\endcsname''%
  \fi
}
% \nm{t:KEY}: the stored long title of any \beT/\beL/\beP/\beR/\beC/\beD/\beE result,
% verbatim and without quotes, so item names follow the macro's first argument.
\newcommand{\nm}[1]{%
  \ifcsname @tlong@#1\endcsname
    \csname @tlong@#1\endcsname%
  \else
    \ref{#1}%
  \fi
}
\makeatother




%% ========================================
%% Other THEOREM ENVIRONMENTS
%% ========================================
\makeatletter
\@ifundefined{lemma}      {\newtheorem{lemma}      {Lemma}}      {}
\@ifundefined{proposition}{\newtheorem{proposition}{Proposition}} {}
\@ifundefined{corollary}  {\newtheorem{corollary}  {Corollary}}  {}
\@ifundefined{conjecture} {\newtheorem{conjecture} {Conjecture}}  {}
\@ifundefined{remark}     {\newtheorem{remark}     {Remark}}     {}
\@ifundefined{question}   {\newtheorem{question}   {Question}}   {}
\@ifundefined{assumption} {\newtheorem{assumption} {Assumption}} {}
\@ifundefined{hypothesis} {\newtheorem{hypothesis} {Hypothesis}} {}
\@ifundefined{fact}       {\newtheorem{fact}       {Fact}}       {}
\makeatother

% \citer: cite the relay-graph threshold theorem of AAHrel. If a document has
% \input the chapter/file containing t:relayRF (so the label resolves), point
% straight at it; otherwise fall back to a bare citation of AAHrel. Shared here
% (not per-paper) so it resolves consistently across every book chapter.
% \citeR: same idea, for the first theorem of AAHRah (t:Pstar).
\makeatletter
\providecommand{\citer}{%
  \@ifundefined{r@t:relayRF}%
    {\cite{AAHrel}}%
    {Theorem~\ref{t:relayRF} of \cite{AAHrel}}%
}
\providecommand{\citeR}{%
  \@ifundefined{r@t:Pstar}%
    {\cite{AAHRah}}%
    {Theorem~\ref{t:Pstar} of \cite{AAHRah}}%
}
\makeatother


% Abbreviations
\def\beQ{\begin{question}}\def\eeQ{\end{question}}
\def\beA{\begin{assumption}}\def\eeA{\end{assumption}}


%% ========================================

%% ========================================
%% COLORS & CUSTOM COMMANDS
%% ========================================
\definecolor{funccolor}{RGB}{25,25,112}
\definecolor{desccolor}{RGB}{64,64,64}
\newcommand{\func}[1]{\texttt{\textcolor{funccolor}{#1}}}
\def\eq{=}
\def\mm{-}

%% ========================================
%% MATHCAL SHORTHANDS
%% ========================================
% Project-wide convention: \mX for a capital X is always \mathcal X.
% Defined for the whole alphabet so the convention never fails to expand.
\newcommand{\mA}{\mathcal A}\newcommand{\mB}{\mathcal B}
\newcommand{\mC}{\mathcal C}\newcommand{\mD}{\mathcal D}
\newcommand{\mE}{\mathcal E}\newcommand{\mF}{\mathcal F}
\newcommand{\mG}{\mathcal G}\newcommand{\mH}{\mathcal H}
\newcommand{\mI}{\mathcal I}\newcommand{\mJ}{\mathcal J}
\newcommand{\mK}{\mathcal K}\newcommand{\mL}{\mathcal L}
\newcommand{\mM}{\mathcal M}\newcommand{\mN}{\mathcal N}
\newcommand{\mO}{\mathcal O}\newcommand{\mP}{\mathcal P}
\newcommand{\mQ}{\mathcal Q}\newcommand{\mR}{\mathcal R}
\newcommand{\mS}{\mathcal S}\newcommand{\mT}{\mathcal T}
\newcommand{\mU}{\mathcal U}\newcommand{\mV}{\mathcal V}
\newcommand{\mW}{\mathcal W}\newcommand{\mX}{\mathcal X}
\newcommand{\mY}{\mathcal Y}\newcommand{\mZ}{\mathcal Z}

%% ========================================
%% BASIC MATH ENVIRONMENTS
%% ========================================
\def\bep{\begin{pmatrix}}\def\eep{\end{pmatrix}}
\def\bev{\begin{vmatrix}}\def\eev{\end{vmatrix}}
\def\bea{\begin{eqnarray*}}\def\eea{\end{eqnarray*}}
\def\bc{\begin{cases}}\def\ec{\end{cases}}
\def\BEN{\begin{enumerate}}\def\EEN{\end{enumerate}}
\def\BI{\begin{itemize}}\def\EI{\end{itemize}}

\newcommand{\be}[1]{\begin{equation}\label{#1}}
\newcommand{\ee}{\end{equation}}
\newcommand{\beq}{\begin{eqnarray}}
\def\eeq{\end{eqnarray}}
\def\eqr{\eqref}\def\fr{\frac}\def\lbl{\label}
\def\Lra{\Longrightarrow}\def\non{\nonumber}\def\cNGM{\cite{Diek,Van,Van08}}
\def\Eq{\Leftrightarrow}\def\b{\beta}\def\uph{\red{UP TO HERE}}
%% ========================================
%% PROBABILITY & STATISTICS
%% ========================================
\renewcommand{\P}{\ensuremath{\mathbb{P}}}  % Note: overrides paragraph symbol
\newcommand{\E}{\ensuremath{\mathbb{E}}}
\newcommand{\ind}[1]{\ensuremath{\mathbbm{1}_{\left\{#1\right\}}}}
\DeclareMathOperator{\Cov}{Cov}
\DeclareMathOperator{\TV}{TV}
\newcommand{\rank}{\operatorname{rank}}

%% ========================================
%% CRN OPERATORS (organizations/ERCs/buffering/siphons framework)
%% ========================================
\def\clos{\mathrm{clos}}
\def\Supply{\mathrm{Supply}}
\def\req{\mathrm{req}}
\def\React{\mathrm{React}}
\def\Prod{\mathrm{Prod}}
\def\Gen{\mathrm{Gen}}
\def\In{\mathrm{In}}
\def\Out{\mathrm{Out}}
\def\Act{\mathrm{Act}}
\def\Min{\mathrm{Min}}

%% ========================================
%% LINE BREAKING
%% ========================================
% Last-resort stretch, used only on paragraphs TeX cannot otherwise break:
% a long \texttt{} token, a comma-separated inline math list (commas are not
% breakpoints in math), or a run of \cite/\ref boxes.  Without it those set as
% overfull \hbox and print into the margin; with it they set slightly loose,
% which is the lesser evil.  Does not affect paragraphs that already break.
\emergencystretch=3em

%% ========================================
%% NUMBER SETS
%% ========================================
\newcommand{\R}{\mathbb{R}}  % Real numbers
\newcommand{\C}{\mathbb{C}}  % Complex numbers
\newcommand{\N}{\mathbb{N}}
\newcommand{\Z}{\mathbb{Z}}
% \providecommand, not \newcommand: some papers define these in their own
% preamble after % Shared preamble for the GLVbook papers.  This is the single canonical copy;
% a driver one level below GLVbook reaches it as
%
%     \input{../basicFiles/preamble}
%
% What stays in the driver, because it differs from paper to paper:
% \documentclass, font choice, \title, \author, \date, and any macro genuinely
% local to one paper.  Everything else belongs here.  The file is
% self-contained: it does not read any other file.

\usepackage[utf8]{inputenc}
\usepackage[novbox]{pdfsync}   % required: SyncTeX alone does not work here

% Page geometry, set once here for every paper that loads this file.  A driver
% that needs different margins says \geometry{...}; a second
% \usepackage[...]{geometry} raises an option clash and must not be used.
\makeatletter
\@ifundefined{AAdefNnogeometry}{\usepackage[a4paper,margin=2cm]{geometry}}{}
\makeatother

% mathematics
\usepackage{amsmath,amssymb,amsthm,mathtools,bm,mathrsfs,bbm}
\usepackage{centernot}

% floats, tables, lists
\usepackage{graphicx,float}
\makeatletter
\@ifpackageloaded{subfig}{}{\usepackage{subcaption}}
\makeatother
\usepackage{booktabs,longtable,threeparttable,array,multirow}
\usepackage{enumitem}

% colour, code, algorithms
\usepackage{xcolor,listings,tcolorbox}
\makeatletter
\@ifundefined{c@algorithm}{\usepackage{algorithm}}{}
\makeatother
\usepackage{algorithmicx,algpseudocode}

% diagrams
\usepackage{tikz,tikz-cd,pgfplots}
\usetikzlibrary{graphs,decorations.pathmorphing,decorations.markings,
arrows.meta,automata,positioning}
\tikzcdset{every label/.append style = {font = \small}}
\pgfplotsset{compat=1.18}

% chemistry
\usepackage[version=4]{mhchem}

\usepackage{microtype}
\usepackage{lineno}
\usepackage{nomencl}
\makenomenclature

% hyperref first, then bookmark, then glossaries: glossaries must follow
% hyperref or the entry links are built against the wrong anchors.
\makeatletter
\@ifundefined{hyperrefpackage}{%
  \@ifundefined{AAdefNnohyperref}{\usepackage{hyperref}}{}%
}{}
\@ifundefined{hyperrefpackage}{\usepackage{bookmark}}{}
\makeatother
\usepackage[toc,sort=use]{glossaries}
\makeglossaries

% Book/standalone switch.  Material that the book already covers in its
% prelude, but that a standalone paper still has to state, is wrapped as
%     \ifbookbuild\else ... \fi
% Drivers default to standalone; bookM.tex says \bookbuildtrue.
\newif\ifbookbuild
\bookbuildfalse

% Part switches for the gradostat paper.  Both default to true; a driver that
% wants the short submission says \madecfalse and/or \slowfastfalse, and every
% cross-reference into the dropped part disappears with it.
%   \ifmadec     Part III, the Lagasquie--Madec model and its break-even geometry
%   \ifslowfast  Part IV, everything carrying the small parameter \eps
\newif\ifmadec     \madectrue
\newif\ifslowfast  \slowfasttrue

% Part switch for the Barrat/simplicial-contagion project.  Defaults to
% true (the hypernetwork-synchrony part is included); the journal-paper
% driver says \hnetpartfalse to drop it and every cross-reference into it.
\newif\ifhnetpart  \hnetparttrue

%% ========================================
%% PACKAGES
%% ========================================
\RequirePackage{etoolbox}
\usepackage{makecell}
\usepackage[novbox]{pdfsync}% obsolete; SyncTeX is built into modern engines
\usepackage{amssymb, amsfonts,mathtools,bbm,eurosym,graphicx,latexsym, amsmath,amsthm,times,nccmath,url}
\usepackage{listings,verbatim,cases,tabularx}
% hyperref is deliberately NOT in the line above.  Journal classes often
% load it themselves, with a long option list, from a begindocument hook
% (mdpi.cls does: \DeclareHookRule{begindocument}{hyperref}{before}{mdpi}).
% Loading it here first, without options, makes that later load raise
% "Option clash for package hyperref".  So load it only if nobody else
% has, or will:
%   - \@ifpackageloaded  catches a class that loaded it already;
%   - \hyperrefpackage    is mdpi.cls's own deferred loader;
%   - \AAdefNnohyperref   lets any other driver opt out by defining it
%                         before % Shared preamble for the GLVbook papers.  This is the single canonical copy;
% a driver one level below GLVbook reaches it as
%
%     \input{../basicFiles/preamble}
%
% What stays in the driver, because it differs from paper to paper:
% \documentclass, font choice, \title, \author, \date, and any macro genuinely
% local to one paper.  Everything else belongs here.  The file is
% self-contained: it does not read any other file.

\usepackage[utf8]{inputenc}
\usepackage[novbox]{pdfsync}   % required: SyncTeX alone does not work here

% Page geometry, set once here for every paper that loads this file.  A driver
% that needs different margins says \geometry{...}; a second
% \usepackage[...]{geometry} raises an option clash and must not be used.
\makeatletter
\@ifundefined{AAdefNnogeometry}{\usepackage[a4paper,margin=2cm]{geometry}}{}
\makeatother

% mathematics
\usepackage{amsmath,amssymb,amsthm,mathtools,bm,mathrsfs,bbm}
\usepackage{centernot}

% floats, tables, lists
\usepackage{graphicx,float}
\makeatletter
\@ifpackageloaded{subfig}{}{\usepackage{subcaption}}
\makeatother
\usepackage{booktabs,longtable,threeparttable,array,multirow}
\usepackage{enumitem}

% colour, code, algorithms
\usepackage{xcolor,listings,tcolorbox}
\makeatletter
\@ifundefined{c@algorithm}{\usepackage{algorithm}}{}
\makeatother
\usepackage{algorithmicx,algpseudocode}

% diagrams
\usepackage{tikz,tikz-cd,pgfplots}
\usetikzlibrary{graphs,decorations.pathmorphing,decorations.markings,
arrows.meta,automata,positioning}
\tikzcdset{every label/.append style = {font = \small}}
\pgfplotsset{compat=1.18}

% chemistry
\usepackage[version=4]{mhchem}

\usepackage{microtype}
\usepackage{lineno}
\usepackage{nomencl}
\makenomenclature

% hyperref first, then bookmark, then glossaries: glossaries must follow
% hyperref or the entry links are built against the wrong anchors.
\makeatletter
\@ifundefined{hyperrefpackage}{%
  \@ifundefined{AAdefNnohyperref}{\usepackage{hyperref}}{}%
}{}
\@ifundefined{hyperrefpackage}{\usepackage{bookmark}}{}
\makeatother
\usepackage[toc,sort=use]{glossaries}
\makeglossaries

% Book/standalone switch.  Material that the book already covers in its
% prelude, but that a standalone paper still has to state, is wrapped as
%     \ifbookbuild\else ... \fi
% Drivers default to standalone; bookM.tex says \bookbuildtrue.
\newif\ifbookbuild
\bookbuildfalse

% Part switches for the gradostat paper.  Both default to true; a driver that
% wants the short submission says \madecfalse and/or \slowfastfalse, and every
% cross-reference into the dropped part disappears with it.
%   \ifmadec     Part III, the Lagasquie--Madec model and its break-even geometry
%   \ifslowfast  Part IV, everything carrying the small parameter \eps
\newif\ifmadec     \madectrue
\newif\ifslowfast  \slowfasttrue

% Part switch for the Barrat/simplicial-contagion project.  Defaults to
% true (the hypernetwork-synchrony part is included); the journal-paper
% driver says \hnetpartfalse to drop it and every cross-reference into it.
\newif\ifhnetpart  \hnetparttrue

%% ========================================
%% PACKAGES
%% ========================================
\RequirePackage{etoolbox}
\usepackage{makecell}
\usepackage[novbox]{pdfsync}% obsolete; SyncTeX is built into modern engines
\usepackage{amssymb, amsfonts,mathtools,bbm,eurosym,graphicx,latexsym, amsmath,amsthm,times,nccmath,url}
\usepackage{listings,verbatim,cases,tabularx}
% hyperref is deliberately NOT in the line above.  Journal classes often
% load it themselves, with a long option list, from a begindocument hook
% (mdpi.cls does: \DeclareHookRule{begindocument}{hyperref}{before}{mdpi}).
% Loading it here first, without options, makes that later load raise
% "Option clash for package hyperref".  So load it only if nobody else
% has, or will:
%   - \@ifpackageloaded  catches a class that loaded it already;
%   - \hyperrefpackage    is mdpi.cls's own deferred loader;
%   - \AAdefNnohyperref   lets any other driver opt out by defining it
%                         before \input{preamble}.
\makeatletter
\@ifpackageloaded{hyperref}{}{%
  \@ifundefined{hyperrefpackage}{%
    \@ifundefined{AAdefNnohyperref}{\usepackage{hyperref}}{}%
  }{}%
}
\makeatother
% linktoc=all makes a table-of-contents page number a link as well as the
% entry text; hyperref's default, linktoc=section, links only the text.
% Deferred to begin-document and guarded: by the block above hyperref may
% be loaded later (mdpi's own hook) or not at all (\AAdefNnohyperref), and
% \hypersetup must exist before it is called.  A driver that wants the
% default back can say \hypersetup{linktoc=section} after \input{preamble}.
\makeatletter
\AtBeginDocument{\@ifundefined{hypersetup}{}{\hypersetup{linktoc=all}}}
\makeatother
\usepackage{graphicx,float,fancybox}
\usepackage{epsfig,epstopdf,tikz}
\usepackage[sans]{dsfont}
\usetikzlibrary{positioning,arrows,arrows.meta,calc}
% Configure listings for Mathematica code
\lstdefinestyle{mathematica}{
 language=Mathematica,
 basicstyle=\small\ttfamily,
 backgroundcolor=\color{gray!10},
 frame=single,
 breaklines=true,
 commentstyle=\color{green!60!black},
 keywordstyle=\color{blue},
 stringstyle=\color{red},
 showstringspaces=false,
 tabsize=2
}
\usepackage{tcolorbox}
% Colored boxes for key results
\newtcolorbox{keyresult}[1][]{
 colback=blue!5!white,
 colframe=blue!75!black,
 title=#1,
 fonttitle=\bfseries
}
\newtcolorbox{examplebox}[1][]{
 colback=green!5!white,
 colframe=green!75!black,
 title=#1,
 fonttitle=\bfseries
}


\makeatletter
\@ifundefined{theorem}   {\newtheorem{theorem}   {Theorem}}   {}
\@ifundefined{definition}{\newtheorem{definition}{Definition}} {}
\@ifundefined{open}      {\newtheorem{open}      {Problem}}   {}
\@ifundefined{example}   {\newtheorem{example}   {Example}}   {}
\makeatother

% COUNTER DEFINITIONS
\newcounter{tabcounter}
\newcounter{defcounter}
\newcounter{opcounter}
\newcounter{thmcounter}
\newcounter{excounter}
\newcounter{ascounter}
\newcounter{cjcounter}
\newcounter{lemcounter}
\newcounter{remcounter}
\newcounter{corcounter}
\newcounter{propcounter}
\newcounter{figcounter}

% UNIFIED GLOSSARY STYLE
\newglossarystyle{unifiedstyle}{%
  \setglossarystyle{list}%
  \renewcommand*{\glossarypreamble}{\gdef\glslastsec{}}%
  \renewcommand{\glossentry}[2]{%
    \edef\glscursec{\glsentryuseriii{##1}}%
    \ifx\glscursec\glslastsec\else
      \item[]\textbf{\S\,\glscursec\enspace\glsentryuseriv{##1}}%
      \global\let\glslastsec\glscursec
    \fi
    \item[\glsentryuserii{##1} \glsentryuseri{##1}:] ``\glossentryname{##1}'' \glsentrydesc{##1}%
  }%
}

\makeatletter
% First-word extractor (used by \beT, \beD and \beXa for the default key)
% \@firstwordof hello world \relax  ->  hello
\def\@firstwordof#1 #2\relax{#1}
% First three letters, used for the heading mnemonic of a theorem given no
% explicit key.  Padding with \@empty keeps it safe for first words shorter
% than three letters: \@empty vanishes under \edef, whereas \relax would
% survive into the label and is illegal inside \csname (it broke every
% two-letter title, e.g. "ME ...", with "Missing \endcsname inserted").
\def\@takethree#1#2#3#4\@nil{#1#2#3}

% EXAMPLE COMMAND -- same signature and behaviour as \beT and \beD.
% Usage: \beXa*[long title][short title][key]
%   star  = suppresses the printed "[key]" tag after the title (as for \beT*
%   and \beD*); the label and glossary entry are unaffected.
%   key = optional label key override; the label becomes e:KEY, and without a
%   key it is e: plus the first three letters of the title's first word.  A
%   mnemonic label lets an example be cited as "Example~\ref{e:KEY}" instead
%   of by section number, which is what the citing prose needs.
\NewDocumentCommand{\beXa}{s O{} O{} o}{%
  \def\@temparg{#2}%
  \ifx\@temparg\@empty
    \begin{example}%
  \else
    \let\@savedSecTitle\@currentsectionname%
    \refstepcounter{excounter}%
    \IfValueTF{#4}{%
      \edef\@exlblkey{#4}%
    }{%
      \edef\@exfirstword{\@firstwordof#2 \relax}%
      \edef\@exlblkey{%
        \expandafter\@takethree\@exfirstword\@empty\@empty\@empty\@nil}%
    }%
    \def\@tempshort{#3}%
    \ifx\@tempshort\@empty
      \expandafter\xdef\csname @xshort@\@exlblkey\endcsname{\@exlblkey}%
    \else
      \expandafter\def\csname @xshort@\@exlblkey\endcsname{#3}%
    \fi
    \IfBooleanTF{#1}{\let\@extag\@empty}{%
      \edef\@extag{\noexpand\ \noexpand\textup{[\noexpand\texttt{e:\@exlblkey}]}}}%
    \begin{example}[#2\ifshowkeys\@extag\fi]%
    \label{e:\@exlblkey}%
    % record the long title so \nref{e:KEY} prints it, as for theorems
    \protected@write\@auxout{}{%
      \string\@thmtitle{e:\@exlblkey}{\unexpanded{#2}}}%
    \expanded{%
      \noexpand\newglossaryentry{e\theexcounter}{%
        name={\unexpanded{#2}%
          \IfBooleanF{#1}{ [\noexpand\texttt{e:\@exlblkey}]}},%
        sort={\theexample},%
        description={p.~\noexpand\pageref{e:\@exlblkey}},%
        user1={\theexample},%
        user2={Example},%
        user3={\thesection},%
        user4={\expandafter\unexpanded\expandafter{\@savedSecTitle}}%
      }%
    }%
    \protected@write\@auxout{}{\string\glsliveentry{e\theexcounter}}%
  \fi
}
\newcommand{\eeXa}{\end{example}}

% DEFINITION COMMAND -- same signature and behaviour as \beT below.
% Usage: \beD[long title][short title][key]
%   long title  = definition heading (plain text only)
%   short title = shown in \dref{word}; defaults to the key
%   key         = optional label key override; the label becomes d:KEY
% Without a key the label is d: plus the first three letters of the title's
% first word, so it is a mnemonic rather than the definition's number: a
% numeric label would silently change whenever a definition is inserted.
\NewDocumentCommand{\beD}{s O{} O{} o}{%
  \def\@temparg{#2}%
  \ifx\@temparg\@empty
    \begin{definition}%
  \else
    \let\@savedSecTitle\@currentsectionname%
    \refstepcounter{defcounter}%
    \IfValueTF{#4}{%
      \edef\@deflblkey{#4}%
    }{%
      \edef\@deffirstword{\@firstwordof#2 \relax}%
      \edef\@deflblkey{%
        \expandafter\@takethree\@deffirstword\@empty\@empty\@empty\@nil}%
    }%
    \def\@tempshort{#3}%
    \ifx\@tempshort\@empty
      \expandafter\xdef\csname @dshort@\@deflblkey\endcsname{\@deflblkey}%
    \else
      \expandafter\def\csname @dshort@\@deflblkey\endcsname{#3}%
    \fi
    \IfBooleanTF{#1}{\let\@deftag\@empty}{%
      \edef\@deftag{\noexpand\ \noexpand\textup{[\noexpand\texttt{d:\@deflblkey}]}}}%
    \begin{definition}[#2\ifshowkeys\@deftag\fi]%
    \label{d:\@deflblkey}%
    % record the long title so \nref{d:KEY} prints it, as for theorems
    \protected@write\@auxout{}{%
      \string\@thmtitle{d:\@deflblkey}{\unexpanded{#2}}}%
    \expanded{%
      \noexpand\newglossaryentry{d\thedefcounter}{%
        name={\unexpanded{#2}%
          \IfBooleanF{#1}{ [\noexpand\texttt{d:\@deflblkey}]}},%
        sort={\thedefinition},%
        description={p.~\noexpand\pageref{d:\@deflblkey}},%
        user1={\thedefinition},%
        user2={Definition},%
        user3={\thesection},%
        user4={\expandafter\unexpanded\expandafter{\@savedSecTitle}}%
      }%
    }%
    \protected@write\@auxout{}{\string\glsliveentry{d\thedefcounter}}%
  \fi
}
\newcommand{\eeD}{\end{definition}}

% \showkeysfalse in a driver hides the [t:KEY] mnemonics in result headings
% (labels and glossary entries are unchanged); used for submission builds.
\newif\ifshowkeys \showkeystrue
% \commonredtrue in a driver that inputs basicFiles/commonRed: tool.tex then
% cites its reduction results instead of restating them.
\newif\ifcommonred \commonredfalse
% TheorEM COMMAND
% Usage: \beT[long title][short title][key]
%   long title  = theorem heading (plain text only)
%   short title = shown in \tref{word}; may contain math; defaults to first word of long title
%   key         = optional label key override; label becomes t:KEY (avoids first-word collisions)
% Label is t:WORD where WORD = first word of long title (or KEY if given); use \tref{WORD} to reference.
\NewDocumentCommand{\beT}{s O{} O{} o}{%
  \def\@temparg{#2}%
  \ifx\@temparg\@empty
    \begin{theorem}%
  \else
    \let\@savedSecTitle\@currentsectionname%
    \refstepcounter{thmcounter}%
    % \@thmtag is the mnemonic shown in the heading; every theorem gets one.
    % With an explicit key the key itself is shown.  Without one the key
    % defaults to the title's first word, and showing that whole word would
    % merely repeat the title ("Dissipativity [Dissipativity]"), so only its
    % first three letters are shown: "Dissipativity [Dis]".  The label is
    % unaffected -- it still uses the full \@thmlblkey.
    \IfValueTF{#4}{%
      \edef\@thmlblkey{#4}%
    }{%
      \edef\@thmfirstword{\@firstwordof#2 \relax}%
      \edef\@thmlblkey{%
        \expandafter\@takethree\@thmfirstword\@empty\@empty\@empty\@nil}%
    }%
    \IfBooleanTF{#1}{\let\@thmtag\@empty}{%
      \edef\@thmtag{\noexpand\ \noexpand\textup{[\noexpand\texttt{t:\@thmlblkey}]}}}%
    \def\@tempshort{#3}%
    \ifx\@tempshort\@empty
      \expandafter\xdef\csname @tshort@\@thmlblkey\endcsname{\@thmlblkey}%
    \else
      \expandafter\def\csname @tshort@\@thmlblkey\endcsname{#3}%
    \fi
    \begin{theorem}[#2\ifshowkeys\@thmtag\fi]%
    \label{t:\@thmlblkey}%
    % record the long title for \nref; via the .aux so forward refs resolve
    \protected@write\@auxout{}{%
      \string\@thmtitle{t:\@thmlblkey}{\unexpanded{#2}}}%
    \expanded{%
      \noexpand\newglossaryentry{t\thethmcounter}{%
        name={\unexpanded{#2}%
          \IfBooleanF{#1}{ [\noexpand\texttt{t:\@thmlblkey}]}},%
        % sort/user1 must be the number the document actually prints, i.e. the
        % theorem environment's counter, not thmcounter.  thmcounter advances
        % only on \beT with a title, so a bare \beT (or \beT with an
        % empty title) makes the two drift and the glossary cites a wrong
        % number.  thmcounter is still fine as the entry key: it is unique.
        sort={\thetheorem},%
        description={p.~\noexpand\pageref{t:\@thmlblkey}},%
        user1={\thetheorem},%
        user2={Theorem},%
        user3={\thesection},%
        user4={\expandafter\unexpanded\expandafter{\@savedSecTitle}}%
      }%
    }%
    \protected@write\@auxout{}{\string\glsliveentry{t\thethmcounter}}%
  \fi
}
\newcommand{\eeT}{\end{theorem}}

% CROSS-PART RELATION PARAGRAPH
% Usage: \xpart{text}  -- place immediately BEFORE \beT, never inside it
%   (anything inside \beT...\eeT would be swallowed into the glossary entry).
% The paragraph is typeset only when Part I is actually present, i.e. in the
% combined book build; in the standalone builds (GranadaDSABNS/A.tex,
% AAHRah-Submission/AAHRah.tex) the label part:glv is undefined and the whole
% paragraph disappears rather than leaving ?? cross-references. Same
% graceful-degradation principle as \GASref in GASt.tex. Resolves on the
% second LaTeX pass, since r@part:glv comes from the .aux file.
% (Already inside the \makeatletter region opened above, so no \makeatletter
% pair here -- adding one would close that region early.)
\newcommand{\xpart}[1]{\@ifundefined{r@part:glv}{}{\paragraph{Relation with Parts~I--II.}#1}}

% \bookref{label}{precise text}{fallback text}: prints the precise cross-reference
% when the label resolves (combined book build), and the fallback text otherwise
% (standalone paper builds, where book-only appendix labels are absent). Same
% graceful-degradation role as \GASref in GASt.tex, but usable in any file since it
% lives in this preamble (loaded by every driver) -- unlike \GASref, which is defined in a
% content file and is therefore undefined in standalone builds that omit that file.
\newcommand{\bookref}[3]{\@ifundefined{r@#1}{#3}{#2}}

% \PVpart and \RBpart: the two parts that a partial build may omit.
% They print "Part~N" in any build that contains the part, and a naming phrase
% in a build that does not, so an excerpt shows no ?? .  Each already supplies
% the word "Part", so write "\PVpart develops ...", not "Part~\PVpart".
\newcommand{\PVpart}{\bookref{part:aahrah}%
  {Part~\ref{part:aahrah}}{the Perron--Volterra part of the full book}}
\newcommand{\RBpart}{\bookref{part:recbrowne}%
  {Part~\ref{part:recbrowne}}{the immuno-virology part of the full book}}

% OPEN PROBLEM COMMANDS
% ASSUMPTION COMMANDS (glossaried)
\newcommand{\beAs}[1][]{%
  \def\@temparg{#1}%
  \ifx\@temparg\@empty
    \begin{assumption}%
  \else
    \let\@savedSecTitle\@currentsectionname%
    \refstepcounter{ascounter}%
    \begin{assumption}[#1]%
    \label{as:\theascounter}%
    \expanded{%
      \noexpand\newglossaryentry{as\theascounter}{%
        name={\unexpanded{#1}},%
        sort={\theascounter},%
        description={p.~\noexpand\pageref{as:\theascounter}},%
        user1={\theascounter},%
        user2={Assumption},%
        user3={\thesection},%
        user4={\expandafter\unexpanded\expandafter{\@savedSecTitle}}%
      }%
    }%
    \protected@write\@auxout{}{\string\glsliveentry{as\theascounter}}%
  \fi
}
\newcommand{\eeAs}{\end{assumption}}
% CONJECTURE COMMANDS (glossaried)
\newcommand{\beCj}[1][]{%
  \def\@temparg{#1}%
  \ifx\@temparg\@empty
    \begin{conjecture}%
  \else
    \let\@savedSecTitle\@currentsectionname%
    \refstepcounter{cjcounter}%
    \begin{conjecture}[#1]%
    \label{cj:\thecjcounter}%
    \expanded{%
      \noexpand\newglossaryentry{cj\thecjcounter}{%
        name={\unexpanded{#1}},%
        sort={\thecjcounter},%
        description={p.~\noexpand\pageref{cj:\thecjcounter}},%
        user1={\thecjcounter},%
        user2={Conjecture},%
        user3={\thesection},%
        user4={\expandafter\unexpanded\expandafter{\@savedSecTitle}}%
      }%
    }%
    \protected@write\@auxout{}{\string\glsliveentry{cj\thecjcounter}}%
  \fi
}
\newcommand{\eeCj}{\end{conjecture}}
% LEMMA COMMANDS
% ---------------------------------------------------------------------
% UNIFORM RESULT WRAPPERS -- \beL, \beC, \beP, \beR share \beT's signature
%
%   \beX[long title][short title][key]  ...  \eeX
%
%   long title  = heading (math allowed)
%   short title = shown in the cross-reference; defaults to the key
%   key         = label key.  With a key the label is X:KEY and the heading
%                 carries the mnemonic [key], exactly as for \beT.  Without
%                 one the label falls back to the running counter, which is
%                 what every pre-existing call site in the book relies on.
%
% A bare \beX -- arguments commented out -- opens the plain environment and
% creates NO label and NO glossary entry.  That is the supported way to keep
% an individual result out of the catalogue.
% ---------------------------------------------------------------------
% #1 wrapper name  #2 environment  #3 label prefix  #4 counter  #5 type name
\newcommand{\@mkresult}[5]{%
  % A leading * suppresses the [key] mnemonic in both the heading and the
  % glossary name, while still making the label X:KEY and the entry:
  %     \beX*[long][short][key]
  \expandafter\NewDocumentCommand\csname be#1\endcsname{s O{} O{} o}{%
    \def\@temparg{##2}%
    \ifx\@temparg\@empty
      \begin{#2}%
    \else
      \let\@savedSecTitle\@currentsectionname%
      \refstepcounter{#4}%
      % With a key: mnemonic label X:KEY and the [key] tag in the heading.
      % Without one: fall back to the running counter, exactly as these
      % wrappers behaved before, so every legacy one-argument call site in
      % the book keeps its label.  The counter is also csname-safe, which a
      % title's first word is not -- titles beginning with math would
      % otherwise build \csname \Omega...\endcsname and fail.
      \IfValueTF{##4}{%
        \edef\@relbl{##4}%
        \IfBooleanTF{##1}{\let\@retag\@empty}{%
          \edef\@retag{\noexpand\ \noexpand\textup{[\noexpand\texttt{#3:##4}]}}}%
      }{%
        \edef\@relbl{\arabic{#4}}%
        \let\@retag\@empty
      }%
      \def\@tempshort{##3}%
      \ifx\@tempshort\@empty
        \expandafter\xdef\csname @sshort@#3\@relbl\endcsname{\@relbl}%
      \else
        \expandafter\def\csname @sshort@#3\@relbl\endcsname{##3}%
      \fi
      \begin{#2}[##2\ifshowkeys\@retag\fi]%
      \edef\@tmplbl{#3:\@relbl}\expandafter\label\expandafter{\@tmplbl}%
      \expanded{%
        \noexpand\newglossaryentry{#3\csname the#4\endcsname}{%
          name={\unexpanded{##2}%
            \IfBooleanF{##1}{\IfValueT{##4}{ [\noexpand\texttt{#3:##4}]}}},%
          sort={\csname the#2\endcsname},%
          description={p.~\noexpand\pageref{#3:\@relbl}},%
          user1={\csname the#2\endcsname},%
          user2={#5},%
          user3={\thesection},%
          user4={\expandafter\unexpanded\expandafter{\@savedSecTitle}}%
        }%
      }%
      \protected@write\@auxout{}{\string\glsliveentry{#3\csname the#4\endcsname}}%
      \protected@write\@auxout{}{\string\@thmtitle{\@tmplbl}{\unexpanded{##2}}}%
    \fi
  }%
  \expandafter\newcommand\csname ee#1\endcsname{\end{#2}}%
}
\@mkresult{L}{lemma}{l}{lemcounter}{Lemma}
% REMARK / COROLLARY / PROPOSITION -- same signature as \beL above
\@mkresult{R}{remark}{r}{remcounter}{Remark}
\@mkresult{C}{corollary}{c}{corcounter}{Corollary}
\@mkresult{P}{proposition}{p}{propcounter}{Proposition}
% OPEN-PROBLEM COMMAND -- generated with the same signature as \beT:
%   \beO[long title][short title][key]   ...   \eeO
% label o:KEY, or o:<counter> when no key is given; a bare \beO makes
% neither a label nor a glossary entry.
\@mkresult{O}{open}{o}{opcounter}{Problem}


% Section-level title tracker: updated only by \section, not \subsection.
% \@currentlabelname is overwritten by every \refstepcounter (including
% \subsection), so we save it separately here.
\newcommand{\@currentsectionname}{}
\AtBeginDocument{%
  \let\@epidcrn@orig@refstepcounter\refstepcounter
  \renewcommand{\refstepcounter}[1]{%
    \@epidcrn@orig@refstepcounter{#1}%
    \ifstrequal{#1}{section}{\global\let\@currentsectionname\@currentlabelname}{}%
  }%
}

% TABLE COMMANDS (glossaried) -- consistent naming with \beT, \beD, etc.
\newcommand{\beTa}[1][]{%
  \def\@temparg{#1}%
  \ifx\@temparg\@empty
    \begin{table}[htbp]%
  \else
    \let\@savedSecTitle\@currentsectionname%
    \refstepcounter{tabcounter}%
    \begin{table}[htbp]%
    \label{A:\thetabcounter}%
    \expanded{%
      \noexpand\newglossaryentry{A\thetabcounter}{%
        name={\unexpanded{#1}},%
        sort={\thetabcounter},%
        description={p.~\noexpand\pageref{A:\thetabcounter}},%
        user1={\thetabcounter},%
        user2={Table},%
        user3={\thesection},%
        user4={\expandafter\unexpanded\expandafter{\@savedSecTitle}}%
      }%
    }%
    \protected@write\@auxout{}{\string\glsliveentry{A\thetabcounter}}%
  \fi
}
\newcommand{\eeTa}{\end{table}}

% FIGURE COMMANDS (glossaried) -- consistent with \beT, \beR, \beTa, etc.
% Usage:
%   \beF[Glossary title]...\caption{Visible caption}\eeF  -> glossaried figure
%   \beF...\caption{Visible caption}\eeF                  -> plain [H] figure, no glossary
\newcommand{\beF}[1][]{%
  \def\@temparg{#1}%
  \ifx\@temparg\@empty
    \begin{figure}[H]%
  \else
    \begin{figure}[H]%
    \let\@savedSecTitle\@currentsectionname%
    \refstepcounter{figcounter}%
    \label{fig:\thefigcounter}%
    \expanded{%
      \noexpand\newglossaryentry{Fig\thefigcounter}{%
        name={\unexpanded{#1}},%
        sort={\thefigcounter},%
        description={p.~\noexpand\pageref{fig:\thefigcounter}},%
        user1={\thefigcounter},%
        user2={Figure},%
        user3={\thesection},%
        user4={\expandafter\unexpanded\expandafter{\@savedSecTitle}}%
      }%
    }%
    \protected@write\@auxout{}{\string\glsliveentry{Fig\thefigcounter}}%
  \fi
}
\newcommand{\eeF}{\end{figure}}
\makeatother
% LIVE-ENTRY REGISTER
% \glsaddall prints every entry the glossaries package currently knows, and
% that set includes entries restored from \jobname.glsdefs.  The package
% rewrites .glsdefs from all *currently defined* entries, so anything it
% restored is written straight back out: a deleted section kept reappearing in
% the catalogue for ever, and the .glsdefs had to be erased by hand.  The
% wrappers therefore record every entry they actually create in the .aux, and
% \printboth adds only those.  The .aux is rebuilt from scratch on each run, so
% this is self-healing: a removed result is gone after one further run.
\makeatletter
\newcommand{\@glslive}{}
\newcommand{\@glsaddlive}[1]{\ifglsentryexists{#1}{\glsadd{#1}}{}}
\newcommand{\glsliveentry}[1]{\g@addto@macro\@glslive{\@glsaddlive{#1}}}
\makeatother
% PRINT COMMAND
% No \clearpage here -- the catalogue continues on the current page.  This has
% been re-added by mistake several times; leave it out.
\makeatletter
\newcommand{\printboth}{%
  \@glslive   % add exactly the entries this document created last run
  \printglossary[style=unifiedstyle, title={0 \hspace{.22em} Definitions, Theorems, Examples and Open Problems}]%
}
\makeatother
% Named theorem reference: \tref{word} -> Thm~N(word-or-short)
% Shows stored short title if available, otherwise shows the key itself.
\newcommand{\tref}[1]{%
  Thm~\ref{t:#1}(\ifcsname @tshort@#1\endcsname\csname @tshort@#1\endcsname\else#1\fi)%
}

% Full-title theorem reference: \nref{t:KEY} prints the number and then the
% theorem's long title, so one writes "Theorem~\nref{t:blocks}" and gets
% "Theorem 3 ``Block factorization and siphons''".  Takes the complete label, like
% \ref, not the bare key like \tref.
%
% \beT writes each title to the .aux (see \@thmtitle below), so \nref resolves
% forward references from the second LaTeX pass, exactly as \ref does.  If the
% title is not yet known -- first pass, or a label that is not a \beT theorem
% -- it degrades to the number rather than producing ??.
% These two use @ in their names, and this point in the file is outside the
% \makeatletter region opened earlier (it closes at \makeatother above), so
% the region has to be reopened here.
\makeatletter
\newcommand{\@thmtitle}[2]{\global\@namedef{@tlong@#1}{#2}}
\newcommand{\nref}[1]{%
  \ref{#1}%
  \ifcsname @tlong@#1\endcsname
    \ ``\csname @tlong@#1\endcsname''%
  \fi
}
% \nm{t:KEY}: the stored long title of any \beT/\beL/\beP/\beR/\beC/\beD/\beE result,
% verbatim and without quotes, so item names follow the macro's first argument.
\newcommand{\nm}[1]{%
  \ifcsname @tlong@#1\endcsname
    \csname @tlong@#1\endcsname%
  \else
    \ref{#1}%
  \fi
}
\makeatother




%% ========================================
%% Other THEOREM ENVIRONMENTS
%% ========================================
\makeatletter
\@ifundefined{lemma}      {\newtheorem{lemma}      {Lemma}}      {}
\@ifundefined{proposition}{\newtheorem{proposition}{Proposition}} {}
\@ifundefined{corollary}  {\newtheorem{corollary}  {Corollary}}  {}
\@ifundefined{conjecture} {\newtheorem{conjecture} {Conjecture}}  {}
\@ifundefined{remark}     {\newtheorem{remark}     {Remark}}     {}
\@ifundefined{question}   {\newtheorem{question}   {Question}}   {}
\@ifundefined{assumption} {\newtheorem{assumption} {Assumption}} {}
\@ifundefined{hypothesis} {\newtheorem{hypothesis} {Hypothesis}} {}
\@ifundefined{fact}       {\newtheorem{fact}       {Fact}}       {}
\makeatother

% \citer: cite the relay-graph threshold theorem of AAHrel. If a document has
% \input the chapter/file containing t:relayRF (so the label resolves), point
% straight at it; otherwise fall back to a bare citation of AAHrel. Shared here
% (not per-paper) so it resolves consistently across every book chapter.
% \citeR: same idea, for the first theorem of AAHRah (t:Pstar).
\makeatletter
\providecommand{\citer}{%
  \@ifundefined{r@t:relayRF}%
    {\cite{AAHrel}}%
    {Theorem~\ref{t:relayRF} of \cite{AAHrel}}%
}
\providecommand{\citeR}{%
  \@ifundefined{r@t:Pstar}%
    {\cite{AAHRah}}%
    {Theorem~\ref{t:Pstar} of \cite{AAHRah}}%
}
\makeatother


% Abbreviations
\def\beQ{\begin{question}}\def\eeQ{\end{question}}
\def\beA{\begin{assumption}}\def\eeA{\end{assumption}}


%% ========================================

%% ========================================
%% COLORS & CUSTOM COMMANDS
%% ========================================
\definecolor{funccolor}{RGB}{25,25,112}
\definecolor{desccolor}{RGB}{64,64,64}
\newcommand{\func}[1]{\texttt{\textcolor{funccolor}{#1}}}
\def\eq{=}
\def\mm{-}

%% ========================================
%% MATHCAL SHORTHANDS
%% ========================================
% Project-wide convention: \mX for a capital X is always \mathcal X.
% Defined for the whole alphabet so the convention never fails to expand.
\newcommand{\mA}{\mathcal A}\newcommand{\mB}{\mathcal B}
\newcommand{\mC}{\mathcal C}\newcommand{\mD}{\mathcal D}
\newcommand{\mE}{\mathcal E}\newcommand{\mF}{\mathcal F}
\newcommand{\mG}{\mathcal G}\newcommand{\mH}{\mathcal H}
\newcommand{\mI}{\mathcal I}\newcommand{\mJ}{\mathcal J}
\newcommand{\mK}{\mathcal K}\newcommand{\mL}{\mathcal L}
\newcommand{\mM}{\mathcal M}\newcommand{\mN}{\mathcal N}
\newcommand{\mO}{\mathcal O}\newcommand{\mP}{\mathcal P}
\newcommand{\mQ}{\mathcal Q}\newcommand{\mR}{\mathcal R}
\newcommand{\mS}{\mathcal S}\newcommand{\mT}{\mathcal T}
\newcommand{\mU}{\mathcal U}\newcommand{\mV}{\mathcal V}
\newcommand{\mW}{\mathcal W}\newcommand{\mX}{\mathcal X}
\newcommand{\mY}{\mathcal Y}\newcommand{\mZ}{\mathcal Z}

%% ========================================
%% BASIC MATH ENVIRONMENTS
%% ========================================
\def\bep{\begin{pmatrix}}\def\eep{\end{pmatrix}}
\def\bev{\begin{vmatrix}}\def\eev{\end{vmatrix}}
\def\bea{\begin{eqnarray*}}\def\eea{\end{eqnarray*}}
\def\bc{\begin{cases}}\def\ec{\end{cases}}
\def\BEN{\begin{enumerate}}\def\EEN{\end{enumerate}}
\def\BI{\begin{itemize}}\def\EI{\end{itemize}}

\newcommand{\be}[1]{\begin{equation}\label{#1}}
\newcommand{\ee}{\end{equation}}
\newcommand{\beq}{\begin{eqnarray}}
\def\eeq{\end{eqnarray}}
\def\eqr{\eqref}\def\fr{\frac}\def\lbl{\label}
\def\Lra{\Longrightarrow}\def\non{\nonumber}\def\cNGM{\cite{Diek,Van,Van08}}
\def\Eq{\Leftrightarrow}\def\b{\beta}\def\uph{\red{UP TO HERE}}
%% ========================================
%% PROBABILITY & STATISTICS
%% ========================================
\renewcommand{\P}{\ensuremath{\mathbb{P}}}  % Note: overrides paragraph symbol
\newcommand{\E}{\ensuremath{\mathbb{E}}}
\newcommand{\ind}[1]{\ensuremath{\mathbbm{1}_{\left\{#1\right\}}}}
\DeclareMathOperator{\Cov}{Cov}
\DeclareMathOperator{\TV}{TV}
\newcommand{\rank}{\operatorname{rank}}

%% ========================================
%% CRN OPERATORS (organizations/ERCs/buffering/siphons framework)
%% ========================================
\def\clos{\mathrm{clos}}
\def\Supply{\mathrm{Supply}}
\def\req{\mathrm{req}}
\def\React{\mathrm{React}}
\def\Prod{\mathrm{Prod}}
\def\Gen{\mathrm{Gen}}
\def\In{\mathrm{In}}
\def\Out{\mathrm{Out}}
\def\Act{\mathrm{Act}}
\def\Min{\mathrm{Min}}

%% ========================================
%% LINE BREAKING
%% ========================================
% Last-resort stretch, used only on paragraphs TeX cannot otherwise break:
% a long \texttt{} token, a comma-separated inline math list (commas are not
% breakpoints in math), or a run of \cite/\ref boxes.  Without it those set as
% overfull \hbox and print into the margin; with it they set slightly loose,
% which is the lesser evil.  Does not affect paragraphs that already break.
\emergencystretch=3em

%% ========================================
%% NUMBER SETS
%% ========================================
\newcommand{\R}{\mathbb{R}}  % Real numbers
\newcommand{\C}{\mathbb{C}}  % Complex numbers
\newcommand{\N}{\mathbb{N}}
\newcommand{\Z}{\mathbb{Z}}
% \providecommand, not \newcommand: some papers define these in their own
% preamble after \input{preamble}, and must not clash.
\providecommand{\Q}{\mathbb{Q}}
\providecommand{\res}{\operatorname{res}}   % resultant; used by epiRes.tex
\newcommand{\MN}{\mathbb{N}}
\newcommand{\MZ}{\mathbb{Z}}
\newcommand{\realpos}{\mathbb{R}_{> 0}}
\newcommand{\realnneg}{\mathbb{R}_{\geq 0}}
\newcommand{\naturalpos}{\mathbb{N}_{>0}}
\newcommand{\Rplus}{\mathbb{R}_+}
\newcommand{\Rp}{\mathbb{R}_{\ge 0}}
\def\Rnne{\mathbb{R}_{\ge 0}}
\def\Se{\bar{S}} \def\Ie{\bar{I}}
\newcommand{\Rpos}{\mathbb{R}_{>0}}
\newcommand{\dd}{\,\mathrm d}
\newcommand{\e}{\mathrm e}
%% ========================================
%% GREEK LETTERS (2+ character abbreviations) greek
%% ========================================
\def\al{\alpha}\def\ba{\bff \alpha}
\def\w{\omega}\def\Lgn{\ga+\mu+\de}\def\om{\omega}
\def\de{\delta}\def\De{\Delta}\def\bde{\bff \delta}
\def\ep{\epsilon}\def\eps{\varepsilon}
\def\f{\varphi}  % Kept as requested
\def\ga{\gamma}\def\Ga{\Gamma}\def\G{\Gamma}\def\bg{\bff \gamma}
\def\up{\upsilon}  % NOT the breve accent: no accent commands anywhere.
                   % Named \up, not \u, so the LaTeX breve stays reachable.
\def\la{\lambda}\def\La{\Lambda}
\def\si{\sigma}\def\ka{\kappa}
\def\th{\theta}\def\vt{\vartheta}
\def\ze{\zeta}\def\Gap{\Gamma^{+}} \def\Gam{\Gamma^{-}}


%% ========================================
%% SPECIAL SHORT COMMANDS (kept as requested)
%% ========================================
\def\I{\infty}
\def\T{\widetilde}
\def\H{\hat}
\newcommand{\relint}{\operatorname{relint}}

%% ========================================
%% BOLD/VECTOR NOTATION
%% ========================================
\newcommand{\bff}[1]{{\mbox{\boldmath$#1$}}}
\newcommand\X{\boldsymbol{X}}
\def\x{\boldsymbol{x}}\def\bx{\bff x}\def\xe{\bar{\x}}
\newcommand\y{\boldsymbol{y}}\def\by{\boldsymbol{y}}\def\ye{\bar{\y}}
\def\yb{\y}\def\xb{\x}
\def\xy{\x-\y}
\newcommand\br{\boldsymbol{r}}\def\re{\bar{\r}}
\newcommand\bd{\boldsymbol{d}}
\newcommand\bQ{\boldsymbol{Q}}
\newcommand{\Yb}{{\mathbf{Y}}}
\newcommand{\Ab}{\mathbf{A}}
\newcommand{\zb}{\mathbf{z}}\def\bz{\bff 0}
\newcommand{\Zb}{\mathbf{Z}}
\newcommand{\bb}{\mathbf{b}}
\def\vbe{\vec {\bff \beta}}
\def\vb{\vec {\bff b}}
\def\v1{\vec {\bff 1}}
\def\vn{\vec \de}
\def\vz{\vec 0}
\def\vi{\overset{\rightarrow}{\mathsf i}}
\def\vv{\overset{\rightarrow}{v}}
\def\vr{\; \vec {\mathsf r}}
\newcommand{\va}{\bff a}
\providecommand{\vect}[1]{\boldsymbol{#1}}   % bold vector (moved from Gradostat.tex)
\providecommand{\one}{\mathbf 1}             % all-ones vector (moved from AliSF)
\providecommand{\spr}{s}                     % spectral bound  (moved from AliSF)
\providecommand{\Ke}{K_{\rm eff}}            % effective carrying capacity
% Critical right/left eigenvectors in the Hopf projection method (AAHOSN/New.tex)
\providecommand{\bq}{\mathbf{q}}
\providecommand{\bp}{\mathbf{p}}


%% ========================================
%% MATHCAL & MATHSCR
%% ========================================
% Mathcal
\newcommand\mc[1]{\mathcal{#1}}
\def\mC{{\mathcal C}}\def\mD{\mathcal D}\def\mE{\mathcal{E}}
\def\mF{\mathcal F}\def\mG{{\mathcal G}}\def\mH{{\mathcal H}}
\def\mI{\mathcal I}\def\mL{\mathcal{L}}\def\mM{{\mathcal M}}
\def\mN{{\mathcal N}}\def\mP{\mathcal{P}}\def\mR{{\mathcal R}}
\def\mS{{\mathcal S}}\def\mT{\mathcal{T}}\def\mV{{\mathcal V}}
\def\mK{\mathcal K}\def\mA{{\mathcal A}}\def\mU{{\mathcal U}}
% Symmetric bilinear and trilinear forms of the Taylor expansion, used in the
% first Lyapunov coefficient (AAHOSN/New.tex)
\providecommand{\Bform}{\mathcal{B}}
\providecommand{\Cform}{\mathcal{C}}

% Blackboard bold.  Kept here, not in the drivers: a \providecommand in a
% driver preamble runs before this file and silently wins, so the same name
% ends up meaning different things in different builds.
\def\IC{\mathbb{C}}\def\IN{\mathbb{N}}\def\IQ{\mathbb{Q}}
\def\IR{\mathbb{R}}\def\IZ{\mathbb{Z}}\def\IG{\mathbb{G}}
\def\II{\mathbb{I}}

% Mathscr
\newcommand{\Rs}{{\mathscr R}}
\newcommand{\Ss}{{\mathscr S}}
\newcommand{\Ns}{{\mathscr N}}
\newcommand{\Cs}{\mathscr C}


%% ========================================
%% OPERATORS & FUNCTIONS
%% ========================================
\DeclareMathOperator{\tr}{tr}
\DeclareMathOperator{\diag}{Diag}
\DeclareMathOperator{\Diag}{Diag}
\DeclareMathOperator{\intr}{\mathrm{Int}}
\DeclareMathOperator{\Inv}{Inv}
\DeclareMathOperator{\spec}{spec}
\DeclareMathOperator{\AUC}{AUC}               % area under the curve
\newcommand{\ii}{\mathrm{i}}                  % imaginary unit
\newcommand\op{\operatorname}
\newcommand{\supp}{\op{supp}}
\newcommand{\diff}{\mathop{}\mathopen{}\mathrm{d}}
\newcommand{\pdd}[1]{\frac{\partial}{\partial{#1}}}
\DeclareMathOperator{\cone}{cone}
\DeclareMathOperator{\proj}{proj}

%% ========================================
%% BRACKETS & NORMS
%% ========================================
\providecommand{\pr}[1]{\left(#1\right)}
\providecommand{\pp}[1]{\left[#1\right]}
\newcommand\croc[1]{\left\langle #1\right\rangle}
\newcommand{\norm}[1]{\|#1\|}
\newcommand{\inner}[2]{\langle #1, #2 \rangle}
% Short synonym of \inner, used by AAHOSN/New.tex as $\ip{u}{v}=\bar u^\top v$
\providecommand{\ip}[2]{\langle #1 , #2 \rangle}


%% ========================================
%% FORMATTING & UTILITIES
%% ========================================
\newcommand*{\Scale}[2][4]{\scalebox{#1}{$#2$}}
\newcommand{\te}[1]{\text{#1}}
\newcommand\steq[1]{\stackrel{\text{\rm #1.}}{=}}
\newcommand\blue[1]{\textcolor{blue}{#1}}
\newcommand{\red}{\textcolor[rgb]{1.00,0.00,0.00}}
\newcommand{\green}{\textcolor[rgb]{0.00,1.00,0.00}}

% Figure command
\newcommand{\figu}[3]{%
\begin{figure}[H]
\centering
\includegraphics[scale=#3]{#1}
\caption{#2\label{f:#1}}
\end{figure}
}

% Footnote
\long\def\symbolfootnote[#1]#2{%
\begingroup
\def\thefootnote{\fnsymbol{footnote}}\footnote[#1]{#2}%
\endgroup}
\def\fn{\symbolfootnote}

\newcommand{\dsum}{\displaystyle\sum}

%% ========================================
%% CRN & EPIDEMIOLOGY SPECIFIC
%% ========================================
% CRN notation
\def\SS{\mathcal S}
\def\CC{\mathcal C}
\def\RR{\mathcal R}

% Reproduction numbers
\def\Rto{R_{1}^{E_2}}
\def\Rot{R_{2}^{E_1}}
\def\Roc{R_{1,c}}
\def\Rtc{R_{2,c}}
\def\rd{r_{0}}
\def\sd{s_{0}}
\def\dd{D_{0}}
\def\vd{v_{0}}
\def\zd{z_{0}}
\def\se{\sE}
\def\xd{x_{0}}
%% ========================================
%% SPECIAL VALUES & EQUILIBRIA
%% ========================================
\def\eE{\bar{e}}
\def\wE{\bar{w}}
\def\sE{s_{*}}
\def\Me{\T M_{ee}}
\def\Rnne{\mathbb{R}_{\ge 0}}

\newcommand{\ones}{\mathbf 1}


%% ========================================
%% ABBREVIATIONS (alphabetical by category)
%% ========================================

% -- General mathematical terms --
\def\adm{admissible}\def\admy{admissibility}
%\def\and{antisymmetric}% disabled 2026-09-30: it overrode LaTeX's author separator \and (printed "antisymmetric" between authors)
\def\ass{assumption}
\def\cadlag{c\`adl\`ag}
\def\ch{characteristic polynomial}
\def\chp{characteristic polynomial}
\def\coe{coefficient}
\def\comp{computation}
\def\con{condition}
\def\coo{coordinate}
\def\corr{corresponding}
\def\dec{decomposition}
\def\den{denominator}
\def\dete{deterministic}
\def\eig{eigenvalue}\def\eval{eigenvalue}
\def\eigv{eigenvector}\def\eiv{eigenvector}\def\evec{eigenvector}
\def\elt{element}
\def\enn{essentially nonnegative/positivity preserving}
\def\equ{equation}
\def\etal{et al.}
\def\fe{for example}
\def\fno{from now on}
\def\for{\forall}
\def\form{formula}
\def\Fr{Furthermore, }
\def\frt{furthermore}
\def\fun{function}
\def\gene{generalization}
\def\heu{heuristic}
\def\how{however}\def\How{However, }
\def\ie{i.e. }
\def\im{\item}
\def\isd{is defined by}
\def\itc{it is easy to check that}
\def\itf{it follows that}\def\Itf{It follows that}
\def\ith{it holds that}
\def\itm{it may be checked that}\def\Itm{It may be checked that}
\def\ito{it turns out that}\def\Ito{It turns out that}\def\Itou{It turns out useful to}
\def\ma{mass action}\def\MAR{mass action representation}
\def\mat{matrix}
\def\mbc{may be checked to be}
\def\mbw{may be written as}\def\mubw{must be written as}
\def\mpr{more precisely, }\def\Mp{More precisely, }
\def\nc{non-catalytic}
\def\nne{non-negative}\def\nny{nonnegativity}
\def\no{\nonumber}
\def\noni{non-infectious}
\def\nota{notation}
\def\np{\newpage}
\def\Nt{Note that}
\def\nt{non trivial}\def\nz{non zero}
\def\oth{otherwise}
\def\Oth{On the other hand}
\def\pa{parameter}\def\para{parameter}
\def\pol{polynomial}
\def\prob{probability}
\def\Prf{{\bf  Proof:}}\def\prf{{\bf  Proof:}}
\def\prop{proportional to}
\def\qu{\quad}
\def\QED{\hfill {$\square$}\goodbreak \medskip}
\def\rep{representation}
\def\resp{respectively}
\def\satd{satisfied}\def\satg{satisfying}\def\sats{satisfies}\def\saty{satisfy}
\def\ssec{\subsection}\def\sssec{\subsubsection}
\def\sef{see for example}
\def\sksy{skew-symmetric}
\def\sub{substrate}
\def\sys{system}
\def\var{variable}
\def\wk{well-known}
\def\wl{without loss of generality}\def\wlo{w.l.o.g.}
\def\wmc{We may check that}\def\Wmc{We may check that}
\def\wmr{we may rewrite}
\def\wms{we must show that}
\def\wrt{with respect to}

% -- CRN & Chemistry --
\newcommand\CRN{chemical reaction network}
\newcommand\CRNs{chemical reaction networks}
\newcommand\RN{reaction network}
\def\ACR{absolute concentration robustness}
\def\AMSD{acyclic siphon decomposition}
\def\DAE{differential algebraic equations}
\def\FHJ{Feinberg-Horn-Jackson}
\def\GMAK{generalized mass action kinetics}
\def\MM{Michaelis-Menten}
\def\sm{stoichiometric matrix}\def\Sm{stoichiometric matrix}
\def\WRZD{weakly reversible, zero deficiency}\def\aut{autocatalytic}

% -- Epidemiology --
\def\brn{basic reproduction number}\def\bzn{basic replacement number}
\def\com{compartment}
\def\dfe{disease free equilibrium}\def\DFE{disease free equilibrium}
\def\dwt{dwell times}\def\edt{expected dwell times}
\def\EM{epidemic model}
\def\endp{endemic fixed point}\def\enp{endemic point}\def\EE{endemic point}
\def\ME{mathematical epidemiology}
\def\MEE{eco-epidemiology}\def\EEP{eco-epidemiological}
\def\ece{eco-epidemiological model}
\def\mgm{multi group model}
\def\mtS{multi-strain model}\def\mSt{multi-strain model}
\def\mSi{minimal siphon}
\def\NGM{next generation matrix}
\def\nim{new infections matrix}
\def\qem{quadratic epidemiological model}
\def\repF{invasion functions}
\def\SPF{SIR-PH-FA model}
\def\sus{susceptible}

% -- Graph theory & Linear algebra --
\def\jac{Jacobi matrix}
\def\jin{Jacobian matrix of the invasion vector field with respect to the invasion variables}
\def\jrn{Jacobian matrix of the resident vector field with respect to the resident variables}
\def\jxn{Jacobian matrix of the $\x$ vector field with respect to the $\x$ variables}
\def\Lm{Laplacian matrix}
\def\MTT{matrix tree theorem}
\def\SGT{spectral graph theory}
\def\SMW{Sherman–Morrison–Woodbury}

% -- Stability & Dynamics --
\def\absC{absorbtion CTMC}
\def\BTb{Bogdanov-Takens bifurcation}
\def\CEP{competitive exclusion partition}
\def\CTMC{continuous time Markov chain}
\def\ds{differential system}
\def\fp{fixed point}\def\bfp{boundary fixed point}\def\fbp{fixed boundary point}
\def\fbl{fixed boundary line}\def\fif{forward invariant face}
\def\gas{global asymptotic stability}\def\GAS{globally asymptotically stable}
\def\GASS{\text{GASS}}
\def\Hb{Hopf bifurcation}
\def\Hur{Hurwitz}\def\MH{Metzler-Hurwitz}\def\RH{Routh-Hurwitz}
\def\IS{\text{IS}}
\def\LAS{locally asymptotically stable}
\def\LdS{Lyapunov diagonally stable}
\def\Lf{Lyapunov function}\def\LyV{Lyapunov Volterra type function}\def\RLF{robust Lyapunov function}
\def\LyM{Lyapunov Malkin stable}
\def\PF{Perron-Frobenius}
\def\spfp{strictly positive fixed point}
\def\sta{stability}
\def\STP{sharp threshold property}\def\STT{sharp threshold theorem}

% -- Models (with \varphi explicitly instead of \f) --
\def\AB{$(A,B)$ Arino-Brauer epidemic models}
\def\ABa{$(A,\bbe,\ba,\varphi)$ balanced bilinear model}
\def\ABP{%$(A,\bbe,P,\varphi)$
\\ balanced bilinear model}
\def\AMG{Anderson May Gupta}
\newcommand\BIN{biological interaction network}
\def\Dt{Descartes type}
\def\gLV{generalized Lotka-Volterra model}
\def\IG{\text{invasion graph}}\def\IGMS{interaction graph on minimal siphons}
\def\LCP{\text{linear complementarity problem}}
\def\LV{Lotka-Volterra}\def\LVK{Lotka-Volterra-Kolmogorov}

% -- Splitting & Structure --
\def\reg{regular splitting}\def\regS{{\bf regular splitting}}
\def\RUR{rational univariate representation}
\def\bds{block-diagonal structure}\def\NGMs{next generation matrices}

% -- Miscellaneous --
\def\Ma{{Mathematica}}
\def\Mm{Max–Min}
\def\num{numerator}
\def\rap{rational approximation} \def\DSR{directed species-reaction bipartite graph}\def\PCP{polynomial complementarity problem}\def\rat{rational substitution}
\def\QSS{quasi steady state} \def\RHS{right hand side} \def\LHS{left hand side}
\def\OSN{online social network}\def\VdD{Van den Driessche}
\def\wli{$\omega$-limit} \def\rk{rich kinetics}\def\sJ{symbolic Jacobian}
\def\cV{\cite{Vas,VasSta,BSV,GSV}} \def\spf{semi-parametric semi-functional}
\def\CRW{Capasso-Ruan-Wang} \def\ABP{$(A,\bbe,P,\f)$ balanced bilinear model}
\def\ABa{%$(A,\bbe,\ba,\f)$ \\
rank-one balanced bilinear model}
\def\bbe{B%{\bff \beta}
} \def\vbe{\vec {\bff \beta}
}
\def\vb{\vec {\bff b}}
\newcommand{\Kh}{\hat{K}}          % K with hat, first mention only
\newcommand{\wS}{w^{\sigma}}       % right Perron J^\perp, full block
\newcommand{\piS}{\vec \pi^{\sigma}}    % left Perron J^\perp, full block
\newcommand{\ws}{w^{\sigma}}       % right Perron J^\perp, sub-block
\newcommand{\pis}{\pi^{\sigma}}    % left Perron J^\perp, sub-block
\newcommand{\hwS}{\hat{w}^{\sigma}}   % right Perron K, full block
\newcommand{\hpiS}{\hat{\pi}^{\sigma}} % left Perron K, full block
\newcommand{\hws}{\hat{w}^{\sigma}}   % right Perron K, sub-block
\newcommand{\hpis}{\hat{\pi}^{\sigma}} % left Perron K, sub-block
\newcommand{\twS}{\tilde{w}^{\sigma}}  % right Perron \tilde K, full block
\newcommand{\tpiS}{\tilde{\pi}^{\sigma}} % left Perron \tilde K
\newcommand{\tws}{\tilde{w}^{\sigma}}
\newcommand{\tpis}{\tilde{\pi}^{\sigma}}
\newcommand{\Mat}[2]{\mathbb{R}_{\ge 0}^{\,#1\times #2}}
\def\tJ{J_\sigma^\perp} \def\lPe{left Perron eigenvector}
% \G is \Gamma (defined at :795).  The Volterra function keeps its own name,
% so that \G no longer has two conflicting meanings in this file.
\def\GV{G}  % Volterra G(u)=u-1-\ln u
\def\PVL{Perron-Volterra Lyapunov minimal siphon ansatz}
\newcommand{\istar}[1]{i_{#1}^{*}}
%\newcommand{\sE}{s_E}
\newcommand{\vE}{v_E}
\def\PV{Perron-Volterra candidate Lyapunov function}
\def\GLV{generalized Lotka-Volterra}
\def\MERN{mathematical epidemiology reaction networks}
\def\CMM{contraction metric method}
\newcommand{\sign}{\operatorname{sign}}
\newcommand{\Var}{\operatorname{Var}} 

% ======================================================================
% IDEAS SHARED BETWEEN THE ABSTRACT AND THE MAIN RESULTS SECTION
%
% Write each idea ONCE, in the abstract; the Main Results section recalls it,
% so the two can never drift apart.
%
%   in the abstract        \ideaitem{1}{The first-order reaction terms ...}
%   in Main Results        \ideahead{1}
%
% \ideaitem stores the text globally and typesets
%   \item \hyperlink{idea:1}{<text>}
% \ideahead typesets
%   \item\hypertarget{idea:1}{}\textbf{<the same text>}
%
% The abstract must come before the Main Results input, which is the normal
% order. If a heading has to read differently from the abstract sentence (for
% instance a long abstract item summarised in one line), override it with the
% optional argument: \ideahead[short heading]{6}.
% ======================================================================
\makeatletter
\newcommand{\ideaitem}[2]{%
  \expandafter\gdef\csname @idea@#1\endcsname{#2}%
  \item\hyperlink{idea:#1}{#2}}
\newcommand{\ideahead}[2][]{%
  \item\hypertarget{idea:#2}{}%
  \def\@ideatmp{#1}%
  \ifx\@ideatmp\@empty
    \@ifundefined{@idea@#2}{%
      \PackageWarning{preamble}{Idea #2 is used in Main Results but was never
        defined by \string\ideaitem\space in the abstract}%
      \textbf{[idea #2 undefined]}%
    }{\textbf{\begingroup\def\ideaterm##1##2{##2}\def\idearest##1{}\@nameuse{@idea@#2}\endgroup}}%
  \else
    \textbf{#1}%
  \fi}
% \ideaterm{<hypertarget>}{<words>}: inside an abstract idea the words become an
% invisible link (the document sets \hypersetup{hidelinks}) jumping to the place
% in Main Results that defines them; inside the Main Results heading itself,
% where the target sits a few lines below, \ideahead strips the link and prints
% the words plainly. Mark the destination with \hypertarget{<name>}{} on the line
% before the display, which adds no equation number and shows nothing.
\newcommand{\ideaterm}[2]{\hyperlink{#1}{#2}}
% \idearest{<text>}: the part of an abstract idea after its first sentence; printed in
% the abstract, omitted from the Main Results heading.
\newcommand{\idearest}[1]{#1}
\makeatother

% ======================================================================
% SHORT-PAPER SWITCH
% \shorttrue in a master that assembles only part of the companion paper
% (Barrat/SipArxiv.tex). \nsh{...} drops text that cites results living in the
% omitted sections, so the short build has no dangling references while the
% companion paper keeps the text verbatim.
% ======================================================================
\newif\ifshort \shortfalse
\newcommand{\nsh}[1]{\ifshort\else#1\fi}
% \shortalt{<short build>}{<companion paper>}: use when a sentence needs
% different wording in the two builds, not merely a deletion.
\newcommand{\shortalt}[2]{\ifshort#1\else#2\fi}
% \Jtr, \JtrS{\Sigma}: transverse Jacobian. J^\perp here; SipArxiv.tex redefines them as J_{\mathrm{tra}}, J_{\mathrm{tra},\Sigma}
\newcommand{\Jtr}{J^\perp}
\newcommand{\JtrS}[1]{J_{#1}^\perp}
.
\makeatletter
\@ifpackageloaded{hyperref}{}{%
  \@ifundefined{hyperrefpackage}{%
    \@ifundefined{AAdefNnohyperref}{\usepackage{hyperref}}{}%
  }{}%
}
\makeatother
% linktoc=all makes a table-of-contents page number a link as well as the
% entry text; hyperref's default, linktoc=section, links only the text.
% Deferred to begin-document and guarded: by the block above hyperref may
% be loaded later (mdpi's own hook) or not at all (\AAdefNnohyperref), and
% \hypersetup must exist before it is called.  A driver that wants the
% default back can say \hypersetup{linktoc=section} after % Shared preamble for the GLVbook papers.  This is the single canonical copy;
% a driver one level below GLVbook reaches it as
%
%     \input{../basicFiles/preamble}
%
% What stays in the driver, because it differs from paper to paper:
% \documentclass, font choice, \title, \author, \date, and any macro genuinely
% local to one paper.  Everything else belongs here.  The file is
% self-contained: it does not read any other file.

\usepackage[utf8]{inputenc}
\usepackage[novbox]{pdfsync}   % required: SyncTeX alone does not work here

% Page geometry, set once here for every paper that loads this file.  A driver
% that needs different margins says \geometry{...}; a second
% \usepackage[...]{geometry} raises an option clash and must not be used.
\makeatletter
\@ifundefined{AAdefNnogeometry}{\usepackage[a4paper,margin=2cm]{geometry}}{}
\makeatother

% mathematics
\usepackage{amsmath,amssymb,amsthm,mathtools,bm,mathrsfs,bbm}
\usepackage{centernot}

% floats, tables, lists
\usepackage{graphicx,float}
\makeatletter
\@ifpackageloaded{subfig}{}{\usepackage{subcaption}}
\makeatother
\usepackage{booktabs,longtable,threeparttable,array,multirow}
\usepackage{enumitem}

% colour, code, algorithms
\usepackage{xcolor,listings,tcolorbox}
\makeatletter
\@ifundefined{c@algorithm}{\usepackage{algorithm}}{}
\makeatother
\usepackage{algorithmicx,algpseudocode}

% diagrams
\usepackage{tikz,tikz-cd,pgfplots}
\usetikzlibrary{graphs,decorations.pathmorphing,decorations.markings,
arrows.meta,automata,positioning}
\tikzcdset{every label/.append style = {font = \small}}
\pgfplotsset{compat=1.18}

% chemistry
\usepackage[version=4]{mhchem}

\usepackage{microtype}
\usepackage{lineno}
\usepackage{nomencl}
\makenomenclature

% hyperref first, then bookmark, then glossaries: glossaries must follow
% hyperref or the entry links are built against the wrong anchors.
\makeatletter
\@ifundefined{hyperrefpackage}{%
  \@ifundefined{AAdefNnohyperref}{\usepackage{hyperref}}{}%
}{}
\@ifundefined{hyperrefpackage}{\usepackage{bookmark}}{}
\makeatother
\usepackage[toc,sort=use]{glossaries}
\makeglossaries

% Book/standalone switch.  Material that the book already covers in its
% prelude, but that a standalone paper still has to state, is wrapped as
%     \ifbookbuild\else ... \fi
% Drivers default to standalone; bookM.tex says \bookbuildtrue.
\newif\ifbookbuild
\bookbuildfalse

% Part switches for the gradostat paper.  Both default to true; a driver that
% wants the short submission says \madecfalse and/or \slowfastfalse, and every
% cross-reference into the dropped part disappears with it.
%   \ifmadec     Part III, the Lagasquie--Madec model and its break-even geometry
%   \ifslowfast  Part IV, everything carrying the small parameter \eps
\newif\ifmadec     \madectrue
\newif\ifslowfast  \slowfasttrue

% Part switch for the Barrat/simplicial-contagion project.  Defaults to
% true (the hypernetwork-synchrony part is included); the journal-paper
% driver says \hnetpartfalse to drop it and every cross-reference into it.
\newif\ifhnetpart  \hnetparttrue

%% ========================================
%% PACKAGES
%% ========================================
\RequirePackage{etoolbox}
\usepackage{makecell}
\usepackage[novbox]{pdfsync}% obsolete; SyncTeX is built into modern engines
\usepackage{amssymb, amsfonts,mathtools,bbm,eurosym,graphicx,latexsym, amsmath,amsthm,times,nccmath,url}
\usepackage{listings,verbatim,cases,tabularx}
% hyperref is deliberately NOT in the line above.  Journal classes often
% load it themselves, with a long option list, from a begindocument hook
% (mdpi.cls does: \DeclareHookRule{begindocument}{hyperref}{before}{mdpi}).
% Loading it here first, without options, makes that later load raise
% "Option clash for package hyperref".  So load it only if nobody else
% has, or will:
%   - \@ifpackageloaded  catches a class that loaded it already;
%   - \hyperrefpackage    is mdpi.cls's own deferred loader;
%   - \AAdefNnohyperref   lets any other driver opt out by defining it
%                         before \input{preamble}.
\makeatletter
\@ifpackageloaded{hyperref}{}{%
  \@ifundefined{hyperrefpackage}{%
    \@ifundefined{AAdefNnohyperref}{\usepackage{hyperref}}{}%
  }{}%
}
\makeatother
% linktoc=all makes a table-of-contents page number a link as well as the
% entry text; hyperref's default, linktoc=section, links only the text.
% Deferred to begin-document and guarded: by the block above hyperref may
% be loaded later (mdpi's own hook) or not at all (\AAdefNnohyperref), and
% \hypersetup must exist before it is called.  A driver that wants the
% default back can say \hypersetup{linktoc=section} after \input{preamble}.
\makeatletter
\AtBeginDocument{\@ifundefined{hypersetup}{}{\hypersetup{linktoc=all}}}
\makeatother
\usepackage{graphicx,float,fancybox}
\usepackage{epsfig,epstopdf,tikz}
\usepackage[sans]{dsfont}
\usetikzlibrary{positioning,arrows,arrows.meta,calc}
% Configure listings for Mathematica code
\lstdefinestyle{mathematica}{
 language=Mathematica,
 basicstyle=\small\ttfamily,
 backgroundcolor=\color{gray!10},
 frame=single,
 breaklines=true,
 commentstyle=\color{green!60!black},
 keywordstyle=\color{blue},
 stringstyle=\color{red},
 showstringspaces=false,
 tabsize=2
}
\usepackage{tcolorbox}
% Colored boxes for key results
\newtcolorbox{keyresult}[1][]{
 colback=blue!5!white,
 colframe=blue!75!black,
 title=#1,
 fonttitle=\bfseries
}
\newtcolorbox{examplebox}[1][]{
 colback=green!5!white,
 colframe=green!75!black,
 title=#1,
 fonttitle=\bfseries
}


\makeatletter
\@ifundefined{theorem}   {\newtheorem{theorem}   {Theorem}}   {}
\@ifundefined{definition}{\newtheorem{definition}{Definition}} {}
\@ifundefined{open}      {\newtheorem{open}      {Problem}}   {}
\@ifundefined{example}   {\newtheorem{example}   {Example}}   {}
\makeatother

% COUNTER DEFINITIONS
\newcounter{tabcounter}
\newcounter{defcounter}
\newcounter{opcounter}
\newcounter{thmcounter}
\newcounter{excounter}
\newcounter{ascounter}
\newcounter{cjcounter}
\newcounter{lemcounter}
\newcounter{remcounter}
\newcounter{corcounter}
\newcounter{propcounter}
\newcounter{figcounter}

% UNIFIED GLOSSARY STYLE
\newglossarystyle{unifiedstyle}{%
  \setglossarystyle{list}%
  \renewcommand*{\glossarypreamble}{\gdef\glslastsec{}}%
  \renewcommand{\glossentry}[2]{%
    \edef\glscursec{\glsentryuseriii{##1}}%
    \ifx\glscursec\glslastsec\else
      \item[]\textbf{\S\,\glscursec\enspace\glsentryuseriv{##1}}%
      \global\let\glslastsec\glscursec
    \fi
    \item[\glsentryuserii{##1} \glsentryuseri{##1}:] ``\glossentryname{##1}'' \glsentrydesc{##1}%
  }%
}

\makeatletter
% First-word extractor (used by \beT, \beD and \beXa for the default key)
% \@firstwordof hello world \relax  ->  hello
\def\@firstwordof#1 #2\relax{#1}
% First three letters, used for the heading mnemonic of a theorem given no
% explicit key.  Padding with \@empty keeps it safe for first words shorter
% than three letters: \@empty vanishes under \edef, whereas \relax would
% survive into the label and is illegal inside \csname (it broke every
% two-letter title, e.g. "ME ...", with "Missing \endcsname inserted").
\def\@takethree#1#2#3#4\@nil{#1#2#3}

% EXAMPLE COMMAND -- same signature and behaviour as \beT and \beD.
% Usage: \beXa*[long title][short title][key]
%   star  = suppresses the printed "[key]" tag after the title (as for \beT*
%   and \beD*); the label and glossary entry are unaffected.
%   key = optional label key override; the label becomes e:KEY, and without a
%   key it is e: plus the first three letters of the title's first word.  A
%   mnemonic label lets an example be cited as "Example~\ref{e:KEY}" instead
%   of by section number, which is what the citing prose needs.
\NewDocumentCommand{\beXa}{s O{} O{} o}{%
  \def\@temparg{#2}%
  \ifx\@temparg\@empty
    \begin{example}%
  \else
    \let\@savedSecTitle\@currentsectionname%
    \refstepcounter{excounter}%
    \IfValueTF{#4}{%
      \edef\@exlblkey{#4}%
    }{%
      \edef\@exfirstword{\@firstwordof#2 \relax}%
      \edef\@exlblkey{%
        \expandafter\@takethree\@exfirstword\@empty\@empty\@empty\@nil}%
    }%
    \def\@tempshort{#3}%
    \ifx\@tempshort\@empty
      \expandafter\xdef\csname @xshort@\@exlblkey\endcsname{\@exlblkey}%
    \else
      \expandafter\def\csname @xshort@\@exlblkey\endcsname{#3}%
    \fi
    \IfBooleanTF{#1}{\let\@extag\@empty}{%
      \edef\@extag{\noexpand\ \noexpand\textup{[\noexpand\texttt{e:\@exlblkey}]}}}%
    \begin{example}[#2\ifshowkeys\@extag\fi]%
    \label{e:\@exlblkey}%
    % record the long title so \nref{e:KEY} prints it, as for theorems
    \protected@write\@auxout{}{%
      \string\@thmtitle{e:\@exlblkey}{\unexpanded{#2}}}%
    \expanded{%
      \noexpand\newglossaryentry{e\theexcounter}{%
        name={\unexpanded{#2}%
          \IfBooleanF{#1}{ [\noexpand\texttt{e:\@exlblkey}]}},%
        sort={\theexample},%
        description={p.~\noexpand\pageref{e:\@exlblkey}},%
        user1={\theexample},%
        user2={Example},%
        user3={\thesection},%
        user4={\expandafter\unexpanded\expandafter{\@savedSecTitle}}%
      }%
    }%
    \protected@write\@auxout{}{\string\glsliveentry{e\theexcounter}}%
  \fi
}
\newcommand{\eeXa}{\end{example}}

% DEFINITION COMMAND -- same signature and behaviour as \beT below.
% Usage: \beD[long title][short title][key]
%   long title  = definition heading (plain text only)
%   short title = shown in \dref{word}; defaults to the key
%   key         = optional label key override; the label becomes d:KEY
% Without a key the label is d: plus the first three letters of the title's
% first word, so it is a mnemonic rather than the definition's number: a
% numeric label would silently change whenever a definition is inserted.
\NewDocumentCommand{\beD}{s O{} O{} o}{%
  \def\@temparg{#2}%
  \ifx\@temparg\@empty
    \begin{definition}%
  \else
    \let\@savedSecTitle\@currentsectionname%
    \refstepcounter{defcounter}%
    \IfValueTF{#4}{%
      \edef\@deflblkey{#4}%
    }{%
      \edef\@deffirstword{\@firstwordof#2 \relax}%
      \edef\@deflblkey{%
        \expandafter\@takethree\@deffirstword\@empty\@empty\@empty\@nil}%
    }%
    \def\@tempshort{#3}%
    \ifx\@tempshort\@empty
      \expandafter\xdef\csname @dshort@\@deflblkey\endcsname{\@deflblkey}%
    \else
      \expandafter\def\csname @dshort@\@deflblkey\endcsname{#3}%
    \fi
    \IfBooleanTF{#1}{\let\@deftag\@empty}{%
      \edef\@deftag{\noexpand\ \noexpand\textup{[\noexpand\texttt{d:\@deflblkey}]}}}%
    \begin{definition}[#2\ifshowkeys\@deftag\fi]%
    \label{d:\@deflblkey}%
    % record the long title so \nref{d:KEY} prints it, as for theorems
    \protected@write\@auxout{}{%
      \string\@thmtitle{d:\@deflblkey}{\unexpanded{#2}}}%
    \expanded{%
      \noexpand\newglossaryentry{d\thedefcounter}{%
        name={\unexpanded{#2}%
          \IfBooleanF{#1}{ [\noexpand\texttt{d:\@deflblkey}]}},%
        sort={\thedefinition},%
        description={p.~\noexpand\pageref{d:\@deflblkey}},%
        user1={\thedefinition},%
        user2={Definition},%
        user3={\thesection},%
        user4={\expandafter\unexpanded\expandafter{\@savedSecTitle}}%
      }%
    }%
    \protected@write\@auxout{}{\string\glsliveentry{d\thedefcounter}}%
  \fi
}
\newcommand{\eeD}{\end{definition}}

% \showkeysfalse in a driver hides the [t:KEY] mnemonics in result headings
% (labels and glossary entries are unchanged); used for submission builds.
\newif\ifshowkeys \showkeystrue
% \commonredtrue in a driver that inputs basicFiles/commonRed: tool.tex then
% cites its reduction results instead of restating them.
\newif\ifcommonred \commonredfalse
% TheorEM COMMAND
% Usage: \beT[long title][short title][key]
%   long title  = theorem heading (plain text only)
%   short title = shown in \tref{word}; may contain math; defaults to first word of long title
%   key         = optional label key override; label becomes t:KEY (avoids first-word collisions)
% Label is t:WORD where WORD = first word of long title (or KEY if given); use \tref{WORD} to reference.
\NewDocumentCommand{\beT}{s O{} O{} o}{%
  \def\@temparg{#2}%
  \ifx\@temparg\@empty
    \begin{theorem}%
  \else
    \let\@savedSecTitle\@currentsectionname%
    \refstepcounter{thmcounter}%
    % \@thmtag is the mnemonic shown in the heading; every theorem gets one.
    % With an explicit key the key itself is shown.  Without one the key
    % defaults to the title's first word, and showing that whole word would
    % merely repeat the title ("Dissipativity [Dissipativity]"), so only its
    % first three letters are shown: "Dissipativity [Dis]".  The label is
    % unaffected -- it still uses the full \@thmlblkey.
    \IfValueTF{#4}{%
      \edef\@thmlblkey{#4}%
    }{%
      \edef\@thmfirstword{\@firstwordof#2 \relax}%
      \edef\@thmlblkey{%
        \expandafter\@takethree\@thmfirstword\@empty\@empty\@empty\@nil}%
    }%
    \IfBooleanTF{#1}{\let\@thmtag\@empty}{%
      \edef\@thmtag{\noexpand\ \noexpand\textup{[\noexpand\texttt{t:\@thmlblkey}]}}}%
    \def\@tempshort{#3}%
    \ifx\@tempshort\@empty
      \expandafter\xdef\csname @tshort@\@thmlblkey\endcsname{\@thmlblkey}%
    \else
      \expandafter\def\csname @tshort@\@thmlblkey\endcsname{#3}%
    \fi
    \begin{theorem}[#2\ifshowkeys\@thmtag\fi]%
    \label{t:\@thmlblkey}%
    % record the long title for \nref; via the .aux so forward refs resolve
    \protected@write\@auxout{}{%
      \string\@thmtitle{t:\@thmlblkey}{\unexpanded{#2}}}%
    \expanded{%
      \noexpand\newglossaryentry{t\thethmcounter}{%
        name={\unexpanded{#2}%
          \IfBooleanF{#1}{ [\noexpand\texttt{t:\@thmlblkey}]}},%
        % sort/user1 must be the number the document actually prints, i.e. the
        % theorem environment's counter, not thmcounter.  thmcounter advances
        % only on \beT with a title, so a bare \beT (or \beT with an
        % empty title) makes the two drift and the glossary cites a wrong
        % number.  thmcounter is still fine as the entry key: it is unique.
        sort={\thetheorem},%
        description={p.~\noexpand\pageref{t:\@thmlblkey}},%
        user1={\thetheorem},%
        user2={Theorem},%
        user3={\thesection},%
        user4={\expandafter\unexpanded\expandafter{\@savedSecTitle}}%
      }%
    }%
    \protected@write\@auxout{}{\string\glsliveentry{t\thethmcounter}}%
  \fi
}
\newcommand{\eeT}{\end{theorem}}

% CROSS-PART RELATION PARAGRAPH
% Usage: \xpart{text}  -- place immediately BEFORE \beT, never inside it
%   (anything inside \beT...\eeT would be swallowed into the glossary entry).
% The paragraph is typeset only when Part I is actually present, i.e. in the
% combined book build; in the standalone builds (GranadaDSABNS/A.tex,
% AAHRah-Submission/AAHRah.tex) the label part:glv is undefined and the whole
% paragraph disappears rather than leaving ?? cross-references. Same
% graceful-degradation principle as \GASref in GASt.tex. Resolves on the
% second LaTeX pass, since r@part:glv comes from the .aux file.
% (Already inside the \makeatletter region opened above, so no \makeatletter
% pair here -- adding one would close that region early.)
\newcommand{\xpart}[1]{\@ifundefined{r@part:glv}{}{\paragraph{Relation with Parts~I--II.}#1}}

% \bookref{label}{precise text}{fallback text}: prints the precise cross-reference
% when the label resolves (combined book build), and the fallback text otherwise
% (standalone paper builds, where book-only appendix labels are absent). Same
% graceful-degradation role as \GASref in GASt.tex, but usable in any file since it
% lives in this preamble (loaded by every driver) -- unlike \GASref, which is defined in a
% content file and is therefore undefined in standalone builds that omit that file.
\newcommand{\bookref}[3]{\@ifundefined{r@#1}{#3}{#2}}

% \PVpart and \RBpart: the two parts that a partial build may omit.
% They print "Part~N" in any build that contains the part, and a naming phrase
% in a build that does not, so an excerpt shows no ?? .  Each already supplies
% the word "Part", so write "\PVpart develops ...", not "Part~\PVpart".
\newcommand{\PVpart}{\bookref{part:aahrah}%
  {Part~\ref{part:aahrah}}{the Perron--Volterra part of the full book}}
\newcommand{\RBpart}{\bookref{part:recbrowne}%
  {Part~\ref{part:recbrowne}}{the immuno-virology part of the full book}}

% OPEN PROBLEM COMMANDS
% ASSUMPTION COMMANDS (glossaried)
\newcommand{\beAs}[1][]{%
  \def\@temparg{#1}%
  \ifx\@temparg\@empty
    \begin{assumption}%
  \else
    \let\@savedSecTitle\@currentsectionname%
    \refstepcounter{ascounter}%
    \begin{assumption}[#1]%
    \label{as:\theascounter}%
    \expanded{%
      \noexpand\newglossaryentry{as\theascounter}{%
        name={\unexpanded{#1}},%
        sort={\theascounter},%
        description={p.~\noexpand\pageref{as:\theascounter}},%
        user1={\theascounter},%
        user2={Assumption},%
        user3={\thesection},%
        user4={\expandafter\unexpanded\expandafter{\@savedSecTitle}}%
      }%
    }%
    \protected@write\@auxout{}{\string\glsliveentry{as\theascounter}}%
  \fi
}
\newcommand{\eeAs}{\end{assumption}}
% CONJECTURE COMMANDS (glossaried)
\newcommand{\beCj}[1][]{%
  \def\@temparg{#1}%
  \ifx\@temparg\@empty
    \begin{conjecture}%
  \else
    \let\@savedSecTitle\@currentsectionname%
    \refstepcounter{cjcounter}%
    \begin{conjecture}[#1]%
    \label{cj:\thecjcounter}%
    \expanded{%
      \noexpand\newglossaryentry{cj\thecjcounter}{%
        name={\unexpanded{#1}},%
        sort={\thecjcounter},%
        description={p.~\noexpand\pageref{cj:\thecjcounter}},%
        user1={\thecjcounter},%
        user2={Conjecture},%
        user3={\thesection},%
        user4={\expandafter\unexpanded\expandafter{\@savedSecTitle}}%
      }%
    }%
    \protected@write\@auxout{}{\string\glsliveentry{cj\thecjcounter}}%
  \fi
}
\newcommand{\eeCj}{\end{conjecture}}
% LEMMA COMMANDS
% ---------------------------------------------------------------------
% UNIFORM RESULT WRAPPERS -- \beL, \beC, \beP, \beR share \beT's signature
%
%   \beX[long title][short title][key]  ...  \eeX
%
%   long title  = heading (math allowed)
%   short title = shown in the cross-reference; defaults to the key
%   key         = label key.  With a key the label is X:KEY and the heading
%                 carries the mnemonic [key], exactly as for \beT.  Without
%                 one the label falls back to the running counter, which is
%                 what every pre-existing call site in the book relies on.
%
% A bare \beX -- arguments commented out -- opens the plain environment and
% creates NO label and NO glossary entry.  That is the supported way to keep
% an individual result out of the catalogue.
% ---------------------------------------------------------------------
% #1 wrapper name  #2 environment  #3 label prefix  #4 counter  #5 type name
\newcommand{\@mkresult}[5]{%
  % A leading * suppresses the [key] mnemonic in both the heading and the
  % glossary name, while still making the label X:KEY and the entry:
  %     \beX*[long][short][key]
  \expandafter\NewDocumentCommand\csname be#1\endcsname{s O{} O{} o}{%
    \def\@temparg{##2}%
    \ifx\@temparg\@empty
      \begin{#2}%
    \else
      \let\@savedSecTitle\@currentsectionname%
      \refstepcounter{#4}%
      % With a key: mnemonic label X:KEY and the [key] tag in the heading.
      % Without one: fall back to the running counter, exactly as these
      % wrappers behaved before, so every legacy one-argument call site in
      % the book keeps its label.  The counter is also csname-safe, which a
      % title's first word is not -- titles beginning with math would
      % otherwise build \csname \Omega...\endcsname and fail.
      \IfValueTF{##4}{%
        \edef\@relbl{##4}%
        \IfBooleanTF{##1}{\let\@retag\@empty}{%
          \edef\@retag{\noexpand\ \noexpand\textup{[\noexpand\texttt{#3:##4}]}}}%
      }{%
        \edef\@relbl{\arabic{#4}}%
        \let\@retag\@empty
      }%
      \def\@tempshort{##3}%
      \ifx\@tempshort\@empty
        \expandafter\xdef\csname @sshort@#3\@relbl\endcsname{\@relbl}%
      \else
        \expandafter\def\csname @sshort@#3\@relbl\endcsname{##3}%
      \fi
      \begin{#2}[##2\ifshowkeys\@retag\fi]%
      \edef\@tmplbl{#3:\@relbl}\expandafter\label\expandafter{\@tmplbl}%
      \expanded{%
        \noexpand\newglossaryentry{#3\csname the#4\endcsname}{%
          name={\unexpanded{##2}%
            \IfBooleanF{##1}{\IfValueT{##4}{ [\noexpand\texttt{#3:##4}]}}},%
          sort={\csname the#2\endcsname},%
          description={p.~\noexpand\pageref{#3:\@relbl}},%
          user1={\csname the#2\endcsname},%
          user2={#5},%
          user3={\thesection},%
          user4={\expandafter\unexpanded\expandafter{\@savedSecTitle}}%
        }%
      }%
      \protected@write\@auxout{}{\string\glsliveentry{#3\csname the#4\endcsname}}%
      \protected@write\@auxout{}{\string\@thmtitle{\@tmplbl}{\unexpanded{##2}}}%
    \fi
  }%
  \expandafter\newcommand\csname ee#1\endcsname{\end{#2}}%
}
\@mkresult{L}{lemma}{l}{lemcounter}{Lemma}
% REMARK / COROLLARY / PROPOSITION -- same signature as \beL above
\@mkresult{R}{remark}{r}{remcounter}{Remark}
\@mkresult{C}{corollary}{c}{corcounter}{Corollary}
\@mkresult{P}{proposition}{p}{propcounter}{Proposition}
% OPEN-PROBLEM COMMAND -- generated with the same signature as \beT:
%   \beO[long title][short title][key]   ...   \eeO
% label o:KEY, or o:<counter> when no key is given; a bare \beO makes
% neither a label nor a glossary entry.
\@mkresult{O}{open}{o}{opcounter}{Problem}


% Section-level title tracker: updated only by \section, not \subsection.
% \@currentlabelname is overwritten by every \refstepcounter (including
% \subsection), so we save it separately here.
\newcommand{\@currentsectionname}{}
\AtBeginDocument{%
  \let\@epidcrn@orig@refstepcounter\refstepcounter
  \renewcommand{\refstepcounter}[1]{%
    \@epidcrn@orig@refstepcounter{#1}%
    \ifstrequal{#1}{section}{\global\let\@currentsectionname\@currentlabelname}{}%
  }%
}

% TABLE COMMANDS (glossaried) -- consistent naming with \beT, \beD, etc.
\newcommand{\beTa}[1][]{%
  \def\@temparg{#1}%
  \ifx\@temparg\@empty
    \begin{table}[htbp]%
  \else
    \let\@savedSecTitle\@currentsectionname%
    \refstepcounter{tabcounter}%
    \begin{table}[htbp]%
    \label{A:\thetabcounter}%
    \expanded{%
      \noexpand\newglossaryentry{A\thetabcounter}{%
        name={\unexpanded{#1}},%
        sort={\thetabcounter},%
        description={p.~\noexpand\pageref{A:\thetabcounter}},%
        user1={\thetabcounter},%
        user2={Table},%
        user3={\thesection},%
        user4={\expandafter\unexpanded\expandafter{\@savedSecTitle}}%
      }%
    }%
    \protected@write\@auxout{}{\string\glsliveentry{A\thetabcounter}}%
  \fi
}
\newcommand{\eeTa}{\end{table}}

% FIGURE COMMANDS (glossaried) -- consistent with \beT, \beR, \beTa, etc.
% Usage:
%   \beF[Glossary title]...\caption{Visible caption}\eeF  -> glossaried figure
%   \beF...\caption{Visible caption}\eeF                  -> plain [H] figure, no glossary
\newcommand{\beF}[1][]{%
  \def\@temparg{#1}%
  \ifx\@temparg\@empty
    \begin{figure}[H]%
  \else
    \begin{figure}[H]%
    \let\@savedSecTitle\@currentsectionname%
    \refstepcounter{figcounter}%
    \label{fig:\thefigcounter}%
    \expanded{%
      \noexpand\newglossaryentry{Fig\thefigcounter}{%
        name={\unexpanded{#1}},%
        sort={\thefigcounter},%
        description={p.~\noexpand\pageref{fig:\thefigcounter}},%
        user1={\thefigcounter},%
        user2={Figure},%
        user3={\thesection},%
        user4={\expandafter\unexpanded\expandafter{\@savedSecTitle}}%
      }%
    }%
    \protected@write\@auxout{}{\string\glsliveentry{Fig\thefigcounter}}%
  \fi
}
\newcommand{\eeF}{\end{figure}}
\makeatother
% LIVE-ENTRY REGISTER
% \glsaddall prints every entry the glossaries package currently knows, and
% that set includes entries restored from \jobname.glsdefs.  The package
% rewrites .glsdefs from all *currently defined* entries, so anything it
% restored is written straight back out: a deleted section kept reappearing in
% the catalogue for ever, and the .glsdefs had to be erased by hand.  The
% wrappers therefore record every entry they actually create in the .aux, and
% \printboth adds only those.  The .aux is rebuilt from scratch on each run, so
% this is self-healing: a removed result is gone after one further run.
\makeatletter
\newcommand{\@glslive}{}
\newcommand{\@glsaddlive}[1]{\ifglsentryexists{#1}{\glsadd{#1}}{}}
\newcommand{\glsliveentry}[1]{\g@addto@macro\@glslive{\@glsaddlive{#1}}}
\makeatother
% PRINT COMMAND
% No \clearpage here -- the catalogue continues on the current page.  This has
% been re-added by mistake several times; leave it out.
\makeatletter
\newcommand{\printboth}{%
  \@glslive   % add exactly the entries this document created last run
  \printglossary[style=unifiedstyle, title={0 \hspace{.22em} Definitions, Theorems, Examples and Open Problems}]%
}
\makeatother
% Named theorem reference: \tref{word} -> Thm~N(word-or-short)
% Shows stored short title if available, otherwise shows the key itself.
\newcommand{\tref}[1]{%
  Thm~\ref{t:#1}(\ifcsname @tshort@#1\endcsname\csname @tshort@#1\endcsname\else#1\fi)%
}

% Full-title theorem reference: \nref{t:KEY} prints the number and then the
% theorem's long title, so one writes "Theorem~\nref{t:blocks}" and gets
% "Theorem 3 ``Block factorization and siphons''".  Takes the complete label, like
% \ref, not the bare key like \tref.
%
% \beT writes each title to the .aux (see \@thmtitle below), so \nref resolves
% forward references from the second LaTeX pass, exactly as \ref does.  If the
% title is not yet known -- first pass, or a label that is not a \beT theorem
% -- it degrades to the number rather than producing ??.
% These two use @ in their names, and this point in the file is outside the
% \makeatletter region opened earlier (it closes at \makeatother above), so
% the region has to be reopened here.
\makeatletter
\newcommand{\@thmtitle}[2]{\global\@namedef{@tlong@#1}{#2}}
\newcommand{\nref}[1]{%
  \ref{#1}%
  \ifcsname @tlong@#1\endcsname
    \ ``\csname @tlong@#1\endcsname''%
  \fi
}
% \nm{t:KEY}: the stored long title of any \beT/\beL/\beP/\beR/\beC/\beD/\beE result,
% verbatim and without quotes, so item names follow the macro's first argument.
\newcommand{\nm}[1]{%
  \ifcsname @tlong@#1\endcsname
    \csname @tlong@#1\endcsname%
  \else
    \ref{#1}%
  \fi
}
\makeatother




%% ========================================
%% Other THEOREM ENVIRONMENTS
%% ========================================
\makeatletter
\@ifundefined{lemma}      {\newtheorem{lemma}      {Lemma}}      {}
\@ifundefined{proposition}{\newtheorem{proposition}{Proposition}} {}
\@ifundefined{corollary}  {\newtheorem{corollary}  {Corollary}}  {}
\@ifundefined{conjecture} {\newtheorem{conjecture} {Conjecture}}  {}
\@ifundefined{remark}     {\newtheorem{remark}     {Remark}}     {}
\@ifundefined{question}   {\newtheorem{question}   {Question}}   {}
\@ifundefined{assumption} {\newtheorem{assumption} {Assumption}} {}
\@ifundefined{hypothesis} {\newtheorem{hypothesis} {Hypothesis}} {}
\@ifundefined{fact}       {\newtheorem{fact}       {Fact}}       {}
\makeatother

% \citer: cite the relay-graph threshold theorem of AAHrel. If a document has
% \input the chapter/file containing t:relayRF (so the label resolves), point
% straight at it; otherwise fall back to a bare citation of AAHrel. Shared here
% (not per-paper) so it resolves consistently across every book chapter.
% \citeR: same idea, for the first theorem of AAHRah (t:Pstar).
\makeatletter
\providecommand{\citer}{%
  \@ifundefined{r@t:relayRF}%
    {\cite{AAHrel}}%
    {Theorem~\ref{t:relayRF} of \cite{AAHrel}}%
}
\providecommand{\citeR}{%
  \@ifundefined{r@t:Pstar}%
    {\cite{AAHRah}}%
    {Theorem~\ref{t:Pstar} of \cite{AAHRah}}%
}
\makeatother


% Abbreviations
\def\beQ{\begin{question}}\def\eeQ{\end{question}}
\def\beA{\begin{assumption}}\def\eeA{\end{assumption}}


%% ========================================

%% ========================================
%% COLORS & CUSTOM COMMANDS
%% ========================================
\definecolor{funccolor}{RGB}{25,25,112}
\definecolor{desccolor}{RGB}{64,64,64}
\newcommand{\func}[1]{\texttt{\textcolor{funccolor}{#1}}}
\def\eq{=}
\def\mm{-}

%% ========================================
%% MATHCAL SHORTHANDS
%% ========================================
% Project-wide convention: \mX for a capital X is always \mathcal X.
% Defined for the whole alphabet so the convention never fails to expand.
\newcommand{\mA}{\mathcal A}\newcommand{\mB}{\mathcal B}
\newcommand{\mC}{\mathcal C}\newcommand{\mD}{\mathcal D}
\newcommand{\mE}{\mathcal E}\newcommand{\mF}{\mathcal F}
\newcommand{\mG}{\mathcal G}\newcommand{\mH}{\mathcal H}
\newcommand{\mI}{\mathcal I}\newcommand{\mJ}{\mathcal J}
\newcommand{\mK}{\mathcal K}\newcommand{\mL}{\mathcal L}
\newcommand{\mM}{\mathcal M}\newcommand{\mN}{\mathcal N}
\newcommand{\mO}{\mathcal O}\newcommand{\mP}{\mathcal P}
\newcommand{\mQ}{\mathcal Q}\newcommand{\mR}{\mathcal R}
\newcommand{\mS}{\mathcal S}\newcommand{\mT}{\mathcal T}
\newcommand{\mU}{\mathcal U}\newcommand{\mV}{\mathcal V}
\newcommand{\mW}{\mathcal W}\newcommand{\mX}{\mathcal X}
\newcommand{\mY}{\mathcal Y}\newcommand{\mZ}{\mathcal Z}

%% ========================================
%% BASIC MATH ENVIRONMENTS
%% ========================================
\def\bep{\begin{pmatrix}}\def\eep{\end{pmatrix}}
\def\bev{\begin{vmatrix}}\def\eev{\end{vmatrix}}
\def\bea{\begin{eqnarray*}}\def\eea{\end{eqnarray*}}
\def\bc{\begin{cases}}\def\ec{\end{cases}}
\def\BEN{\begin{enumerate}}\def\EEN{\end{enumerate}}
\def\BI{\begin{itemize}}\def\EI{\end{itemize}}

\newcommand{\be}[1]{\begin{equation}\label{#1}}
\newcommand{\ee}{\end{equation}}
\newcommand{\beq}{\begin{eqnarray}}
\def\eeq{\end{eqnarray}}
\def\eqr{\eqref}\def\fr{\frac}\def\lbl{\label}
\def\Lra{\Longrightarrow}\def\non{\nonumber}\def\cNGM{\cite{Diek,Van,Van08}}
\def\Eq{\Leftrightarrow}\def\b{\beta}\def\uph{\red{UP TO HERE}}
%% ========================================
%% PROBABILITY & STATISTICS
%% ========================================
\renewcommand{\P}{\ensuremath{\mathbb{P}}}  % Note: overrides paragraph symbol
\newcommand{\E}{\ensuremath{\mathbb{E}}}
\newcommand{\ind}[1]{\ensuremath{\mathbbm{1}_{\left\{#1\right\}}}}
\DeclareMathOperator{\Cov}{Cov}
\DeclareMathOperator{\TV}{TV}
\newcommand{\rank}{\operatorname{rank}}

%% ========================================
%% CRN OPERATORS (organizations/ERCs/buffering/siphons framework)
%% ========================================
\def\clos{\mathrm{clos}}
\def\Supply{\mathrm{Supply}}
\def\req{\mathrm{req}}
\def\React{\mathrm{React}}
\def\Prod{\mathrm{Prod}}
\def\Gen{\mathrm{Gen}}
\def\In{\mathrm{In}}
\def\Out{\mathrm{Out}}
\def\Act{\mathrm{Act}}
\def\Min{\mathrm{Min}}

%% ========================================
%% LINE BREAKING
%% ========================================
% Last-resort stretch, used only on paragraphs TeX cannot otherwise break:
% a long \texttt{} token, a comma-separated inline math list (commas are not
% breakpoints in math), or a run of \cite/\ref boxes.  Without it those set as
% overfull \hbox and print into the margin; with it they set slightly loose,
% which is the lesser evil.  Does not affect paragraphs that already break.
\emergencystretch=3em

%% ========================================
%% NUMBER SETS
%% ========================================
\newcommand{\R}{\mathbb{R}}  % Real numbers
\newcommand{\C}{\mathbb{C}}  % Complex numbers
\newcommand{\N}{\mathbb{N}}
\newcommand{\Z}{\mathbb{Z}}
% \providecommand, not \newcommand: some papers define these in their own
% preamble after \input{preamble}, and must not clash.
\providecommand{\Q}{\mathbb{Q}}
\providecommand{\res}{\operatorname{res}}   % resultant; used by epiRes.tex
\newcommand{\MN}{\mathbb{N}}
\newcommand{\MZ}{\mathbb{Z}}
\newcommand{\realpos}{\mathbb{R}_{> 0}}
\newcommand{\realnneg}{\mathbb{R}_{\geq 0}}
\newcommand{\naturalpos}{\mathbb{N}_{>0}}
\newcommand{\Rplus}{\mathbb{R}_+}
\newcommand{\Rp}{\mathbb{R}_{\ge 0}}
\def\Rnne{\mathbb{R}_{\ge 0}}
\def\Se{\bar{S}} \def\Ie{\bar{I}}
\newcommand{\Rpos}{\mathbb{R}_{>0}}
\newcommand{\dd}{\,\mathrm d}
\newcommand{\e}{\mathrm e}
%% ========================================
%% GREEK LETTERS (2+ character abbreviations) greek
%% ========================================
\def\al{\alpha}\def\ba{\bff \alpha}
\def\w{\omega}\def\Lgn{\ga+\mu+\de}\def\om{\omega}
\def\de{\delta}\def\De{\Delta}\def\bde{\bff \delta}
\def\ep{\epsilon}\def\eps{\varepsilon}
\def\f{\varphi}  % Kept as requested
\def\ga{\gamma}\def\Ga{\Gamma}\def\G{\Gamma}\def\bg{\bff \gamma}
\def\up{\upsilon}  % NOT the breve accent: no accent commands anywhere.
                   % Named \up, not \u, so the LaTeX breve stays reachable.
\def\la{\lambda}\def\La{\Lambda}
\def\si{\sigma}\def\ka{\kappa}
\def\th{\theta}\def\vt{\vartheta}
\def\ze{\zeta}\def\Gap{\Gamma^{+}} \def\Gam{\Gamma^{-}}


%% ========================================
%% SPECIAL SHORT COMMANDS (kept as requested)
%% ========================================
\def\I{\infty}
\def\T{\widetilde}
\def\H{\hat}
\newcommand{\relint}{\operatorname{relint}}

%% ========================================
%% BOLD/VECTOR NOTATION
%% ========================================
\newcommand{\bff}[1]{{\mbox{\boldmath$#1$}}}
\newcommand\X{\boldsymbol{X}}
\def\x{\boldsymbol{x}}\def\bx{\bff x}\def\xe{\bar{\x}}
\newcommand\y{\boldsymbol{y}}\def\by{\boldsymbol{y}}\def\ye{\bar{\y}}
\def\yb{\y}\def\xb{\x}
\def\xy{\x-\y}
\newcommand\br{\boldsymbol{r}}\def\re{\bar{\r}}
\newcommand\bd{\boldsymbol{d}}
\newcommand\bQ{\boldsymbol{Q}}
\newcommand{\Yb}{{\mathbf{Y}}}
\newcommand{\Ab}{\mathbf{A}}
\newcommand{\zb}{\mathbf{z}}\def\bz{\bff 0}
\newcommand{\Zb}{\mathbf{Z}}
\newcommand{\bb}{\mathbf{b}}
\def\vbe{\vec {\bff \beta}}
\def\vb{\vec {\bff b}}
\def\v1{\vec {\bff 1}}
\def\vn{\vec \de}
\def\vz{\vec 0}
\def\vi{\overset{\rightarrow}{\mathsf i}}
\def\vv{\overset{\rightarrow}{v}}
\def\vr{\; \vec {\mathsf r}}
\newcommand{\va}{\bff a}
\providecommand{\vect}[1]{\boldsymbol{#1}}   % bold vector (moved from Gradostat.tex)
\providecommand{\one}{\mathbf 1}             % all-ones vector (moved from AliSF)
\providecommand{\spr}{s}                     % spectral bound  (moved from AliSF)
\providecommand{\Ke}{K_{\rm eff}}            % effective carrying capacity
% Critical right/left eigenvectors in the Hopf projection method (AAHOSN/New.tex)
\providecommand{\bq}{\mathbf{q}}
\providecommand{\bp}{\mathbf{p}}


%% ========================================
%% MATHCAL & MATHSCR
%% ========================================
% Mathcal
\newcommand\mc[1]{\mathcal{#1}}
\def\mC{{\mathcal C}}\def\mD{\mathcal D}\def\mE{\mathcal{E}}
\def\mF{\mathcal F}\def\mG{{\mathcal G}}\def\mH{{\mathcal H}}
\def\mI{\mathcal I}\def\mL{\mathcal{L}}\def\mM{{\mathcal M}}
\def\mN{{\mathcal N}}\def\mP{\mathcal{P}}\def\mR{{\mathcal R}}
\def\mS{{\mathcal S}}\def\mT{\mathcal{T}}\def\mV{{\mathcal V}}
\def\mK{\mathcal K}\def\mA{{\mathcal A}}\def\mU{{\mathcal U}}
% Symmetric bilinear and trilinear forms of the Taylor expansion, used in the
% first Lyapunov coefficient (AAHOSN/New.tex)
\providecommand{\Bform}{\mathcal{B}}
\providecommand{\Cform}{\mathcal{C}}

% Blackboard bold.  Kept here, not in the drivers: a \providecommand in a
% driver preamble runs before this file and silently wins, so the same name
% ends up meaning different things in different builds.
\def\IC{\mathbb{C}}\def\IN{\mathbb{N}}\def\IQ{\mathbb{Q}}
\def\IR{\mathbb{R}}\def\IZ{\mathbb{Z}}\def\IG{\mathbb{G}}
\def\II{\mathbb{I}}

% Mathscr
\newcommand{\Rs}{{\mathscr R}}
\newcommand{\Ss}{{\mathscr S}}
\newcommand{\Ns}{{\mathscr N}}
\newcommand{\Cs}{\mathscr C}


%% ========================================
%% OPERATORS & FUNCTIONS
%% ========================================
\DeclareMathOperator{\tr}{tr}
\DeclareMathOperator{\diag}{Diag}
\DeclareMathOperator{\Diag}{Diag}
\DeclareMathOperator{\intr}{\mathrm{Int}}
\DeclareMathOperator{\Inv}{Inv}
\DeclareMathOperator{\spec}{spec}
\DeclareMathOperator{\AUC}{AUC}               % area under the curve
\newcommand{\ii}{\mathrm{i}}                  % imaginary unit
\newcommand\op{\operatorname}
\newcommand{\supp}{\op{supp}}
\newcommand{\diff}{\mathop{}\mathopen{}\mathrm{d}}
\newcommand{\pdd}[1]{\frac{\partial}{\partial{#1}}}
\DeclareMathOperator{\cone}{cone}
\DeclareMathOperator{\proj}{proj}

%% ========================================
%% BRACKETS & NORMS
%% ========================================
\providecommand{\pr}[1]{\left(#1\right)}
\providecommand{\pp}[1]{\left[#1\right]}
\newcommand\croc[1]{\left\langle #1\right\rangle}
\newcommand{\norm}[1]{\|#1\|}
\newcommand{\inner}[2]{\langle #1, #2 \rangle}
% Short synonym of \inner, used by AAHOSN/New.tex as $\ip{u}{v}=\bar u^\top v$
\providecommand{\ip}[2]{\langle #1 , #2 \rangle}


%% ========================================
%% FORMATTING & UTILITIES
%% ========================================
\newcommand*{\Scale}[2][4]{\scalebox{#1}{$#2$}}
\newcommand{\te}[1]{\text{#1}}
\newcommand\steq[1]{\stackrel{\text{\rm #1.}}{=}}
\newcommand\blue[1]{\textcolor{blue}{#1}}
\newcommand{\red}{\textcolor[rgb]{1.00,0.00,0.00}}
\newcommand{\green}{\textcolor[rgb]{0.00,1.00,0.00}}

% Figure command
\newcommand{\figu}[3]{%
\begin{figure}[H]
\centering
\includegraphics[scale=#3]{#1}
\caption{#2\label{f:#1}}
\end{figure}
}

% Footnote
\long\def\symbolfootnote[#1]#2{%
\begingroup
\def\thefootnote{\fnsymbol{footnote}}\footnote[#1]{#2}%
\endgroup}
\def\fn{\symbolfootnote}

\newcommand{\dsum}{\displaystyle\sum}

%% ========================================
%% CRN & EPIDEMIOLOGY SPECIFIC
%% ========================================
% CRN notation
\def\SS{\mathcal S}
\def\CC{\mathcal C}
\def\RR{\mathcal R}

% Reproduction numbers
\def\Rto{R_{1}^{E_2}}
\def\Rot{R_{2}^{E_1}}
\def\Roc{R_{1,c}}
\def\Rtc{R_{2,c}}
\def\rd{r_{0}}
\def\sd{s_{0}}
\def\dd{D_{0}}
\def\vd{v_{0}}
\def\zd{z_{0}}
\def\se{\sE}
\def\xd{x_{0}}
%% ========================================
%% SPECIAL VALUES & EQUILIBRIA
%% ========================================
\def\eE{\bar{e}}
\def\wE{\bar{w}}
\def\sE{s_{*}}
\def\Me{\T M_{ee}}
\def\Rnne{\mathbb{R}_{\ge 0}}

\newcommand{\ones}{\mathbf 1}


%% ========================================
%% ABBREVIATIONS (alphabetical by category)
%% ========================================

% -- General mathematical terms --
\def\adm{admissible}\def\admy{admissibility}
%\def\and{antisymmetric}% disabled 2026-09-30: it overrode LaTeX's author separator \and (printed "antisymmetric" between authors)
\def\ass{assumption}
\def\cadlag{c\`adl\`ag}
\def\ch{characteristic polynomial}
\def\chp{characteristic polynomial}
\def\coe{coefficient}
\def\comp{computation}
\def\con{condition}
\def\coo{coordinate}
\def\corr{corresponding}
\def\dec{decomposition}
\def\den{denominator}
\def\dete{deterministic}
\def\eig{eigenvalue}\def\eval{eigenvalue}
\def\eigv{eigenvector}\def\eiv{eigenvector}\def\evec{eigenvector}
\def\elt{element}
\def\enn{essentially nonnegative/positivity preserving}
\def\equ{equation}
\def\etal{et al.}
\def\fe{for example}
\def\fno{from now on}
\def\for{\forall}
\def\form{formula}
\def\Fr{Furthermore, }
\def\frt{furthermore}
\def\fun{function}
\def\gene{generalization}
\def\heu{heuristic}
\def\how{however}\def\How{However, }
\def\ie{i.e. }
\def\im{\item}
\def\isd{is defined by}
\def\itc{it is easy to check that}
\def\itf{it follows that}\def\Itf{It follows that}
\def\ith{it holds that}
\def\itm{it may be checked that}\def\Itm{It may be checked that}
\def\ito{it turns out that}\def\Ito{It turns out that}\def\Itou{It turns out useful to}
\def\ma{mass action}\def\MAR{mass action representation}
\def\mat{matrix}
\def\mbc{may be checked to be}
\def\mbw{may be written as}\def\mubw{must be written as}
\def\mpr{more precisely, }\def\Mp{More precisely, }
\def\nc{non-catalytic}
\def\nne{non-negative}\def\nny{nonnegativity}
\def\no{\nonumber}
\def\noni{non-infectious}
\def\nota{notation}
\def\np{\newpage}
\def\Nt{Note that}
\def\nt{non trivial}\def\nz{non zero}
\def\oth{otherwise}
\def\Oth{On the other hand}
\def\pa{parameter}\def\para{parameter}
\def\pol{polynomial}
\def\prob{probability}
\def\Prf{{\bf  Proof:}}\def\prf{{\bf  Proof:}}
\def\prop{proportional to}
\def\qu{\quad}
\def\QED{\hfill {$\square$}\goodbreak \medskip}
\def\rep{representation}
\def\resp{respectively}
\def\satd{satisfied}\def\satg{satisfying}\def\sats{satisfies}\def\saty{satisfy}
\def\ssec{\subsection}\def\sssec{\subsubsection}
\def\sef{see for example}
\def\sksy{skew-symmetric}
\def\sub{substrate}
\def\sys{system}
\def\var{variable}
\def\wk{well-known}
\def\wl{without loss of generality}\def\wlo{w.l.o.g.}
\def\wmc{We may check that}\def\Wmc{We may check that}
\def\wmr{we may rewrite}
\def\wms{we must show that}
\def\wrt{with respect to}

% -- CRN & Chemistry --
\newcommand\CRN{chemical reaction network}
\newcommand\CRNs{chemical reaction networks}
\newcommand\RN{reaction network}
\def\ACR{absolute concentration robustness}
\def\AMSD{acyclic siphon decomposition}
\def\DAE{differential algebraic equations}
\def\FHJ{Feinberg-Horn-Jackson}
\def\GMAK{generalized mass action kinetics}
\def\MM{Michaelis-Menten}
\def\sm{stoichiometric matrix}\def\Sm{stoichiometric matrix}
\def\WRZD{weakly reversible, zero deficiency}\def\aut{autocatalytic}

% -- Epidemiology --
\def\brn{basic reproduction number}\def\bzn{basic replacement number}
\def\com{compartment}
\def\dfe{disease free equilibrium}\def\DFE{disease free equilibrium}
\def\dwt{dwell times}\def\edt{expected dwell times}
\def\EM{epidemic model}
\def\endp{endemic fixed point}\def\enp{endemic point}\def\EE{endemic point}
\def\ME{mathematical epidemiology}
\def\MEE{eco-epidemiology}\def\EEP{eco-epidemiological}
\def\ece{eco-epidemiological model}
\def\mgm{multi group model}
\def\mtS{multi-strain model}\def\mSt{multi-strain model}
\def\mSi{minimal siphon}
\def\NGM{next generation matrix}
\def\nim{new infections matrix}
\def\qem{quadratic epidemiological model}
\def\repF{invasion functions}
\def\SPF{SIR-PH-FA model}
\def\sus{susceptible}

% -- Graph theory & Linear algebra --
\def\jac{Jacobi matrix}
\def\jin{Jacobian matrix of the invasion vector field with respect to the invasion variables}
\def\jrn{Jacobian matrix of the resident vector field with respect to the resident variables}
\def\jxn{Jacobian matrix of the $\x$ vector field with respect to the $\x$ variables}
\def\Lm{Laplacian matrix}
\def\MTT{matrix tree theorem}
\def\SGT{spectral graph theory}
\def\SMW{Sherman–Morrison–Woodbury}

% -- Stability & Dynamics --
\def\absC{absorbtion CTMC}
\def\BTb{Bogdanov-Takens bifurcation}
\def\CEP{competitive exclusion partition}
\def\CTMC{continuous time Markov chain}
\def\ds{differential system}
\def\fp{fixed point}\def\bfp{boundary fixed point}\def\fbp{fixed boundary point}
\def\fbl{fixed boundary line}\def\fif{forward invariant face}
\def\gas{global asymptotic stability}\def\GAS{globally asymptotically stable}
\def\GASS{\text{GASS}}
\def\Hb{Hopf bifurcation}
\def\Hur{Hurwitz}\def\MH{Metzler-Hurwitz}\def\RH{Routh-Hurwitz}
\def\IS{\text{IS}}
\def\LAS{locally asymptotically stable}
\def\LdS{Lyapunov diagonally stable}
\def\Lf{Lyapunov function}\def\LyV{Lyapunov Volterra type function}\def\RLF{robust Lyapunov function}
\def\LyM{Lyapunov Malkin stable}
\def\PF{Perron-Frobenius}
\def\spfp{strictly positive fixed point}
\def\sta{stability}
\def\STP{sharp threshold property}\def\STT{sharp threshold theorem}

% -- Models (with \varphi explicitly instead of \f) --
\def\AB{$(A,B)$ Arino-Brauer epidemic models}
\def\ABa{$(A,\bbe,\ba,\varphi)$ balanced bilinear model}
\def\ABP{%$(A,\bbe,P,\varphi)$
\\ balanced bilinear model}
\def\AMG{Anderson May Gupta}
\newcommand\BIN{biological interaction network}
\def\Dt{Descartes type}
\def\gLV{generalized Lotka-Volterra model}
\def\IG{\text{invasion graph}}\def\IGMS{interaction graph on minimal siphons}
\def\LCP{\text{linear complementarity problem}}
\def\LV{Lotka-Volterra}\def\LVK{Lotka-Volterra-Kolmogorov}

% -- Splitting & Structure --
\def\reg{regular splitting}\def\regS{{\bf regular splitting}}
\def\RUR{rational univariate representation}
\def\bds{block-diagonal structure}\def\NGMs{next generation matrices}

% -- Miscellaneous --
\def\Ma{{Mathematica}}
\def\Mm{Max–Min}
\def\num{numerator}
\def\rap{rational approximation} \def\DSR{directed species-reaction bipartite graph}\def\PCP{polynomial complementarity problem}\def\rat{rational substitution}
\def\QSS{quasi steady state} \def\RHS{right hand side} \def\LHS{left hand side}
\def\OSN{online social network}\def\VdD{Van den Driessche}
\def\wli{$\omega$-limit} \def\rk{rich kinetics}\def\sJ{symbolic Jacobian}
\def\cV{\cite{Vas,VasSta,BSV,GSV}} \def\spf{semi-parametric semi-functional}
\def\CRW{Capasso-Ruan-Wang} \def\ABP{$(A,\bbe,P,\f)$ balanced bilinear model}
\def\ABa{%$(A,\bbe,\ba,\f)$ \\
rank-one balanced bilinear model}
\def\bbe{B%{\bff \beta}
} \def\vbe{\vec {\bff \beta}
}
\def\vb{\vec {\bff b}}
\newcommand{\Kh}{\hat{K}}          % K with hat, first mention only
\newcommand{\wS}{w^{\sigma}}       % right Perron J^\perp, full block
\newcommand{\piS}{\vec \pi^{\sigma}}    % left Perron J^\perp, full block
\newcommand{\ws}{w^{\sigma}}       % right Perron J^\perp, sub-block
\newcommand{\pis}{\pi^{\sigma}}    % left Perron J^\perp, sub-block
\newcommand{\hwS}{\hat{w}^{\sigma}}   % right Perron K, full block
\newcommand{\hpiS}{\hat{\pi}^{\sigma}} % left Perron K, full block
\newcommand{\hws}{\hat{w}^{\sigma}}   % right Perron K, sub-block
\newcommand{\hpis}{\hat{\pi}^{\sigma}} % left Perron K, sub-block
\newcommand{\twS}{\tilde{w}^{\sigma}}  % right Perron \tilde K, full block
\newcommand{\tpiS}{\tilde{\pi}^{\sigma}} % left Perron \tilde K
\newcommand{\tws}{\tilde{w}^{\sigma}}
\newcommand{\tpis}{\tilde{\pi}^{\sigma}}
\newcommand{\Mat}[2]{\mathbb{R}_{\ge 0}^{\,#1\times #2}}
\def\tJ{J_\sigma^\perp} \def\lPe{left Perron eigenvector}
% \G is \Gamma (defined at :795).  The Volterra function keeps its own name,
% so that \G no longer has two conflicting meanings in this file.
\def\GV{G}  % Volterra G(u)=u-1-\ln u
\def\PVL{Perron-Volterra Lyapunov minimal siphon ansatz}
\newcommand{\istar}[1]{i_{#1}^{*}}
%\newcommand{\sE}{s_E}
\newcommand{\vE}{v_E}
\def\PV{Perron-Volterra candidate Lyapunov function}
\def\GLV{generalized Lotka-Volterra}
\def\MERN{mathematical epidemiology reaction networks}
\def\CMM{contraction metric method}
\newcommand{\sign}{\operatorname{sign}}
\newcommand{\Var}{\operatorname{Var}} 

% ======================================================================
% IDEAS SHARED BETWEEN THE ABSTRACT AND THE MAIN RESULTS SECTION
%
% Write each idea ONCE, in the abstract; the Main Results section recalls it,
% so the two can never drift apart.
%
%   in the abstract        \ideaitem{1}{The first-order reaction terms ...}
%   in Main Results        \ideahead{1}
%
% \ideaitem stores the text globally and typesets
%   \item \hyperlink{idea:1}{<text>}
% \ideahead typesets
%   \item\hypertarget{idea:1}{}\textbf{<the same text>}
%
% The abstract must come before the Main Results input, which is the normal
% order. If a heading has to read differently from the abstract sentence (for
% instance a long abstract item summarised in one line), override it with the
% optional argument: \ideahead[short heading]{6}.
% ======================================================================
\makeatletter
\newcommand{\ideaitem}[2]{%
  \expandafter\gdef\csname @idea@#1\endcsname{#2}%
  \item\hyperlink{idea:#1}{#2}}
\newcommand{\ideahead}[2][]{%
  \item\hypertarget{idea:#2}{}%
  \def\@ideatmp{#1}%
  \ifx\@ideatmp\@empty
    \@ifundefined{@idea@#2}{%
      \PackageWarning{preamble}{Idea #2 is used in Main Results but was never
        defined by \string\ideaitem\space in the abstract}%
      \textbf{[idea #2 undefined]}%
    }{\textbf{\begingroup\def\ideaterm##1##2{##2}\def\idearest##1{}\@nameuse{@idea@#2}\endgroup}}%
  \else
    \textbf{#1}%
  \fi}
% \ideaterm{<hypertarget>}{<words>}: inside an abstract idea the words become an
% invisible link (the document sets \hypersetup{hidelinks}) jumping to the place
% in Main Results that defines them; inside the Main Results heading itself,
% where the target sits a few lines below, \ideahead strips the link and prints
% the words plainly. Mark the destination with \hypertarget{<name>}{} on the line
% before the display, which adds no equation number and shows nothing.
\newcommand{\ideaterm}[2]{\hyperlink{#1}{#2}}
% \idearest{<text>}: the part of an abstract idea after its first sentence; printed in
% the abstract, omitted from the Main Results heading.
\newcommand{\idearest}[1]{#1}
\makeatother

% ======================================================================
% SHORT-PAPER SWITCH
% \shorttrue in a master that assembles only part of the companion paper
% (Barrat/SipArxiv.tex). \nsh{...} drops text that cites results living in the
% omitted sections, so the short build has no dangling references while the
% companion paper keeps the text verbatim.
% ======================================================================
\newif\ifshort \shortfalse
\newcommand{\nsh}[1]{\ifshort\else#1\fi}
% \shortalt{<short build>}{<companion paper>}: use when a sentence needs
% different wording in the two builds, not merely a deletion.
\newcommand{\shortalt}[2]{\ifshort#1\else#2\fi}
% \Jtr, \JtrS{\Sigma}: transverse Jacobian. J^\perp here; SipArxiv.tex redefines them as J_{\mathrm{tra}}, J_{\mathrm{tra},\Sigma}
\newcommand{\Jtr}{J^\perp}
\newcommand{\JtrS}[1]{J_{#1}^\perp}
.
\makeatletter
\AtBeginDocument{\@ifundefined{hypersetup}{}{\hypersetup{linktoc=all}}}
\makeatother
\usepackage{graphicx,float,fancybox}
\usepackage{epsfig,epstopdf,tikz}
\usepackage[sans]{dsfont}
\usetikzlibrary{positioning,arrows,arrows.meta,calc}
% Configure listings for Mathematica code
\lstdefinestyle{mathematica}{
 language=Mathematica,
 basicstyle=\small\ttfamily,
 backgroundcolor=\color{gray!10},
 frame=single,
 breaklines=true,
 commentstyle=\color{green!60!black},
 keywordstyle=\color{blue},
 stringstyle=\color{red},
 showstringspaces=false,
 tabsize=2
}
\usepackage{tcolorbox}
% Colored boxes for key results
\newtcolorbox{keyresult}[1][]{
 colback=blue!5!white,
 colframe=blue!75!black,
 title=#1,
 fonttitle=\bfseries
}
\newtcolorbox{examplebox}[1][]{
 colback=green!5!white,
 colframe=green!75!black,
 title=#1,
 fonttitle=\bfseries
}


\makeatletter
\@ifundefined{theorem}   {\newtheorem{theorem}   {Theorem}}   {}
\@ifundefined{definition}{\newtheorem{definition}{Definition}} {}
\@ifundefined{open}      {\newtheorem{open}      {Problem}}   {}
\@ifundefined{example}   {\newtheorem{example}   {Example}}   {}
\makeatother

% COUNTER DEFINITIONS
\newcounter{tabcounter}
\newcounter{defcounter}
\newcounter{opcounter}
\newcounter{thmcounter}
\newcounter{excounter}
\newcounter{ascounter}
\newcounter{cjcounter}
\newcounter{lemcounter}
\newcounter{remcounter}
\newcounter{corcounter}
\newcounter{propcounter}
\newcounter{figcounter}

% UNIFIED GLOSSARY STYLE
\newglossarystyle{unifiedstyle}{%
  \setglossarystyle{list}%
  \renewcommand*{\glossarypreamble}{\gdef\glslastsec{}}%
  \renewcommand{\glossentry}[2]{%
    \edef\glscursec{\glsentryuseriii{##1}}%
    \ifx\glscursec\glslastsec\else
      \item[]\textbf{\S\,\glscursec\enspace\glsentryuseriv{##1}}%
      \global\let\glslastsec\glscursec
    \fi
    \item[\glsentryuserii{##1} \glsentryuseri{##1}:] ``\glossentryname{##1}'' \glsentrydesc{##1}%
  }%
}

\makeatletter
% First-word extractor (used by \beT, \beD and \beXa for the default key)
% \@firstwordof hello world \relax  ->  hello
\def\@firstwordof#1 #2\relax{#1}
% First three letters, used for the heading mnemonic of a theorem given no
% explicit key.  Padding with \@empty keeps it safe for first words shorter
% than three letters: \@empty vanishes under \edef, whereas \relax would
% survive into the label and is illegal inside \csname (it broke every
% two-letter title, e.g. "ME ...", with "Missing \endcsname inserted").
\def\@takethree#1#2#3#4\@nil{#1#2#3}

% EXAMPLE COMMAND -- same signature and behaviour as \beT and \beD.
% Usage: \beXa*[long title][short title][key]
%   star  = suppresses the printed "[key]" tag after the title (as for \beT*
%   and \beD*); the label and glossary entry are unaffected.
%   key = optional label key override; the label becomes e:KEY, and without a
%   key it is e: plus the first three letters of the title's first word.  A
%   mnemonic label lets an example be cited as "Example~\ref{e:KEY}" instead
%   of by section number, which is what the citing prose needs.
\NewDocumentCommand{\beXa}{s O{} O{} o}{%
  \def\@temparg{#2}%
  \ifx\@temparg\@empty
    \begin{example}%
  \else
    \let\@savedSecTitle\@currentsectionname%
    \refstepcounter{excounter}%
    \IfValueTF{#4}{%
      \edef\@exlblkey{#4}%
    }{%
      \edef\@exfirstword{\@firstwordof#2 \relax}%
      \edef\@exlblkey{%
        \expandafter\@takethree\@exfirstword\@empty\@empty\@empty\@nil}%
    }%
    \def\@tempshort{#3}%
    \ifx\@tempshort\@empty
      \expandafter\xdef\csname @xshort@\@exlblkey\endcsname{\@exlblkey}%
    \else
      \expandafter\def\csname @xshort@\@exlblkey\endcsname{#3}%
    \fi
    \IfBooleanTF{#1}{\let\@extag\@empty}{%
      \edef\@extag{\noexpand\ \noexpand\textup{[\noexpand\texttt{e:\@exlblkey}]}}}%
    \begin{example}[#2\ifshowkeys\@extag\fi]%
    \label{e:\@exlblkey}%
    % record the long title so \nref{e:KEY} prints it, as for theorems
    \protected@write\@auxout{}{%
      \string\@thmtitle{e:\@exlblkey}{\unexpanded{#2}}}%
    \expanded{%
      \noexpand\newglossaryentry{e\theexcounter}{%
        name={\unexpanded{#2}%
          \IfBooleanF{#1}{ [\noexpand\texttt{e:\@exlblkey}]}},%
        sort={\theexample},%
        description={p.~\noexpand\pageref{e:\@exlblkey}},%
        user1={\theexample},%
        user2={Example},%
        user3={\thesection},%
        user4={\expandafter\unexpanded\expandafter{\@savedSecTitle}}%
      }%
    }%
    \protected@write\@auxout{}{\string\glsliveentry{e\theexcounter}}%
  \fi
}
\newcommand{\eeXa}{\end{example}}

% DEFINITION COMMAND -- same signature and behaviour as \beT below.
% Usage: \beD[long title][short title][key]
%   long title  = definition heading (plain text only)
%   short title = shown in \dref{word}; defaults to the key
%   key         = optional label key override; the label becomes d:KEY
% Without a key the label is d: plus the first three letters of the title's
% first word, so it is a mnemonic rather than the definition's number: a
% numeric label would silently change whenever a definition is inserted.
\NewDocumentCommand{\beD}{s O{} O{} o}{%
  \def\@temparg{#2}%
  \ifx\@temparg\@empty
    \begin{definition}%
  \else
    \let\@savedSecTitle\@currentsectionname%
    \refstepcounter{defcounter}%
    \IfValueTF{#4}{%
      \edef\@deflblkey{#4}%
    }{%
      \edef\@deffirstword{\@firstwordof#2 \relax}%
      \edef\@deflblkey{%
        \expandafter\@takethree\@deffirstword\@empty\@empty\@empty\@nil}%
    }%
    \def\@tempshort{#3}%
    \ifx\@tempshort\@empty
      \expandafter\xdef\csname @dshort@\@deflblkey\endcsname{\@deflblkey}%
    \else
      \expandafter\def\csname @dshort@\@deflblkey\endcsname{#3}%
    \fi
    \IfBooleanTF{#1}{\let\@deftag\@empty}{%
      \edef\@deftag{\noexpand\ \noexpand\textup{[\noexpand\texttt{d:\@deflblkey}]}}}%
    \begin{definition}[#2\ifshowkeys\@deftag\fi]%
    \label{d:\@deflblkey}%
    % record the long title so \nref{d:KEY} prints it, as for theorems
    \protected@write\@auxout{}{%
      \string\@thmtitle{d:\@deflblkey}{\unexpanded{#2}}}%
    \expanded{%
      \noexpand\newglossaryentry{d\thedefcounter}{%
        name={\unexpanded{#2}%
          \IfBooleanF{#1}{ [\noexpand\texttt{d:\@deflblkey}]}},%
        sort={\thedefinition},%
        description={p.~\noexpand\pageref{d:\@deflblkey}},%
        user1={\thedefinition},%
        user2={Definition},%
        user3={\thesection},%
        user4={\expandafter\unexpanded\expandafter{\@savedSecTitle}}%
      }%
    }%
    \protected@write\@auxout{}{\string\glsliveentry{d\thedefcounter}}%
  \fi
}
\newcommand{\eeD}{\end{definition}}

% \showkeysfalse in a driver hides the [t:KEY] mnemonics in result headings
% (labels and glossary entries are unchanged); used for submission builds.
\newif\ifshowkeys \showkeystrue
% \commonredtrue in a driver that inputs basicFiles/commonRed: tool.tex then
% cites its reduction results instead of restating them.
\newif\ifcommonred \commonredfalse
% TheorEM COMMAND
% Usage: \beT[long title][short title][key]
%   long title  = theorem heading (plain text only)
%   short title = shown in \tref{word}; may contain math; defaults to first word of long title
%   key         = optional label key override; label becomes t:KEY (avoids first-word collisions)
% Label is t:WORD where WORD = first word of long title (or KEY if given); use \tref{WORD} to reference.
\NewDocumentCommand{\beT}{s O{} O{} o}{%
  \def\@temparg{#2}%
  \ifx\@temparg\@empty
    \begin{theorem}%
  \else
    \let\@savedSecTitle\@currentsectionname%
    \refstepcounter{thmcounter}%
    % \@thmtag is the mnemonic shown in the heading; every theorem gets one.
    % With an explicit key the key itself is shown.  Without one the key
    % defaults to the title's first word, and showing that whole word would
    % merely repeat the title ("Dissipativity [Dissipativity]"), so only its
    % first three letters are shown: "Dissipativity [Dis]".  The label is
    % unaffected -- it still uses the full \@thmlblkey.
    \IfValueTF{#4}{%
      \edef\@thmlblkey{#4}%
    }{%
      \edef\@thmfirstword{\@firstwordof#2 \relax}%
      \edef\@thmlblkey{%
        \expandafter\@takethree\@thmfirstword\@empty\@empty\@empty\@nil}%
    }%
    \IfBooleanTF{#1}{\let\@thmtag\@empty}{%
      \edef\@thmtag{\noexpand\ \noexpand\textup{[\noexpand\texttt{t:\@thmlblkey}]}}}%
    \def\@tempshort{#3}%
    \ifx\@tempshort\@empty
      \expandafter\xdef\csname @tshort@\@thmlblkey\endcsname{\@thmlblkey}%
    \else
      \expandafter\def\csname @tshort@\@thmlblkey\endcsname{#3}%
    \fi
    \begin{theorem}[#2\ifshowkeys\@thmtag\fi]%
    \label{t:\@thmlblkey}%
    % record the long title for \nref; via the .aux so forward refs resolve
    \protected@write\@auxout{}{%
      \string\@thmtitle{t:\@thmlblkey}{\unexpanded{#2}}}%
    \expanded{%
      \noexpand\newglossaryentry{t\thethmcounter}{%
        name={\unexpanded{#2}%
          \IfBooleanF{#1}{ [\noexpand\texttt{t:\@thmlblkey}]}},%
        % sort/user1 must be the number the document actually prints, i.e. the
        % theorem environment's counter, not thmcounter.  thmcounter advances
        % only on \beT with a title, so a bare \beT (or \beT with an
        % empty title) makes the two drift and the glossary cites a wrong
        % number.  thmcounter is still fine as the entry key: it is unique.
        sort={\thetheorem},%
        description={p.~\noexpand\pageref{t:\@thmlblkey}},%
        user1={\thetheorem},%
        user2={Theorem},%
        user3={\thesection},%
        user4={\expandafter\unexpanded\expandafter{\@savedSecTitle}}%
      }%
    }%
    \protected@write\@auxout{}{\string\glsliveentry{t\thethmcounter}}%
  \fi
}
\newcommand{\eeT}{\end{theorem}}

% CROSS-PART RELATION PARAGRAPH
% Usage: \xpart{text}  -- place immediately BEFORE \beT, never inside it
%   (anything inside \beT...\eeT would be swallowed into the glossary entry).
% The paragraph is typeset only when Part I is actually present, i.e. in the
% combined book build; in the standalone builds (GranadaDSABNS/A.tex,
% AAHRah-Submission/AAHRah.tex) the label part:glv is undefined and the whole
% paragraph disappears rather than leaving ?? cross-references. Same
% graceful-degradation principle as \GASref in GASt.tex. Resolves on the
% second LaTeX pass, since r@part:glv comes from the .aux file.
% (Already inside the \makeatletter region opened above, so no \makeatletter
% pair here -- adding one would close that region early.)
\newcommand{\xpart}[1]{\@ifundefined{r@part:glv}{}{\paragraph{Relation with Parts~I--II.}#1}}

% \bookref{label}{precise text}{fallback text}: prints the precise cross-reference
% when the label resolves (combined book build), and the fallback text otherwise
% (standalone paper builds, where book-only appendix labels are absent). Same
% graceful-degradation role as \GASref in GASt.tex, but usable in any file since it
% lives in this preamble (loaded by every driver) -- unlike \GASref, which is defined in a
% content file and is therefore undefined in standalone builds that omit that file.
\newcommand{\bookref}[3]{\@ifundefined{r@#1}{#3}{#2}}

% \PVpart and \RBpart: the two parts that a partial build may omit.
% They print "Part~N" in any build that contains the part, and a naming phrase
% in a build that does not, so an excerpt shows no ?? .  Each already supplies
% the word "Part", so write "\PVpart develops ...", not "Part~\PVpart".
\newcommand{\PVpart}{\bookref{part:aahrah}%
  {Part~\ref{part:aahrah}}{the Perron--Volterra part of the full book}}
\newcommand{\RBpart}{\bookref{part:recbrowne}%
  {Part~\ref{part:recbrowne}}{the immuno-virology part of the full book}}

% OPEN PROBLEM COMMANDS
% ASSUMPTION COMMANDS (glossaried)
\newcommand{\beAs}[1][]{%
  \def\@temparg{#1}%
  \ifx\@temparg\@empty
    \begin{assumption}%
  \else
    \let\@savedSecTitle\@currentsectionname%
    \refstepcounter{ascounter}%
    \begin{assumption}[#1]%
    \label{as:\theascounter}%
    \expanded{%
      \noexpand\newglossaryentry{as\theascounter}{%
        name={\unexpanded{#1}},%
        sort={\theascounter},%
        description={p.~\noexpand\pageref{as:\theascounter}},%
        user1={\theascounter},%
        user2={Assumption},%
        user3={\thesection},%
        user4={\expandafter\unexpanded\expandafter{\@savedSecTitle}}%
      }%
    }%
    \protected@write\@auxout{}{\string\glsliveentry{as\theascounter}}%
  \fi
}
\newcommand{\eeAs}{\end{assumption}}
% CONJECTURE COMMANDS (glossaried)
\newcommand{\beCj}[1][]{%
  \def\@temparg{#1}%
  \ifx\@temparg\@empty
    \begin{conjecture}%
  \else
    \let\@savedSecTitle\@currentsectionname%
    \refstepcounter{cjcounter}%
    \begin{conjecture}[#1]%
    \label{cj:\thecjcounter}%
    \expanded{%
      \noexpand\newglossaryentry{cj\thecjcounter}{%
        name={\unexpanded{#1}},%
        sort={\thecjcounter},%
        description={p.~\noexpand\pageref{cj:\thecjcounter}},%
        user1={\thecjcounter},%
        user2={Conjecture},%
        user3={\thesection},%
        user4={\expandafter\unexpanded\expandafter{\@savedSecTitle}}%
      }%
    }%
    \protected@write\@auxout{}{\string\glsliveentry{cj\thecjcounter}}%
  \fi
}
\newcommand{\eeCj}{\end{conjecture}}
% LEMMA COMMANDS
% ---------------------------------------------------------------------
% UNIFORM RESULT WRAPPERS -- \beL, \beC, \beP, \beR share \beT's signature
%
%   \beX[long title][short title][key]  ...  \eeX
%
%   long title  = heading (math allowed)
%   short title = shown in the cross-reference; defaults to the key
%   key         = label key.  With a key the label is X:KEY and the heading
%                 carries the mnemonic [key], exactly as for \beT.  Without
%                 one the label falls back to the running counter, which is
%                 what every pre-existing call site in the book relies on.
%
% A bare \beX -- arguments commented out -- opens the plain environment and
% creates NO label and NO glossary entry.  That is the supported way to keep
% an individual result out of the catalogue.
% ---------------------------------------------------------------------
% #1 wrapper name  #2 environment  #3 label prefix  #4 counter  #5 type name
\newcommand{\@mkresult}[5]{%
  % A leading * suppresses the [key] mnemonic in both the heading and the
  % glossary name, while still making the label X:KEY and the entry:
  %     \beX*[long][short][key]
  \expandafter\NewDocumentCommand\csname be#1\endcsname{s O{} O{} o}{%
    \def\@temparg{##2}%
    \ifx\@temparg\@empty
      \begin{#2}%
    \else
      \let\@savedSecTitle\@currentsectionname%
      \refstepcounter{#4}%
      % With a key: mnemonic label X:KEY and the [key] tag in the heading.
      % Without one: fall back to the running counter, exactly as these
      % wrappers behaved before, so every legacy one-argument call site in
      % the book keeps its label.  The counter is also csname-safe, which a
      % title's first word is not -- titles beginning with math would
      % otherwise build \csname \Omega...\endcsname and fail.
      \IfValueTF{##4}{%
        \edef\@relbl{##4}%
        \IfBooleanTF{##1}{\let\@retag\@empty}{%
          \edef\@retag{\noexpand\ \noexpand\textup{[\noexpand\texttt{#3:##4}]}}}%
      }{%
        \edef\@relbl{\arabic{#4}}%
        \let\@retag\@empty
      }%
      \def\@tempshort{##3}%
      \ifx\@tempshort\@empty
        \expandafter\xdef\csname @sshort@#3\@relbl\endcsname{\@relbl}%
      \else
        \expandafter\def\csname @sshort@#3\@relbl\endcsname{##3}%
      \fi
      \begin{#2}[##2\ifshowkeys\@retag\fi]%
      \edef\@tmplbl{#3:\@relbl}\expandafter\label\expandafter{\@tmplbl}%
      \expanded{%
        \noexpand\newglossaryentry{#3\csname the#4\endcsname}{%
          name={\unexpanded{##2}%
            \IfBooleanF{##1}{\IfValueT{##4}{ [\noexpand\texttt{#3:##4}]}}},%
          sort={\csname the#2\endcsname},%
          description={p.~\noexpand\pageref{#3:\@relbl}},%
          user1={\csname the#2\endcsname},%
          user2={#5},%
          user3={\thesection},%
          user4={\expandafter\unexpanded\expandafter{\@savedSecTitle}}%
        }%
      }%
      \protected@write\@auxout{}{\string\glsliveentry{#3\csname the#4\endcsname}}%
      \protected@write\@auxout{}{\string\@thmtitle{\@tmplbl}{\unexpanded{##2}}}%
    \fi
  }%
  \expandafter\newcommand\csname ee#1\endcsname{\end{#2}}%
}
\@mkresult{L}{lemma}{l}{lemcounter}{Lemma}
% REMARK / COROLLARY / PROPOSITION -- same signature as \beL above
\@mkresult{R}{remark}{r}{remcounter}{Remark}
\@mkresult{C}{corollary}{c}{corcounter}{Corollary}
\@mkresult{P}{proposition}{p}{propcounter}{Proposition}
% OPEN-PROBLEM COMMAND -- generated with the same signature as \beT:
%   \beO[long title][short title][key]   ...   \eeO
% label o:KEY, or o:<counter> when no key is given; a bare \beO makes
% neither a label nor a glossary entry.
\@mkresult{O}{open}{o}{opcounter}{Problem}


% Section-level title tracker: updated only by \section, not \subsection.
% \@currentlabelname is overwritten by every \refstepcounter (including
% \subsection), so we save it separately here.
\newcommand{\@currentsectionname}{}
\AtBeginDocument{%
  \let\@epidcrn@orig@refstepcounter\refstepcounter
  \renewcommand{\refstepcounter}[1]{%
    \@epidcrn@orig@refstepcounter{#1}%
    \ifstrequal{#1}{section}{\global\let\@currentsectionname\@currentlabelname}{}%
  }%
}

% TABLE COMMANDS (glossaried) -- consistent naming with \beT, \beD, etc.
\newcommand{\beTa}[1][]{%
  \def\@temparg{#1}%
  \ifx\@temparg\@empty
    \begin{table}[htbp]%
  \else
    \let\@savedSecTitle\@currentsectionname%
    \refstepcounter{tabcounter}%
    \begin{table}[htbp]%
    \label{A:\thetabcounter}%
    \expanded{%
      \noexpand\newglossaryentry{A\thetabcounter}{%
        name={\unexpanded{#1}},%
        sort={\thetabcounter},%
        description={p.~\noexpand\pageref{A:\thetabcounter}},%
        user1={\thetabcounter},%
        user2={Table},%
        user3={\thesection},%
        user4={\expandafter\unexpanded\expandafter{\@savedSecTitle}}%
      }%
    }%
    \protected@write\@auxout{}{\string\glsliveentry{A\thetabcounter}}%
  \fi
}
\newcommand{\eeTa}{\end{table}}

% FIGURE COMMANDS (glossaried) -- consistent with \beT, \beR, \beTa, etc.
% Usage:
%   \beF[Glossary title]...\caption{Visible caption}\eeF  -> glossaried figure
%   \beF...\caption{Visible caption}\eeF                  -> plain [H] figure, no glossary
\newcommand{\beF}[1][]{%
  \def\@temparg{#1}%
  \ifx\@temparg\@empty
    \begin{figure}[H]%
  \else
    \begin{figure}[H]%
    \let\@savedSecTitle\@currentsectionname%
    \refstepcounter{figcounter}%
    \label{fig:\thefigcounter}%
    \expanded{%
      \noexpand\newglossaryentry{Fig\thefigcounter}{%
        name={\unexpanded{#1}},%
        sort={\thefigcounter},%
        description={p.~\noexpand\pageref{fig:\thefigcounter}},%
        user1={\thefigcounter},%
        user2={Figure},%
        user3={\thesection},%
        user4={\expandafter\unexpanded\expandafter{\@savedSecTitle}}%
      }%
    }%
    \protected@write\@auxout{}{\string\glsliveentry{Fig\thefigcounter}}%
  \fi
}
\newcommand{\eeF}{\end{figure}}
\makeatother
% LIVE-ENTRY REGISTER
% \glsaddall prints every entry the glossaries package currently knows, and
% that set includes entries restored from \jobname.glsdefs.  The package
% rewrites .glsdefs from all *currently defined* entries, so anything it
% restored is written straight back out: a deleted section kept reappearing in
% the catalogue for ever, and the .glsdefs had to be erased by hand.  The
% wrappers therefore record every entry they actually create in the .aux, and
% \printboth adds only those.  The .aux is rebuilt from scratch on each run, so
% this is self-healing: a removed result is gone after one further run.
\makeatletter
\newcommand{\@glslive}{}
\newcommand{\@glsaddlive}[1]{\ifglsentryexists{#1}{\glsadd{#1}}{}}
\newcommand{\glsliveentry}[1]{\g@addto@macro\@glslive{\@glsaddlive{#1}}}
\makeatother
% PRINT COMMAND
% No \clearpage here -- the catalogue continues on the current page.  This has
% been re-added by mistake several times; leave it out.
\makeatletter
\newcommand{\printboth}{%
  \@glslive   % add exactly the entries this document created last run
  \printglossary[style=unifiedstyle, title={0 \hspace{.22em} Definitions, Theorems, Examples and Open Problems}]%
}
\makeatother
% Named theorem reference: \tref{word} -> Thm~N(word-or-short)
% Shows stored short title if available, otherwise shows the key itself.
\newcommand{\tref}[1]{%
  Thm~\ref{t:#1}(\ifcsname @tshort@#1\endcsname\csname @tshort@#1\endcsname\else#1\fi)%
}

% Full-title theorem reference: \nref{t:KEY} prints the number and then the
% theorem's long title, so one writes "Theorem~\nref{t:blocks}" and gets
% "Theorem 3 ``Block factorization and siphons''".  Takes the complete label, like
% \ref, not the bare key like \tref.
%
% \beT writes each title to the .aux (see \@thmtitle below), so \nref resolves
% forward references from the second LaTeX pass, exactly as \ref does.  If the
% title is not yet known -- first pass, or a label that is not a \beT theorem
% -- it degrades to the number rather than producing ??.
% These two use @ in their names, and this point in the file is outside the
% \makeatletter region opened earlier (it closes at \makeatother above), so
% the region has to be reopened here.
\makeatletter
\newcommand{\@thmtitle}[2]{\global\@namedef{@tlong@#1}{#2}}
\newcommand{\nref}[1]{%
  \ref{#1}%
  \ifcsname @tlong@#1\endcsname
    \ ``\csname @tlong@#1\endcsname''%
  \fi
}
% \nm{t:KEY}: the stored long title of any \beT/\beL/\beP/\beR/\beC/\beD/\beE result,
% verbatim and without quotes, so item names follow the macro's first argument.
\newcommand{\nm}[1]{%
  \ifcsname @tlong@#1\endcsname
    \csname @tlong@#1\endcsname%
  \else
    \ref{#1}%
  \fi
}
\makeatother




%% ========================================
%% Other THEOREM ENVIRONMENTS
%% ========================================
\makeatletter
\@ifundefined{lemma}      {\newtheorem{lemma}      {Lemma}}      {}
\@ifundefined{proposition}{\newtheorem{proposition}{Proposition}} {}
\@ifundefined{corollary}  {\newtheorem{corollary}  {Corollary}}  {}
\@ifundefined{conjecture} {\newtheorem{conjecture} {Conjecture}}  {}
\@ifundefined{remark}     {\newtheorem{remark}     {Remark}}     {}
\@ifundefined{question}   {\newtheorem{question}   {Question}}   {}
\@ifundefined{assumption} {\newtheorem{assumption} {Assumption}} {}
\@ifundefined{hypothesis} {\newtheorem{hypothesis} {Hypothesis}} {}
\@ifundefined{fact}       {\newtheorem{fact}       {Fact}}       {}
\makeatother

% \citer: cite the relay-graph threshold theorem of AAHrel. If a document has
% \input the chapter/file containing t:relayRF (so the label resolves), point
% straight at it; otherwise fall back to a bare citation of AAHrel. Shared here
% (not per-paper) so it resolves consistently across every book chapter.
% \citeR: same idea, for the first theorem of AAHRah (t:Pstar).
\makeatletter
\providecommand{\citer}{%
  \@ifundefined{r@t:relayRF}%
    {\cite{AAHrel}}%
    {Theorem~\ref{t:relayRF} of \cite{AAHrel}}%
}
\providecommand{\citeR}{%
  \@ifundefined{r@t:Pstar}%
    {\cite{AAHRah}}%
    {Theorem~\ref{t:Pstar} of \cite{AAHRah}}%
}
\makeatother


% Abbreviations
\def\beQ{\begin{question}}\def\eeQ{\end{question}}
\def\beA{\begin{assumption}}\def\eeA{\end{assumption}}


%% ========================================

%% ========================================
%% COLORS & CUSTOM COMMANDS
%% ========================================
\definecolor{funccolor}{RGB}{25,25,112}
\definecolor{desccolor}{RGB}{64,64,64}
\newcommand{\func}[1]{\texttt{\textcolor{funccolor}{#1}}}
\def\eq{=}
\def\mm{-}

%% ========================================
%% MATHCAL SHORTHANDS
%% ========================================
% Project-wide convention: \mX for a capital X is always \mathcal X.
% Defined for the whole alphabet so the convention never fails to expand.
\newcommand{\mA}{\mathcal A}\newcommand{\mB}{\mathcal B}
\newcommand{\mC}{\mathcal C}\newcommand{\mD}{\mathcal D}
\newcommand{\mE}{\mathcal E}\newcommand{\mF}{\mathcal F}
\newcommand{\mG}{\mathcal G}\newcommand{\mH}{\mathcal H}
\newcommand{\mI}{\mathcal I}\newcommand{\mJ}{\mathcal J}
\newcommand{\mK}{\mathcal K}\newcommand{\mL}{\mathcal L}
\newcommand{\mM}{\mathcal M}\newcommand{\mN}{\mathcal N}
\newcommand{\mO}{\mathcal O}\newcommand{\mP}{\mathcal P}
\newcommand{\mQ}{\mathcal Q}\newcommand{\mR}{\mathcal R}
\newcommand{\mS}{\mathcal S}\newcommand{\mT}{\mathcal T}
\newcommand{\mU}{\mathcal U}\newcommand{\mV}{\mathcal V}
\newcommand{\mW}{\mathcal W}\newcommand{\mX}{\mathcal X}
\newcommand{\mY}{\mathcal Y}\newcommand{\mZ}{\mathcal Z}

%% ========================================
%% BASIC MATH ENVIRONMENTS
%% ========================================
\def\bep{\begin{pmatrix}}\def\eep{\end{pmatrix}}
\def\bev{\begin{vmatrix}}\def\eev{\end{vmatrix}}
\def\bea{\begin{eqnarray*}}\def\eea{\end{eqnarray*}}
\def\bc{\begin{cases}}\def\ec{\end{cases}}
\def\BEN{\begin{enumerate}}\def\EEN{\end{enumerate}}
\def\BI{\begin{itemize}}\def\EI{\end{itemize}}

\newcommand{\be}[1]{\begin{equation}\label{#1}}
\newcommand{\ee}{\end{equation}}
\newcommand{\beq}{\begin{eqnarray}}
\def\eeq{\end{eqnarray}}
\def\eqr{\eqref}\def\fr{\frac}\def\lbl{\label}
\def\Lra{\Longrightarrow}\def\non{\nonumber}\def\cNGM{\cite{Diek,Van,Van08}}
\def\Eq{\Leftrightarrow}\def\b{\beta}\def\uph{\red{UP TO HERE}}
%% ========================================
%% PROBABILITY & STATISTICS
%% ========================================
\renewcommand{\P}{\ensuremath{\mathbb{P}}}  % Note: overrides paragraph symbol
\newcommand{\E}{\ensuremath{\mathbb{E}}}
\newcommand{\ind}[1]{\ensuremath{\mathbbm{1}_{\left\{#1\right\}}}}
\DeclareMathOperator{\Cov}{Cov}
\DeclareMathOperator{\TV}{TV}
\newcommand{\rank}{\operatorname{rank}}

%% ========================================
%% CRN OPERATORS (organizations/ERCs/buffering/siphons framework)
%% ========================================
\def\clos{\mathrm{clos}}
\def\Supply{\mathrm{Supply}}
\def\req{\mathrm{req}}
\def\React{\mathrm{React}}
\def\Prod{\mathrm{Prod}}
\def\Gen{\mathrm{Gen}}
\def\In{\mathrm{In}}
\def\Out{\mathrm{Out}}
\def\Act{\mathrm{Act}}
\def\Min{\mathrm{Min}}

%% ========================================
%% LINE BREAKING
%% ========================================
% Last-resort stretch, used only on paragraphs TeX cannot otherwise break:
% a long \texttt{} token, a comma-separated inline math list (commas are not
% breakpoints in math), or a run of \cite/\ref boxes.  Without it those set as
% overfull \hbox and print into the margin; with it they set slightly loose,
% which is the lesser evil.  Does not affect paragraphs that already break.
\emergencystretch=3em

%% ========================================
%% NUMBER SETS
%% ========================================
\newcommand{\R}{\mathbb{R}}  % Real numbers
\newcommand{\C}{\mathbb{C}}  % Complex numbers
\newcommand{\N}{\mathbb{N}}
\newcommand{\Z}{\mathbb{Z}}
% \providecommand, not \newcommand: some papers define these in their own
% preamble after % Shared preamble for the GLVbook papers.  This is the single canonical copy;
% a driver one level below GLVbook reaches it as
%
%     \input{../basicFiles/preamble}
%
% What stays in the driver, because it differs from paper to paper:
% \documentclass, font choice, \title, \author, \date, and any macro genuinely
% local to one paper.  Everything else belongs here.  The file is
% self-contained: it does not read any other file.

\usepackage[utf8]{inputenc}
\usepackage[novbox]{pdfsync}   % required: SyncTeX alone does not work here

% Page geometry, set once here for every paper that loads this file.  A driver
% that needs different margins says \geometry{...}; a second
% \usepackage[...]{geometry} raises an option clash and must not be used.
\makeatletter
\@ifundefined{AAdefNnogeometry}{\usepackage[a4paper,margin=2cm]{geometry}}{}
\makeatother

% mathematics
\usepackage{amsmath,amssymb,amsthm,mathtools,bm,mathrsfs,bbm}
\usepackage{centernot}

% floats, tables, lists
\usepackage{graphicx,float}
\makeatletter
\@ifpackageloaded{subfig}{}{\usepackage{subcaption}}
\makeatother
\usepackage{booktabs,longtable,threeparttable,array,multirow}
\usepackage{enumitem}

% colour, code, algorithms
\usepackage{xcolor,listings,tcolorbox}
\makeatletter
\@ifundefined{c@algorithm}{\usepackage{algorithm}}{}
\makeatother
\usepackage{algorithmicx,algpseudocode}

% diagrams
\usepackage{tikz,tikz-cd,pgfplots}
\usetikzlibrary{graphs,decorations.pathmorphing,decorations.markings,
arrows.meta,automata,positioning}
\tikzcdset{every label/.append style = {font = \small}}
\pgfplotsset{compat=1.18}

% chemistry
\usepackage[version=4]{mhchem}

\usepackage{microtype}
\usepackage{lineno}
\usepackage{nomencl}
\makenomenclature

% hyperref first, then bookmark, then glossaries: glossaries must follow
% hyperref or the entry links are built against the wrong anchors.
\makeatletter
\@ifundefined{hyperrefpackage}{%
  \@ifundefined{AAdefNnohyperref}{\usepackage{hyperref}}{}%
}{}
\@ifundefined{hyperrefpackage}{\usepackage{bookmark}}{}
\makeatother
\usepackage[toc,sort=use]{glossaries}
\makeglossaries

% Book/standalone switch.  Material that the book already covers in its
% prelude, but that a standalone paper still has to state, is wrapped as
%     \ifbookbuild\else ... \fi
% Drivers default to standalone; bookM.tex says \bookbuildtrue.
\newif\ifbookbuild
\bookbuildfalse

% Part switches for the gradostat paper.  Both default to true; a driver that
% wants the short submission says \madecfalse and/or \slowfastfalse, and every
% cross-reference into the dropped part disappears with it.
%   \ifmadec     Part III, the Lagasquie--Madec model and its break-even geometry
%   \ifslowfast  Part IV, everything carrying the small parameter \eps
\newif\ifmadec     \madectrue
\newif\ifslowfast  \slowfasttrue

% Part switch for the Barrat/simplicial-contagion project.  Defaults to
% true (the hypernetwork-synchrony part is included); the journal-paper
% driver says \hnetpartfalse to drop it and every cross-reference into it.
\newif\ifhnetpart  \hnetparttrue

%% ========================================
%% PACKAGES
%% ========================================
\RequirePackage{etoolbox}
\usepackage{makecell}
\usepackage[novbox]{pdfsync}% obsolete; SyncTeX is built into modern engines
\usepackage{amssymb, amsfonts,mathtools,bbm,eurosym,graphicx,latexsym, amsmath,amsthm,times,nccmath,url}
\usepackage{listings,verbatim,cases,tabularx}
% hyperref is deliberately NOT in the line above.  Journal classes often
% load it themselves, with a long option list, from a begindocument hook
% (mdpi.cls does: \DeclareHookRule{begindocument}{hyperref}{before}{mdpi}).
% Loading it here first, without options, makes that later load raise
% "Option clash for package hyperref".  So load it only if nobody else
% has, or will:
%   - \@ifpackageloaded  catches a class that loaded it already;
%   - \hyperrefpackage    is mdpi.cls's own deferred loader;
%   - \AAdefNnohyperref   lets any other driver opt out by defining it
%                         before \input{preamble}.
\makeatletter
\@ifpackageloaded{hyperref}{}{%
  \@ifundefined{hyperrefpackage}{%
    \@ifundefined{AAdefNnohyperref}{\usepackage{hyperref}}{}%
  }{}%
}
\makeatother
% linktoc=all makes a table-of-contents page number a link as well as the
% entry text; hyperref's default, linktoc=section, links only the text.
% Deferred to begin-document and guarded: by the block above hyperref may
% be loaded later (mdpi's own hook) or not at all (\AAdefNnohyperref), and
% \hypersetup must exist before it is called.  A driver that wants the
% default back can say \hypersetup{linktoc=section} after \input{preamble}.
\makeatletter
\AtBeginDocument{\@ifundefined{hypersetup}{}{\hypersetup{linktoc=all}}}
\makeatother
\usepackage{graphicx,float,fancybox}
\usepackage{epsfig,epstopdf,tikz}
\usepackage[sans]{dsfont}
\usetikzlibrary{positioning,arrows,arrows.meta,calc}
% Configure listings for Mathematica code
\lstdefinestyle{mathematica}{
 language=Mathematica,
 basicstyle=\small\ttfamily,
 backgroundcolor=\color{gray!10},
 frame=single,
 breaklines=true,
 commentstyle=\color{green!60!black},
 keywordstyle=\color{blue},
 stringstyle=\color{red},
 showstringspaces=false,
 tabsize=2
}
\usepackage{tcolorbox}
% Colored boxes for key results
\newtcolorbox{keyresult}[1][]{
 colback=blue!5!white,
 colframe=blue!75!black,
 title=#1,
 fonttitle=\bfseries
}
\newtcolorbox{examplebox}[1][]{
 colback=green!5!white,
 colframe=green!75!black,
 title=#1,
 fonttitle=\bfseries
}


\makeatletter
\@ifundefined{theorem}   {\newtheorem{theorem}   {Theorem}}   {}
\@ifundefined{definition}{\newtheorem{definition}{Definition}} {}
\@ifundefined{open}      {\newtheorem{open}      {Problem}}   {}
\@ifundefined{example}   {\newtheorem{example}   {Example}}   {}
\makeatother

% COUNTER DEFINITIONS
\newcounter{tabcounter}
\newcounter{defcounter}
\newcounter{opcounter}
\newcounter{thmcounter}
\newcounter{excounter}
\newcounter{ascounter}
\newcounter{cjcounter}
\newcounter{lemcounter}
\newcounter{remcounter}
\newcounter{corcounter}
\newcounter{propcounter}
\newcounter{figcounter}

% UNIFIED GLOSSARY STYLE
\newglossarystyle{unifiedstyle}{%
  \setglossarystyle{list}%
  \renewcommand*{\glossarypreamble}{\gdef\glslastsec{}}%
  \renewcommand{\glossentry}[2]{%
    \edef\glscursec{\glsentryuseriii{##1}}%
    \ifx\glscursec\glslastsec\else
      \item[]\textbf{\S\,\glscursec\enspace\glsentryuseriv{##1}}%
      \global\let\glslastsec\glscursec
    \fi
    \item[\glsentryuserii{##1} \glsentryuseri{##1}:] ``\glossentryname{##1}'' \glsentrydesc{##1}%
  }%
}

\makeatletter
% First-word extractor (used by \beT, \beD and \beXa for the default key)
% \@firstwordof hello world \relax  ->  hello
\def\@firstwordof#1 #2\relax{#1}
% First three letters, used for the heading mnemonic of a theorem given no
% explicit key.  Padding with \@empty keeps it safe for first words shorter
% than three letters: \@empty vanishes under \edef, whereas \relax would
% survive into the label and is illegal inside \csname (it broke every
% two-letter title, e.g. "ME ...", with "Missing \endcsname inserted").
\def\@takethree#1#2#3#4\@nil{#1#2#3}

% EXAMPLE COMMAND -- same signature and behaviour as \beT and \beD.
% Usage: \beXa*[long title][short title][key]
%   star  = suppresses the printed "[key]" tag after the title (as for \beT*
%   and \beD*); the label and glossary entry are unaffected.
%   key = optional label key override; the label becomes e:KEY, and without a
%   key it is e: plus the first three letters of the title's first word.  A
%   mnemonic label lets an example be cited as "Example~\ref{e:KEY}" instead
%   of by section number, which is what the citing prose needs.
\NewDocumentCommand{\beXa}{s O{} O{} o}{%
  \def\@temparg{#2}%
  \ifx\@temparg\@empty
    \begin{example}%
  \else
    \let\@savedSecTitle\@currentsectionname%
    \refstepcounter{excounter}%
    \IfValueTF{#4}{%
      \edef\@exlblkey{#4}%
    }{%
      \edef\@exfirstword{\@firstwordof#2 \relax}%
      \edef\@exlblkey{%
        \expandafter\@takethree\@exfirstword\@empty\@empty\@empty\@nil}%
    }%
    \def\@tempshort{#3}%
    \ifx\@tempshort\@empty
      \expandafter\xdef\csname @xshort@\@exlblkey\endcsname{\@exlblkey}%
    \else
      \expandafter\def\csname @xshort@\@exlblkey\endcsname{#3}%
    \fi
    \IfBooleanTF{#1}{\let\@extag\@empty}{%
      \edef\@extag{\noexpand\ \noexpand\textup{[\noexpand\texttt{e:\@exlblkey}]}}}%
    \begin{example}[#2\ifshowkeys\@extag\fi]%
    \label{e:\@exlblkey}%
    % record the long title so \nref{e:KEY} prints it, as for theorems
    \protected@write\@auxout{}{%
      \string\@thmtitle{e:\@exlblkey}{\unexpanded{#2}}}%
    \expanded{%
      \noexpand\newglossaryentry{e\theexcounter}{%
        name={\unexpanded{#2}%
          \IfBooleanF{#1}{ [\noexpand\texttt{e:\@exlblkey}]}},%
        sort={\theexample},%
        description={p.~\noexpand\pageref{e:\@exlblkey}},%
        user1={\theexample},%
        user2={Example},%
        user3={\thesection},%
        user4={\expandafter\unexpanded\expandafter{\@savedSecTitle}}%
      }%
    }%
    \protected@write\@auxout{}{\string\glsliveentry{e\theexcounter}}%
  \fi
}
\newcommand{\eeXa}{\end{example}}

% DEFINITION COMMAND -- same signature and behaviour as \beT below.
% Usage: \beD[long title][short title][key]
%   long title  = definition heading (plain text only)
%   short title = shown in \dref{word}; defaults to the key
%   key         = optional label key override; the label becomes d:KEY
% Without a key the label is d: plus the first three letters of the title's
% first word, so it is a mnemonic rather than the definition's number: a
% numeric label would silently change whenever a definition is inserted.
\NewDocumentCommand{\beD}{s O{} O{} o}{%
  \def\@temparg{#2}%
  \ifx\@temparg\@empty
    \begin{definition}%
  \else
    \let\@savedSecTitle\@currentsectionname%
    \refstepcounter{defcounter}%
    \IfValueTF{#4}{%
      \edef\@deflblkey{#4}%
    }{%
      \edef\@deffirstword{\@firstwordof#2 \relax}%
      \edef\@deflblkey{%
        \expandafter\@takethree\@deffirstword\@empty\@empty\@empty\@nil}%
    }%
    \def\@tempshort{#3}%
    \ifx\@tempshort\@empty
      \expandafter\xdef\csname @dshort@\@deflblkey\endcsname{\@deflblkey}%
    \else
      \expandafter\def\csname @dshort@\@deflblkey\endcsname{#3}%
    \fi
    \IfBooleanTF{#1}{\let\@deftag\@empty}{%
      \edef\@deftag{\noexpand\ \noexpand\textup{[\noexpand\texttt{d:\@deflblkey}]}}}%
    \begin{definition}[#2\ifshowkeys\@deftag\fi]%
    \label{d:\@deflblkey}%
    % record the long title so \nref{d:KEY} prints it, as for theorems
    \protected@write\@auxout{}{%
      \string\@thmtitle{d:\@deflblkey}{\unexpanded{#2}}}%
    \expanded{%
      \noexpand\newglossaryentry{d\thedefcounter}{%
        name={\unexpanded{#2}%
          \IfBooleanF{#1}{ [\noexpand\texttt{d:\@deflblkey}]}},%
        sort={\thedefinition},%
        description={p.~\noexpand\pageref{d:\@deflblkey}},%
        user1={\thedefinition},%
        user2={Definition},%
        user3={\thesection},%
        user4={\expandafter\unexpanded\expandafter{\@savedSecTitle}}%
      }%
    }%
    \protected@write\@auxout{}{\string\glsliveentry{d\thedefcounter}}%
  \fi
}
\newcommand{\eeD}{\end{definition}}

% \showkeysfalse in a driver hides the [t:KEY] mnemonics in result headings
% (labels and glossary entries are unchanged); used for submission builds.
\newif\ifshowkeys \showkeystrue
% \commonredtrue in a driver that inputs basicFiles/commonRed: tool.tex then
% cites its reduction results instead of restating them.
\newif\ifcommonred \commonredfalse
% TheorEM COMMAND
% Usage: \beT[long title][short title][key]
%   long title  = theorem heading (plain text only)
%   short title = shown in \tref{word}; may contain math; defaults to first word of long title
%   key         = optional label key override; label becomes t:KEY (avoids first-word collisions)
% Label is t:WORD where WORD = first word of long title (or KEY if given); use \tref{WORD} to reference.
\NewDocumentCommand{\beT}{s O{} O{} o}{%
  \def\@temparg{#2}%
  \ifx\@temparg\@empty
    \begin{theorem}%
  \else
    \let\@savedSecTitle\@currentsectionname%
    \refstepcounter{thmcounter}%
    % \@thmtag is the mnemonic shown in the heading; every theorem gets one.
    % With an explicit key the key itself is shown.  Without one the key
    % defaults to the title's first word, and showing that whole word would
    % merely repeat the title ("Dissipativity [Dissipativity]"), so only its
    % first three letters are shown: "Dissipativity [Dis]".  The label is
    % unaffected -- it still uses the full \@thmlblkey.
    \IfValueTF{#4}{%
      \edef\@thmlblkey{#4}%
    }{%
      \edef\@thmfirstword{\@firstwordof#2 \relax}%
      \edef\@thmlblkey{%
        \expandafter\@takethree\@thmfirstword\@empty\@empty\@empty\@nil}%
    }%
    \IfBooleanTF{#1}{\let\@thmtag\@empty}{%
      \edef\@thmtag{\noexpand\ \noexpand\textup{[\noexpand\texttt{t:\@thmlblkey}]}}}%
    \def\@tempshort{#3}%
    \ifx\@tempshort\@empty
      \expandafter\xdef\csname @tshort@\@thmlblkey\endcsname{\@thmlblkey}%
    \else
      \expandafter\def\csname @tshort@\@thmlblkey\endcsname{#3}%
    \fi
    \begin{theorem}[#2\ifshowkeys\@thmtag\fi]%
    \label{t:\@thmlblkey}%
    % record the long title for \nref; via the .aux so forward refs resolve
    \protected@write\@auxout{}{%
      \string\@thmtitle{t:\@thmlblkey}{\unexpanded{#2}}}%
    \expanded{%
      \noexpand\newglossaryentry{t\thethmcounter}{%
        name={\unexpanded{#2}%
          \IfBooleanF{#1}{ [\noexpand\texttt{t:\@thmlblkey}]}},%
        % sort/user1 must be the number the document actually prints, i.e. the
        % theorem environment's counter, not thmcounter.  thmcounter advances
        % only on \beT with a title, so a bare \beT (or \beT with an
        % empty title) makes the two drift and the glossary cites a wrong
        % number.  thmcounter is still fine as the entry key: it is unique.
        sort={\thetheorem},%
        description={p.~\noexpand\pageref{t:\@thmlblkey}},%
        user1={\thetheorem},%
        user2={Theorem},%
        user3={\thesection},%
        user4={\expandafter\unexpanded\expandafter{\@savedSecTitle}}%
      }%
    }%
    \protected@write\@auxout{}{\string\glsliveentry{t\thethmcounter}}%
  \fi
}
\newcommand{\eeT}{\end{theorem}}

% CROSS-PART RELATION PARAGRAPH
% Usage: \xpart{text}  -- place immediately BEFORE \beT, never inside it
%   (anything inside \beT...\eeT would be swallowed into the glossary entry).
% The paragraph is typeset only when Part I is actually present, i.e. in the
% combined book build; in the standalone builds (GranadaDSABNS/A.tex,
% AAHRah-Submission/AAHRah.tex) the label part:glv is undefined and the whole
% paragraph disappears rather than leaving ?? cross-references. Same
% graceful-degradation principle as \GASref in GASt.tex. Resolves on the
% second LaTeX pass, since r@part:glv comes from the .aux file.
% (Already inside the \makeatletter region opened above, so no \makeatletter
% pair here -- adding one would close that region early.)
\newcommand{\xpart}[1]{\@ifundefined{r@part:glv}{}{\paragraph{Relation with Parts~I--II.}#1}}

% \bookref{label}{precise text}{fallback text}: prints the precise cross-reference
% when the label resolves (combined book build), and the fallback text otherwise
% (standalone paper builds, where book-only appendix labels are absent). Same
% graceful-degradation role as \GASref in GASt.tex, but usable in any file since it
% lives in this preamble (loaded by every driver) -- unlike \GASref, which is defined in a
% content file and is therefore undefined in standalone builds that omit that file.
\newcommand{\bookref}[3]{\@ifundefined{r@#1}{#3}{#2}}

% \PVpart and \RBpart: the two parts that a partial build may omit.
% They print "Part~N" in any build that contains the part, and a naming phrase
% in a build that does not, so an excerpt shows no ?? .  Each already supplies
% the word "Part", so write "\PVpart develops ...", not "Part~\PVpart".
\newcommand{\PVpart}{\bookref{part:aahrah}%
  {Part~\ref{part:aahrah}}{the Perron--Volterra part of the full book}}
\newcommand{\RBpart}{\bookref{part:recbrowne}%
  {Part~\ref{part:recbrowne}}{the immuno-virology part of the full book}}

% OPEN PROBLEM COMMANDS
% ASSUMPTION COMMANDS (glossaried)
\newcommand{\beAs}[1][]{%
  \def\@temparg{#1}%
  \ifx\@temparg\@empty
    \begin{assumption}%
  \else
    \let\@savedSecTitle\@currentsectionname%
    \refstepcounter{ascounter}%
    \begin{assumption}[#1]%
    \label{as:\theascounter}%
    \expanded{%
      \noexpand\newglossaryentry{as\theascounter}{%
        name={\unexpanded{#1}},%
        sort={\theascounter},%
        description={p.~\noexpand\pageref{as:\theascounter}},%
        user1={\theascounter},%
        user2={Assumption},%
        user3={\thesection},%
        user4={\expandafter\unexpanded\expandafter{\@savedSecTitle}}%
      }%
    }%
    \protected@write\@auxout{}{\string\glsliveentry{as\theascounter}}%
  \fi
}
\newcommand{\eeAs}{\end{assumption}}
% CONJECTURE COMMANDS (glossaried)
\newcommand{\beCj}[1][]{%
  \def\@temparg{#1}%
  \ifx\@temparg\@empty
    \begin{conjecture}%
  \else
    \let\@savedSecTitle\@currentsectionname%
    \refstepcounter{cjcounter}%
    \begin{conjecture}[#1]%
    \label{cj:\thecjcounter}%
    \expanded{%
      \noexpand\newglossaryentry{cj\thecjcounter}{%
        name={\unexpanded{#1}},%
        sort={\thecjcounter},%
        description={p.~\noexpand\pageref{cj:\thecjcounter}},%
        user1={\thecjcounter},%
        user2={Conjecture},%
        user3={\thesection},%
        user4={\expandafter\unexpanded\expandafter{\@savedSecTitle}}%
      }%
    }%
    \protected@write\@auxout{}{\string\glsliveentry{cj\thecjcounter}}%
  \fi
}
\newcommand{\eeCj}{\end{conjecture}}
% LEMMA COMMANDS
% ---------------------------------------------------------------------
% UNIFORM RESULT WRAPPERS -- \beL, \beC, \beP, \beR share \beT's signature
%
%   \beX[long title][short title][key]  ...  \eeX
%
%   long title  = heading (math allowed)
%   short title = shown in the cross-reference; defaults to the key
%   key         = label key.  With a key the label is X:KEY and the heading
%                 carries the mnemonic [key], exactly as for \beT.  Without
%                 one the label falls back to the running counter, which is
%                 what every pre-existing call site in the book relies on.
%
% A bare \beX -- arguments commented out -- opens the plain environment and
% creates NO label and NO glossary entry.  That is the supported way to keep
% an individual result out of the catalogue.
% ---------------------------------------------------------------------
% #1 wrapper name  #2 environment  #3 label prefix  #4 counter  #5 type name
\newcommand{\@mkresult}[5]{%
  % A leading * suppresses the [key] mnemonic in both the heading and the
  % glossary name, while still making the label X:KEY and the entry:
  %     \beX*[long][short][key]
  \expandafter\NewDocumentCommand\csname be#1\endcsname{s O{} O{} o}{%
    \def\@temparg{##2}%
    \ifx\@temparg\@empty
      \begin{#2}%
    \else
      \let\@savedSecTitle\@currentsectionname%
      \refstepcounter{#4}%
      % With a key: mnemonic label X:KEY and the [key] tag in the heading.
      % Without one: fall back to the running counter, exactly as these
      % wrappers behaved before, so every legacy one-argument call site in
      % the book keeps its label.  The counter is also csname-safe, which a
      % title's first word is not -- titles beginning with math would
      % otherwise build \csname \Omega...\endcsname and fail.
      \IfValueTF{##4}{%
        \edef\@relbl{##4}%
        \IfBooleanTF{##1}{\let\@retag\@empty}{%
          \edef\@retag{\noexpand\ \noexpand\textup{[\noexpand\texttt{#3:##4}]}}}%
      }{%
        \edef\@relbl{\arabic{#4}}%
        \let\@retag\@empty
      }%
      \def\@tempshort{##3}%
      \ifx\@tempshort\@empty
        \expandafter\xdef\csname @sshort@#3\@relbl\endcsname{\@relbl}%
      \else
        \expandafter\def\csname @sshort@#3\@relbl\endcsname{##3}%
      \fi
      \begin{#2}[##2\ifshowkeys\@retag\fi]%
      \edef\@tmplbl{#3:\@relbl}\expandafter\label\expandafter{\@tmplbl}%
      \expanded{%
        \noexpand\newglossaryentry{#3\csname the#4\endcsname}{%
          name={\unexpanded{##2}%
            \IfBooleanF{##1}{\IfValueT{##4}{ [\noexpand\texttt{#3:##4}]}}},%
          sort={\csname the#2\endcsname},%
          description={p.~\noexpand\pageref{#3:\@relbl}},%
          user1={\csname the#2\endcsname},%
          user2={#5},%
          user3={\thesection},%
          user4={\expandafter\unexpanded\expandafter{\@savedSecTitle}}%
        }%
      }%
      \protected@write\@auxout{}{\string\glsliveentry{#3\csname the#4\endcsname}}%
      \protected@write\@auxout{}{\string\@thmtitle{\@tmplbl}{\unexpanded{##2}}}%
    \fi
  }%
  \expandafter\newcommand\csname ee#1\endcsname{\end{#2}}%
}
\@mkresult{L}{lemma}{l}{lemcounter}{Lemma}
% REMARK / COROLLARY / PROPOSITION -- same signature as \beL above
\@mkresult{R}{remark}{r}{remcounter}{Remark}
\@mkresult{C}{corollary}{c}{corcounter}{Corollary}
\@mkresult{P}{proposition}{p}{propcounter}{Proposition}
% OPEN-PROBLEM COMMAND -- generated with the same signature as \beT:
%   \beO[long title][short title][key]   ...   \eeO
% label o:KEY, or o:<counter> when no key is given; a bare \beO makes
% neither a label nor a glossary entry.
\@mkresult{O}{open}{o}{opcounter}{Problem}


% Section-level title tracker: updated only by \section, not \subsection.
% \@currentlabelname is overwritten by every \refstepcounter (including
% \subsection), so we save it separately here.
\newcommand{\@currentsectionname}{}
\AtBeginDocument{%
  \let\@epidcrn@orig@refstepcounter\refstepcounter
  \renewcommand{\refstepcounter}[1]{%
    \@epidcrn@orig@refstepcounter{#1}%
    \ifstrequal{#1}{section}{\global\let\@currentsectionname\@currentlabelname}{}%
  }%
}

% TABLE COMMANDS (glossaried) -- consistent naming with \beT, \beD, etc.
\newcommand{\beTa}[1][]{%
  \def\@temparg{#1}%
  \ifx\@temparg\@empty
    \begin{table}[htbp]%
  \else
    \let\@savedSecTitle\@currentsectionname%
    \refstepcounter{tabcounter}%
    \begin{table}[htbp]%
    \label{A:\thetabcounter}%
    \expanded{%
      \noexpand\newglossaryentry{A\thetabcounter}{%
        name={\unexpanded{#1}},%
        sort={\thetabcounter},%
        description={p.~\noexpand\pageref{A:\thetabcounter}},%
        user1={\thetabcounter},%
        user2={Table},%
        user3={\thesection},%
        user4={\expandafter\unexpanded\expandafter{\@savedSecTitle}}%
      }%
    }%
    \protected@write\@auxout{}{\string\glsliveentry{A\thetabcounter}}%
  \fi
}
\newcommand{\eeTa}{\end{table}}

% FIGURE COMMANDS (glossaried) -- consistent with \beT, \beR, \beTa, etc.
% Usage:
%   \beF[Glossary title]...\caption{Visible caption}\eeF  -> glossaried figure
%   \beF...\caption{Visible caption}\eeF                  -> plain [H] figure, no glossary
\newcommand{\beF}[1][]{%
  \def\@temparg{#1}%
  \ifx\@temparg\@empty
    \begin{figure}[H]%
  \else
    \begin{figure}[H]%
    \let\@savedSecTitle\@currentsectionname%
    \refstepcounter{figcounter}%
    \label{fig:\thefigcounter}%
    \expanded{%
      \noexpand\newglossaryentry{Fig\thefigcounter}{%
        name={\unexpanded{#1}},%
        sort={\thefigcounter},%
        description={p.~\noexpand\pageref{fig:\thefigcounter}},%
        user1={\thefigcounter},%
        user2={Figure},%
        user3={\thesection},%
        user4={\expandafter\unexpanded\expandafter{\@savedSecTitle}}%
      }%
    }%
    \protected@write\@auxout{}{\string\glsliveentry{Fig\thefigcounter}}%
  \fi
}
\newcommand{\eeF}{\end{figure}}
\makeatother
% LIVE-ENTRY REGISTER
% \glsaddall prints every entry the glossaries package currently knows, and
% that set includes entries restored from \jobname.glsdefs.  The package
% rewrites .glsdefs from all *currently defined* entries, so anything it
% restored is written straight back out: a deleted section kept reappearing in
% the catalogue for ever, and the .glsdefs had to be erased by hand.  The
% wrappers therefore record every entry they actually create in the .aux, and
% \printboth adds only those.  The .aux is rebuilt from scratch on each run, so
% this is self-healing: a removed result is gone after one further run.
\makeatletter
\newcommand{\@glslive}{}
\newcommand{\@glsaddlive}[1]{\ifglsentryexists{#1}{\glsadd{#1}}{}}
\newcommand{\glsliveentry}[1]{\g@addto@macro\@glslive{\@glsaddlive{#1}}}
\makeatother
% PRINT COMMAND
% No \clearpage here -- the catalogue continues on the current page.  This has
% been re-added by mistake several times; leave it out.
\makeatletter
\newcommand{\printboth}{%
  \@glslive   % add exactly the entries this document created last run
  \printglossary[style=unifiedstyle, title={0 \hspace{.22em} Definitions, Theorems, Examples and Open Problems}]%
}
\makeatother
% Named theorem reference: \tref{word} -> Thm~N(word-or-short)
% Shows stored short title if available, otherwise shows the key itself.
\newcommand{\tref}[1]{%
  Thm~\ref{t:#1}(\ifcsname @tshort@#1\endcsname\csname @tshort@#1\endcsname\else#1\fi)%
}

% Full-title theorem reference: \nref{t:KEY} prints the number and then the
% theorem's long title, so one writes "Theorem~\nref{t:blocks}" and gets
% "Theorem 3 ``Block factorization and siphons''".  Takes the complete label, like
% \ref, not the bare key like \tref.
%
% \beT writes each title to the .aux (see \@thmtitle below), so \nref resolves
% forward references from the second LaTeX pass, exactly as \ref does.  If the
% title is not yet known -- first pass, or a label that is not a \beT theorem
% -- it degrades to the number rather than producing ??.
% These two use @ in their names, and this point in the file is outside the
% \makeatletter region opened earlier (it closes at \makeatother above), so
% the region has to be reopened here.
\makeatletter
\newcommand{\@thmtitle}[2]{\global\@namedef{@tlong@#1}{#2}}
\newcommand{\nref}[1]{%
  \ref{#1}%
  \ifcsname @tlong@#1\endcsname
    \ ``\csname @tlong@#1\endcsname''%
  \fi
}
% \nm{t:KEY}: the stored long title of any \beT/\beL/\beP/\beR/\beC/\beD/\beE result,
% verbatim and without quotes, so item names follow the macro's first argument.
\newcommand{\nm}[1]{%
  \ifcsname @tlong@#1\endcsname
    \csname @tlong@#1\endcsname%
  \else
    \ref{#1}%
  \fi
}
\makeatother




%% ========================================
%% Other THEOREM ENVIRONMENTS
%% ========================================
\makeatletter
\@ifundefined{lemma}      {\newtheorem{lemma}      {Lemma}}      {}
\@ifundefined{proposition}{\newtheorem{proposition}{Proposition}} {}
\@ifundefined{corollary}  {\newtheorem{corollary}  {Corollary}}  {}
\@ifundefined{conjecture} {\newtheorem{conjecture} {Conjecture}}  {}
\@ifundefined{remark}     {\newtheorem{remark}     {Remark}}     {}
\@ifundefined{question}   {\newtheorem{question}   {Question}}   {}
\@ifundefined{assumption} {\newtheorem{assumption} {Assumption}} {}
\@ifundefined{hypothesis} {\newtheorem{hypothesis} {Hypothesis}} {}
\@ifundefined{fact}       {\newtheorem{fact}       {Fact}}       {}
\makeatother

% \citer: cite the relay-graph threshold theorem of AAHrel. If a document has
% \input the chapter/file containing t:relayRF (so the label resolves), point
% straight at it; otherwise fall back to a bare citation of AAHrel. Shared here
% (not per-paper) so it resolves consistently across every book chapter.
% \citeR: same idea, for the first theorem of AAHRah (t:Pstar).
\makeatletter
\providecommand{\citer}{%
  \@ifundefined{r@t:relayRF}%
    {\cite{AAHrel}}%
    {Theorem~\ref{t:relayRF} of \cite{AAHrel}}%
}
\providecommand{\citeR}{%
  \@ifundefined{r@t:Pstar}%
    {\cite{AAHRah}}%
    {Theorem~\ref{t:Pstar} of \cite{AAHRah}}%
}
\makeatother


% Abbreviations
\def\beQ{\begin{question}}\def\eeQ{\end{question}}
\def\beA{\begin{assumption}}\def\eeA{\end{assumption}}


%% ========================================

%% ========================================
%% COLORS & CUSTOM COMMANDS
%% ========================================
\definecolor{funccolor}{RGB}{25,25,112}
\definecolor{desccolor}{RGB}{64,64,64}
\newcommand{\func}[1]{\texttt{\textcolor{funccolor}{#1}}}
\def\eq{=}
\def\mm{-}

%% ========================================
%% MATHCAL SHORTHANDS
%% ========================================
% Project-wide convention: \mX for a capital X is always \mathcal X.
% Defined for the whole alphabet so the convention never fails to expand.
\newcommand{\mA}{\mathcal A}\newcommand{\mB}{\mathcal B}
\newcommand{\mC}{\mathcal C}\newcommand{\mD}{\mathcal D}
\newcommand{\mE}{\mathcal E}\newcommand{\mF}{\mathcal F}
\newcommand{\mG}{\mathcal G}\newcommand{\mH}{\mathcal H}
\newcommand{\mI}{\mathcal I}\newcommand{\mJ}{\mathcal J}
\newcommand{\mK}{\mathcal K}\newcommand{\mL}{\mathcal L}
\newcommand{\mM}{\mathcal M}\newcommand{\mN}{\mathcal N}
\newcommand{\mO}{\mathcal O}\newcommand{\mP}{\mathcal P}
\newcommand{\mQ}{\mathcal Q}\newcommand{\mR}{\mathcal R}
\newcommand{\mS}{\mathcal S}\newcommand{\mT}{\mathcal T}
\newcommand{\mU}{\mathcal U}\newcommand{\mV}{\mathcal V}
\newcommand{\mW}{\mathcal W}\newcommand{\mX}{\mathcal X}
\newcommand{\mY}{\mathcal Y}\newcommand{\mZ}{\mathcal Z}

%% ========================================
%% BASIC MATH ENVIRONMENTS
%% ========================================
\def\bep{\begin{pmatrix}}\def\eep{\end{pmatrix}}
\def\bev{\begin{vmatrix}}\def\eev{\end{vmatrix}}
\def\bea{\begin{eqnarray*}}\def\eea{\end{eqnarray*}}
\def\bc{\begin{cases}}\def\ec{\end{cases}}
\def\BEN{\begin{enumerate}}\def\EEN{\end{enumerate}}
\def\BI{\begin{itemize}}\def\EI{\end{itemize}}

\newcommand{\be}[1]{\begin{equation}\label{#1}}
\newcommand{\ee}{\end{equation}}
\newcommand{\beq}{\begin{eqnarray}}
\def\eeq{\end{eqnarray}}
\def\eqr{\eqref}\def\fr{\frac}\def\lbl{\label}
\def\Lra{\Longrightarrow}\def\non{\nonumber}\def\cNGM{\cite{Diek,Van,Van08}}
\def\Eq{\Leftrightarrow}\def\b{\beta}\def\uph{\red{UP TO HERE}}
%% ========================================
%% PROBABILITY & STATISTICS
%% ========================================
\renewcommand{\P}{\ensuremath{\mathbb{P}}}  % Note: overrides paragraph symbol
\newcommand{\E}{\ensuremath{\mathbb{E}}}
\newcommand{\ind}[1]{\ensuremath{\mathbbm{1}_{\left\{#1\right\}}}}
\DeclareMathOperator{\Cov}{Cov}
\DeclareMathOperator{\TV}{TV}
\newcommand{\rank}{\operatorname{rank}}

%% ========================================
%% CRN OPERATORS (organizations/ERCs/buffering/siphons framework)
%% ========================================
\def\clos{\mathrm{clos}}
\def\Supply{\mathrm{Supply}}
\def\req{\mathrm{req}}
\def\React{\mathrm{React}}
\def\Prod{\mathrm{Prod}}
\def\Gen{\mathrm{Gen}}
\def\In{\mathrm{In}}
\def\Out{\mathrm{Out}}
\def\Act{\mathrm{Act}}
\def\Min{\mathrm{Min}}

%% ========================================
%% LINE BREAKING
%% ========================================
% Last-resort stretch, used only on paragraphs TeX cannot otherwise break:
% a long \texttt{} token, a comma-separated inline math list (commas are not
% breakpoints in math), or a run of \cite/\ref boxes.  Without it those set as
% overfull \hbox and print into the margin; with it they set slightly loose,
% which is the lesser evil.  Does not affect paragraphs that already break.
\emergencystretch=3em

%% ========================================
%% NUMBER SETS
%% ========================================
\newcommand{\R}{\mathbb{R}}  % Real numbers
\newcommand{\C}{\mathbb{C}}  % Complex numbers
\newcommand{\N}{\mathbb{N}}
\newcommand{\Z}{\mathbb{Z}}
% \providecommand, not \newcommand: some papers define these in their own
% preamble after \input{preamble}, and must not clash.
\providecommand{\Q}{\mathbb{Q}}
\providecommand{\res}{\operatorname{res}}   % resultant; used by epiRes.tex
\newcommand{\MN}{\mathbb{N}}
\newcommand{\MZ}{\mathbb{Z}}
\newcommand{\realpos}{\mathbb{R}_{> 0}}
\newcommand{\realnneg}{\mathbb{R}_{\geq 0}}
\newcommand{\naturalpos}{\mathbb{N}_{>0}}
\newcommand{\Rplus}{\mathbb{R}_+}
\newcommand{\Rp}{\mathbb{R}_{\ge 0}}
\def\Rnne{\mathbb{R}_{\ge 0}}
\def\Se{\bar{S}} \def\Ie{\bar{I}}
\newcommand{\Rpos}{\mathbb{R}_{>0}}
\newcommand{\dd}{\,\mathrm d}
\newcommand{\e}{\mathrm e}
%% ========================================
%% GREEK LETTERS (2+ character abbreviations) greek
%% ========================================
\def\al{\alpha}\def\ba{\bff \alpha}
\def\w{\omega}\def\Lgn{\ga+\mu+\de}\def\om{\omega}
\def\de{\delta}\def\De{\Delta}\def\bde{\bff \delta}
\def\ep{\epsilon}\def\eps{\varepsilon}
\def\f{\varphi}  % Kept as requested
\def\ga{\gamma}\def\Ga{\Gamma}\def\G{\Gamma}\def\bg{\bff \gamma}
\def\up{\upsilon}  % NOT the breve accent: no accent commands anywhere.
                   % Named \up, not \u, so the LaTeX breve stays reachable.
\def\la{\lambda}\def\La{\Lambda}
\def\si{\sigma}\def\ka{\kappa}
\def\th{\theta}\def\vt{\vartheta}
\def\ze{\zeta}\def\Gap{\Gamma^{+}} \def\Gam{\Gamma^{-}}


%% ========================================
%% SPECIAL SHORT COMMANDS (kept as requested)
%% ========================================
\def\I{\infty}
\def\T{\widetilde}
\def\H{\hat}
\newcommand{\relint}{\operatorname{relint}}

%% ========================================
%% BOLD/VECTOR NOTATION
%% ========================================
\newcommand{\bff}[1]{{\mbox{\boldmath$#1$}}}
\newcommand\X{\boldsymbol{X}}
\def\x{\boldsymbol{x}}\def\bx{\bff x}\def\xe{\bar{\x}}
\newcommand\y{\boldsymbol{y}}\def\by{\boldsymbol{y}}\def\ye{\bar{\y}}
\def\yb{\y}\def\xb{\x}
\def\xy{\x-\y}
\newcommand\br{\boldsymbol{r}}\def\re{\bar{\r}}
\newcommand\bd{\boldsymbol{d}}
\newcommand\bQ{\boldsymbol{Q}}
\newcommand{\Yb}{{\mathbf{Y}}}
\newcommand{\Ab}{\mathbf{A}}
\newcommand{\zb}{\mathbf{z}}\def\bz{\bff 0}
\newcommand{\Zb}{\mathbf{Z}}
\newcommand{\bb}{\mathbf{b}}
\def\vbe{\vec {\bff \beta}}
\def\vb{\vec {\bff b}}
\def\v1{\vec {\bff 1}}
\def\vn{\vec \de}
\def\vz{\vec 0}
\def\vi{\overset{\rightarrow}{\mathsf i}}
\def\vv{\overset{\rightarrow}{v}}
\def\vr{\; \vec {\mathsf r}}
\newcommand{\va}{\bff a}
\providecommand{\vect}[1]{\boldsymbol{#1}}   % bold vector (moved from Gradostat.tex)
\providecommand{\one}{\mathbf 1}             % all-ones vector (moved from AliSF)
\providecommand{\spr}{s}                     % spectral bound  (moved from AliSF)
\providecommand{\Ke}{K_{\rm eff}}            % effective carrying capacity
% Critical right/left eigenvectors in the Hopf projection method (AAHOSN/New.tex)
\providecommand{\bq}{\mathbf{q}}
\providecommand{\bp}{\mathbf{p}}


%% ========================================
%% MATHCAL & MATHSCR
%% ========================================
% Mathcal
\newcommand\mc[1]{\mathcal{#1}}
\def\mC{{\mathcal C}}\def\mD{\mathcal D}\def\mE{\mathcal{E}}
\def\mF{\mathcal F}\def\mG{{\mathcal G}}\def\mH{{\mathcal H}}
\def\mI{\mathcal I}\def\mL{\mathcal{L}}\def\mM{{\mathcal M}}
\def\mN{{\mathcal N}}\def\mP{\mathcal{P}}\def\mR{{\mathcal R}}
\def\mS{{\mathcal S}}\def\mT{\mathcal{T}}\def\mV{{\mathcal V}}
\def\mK{\mathcal K}\def\mA{{\mathcal A}}\def\mU{{\mathcal U}}
% Symmetric bilinear and trilinear forms of the Taylor expansion, used in the
% first Lyapunov coefficient (AAHOSN/New.tex)
\providecommand{\Bform}{\mathcal{B}}
\providecommand{\Cform}{\mathcal{C}}

% Blackboard bold.  Kept here, not in the drivers: a \providecommand in a
% driver preamble runs before this file and silently wins, so the same name
% ends up meaning different things in different builds.
\def\IC{\mathbb{C}}\def\IN{\mathbb{N}}\def\IQ{\mathbb{Q}}
\def\IR{\mathbb{R}}\def\IZ{\mathbb{Z}}\def\IG{\mathbb{G}}
\def\II{\mathbb{I}}

% Mathscr
\newcommand{\Rs}{{\mathscr R}}
\newcommand{\Ss}{{\mathscr S}}
\newcommand{\Ns}{{\mathscr N}}
\newcommand{\Cs}{\mathscr C}


%% ========================================
%% OPERATORS & FUNCTIONS
%% ========================================
\DeclareMathOperator{\tr}{tr}
\DeclareMathOperator{\diag}{Diag}
\DeclareMathOperator{\Diag}{Diag}
\DeclareMathOperator{\intr}{\mathrm{Int}}
\DeclareMathOperator{\Inv}{Inv}
\DeclareMathOperator{\spec}{spec}
\DeclareMathOperator{\AUC}{AUC}               % area under the curve
\newcommand{\ii}{\mathrm{i}}                  % imaginary unit
\newcommand\op{\operatorname}
\newcommand{\supp}{\op{supp}}
\newcommand{\diff}{\mathop{}\mathopen{}\mathrm{d}}
\newcommand{\pdd}[1]{\frac{\partial}{\partial{#1}}}
\DeclareMathOperator{\cone}{cone}
\DeclareMathOperator{\proj}{proj}

%% ========================================
%% BRACKETS & NORMS
%% ========================================
\providecommand{\pr}[1]{\left(#1\right)}
\providecommand{\pp}[1]{\left[#1\right]}
\newcommand\croc[1]{\left\langle #1\right\rangle}
\newcommand{\norm}[1]{\|#1\|}
\newcommand{\inner}[2]{\langle #1, #2 \rangle}
% Short synonym of \inner, used by AAHOSN/New.tex as $\ip{u}{v}=\bar u^\top v$
\providecommand{\ip}[2]{\langle #1 , #2 \rangle}


%% ========================================
%% FORMATTING & UTILITIES
%% ========================================
\newcommand*{\Scale}[2][4]{\scalebox{#1}{$#2$}}
\newcommand{\te}[1]{\text{#1}}
\newcommand\steq[1]{\stackrel{\text{\rm #1.}}{=}}
\newcommand\blue[1]{\textcolor{blue}{#1}}
\newcommand{\red}{\textcolor[rgb]{1.00,0.00,0.00}}
\newcommand{\green}{\textcolor[rgb]{0.00,1.00,0.00}}

% Figure command
\newcommand{\figu}[3]{%
\begin{figure}[H]
\centering
\includegraphics[scale=#3]{#1}
\caption{#2\label{f:#1}}
\end{figure}
}

% Footnote
\long\def\symbolfootnote[#1]#2{%
\begingroup
\def\thefootnote{\fnsymbol{footnote}}\footnote[#1]{#2}%
\endgroup}
\def\fn{\symbolfootnote}

\newcommand{\dsum}{\displaystyle\sum}

%% ========================================
%% CRN & EPIDEMIOLOGY SPECIFIC
%% ========================================
% CRN notation
\def\SS{\mathcal S}
\def\CC{\mathcal C}
\def\RR{\mathcal R}

% Reproduction numbers
\def\Rto{R_{1}^{E_2}}
\def\Rot{R_{2}^{E_1}}
\def\Roc{R_{1,c}}
\def\Rtc{R_{2,c}}
\def\rd{r_{0}}
\def\sd{s_{0}}
\def\dd{D_{0}}
\def\vd{v_{0}}
\def\zd{z_{0}}
\def\se{\sE}
\def\xd{x_{0}}
%% ========================================
%% SPECIAL VALUES & EQUILIBRIA
%% ========================================
\def\eE{\bar{e}}
\def\wE{\bar{w}}
\def\sE{s_{*}}
\def\Me{\T M_{ee}}
\def\Rnne{\mathbb{R}_{\ge 0}}

\newcommand{\ones}{\mathbf 1}


%% ========================================
%% ABBREVIATIONS (alphabetical by category)
%% ========================================

% -- General mathematical terms --
\def\adm{admissible}\def\admy{admissibility}
%\def\and{antisymmetric}% disabled 2026-09-30: it overrode LaTeX's author separator \and (printed "antisymmetric" between authors)
\def\ass{assumption}
\def\cadlag{c\`adl\`ag}
\def\ch{characteristic polynomial}
\def\chp{characteristic polynomial}
\def\coe{coefficient}
\def\comp{computation}
\def\con{condition}
\def\coo{coordinate}
\def\corr{corresponding}
\def\dec{decomposition}
\def\den{denominator}
\def\dete{deterministic}
\def\eig{eigenvalue}\def\eval{eigenvalue}
\def\eigv{eigenvector}\def\eiv{eigenvector}\def\evec{eigenvector}
\def\elt{element}
\def\enn{essentially nonnegative/positivity preserving}
\def\equ{equation}
\def\etal{et al.}
\def\fe{for example}
\def\fno{from now on}
\def\for{\forall}
\def\form{formula}
\def\Fr{Furthermore, }
\def\frt{furthermore}
\def\fun{function}
\def\gene{generalization}
\def\heu{heuristic}
\def\how{however}\def\How{However, }
\def\ie{i.e. }
\def\im{\item}
\def\isd{is defined by}
\def\itc{it is easy to check that}
\def\itf{it follows that}\def\Itf{It follows that}
\def\ith{it holds that}
\def\itm{it may be checked that}\def\Itm{It may be checked that}
\def\ito{it turns out that}\def\Ito{It turns out that}\def\Itou{It turns out useful to}
\def\ma{mass action}\def\MAR{mass action representation}
\def\mat{matrix}
\def\mbc{may be checked to be}
\def\mbw{may be written as}\def\mubw{must be written as}
\def\mpr{more precisely, }\def\Mp{More precisely, }
\def\nc{non-catalytic}
\def\nne{non-negative}\def\nny{nonnegativity}
\def\no{\nonumber}
\def\noni{non-infectious}
\def\nota{notation}
\def\np{\newpage}
\def\Nt{Note that}
\def\nt{non trivial}\def\nz{non zero}
\def\oth{otherwise}
\def\Oth{On the other hand}
\def\pa{parameter}\def\para{parameter}
\def\pol{polynomial}
\def\prob{probability}
\def\Prf{{\bf  Proof:}}\def\prf{{\bf  Proof:}}
\def\prop{proportional to}
\def\qu{\quad}
\def\QED{\hfill {$\square$}\goodbreak \medskip}
\def\rep{representation}
\def\resp{respectively}
\def\satd{satisfied}\def\satg{satisfying}\def\sats{satisfies}\def\saty{satisfy}
\def\ssec{\subsection}\def\sssec{\subsubsection}
\def\sef{see for example}
\def\sksy{skew-symmetric}
\def\sub{substrate}
\def\sys{system}
\def\var{variable}
\def\wk{well-known}
\def\wl{without loss of generality}\def\wlo{w.l.o.g.}
\def\wmc{We may check that}\def\Wmc{We may check that}
\def\wmr{we may rewrite}
\def\wms{we must show that}
\def\wrt{with respect to}

% -- CRN & Chemistry --
\newcommand\CRN{chemical reaction network}
\newcommand\CRNs{chemical reaction networks}
\newcommand\RN{reaction network}
\def\ACR{absolute concentration robustness}
\def\AMSD{acyclic siphon decomposition}
\def\DAE{differential algebraic equations}
\def\FHJ{Feinberg-Horn-Jackson}
\def\GMAK{generalized mass action kinetics}
\def\MM{Michaelis-Menten}
\def\sm{stoichiometric matrix}\def\Sm{stoichiometric matrix}
\def\WRZD{weakly reversible, zero deficiency}\def\aut{autocatalytic}

% -- Epidemiology --
\def\brn{basic reproduction number}\def\bzn{basic replacement number}
\def\com{compartment}
\def\dfe{disease free equilibrium}\def\DFE{disease free equilibrium}
\def\dwt{dwell times}\def\edt{expected dwell times}
\def\EM{epidemic model}
\def\endp{endemic fixed point}\def\enp{endemic point}\def\EE{endemic point}
\def\ME{mathematical epidemiology}
\def\MEE{eco-epidemiology}\def\EEP{eco-epidemiological}
\def\ece{eco-epidemiological model}
\def\mgm{multi group model}
\def\mtS{multi-strain model}\def\mSt{multi-strain model}
\def\mSi{minimal siphon}
\def\NGM{next generation matrix}
\def\nim{new infections matrix}
\def\qem{quadratic epidemiological model}
\def\repF{invasion functions}
\def\SPF{SIR-PH-FA model}
\def\sus{susceptible}

% -- Graph theory & Linear algebra --
\def\jac{Jacobi matrix}
\def\jin{Jacobian matrix of the invasion vector field with respect to the invasion variables}
\def\jrn{Jacobian matrix of the resident vector field with respect to the resident variables}
\def\jxn{Jacobian matrix of the $\x$ vector field with respect to the $\x$ variables}
\def\Lm{Laplacian matrix}
\def\MTT{matrix tree theorem}
\def\SGT{spectral graph theory}
\def\SMW{Sherman–Morrison–Woodbury}

% -- Stability & Dynamics --
\def\absC{absorbtion CTMC}
\def\BTb{Bogdanov-Takens bifurcation}
\def\CEP{competitive exclusion partition}
\def\CTMC{continuous time Markov chain}
\def\ds{differential system}
\def\fp{fixed point}\def\bfp{boundary fixed point}\def\fbp{fixed boundary point}
\def\fbl{fixed boundary line}\def\fif{forward invariant face}
\def\gas{global asymptotic stability}\def\GAS{globally asymptotically stable}
\def\GASS{\text{GASS}}
\def\Hb{Hopf bifurcation}
\def\Hur{Hurwitz}\def\MH{Metzler-Hurwitz}\def\RH{Routh-Hurwitz}
\def\IS{\text{IS}}
\def\LAS{locally asymptotically stable}
\def\LdS{Lyapunov diagonally stable}
\def\Lf{Lyapunov function}\def\LyV{Lyapunov Volterra type function}\def\RLF{robust Lyapunov function}
\def\LyM{Lyapunov Malkin stable}
\def\PF{Perron-Frobenius}
\def\spfp{strictly positive fixed point}
\def\sta{stability}
\def\STP{sharp threshold property}\def\STT{sharp threshold theorem}

% -- Models (with \varphi explicitly instead of \f) --
\def\AB{$(A,B)$ Arino-Brauer epidemic models}
\def\ABa{$(A,\bbe,\ba,\varphi)$ balanced bilinear model}
\def\ABP{%$(A,\bbe,P,\varphi)$
\\ balanced bilinear model}
\def\AMG{Anderson May Gupta}
\newcommand\BIN{biological interaction network}
\def\Dt{Descartes type}
\def\gLV{generalized Lotka-Volterra model}
\def\IG{\text{invasion graph}}\def\IGMS{interaction graph on minimal siphons}
\def\LCP{\text{linear complementarity problem}}
\def\LV{Lotka-Volterra}\def\LVK{Lotka-Volterra-Kolmogorov}

% -- Splitting & Structure --
\def\reg{regular splitting}\def\regS{{\bf regular splitting}}
\def\RUR{rational univariate representation}
\def\bds{block-diagonal structure}\def\NGMs{next generation matrices}

% -- Miscellaneous --
\def\Ma{{Mathematica}}
\def\Mm{Max–Min}
\def\num{numerator}
\def\rap{rational approximation} \def\DSR{directed species-reaction bipartite graph}\def\PCP{polynomial complementarity problem}\def\rat{rational substitution}
\def\QSS{quasi steady state} \def\RHS{right hand side} \def\LHS{left hand side}
\def\OSN{online social network}\def\VdD{Van den Driessche}
\def\wli{$\omega$-limit} \def\rk{rich kinetics}\def\sJ{symbolic Jacobian}
\def\cV{\cite{Vas,VasSta,BSV,GSV}} \def\spf{semi-parametric semi-functional}
\def\CRW{Capasso-Ruan-Wang} \def\ABP{$(A,\bbe,P,\f)$ balanced bilinear model}
\def\ABa{%$(A,\bbe,\ba,\f)$ \\
rank-one balanced bilinear model}
\def\bbe{B%{\bff \beta}
} \def\vbe{\vec {\bff \beta}
}
\def\vb{\vec {\bff b}}
\newcommand{\Kh}{\hat{K}}          % K with hat, first mention only
\newcommand{\wS}{w^{\sigma}}       % right Perron J^\perp, full block
\newcommand{\piS}{\vec \pi^{\sigma}}    % left Perron J^\perp, full block
\newcommand{\ws}{w^{\sigma}}       % right Perron J^\perp, sub-block
\newcommand{\pis}{\pi^{\sigma}}    % left Perron J^\perp, sub-block
\newcommand{\hwS}{\hat{w}^{\sigma}}   % right Perron K, full block
\newcommand{\hpiS}{\hat{\pi}^{\sigma}} % left Perron K, full block
\newcommand{\hws}{\hat{w}^{\sigma}}   % right Perron K, sub-block
\newcommand{\hpis}{\hat{\pi}^{\sigma}} % left Perron K, sub-block
\newcommand{\twS}{\tilde{w}^{\sigma}}  % right Perron \tilde K, full block
\newcommand{\tpiS}{\tilde{\pi}^{\sigma}} % left Perron \tilde K
\newcommand{\tws}{\tilde{w}^{\sigma}}
\newcommand{\tpis}{\tilde{\pi}^{\sigma}}
\newcommand{\Mat}[2]{\mathbb{R}_{\ge 0}^{\,#1\times #2}}
\def\tJ{J_\sigma^\perp} \def\lPe{left Perron eigenvector}
% \G is \Gamma (defined at :795).  The Volterra function keeps its own name,
% so that \G no longer has two conflicting meanings in this file.
\def\GV{G}  % Volterra G(u)=u-1-\ln u
\def\PVL{Perron-Volterra Lyapunov minimal siphon ansatz}
\newcommand{\istar}[1]{i_{#1}^{*}}
%\newcommand{\sE}{s_E}
\newcommand{\vE}{v_E}
\def\PV{Perron-Volterra candidate Lyapunov function}
\def\GLV{generalized Lotka-Volterra}
\def\MERN{mathematical epidemiology reaction networks}
\def\CMM{contraction metric method}
\newcommand{\sign}{\operatorname{sign}}
\newcommand{\Var}{\operatorname{Var}} 

% ======================================================================
% IDEAS SHARED BETWEEN THE ABSTRACT AND THE MAIN RESULTS SECTION
%
% Write each idea ONCE, in the abstract; the Main Results section recalls it,
% so the two can never drift apart.
%
%   in the abstract        \ideaitem{1}{The first-order reaction terms ...}
%   in Main Results        \ideahead{1}
%
% \ideaitem stores the text globally and typesets
%   \item \hyperlink{idea:1}{<text>}
% \ideahead typesets
%   \item\hypertarget{idea:1}{}\textbf{<the same text>}
%
% The abstract must come before the Main Results input, which is the normal
% order. If a heading has to read differently from the abstract sentence (for
% instance a long abstract item summarised in one line), override it with the
% optional argument: \ideahead[short heading]{6}.
% ======================================================================
\makeatletter
\newcommand{\ideaitem}[2]{%
  \expandafter\gdef\csname @idea@#1\endcsname{#2}%
  \item\hyperlink{idea:#1}{#2}}
\newcommand{\ideahead}[2][]{%
  \item\hypertarget{idea:#2}{}%
  \def\@ideatmp{#1}%
  \ifx\@ideatmp\@empty
    \@ifundefined{@idea@#2}{%
      \PackageWarning{preamble}{Idea #2 is used in Main Results but was never
        defined by \string\ideaitem\space in the abstract}%
      \textbf{[idea #2 undefined]}%
    }{\textbf{\begingroup\def\ideaterm##1##2{##2}\def\idearest##1{}\@nameuse{@idea@#2}\endgroup}}%
  \else
    \textbf{#1}%
  \fi}
% \ideaterm{<hypertarget>}{<words>}: inside an abstract idea the words become an
% invisible link (the document sets \hypersetup{hidelinks}) jumping to the place
% in Main Results that defines them; inside the Main Results heading itself,
% where the target sits a few lines below, \ideahead strips the link and prints
% the words plainly. Mark the destination with \hypertarget{<name>}{} on the line
% before the display, which adds no equation number and shows nothing.
\newcommand{\ideaterm}[2]{\hyperlink{#1}{#2}}
% \idearest{<text>}: the part of an abstract idea after its first sentence; printed in
% the abstract, omitted from the Main Results heading.
\newcommand{\idearest}[1]{#1}
\makeatother

% ======================================================================
% SHORT-PAPER SWITCH
% \shorttrue in a master that assembles only part of the companion paper
% (Barrat/SipArxiv.tex). \nsh{...} drops text that cites results living in the
% omitted sections, so the short build has no dangling references while the
% companion paper keeps the text verbatim.
% ======================================================================
\newif\ifshort \shortfalse
\newcommand{\nsh}[1]{\ifshort\else#1\fi}
% \shortalt{<short build>}{<companion paper>}: use when a sentence needs
% different wording in the two builds, not merely a deletion.
\newcommand{\shortalt}[2]{\ifshort#1\else#2\fi}
% \Jtr, \JtrS{\Sigma}: transverse Jacobian. J^\perp here; SipArxiv.tex redefines them as J_{\mathrm{tra}}, J_{\mathrm{tra},\Sigma}
\newcommand{\Jtr}{J^\perp}
\newcommand{\JtrS}[1]{J_{#1}^\perp}
, and must not clash.
\providecommand{\Q}{\mathbb{Q}}
\providecommand{\res}{\operatorname{res}}   % resultant; used by epiRes.tex
\newcommand{\MN}{\mathbb{N}}
\newcommand{\MZ}{\mathbb{Z}}
\newcommand{\realpos}{\mathbb{R}_{> 0}}
\newcommand{\realnneg}{\mathbb{R}_{\geq 0}}
\newcommand{\naturalpos}{\mathbb{N}_{>0}}
\newcommand{\Rplus}{\mathbb{R}_+}
\newcommand{\Rp}{\mathbb{R}_{\ge 0}}
\def\Rnne{\mathbb{R}_{\ge 0}}
\def\Se{\bar{S}} \def\Ie{\bar{I}}
\newcommand{\Rpos}{\mathbb{R}_{>0}}
\newcommand{\dd}{\,\mathrm d}
\newcommand{\e}{\mathrm e}
%% ========================================
%% GREEK LETTERS (2+ character abbreviations) greek
%% ========================================
\def\al{\alpha}\def\ba{\bff \alpha}
\def\w{\omega}\def\Lgn{\ga+\mu+\de}\def\om{\omega}
\def\de{\delta}\def\De{\Delta}\def\bde{\bff \delta}
\def\ep{\epsilon}\def\eps{\varepsilon}
\def\f{\varphi}  % Kept as requested
\def\ga{\gamma}\def\Ga{\Gamma}\def\G{\Gamma}\def\bg{\bff \gamma}
\def\up{\upsilon}  % NOT the breve accent: no accent commands anywhere.
                   % Named \up, not \u, so the LaTeX breve stays reachable.
\def\la{\lambda}\def\La{\Lambda}
\def\si{\sigma}\def\ka{\kappa}
\def\th{\theta}\def\vt{\vartheta}
\def\ze{\zeta}\def\Gap{\Gamma^{+}} \def\Gam{\Gamma^{-}}


%% ========================================
%% SPECIAL SHORT COMMANDS (kept as requested)
%% ========================================
\def\I{\infty}
\def\T{\widetilde}
\def\H{\hat}
\newcommand{\relint}{\operatorname{relint}}

%% ========================================
%% BOLD/VECTOR NOTATION
%% ========================================
\newcommand{\bff}[1]{{\mbox{\boldmath$#1$}}}
\newcommand\X{\boldsymbol{X}}
\def\x{\boldsymbol{x}}\def\bx{\bff x}\def\xe{\bar{\x}}
\newcommand\y{\boldsymbol{y}}\def\by{\boldsymbol{y}}\def\ye{\bar{\y}}
\def\yb{\y}\def\xb{\x}
\def\xy{\x-\y}
\newcommand\br{\boldsymbol{r}}\def\re{\bar{\r}}
\newcommand\bd{\boldsymbol{d}}
\newcommand\bQ{\boldsymbol{Q}}
\newcommand{\Yb}{{\mathbf{Y}}}
\newcommand{\Ab}{\mathbf{A}}
\newcommand{\zb}{\mathbf{z}}\def\bz{\bff 0}
\newcommand{\Zb}{\mathbf{Z}}
\newcommand{\bb}{\mathbf{b}}
\def\vbe{\vec {\bff \beta}}
\def\vb{\vec {\bff b}}
\def\v1{\vec {\bff 1}}
\def\vn{\vec \de}
\def\vz{\vec 0}
\def\vi{\overset{\rightarrow}{\mathsf i}}
\def\vv{\overset{\rightarrow}{v}}
\def\vr{\; \vec {\mathsf r}}
\newcommand{\va}{\bff a}
\providecommand{\vect}[1]{\boldsymbol{#1}}   % bold vector (moved from Gradostat.tex)
\providecommand{\one}{\mathbf 1}             % all-ones vector (moved from AliSF)
\providecommand{\spr}{s}                     % spectral bound  (moved from AliSF)
\providecommand{\Ke}{K_{\rm eff}}            % effective carrying capacity
% Critical right/left eigenvectors in the Hopf projection method (AAHOSN/New.tex)
\providecommand{\bq}{\mathbf{q}}
\providecommand{\bp}{\mathbf{p}}


%% ========================================
%% MATHCAL & MATHSCR
%% ========================================
% Mathcal
\newcommand\mc[1]{\mathcal{#1}}
\def\mC{{\mathcal C}}\def\mD{\mathcal D}\def\mE{\mathcal{E}}
\def\mF{\mathcal F}\def\mG{{\mathcal G}}\def\mH{{\mathcal H}}
\def\mI{\mathcal I}\def\mL{\mathcal{L}}\def\mM{{\mathcal M}}
\def\mN{{\mathcal N}}\def\mP{\mathcal{P}}\def\mR{{\mathcal R}}
\def\mS{{\mathcal S}}\def\mT{\mathcal{T}}\def\mV{{\mathcal V}}
\def\mK{\mathcal K}\def\mA{{\mathcal A}}\def\mU{{\mathcal U}}
% Symmetric bilinear and trilinear forms of the Taylor expansion, used in the
% first Lyapunov coefficient (AAHOSN/New.tex)
\providecommand{\Bform}{\mathcal{B}}
\providecommand{\Cform}{\mathcal{C}}

% Blackboard bold.  Kept here, not in the drivers: a \providecommand in a
% driver preamble runs before this file and silently wins, so the same name
% ends up meaning different things in different builds.
\def\IC{\mathbb{C}}\def\IN{\mathbb{N}}\def\IQ{\mathbb{Q}}
\def\IR{\mathbb{R}}\def\IZ{\mathbb{Z}}\def\IG{\mathbb{G}}
\def\II{\mathbb{I}}

% Mathscr
\newcommand{\Rs}{{\mathscr R}}
\newcommand{\Ss}{{\mathscr S}}
\newcommand{\Ns}{{\mathscr N}}
\newcommand{\Cs}{\mathscr C}


%% ========================================
%% OPERATORS & FUNCTIONS
%% ========================================
\DeclareMathOperator{\tr}{tr}
\DeclareMathOperator{\diag}{Diag}
\DeclareMathOperator{\Diag}{Diag}
\DeclareMathOperator{\intr}{\mathrm{Int}}
\DeclareMathOperator{\Inv}{Inv}
\DeclareMathOperator{\spec}{spec}
\DeclareMathOperator{\AUC}{AUC}               % area under the curve
\newcommand{\ii}{\mathrm{i}}                  % imaginary unit
\newcommand\op{\operatorname}
\newcommand{\supp}{\op{supp}}
\newcommand{\diff}{\mathop{}\mathopen{}\mathrm{d}}
\newcommand{\pdd}[1]{\frac{\partial}{\partial{#1}}}
\DeclareMathOperator{\cone}{cone}
\DeclareMathOperator{\proj}{proj}

%% ========================================
%% BRACKETS & NORMS
%% ========================================
\providecommand{\pr}[1]{\left(#1\right)}
\providecommand{\pp}[1]{\left[#1\right]}
\newcommand\croc[1]{\left\langle #1\right\rangle}
\newcommand{\norm}[1]{\|#1\|}
\newcommand{\inner}[2]{\langle #1, #2 \rangle}
% Short synonym of \inner, used by AAHOSN/New.tex as $\ip{u}{v}=\bar u^\top v$
\providecommand{\ip}[2]{\langle #1 , #2 \rangle}


%% ========================================
%% FORMATTING & UTILITIES
%% ========================================
\newcommand*{\Scale}[2][4]{\scalebox{#1}{$#2$}}
\newcommand{\te}[1]{\text{#1}}
\newcommand\steq[1]{\stackrel{\text{\rm #1.}}{=}}
\newcommand\blue[1]{\textcolor{blue}{#1}}
\newcommand{\red}{\textcolor[rgb]{1.00,0.00,0.00}}
\newcommand{\green}{\textcolor[rgb]{0.00,1.00,0.00}}

% Figure command
\newcommand{\figu}[3]{%
\begin{figure}[H]
\centering
\includegraphics[scale=#3]{#1}
\caption{#2\label{f:#1}}
\end{figure}
}

% Footnote
\long\def\symbolfootnote[#1]#2{%
\begingroup
\def\thefootnote{\fnsymbol{footnote}}\footnote[#1]{#2}%
\endgroup}
\def\fn{\symbolfootnote}

\newcommand{\dsum}{\displaystyle\sum}

%% ========================================
%% CRN & EPIDEMIOLOGY SPECIFIC
%% ========================================
% CRN notation
\def\SS{\mathcal S}
\def\CC{\mathcal C}
\def\RR{\mathcal R}

% Reproduction numbers
\def\Rto{R_{1}^{E_2}}
\def\Rot{R_{2}^{E_1}}
\def\Roc{R_{1,c}}
\def\Rtc{R_{2,c}}
\def\rd{r_{0}}
\def\sd{s_{0}}
\def\dd{D_{0}}
\def\vd{v_{0}}
\def\zd{z_{0}}
\def\se{\sE}
\def\xd{x_{0}}
%% ========================================
%% SPECIAL VALUES & EQUILIBRIA
%% ========================================
\def\eE{\bar{e}}
\def\wE{\bar{w}}
\def\sE{s_{*}}
\def\Me{\T M_{ee}}
\def\Rnne{\mathbb{R}_{\ge 0}}

\newcommand{\ones}{\mathbf 1}


%% ========================================
%% ABBREVIATIONS (alphabetical by category)
%% ========================================

% -- General mathematical terms --
\def\adm{admissible}\def\admy{admissibility}
%\def\and{antisymmetric}% disabled 2026-09-30: it overrode LaTeX's author separator \and (printed "antisymmetric" between authors)
\def\ass{assumption}
\def\cadlag{c\`adl\`ag}
\def\ch{characteristic polynomial}
\def\chp{characteristic polynomial}
\def\coe{coefficient}
\def\comp{computation}
\def\con{condition}
\def\coo{coordinate}
\def\corr{corresponding}
\def\dec{decomposition}
\def\den{denominator}
\def\dete{deterministic}
\def\eig{eigenvalue}\def\eval{eigenvalue}
\def\eigv{eigenvector}\def\eiv{eigenvector}\def\evec{eigenvector}
\def\elt{element}
\def\enn{essentially nonnegative/positivity preserving}
\def\equ{equation}
\def\etal{et al.}
\def\fe{for example}
\def\fno{from now on}
\def\for{\forall}
\def\form{formula}
\def\Fr{Furthermore, }
\def\frt{furthermore}
\def\fun{function}
\def\gene{generalization}
\def\heu{heuristic}
\def\how{however}\def\How{However, }
\def\ie{i.e. }
\def\im{\item}
\def\isd{is defined by}
\def\itc{it is easy to check that}
\def\itf{it follows that}\def\Itf{It follows that}
\def\ith{it holds that}
\def\itm{it may be checked that}\def\Itm{It may be checked that}
\def\ito{it turns out that}\def\Ito{It turns out that}\def\Itou{It turns out useful to}
\def\ma{mass action}\def\MAR{mass action representation}
\def\mat{matrix}
\def\mbc{may be checked to be}
\def\mbw{may be written as}\def\mubw{must be written as}
\def\mpr{more precisely, }\def\Mp{More precisely, }
\def\nc{non-catalytic}
\def\nne{non-negative}\def\nny{nonnegativity}
\def\no{\nonumber}
\def\noni{non-infectious}
\def\nota{notation}
\def\np{\newpage}
\def\Nt{Note that}
\def\nt{non trivial}\def\nz{non zero}
\def\oth{otherwise}
\def\Oth{On the other hand}
\def\pa{parameter}\def\para{parameter}
\def\pol{polynomial}
\def\prob{probability}
\def\Prf{{\bf  Proof:}}\def\prf{{\bf  Proof:}}
\def\prop{proportional to}
\def\qu{\quad}
\def\QED{\hfill {$\square$}\goodbreak \medskip}
\def\rep{representation}
\def\resp{respectively}
\def\satd{satisfied}\def\satg{satisfying}\def\sats{satisfies}\def\saty{satisfy}
\def\ssec{\subsection}\def\sssec{\subsubsection}
\def\sef{see for example}
\def\sksy{skew-symmetric}
\def\sub{substrate}
\def\sys{system}
\def\var{variable}
\def\wk{well-known}
\def\wl{without loss of generality}\def\wlo{w.l.o.g.}
\def\wmc{We may check that}\def\Wmc{We may check that}
\def\wmr{we may rewrite}
\def\wms{we must show that}
\def\wrt{with respect to}

% -- CRN & Chemistry --
\newcommand\CRN{chemical reaction network}
\newcommand\CRNs{chemical reaction networks}
\newcommand\RN{reaction network}
\def\ACR{absolute concentration robustness}
\def\AMSD{acyclic siphon decomposition}
\def\DAE{differential algebraic equations}
\def\FHJ{Feinberg-Horn-Jackson}
\def\GMAK{generalized mass action kinetics}
\def\MM{Michaelis-Menten}
\def\sm{stoichiometric matrix}\def\Sm{stoichiometric matrix}
\def\WRZD{weakly reversible, zero deficiency}\def\aut{autocatalytic}

% -- Epidemiology --
\def\brn{basic reproduction number}\def\bzn{basic replacement number}
\def\com{compartment}
\def\dfe{disease free equilibrium}\def\DFE{disease free equilibrium}
\def\dwt{dwell times}\def\edt{expected dwell times}
\def\EM{epidemic model}
\def\endp{endemic fixed point}\def\enp{endemic point}\def\EE{endemic point}
\def\ME{mathematical epidemiology}
\def\MEE{eco-epidemiology}\def\EEP{eco-epidemiological}
\def\ece{eco-epidemiological model}
\def\mgm{multi group model}
\def\mtS{multi-strain model}\def\mSt{multi-strain model}
\def\mSi{minimal siphon}
\def\NGM{next generation matrix}
\def\nim{new infections matrix}
\def\qem{quadratic epidemiological model}
\def\repF{invasion functions}
\def\SPF{SIR-PH-FA model}
\def\sus{susceptible}

% -- Graph theory & Linear algebra --
\def\jac{Jacobi matrix}
\def\jin{Jacobian matrix of the invasion vector field with respect to the invasion variables}
\def\jrn{Jacobian matrix of the resident vector field with respect to the resident variables}
\def\jxn{Jacobian matrix of the $\x$ vector field with respect to the $\x$ variables}
\def\Lm{Laplacian matrix}
\def\MTT{matrix tree theorem}
\def\SGT{spectral graph theory}
\def\SMW{Sherman–Morrison–Woodbury}

% -- Stability & Dynamics --
\def\absC{absorbtion CTMC}
\def\BTb{Bogdanov-Takens bifurcation}
\def\CEP{competitive exclusion partition}
\def\CTMC{continuous time Markov chain}
\def\ds{differential system}
\def\fp{fixed point}\def\bfp{boundary fixed point}\def\fbp{fixed boundary point}
\def\fbl{fixed boundary line}\def\fif{forward invariant face}
\def\gas{global asymptotic stability}\def\GAS{globally asymptotically stable}
\def\GASS{\text{GASS}}
\def\Hb{Hopf bifurcation}
\def\Hur{Hurwitz}\def\MH{Metzler-Hurwitz}\def\RH{Routh-Hurwitz}
\def\IS{\text{IS}}
\def\LAS{locally asymptotically stable}
\def\LdS{Lyapunov diagonally stable}
\def\Lf{Lyapunov function}\def\LyV{Lyapunov Volterra type function}\def\RLF{robust Lyapunov function}
\def\LyM{Lyapunov Malkin stable}
\def\PF{Perron-Frobenius}
\def\spfp{strictly positive fixed point}
\def\sta{stability}
\def\STP{sharp threshold property}\def\STT{sharp threshold theorem}

% -- Models (with \varphi explicitly instead of \f) --
\def\AB{$(A,B)$ Arino-Brauer epidemic models}
\def\ABa{$(A,\bbe,\ba,\varphi)$ balanced bilinear model}
\def\ABP{%$(A,\bbe,P,\varphi)$
\\ balanced bilinear model}
\def\AMG{Anderson May Gupta}
\newcommand\BIN{biological interaction network}
\def\Dt{Descartes type}
\def\gLV{generalized Lotka-Volterra model}
\def\IG{\text{invasion graph}}\def\IGMS{interaction graph on minimal siphons}
\def\LCP{\text{linear complementarity problem}}
\def\LV{Lotka-Volterra}\def\LVK{Lotka-Volterra-Kolmogorov}

% -- Splitting & Structure --
\def\reg{regular splitting}\def\regS{{\bf regular splitting}}
\def\RUR{rational univariate representation}
\def\bds{block-diagonal structure}\def\NGMs{next generation matrices}

% -- Miscellaneous --
\def\Ma{{Mathematica}}
\def\Mm{Max–Min}
\def\num{numerator}
\def\rap{rational approximation} \def\DSR{directed species-reaction bipartite graph}\def\PCP{polynomial complementarity problem}\def\rat{rational substitution}
\def\QSS{quasi steady state} \def\RHS{right hand side} \def\LHS{left hand side}
\def\OSN{online social network}\def\VdD{Van den Driessche}
\def\wli{$\omega$-limit} \def\rk{rich kinetics}\def\sJ{symbolic Jacobian}
\def\cV{\cite{Vas,VasSta,BSV,GSV}} \def\spf{semi-parametric semi-functional}
\def\CRW{Capasso-Ruan-Wang} \def\ABP{$(A,\bbe,P,\f)$ balanced bilinear model}
\def\ABa{%$(A,\bbe,\ba,\f)$ \\
rank-one balanced bilinear model}
\def\bbe{B%{\bff \beta}
} \def\vbe{\vec {\bff \beta}
}
\def\vb{\vec {\bff b}}
\newcommand{\Kh}{\hat{K}}          % K with hat, first mention only
\newcommand{\wS}{w^{\sigma}}       % right Perron J^\perp, full block
\newcommand{\piS}{\vec \pi^{\sigma}}    % left Perron J^\perp, full block
\newcommand{\ws}{w^{\sigma}}       % right Perron J^\perp, sub-block
\newcommand{\pis}{\pi^{\sigma}}    % left Perron J^\perp, sub-block
\newcommand{\hwS}{\hat{w}^{\sigma}}   % right Perron K, full block
\newcommand{\hpiS}{\hat{\pi}^{\sigma}} % left Perron K, full block
\newcommand{\hws}{\hat{w}^{\sigma}}   % right Perron K, sub-block
\newcommand{\hpis}{\hat{\pi}^{\sigma}} % left Perron K, sub-block
\newcommand{\twS}{\tilde{w}^{\sigma}}  % right Perron \tilde K, full block
\newcommand{\tpiS}{\tilde{\pi}^{\sigma}} % left Perron \tilde K
\newcommand{\tws}{\tilde{w}^{\sigma}}
\newcommand{\tpis}{\tilde{\pi}^{\sigma}}
\newcommand{\Mat}[2]{\mathbb{R}_{\ge 0}^{\,#1\times #2}}
\def\tJ{J_\sigma^\perp} \def\lPe{left Perron eigenvector}
% \G is \Gamma (defined at :795).  The Volterra function keeps its own name,
% so that \G no longer has two conflicting meanings in this file.
\def\GV{G}  % Volterra G(u)=u-1-\ln u
\def\PVL{Perron-Volterra Lyapunov minimal siphon ansatz}
\newcommand{\istar}[1]{i_{#1}^{*}}
%\newcommand{\sE}{s_E}
\newcommand{\vE}{v_E}
\def\PV{Perron-Volterra candidate Lyapunov function}
\def\GLV{generalized Lotka-Volterra}
\def\MERN{mathematical epidemiology reaction networks}
\def\CMM{contraction metric method}
\newcommand{\sign}{\operatorname{sign}}
\newcommand{\Var}{\operatorname{Var}} 

% ======================================================================
% IDEAS SHARED BETWEEN THE ABSTRACT AND THE MAIN RESULTS SECTION
%
% Write each idea ONCE, in the abstract; the Main Results section recalls it,
% so the two can never drift apart.
%
%   in the abstract        \ideaitem{1}{The first-order reaction terms ...}
%   in Main Results        \ideahead{1}
%
% \ideaitem stores the text globally and typesets
%   \item \hyperlink{idea:1}{<text>}
% \ideahead typesets
%   \item\hypertarget{idea:1}{}\textbf{<the same text>}
%
% The abstract must come before the Main Results input, which is the normal
% order. If a heading has to read differently from the abstract sentence (for
% instance a long abstract item summarised in one line), override it with the
% optional argument: \ideahead[short heading]{6}.
% ======================================================================
\makeatletter
\newcommand{\ideaitem}[2]{%
  \expandafter\gdef\csname @idea@#1\endcsname{#2}%
  \item\hyperlink{idea:#1}{#2}}
\newcommand{\ideahead}[2][]{%
  \item\hypertarget{idea:#2}{}%
  \def\@ideatmp{#1}%
  \ifx\@ideatmp\@empty
    \@ifundefined{@idea@#2}{%
      \PackageWarning{preamble}{Idea #2 is used in Main Results but was never
        defined by \string\ideaitem\space in the abstract}%
      \textbf{[idea #2 undefined]}%
    }{\textbf{\begingroup\def\ideaterm##1##2{##2}\def\idearest##1{}\@nameuse{@idea@#2}\endgroup}}%
  \else
    \textbf{#1}%
  \fi}
% \ideaterm{<hypertarget>}{<words>}: inside an abstract idea the words become an
% invisible link (the document sets \hypersetup{hidelinks}) jumping to the place
% in Main Results that defines them; inside the Main Results heading itself,
% where the target sits a few lines below, \ideahead strips the link and prints
% the words plainly. Mark the destination with \hypertarget{<name>}{} on the line
% before the display, which adds no equation number and shows nothing.
\newcommand{\ideaterm}[2]{\hyperlink{#1}{#2}}
% \idearest{<text>}: the part of an abstract idea after its first sentence; printed in
% the abstract, omitted from the Main Results heading.
\newcommand{\idearest}[1]{#1}
\makeatother

% ======================================================================
% SHORT-PAPER SWITCH
% \shorttrue in a master that assembles only part of the companion paper
% (Barrat/SipArxiv.tex). \nsh{...} drops text that cites results living in the
% omitted sections, so the short build has no dangling references while the
% companion paper keeps the text verbatim.
% ======================================================================
\newif\ifshort \shortfalse
\newcommand{\nsh}[1]{\ifshort\else#1\fi}
% \shortalt{<short build>}{<companion paper>}: use when a sentence needs
% different wording in the two builds, not merely a deletion.
\newcommand{\shortalt}[2]{\ifshort#1\else#2\fi}
% \Jtr, \JtrS{\Sigma}: transverse Jacobian. J^\perp here; SipArxiv.tex redefines them as J_{\mathrm{tra}}, J_{\mathrm{tra},\Sigma}
\newcommand{\Jtr}{J^\perp}
\newcommand{\JtrS}[1]{J_{#1}^\perp}
, and must not clash.
\providecommand{\Q}{\mathbb{Q}}
\providecommand{\res}{\operatorname{res}}   % resultant; used by epiRes.tex
\newcommand{\MN}{\mathbb{N}}
\newcommand{\MZ}{\mathbb{Z}}
\newcommand{\realpos}{\mathbb{R}_{> 0}}
\newcommand{\realnneg}{\mathbb{R}_{\geq 0}}
\newcommand{\naturalpos}{\mathbb{N}_{>0}}
\newcommand{\Rplus}{\mathbb{R}_+}
\newcommand{\Rp}{\mathbb{R}_{\ge 0}}
\def\Rnne{\mathbb{R}_{\ge 0}}
\def\Se{\bar{S}} \def\Ie{\bar{I}}
\newcommand{\Rpos}{\mathbb{R}_{>0}}
\newcommand{\dd}{\,\mathrm d}
\newcommand{\e}{\mathrm e}
%% ========================================
%% GREEK LETTERS (2+ character abbreviations) greek
%% ========================================
\def\al{\alpha}\def\ba{\bff \alpha}
\def\w{\omega}\def\Lgn{\ga+\mu+\de}\def\om{\omega}
\def\de{\delta}\def\De{\Delta}\def\bde{\bff \delta}
\def\ep{\epsilon}\def\eps{\varepsilon}
\def\f{\varphi}  % Kept as requested
\def\ga{\gamma}\def\Ga{\Gamma}\def\G{\Gamma}\def\bg{\bff \gamma}
\def\up{\upsilon}  % NOT the breve accent: no accent commands anywhere.
                   % Named \up, not \u, so the LaTeX breve stays reachable.
\def\la{\lambda}\def\La{\Lambda}
\def\si{\sigma}\def\ka{\kappa}
\def\th{\theta}\def\vt{\vartheta}
\def\ze{\zeta}\def\Gap{\Gamma^{+}} \def\Gam{\Gamma^{-}}


%% ========================================
%% SPECIAL SHORT COMMANDS (kept as requested)
%% ========================================
\def\I{\infty}
\def\T{\widetilde}
\def\H{\hat}
\newcommand{\relint}{\operatorname{relint}}

%% ========================================
%% BOLD/VECTOR NOTATION
%% ========================================
\newcommand{\bff}[1]{{\mbox{\boldmath$#1$}}}
\newcommand\X{\boldsymbol{X}}
\def\x{\boldsymbol{x}}\def\bx{\bff x}\def\xe{\bar{\x}}
\newcommand\y{\boldsymbol{y}}\def\by{\boldsymbol{y}}\def\ye{\bar{\y}}
\def\yb{\y}\def\xb{\x}
\def\xy{\x-\y}
\newcommand\br{\boldsymbol{r}}\def\re{\bar{\r}}
\newcommand\bd{\boldsymbol{d}}
\newcommand\bQ{\boldsymbol{Q}}
\newcommand{\Yb}{{\mathbf{Y}}}
\newcommand{\Ab}{\mathbf{A}}
\newcommand{\zb}{\mathbf{z}}\def\bz{\bff 0}
\newcommand{\Zb}{\mathbf{Z}}
\newcommand{\bb}{\mathbf{b}}
\def\vbe{\vec {\bff \beta}}
\def\vb{\vec {\bff b}}
\def\v1{\vec {\bff 1}}
\def\vn{\vec \de}
\def\vz{\vec 0}
\def\vi{\overset{\rightarrow}{\mathsf i}}
\def\vv{\overset{\rightarrow}{v}}
\def\vr{\; \vec {\mathsf r}}
\newcommand{\va}{\bff a}
\providecommand{\vect}[1]{\boldsymbol{#1}}   % bold vector (moved from Gradostat.tex)
\providecommand{\one}{\mathbf 1}             % all-ones vector (moved from AliSF)
\providecommand{\spr}{s}                     % spectral bound  (moved from AliSF)
\providecommand{\Ke}{K_{\rm eff}}            % effective carrying capacity
% Critical right/left eigenvectors in the Hopf projection method (AAHOSN/New.tex)
\providecommand{\bq}{\mathbf{q}}
\providecommand{\bp}{\mathbf{p}}


%% ========================================
%% MATHCAL & MATHSCR
%% ========================================
% Mathcal
\newcommand\mc[1]{\mathcal{#1}}
\def\mC{{\mathcal C}}\def\mD{\mathcal D}\def\mE{\mathcal{E}}
\def\mF{\mathcal F}\def\mG{{\mathcal G}}\def\mH{{\mathcal H}}
\def\mI{\mathcal I}\def\mL{\mathcal{L}}\def\mM{{\mathcal M}}
\def\mN{{\mathcal N}}\def\mP{\mathcal{P}}\def\mR{{\mathcal R}}
\def\mS{{\mathcal S}}\def\mT{\mathcal{T}}\def\mV{{\mathcal V}}
\def\mK{\mathcal K}\def\mA{{\mathcal A}}\def\mU{{\mathcal U}}
% Symmetric bilinear and trilinear forms of the Taylor expansion, used in the
% first Lyapunov coefficient (AAHOSN/New.tex)
\providecommand{\Bform}{\mathcal{B}}
\providecommand{\Cform}{\mathcal{C}}

% Blackboard bold.  Kept here, not in the drivers: a \providecommand in a
% driver preamble runs before this file and silently wins, so the same name
% ends up meaning different things in different builds.
\def\IC{\mathbb{C}}\def\IN{\mathbb{N}}\def\IQ{\mathbb{Q}}
\def\IR{\mathbb{R}}\def\IZ{\mathbb{Z}}\def\IG{\mathbb{G}}
\def\II{\mathbb{I}}

% Mathscr
\newcommand{\Rs}{{\mathscr R}}
\newcommand{\Ss}{{\mathscr S}}
\newcommand{\Ns}{{\mathscr N}}
\newcommand{\Cs}{\mathscr C}


%% ========================================
%% OPERATORS & FUNCTIONS
%% ========================================
\DeclareMathOperator{\tr}{tr}
\DeclareMathOperator{\diag}{Diag}
\DeclareMathOperator{\Diag}{Diag}
\DeclareMathOperator{\intr}{\mathrm{Int}}
\DeclareMathOperator{\Inv}{Inv}
\DeclareMathOperator{\spec}{spec}
\DeclareMathOperator{\AUC}{AUC}               % area under the curve
\newcommand{\ii}{\mathrm{i}}                  % imaginary unit
\newcommand\op{\operatorname}
\newcommand{\supp}{\op{supp}}
\newcommand{\diff}{\mathop{}\mathopen{}\mathrm{d}}
\newcommand{\pdd}[1]{\frac{\partial}{\partial{#1}}}
\DeclareMathOperator{\cone}{cone}
\DeclareMathOperator{\proj}{proj}

%% ========================================
%% BRACKETS & NORMS
%% ========================================
\providecommand{\pr}[1]{\left(#1\right)}
\providecommand{\pp}[1]{\left[#1\right]}
\newcommand\croc[1]{\left\langle #1\right\rangle}
\newcommand{\norm}[1]{\|#1\|}
\newcommand{\inner}[2]{\langle #1, #2 \rangle}
% Short synonym of \inner, used by AAHOSN/New.tex as $\ip{u}{v}=\bar u^\top v$
\providecommand{\ip}[2]{\langle #1 , #2 \rangle}


%% ========================================
%% FORMATTING & UTILITIES
%% ========================================
\newcommand*{\Scale}[2][4]{\scalebox{#1}{$#2$}}
\newcommand{\te}[1]{\text{#1}}
\newcommand\steq[1]{\stackrel{\text{\rm #1.}}{=}}
\newcommand\blue[1]{\textcolor{blue}{#1}}
\newcommand{\red}{\textcolor[rgb]{1.00,0.00,0.00}}
\newcommand{\green}{\textcolor[rgb]{0.00,1.00,0.00}}

% Figure command
\newcommand{\figu}[3]{%
\begin{figure}[H]
\centering
\includegraphics[scale=#3]{#1}
\caption{#2\label{f:#1}}
\end{figure}
}

% Footnote
\long\def\symbolfootnote[#1]#2{%
\begingroup
\def\thefootnote{\fnsymbol{footnote}}\footnote[#1]{#2}%
\endgroup}
\def\fn{\symbolfootnote}

\newcommand{\dsum}{\displaystyle\sum}

%% ========================================
%% CRN & EPIDEMIOLOGY SPECIFIC
%% ========================================
% CRN notation
\def\SS{\mathcal S}
\def\CC{\mathcal C}
\def\RR{\mathcal R}

% Reproduction numbers
\def\Rto{R_{1}^{E_2}}
\def\Rot{R_{2}^{E_1}}
\def\Roc{R_{1,c}}
\def\Rtc{R_{2,c}}
\def\rd{r_{0}}
\def\sd{s_{0}}
\def\dd{D_{0}}
\def\vd{v_{0}}
\def\zd{z_{0}}
\def\se{\sE}
\def\xd{x_{0}}
%% ========================================
%% SPECIAL VALUES & EQUILIBRIA
%% ========================================
\def\eE{\bar{e}}
\def\wE{\bar{w}}
\def\sE{s_{*}}
\def\Me{\T M_{ee}}
\def\Rnne{\mathbb{R}_{\ge 0}}

\newcommand{\ones}{\mathbf 1}


%% ========================================
%% ABBREVIATIONS (alphabetical by category)
%% ========================================

% -- General mathematical terms --
\def\adm{admissible}\def\admy{admissibility}
%\def\and{antisymmetric}% disabled 2026-09-30: it overrode LaTeX's author separator \and (printed "antisymmetric" between authors)
\def\ass{assumption}
\def\cadlag{c\`adl\`ag}
\def\ch{characteristic polynomial}
\def\chp{characteristic polynomial}
\def\coe{coefficient}
\def\comp{computation}
\def\con{condition}
\def\coo{coordinate}
\def\corr{corresponding}
\def\dec{decomposition}
\def\den{denominator}
\def\dete{deterministic}
\def\eig{eigenvalue}\def\eval{eigenvalue}
\def\eigv{eigenvector}\def\eiv{eigenvector}\def\evec{eigenvector}
\def\elt{element}
\def\enn{essentially nonnegative/positivity preserving}
\def\equ{equation}
\def\etal{et al.}
\def\fe{for example}
\def\fno{from now on}
\def\for{\forall}
\def\form{formula}
\def\Fr{Furthermore, }
\def\frt{furthermore}
\def\fun{function}
\def\gene{generalization}
\def\heu{heuristic}
\def\how{however}\def\How{However, }
\def\ie{i.e. }
\def\im{\item}
\def\isd{is defined by}
\def\itc{it is easy to check that}
\def\itf{it follows that}\def\Itf{It follows that}
\def\ith{it holds that}
\def\itm{it may be checked that}\def\Itm{It may be checked that}
\def\ito{it turns out that}\def\Ito{It turns out that}\def\Itou{It turns out useful to}
\def\ma{mass action}\def\MAR{mass action representation}
\def\mat{matrix}
\def\mbc{may be checked to be}
\def\mbw{may be written as}\def\mubw{must be written as}
\def\mpr{more precisely, }\def\Mp{More precisely, }
\def\nc{non-catalytic}
\def\nne{non-negative}\def\nny{nonnegativity}
\def\no{\nonumber}
\def\noni{non-infectious}
\def\nota{notation}
\def\np{\newpage}
\def\Nt{Note that}
\def\nt{non trivial}\def\nz{non zero}
\def\oth{otherwise}
\def\Oth{On the other hand}
\def\pa{parameter}\def\para{parameter}
\def\pol{polynomial}
\def\prob{probability}
\def\Prf{{\bf  Proof:}}\def\prf{{\bf  Proof:}}
\def\prop{proportional to}
\def\qu{\quad}
\def\QED{\hfill {$\square$}\goodbreak \medskip}
\def\rep{representation}
\def\resp{respectively}
\def\satd{satisfied}\def\satg{satisfying}\def\sats{satisfies}\def\saty{satisfy}
\def\ssec{\subsection}\def\sssec{\subsubsection}
\def\sef{see for example}
\def\sksy{skew-symmetric}
\def\sub{substrate}
\def\sys{system}
\def\var{variable}
\def\wk{well-known}
\def\wl{without loss of generality}\def\wlo{w.l.o.g.}
\def\wmc{We may check that}\def\Wmc{We may check that}
\def\wmr{we may rewrite}
\def\wms{we must show that}
\def\wrt{with respect to}

% -- CRN & Chemistry --
\newcommand\CRN{chemical reaction network}
\newcommand\CRNs{chemical reaction networks}
\newcommand\RN{reaction network}
\def\ACR{absolute concentration robustness}
\def\AMSD{acyclic siphon decomposition}
\def\DAE{differential algebraic equations}
\def\FHJ{Feinberg-Horn-Jackson}
\def\GMAK{generalized mass action kinetics}
\def\MM{Michaelis-Menten}
\def\sm{stoichiometric matrix}\def\Sm{stoichiometric matrix}
\def\WRZD{weakly reversible, zero deficiency}\def\aut{autocatalytic}

% -- Epidemiology --
\def\brn{basic reproduction number}\def\bzn{basic replacement number}
\def\com{compartment}
\def\dfe{disease free equilibrium}\def\DFE{disease free equilibrium}
\def\dwt{dwell times}\def\edt{expected dwell times}
\def\EM{epidemic model}
\def\endp{endemic fixed point}\def\enp{endemic point}\def\EE{endemic point}
\def\ME{mathematical epidemiology}
\def\MEE{eco-epidemiology}\def\EEP{eco-epidemiological}
\def\ece{eco-epidemiological model}
\def\mgm{multi group model}
\def\mtS{multi-strain model}\def\mSt{multi-strain model}
\def\mSi{minimal siphon}
\def\NGM{next generation matrix}
\def\nim{new infections matrix}
\def\qem{quadratic epidemiological model}
\def\repF{invasion functions}
\def\SPF{SIR-PH-FA model}
\def\sus{susceptible}

% -- Graph theory & Linear algebra --
\def\jac{Jacobi matrix}
\def\jin{Jacobian matrix of the invasion vector field with respect to the invasion variables}
\def\jrn{Jacobian matrix of the resident vector field with respect to the resident variables}
\def\jxn{Jacobian matrix of the $\x$ vector field with respect to the $\x$ variables}
\def\Lm{Laplacian matrix}
\def\MTT{matrix tree theorem}
\def\SGT{spectral graph theory}
\def\SMW{Sherman–Morrison–Woodbury}

% -- Stability & Dynamics --
\def\absC{absorbtion CTMC}
\def\BTb{Bogdanov-Takens bifurcation}
\def\CEP{competitive exclusion partition}
\def\CTMC{continuous time Markov chain}
\def\ds{differential system}
\def\fp{fixed point}\def\bfp{boundary fixed point}\def\fbp{fixed boundary point}
\def\fbl{fixed boundary line}\def\fif{forward invariant face}
\def\gas{global asymptotic stability}\def\GAS{globally asymptotically stable}
\def\GASS{\text{GASS}}
\def\Hb{Hopf bifurcation}
\def\Hur{Hurwitz}\def\MH{Metzler-Hurwitz}\def\RH{Routh-Hurwitz}
\def\IS{\text{IS}}
\def\LAS{locally asymptotically stable}
\def\LdS{Lyapunov diagonally stable}
\def\Lf{Lyapunov function}\def\LyV{Lyapunov Volterra type function}\def\RLF{robust Lyapunov function}
\def\LyM{Lyapunov Malkin stable}
\def\PF{Perron-Frobenius}
\def\spfp{strictly positive fixed point}
\def\sta{stability}
\def\STP{sharp threshold property}\def\STT{sharp threshold theorem}

% -- Models (with \varphi explicitly instead of \f) --
\def\AB{$(A,B)$ Arino-Brauer epidemic models}
\def\ABa{$(A,\bbe,\ba,\varphi)$ balanced bilinear model}
\def\ABP{%$(A,\bbe,P,\varphi)$
\\ balanced bilinear model}
\def\AMG{Anderson May Gupta}
\newcommand\BIN{biological interaction network}
\def\Dt{Descartes type}
\def\gLV{generalized Lotka-Volterra model}
\def\IG{\text{invasion graph}}\def\IGMS{interaction graph on minimal siphons}
\def\LCP{\text{linear complementarity problem}}
\def\LV{Lotka-Volterra}\def\LVK{Lotka-Volterra-Kolmogorov}

% -- Splitting & Structure --
\def\reg{regular splitting}\def\regS{{\bf regular splitting}}
\def\RUR{rational univariate representation}
\def\bds{block-diagonal structure}\def\NGMs{next generation matrices}

% -- Miscellaneous --
\def\Ma{{Mathematica}}
\def\Mm{Max–Min}
\def\num{numerator}
\def\rap{rational approximation} \def\DSR{directed species-reaction bipartite graph}\def\PCP{polynomial complementarity problem}\def\rat{rational substitution}
\def\QSS{quasi steady state} \def\RHS{right hand side} \def\LHS{left hand side}
\def\OSN{online social network}\def\VdD{Van den Driessche}
\def\wli{$\omega$-limit} \def\rk{rich kinetics}\def\sJ{symbolic Jacobian}
\def\cV{\cite{Vas,VasSta,BSV,GSV}} \def\spf{semi-parametric semi-functional}
\def\CRW{Capasso-Ruan-Wang} \def\ABP{$(A,\bbe,P,\f)$ balanced bilinear model}
\def\ABa{%$(A,\bbe,\ba,\f)$ \\
rank-one balanced bilinear model}
\def\bbe{B%{\bff \beta}
} \def\vbe{\vec {\bff \beta}
}
\def\vb{\vec {\bff b}}
\newcommand{\Kh}{\hat{K}}          % K with hat, first mention only
\newcommand{\wS}{w^{\sigma}}       % right Perron J^\perp, full block
\newcommand{\piS}{\vec \pi^{\sigma}}    % left Perron J^\perp, full block
\newcommand{\ws}{w^{\sigma}}       % right Perron J^\perp, sub-block
\newcommand{\pis}{\pi^{\sigma}}    % left Perron J^\perp, sub-block
\newcommand{\hwS}{\hat{w}^{\sigma}}   % right Perron K, full block
\newcommand{\hpiS}{\hat{\pi}^{\sigma}} % left Perron K, full block
\newcommand{\hws}{\hat{w}^{\sigma}}   % right Perron K, sub-block
\newcommand{\hpis}{\hat{\pi}^{\sigma}} % left Perron K, sub-block
\newcommand{\twS}{\tilde{w}^{\sigma}}  % right Perron \tilde K, full block
\newcommand{\tpiS}{\tilde{\pi}^{\sigma}} % left Perron \tilde K
\newcommand{\tws}{\tilde{w}^{\sigma}}
\newcommand{\tpis}{\tilde{\pi}^{\sigma}}
\newcommand{\Mat}[2]{\mathbb{R}_{\ge 0}^{\,#1\times #2}}
\def\tJ{J_\sigma^\perp} \def\lPe{left Perron eigenvector}
% \G is \Gamma (defined at :795).  The Volterra function keeps its own name,
% so that \G no longer has two conflicting meanings in this file.
\def\GV{G}  % Volterra G(u)=u-1-\ln u
\def\PVL{Perron-Volterra Lyapunov minimal siphon ansatz}
\newcommand{\istar}[1]{i_{#1}^{*}}
%\newcommand{\sE}{s_E}
\newcommand{\vE}{v_E}
\def\PV{Perron-Volterra candidate Lyapunov function}
\def\GLV{generalized Lotka-Volterra}
\def\MERN{mathematical epidemiology reaction networks}
\def\CMM{contraction metric method}
\newcommand{\sign}{\operatorname{sign}}
\newcommand{\Var}{\operatorname{Var}} 

% ======================================================================
% IDEAS SHARED BETWEEN THE ABSTRACT AND THE MAIN RESULTS SECTION
%
% Write each idea ONCE, in the abstract; the Main Results section recalls it,
% so the two can never drift apart.
%
%   in the abstract        \ideaitem{1}{The first-order reaction terms ...}
%   in Main Results        \ideahead{1}
%
% \ideaitem stores the text globally and typesets
%   \item \hyperlink{idea:1}{<text>}
% \ideahead typesets
%   \item\hypertarget{idea:1}{}\textbf{<the same text>}
%
% The abstract must come before the Main Results input, which is the normal
% order. If a heading has to read differently from the abstract sentence (for
% instance a long abstract item summarised in one line), override it with the
% optional argument: \ideahead[short heading]{6}.
% ======================================================================
\makeatletter
\newcommand{\ideaitem}[2]{%
  \expandafter\gdef\csname @idea@#1\endcsname{#2}%
  \item\hyperlink{idea:#1}{#2}}
\newcommand{\ideahead}[2][]{%
  \item\hypertarget{idea:#2}{}%
  \def\@ideatmp{#1}%
  \ifx\@ideatmp\@empty
    \@ifundefined{@idea@#2}{%
      \PackageWarning{preamble}{Idea #2 is used in Main Results but was never
        defined by \string\ideaitem\space in the abstract}%
      \textbf{[idea #2 undefined]}%
    }{\textbf{\begingroup\def\ideaterm##1##2{##2}\def\idearest##1{}\@nameuse{@idea@#2}\endgroup}}%
  \else
    \textbf{#1}%
  \fi}
% \ideaterm{<hypertarget>}{<words>}: inside an abstract idea the words become an
% invisible link (the document sets \hypersetup{hidelinks}) jumping to the place
% in Main Results that defines them; inside the Main Results heading itself,
% where the target sits a few lines below, \ideahead strips the link and prints
% the words plainly. Mark the destination with \hypertarget{<name>}{} on the line
% before the display, which adds no equation number and shows nothing.
\newcommand{\ideaterm}[2]{\hyperlink{#1}{#2}}
% \idearest{<text>}: the part of an abstract idea after its first sentence; printed in
% the abstract, omitted from the Main Results heading.
\newcommand{\idearest}[1]{#1}
\makeatother

% ======================================================================
% SHORT-PAPER SWITCH
% \shorttrue in a master that assembles only part of the companion paper
% (Barrat/SipArxiv.tex). \nsh{...} drops text that cites results living in the
% omitted sections, so the short build has no dangling references while the
% companion paper keeps the text verbatim.
% ======================================================================
\newif\ifshort \shortfalse
\newcommand{\nsh}[1]{\ifshort\else#1\fi}
% \shortalt{<short build>}{<companion paper>}: use when a sentence needs
% different wording in the two builds, not merely a deletion.
\newcommand{\shortalt}[2]{\ifshort#1\else#2\fi}
% \Jtr, \JtrS{\Sigma}: transverse Jacobian. J^\perp here; SipArxiv.tex redefines them as J_{\mathrm{tra}}, J_{\mathrm{tra},\Sigma}
\newcommand{\Jtr}{J^\perp}
\newcommand{\JtrS}[1]{J_{#1}^\perp}
, and must not clash.
\providecommand{\Q}{\mathbb{Q}}
\providecommand{\res}{\operatorname{res}}   % resultant; used by epiRes.tex
\newcommand{\MN}{\mathbb{N}}
\newcommand{\MZ}{\mathbb{Z}}
\newcommand{\realpos}{\mathbb{R}_{> 0}}
\newcommand{\realnneg}{\mathbb{R}_{\geq 0}}
\newcommand{\naturalpos}{\mathbb{N}_{>0}}
\newcommand{\Rplus}{\mathbb{R}_+}
\newcommand{\Rp}{\mathbb{R}_{\ge 0}}
\def\Rnne{\mathbb{R}_{\ge 0}}
\def\Se{\bar{S}} \def\Ie{\bar{I}}
\newcommand{\Rpos}{\mathbb{R}_{>0}}
\newcommand{\dd}{\,\mathrm d}
\newcommand{\e}{\mathrm e}
%% ========================================
%% GREEK LETTERS (2+ character abbreviations) greek
%% ========================================
\def\al{\alpha}\def\ba{\bff \alpha}
\def\w{\omega}\def\Lgn{\ga+\mu+\de}\def\om{\omega}
\def\de{\delta}\def\De{\Delta}\def\bde{\bff \delta}
\def\ep{\epsilon}\def\eps{\varepsilon}
\def\f{\varphi}  % Kept as requested
\def\ga{\gamma}\def\Ga{\Gamma}\def\G{\Gamma}\def\bg{\bff \gamma}
\def\up{\upsilon}  % NOT the breve accent: no accent commands anywhere.
                   % Named \up, not \u, so the LaTeX breve stays reachable.
\def\la{\lambda}\def\La{\Lambda}
\def\si{\sigma}\def\ka{\kappa}
\def\th{\theta}\def\vt{\vartheta}
\def\ze{\zeta}\def\Gap{\Gamma^{+}} \def\Gam{\Gamma^{-}}


%% ========================================
%% SPECIAL SHORT COMMANDS (kept as requested)
%% ========================================
\def\I{\infty}
\def\T{\widetilde}
\def\H{\hat}
\newcommand{\relint}{\operatorname{relint}}

%% ========================================
%% BOLD/VECTOR NOTATION
%% ========================================
\newcommand{\bff}[1]{{\mbox{\boldmath$#1$}}}
\newcommand\X{\boldsymbol{X}}
\def\x{\boldsymbol{x}}\def\bx{\bff x}\def\xe{\bar{\x}}
\newcommand\y{\boldsymbol{y}}\def\by{\boldsymbol{y}}\def\ye{\bar{\y}}
\def\yb{\y}\def\xb{\x}
\def\xy{\x-\y}
\newcommand\br{\boldsymbol{r}}\def\re{\bar{\r}}
\newcommand\bd{\boldsymbol{d}}
\newcommand\bQ{\boldsymbol{Q}}
\newcommand{\Yb}{{\mathbf{Y}}}
\newcommand{\Ab}{\mathbf{A}}
\newcommand{\zb}{\mathbf{z}}\def\bz{\bff 0}
\newcommand{\Zb}{\mathbf{Z}}
\newcommand{\bb}{\mathbf{b}}
\def\vbe{\vec {\bff \beta}}
\def\vb{\vec {\bff b}}
\def\v1{\vec {\bff 1}}
\def\vn{\vec \de}
\def\vz{\vec 0}
\def\vi{\overset{\rightarrow}{\mathsf i}}
\def\vv{\overset{\rightarrow}{v}}
\def\vr{\; \vec {\mathsf r}}
\newcommand{\va}{\bff a}
\providecommand{\vect}[1]{\boldsymbol{#1}}   % bold vector (moved from Gradostat.tex)
\providecommand{\one}{\mathbf 1}             % all-ones vector (moved from AliSF)
\providecommand{\spr}{s}                     % spectral bound  (moved from AliSF)
\providecommand{\Ke}{K_{\rm eff}}            % effective carrying capacity
% Critical right/left eigenvectors in the Hopf projection method (AAHOSN/New.tex)
\providecommand{\bq}{\mathbf{q}}
\providecommand{\bp}{\mathbf{p}}


%% ========================================
%% MATHCAL & MATHSCR
%% ========================================
% Mathcal
\newcommand\mc[1]{\mathcal{#1}}
\def\mC{{\mathcal C}}\def\mD{\mathcal D}\def\mE{\mathcal{E}}
\def\mF{\mathcal F}\def\mG{{\mathcal G}}\def\mH{{\mathcal H}}
\def\mI{\mathcal I}\def\mL{\mathcal{L}}\def\mM{{\mathcal M}}
\def\mN{{\mathcal N}}\def\mP{\mathcal{P}}\def\mR{{\mathcal R}}
\def\mS{{\mathcal S}}\def\mT{\mathcal{T}}\def\mV{{\mathcal V}}
\def\mK{\mathcal K}\def\mA{{\mathcal A}}\def\mU{{\mathcal U}}
% Symmetric bilinear and trilinear forms of the Taylor expansion, used in the
% first Lyapunov coefficient (AAHOSN/New.tex)
\providecommand{\Bform}{\mathcal{B}}
\providecommand{\Cform}{\mathcal{C}}

% Blackboard bold.  Kept here, not in the drivers: a \providecommand in a
% driver preamble runs before this file and silently wins, so the same name
% ends up meaning different things in different builds.
\def\IC{\mathbb{C}}\def\IN{\mathbb{N}}\def\IQ{\mathbb{Q}}
\def\IR{\mathbb{R}}\def\IZ{\mathbb{Z}}\def\IG{\mathbb{G}}
\def\II{\mathbb{I}}

% Mathscr
\newcommand{\Rs}{{\mathscr R}}
\newcommand{\Ss}{{\mathscr S}}
\newcommand{\Ns}{{\mathscr N}}
\newcommand{\Cs}{\mathscr C}


%% ========================================
%% OPERATORS & FUNCTIONS
%% ========================================
\DeclareMathOperator{\tr}{tr}
\DeclareMathOperator{\diag}{Diag}
\DeclareMathOperator{\Diag}{Diag}
\DeclareMathOperator{\intr}{\mathrm{Int}}
\DeclareMathOperator{\Inv}{Inv}
\DeclareMathOperator{\spec}{spec}
\DeclareMathOperator{\AUC}{AUC}               % area under the curve
\newcommand{\ii}{\mathrm{i}}                  % imaginary unit
\newcommand\op{\operatorname}
\newcommand{\supp}{\op{supp}}
\newcommand{\diff}{\mathop{}\mathopen{}\mathrm{d}}
\newcommand{\pdd}[1]{\frac{\partial}{\partial{#1}}}
\DeclareMathOperator{\cone}{cone}
\DeclareMathOperator{\proj}{proj}

%% ========================================
%% BRACKETS & NORMS
%% ========================================
\providecommand{\pr}[1]{\left(#1\right)}
\providecommand{\pp}[1]{\left[#1\right]}
\newcommand\croc[1]{\left\langle #1\right\rangle}
\newcommand{\norm}[1]{\|#1\|}
\newcommand{\inner}[2]{\langle #1, #2 \rangle}
% Short synonym of \inner, used by AAHOSN/New.tex as $\ip{u}{v}=\bar u^\top v$
\providecommand{\ip}[2]{\langle #1 , #2 \rangle}


%% ========================================
%% FORMATTING & UTILITIES
%% ========================================
\newcommand*{\Scale}[2][4]{\scalebox{#1}{$#2$}}
\newcommand{\te}[1]{\text{#1}}
\newcommand\steq[1]{\stackrel{\text{\rm #1.}}{=}}
\newcommand\blue[1]{\textcolor{blue}{#1}}
\newcommand{\red}{\textcolor[rgb]{1.00,0.00,0.00}}
\newcommand{\green}{\textcolor[rgb]{0.00,1.00,0.00}}

% Figure command
\newcommand{\figu}[3]{%
\begin{figure}[H]
\centering
\includegraphics[scale=#3]{#1}
\caption{#2\label{f:#1}}
\end{figure}
}

% Footnote
\long\def\symbolfootnote[#1]#2{%
\begingroup
\def\thefootnote{\fnsymbol{footnote}}\footnote[#1]{#2}%
\endgroup}
\def\fn{\symbolfootnote}

\newcommand{\dsum}{\displaystyle\sum}

%% ========================================
%% CRN & EPIDEMIOLOGY SPECIFIC
%% ========================================
% CRN notation
\def\SS{\mathcal S}
\def\CC{\mathcal C}
\def\RR{\mathcal R}

% Reproduction numbers
\def\Rto{R_{1}^{E_2}}
\def\Rot{R_{2}^{E_1}}
\def\Roc{R_{1,c}}
\def\Rtc{R_{2,c}}
\def\rd{r_{0}}
\def\sd{s_{0}}
\def\dd{D_{0}}
\def\vd{v_{0}}
\def\zd{z_{0}}
\def\se{\sE}
\def\xd{x_{0}}
%% ========================================
%% SPECIAL VALUES & EQUILIBRIA
%% ========================================
\def\eE{\bar{e}}
\def\wE{\bar{w}}
\def\sE{s_{*}}
\def\Me{\T M_{ee}}
\def\Rnne{\mathbb{R}_{\ge 0}}

\newcommand{\ones}{\mathbf 1}


%% ========================================
%% ABBREVIATIONS (alphabetical by category)
%% ========================================

% -- General mathematical terms --
\def\adm{admissible}\def\admy{admissibility}
%\def\and{antisymmetric}% disabled 2026-09-30: it overrode LaTeX's author separator \and (printed "antisymmetric" between authors)
\def\ass{assumption}
\def\cadlag{c\`adl\`ag}
\def\ch{characteristic polynomial}
\def\chp{characteristic polynomial}
\def\coe{coefficient}
\def\comp{computation}
\def\con{condition}
\def\coo{coordinate}
\def\corr{corresponding}
\def\dec{decomposition}
\def\den{denominator}
\def\dete{deterministic}
\def\eig{eigenvalue}\def\eval{eigenvalue}
\def\eigv{eigenvector}\def\eiv{eigenvector}\def\evec{eigenvector}
\def\elt{element}
\def\enn{essentially nonnegative/positivity preserving}
\def\equ{equation}
\def\etal{et al.}
\def\fe{for example}
\def\fno{from now on}
\def\for{\forall}
\def\form{formula}
\def\Fr{Furthermore, }
\def\frt{furthermore}
\def\fun{function}
\def\gene{generalization}
\def\heu{heuristic}
\def\how{however}\def\How{However, }
\def\ie{i.e. }
\def\im{\item}
\def\isd{is defined by}
\def\itc{it is easy to check that}
\def\itf{it follows that}\def\Itf{It follows that}
\def\ith{it holds that}
\def\itm{it may be checked that}\def\Itm{It may be checked that}
\def\ito{it turns out that}\def\Ito{It turns out that}\def\Itou{It turns out useful to}
\def\ma{mass action}\def\MAR{mass action representation}
\def\mat{matrix}
\def\mbc{may be checked to be}
\def\mbw{may be written as}\def\mubw{must be written as}
\def\mpr{more precisely, }\def\Mp{More precisely, }
\def\nc{non-catalytic}
\def\nne{non-negative}\def\nny{nonnegativity}
\def\no{\nonumber}
\def\noni{non-infectious}
\def\nota{notation}
\def\np{\newpage}
\def\Nt{Note that}
\def\nt{non trivial}\def\nz{non zero}
\def\oth{otherwise}
\def\Oth{On the other hand}
\def\pa{parameter}\def\para{parameter}
\def\pol{polynomial}
\def\prob{probability}
\def\Prf{{\bf  Proof:}}\def\prf{{\bf  Proof:}}
\def\prop{proportional to}
\def\qu{\quad}
\def\QED{\hfill {$\square$}\goodbreak \medskip}
\def\rep{representation}
\def\resp{respectively}
\def\satd{satisfied}\def\satg{satisfying}\def\sats{satisfies}\def\saty{satisfy}
\def\ssec{\subsection}\def\sssec{\subsubsection}
\def\sef{see for example}
\def\sksy{skew-symmetric}
\def\sub{substrate}
\def\sys{system}
\def\var{variable}
\def\wk{well-known}
\def\wl{without loss of generality}\def\wlo{w.l.o.g.}
\def\wmc{We may check that}\def\Wmc{We may check that}
\def\wmr{we may rewrite}
\def\wms{we must show that}
\def\wrt{with respect to}

% -- CRN & Chemistry --
\newcommand\CRN{chemical reaction network}
\newcommand\CRNs{chemical reaction networks}
\newcommand\RN{reaction network}
\def\ACR{absolute concentration robustness}
\def\AMSD{acyclic siphon decomposition}
\def\DAE{differential algebraic equations}
\def\FHJ{Feinberg-Horn-Jackson}
\def\GMAK{generalized mass action kinetics}
\def\MM{Michaelis-Menten}
\def\sm{stoichiometric matrix}\def\Sm{stoichiometric matrix}
\def\WRZD{weakly reversible, zero deficiency}\def\aut{autocatalytic}

% -- Epidemiology --
\def\brn{basic reproduction number}\def\bzn{basic replacement number}
\def\com{compartment}
\def\dfe{disease free equilibrium}\def\DFE{disease free equilibrium}
\def\dwt{dwell times}\def\edt{expected dwell times}
\def\EM{epidemic model}
\def\endp{endemic fixed point}\def\enp{endemic point}\def\EE{endemic point}
\def\ME{mathematical epidemiology}
\def\MEE{eco-epidemiology}\def\EEP{eco-epidemiological}
\def\ece{eco-epidemiological model}
\def\mgm{multi group model}
\def\mtS{multi-strain model}\def\mSt{multi-strain model}
\def\mSi{minimal siphon}
\def\NGM{next generation matrix}
\def\nim{new infections matrix}
\def\qem{quadratic epidemiological model}
\def\repF{invasion functions}
\def\SPF{SIR-PH-FA model}
\def\sus{susceptible}

% -- Graph theory & Linear algebra --
\def\jac{Jacobi matrix}
\def\jin{Jacobian matrix of the invasion vector field with respect to the invasion variables}
\def\jrn{Jacobian matrix of the resident vector field with respect to the resident variables}
\def\jxn{Jacobian matrix of the $\x$ vector field with respect to the $\x$ variables}
\def\Lm{Laplacian matrix}
\def\MTT{matrix tree theorem}
\def\SGT{spectral graph theory}
\def\SMW{Sherman–Morrison–Woodbury}

% -- Stability & Dynamics --
\def\absC{absorbtion CTMC}
\def\BTb{Bogdanov-Takens bifurcation}
\def\CEP{competitive exclusion partition}
\def\CTMC{continuous time Markov chain}
\def\ds{differential system}
\def\fp{fixed point}\def\bfp{boundary fixed point}\def\fbp{fixed boundary point}
\def\fbl{fixed boundary line}\def\fif{forward invariant face}
\def\gas{global asymptotic stability}\def\GAS{globally asymptotically stable}
\def\GASS{\text{GASS}}
\def\Hb{Hopf bifurcation}
\def\Hur{Hurwitz}\def\MH{Metzler-Hurwitz}\def\RH{Routh-Hurwitz}
\def\IS{\text{IS}}
\def\LAS{locally asymptotically stable}
\def\LdS{Lyapunov diagonally stable}
\def\Lf{Lyapunov function}\def\LyV{Lyapunov Volterra type function}\def\RLF{robust Lyapunov function}
\def\LyM{Lyapunov Malkin stable}
\def\PF{Perron-Frobenius}
\def\spfp{strictly positive fixed point}
\def\sta{stability}
\def\STP{sharp threshold property}\def\STT{sharp threshold theorem}

% -- Models (with \varphi explicitly instead of \f) --
\def\AB{$(A,B)$ Arino-Brauer epidemic models}
\def\ABa{$(A,\bbe,\ba,\varphi)$ balanced bilinear model}
\def\ABP{%$(A,\bbe,P,\varphi)$
\\ balanced bilinear model}
\def\AMG{Anderson May Gupta}
\newcommand\BIN{biological interaction network}
\def\Dt{Descartes type}
\def\gLV{generalized Lotka-Volterra model}
\def\IG{\text{invasion graph}}\def\IGMS{interaction graph on minimal siphons}
\def\LCP{\text{linear complementarity problem}}
\def\LV{Lotka-Volterra}\def\LVK{Lotka-Volterra-Kolmogorov}

% -- Splitting & Structure --
\def\reg{regular splitting}\def\regS{{\bf regular splitting}}
\def\RUR{rational univariate representation}
\def\bds{block-diagonal structure}\def\NGMs{next generation matrices}

% -- Miscellaneous --
\def\Ma{{Mathematica}}
\def\Mm{Max–Min}
\def\num{numerator}
\def\rap{rational approximation} \def\DSR{directed species-reaction bipartite graph}\def\PCP{polynomial complementarity problem}\def\rat{rational substitution}
\def\QSS{quasi steady state} \def\RHS{right hand side} \def\LHS{left hand side}
\def\OSN{online social network}\def\VdD{Van den Driessche}
\def\wli{$\omega$-limit} \def\rk{rich kinetics}\def\sJ{symbolic Jacobian}
\def\cV{\cite{Vas,VasSta,BSV,GSV}} \def\spf{semi-parametric semi-functional}
\def\CRW{Capasso-Ruan-Wang} \def\ABP{$(A,\bbe,P,\f)$ balanced bilinear model}
\def\ABa{%$(A,\bbe,\ba,\f)$ \\
rank-one balanced bilinear model}
\def\bbe{B%{\bff \beta}
} \def\vbe{\vec {\bff \beta}
}
\def\vb{\vec {\bff b}}
\newcommand{\Kh}{\hat{K}}          % K with hat, first mention only
\newcommand{\wS}{w^{\sigma}}       % right Perron J^\perp, full block
\newcommand{\piS}{\vec \pi^{\sigma}}    % left Perron J^\perp, full block
\newcommand{\ws}{w^{\sigma}}       % right Perron J^\perp, sub-block
\newcommand{\pis}{\pi^{\sigma}}    % left Perron J^\perp, sub-block
\newcommand{\hwS}{\hat{w}^{\sigma}}   % right Perron K, full block
\newcommand{\hpiS}{\hat{\pi}^{\sigma}} % left Perron K, full block
\newcommand{\hws}{\hat{w}^{\sigma}}   % right Perron K, sub-block
\newcommand{\hpis}{\hat{\pi}^{\sigma}} % left Perron K, sub-block
\newcommand{\twS}{\tilde{w}^{\sigma}}  % right Perron \tilde K, full block
\newcommand{\tpiS}{\tilde{\pi}^{\sigma}} % left Perron \tilde K
\newcommand{\tws}{\tilde{w}^{\sigma}}
\newcommand{\tpis}{\tilde{\pi}^{\sigma}}
\newcommand{\Mat}[2]{\mathbb{R}_{\ge 0}^{\,#1\times #2}}
\def\tJ{J_\sigma^\perp} \def\lPe{left Perron eigenvector}
% \G is \Gamma (defined at :795).  The Volterra function keeps its own name,
% so that \G no longer has two conflicting meanings in this file.
\def\GV{G}  % Volterra G(u)=u-1-\ln u
\def\PVL{Perron-Volterra Lyapunov minimal siphon ansatz}
\newcommand{\istar}[1]{i_{#1}^{*}}
%\newcommand{\sE}{s_E}
\newcommand{\vE}{v_E}
\def\PV{Perron-Volterra candidate Lyapunov function}
\def\GLV{generalized Lotka-Volterra}
\def\MERN{mathematical epidemiology reaction networks}
\def\CMM{contraction metric method}
\newcommand{\sign}{\operatorname{sign}}
\newcommand{\Var}{\operatorname{Var}} 

% ======================================================================
% IDEAS SHARED BETWEEN THE ABSTRACT AND THE MAIN RESULTS SECTION
%
% Write each idea ONCE, in the abstract; the Main Results section recalls it,
% so the two can never drift apart.
%
%   in the abstract        \ideaitem{1}{The first-order reaction terms ...}
%   in Main Results        \ideahead{1}
%
% \ideaitem stores the text globally and typesets
%   \item \hyperlink{idea:1}{<text>}
% \ideahead typesets
%   \item\hypertarget{idea:1}{}\textbf{<the same text>}
%
% The abstract must come before the Main Results input, which is the normal
% order. If a heading has to read differently from the abstract sentence (for
% instance a long abstract item summarised in one line), override it with the
% optional argument: \ideahead[short heading]{6}.
% ======================================================================
\makeatletter
\newcommand{\ideaitem}[2]{%
  \expandafter\gdef\csname @idea@#1\endcsname{#2}%
  \item\hyperlink{idea:#1}{#2}}
\newcommand{\ideahead}[2][]{%
  \item\hypertarget{idea:#2}{}%
  \def\@ideatmp{#1}%
  \ifx\@ideatmp\@empty
    \@ifundefined{@idea@#2}{%
      \PackageWarning{preamble}{Idea #2 is used in Main Results but was never
        defined by \string\ideaitem\space in the abstract}%
      \textbf{[idea #2 undefined]}%
    }{\textbf{\begingroup\def\ideaterm##1##2{##2}\def\idearest##1{}\@nameuse{@idea@#2}\endgroup}}%
  \else
    \textbf{#1}%
  \fi}
% \ideaterm{<hypertarget>}{<words>}: inside an abstract idea the words become an
% invisible link (the document sets \hypersetup{hidelinks}) jumping to the place
% in Main Results that defines them; inside the Main Results heading itself,
% where the target sits a few lines below, \ideahead strips the link and prints
% the words plainly. Mark the destination with \hypertarget{<name>}{} on the line
% before the display, which adds no equation number and shows nothing.
\newcommand{\ideaterm}[2]{\hyperlink{#1}{#2}}
% \idearest{<text>}: the part of an abstract idea after its first sentence; printed in
% the abstract, omitted from the Main Results heading.
\newcommand{\idearest}[1]{#1}
\makeatother

% ======================================================================
% SHORT-PAPER SWITCH
% \shorttrue in a master that assembles only part of the companion paper
% (Barrat/SipArxiv.tex). \nsh{...} drops text that cites results living in the
% omitted sections, so the short build has no dangling references while the
% companion paper keeps the text verbatim.
% ======================================================================
\newif\ifshort \shortfalse
\newcommand{\nsh}[1]{\ifshort\else#1\fi}
% \shortalt{<short build>}{<companion paper>}: use when a sentence needs
% different wording in the two builds, not merely a deletion.
\newcommand{\shortalt}[2]{\ifshort#1\else#2\fi}
% \Jtr, \JtrS{\Sigma}: transverse Jacobian. J^\perp here; SipArxiv.tex redefines them as J_{\mathrm{tra}}, J_{\mathrm{tra},\Sigma}
\newcommand{\Jtr}{J^\perp}
\newcommand{\JtrS}[1]{J_{#1}^\perp}
